\documentclass[11pt,a4paper]{article}
\pdfextension mapfile{+cm.map}
\pdfextension mapfile{+cmextra.map}
\pdfextension mapfile{+symbols.map}
\pdfextension mapfile{+euler.map}
\pdfextension mapfile{+latxfont.map}
\pdfextension mapfile{+rsfs.map}
\usepackage[margin=25mm,headheight=15pt]{geometry}
\usepackage{fontspec}
\setmainfont{Latin Modern Roman}
\usepackage{amsmath,amssymb,amsthm,amscd,mathtools,mathrsfs}
\usepackage{microtype,booktabs,longtable,array,xcolor,graphicx,enumitem,needspace,fancyhdr}
\usepackage[round,authoryear]{natbib}
\usepackage[unicode,colorlinks=true,linkcolor=blue!40!black,citecolor=blue!40!black,urlcolor=blue!40!black]{hyperref}
\usepackage{xurl}
\usepackage{listings}
\usepackage{tabularx}
\lstset{basicstyle=\ttfamily\small,breaklines=true,columns=fullflexible,keepspaces=true}
\usepackage{tikz}
\usepackage{tikz-cd}
\usetikzlibrary{arrows.meta,calc,patterns,positioning}
\usepackage{aliascnt}
\usepackage[nameinlink,noabbrev]{cleveref}
\allowdisplaybreaks[2]
\setcounter{MaxMatrixCols}{32}
\setlength{\emergencystretch}{3em}
\setlength{\parskip}{2pt}
\newtheorem{theorem}{Theorem}[section]
\newaliascnt{proposition}{theorem}
\newtheorem{proposition}[proposition]{Proposition}
\aliascntresetthe{proposition}
\crefname{proposition}{proposition}{propositions}
\newaliascnt{lemma}{theorem}
\newtheorem{lemma}[lemma]{Lemma}
\aliascntresetthe{lemma}
\crefname{lemma}{lemma}{lemmas}
\newaliascnt{corollary}{theorem}
\newtheorem{corollary}[corollary]{Corollary}
\aliascntresetthe{corollary}
\crefname{corollary}{corollary}{corollarys}
\newaliascnt{claim}{theorem}
\newtheorem{claim}[claim]{Claim}
\aliascntresetthe{claim}
\crefname{claim}{claim}{claims}
\newaliascnt{statement}{theorem}
\newtheorem{statement}[statement]{Statement}
\aliascntresetthe{statement}
\crefname{statement}{statement}{statements}
\theoremstyle{definition}
\newaliascnt{definition}{theorem}
\newtheorem{definition}[definition]{Definition}
\aliascntresetthe{definition}
\crefname{definition}{definition}{definitions}
\newaliascnt{remark}{theorem}
\newtheorem{remark}[remark]{Remark}
\aliascntresetthe{remark}
\crefname{remark}{remark}{remarks}
\newaliascnt{example}{theorem}
\newtheorem{example}[example]{Example}
\aliascntresetthe{example}
\crefname{example}{example}{examples}
\newaliascnt{problem}{theorem}
\newtheorem{problem}[problem]{Open problem}
\aliascntresetthe{problem}
\crefname{problem}{open problem}{open problems}
\newaliascnt{question}{theorem}
\newtheorem{question}[question]{Question}
\aliascntresetthe{question}
\crefname{question}{question}{questions}
\newaliascnt{auditrecord}{theorem}
\newtheorem{auditrecord}[auditrecord]{Verification record}
\aliascntresetthe{auditrecord}
\crefname{auditrecord}{verification record}{verification records}
\newaliascnt{test}{theorem}
\newtheorem{test}[test]{Verification}
\aliascntresetthe{test}
\crefname{test}{verification}{verifications}
\newaliascnt{conjecture}{theorem}
\newtheorem{conjecture}[conjecture]{Conjecture}
\aliascntresetthe{conjecture}
\crefname{conjecture}{conjecture}{conjectures}
\newcommand{\C}{\mathbb C}\newcommand{\R}{\mathbb R}\newcommand{\Z}{\mathbb Z}\newcommand{\Q}{\mathbb Q}\newcommand{\N}{\mathbb N}\newcommand{\F}{\mathbb F}
\newcommand{\CC}{\mathbb C}\newcommand{\RR}{\mathbb R}\newcommand{\ZZ}{\mathbb Z}\newcommand{\QQ}{\mathbb Q}\newcommand{\NN}{\mathbb N}
\newcommand{\PP}{\mathbb P}\newcommand{\Pp}{\mathbb P}\newcommand{\Pone}{\mathbb P^1}\newcommand{\HH}{\mathbb H}
\newcommand{\RP}{\mathbb{RP}}\newcommand{\CP}{\mathbb{CP}}
\newcommand{\OO}{\mathcal O}\newcommand{\cO}{\mathcal O}\newcommand{\cM}{\mathcal M}\newcommand{\cF}{\mathcal F}
\newcommand{\h}{\mathfrak h}\newcommand{\hz}{\mathfrak h_z}\newcommand{\Acon}{\mathfrak A}\newcommand{\wt}[1]{\widehat{#1}}
\newcommand{\red}{\mathrm{red}}\newcommand{\tor}{\mathrm{tor}}\newcommand{\dP}{\mathrm{dP}}
\newcommand{\gh}{\widehat\gamma}\newcommand{\uh}{\widehat u}\newcommand{\wh}{\widehat w}\newcommand{\dhat}{\widehat\delta}\renewcommand{\dh}{\widehat\delta}
\newcommand{\dd}{\mathrm d}\newcommand{\ii}{\mathrm i}\newcommand{\eps}{\varepsilon}\newcommand{\partt}{\mathrm{part}}\newcommand{\Bcirc}{B^\circ}\newcommand{\ddbar}{\partial\bar\partial}
\DeclareMathOperator{\Pic}{Pic}\DeclareMathOperator{\Alb}{Alb}\DeclareMathOperator{\Aut}{Aut}\DeclareMathOperator{\NS}{NS}
\DeclareMathOperator{\Hom}{Hom}\DeclareMathOperator{\End}{End}\DeclareMathOperator{\Ext}{Ext}\DeclareMathOperator{\Tor}{Tor}
\DeclareMathOperator{\Spec}{Spec}\DeclareMathOperator{\Proj}{Proj}\DeclareMathOperator{\Sym}{Sym}\DeclareMathOperator{\Diff}{Diff}
\DeclareMathOperator{\GL}{GL}\DeclareMathOperator{\SL}{SL}\DeclareMathOperator{\PSL}{PSL}\DeclareMathOperator{\Sp}{Sp}
\DeclareMathOperator{\im}{im}\DeclareMathOperator{\coker}{coker}\DeclareMathOperator{\rank}{rank}\DeclareMathOperator{\diag}{diag}
\DeclareMathOperator{\ord}{ord}\DeclareMathOperator{\vol}{vol}\DeclareMathOperator{\pr}{pr}\DeclareMathOperator{\id}{id}
\DeclareMathOperator{\divisor}{div}\DeclareMathOperator{\Tr}{Tr}\DeclareMathOperator{\Nm}{Nm}\DeclareMathOperator{\Sha}{Sha}
\DeclareMathOperator{\Imn}{Im}\DeclareMathOperator{\Ren}{Re}\DeclareMathOperator{\trdeg}{tr.deg}
\newcommand{\Uses}[1]{\par\smallskip\noindent\textit{Uses:} #1.\par\smallskip}
\newcommand{\Verified}{}\newcommand{\AuthorSelfCheck}{}\newcommand{\NeedsCheck}{}\newcommand{\Open}{\textbf{Open.}\ }

\newcommand{\StatementRecord}[4]{\par\smallskip\begingroup\upshape\raggedright\noindent\textbf{Label:} \nolinkurl{#1}.\par\noindent\textbf{Status:} #2.\par\noindent\textbf{Check record:} #3.\par\noindent\textbf{Uses:} #4.\par\endgroup\smallskip}
\newcommand{\RungZero}[1]{\par\noindent\textbf{Result and mechanism.} #1\par\smallskip}
\newcommand{\RungOne}{\par\Needspace{6\baselineskip}\noindent\textbf{Proof at the first rung.}\par\nobreak\smallskip}
\newcommand{\CurrentChecker}{GPT-6 Astra (OpenAI, ChatGPT Pro mode, effort Pro), 18 November 2026}
\newaliascnt{cited}{theorem}
\newtheorem{cited}[cited]{Cited result}
\aliascntresetthe{cited}
\crefname{cited}{cited result}{cited results}
\newaliascnt{refuted}{theorem}
\newtheorem{refuted}[refuted]{Correction}
\aliascntresetthe{refuted}
\crefname{refuted}{correction}{corrections}
\newaliascnt{computed}{theorem}
\newtheorem{computed}[computed]{Computation}
\aliascntresetthe{computed}
\crefname{computed}{computation}{computations}
\newaliascnt{openq}{theorem}
\newtheorem{openq}[openq]{Open question}
\aliascntresetthe{openq}
\crefname{openq}{open question}{open questions}

\providecommand{\rk}{\operatorname{rk}}
\providecommand{\Ann}{\operatorname{Ann}}
\providecommand{\Log}{\operatorname{Log}}
\providecommand{\vir}{\mathrm{vir}}

\begin{document}
\title{Paper VI}\author{AI-produced mathematical text}\date{19 November 2026}\maketitle
\begin{abstract}
Let $f:Y\to B$ be a proper holomorphic map with connected fibres
from a smooth compact complex threefold with $b_1(Y)=b_2(Y)=0$,
and suppose that its smooth general fibre is a complex two-torus.
We prove that every reduced fibre support has nonnegative Euler
number, with arbitrary component multiplicities and arbitrary
support singularities. Consequently $B\simeq\mathbb P^1$ and
$c_3(Y)\in\{0,2\}$. Finite-character support cohomology computes
the Euler number of an irreducible support with rational or ruled
normalization resolution; local derivative and conductor
comparisons identify the zero-Euler normalizations. For reducible
supports, the actual Picard restrictions exclude elliptic--Enoki
interfaces, and a local vanishing sphere excludes the remaining
all-Enoki case. The conductor Euler formula then gives both
nonnegativity and the complete equality criterion. The supporting
occurrence constructions, multiplicity calculations, contractions
and comparison examples are included with their precise maps.
\end{abstract}

\paragraph{The proof's scope.}
The irreducible argument uses the local comparisons of Paper~II
for the single support in question, together with the
finite-character mixed-Hodge calculation proved here. Its use
does not require the other supports to be irreducible or to have
normal crossings. For a singular irreducible support of Euler
number zero, the actual normalization is smooth minimal elliptic
ruled or smooth minimal Hopf. For a reducible zero-Euler support,
the actual normalization components are smooth minimal Hopf or
non-K\"ahler elliptic surfaces and their reduced conductor maps
are finite unramified covers of disjoint smooth elliptic curves.

\paragraph{Known inputs and the additional analytic form.}
The value $e(Y)=2$ for the supplied completed torus family
is stated in the original manuscript
\cite[Theorem on p.~4; Theorem~7.22, p.~60]{AC26}.
Engel computes this value from the two zero-dimensional
toric strata in \cite[Proposition~7.1]{Engel2026}.
The relative cusp construction is also stated in
\cite[Section~2, relative cusp theorem]{DGKX2026v1,DGKX2026v2}.
These existing constructions and Euler computations retain
their credit.  The nonnegativity and equality conclusions
of \ref{cdp:st:C.2.G.1} are proved for every family satisfying
its displayed hypotheses by the local comparisons below.
The classification in
\cite[Section~1 and Theorem~3.1]{KamenovaTori} assumes
a semistable K\"ahler total space with trivial canonical
bundle; each use of that classification retains those
hypotheses.
The surface-classification facts are cited at
\cite[Proposition~2.1 and Theorem~2.2, pp.~842--843]{Peters2021}.
The logarithmic support complex uses
\cite[Section~3, Theorem~1 and Corollary~2; Section~4, equation~(3)]{Steenbrink1999}.
That source works with a complete algebraic ambient space.
The compact-projective-strata construction required for the
analytic neighbourhood here is proved explicitly in
\ref{cdp:st:C.2.MH.1}; the source is not cited as already
containing that broader form. The class-VII foliation and
negative-definite curve inputs are cited at the exact loci
opened in Favre and Dloussky--Oeljeklaus.

The cyclic symplectic torus orders used in
\ref{cdp:st:C.2.UN.24} and \ref{cdp:st:C.2.UN.32}
are classical \cite[Lemmas~3.2--3.3, pp.~30--31]{Fujiki1988Tori}.
The four cyclic quotient configurations are listed in
\cite[table on p.~7659; Theorem~27, p.~7663]{Roulleau2019Kummer},
with their credit to Nikulin and Bertin retained.
For an algebraic-dimension-zero K3 surface, the cyclic
torus quotient associated with a foliation is already proved in
\cite[Section~10, pp.~29--30]{Brunella2002Foliations}.
The stabilizer action on the saturated limiting-field space
in \ref{cdp:st:C.2.UN.16} has an eigenline.  Its tangent line
descends through the finite quotient on the unramified locus,
then extends by saturation on the K3 resolution, giving the
foliation to which that result applies.  The elliptic-cover
case in \ref{cdp:st:C.2.UN.32} uses the explicit translation
quotient and principal-bundle characteristic-class calculation
to construct its torus.  The occurrence statements retain the
specified quotient component through the actual vertical
modifications.

\paragraph{Reading the verification record.}
Historical checks, current re-derivations, partial current
reviews and source openings are recorded separately. The long
occurrence catalogue has its own per-statement review record;
the assembly of the main theorem supplies no blanket check for
those constructions.

\section*{The sign and its proof at the first rung}
We consider a proper surjective holomorphic map
$f:Y\to B$ with connected fibres, where $Y$ is a connected
smooth compact complex threefold, $B$ is a smooth compact curve,
$b_1(Y)=b_2(Y)=0$, and the general fibre is a complex two-torus.
We retain each full fibre $F_b=\sum_i m_iS_i$ and its
reduced support $S_b$.  The multiplicities $m_i$ are arbitrary
positive integers and the supports have arbitrary singularities.

\RungZero{Every reduced fibre support has nonnegative Euler
number because, on a reducible support, the normalization and conductor
give an exact sum of nonnegative contributions; on an irreducible
support finite-character duality and the local conductor calculation
give the value zero or two.
Euler additivity and the skew middle pairing therefore give
$B=\PP^1$, $c_3(Y)\in\{0,2\}$, and $b_3(Y)=2-c_3(Y)$
(\ref{cdp:st:C.2.G.1}).}
\RungOne
\begin{enumerate}[label=\arabic*.,leftmargin=*]
\item Connected-fibre Leray injects $H^1(B,\Q)$ into
$H^1(Y,\Q)$, so $B=\PP^1$.  The finitely many critical
values isolate every contribution to $e(Y)$ because the
smooth torus has Euler number zero.
\item Finite-character support duality gives $e(S)=2$ for a
rational normalization resolution and $e(S)=2-2g$ for a ruled
one. The local irreducible nonnegativity of Paper~II,
\ref{night:P4.26}, forces $g=1$ in the latter case.
The equality comparison and the saturated-foliation Chern
identities give smooth minimal elliptic ruled or Hopf
normalization (\ref{cdp:st:C.2.MH.2}, \ref{night:P4.32},
\ref{night:P4.36}, \ref{night:P4.37}).
\item On every component of a reducible support, an adjacent
component pulls back to a nonzero effective divisor with
zero rational Chern class.  Integration excludes a K\"ahler
resolution.  The class-VII and elliptic-surface inputs then
leave Hopf, non-K\"ahler elliptic, or Enoki minimal models
(\ref{cdp:st:C.2.INT.1}).
\item The actual Picard restrictions prevent an Enoki
component from meeting a Hopf or elliptic component.
The full-fibre multiplicities and canonical divisor
coefficients reduce an all-Enoki fibre to a nodal
two-component configuration.  Its local equation
$z^a(xy+z^n)^b$ has a vanishing sphere of self-intersection
$-2$ in a nearby torus, contradicting the torus's
contractible universal cover (\ref{cdp:st:C.2.INT.2}).
\item For the surviving components, the actual normalization
maps factor their minimal-model maps.  Coherent conductor
duality gives $\chi(\OO_{C})\ge0$ on the reduced conductor
image $C$.  Curve normalization and the exact blowup counts
then give $e(C)\ge0$ and
$e(T_i\setminus A_i)\ge0$; their sum equals $e(S_b)$
(\ref{cdp:st:C.2.INT.3}).
\item Thus $c_3(Y)=e(Y)=\sum_b e(S_b)\ge0$.
Poincar\'e duality gives $e(Y)=2-b_3(Y)$, and the
nondegenerate skew pairing on $H^3(Y,\Q)$ makes $b_3(Y)$
even.  The only resulting Euler values are zero and two.
\end{enumerate}

\subsection*{The equality criterion}
\RungZero{The equality $c_3(Y)=0$ holds exactly when every
support has its zero-Euler normalization and conductor form.
For reducible supports, equality in each nonnegative summand
forces all blowups and conductor singularities to disappear:
the normalizations are smooth minimal Hopf or non-K\"ahler
elliptic surfaces, and the conductor maps are finite
unramified covers of smooth elliptic curves.}
\RungOne
\begin{enumerate}[label=\arabic*.,leftmargin=*]
\item On an irreducible support, Paper~II identifies the
zero-Euler possibilities as a smooth support with a finite
unramified torus cover, or a singular support with smooth
minimal elliptic ruled or Hopf normalization.
\item For a reducible support, curve normalization writes
$e(C)$ as $\chi(\OO_C)$ plus the number of rational-normalized
components plus the nonnegative defects
$\delta_x-(s_x-1)$.  A zero sum makes $C$ a disjoint union
of smooth elliptic curves.
\item Every rational curve upstairs is then contracted by
the normalization resolution.  A last blowup would produce
a contracted $(-1)$-curve, contradicting minimality of
that resolution.  Thus each normalization is already its
smooth minimal model.
\item Its conductor curves are smooth elliptic curves.
Riemann--Hurwitz makes their finite maps unramified.
Conversely these exact normalization and conductor data
have zero Euler number by the normalization formula.
\item The global Euler number is the sum of nonnegative
support Euler numbers, so it vanishes exactly when every
support meets the corresponding local equality criterion.
\end{enumerate}

\subsection*{The local comparison retained with the global theorem}
\RungZero{For a torus-degeneration support $S$, with nearby-fibre
monodromy $T:H^1(F,\Q)\to H^1(F,\Q)$, the tube and
mapping-torus exact sequences give
$e(S)\equiv\dim\ker(\Lambda^2T-I)\pmod2$; the symmetric
wedge pairing and its Cayley transform make this dimension even,
without the global Betti conditions.
All even integers occur after explicit curve and point blowups
of a product torus fibration
(\ref{cdp:st:C.2.UN.1}, \ref{cdp:st:C.2.UN.2}).}
\RungOne
\begin{enumerate}[label=\arabic*.,leftmargin=*]
\item A small full fibre tube $U$ retracts onto the support
and has boundary the mapping torus $M$ of an orientation-
preserving automorphism of the real four-torus.
\item The pair sequence and Poincar\'e--Lefschetz duality
give $\chi(U)\equiv b_0(M)+b_1(M)+b_2(M)\pmod2$;
the omitted middle intersection rank is even.
\item Wang's sequence identifies this parity with the
dimension of the fixed subspace of $\Lambda^2T$.
The symmetric wedge pairing and its Cayley transform
make $\operatorname{rank}(\Lambda^2T-I)$ even, proving
the parity assertion.
\item A blowup along a genus-two curve in a special
torus fibre changes its support Euler number by $-2$;
a point blowup changes it by $2$.  The actual iterated
centres in \ref{cdp:st:C.2.UN.2} realize every integer
$2(m-k)$ while leaving every nearby fibre the original
torus.
\end{enumerate}

\section{The degeneration and its normalization maps}
Throughout this paper, $f:Y\to B$ is a proper holomorphic map with connected fibres from a smooth compact complex threefold. The sign theorem assumes $b_1(Y)=b_2(Y)=0$ and a smooth general fibre that is a complex two-torus. Local comparison constructions state their own hypotheses; they are used with precisely the displayed maps. A full fibre is its Cartier scheme. Its reduced support is denoted $S$, its normalization by $\nu:T\to S$, and a resolution by $\pi:Z\to T$. Euler characteristic means topological Euler characteristic unless $\chi(\mathcal O)$ is written.

\section{Unrestricted Euler topology and component covers}\label{sec:II:06}


\begin{proposition}[Every torus-degeneration fibre has even Euler number]\label{cdp:st:C.2.UN.1}
\StatementRecord{cdp:st:C.2.UN.1}{proved}{Historical events: Claude Fable 5.2 (Anthropic), effort high, 16 November 2026: carried as background without fault found. GPT-6 Astra (OpenAI), ChatGPT Pro mode, effort Pro, S6\_09\_PAPERS, 17 November 2026: targeted proof check. Also checked in this writing. Pre-continuation writing events: GPT-6 Astra (OpenAI), ChatGPT Pro mode, effort Pro, S6\_10\_PAPERS, 18 November 2026: targeted proof check. Continuation mathematical and source checks: GPT-6 Astra (OpenAI), ChatGPT Pro mode, effort Pro, S6\_10\_PAPERS continuation, 19 November 2026: computation or formal certificate; targeted proof check. Continuation metadata events: GPT-6 Astra (OpenAI), ChatGPT Pro mode, effort Pro, S6\_10\_PAPERS continuation, 19 November 2026: used and rechecked in this run; status-vocabulary reconciliation; metadata reconciliation only; no new mathematical check. \emph{Pre-continuation recorded proof scope.} The pair sequence, skew pairing in degree three, exterior-square monodromy pairing, and even rank of G-I were re-derived}{the hypotheses, notation and definitions stated here; \cref{cdp:eq:C.UN.3}; \cref{cdp:eq:C.UN.4}}

Let $f:Y\to B$ be a proper holomorphic map from a smooth complex
threefold to a smooth complex curve, with connected fibres and smooth
general fibre a complex two-torus. For every $b\in B$,
\[
                    e((Y_b)_{\rm red})\in2\Z.       \tag{C.UN.1}\label{cdp:eq:C.UN.1}
\]
This assertion allows arbitrary multiplicities and singularities and
imposes no K\"ahler or Betti-number hypothesis on $Y$.
\end{proposition}


\begin{proof}
\noindent\textbf{Step 1.} Take a sufficiently small closed coordinate disc $\Delta$ about $b$,
containing no other critical value. Put $U=f^{-1}(\Delta)$ and
$M=\partial U$. Then $U$ is a compact oriented smooth real six-manifold
with boundary, and $M$ is the mapping torus of the orientation-preserving
monodromy of a real four-torus. The analytic central support $S$ is a
deformation retract of $U$. One precise construction uses a controlled
gradient for the real analytic function $|f|^2$ on a Whitney
stratification: it decreases $|f|$, is tangent to the strata over zero,
and has trajectories with a limit on the zero set by the local
analytic gradient inequality. Compactness of the zero set gives a
common small disc, and the resulting continuous limit homotopy fixes
the zero set. Equivalently the triangulated proper local tube is a
regular neighbourhood of the zero fibre. Nilpotents do not change this
retraction. Thus $e(S)=\chi(U)$.

\noindent\textbf{Step 2.} We give the parity calculation entirely over $\Q$. Write
$s_i=\operatorname{rank}(H_i(U)\to H_i(U,M))$.
Exactness of the pair sequence, with Poincar\'e--Lefschetz duality,
gives
\[
 b_0(M)+b_1(M)+b_2(M)
       =\sum_{i=0}^{5}b_i(U)-s_0-s_3-2s_1-2s_2.
                                                        \tag{C.UN.2}\label{cdp:eq:C.UN.2}
\]
Indeed, the summand $b_i(M)$ is the sum of the ranks of
$H_i(M)\to H_i(U)$ and $H_{i+1}(U,M)\to H_i(M)$;
substitution of $b_{i+1}(U,M)=b_{5-i}(U)$ gives the formula.
Here $s_0=0$ and $b_6(U)=0$. The map defining $s_3$ is represented,
after duality, by the middle-dimensional intersection pairing of $U$.
That pairing is skew-symmetric because the middle dimension is three,
so $s_3$ is even. Consequently
\[
 \chi(U)\equiv b_0(M)+b_1(M)+b_2(M)\pmod2.            \tag{C.UN.3}\label{cdp:eq:C.UN.3}
\]

\noindent\textbf{Step 3.} Let $T$ act on $H^1(F,\Q)$ for the four-torus fibre $F$, and set
$a_i=\dim\ker(\Lambda^iT-I)$. Cutting the mapping torus along one
fibre gives the Wang exact sequence; hence
\[
 b_0(M)=1,\quad b_1(M)=1+a_1,\quad b_2(M)=a_1+a_2.
                                                        \tag{C.UN.4}\label{cdp:eq:C.UN.4}
\]
Orientation preservation says $\det T=1$. The transformation
$G=\Lambda^2T$ preserves the nondegenerate symmetric wedge pairing
on $\Lambda^2\Q^4$ and has determinant $(\det T)^3=1$.
We verify that $\operatorname{rank}(G-I)$ is even without a
semisimplicity assumption. Complexify and split off the generalized
$(-1)$-eigenspace. Its dimension is even: all eigenvalues outside
$\{1,-1\}$ occur in reciprocal pairs with equal algebraic
multiplicity, and $\det G=1$. This generalized eigenspace is
nondegenerate and orthogonal to its invariant complement. On it
$G-I$ is invertible, of even rank. On the complement $G+I$ is
invertible, and
\[
                         (G-I)(G+I)^{-1}
\]
is skew-adjoint for the symmetric wedge pairing. Multiplication by
that nondegenerate pairing represents it as a skew-symmetric matrix,
whose rank is even. Its rank equals that of $G-I$ on the complement.
This proves the assertion. Since $\dim\Lambda^2\Q^4=6$, $a_2$ is
even. Equations \eqref{cdp:eq:C.UN.3}--\eqref{cdp:eq:C.UN.4} prove \eqref{cdp:eq:C.UN.1}.
\end{proof}






\begin{proposition}[All even Euler numbers and an unbounded zero-Euler family]\label{cdp:st:C.2.UN.2}
\StatementRecord{cdp:st:C.2.UN.2}{proved}{Historical events: Claude Fable 5.2 (Anthropic), effort high, 16 November 2026: carried as background without fault found. GPT-6 Astra (OpenAI), ChatGPT Pro mode, effort Pro, S6\_09\_PAPERS, 17 November 2026: targeted proof check. Also checked in this writing. Pre-continuation writing events: GPT-6 Astra (OpenAI), ChatGPT Pro mode, effort Pro, S6\_10\_PAPERS, 18 November 2026: targeted proof check. Continuation mathematical and source checks: GPT-6 Astra (OpenAI), ChatGPT Pro mode, effort Pro, S6\_10\_PAPERS continuation, 19 November 2026: computation or formal certificate; targeted proof check. Continuation metadata events: GPT-6 Astra (OpenAI), ChatGPT Pro mode, effort Pro, S6\_10\_PAPERS continuation, 19 November 2026: used and rechecked in this run; status-vocabulary reconciliation; metadata reconciliation only; no new mathematical check. \emph{Pre-continuation recorded proof scope.} The changes -2 for a genus-two curve blowup and +2 for a point blowup, and the complete even-integer realization, were re-derived using the stated compact genus-two curve input}{the hypotheses, notation and definitions stated here; \cref{cdp:st:C2b}; \cref{cdp:st:C.2.UN.1}}

The set of Euler numbers of reduced fibres among compact smooth
threefolds fibred in two-tori is exactly $2\Z$. For each $k\ge1$
there is a smooth projective such threefold with one singular reduced
normal-crossing fibre of Euler number zero consisting of a blown-up
abelian surface, $k$ ruled surfaces over a genus-two curve, and $k$
projective planes. Their incidence graph is a chain of the ruled
components attached to the abelian component, together with $k$
projective-plane leaves attached to that component.
\end{proposition}


\begin{proof}
\noindent\textbf{Step 1.} Use the explicit embedded genus-two curve $C_2\subset A$ of \ref{cdp:st:C2b}.
In $A\times\PP^1$ retain the smooth horizontal surface
$R=C_2\times\PP^1$. First blow up $C_2\times\{0\}$, as in \ref{cdp:st:C2b}.
The strict transform of $R$ is still $R$, because its centre is a
Cartier divisor in $R$. It meets the exceptional ruled surface in
the section corresponding to the base-normal direction; this
section is disjoint from the intersection with the strict abelian
component. The map to the base is a submersion near this section.

\noindent\textbf{Step 2.} Blow up that section and repeat, always using the intersection of
the strict transform of $R$ with the newest exceptional component.
The next centre is a smooth genus-two curve in the smooth locus of
the reduced central fibre, disjoint from its double curves. Its
normal sequence splits: the normal line to the curve inside $R$
maps isomorphically under $df$ to the tangent line of the base,
and supplies a splitting of the normal sequence of the central
surface. The two line summands give two disjoint sections on the
new exceptional projective bundle. One meets the previous
component; the other is the next centre. This proves by induction
that all $k$ curve blowups are defined and that the reduced central
support has the stated chain. In local coordinates the comparison
is always $t=u\tau$ as in \ref{cdp:st:C2b}, so every new component has
multiplicity one and no triple point is introduced.

\noindent\textbf{Step 3.} Finally blow up $m\ge0$ distinct points of the strict abelian
component away from the genus-two double curve. They introduce
$m$ projective-plane leaves, each meeting the blown-up abelian
component in one projective line. All maps are proper holomorphic
blowdowns followed by the original projection. Every nonzero
fibre is biholomorphic to $A$.

\noindent\textbf{Step 4.} Each curve blowup increases the Euler number of the central
support by $e(C_2)=-2$: the exceptional projective bundle has Euler
number $2e(C_2)$ and its intersection with the old component has
Euler number $e(C_2)$. Each point blowup increases it by two,
because the old surface gains one, the exceptional plane adds
three, and their intersection subtracts two. Thus
\[
                         e(S)=2(m-k).                 \tag{C.UN.5}\label{cdp:eq:C.UN.5}
\]
Given any integer $n$, take $k\ge\max(1,-n)$ and $m=k+n$.
This realizes $2n$; \ref{cdp:st:C.2.UN.1} supplies the opposite inclusion.
Taking $m=k$ gives the asserted zero-Euler family.

\noindent\textbf{Step 5.} The integral cohomology of the total space is torsion-free. The
excision and projective-bundle comparison of \ref{cdp:st:C2b} gives
\[
 (b_0,\ldots,b_6)
      =(1,4,7+k+m,8+4k,7+k+m,4,1).                  \tag{C.UN.6}\label{cdp:eq:C.UN.6}
\]
Its Euler characteristic is independently $2(m-k)$. The structure
sheaf direct images of each blowup are $\OO$ in degree zero and
zero in positive degrees, so
$(h^0(\OO),h^1(\OO),h^2(\OO),h^3(\OO))=(1,2,1,0)$.
The exact arithmetic and exterior-square checks are in
\path{code/root_unrestricted_euler.py} and its output.
This family supplies further members of the requested unrestricted
zero-Euler classification; it does not claim that every zero-Euler
fibre has this construction.
\end{proof}






\begin{proposition}[The same punctured torus family has non-bimeromorphic zero-Euler fillings]\label{cdp:st:C.2.UN.3}
\StatementRecord{cdp:st:C.2.UN.3}{proved}{Historical events: Claude Fable 5.2 (Anthropic), effort high, 16 November 2026: carried as background without fault found. GPT-6 Astra (OpenAI), ChatGPT Pro mode, effort Pro, S6\_09\_PAPERS, 17 November 2026: targeted proof check. Also checked in this writing. Pre-continuation writing events: GPT-6 Astra (OpenAI), ChatGPT Pro mode, effort Pro, S6\_10\_PAPERS, 18 November 2026: targeted proof check; used and rechecked in this run. Continuation mathematical and source checks: GPT-6 Astra (OpenAI), ChatGPT Pro mode, effort Pro, S6\_10\_PAPERS continuation, 19 November 2026: targeted proof check. Continuation metadata events: GPT-6 Astra (OpenAI), ChatGPT Pro mode, effort Pro, S6\_10\_PAPERS continuation, 19 November 2026: used and rechecked in this run; status-vocabulary reconciliation; metadata reconciliation only; no new mathematical check. \emph{Pre-continuation recorded proof scope.} The nonzero-line quotient, Kodaira and Hopf components, integral topology and finite-base-change equation were re-derived. AKMW2002 section1.2.4 p539, Lemma1.3.1 p540 and the analytic extension in section5.4 p564 were opened for the actual proper graph elimination; compactness permits shrinking its family germ after the finite local stages}{the hypotheses, notation and definitions stated here; \cref{cdp:st:D.2.7}; \cref{cdp:eq:C.UC.2}; \cite[Section~1.2.4 and Lemma~1.3.1]{AKMW2002}; \cref{cdp:eq:C.UC.1}; \cref{cdp:eq:C.UC.4}; \cref{cdp:eq:C.UC.3}; \cref{cdp:eq:C.UC.6}; \cref{cdp:eq:C.UC.8}; \cref{cdp:eq:C.UC.9}; \cref{cdp:eq:C.UC.10}; \cref{cdp:eq:C.UC.7}}

Let $A=E\times E$, $E=\C/(\Z+\mathrm i\Z)$.
There are two proper smooth-total-space fillings of
$A\times\Delta^*$ over a disc whose reduced central
supports both have Euler number zero, but for which
no bimeromorphic map over the disc exists. One is
$A\times\Delta$. The other is the local filling
(D.23), whose central support consists of a primary
Kodaira surface and a primary Hopf surface meeting
along one elliptic curve. The nonexistence persists
after every finite base change and resolution of the
normalized base change, with the given punctured
family retained.

These fillings occur in compact smooth threefolds:
the principal $A$-bundle $Y_0$ and the threefold $Y$
constructed in \ref{cdp:st:D.2.7} are isomorphic away from the
three specified fibres. That isomorphism does not
extend bimeromorphically, even after a finite base
change of their base curve. In particular smooth
torus fillings, together with modifications by
blowups and blowdowns over the central fibre, do
not generate all zero-Euler fillings of this
punctured family.

More strongly, the same punctured product has
infinitely many pairwise non-bimeromorphic
zero-Euler fillings. They are obtained by replacing
$\OO_V(R)$ in \eqref{cdp:eq:C.UC.2} by $\OO_V(nR)$,
$n\ge1$. This distinction also persists after
arbitrary finite base change and resolution.
\end{proposition}


\begin{proof}
\noindent\textbf{Step 1.} We first record the precise birational comparison
used in the proof. Let $\phi:M\dashrightarrow N$
be a bimeromorphic map between smooth complex
threefolds, and let $P\subset M$ be a compact
irreducible divisor whose smooth resolution is
not covered by rational curves. Then $\phi$
cannot contract $P$ to a subset of dimension
at most one. Resolve the inverse meromorphic
map by smooth-centre blowups:
\[
 \widehat N\xrightarrow{\ a\ }N,
 \qquad b:\widehat N\longrightarrow M,
 \qquad b=\phi^{-1}a.
\tag{C.UC.1}\label{cdp:eq:C.UC.1}
\]
The elimination-of-indeterminacy construction
has this form in the analytic category as
well as the algebraic one; see
\cite[Section~1.2.4 and Lemma~1.3.1]{AKMW2002}.
For maps of proper disc models, the same
construction is made on a smaller disc; its
centres above the compact central fibre are
compact. For the compact threefolds in the
statement the cited construction applies
directly.

\noindent\textbf{Step 2.} Since $b$ is a proper modification of the
smooth target $M$, the strict transform
$\widehat P$ maps birationally onto $P$.
Here is the analytic justification, without
using a global function field. The locus in
$M$ over which $b$ has positive-dimensional
fibres is analytic by properness. It cannot
contain a divisor: its preimage would then
have dimension at least three while being
a proper analytic subset of the irreducible
threefold $\widehat N$. Thus $b$ is finite
over a dense open subset of $P$. A finite
modification of the normal manifold $M$
is an isomorphism, so its degree there is
one. If $\phi$ contracted
$P$, then $a$ would contract $\widehat P$.
It would be an exceptional divisor of the
composition of smooth-centre blowups $a$.
Every such divisor is birational to the
exceptional divisor created at one of those
steps. In dimension three that divisor is
$\PP^2$ for a point centre, or a
$\PP^1$-bundle over a smooth curve for a
curve centre. A divisor centre gives an
isomorphism and creates no exceptional
divisor. Hence $\widehat P$, and therefore
$P$, would be covered by rational curves,
a contradiction. Applying the same argument
to the inverse shows that the surviving
divisors correspond birationally, rather
than only by a generically finite map.

\noindent\textbf{Step 3.} We now give the actual local models. Put
$q=e^{-2\pi}$ and
\[
 V=\operatorname{Bl}_{(0,0)}(E\times\Delta),
 \qquad R=\operatorname{Exc}(V\to E\times\Delta),
 \qquad U=\OO_V(R)^\times/q^{\Z}.
\tag{C.UC.2}\label{cdp:eq:C.UC.2}
\]
The quotient map $U\to V$ is a principal
elliptic bundle with fibre $\C^*/q^{\Z}=E$.
It is smooth and proper, since each bundle
chart is the product with this compact
elliptic curve. Composition with $V\to\Delta$
is therefore a proper map from a smooth
threefold. The canonical section of
$\OO_V(R)$ trivializes it on the punctured
disc, giving the specified isomorphism
\[
                         U|_{\Delta^*}=A\times\Delta^*.
\tag{C.UC.3}\label{cdp:eq:C.UC.3}
\]
This is exactly the construction in \ref{cdp:st:D.2.7},
including its holomorphic gluing to the
compact threefold.

\noindent\textbf{Step 4.} The strict transform $E_0$ of $E\times\{0\}$
and the exceptional $R$ meet transversely at
one point. The restrictions of $\OO_V(R)$
to them have degree $+1$ and $-1$,
respectively. Their inverse images in $U$
are
\[
 P=\OO_E(0)^\times/q^{\Z},\qquad
 Q=\OO_{\PP^1}(-1)^\times/q^{\Z}
       =(\C^2\setminus\{0\})/\langle q\rangle.
\tag{C.UC.4}\label{cdp:eq:C.UC.4}
\]
They meet in a single elliptic fibre $J$,
and the local equation of the map is
$xy=t$ along $J$. The first surface is a
primary Kodaira surface with $b_1(P)=3$:
in its principal elliptic bundle the primitive
Chern class kills one of the two elliptic
fibre cycles. The second is the indicated
primary Hopf surface, with $b_1(Q)=1$.
These identifications and the integral
transgression are also calculated in
(D.23)--(D.25).

\noindent\textbf{Step 5.} Neither surface contains a rational curve.
A holomorphic map $\PP^1\to P$ projects
constantly to its elliptic base and then
constantly to its elliptic fibre. A map
$\PP^1\to Q$ lifts to $\C^2\setminus\{0\}$,
since $\PP^1$ is simply connected; its two
holomorphic coordinate functions are constant.
Thus their resolutions are not covered by
rational curves. A bimeromorphic comparison
from $U$ to $A\times\Delta$ over $\Delta$
would, by \eqref{cdp:eq:C.UC.1}, have to carry each of
$P,Q$ birationally onto a central divisor.
The product has only the one central
divisor $A$, and neither $P$ nor $Q$ is
birational to it. Indeed the first Betti
number of a smooth compact surface is
unchanged by point blowups: a ball is
replaced by the simply connected exceptional
neighbourhood, and van Kampen gives the
same fundamental group. Resolution of a
surface bimeromorphism therefore preserves
$b_1$, whereas $b_1(A)=4$ differs from
both values in \eqref{cdp:eq:C.UC.4}. This proves that
there is no such bimeromorphic comparison.

\noindent\textbf{Step 6.} All these supports have Euler number zero.
For $P$ and $Q$, their principal elliptic
bundle projections are locally trivial
proper bundles, so their Euler numbers are
the product of the Euler numbers of base
and fibre, hence zero. The intersection
$J$ also has Euler number zero. Additivity
gives
\[
                    e(P\cup Q)=e(P)+e(Q)-e(J)=0=e(A).
\tag{C.UC.5}\label{cdp:eq:C.UC.5}
\]

\noindent\textbf{Step 7.} Here is the comparison after a finite base
change, with its ramification retained.
Locally such a change has equation
$t=u(s)s^m$, $u(0)\ne0$. Taking a local
$m$th root of the unit replaces its
parameter by one with $t=s^m$.
At a general point of either $P$ or $Q$
the original map is a submersion with
coordinate $x=t$. Its base change is the
smooth chart
\[
                          x=s^m,
\tag{C.UC.6}\label{cdp:eq:C.UC.6}
\]
with the two surface coordinates unchanged.
The component over $s=0$ thus maps with
degree one to the original component.
At their intersection the normalized base
change has equation $xy=s^m$, with one
additional smooth elliptic coordinate.
This is normal, with transverse rational
double point $A_{m-1}$ when $m>1$.
A resolution adds ruled exceptional
surfaces above $J$; the strict transforms
of $P$ and $Q$ remain birational to the
surfaces in \eqref{cdp:eq:C.UC.4}. No choice of further
resolution changes those birational
classes. The target base change remains
the smooth product $A\times\Delta_s$.
The argument using \eqref{cdp:eq:C.UC.1} and their
first Betti numbers again excludes every
bimeromorphic comparison over $\Delta_s$.

\noindent\textbf{Step 8.} Finally use the compact gluing of \ref{cdp:st:D.2.7}.
The original principal $A$-bundle $Y_0$
has a smooth abelian central fibre at
each of the three chosen base points.
The replacement $Y$ has the local model
\eqref{cdp:eq:C.UC.2} at each point and is identified
with $Y_0$ elsewhere by \eqref{cdp:eq:C.UC.3}.
An extension of this isomorphism would
respect the base, since the two base
maps agree on its dense open graph.
Its restriction near any of the three
fibres would be the already excluded
comparison. The same argument applies
above every preimage of these points
under a finite cover of the compact base,
using \eqref{cdp:eq:C.UC.6}. This proves the compact
and base-change assertions.

\noindent\textbf{Step 9.} For the infinite family put
\[
 U_n=\OO_V(nR)^\times/q^{\Z},\qquad
 P_n=\OO_E(n[0])^\times/q^{\Z},\qquad
 Q_n=\OO_{\PP^1}(-n)^\times/q^{\Z}.
\tag{C.UC.7}\label{cdp:eq:C.UC.7}
\]
The canonical section again trivializes the
bundle off the central fibre. Thus the punctured
isomorphism is the same exact product \eqref{cdp:eq:C.UC.3}.
The central support is $P_n\cup_J Q_n$;
all three terms have Euler number zero by their
principal elliptic bundle projections. Each
$P_n$ contains no rational curve by its
projection to $E$ and then to its elliptic
fibre. For $Q_n$, the usual quotient
$\OO_{\PP^1}(-n)^\times
=(\C^2\setminus\{0\})/\boldsymbol\mu_n$
is explicit: a nonzero vector $v$ represents
$v^{\otimes n}$ in the $n$th tensor power of
the tautological line. Multiplication by $q$
on that tensor power lifts to multiplication
by the positive real number $q^{1/n}$ on
$\C^2\setminus\{0\}$. Hence
\[
 Q_n=(\C^2\setminus\{0\})/
       \langle\boldsymbol\mu_n,q^{1/n}\rangle,
 \qquad
 \pi_1(Q_n)=\Z\oplus\Z/n\Z.
\tag{C.UC.8}\label{cdp:eq:C.UC.8}
\]
The two commuting generators are independent,
because a nonzero power of $q^{1/n}$ has
absolute value different from one. The action
is free and properly discontinuous, and its
covering space retracts to the simply connected
$S^3$. This proves the fundamental-group
calculation and also proves that every map
$\PP^1\to Q_n$ is constant by lifting it to
the cover.

\noindent\textbf{Step 10.} For completeness the corresponding calculation
for $P_n$ retains the extension class. A Hermitian
metric on $\OO_E(n[0])$ identifies its nonzero
total space smoothly with its unit-circle bundle
times $\R_{>0}$. As $q$ is positive real, its
quotient is the unit-circle bundle times
$\R_{>0}/q^{\Z}\simeq S^1$. The circle bundle
has Chern number $n$. Lifting the two generators
of $\pi_1(E)$, the attaching map of the base
two-cell has winding number $n$ in the circle
fibre. With a consistent orientation this gives
the exact presentation
\[
 \pi_1(P_n)=
 \langle a,b,c,d\mid
 [a,b]=c^n,\ [a,c]=[b,c]=1,
 [a,d]=[b,d]=[c,d]=1\rangle.
\tag{C.UC.9}\label{cdp:eq:C.UC.9}
\]
There are no further relations: the homotopy
sequence of the circle bundle has injective
fibre group because $\pi_2(E)=0$, and the last
factor is the displayed direct $S^1$ factor.
Abelianizing \eqref{cdp:eq:C.UC.8}--\eqref{cdp:eq:C.UC.9} yields
\[
 H_1(P_n,\Z)=\Z^3\oplus\Z/n\Z,
 \qquad H_1(Q_n,\Z)=\Z\oplus\Z/n\Z.
\tag{C.UC.10}\label{cdp:eq:C.UC.10}
\]
Integral first homology, including its torsion,
is invariant under a bimeromorphism of smooth
compact surfaces, by the same point-blowup
comparison already used above. By \eqref{cdp:eq:C.UC.1},
a bimeromorphism $U_n\dashrightarrow U_{n'}$
over the disc must carry the two non-uniruled
central divisors birationally to the two such
divisors on the target. The ranks in \eqref{cdp:eq:C.UC.10}
force $P_n$ to correspond to $P_{n'}$ and
$Q_n$ to $Q_{n'}$. The torsion orders then
force $n=n'$. After any finite base change,
\eqref{cdp:eq:C.UC.6} still leaves these components
birationally unchanged, and the resolution of
$xy=s^m$ adds only ruled surfaces over $J$.
Further resolutions add only uniruled
exceptional divisors. They cannot replace a
non-uniruled component under \eqref{cdp:eq:C.UC.1}, so
the same torsion argument proves persistence.
Finally all $U_n$ have compact realizations:
replace $V$ by
$\operatorname{Bl}_{(0,0)}(E\times\PP^1)$
and use the same line bundle $\OO(nR)$.
Its elliptic quotient is a compact smooth
threefold, and the projection to $\PP^1$
restricts to \eqref{cdp:eq:C.UC.7} near zero.
\end{proof}



The comparison obstruction specifies the limit
of a particular proposed classification route.
In \cite[Sections~5.1--5.3]{KKN2008}, a
nondegenerate logarithmic torus pairing is
given first, and a suitable complete fan
produces a proper model of that logarithmic
complex torus. The theorem there supplies
those models; an extension of an arbitrary
punctured-family isomorphism to a
bimeromorphic map of its chosen fillings
is an additional requirement. Statement
\ref{cdp:st:C.2.UN.3} refutes that additional requirement
for the displayed fillings, even though
the punctured family is a product.




\begin{proposition}[Every even Euler number occurs above each fixed torus filling]\label{cdp:st:C.2.UN.4}
\StatementRecord{cdp:st:C.2.UN.4}{proved}{Historical events: Claude Fable 5.2 (Anthropic), effort high, 16 November 2026: carried as background without fault found. GPT-6 Astra (OpenAI), ChatGPT Pro mode, effort Pro, S6\_09\_PAPERS, 17 November 2026: targeted proof check. Also checked in this writing. Pre-continuation writing events: GPT-6 Astra (OpenAI), ChatGPT Pro mode, effort Pro, S6\_10\_PAPERS, 18 November 2026: targeted proof check. Continuation mathematical and source checks: GPT-6 Astra (OpenAI), ChatGPT Pro mode, effort Pro, S6\_10\_PAPERS continuation, 19 November 2026: computation or formal certificate; targeted proof check. Continuation metadata events: GPT-6 Astra (OpenAI), ChatGPT Pro mode, effort Pro, S6\_10\_PAPERS continuation, 19 November 2026: used and rechecked in this run; status-vocabulary reconciliation; metadata reconciliation only; no new mathematical check. \emph{Pre-continuation recorded proof scope.} The point and smooth plane-quartic blowups change the support Euler number by plus two and minus two, with full multiplicities tracked in actual charts. The integral projective-bundle and relative-pair map gives the stated cohomology summands; the net-zero cycle adds the asserted components and Betti numbers}{the hypotheses, notation and definitions stated here; \cref{cdp:eq:C.VR.4}; \cref{cdp:eq:C.VR.6}; \cref{cdp:eq:C.VR.8}; \cref{cdp:eq:C.VR.9}; \cref{cdp:st:C.2.UN.1}}

Let $f:Y\to B$ be a proper holomorphic map from a smooth complex
threefold to a smooth complex curve, with connected fibres and smooth
general fibre a complex two-torus. Fix $b\in B$ and put
$S=(Y_b)_{\rm red}$. For every integer $k$ there is a finite composition
of blowups with smooth compact centres in successive fibres over $b$,
\[
 \pi:\widetilde Y\longrightarrow Y,\qquad
 \widetilde f=f\pi,\qquad
 e((\widetilde Y_b)_{\rm red})=2k.                 \tag{C.VR.1}\label{cdp:eq:C.VR.1}
\]
The map $\pi$ is an isomorphism over $B\setminus\{b\}$.
In particular every such filling is dominated by a smooth filling
whose central reduced support has Euler number zero. More precisely,
every further smooth model obtained by vertical blowups admits a
zero-Euler refinement by additional vertical blowups. Thus zero-Euler
models are cofinal for this domination relation, in the explicit sense
that every object is dominated by one of them.
Above every given filling there are such zero-Euler refinements with
arbitrarily many irreducible components of their central support.

An elementary operation lowering the Euler number by two consists of
one point blowup and one smooth genus-three curve blowup. On integral
cohomology its additional summands are $\Z^2$ in degrees two and four
and $\Z^6$ in degree three. In particular, whenever the Betti numbers
are finite,
\[
 (\Delta b_2,\Delta b_3,\Delta b_4)=(2,6,2),
 \qquad \Delta b_j=0\quad(j\notin\{2,3,4\}).       \tag{C.VR.2}\label{cdp:eq:C.VR.2}
\]
All the original positive full-fibre multiplicities are retained on
their strict transforms. The construction requires no normal crossings
condition on the original support.
\end{proposition}


\begin{proof}
\noindent\textbf{Step 1.} We first specify the comparison for a blowup whose centre meets
singularities of the support. Let $Z\subset S$ be a smooth compact
centre of complex codimension $r$ in the smooth ambient threefold,
and let $\pi_Z:Y'\to Y$ be its blowup. Write $S'=(f\pi_Z)^{-1}(b)_{\rm red}$.
The equality of underlying sets is $S'=\pi_Z^{-1}(S)$, and the actual
restrictions of the blowup give
\[
 S'\setminus\pi_Z^{-1}(Z)\simeq S\setminus Z,
 \qquad \pi_Z^{-1}(Z)=\PP(N_{Z/Y}).                \tag{C.VR.3}\label{cdp:eq:C.VR.3}
\]
Compact complex analytic spaces admit finite triangulations, compatible
with these closed analytic subsets. Euler additivity, applied with
compact supports on the complements, and the projective-bundle
formula therefore give
\[
 \begin{split}
 e(S')-e(S)
 &=e(\PP(N_{Z/Y}))-e(Z)\\
 &=(r-1)e(Z).
 \end{split}                                      \tag{C.VR.4}\label{cdp:eq:C.VR.4}
\]
The projective-bundle Euler equality also follows from its locally
trivial fibre $\PP^{r-1}$, whose Euler number is $r$.
Thus \eqref{cdp:eq:C.VR.4} retains any intersections of the centre with other
components or with the singular support. It makes no replacement of
the full inverse image by a strict transform.

\noindent\textbf{Step 2.} Choose a smooth point $p$ of one component of $S$, away from every
other component. Such points form a nonempty open subset of each
irreducible component. If the full fibre has multiplicity $m>0$ on
that component, local holomorphic coordinates may be chosen with
\[
                    t\circ f=x^m.                \tag{C.VR.5}\label{cdp:eq:C.VR.5}
\]
Indeed the local equation initially has the form $x^m a$, where $a$
is a holomorphic unit; a local holomorphic $m$th root of $a$ replaces
$x$ by $x a^{1/m}$. Blow up $p$. The exceptional surface
$E\simeq\PP^2$ meets the strict component in the line
$L=\{X=0\}\subset\PP(T_pY)$, with homogeneous coordinates
$[X:Y:Z]$ induced by $(x,y,z)$. It meets no other central component.
In the chart $y=r$, $x=ru$, $z=rv$, the actual map is
\[
                         t=(ru)^m.                \tag{C.VR.6}\label{cdp:eq:C.VR.6}
\]
Consequently both $E=(r=0)$ and the strict component $(u=0)$
have multiplicity $m$. Formula \eqref{cdp:eq:C.VR.4}, with $r=3$ and centre a
point, shows that this step increases the central Euler number by two.

\noindent\textbf{Step 3.} Inside $E$ take the Fermat quartic
\[
 C=\{X^4+Y^4+Z^4=0\}\subset\PP^2.                \tag{C.VR.7}\label{cdp:eq:C.VR.7}
\]
Its three partial derivatives cannot vanish simultaneously at a
projective point, so it is smooth. It is connected: if the quartic
were a disjoint union of distinct plane components, two such
components would meet by their positive intersection number; their
intersection would be a singular point. Adjunction gives
$K_C=\OO_C(1)$ and $2g(C)-2=4$, hence $g(C)=3$ and $e(C)=-4$.
On $L$ its equation is $Y^4+Z^4=0$. It has four distinct roots,
none at $Z=0$, and the two gradients of the quartic and $X$ are
independent there. Thus $C$ meets $L$ transversely in exactly four
points.

\noindent\textbf{Step 4.} Now blow up the smooth curve $C$ in the threefold obtained after the
first blowup. Its exceptional divisor is the projective line bundle
$\PP(N_{C/Y_1})\to C$. Equation \eqref{cdp:eq:C.VR.4}, with $r=2$, gives Euler
change $e(C)=-4$. The two successive changes sum to $-2$.
For clarity, the four points of $C\cap L$ introduce no singularities
of the ambient space and do not alter this Euler comparison. Near any
one of them, preserve the local equations $r=0$ of $E$ and $u=0$ of
the strict old component. Transversality gives a third coordinate $v$
such that $C=(r=v=0)$. Absorbing a remaining unit in $r$ preserves
these equations and gives $t=(ru)^m$ as in \eqref{cdp:eq:C.VR.6}. The two charts of
the curve blowup are
\[
 \begin{array}{c|c|c}
 \text{chart}&\text{blowdown substitution}&\text{base map}\\ \hline
 \text{$r$ chart}&v=rw&t=(ru)^m\\
 \text{$v$ chart}&r=vw&t=(uvw)^m .
 \end{array}                                      \tag{C.VR.8}\label{cdp:eq:C.VR.8}
\]
In the second chart the new exceptional divisor, the strict transform
of $E$, and the strict old component have equations $v=0$, $w=0$,
and $u=0$, respectively. They meet with normal crossings and all
have multiplicity $m$. Away from $C\cap L$ the comparison is the
same two-coordinate blowup of $t=r^m$. The strict transform of $E$
is isomorphic to $E$, since $C$ is Cartier in $E$; the strict old
component is its point blowup at the four transverse intersection
points. These identifications follow directly from the charts in
\eqref{cdp:eq:C.VR.8}. Outside the specified neighbourhoods the maps are
isomorphisms, so any other singularities of the original support
are simply retained.

\noindent\textbf{Step 5.} We include the integral cohomology comparison giving \eqref{cdp:eq:C.VR.2}.
For a smooth centre of codimension $r$, let $j:E_Z\hookrightarrow Y'$
and $p:E_Z\to Z$ denote the exceptional inclusion and projection,
and put $\xi=c_1(\OO_{E_Z}(1))$. The actual blowup homomorphism is
\[
 \begin{split}
 H^q(Y,\Z)\oplus\bigoplus_{i=1}^{r-1}H^{q-2i}(Z,\Z)
   &\longrightarrow H^q(Y',\Z),\\
 (\alpha,\beta_1,\ldots,\beta_{r-1})
   &\longmapsto \pi_Z^*\alpha+
       \sum_{i=1}^{r-1}j_*(\xi^{i-1}p^*\beta_i).
 \end{split}                                      \tag{C.VR.9}\label{cdp:eq:C.VR.9}
\]
It is an isomorphism. To check it integrally, excision identifies the
relative cohomology of $(Y',E_Z)$ with that of $(Y,Z)$. The
projective-bundle theorem makes $H^q(E_Z,\Z)$ the direct sum of
$\xi^i p^*H^{q-2i}(Z,\Z)$ for $0\le i<r$, and the self-intersection
formula gives $j^*j_*\eta=-\xi\eta$. Thus the extra terms of
\eqref{cdp:eq:C.VR.9} map isomorphically, with sign minus, onto the summands with
$i>0$. Comparing the relative long exact sequences proves
surjectivity; their remaining $i=0$ term is exactly the pullback from
$Y$. The proper degree-one trace makes $\pi_Z^*$ injective, and
restriction to $E_Z$ then proves independence of all the displayed
extra summands. This proves \eqref{cdp:eq:C.VR.9}. The same argument applies on a
compact fibre tube with boundary, since the centre is in its interior
and the blowup is unchanged on the boundary.

\noindent\textbf{Step 6.} The point blowup adds $\Z$ in each of degrees two and four. The
genus-three curve blowup adds $H^{q-2}(C,\Z)$, whose ranks in these
three successive degrees are $1,6,1$. This proves \eqref{cdp:eq:C.VR.2} and shows
that pre-existing integral torsion is unchanged.

\noindent\textbf{Step 7.} Both steps are proper holomorphic maps from smooth threefolds. A
blowup has connected fibres and pushes its structure sheaf to the
structure sheaf of the original smooth space. Stein factorization
therefore proves that their composition with $f$ again has connected
fibres. The centres lie over $b$, so the resulting map is the same
proper torus family over $B\setminus\{b\}$.

\noindent\textbf{Step 8.} Finally \ref{cdp:st:C.2.UN.1} writes $e(S)=2n$ for an integer $n$. To obtain $2k$
with $k\ge n$, use $k-n$ point blowups. To obtain $2k$ with $k<n$,
use $n-k$ of the two-step operations just proved. At each stage a
fresh smooth point of a single component exists, and the construction
can be performed there; only finitely many centres are needed.
Applying this same procedure to any given vertical blowup model
proves the asserted cofinality with its specified domination maps.
After reaching Euler number zero, compose one further point blowup
with one of the two-step decreasing operations, at fresh smooth
points. Their net Euler change is zero, while they add three
irreducible central components. Repetition gives arbitrarily many
components above the original filling. Each repetition adds
$(3,6,3)$ to $(b_2,b_3,b_4)$ when those numbers are finite; the
actual maps remain the successive blowdowns already described.
The exact chart, quartic, Euler, and cohomology calculations are
retained in \path{code/E_vertical_refinements.py} and its output.
\end{proof}






\begin{proposition}[The same non-uniruled components occur with every even Euler number]\label{cdp:st:C.2.UN.14}
\StatementRecord{cdp:st:C.2.UN.14}{proved}{Historical events: Claude Fable 5.2 (Anthropic), effort high, 16 November 2026: carried as background without fault found. GPT-6 Astra (OpenAI), ChatGPT Pro mode, effort Pro, S6\_09\_PAPERS, 17 November 2026: targeted proof check. Also checked in this writing. Pre-continuation writing events: GPT-6 Astra (OpenAI), ChatGPT Pro mode, effort Pro, S6\_10\_PAPERS, 18 November 2026: targeted proof check; used and rechecked in this run. Continuation mathematical and source checks: GPT-6 Astra (OpenAI), ChatGPT Pro mode, effort Pro, S6\_10\_PAPERS continuation, 19 November 2026: independent re-proof. Continuation metadata events: GPT-6 Astra (OpenAI), ChatGPT Pro mode, effort Pro, S6\_10\_PAPERS continuation, 19 November 2026: used and rechecked in this run; status-vocabulary reconciliation; metadata reconciliation only; no new mathematical check. \emph{Pre-continuation recorded proof scope.} The normalization of the actual blowup graph preserves every strict nonuniruled component and its generic multiplicity. The exceptional projective plane or ruled surface is explicitly uniruled. Applying the already checked local Euler operations gives all stated even values, independently at each of the finitely many critical fibres}{the hypotheses, notation and definitions stated here; \cref{cdp:st:C.2.UN.4}; \cref{cdp:eq:C.EC.2}}

Fix a proper holomorphic family $p^\times:\mathcal T^\times\to\Delta^\times$
of connected complex two-tori. A marked smooth filling is a proper
holomorphic map $f:Y\to\Delta$ from a smooth complex threefold with
connected fibres, together with a biholomorphism
$f^{-1}(\Delta^\times)\simeq\mathcal T^\times$ over $\Delta^\times$.
For each irreducible component $S_i$ of $(Y_0)_{\rm red}$, let
$Z_i\to S_i^\nu\to S_i$ be a smooth compact resolution of its
normalization. Write $\mathcal E(f)$ for the finite multiset of
bimeromorphism classes $[Z_i]$ for which $Z_i$ is not uniruled.
Here a smooth compact surface is uniruled when it admits a dominant
meromorphic map $C\times\PP^1\dashrightarrow Z$ for a smooth compact
curve $C$. Repetitions in the multiset record distinct components.

For every integer $k$, the following equality holds between sets of
finite multisets:
\[
 \begin{split}
 \{\mathcal E(f): f\text{ is a marked smooth filling of }p^\times\}
 ={}&\{\mathcal E(f): f\text{ is such a filling},\\[-2pt]
 &\hspace{48mm}e((Y_0)_{\rm red})=2k\}.
 \end{split}                                                     \tag{C.EC.1}\label{cdp:eq:C.EC.1}
\]
For each filling in the left-hand set, the representative on the
right is obtained by the explicit proper holomorphic vertical
blowup map of \ref{cdp:st:C.2.UN.4}. That map induces a bijection between the
original non-uniruled components and the resulting non-uniruled
components, and a bimeromorphism of their normalization resolutions.
Every original full-fibre multiplicity is preserved by this bijection.

The same assertion holds for compact smooth total spaces over a
fixed compact base curve, with the marked family fixed away from a
finite set $Q$. One can prescribe an arbitrary even Euler number at
each point of $Q$ simultaneously, preserving the multiset and
multiplicity at each point. In particular, the collection of
non-uniruled surface classes which occur in arbitrary compact
torus-fibred fillings is exactly the collection which occur in
compact fillings all of whose reduced fibres have Euler number zero.
\end{proposition}


\begin{proof}
\noindent\textbf{Step 1.} We specify the component comparison for one of the smooth blowups
$\pi:Y'\to Y$ used in \ref{cdp:st:C.2.UN.4}. Its centre $C$ is a point or a
smooth compact curve in the central support. The inverse image of
that support consists of the strict transforms $S_i'$ of its
original components together with the exceptional divisor
\[
 E=\PP(N_{C/Y})\longrightarrow C.                    \tag{C.EC.2}\label{cdp:eq:C.EC.2}
\]
Each restriction $S_i'\to S_i$ is a proper holomorphic map and is an
isomorphism over the complement of $C$. Since $C$ has smaller
dimension than $S_i$, this complement is dense. Thus this actual
map is bimeromorphic. Normalizing both surfaces and resolving them
gives a bimeromorphism $Z_i'\dashrightarrow Z_i$: explicitly, take
the closure of the graph of the isomorphism on the dense normal
open sets, normalize it, and resolve the resulting compact analytic
surface. Its two proper holomorphic projections resolve the stated
meromorphic map. The pullback of the defining equation of the full
fibre has its original vanishing order at the generic point of
$S_i'$, because $\pi$ is an isomorphism there. Hence its
multiplicity is unchanged.

\noindent\textbf{Step 2.} Uniruledness is preserved by these bimeromorphisms. Indeed, composing
a dominant meromorphic map $C_0\times\PP^1\dashrightarrow Z_i$
with either bimeromorphic inverse gives such a map for $Z_i'$; its
dominance follows on the dense open set where that inverse is an
isomorphism. Resolving the graph makes the composition a
well-defined meromorphic map. The same argument in the reverse
direction proves the equivalence.

\noindent\textbf{Step 3.} The newly introduced divisor in \eqref{cdp:eq:C.EC.2} is uniruled. For a point
centre it is $\PP^2$, which is rational: the map
$\PP^1\times\PP^1\dashrightarrow\PP^2$ given on the affine chart
by $(u,v)\mapsto[1:u:v]$ is bimeromorphic. For a curve centre it is
a projective line bundle over $C$. Its rank-two bundle is
meromorphically trivial on $C$: tensoring by a sufficiently positive
line bundle supplies two global sections independent at the
generic point, and their ratio gives the meromorphic
trivialization. Projectivizing it gives a bimeromorphism
$C\times\PP^1\dashrightarrow\PP(N_{C/Y})$.
Its subsequent strict transforms remain uniruled by the same
composition argument. Therefore a finite sequence of these blowups
introduces no new non-uniruled class and loses none of the original
classes. It preserves their repetitions and multiplicities as well.

\noindent\textbf{Step 4.} Starting from any marked filling, \ref{cdp:st:C.2.UN.4} constructs a finite
sequence of exactly these maps whose final reduced Euler number is
$2k$. Every centre lies over zero, so the blowdown is a
biholomorphism over $\Delta^\times$ and transports the given marking
by composition. We have just proved that its multiset is the
original $\mathcal E(f)$. This gives inclusion of the left side of
\eqref{cdp:eq:C.EC.1} in the right. Inclusion in the other direction follows
because each filling with Euler number $2k$ is among the same
marked fillings defining the left side. This also proves the
assertion when there are no fillings: both sets are empty.

\noindent\textbf{Step 5.} For the compact version, apply the same construction successively
at the finitely many points of $Q$. Every centre is a compact
analytic submanifold of the current compact threefold. Its global
blowup is compact and smooth, and the composite with the fixed base
map is proper holomorphic with connected fibres, by the
structure-sheaf and Stein-factor comparison in \ref{cdp:st:C.2.UN.4}. Its
restriction off that point is an isomorphism. Thus a construction
at one point preserves every fibre and every marking already
treated at another point. There are only finitely many critical
values of a compact proper map to a curve which is smooth
generically: the image of its compact analytic critical locus is
a proper analytic subset of the curve. Taking $Q$ to contain them
all and taking $k=0$ at each of them proves the final assertion;
the remaining smooth torus fibres already have Euler number zero.
\end{proof}



For completeness, a surface dominated meromorphically by
$C\times\PP^1$ has no nonzero pluricanonical section. The pullback
of such a section is nonzero, because a dominant meromorphic map
between smooth surfaces is generically of full rank. It is
holomorphic away from the finitely many indeterminacy points and
extends over them by Hartogs' theorem. Its restriction to any
$\{c\}\times\PP^1$ has degree $-2m$ in pluricanonical degree $m>0$
and must vanish. This contradicts its being nonzero. Thus the
equality applies in particular to a specified general-type
normalization class. It makes the occurrence problem for such a
component part of the unrestricted zero-Euler problem through the
explicit maps just proved. It does not supply a comparison between
two independently specified marked fillings. The construction
uses \ref{cdp:st:C.2.UN.4} only and invokes no assertion that such a comparison
extends across the central fibre.




\begin{proposition}[No component of a torus degeneration has general-type normalization]\label{cdp:st:C.2.UN.15}
\StatementRecord{cdp:st:C.2.UN.15}{proved}{Historical events: Claude Fable 5.2 (Anthropic), effort high, 16 November 2026: carried as background without fault found. GPT-6 Astra (OpenAI), ChatGPT Pro mode, effort Pro, S6\_09\_PAPERS, 17 November 2026: targeted proof check. Also checked in this writing. Pre-continuation writing events: GPT-6 Astra (OpenAI), ChatGPT Pro mode, effort Pro, S6\_10\_PAPERS, 18 November 2026: targeted proof check; used and rechecked in this run. Continuation mathematical and source checks: GPT-6 Astra (OpenAI), ChatGPT Pro mode, effort Pro, S6\_10\_PAPERS continuation, 19 November 2026: targeted proof check; used and rechecked in this run. Continuation metadata events: GPT-6 Astra (OpenAI), ChatGPT Pro mode, effort Pro, S6\_10\_PAPERS continuation, 19 November 2026: used and rechecked in this run; status-vocabulary reconciliation; metadata reconciliation only; no new mathematical check. \emph{Pre-continuation recorded proof scope.} The actual root base change makes the chosen covering component multiplicity one. Rescaling the coherent vertical vector field gives a nonzero meromorphic restriction, and Guillot2024 Theorem3.3 p9 supplies the exact closed-univalence comparison. The dense-leaf branch matches Guillot2019 Main Theorem p2 and Proposition22 p25; the algebraic-leaf branch constructs its Hilbert-incidence first integral and excludes genus at least two by the time-form pole calculation}{the hypotheses, notation and definitions stated here; \cref{cdp:eq:C.SC.3}; \cite[Theorem~3.3, with its proof]{Guillot2024}; \cite{Guillot2019}; \cref{cdp:eq:C.SC.6}}

Let $f:U\to\Delta$ be a proper holomorphic map from
a smooth complex threefold to a disc, with connected
fibres and smooth fibres complex two-tori over
$\Delta^*$. For every irreducible component $S_i$
of the reduced central support, its smooth
normalization resolution $Z_i$ satisfies
\[
                            \kappa(Z_i)\le1.
 \tag{C.SC.1}\label{cdp:eq:C.SC.1}
\]
No K\"ahler, projectivity, Betti-number or Euler-number
assumption is imposed on $U$ or its central fibre.
\end{proposition}


\begin{proof}
\noindent\textbf{Step 1.} We first make the chosen component reduced by an
actual finite base change. If its full-fibre
multiplicity is $m_i$, put $t=s^{m_i}$ and normalize
the fibre product. At a generic smooth point of
$S_i$, away from the other components, the original
equation is $t=u z^{m_i}$ with $u$ a unit. The
normalized base change has the explicit local maps
\[
                  w^{m_i}=u,\qquad z=s/w.
 \tag{C.SC.2}\label{cdp:eq:C.SC.2}
\]
It is therefore smooth at those points, and every
central component above $S_i$ has multiplicity one
and maps generically finitely onto $S_i$. Choose one
such component and resolve the normalized total
space and central support, isomorphically over
$\Delta^*$ and generically along that component.
Write $f':U'\to\Delta_s$ for the resulting proper
map and $E$ for the smooth strict component. The
space $U'$ is smooth, $E$ still has multiplicity
one, and the smooth nearby fibres are the original
connected tori. Connectedness of all full fibres
follows also from Stein factorization over the
normal disc, since its generic degree is one.

\noindent\textbf{Step 2.} If $Z_i$ were of general type, so would be $E$.
Here is the precise comparison. The map $E\to S_i$
factors through normalization and gives a
generically finite meromorphic map $E\dashrightarrow
Z_i$. Resolve its graph to obtain a birational
map $a:\widehat E\to E$ and a generically finite
holomorphic map $b:\widehat E\to Z_i$ between
smooth compact surfaces. The actual Jacobian map
gives
\[
                 K_{\widehat E}=b^*K_{Z_i}+R,
                 \qquad R\ge0.                 \tag{C.SC.3}\label{cdp:eq:C.SC.3}
\]
Pullback of pluricanonical sections is injective.
The generically finite pullback of a big canonical
line is big; equivalently a generically finite
pluricanonical map of $Z_i$ remains generically
finite after composition with $b$. Equation
\eqref{cdp:eq:C.SC.3} therefore makes $K_{\widehat E}$ big.
Birational invariance gives $\kappa(E)=2$.
In particular $E$ is projective, as is every
smooth compact complex surface of general type.

\noindent\textbf{Step 3.} We construct a nonzero meromorphic vector field
on $E$ with single-valued solutions. Define the
coherent relative tangent sheaf by the actual
kernel
\[
 \mathcal T_{U'/\Delta_s}
 =\ker\bigl[df':T_{U'}\longrightarrow f'^*T_{\Delta_s}\bigr].
 \tag{C.SC.4}\label{cdp:eq:C.SC.4}
\]
Its proper direct image is torsion-free on the
disc: a section annihilated by a power of $s$
vanishes on the dense open punctured tube and
hence vanishes identically. Its rank is two,
because the global holomorphic vector fields
of a two-torus are its two translation fields.
It is thus a locally free rank-two sheaf after
shrinking the disc. Choose a nonzero section
$v$, viewed by adjunction as an actual
holomorphic vertical vector field on $U'$.
At the generic point of $E$ its order of
vanishing is a finite integer $k\ge0$.
Set
\[
                              w=s^{-k}v.
 \tag{C.SC.5}\label{cdp:eq:C.SC.5}
\]
This is a global meromorphic vector field on
$U'$, with poles confined to the central
divisor. At a generic point of $E$, the
multiplicity-one assertion makes $s$ a local
transverse coordinate. Since $v(s)=0$, write
$v=a(x,y,s)\partial_x+b(x,y,s)\partial_y$.
The integer $k$ is the minimum of the two
coefficient orders along $s=0$. Thus $w$
is holomorphic and nonzero generically on
$E$ and has a well-defined nonzero tangent
meromorphic restriction $w_E$ to $E$.

\noindent\textbf{Step 4.} We use the precise closed-univalence theorem
\cite[Theorem~3.3, with its proof]{Guillot2024}:
for a holomorphic vector field on a complex
manifold of any dimension, the locus where
its maximal solutions are single-valued is
closed. Such a field is called semicomplete
when that locus is the whole manifold.
The closure assertion is a path-continuation
statement. A failure at a point is witnessed
by a nonclosed path in a solution leaf whose
time-form integral is zero. Lifting the
time parametrization of this compact path
through finitely many flow charts produces
the same nonclosed zero-time path through
all sufficiently nearby initial points.
This proves openness of the failure locus
and records exactly the comparison used here.

\noindent\textbf{Step 5.} Let $P$ be the polar divisor of $w$ and put
$W=U'\setminus P$. On the dense open subset
$U'^*\subset W$, $w$ restricts to a
holomorphic vector field on each compact
torus, multiplied by the constant $s^{-k}$
on that fibre. Each such field is complete:
its complex-time solution stays in the
compact fibre, so local existence extends
it for all times. Consequently $w$ is
semicomplete on all of $W$ by the stated
closedness theorem.

\noindent\textbf{Step 6.} The hypersurface $E\cap W$ is invariant.
Indeed $w(s)=0$, and at any point where
$w$ is holomorphic, factoring the local
equation of $s$ into its irreducible
components shows that $w$ preserves the
ideal of each component. A solution
starting in $E\cap W$ therefore stays
there, and its single-valuedness is the
one inherited from $w$. The curve removed
from $E$ by the ambient polar divisor may
contain removable poles of $w_E$. Apply
the same closedness theorem once more on
$E$ minus the actual polar divisor of
$w_E$. Its dense open subset $E\cap W$
is already semicomplete. Hence $w_E$ is
a nonzero semicomplete meromorphic vector
field on the entire compact surface $E$,
in the precise sense of single-valuedness
on its holomorphic locus.

\noindent\textbf{Step 7.} We finish by proving that a projective
surface of general type admits no such
field. Suppose first that one maximal
solution of $w_E$ has Zariski-dense image.
The Main Theorem of \cite{Guillot2019},
including the clarification immediately
following it, gives a birational model
which is an abelian surface or admits a
rational or elliptic fibration. The
holomorphic alternative is covered there
by Proposition~22; the dense leaf excludes
a first integral in that alternative.
All these models have Kodaira dimension
at most one, contradicting $\kappa(E)=2$.
Only single-valued maximality, already
proved, is needed in this application;
the solution need not be entire.

\noindent\textbf{Step 8.} Here is the other case, with its actual
fibration map. If every solution leaf is
contained in a proper analytic curve,
the closures of the regular leaves are
irreducible algebraic invariant curves
on the projective surface $E$. They
cover its open regular foliation locus.
Use the countably many Hilbert polynomials
and flattening strata of their universal
families to parametrize reduced irreducible
curves. The locus of such integral curves
is constructible. After clearing the
polar divisor, invariance is the vanishing
of the induced normal-component map of
the vector field on the universal curve;
on these strata it is a constructible
condition. Thus we have countably many
irreducible locally closed parameter
strata of integral invariant curves.
Their incidence images in $E$ are
constructible. If none were dense, their
Zariski closures would be countably many
proper closed analytic subsets covering
the regular foliation locus, contrary
to the Baire theorem. Choose a stratum
whose incidence dominates $E$. Its
generic curves are reduced and
irreducible. At a regular point of the
foliation exactly one such irreducible
curve can pass: two would contain the
same local solution germ and would
coincide. Thus the parameter has dimension
one and the incidence map to $E$ is
generically one-to-one. Compactify and
normalize this parameter curve and
close the incidence graph. Inverting
its birational projection to $E$ and
projecting to the parameter gives an
actual meromorphic first integral.
Resolve its indeterminacies and take its
Stein factor to obtain a morphism
\[
                         q:E'\longrightarrow C
 \tag{C.SC.6}\label{cdp:eq:C.SC.6}
\]
with connected generic fibre a smooth
compact curve. The lifted meromorphic
field remains semicomplete: it agrees
with the old field on the birational
open set, and the same closedness theorem
extends single-valuedness over its full
holomorphic locus. It is a nonzero
tangent meromorphic field on a generic
fibre $G$ of \eqref{cdp:eq:C.SC.6}.

\noindent\textbf{Step 9.} Such a field on a compact curve has no
pole. A pole of order $r\ge1$ would
have the local form $z^{-r}a(z)\partial_z$
with $a(0)\ne0$. Its time form is
$z^r a(z)^{-1}dz$, whose primitive has
a zero of order $r+1$. The inverse
solution therefore has $r+1$ branches
around that time value, all lying in
the punctured coordinate disc. This
contradicts semicompleteness. The field
on $G$ is consequently holomorphic
and nonzero, so
$\deg T_G=2-2g(G)\ge0$ and $g(G)\le1$.

\noindent\textbf{Step 10.} This also contradicts general type
directly through pluricanonical maps.
Adjunction on the smooth generic fibre
gives $K_{E'}|_G=K_G$, of degree at
most zero. Every linear system
$|nK_{E'}|$ therefore restricts to a
constant map on $G$: a negative-degree
line has no sections, and a degree-zero
line on a connected compact curve has
at most one independent section. A
generically finite pluricanonical map
of the general-type surface $E'$ cannot
contract its generic fibres in this
way. This excludes the all-algebraic-
leaves case as well.

\noindent\textbf{Step 11.} Both cases contradict the general-type
consequence of \eqref{cdp:eq:C.SC.3}. The chosen
$Z_i$ cannot have Kodaira dimension
two. Since $S_i$ was arbitrary,
\eqref{cdp:eq:C.SC.1} follows for every component.
\end{proof}






\begin{proposition}[Two commuting semicomplete fields on the actual component cover]\label{cdp:st:C.2.UN.16}
\StatementRecord{cdp:st:C.2.UN.16}{proved}{Historical events: Claude Fable 5.2 (Anthropic), effort high, 16 November 2026: carried as background without fault found. GPT-6 Astra (OpenAI), ChatGPT Pro mode, effort Pro, S6\_09\_PAPERS, 17 November 2026: targeted proof check. Also checked in this writing. Pre-continuation writing events: GPT-6 Astra (OpenAI), ChatGPT Pro mode, effort Pro, S6\_10\_PAPERS, 18 November 2026: targeted proof check; used and rechecked in this run. Continuation mathematical and source checks: GPT-6 Astra (OpenAI), ChatGPT Pro mode, effort Pro, S6\_10\_PAPERS continuation, 19 November 2026: targeted proof check; used and rechecked in this run. Continuation metadata events: GPT-6 Astra (OpenAI), ChatGPT Pro mode, effort Pro, S6\_10\_PAPERS continuation, 19 November 2026: used and rechecked in this run; status-vocabulary reconciliation; metadata reconciliation only; no new mathematical check. \emph{Pre-continuation recorded proof scope.} The determinant valuation bounds the actual saturation between M and s-to-the-minus-d M, and its restriction kernel is exactly sL. This produces two commuting independent meromorphic fields, with semicompleteness for every constant combination. Their rank-one ratio constructs the actual first integral. Guillot2024 Theorems3.3/9.2 and DOT2001 Theorem0.3 were opened; DOT2000 pp256/260 supplies the positive-b2 one-field bound. The Inoue1974 exact-proposition scan remains a separate reference gate}{the hypotheses, notation and definitions stated here; \cref{cdp:st:C.2.UN.15}; \cref{cdp:eq:C.CF.3}; \cref{cdp:eq:C.CF.5}; \cref{cdp:eq:C.CF.7}; \cite[Theorem~3.3]{Guillot2024}; \cref{cdp:eq:C.CF.8}; \cite[Theorem~9.2]{Guillot2024}; \cite[Theorem~0.3]{DOT2001}; \cite[pp.~256 and~260]{DOT2000}; \cite[Propositions~2(ii), 3(ii) and~5]{Inoue1974}}

Let $f:U\to\Delta$ be a proper holomorphic map from a smooth
complex threefold, with connected fibres and with smooth complex
two-tori over $\Delta^*$. Fix an irreducible central component $S_i$
of multiplicity $m_i$. Make the actual base change $t=s^{m_i}$,
normalize, and resolve as in \ref{cdp:st:C.2.UN.15}. Every chosen smooth strict
component $E$ above $S_i$ has multiplicity one and carries a
two-dimensional complex vector space
\[
 \mathfrak a_E\subset
 H^0(E,T_E\otimes_{\mathcal O_E}\mathcal M_E)
 \tag{C.CF.1}\label{cdp:eq:C.CF.1}
\]
of commuting meromorphic vector fields. Every member of
$\mathfrak a_E$ is semicomplete on its holomorphic locus.
The map $E\to S_i$ factors through normalization and is
proper and generically finite.

There are two exhaustive geometric alternatives for this particular
space of fields. Its generic rank is two, in which case any complex
basis is a commuting meromorphic frame; or its generic rank is one,
in which case an actual resolution of a meromorphic first integral
and its Stein factor give
\[
 \widehat E\xrightarrow{\ q\ }B,
 \qquad \widehat E\xrightarrow{\rho}E,
 \tag{C.CF.2}\label{cdp:eq:C.CF.2}
\]
where $\rho$ is proper birational, $B$ is a smooth compact curve,
and the connected general fibre of $q$ is rational or elliptic.

When $a(S_i^\nu)=0$, the unique minimal model of $E$ is a
two-torus or a Hopf surface. This last conclusion concerns the
specified generically finite covering component; the maps giving
it are part of the assertion.
\end{proposition}



\begin{proof}
\noindent\textbf{Step 1.} Write $f':U'\to\Delta_s$ for the resolved base change.
The local equations
$w^{m_i}=u$, $z=s/w$ in (C.SC.2), starting from
$t=u z^{m_i}$, prove both multiplicity one and generic
finiteness. Properness and the universal property of normalization
give the asserted factorization $E\to S_i^\nu\to S_i$.
These are the actual maps used below.

\noindent\textbf{Step 2.} Put $\mathcal O=\mathbb C\{s\}$ and
$K=\operatorname{Frac}(\mathcal O)$. Let
\[
 M=(f'_*\mathcal T_{U'/\Delta_s})_0,
 \qquad
 L=\{v\in M\otimes_{\mathcal O}K:
                     \operatorname{ord}_E(v)\ge0\}.
 \tag{C.CF.3}\label{cdp:eq:C.CF.3}
\]
The order in this formula is the minimum of the coefficient
orders in a local tangent frame at the generic point of $E$.
It is independent of that frame. The direct image is coherent
by properness, is torsion-free because its sections are actual
vector fields on $U'$, and has rank two by base change on the
smooth torus locus. Consequently $M$ is a free rank-two
$\mathcal O$-module. Its sections are identified with actual
holomorphic vertical fields by the definition of direct image.
Any two of them commute: their restrictions to a smooth torus
are translation fields, so their bracket vanishes on the dense
punctured tube and hence everywhere.

\noindent\textbf{Step 3.} Here is a finite bound for the saturation in \eqref{cdp:eq:C.CF.3}.
Choose an $\mathcal O$-basis $v_1,v_2$ of $M$.
At a generic point of $E$, multiplicity one makes $s$ a
coordinate, so the two vertical tangent coordinates $x,y$
give
\[
 (v_1\ v_2)=(\partial_x\ \partial_y)A,
 \qquad d=\operatorname{ord}_E(\det A)<\infty.
 \tag{C.CF.4}\label{cdp:eq:C.CF.4}
\]
All entries of $A$ are holomorphic there. The determinant is
nonzero because the evaluation of the two translation fields
is an isomorphism on each nearby torus. In particular $d\ge0$.
For $v=a_1(s)v_1+a_2(s)v_2\in L$, the column $Aa$ has
nonnegative order along $E$. The identity
$a=(\det A)^{-1}\operatorname{adj}(A)(Aa)$ shows that
each $a_j$ has order at least $-d$. The order of a function
of $s$ along $E$ equals its $s$-adic order, again by
multiplicity one. We obtain actual inclusions
\[
                        M\subset L\subset s^{-d}M.
 \tag{C.CF.5}\label{cdp:eq:C.CF.5}
\]
The set $L$ is an $\mathcal O$-submodule. Since $\mathcal O$
is noetherian, \eqref{cdp:eq:C.CF.5} makes it finite; since it is torsion-free
of rank two over the discrete valuation ring $\mathcal O$,
it is free of rank two.

\noindent\textbf{Step 4.} Every member of $L$ is a global meromorphic vertical field
on a sufficiently small representative of $U'$, regular at
the generic point of $E$. Its restriction is therefore a
global meromorphic section of $T_E$. Tangency follows from
$v(s)=0$ at the generic smooth point and then as an identity
of meromorphic sections. Restriction defines the exact map
\[
 \operatorname{res}_E:L\longrightarrow
 H^0(E,T_E\otimes\mathcal M_E),
 \qquad
 \ker(\operatorname{res}_E)=sL.
 \tag{C.CF.6}\label{cdp:eq:C.CF.6}
\]
Indeed, zero restriction means order at least one along $E$;
division by $s$ then belongs to $L$. Conversely, every member
of $sL$ restricts to zero. Thus
\[
                  L/sL\simeq\mathbb C^2
                    \hookrightarrow
                  H^0(E,T_E\otimes\mathcal M_E).
 \tag{C.CF.7}\label{cdp:eq:C.CF.7}
\]
Define $\mathfrak a_E$ to be its image. Since vertical fields
annihilate functions of $s$, brackets of members of
$M\otimes K$ vanish. Hence the fields in \eqref{cdp:eq:C.CF.7} commute.

\noindent\textbf{Step 5.} For completeness, semicompleteness applies to the entire
space, not only to a selected nonzero element. Given
$V\in\mathfrak a_E$, choose its representative $v\in L$.
Its poles are confined to the central fibre. On every
punctured fibre it is a holomorphic translation field
multiplied by constants, and is complete. Remove the
ambient polar divisor. The closedness of the
uniformizability locus for a holomorphic vector field
\cite[Theorem~3.3]{Guillot2024} makes $v$ semicomplete
there, since the punctured tube is dense. The component
$E$ is invariant: locally $v(s)=0$ and factorization of
the local equation of $s$ show that each irreducible
component ideal is invariant wherever $v$ is holomorphic.
Solutions in $E$ inherit single-valuedness. Some removed
curves may be removable poles for $V$; the same closedness
theorem on $E$ minus its actual polar divisor extends
semicompleteness across those curves. This proves the
assertion for every constant linear combination in
$\mathfrak a_E$.

\noindent\textbf{Step 6.} Take a complex basis $V_1,V_2$ of $\mathfrak a_E$.
When $V_1\wedge V_2$ is nonzero, it is the announced
meromorphic frame. Otherwise there is a globally defined
meromorphic function $h$ with
\[
                  V_2=hV_1,
                  \qquad V_1(h)=0.
 \tag{C.CF.8}\label{cdp:eq:C.CF.8}
\]
The local coefficient ratios agree on their overlaps,
which proves that $h$ is global. It is nonconstant by
the complex-linear independence in \eqref{cdp:eq:C.CF.7}. The second
identity follows from the vanishing bracket.
Resolve the graph of $h:E\dashrightarrow\mathbb P^1$
and take the Stein factor of its projection to
$\mathbb P^1$. This gives exactly \eqref{cdp:eq:C.CF.2}, with a
finite map $B\to\mathbb P^1$. The lifted field is
tangent to $q$ and nonzero on a general fibre.
Its semicompleteness passes through the birational
open isomorphism and then over its full holomorphic
locus by the same closedness theorem.

\noindent\textbf{Step 7.} A nonzero semicomplete meromorphic field on a smooth
compact curve has no pole. At a pole of order $r\ge1$
its expression is $z^{-r}a(z)\partial_z$, $a(0)\ne0$.
The time form $z^r a(z)^{-1}dz$ has a primitive
vanishing to order $r+1$; its inverse solution has
distinct analytic continuations around that time
value. They all lie in the punctured coordinate
disc, contradicting single-valuedness. The restricted
field is therefore holomorphic, so
$\deg T_G=2-2g(G)\ge0$ for a general fibre $G$.
This proves the rational-or-elliptic assertion.

\noindent\textbf{Step 8.} Finally suppose $a(S_i^\nu)=0$. The map
$E\to S_i^\nu$ has a finite Stein factor
$E\to T'\to S_i^\nu$, whose first map is
bimeromorphic. For a meromorphic function on $T'$, the
elementary symmetric functions of its values over
the finite map are meromorphic on $S_i^\nu$.
They give a monic algebraic equation over
$\mathcal M(S_i^\nu)$. Thus the meromorphic function
field extension is algebraic and
\[
                            a(E)=a(S_i^\nu)=0.
 \tag{C.CF.9}\label{cdp:eq:C.CF.9}
\]
In particular \eqref{cdp:eq:C.CF.8} cannot occur.

\noindent\textbf{Step 9.} Apply the compact-surface classification of semicomplete
meromorphic fields \cite[Theorem~9.2]{Guillot2024}
separately to $V_1$ and $V_2$. Its fibration alternative
would give a nonconstant meromorphic function on $E$;
therefore each field is holomorphic on some smooth
bimeromorphic model. An algebraic-dimension-zero
surface is not ruled. Its smooth minimal model $E_0$
is unique, and every smooth bimeromorphic model maps
to it by successive point blowdowns. Holomorphic
fields descend under these actual maps: for a
blowdown $\pi$ the equality
$\pi_*\mathcal O=\mathcal O$ sends the derivation
$g\mapsto V(\pi^*g)$ to a holomorphic derivation on
the base. On the common open isomorphism it is the
original field. The two resulting holomorphic fields
on the same $E_0$ remain complex-linearly independent
and commute.

\noindent\textbf{Step 10.} The complete list of minimal compact surfaces
admitting a nonzero holomorphic field is
\cite[Theorem~0.3]{DOT2001}. In algebraic dimension
zero, its possibilities are two-tori, Hopf surfaces,
Inoue surfaces, and minimal class-VII surfaces with
positive second Betti number. The last class has
at most one independent holomorphic field
\cite[pp.~256 and~260]{DOT2000}; Inoue surfaces
also have at most one
\cite[Propositions~2(ii), 3(ii) and~5]{Inoue1974}.
The two independent fields on $E_0$ exclude both
of those alternatives. Hence $E_0$ is a two-torus
or a Hopf surface, as claimed.
\end{proof}






\begin{proposition}[Canonical support forces the numerically anticanonical class VII case]\label{cdp:st:C.2.UN.17}
\StatementRecord{cdp:st:C.2.UN.17}{proved}{Historical events: Claude Fable 5.2 (Anthropic), effort high, 16 November 2026: carried as background without fault found. GPT-6 Astra (OpenAI), ChatGPT Pro mode, effort Pro, S6\_09\_PAPERS, 17 November 2026: targeted proof check. Also checked in this writing. Pre-continuation writing events: GPT-6 Astra (OpenAI), ChatGPT Pro mode, effort Pro, S6\_10\_PAPERS, 18 November 2026: targeted proof check; used and rechecked in this run. Continuation mathematical and source checks: GPT-6 Astra (OpenAI), ChatGPT Pro mode, effort Pro, S6\_10\_PAPERS continuation, 19 November 2026: targeted proof check; used and rechecked in this run. Continuation metadata events: GPT-6 Astra (OpenAI), ChatGPT Pro mode, effort Pro, S6\_10\_PAPERS continuation, 19 November 2026: used and rechecked in this run; status-vocabulary reconciliation; metadata reconciliation only; no new mathematical check. \emph{Pre-continuation recorded proof scope.} Canonical evaluation and every effective Cartier pullback give the exact full-multiplicity line identity. On the minimal model, the negative-semidefinite curve matrix is handled blockwise, including a strictly positive kernel vector for each singular block. Clearing rational coefficients and integral torsion gives the actual effective anticanonical power with flat line; the flat-character surjection was re-derived from the real-to-coherent H1 map}{the hypotheses, notation and definitions stated here; \cref{cdp:st:A.5.3}; \cref{cdp:eq:C.CA.4}; \cref{cdp:eq:C.CA.3}; \cref{cdp:eq:C.CA.5}; \cite[Remark~2.7]{DlousskyTeleman2012}; \cite[\S1, pp.~4--5]{Dloussky2006}; \cref{cdp:eq:C.CA.7}; \cref{cdp:eq:C.CA.9}; \cite[Theorem~0.2]{Dloussky2006}; \cref{cdp:eq:C.CA.6}}

Let $f:U\to\Delta$ be a proper holomorphic map from a smooth
complex threefold, with connected fibres and smooth two-torus
fibres over $\Delta^*$. Write its full central fibre as
$F=\sum_i m_iS_i$. For each component let
$h_i:Z_i\to S_i$ be its normalization followed by minimal
resolution, and let $D_i\ge0$ be the intrinsic resolution
conductor of \ref{cdp:st:A.5.3}. There are nonnegative integers $a_j$ with
$K_U\simeq\OO_U(\sum_j a_jS_j)$. Put
$E_{ij}=h_i^*(S_j)$ for $j\ne i$, using the effective
Cartier pullback. The exact line-bundle identity is
\[
 K_{Z_i}^{\otimes m_i}\otimes\OO_{Z_i}(m_iD_i)
 \simeq
 \OO_{Z_i}\!\left(
  \sum_{j\ne i}\bigl[m_i a_j-m_j(a_i+1)\bigr]E_{ij}
 \right).                                                \tag{C.CA.1}\label{cdp:eq:C.CA.1}
\]
In particular a positive tensor power of the canonical line on the minimal
model of $Z_i$ has a meromorphic section whose divisor is
supported on the images of these interface and conductor
curves.

When this minimal model $M$ is of class VII, it has the
following restrictions. If $b_2(M)=0$, it is a Hopf surface.
If $b_2(M)>0$, there are an integer $n>0$, a nonzero
effective integral divisor $A$ and a complex-flat line
bundle $L$ such that
\[
                 K_M^{\otimes n}\otimes\OO_M(A)
                         \simeq L.                       \tag{C.CA.2}\label{cdp:eq:C.CA.2}
\]
Consequently $M$ contains a global spherical shell by the
proved numerically anticanonical theorem. In particular,
the unrestricted normalization problem does not require
the classification of class VII surfaces without a
numerically anticanonical divisor. These are restrictions
on the components of the actual map $f$; no occurrence of
every surface with \eqref{cdp:eq:C.CA.2} is asserted.
\end{proposition}


\begin{proof}
\noindent\textbf{Step 1.} After shrinking the disc, $f_*K_U$ is a locally free
rank-one sheaf. Coherence follows from properness, and
torsion-freeness follows because a holomorphic canonical
section killed by a power of the base coordinate vanishes
on a dense open subset. Its rank is one since every nearby
torus has a one-dimensional space of canonical forms.
Choose a frame and evaluate it on $U$. Over $\Delta^*$,
Grauert base change identifies its restriction with a
nonzero torus canonical form, which is nowhere zero.
The evaluation is therefore a nonzero holomorphic section
of $K_U$ with divisor $\sum_j a_jS_j$, where $a_j\ge0$.
This constructs the claimed line-bundle isomorphism.

\noindent\textbf{Step 2.} Let $N_i=h_i^*\OO_U(S_i)$. The canonical sections defining
$S_j$ do not vanish identically on $S_i$ for $j\ne i$.
Their pullbacks define the actual effective divisors
$E_{ij}$, including any exceptional orders. Pullback of
the principal full fibre gives
\[
 N_i^{\otimes m_i}\otimes
 \OO_{Z_i}\!\left(\sum_{j\ne i}m_jE_{ij}\right)
                         \simeq\OO_{Z_i}.               \tag{C.CA.3}\label{cdp:eq:C.CA.3}
\]
Adjunction for the individual reduced hypersurface $S_i$,
followed by the conductor comparison of \ref{cdp:st:A.5.3}, gives
\[
 \begin{split}
 K_{Z_i}\otimes\OO_{Z_i}(D_i)
   &\simeq h_i^*\bigl((K_U\otimes\OO_U(S_i))|_{S_i}\bigr)\\
   &\simeq N_i^{\otimes(a_i+1)}\otimes
      \OO_{Z_i}\!\left(\sum_{j\ne i}a_jE_{ij}\right).
 \end{split}                                             \tag{C.CA.4}\label{cdp:eq:C.CA.4}
\]
Raise \eqref{cdp:eq:C.CA.4} to the $m_i$-th power and substitute
\eqref{cdp:eq:C.CA.3}. This proves \eqref{cdp:eq:C.CA.1}, with its signs and every
full-fibre multiplicity retained.

\noindent\textbf{Step 3.} Let $\tau:Z_i\to M$ be the sequence of point blowdowns
to a minimal model. Identity \eqref{cdp:eq:C.CA.1} provides a
meromorphic section of $K_{Z_i}^{\otimes m_i}$ with the
displayed signed divisor, less $m_iD_i$. On the open set
where $\tau$ is an isomorphism this is a meromorphic
$m_i$-canonical section on $M$. In local coordinates at
each omitted point, the coefficient of that section is a
meromorphic function on a punctured smooth surface and
extends meromorphically. Thus the section extends across
those points. Its divisor is precisely the pushforward:
exceptional curves disappear, and every other curve
keeps its order. Equivalently, at an individual blowup
the exact canonical formula
$K_{\widehat M}=\tau^*K_M+E$ removes the corresponding
$m_iE$ from the pulled-back pluricanonical divisor.
We have therefore constructed a signed integral divisor
$B$ on $M$ and an actual isomorphism
\[
                         K_M^{\otimes m_i}\simeq\OO_M(B).
                                                               \tag{C.CA.5}\label{cdp:eq:C.CA.5}
\]
Every component of $B$ is the image of an interface or
intrinsic conductor curve from \eqref{cdp:eq:C.CA.1}.

\noindent\textbf{Step 4.} Suppose now that $M$ is of class VII. If it contained no
curve, \eqref{cdp:eq:C.CA.5} would have $B=0$ and would trivialize a
positive power of $K_M$. This contradicts
$\kappa(M)=-\infty$. When $b_2(M)=0$, choose a nonempty
curve. A positive multiple of it is homologically trivial,
since its integral class is torsion. The classification
of such effective divisors in
\cite[Remark~2.7]{DlousskyTeleman2012} gives a Hopf
surface in this zero-Betti case. This proves the first
alternative without an assertion about unknown class VII
surfaces.

\noindent\textbf{Step 5.} For $b_2(M)>0$, let $C_1,\ldots,C_r$ be the distinct
curves in the support of $B$, and write
\[
 c_1(K_M)_{\mathbb R}=\sum_{\ell=1}^r c_\ell[C_\ell],
 \qquad c_\ell\in\Q.                                  \tag{C.CA.6}\label{cdp:eq:C.CA.6}
\]
The intersection form on $H^2(M,\mathbb R)$ is negative
definite and $K_M^2=-b_2(M)$. Moreover
$K_M\cdot C_\ell\ge0$ for every curve on the minimal
surface; these standard class VII facts are stated with
their adjunction argument in \cite[\S1, pp.~4--5]{Dloussky2006}.
For completeness, adjunction gives
$K_M\cdot C=2p_a(C)-2-C^2$. A curve of arithmetic genus
zero is a smooth rational curve; minimality excludes
self-intersection $-1$, and self-intersection zero would
give zero real class and hence $K_M\cdot C=0$, contradicting
adjunction. Every other curve has arithmetic genus at
least one and $C^2\le0$, so also has nonnegative canonical
intersection.

\noindent\textbf{Step 6.} Let $Q=(C_\ell\cdot C_k)$ and
$v_\ell=K_M\cdot C_\ell$. Then
\[
                         Qc=v,\qquad v\ge0.             \tag{C.CA.7}\label{cdp:eq:C.CA.7}
\]
The matrix $Q$ is negative semidefinite, and all its
off-diagonal entries are nonnegative. Separate it into
blocks according to the connected components of the
intersection graph. For a negative-definite block, write
$c=c_+-c_-$ with disjoint nonnegative coordinate supports.
If $c_+\ne0$, then
\[
 c_+^tQc=c_+^tQc_+-c_+^tQc_-<0,
 \qquad c_+^tv\ge0,                                    \tag{C.CA.8}\label{cdp:eq:C.CA.8}
\]
contradicting \eqref{cdp:eq:C.CA.7}. Thus every coefficient in this
block is nonpositive.

\noindent\textbf{Step 7.} We also treat every singular block, rather than assuming
linear independence of the curve classes. Choose a nonzero
kernel vector $u$. Replacing it by its coordinatewise
absolute value cannot decrease its quadratic form, because
the off-diagonal entries are nonnegative. Both quadratic
forms must be zero by negative semidefiniteness. Hence
$|u|$ is again in the kernel. A zero coordinate of $|u|$
would, in the corresponding kernel equation, force every
neighboring coordinate to be zero. Connectedness would
then force the whole vector to vanish. Thus the block has
a strictly positive kernel vector. Pairing \eqref{cdp:eq:C.CA.7} with
this vector gives $v=0$ on the block, and then $Qc=0$.
The real divisor class $\beta=\sum c_\ell[C_\ell]$ of
that block satisfies $\beta^2=0$. The ambient negative
definite form forces $\beta=0$.

\noindent\textbf{Step 8.} Discard the singular blocks and negate the coefficients
on the negative-definite blocks. This gives an effective
rational divisor $A_{\Q}$ on the actual curves of $M$
such that
\[
                   c_1(K_M+A_{\Q})_{\mathbb R}=0.        \tag{C.CA.9}\label{cdp:eq:C.CA.9}
\]
It is nonzero, since otherwise $c_1(K_M)_{\mathbb R}=0$
would contradict $K_M^2=-b_2(M)<0$.

\noindent\textbf{Step 9.} Choose a positive multiple clearing the denominators in
$A_{\Q}$ and then another positive multiple killing the
torsion integral Chern class in \eqref{cdp:eq:C.CA.9}. We obtain an
integer $n>0$, a nonzero effective integral divisor $A$,
and a topologically trivial holomorphic line
\[
                       L=K_M^{\otimes n}\otimes\OO_M(A).
                                                               \tag{C.CA.10}\label{cdp:eq:C.CA.10}
\]
We specify why this line is complex-flat. On a class VII
surface $p_g=0$, $c_2(M)=b_2(M)$ and $K_M^2=-b_2(M)$,
so Noether's formula gives $\chi(\OO_M)=0$ and
$h^1(\OO_M)=1$. Also $\dim_{\C}H^1(M,\C)=1$.
The natural map $H^1(M,\R)\to H^1(M,\OO_M)$ is
injective: if a real closed one-form $\alpha$ has
$\alpha^{0,1}=\bar\partial(p+\mathrm i q)$, with $p,q$
real, then
$\alpha=dp+\mathrm i(\bar\partial q-\partial q)$ and
closedness gives $\partial\bar\partial q=0$. The
pluriharmonic maximum principle on the compact connected
surface makes $q$ constant, so $\alpha$ is exact.
Consequently the complex-linear map
$H^1(M,\C)\to H^1(M,\OO_M)$ is an isomorphism.
The exponential sequences for constant and holomorphic
coefficients therefore give the surjection
\[
             H^1(M,\C^*)\longrightarrow\Pic^0(M).       \tag{C.CA.11}\label{cdp:eq:C.CA.11}
\]
This is the actual flat-character map, taking constant
transition functions to the same holomorphic transition
functions. It supplies a complex-flat realization of $L$.

\noindent\textbf{Step 10.} The canonical section of $\OO_M(A)$ now becomes a
nonzero section of $K_M^{-n}\otimes L$. These are exactly
the hypotheses of \cite[Theorem~0.2]{Dloussky2006},
which proves existence of a global spherical shell.
This proves \eqref{cdp:eq:C.CA.2} and the stated normalization
restriction. The argument also shows directly why a
positive-Betti class VII surface whose every curve class
is zero in real cohomology cannot be a minimal
normalization model here: \eqref{cdp:eq:C.CA.6} would contradict
$K_M^2=-b_2(M)$.
\end{proof}



\ref{cdp:st:C.2.UN.19} subsequently excludes the positive-$b_2$ alternative
altogether by the actual cyclic base-change quotient and its
descent to a Hopf surface. The canonical-support and
numerically anticanonical comparison above is retained as an
independent necessary reduction. It does not supply an example
of a positive-$b_2$ class VII component.




\begin{proposition}[Every class VII normalization component has Hopf minimal model]\label{cdp:st:C.2.UN.19}
\StatementRecord{cdp:st:C.2.UN.19}{proved}{Historical events: Claude Fable 5.2 (Anthropic), effort high, 16 November 2026: carried as background without fault found. GPT-6 Astra (OpenAI), ChatGPT Pro mode, effort Pro, S6\_09\_PAPERS, 17 November 2026: targeted proof check. Also checked in this writing. Pre-continuation writing events: GPT-6 Astra (OpenAI), ChatGPT Pro mode, effort Pro, S6\_10\_PAPERS, 18 November 2026: targeted proof check; used and rechecked in this run. Continuation mathematical and source checks: GPT-6 Astra (OpenAI), ChatGPT Pro mode, effort Pro, S6\_10\_PAPERS continuation, 19 November 2026: targeted proof check. Continuation metadata events: GPT-6 Astra (OpenAI), ChatGPT Pro mode, effort Pro, S6\_10\_PAPERS continuation, 19 November 2026: used and rechecked in this run; status-vocabulary reconciliation; metadata reconciliation only; no new mathematical check. \emph{Pre-continuation recorded proof scope.} The stabilizer of the chosen root-cover component gives the precise generically finite quotient of the original normalization. Torus quotients have even first Betti number through the invariant-cohomology and cyclic-resolution maps. For Hopf quotients the derivative determinant character fixes H1, the Lefschetz calculation excludes isolated fixed points, and the explicit pseudoreflection invariant ring makes the quotient smooth. Its Betti numbers, canonical pullback and actual curve then identify the Hopf minimal model}{the hypotheses, notation and definitions stated here; \cite[\S1, universal and canonical desingularization]{BierstoneMilman1997}; \cref{cdp:eq:C.HQ.2}; \cref{cdp:st:C.2.UN.17}; \cite[p.~256]{DOT2000}; \cref{cdp:st:C.2.UN.16}; \cref{cdp:eq:C.HQ.3}; \cref{cdp:eq:C.HQ.4}; \cite[\S1.1]{Nakamura2008VII}; \cite[Theorem~2.1]{Khaled2024Hopf}; \cref{cdp:eq:C.HQ.6}; \cref{cdp:eq:C.HQ.7}; \cref{cdp:eq:C.HQ.8}; \cref{cdp:eq:C.HQ.10}; \cite[Remark~2.7]{DlousskyTeleman2012}}

For every proper holomorphic degeneration of smooth complex
two-tori with smooth three-dimensional total space, each
central component whose normalization resolution is of class
VII has a Hopf surface as its minimal model. This holds for
arbitrary multiplicities, arbitrary central singularities,
and arbitrary Euler number. In particular no Inoue surface
and no minimal class VII surface with positive second Betti
number is a normalization minimal model of such a component.

The comparison proving this assertion is the cyclic cover
obtained from $t=s^{m_i}$ for that component. Its stabilizer
acts on the minimal model of the covering component, and
the quotient is bimeromorphic to the original normalization.
\end{proposition}



\begin{proof}
\noindent\textbf{Step 1.} We first identify the quotient in the actual base change.
Normalize $U\times_{\Delta_t}\Delta_s$, where
$t=s^{m_i}$. Its cyclic group $\mu_{m_i}$ acts over $U$.
Choose a resolution of the total space and central divisor
equivariant for this finite group. Canonical analytic
resolution has this property: the invariants defining its
centres are preserved by every local isomorphism, so the
action lifts through every blowup
\cite[\S1, universal and canonical desingularization]{BierstoneMilman1997}.
The algorithm is an isomorphism over the punctured tube
and at the generic smooth points used in (C.SC.2).

\noindent\textbf{Step 2.} Let $E$ be a chosen smooth strict component above $S_i$,
and let $G\subset\mu_{m_i}$ be its stabilizer. At a
generic point of $S_i$ the equations are exactly
\[
                 w^{m_i}=u,\qquad z=s/w.
 \tag{C.HQ.1}\label{cdp:eq:C.HQ.1}
\]
Consequently the reduced inverse image over a dense
smooth open subset of $S_i$ is a $\mu_{m_i}$-torsor.
The cyclic group permutes its connected components
transitively. The stabilizer of a selected component
acts freely and transitively on each fibre of that
component. Equivalently, the degree of $E\to S_i$
at the generic point is $|G|$, and quotienting by
$G$ gives degree one. The proper map to normalization
therefore descends to a proper bimeromorphic map
\[
                         E/G\longrightarrow S_i^\nu.
 \tag{C.HQ.2}\label{cdp:eq:C.HQ.2}
\]
This argument remains valid when the root cover is
disconnected over $S_i$; the stabilizer, rather than
the whole base-change group, is the group in \eqref{cdp:eq:C.HQ.2}.

\noindent\textbf{Step 3.} Write $Z_i$ for a smooth normalization resolution and
$M_i$ for its minimal model. The case $b_2(M_i)=0$
already has Hopf minimal model by \ref{cdp:st:C.2.UN.17}. Every
minimal class VII surface with $b_2>0$ has algebraic
dimension zero \cite[p.~256]{DOT2000}. It therefore
suffices to treat $a(M_i)=0$, by an argument which
also covers the zero-Betti case. The field comparison
in (C.CF.9) gives
$a(E)=0$, and \ref{cdp:st:C.2.UN.16} says that the unique smooth
minimal model $E_0$ is a torus or a Hopf surface.

\noindent\textbf{Step 4.} Every element of $G$ induces a birational automorphism
of $E_0$. On a nonruled minimal compact surface this
is holomorphic, by uniqueness of the minimal model.
Thus the contraction $\rho:E\to E_0$ is equivariant:
the identity $\rho g=g_0\rho$ holds on the common
open isomorphism and hence everywhere. It descends
to a proper bimeromorphic map
\[
                    E/G\longrightarrow E_0/G.
 \tag{C.HQ.3}\label{cdp:eq:C.HQ.3}
\]
The actions on $E_0$ are faithful after removing
their kernel, which does not change the quotient.
Resolving the graph of \eqref{cdp:eq:C.HQ.2}--\eqref{cdp:eq:C.HQ.3} gives
actual proper bimeromorphic comparisons between
$Z_i$ and a resolution of $E_0/G$.

\noindent\textbf{Step 5.} We exclude the torus alternative by its first
cohomology. For a finite holomorphic quotient
$q:A\to A/G$, proper base change gives
$R^jq_*\mathbb Q=0$ for $j>0$, and averaging gives
$(q_*\mathbb Q)^G=\mathbb Q$. Taking invariants is
exact over $\mathbb Q$, so the actual pullback map
identifies
\[
                     H^1(A/G,\mathbb Q)
                          \simeq H^1(A,\mathbb Q)^G.
 \tag{C.HQ.4}\label{cdp:eq:C.HQ.4}
\]
For a complex torus the holomorphic action preserves
$H^{1,0}\oplus H^{0,1}$ and conjugation interchanges
the two invariant summands. Hence the dimension in
\eqref{cdp:eq:C.HQ.4} is even. Each singularity of a cyclic surface
quotient has a resolution with a connected chain
of rational curves: after linearization, the
quotient cone is the first quadrant in
$\mathbb Z^2+\frac1n(a,b)\mathbb Z$, and a regular
fan subdivision gives that resolution. Removing
pseudoreflections first only replaces the local
coordinates by their invariant powers. Thus every
exceptional fibre has $H^1=0$ and is connected.
Proper base change and the degree-one part of
Leray show that resolving $A/G$ leaves $H^1$
unchanged. Birational maps between smooth compact
surfaces preserve $H^1$, as follows from their
factorization into point blowups. The comparisons
\eqref{cdp:eq:C.HQ.2}--\eqref{cdp:eq:C.HQ.3} would therefore make $b_1(Z_i)$
even, contradicting $b_1(Z_i)=1$.

\noindent\textbf{Step 6.} It remains to prove that the cyclic quotient of
the Hopf surface $H=E_0$ is again Hopf. We give
the fixed-point and quotient argument explicitly.
Let
\[
            H=(\mathbb C^2\setminus\{0\})/\Gamma.
 \tag{C.HQ.5}\label{cdp:eq:C.HQ.5}
\]
The deck group contains an infinite cyclic
subgroup of finite index generated by a
holomorphic contraction. Equivalently $H$ has
a finite unramified primary Hopf cover; the
normal forms and the finite-cover description
are recalled in \cite[\S1.1]{Nakamura2008VII}
and \cite[Theorem~2.1]{Khaled2024Hopf}.
Every deck transformation extends across the
origin by Hartogs, as does its inverse. These
extensions fix the origin and are biholomorphic.
The same is true for any lift $\widetilde g$
of a holomorphic automorphism $g$ of $H$.

\noindent\textbf{Step 7.} Define the real character
\[
 \delta:\Gamma\longrightarrow\mathbb R,
 \qquad \delta(\gamma)=\log|\det(d\gamma)_0|.
 \tag{C.HQ.6}\label{cdp:eq:C.HQ.6}
\]
It is nonzero, since its value on a contraction
is negative. For any automorphism lift,
conjugacy of derivatives gives
\[
 \delta(\widetilde g\gamma\widetilde g^{-1})
                         =\delta(\gamma).
 \tag{C.HQ.7}\label{cdp:eq:C.HQ.7}
\]
A primary Hopf cover has rational Betti numbers
$(1,1,0,1,1)$: this follows directly from the
covering spectral sequence for
$\mathbb C^2\setminus\{0\}\simeq S^3$ and the
infinite cyclic deck group, whose action on
$H^3(S^3,\mathbb Q)$ preserves orientation.
Transfer from the finite primary cover gives
$b_1(H)\le1$ and $b_2(H)=0$, while \eqref{cdp:eq:C.HQ.6}
gives $b_1(H)\ge1$. Hence $b_1(H)=1$.
Equation \eqref{cdp:eq:C.HQ.7} says that every holomorphic
automorphism acts identically on $H^1(H,\mathbb R)$.
Poincar\'e duality and its complex orientation
then give the identity action on $H^3$ as well.
Its Lefschetz number is therefore
\[
                 L(g)=1-1+0-1+1=0.
 \tag{C.HQ.8}\label{cdp:eq:C.HQ.8}
\]

\noindent\textbf{Step 8.} For a nonidentity finite-order automorphism,
the fixed locus is a disjoint union of smooth
compact curves and isolated points, by finite
holomorphic linearization. Each fixed curve $C$
is smooth. Since $H^2(H,\mathbb Q)=0$, one has
$C^2=K_H\cdot C=0$. The actual adjunction
isomorphism $K_C=(K_H\otimes\mathcal O_H(C))|_C$
therefore gives $2g(C)-2=0$. Thus all the
curve components of the fixed locus have Euler
number zero. For a finite-order
map the topological Lefschetz number equals
the Euler number of its fixed locus. Indeed,
use an equivariant triangulation and barycentric
subdivision; each invariant simplex is then
fixed pointwise, so the alternating trace
on simplicial cochains counts exactly the
fixed simplices. Equation \eqref{cdp:eq:C.HQ.8} consequently
shows that there are no isolated fixed points.

\noindent\textbf{Step 9.} At any point of $H$, the stabilizer in the
effective cyclic group $G$ is cyclic. Its
generator has a fixed curve through that
point and, in linearizing coordinates, acts
by
\[
                (x,y)\longmapsto(x,\zeta y),
 \qquad
 \mathbb C\{x,y\}^{\langle\zeta\rangle}
                    =\mathbb C\{x,y^{|G_x|}\}.
 \tag{C.HQ.9}\label{cdp:eq:C.HQ.9}
\]
Here $G_x$ denotes the stabilizer, and
$\zeta$ has its full order; a finite-order
germ with identity derivative is the identity.
This proves that the quotient $Q=H/G$ is
smooth, with the displayed local quotient
map even at its branch curves.

\noindent\textbf{Step 10.} The invariant-cohomology comparison in
\eqref{cdp:eq:C.HQ.4}, now in every degree and combined
with \eqref{cdp:eq:C.HQ.7}, gives
\[
                         b_1(Q)=1,
                         \qquad b_2(Q)=0.
 \tag{C.HQ.10}\label{cdp:eq:C.HQ.10}
\]
The finite map $H\to Q$ injects every space
of pluricanonical sections by the actual
Jacobian pullback. Since a Hopf surface
has $\kappa=-\infty$, this proves
$\kappa(Q)=-\infty$. Also $Q$ is minimal:
an exceptional curve of self-intersection
$-1$ would have a nonzero real $H^2$ class,
contrary to \eqref{cdp:eq:C.HQ.10}.

\noindent\textbf{Step 11.} There is a curve on $Q$. Explicitly, the
primary Hopf normal form
$(z_1,z_2)\mapsto(\alpha z_1+a z_2^r,\beta z_2)$
contains the elliptic curve given by $z_2=0$;
its finite image in $H$, and then in $Q$,
is a nonempty curve. A positive multiple
of that curve is homologically trivial by
\eqref{cdp:eq:C.HQ.10}. The divisor classification
\cite[Remark~2.7]{DlousskyTeleman2012} applied
to the minimal class VII surface $Q$ now
gives a Hopf surface: its Enoki alternative
has positive second Betti number and is
excluded by \eqref{cdp:eq:C.HQ.10}. We have proved the
cyclic-quotient assertion by actual quotient
maps, including their branch points.

\noindent\textbf{Step 12.} Finally \eqref{cdp:eq:C.HQ.2}--\eqref{cdp:eq:C.HQ.3} make $Z_i$
bimeromorphic to this smooth Hopf surface
$Q$. Uniqueness of the class VII minimal
model proves that $M_i\simeq Q$. This
proves the assertion for every component
in the original, unaltered degeneration.
\end{proof}






\begin{proposition}[The rank-one component fibration descends to the original support]\label{cdp:st:C.2.UN.20}
\StatementRecord{cdp:st:C.2.UN.20}{proved}{Historical events: Claude Fable 5.2 (Anthropic), effort high, 16 November 2026: carried as background without fault found. GPT-6 Astra (OpenAI), ChatGPT Pro mode, effort Pro, S6\_09\_PAPERS, 17 November 2026: targeted proof check. Also checked in this writing. Pre-continuation writing events: GPT-6 Astra (OpenAI), ChatGPT Pro mode, effort Pro, S6\_10\_PAPERS, 18 November 2026: targeted proof check; used and rechecked in this run. Continuation mathematical and source checks: GPT-6 Astra (OpenAI), ChatGPT Pro mode, effort Pro, S6\_10\_PAPERS continuation, 19 November 2026: independent re-proof. Continuation metadata events: GPT-6 Astra (OpenAI), ChatGPT Pro mode, effort Pro, S6\_10\_PAPERS continuation, 19 November 2026: used and rechecked in this run; status-vocabulary reconciliation; metadata reconciliation only; no new mathematical check. \emph{Pre-continuation recorded proof scope.} The stabilizer preserves the valuation lattice semilinearly and induces the exact Mobius action on the meromorphic field ratio. Equivariant graph resolution and Stein factorization yield the quotient diagram. The general quotient fibre is exactly G\_b modulo the kernel of the base action, is connected and smooth, and shows that the quotient base already is the Stein base. Adjunction excludes rational fibres for nonnegative Kodaira dimension}{the hypotheses, notation and definitions stated here; \cref{cdp:st:C.2.UN.16}; \cref{cdp:st:C.2.UN.19}; \cref{cdp:eq:C.FD.6}}

In the setting and with the actual component cover of \ref{cdp:st:C.2.UN.16},
the generic-rank-one branch gives a fibration on a resolution of
the original normalization. More precisely, for the stabilizer
$G$ of the chosen covering component $E$, there are proper
birational maps and quotient maps in a commutative diagram
\[
\begin{array}{ccc}
 \widehat E&\xrightarrow{\ q\ }&B\\
 \downarrow&&\downarrow\\
 \widehat E/G&\longrightarrow&B/G,
\end{array}                                                    \tag{C.FD.1}\label{cdp:eq:C.FD.1}
\]
where $B$ is a smooth compact curve, $q$ has connected general
fibre, and $\widehat E/G$ is bimeromorphic to $S_i^\nu$.
Resolving this last quotient and taking its Stein factor gives
an actual proper fibration
\[
        Z'_i\longrightarrow B/G
        \quad\hbox{with connected rational or elliptic general fibre}.
                                                               \tag{C.FD.2}\label{cdp:eq:C.FD.2}
\]
The base in \eqref{cdp:eq:C.FD.2} is already the Stein base; in particular
there is no additional finite cover of $B/G$ implicit in this
description. For every original component with nonnegative
Kodaira dimension, its general fibre in \eqref{cdp:eq:C.FD.2} is elliptic.
The same conclusion holds for every original component that
is not uniruled.
\end{proposition}


\begin{proof}
\noindent\textbf{Step 1.} Use a resolution equivariant for the finite cyclic group of the
normalized base change, as in \ref{cdp:st:C.2.UN.19}. Its stabilizer $G$
preserves $E$. Recall the actual modules
\[
 M=(f'_*\mathcal T_{U'/\Delta_s})_0,
 \qquad
 L=\{v\in M\otimes_{\mathbb C\{s\}}\mathbb C(\{s\}):
                  \operatorname{ord}_E(v)\ge0\}.
                                                               \tag{C.FD.3}\label{cdp:eq:C.FD.3}
\]
An element $g\in G$ covers $s\mapsto\zeta_gs$. Pushforward of
vector fields is therefore a semilinear automorphism of $M$
and of its meromorphic localization. Since $g(E)=E$, it
preserves the valuation defining $L$. It also preserves $sL$,
because $s$ changes by a nonzero constant. Reduction gives
a complex-linear representation on $L/sL$, and the restriction
map of (C.CF.6) is equivariant:
\[
 g_*\bigl(v|_E\bigr)=(g_*v)|_E.
                                                               \tag{C.FD.4}\label{cdp:eq:C.FD.4}
\]
Consequently the two-dimensional space $\mathfrak a_E$ and
its rank-one meromorphic tangent distribution are preserved
by $G$.

\noindent\textbf{Step 2.} For a basis $V_1,V_2$ write $V_2=hV_1$. The function $h$ is
nonconstant, and both fields annihilate it. If the matrix of
$g_*$ in this basis is
\[
       g_*V_1=aV_1+bV_2,\qquad
       g_*V_2=cV_1+dV_2,
\]
then the actual coefficient comparison gives
\[
              h\circ g^{-1}=\frac{c+dh}{a+bh}.
                                                               \tag{C.FD.5}\label{cdp:eq:C.FD.5}
\]
This identity holds where the displayed denominator is nonzero
and hence as an identity of meromorphic functions. It supplies
the finite projective action on the target of
$h:E\dashrightarrow\mathbb P^1$. Resolve its graph
equivariantly, and then take the Stein factor. Uniqueness of
the Stein factor lifts the action on $\mathbb P^1$ to an action
on the smooth compact curve $B$, and gives an equivariant map
$q:\widehat E\to B$. This is the same connected-leaf
fibration as in \ref{cdp:st:C.2.UN.16}. Its general fibre is rational or
elliptic, by the semicomplete one-dimensional field calculation
there. Taking the finite quotient produces \eqref{cdp:eq:C.FD.1}.

\noindent\textbf{Step 3.} We spell out the general quotient fibre. Let
$K\subset G$ be the kernel of the action on $B$. Outside
the finitely many fixed points of the nonidentity elements
of $G/K$, a point $b\in B$ has stabilizer exactly $K$.
The inverse image of its orbit in $B/G$ is the disjoint union
of the fibres $q^{-1}(gb)$, and $G$ acts transitively on these
fibres. Hence the quotient fibre is exactly
\[
                       q^{-1}(b)/K.
                                                               \tag{C.FD.6}\label{cdp:eq:C.FD.6}
\]
It is connected. A finite quotient of a smooth compact curve
is smooth, as follows locally from the invariant coordinate
$z\mapsto z^r$ at a cyclic stabilizer. Riemann--Hurwitz
shows that the genus of \eqref{cdp:eq:C.FD.6} is at most the genus of
$q^{-1}(b)$, and is therefore zero or one.

\noindent\textbf{Step 4.} For clarity, the generic fibre remains unchanged when the
surface quotient is resolved. A stabilizer of a point in a
general fibre is contained in $K$. It fixes the base coordinate
and, after finite holomorphic linearization, acts by
$(x,y)\mapsto(x,\xi y)$. Its invariant ring is
$\mathbb C\{x,y^r\}$, so the surface quotient is smooth
along a general fibre. A resolution can be an isomorphism
there. The resulting map to $B/G$ has connected general
fibre by \eqref{cdp:eq:C.FD.6}. Its Stein factor is a finite map of normal
compact curves of degree one, hence an isomorphism. This
proves the assertion about the base in \eqref{cdp:eq:C.FD.2}.

\noindent\textbf{Step 5.} The quotient comparisons (C.HQ.2)--(C.HQ.3), or directly
the generic $G$-torsor in (C.HQ.1), identify
$\widehat E/G$ bimeromorphically with $S_i^\nu$.
Resolving their graph gives the proper birational comparison
of \eqref{cdp:eq:C.FD.2} with a smooth normalization resolution of the
original component. Thus the fibration has been constructed
on the original birational component, through the specified
finite quotient, rather than inferred from images of a
covering family of curves.

\noindent\textbf{Step 6.} Finally a rational general fibre $F$ of \eqref{cdp:eq:C.FD.2} has trivial
normal line in $Z'_i$, and adjunction gives
$K_{Z'_i}|_F=K_F=\mathcal O_{\mathbb P^1}(-2)$.
Every positive pluricanonical section restricts to zero on
every general fibre and therefore vanishes identically.
This makes $\kappa(Z'_i)=-\infty$, so the nonnegative
Kodaira-dimension cases have elliptic general fibre.
More generally the rational fibration itself is a family
of rational curves covering the surface; it is precisely
an actual ruling on a proper birational model and is
excluded when the original component is not uniruled.
\end{proof}






\begin{proposition}[A component cover with a rank-two translation frame has coherent Euler characteristic zero]\label{cdp:st:C.2.UN.22}
\StatementRecord{cdp:st:C.2.UN.22}{proved}{Historical events: Claude Fable 5.2 (Anthropic), effort high, 16 November 2026: carried as background without fault found. GPT-6 Astra (OpenAI), ChatGPT Pro mode, effort Pro, S6\_09\_PAPERS, 17 November 2026: targeted proof check. Also checked in this writing. Pre-continuation writing events: GPT-6 Astra (OpenAI), ChatGPT Pro mode, effort Pro, S6\_10\_PAPERS, 18 November 2026: targeted proof check; used and rechecked in this run. Continuation mathematical and source checks: GPT-6 Astra (OpenAI), ChatGPT Pro mode, effort Pro, S6\_10\_PAPERS continuation, 19 November 2026: targeted proof check; used and rechecked in this run. Continuation metadata events: GPT-6 Astra (OpenAI), ChatGPT Pro mode, effort Pro, S6\_10\_PAPERS continuation, 19 November 2026: used and rechecked in this run; status-vocabulary reconciliation; metadata reconciliation only; no new mathematical check. \emph{Pre-continuation recorded proof scope.} The time-form pole obstruction, exact projection of the commuting field, period-lattice transport and root-cover removal of multiple fibres were re-derived. Each algebraic-dimension case is treated through its actual meromorphic field and quotient fibration. The dense-leaf case uses the explicit finite isotrivial product cover to rule out a horizontal pole. Kodaira1963II Theorem6.2 pp564-566 and section8(iv)(11) p585 were opened in reproduced original journal text; Ulmer-Urzua section2.3 pp5-6 supplies a separately opened primary author-preprint table}{the hypotheses, notation and definitions stated here; \cref{cdp:st:C.2.UN.16}; \cite[Theorem~3.3]{Guillot2024}; \cite[Theorem~6.2, pp.~564--566; Section~8(iv)(11), p.~585]{Kodaira1963II}; \citet[Section~2.3, pp.~5--6]{UlmerUrzua2021}}

Let $X$ be a smooth compact complex surface with
$\kappa(X)\ge0$. Suppose that it carries a two-dimensional
complex vector space $\mathfrak a$ of commuting meromorphic
vector fields, of generic rank two, every member of which
is semicomplete on its holomorphic locus. Then
\[
                         \chi(\mathcal O_X)=0.
                                                               \tag{C.RF.5}\label{cdp:eq:C.RF.5}
\]
More precisely, $X$ has a two-torus birational model,
or it has an actual relatively minimal elliptic fibration
on a proper birational model, over a smooth compact curve
of genus at most one, with identity integral $H_1$ monodromy
and with every fibre a positive multiple of a smooth elliptic
curve. Its smooth fibres are torsors under one fixed elliptic
curve. All actual multiple fibres and their multiplicities
are retained in this conclusion.

This applies in particular to every covering component
$E$ of \ref{cdp:st:C.2.UN.16} with $\kappa(E)\ge0$ in its
generic-rank-two branch. Such a covering component cannot
have a K3 or Enriques minimal model. The assertion concerns
the actual specified root-cover component; quotients of that
component are compared separately.
\end{proposition}



\begin{proof}[The elliptic-first-integral step in \ref{cdp:st:C.2.UN.22}]
\noindent\textbf{Step 1.} Let $X$ be a smooth compact complex surface and let $V,Y$ be
commuting meromorphic vector fields with $V\wedge Y\ne0$.
Assume that both fields are semicomplete on their holomorphic
loci and that an actual resolution of a meromorphic first
integral of $V$, followed by its Stein factorization, gives
a map
\[
 q:\widehat X\longrightarrow B
 \tag{C.RF.1}\label{cdp:eq:C.RF.1}
\]
with connected elliptic general fibre.  The proof below does
not require $X$ to be K\"ahler.  We use the meromorphic lifts
of the two fields on $\widehat X$; their brackets still vanish
on the common open isomorphism and hence everywhere, and
their semicompleteness extends over their holomorphic loci
by \cite[Theorem~3.3]{Guillot2024}.

\noindent\textbf{Step 2.} First, $V$ has no horizontal polar curve.  On a general
smooth fibre $F$ of $q$, its restriction is a nonzero
semicomplete meromorphic field on the compact elliptic
curve.  The time-form argument in \ref{cdp:st:C.2.UN.16} excludes poles
on $F$.  Thus this restriction is holomorphic and nowhere
zero, and its flow acts by all translations of $F$.
The commutation relation implies
$(\Phi_V^a)_*Y=Y$ as an identity of meromorphic fields.
Consequently the intersection of the polar divisor of $Y$
with $F$ is invariant under all translations of $F$.
It is a finite set, so it is empty.  This proves that
$Y$ also has no horizontal polar curve.

\noindent\textbf{Step 3.} The field $Y$ projects to a nonzero meromorphic field $Z$
on $B$.  In fact, for each local meromorphic function $g$
on $B$,
\[
 V\bigl(Y(q^*g)\bigr)
   =Y\bigl(V(q^*g)\bigr)=0.
\]
The meromorphic function $Y(q^*g)$ is therefore constant
on the connected general fibres; proper Stein descent
makes it the pullback of a meromorphic function on $B$.
These functions define $Z$.  Were $Z$ zero, both fields
would be tangent to the one-dimensional fibres, contrary
to $V\wedge Y\ne0$.

\noindent\textbf{Step 4.} The projected field has no pole.  To verify this without
an unproved descent assertion about semicompleteness,
suppose that $Z$ has a pole of order $r\ge1$ at $b$.
The primitive of its time form has a zero of order $r+1$
there.  Its inverse solution consequently has distinct
continuations after a loop about the corresponding time
value.  Choose that loop and its lifted base paths in a
small punctured disc about $b$ avoiding all other bad
values.  The inverse image of the compact union of these
paths is compact.  Both $q$ and $Y$ are holomorphic there,
and $Y$ has no pole there, since all its polar curves are
vertical and the disc was chosen accordingly.  The usual
continuation theorem for holomorphic ordinary differential
equations therefore lifts the whole time loop to a
solution of $Y$.  Its endpoint has a different image in
$B$.  This contradicts semicompleteness of $Y$.
Hence $Z$ is a nonzero holomorphic field on the compact
curve $B$, and $g(B)\le1$.

\noindent\textbf{Step 5.} Remove finitely many values so that $q$ is a smooth
elliptic fibration, $V,Y$ are holomorphic, $V$ is nowhere
zero vertically, and $Z$ is nonzero.  On this open base
$B^\circ$, let $\alpha_b$ be the time form of $V|_{F_b}$.
The flow of $Y$, defined for a uniform small time along
each compact fibre, carries whole fibres biholomorphically
to whole fibres and carries $V$ to itself.  It therefore
preserves their time forms.  Integration gives the exact
horizontal map from the integral homology local system
\[
 \mathcal H_1(q)\big|_{B^\circ}
   \longrightarrow\underline{\mathbb C},
 \qquad
 \gamma\longmapsto\int_\gamma\alpha_b.
 \tag{C.RF.2}\label{cdp:eq:C.RF.2}
\]
The map is injective on each stalk: the periods of a
nonzero holomorphic one-form on an elliptic curve are a
rank-two lattice in $\mathbb C$.  The preceding flow
comparison proves that its values are locally constant.
Analytic continuation of any integral cycle around a
loop preserves its complex period; injectivity then
proves that every monodromy operator on $H_1(F_b,\mathbb Z)$
is the identity.  The image lattice is fixed, say
$\Lambda\subset\mathbb C$, and the fibres over
$B^\circ$ are torsors under the fixed elliptic curve
$\mathbb C/\Lambda$.

\noindent\textbf{Step 6.} Contract the vertical exceptional curves of the first
kind to obtain the actual relatively minimal elliptic
map
\[
 \widehat X\xrightarrow{\ \mu\ }X_0
       \xrightarrow{\ q_0\ }B.
 \tag{C.RF.3}\label{cdp:eq:C.RF.3}
\]
These contractions preserve the preceding integral
monodromy representation.  Kodaira's classification
\cite[Theorem~6.2, pp.~564--566; Section~8(iv)(11), p.~585]{Kodaira1963II}, together with the explicit positive-loop matrices in \citet[Section~2.3, pp.~5--6]{UlmerUrzua2021},
now shows that every fibre has smooth elliptic support.
Here is the comparison for multiplicities, to specify
the use of the table.  The only multiple types in that
classification are $mI_b$.  Locally write such a full
fibre as $t=u h^m$, where $h=0$ is its reduced Cartier
support and $u$ is a unit.  In the normalization of the
actual base change $t=s^m$, the charts
$w^m=u$, $s=wh$ are \emph{\'etale} over the original
smooth surface.  The covering family has a reduced
central fibre and its linear monodromy is the $m$th
power of the original one, hence the identity.  It
remains relatively minimal: every central component
downstairs has canonical degree zero by the displayed
Kodaira types; the \emph{\'etale} canonical pullback
gives the same zero degree on every central prime
curve upstairs, and adjunction excludes an exceptional
rational curve of self-intersection $-1$ upstairs.
Among the reduced relatively minimal types in the
cited table, identity monodromy is the smooth type
$I_0$.  Hence the original support is smooth as well.
The argument applies to every connected covering
component and makes no connectedness assumption on
the root cover.

\noindent\textbf{Step 7.} In particular, for pairwise distinct base values the
fibre supports $C_b$ are disjoint smooth elliptic
curves, and $C_b^2=0$, since some positive multiple
of $C_b$ is a full fibre.  The Euler characteristic
is zero by additivity over the smooth base stratum
and the finitely many special fibres.  To compute
the canonical square, choose a nonzero meromorphic
one-form $\eta$ on $B$.  The meromorphic field $V$
descends through $\mu$, and at the generic point
there is a unique meromorphic two-form $\Omega$
satisfying
\[
 \iota_V\Omega=q_0^*\eta.
 \tag{C.RF.4}\label{cdp:eq:C.RF.4}
\]
Both sides are interpreted as actual meromorphic
differentials.  Since $V$ is tangent to $q_0$, the
right-hand side annihilates $V$, which proves
existence and uniqueness of $\Omega$ by a local
tangent-frame calculation.  Its divisor is
vertical: on a general compact fibre $V$ is
holomorphic and nowhere zero, and $\eta$ is
holomorphic and nonzero at its base value.
Thus a canonical divisor is a finite integral
combination of the disjoint square-zero curves
$C_b$, and $K_{X_0}^2=0$.  The surface
Riemann--Roch formula gives
$\chi(\mathcal O_{X_0})=(K_{X_0}^2+e(X_0))/12=0$.
Proper birational invariance of the coherent
Euler characteristic under the displayed point
blowups and blowdowns yields
$\chi(\mathcal O_X)=0$.

\noindent\textbf{Step 8.} No vanishing of $Z$ at a multiple value is asserted.
Indeed, a logarithmic elliptic translation can absorb
the cyclic root monodromy.  Multiple smooth fibres
are included in the conclusion.
\end{proof}

\begin{proof}[Completion of \ref{cdp:st:C.2.UN.22}]
\noindent\textbf{Step 1.} We first treat the nonprojective cases by actual
meromorphic maps. If $a(X)=1$, choose a nonconstant
$h\in\mathcal M(X)$, resolve its graph, and take the
Stein factor $q:\widehat X\to B$. Every
$g\in\mathcal M(X)$ is algebraic over $\mathbb C(h)$.
On a connected general fibre of $q$ it therefore takes
finitely many values and is constant. Proper meromorphic
descent consequently gives the exact field identity
\[
                       \mathcal M(X)=q^*\mathcal M(B).
                                                               \tag{C.RF.9}\label{cdp:eq:C.RF.9}
\]
For example, this descent can be seen by mapping the
proper graph of $g$ to $B\times\mathbb P^1$; its image
is the graph of the descended meromorphic function.
Each field in $\mathfrak a$ differentiates this field
of functions and therefore induces a meromorphic vector
field on $B$. The induced fields commute. They cannot
both vanish, since that would make both original fields
tangent to the one-dimensional fibres and violate
generic rank two. On a curve any two commuting fields
with one nonzero are constant multiples, by \eqref{cdp:eq:C.RF.8}.
A nonzero constant linear combination in $\mathfrak a$
is thus vertical. On the compact general fibre its
semicompleteness excludes poles, so the genus is at most
one; genus zero again contradicts $\kappa(X)\ge0$ by
adjunction. This gives an elliptic fibration and reduces
the case $a(X)=1$ to \eqref{cdp:eq:C.RF.1}--\eqref{cdp:eq:C.RF.4}.

\noindent\textbf{Step 2.} If $a(X)=0$, apply the final argument of \ref{cdp:st:C.2.UN.16}.
That argument uses only the given two-dimensional
space of semicomplete fields: the fibration alternative
of \cite[Theorem~9.2]{Guillot2024} is impossible in
algebraic dimension zero, and the fields on their
respective holomorphic birational models descend to
two independent holomorphic fields on the same unique
nonruled minimal model. The compact-surface field
classification and dimension bounds used there give
a two-torus or a Hopf surface. Kodaira dimension
excludes the latter. This gives the two-torus alternative
and \eqref{cdp:eq:C.RF.5}.

\noindent\textbf{Step 3.} It remains to treat $a(X)=2$. A compact complex surface
of algebraic dimension two has a smooth projective
birational model. Passing through the common resolution
preserves the bracket, generic rank, and coherent Euler
characteristic; semicompleteness on the full holomorphic
locus follows from the same closedness comparison as in
\ref{cdp:st:C.2.UN.16}. We may therefore work on that projective model.
Choose a basis $V_1,V_2$ of $\mathfrak a$. The nonzero
meromorphic field $V_1$ has either a leaf with Zariski-dense
image or only algebraic leaves. In the latter case the
constructible integral-curve incidence argument in
\ref{cdp:st:C.2.UN.15} gives an actual resolution of its first integral
and a Stein fibration. The semicomplete time-form calculation
on its compact general fibre gives genus at most one.
A rational general fibre would force every pluricanonical
section to vanish, by adjunction as in \ref{cdp:st:C.2.UN.20}, contrary
to $\kappa(X)\ge0$. Thus its general fibre is elliptic,
and the preceding proof applies to $V_1,V_2$.

\noindent\textbf{Step 4.} It remains to treat a Zariski-dense leaf of $V_1$.
Semicompleteness supplies the univalent maximal solution
required by \cite[Main Theorem and the clarification immediately
following it, p.~2]{Guillot2019}. On an actual smooth
birational model, either the surface is Abelian, when
\eqref{cdp:eq:C.RF.5} follows at once, or there is a fibration
\[
                         q:X'\longrightarrow B
                                                               \tag{C.RF.6}\label{cdp:eq:C.RF.6}
\]
over a rational or elliptic curve, with rational or elliptic
general fibre, which $V_1$ preserves and whose fibres
contain every polar curve of $V_1$. Again the rational-fibre
case contradicts $\kappa(X')\ge0$. Thus \eqref{cdp:eq:C.RF.6} is an
elliptic fibration. The projection $Z_1=q_*V_1$ is nonzero:
otherwise every leaf of $V_1$ would lie in a compact fibre,
contrary to the selected dense leaf. Both lifted fields
remain semicomplete on their holomorphic loci by the same
birational-open comparison and closedness argument used
in \ref{cdp:st:C.2.UN.16}.

\noindent\textbf{Step 5.} We prove that the polar divisor of $V_2$ is also vertical.
At a generic point where $V_1$ is holomorphic and nonzero,
the identity $[V_1,V_2]=0$ says that its local flow preserves
$V_2$ as a meromorphic tensor and therefore preserves its
polar divisor. Equivalently, in a flow coordinate
$V_1=\partial_z$, the coefficients of $V_2$ are independent
of $z$. Thus a horizontal polar component would be a
compact invariant multisection of the foliation of $V_1$.
We now rule out such a multisection with the actual finite
maps needed for this conclusion.

\noindent\textbf{Step 6.} Shrink to a nonempty open base $B^\circ$ on which the
fibration is smooth, $V_1$ is holomorphic on each whole
fibre, and $Z_1$ is nonzero. Its flow is defined uniformly
for small time on any one compact fibre, and maps that
whole fibre biholomorphically to a neighboring fibre.
Consequently the elliptic curves are isotrivial over
$B^\circ$. Suppose that $C\subset X'$ were a compact
horizontal invariant curve. Its normalization $\widetilde C$
maps finitely to $B$. In the normalized fibre product with
$\widetilde C$, the diagonal curve supplies an actual
section; it is horizontal for the lifted field. Resolve
this fibre product, taking a resolution which is an
isomorphism above the smooth open locus.

\noindent\textbf{Step 7.} Fix an elliptic curve $A$ isomorphic to a smooth fibre.
Over a smaller open subset of $\widetilde C$, the
origin-preserving isomorphisms from $A$ to the fibres
form a finite \'etale torsor with structure group
$\operatorname{Aut}(A,0)$. This is a concrete finite
cover: in a Weierstrass equation with constant $j$, the
change of coordinates is obtained by adjoining the
scalars $\lambda$ with
$g_2=\lambda^4g_{2,A}$ and
$g_3=\lambda^6g_{3,A}$, omitting the equation with
zero constant coefficient when $j=0$ or $1728$.
Its ambiguity is exactly the finite origin-preserving
automorphism group. Normalize the compact curve
$\widetilde C$ in a connected component of this finite
cover. On the resulting compact curve $B'$ there is
an actual product isomorphism over a nonempty open set,
and resolving its graph gives proper maps
\[
 \begin{array}{ccc}
 \widetilde X'&\longrightarrow&X'\\
 \downarrow&&\downarrow q\\
 B'&\longrightarrow&B,
 \end{array}
 \qquad
 \widetilde X'\dashrightarrow A\times B'
 \quad\hbox{birational}.
                                                               \tag{C.RF.7}\label{cdp:eq:C.RF.7}
\]
Here all curves and surfaces are projective, so the
isomorphism over the algebraic open set defines the
displayed rational map; its graph closure is proper.
The section is carried to the zero section.

\noindent\textbf{Step 8.} On the product open set the lifted field has the form
$a(b)\partial_b+c(b)\partial_u$, where $\partial_u$
is a fixed translation field on $A$. Indeed, its vertical
coefficient is holomorphic on each whole compact elliptic
fibre, so is constant on that fibre. Tangency to the zero
section gives $c=0$. Every local leaf therefore lies in
$\{u\}\times B'$. The compact closure of this section
in the proper graph model maps to an algebraic invariant
curve in $X'$. Through a general point this contains
the entire leaf by uniqueness of analytic continuation.
This contradicts the dense leaf of $V_1$. This argument
also explains the invariant-multisection step in
\cite[\S6.3]{Guillot2019}; no extension of an arbitrary
transcendental horizontal trivialization is being used.

\noindent\textbf{Step 9.} It follows that $V_2$ has no horizontal polar curve.
Hence it is holomorphic on every fibre over a suitable
nonempty open subset of $B^\circ$. Its normal component
there is a holomorphic function on a compact elliptic
curve and is constant. It consequently projects to a
meromorphic field $Z_2$ on $B$; the coefficients descend
through the proper connected-fibre map, or can be
recovered by the trace along any finite multisection.
The bracket identity gives $[Z_1,Z_2]=0$. For nonzero
$Z_1$ on a connected curve this is the precise equation
\[
              Z_1(Z_2/Z_1)=0,
              \qquad Z_2=cZ_1\quad(c\in\mathbb C).
                                                               \tag{C.RF.8}\label{cdp:eq:C.RF.8}
\]
Thus $W=V_2-cV_1$ is vertical. It is nonzero and
$W\wedge V_1\ne0$ by the generic-rank-two hypothesis.
It is semicomplete because it is a constant member of
$\mathfrak a$, and it commutes with $V_1$. Apply the
preceding elliptic-first-integral proof to $W,V_1$
and \eqref{cdp:eq:C.RF.6}. It gives the stated relatively minimal
model, fixed period lattice, identity integral monodromy,
and \eqref{cdp:eq:C.RF.5}.

\noindent\textbf{Step 10.} Finally $\chi(\mathcal O)=2$ on a K3 surface and
$\chi(\mathcal O)=1$ on an Enriques surface. These
values and coherent Euler characteristic are preserved
under the actual point blowups and blowdowns comparing
smooth birational surface models. They exclude those
two minimal models for the rank-two covering component.
\end{proof}






\begin{proposition}[A cyclic torus quotient with nonnegative Kodaira dimension is a torus, bielliptic, or K3]\label{cdp:st:C.2.UN.23}
\StatementRecord{cdp:st:C.2.UN.23}{proved}{Historical events: Claude Fable 5.2 (Anthropic), effort high, 16 November 2026: carried as background without fault found. GPT-6 Astra (OpenAI), ChatGPT Pro mode, effort Pro, S6\_09\_PAPERS, 17 November 2026: targeted proof check. Also checked in this writing. Pre-continuation writing events: GPT-6 Astra (OpenAI), ChatGPT Pro mode, effort Pro, S6\_10\_PAPERS, 18 November 2026: targeted proof check; used and rechecked in this run. Continuation mathematical and source checks: GPT-6 Astra (OpenAI), ChatGPT Pro mode, effort Pro, S6\_10\_PAPERS continuation, 19 November 2026: independent re-proof. Continuation metadata events: GPT-6 Astra (OpenAI), ChatGPT Pro mode, effort Pro, S6\_10\_PAPERS continuation, 19 November 2026: used and rechecked in this run; status-vocabulary reconciliation; metadata reconciliation only; no new mathematical check. \emph{Pre-continuation recorded proof scope.} The affine action, translated fixed point, exact quotient lattice and toric pluricanonical orders were re-derived. A non-symplectic cyclic quotient forces the displayed pole or branch zero. The symplectic case has crepant height-one rays; the actual fundamental-group quotient by all fixed-point stabilizers is trivial, yielding the K3 resolution. The equivariant minimal contraction transfers this conclusion to the covering component}{the hypotheses, notation and definitions stated here; \cref{cdp:st:C.2.UN.16}; \cref{cdp:eq:C.CQ.3}; \cref{cdp:eq:C.CQ.4}; \cref{cdp:eq:C.CQ.6}; \cref{cdp:eq:C.CQ.9}; \cref{cdp:eq:C.CQ.11}; \cref{cdp:st:C.2.UN.20}}

Let $A$ be a compact complex two-torus, let $G$ be a finite
cyclic group of holomorphic automorphisms acting effectively
on $A$, and let $Z\to A/G$ be a smooth resolution.
Every positive plurigenus of $Z$ is at most one. There are
the following exhaustive alternatives for a generator
$g(z)=Lz+b$:
\[
\begin{array}{c|c}
\text{linear part and action}&\text{minimal resolution type}\\ \hline
L=I&\text{a two-torus}\\
\dim_{\C}\ker(L-I)=1,\ G\text{ acts freely}
   &\text{a bielliptic surface}\\
\ker(L-I)=0,\ \det L=1&\text{a K3 surface}\\
\text{all other cases}&P_m(Z)=0\text{ for every }m>0 .
\end{array}                                                    \tag{C.CQ.1}\label{cdp:eq:C.CQ.1}
\]
In particular an Enriques surface is not bimeromorphic to
a cyclic quotient of a complex two-torus.

For the actual component cover in \ref{cdp:st:C.2.UN.16}, suppose its
minimal model is a two-torus. The quotient by its cyclic
stabilizer, followed by the specified proper birational
comparison to the original normalization, satisfies
\eqref{cdp:eq:C.CQ.1}. Thus an original component of nonnegative
Kodaira dimension in this branch has minimal model a
torus, a bielliptic surface, or a K3 surface.
\end{proposition}


\begin{proof}
\noindent\textbf{Step 1.} The derivative of a holomorphic torus automorphism is
constant in the translation trivialization: its entries
are holomorphic functions on a compact torus. Lifting to
the universal cover consequently writes the automorphism
as $z\mapsto Lz+b$, with $L\Lambda=\Lambda$ for the
lattice $\Lambda$. For a finite-order automorphism, $L$
is diagonalizable with roots of unity as eigenvalues.

\noindent\textbf{Step 2.} We first specify the canonical comparison used in every
case. Resolve the rational lift from $A$ to $Z$:
\[
\begin{array}{ccc}
A'&\xrightarrow{\ v\ }&Z\\
\downarrow{\scriptstyle u}&&\downarrow\\
A&\longrightarrow&A/G .
\end{array}                                                    \tag{C.CQ.2}\label{cdp:eq:C.CQ.2}
\]
Here $u$ is a sequence of point blowups and $v$ is
generically finite. Pullback of a pluricanonical form
through $v$ is holomorphic. On the open set where $u$
is an isomorphism it gives a pluricanonical form on
$A$; its coefficients extend across the finitely many
missing points by Hartogs. Hence there is an injective
map
\[
 H^0(Z,mK_Z)\longrightarrow H^0(A,mK_A)
                  =\C\,\sigma^{\otimes m},
 \qquad \sigma=dz_1\wedge dz_2 .
                                                               \tag{C.CQ.3}\label{cdp:eq:C.CQ.3}
\]
Injectivity follows on the dense locus of the dominant
map. The image is $G$-invariant, so a nonzero image
also requires $(\det L)^m=1$.

\noindent\textbf{Step 3.} If the quotient map $A\to A/G$ is ramified along a
curve, its generic local form is
$(x,y)\mapsto(x^e,y)$ with $e>1$. The Jacobian
pullback of a holomorphic $m$-canonical form vanishes
to order at least $m(e-1)$ along that curve. This is
incompatible with the nonzero constant form in
\eqref{cdp:eq:C.CQ.3}. We have proved
\[
 \text{a branch curve in }A\longrightarrow A/G
       \quad\Longrightarrow\quad P_m(Z)=0\quad(m>0).
                                                               \tag{C.CQ.4}\label{cdp:eq:C.CQ.4}
\]
This conclusion uses the actual ramification map,
including when the quotient is smooth along the
branch image.

\noindent\textbf{Step 4.} When $L=I$, the group consists of translations. The
action is free and its quotient is the torus obtained
by enlarging $\Lambda$ by the finitely many translation
vectors. This is the first row of \eqref{cdp:eq:C.CQ.1}.

\noindent\textbf{Step 5.} Suppose that $L$ has exactly one eigenvalue equal to
one. Every power of $L$ has that eigenvalue. A
nonidentity element of $G$ with a fixed point therefore
has a fixed curve, unless its derivative is the
identity; in the latter case it is a translation and
has no fixed point. Thus a nonfree action gives
\eqref{cdp:eq:C.CQ.4}. A free action gives a smooth finite
\'etale quotient. Its canonical line has finite order,
since the finite-order character $\det L$ is the
action on $\sigma$; it has no nonzero canonical
form. Averaging over the finite cover gives
\[
 h^1(Z,\OO_Z)=1,\qquad h^0(Z,K_Z)=0,\qquad
                  b_1(Z)=2 .
                                                               \tag{C.CQ.5}\label{cdp:eq:C.CQ.5}
\]
There is also an explicit product-elliptic comparison.
Let
$A_+=\ker(L-I)^0$ and $A_-=\operatorname{im}(L-I)$.
They are elliptic subtori, their tangent spaces are
complementary, and the addition map
\[
                  A_+\times A_-\longrightarrow A
                                                               \tag{C.CQ.6}\label{cdp:eq:C.CQ.6}
\]
is a finite \'etale isogeny. The finite group of lifts
of $G$ to this product acts freely, by translations
on the first factor and affine automorphisms on the
second; its quotient is exactly $A/G$. To see the
lift explicitly, lift the translation vector $b$
through \eqref{cdp:eq:C.CQ.6} and use the linear action
$(1,L|_{A_-})$. Its power equal to the order of $g$
is a deck translation of \eqref{cdp:eq:C.CQ.6}, so the group
generated with the finite deck group is finite.
This proves the stated finite free product quotient,
the usual bielliptic presentation, without inferring
it only from the numerical invariants. In particular
this quotient is projective. This proves the second
row of \eqref{cdp:eq:C.CQ.1}.

\noindent\textbf{Step 6.} It remains to treat $\ker(L-I)=0$. The endomorphism
$I-L$ of $A$ is an isogeny, because its derivative
is invertible. Solve $(I-L)p=b$ and translate the
origin to $p$. The generator becomes the linear
map $L$ and has a fixed point. Its order equals
the order of $L$.

\noindent\textbf{Step 7.} Assume first that there is no branch curve. At the
fixed origin, no nontrivial power of $L$ can have
exactly one eigenvalue one; such a power would fix
a divisor. If the order is $n$, each eigenvalue
of $L$ is therefore a primitive $n$th root.
After changing the cyclic generator, its local
representation is
\[
             (z_1,z_2)\longmapsto
                     (\zeta z_1,\zeta^r z_2),
       \quad 1\le r<n,\quad \gcd(r,n)=1 .
                                                               \tag{C.CQ.7}\label{cdp:eq:C.CQ.7}
\]
We give the precise local pluricanonical test. The
quotient is the affine toric surface for the first
quadrant in
\[
 N=\Z^2+\Z(1/n,r/n),\qquad
 M=\{(a,b)\in\Z^2:a+rb\equiv0\pmod n\}.
                                                               \tag{C.CQ.8}\label{cdp:eq:C.CQ.8}
\]
These are the cocharacter and character lattices:
the invariant monomials are exactly
$z_1^az_2^b$ with $(a,b)\in M$. If
$(\det L)^m=1$, the form in \eqref{cdp:eq:C.CQ.3} descends on
the smooth locus as
\[
 (z_1z_2)^m
       \left(\frac{dz_1}{z_1}\wedge
             \frac{dz_2}{z_2}\right)^{\otimes m}.
\]
Its order on the divisor of a primitive toric ray
$v=(v_1,v_2)\in N$ is
\[
                       m(v_1+v_2-1).
                                                               \tag{C.CQ.9}\label{cdp:eq:C.CQ.9}
\]
Indeed a unimodular toric chart expresses the
monomial's order as its character pairing with $v$,
while the logarithmic volume form has one simple
pole along that divisor. The same formula holds
on a refinement, by its explicit monomial map.

\noindent\textbf{Step 8.} The vector $(1/n,r/n)$ is primitive in $N$ and
can be included in a smooth fan subdivision of
the quadrant. To obtain such a subdivision, first
insert this ray and then insert primitive lattice
rays until every adjacent cone has determinant
one in a lattice basis; the Euclidean algorithm
terminates this two-dimensional procedure.
When $r<n-1$, \eqref{cdp:eq:C.CQ.9} is negative on this ray.
The descended form cannot extend to that smooth
resolution. Smooth resolutions have the same
pluricanonical sections under their actual common
point-blowup resolution, by the Hartogs comparison
already used in \eqref{cdp:eq:C.CQ.3}. Therefore
\[
 \ker(L-I)=0,\quad \det L\ne1
      \quad\Longrightarrow\quad P_m(Z)=0\quad(m>0),
                                                               \tag{C.CQ.10}\label{cdp:eq:C.CQ.10}
\]
also including the branch-curve case by \eqref{cdp:eq:C.CQ.4}.

\noindent\textbf{Step 9.} Finally let $\det L=1$. Every stabilizer acts with
weights $(1,-1)$ after choosing its generator, so
there are no branch curves. For a stabilizer of
order $n$, the rays
\[
 (1,0),\ ((n-1)/n,1/n),\ldots,
                 (1/n,(n-1)/n),\ (0,1)
                                                               \tag{C.CQ.11}\label{cdp:eq:C.CQ.11}
\]
give a smooth subdivision. Consecutive determinants
are $1/n$, the covolume of $N$, and all these rays
have height $v_1+v_2=1$. Formula \eqref{cdp:eq:C.CQ.9} shows
that the invariant two-form extends without a
zero or pole across every exceptional divisor.
Thus the simultaneous resolution $Z_0$ has
$K_{Z_0}\simeq\OO_{Z_0}$. Its exceptional curves
are the chains of rational curves given by
\eqref{cdp:eq:C.CQ.11}.

\noindent\textbf{Step 10.} For clarity, its simple connectedness follows
through the quotient map as well. Remove the
discrete set of points in $\C^2$ fixed by
nontrivial elements of
$\Gamma=\Lambda\rtimes\langle L\rangle$.
The complement is simply connected: a compact
loop bounds a disc in $\R^4$, and a perturbation
relative to its boundary avoids the finitely
many removed points met in a compact region.
It is the universal cover of the smooth
quotient locus. Filling in a quotient point
kills the normal closure of its finite
stabilizer in $\Gamma$, by van Kampen.
For every $\lambda\in\Lambda$, the element
$(\lambda,L)$ has a fixed point, since $I-L$
is invertible. Thus all $(\lambda,L)$ are
killed, including $(0,L)$; their products
kill every translation $(\lambda,I)$ as well.
The resulting fundamental group is trivial.
Replacing a quotient point by the rational
chain in \eqref{cdp:eq:C.CQ.11} imposes the same relation:
its plumbing neighbourhood is simply connected
and has the same punctured boundary as the
quotient ball. Hence $\pi_1(Z_0)=1$.
Together with $K_{Z_0}=\OO$, this identifies
$Z_0$ as a K3 surface. It is minimal, since
adjunction excludes a rational $(-1)$-curve
on a surface with trivial canonical line.
This proves the third row.

\noindent\textbf{Step 11.} An Enriques surface has a nonzero even
pluricanonical form, but has no nonzero
canonical form and has $b_1=0$. None of
the three nonnegative cases in \eqref{cdp:eq:C.CQ.1}
has those properties: the torus and K3
have a canonical form, and the bielliptic
case has $b_1=2$. All the other cases
have no positive pluricanonical form.
This proves the exclusion.

\noindent\textbf{Step 12.} We finish with the component comparison.
The stabilizer of a component in the
cyclic base change of \ref{cdp:st:C.2.UN.16} is cyclic,
and its quotient is properly bimeromorphic
to the original normalization, as in
(C.HQ.1)--(C.HQ.3). If the minimal model
of the covering component $E$ is a torus
$A$, its actual minimal contraction
$E\to A$ is equivariant. One way to
verify the descent is to factor all the
exceptional rational curves: their maps
to the torus are constant, and the
equality $\rho_*\OO_E=\OO_A$ descends
the coordinate functions and hence
every conjugate action. Thus the induced
action on $A$ is holomorphic, and
\[
 \begin{array}{ccc}
 E&\longrightarrow&A\\
 \downarrow&&\downarrow\\
 E/G&\longrightarrow&A/G
 \end{array}
 \qquad
             E/G\dashrightarrow S_i^\nu
                                                               \tag{C.CQ.12}\label{cdp:eq:C.CQ.12}
\]
has a proper birational horizontal quotient
map and the specified proper birational
comparison on the right. Resolving their
graph identifies the pluricanonical
spaces and minimal birational type.
Applying \eqref{cdp:eq:C.CQ.1} gives exactly the last
assertion. This statement concerns this
actual torus-cover branch; the separate
rank-one fibration branch is handled by
the different comparison of \ref{cdp:st:C.2.UN.20}.
\end{proof}



For comparison with the primary literature, the projective
quotient classification is already given by
\cite[Theorem~2.1 and Lemma~3.4]{Yoshihara1995}.
Its Enriques cases have noncyclic holonomy: the Klein
four group or a dihedral group of order eight
\cite[Example~2.3 and Lemma~3.6]{Yoshihara1995}.
The proof above includes nonprojective tori and supplies
the actual cyclic component-cover comparison needed here.




\begin{proposition}[Every compact elliptic surface satisfies the intrinsic rank-one field condition]\label{cdp:st:C.2.UN.26}
\StatementRecord{cdp:st:C.2.UN.26}{proved}{Historical events: Claude Fable 5.2 (Anthropic), effort high, 16 November 2026: carried as background without fault found. GPT-6 Astra (OpenAI), ChatGPT Pro mode, effort Pro, S6\_09\_PAPERS, 17 November 2026: targeted proof check. Also checked in this writing. Pre-continuation writing events: GPT-6 Astra (OpenAI), ChatGPT Pro mode, effort Pro, S6\_10\_PAPERS, 18 November 2026: targeted proof check. Continuation mathematical and source checks: GPT-6 Astra (OpenAI), ChatGPT Pro mode, effort Pro, S6\_10\_PAPERS continuation, 19 November 2026: independent re-proof; used and rechecked in this run. Continuation metadata events: GPT-6 Astra (OpenAI), ChatGPT Pro mode, effort Pro, S6\_10\_PAPERS continuation, 19 November 2026: used and rechecked in this run; status-vocabulary reconciliation; metadata reconciliation only; no new mathematical check. \emph{Pre-continuation recorded proof scope.} The rank-one direct image of the relative canonical line supplies a global meromorphic time form and its uniquely contracted vertical field. All constant combinations are complete on the dense compact elliptic fibres; Guillot's exact closed-univalence theorem gives semicompleteness on each full holomorphic locus. The finite base map has precisely the original connected-fibre Stein factor. The cyclic eigenvector and quotient-base function give the stated equivariant space}{the hypotheses, notation and definitions stated here; \cref{cdp:st:C.2.UN.16}; \cref{cdp:eq:C.FL.3}; \cite[Theorem~3.3]{Guillot2024}}

Let $q:X\to B$ be a proper holomorphic map from a
smooth compact connected complex surface to a smooth
compact curve, with connected elliptic general fibre.
There are commuting meromorphic vector fields $V_1,V_2$
which are linearly independent over $\mathbb C$, have
generic rank one, and whose every constant linear
combination is semicomplete on its full holomorphic
locus. Their ratio has an actual graph resolution and
Stein factor birationally equal to $q$.

If a finite cyclic group $G$ acts holomorphically on
$X$ and preserves $q$, the two-dimensional space can
be chosen $G$-stable. Thus the rank-one field-space
condition of \ref{cdp:st:C.2.UN.16}, even with this equivariance,
is satisfied by every specified compact elliptic
surface and cyclic fibration action. Any additional
restriction on its occurrence as a degeneration
component must use further data of that degeneration.
\end{proposition}



\begin{proof}
\noindent\textbf{Step 1.} The relative canonical line is the actual invertible
sheaf
\[
 \omega_{X/B}=K_X\otimes q^*K_B^{-1},
 \qquad \mathcal E=q_*\omega_{X/B}.
                                                               \tag{C.FL.1}\label{cdp:eq:C.FL.1}
\]
Proper coherence makes $\mathcal E$ coherent. It is
torsion-free: a section of a line bundle annihilated
by a nonzero base function vanishes over the dense
complement of its zero fibre and hence vanishes
identically. On a smooth elliptic fibre the canonical
space has dimension one, so $\mathcal E$ has rank one.
It is therefore a line bundle on $B$.

\noindent\textbf{Step 2.} Choose a nonzero meromorphic section $\sigma$ of
$\mathcal E$. This is an actual global choice:
Riemann--Roch on $B$ gives a nonzero holomorphic
section of $\mathcal E(D)$ for an effective divisor
$D$ of sufficiently large degree, which is the
required meromorphic section of $\mathcal E$.
Evaluation gives a nonzero meromorphic section
\[
 a:\mathcal O_X\dashrightarrow\omega_{X/B}
       \quad\hbox{from}\quad q^*\mathcal E\longrightarrow
                                      \omega_{X/B}.
                                                               \tag{C.FL.2}\label{cdp:eq:C.FL.2}
\]
Here the dashed notation records poles of the section,
not a replacement of the displayed sheaf evaluation.
On the smooth part of $q$, the canonical differential
identification $\omega_{X/B}=\Omega^1_{X/B}$ makes
$a$ an actual relative one-form.

\noindent\textbf{Step 3.} Its dual defines a global meromorphic vertical field
$V$. The coordinate comparison at critical fibres
can be specified as follows. On a coordinate disc
of $B$ with coordinate $z$, set
$\Omega_z=a\otimes q^*(dz)$, a meromorphic section
of $K_X$. Define $V$ by
\[
                   \iota_V\Omega_z=-q^*(dz).
                                                               \tag{C.FL.3}\label{cdp:eq:C.FL.3}
\]
Contraction by a nonzero meromorphic two-form is an
isomorphism between the meromorphic tangent and
cotangent bundles, so this has a unique meromorphic
solution. Under $z'=\phi(z)$ both sides are multiplied
by $q^*\phi'(z)$; the solution is unchanged and
therefore glues. The alternating identity in
\eqref{cdp:eq:C.FL.3} gives $dq(V)=0$. On a regular fibre its
relative time form is $a$, so $a(V)=1$ with the
usual convention $\Omega_z=q^*(dz)\wedge a$.

\noindent\textbf{Step 4.} Remove finitely many base values so that $q$ is
smooth and the evaluated section is holomorphic
and nonzero on every remaining fibre. Such a
removal is possible by generic base change for
$q_*\omega_{X/B}$ and the finite zero and pole
divisors of $\sigma$. A nonzero canonical form
on a compact elliptic curve is nowhere zero.
Thus $V$ restricts on each of these fibres to
a nonzero holomorphic translation field and is
complete there.

\noindent\textbf{Step 5.} Choose any nonconstant meromorphic function
$h$ on $B$, and put
\[
              V_1=V,\qquad V_2=(q^*h)V.
                                                               \tag{C.FL.4}\label{cdp:eq:C.FL.4}
\]
Their bracket vanishes because $V(q^*h)=0$.
They are independent over $\mathbb C$ since $h$
is nonconstant, and have generic rank one by
the formula. For $c,d\in\mathbb C$, the field
$W=(c+dq^*h)V$ is complete on the dense union
of the preceding compact smooth fibres away
from the poles of $h$. Its coefficient is
constant on each such fibre; zero coefficients
give constant complete solutions. Apply
\cite[Theorem~3.3]{Guillot2024} to the complex
manifold obtained by deleting the actual polar
divisor of $W$. The univalence locus is closed
there and contains this dense union. It is
therefore the whole holomorphic locus, proving
the asserted semicompleteness for every $W$.

\noindent\textbf{Step 6.} The ratio is $q^*h$, so the resulting map is
the actual composite $X\xrightarrow q B
\xrightarrow h\mathbb P^1$. The second arrow
is finite. Its Stein factor is $q$ itself:
$q_*\mathcal O_X=\mathcal O_B$, and $B$ is
the normalization in the field of functions
constant on the connected general fibres.
This proves the precise first-integral comparison.

\noindent\textbf{Step 7.} For the equivariant assertion, the finite orbit
of $V$ spans a nonzero finite-dimensional
representation of $G$ by meromorphic vertical
fields. A cyclic complex representation has
an eigenvector $V'\ne0$. On a general smooth
fibre each orbit field is holomorphic and
vertical, and therefore a translation field.
The same is true of $V'$; after removing
finitely many further values it is nonzero
on every remaining fibre. Choose a nonconstant
meromorphic function $h$ on the quotient
curve $B/G$, and pull it back to $B$.
Then the space
\[
                  \mathbb C V'+\mathbb C(q^*h)V'
                                                               \tag{C.FL.5}\label{cdp:eq:C.FL.5}
\]
is $G$-stable, with the same eigencharacter
on its two displayed generators. Its members
are semicomplete by the dense compact-fibre
and closedness argument just proved, and its
ratio has the same connected-fibre Stein
factor $q$. This proves the full statement.
\end{proof}






\begin{proposition}[The sum of the normalization geometric genera in one torus fibre is at most one]\label{cdp:st:C.2.UN.28}
\StatementRecord{cdp:st:C.2.UN.28}{proved}{Historical events: Claude Fable 5.2 (Anthropic), effort high, 16 November 2026: carried as background without fault found. GPT-6 Astra (OpenAI), ChatGPT Pro mode, effort Pro, S6\_09\_PAPERS, 17 November 2026: targeted proof check. Also checked in this writing. Pre-continuation writing events: GPT-6 Astra (OpenAI), ChatGPT Pro mode, effort Pro, S6\_10\_PAPERS, 18 November 2026: targeted proof check. Continuation mathematical and source checks: GPT-6 Astra (OpenAI), ChatGPT Pro mode, effort Pro, S6\_10\_PAPERS continuation, 19 November 2026: independent re-proof. Continuation metadata events: GPT-6 Astra (OpenAI), ChatGPT Pro mode, effort Pro, S6\_10\_PAPERS continuation, 19 November 2026: used and rechecked in this run; status-vocabulary reconciliation; metadata reconciliation only; no new mathematical check. \emph{Pre-continuation recorded proof scope.} The normalized root charts make each chosen strict covering component multiplicity one. The Poincare-Alexander map is exactly intersection with the boundary curves; positive real parts of holomorphic two-forms lift to compact cycles in its kernel. Actual horizontal flows embed disjoint cores into one nearby torus while preserving all intersection pairings. The positive index three of its exterior-square lattice gives twice the total geometric genus at most three}{the hypotheses, notation and definitions stated here; \cref{cdp:eq:C.PG.2}; \cref{cdp:eq:C.PG.6}; \cref{cdp:eq:C.PG.5}}

Let $f:U\to\Delta$ be a proper holomorphic map from a
smooth complex threefold, with connected two-torus
fibres over $\Delta^*$. Write its full central fibre
as $\sum_i m_iS_i$, and let $Z_i\to S_i^\nu$ be
smooth resolutions of the normalizations. Then
\[
                         \sum_i p_g(Z_i)\le1.
                                                               \tag{C.PG.1}\label{cdp:eq:C.PG.1}
\]
The assertion has no K\"ahler, Betti-number, reducedness,
or normal-crossing hypothesis. In particular at most
one normalization component in a given fibre has a
nonzero holomorphic two-form, and that component has
geometric genus one.

The same bound holds for any collection of multiplicity-one
components in a smooth resolved finite base change of
this family. The proof compares their compact homology
cores with one actual nearby torus, preserving their
intersection forms.
\end{proposition}



\begin{proof}
\noindent\textbf{Step 1.} Choose $N$ divisible by every $m_i$, make the actual
base change $t=s^N$, normalize, and resolve the total
space and its central support, isomorphically over
$s\ne0$ and at a generic point of every strict
component. This produces a proper map from a smooth
threefold
\[
                         g:\widetilde U\to\Delta_s
                                                               \tag{C.PG.2}\label{cdp:eq:C.PG.2}
\]
whose nearby fibre is exactly $f^{-1}(s^N)$.
Every strict component above an $S_i$ has multiplicity
one. To see the map explicitly at a generic point,
write $t=u x^{m_i}$ with $u$ a unit. The normalized
root charts are
\[
       w^{m_i}=u,\qquad x=s^{N/m_i}/w.
                                                               \tag{C.PG.3}\label{cdp:eq:C.PG.3}
\]
The unit-root equation is nonsingular at $s=0$;
there $s$ is a normal parameter. Subsequent
resolution is an isomorphism at this generic
point and preserves the multiplicity one.

\noindent\textbf{Step 2.} Select one smooth strict component $E_i$ above each
$S_i$. Its map to $S_i^\nu$ is generically finite.
Resolve the actual rational comparison to $Z_i$:
\[
              E_i\xleftarrow{\rho_i}W_i
                            \xrightarrow{v_i}Z_i.
                                                               \tag{C.PG.4}\label{cdp:eq:C.PG.4}
\]
The first map is a sequence of point blowups,
and the second is generically finite. Pullback
of holomorphic two-forms through $v_i$ is
injective, since its Jacobian is nonzero on
a dense open set. Point-blowup invariance of
holomorphic two-forms follows by restriction
and Hartogs extension across the blown-down
points. Hence
\[
                         p_g(E_i)\ge p_g(Z_i).
                                                               \tag{C.PG.5}\label{cdp:eq:C.PG.5}
\]
It suffices to prove the bound for the selected
components $E_i$ of \eqref{cdp:eq:C.PG.2}.

\noindent\textbf{Step 3.} Let $D_i$ be the reduced intersection of $E_i$
with all the other central components of the
resolved support. It is a compact normal-crossing
curve, or is empty. The restriction of $g$ is
a submersion near $E_i\setminus D_i$: in the
normal-crossing local equation only this
multiplicity-one component remains, so $g$
is a unit times its local defining coordinate.

\noindent\textbf{Step 4.} We first specify the relevant homology map on
any such compact smooth surface $E=E_i$ and
boundary curve $D=D_i$. Poincar\'e--Alexander
duality, applied to the triangulable compact
curve $D$ in the oriented real four-manifold
$E$, identifies the following portion of the
relative homology sequence:
\[
 H_2(E\setminus D,\mathbb R)\longrightarrow
 H_2(E,\mathbb R)\xrightarrow{\ I_D\ }
 H_2(E,E\setminus D;\mathbb R)
       \simeq H^2(D,\mathbb R).
                                                               \tag{C.PG.6}\label{cdp:eq:C.PG.6}
\]
There is one basis vector in $H^2(D,\mathbb R)$
for each irreducible curve of $D$. In these
coordinates $I_D$ is precisely the vector of
intersection numbers with those curves: under
Poincar\'e duality the map is restriction of
$H^2(E)$ to $D$. Thus its kernel is exactly
the image from the open complement in
\eqref{cdp:eq:C.PG.6}. This statement includes all
intersections between boundary curves; no
disjointness of those curves is assumed.

\noindent\textbf{Step 5.} Every holomorphic two-form $\alpha$ on $E$
is closed, because both its $(3,0)$ derivative
and its antiholomorphic derivative vanish.
The real classes of these forms give a
positive subspace
\[
 P_E=\{[\operatorname{Re}\alpha]:
                       \alpha\in H^0(E,K_E)\}
       \subset H^2(E,\mathbb R),
 \qquad \dim_{\mathbb R}P_E=2p_g(E).
                                                               \tag{C.PG.7}\label{cdp:eq:C.PG.7}
\]
Indeed, with the complex orientation,
\[
 \int_E(\operatorname{Re}\alpha)^2
       =\tfrac12\int_E\alpha\wedge\overline\alpha>0
                      \quad (\alpha\ne0).
                                                               \tag{C.PG.8}\label{cdp:eq:C.PG.8}
\]
This also proves injectivity of the real-part
map on the underlying real vector space of
holomorphic forms. Every such class has zero
integral on every boundary curve, since the
pullback of a two-form to the normalization
of a curve is zero. Its Poincar\'e-dual
homology class therefore lies in
$\ker I_D$ and lifts to
$H_2(E\setminus D,\mathbb R)$ by \eqref{cdp:eq:C.PG.6}.
No Hodge decomposition or K\"ahler form has
been used.

\noindent\textbf{Step 6.} Choose a real basis of each $P_{E_i}$,
choose lifts of its Poincar\'e-dual classes,
and represent those lifts by finite compact
smooth singular cycles in $E_i\setminus D_i$.
Include their transverse perturbations needed
to compute all intersection numbers in a
compact set $K_i$ of that complement. These
are homology-cycle representatives; they
need not be embedded surfaces.

\noindent\textbf{Step 7.} There is an actual orientation-preserving
smooth embedding of a neighbourhood of
$K_i$ into a nearby fibre of \eqref{cdp:eq:C.PG.2}.
Here is the local construction. On a
neighbourhood of $K_i$, $dg$ is surjective.
Choose smooth horizontal lifts of the two
coordinate fields of the real base disc;
a partition of unity preserves their
specified images under $dg$. Multiply by
a cutoff outside a slightly larger compact
neighbourhood and flow successively for
the real and imaginary parts of $s$.
For sufficiently small $s$, the trajectories
stay where the specified horizontal equations
hold. Their composition maps a neighbourhood
of $K_i$ in $E_i$ diffeomorphically to an
open subset of $g^{-1}(s)$. The image is
in that exact fibre by the two equations
for $dg$, and its orientation agrees with
the complex orientation by continuity
from $s=0$.

\noindent\textbf{Step 8.} The finitely many compact sets $K_i$ lie
in different central components away
from all intersections. Choose their
ambient neighbourhoods disjoint, and
then choose the same sufficiently small
$s\ne0$ for all the flows. Their images
in the single torus $F=g^{-1}(s)$ remain
disjoint. Transverse intersection numbers
of the chosen cycles are preserved by
the local embeddings, and cycles from
different components have zero intersection.
Consequently their images give mutually
orthogonal positive subspaces of
$H_2(F,\mathbb R)$, of dimensions
$2p_g(E_i)$. Injectivity on each positive
subspace, and on their direct sum,
follows from preservation of its positive
intersection form.

\noindent\textbf{Step 9.} The intersection form of an oriented
real four-torus is $3H$, where $H$ is
the hyperbolic plane. This follows
directly from the exterior-algebra
identification of its integral
cohomology: the three complementary
degree-two pairs give three hyperbolic
blocks. Its positive index is therefore
three. We have proved
\[
                    2\sum_i p_g(E_i)\le3.
                                                               \tag{C.PG.9}\label{cdp:eq:C.PG.9}
\]
The sum is integral, so it is at most
one. Together with \eqref{cdp:eq:C.PG.5} this is
\eqref{cdp:eq:C.PG.1}. The compact-core argument
used only multiplicity one, and hence
also proves the stated assertion for
any collection of such components in
a smooth finite-base-change model.
\end{proof}






\begin{proposition}[Rank-one elliptic component covers have only smooth multiple fibres]\label{cdp:st:C.2.UN.29}
\StatementRecord{cdp:st:C.2.UN.29}{proved}{Historical events: Claude Fable 5.2 (Anthropic), effort high, 16 November 2026: carried as background without fault found. GPT-6 Astra (OpenAI), ChatGPT Pro mode, effort Pro, S6\_09\_PAPERS, 17 November 2026: targeted proof check. Also checked in this writing. Pre-continuation writing events: GPT-6 Astra (OpenAI), ChatGPT Pro mode, effort Pro, S6\_10\_PAPERS, 18 November 2026: targeted proof check; used and rechecked in this run. Continuation mathematical and source checks: GPT-6 Astra (OpenAI), ChatGPT Pro mode, effort Pro, S6\_10\_PAPERS continuation, 19 November 2026: computation or formal certificate; targeted proof check. Continuation metadata events: GPT-6 Astra (OpenAI), ChatGPT Pro mode, effort Pro, S6\_10\_PAPERS continuation, 19 November 2026: used and rechecked in this run; status-vocabulary reconciliation; metadata reconciliation only; no new mathematical check. \emph{Pre-continuation recorded proof scope.} The continued compact time paths construct actual return germs preserving both ambient fields. Their exact pushforward equations at orders d and 2d force each period to be affine in the global meromorphic ratio. Period injectivity and the bounded Cayley transform force identity integral monodromy and fixed modulus. The root-disc invariant calculation proves the global R1 line and coherent genus formula. Kodaira1963II section8(iv)(11) p585 and Theorem6.2 pp564-566, together with Ulmer-Urzua section2.3 pp5-6, were opened at the exact smooth-fibre and monodromy loci}{the hypotheses, notation and definitions stated here; \cref{cdp:st:C.2.UN.16}; \cref{cdp:st:C.2.UN.22}; \cref{cdp:eq:C.RG.6}; \cref{cdp:eq:C.RG.8}; \cref{cdp:eq:C.RG.5}; \cref{cdp:eq:C.RG.12}; \cite[Theorem~6.2; Section~8(iv)(11), p.~585]{Kodaira1963II}; \citet[Section~2.3, pp.~5--6]{UlmerUrzua2021}; \cref{cdp:st:C.2.UN.28}; \cref{cdp:st:C.2.UN.20}; \cref{cdp:st:C.2.UN.15}; \cref{cdp:st:C.2.UN.19}}

Let $E$ be an actual multiplicity-one covering component
constructed in \ref{cdp:st:C.2.UN.16}, with its saturated two-dimensional
space of limiting translation fields. In the generic-rank-one
branch, take the actual graph resolution and Stein factor
of their ratio. When its connected general fibre is elliptic,
the resulting relatively minimal elliptic surface
\[
                         q:E_0\longrightarrow B
                                                               \tag{C.RG.1}\label{cdp:eq:C.RG.1}
\]
has identity integral $H_1$ monodromy. Every fibre is a
positive multiple of a smooth elliptic curve, and all smooth
fibres are torsors under one fixed elliptic curve. Moreover
\[
 R^1q_*\OO_{E_0}\simeq\OO_B,\qquad
 p_g(E)=g(B)\le1,\qquad \chi(\OO_E)=0.
                                                               \tag{C.RG.2}\label{cdp:eq:C.RG.2}
\]
The maps to $E_0$ are the specified point blowups resolving
the first integral followed by vertical exceptional-curve
contractions. No extension of that first integral to an
ambient neighbourhood is assumed.

Together with \ref{cdp:st:C.2.UN.22}, this proves that every such root-cover
component of nonnegative Kodaira dimension has coherent
Euler characteristic zero. In particular its minimal model
is neither K3 nor Enriques. Every original component of
Kodaira dimension one has an isotrivial elliptic Iitaka
fibration over a curve of genus at most one; its possible
singular fibres are retained by the actual cyclic quotient
comparison, and are not asserted to be smooth.
\end{proposition}


\begin{proof}
\noindent\textbf{Step 1.} We work on the actual root-change and resolution model
$f:U\to\Delta_s$ in \ref{cdp:st:C.2.UN.16}. Its smooth general fibre
is still a complex two-torus. The chosen component $E$
has multiplicity one. Its two ambient meromorphic fields
$V_1,V_2$ are vertical, commute, are holomorphic and
independent on each general torus, and have no pole along
the generic point of $E$. Their restrictions are independent
over $\C$ and have rank one over the meromorphic function
field of $E$.

\noindent\textbf{Step 2.} Resolve the meromorphic ratio on $E$ by point blowups
and take its connected-fibre Stein factor $q:E\to B$.
Perform the same point blowups in the smooth ambient
threefold. This gives the actual ambient extension of
these modifications, retains multiplicity one on the
strict transform of $E$, and lifts both fields
meromorphically. Their generic restrictions and bracket
are unchanged through the common open isomorphism.
We use the same letter $E$ for this modified component.
There is a nonconstant meromorphic function $H$ on $B$
such that
\[
                    V_2|_E=(q^*H)V_1|_E.
                                                               \tag{C.RG.3}\label{cdp:eq:C.RG.3}
\]
This is precisely the function and Stein map given by
the original ratio, rather than a chosen ambient extension
of it.

\noindent\textbf{Step 3.} Remove finitely many values of $B$ so that $q$ has smooth
connected elliptic fibres and $V_1|_E$ is a nonzero
holomorphic translation field on each of them. The latter
assertion follows from the compact-curve time-form argument
in C.2.UN.16: a semicomplete meromorphic field on an
elliptic curve has no pole. After excluding the finitely
many fibres on which its restriction is identically zero,
it is nowhere zero on every remaining fibre. Remove also
the poles and critical points of $H$.

\noindent\textbf{Step 4.} Let $D$ be the union on $E$ of its intersections with
the other central divisors and with the actual ambient
polar divisors of the two fields. Each remaining general
elliptic fibre meets $D$ in finitely many points, after
discarding any finitely many fibres contained in $D$.
The fields are holomorphic on $U^*$; hence all their
ambient polar divisors lie over zero. In particular
outside this explicitly removed set the ambient fields
and the submersion $f$ are holomorphic. At a suitable
point $p$ of a remaining fibre $F_{b_0}$ there are
holomorphic coordinates $(u,z,s)$ with
\[
 f=s,\qquad V_1=\partial_u,\qquad
 V_2=h(z,s)\partial_u+A(z,s)\partial_z,
 \qquad h(z,0)=H(z).
                                                               \tag{C.RG.4}\label{cdp:eq:C.RG.4}
\]
Here $z$ is a local coordinate on $B$ along $E$.
Indeed the holomorphic flow-box theorem gives $V_1=\partial_u$,
verticality preserves $s$, and commutation makes both
coefficients of $V_2$ independent of $u$. The local coordinate
$z$ is only used in this one flow box. It need not extend
around the compact elliptic fibre.

\noindent\textbf{Step 5.} Since the ambient fields are independent on the nearby
tori, $A$ is not identically zero. Since their restrictions
to $E$ have rank one, $A(z,0)=0$. Choosing $p$ away
from the zeros of its leading coefficient gives
\[
 A(z,s)=s^d\bigl(a(z)+O(s)\bigr),\qquad d\ge1,
 \qquad a(z)\ne0,\qquad H'(z)\ne0
                                                               \tag{C.RG.5}\label{cdp:eq:C.RG.5}
\]
on a smaller disc. The order $d$ is finite; it is the
order of the ambient wedge along this component in the
chosen nonvanishing frame.

\par\medskip\noindent\textbf{Actual elliptic return maps.}
Let $\lambda_1(z),\lambda_2(z)$ be the periods of the
time form of $V_1|_{F_z}$, for an oriented integral
homology basis continued over a small disc about $b_0$.
These are holomorphic functions; a local holomorphic
section and the ordinary integration construction of
an elliptic curve supply the marked periods. Put
$c_i=\lambda_i(b_0)$.

\noindent\textbf{Step 6.} The time covering $\C\to F_{b_0}$ of the translation
field pulls the finite bad set $D\cap F_{b_0}$ back to
a discrete subset of $\C$. Choose a piecewise smooth
time path from $0$ to $c_i$ avoiding this discrete set.
Its image is a compact path on which both ambient fields
are holomorphic and $V_1$ is nonzero. A finite chain of
holomorphic flow charts therefore continues the solution
along this entire path for every initial point in a
sufficiently small ambient neighbourhood of $p$.
Its endpoint is again near $p$, since the central path
closes. This constructs a genuine holomorphic return
germ $\Phi_i$; it does not deform a compact elliptic
orbit or assert that such an orbit exists on a nearby
torus.

\noindent\textbf{Step 7.} Every local continuation of the flow of $V_1$ preserves
$f,V_1,V_2$, by verticality and the bracket identity.
Thus the return germ in the single flow box has the form
\[
 \Phi_i(u,z,s)=\bigl(u+\delta_i(z,s),\phi_i(z,s),s\bigr).
                                                               \tag{C.RG.6}\label{cdp:eq:C.RG.6}
\]
On $E$ the first integral is constant along its elliptic
orbits, and comparison with the continued period gives
\[
 \phi_i(z,0)=z,\qquad
 \delta_i(z,0)=c_i-\lambda_i(z).
                                                               \tag{C.RG.7}\label{cdp:eq:C.RG.7}
\]
To see the second equality directly, follow the fixed
time path of total time $c_i$ on $F_z$, then return
to the local section by the small flow segment of
time $-\delta_i(z,0)$. The resulting closed path is
the continued integral cycle, so its time integral is
$c_i-\delta_i(z,0)=\lambda_i(z)$. This also specifies
which integral cycle and period enter the return map.

\noindent\textbf{Step 8.} Pushing $V_2$ forward by \eqref{cdp:eq:C.RG.6} and equating it to
$V_2$ at the endpoint gives the two exact identities
\[
 \begin{split}
 A(\phi_i(z,s),s)&=A(z,s)\,\partial_z\phi_i(z,s),\\
 h(\phi_i(z,s),s)&=h(z,s)+A(z,s)\,\partial_z\delta_i(z,s).
 \end{split}                                                   \tag{C.RG.8}\label{cdp:eq:C.RG.8}
\]
These are identities of holomorphic germs on the actual
ambient threefold.

\par\medskip\noindent\textbf{The leading normal coefficients.}
The second identity of \eqref{cdp:eq:C.RG.8}, together with $H'$ being
a unit and \eqref{cdp:eq:C.RG.5}, first gives
$\phi_i-z=O(s^d)$. For instance, factor
$h(\phi_i,s)-h(z,s)=(\phi_i-z)h_1(z,s)$;
the factor $h_1$ is a unit because $h_1(z,0)=H'(z)$.
Write
\[
                \phi_i=z+s^d\psi_i(z)+O(s^{d+1}).
                                                               \tag{C.RG.9}\label{cdp:eq:C.RG.9}
\]
In the first identity of \eqref{cdp:eq:C.RG.8}, all coefficients
of $A$ below order $2d$ cancel. At order $2d$ the
remaining equality is
\[
                         a'\psi_i=a\psi_i'.
                                                               \tag{C.RG.10}\label{cdp:eq:C.RG.10}
\]
Terms $s^{d+k}a_k(z)$ with $k\ge1$ first contribute
to the difference at order $2d+k$, so they do not
enter this calculation. Since $a$ is a unit,
$(\psi_i/a)'=0$. Hence $\psi_i=k_i a$ for a constant
$k_i\in\C$. At order $d$, the second identity gives
\[
 H'\psi_i=a\,\partial_z\delta_i(z,0)
          =-a\lambda_i',\qquad
             \lambda_i'= -k_i H'.
                                                               \tag{C.RG.11}\label{cdp:eq:C.RG.11}
\]
Consequently there are constants $d_i\in\C$ with
\[
                         \lambda_i=d_i-k_iH
                    \quad(i=1,2)
                                                               \tag{C.RG.12}\label{cdp:eq:C.RG.12}
\]
on the small disc. This uses only the two return germs
and one flow box. In particular it does not require
an ambient extension of $q$ or of $H$ around a fibre.
The formal coefficient extraction is independently
certified in \path{code/A_return_germs.py}; its executed
output records the higher-coefficient cancellations.

\par\medskip\noindent\textbf{Global periods and the singular fibres.}
Let $B^\circ$ be a connected complement of finitely
many values containing the disc used above, on which
the elliptic fibration, time form, and $H$ are regular.
Continue the chosen homology basis and periods on its
universal cover. Both sides of \eqref{cdp:eq:C.RG.12} are holomorphic
there, so the identity theorem gives the same formula
everywhere, with the same constants. Since $H$ is an
actual single-valued meromorphic function on $B$, each
continued period is single-valued on $B^\circ$.
The integration map from $H_1(F_b,\Z)$ to $\C$ is
injective for the time form of a nonzero elliptic
translation field. Thus returning both periods
unchanged forces every integral monodromy matrix to
be the identity.

\noindent\textbf{Step 9.} The oriented ratio $\tau=\lambda_2/\lambda_1$ is
therefore a single-valued holomorphic map
$B^\circ\to\mathbb H$. The denominator never vanishes,
because a nonzero integral elliptic cycle has nonzero
period. The Cayley transform $(\tau-\mathrm i)/(\tau+\mathrm i)$
is bounded, extends holomorphically across all finitely
many omitted points, and is constant on the compact
curve $B$. Hence $\tau$ is constant. Integrating a
locally normalized time form identifies the smooth
fibres with torsors under one fixed elliptic curve.

\noindent\textbf{Step 10.} Contract the vertical exceptional curves to obtain
\eqref{cdp:eq:C.RG.1}. The integral monodromy is unchanged. The
exact Kodaira-table and multiple-root comparison in
(C.RF.3) applies: for a full fibre $mC$, write
$t=u h^m$, make the actual change $t=s^m$, and
normalize by $w^m=u$, $s=wh$. This covering surface
is \'etale over the original one. Its reduced
relatively minimal fibre has identity monodromy,
so is the smooth type $I_0$. It follows that every
original reduced support $C$ is smooth elliptic.
This is precisely the comparison established from
\cite[Theorem~6.2; Section~8(iv)(11), p.~585]{Kodaira1963II} and \citet[Section~2.3, pp.~5--6]{UlmerUrzua2021} in
\ref{cdp:st:C.2.UN.22}; it includes all multiple fibres.

\par\medskip\noindent\textbf{The direct image and the genus bound.}
Over a small smooth base disc, a local holomorphic
section and integration of the marked differential
give an actual product with the fixed elliptic curve.
Since the integral monodromy is the identity, the
origin-preserving linear part of every overlap map
is the identity. The overlap maps are translations,
which act trivially on $H^1(\OO)$.
At a multiple value use the actual normalized root
family just described. It is smooth with constant
elliptic modulus. The cyclic action on its central
elliptic fibre is free, since the map to $E_0$ is
\'etale, and is therefore a translation. Its linear
part on the neighbouring marked elliptic fibres
is the same discrete automorphism and hence is
the identity. On $H^1(\OO)$ it acts trivially.
Taking invariants in the local root disc consequently
gives $R^1q_*\OO=\OO$ at the multiple value as well:
the coefficient functions invariant under
$s\mapsto\zeta_m s$ are exactly the holomorphic
functions of $t=s^m$. These local generators agree
with the preceding translation-invariant generators
over the punctured discs. This proves the first
assertion of \eqref{cdp:eq:C.RG.2} with its actual descent map.

\noindent\textbf{Step 11.} Connected fibres give $q_*\OO_{E_0}=\OO_B$.
The coherent Leray sequence on the curve now gives
\[
 H^2(E_0,\OO_{E_0})\simeq H^1(B,\OO_B),\qquad
 h^1(E_0,\OO_{E_0})=g(B)+1.
                                                               \tag{C.RG.13}\label{cdp:eq:C.RG.13}
\]
Serre duality and birational invariance give
$p_g(E)=g(B)$ and $\chi(\OO_E)=0$.
The independent compact-core inequality
\ref{cdp:st:C.2.UN.28} gives $p_g(E)\le1$, and therefore
$g(B)\le1$.

\noindent\textbf{Step 12.} Every root-cover component of nonnegative Kodaira
dimension has either generic rank two, covered by
\ref{cdp:st:C.2.UN.22}, or rank one. In the latter case a rational
general first-integral fibre would make the surface
uniruled and give Kodaira dimension $-\infty$ by
the actual adjunction argument in \ref{cdp:st:C.2.UN.20}.
Its general fibre is therefore elliptic and the
present proof applies. This proves the stated
coherent-Euler and K3/Enriques exclusions on the
covering components.

\noindent\textbf{Step 13.} Finally let an original component have Kodaira
dimension one. Its actual generically finite
component cover has Kodaira dimension at least
one by pluricanonical pullback, and at most one
by \ref{cdp:st:C.2.UN.15}. Its elliptic fibration is therefore
its connected-fibre Iitaka fibration. The cyclic
stabilizer acts on this fibration, since it acts
on the pluricanonical field, and the proper
quotient and graph comparisons in \ref{cdp:st:C.2.UN.19}--20
give the descended fibration on a smooth
birational model of the original component.
Its base is the quotient of $B$, so has genus
at most one. On a general fibre, the subgroup
fixing a general base point must act by
translations: an origin-preserving nontrivial
linear part would give a rational general quotient
fibre. Its normal line is trivial, so adjunction
restricts every positive pluricanonical line to
$\OO_{\mathbb P^1}(-2m)$; all pluricanonical sections
would then vanish on the dense family of these
fibres, contrary to Kodaira dimension one. The finite translation
subgroup of the fixed elliptic curve is locally
constant, because it consists of points of
bounded finite order. Thus the general quotient
fibre is one fixed elliptic curve as well.
This proves isotriviality through the actual
finite quotient, while leaving its exceptional
and singular fibres in place.
\end{proof}






\begin{proposition}[Nonnegative-Kodaira-dimension components have tightly specified cyclic covers]\label{cdp:st:C.2.UN.31}
\StatementRecord{cdp:st:C.2.UN.31}{proved}{Historical events: Claude Fable 5.2 (Anthropic), effort high, 16 November 2026: carried as background without fault found. GPT-6 Astra (OpenAI), ChatGPT Pro mode, effort Pro, S6\_09\_PAPERS, 17 November 2026: targeted proof check. Also checked in this writing. Pre-continuation writing events: GPT-6 Astra (OpenAI), ChatGPT Pro mode, effort Pro, S6\_10\_PAPERS, 18 November 2026: targeted proof check. Continuation mathematical and source checks: GPT-6 Astra (OpenAI), ChatGPT Pro mode, effort Pro, S6\_10\_PAPERS continuation, 19 November 2026: targeted proof check. Continuation metadata events: GPT-6 Astra (OpenAI), ChatGPT Pro mode, effort Pro, S6\_10\_PAPERS continuation, 19 November 2026: used and rechecked in this run; status-vocabulary reconciliation; metadata reconciliation only; no new mathematical check. \emph{Pre-continuation recorded proof scope.} The unique minimal contraction descends the cyclic action, and each original fibration is formed through its actual quotient and connected-fibre Stein map. The pluricanonical restriction proves that the generic kernel acts by translations. Canonical nefness makes every exceptional curve vertical, so the map descends to the original minimal surface. The bielliptic connected-automorphism orbit map and K3 form pullback treat the remaining cases; the coherent Euler range follows from the compact-core bound and surface invariants}{the hypotheses, notation and definitions stated here; \cref{cdp:st:C.2.UN.24}; \cref{cdp:st:C.2.UN.16}; \cref{cdp:st:C.2.UN.15}; \cref{cdp:st:C.2.UN.22}; \cref{cdp:st:C.2.UN.20}; \cref{cdp:eq:C.CS.2}; \cref{cdp:st:C.2.UN.29}; \cref{cdp:st:C.2.UN.28}; \cref{cdp:eq:C.CS.4}; \cref{cdp:st:C.3.1}; \cref{cdp:eq:C.CS.5}; \cref{cdp:st:C.2.UN.23}; \cref{cdp:st:C.2.UN.32}}

Let $f:U\to\Delta$ be an unrestricted proper degeneration
of connected complex two-tori from a smooth threefold.
Let $Z$ be the smooth minimal model of the normalization
of one central component, with $\kappa(Z)\ge0$.
There is an actual generically finite component cover
$E\to S^\nu$, obtained by the normalized root base change
and resolution, with cyclic stabilizer $G$, such that
$E/G$ is properly birational to $S^\nu$. Its smooth
minimal model $E_0$ satisfies one of these descriptions:

\begin{itemize}
\item $E_0$ is a complex two-torus;
\item $E_0$ has an isotrivial elliptic fibration over a
curve $B$ with $g(B)\le1$, identity integral $H_1$
monodromy, and only positive multiples of smooth
elliptic fibres. In this case
$R^1q_*\mathcal O_{E_0}=\mathcal O_B$,
$p_g(E_0)=g(B)$, and $\chi(\mathcal O_{E_0})=0$.
\end{itemize}
In both cases $\chi(\mathcal O_{E_0})=0$, and the
action of $G$ descends holomorphically to $E_0$.

For the original minimal component the consequences are
\[
 p_g(Z)\le1,
 \qquad \chi(\mathcal O_Z)\in\{0,1,2\}.
                                                               \tag{C.CS.1}\label{cdp:eq:C.CS.1}
\]
Every component with $\kappa(Z)=1$ has an actual
isotrivial elliptic Iitaka fibration on $Z$, over a
curve of genus at most one. An Enriques component
has an isotrivial elliptic fibration over $\mathbb P^1$.
A K3 component is either the cyclic torus quotient
of \ref{cdp:st:C.2.UN.24}, with the exact orders and configurations
listed there, or has an isotrivial elliptic fibration
over $\mathbb P^1$ through the non-torus-cover branch.
The singular and multiple fibres on these original
quotient surfaces are retained by the quotient maps;
the smooth-multiple assertion concerns the specified
cover $E_0$.
\end{proposition}



\begin{proof}
\noindent\textbf{Step 1.} Take the actual component, cyclic stabilizer, and
proper birational quotient comparison in \ref{cdp:st:C.2.UN.16}
and (C.CQ.12). Pluricanonical pullback through a
resolution of the generically finite graph gives
\[
             \kappa(E)\ge\kappa(Z)\ge0.
                                                               \tag{C.CS.2}\label{cdp:eq:C.CS.2}
\]
The compact-component theorem \ref{cdp:st:C.2.UN.15} applies
also to this finite base-change family, so
$\kappa(E)\le1$. For the saturated limiting field
space, the generic-rank-two branch is covered
by \ref{cdp:st:C.2.UN.22}. In rank one the first integral
has rational or elliptic general fibre. Rational
fibres would force every positive pluricanonical
section to vanish, by the exact adjunction
argument of \ref{cdp:st:C.2.UN.20}, contradicting \eqref{cdp:eq:C.CS.2}.
Thus \ref{cdp:st:C.2.UN.29} applies. These two statements
give precisely the asserted torus or smooth-multiple
elliptic alternatives and coherent Euler characteristic
zero.

\noindent\textbf{Step 2.} In the elliptic alternative the actual root-disc
invariant calculation of (C.RG.2)--(C.RG.13) gives
$R^1q_*\mathcal O=\mathcal O_B$ and
$p_g(E)=g(B)$. That calculation requires only
the fixed elliptic modulus, identity integral
monodromy and the smooth-multiple local charts,
so it applies to the elliptic alternative of
\ref{cdp:st:C.2.UN.22} as well. The independent compact-core
bound \ref{cdp:st:C.2.UN.28} gives $p_g(E)\le1$, and hence
$g(B)\le1$ in both branches.

\noindent\textbf{Step 3.} All proper birational maps between the displayed
smooth surface models factor through actual
point blowups and contractions. On a nonruled
surface the minimal model is unique. Equivalently,
the exceptional curves of a minimal contraction
form a negative-definite configuration; its
$(-1)$-curves are disjoint, so their finite
$G$-orbits can be contracted equivariantly,
one stage at a time. This gives the holomorphic
action on $E_0$ and the exact comparison
\[
 \begin{array}{ccc}
 E&\longrightarrow&E_0\\
 \downarrow&&\downarrow\\
 E/G&\longrightarrow&E_0/G
 \end{array}
 \qquad E_0/G\dashrightarrow Z
       \quad\hbox{properly birational}.
                                                               \tag{C.CS.3}\label{cdp:eq:C.CS.3}
\]
It is this action and quotient which are used
below.

\noindent\textbf{Step 4.} Suppose first that $\kappa(Z)=1$. Then
\eqref{cdp:eq:C.CS.2} and \ref{cdp:st:C.2.UN.15} give $\kappa(E_0)=1$,
so the torus alternative is impossible. The
elliptic fibration is the connected-fibre
Iitaka fibration. To specify the invariance,
the group acts on every space of pluricanonical
sections. An equivariant graph resolution of
a sufficiently large pluricanonical map and
its Stein factor therefore give a $G$-equivariant
map to $B$. The same Stein base is the one
of the preceding elliptic fibration: pluricanonical
ratios restrict constantly to its elliptic
fibres, and Kodaira dimension one makes the
induced map of bases finite. Taking the
connected-fibre Stein factor removes that
finite map.

\noindent\textbf{Step 5.} The actual quotient fibration is
\[
                   E'/G\longrightarrow B/G,
                                                               \tag{C.CS.4}\label{cdp:eq:C.CS.4}
\]
after this equivariant graph resolution.
For a general base point, let
$K=\ker(G\to\operatorname{Aut}B)$.
Its fibre is exactly $F/K$, as in \ref{cdp:st:C.2.UN.20},
and is connected. Every $F$ is isomorphic
to the same elliptic curve $A$.
The subgroup $K$ acts by translations:
a finite automorphism of an elliptic curve
with nontrivial linear part has a fixed
point, and its quotient is rational by
Riemann--Hurwitz. A rational general fibre
in \eqref{cdp:eq:C.CS.4} would annihilate every positive
pluricanonical section, contrary to
$\kappa(Z)=1$. The translation points
belong to the finite subgroup $A[|K|]$,
so they are locally constant under the
marked product identifications. Therefore
the quotient general fibre is the one
fixed elliptic curve $A/K$.
The quotient base has genus at most one
by Riemann--Hurwitz and $g(B)\le1$.

\noindent\textbf{Step 6.} This fibration is the Iitaka fibration
of $Z$. Its map is holomorphic on the
minimal surface. One can also check the
descent on an actual common resolution
$Z\xleftarrow{\rho}X\xrightarrow q B/G$:
for a smooth general elliptic fibre $F$,
\[
 0=K_X\cdot F
   =\rho^*K_Z\cdot F+\sum_j a_jE_j\cdot F,
                   \qquad a_j>0.
                                                               \tag{C.CS.5}\label{cdp:eq:C.CS.5}
\]
On a minimal surface of Kodaira dimension
one the canonical class is nef; this is
the elliptic canonical-bundle formula
in the surface classification used in
\ref{cdp:st:C.3.1}. Every term in \eqref{cdp:eq:C.CS.5} is
nonnegative. Hence every exceptional
curve is vertical. Connected exceptional
fibres and $\rho_*\mathcal O_X=\mathcal O_Z$
then descend the actual proper graph to
a holomorphic map on $Z$. This proves
the claimed isotrivial Iitaka fibration.

\noindent\textbf{Step 7.} Now let $Z$ be Enriques. Algebraic
dimension is preserved by the actual
generically finite component map, by
the symmetric-function field argument
in (C.K3.2)--(C.K3.3). Thus $E_0$ is
projective. If it has Kodaira dimension
one, the preceding intrinsic Iitaka
quotient argument again gives an
isotrivial elliptic fibration on a
birational model of $Z$. If it has
Kodaira dimension zero, the surface
classification and $\chi(\mathcal O_{E_0})=0$
make it Abelian or bielliptic.
The Abelian case is excluded by the
exact cyclic quotient theorem \ref{cdp:st:C.2.UN.23}.
For the bielliptic case, use the actual
orbit fibration of its connected
automorphism group. In its finite free
product presentation $(A_1\times A_2)/H$,
$H$ acts by translations on $A_1$ and
has nontrivial linear action on $A_2$.
The invariant tangent fields are exactly
the one-dimensional translation space
of $A_1$. Thus the image of the compact
elliptic translation group $A_1$ is
the whole connected automorphism group,
and its orbit map is the explicit
elliptic fibration to $A_2/H\simeq\mathbb P^1$.
Its fibres are smooth elliptic multiples;
its integral monodromy is the identity,
since all actions on the fibre direction
are translations. Every automorphism
normalizes the connected automorphism
group, so preserves this orbit map and
acts on its connected-fibre base. It
is therefore $G$-equivariant. The quotient
\eqref{cdp:eq:C.CS.4} is available again. Its
general fibre must be elliptic, since
a rational general fibre would force
the nonzero even pluricanonical form
of an Enriques surface to vanish.

\noindent\textbf{Step 8.} In either case, \eqref{cdp:eq:C.CS.5} descends the
elliptic map to the minimal Enriques
surface: now $K_Z$ is torsion, so its
intersection term is zero. The Leray
injection of $H^1$ of the base into
$H^1(Z,\mathcal O_Z)=0$ makes the
base $\mathbb P^1$. This proves the
Enriques assertion through the actual
quotient fibration.

\noindent\textbf{Step 9.} For a K3 target, the nonzero canonical
form gives $p_g(E_0)\ge1$. If $E_0$
is a torus, \ref{cdp:st:C.2.UN.23}--24 give exactly
the cyclic symplectic torus quotient
already stated. Otherwise it has an
elliptic fibration of the specified
fixed-modulus type. If $\kappa(E_0)=1$,
use the equivariant Iitaka fibration.
If $a(E_0)=1$, use the intrinsic
field identity $\mathcal M(E_0)=q^*\mathcal M(B)$
and the equivariant Stein factor in
(C.K3.4). The remaining projective
Kodaira-dimension-zero case would be
bielliptic, but has $p_g=0$ and is
excluded. Algebraic dimension zero
has only the torus alternative by
\ref{cdp:st:C.2.UN.16}. Thus these are all cases.
Taking the quotient gives an isotrivial
elliptic fibration, since its general
fibre cannot be rational on a K3
birational model. Formula \eqref{cdp:eq:C.CS.5}
with $K_Z=0$ descends it to $Z$,
and $H^1(Z,\mathcal O_Z)=0$ again
makes its base $\mathbb P^1$.

\noindent\textbf{Step 10.} Finally \ref{cdp:st:C.2.UN.28} gives $p_g(Z)\le1$.
The compact surface classification
gives $\chi(\mathcal O_Z)\ge0$ when
$\kappa(Z)\ge0$: in Kodaira dimension
zero its values are the explicit
torus, bielliptic, Kodaira, K3 and
Enriques values; in Kodaira dimension
one the relatively minimal elliptic
surface has $K_Z^2=0$ and nonnegative
Euler number by Kodaira's fibre list,
so surface Riemann--Roch gives the
inequality. On the other hand
\[
        \chi(\mathcal O_Z)=1-h^1(Z,\mathcal O_Z)+p_g(Z)
                              \le2.
                                                               \tag{C.CS.6}\label{cdp:eq:C.CS.6}
\]
It is an integer, giving \eqref{cdp:eq:C.CS.1}.
The later exact quotient theorem
\ref{cdp:st:C.2.UN.32} further identifies every K3
component, in all algebraic dimensions,
with the cyclic torus quotient family.
The preceding non-torus-cover comparison
is retained as its explicit intermediate
fibration map.

\noindent\textbf{Step 11.} All these comparisons retain the
finite quotient's singular fibres;
they do not identify them with the
smooth fibres of its root cover.
\end{proof}






\begin{proposition}[A properly elliptic normalization component has one of eleven positive-coherent-Euler fibre configurations]\label{cdp:st:C.2.UN.33}
\StatementRecord{cdp:st:C.2.UN.33}{proved}{Historical events: Claude Fable 5.2 (Anthropic), effort high, 16 November 2026: carried as background without fault found. GPT-6 Astra (OpenAI), ChatGPT Pro mode, effort Pro, S6\_09\_PAPERS, 17 November 2026: targeted proof check. Also checked in this writing. Pre-continuation writing events: GPT-6 Astra (OpenAI), ChatGPT Pro mode, effort Pro, S6\_10\_PAPERS, 18 November 2026: targeted proof check; used and rechecked in this run. Continuation mathematical and source checks: GPT-6 Astra (OpenAI), ChatGPT Pro mode, effort Pro, S6\_10\_PAPERS continuation, 19 November 2026: computation or formal certificate; targeted proof check. Continuation metadata events: GPT-6 Astra (OpenAI), ChatGPT Pro mode, effort Pro, S6\_10\_PAPERS continuation, 19 November 2026: used and rechecked in this run; status-vocabulary reconciliation; metadata reconciliation only; no new mathematical check. \emph{Pre-continuation recorded proof scope.} The exact root-chart quotient by the fibrewise translation kernel, the two divisibility conditions excluding a multiple branch fibre with nontrivial linear action, and the invariant Weierstrass coordinates were re-derived. All cyclic base-character cases and their fibre Euler numbers give the eleven necessary configurations. The invariant canonical character and compact-core bound give the coherent invariants. Kodaira's complete fibre table and the exact period-extension case were opened in reproduced original journal text; the primary author-preprint monodromy table was opened separately. No occurrence assertion is inferred from this necessary list}{the hypotheses, notation and definitions stated here; \cref{cdp:st:C.2.UN.16}; \cref{cdp:st:C.2.UN.15}; \cref{cdp:st:C.2.UN.22}; \cref{cdp:st:C.2.UN.29}; \cref{cdp:st:C.2.UN.19}; \cref{cdp:eq:C.EQ.3}; \cref{cdp:st:C.2.UN.32}; \cref{cdp:eq:C.EQ.4}; \cref{cdp:eq:C.EQ.5}; \citet[Section~2.3, pp.~5--6]{UlmerUrzua2021}; \citet[Theorem~6.2 and Section~8(iv)]{Kodaira1963II}; \cref{cdp:eq:C.EQ.7}; \cref{cdp:eq:C.EQ.8}; \cref{cdp:eq:C.EQ.10}; \cref{cdp:eq:C.EQ.11}; \cref{cdp:st:C.2.UN.28}}

Let $M$ be the smooth minimal model of a normalization
component of a fibre of an unrestricted two-torus degeneration,
and suppose $\kappa(M)=1$. Its connected elliptic Iitaka map
$p:M\to B_0$ is isotrivial, and $g(B_0)\le1$.
Its coherent Euler characteristic belongs to $\{0,1,2\}$.
When it is zero, every fibre of $p$ is a positive multiple
of a smooth elliptic curve. When it is positive,
$B_0\simeq\PP^1$, $q(M)=0$, and the complete configuration
of the fibres of $p$ with singular reduced support is one
of the following rows:
\[
\begin{array}{c|l|c}
\chi(\OO_M)&\text{singular reduced fibres of }p&e(M)\\ \hline
1&\mathrm{II}+\mathrm{II}^*&12\\
1&\mathrm{III}+\mathrm{III}^*&12\\
1&\mathrm{IV}+\mathrm{IV}^*&12\\
1&2\mathrm I_0^*&12\\
1&3\mathrm{IV}&12\\
1&2\mathrm{III}+\mathrm I_0^*&12\\
1&\mathrm{II}+\mathrm{IV}+\mathrm I_0^*&12\\ \hline
2&4\mathrm I_0^*&24\\
2&3\mathrm{IV}^*&24\\
2&2\mathrm{III}^*+\mathrm I_0^*&24\\
2&\mathrm{II}^*+\mathrm{IV}^*+\mathrm I_0^*&24 .
\end{array}                                                   \tag{C.EQ.1}\label{cdp:eq:C.EQ.1}
\]
In each row $p_g(M)=\chi(\OO_M)-1$. All additional
smooth multiple fibres are retained as part of the specified
surface; the table lists its nonsmooth reduced fibres.
This is a necessary classification of the displayed
invariant for actual components. The statement does not
assert occurrence of every abstract surface with these data.
The map $p$ is the elliptic map on the component surface;
the original threefold map has the degeneration disc as target.
\end{proposition}



\begin{proof}
\noindent\textbf{Step 1.} Take the actual cyclic component cover of \ref{cdp:st:C.2.UN.16}.
Pluricanonical pullback gives Kodaira dimension at least
one, and \ref{cdp:st:C.2.UN.15} excludes dimension two. Its minimal model
$E$ therefore has Kodaira dimension one. The results
\ref{cdp:st:C.2.UN.22} and \ref{cdp:st:C.2.UN.29} give its connected elliptic Iitaka map
\[
 q:E\longrightarrow B,\qquad g(B)\le1,                 \tag{C.EQ.2}\label{cdp:eq:C.EQ.2}
\]
with a fixed elliptic fibre, identity integral monodromy,
and only smooth multiple fibres. The cyclic stabilizer
$G$ preserves $q$, because it acts on the pluricanonical
field. The normalization of the original component is
birational to the actual finite quotient $E/G$, through
the equivariant graph and minimal-model comparisons of
\ref{cdp:st:C.2.UN.19}--20. Thus its relatively minimal elliptic model is
the model of that quotient over $B/G$.

\noindent\textbf{Step 2.} Let $H$ be the kernel of $G\to\operatorname{Aut}(B)$.
On a general elliptic fibre it acts by translations.
Indeed a nontrivial linear elliptic automorphism has a
fixed point and its quotient is rational by
Riemann--Hurwitz. A rational general fibre on the
quotient would make every positive pluricanonical
section vanish by adjunction, contrary to $\kappa(M)=1$.
The finite translations form a fixed subgroup of the
constant elliptic curve, since points of bounded order
are discrete. We must retain the action at multiple
fibres when taking the quotient by $H$.

\par\medskip\noindent\textbf{The quotient by fibre translations.}
At a full fibre $mF$ of $q$, the actual normalized root
chart is
\[
 (F'\times\Delta_z)/\langle\delta\rangle,
 \qquad \delta(v,z)=(v+\tau,\zeta_m z),
 \qquad \operatorname{ord}(\tau)=m,
 \qquad t=z^m.                                        \tag{C.EQ.3}\label{cdp:eq:C.EQ.3}
\]
The root family is a smooth fixed-elliptic family.
The cyclic deck action is a translation on every fibre,
and its free action on the central fibre gives exactly
the indicated order of $\tau$. This is the normalized
root comparison proved in \ref{cdp:st:C.2.UN.22} and \ref{cdp:st:C.2.UN.29}, with the
constant-translation gauge written in (C.KX.4).

\noindent\textbf{Step 3.} Lift a base-fixing element of $H$ to this chart.
Its multiplier on $z$ is a power of $\zeta_m$;
composition with that inverse power of $\delta$
makes it fix $z$. The appropriate finite power
is a deck transformation fixing $z$, and hence
the identity. Its fibre action is therefore a
translation by a finite-order constant point.
The lifts so chosen generate a finite translation
subgroup $T\subset F'$, commuting with $\delta$.
First quotient by $T$. Write $\bar F=F'/T$ and
$m'=\operatorname{ord}(\bar\tau)$, so $m'\mid m$.
The remaining action is
$(\bar v,z)\mapsto(\bar v+\bar\tau,\zeta_m z)$.
Its $m'$-th power is a pure rotation of the disc.
Quotienting that subgroup replaces $z$ by
$w=z^{m/m'}$; the remaining free action is
\[
 (\bar v,w)\longmapsto(\bar v+\bar\tau,\zeta_{m'}w),
 \qquad t=w^{m'}.                                    \tag{C.EQ.4}\label{cdp:eq:C.EQ.4}
\]
This proves directly that $E_1=E/H$ is smooth near
the fibre and has the full smooth multiple fibre
$m'(\bar F/\langle\bar\tau\rangle)$. It is smooth
elsewhere by the same local translation quotient.
Thus $E_1\to B$ still has a fixed elliptic fibre,
identity integral monodromy, and only smooth multiples.
The remaining cyclic group $Q=G/H$ acts faithfully
on $B$.

\par\medskip\noindent\textbf{The local action at a base branch value.}
Let a stabilizer of order $s$ act at $b\in B$, and
let $\beta$ be its linear action on the fixed elliptic
curve. The linear action is constant along the regular
base, because it belongs to a discrete finite automorphism
group. Suppose first that $\beta\ne1$.
We prove that the fibre of $E_1$ at $b$ has multiplicity
one; this prevents a hidden change of the singular-fibre
type in the quotient.

\noindent\textbf{Step 4.} Use \eqref{cdp:eq:C.EQ.3} for a putative multiplicity $m$.
The generator lifts to the normalized root chart,
with normal multiplier $\rho$ satisfying
$\rho^m=\zeta_s$, and affine fibre derivative $\beta$.
It commutes with the deck action: its conjugation
preserves the multiplier $\zeta_m$ on the root
coordinate, and hence preserves the corresponding
deck generator. Restriction to the central elliptic
curve gives $\beta\tau=\tau$. The finite-order
endomorphism $\beta\ne1$ satisfies
$1+\beta+\cdots+\beta^{s-1}=0$. Applying it to
$\tau$ yields $s\tau=0$, so $m\mid s$.

\noindent\textbf{Step 5.} The induced affine map on the reduced elliptic
fibre has a fixed point, because $\beta-1$ is a
surjective elliptic isogeny. Choose a lift of that
point, and compose the lifted generator with a
deck transformation to make it fix the chosen
lift. Its $s$-th power is a deck transformation
fixing a point of the central elliptic curve;
the deck action there is free, so that power is
the identity. Therefore $\rho^s=1$.
Together with $\rho^m$ having order $s$, this
gives $\gcd(m,s)=1$. Since $m\mid s$, we have
\[
                 m=1\qquad(\beta\ne1).               \tag{C.EQ.5}\label{cdp:eq:C.EQ.5}
\]
This is the same explicit local comparison used
in \ref{cdp:st:C.2.UN.32}, and the argument did not use a canonical
form or a K3 hypothesis.

\noindent\textbf{Step 6.} When $\beta=1$, every element with a fixed point
on a reduced elliptic fibre fixes that whole fibre:
an elliptic translation with a fixed point is the
identity. Its stabilizer acts only on the normal
coordinate. A finite holomorphic action with
identity derivative is locally the identity,
by the averaging linearization, so every nontrivial
stabilizer is a normal reflection. The quotient
surface is smooth and the reduced fibre remains
smooth. Equivalently, the translation and root
quotients \eqref{cdp:eq:C.EQ.3}--\eqref{cdp:eq:C.EQ.4} display that fibre
as a smooth multiple. Hence these values have
zero topological Euler contribution.

\noindent\textbf{Step 7.} For $\beta\ne1$, \eqref{cdp:eq:C.EQ.5} gives a smooth product
near the source fibre. An equivariant local section
through a fixed point removes the translation part
of the finite action. To construct it, lift its
translation cocycle to small holomorphic complex-valued
functions near the fixed point and average over the
finite group. The cocycle sum is zero because the
action has finite order and the lift vanishes at
the fixed point. In additive notation, for an operator
$P$ of order $s$ and a cocycle $c$ with
$\sum_{k=0}^{s-1}P^kc=0$, the function
$s^{-1}\sum_{k=0}^{s-1}kP^kc$ solves
$(P-1)b=c$. This supplies the claimed holomorphic
section and gives a product action.
Quotient the subgroup acting trivially on the
elliptic factor. It consists of pure normal
reflections. The remaining action has order
$d\in\{2,3,4,6\}$ and the exact form
\[
             (v,z)\longmapsto(\zeta_d^r v,\zeta_d z),
             \qquad r\in\{1,d-1\}.                   \tag{C.EQ.6}\label{cdp:eq:C.EQ.6}
\]
Translations of an origin have already been removed.

\noindent\textbf{Step 8.} Here is an explicit comparison with the Kodaira
types, including the direction of the character.
Write the fixed elliptic curve as
$y^2=x^3+Ax+B$, so multiplication by $\zeta_d^r$
acts on $x,y$ by $\zeta_d^{-2r},\zeta_d^{-3r}$.
The invariant meromorphic functions
\[
          T=z^d,\qquad X=z^{2r}x,\qquad Y=z^{3r}y
                                                               \tag{C.EQ.7}\label{cdp:eq:C.EQ.7}
\]
generate the quotient function field. They give
the following minimal Weierstrass equations:
\[
\begin{array}{c|c|c|c}
d&r&\text{equation}&\text{Kodaira type}\\ \hline
2&1&Y^2=X^3+AT^2X+BT^3&\mathrm I_0^*\\
3&1&Y^2=X^3+BT^2&\mathrm{IV}\\
3&2&Y^2=X^3+BT^4&\mathrm{IV}^*\\
4&1&Y^2=X^3+ATX&\mathrm{III}\\
4&3&Y^2=X^3+AT^3X&\mathrm{III}^*\\
6&1&Y^2=X^3+BT&\mathrm{II}\\
6&5&Y^2=X^3+BT^5&\mathrm{II}^* .
\end{array}                                                   \tag{C.EQ.8}\label{cdp:eq:C.EQ.8}
\]
For $d=3,6$ the curve has $A=0$ and $B\ne0$;
for $d=4$ it has $B=0$ and $A\ne0$.
The equations are minimal, since their indicated
orders are below the Weierstrass rescaling threshold.
The primitive monodromy matrices and their characters distinguish the
unstarred and starred types; see \citet[Section~2.3, pp.~5--6]{UlmerUrzua2021}
and \citet[Theorem~6.2 and Section~8(iv)]{Kodaira1963II}.
For the displayed minimal equations the discriminant orders are computed
directly as $12r/d$, and the corresponding Kodaira configurations have
the same Euler numbers.
The map \eqref{cdp:eq:C.EQ.7} is birational over the base
after taking the finite quotient. Resolving its
graph and contracting vertical exceptional curves
gives the same relatively minimal elliptic model;
therefore these are the fibre types of the actual
quotient. This verifies the reduction to that table,
including the two different types of each primitive
order. Their Euler numbers are $12r/d$.

\par\medskip\noindent\textbf{All faithful cyclic base actions.}
On $B=\PP^1$, a nontrivial cyclic action is conjugate
to $z\mapsto\zeta_nz$. Its only branch values are
zero and infinity, where positive local generators
are inverse. A nontrivial fibre character of order
$d$ consequently gives the two opposite exponents
in \eqref{cdp:eq:C.EQ.8}. The resulting pairs are
\[
 (\mathrm{II},\mathrm{II}^*),\quad
 (\mathrm{III},\mathrm{III}^*),\quad
 (\mathrm{IV},\mathrm{IV}^*),\quad
 (\mathrm I_0^*,\mathrm I_0^*).                         \tag{C.EQ.9}\label{cdp:eq:C.EQ.9}
\]
Every pair has Euler sum twelve. A trivial fibre
character gives only smooth multiples.

\noindent\textbf{Step 9.} On an elliptic $B$, a finite cyclic group either
acts by translations, in which case it is free,
or has a fixed point. In the latter case translate
that point to the origin. The action is linear of
order $n=2,3,4,6$, as follows from its action on
the rank-two integral period lattice. The branch
indices of $B\to B/Q\simeq\PP^1$ are
\[
\begin{array}{c|c}
n&\text{branch indices}\\ \hline
2&(2,2,2,2)\\
3&(3,3,3)\\
4&(4,4,2)\\
6&(6,3,2).
\end{array}                                                   \tag{C.EQ.10}\label{cdp:eq:C.EQ.10}
\]
For completeness, the number fixed by a nontrivial
power is $|1-\zeta_n^k|^2$, the determinant of
that power minus one on the lattice. Removing
the points with larger stabilizer and dividing
by the corresponding orbit sizes gives exactly
\eqref{cdp:eq:C.EQ.10}. Each row also satisfies the degree-$n$
Riemann--Hurwitz equation.

\noindent\textbf{Step 10.} Write the fibre character as $\zeta_n^a$,
$0\le a<n$. At a branch index $e$, its exponent
is $a$ modulo $e$. Exponent zero gives a smooth
multiple. Otherwise put $c=\gcd(a,e)$; quotienting
the normal reflection kernel gives the entry of
\eqref{cdp:eq:C.EQ.8} with $d=e/c$ and $r=(a\bmod e)/c$.
Enumerating these actual characters gives
\[
\begin{array}{c|c|l}
n&a&\text{nonsmooth fibre configuration}\\ \hline
2&1&4\mathrm I_0^*\\
3&1&3\mathrm{IV}\\
3&2&3\mathrm{IV}^*\\
4&1&2\mathrm{III}+\mathrm I_0^*\\
4&2&2\mathrm I_0^*\\
4&3&2\mathrm{III}^*+\mathrm I_0^*\\
6&1&\mathrm{II}+\mathrm{IV}+\mathrm I_0^*\\
6&2,4&\mathrm{IV}+\mathrm{IV}^*\\
6&3&2\mathrm I_0^*\\
6&5&\mathrm{II}^*+\mathrm{IV}^*+\mathrm I_0^* .
\end{array}                                                   \tag{C.EQ.11}\label{cdp:eq:C.EQ.11}
\]
The zero character gives no nonsmooth fibre.
These finite counts, the local reflection reductions,
and the Euler sums are independently executed in
\path{code/root_cyclic_elliptic_types.py} and its
saved output.

\noindent\textbf{Step 11.} We finish by identifying the coherent invariants
without adding any unproved occurrence assertion.
The singular values just listed are all the possible
ones: a free base orbit is locally unchanged, and
the source had only smooth multiple fibres.
A free cyclic action on an elliptic base also leaves
only smooth multiples. Thus positive Euler occurs
only over $B_0=\PP^1$ and with nontrivial fibre
linear character. The exact smooth-multiple comparison
of \ref{cdp:st:C.2.UN.29} applies to $E_1$ and gives $p_g(E_1)=g(B)$.
The generically finite quotient comparison, resolved
by its actual graph, injects holomorphic two-forms
on $M$ into the invariant forms on $E_1$, as in
(C.PG.4)--(C.PG.5).
For $B=\PP^1$ this immediately gives $p_g(M)=0$.
For elliptic $B$, the unique two-form on $E_1$ is
the wedge of a base differential and the fibre
differential with fixed normalized periods. At a
smooth multiple fibre the same form extends with
the zero of order $m-1$ supplied by the root
coordinate; hence it is a global holomorphic form.
Its character is $\zeta_n^{a+1}$. Consequently
$p_g(M)=0$ whenever $a+1\not\equiv0\pmod n$.
Every Euler-twelve row of \eqref{cdp:eq:C.EQ.11} has precisely
that nontrivial character. This argument uses
pullback through the actual graph and Hartogs
extension, and does not infer extension of a
form from numerical invariance alone.

\noindent\textbf{Step 12.} On a relatively minimal elliptic surface the canonical
divisor is vertical. An actual meromorphic differential
on the base and a meromorphic relative differential
give such a divisor; restriction to a general compact
elliptic fibre has no zero or pole. The components
of the displayed reducible singular fibres are $(-2)$
rational curves. A type $\mathrm{II}$ fibre has arithmetic
genus one and square zero, and a smooth reduced multiple
fibre is an elliptic curve of square zero. Adjunction gives
$K_M\cdot C=0$ for each of them, and hence $K_M^2=0$.
Analytic Riemann--Roch gives
$12\chi(\OO_M)=e(M)$. Euler localization for the
elliptic map $p$ and \eqref{cdp:eq:C.EQ.8}--\eqref{cdp:eq:C.EQ.11} give
exactly the rows of \eqref{cdp:eq:C.EQ.1}, or zero. For Euler
twelve we proved $p_g=0$, so $\chi=1$ gives $q=0$.
For Euler twenty-four, \ref{cdp:st:C.2.UN.28} gives $p_g\le1$, while
$\chi=2=1-q+p_g$ and $q\ge0$ force $p_g=1,q=0$.
In the zero case
the local list has no nonsmooth reduced fibre,
since every such entry has positive Euler number.
This proves all assertions while preserving every
smooth multiple fibre and every actual quotient map.
\end{proof}






\begin{proposition}[Compact component cores give a multiplicity-weighted invariant lattice in one nearby torus]\label{cdp:st:C.2.UN.36}
\StatementRecord{cdp:st:C.2.UN.36}{proved}{Historical events: Claude Fable 5.2 (Anthropic), effort high, 16 November 2026: carried as background without fault found. GPT-6 Astra (OpenAI), ChatGPT Pro mode, effort Pro, S6\_09\_PAPERS, 17 November 2026: targeted proof check. Also checked in this writing. Pre-continuation writing events: GPT-6 Astra (OpenAI), ChatGPT Pro mode, effort Pro, S6\_10\_PAPERS, 18 November 2026: targeted proof check. Continuation mathematical and source checks: GPT-6 Astra (OpenAI), ChatGPT Pro mode, effort Pro, S6\_10\_PAPERS continuation, 19 November 2026: computation or formal certificate; independent re-proof. Continuation metadata events: GPT-6 Astra (OpenAI), ChatGPT Pro mode, effort Pro, S6\_10\_PAPERS continuation, 19 November 2026: used and rechecked in this run; status-vocabulary reconciliation; metadata reconciliation only; no new mathematical check. \emph{Pre-continuation recorded proof scope.} The integral Alexander map is exactly intersection with boundary curves, so its kernel classes have compact integral representatives. Normalized root-cover transfer multiplies each intersection form by the component multiplicity. Actual simultaneous horizontal flows map the cores into one nearby torus, with kernel in the radical and monodromy-invariant image. The determinant-square obstruction, every inertia bound and all four displayed exterior-square Gram cases were re-derived and independently replayed. The square-zero input UN.11 now has its own separately completed current review}{the hypotheses, notation and definitions stated here; \cref{cdp:eq:C.CL.8}; \cref{cdp:eq:C.CL.9}; \cref{cdp:eq:C.CL.7}; \cref{cdp:st:C.2.UN.33}; \cref{cdp:st:C.2.UN.11}; \cref{cdp:eq:C.CL.10}}

Let $f:U\to\Delta$ be a proper holomorphic map from
a smooth complex threefold, with connected two-torus
fibres over $\Delta^*$. Write its full central fibre
as $\sum_i m_iS_i$, and put $S=\bigcup_iS_i$ with
its reduced structure. For each component choose a
resolution of its normalization
$h_i:Z_i\to S_i\subset S$ which is an isomorphism
over $S_i\setminus\operatorname{Sing}S$. Set
\[
 \begin{gathered}
 D_i=h_i^{-1}(\operatorname{Sing}S)_{\rm red},\qquad
 R_i=\operatorname{span}_{\mathbb R}
       \{[C]:C\text{ is a curve component of }D_i\},
 \\
 K_i=\{v\in H_2(Z_i,\Z)/\operatorname{tors}:v\cdot R_i=0\}.
 \end{gathered}
                                                               \tag{C.CL.1}\label{cdp:eq:C.CL.1}
\]
The set $D_i$ may also have isolated points; they
contribute no curve class to $R_i$. The pairing on
$K_i$ is the restriction of the ordinary oriented
surface intersection form and may have a radical.

Choose a common base-change order $N$ divisible by
all $m_i$ and by the semisimple monodromy order.
Let $U_1$ be the monodromy on the integral first
homology of the actual nearby torus after this
base change. There is a homomorphism, obtained by
choosing compact representatives and their transfers,
\[
 \Psi:\bigoplus_iK_i\longrightarrow
       H_2(F,\Z)^{\Lambda^2U_1},\qquad
 \Psi(v)\cdot\Psi(w)=\sum_i m_i(v_i\cdot w_i).
                                                               \tag{C.CL.2}\label{cdp:eq:C.CL.2}
\]
Thus every multiplicity appears in the preserved
bilinear form. The kernel is contained in the
radical of that form. If its source is nondegenerate
of rank six, its image has finite integral index
in $H_2(F,\Z)$ and
\[
 \left|\det\left(\bigoplus_i m_iQ_{K_i}\right)\right|
                       \text{ is an integer square}.
                                                               \tag{C.CL.3}\label{cdp:eq:C.CL.3}
\]

Let $(p_i,q_i,z_i)$ be the positive, negative and
radical dimensions of the real form on $K_i$.
Let $(p_U,q_U,z_U)$ be those of the intersection
form restricted to
$\ker(\Lambda^2U_1-I)\subset H_2(F,\mathbb R)$.
Then
\[
             \sum_i p_i\le p_U\le3,\qquad
             \sum_i q_i\le q_U\le3.                 \tag{C.CL.4}\label{cdp:eq:C.CL.4}
\]
In particular, if $r_i=\dim R_i$, then
\[
 \sum_i\max\{b_2^-(Z_i)-r_i,0\}\le q_U\le3,
 \qquad
                 2\sum_i p_g(Z_i)\le p_U.           \tag{C.CL.5}\label{cdp:eq:C.CL.5}
\]
A K3 component whose chosen resolution is obtained
from its minimal model by $n$ point blowups must
therefore have at least $16+n$ independent boundary
curve classes. For an Enriques component or a
properly elliptic component of coherent Euler
characteristic one, the corresponding lower bound
is $6+n$.

The existing monodromy theorem gives
$(U_1-I)^2=0$. In an oriented real basis of
$H_1(F,\mathbb R)$, all its possible Jordan forms
give the following exact invariant forms:
\[
\begin{array}{c|c|c|c}
\operatorname{rank}(U_1-I)&\text{nonzero values of }U_1-I
 &\dim\ker(\Lambda^2U_1-I)&(p_U,q_U,z_U)\\ \hline
0&\text{none}&6&(3,3,0)\\
1&e_2\mapsto e_1&4&(1,1,2)\\
2&e_2\mapsto e_1,\ e_4\mapsto e_3&4&(1,2,1)\\
2&e_2\mapsto e_1,\ e_4\mapsto-e_3&4&(2,1,1).
\end{array}                                                   \tag{C.CL.6}\label{cdp:eq:C.CL.6}
\]
The two last rows record the two oriented real
forms of that Jordan type. In particular rank-one
unipotent monodromy forces $p_g(Z_i)=0$ for every
normalization component. All these assertions
apply to arbitrary actual central fibres, with
no normal-crossing, reducedness or K\"ahler
hypothesis on the original fibre.
\end{proposition}



\begin{proof}
\noindent\textbf{Step 1.} The open set
$Z_i^\circ=Z_i\setminus D_i$ maps isomorphically
to $S_i\setminus\operatorname{Sing}S$. It is
therefore precisely the part of this component
where the reduced support is one smooth embedded
surface. The compact analytic set $D_i$ consists
of curves and possibly isolated points. Its
degree-two integral cohomology is free, with
one generator for each curve component. Indeed
normalization of a compact reduced curve changes
only the identifications at finitely many points,
which do not change its degree-two cohomology;
isolated points have none.

\noindent\textbf{Step 2.} Poincar\'e--Alexander duality and the relative
homology sequence give the exact map
\[
 H_2(Z_i^\circ,\Z)\longrightarrow H_2(Z_i,\Z)
 \longrightarrow H_2(Z_i,Z_i^\circ;\Z)
       \simeq H^2(D_i,\Z).                          \tag{C.CL.7}\label{cdp:eq:C.CL.7}
\]
The second arrow, under these identifications,
is the vector of intersections with the curve
components of $D_i$. This follows by applying
Poincar\'e duality to the restriction map to
that compact analytic set. Torsion maps to
zero since the target is free. Consequently
every element of $K_i$ has an integral lift
represented by a compact finite cycle in
$Z_i^\circ$. Choose such lifts for an integral
basis of each $K_i$. Their mutual intersections
can be computed by compact transverse smooth
chain representatives in that same open set.
Include all the required perturbations in
one compact subset for each component.

\noindent\textbf{Step 3.} Make the actual change $t=s^N$, normalize,
and resolve the total space and central support,
isomorphically over the smooth open sets just
specified and over $s\ne0$. Let
$g:\widetilde U\to\Delta_s$ be the resulting
proper map from a smooth threefold. Over a
point of $S_i\setminus\operatorname{Sing}S$
the original local equation is
$t=u x^{m_i}$, with $u$ a unit. Its normalized
root charts are
\[
       w^{m_i}=u,\qquad x=s^{N/m_i}/w.               \tag{C.CL.8}\label{cdp:eq:C.CL.8}
\]
At $s=0$ the unit-root equation has exactly
$m_i$ local solutions, each nonsingular.
The full collection of strict components
above $S_i$ therefore restricts to a finite
\'etale covering
\[
        \pi_i:\widetilde Z_i^\circ
                       \longrightarrow Z_i^\circ
                  \quad\text{of total degree }m_i.  \tag{C.CL.9}\label{cdp:eq:C.CL.9}
\]
It may be disconnected. Every one of these
strict components has multiplicity one;
the root normal coordinate is $s$.
No modification used in the resolution
changes \eqref{cdp:eq:C.CL.8}--\eqref{cdp:eq:C.CL.9}.

\noindent\textbf{Step 4.} Transfer an integral compact cycle by taking
the sum of all its lifts under \eqref{cdp:eq:C.CL.9},
subdividing its finite singular chain into
evenly covered pieces. This is the ordinary
integral covering transfer and is a chain
map. For two transverse compact cycles,
every signed intersection point downstairs
has exactly $m_i$ lifts with the same local
orientation. Hence their transferred
intersection is exactly $m_i$ times the
original intersection. This proves the
multiplicity factor directly, including
disconnected root covers.

\noindent\textbf{Step 5.} The transferred cycles lie in the smooth
single-component locus of the central
support of $\widetilde U$. The map $g$ is
a submersion on a neighbourhood of their
compact union. Choose smooth horizontal
lifts of the two real coordinate fields
on the base disc. A partition of unity
preserves their specified images under
$dg$. After a cutoff outside a slightly
larger neighbourhood, their successive
real and imaginary flows give, for every
sufficiently small $s$, an actual smooth
embedding of a neighbourhood of the compact
cycles into $g^{-1}(s)$. The image lies in
that exact fibre, and the embedding preserves
the complex orientation by continuity.
Choose disjoint ambient neighbourhoods for
the compact subsets in distinct strict
components, and the same small disc for
all flows. Their images then remain
disjoint in one nearby torus $F$.

\noindent\textbf{Step 6.} Transport the selected transferred basis
cycles and extend by integer linearity.
This defines the homomorphism $\Psi$.
The transfer intersection calculation and
the disjoint neighbourhoods give the
bilinear equality of \eqref{cdp:eq:C.CL.2}. The choice
of lifts is part of this construction;
no canonical splitting of \eqref{cdp:eq:C.CL.7} is
needed.

\noindent\textbf{Step 7.} The image is invariant under the actual
nearby monodromy. In fact the two-flow
construction defines a smooth family of
each finite cycle over the \emph{whole}
small $s$-disc. Its restriction to a
circle in that disc returns to the
identical parametrized cycle. In any
local differentiable trivialization of
the proper smooth family over $s\ne0$,
this continuous family has a locally
constant integral homology class.
Thus the class returns unchanged under
Gauss--Manin transport. The actual
identification
$H_2(F,\Z)\simeq\Lambda^2H_1(F,\Z)$
intertwines transport, so its image is
in the group written in \eqref{cdp:eq:C.CL.2}.
Only these explicitly transported
cycles enter the assertion.

\noindent\textbf{Step 8.} If $\Psi(v)=0$, the preserved pairing
with every $\Psi(w)$ shows that $v$
is in the radical of the source form.
Choose a real complement of that radical
on which the restricted form is
nondegenerate. Its image is injective
and has the same positive and negative
indices. Maximal positive and negative
subspaces of the target restricted form
therefore give \eqref{cdp:eq:C.CL.4}. If the integral
source is nondegenerate of rank six,
$\Psi$ is an injection of full-rank
free lattices into $H_2(F,\Z)$.
The latter has the unimodular form
$3H$, of determinant $-1$. If the
image has index $c$, the change-of-basis
determinant gives source determinant
$-c^2$. This proves \eqref{cdp:eq:C.CL.3}.

\noindent\textbf{Step 9.} Intersect a maximal negative definite
subspace of $H_2(Z_i,\mathbb R)$ with
the codimension-$r_i$ subspace $K_i\otimes\mathbb R$.
Its dimension is at least
$b_2^-(Z_i)-r_i$, so its negative index
is at least the corresponding nonnegative
maximum in \eqref{cdp:eq:C.CL.5}. For the second
inequality there, the real and imaginary
parts of holomorphic two-forms give the
positive $2p_g(Z_i)$-dimensional subspace
proved in (C.PG.7)--(C.PG.8). Their
integrals on every boundary curve
vanish by type; their dual homology
classes therefore belong to
$K_i\otimes\mathbb R$ by \eqref{cdp:eq:C.CL.7}.
This proves \eqref{cdp:eq:C.CL.5}.

\noindent\textbf{Step 10.} A K3 minimal surface has intersection
signature $(3,19)$; $n$ point blowups
add $n$ orthogonal classes of square
$-1$. For Enriques surfaces the
signature is $(1,9)$. A properly
elliptic component with $\chi(\OO)=1$
has $q=p_g=0$ and $e=12$ by \ref{cdp:st:C.2.UN.33},
and signature $(1,9)$ as well:
$b_1=0$, $b_2=10$, and the signature
formula with $K^2=0$ gives signature
$-8$. Applying the negative-index
bound to $19+n$ or $9+n$ gives
the asserted boundary ranks.

\noindent\textbf{Step 11.} It remains to compute the exact
monodromy bound. \ref{cdp:st:C.2.UN.11} proves
that $N_0=U_1-I$ has square zero
and rank at most two. In rank one,
choose a vector and its nonzero
image, then complete the basis in
the kernel. Changing the sign of
one unused basis vector makes the
basis positively oriented, giving
the second row of \eqref{cdp:eq:C.CL.6}.
In rank two, choose two image
vectors and their preimages. If
their resulting ordered basis is
negative, reverse one preimage.
This gives the two oriented forms
\[
          N_0e_2=e_1,\qquad N_0e_4=\epsilon e_3,
                     \qquad\epsilon\in\{1,-1\}.     \tag{C.CL.10}\label{cdp:eq:C.CL.10}
\]
Write $e_{ij}=e_i\wedge e_j$.
In the oriented wedge basis
$(e_{12},e_{13},e_{14},e_{23},e_{24},e_{34})$
the only nonzero complementary
pairings are
\[
 e_{12}\cdot e_{34}=1,\qquad
 e_{13}\cdot e_{24}=-1,\qquad
 e_{14}\cdot e_{23}=1.                              \tag{C.CL.11}\label{cdp:eq:C.CL.11}
\]
For rank one the invariant space
has basis $(e_{12},e_{13},e_{14},e_{34})$;
its form is one hyperbolic plane
and a two-dimensional radical.
For \eqref{cdp:eq:C.CL.10} an invariant basis is
\[
       e_{12},\ e_{34},\ e_{13},\ e_{14}-\epsilon e_{23}.
                                                               \tag{C.CL.12}\label{cdp:eq:C.CL.12}
\]
Its Gram matrix is
\[
 \begin{pmatrix}
 0&1&0&0\\
 1&0&0&0\\
 0&0&0&0\\
 0&0&0&-2\epsilon
 \end{pmatrix}.                                               \tag{C.CL.13}\label{cdp:eq:C.CL.13}
\]
These are exactly the rows of
\eqref{cdp:eq:C.CL.6}; the identity row is
the whole form $3H$.
The exterior-square matrices,
their exact kernels, and all
four invariant inertias are
executed over the integers and
rationals in
\path{code/A_compact_core_lattices.py}
and its saved output. Finally
the rank-one row has $p_U=1$,
so \eqref{cdp:eq:C.CL.5} forces every
$p_g(Z_i)$ to vanish, as claimed.
\end{proof}






\begin{proposition}[Every curve of a smooth minimal algebraic-dimension-zero Hopf normalization lies in the fibre boundary]\label{cdp:st:C.2.UN.38}
\StatementRecord{cdp:st:C.2.UN.38}{proved}{Historical events: Claude Fable 5.2 (Anthropic), effort high, 16 November 2026: carried as background without fault found. GPT-6 Astra (OpenAI), ChatGPT Pro mode, effort Pro, S6\_09\_PAPERS, 17 November 2026: targeted proof check. Also checked in this writing. Pre-continuation writing events: GPT-6 Astra (OpenAI), ChatGPT Pro mode, effort Pro, S6\_10\_PAPERS, 18 November 2026: targeted proof check; used and rechecked in this run. Continuation mathematical and source checks: GPT-6 Astra (OpenAI), ChatGPT Pro mode, effort Pro, S6\_10\_PAPERS continuation, 19 November 2026: targeted proof check. Continuation metadata events: GPT-6 Astra (OpenAI), ChatGPT Pro mode, effort Pro, S6\_10\_PAPERS continuation, 19 November 2026: used and rechecked in this run; status-vocabulary reconciliation; metadata reconciliation only; no new mathematical check. \emph{Pre-continuation recorded proof scope.} The actual vertical canonical section and full fibre equation give the neighborhood torsion isomorphism. The diagonal normal character is non-torsion by its Laurent equation. In the non-diagonal case the coefficient of z to the minus k minus one at order j plus r gives the explicit nonzero obstruction. Remmert-Stein closure, the convergent unit cocycle and weighted eigenfunction calculation identify all curves, so every absent axis would be disjoint from the full boundary. The actual finite unramified primary cover pulls back the same canonical trivialization. Nakamura section1.1 Theorem1.1.2 p2 and Khaled's primary normal forms were matched}{the hypotheses, notation and definitions stated here; \cref{cdp:eq:C.HB.2}; \cite[Theorem~2.1]{Khaled2024Hopf}; \cite[\S1.1]{Nakamura2008VII}; \cref{cdp:eq:C.HB.3}; \cref{cdp:eq:C.HB.5}; \cref{cdp:eq:C.HB.4}}

Let $f:U\to\Delta$ be a proper degeneration of
connected complex two-tori from a smooth threefold,
and write its full central fibre as $\sum_i m_iS_i$.
On the open subset of $S_i$ where the reduced
support is smooth, its canonical line has an
actual trivial $m_i$-th power in a neighborhood
of every compact subset. In particular, a
smooth compact elliptic curve $C$ in that open
subset satisfies
\[
                         N_{C/S_i}^{\otimes m_i}
                                      \simeq\OO_C.    \tag{C.HB.1}\label{cdp:eq:C.HB.1}
\]

Suppose the normalization of a central component
is itself a smooth minimal Hopf surface of
algebraic dimension zero. Then every elliptic
curve on that normalization is contained in
the inverse image of the singular locus of
the reduced support. For a diagonal primary
Hopf surface this includes both axes. For a
non-diagonal primary Hopf surface it includes
its unique axis; in this case the conclusion
uses the full canonical neighborhood and
the explicit higher-jet calculation below.
The same conclusion descends through every
actual finite primary Hopf covering.
\end{proposition}



\begin{proof}
\noindent\textbf{Step 1.} The proper direct image $f_*K_U$ is a
torsion-free coherent sheaf of generic rank
one on the disc and hence is a line bundle.
Choose a generator after shrinking the disc.
Its evaluation is a holomorphic section
$\omega$ of $K_U$. On each nonzero fibre,
proper base change identifies this evaluation
with a nonzero canonical form on that torus,
up to the fixed base differential. It has
no zero there. Consequently its exact divisor
is vertical:
\[
                   \operatorname{div}(\omega)
                               =\sum_i k_iS_i,
                        \qquad k_i\ge0.              \tag{C.HB.2}\label{cdp:eq:C.HB.2}
\]
Let $W\subset U$ be a neighborhood of a
compact subset in the indicated regular
part of $S=S_i$, small enough to avoid all
other components and the singularities of
$S$. On $W$ the actual fibre equation has
divisor $mS$, where $m=m_i$. Thus it gives
the trivialization of $\OO_W(mS)$, and
\eqref{cdp:eq:C.HB.2} gives
\[
 \begin{gathered}
 K_U|_W\simeq\OO_W(k_iS),\\
 K_S|_{S\cap W}\simeq
                  \OO_W((k_i+1)S)|_{S\cap W},\\
 (K_S|_{S\cap W})^{\otimes m}\simeq\OO_{S\cap W}.
 \end{gathered}
                                                               \tag{C.HB.3}\label{cdp:eq:C.HB.3}
\]
These are line-bundle isomorphisms on the
neighborhood, obtained from the actual
canonical and fibre sections. Adjunction
for an elliptic $C\subset S\cap W$ gives
$N_{C/S}=K_C\otimes K_S^{-1}|_C$.
Since $K_C=\OO_C$, \eqref{cdp:eq:C.HB.1} follows.

\noindent\textbf{Step 2.} Consider first the diagonal primary
Hopf surface
\[
 H_{\alpha,\beta}=(\C^2\setminus\{0\})/
       \langle(z,w)\mapsto(\alpha z,\beta w)\rangle,
                    \qquad0<|\alpha|,|\beta|<1.
                                                               \tag{C.HB.4}\label{cdp:eq:C.HB.4}
\]
The axis $C_\alpha$ is $\C^*/\alpha^{\Z}$,
and its actual normal line is the quotient
of $\C^*\times\C$ by
$(z,v)\mapsto(\alpha z,\beta v)$.
If a positive tensor power $n$ of this
line is trivial, a nowhere-zero holomorphic
function $a$ on $\C^*$ satisfies
$a(\alpha z)=\beta^n a(z)$.
In its Laurent expansion each nonzero
coefficient of $z^j$ requires
$\alpha^j=\beta^n$. Since $|\alpha|<1$,
there is at most one such exponent, and
there must be one. Its exponent is positive
because $|\beta|<1$ and $n>0$.
Conversely that equality gives the
explicit trivialization $a(z)=z^j$.
It also gives the nonconstant invariant
meromorphic function $w^n/z^j$ on the
Hopf surface. Algebraic dimension zero
therefore makes the normal of each axis
non-torsion. Formula \eqref{cdp:eq:C.HB.1} excludes
an entire such axis from the regular
part of the reduced support.

\noindent\textbf{Step 3.} The non-diagonal primary normal form
requires a further argument. The primary
normal forms and finite unramified cover
used here are the explicit ones recalled
in \cite[Theorem~2.1]{Khaled2024Hopf}
and \cite[\S1.1]{Nakamura2008VII}.
Write the non-diagonal contraction
as
\[
 H=(\C^2\setminus\{0\})/\langle F\rangle,
 \qquad F(z,w)=(\beta^r z+w^r,\beta w),
                  \quad r\ge1,\quad0<|\beta|<1,
                                                               \tag{C.HB.5}\label{cdp:eq:C.HB.5}
\]
and let $C$ be the image of $w=0$.
The invariant meromorphic two-form
$dz\wedge dw/w^{r+1}$ gives the actual
canonical divisor
\[
                         K_H\simeq\OO_H(-(r+1)C).
                                                               \tag{C.HB.6}\label{cdp:eq:C.HB.6}
\]
Its normal line has torsion, so \eqref{cdp:eq:C.HB.1}
alone gives no contradiction. Instead,
\eqref{cdp:eq:C.HB.3} on a neighborhood of $C$ would
make $\OO(m(r+1)C)$ trivial there.
Its canonical divisor section would
then be a holomorphic function with
positive divisor $m(r+1)C$ on that
neighborhood. We show directly that
no such function exists.

\noindent\textbf{Step 4.} Lift it to the invariant neighborhood
of $\C^*\times\{0\}$ in the covering
space and expand in the normal variable:
\[
 h(z,w)=\sum_{l\ge j}h_l(z)w^l,
       \qquad j>0,\qquad h_j\ne0,
       \qquad h_l\in\OO(\C^*),
       \qquad h\circ F=h.                           \tag{C.HB.7}\label{cdp:eq:C.HB.7}
\]
The radius of convergence may depend
on $z$; the Taylor coefficients are
nevertheless holomorphic on all of
$\C^*$ because the lifted neighborhood
contains that entire axis. The leading
coefficient satisfies
$\beta^j h_j(\beta^r z)=h_j(z)$.
Its Laurent expansion forces
\[
                      j=rk,\qquad
                      h_j(z)=c z^{-k},
                      \qquad k\ge1,\quad c\ne0.
                                                               \tag{C.HB.8}\label{cdp:eq:C.HB.8}
\]
At order $j+r$, the same invariant
function equation is exactly
\[
 \beta^{j+r}h_{j+r}(\beta^r z)-h_{j+r}(z)
                    =k c\beta^{-r}z^{-k-1}.          \tag{C.HB.9}\label{cdp:eq:C.HB.9}
\]
Indeed the only extra contribution
at that order is the first derivative
of $h_j(\beta^r z+w^r)\beta^j$;
all other Taylor terms start at
higher order. In the Laurent
expansion of the left side, the
coefficient multiplying $z^{-k-1}$
is zero, since
$\beta^{j+r+r(-k-1)}-1=0$.
The right side has the nonzero
coefficient $k c\beta^{-r}$.
This contradiction proves the
neighborhood obstruction for
the non-diagonal axis as well.

\noindent\textbf{Step 5.} Finally let $\nu:H\to S$ be the
given smooth minimal Hopf
normalization. Its conductor
and the inverse images of
the other central components
have curve support. Outside
this support $\nu$ is the
actual isomorphism onto the
regular locus of the reduced
fibre. For clarity, there
is no extra isolated
normalization defect: a
hypersurface is Cohen--Macaulay,
so its nonnormal locus is
the conductor curve locus;
where it is normal, the
normalization is an isomorphism,
and smoothness of $H$ then
gives smoothness there.
An intersection with another
central component is a
Cartier curve divisor and
its finite pullback also
has curve support.

\noindent\textbf{Step 6.} We also verify the curve assertion
used in applying the obstruction.
The inverse image of a compact
curve in the covering punctured
plane extends analytically across
the origin by the removable-set
theorem for analytic sets, since
its dimension is one and the
added set has dimension zero.
Its germ has a reduced equation
$a$ satisfying $a\circ F=v a$
for a holomorphic unit $v$.
After shrinking to a contracting
ball, the product solving the
unit cocycle converges: if
$u=v/v(0)$, then
$b=\prod_{j\ge0}u\circ F^j$
is a holomorphic unit and
$b\circ F=b/u$.
Thus $h=ba$ satisfies
$h\circ F=v(0)h$.
For the diagonal contraction,
two distinct nonzero monomials
in $h$ would give
$\alpha^i\beta^j=1$ for a
nonzero integer pair; because
both moduli are less than one,
this is a positive-power
relation between $\alpha$
and $\beta$ and gives a
nonconstant meromorphic
invariant. In algebraic
dimension zero there is
therefore just one monomial,
so the curve consists of
axes. For \eqref{cdp:eq:C.HB.5}, decompose
the convergent power series
into weighted degrees with
weights $r,1$ on $z,w$.
The contraction preserves
each finite-dimensional
weighted piece and acts
there with sole eigenvalue
$\beta^d$ in degree $d$.
Their distinct moduli force
an eigenfunction into one
degree. Write it as
$w^dP(z/w^r)$; after taking
out $\beta^d$, invariance
requires the polynomial
$P$ to be periodic under
translation by $\beta^{-r}$.
It is constant. Thus only
the axis $w=0$ occurs.
Analytic continuation along
the contraction orbits
gives this assertion on
the entire punctured plane.
Hence the only curves
are the two disjoint axes
in \eqref{cdp:eq:C.HB.4}, or the single
axis in \eqref{cdp:eq:C.HB.5}.
Therefore an axis not
contained in that boundary
is disjoint from it. The
preceding neighborhood
contradictions apply and
force every axis into the
boundary. For a secondary
Hopf surface, pull back
the neighborhood and its
canonical trivialization
along its finite unramified
primary covering. The
covering has the same
algebraic dimension:
meromorphic functions
descend through their
elementary symmetric
polynomials. A curve
outside the boundary
would lift to an axis
outside the pulled-back
boundary, contradicting
the already proved
primary calculation.
This proves the stated
descent and the theorem.
\end{proof}






\begin{proposition}[Curve genera, triple points and the finite monodromy constant decide zero Euler number]\label{cdp:st:C.2.UN.39}
\StatementRecord{cdp:st:C.2.UN.39}{proved}{Historical events: Claude Fable 5.2 (Anthropic), effort high, 16 November 2026: carried as background without fault found. GPT-6 Astra (OpenAI), ChatGPT Pro mode, effort Pro, S6\_09\_PAPERS, 17 November 2026: targeted proof check. Also checked in this writing. Pre-continuation writing events: GPT-6 Astra (OpenAI), ChatGPT Pro mode, effort Pro, S6\_10\_PAPERS, 18 November 2026: targeted proof check. Continuation mathematical and source checks: GPT-6 Astra (OpenAI), ChatGPT Pro mode, effort Pro, S6\_10\_PAPERS continuation, 19 November 2026: computation or formal certificate; independent re-proof. Continuation metadata events: GPT-6 Astra (OpenAI), ChatGPT Pro mode, effort Pro, S6\_10\_PAPERS continuation, 19 November 2026: used and rechecked in this run; status-vocabulary reconciliation; metadata reconciliation only; no new mathematical check. \emph{Pre-continuation recorded proof scope.} Averaging the exact nearby-cycle trace over a common multiple of all component multiplicities gives the open-stratum sum. Double-curve inclusion-exclusion and the actual projective fibres of each point or curve blowup give every sign and centre correction. This identifies the index with (j(T)-e(S))/2 for each actual resolution. All nineteen finite monodromy averages were independently recomputed from exact companion matrices; the point-plus-quartic example verifies the geometric correction}{the hypotheses, notation and definitions stated here; \cref{cdp:st:C.2.UN.11}; \cref{cdp:eq:C.ZC.5}; \cref{cdp:eq:C.ZC.6}; \cref{cdp:eq:C.ZC.8}; \cref{cdp:eq:C.ZC.7}; \cref{cdp:eq:C.ZC.9}; \cref{cdp:st:C.2.UN.4}}

Let $f:U\to\Delta$ be an actual proper holomorphic
two-torus degeneration with smooth total threefold,
and let $S=f^{-1}(0)_{\rm red}$. Choose an embedded
resolution $\pi:\widetilde U\to U$ as a sequence
of blowups with smooth centres contained in the
successive central supports. Write its full
central divisor as $\sum_jN_jE_j$ with simple
normal crossings. Let $\mathcal D$ be the set
of connected double curves of $\bigcup_jE_j$,
and let $t_3$ be its number of triple points.
Count every point centre at every stage of
$\pi$, obtaining an integer $p_\pi$, and list
every connected smooth curve centre as $C_\alpha$.
Define the integer
\[
 I_\pi=
     \sum_{D\in\mathcal D}\bigl(g(D)-1\bigr)+t_3
       +p_\pi+\sum_\alpha\bigl(1-g(C_\alpha)\bigr).
                                                               \tag{C.ZC.1}\label{cdp:eq:C.ZC.1}
\]
It is independent of the chosen embedded
resolution, including additional vertical
blowups used to obtain another such resolution.

Let $T$ be the actual monodromy on $H^1(F,\Z)$
of a nearby torus, and let $M$ be the order of
its semisimple part. Put
\[
             j(T)=\frac1M\sum_{k=1}^{M}\det(I-T^k).
                                                               \tag{C.ZC.2}\label{cdp:eq:C.ZC.2}
\]
Then
\[
                e(S)=j(T)-2I_\pi,\qquad
                e(S)=0\quad\Longleftrightarrow\quad
                         I_\pi=\frac{j(T)}2 .        \tag{C.ZC.3}\label{cdp:eq:C.ZC.3}
\]
The constant $j(T)$ belongs to $\{0,4,6,8\}$.
For all nineteen monodromy polynomials its
value and the required zero index are
\[
 \begin{array}{c|r|r|r}
\det(xI-T)&M&j(T)&I_\pi\text{ for }e(S)=0\\ \hline
\Phi_{1}^{4}&1&0&0\\
\Phi_{1}^{2} \Phi_{2}^{2}&2&0&0\\
\Phi_{1}^{2} \Phi_{3}&3&0&0\\
\Phi_{1}^{2} \Phi_{4}&4&0&0\\
\Phi_{1}^{2} \Phi_{6}&6&0&0\\
\Phi_{2}^{4}&2&8&4\\
\Phi_{3}^{2}&3&6&3\\
\Phi_{2}^{2} \Phi_{4}&4&4&2\\
\Phi_{4}^{2}&4&6&3\\
\Phi_{5}&5&4&2\\
\Phi_{2}^{2} \Phi_{3}&6&4&2\\
\Phi_{2}^{2} \Phi_{6}&6&4&2\\
\Phi_{3} \Phi_{6}&6&4&2\\
\Phi_{6}^{2}&6&6&3\\
\Phi_{8}&8&4&2\\
\Phi_{10}&10&4&2\\
\Phi_{3} \Phi_{4}&12&4&2\\
\Phi_{4} \Phi_{6}&12&4&2\\
\Phi_{12}&12&4&2\\
\end{array}
                                                               \tag{C.ZC.4}\label{cdp:eq:C.ZC.4}
\]
Thus this is a necessary and sufficient
geometric test for the zero Euler number of
every actual fibre, using only the indicated
curve genera, triple points, resolution centres
and finite monodromy constant. Arbitrary
component and centre multiplicities are
included through the full divisor and the
actual blowup maps. The statement classifies
this numerical property of actual fibres;
it makes no occurrence assertion for abstract
genus or monodromy data and does not enumerate
analytic isomorphism classes of supports.
\end{proposition}



\begin{proof}
\noindent\textbf{Step 1.} An embedded resolution of the specified kind
exists by the analytic resolution theorem used
in \ref{cdp:st:C.2.UN.11}, with all centres over the central
support and with an isomorphism over
$\Delta^*$. Smooth divisor centres need
not be included, since their blowups are
isomorphisms. The remaining connected centres
in a threefold are compact smooth curves or
points. A disconnected centre can be counted
by its disjoint connected components; their
blowups commute. In particular every genus
and count in \eqref{cdp:eq:C.ZC.1} is attached to an actual
compact centre or curve in this resolution.

\noindent\textbf{Step 2.} We first compute the resolved support without
any reduction of its multiplicities. Put
\[
             E_j^\circ=E_j\setminus\bigcup_{i\ne j}E_i.
\]
The proper analytic A'Campo comparison, with
its actual nearby-cycle maps as established
in (C.MS.6), says for every positive $k$ that
\[
         \det(I-T^k)=\sum_{N_j\mid k}N_j e(E_j^\circ).
                                                               \tag{C.ZC.5}\label{cdp:eq:C.ZC.5}
\]
Here the full integer $N_j$ is used for each
component. Choose $L$ divisible by $M$ and
by every $N_j$. The determinant on the left
depends only on the eigenvalues of $T$,
so is periodic with period $M$ even when
its unipotent part is nontrivial. Average
\eqref{cdp:eq:C.ZC.5} over $1\le k\le L$. Exactly $L/N_j$
of these integers are divisible by $N_j$.
It follows that
\[
 \begin{split}
 j(T)
 &=\frac1L\sum_{k=1}^{L}\det(I-T^k)\\
 &=\sum_j\frac{N_j}{L}\frac{L}{N_j}e(E_j^\circ)
   =\sum_j e(E_j^\circ).
 \end{split}                                                   \tag{C.ZC.6}\label{cdp:eq:C.ZC.6}
\]
This calculation includes components whose
multiplicity does not divide $M$; choosing
$L$ retains them before the averaging.

\noindent\textbf{Step 3.} Let $\widetilde S=\bigcup_jE_j$. Its
remaining strata are the double curves
with their triple points removed, and
the triple points themselves. Every
triple point lies on exactly three
double curves. Since each connected
smooth compact double curve $D$ has
Euler number $2-2g(D)$, additivity gives
\[
 \begin{split}
 e(\widetilde S)
 &=\sum_j e(E_j^\circ)
       +\sum_{D\in\mathcal D}(2-2g(D))-3t_3+t_3\\
 &=j(T)-2\left(
           \sum_{D\in\mathcal D}(g(D)-1)+t_3\right).
 \end{split}                                                   \tag{C.ZC.7}\label{cdp:eq:C.ZC.7}
\]
The surface Euler numbers have disappeared:
the open surface contribution is the finite
monodromy average computed in \eqref{cdp:eq:C.ZC.6}.

\noindent\textbf{Step 4.} We now compare the original support and
the resolved one through each actual
blowup. Let $\sigma:U'\to U_0$ blow up
a connected smooth centre $Z$ contained
in a current reduced central support
$S_0$. Its reduced inverse image is
the next central support $S'$. Away
from $Z$ the map is an isomorphism;
over every point of $Z$ its fibre is
$\mathbb P^{c-1}$, where $c$ is the
complex codimension of $Z$ in $U_0$.
The entire projective fibre belongs
to $S'$, because the centre is contained
in $S_0$. Therefore the actual proper
map, stratified by $Z$ and its complement,
gives
\[
                    e(S')-e(S_0)=(c-1)e(Z).
                                                               \tag{C.ZC.8}\label{cdp:eq:C.ZC.8}
\]
For a point this is $2$; for a smooth
connected curve it is $2-2g(Z)$.
This computation is valid for every
multiplicity of the full fibre along
the centre. More explicitly, if the
generic order of its local equation
along the centre is $\mu$, its pullback
has the exceptional factor of order
$\mu$. That factor changes the full
divisor coefficient but not the
projective fibre in its reduced
inverse image. No value of $\mu$
was assumed in \eqref{cdp:eq:C.ZC.8}.

\noindent\textbf{Step 5.} Iterating the actual proper maps in
$\pi$ gives
\[
 e(\widetilde S)=e(S)+2p_\pi
                        +2\sum_\alpha(1-g(C_\alpha)).
                                                               \tag{C.ZC.9}\label{cdp:eq:C.ZC.9}
\]
Combining \eqref{cdp:eq:C.ZC.7} and \eqref{cdp:eq:C.ZC.9}
proves \eqref{cdp:eq:C.ZC.3} with exactly the
integer in \eqref{cdp:eq:C.ZC.1}. Every such
resolution is an isomorphism over
the punctured disc and therefore
has the same actual monodromy.
The formula consequently identifies
$I_\pi$ with $(j(T)-e(S))/2$ for
each choice. This proves independence
by the actual blowup comparisons;
no comparison between unrelated
fillings is required.

\noindent\textbf{Step 6.} The nineteen possible characteristic
polynomials and their one-period
determinants were computed in \ref{cdp:st:C.2.UN.11}.
Applying the finite average \eqref{cdp:eq:C.ZC.2}
gives the entire table \eqref{cdp:eq:C.ZC.4}.
Equivalently, $j=0$ when the polynomial
has a factor $\Phi_1$; $j=8$ for
$\Phi_2^4$; $j=6$ for
$\Phi_3^2,\Phi_4^2,\Phi_6^2$;
and $j=4$ for the ten other rows.
The exact independent calculation
in \path{code/A_zero_curve_index.py}
checks all nineteen finite averages
against the sums of the A'Campo
signature coefficients, and writes
the displayed table and
\path{records/A_zero_curve_index.csv}.
All arithmetic is integral or rational.

\noindent\textbf{Step 7.} For a concrete check of the geometric
terms, start from a smooth torus support
with identity monodromy. A point blowup
creates one rational double curve and
no triple point. The new support has
index $-1$ and Euler number $2$.
Follow it by the genus-three quartic
blowup of \ref{cdp:st:C.2.UN.4}. There are then five
rational double curves, one double
curve of genus three, and four triple
points. Their index is
$-5+(3-1)+4=1$, giving Euler number
$-2$. Viewed instead as a resolution
of the original smooth support, these
two centres contribute
$1+(1-3)=-1$ to the centre term,
so \eqref{cdp:eq:C.ZC.1} returns index zero for
that original support. These exact
charts verify both the sign of the
centre correction and its compatibility
with the universal balancing operations.
\end{proof}






\begin{proposition}[The fundamental group and the component-incidence constraint of every torus fibre]\label{cdp:st:C.2.UN.44}
\StatementRecord{cdp:st:C.2.UN.44}{proved}{Historical events: Claude Fable 5.2 (Anthropic), effort high, 16 November 2026: carried as background without fault found. GPT-6 Astra (OpenAI), ChatGPT Pro mode, effort Pro, S6\_09\_PAPERS, 17 November 2026: targeted proof check. Also checked in this writing. Pre-continuation writing events: GPT-6 Astra (OpenAI), ChatGPT Pro mode, effort Pro, S6\_10\_PAPERS, 18 November 2026: targeted proof check; used and rechecked in this run. Continuation mathematical and source checks: GPT-6 Astra (OpenAI), ChatGPT Pro mode, effort Pro, S6\_10\_PAPERS continuation, 19 November 2026: computation or formal certificate; targeted proof check. Continuation metadata events: GPT-6 Astra (OpenAI), ChatGPT Pro mode, effort Pro, S6\_10\_PAPERS continuation, 19 November 2026: used and rechecked in this run; status-vocabulary reconciliation; metadata reconciliation only; no new mathematical check. \emph{Pre-continuation recorded proof scope.} The local analytic gradient inequality and finite-length flow produce the actual strong tube retraction. Transverse discs give precisely the normal closure of component meridians; the normalized root cover is connected and etale in the explicit equation s=vh. Reidemeister-Schreier conjugates and the Bezout combination give the exact lattice quotient. Actual blowdown maps preserve fundamental groups. The triangulated mapping cylinders and triple-point triangles give the dual-complex quotient. Its degree-one map into torus cohomology is injective and preserves products; pulling back the finite covering bounds the rank of its finite-index abelian subgroup by two. Kurdyka Theorem(LI) p775 and Lojasiewicz Theorem2 p464 match the required analytic hypotheses}{the hypotheses, notation and definitions stated here; \cref{cdp:st:C.2.UN.39}; \cite[Theorem (LI), p.~775]{Kurdyka1998}; \cref{cdp:eq:C.FG.9}; \cref{cdp:eq:C.FG.10}; \cite[\S3, Theorem~2]{Lojasiewicz1964}; \cref{cdp:eq:C.FG.14}; \cref{cdp:eq:C.FG.15}; \cref{cdp:eq:C.FG.12}; \cref{cdp:eq:C.FG.17}; \cref{cdp:eq:C.FG.18}}

Let $f:U\to\Delta$ be a proper holomorphic degeneration of connected
complex two-tori, with smooth total threefold. Write
\[
 f^{-1}(0)=\sum_{i=1}^r m_iS_i,\qquad
 S=\bigcup_iS_i,\qquad d=\gcd(m_1,\ldots,m_r).
                                                        \tag{C.FG.1}\label{cdp:eq:C.FG.1}
\]
All fundamental groups below are based; changing the paths to the
chosen base point changes the indicated presentations by their actual
change-of-base-point isomorphisms. Choose a lift $t$ of the positive
generator of $\pi_1(\Delta^*)$, put $\Lambda=\pi_1(F)=\Z^4$, and
write $T$ for conjugation by $t$ on $\Lambda$. There are vectors
$v_i\in\Lambda$, obtained from positively oriented meridians of the
actual components, for which
\[
 \Gamma:=\pi_1(S)
  \simeq
  \bigl(\Lambda\rtimes_T\langle t\rangle\bigr)
       \big/\langle\!\langle x^{v_i}t^{m_i}:1\le i\le r
                    \rangle\!\rangle .                \tag{C.FG.2}\label{cdp:eq:C.FG.2}
\]
Here $x^v$ denotes the element $v$ of the additive lattice written
multiplicatively. Its image $A$ in $\Gamma$ is abelian, and there is
an exact sequence
\[
       1\longrightarrow A\longrightarrow\Gamma
          \longrightarrow\Z/d\Z\longrightarrow1 .    \tag{C.FG.3}\label{cdp:eq:C.FG.3}
\]
In particular the index is exactly $d$ and $A$ is generated by at most
four elements. There is an exact integral description as follows.
Set $n_i=m_i/d$, choose integers $c_i$ with $\sum_i c_in_i=1$, and
put $v_0=\sum_i c_iv_i$. Then
\[
 A\simeq\Lambda/R_d,\qquad
 R_d=(T^d-I)\Lambda+
       \left\langle T^jv_i-n_iv_0:
                         1\le i\le r,\ 0\le j<d\right\rangle_{\Z}.
                                                        \tag{C.FG.4}\label{cdp:eq:C.FG.4}
\]
For $d=1$, the whole group is abelian and this specializes to
\[
 \pi_1(S)\simeq
   \Lambda\big/\left((T-I)\Lambda+
                   \langle v_i-m_iv_0\rangle_{\Z}\right).
                                                        \tag{C.FG.5}\label{cdp:eq:C.FG.5}
\]
The normalized base change $s^d=f$ is an actual connected degree-$d$
\emph{\'etale} cover of the total space near $S$. Its reduced central
support $S_d$ is connected, its restriction $S_d\to S$ is the
connected cyclic covering belonging to the subgroup $A$, and
$\pi_1(S_d)=A$ under that covering injection.

For every embedded resolution of the central divisor by smooth
vertical blowups, let $\mathcal D$ be the dual complex of its simple
normal-crossing support. There is a surjection
\[
                  \pi_1(S)\longrightarrow\pi_1(\mathcal D).
                                                        \tag{C.FG.6}\label{cdp:eq:C.FG.6}
\]
Its target has an abelian normal subgroup of index dividing $d$,
generated by at most four elements and of rank at most two.
Moreover the cup-product map
$\bigwedge^*H^1(\mathcal D,\Q)\to H^*(\mathcal D,\Q)$ is injective.
When the resolved support has
no triple points, its connected incidence graph is a tree or has
exactly one independent cycle. In particular two components of an
actual simple normal-crossing fibre with no triple points can meet
in at most two connected curves.

For an actual simple normal-crossing support with no triple points,
write $e_1$ for the number of its connected double curves. The zero
Euler criterion of \ref{cdp:st:C.2.UN.39} takes the following joint form:
\[
 \begin{gathered}
 \#\{
      \text{independent incidence cycles}\}\le1,\\
 e(S)=0\ \Longleftrightarrow\
       \sum_{D\ {\rm double}}g(D)=e_1+\frac{j(T)}2 .
 \end{gathered}
                                                        \tag{C.FG.7}\label{cdp:eq:C.FG.7}
\]
Here $j(T)$ is unchanged by replacing the action on $H^1(F)$ by
the corresponding action on $H_1(F)$ or its inverse. If all $m_i$
are equal, $j(T)=0$. If $1$ is not an eigenvalue of $T$, a
no-triple incidence graph must be a tree.
These are constraints and exact group formulas
for actual degenerations; they include arbitrary multiplicities
and do not assert the existence of a smoothing of abstract gluing data.
\end{proposition}



\begin{proof}
\noindent\textbf{Step 1.} \par\medskip\noindent\textbf{The actual small-tube retraction.}
Shrink the disc so that zero is its only critical value, and take
a compact closed tube $U_\epsilon=f^{-1}(|z|\le\epsilon)$ inside
the original open total space. Put $\varphi=|f|^2$ and choose a
smooth Riemannian metric near this tube. The local analytic gradient
inequality, in the form of
\cite[Theorem (LI), p.~775]{Kurdyka1998}, gives near each point of
$S$ an exponent smaller than one. Compactness, a finite coordinate
cover, and uniform equivalence of the coordinate Euclidean metrics
with the chosen metric give constants $c>0$ and $0<\theta<1$ with
\[
                 \|\nabla\varphi\|\ge c\varphi^\theta
                   \quad\hbox{on }U_\epsilon\setminus S.
                                                        \tag{C.FG.8}\label{cdp:eq:C.FG.8}
\]
Indeed take the maximum of the finitely many exponents, then shrink
to $0\le\varphi<1$; properness permits a whole smaller tube to be
contained in the resulting neighbourhood of $S$. There is no
critical point of $\varphi$ off $S$, since $df$ is surjective there.

\noindent\textbf{Step 2.} On that complement use the vector field
$V=-\nabla\varphi/\|\nabla\varphi\|^2$. Along its flow the value
of $\varphi$ decreases at speed one. A trajectory starting at level
$a$ therefore exists up to every time smaller than $a$: each
positive closed level band is compact and the vector field is
smooth there. Its remaining length below level $b$ is at most
\[
          \int_0^b\frac{du}{c u^\theta}
                 =\frac{b^{1-\theta}}{c(1-\theta)}.    \tag{C.FG.9}\label{cdp:eq:C.FG.9}
\]
It consequently has a limit in $S$. Let $H(x,q)$ be the flow for
time $q\varphi(x)$ when $q<1$, and its limiting point when $q=1$;
put $H(x,q)=x$ on $S$. The estimate \eqref{cdp:eq:C.FG.9} proves continuity
at points of $S$ and uniform convergence as $q\to1$. For $q<1$
away from $S$, ordinary smooth dependence of flows proves
continuity. Thus $H$ is an actual strong deformation retraction
of $U_\epsilon$ onto $S$. This also proves the needed retraction
for every later resolved or covered tube. No invariant-cycle
surjectivity is involved.

\par\medskip\noindent\textbf{The exact normal subgroup of meridians.}
Write $U^*=U_\epsilon\setminus S$. Proper horizontal transport
over the punctured disc gives a smooth torus bundle. Its homotopy
sequence, and $\pi_2(\Delta^*)=0$, identify
\[
       \pi_1(U^*)=\Lambda\rtimes_T\langle t\rangle,
       \qquad \pi_1(U^*)\longrightarrow\pi_1(\Delta^*)=\Z
       \text{ sends }x^v t^a\text{ to }a .             \tag{C.FG.10}\label{cdp:eq:C.FG.10}
\]
The splitting exists because an extension onto the infinite cyclic
group splits after choosing a lift of its generator.

\noindent\textbf{Step 3.} General position makes every loop in the smooth six-manifold
$U_\epsilon$ avoid the divisor $S$. A null-homotopy of a loop from
$U^*$ can be made transverse to the smooth part of $S$ and disjoint
from its singular set: the latter has real dimension at most two,
so a transverse real two-disc misses it. Removing little discs
around the finitely many transverse intersections expresses the
original loop as a product of conjugates of positive or negative
meridians. Conversely each such meridian bounds its normal disc
in $U_\epsilon$.

\noindent\textbf{Step 4.} One meridian for each global irreducible component suffices. The
smooth locus of an irreducible component, with its intersections
with other components removed, is path connected. This follows
also by taking its connected normalization resolution and moving
paths off the proper complex analytic curves and points. A path
there moves positive normal circles through positive normal circles,
giving conjugate based meridians in $U^*$. Near a chosen smooth
point of $S_i$ the actual equation is $f=u z^{m_i}$ with $u$ a
unit. The meridian has winding $m_i$, hence is exactly
$x^{v_i}t^{m_i}$ in \eqref{cdp:eq:C.FG.10}, for a uniquely determined $v_i$
after the indicated choices. The retraction and this kernel
calculation prove \eqref{cdp:eq:C.FG.2}.

\noindent\textbf{Step 5.} The image $A$ of $\Lambda$ is normal and abelian. Quotienting
\eqref{cdp:eq:C.FG.2} by this image leaves precisely the presentation
$\langle t\mid t^{m_i}=1\text{ for all }i\rangle$.
This is $\Z/d\Z$, which proves the exact index in \eqref{cdp:eq:C.FG.3}.
For later comparison, abelianizing the same actual presentation
also gives
\[
 H_1(S,\Z)\simeq
 \frac{\bigl(\Lambda/(T-I)\Lambda\bigr)\oplus\Z}
      {\langle ([v_i],m_i):1\le i\le r\rangle_{\Z}}.
                                                        \tag{C.FG.11}\label{cdp:eq:C.FG.11}
\]

\par\medskip\noindent\textbf{The connected normalized root cover.}
Work momentarily in the open complex neighbourhood of the closed
tube. Since all divisor multiplicities of $f$ are divisible by
$d$, unique factorization in the regular local analytic ring
gives $f=u h^d$ near every point, including singular points of the
support. The normalization of $s^d=u h^d$ is the explicit map
\[
       \{v^d=u\}\longrightarrow\{s^d=u h^d\},
                         \qquad s=vh .                \tag{C.FG.12}\label{cdp:eq:C.FG.12}
\]
Its source is finite and \'etale over the smooth original chart:
$v$ is a unit and $d v^{d-1}$ never vanishes. It is normal and is
isomorphic to the base change over $h\ne0$, hence is exactly that
base change's normalization. These normalization maps glue.

\noindent\textbf{Step 6.} Over $U^*$ this is the cyclic covering given by the winding number
modulo $d$ in \eqref{cdp:eq:C.FG.10}. That homomorphism is onto, so the cover
is connected. Its complement in the normalization is a divisor;
the full normalized cover is therefore connected as well.
The tube retraction $H$ lifts uniquely to this actual covering.
Its lift starts at the identity, fixes the inverse image of $S$
pointwise, and at time one has image in that inverse image. It is
therefore a strong deformation retraction onto $S_d$, proving
connectedness of $S_d$ and the asserted covering injection.
The covering subgroup in $\Gamma$ is the kernel of the winding
modulo $d$ map, which is exactly $A$ by \eqref{cdp:eq:C.FG.3}. In particular
this construction does not presume that the cover of an individual
$S_i$ is connected or that it has degree one.

\noindent\textbf{Step 7.} There is also a useful integral calculation using the punctured
part of this same cover. Its group is
$\Lambda\rtimes_{T^d}\langle w\rangle$, where $w=t^d$.
The kernel of its map to $\pi_1(S_d)$ is normally generated by
all the lifted meridians
\[
                x^{T^jv_i}w^{n_i},
                  \qquad 1\le i\le r,\quad0\le j<d .  \tag{C.FG.13}\label{cdp:eq:C.FG.13}
\]
To see completeness of this list, use the coset representatives
$1,t,\ldots,t^{d-1}$ for the winding cover in \eqref{cdp:eq:C.FG.10}.
Conjugating the original normal generators by these representatives,
and then by elements of the covering subgroup, gives exactly its
kernel. Thus no choice of the connected components over $S_i$
has omitted a meridian.

\noindent\textbf{Step 8.} The relations with $j=0$ imply $w^{n_i}=x^{-v_i}$.
Multiplying their $c_i$th powers gives $w=x^{-v_0}$. Consequently
the covering group is abelian, $(T^d-I)\Lambda$ is killed, and
the remaining relations are precisely
$T^jv_i-n_iv_0=0$. Conversely the map $x^v\mapsto[v]$,
$w\mapsto-[v_0]$ from this presentation to $\Lambda/R_d$
respects every semidirect and meridian relation. It is inverse to
the map induced by $\Lambda$. This proves \eqref{cdp:eq:C.FG.4}. For $d=1$
the same argument proves \eqref{cdp:eq:C.FG.5}; in particular it first proves
$t\in A$ and only then imposes $T=I$ on that quotient.
The Smith-normal-form comparisons in
\path{code/A_meridian_groups.py} check the two integral
presentations and their torsion in explicit test matrices; the
presentation is proved here for every actual monodromy.

\par\medskip\noindent\textbf{Resolution and the actual dual-complex map.}
A smooth vertical point or curve blowup induces an isomorphism on
the fundamental group of the total tube. For a curve centre, cut
out a small normal disc bundle. Its boundary is an $S^3$-bundle
over the curve, and the exceptional neighbourhood retracts to the
$\PP^1$-bundle over that same curve. The boundary inclusion
induces the same isomorphism on fundamental groups as the old
disc bundle, since both fibre fundamental groups are zero. For
a point centre the boundary is $S^5$ and the exceptional fibre
is $\PP^2$, again both simply connected. Van Kampen with the
unchanged complement proves the assertion for the actual blowdown
map. Iteration, followed by the two tube retractions, gives
\[
              \pi_1(S^{\rm res})\xrightarrow{\ \simeq\ }\pi_1(S).
                                                        \tag{C.FG.14}\label{cdp:eq:C.FG.14}
\]
The old multiplicities are retained on strict transforms. Every
new multiplicity is the order of the old fibre equation along its
centre and is divisible by $d$, so their gcd remains $d$.

\noindent\textbf{Step 9.} For an SNC support $E=\bigcup E_i$, form the space $P$ from the
disjoint components $E_i$, one $D\times[0,1]$ for each connected
double curve, and one two-simplex for each triple point. The faces
are glued by the actual intersection inclusions. There are actual
continuous maps
\[
                   E\ \xleftarrow{\ p\ }\ P
                       \xrightarrow{\ q\ }\mathcal D,             
                                                        \tag{C.FG.15}\label{cdp:eq:C.FG.15}
\]
where $p$ forgets the simplex coordinate and $q$ collapses every
connected stratum to its associated cell. Compatible triangulation
of the finite analytic family of components and intersections makes
all the inclusions closed subcomplex inclusions; see
\cite[\S3, Theorem~2]{Lojasiewicz1964}. The map $p$ is a homotopy
equivalence. Explicitly, for two subcomplexes the double mapping
cylinder projects to their union by the homotopy-pushout comparison
for a cofibration. Inductively add one component, applying that
comparison also to its intersection with the previous union.
The intersection's mapping cylinders and the triple-point
triangles are precisely the pieces already included in $P$.
This proves the comparison for the finite union and supplies a
homotopy inverse of $p$.

\noindent\textbf{Step 10.} Van Kampen on these mapping cylinders gives a particularly precise
description of $q_*$. Start with the vertex fundamental groups and
edge paths, impose the relations induced by the two inclusions of
each $\pi_1(D)$, and then the triangle boundary relations.
Killing the vertex groups leaves exactly the presentation of
$\pi_1(\mathcal D)$. Thus
\[
 \pi_1(\mathcal D)\simeq
 \pi_1(E)\big/
       \langle\!\langle\operatorname{im}\pi_1(E_i):i
                              \rangle\!\rangle .       \tag{C.FG.16}\label{cdp:eq:C.FG.16}
\]
The map $q\circ p^{-1}$, where $p^{-1}$ is this homotopy inverse,
and \eqref{cdp:eq:C.FG.14} give \eqref{cdp:eq:C.FG.6}. In particular it is an actual
topological quotient map on fundamental groups, including when
two components have several connected double curves.

\noindent\textbf{Step 11.} The image of $A$ is an abelian normal subgroup whose quotient is
a quotient of $\Z/d\Z$. If there are no triple points,
$\mathcal D$ is a finite connected graph. Let its free-group rank
be $b$ and let the index of that abelian subgroup be $a$.
The corresponding connected graph covering has $a$ times as many
vertices and edges, hence free fundamental group of rank
$1+a(b-1)$. An abelian free group has rank at most one. Thus
$b\le1$, proving the tree-or-one-cycle assertion and its two-vertex
consequence.

\par\medskip\noindent\textbf{The rank-two bound, with the actual finite-cover comparison.}
For any of these tubes, inclusion of a nearby fibre followed by
the constructed retraction is a continuous map
$\sigma:F\to S$. Restriction to its fundamental group has image
$A$ of finite index. Therefore
\[
 \sigma^*:H^1(S,\Q)\longrightarrow H^1(F,\Q)
           \quad\hbox{is injective}.                  \tag{C.FG.17}\label{cdp:eq:C.FG.17}
\]
Indeed a homomorphism $\Gamma\to\Q$ vanishing on $A$ factors through
the finite group $\Gamma/A$ and is zero. This proves the stated
direction of the cohomology map. The map $\sigma^*$ in all degrees
is a ring homomorphism because it comes from that actual continuous
map. Its degree-one image is $T$-invariant, as is also immediate
from the semidirect presentation.

\noindent\textbf{Step 12.} Apply this to the resolved tube $W$ with central support
$E=S^{\rm res}$. The map $q\circ p^{-1}:E\to\mathcal D$ of
\eqref{cdp:eq:C.FG.15} induces an injection on $H^1$ because its fundamental
group map is onto. The composite
\[
 H^*(\mathcal D,\Q)\longrightarrow H^*(E,\Q)
               \xrightarrow{\ \sigma^*\ }H^*(F,\Q)
                  =\bigwedge^*H^1(F,\Q)               \tag{C.FG.18}\label{cdp:eq:C.FG.18}
\]
is again a ring homomorphism and is injective in degree one.
The exterior algebra on any linear subspace of $H^1(F,\Q)$
injects into this exterior algebra. Consequently cup product
on $\mathcal D$ is injective on
$\bigwedge^*H^1(\mathcal D,\Q)$. Since $\mathcal D$ is a
two-dimensional cell complex, its third cohomology vanishes,
and three independent degree-one classes would give a
contradiction. Thus $b_1(\mathcal D)\le2$.

\noindent\textbf{Step 13.} This also bounds the rank of its finite-index abelian subgroup,
rather than only the rank of its abelianization. Put
$\Pi=\pi_1(\mathcal D)$ and let $B$ be the image of $A$ in $\Pi$.
The index $a=[\Pi:B]$ divides $d$. Take the actual connected
degree-$a$ covering $\mathcal D'\to\mathcal D$ belonging to $B$,
and pull it back along the continuous map
$W\to E\to\mathcal D$ just constructed. The pullback is connected
because $\pi_1(W)\to\Pi$ is onto. Its covering subgroup is the
inverse image of $B$ in $\Gamma$. The quotient $\Pi/B$ is cyclic
of order $a$, generated by the image of $t$, so this inverse
image is precisely the kernel of the winding number modulo $a$.
The uniqueness of a based covering with a given subgroup
identifies this pullback, over $W$, with the actual normalized
root cover
\[
                  s^a=f\circ\pi:W_a\longrightarrow W. \tag{C.FG.19}\label{cdp:eq:C.FG.19}
\]
It is \'{e}tale by \eqref{cdp:eq:C.FG.12}, since $a$ divides every resolved
component multiplicity. Its lifted tube retraction has connected
central support $E_a$. The lifted map
$E_a\to\mathcal D'$ is onto on fundamental groups: its image is
exactly $B$, by the inverse-image description of the covering
subgroup. Regarded over the $s$-disc, $W_a$ has connected nearby
fibre equal to the actual torus $F_{s^a}$. Hence \eqref{cdp:eq:C.FG.17}--\eqref{cdp:eq:C.FG.18}
apply to this map and give $b_1(\mathcal D')\le2$. Since
$\pi_1(\mathcal D')=B$ is finitely generated abelian, this is
exactly $\operatorname{rank}B\le2$, as claimed.

\noindent\textbf{Step 14.} If a no-triple incidence graph had one cycle, its nonzero
$H^1$ would inject by \eqref{cdp:eq:C.FG.18} into the $T$-invariant part of
$H^1(F,\Q)$. Thus $1$ must be an eigenvalue of $T$ in that case.

\noindent\textbf{Step 15.} Finally, for the original support already SNC with no triple
points, \ref{cdp:st:C.2.UN.39} has no resolution-centre terms and gives
$e(S)=j(T)-2\sum_D(g(D)-1)$, proving \eqref{cdp:eq:C.FG.7}.
If all multiplicities are $m$, the full trace formula is
$\det(I-T^k)=m\sum_i e(S_i^\circ)$ when $m\mid k$ and zero
otherwise. Taking $k$ also divisible by the semisimple order makes
$T^k$ unipotent, so this determinant is zero. Hence the common
sum is zero and all traces vanish; in particular $j(T)=0$.
\end{proof}






\begin{proposition}[Positive geometric genus controls the whole support group]\label{cdp:st:C.2.UN.47}
\StatementRecord{cdp:st:C.2.UN.47}{proved}{Historical events: Claude Fable 5.2 (Anthropic), effort high, 16 November 2026: carried as background without fault found. GPT-6 Astra (OpenAI), ChatGPT Pro mode, effort Pro, S6\_09\_PAPERS, 17 November 2026: targeted proof check. Also checked in this writing. Pre-continuation writing events: GPT-6 Astra (OpenAI), ChatGPT Pro mode, effort Pro, S6\_10\_PAPERS, 18 November 2026: targeted proof check. Continuation mathematical and source checks: GPT-6 Astra (OpenAI), ChatGPT Pro mode, effort Pro, S6\_10\_PAPERS continuation, 19 November 2026: independent re-proof. Continuation metadata events: GPT-6 Astra (OpenAI), ChatGPT Pro mode, effort Pro, S6\_10\_PAPERS continuation, 19 November 2026: used and rechecked in this run; status-vocabulary reconciliation; metadata reconciliation only; no new mathematical check. \emph{Pre-continuation recorded proof scope.} Every finite quotient of the tube group yields an actual finite etale cover and a connected Stein-disc torus degeneration. Pullback normalization resolutions have exactly the subgroup-index number of components. The compact-core genus bound makes a positive-genus component surjective in every finite quotient, and the explicit nA quotient separates every proper subgroup of the virtually abelian support group. Etale Riemann-Roch gives the exact Euler-characteristic degree inequality, forcing trivial support group for chi=2 and only the stated order-two exception for nontrivial chi=1 image}{the hypotheses, notation and definitions stated here; \cref{cdp:st:C.2.UN.33}; \cref{cdp:st:C.2.UN.44}; \cref{cdp:eq:C.PF.5}; \cref{cdp:st:C.2.UN.28}; \cref{cdp:eq:C.PF.4}; \cref{cdp:eq:C.PF.7}; \cref{cdp:eq:C.PF.8}; \cref{cdp:eq:C.PF.6}; \cref{cdp:eq:C.PF.12}; \cref{cdp:eq:C.PF.10}}

Let $f:U\to\Delta$ be a proper holomorphic two-torus degeneration
with smooth connected total threefold, let $S=f^{-1}(0)_{\rm red}$,
and put $\Gamma=\pi_1(S)$. For an irreducible component $S_i$ let
$h_i:Z_i\to S_i\hookrightarrow U$ be a resolution of its normalization.
The following are simultaneous constraints on this actual fibre.

Every $Z_i$ with $p_g(Z_i)>0$ has $p_g(Z_i)=1$, and its actual
map induces a surjection
\[
                         \pi_1(Z_i)\longrightarrow\Gamma .
                                                        \tag{C.PF.1}\label{cdp:eq:C.PF.1}
\]
Consequently every SNC dual complex of an embedded resolution of
the fibre is simply connected whenever such a component occurs.
A component with K3 minimal model forces $\Gamma=1$.
More generally,
\[
                    \chi(\OO_{Z_i})=2
                         \quad\Longrightarrow\quad\Gamma=1 .
                                                        \tag{C.PF.2}\label{cdp:eq:C.PF.2}
\]
This includes the Kodaira-dimension-one components of holomorphic
Euler characteristic two in \ref{cdp:st:C.2.UN.33}.

For a component with $\chi(\OO_{Z_i})=1$, its fundamental-group
image in $\Gamma$ is zero, with just the following possible
group configuration:
\[
 \Gamma\simeq\Z/2\Z,\qquad
 \operatorname{im}\bigl(\pi_1(Z_i)\to\Gamma\bigr)=\Gamma .
                                                        \tag{C.PF.3}\label{cdp:eq:C.PF.3}
\]
In that configuration the actual connected double cover of $Z_i$
induced from the tube has
$\chi(\OO)=2$, $p_g=1$ and $q=0$. At most one component with
$\chi(\OO)=1$ can have nontrivial fundamental-group image.
The assertion \eqref{cdp:eq:C.PF.3} is a restriction on groups of actual fibres,
and is not an assertion that every surface or gluing with these
invariants occurs.
\end{proposition}



\begin{proof}
\noindent\textbf{Step 1.} \par\medskip\noindent\textbf{Every finite quotient gives another actual torus degeneration.}
Use the strong deformation retraction of \ref{cdp:st:C.2.UN.44} on a sufficiently
small tube; its annular complement also retracts radially by proper
horizontal transport. Thus an open tube and its closed smaller
tube both have group $\Gamma$. Given a finite quotient
$\rho:\Gamma\twoheadrightarrow Q$ of order $N$, its connected
regular topological covering carries the unique complex structure
obtained by lifting the holomorphic coordinate charts. This gives
an actual finite \'{e}tale map
\[
                     p:U_Q\longrightarrow U,\qquad \deg p=N.
                                                        \tag{C.PF.4}\label{cdp:eq:C.PF.4}
\]
It is proper: a finite covering of a locally compact Hausdorff
space is proper, as follows by covering a compact set with
finitely many relatively compact evenly covered neighbourhoods.
Consequently $f\circ p$ is proper. Its total space is smooth and
connected.

\noindent\textbf{Step 2.} The punctured covering $U_Q^*$ is connected, since
$\pi_1(U^*)\to\Gamma\to Q$ is onto by \ref{cdp:st:C.2.UN.44}. Let $A$ be the image
of the nearby fibre group in $\Gamma$ and let $B=\rho(A)$.
The connected components over a fixed nearby fibre are indexed
by $Q/B$. The base generator permutes them transitively, and
$Q/B$ is cyclic, because $\Gamma/A$ is cyclic. If
$a=[Q:B]$, the connected base covering in the Stein factor of
$f\circ p$ is therefore the punctured degree-$a$ cover of the disc.
Its normal completion has one point above zero and is the disc
with actual coordinate map $t=s^a$. One can obtain this
identification by the classification of coverings of the
punctured disc and then extend the two inverse maps across zero:
each is bounded over a smaller disc, and the completed normal
one-dimensional spaces have the removable-singularity property.
Equivalently it is the normalization of this finite disc cover.
The proper Stein factor is this completion, because it is finite,
normal and agrees over the dense punctured disc. We have thus
constructed a proper map with connected fibres
\[
                        f_Q:U_Q\longrightarrow\Delta_s .
                                                        \tag{C.PF.5}\label{cdp:eq:C.PF.5}
\]
Its reduced central support is exactly $(p^{-1}S)_{\rm red}$,
because the completed Stein disc has just one point over zero.
Every smooth general fibre of $f_Q$ is a connected finite \'{e}tale
cover of $F$. Writing $F=\C^2/\Lambda$, its actual lifted universal
cover identifies it with $\C^2/\Lambda'$ for a finite-index
sublattice $\Lambda'\subset\Lambda$. Hence \eqref{cdp:eq:C.PF.5} is another
two-torus degeneration, to which \ref{cdp:st:C.2.UN.28} applies. This comparison
does not require the original cover to be a pure base cover.

\par\medskip\noindent\textbf{The component pullbacks and their precise degree.}
Fix $Z=Z_i$ and take its fibre product with \eqref{cdp:eq:C.PF.4}. It is a
smooth compact, possibly disconnected surface
\[
               Z_Q=Z\times_U U_Q\longrightarrow Z
                       \quad\hbox{finite \'{e}tale of degree }N .
                                                        \tag{C.PF.6}\label{cdp:eq:C.PF.6}
\]
Let $H_Q=\rho(h_{i*}\pi_1(Z))$, let $e=|H_Q|$, and put
$k=[Q:H_Q]=N/e$. Path lifting gives exactly $k$ connected
components $Z_1',\ldots,Z_k'$, each a connected cover of $Z$
of degree $e$.

\noindent\textbf{Step 3.} These are resolutions of distinct actual normalization components
of the central support of $U_Q$. To verify this comparison, remove
the singular set of the original support and the exceptional
curves of $Z\to S_i$. Over the remaining dense open set the map
$Z\to U$ is its embedding as the smooth single-component locus.
Its base change is therefore an isomorphism onto the corresponding
open subset of the central support of $U_Q$. Removing those proper
analytic curves and points from a connected smooth $Z_j'$ leaves
it connected. Its proper image is consequently an irreducible
central component, and the displayed isomorphism over the dense
open set proves that $Z_j'$ maps birationally onto its normalization.
Different $Z_j'$ cannot give the same component, since their dense
open isomorphisms have disjoint images. Thus the compact-core
bound of \ref{cdp:st:C.2.UN.28} gives the exact useful inequality
\[
                         \sum_{j=1}^k p_g(Z_j')\le1 .
                                                        \tag{C.PF.7}\label{cdp:eq:C.PF.7}
\]
Other central components can only add nonnegative terms to that
bound.

\noindent\textbf{Step 4.} Pullback of a holomorphic two-form under a finite \'{e}tale
cover is injective, by restriction to any open set on which
that cover is biholomorphic. If $p_g(Z)>0$, \ref{cdp:st:C.2.UN.28} on the
original degeneration already gives $p_g(Z)=1$. Every
$Z_j'$ then contributes at least one in \eqref{cdp:eq:C.PF.7}, so $k=1$.
We have proved that
\[
       \rho(h_{i*}\pi_1(Z))=Q
                 \quad\hbox{for every finite quotient }\rho
                         \hbox{ of }\Gamma .           \tag{C.PF.8}\label{cdp:eq:C.PF.8}
\]

\par\medskip\noindent\textbf{The group-theoretic passage retains the actual subgroup.}
\ref{cdp:st:C.2.UN.44} supplies a normal finitely generated abelian subgroup
$A\subset\Gamma$ of finite index. Every subgroup $H$ of $\Gamma$
is separable by finite quotients in the sense needed here.
If $HA\ne\Gamma$, its image already fails to be all of
the finite quotient $\Gamma/A$. Otherwise, for a proper
$H$, choose $b\in A\setminus H$. The finitely generated
abelian group $A/(H\cap A)$ has a finite quotient detecting
the class of $b$; equivalently some positive integer $n$
satisfies
\[
                         b\notin(H\cap A)+nA .
                                                        \tag{C.PF.9}\label{cdp:eq:C.PF.9}
\]
The subgroup $nA$ is invariant under every automorphism of $A$,
hence is normal in $\Gamma$, and $\Gamma/nA$ is finite.
Moreover $HN\cap A=(H\cap A)+nA$ for $N=nA$.
Thus the image of $H$ in this finite quotient misses $b$.
This proves directly that \eqref{cdp:eq:C.PF.8} implies
$h_{i*}\pi_1(Z)=\Gamma$, establishing \eqref{cdp:eq:C.PF.1}.

\noindent\textbf{Step 5.} The same elementary argument gives residual finiteness of
$\Gamma$. An element outside $A$ is detected in $\Gamma/A$.
A nonzero element of $A$ is detected in $A/nA$ for some $n$,
and therefore in the finite quotient $\Gamma/nA$.
No residual-finiteness assertion for a general degeneration
group is being imported.

\par\medskip\noindent\textbf{The dual complex and the K3 consequence.}
Choose any embedded resolution of the central divisor.
The strict transform of $S_i$ is a smooth resolution of its
normalization, with the same geometric genus as $Z$.
For completeness, the graph of the birational correspondence
between these two smooth resolutions has a common resolution
by point blowups. Pullback and Hartogs extension of holomorphic
two-form coefficients across a point show that each point blowup
preserves $p_g$. The same blowup maps induce isomorphisms on
$\pi_1$, by the $S^3$ and $\PP^1$ neighbourhood comparison.
The two maps from the common resolution to $U$ agree on a
dense open set and hence everywhere. Thus the image of the
strict-transform vertex group in $\Gamma$ is exactly the
image in \eqref{cdp:eq:C.PF.1}.

\noindent\textbf{Step 6.} By (C.FG.14), the resolved support has group $\Gamma$.
By (C.FG.16), passage to its dual complex kills the image of
every component group. In particular it kills this surjective
vertex group, so the dual-complex fundamental group is trivial.
If $Z$ has K3 minimal model, its fundamental group is zero:
the K3 is simply connected and the point blowups do not change
its fundamental group. Equation \eqref{cdp:eq:C.PF.1} then gives $\Gamma=1$.

\par\medskip\noindent\textbf{Holomorphic Euler characteristic gives stronger restrictions.}
The degree in \eqref{cdp:eq:C.PF.6} also gives
\[
       \sum_{j=1}^k\chi(\OO_{Z_j'})=N\chi(\OO_Z).
                                                        \tag{C.PF.10}\label{cdp:eq:C.PF.10}
\]
This follows, for example, by applying surface Riemann--Roch to
the actual \'{e}tale tangent-bundle pullback:
$12\chi(\OO)=\int(c_1^2+c_2)$, and integration of a pulled-back
top-degree class multiplies by the covering degree. Combining
this equality with
$p_g(Z_j')=\chi(\OO_{Z_j'})-1+q(Z_j')$ gives
\[
       1\ \ge\ \sum_j p_g(Z_j')
           \ \ge\ N\chi(\OO_Z)-k,\qquad k\le N .
                                                        \tag{C.PF.11}\label{cdp:eq:C.PF.11}
\]
For $\chi(\OO_Z)=2$, the right side is at least $N$.
Every finite quotient of $\Gamma$ therefore has order one.
Residual finiteness proves \eqref{cdp:eq:C.PF.2}.

\noindent\textbf{Step 7.} For $\chi(\OO_Z)=1$, a finite quotient in which $H_Q$ is
nontrivial has $e\ge2$ and satisfies
\[
                    1\ge N-\frac Ne\ge\frac N2 .
                                                        \tag{C.PF.12}\label{cdp:eq:C.PF.12}
\]
It follows that $N=e=2$, so the only such finite quotient
is a quotient of order two. If the image
$H=h_{i*}\pi_1(Z)$ in $\Gamma$ is nontrivial, residual
finiteness provides one such quotient
$\rho:\Gamma\to\Z/2\Z$ detecting a nonidentity element of $H$.
If $\ker\rho$ contained a nonidentity element, another finite
quotient would detect that element. The image of the product
of these two quotient maps would then be a finite quotient
of order greater than two still detecting the original
element of $H$, contradicting \eqref{cdp:eq:C.PF.12}. Thus $\ker\rho=1$,
which proves precisely \eqref{cdp:eq:C.PF.3}.

\noindent\textbf{Step 8.} In this case \eqref{cdp:eq:C.PF.6} is a connected double cover of $Z$,
and \eqref{cdp:eq:C.PF.10} gives $\chi=2$. Equation \eqref{cdp:eq:C.PF.7} and
$p_g=\chi-1+q$ give $p_g=1$ and $q=0$. If two different
components with $\chi=1$ both had nontrivial image, the
same actual double cover of the tube would contain their
two distinct normalization resolutions, each of geometric
genus one. This contradicts \ref{cdp:st:C.2.UN.28}. It proves the final
joint assertion.
\end{proof}






\section{Irreducible and reducible torus-fibre Euler criteria}\label{sec:II:07}


\begin{proposition}[Every normal irreducible support is a smooth torus quotient]\label{cdp:st:C.2.UN.5}
\StatementRecord{cdp:st:C.2.UN.5}{proved}{Historical events: Claude Fable 5.2 (Anthropic), effort high, 16 November 2026: carried as background without fault found. GPT-6 Astra (OpenAI), ChatGPT Pro mode, effort Pro, S6\_09\_PAPERS, 17 November 2026: targeted proof check. Also checked in this writing. Pre-continuation writing events: GPT-6 Astra (OpenAI), ChatGPT Pro mode, effort Pro, S6\_10\_PAPERS, 18 November 2026: targeted proof check; used and rechecked in this run. Continuation mathematical and source checks: GPT-6 Astra (OpenAI), ChatGPT Pro mode, effort Pro, S6\_10\_PAPERS continuation, 19 November 2026: computation or formal certificate; independent re-proof; targeted proof check; used and rechecked in this run. Continuation metadata events: GPT-6 Astra (OpenAI), ChatGPT Pro mode, effort Pro, S6\_10\_PAPERS continuation, 19 November 2026: used and rechecked in this run; status-vocabulary reconciliation; metadata reconciliation only; no new mathematical check. \emph{Continuation recorded proof scope.} Complete written derivation before reading the target proof, followed by comparison of all seven printed steps. The normalized cover, Albanese map, affine cyclic action and canonical-character factorization were checked. The fresh program \path{code/continuation_root/UN5_characters_fresh.py} checks all 47 primitive characters and the integral elliptic unit orders; the cited Buchdahl and Catanese statements were opened.}{the hypotheses, notation and definitions stated here; \cref{cdp:st:A.5.17}; \cite[Theorem~11]{Buchdahl1999}; \cite[Theorem~4.6]{Catanese2002}; \cite[Definition~3.1 and Theorems~3.2--3.3]{Catanese2002}; \cref{cdp:eq:C.UNN.2}}

Let $f:Y\to B$ be a proper holomorphic map from a smooth
complex threefold to a smooth curve, with connected fibres
and smooth general fibre a complex two-torus. Let $mS$ be
a full fibre whose reduced support is irreducible and normal.
Then $S$ is smooth and $e(S)=0$. More precisely, the
normalized base change $t=s^m$ gives a connected two-torus
$A$ and a finite unramified cyclic map
\[
                         A\longrightarrow S
\]
of degree $m$. Thus $S$ is a two-torus or a bielliptic
surface. In the bielliptic case the order $n$ of $K_S$ is
$2,3,4$, or $6$, and $n\mid m$; in particular a reduced
normal fibre is a torus.

Every two-torus occurs with every positive multiplicity.
Every bielliptic surface of canonical order $n$ occurs
with every multiplicity divisible by $n$. These
realizations can be made in smooth compact threefolds
over $\PP^1$.
\end{proposition}


\begin{proof}
\noindent\textbf{Step 1.} A normal surface hypersurface has only isolated
singularities. At a singular point of $S$, write the
actual map as $t=u h^m$, with $u$ a holomorphic unit
and $h$ a reduced local equation of $S$. A local
$m$th root of $u$ gives $t=(u^{1/m}h)^m$. The
isolated-hypersurface assertion of \ref{cdp:st:A.5.17} supplies an
embedded vanishing two-sphere of self-intersection
$-2$ in a small nonzero Milnor fibre of $u^{1/m}h$.
This fibre is an open part of the actual nearby
torus fibre of $f$. A sphere in a torus is nullhomotopic,
and hence has zero homology class and zero
self-intersection. The contradiction proves that $S$
is smooth. This argument uses no Betti-number or
K\"ahler hypothesis on the ambient threefold.

\noindent\textbf{Step 2.} Shrink the base to a disc with coordinate $t$. Near
any point of $S$ write $t=u z^m$ with $z$ a smooth
equation of $S$. The normalization $U'$ of the base
change $t=s^m$ is given locally by the explicit maps
\[
        w^m=u,\qquad s=wz,
        \qquad p:U'\longrightarrow U.
 \tag{C.UNN.1}\label{cdp:eq:C.UNN.1}
\]
Since $u$ never vanishes, $p$ is finite and unramified
of degree $m$; these local normalizations glue to the
normalization of the actual fibre product. The map
$h:U'\to\Delta_s$ has $ds=w\,dz$ along its central
support, so that support is smooth and reduced.
For each $s\ne0$, the map $p$ identifies $h^{-1}(s)$
with $f^{-1}(s^m)$, a connected torus. The Stein
factor of $h$ has generic degree one over the normal
disc and is therefore the disc. In particular $U'$
and its central fibre $A$ are connected. After
shrinking the disc, $h$ is a proper submersion.
The deck action is $w\mapsto\zeta_m w$, so it is
free also on $A$, and its quotient there is exactly
$S$.

\noindent\textbf{Step 3.} We verify that $A$ is a torus through its actual
Albanese map. Ehresmann's theorem identifies its
integral cohomology ring with that of a real
four-torus. Thus $b_1(A)=4$, and
\cite[Theorem~11]{Buchdahl1999} makes $A$ K\"ahler.
The period-integral map
$a:A\to\operatorname{Alb}(A)$ is consequently
defined by its two independent closed holomorphic
one-forms and induces an isomorphism on integral
$H^1$. Both surfaces have the exterior algebra on
this lattice as their integral cohomology ring.
The pullback of a top-degree generator is therefore
a generator, so $a$ has degree one; holomorphicity
fixes the positive sign. Its pullback on $H^2$ is
also an isomorphism. A curve contracted by $a$
would pair to zero with every pullback class and
hence with a K\"ahler class on $A$, contradicting
positive area. Thus $a$ is finite of degree one
and is an isomorphism. This proves the torus
identification, without imposing a metric on $U'$.
It also agrees with the deformation theorem
\cite[Theorem~4.6]{Catanese2002}.

\noindent\textbf{Step 4.} Write a generator of the deck group on
$A=\mathbb C^2/\Lambda$ as $g(z)=Lz+b$.
Every holomorphic torus automorphism has this affine
form: the derivative of a lift is periodic and hence
constant on the compact quotient. The matrix $L$
preserves $\Lambda$ and has finite order. If $1$
were not an eigenvalue of $L$, the equation
$(I-L)z=b$ in $\mathbb C^2$ would give a fixed
point of $g$, contrary to the free deck action.
If $L=I$, the deck group acts by translations and
$S$ is the torus quotient by the resulting finite
subgroup.

\noindent\textbf{Step 5.} Otherwise the eigenvalues on $\mathbb C^2$ are
$1$ and a nontrivial root of unity $\zeta$.
The corresponding integral real representation
has characteristic polynomial
$(x-1)^2(x-\zeta)(x-\bar\zeta)$, so the order
$n$ of $\zeta$ is $2,3,4$, or $6$. The images of
$I-L$ and of $I+L+\cdots+L^{n-1}$ are elliptic
subtori with complementary tangent lines. Their
addition map is an isogeny from a product of
elliptic curves to $A$. The kernel of the deck
group's linear representation consists of
translations. Quotienting by it first gives a
torus $A_1$ with a free cyclic action of order
$n$ and quotient $S$. On the canonical form
this action has character $\zeta$, so $K_S$
has exact order $n$, while the invariant
holomorphic one-form space has dimension one.
Equivalently, lifting to the product isogeny
and removing the pure translation kernels gives
the bielliptic presentation
$(E\times F)/G$, with $G$ acting faithfully
by translations on $E$ and with $F/G\simeq\PP^1$.
These are precisely the bielliptic surfaces
of \cite[Definition~3.1 and Theorems~3.2--3.3]{Catanese2002}.
The linear order divides the deck order $m$.
In both cases the unramified cover gives
$m\,e(S)=e(A)=0$.

\noindent\textbf{Step 6.} We finish with realizations and their multiplicities.
Given any free cyclic torus quotient
$A\to S=A/\langle g\rangle$ of degree $m$, set
\[
 Y=(A\times\PP^1)/\langle(a,z)\mapsto
                         (g(a),\zeta_m z)\rangle,
 \qquad f([a,z])=z^m.
 \tag{C.UNN.2}\label{cdp:eq:C.UNN.2}
\]
The action is free, hence $Y$ is smooth and compact.
For $t\ne0,\infty$ its $m$ slices are permuted
transitively, and the quotient fibre is exactly
$A$. The full fibres at $0$ and $\infty$ are
$mS$. Thus these are actual compact realizations
of the displayed support and multiplicity.
For a prescribed torus $S$, an index-$m$ sublattice
with cyclic quotient gives the required cyclic
isogeny $A\to S$ for every $m\ge1$.

\noindent\textbf{Step 7.} For a prescribed bielliptic surface, use its
product presentation $(E\times F)/G$. The linear
character $G\to\mu_n$ is surjective, and its
kernel acts by translations on both factors.
Consequently its canonical cyclic cover is a
torus. Moreover the character factors through
$\pi_1(E/G)\simeq\mathbb Z^2$, because $G$
acts on $E$ by translations. A surjective
character $\mathbb Z^2\to\mathbb Z/n$ lifts
to a primitive integral homomorphism
$\mathbb Z^2\to\mathbb Z$: choose integer
lifts of its two coefficients, and change one
by a multiple of $n$ to make their gcd one.
The possible prime divisors of the other
coefficient are avoided by the Chinese
remainder theorem; primes dividing $n$ already
cannot divide both coefficients. Pull this
homomorphism back along the connected elliptic
fibration $S\to E/G$ and reduce it modulo
$m=kn$. Its kernel defines a connected cyclic
unramified cover of degree $m$ that factors
through the canonical torus cover. A connected
finite unramified cover of a torus is the
quotient of $\mathbb C^2$ by a finite-index
sublattice, so this new cover is again a
torus. Applying \eqref{cdp:eq:C.UNN.2} proves every claimed
bielliptic multiplicity.
\end{proof}






The statements \ref{cdp:st:C.2.CR.1}--\ref{cdp:st:C.2.CR.6} below retain the
proved local and global reductions. The later
\ref{cdp:st:C.2.UN.7} proves that the contracted irrational
configuration has $e(S)\ge2$ and $\pi_1(S)=1$.
Together with the CDP transgression, or the zero
Euler conclusion \ref{cdp:st:C.2.MH.3}, this excludes the
configuration entirely under the CDP hypotheses.
The unrestricted genus assertion is separately
proved in \ref{cdp:st:C.2.UN.10}.

\begin{proposition}[Topology and coefficients in the contracted-section case]\label{cdp:st:C.2.CR.1}
\StatementRecord{cdp:st:C.2.CR.1}{proved}{Historical events: Claude Fable 5.2 (Anthropic), effort high, 16 November 2026: carried as background without fault found. GPT-6 Astra (OpenAI), ChatGPT Pro mode, effort Pro, S6\_09\_PAPERS, 17 November 2026: targeted proof check. Also checked in this writing. Pre-continuation writing events: GPT-6 Astra (OpenAI), ChatGPT Pro mode, effort Pro, S6\_10\_PAPERS, 18 November 2026: targeted proof check. Continuation metadata events: GPT-6 Astra (OpenAI), ChatGPT Pro mode, effort Pro, S6\_10\_PAPERS continuation, 19 November 2026: used and rechecked in this run; status-vocabulary reconciliation; metadata reconciliation only; no new mathematical check. \emph{Pre-continuation recorded proof scope.} The formal-function restriction to the contracted section gives exact exceptional length g and both degree-one vanishings. Every reverse ruling blowup has coefficient d minus one; its tree, shifted gcd and Euler changes were checked. The minimal-resolution adjunction calculation identifies the complete conductor preimage, and the actual pushout gives its fundamental group. The punctured torus-bundle surjection and global nonzero character force rank one. The conductor-duality and normalization sequences give each coherent and delta inequality}{the hypotheses, notation and definitions stated here; \cref{cdp:st:A.5.3}; \cref{cdp:st:C.2.UN.7}; \cref{cdp:eq:C.CR.4}; \cref{cdp:st:A.5.17}; \cref{cdp:eq:C.CR.6}; \cref{cdp:eq:C.CR.7}; \cref{cdp:st:D.1.9}; \cref{cdp:eq:C.CR.8}}

Let $Y$ be a smooth compact complex threefold with
$b_1(Y)=b_2(Y)=0$, and let $f:Y\to\PP^1$ be a proper
holomorphic surjection with connected fibres and smooth
general fibre a complex two-torus. Let $mS$
be a full fibre with irreducible nonnormal reduced support, and
write $h=\eta\pi:Z\to T\to S$ for its normalization and minimal
normalization resolution. Suppose that $p:Z\to B$ is an
irrational ruling, that the effective divisor in \ref{cdp:st:A.5.3} is
$D=2C+V$, where $C$ is a section and $V$ is vertical, and
that $\pi$ contracts $C$. Put $g=g(B)$, let $\Gamma$ be the
conductor-image scheme, and let $n$ be the number of irreducible
components of $\Gamma_{\rm red}$.

Then the following are exact consequences of these hypotheses:
\[
\begin{gathered}
 H^1(T;\mathbb R)=H^1(T,\OO_T)=H^2(T,\OO_T)=0,
 \qquad \operatorname{length}R^1\pi_*\OO_Z=g,\\
 \pi_1(S)\simeq\pi_1(\Gamma_{\rm red})\simeq\mathbb Z,
 \qquad e(\Gamma_{\rm red})=n,\\
 \chi(\OO_\Gamma)=-1,\qquad
 \chi(\OO_{\Gamma_{\rm red}})\in\{-1,0\},\qquad
 \sum_{x\in\Gamma_{\rm red}}\bigl(\delta_x-r_x+1\bigr)\le1.
\end{gathered}
\tag{C.CR.1}\label{cdp:eq:C.CR.1}
\]
Here $r_x$ is the number of normalization branches at $x$.
The full inverse image $D_{\rm red}$ is a smooth genus-$g$
section with trees of smooth rational curves attached; all
its intersections are ordinary double points. For adjacent
components of coefficients $a,b$ one has
\[
                         \gcd(a-1,b-1)=1.          \tag{C.CR.2}\label{cdp:eq:C.CR.2}
\]
Choose a specified sequence of ruling blowdowns to a bundle
$Z_0\to B$. Write its pushed-forward divisor as
$2C_0+\sum_{j=1}^{v}a_jF_j$, with distinct fibres $F_j$
and $a_j>0$, and let $N_0$ count the reverse blowups whose
centre has multiplicity one in the current divisor. Then
\[
                     e(S)=2-2g-v+N_0+n.           \tag{C.CR.3}\label{cdp:eq:C.CR.3}
\]
These necessary equations are retained in the Euler
calculation. \ref{cdp:st:C.2.UN.7} later excludes their simultaneous
geometric occurrence under the stated CDP hypotheses.
\end{proposition}


\begin{proof}
\noindent\textbf{Step 1.} The finite-duality calculation (A.RT.5), whose proof does
not assume rationality, gives $H^2(T,\OO_T)=0$. The
degree-one Leray sequence is consequently
\[
 0\longrightarrow H^1(T,\OO_T)\longrightarrow H^1(Z,\OO_Z)
 \longrightarrow H^0(T,R^1\pi_*\OO_Z)\longrightarrow0.
\tag{C.CR.4}\label{cdp:eq:C.CR.4}
\]
The middle vector space has dimension $g$. At the point
to which $C$ is contracted, formal functions and the
successive infinitesimal exceptional fibres give a
surjection from the completed stalk of $R^1\pi_*\OO_Z$
onto the $H^1$ of the scheme-theoretic exceptional fibre.
Indeed the restriction maps on $H^1$ of these thickenings
are surjective, since their kernels are supported in
dimension at most one. The completed stalk equals the
stalk, because this higher direct image has finite support.
Restriction from the exceptional fibre to $C$ is likewise
surjective on $H^1$. Thus its length is at least
$h^1(C,\OO_C)=g$, whereas \eqref{cdp:eq:C.CR.4} bounds the total length
by $g$. Equality proves the coherent assertions in the
first line of \eqref{cdp:eq:C.CR.1}.

\noindent\textbf{Step 2.} Connectedness of resolution fibres makes the constant-sheaf
Leray map $H^1(T;\mathbb R)\to H^1(Z;\mathbb R)$ injective.
Restriction from $Z$ to its section $C$ is an isomorphism
on $H^1$, because a ruled bundle has simply connected
fibres and point blowups do not change $H^1$. The image of
$H^1(T;\mathbb R)$ restricts to zero on the contracted
curve $C$. This proves its vanishing.

\noindent\textbf{Step 3.} We describe the actual divisor comparisons. At every ruling
blowdown $\sigma:W\to W'$ with exceptional curve $E$, let
$a$ be the coefficient of $E$ and let $d$ be the multiplicity
of the pushed-forward boundary at the centre. The identity
\[
 K_W+D_W=\sigma^*(K_{W'}+D_{W'})+(1+a-d)E
\tag{C.CR.5}\label{cdp:eq:C.CR.5}
\]
and numerical triviality give $a=d-1$. The minimal model
boundary is therefore $2C_0+\sum a_jF_j$. Its reduced
support is the section with disjoint rational fibre leaves.
Every reverse blowup preserves the asserted simple-normal-
crossings tree structure: a positive exceptional coefficient
adds a leaf at a smooth boundary point or subdivides an edge
at a boundary node; a zero exceptional coefficient leaves
the reduced boundary isomorphic to its predecessor. A centre
off the current boundary would create coefficient $-1$,
which persists on its strict transform and contradicts the
effectiveness of the original divisor. In particular
$D_{\rm red}$ is connected, and its fundamental group maps
isomorphically to that of $Z$ by inclusion, with both
identified with $\pi_1(B)$ through the section.

\noindent\textbf{Step 4.} This also proves \eqref{cdp:eq:C.CR.2} without discarding coefficients.
Initially every edge has coefficients $(2,a_j)$ and hence
coprime shifted coefficients. A node blowup inserts the
coefficient $a+b-1$, so its shifted coefficient is
$(a-1)+(b-1)$; the two new edge gcds equal the old gcd.
A positive-coefficient blowup at a smooth point of a
coefficient-$a$ boundary creates coefficient $a-1$ and
an edge with shifted coefficients $a-1,a-2$, whose gcd
is one. These are all the possible operations.

\noindent\textbf{Step 5.} On a minimal normalization resolution, a coefficient-zero
exceptional curve $E$ has $K_Z\cdot E\ge0$. Since
$(K_Z+D)\cdot E=0$, it follows that $D\cdot E=0$.
Such a curve cannot meet $D_{\rm red}$. Connectedness of
exceptional fibres and of $D_{\rm red}$ now gives the same
support equality as in A.5.16:
\[
 D_{\rm red}=h^{-1}(\Gamma_{\rm red}).
\tag{C.CR.6}\label{cdp:eq:C.CR.6}
\]
Any exceptional fibre outside this set would correspond to
an isolated normal hypersurface singularity of $S$, because
the normal conductor ideal is reflexive and has no isolated
support. \ref{cdp:st:A.5.17} excludes every such point in the actual
torus smoothing. Thus $h$ is an isomorphism off the two
closed sets in \eqref{cdp:eq:C.CR.6}.

\noindent\textbf{Step 6.} It follows that $S$ is the topological pushout of
$D_{\rm red}\hookrightarrow Z$ and
$D_{\rm red}\to\Gamma_{\rm red}$. This is the actual
pushout: its map to $S$ is a continuous bijection from a
compact space to a Hausdorff space. Triangulation of the
pair makes the first map a cofibration. Van Kampen applied
to its double mapping cylinder therefore gives
\[
 \pi_1(S)=\pi_1(Z)*_{\pi_1(D_{\rm red})}
                         \pi_1(\Gamma_{\rm red})
          \simeq\pi_1(\Gamma_{\rm red}).
\tag{C.CR.7}\label{cdp:eq:C.CR.7}
\]
Indeed the first amalgamating map is an isomorphism, and
the second map is trivial, since the section $C$ generates
$\pi_1(D_{\rm red})$ and $h(C)$ is a point. Every component
of $\Gamma_{\rm red}$ is dominated by a vertical rational
curve, so its normalization is $\PP^1$. Its fundamental
group is the free group of rank $b_1(\Gamma_{\rm red})$.

\noindent\textbf{Step 7.} For the exact local comparison with the nearby torus, take
a triangulated regular neighbourhood retracting onto $S$
and put a sufficiently small full fibre tube inside it.
Perturbing loops off the divisor gives a surjection from
the fundamental group of its punctured torus bundle onto
that of $S$. The former is $\mathbb Z^4\rtimes\mathbb Z$,
which is metabelian. A free group of rank at least two
is not metabelian, since it surjects onto $S_4$.
Consequently the rank in \eqref{cdp:eq:C.CR.7} is zero or one. The
nonzero integral class on every support supplied by \ref{cdp:st:D.1.9}
excludes zero. This proves the fundamental-group assertion.

\noindent\textbf{Step 8.} Normalization of the rational image components gives the
exact identities
\[
 b_1(\Gamma_{\rm red})
   =\sum_x(r_x-1)-n+1=1,\qquad
 e(\Gamma_{\rm red})=2n-\sum_x(r_x-1)=n.
\tag{C.CR.8}\label{cdp:eq:C.CR.8}
\]
Put $\epsilon=h^0(S,\omega_S)\in\{0,1\}$.
The vanishings already proved, the normalization exact
sequence and conductor duality give
\[
 h^0(\eta_*\OO_T/\OO_S)=1+\epsilon,\qquad
 h^1(\eta_*\OO_T/\OO_S)=\epsilon,
 \qquad\chi(\OO_\Gamma)=-1.
\tag{C.CR.9}\label{cdp:eq:C.CR.9}
\]
The nilradical sequence and curve duality give
$\chi(\OO_{\Gamma_{\rm red}})\ge-1$ exactly as in
C.2.VII.2: for $\epsilon=0$ the relevant twisted $H^1$
has dimension at most one, and for $\epsilon=1$ the twist
is trivial and connectedness supplies $h^0=1$.
On a connected curve with rational normalizations,
$\chi(\OO_{\Gamma_{\rm red}})=n-\sum_x\delta_x$.
Together with \eqref{cdp:eq:C.CR.8} and $\delta_x\ge r_x-1$, this
proves the remaining assertions of \eqref{cdp:eq:C.CR.1}.

\noindent\textbf{Step 9.} Finally $e(Z_0)=4-4g$ and
$e((D_0)_{\rm red})=2-2g+v$. A reverse blowup of
multiplicity at least two increases the Euler numbers
of both the surface and the reduced boundary by one,
leaving their difference unchanged. A multiplicity-one
centre leaves the boundary isomorphic to its predecessor
and increases the complement Euler number by one.
Therefore
\[
 e(Z\setminus D_{\rm red})=2-2g-v+N_0.
\]
Equation \eqref{cdp:eq:C.CR.6}, the excluded off-conductor exceptionals,
and additivity of Euler number now give \eqref{cdp:eq:C.CR.3}.
\end{proof}






\begin{proposition}[The ruling genus is at most two]\label{cdp:st:C.2.CR.2}
\StatementRecord{cdp:st:C.2.CR.2}{proved}{Historical events: Claude Fable 5.2 (Anthropic), effort high, 16 November 2026: carried as background without fault found. GPT-6 Astra (OpenAI), ChatGPT Pro mode, effort Pro, S6\_09\_PAPERS, 17 November 2026: targeted proof check. Also checked in this writing. Pre-continuation writing events: GPT-6 Astra (OpenAI), ChatGPT Pro mode, effort Pro, S6\_10\_PAPERS, 18 November 2026: targeted proof check; used and rechecked in this run. Continuation metadata events: GPT-6 Astra (OpenAI), ChatGPT Pro mode, effort Pro, S6\_10\_PAPERS continuation, 19 November 2026: used and rechecked in this run; status-vocabulary reconciliation; metadata reconciliation only; no new mathematical check. \emph{Pre-continuation recorded proof scope.} Proper duality and the residue sequence identify the branch canonical ideal modulo pi\_*omega with H0(C,omega\_C). The containment of its intersection with m squared preserves every vertical divisor order. Nakayama and purity exclude cotangent dimension three; dimension two forces the exact full conductor ideal (x,y), a unique normal branch and transverse delta one. Van Straten sections2.2.2-2.2.5 pp60-61 were visually opened and matched to the proper analytic resolution}{the hypotheses, notation and definitions stated here; \cref{cdp:st:C.2.CR.1}; \cite[Sections~2.2.2--2.2.5]{vanStraten1987}; \cref{cdp:eq:C.CR.10}; \cref{cdp:st:A.5.3}; \cref{cdp:eq:C.CR.12}}

In the situation of \ref{cdp:st:C.2.CR.1}, the ruling genus is $g=1$ or
$g=2$. Let $p=\pi(C)$ and $q=\eta(p)$. When $g=2$, the
normalization has just the one point $p$ over $q$, and the
full conductor-image scheme $\Gamma$ is smooth and reduced
at $q$. Its component through $q$ has generic transversal
either an ordinary node or an ordinary cusp.
\end{proposition}


\begin{proof}
\noindent\textbf{Step 1.} Work over a small neighbourhood of $q$ and fix a generator
of the invertible dualizing sheaf $\omega_S$. Finite
duality then identifies the conductor on the normal branch
$(T,p)$ with its canonical module:
\[
 I_p=\omega_{T,p}\subset\OO_{T,p}.
\]
The conductor ideal on $S$ is the direct sum of these
normalization-branch ideals; thus $I_p$ itself embeds
as an ideal of $R=\OO_{S,q}$ by putting zero on the other
normalization branches.

\noindent\textbf{Step 2.} Put $J=(\pi_*\omega_Z)_p$. Proper duality and
Grauert--Riemenschneider vanishing give
\[
 \operatorname{length}(I_p/J)
   =\operatorname{length}(R^1\pi_*\OO_Z)_p=g.
 \tag{C.CR.10}\label{cdp:eq:C.CR.10}
\]
These are the analytic duality and vanishing maps described
in \cite[Sections~2.2.2--2.2.5]{vanStraten1987}.
For precision, the defect of $\OO_T\to R\pi_*\OO_Z$ is
the finite-length module $R^1\pi_*\OO_Z$ in degree one;
duality turns it into the quotient
$\omega_T/\pi_*\omega_Z$, since the positive direct
images of $\omega_Z$ vanish. Finite-length duality
preserves length.

\noindent\textbf{Step 3.} The residue sequence for the smooth exceptional section is
\[
 0\longrightarrow\omega_Z\longrightarrow\omega_Z(C)
   \longrightarrow\omega_C\longrightarrow0 .
 \tag{C.CR.11}\label{cdp:eq:C.CR.11}
\]
The same vanishing theorem makes its pushforward quotient
at $p$ equal to $H^0(C,\omega_C)$, of dimension $g$.
The sheaf $\pi_*\omega_Z(C)$ embeds in $\omega_T$:
away from the exceptional set its sections are holomorphic
two-forms, and the canonical module of a normal surface
extends these across its isolated singular points.
Comparing lengths with \eqref{cdp:eq:C.CR.10} proves
\[
 I_p=(\pi_*\omega_Z(C))_p,\qquad
 I_p/J\simeq H^0(C,\omega_C).
 \tag{C.CR.12}\label{cdp:eq:C.CR.12}
\]
In the fixed dualizing trivialization the exact comparison
of \ref{cdp:st:A.5.3} reads
$\omega_Z=\OO_Z(-2C-V)$, so that
$\omega_Z(C)=\OO_Z(-C-V)$. The quotient in \eqref{cdp:eq:C.CR.12}
is consequently the leading coefficient of order one
along $C$, with the vertical vanishing conditions retained.

\noindent\textbf{Step 4.} Let $\mathfrak m$ be the maximal ideal of $R$. Every
element of $\mathfrak m$ pulls back to a function vanishing
on $C$, because $h(C)=q$. Its square therefore vanishes
to order at least two on $C$. Formula \eqref{cdp:eq:C.CR.12}, including
all vertical conditions in $V$, implies
\[
                 I_p\cap\mathfrak m^2\subset J.
 \tag{C.CR.13}\label{cdp:eq:C.CR.13}
\]
Indeed an element of $I_p$ is already a section of
$\OO_Z(-C-V)$; an additional zero along $C$ puts it
in $\OO_Z(-2C-V)$.

\noindent\textbf{Step 5.} It follows that
\[
 g=\dim(I_p/J)
 \le \dim\bigl(I_p/(I_p\cap\mathfrak m^2)\bigr)
 \le \dim(\mathfrak m/\mathfrak m^2)\le3.
 \tag{C.CR.14}\label{cdp:eq:C.CR.14}
\]
If the middle dimension were three, the ideal $I_p$
would generate the whole cotangent space. Nakayama's
lemma applied to $\mathfrak m/I_p$ would give
$I_p=\mathfrak m$. The full conductor ideal, which
contains $I_p$, would then be $\mathfrak m$, making its
image zero-dimensional. This contradicts the pure
one-dimensional Cohen--Macaulay conductor image.
Thus the middle dimension is at most two, proving
$g\le2$.

\noindent\textbf{Step 6.} When $g=2$, choose $x,y\in I_p$ whose linear parts
are independent, and complete them to coordinates
$(x,y,z)$ on the smooth ambient threefold. The full
conductor ideal contains $(x,y)$ in $R$. The defining
equation of $S$ belongs to $(x,y)$, since the conductor
has dimension one. The quotient by $(x,y)$ is therefore
$\mathbb C\{z\}$. A strictly larger proper conductor
ideal would give a zero-dimensional quotient of this
ring. Purity hence forces the full conductor ideal
to be exactly $(x,y)$, proving smoothness and
reducedness of $\Gamma$ at $q$.

\noindent\textbf{Step 7.} There is no other normalization point over $q$.
The two elements $x,y$ were chosen in the single
summand $I_p$, so both vanish identically on any other
local normalization branch. Such a two-dimensional
hypersurface branch cannot be contained in the
one-dimensional smooth germ $\{x=y=0\}$.

\noindent\textbf{Step 8.} Finally conductor duality gives
$Q\simeq\omega_\Gamma\otimes\omega_S^{-1}|_\Gamma$.
Near the smooth reduced curve $\Gamma$ this is a
line bundle, so the transverse normalization quotient
has length one. A reduced plane curve of delta
invariant one is a node or a cusp. One can see this
directly: multiple branches require exactly two
smooth branches with intersection multiplicity one;
a single branch has normalization
$\mathbb C\{u\}$ and a colength-one subring, necessarily
$\mathbb C\{u^2,u^3\}$.
\end{proof}






\begin{proposition}[The genus-two node alternative is excluded]\label{cdp:st:C.2.CR.3}
\StatementRecord{cdp:st:C.2.CR.3}{proved}{Historical events: Claude Fable 5.2 (Anthropic), effort high, 16 November 2026: carried as background without fault found. GPT-6 Astra (OpenAI), ChatGPT Pro mode, effort Pro, S6\_09\_PAPERS, 17 November 2026: targeted proof check. Also checked in this writing. Pre-continuation writing events: GPT-6 Astra (OpenAI), ChatGPT Pro mode, effort Pro, S6\_10\_PAPERS, 18 November 2026: targeted proof check. Continuation metadata events: GPT-6 Astra (OpenAI), ChatGPT Pro mode, effort Pro, S6\_10\_PAPERS continuation, 19 November 2026: used and rechecked in this run; status-vocabulary reconciliation; metadata reconciliation only; no new mathematical check. \emph{Pre-continuation recorded proof scope.} The coefficient-one leaf and shifted-gcd calculation identifies every strict conductor curve attached to the contracted component. The actual exceptional contraction makes the image conductor rational and unibranch. A second hypersurface branch would share a curve of intersection and contradict the unique normalized conductor branch. Its normalization is therefore a homeomorphism from P1, contradicting the previously derived infinite cyclic conductor group}{the hypotheses, notation and definitions stated here; \cref{cdp:st:C.2.CR.2}}

In the genus-two situation of \ref{cdp:st:C.2.CR.2}, the component
of $\Gamma$ through $q=h(C)$ cannot have generic
transversal an ordinary node. Consequently the
remaining genus-two alternative has an ordinary cusp
at its generic point, and exactly one strict conductor
curve of coefficient two and degree one lies above
that component.
\end{proposition}


\begin{proof}
\noindent\textbf{Step 1.} We first record an additional consequence of the
specified boundary blowups in C.2.CR.1: every
coefficient-one component of $D_{\rm red}$ is a leaf.
This holds on the initial section-and-fibres star.
A smooth boundary blowup creates coefficient one
only when the old coefficient is two, and creates
it as a leaf. A node blowup could create coefficient
one only from an edge $(1,1)$, which (C.CR.2)
excludes. Subsequent blowups preserve the valence
one of an existing coefficient-one vertex: at its
unique boundary node they merely subdivide that edge,
and at a smooth point its new exceptional coefficient
is zero.

\noindent\textbf{Step 2.} Let $E_0$ be the connected exceptional fibre containing
$C$. In the boundary tree, the first nonexceptional
vertex on every path from $C$ meets $E_0$, so its image
passes through $p$. Its further image is therefore
a component of $\Gamma$ through $q$. By \ref{cdp:st:C.2.CR.2}
there is only one such component, denoted $G$.

\noindent\textbf{Step 3.} Suppose its generic transversal were a node.
Every strict conductor curve above $G$ then has
coefficient one. Thus each first nonexceptional
vertex just found is a leaf. Every strict conductor
curve is consequently such a leaf, attached at a
single point to $E_0$. No additional exceptional
tail can meet it, since that would raise its
boundary valence. The reduced conductor on $T$
is therefore a union of rational curves whose only
possible multibranch point is $p$; every component
passes through $p$, is smooth away from $p$, and
has only one branch at $p$.

\noindent\textbf{Step 4.} All their images pass through $q$, where $\Gamma$
is smooth, so the conductor image has just the one
global component $G$. It also has no multibranch
point elsewhere. Indeed, at $x\ne q$ each point
of the normalization over $x$ receives just one
local branch of $G$. If $G$ had two distinct local
branches there, their receiving normalization
points would be distinct. The corresponding
irreducible hypersurface branches of $S$ must
intersect in a curve, because their distinct
local equations generate a height-two ideal in
the regular local ring of the threefold. That
intersection curve is a conductor-image branch
received by both normalization points, contradicting
their different singleton branch sets.

\noindent\textbf{Step 5.} The normalization of $G$ is rational and all its
germs are unibranch. The normalization map is a
continuous bijection from a compact space to a
Hausdorff space and hence a homeomorphism.
Thus $\pi_1(G)=1$, contradicting (C.CR.1).
The generic cusp alternative is the only one left.
Its transverse normalization is unibranch with
conductor exponent two, so it contributes exactly
one strict coefficient-two curve, of degree one
over $G$.
\end{proof}






\begin{proposition}[A smooth critical curve with generic cusp is a constant cusp]\label{cdp:st:C.2.CR.4}
\StatementRecord{cdp:st:C.2.CR.4}{proved}{Historical events: Claude Fable 5.2 (Anthropic), effort high, 16 November 2026: carried as background without fault found. GPT-6 Astra (OpenAI), ChatGPT Pro mode, effort Pro, S6\_09\_PAPERS, 17 November 2026: targeted proof check. Also checked in this writing. Pre-continuation writing events: GPT-6 Astra (OpenAI), ChatGPT Pro mode, effort Pro, S6\_10\_PAPERS, 18 November 2026: targeted proof check. Continuation metadata events: GPT-6 Astra (OpenAI), ChatGPT Pro mode, effort Pro, S6\_10\_PAPERS continuation, 19 November 2026: used and rechecked in this run; status-vocabulary reconciliation; metadata reconciliation only; no new mathematical check. \emph{Pre-continuation recorded proof scope.} Factoring the primitive rank-one quadratic matrix uses an actual holomorphic normal-coordinate change. Parameter Morse reduction treats the nonzero quadratic coefficient case; the two-parameter perturbation, Weierstrass coefficient divisibilities and Tschirnhaus substitution give the exact homogeneous germ Y cubed plus x squared Z in the remaining case. The uniform product-cusp boundary charts transport its spherical class of square minus four into the original torus, giving the contradiction. The common unit and multiplicity are retained by a local root}{the hypotheses, notation and definitions stated here; \cref{cdp:st:A.5.17}; \cref{cdp:eq:C.CR.17}; \cref{cdp:st:A.5.15}; \cref{cdp:eq:C.CR.18}}

Let a reduced hypersurface germ $\{f=0\}$ in a smooth
complex threefold occur in a fibre whose nearby smooth
fibres are complex two-tori. Suppose that the full reduced
critical locus of $f$ near the point is a smooth curve
and that the generic transverse singularity is an ordinary
cusp. Then there are local holomorphic coordinates in which
\[
                            f=X^2+Y^3.             \tag{C.CR.15}\label{cdp:eq:C.CR.15}
\]
An analytic unit and a common full-fibre multiplicity may
be retained in the original map: they do not change this
assertion about the reduced hypersurface.
\end{proposition}


\begin{proof}
\noindent\textbf{Step 1.} Choose coordinates $(x,y,z)$ with critical curve
$\{x=y=0\}$. The rank-one quadratic form along this curve
can be written
\[
 a(z)x^2.
\]
To obtain these coordinates without a meromorphic change,
divide the matrix of its coefficients by its common
vanishing order at $z=0$. A primitive symmetric rank-one
matrix has a unit entry, factors as a unit times the
square of a primitive linear form, and that linear form
extends to an invertible holomorphic change of the normal
coordinates. The intrinsic cubic coefficient on its
kernel is then $b(z)y^3$. The generic cusp hypothesis says
that both $a$ and $b$ are nonzero functions; shrink the
disc so that their only possible zero is at $z=0$.

\noindent\textbf{Step 2.} If $a(0)\ne0$, the Morse lemma with parameters gives
\[
 f=X^2+y^3u(y,z),\qquad u(0,z)=b(z)
\]
up to multiplication of $b$ by a holomorphic unit.
If $u(0,0)=0$, the reduced support of $y^3u(y,z)$ is
singular: $u$ has a component through the origin
different from $y=0$, because $u(0,z)$ is not identically
zero. The arbitrary-suspension clause of \ref{cdp:st:A.5.17} then
puts a sphere of nonzero self-intersection into a
nearby actual torus fibre, which is impossible.
Thus $u$ is a unit, and replacing $y$ by a holomorphic
cube root of $u$ times $y$ gives \eqref{cdp:eq:C.CR.15}.

\noindent\textbf{Step 3.} It remains to exclude $a(0)=0$. Consider the actual
holomorphic perturbation
\[
                  f_{\epsilon,\delta}
                      =f+\epsilon x^2+\delta y^3.
 \tag{C.CR.16}\label{cdp:eq:C.CR.16}
\]
Choose small parameters so that $a(z)+\epsilon$ has a
simple zero $z_\epsilon$ and
$b(z_\epsilon)+\delta\ne0$. Such choices exist:
regular values of the nonconstant function $a$ are
dense near zero, and finitely many values of $\delta$
are then excluded.

\noindent\textbf{Step 4.} We specify the germ produced at
$(0,0,z_\epsilon)$. Weierstrass preparation in $y$,
followed by a Tschirnhaus substitution, writes it as
a holomorphic unit times
\[
             Y^3+x^2\bigl(\beta(x,z)Y+\gamma(x,z)\bigr).
 \tag{C.CR.17}\label{cdp:eq:C.CR.17}
\]
Here the absence of a linear transverse jet gives
the divisibility of the constant term by $x^2$,
and the rank-one quadratic form gives the divisibility
of the coefficient of $Y$ by $x^2$. More explicitly,
before the Tschirnhaus substitution the coefficient
of $y^2$ is divisible by $x$, while those of $y$
and of $1$ are divisible by $x^2$; these divisibilities
are preserved by that substitution.
The function $\gamma(0,z)$ is a unit multiple of
$a(z)+\epsilon$. Consequently
\[
             Z=\beta(x,z)Y+\gamma(x,z)
\]
is an invertible change of the curve coordinate at
the chosen simple zero. Equation \eqref{cdp:eq:C.CR.17} is thus
a unit times $Y^3+x^2Z$. This remaining unit can
also be removed: this polynomial is homogeneous
of degree three, so scaling all three coordinates
by a holomorphic cube root of the unit gives the
exact equation
\[
                            Y^3+x^2Z.
 \tag{C.CR.18}\label{cdp:eq:C.CR.18}
\]
\ref{cdp:st:A.5.15} supplies a compact spherical cycle of square $-4$
in an arbitrarily small nonzero level of this germ.
Indeed its Hirzebruch formula is
$-2(b/\gcd(b,c))(c/\gcd(b,c))$, here with
$(b,c)=(2,1)$. The assertion is about the homology
class of the explicit continuous map from $S^2$;
no embedded representative is needed.

\noindent\textbf{Step 5.} For completeness we prove the boundary comparison
needed to put this spherical cycle in the original fibre.
Choose a sufficiently small closed ball whose
boundary meets the smooth critical curve transversely
and is transverse to the smooth stratum of $f=0$.
On the compact link of the critical curve, both
$a$ and $b$ are units. For all sufficiently small
$(\epsilon,\delta)$ the parametrized Morse and
cusp-coordinate changes therefore give holomorphic
product-cusp charts there, retaining $z$ as curve
coordinate. In each such chart the vector
$\partial/\partial z$ is tangent to every level
of $f_{\epsilon,\delta}$. Along the critical link,
an appropriate real radial vector in this direction
is transverse to the fixed sphere. This remains
true on a neighbourhood of that link, uniformly
for the small perturbations. Finitely many charts
suffice. On the rest of the boundary ordinary
compactness and transversality at parameter zero
give the same conclusion for small target values.
Thus all sufficiently small nonzero levels have
uniformly transverse boundary. The full-critical-
locus hypothesis ensures that no other boundary
singular curve has been omitted from this argument.

\noindent\textbf{Step 6.} For fixed small parameters, the nonzero critical
values inside the ball are finite. The local
spherical cycle from \eqref{cdp:eq:C.CR.18} can therefore be transported
through regular target values to a fixed small
value $t_1$, avoiding those critical values.
Choose the parameters small enough that every
critical value throughout the segment joining
them to zero has modulus less than $|t_1|/2$.
This is possible by compactness: every limit of
critical points is a critical point of $f$, and
its value is zero; the just-proved product charts
control the boundary as well. The family of
$t_1$-levels along this parameter segment is a
proper submersion with transverse boundary.
Ehresmann transport identifies it with the
original $f=t_1$ level and preserves the square
of the transported spherical class.

\noindent\textbf{Step 7.} The original level lies in a smooth nearby torus
fibre. A two-torus is aspherical, so every sphere
in it is null-homotopic and has zero homology
class and self-intersection. This contradicts
the nonzero square obtained above. Hence $a(0)\ne0$,
and the first part of the proof gives \eqref{cdp:eq:C.CR.15}.
If the original map is a unit times $f^m$, absorb
the unit by its local $m$th root before this
comparison; the resulting sphere lies in the
actual level with parameter $t_1^m$.
\end{proof}






\begin{proposition}[A contracted horizontal section is elliptic]\label{cdp:st:C.2.CR.5}
\StatementRecord{cdp:st:C.2.CR.5}{proved}{Historical events: Claude Fable 5.2 (Anthropic), effort high, 16 November 2026: carried as background without fault found. GPT-6 Astra (OpenAI), ChatGPT Pro mode, effort Pro, S6\_09\_PAPERS, 17 November 2026: targeted proof check. Also checked in this writing. Pre-continuation writing events: GPT-6 Astra (OpenAI), ChatGPT Pro mode, effort Pro, S6\_10\_PAPERS, 18 November 2026: targeted proof check. Continuation metadata events: GPT-6 Astra (OpenAI), ChatGPT Pro mode, effort Pro, S6\_10\_PAPERS continuation, 19 November 2026: used and rechecked in this run; status-vocabulary reconciliation; metadata reconciliation only; no new mathematical check. \emph{Pre-continuation recorded proof scope.} The residue cotangent bound leaves genus at most two. In genus two the exact conductor is smooth and has transverse delta one: the node is excluded by CR.3, while the product-cusp comparison CR.4 makes its normalization smooth. A resolution of that smooth germ cannot contract the positive-genus section. The remaining genus-one case gives the stated exact Euler formula}{the hypotheses, notation and definitions stated here; \cref{cdp:st:C.2.CR.1}; \cref{cdp:st:C.2.CR.2}; \cref{cdp:st:C.2.CR.3}; \cref{cdp:st:C.2.CR.4}}

In the situation of \ref{cdp:st:C.2.CR.1} one has
\[
                              g(B)=1.              \tag{C.CR.19}\label{cdp:eq:C.CR.19}
\]
In particular the remaining exact Euler expression is
\[
                            e(S)=n-v+N_0.          \tag{C.CR.20}\label{cdp:eq:C.CR.20}
\]
\end{proposition}


\begin{proof}
\noindent\textbf{Step 1.} By \ref{cdp:st:C.2.CR.2} the only other possibility is $g=2$.
The conductor image is then smooth and reduced at
$q=h(C)$, and its generic transversal is a node or
a cusp. \ref{cdp:st:C.2.CR.3} excludes the node. For the cusp,
the reduced critical locus near $q$ is precisely
that smooth conductor curve: every positive-dimensional
singular locus of a reduced hypersurface surface
belongs to its nonnormal locus, and there are no
other isolated points after the neighbourhood is
shrunk. \ref{cdp:st:C.2.CR.4} makes the germ a product cusp.
Its normalization is smooth, contrary to the
exceptional genus-two curve contracted to the
normalization point $p$. Thus $g=2$ is excluded.
Substitution in (C.CR.3) gives \eqref{cdp:eq:C.CR.20}.
\end{proof}






\begin{proposition}[The elliptic contraction has a planar conductor and only double transverse branches at its image]\label{cdp:st:C.2.CR.6}
\StatementRecord{cdp:st:C.2.CR.6}{proved}{Historical events: Claude Fable 5.2 (Anthropic), effort high, 16 November 2026: carried as background without fault found. GPT-6 Astra (OpenAI), ChatGPT Pro mode, effort Pro, S6\_09\_PAPERS, 17 November 2026: targeted proof check. Also checked in this writing. Pre-continuation writing events: GPT-6 Astra (OpenAI), ChatGPT Pro mode, effort Pro, S6\_10\_PAPERS, 18 November 2026: targeted proof check. Continuation metadata events: GPT-6 Astra (OpenAI), ChatGPT Pro mode, effort Pro, S6\_10\_PAPERS continuation, 19 November 2026: used and rechecked in this run; status-vocabulary reconciliation; metadata reconciliation only; no new mathematical check. \emph{Pre-continuation recorded proof scope.} The residue quotient supplies a linear conductor equation; Cohen-Macaulay purity makes the remaining height-one ideal principal. Pulling the residue generator to the first point blowup gives the order bound ord(a)>=mult(D)-1, so its linear conductor forces a double plane germ. Weierstrass preparation, square completion and a root of a unit give exactly A\_k. The individual-branch intersection and numerical-semigroup calculations give every stated conductor exponent and covering degree}{the hypotheses, notation and definitions stated here; \cref{cdp:st:C.2.CR.1}; \cref{cdp:eq:C.CR.22}}

In the remaining case $g=1$ of \ref{cdp:st:C.2.CR.1}, let
$p=\pi(C)$ and $q=\eta(p)$. The conductor-image
scheme is a local complete intersection at $q$.
More precisely, ambient coordinates can be chosen
so that its ideal is
\[
                         (x,g(y,z)).               \tag{C.CR.21}\label{cdp:eq:C.CR.21}
\]
Every global conductor component passing through
$q$ has generic transverse type $A_k$ for some
$k\ge1$. For type $A_{2s-1}$, its two transverse
normalization branches have conductor coefficient
$s$ each; their global covering degrees sum to two.
For type $A_{2s}$ there is one coefficient-$2s$
normalization branch, of covering degree one.
\end{proposition}


\begin{proof}
\noindent\textbf{Step 1.} The residue comparison (C.CR.12) is valid also
when $g=1$. By (C.CR.13) its nonzero one-dimensional
quotient requires an element of the conductor
summand $I_p$ with nonzero ambient linear part.
Call this element $x$ and extend it to coordinates
$(x,y,z)$ on the smooth threefold.

\noindent\textbf{Step 2.} The full conductor ideal contains $x$ and has a
pure one-dimensional Cohen--Macaulay quotient.
After quotienting the ambient local ring by $x$,
its remaining ideal is an unmixed height-one
ideal in the regular two-dimensional ring
$\mathbb C\{y,z\}$. Such an ideal is principal:
it is the intersection of its height-one primary
localizations, and each is a power of a principal
prime in this factorial ring. Equivalently its
quotient has depth one, so the ideal is a rank-one
reflexive module on this smooth surface and
is invertible. This proves \eqref{cdp:eq:C.CR.21}, retaining
all multiplicities in $g$.

\noindent\textbf{Step 3.} We give the precise adjunction comparison behind
the transverse multiplicity assertion. For a
reduced plane curve germ $D\subset W$, let
$\nu:\widetilde D\to D$ be its normalization,
and let $\rho:\widetilde W\to W$ be an embedded
resolution with smooth disjoint strict transform
$\widehat D$. Finite curve duality identifies
\[
 \mathfrak c_D\,\omega_D=\nu_*\omega_{\widetilde D}.
\]
The residue sequence on $\widetilde W$ and
Grauert--Riemenschneider vanishing identify the
right side with the residue of
$\rho_*\omega_{\widetilde W}(\widehat D)$.
In a local generator $du\wedge dv/f_D$ of the
ambient meromorphic residue description, a
conductor element $a$ therefore lifts to a form
\[
                     a\,du\wedge dv/f_D
\]
whose pullback has no exceptional pole. On the
first plane blowup, with exceptional curve $E$,
the canonical divisor has coefficient one and
the total transform of $D$ has coefficient
$m=\operatorname{mult}_0 D$. The coefficient
of this pulled-back form along $E$ is
$\operatorname{ord}_0(a)+1-m$. Hence
\[
 \mathfrak c_D\subset\mathfrak m_D^{\,m-1}
 \quad\hbox{in the sense that every ambient
 representative has order at least }m-1.
 \tag{C.CR.22}\label{cdp:eq:C.CR.22}
\]
Adding a multiple of $f_D$ does not change this
order bound, since $f_D$ has order $m$.

\noindent\textbf{Step 4.} At a general smooth point of any conductor
component through $q$, the ambient coordinate
$x$ remains in its conductor ideal and has
a nonzero transverse linear part: its zero
set is a smooth ambient plane containing the
conductor curve. Applying \eqref{cdp:eq:C.CR.22} to a
transverse plane slice gives $m\le2$.
The slice is singular, so $m=2$.
Weierstrass preparation, completion of the
square, and absorption of a unit then put
a reduced multiplicity-two plane curve in
the form $u^2-v^{k+1}=0$.

\noindent\textbf{Step 5.} For $k=2s-1$, the two branches are
$u=\pm v^s$; their intersection multiplicity
is $s$, and the conductor on each normalized
branch is $(v^s)$. For $k=2s$, the normalization
is $v=t^2$, $u=t^{2s+1}$, whose semigroup
has conductor $2s$, giving $(t^{2s})$.
These calculations specify both the coefficients
and the numbers of transverse branches; finite
normalization over the global image component
expresses those numbers as the asserted sums
of covering degrees.
\end{proof}






\begin{proposition}[A smooth full critical curve with generic node is a product node]\label{cdp:st:C.2.UN.9}
\StatementRecord{cdp:st:C.2.UN.9}{proved}{Historical events: Claude Fable 5.2 (Anthropic), effort high, 16 November 2026: carried as background without fault found. GPT-6 Astra (OpenAI), ChatGPT Pro mode, effort Pro, S6\_09\_PAPERS, 17 November 2026: targeted proof check. Also checked in this writing. Pre-continuation writing events: GPT-6 Astra (OpenAI), ChatGPT Pro mode, effort Pro, S6\_10\_PAPERS, 18 November 2026: targeted proof check. Continuation metadata events: GPT-6 Astra (OpenAI), ChatGPT Pro mode, effort Pro, S6\_10\_PAPERS continuation, 19 November 2026: used and rechecked in this run; status-vocabulary reconciliation; metadata reconciliation only; no new mathematical check. \emph{Pre-continuation recorded proof scope.} The rank-one symmetric perturbation with simple determinant zero, parameter Morse reduction and exact coordinate Z=h(y,z) produce X squared plus y squared Z. The compact spherical class of square minus four is transported through the fixed regular target level using the explicitly uniform boundary charts. Its image in the actual aspherical torus is nullhomotopic, contradicting its preserved square. The unit-root and full-multiplicity comparison were checked}{the hypotheses, notation and definitions stated here; \cref{cdp:eq:C.UNNODE.2}; \cref{cdp:st:A.5.15}; \cref{cdp:eq:C.UNNODE.3}}

Let a reduced hypersurface germ $\{f=0\}$ in a smooth
complex threefold occur in a fibre of a proper torus-surface
degeneration. Suppose that the full reduced critical locus
of $f$ is a smooth curve and its generic transverse
singularity is an ordinary node. There are local
holomorphic coordinates in which
\[
                            f=XY.                  \tag{C.UNNODE.1}\label{cdp:eq:C.UNNODE.1}
\]
The conclusion about the reduced hypersurface also holds
when the actual map has a holomorphic unit and a common
positive multiplicity.
\end{proposition}


\begin{proof}
\noindent\textbf{Step 1.} Choose coordinates $(x,y,z)$ with critical curve
$\Gamma=\{x=y=0\}$. Since $f$ vanishes on $\Gamma$ and
its first derivatives vanish there, its transverse
quadratic term is
\[
                  (x,y)Q(z)(x,y)^{\mathsf t},
\]
where $Q$ is a holomorphic symmetric $2\times2$ matrix.
Its determinant is not identically zero. Shrink the
curve disc so that $\det Q(z)\ne0$ for $z\ne0$.
If $Q(0)$ is invertible, the holomorphic Morse lemma
with parameter $z$ gives $f=X^2+Y^2$, and a linear
change gives \eqref{cdp:eq:C.UNNODE.1}.

\noindent\textbf{Step 2.} Suppose instead that $Q(0)$ is singular. We construct
arbitrarily small constant symmetric matrices $E$
for which $Q(z)+E$ has a rank-one value at a simple
zero of its determinant near the origin. Choose
$z_0\to0$ with $Q'(z_0)\ne0$; such points exist
because the nonconstant determinant makes $Q$
nonconstant. Choose a rank-one symmetric matrix
$R$ close to $Q(0)$ with
\[
                  \operatorname{tr}
                      (\operatorname{adj}(R)Q'(z_0))\ne0.
\]
This choice is possible also when $Q(0)=0$: for a
rank-one symmetric matrix with kernel spanned by
$v$, its adjugate is a nonzero scalar multiple of
$vv^{\mathsf t}$, and the polynomial
$v^{\mathsf t}Q'(z_0)v$ is not identically zero.
Thus the forbidden kernels form a proper finite
set in the projective line, and the rank-one
matrix can be chosen as close to $Q(0)$ as desired.
Set $E=R-Q(z_0)$. Then $E\to0$, and
\[
 \det(Q(z_0)+E)=0,\qquad
 \left.\frac{d}{dz}\det(Q(z)+E)\right|_{z=z_0}\ne0.
 \tag{C.UNNODE.2}\label{cdp:eq:C.UNNODE.2}
\]

\noindent\textbf{Step 3.} Consider the actual perturbation
$f_E=f+(x,y)E(x,y)^{\mathsf t}$ at $(0,0,z_0)$.
The rank-one Morse lemma, retaining $(y,z)$ as
parameters after a linear normal-coordinate change,
writes
\[
                         f_E=X^2+y^2h(y,z).
\]
The factor $y^2$ follows from $f_E(0,0,z)=0$ and
the vanishing of its normal first derivative along
the same curve. Equation \eqref{cdp:eq:C.UNNODE.2} says exactly
that $h(0,z_0)=0$ and $\partial h/\partial z(0,z_0)
\ne0$, up to a nowhere-zero coefficient from the
nondegenerate quadratic direction. Hence the actual
coordinate substitution $Z=h(y,z)$ gives
\[
                            f_E=X^2+y^2Z.          \tag{C.UNNODE.3}\label{cdp:eq:C.UNNODE.3}
\]
The explicit construction in \ref{cdp:st:A.5.15} gives a compact
spherical class in a small nonzero level of
\eqref{cdp:eq:C.UNNODE.3}. Its square is
$-2(2/\gcd(2,1))(1/\gcd(2,1))=-4$.
The construction supplies a continuous map from
$S^2$ representing this class; embeddedness is
not required.

\noindent\textbf{Step 4.} Here is the precise transport into an original
fibre. Choose a small closed ball whose boundary
is transverse to $\Gamma$ and to the smooth
stratum of $f=0$. Along the compact boundary link
of $\Gamma$, the matrix $Q$ is invertible.
For all sufficiently small $E$, the parameter
Morse lemma therefore gives product-node charts
there, retaining the curve coordinate. In these
charts $\partial/\partial z$ is tangent to every
level of $f_E$. A real radial curve-direction
vector is transverse to the fixed boundary at
the critical link and remains transverse on its
neighbourhood, uniformly for small $E$. On the
remaining compact boundary, ordinary transversality
gives the same uniform conclusion for small target
values. The full-critical-locus hypothesis ensures
that these charts cover every boundary singular
point requiring this argument.

\noindent\textbf{Step 5.} For a fixed small $E$, transport the spherical
class from \eqref{cdp:eq:C.UNNODE.3} through regular target values
to a fixed small value $t_1$, avoiding the finitely
many nonzero critical values in the ball. Choose
$E$ small enough that every critical value along
the segment from $E$ to $0$ has modulus less than
$|t_1|/2$. This follows by compactness: a limit of
critical points is a critical point of $f$ with
value zero; the just-proved boundary charts rule
out loss of transversality at the boundary.
The $t_1$-levels along that segment form a proper
submersion with boundary. Ehresmann transport
therefore maps the actual spherical class to a
class of square $-4$ in the original $f=t_1$ level.
Its inclusion into the nearby torus preserves
the square, since both the compact representative
and its intersection push-off lie in that open
level. But this inclusion also maps its representing
$S^2$ to an aspherical torus, where it is
nullhomotopic and has zero homology class. This
contradiction proves $\det Q(0)\ne0$.

\noindent\textbf{Step 6.} If the original map is $u f^m$, replace the reduced
equation by $u^{1/m}f$ using a local holomorphic
root of the unit. Its critical-set germ is still
$\Gamma$: every critical component through the
point has constant value zero, and the singular
locus of the reduced zero hypersurface is unchanged.
The argument then places the
same class in the actual nearby level with base
parameter $t_1^m$. This proves the last assertion.
\end{proof}






\begin{proposition}[A contracted irrational horizontal section is elliptic without ambient Betti assumptions]\label{cdp:st:C.2.UN.10}
\StatementRecord{cdp:st:C.2.UN.10}{proved}{Historical events: Claude Fable 5.2 (Anthropic), effort high, 16 November 2026: carried as background without fault found. GPT-6 Astra (OpenAI), ChatGPT Pro mode, effort Pro, S6\_09\_PAPERS, 17 November 2026: targeted proof check. Also checked in this writing. Pre-continuation writing events: GPT-6 Astra (OpenAI), ChatGPT Pro mode, effort Pro, S6\_10\_PAPERS, 18 November 2026: targeted proof check; used and rechecked in this run. Continuation metadata events: GPT-6 Astra (OpenAI), ChatGPT Pro mode, effort Pro, S6\_10\_PAPERS continuation, 19 November 2026: used and rechecked in this run; status-vocabulary reconciliation; metadata reconciliation only; no new mathematical check. \emph{Pre-continuation recorded proof scope.} Canonical evaluation on the actual irreducible full tube supplies the torsion dualizing line and coherent Euler zero. Formal functions gives the exact exceptional H1 length, and the residue map identifies the branch conductor with pi\_*omega(C). Its intersection with the square of the ambient maximal ideal retains all vertical valuations and bounds the cotangent dimension. In genus two, the exact ideal (x,y), unique normal branch and node/cusp normalizations give the contradiction. Van Straten sections2.2.2-2.2.5 pp60-61 were visually opened for the proper analytic duality and canonical higher-direct-image vanishing}{the hypotheses, notation and definitions stated here; \cref{cdp:st:C.2.CR.1}; \cref{cdp:st:C.2.CR.2}; \cref{cdp:st:C.2.UN.9}; \cref{cdp:st:C.2.CR.4}}

Let $f:Y\to B$ be a proper holomorphic map from a
smooth complex threefold to a smooth curve, with
connected fibres and smooth general fibre a complex
two-torus. Let $mS$ be an irreducible nonnormal full
fibre. Write $h=\eta\pi:Z\to T\to S$ for the
normalization and its minimal resolution. Suppose
that $Z$ is ruled over a curve of positive genus,
that its effective adjunction divisor is
$D=2C+V$, with $C$ a section and $V$ vertical,
and that $\pi$ contracts $C$. Then the ruling
base has genus one.
\end{proposition}


\begin{proof}
\noindent\textbf{Step 1.} We first supply the comparisons that previously
came from the global Betti assumptions. On a
small full fibre tube $U$, the rank-one torsion-free
sheaf $f_*K_U$ is a line bundle. A generator has
no horizontal zeros, since its restriction to a
smooth torus is a nonzero canonical form. Its
zero divisor is consequently $kS$ for some
$k\ge0$, after shrinking away other special
fibres. Thus the evaluation map gives the actual
isomorphism
\[
 K_U\simeq\OO_U(kS),\qquad
 \omega_S\simeq\OO_U((k+1)S)|_S.                \tag{C.UNCR.1}\label{cdp:eq:C.UNCR.1}
\]
The divisor $mS$ is principal, so both these line
bundles are torsion. The map to the smooth curve
is flat, and the holomorphic Euler characteristic
of its full fibres is constant. The filtration
of $\OO_{mS}$ has quotients
$\OO_S(-jS)$ for $0\le j<m$; each is a torsion
line-bundle twist of $\OO_S$. Analytic
Riemann--Roch gives the same Euler characteristic
for these twists, since their rational Chern
characters are one. It follows that
\[
            0=\chi(\OO_{mS})=m\,\chi(\OO_S).
 \tag{C.UNCR.2}\label{cdp:eq:C.UNCR.2}
\]
The exact normalization and resolution duality
maps therefore give $K_Z+D=h^*\omega_S$, with
the same effective divisor and torsion right-hand
side as in \ref{cdp:st:C.2.CR.1}--\ref{cdp:st:C.2.CR.2}.

\noindent\textbf{Step 2.} Here are the parts of that argument needed here,
with no transgression-class assumption. Finite
duality embeds $\omega_T$ into the torsion line
$\eta^*\omega_S$, with its sections vanishing on
the nonempty conductor. A nonzero section of a
torsion line on a connected compact normal space
is nowhere zero: taking a trivializing power
makes it a nonzero constant. Consequently
$H^2(T,\OO_T)=0$. The coherent Leray sequence
then bounds $\operatorname{length}R^1\pi_*\OO_Z$
by $h^1(Z,\OO_Z)=g$, whereas formal functions
and restriction to the contracted section $C$
give length at least $h^1(C,\OO_C)=g$ at its
image point $p$. Thus equality holds at $p$ and
there are no further contributions.

\noindent\textbf{Step 3.} At $q=\eta(p)$, fix a local generator of $\omega_S$.
Finite duality identifies the branch conductor
ideal $I_p$ with $\omega_{T,p}$, embedded as an
actual ideal summand in $\OO_{S,q}$. Proper
duality and the residue sequence of $C$ give
exactly
\[
 I_p=(\pi_*\omega_Z(C))_p,\qquad
 I_p/(\pi_*\omega_Z)_p\simeq H^0(C,\omega_C).
 \tag{C.UNCR.3}\label{cdp:eq:C.UNCR.3}
\]
The equality of lengths follows from the preceding
formal-functions computation. In the fixed local
dualizing trivialization these sheaves are
$\OO_Z(-C-V)$ and $\OO_Z(-2C-V)$.
The maximal ideal of $\OO_{S,q}$ vanishes on $C$;
its square supplies an additional order along
$C$ without losing any prescribed order along
$V$. Hence
$I_p\cap\mathfrak m_q^2\subset(\pi_*\omega_Z)_p$.
The embedding dimension three and Nakayama
argument of \ref{cdp:st:C.2.CR.2} now give $g\le2$. When
$g=2$ they give the actual conductor ideal
$(x,y)$, a unique normalization point $p$ above
$q$, and generic transverse delta invariant
one. Thus the full conductor-image curve is
smooth and reduced at $q$, and the generic
transversal is a node or a cusp.

\noindent\textbf{Step 4.} Every positive-dimensional critical component
of the reduced surface is part of that conductor
curve; after shrinking, its full reduced critical
locus is exactly the smooth curve just obtained.
The node case is a product node by \ref{cdp:st:C.2.UN.9}.
Its normalization has two smooth points above
$q$, contradicting the unique point $p$.
The cusp case is a product cusp by \ref{cdp:st:C.2.CR.4}.
Its explicit normalization
$(u,z)\mapsto(X,Y,z)=(u^3,-u^2,z)$, after
adjusting the harmless sign in $X^2+Y^3$,
is smooth. This contradicts a minimal
normalization resolution contracting a
positive-genus curve to $p$. Both genus-two
possibilities are excluded, and $g=1$ follows.
\end{proof}






\begin{proposition}[The reduced horizontal case is a minimal elliptic ruled double-curve fibre]\label{cdp:st:C.2.IRR.1}
\StatementRecord{cdp:st:C.2.IRR.1}{proved}{Historical events: Claude Fable 5.2 (Anthropic), effort high, 16 November 2026: carried as background without fault found. GPT-6 Astra (OpenAI), ChatGPT Pro mode, effort Pro, S6\_09\_PAPERS, 17 November 2026: targeted proof check. Also checked in this writing. Pre-continuation writing events: GPT-6 Astra (OpenAI), ChatGPT Pro mode, effort Pro, S6\_10\_PAPERS, 18 November 2026: targeted proof check; used and rechecked in this run. Continuation metadata events: GPT-6 Astra (OpenAI), ChatGPT Pro mode, effort Pro, S6\_10\_PAPERS continuation, 19 November 2026: used and rechecked in this run; status-vocabulary reconciliation; metadata reconciliation only; no new mathematical check. \emph{Pre-continuation recorded proof scope.} The adjunction and blowdown identities, exclusion of coefficient-zero exceptional fibres outside the conductor, exact conductor double cover, pinched-point obstruction on a rational image and elliptic Hessian degree were re-derived. The resulting constant double curve and actual nearby-cycle Euler identity give smooth minimal normalization and Euler zero. A'Campo1975 sections2-3 pp237-245 were opened in primary-host indexed text; no visual-scan credit is claimed}{the hypotheses, notation and definitions stated here; \cref{cdp:st:A.5.3}; \cref{cdp:eq:C.IR.3}; \cref{cdp:eq:C.IR.2}; \cref{cdp:eq:C.IR.4}; \cref{cdp:eq:C.IR.5}; \cref{cdp:st:A.5.17}; \cref{cdp:st:C4e}; \cref{cdp:st:A.5.14}; \cref{cdp:eq:C.IR.6}; \cref{cdp:st:A.5.15}; \cref{cdp:st:A.5.5}; \cref{cdp:st:C.2.NC.1}}

Let $f:Y\to C_0$ be a proper holomorphic map from a smooth compact
complex threefold with $H^2(Y,\mathbb Q)=0$, with connected fibres
and smooth general fibre a complex two-torus. Let $mS$ be a full
fibre with irreducible nonnormal reduced support. Let
$h=\eta\pi:Z\to T\to S$ be the minimal resolution of its
normalization, and suppose that $p:Z\to B$ is an irrational
ruling. Write the effective divisor of \ref{cdp:st:A.5.3} as
$D_Z=H+V$, where $V$ is vertical and the horizontal part $H$
is reduced, of total degree two over $B$.

Then $B$ is elliptic, $T=Z$ is a smooth minimal ruled surface
over $B$, and $D_Z=H$ is a disjoint union of smooth elliptic
curves. The reduced conductor image is a disjoint union of
smooth elliptic curves. Its inverse image $H$ maps onto it
by a finite etale covering of total degree two over each
component. The support $S$ has only ordinary double-curve
normal crossings and
\[
                         e(S)=0.                 \tag{C.IR.1}\label{cdp:eq:C.IR.1}
\]
No assumption on the initial normalization singularities or
on the ambient Picard pullback is made.
\end{proposition}


\begin{proof}
\noindent\textbf{Step 1.} All divisor intersections in this proof use rational Chern
classes. The actual canonical comparison in \ref{cdp:st:A.5.3} gives
\[
 K_Z+D_Z=h^*\omega_S,\qquad c_1(K_Z+D_Z)=0.
                                                     \tag{C.IR.2}\label{cdp:eq:C.IR.2}
\]
For an irreducible component $H_i$ of the reduced horizontal
part, adjunction therefore gives
\[
 2p_a(H_i)-2+(D_Z-H_i)\cdot H_i=0.                 \tag{C.IR.3}\label{cdp:eq:C.IR.3}
\]
The second term is nonnegative. The normalization of $H_i$
maps finitely onto the positive-genus curve $B$; Riemann--Hurwitz
makes its genus at least one. Equality in \eqref{cdp:eq:C.IR.3} thus makes
$H_i$ a smooth elliptic curve, makes $B$ elliptic, and makes
all horizontal components disjoint from each other and from
$V$.

\noindent\textbf{Step 2.} We give the required comparison under the ruling blowdowns.
For one $\sigma:W\to W'$ with exceptional curve $E$, put
$D'=\sigma_*D_W$, let $a$ be the coefficient of $E$ in $D_W$,
and let $d$ be the multiplicity of $D'$ at the centre. The
exact divisor identity is
\[
 K_W+D_W=\sigma^*(K_{W'}+D')+(1+a-d)E.
                                                     \tag{C.IR.4}\label{cdp:eq:C.IR.4}
\]
Intersection with $E$ gives $a=d-1$, so \eqref{cdp:eq:C.IR.2} descends.
An exceptional curve in the vertical conductor is disjoint
from $H$. An exceptional curve outside the conductor has
$D_W\cdot E=1$, and therefore meets the horizontal part in
at most one simple point; it cannot identify two horizontal
points or join a horizontal component to the vertical
conductor. Thus the horizontal curves stay smooth, disjoint,
and disjoint from the remaining vertical conductor after
each blowdown. On the resulting minimal ruled bundle every
nonzero effective vertical divisor meets the horizontal
part. Its vertical conductor is consequently zero.

\noindent\textbf{Step 3.} In the reverse blowup sequence a centre away from the current
horizontal divisor would have $d=0$ in \eqref{cdp:eq:C.IR.4} and create
coefficient $-1$ on the new exceptional curve. That coefficient
persists on its strict transform under later blowups, contrary
to effectiveness of the original $D_Z$. Every centre
therefore lies on the smooth coefficient-one horizontal
divisor and creates coefficient zero. Induction gives
\[
                             D_Z=H.               \tag{C.IR.5}\label{cdp:eq:C.IR.5}
\]
This argument did not require the map $h$ to factor through
any ruling blowdown.

\noindent\textbf{Step 4.} We next prove that $\pi$ has no exceptional curves, including
horizontal ones. An exceptional curve $E$ of coefficient
zero in \eqref{cdp:eq:C.IR.5} satisfies $K_Z\cdot E\ge0$ on the minimal
resolution: adjunction and negative self-intersection give
this inequality, the sole contrary possibility being a
smooth rational $(-1)$-curve, which minimality excludes.
Equation \eqref{cdp:eq:C.IR.2} then forces $K_Z\cdot E=H\cdot E=0$.
Such a curve is disjoint from $H$. A positive-coefficient
exceptional curve is one of the mutually disjoint elliptic
curves $H_i$. Every connected exceptional fibre is therefore
disjoint from the strict conductor curves.

\noindent\textbf{Step 5.} The normal conductor ideal on $T$ is the rank-one reflexive
sheaf $\omega_T\otimes\eta^*\omega_S^{-1}$, embedded in
$\mathcal O_T$ by finite duality. On an open set without
its conductor divisors it equals $\mathcal O_T$: equality
off a codimension-two set extends by reflexivity on a normal
surface. Thus $\eta$ is an isomorphism near the image of
every exceptional fibre just described. Such an image is
an isolated normal hypersurface singularity of $S$.
The isolated-singularity clause of \ref{cdp:st:A.5.17} excludes it from
this actual torus smoothing, including its full multiplicity
$m$ and its local unit. Hence $\pi$ is an isomorphism.
In particular $T=Z$ is smooth and its conductor is exactly
the reduced divisor $H$.

\noindent\textbf{Step 6.} Let $G$ be the conductor image. Its ideal is the contraction
of the radical ideal of $H$ under the finite normalization,
so it is radical. Thus $G$ itself is reduced. Every local
point of $H$ has a single smooth branch, whose image is one
local branch of $G$. We claim that $G$ has only unibranch
germs and that its distinct global components are disjoint.
Otherwise choose two local branches of $G$ at a common point
and points of $T$ receiving the respective branches. These
points are different because $H$ is smooth at each point.
They correspond to different irreducible hypersurface
branches of $S$ in the smooth threefold. Two such branches
meet in a curve: their distinct equations generate an ideal
of height two in the regular three-dimensional local ring.
That intersection is a conductor-image branch received by
both points of $T$. Each of those points receives only one
branch, contradicting the chosen distinct images. This
proves the claim.

\noindent\textbf{Step 7.} The generic transverse curve has conductor exponent one on
every normalized branch. The branchwise conductor formula
in the proof of \ref{cdp:st:C4e} shows that it has exactly two smooth
transverse branches: a singular branch has intrinsic
conductor exponent at least two, and for smooth branches
the exponent is the sum of their positive pairwise
intersection multiplicities. Consequently the finite map
$H\to G$ has total degree two over every irreducible component.

\noindent\textbf{Step 8.} Fix such a component $G_j$, with normalization
$\nu:E_j\to G_j$. Each component of its inverse image is
elliptic. Riemann--Hurwitz gives $g(E_j)\le1$.
We exclude $E_j=\mathbf P^1$ by an actual local comparison.
In that case its inverse image must be one elliptic curve
of degree two; two degree-one maps from elliptic curves to
$\mathbf P^1$ cannot occur. Its four ramification points
have four distinct images on $\mathbf P^1$. \ref{cdp:st:A.5.14} gives
$p_a(G_j)\le2$, and all singularities of $G_j$ are unibranch.
There are therefore at most two normalization points above
its singular locus. At least one of the four branch points
lies over a smooth point of $G_j$.

\noindent\textbf{Step 9.} At that point choose coordinates $(u,v)$ on $T$, with
$H=\{u=0\}$, so that the degree-two conductor map is
$z=v^2$. The normalization-conductor square is cartesian
as a square of local analytic rings. Because its conductor
is $(u)$, the surface local ring is exactly
\[
 \{a\in\mathbb C\{u,v\}:a(0,v)\in\mathbb C\{v^2\}\}
   =\mathbb C\{u,uv,v^2\}
   \simeq\mathbb C\{x,y,z\}/(x^2-y^2z).       \tag{C.IR.6}\label{cdp:eq:C.IR.6}
\]
Here $(x,y,z)=(uv,u,v^2)$, so every map in the comparison
is specified. The pinch germ in \eqref{cdp:eq:C.IR.6} is excluded by
\ref{cdp:st:A.5.15}. It follows that $E_j$ is elliptic.

\noindent\textbf{Step 10.} Now \ref{cdp:st:A.5.14} bounds the total delta invariant of $G_j$ by one
and proves that its embedding in $Y$ is a local complete
intersection. On $E_j$ let $N$ be the dual of its pulled-back
conormal bundle, and let $\Xi$ be its normalization conductor
divisor. The determinant comparison from \ref{cdp:st:A.5.5} is
\[
 \det N=K_{E_j}(\Xi)\otimes\nu^*K_Y^{-1},
 \qquad \deg N=2\delta_j.                       \tag{C.IR.7}\label{cdp:eq:C.IR.7}
\]
The reduced equation of $S$ belongs to the ordinary square
of the ideal of $G_j$: the ideal is a regular sequence,
its powers are unmixed, and it has order at least two at
the generic point. Its quadratic jet is therefore a
holomorphic section of $\operatorname{Sym}^2N^*\otimes M$,
where $M=\nu^*\mathcal O_Y(S)$ has degree zero. Its determinant
is nonzero generically and has degree
\[
 \deg\bigl((\det N)^{-2}\otimes M^2\bigr)
                         =-4\delta_j.           \tag{C.IR.8}\label{cdp:eq:C.IR.8}
\]
Thus $\delta_j=0$. The curve $G_j$ is smooth elliptic and
the determinant, now a degree-zero holomorphic section,
is nowhere zero. The holomorphic Morse lemma with the
curve parameter gives the local reduced equation $xy=0$
at every point of $G_j$. This also proves that the
degree-two map $H\to G$ is etale everywhere.

\noindent\textbf{Step 11.} There are no triple points, because $H$ is smooth and its
components are disjoint. Apply the already proved
normal-crossings comparison \ref{cdp:st:C.2.NC.1} to the actual full
fibre $mS$. It gives
\[
 0=e(Y_t)=m\,e(T\setminus H)=m\,e(T),            \tag{C.IR.9}\label{cdp:eq:C.IR.9}
\]
because $e(H)=0$. If $N_0$ is the number of ruling blowups
above a minimal bundle over the elliptic curve $B$, the
point-blowup formula gives $e(T)=N_0$. Hence $N_0=0$ and
$T$ is minimal. Finally normalization is an isomorphism
off $H$ and $G$, so
$e(S)=e(T)-e(H)+e(G)=0$. This proves every assertion,
with arbitrary full-fibre multiplicity retained.
\end{proof}



The elliptic ruled geometry in \ref{cdp:st:C.2.IRR.1} has a direct literature
comparison. Kamenova's Theorem 3.1 in \cite{KamenovaTori} treats
semistable Kahler degenerations with trivial total-space canonical
bundle; its elliptic ruled cases have elliptic double curves arranged
along a path or a cycle. Our finite normalization $T\to S$ specifies
the self-identification of the double curves in the irreducible case.
The proof above establishes this geometry from the actual conductor
and the full-fibre smoothing, retaining the multiplicity and proving
the necessary local pinch exclusion. We have not supplied a Kahler
form or a trivialization of the total-space canonical bundle needed
to apply that theorem directly to our tube. 




\begin{proposition}[Proved anticanonical support comparison]\label{cdp:st:C.2.VII.1}
\StatementRecord{cdp:st:C.2.VII.1}{proved}{Historical events: Claude Fable 5.2 (Anthropic), effort high, 16 November 2026: carried as background without fault found. GPT-6 Astra (OpenAI), ChatGPT Pro mode, effort Pro, S6\_09\_PAPERS, 17 November 2026: targeted proof check. Also checked in this writing. Pre-continuation writing events: GPT-6 Astra (OpenAI), ChatGPT Pro mode, effort Pro, S6\_10\_PAPERS, 18 November 2026: targeted proof check; used and rechecked in this run. Continuation metadata events: GPT-6 Astra (OpenAI), ChatGPT Pro mode, effort Pro, S6\_10\_PAPERS continuation, 19 November 2026: used and rechecked in this run; status-vocabulary reconciliation; metadata reconciliation only; no new mathematical check. \emph{Pre-continuation recorded proof scope.} The torsion-line descent through every point blowdown, exact divisor transformation, nodal-cycle and leaf Euler changes, and rational-curve count were re-derived. The parabolic Inoue coefficient is fixed to one by the actual canonical relation and positive Gauduchon degree. The cited anticanonical graph and support propositions remain exact external inputs, with their remaining opening gates recorded separately}{the hypotheses, notation and definitions stated here; \cite[Theorem~0.2, Proposition~2.29 and Theorem~2.38]{Dloussky2006}; \cite[\S1, pp.~284--285, and \S3, p.~296]{DOT2003}; \cite[Corollary~2.25 and Proposition~2.27]{Dloussky2006}; \cite[Introduction, p.~283]{DOT2003}; \cite[Lemma~2.22 and the following paragraph; Proposition~2.27]{Dloussky2006}; \cref{cdp:st:C.3.8}}


Let $Z$ be a smooth compact surface of class VII, not
necessarily minimal, and let $D>0$ be an effective integral
divisor such that $K_Z+D$ is torsion in $\Pic(Z)$.
Write $R=D_{\mathrm{red}}$, and let $r(R)$ be its number
of rational irreducible components. Then
\[
 e(R)=r(R),\qquad r(Z)\le b_2(Z),
\tag{C.VII.1}\label{cdp:eq:C.VII.1}
\]
where $r(Z)$ counts all rational curves of $Z$.
Every connected component of $R$ has arithmetic genus one.
In particular a connected component containing an elliptic
curve consists of that single elliptic curve with trees of
rational curves attached; it cannot contain a second
elliptic curve or a path joining two distinct points of the
first one.
\end{proposition}


\begin{proof}
\noindent\textbf{Step 1.} We specify the established surface results used in this
comparison. For a minimal class VII surface with positive
second Betti number, existence of a numerically
anticanonical divisor implies existence of a global
spherical shell by
\cite[Theorem~0.2, Proposition~2.29 and Theorem~2.38]{Dloussky2006}.
The reduced support is one rational cycle with attached
trees, one or two rational cycles, or, in the parabolic
Inoue case, a disjoint rational cycle and a smooth elliptic
curve. The intermediate graph and its actual curve
configuration are described in
\cite[\S1, pp.~284--285, and \S3, p.~296]{DOT2003}.
The fact that no connected piece of the anticanonical
support is an isolated tree is
\cite[Corollary~2.25 and Proposition~2.27]{Dloussky2006};
Lemma~2.22 there shows that a curve meeting that support
belongs to it. These statements apply to our divisor with
index one, since its twisting line is torsion and therefore
flat. They give, on the minimal model, arithmetic genus
one for each connected component and Euler number equal
to the number of its rational components. The rational
curve count is at most $b_2$, as stated and used in
\cite[Introduction, p.~283]{DOT2003}.
These are consequences of the proved theorem about
anticanonical divisors; no classification assertion about
class VII surfaces without such a divisor is used.

\noindent\textbf{Step 2.} There is a coefficient issue in the parabolic Inoue case
which the torsion hypothesis resolves explicitly. Write
$E$ for its elliptic curve and $C$ for its reduced rational
cycle. The exact canonical relation is
$K_{Z_0}+E+C=0$; the curves $E$ and $C$ are disjoint,
$E^2=-b_2(Z_0)$, and the intersection matrix of the
components of $C$ has kernel generated by their sum
\cite[Lemma~2.22 and the following paragraph; Proposition~2.27]{Dloussky2006}.
Consequently an effective integral divisor with
$c_1(K_{Z_0}+D_0)_{\mathbb R}=0$ has the form
$D_0=E+kC$, $k\ge0$: intersection with $E$ fixes its
coefficient at one, and the cycle intersection matrix
makes all cycle coefficients equal. Our stronger
torsion relation then gives
$K_{Z_0}+D_0=(k-1)C$ torsion. For any Gauduchon metric
$\omega$, the degree of $\OO(C)$ is
$\int_C\omega>0$, whereas a torsion line has degree
zero. Hence $k=1$. In particular both connected
pieces of the stated reduced support really occur.
This deduction uses the torsion relation, not merely
the vanishing of its real first Chern class.

\noindent\textbf{Step 3.} For a minimal model with $b_2=0$, the nonzero effective
divisor excludes the curve-free Inoue surfaces. The
Hopf alternative has only smooth elliptic curves, which
are pairwise disjoint. This follows as well from the
explicit contraction description used in \ref{cdp:st:C.3.8} and its
finite unramified quotient: the preimages of two distinct
curves are disjoint elliptic curves, and the quotient
is locally biholomorphic. All assertions are then
immediate on that model.

\noindent\textbf{Step 4.} We verify the transfer through the actual blowdowns.
Let $\sigma:Z\to Z_0$ be a sequence of point blowups
from a minimal model. A torsion line is represented by
a finite unitary character of the fundamental group,
and a point blowup induces an isomorphism on fundamental
groups. Thus the torsion line $K_Z+D$ descends to $Z_0$.
Pushing forward a meromorphic canonical section with pole
divisor $D$, or its torsion-trivializing power, gives an
effective divisor $D_0$ with $K_{Z_0}+D_0$ torsion.
It is nonzero, since a torsion canonical bundle would
give Kodaira dimension zero. At each individual blowup
the canonical divisor formula and the pushforward section
give the exact divisor equality
\[
 D_{\mathrm{new}}=\sigma^*D_{\mathrm{old}}-E.
\tag{C.VII.2}\label{cdp:eq:C.VII.2}
\]
The coefficient of $E$ is
$\operatorname{mult}_pD_{\mathrm{old}}-1$.
Effectivity forces the centre to lie on $D_{\mathrm{old}}$.
If this coefficient is positive, the reduced support
is its entire reduced inverse image: a smooth point
acquires a rational leaf, or an ordinary intersection
is subdivided by the exceptional curve. These operations
preserve arithmetic genus and increase both the Euler
number and the rational-component count by one.
If the coefficient is zero, the centre is a smooth point
of a multiplicity-one component and the new support is
its strict transform, isomorphic to the old support.
Thus all the stated support properties transfer in the
required direction. Finally each blowup increases $b_2$
by one and adds precisely its exceptional rational
curve; every other rational curve is the strict transform
of one on the preceding surface. The curve-count
inequality in \eqref{cdp:eq:C.VII.1} follows.
\end{proof}






\begin{proposition}[Proved global conductor Euler bound]\label{cdp:st:C.2.VII.2}
\StatementRecord{cdp:st:C.2.VII.2}{proved}{Historical events: Claude Fable 5.2 (Anthropic), effort high, 16 November 2026: carried as background without fault found. GPT-6 Astra (OpenAI), ChatGPT Pro mode, effort Pro, S6\_09\_PAPERS, 17 November 2026: targeted proof check. Also checked in this writing. Pre-continuation writing events: GPT-6 Astra (OpenAI), ChatGPT Pro mode, effort Pro, S6\_10\_PAPERS, 18 November 2026: targeted proof check. Continuation metadata events: GPT-6 Astra (OpenAI), ChatGPT Pro mode, effort Pro, S6\_10\_PAPERS continuation, 19 November 2026: used and rechecked in this run; status-vocabulary reconciliation; metadata reconciliation only; no new mathematical check. \emph{Pre-continuation recorded proof scope.} Finite duality retains the full conductor scheme and gives coherent Euler at least minus one. The exact topological normalization map then leaves only an elliptic node as a possible negative case. Height-two intersections of actual hypersurface branches force a common normalization point, and its connected resolution fibre contradicts the proved genus-one conductor-support graph}{the hypotheses, notation and definitions stated here; \cref{cdp:st:A.5.12}; \cref{cdp:st:A.5.3}; \cref{cdp:st:C.2.VII.1}; \cref{cdp:eq:C.VII.5}; \cref{cdp:eq:C.VII.6}}


Let $S$ be a compact irreducible reduced hypersurface in
a smooth complex threefold, with $\omega_S$ torsion and
$\chi(S,\OO_S)=0$. Let $\eta:T\to S$ be normalization,
let $\Gamma$ be its conductor image scheme, and put
$C=\Gamma_{\mathrm{red}}$. Then
\[
 \chi(C,\OO_C)\ge-1.
\tag{C.VII.3}\label{cdp:eq:C.VII.3}
\]
Suppose in addition that the minimal resolution
$\pi:Z\to T$ is a class VII surface. Then
\[
 e(C)\ge0.
\tag{C.VII.4}\label{cdp:eq:C.VII.4}
\]
The conductor scheme may be nonreduced and its generic
transverse exponent is unrestricted.
\end{proposition}


\begin{proof}
\noindent\textbf{Step 1.} If the conductor is empty, both conclusions are immediate.
We therefore assume that it is nonempty.
Put $L=\omega_S$ and
$\varepsilon=h^0(S,L)\in\{0,1\}$.
A nonzero section of a torsion line on a compact
connected reduced space trivializes it, by taking a
torsion-trivializing power. Serre duality and
$\chi(\OO_S)=0$ therefore give
$h^1(S,\OO_S)=1+\varepsilon$, with
$\varepsilon=1$ exactly when $L=\OO_S$.
The normalization quotient
$\mathcal Q=\eta_*\OO_T/\OO_S$ has
$H^0(\mathcal Q)\hookrightarrow H^1(S,\OO_S)$,
because $S$ and $T$ have the same constant holomorphic
functions. The actual conductor duality of \ref{cdp:st:A.5.12} is
\[
 \mathcal Q=\omega_\Gamma\otimes L^{-1}|_\Gamma .
\]
Serre duality on the Cohen--Macaulay conductor curve
gives $h^1(\Gamma,L|_\Gamma)\le1+\varepsilon$.
The restriction onto $C$ is a quotient of coherent
sheaves on a one-dimensional space. Its kernel has
no second cohomology, so
$h^1(C,L|_C)\le1+\varepsilon$.
Since $L|_C$ is torsion, it has degree zero on every
irreducible component and
$\chi(C,L|_C)=\chi(C,\OO_C)$.
If $\varepsilon=0$, Riemann--Roch gives
$\chi(\OO_C)=h^0(C,L|_C)-h^1(C,L|_C)\ge-1$.
If $\varepsilon=1$, then $L=\OO_S$ and
$h^0(C,\OO_C)\ge1$, while $h^1(C,\OO_C)\le2$.
This proves \eqref{cdp:eq:C.VII.3}, including all nilpotents of
the original conductor scheme in the comparison.

\noindent\textbf{Step 2.} Now assume that $Z$ is class VII. \ref{cdp:st:A.5.3} gives a
nonzero effective integral divisor $D$ with
$K_Z+D=(\eta\pi)^*\omega_S$.
Each curve in the reduced strict conductor on $Z$
is rational or elliptic by \ref{cdp:st:C.2.VII.1}. Its finite image
in $C$ has normalization of genus at most one, by
Riemann--Hurwitz. Let $n$ be the number of irreducible
components of $C$, let $g_i\in\{0,1\}$ be the genera
of their normalizations, and let $r_x$ and $\delta_x$
be the branch number and delta invariant at a singular
point of $C$. The exact normalization sequence and
the topological normalization map give
\[
 \chi(\OO_C)=\sum_i(1-g_i)-\sum_x\delta_x,\qquad
 e(C)=\sum_i(2-2g_i)-\sum_x(r_x-1).
\tag{C.VII.5}\label{cdp:eq:C.VII.5}
\]
For every reduced curve germ,
$\delta_x\ge r_x-1$: its normalization permits
independent constants on the branches, whereas a
function on the original germ has one common constant.
If $n_0$ is the number of rational-normalization
components, \eqref{cdp:eq:C.VII.3}--\eqref{cdp:eq:C.VII.5} imply
$e(C)\ge n_0-1$.

\noindent\textbf{Step 3.} The only possible negative value is consequently
$e(C)=-1$, with every normalized component elliptic.
Then \eqref{cdp:eq:C.VII.3} says $\sum\delta_x\le1$.
Negativity requires two branches, so there is exactly
one singular point, it has two branches and
$\delta=1$, and the local curve ring is
\[
 \{(a(u),b(v))\in\C\{u\}\oplus\C\{v\}:a(0)=b(0)\}.
\tag{C.VII.6}\label{cdp:eq:C.VII.6}
\]
Indeed that subring already has colength one in the
normalization, and the curve ring is contained in it
with the same colength. Thus the point is an ordinary
node. Its two branches come either from distinct
elliptic normalizations or from two distinct points
of one elliptic normalization.

\noindent\textbf{Step 4.} We show that this node is incompatible with the
specified hypersurface and resolution maps. At least
one normalization point $t\in T$ over the node receives
preimages of both local branches of $C$.
Otherwise partition the irreducible local branches
of the hypersurface germ $S$ according to the one
conductor branch their normalization receives. Both
parts are nonempty. Two hypersurface branches from
opposite parts intersect in a curve: in the regular
three-dimensional local ambient ring, the ideal of
their two distinct equations has height two.
That intersection curve is contained in the conductor
image and therefore is one of its two branches.
Both surface branches must receive it, contrary to
the partition. With only one local surface branch
the assertion is immediate.

\noindent\textbf{Step 5.} Every $\pi$-exceptional curve not contained in
$D_{\mathrm{red}}$ is disjoint from it. To see this,
write $E$ for such a curve. Minimality of the
resolution and adjunction give $K_Z\cdot E\ge0$,
while $D\cdot E\ge0$, and
$(K_Z+D)\cdot E=0$ forces equality. In particular
$D\cdot E=0$. A resolution fibre over a point of
the conductor must meet its strict transform, by
properness. Since its fibre is connected, this
disjointness proves
$\pi^{-1}(t)\subset D_{\mathrm{red}}$ for the
normalization point just found.

\noindent\textbf{Step 6.} The strict curves over the two elliptic-normalized
branches are elliptic: a rational curve cannot map
nonconstantly onto an elliptic curve. They meet the
connected fibre $\pi^{-1}(t)$. If the curves are
distinct, one connected component of $D_{\mathrm{red}}$
therefore contains two elliptic curves. If they are
the same elliptic curve, their meeting points with
the fibre are distinct. This follows from the exact
factorization of a map from a normal curve through
the normalization of its image: one point cannot
map to both points over the node in \eqref{cdp:eq:C.VII.6}.
The connected resolution fibre consequently gives
a path in $D_{\mathrm{red}}$ joining these distinct
points. Both possibilities contradict C.2.VII.1:
the first has two positive-genus components, and
the second adds a graph cycle to an elliptic
component, giving arithmetic genus at least two.
This excludes the only negative possibility and
proves \eqref{cdp:eq:C.VII.4}.
\end{proof}






\begin{proposition}[Proved sign for class VII normalization]\label{cdp:st:C.2.VII.3}
\StatementRecord{cdp:st:C.2.VII.3}{proved}{Historical events: Claude Fable 5.2 (Anthropic), effort high, 16 November 2026: carried as background without fault found. GPT-6 Astra (OpenAI), ChatGPT Pro mode, effort Pro, S6\_09\_PAPERS, 17 November 2026: targeted proof check. Also checked in this writing. Pre-continuation writing events: GPT-6 Astra (OpenAI), ChatGPT Pro mode, effort Pro, S6\_10\_PAPERS, 18 November 2026: targeted proof check; used and rechecked in this run. Continuation metadata events: GPT-6 Astra (OpenAI), ChatGPT Pro mode, effort Pro, S6\_10\_PAPERS continuation, 19 November 2026: used and rechecked in this run; status-vocabulary reconciliation; metadata reconciliation only; no new mathematical check. \emph{Pre-continuation recorded proof scope.} The full-fibre canonical evaluation and rational divisor Riemann-Roch give the precise torsion canonical and coherent Euler inputs. A zero-coefficient minimal-resolution exceptional curve is a disjoint smooth rational minus-two curve; negative definiteness makes each such configuration a transverse tree. The actual resolution replaces those trees by points and restores exactly the omitted positive fibres, giving e(T minus A)=b2-r-a+b. The rational-curve count and the exact conductor Euler bound then give the sign for the original support, with no condition on its conductor multiplicities. The bounded Dloussky and DOT surface-configuration source loci are separately recorded}{the hypotheses, notation and definitions stated here; \cref{cdp:st:A.5.12}; \cref{cdp:st:A.5.3}; \cref{cdp:st:C.2.VII.1}; \cref{cdp:st:C.2.VII.2}}


Let $f:Y\to B$ be a proper connected-fibre holomorphic
map from a smooth compact complex threefold with
$b_2(Y)=0$ to a smooth curve, whose general fibre is a
complex two-torus. Let $mS$ be a full fibre with
irreducible reduced nonnormal support, and suppose
that the minimal resolution $Z$ of its normalization
is class VII. Then
\[
 e(S)\ge0.
\tag{C.VII.7}\label{cdp:eq:C.VII.7}
\]
No restriction on the singularity type or conductor
multiplicity of $S$ is made.
\end{proposition}


\begin{proof}
\noindent\textbf{Step 1.} The canonical evaluation and adjunction comparison
(A.CD.1) give $\omega_S^{\otimes m}=\OO_S$.
The divisor sequence and Hirzebruch--Riemann--Roch
give $\chi(\OO_S)=0$, because $c_1(\OO_Y(S))=0$
rationally. These arguments require only $b_2(Y)=0$;
the integral condition and Picard-variation assumption
in the later part of \ref{cdp:st:A.5.12} are not needed here.
Let $\eta:T\to S$, $\pi:Z\to T$, and
$D>0$ be the maps and divisor of \ref{cdp:st:A.5.3}. Put
$C=\Gamma_{\mathrm{red}}$ and
$A=(\eta^{-1}C)_{\mathrm{red}}\subset T$.
Thus $K_Z+D$ is torsion, and \ref{cdp:st:C.2.VII.1}--\ref{cdp:st:C.2.VII.2}
apply.

\noindent\textbf{Step 2.} We compute $e(T\setminus A)$ through the actual
resolution. As in \ref{cdp:st:C.2.VII.2}, an exceptional curve
$E$ with coefficient zero in $D$ has
$K_Z\cdot E=D\cdot E=0$.
Adjunction and negative self-intersection then give
$p_a(E)=0$ and $E^2=-2$.
These curves are smooth rational and disjoint
from $D_{\mathrm{red}}$.
Their negative-definite intersection matrix makes
each connected configuration a tree: a repeated
edge or a cycle would give a nonnegative square
for a nonzero vector in the matrix. The intersections
are transverse, since their local intersection
multiplicities are one; three curves cannot meet
at a point, since that would give a triangle.
Connectedness of resolution fibres shows that a
fibre either consists entirely of these curves or
lies entirely in $D_{\mathrm{red}}$.
The zero-coefficient fibres lie outside $A$, since
a fibre over $A$ must meet the strict conductor.

\noindent\textbf{Step 3.} Let $a$ be the total number of zero-coefficient
exceptional curves, and let $b$ be the number of
exceptional fibres lying in $D_{\mathrm{red}}$
whose images in $T$ lie outside $A$. On
$Z\setminus D_{\mathrm{red}}$ the former fibres
are still present and the latter are absent.
A tree with $a_j$ rational curves has Euler
number $a_j+1$, so replacing it by its image
point subtracts $a_j$. Adding back the image
point of an absent positive-coefficient fibre
adds one. Away from these finitely many fibres,
$\pi$ is the required isomorphism of complements.
Additivity of compactly supported Euler
characteristic therefore gives the exact equality
\[
 \begin{aligned}
 e(T\setminus A)
   &=e(Z)-e(D_{\mathrm{red}})-a+b\\
   &=b_2(Z)-r(D_{\mathrm{red}})-a+b
   \ \ge0 .
 \end{aligned}
\tag{C.VII.8}\label{cdp:eq:C.VII.8}
\]
The inequality follows because the rational curves
counted by $r(D_{\mathrm{red}})$ and by $a$ are
distinct, and their total is at most $b_2(Z)$ by
\ref{cdp:st:C.2.VII.1}.
Finally $\eta$ is a biholomorphism
$T\setminus A\to S\setminus C$. Hence
\[
 e(S)=e(T\setminus A)+e(C)\ge0
\]
by \ref{cdp:st:C.2.VII.2}. This proves the assertion for the
original reduced support and all full multiplicities.
\end{proof}






\begin{proposition}[Proved equality comparison]\label{cdp:st:C.2.VII.4}
\StatementRecord{cdp:st:C.2.VII.4}{proved}{Historical events: Claude Fable 5.2 (Anthropic), effort high, 16 November 2026: carried as background without fault found. GPT-6 Astra (OpenAI), ChatGPT Pro mode, effort Pro, S6\_09\_PAPERS, 17 November 2026: targeted proof check. Also checked in this writing. Pre-continuation writing events: GPT-6 Astra (OpenAI), ChatGPT Pro mode, effort Pro, S6\_10\_PAPERS, 18 November 2026: targeted proof check; used and rechecked in this run. Continuation metadata events: GPT-6 Astra (OpenAI), ChatGPT Pro mode, effort Pro, S6\_10\_PAPERS continuation, 19 November 2026: used and rechecked in this run; status-vocabulary reconciliation; metadata reconciliation only; no new mathematical check. \emph{Pre-continuation recorded proof scope.} The nonnegative terms in the exact resolution Euler identity give the stated equality conditions. The curve normalization sequence and branch defect bound give precisely the two rational-component alternatives and their delta counts. A nonconstant actual Picard pullback has nonzero differential, sharpening the conductor-duality dimension bound and forcing zero-Euler conductor curves to be disjoint smooth elliptic curves. The established anticanonical class-VII configurations exclude the positive-Betti alternatives; their square-zero rational cycle cannot lie in a negative-definite exceptional matrix. The remaining point-blown-up Hopf case would have an exceptional last minus-one curve, contradicting the minimal resolution. Riemann-Hurwitz supplies the exact unramified conductor maps and normalization gives the converse}{the hypotheses, notation and definitions stated here; \cref{cdp:st:C.2.VII.3}; \cref{cdp:st:C.2.VII.2}; \cref{cdp:st:A.5.12}; \cref{cdp:st:C.2.VII.1}; \cref{cdp:st:C.2.INT.3}}


Retain the irreducible class VII situation of
\ref{cdp:st:C.2.VII.3}. Let $C=\Gamma_{\rm red}$, let $n_0$
be the number of its rational-normalized components,
and put $v(C)=\sum_x(s_x-1)$, where $s_x$ is the
number of branches at a singular point of $C$.
On the minimal resolution $Z$ let $D$ be the
resolution conductor, let $r$ count the rational
curves in $D_{\rm red}$, and let $a$ count the
zero-coefficient exceptional curves lying over
points outside the conductor. Let $b$ count the
positive-coefficient exceptional fibres over
such points, as in \ref{cdp:st:C.2.VII.3}. Then $e(S)=0$
is equivalent to the following explicit
normalization and curve data:
\[
 b=0,\qquad b_2(Z)=r+a,
 \qquad
 \begin{cases}
 n_0=0, v(C)=0,\quad\text{or}\\
 n_0=1, v(C)=2.
 \end{cases}
 \tag{C.VII.9}\label{cdp:eq:C.VII.9}
\]
In the first curve alternative all components
have elliptic normalization and the sum of
their delta invariants is at most one. In the
second the total delta invariant is exactly
two, and every singular point satisfies
$\delta_x=s_x-1$.

If the actual homomorphism
$\Pic^0(Y)\to\Pic^0(Z)$ is nonconstant,
there is the sharper geometric equivalence:
$e(S)=0$ exactly when its normalization is a
smooth minimal Hopf surface and the reduced
conductor image and inverse image are disjoint
unions of smooth elliptic curves, with finite
unramified maps between them.
\end{proposition}


\begin{proof}
\noindent\textbf{Step 1.} The exact comparison already proved in
\ref{cdp:st:C.2.VII.3} is
\[
 e(S)=b_2(Z)-r-a+b+e(C),
 \qquad b_2(Z)-r-a\ge0,\quad b\ge0,\quad e(C)\ge0.
 \tag{C.VII.10}\label{cdp:eq:C.VII.10}
\]
Every normalized component of $C$ has genus
zero or one. Its topological normalization
formula is $e(C)=2n_0-v(C)$. The stronger
bound in \ref{cdp:st:C.2.VII.2} is $e(C)\ge n_0-1$.
Thus $e(C)=0$ permits precisely $n_0=0,1$,
with $v(C)=0,2$, respectively. This proves
both directions of \eqref{cdp:eq:C.VII.9}, using actual
branch numbers rather than an assumption
about a local singularity type.
The same proof gives
$\chi(\OO_C)=n_0-\sum\delta_x\ge-1$.
For $n_0=0$ this bounds the total delta
invariant by one. For $n_0=1$, the value
$v(C)=2$ and $\delta_x\ge s_x-1$ force
$\sum\delta_x=2$ and equality at every point.
These are exact consequences for fibres which
satisfy the criterion; no separate claim
that all such abstract curve patterns admit
a torus smoothing is used.

\noindent\textbf{Step 2.} Now suppose that the Picard pullback is
nonconstant. The argument of \ref{cdp:st:A.5.12} applies
to $\Pic^0(Y)$ even with only rational
$H^2(Y,\Q)=0$. Namely its nonzero differential
is the actual composite
\[
 H^1(Y,\OO_Y)\longrightarrow H^1(S,\OO_S)
 \longrightarrow H^1(T,\OO_T)
 \longrightarrow H^1(Z,\OO_Z).
 \tag{C.VII.11}\label{cdp:eq:C.VII.11}
\]
The middle arrow consequently has positive
rank. Put $\varepsilon=h^0(S,\omega_S)$;
the torsion canonical relation makes
$\varepsilon\in\{0,1\}$ and
$h^1(S,\OO_S)=1+\varepsilon$.
The normalization sequence then gives
$h^0(\mathcal Q)\le\varepsilon$.
Conductor duality and restriction onto $C$
therefore give
\[
 h^1(C,\omega_S|_C)\le\varepsilon.
 \tag{C.VII.12}\label{cdp:eq:C.VII.12}
\]
If $\varepsilon=0$, degree-zero
Riemann--Roch gives $\chi(\OO_C)\ge0$.
If $\varepsilon=1$, the torsion line
$\omega_S$ is trivial, and
$h^0(C,\OO_C)\ge1$ gives the same inequality.
The exact formula (C.INT.26) now implies
that $e(C)=0$ holds precisely when $C$ is
a disjoint union of smooth elliptic curves.

\noindent\textbf{Step 3.} Assume $e(S)=0$. Every rational curve in
$D_{\rm red}$ must be resolution-exceptional,
because its image lies in $C$ and cannot
dominate an elliptic curve. A minimal class
VII model with positive second Betti number
and a torsion anticanonical divisor is one
of the proved configurations used in
\ref{cdp:st:C.2.VII.1}. An intermediate or
Inoue--Hirzebruch model has no elliptic
curves, and therefore could not dominate
the nonempty elliptic conductor. A
parabolic Inoue model has a rational cycle
of square zero and exactly $b_2$ rational
curves. Its point blowups therefore still
have exactly $b_2(Z)$ rational curves.
The equality $b_2(Z)=r+a$ makes every one
of them resolution-exceptional, including
any zero-coefficient blowup curves. Thus
the full total transform of the minimal
rational cycle, whose square remains zero,
would be exceptional. This contradicts
negative definiteness of a resolution's
exceptional intersection matrix.
Consequently the minimal model is Hopf.
All rational curves on $Z$ are now the
exceptional curves of its point blowups,
and their number is $b_2(Z)$.
The equality $b_2(Z)=r+a$ makes all of
them resolution-exceptional. As in the
last paragraph of \ref{cdp:st:C.2.INT.3}, a nonempty
blowup sequence would have a last
$(-1)$ curve, contradicting minimality
of the resolution. Thus the normalization
is already a smooth minimal Hopf surface.
Its reduced conductor curves are elliptic;
their finite maps to the smooth elliptic
curves of $C$ are unramified by
Riemann--Hurwitz.
Conversely those normalization and
conductor data have Euler number zero
by the exact normalization formula.
This proves the sharper equivalence.
\end{proof}






\begin{proposition}[Euler signs for two of the rational boundary classes]\label{cdp:st:C.2.RAT.1}
\StatementRecord{cdp:st:C.2.RAT.1}{proved}{Historical events: Claude Fable 5.2 (Anthropic), effort high, 16 November 2026: carried as background without fault found. GPT-6 Astra (OpenAI), ChatGPT Pro mode, effort Pro, S6\_09\_PAPERS, 17 November 2026: targeted proof check. Also checked in this writing. Pre-continuation writing events: GPT-6 Astra (OpenAI), ChatGPT Pro mode, effort Pro, S6\_10\_PAPERS, 18 November 2026: targeted proof check. Continuation metadata events: GPT-6 Astra (OpenAI), ChatGPT Pro mode, effort Pro, S6\_10\_PAPERS continuation, 19 November 2026: used and rechecked in this run; status-vocabulary reconciliation; metadata reconciliation only; no new mathematical check. \emph{Pre-continuation recorded proof scope.} The complete delta-at-most-two curve list, the two non-lci symbolic-square ideals, conormal determinant defects, uncancellable Hessian poles and radical cubic bounds were re-derived. A fresh execution of code/replay\_low\_delta\_cusp.py passed every exact saturation and jet identity. The lci normal-jet section has degree 12-4*delta; the actual Hirzebruch coordinate germ excludes a simple Hessian zero. Its zero budget and the connected positive-boundary graph yield all stated r and Euler bounds}{the hypotheses, notation and definitions stated here; \cref{cdp:st:A.5.16}; \cref{cdp:eq:C.RA.7}; \cref{cdp:eq:C.RA.8}; \cref{cdp:eq:C.RA.9}; \cref{cdp:eq:C.RA.10}; \cref{cdp:st:A.5.5}; \cref{cdp:st:A.5.15}; \cref{cdp:st:A.5.17}; \cref{cdp:eq:C.RA.14}; \cref{cdp:eq:C.RA.15}; \cref{cdp:eq:C.RA.6}; \cref{cdp:eq:C.RA.13}; \cref{cdp:eq:C.RA.2}}

Let $f:Y\to B$ be a proper holomorphic map from a smooth compact
complex threefold with $b_2(Y)=0$, with connected fibres and smooth
general fibre a complex two-torus. Let $mS$ be a full fibre with
irreducible nonnormal reduced support. Write
$h=\eta\pi:Z\to T\to S$ for the minimal resolution of its
normalization and suppose that $Z$ is rational. Use the effective
divisor $K_Z+D=0$ of \ref{cdp:st:A.5.16}.

If a contraction to $\mathbf P^2$ sends $D$ to three concurrent
reduced lines, then $e(S)\ge1$. More precisely,
$e(S)=-1+N_0+e(\Gamma_{\rm red})\ge1+N_0$, with $N_0$ as in
(A.RT.12).

Consider also a specified morphism $\rho:Z\to\mathbb F_n$,
$n\ge1$, which is a sequence of point blowups in reverse and has
\[
 \rho_*D=2C_0+\sum_{j=1}^r a_jF_j,\qquad
 C_0^2=-n,\qquad
 a_j\ge1,\quad \sum_j a_j=n+2,                  \tag{C.RA.1}\label{cdp:eq:C.RA.1}
\]
where the $F_j$ are distinct ruling fibres. Let $C$ be the
strict transform of $C_0$. Suppose $h(C)$ is a curve and the
generic transverse singularity there is a cusp $A_2$.
Then $r\le3$ and $e(S)\ge1$. If $h(C)$ is smooth, then
$r\le2$. These assertions allow singular normalization,
arbitrary vertical conductor multiplicities and a nonreduced
conductor-image scheme.
\end{proposition}


\begin{proof}
\noindent\textbf{Step 1.} Put $G=\Gamma_{\rm red}$. \ref{cdp:st:A.5.16} has already proved
\[
 e(S)=e(Z\setminus D_{\rm red})+e(G).             \tag{C.RA.2}\label{cdp:eq:C.RA.2}
\]
In either boundary class under consideration, $D_{\rm red}$
is simply connected, all its components are rational, and
(A.RT.18)--(A.RT.20) give
\[
 \pi_1(S)=\pi_1(G)\in\{1,\mathbb Z\}.
 \tag{C.RA.3}\label{cdp:eq:C.RA.3}
\]
Write $n_G$ for the number of irreducible components of $G$
and $\beta=b_1(G)\le1$. The normalization maps from their
rational normalizations are homeomorphisms except for the
identification of finitely many points. Consequently
\[
                   e(G)=n_G+1-\beta\ge n_G\ge1.
 \tag{C.RA.4}\label{cdp:eq:C.RA.4}
\]
The exact two-branch incidence argument (A.RT.21) strengthens
this to $e(G)\ge2$ in the present boundary classes.
For three concurrent lines, (A.RT.12)--(A.RT.13) therefore
prove the first assertion.

\noindent\textbf{Step 2.} We prove the cusp assertion. Put $H=h(C)_{\rm red}$.
The transverse $A_2$ curve has a single normalized branch.
Thus $C\to H$ is the normalization map of degree one, and
no other strict conductor component maps onto $H$.
In particular $H^\nu=C=\mathbf P^1$. The restriction
$H^1(G,\OO_G)\to H^1(H,\OO_H)$ is surjective, since its
kernel has one-dimensional support and no second cohomology.
\ref{cdp:st:A.5.16} therefore gives
\[
              \delta:=p_a(H)\le2.              \tag{C.RA.5}\label{cdp:eq:C.RA.5}
\]
Since $H^\nu$ is rational, $\delta$ is the sum of its local
delta invariants. Its topological incidence graph is a
subgraph of the incidence graph of $G$, so $b_1(H)\le1$.
If $s$ denotes the number of points of $C$ above singular
points of $H$, normalization of the reduced curve gives
\[
 s=\sum_{p\in\operatorname{Sing}H}r_p
   =b_1(H)+\#\operatorname{Sing}H\le1+\delta
 \quad(\delta>0),\qquad s=0\quad(\delta=0).       \tag{C.RA.6}\label{cdp:eq:C.RA.6}
\]

\noindent\textbf{Step 3.} We need the precise low-delta curve analysis, including its
two possible non-lci cases. A reduced curve germ of delta
invariant at most two has at most three branches.
Three branches require delta two, and the ring is exactly
the subring of its normalization consisting of triples
with equal constant value. This follows by equality of
colengths, both equal to two. Its topological branch
identification alone contributes two loops to a globally
irreducible curve, so $b_1(H)\le1$ excludes it here.

\noindent\textbf{Step 4.} For two branches, the exact sequence for their union gives
$\delta=\delta_1+\delta_2+\ell(\OO_{H_1\cap H_2})$.
Two smooth branches therefore give a node, or, at delta
two, the fibre product matching their constant and linear
terms. The latter has the plane tacnode ring
$\mathbb C\{x,y\}/(y^2-x^4)$.
The only remaining two-branch possibility is a cusp and
a transverse smooth branch, whose ring is
$\mathbb C\{t^2,t^3\}\times_{\mathbb C}\mathbb C\{u\}$.
In embedding dimension three its ideal can be written
\[
                  I=(xz,yz,y^2-x^3).            \tag{C.RA.7}\label{cdp:eq:C.RA.7}
\]
For one branch, the delta invariant equals the number of
gaps of its valuation semigroup. With at most two gaps,
the singular possibilities are
$\langle2,3\rangle$, $\langle2,5\rangle$, and
$\langle3,4,5\rangle$. The first two are plane cusps.
For the last, the conductor is $t^3\mathbb C\{t\}$,
and colength two makes the ring
$\mathbb C+t^3\mathbb C\{t\}$ exactly. Its ideal is
\[
 I=(g_1,g_2,g_3)
   =(y^2-xz,\ yz-x^3,\ z^2-x^2y),\qquad
 (x,y,z)=(t^3,t^4,t^5).                         \tag{C.RA.8}\label{cdp:eq:C.RA.8}
\]
For completeness, the asserted plane form at semigroup
$\langle2,5\rangle$ follows by making an order-two
generator equal to $t^2$ by a change of parameter.
Its conductor is then $t^4\mathbb C\{t\}$, and the ring
is $\mathbb C+\mathbb Ct^2+t^4\mathbb C\{t\}$.
These conductor descriptions also follow in the completed
local rings by eliminating successive higher terms; finite
normalization and faithful flatness of completion give the
same conductor in the convergent rings. Thus there are no
unlisted cases. A germ of embedding dimension two is lci
in the smooth threefold: choose one regular transverse
ambient coordinate in its ideal and then its single plane
equation. We now exclude \eqref{cdp:eq:C.RA.7}--\eqref{cdp:eq:C.RA.8} for our cusp
image.

\noindent\textbf{Step 5.} For \eqref{cdp:eq:C.RA.8}, put
\[
 F=x^5+xy^3+z^3-3x^2yz .
\]
The symbolic square is $I^{(2)}=I^2+(F)$, with
\[
 xF=g_2^2-g_1g_3,\quad
 yF=g_2g_3+xg_1^2,\quad
 zF=g_3^2+xg_1g_2.                              \tag{C.RA.9}\label{cdp:eq:C.RA.9}
\]
The saturation computation $I^2:x^\infty=I^2+(F)$,
including the complete exact Groebner certificate, is
given in \path{code/A_low_delta_cusp.py} and its output.
It computes the symbolic square because the curve is lci
away from the origin and $x$ vanishes on it only at the
origin. The same equality in the analytic local ring
follows by flatness of the convergent extension and the
displayed finite saturation.

\noindent\textbf{Step 6.} Let $N$ be the dual of the torsion-free pullback to $C$ of
the ambient conormal module of $H$. At \eqref{cdp:eq:C.RA.8} the
normalized conormal basis and its dual lifts are
\[
\begin{gathered}
 a=dg_1,\qquad b=dg_2+t\,dg_1,\qquad
 dg_3=-2t^2a+tb,\\
 A=-t^{-3}\partial_z,\qquad
 B=(t^{-5}\partial_y+2t^{-4}\partial_z)/3 .
\end{gathered}                                  \tag{C.RA.10}\label{cdp:eq:C.RA.10}
\]
Their wedge is $a\wedge b=t^6\,\iota_{\gamma'}(dx\wedge
dy\wedge dz)$, where $\gamma=(t^3,t^4,t^5)$.
Thus the determinant formula is
$\det N=K_C\otimes h_C^*K_Y^{-1}(6p)$.
Since every ambient line has degree zero on every compact
curve when $b_2(Y)=0$, its degree here is four.

\noindent\textbf{Step 7.} Write the local equation of $S$ as $cF+f_2$, $f_2\in I^2$.
Its normalized quadratic form is the sum of a regular
quadratic form and $c$ times
\[
 q_F=\frac3t a^2-\frac3{t^2}ab+\frac1{t^3}b^2,
 \qquad \det(q_F)=\frac3{4t^4}.                  \tag{C.RA.11}\label{cdp:eq:C.RA.11}
\]
If $c(0)\ne0$, the determinant has an uncancellable
pole of order four: a cross term with the regular form
has pole at most three. The generic cusp requires rank
one, so $c(0)=0$. Equations \eqref{cdp:eq:C.RA.9} then give
$cF\in I^2$, and the quadratic form is regular.

\noindent\textbf{Step 8.} The intrinsic cubic on its radical has pole order at most
ten at $p$. Here is a direct bound for arbitrary coefficients,
not only for the polynomial $F$. In the lifts \eqref{cdp:eq:C.RA.10},
the differentials $dg_i$ are regular, the second derivatives
$D^2g_i(Au+Bv)^2$ have poles at most ten, and a derivative
of a holomorphic coefficient has pole at most five.
The formula for $D^3(b_{ij}g_i g_j)$ therefore has pole
at most ten. Terms containing an undifferentiated $g_i$
vanish on the curve. The exact derivatives, with equality
in the pole-ten bound, are in the cited computation.
The radical cubic is independent of a tangent component
of the lift: differentiating the identically zero quadratic
form on its radical along the smooth part of the curve
gives this invariance, and equality extends meromorphically.

\noindent\textbf{Step 9.} The same argument is shorter for \eqref{cdp:eq:C.RA.7}. Put
$g_1=xz$, $g_2=yz$, $g_3=y^2-x^3$. Then
$I^{(2)}=I^2+(z g_3)$ and
\[
 x(zg_3)=g_1g_3,\quad y(zg_3)=g_2g_3,\quad
 z(zg_3)=g_2^2-xg_1^2.
 \tag{C.RA.12}\label{cdp:eq:C.RA.12}
\]
On the cusp branch $(t^2,t^3,0)$ use $a=dg_1$, $b=dg_3$.
The quadratic form of $z g_3$ is $t^{-2}ab$, whose
determinant has pole four and cannot be cancelled by a
regular quadratic form. Thus its coefficient again
vanishes at the origin and the surface equation belongs
to $I^2$. The determinant corrections are four on the
cusp branch and two on the smooth branch $(0,0,t)$.
The total determinant degree is again four.
Dual lifts on the cusp branch are
$t^{-2}\partial_z$ and $\frac12t^{-3}\partial_y$;
on the smooth branch they are
$t^{-1}\partial_x,t^{-1}\partial_y$.
The same product derivative formula bounds the cubic
poles by six and two respectively, of total degree eight.

\noindent\textbf{Step 10.} In either non-lci case let $R\subset N$ be the saturated
radical of the regular rank-one quadratic form and let
$Q=N/R$. Its nonzero square coefficient is a section of
$Q^{-2}\otimes M$, where $M=\OO_Y(S)|_C$ has degree
zero. Hence $\deg Q\le0$ and $\deg R=4-\deg Q\ge4$.
The nonzero intrinsic cubic of a generic $A_2$ is a
section of $R^{-3}\otimes M$ with poles of total degree
at most ten, or eight. It would give respectively
$3\deg R\le10$ or $3\deg R\le8$, both impossible.
This excludes the two non-lci cases, without assuming
that the conductor-image scheme is reduced.

\noindent\textbf{Step 11.} We may consequently use the lci conormal bundle along $H$.
Let $\Xi$ be its normalization conductor divisor, of degree
$2\delta$. The determinant formula and the
conductor-cleared intrinsic cubic from \ref{cdp:st:A.5.5} give
\[
\begin{gathered}
 \det N=K_C(\Xi)\otimes h_C^*K_Y^{-1},\\
 q\in H^0(C,Q^{-2}\otimes M),\qquad
 k\in H^0(C,R^{-3}\otimes M(2\Xi)),\\
 0\ne q^3k^2\in
 H^0\bigl(C,(\det N)^{-6}\otimes M^5(4\Xi)\bigr),\\
 \deg\bigl((\det N)^{-6}\otimes M^5(4\Xi)\bigr)
                         =12-4\delta .
\end{gathered}                                  \tag{C.RA.13}\label{cdp:eq:C.RA.13}
\]
The section is nonzero because the generic transverse
singularity is a cusp.

\noindent\textbf{Step 12.} We give the local comparison that strengthens the zero
count at a smooth point of $H$. A simple zero of $q$ at
which $k$ is nonzero is impossible in a torus smoothing.
Choose coordinates with $H=\{x=y=0\}$, radical direction
$\partial_x$, and coordinate $z$ along $H$.
The nonzero cubic permits Weierstrass preparation in $x$:
\[
 f_S=U\bigl(x^3+a_2(y,z)x^2+a_1(y,z)x+a_0(y,z)\bigr).
 \tag{C.RA.14}\label{cdp:eq:C.RA.14}
\]
The critical-curve equations and the rank-one Hessian
along it give $a_2\in(y)$ and $a_1,a_0\in(y^2)$.
After depressing the cube by $x'=x+a_2/3$, this is
\[
 f_S=U'\bigl(x'^3+y^2(B_1(y,z)x'+B_0(y,z))\bigr).
 \tag{C.RA.15}\label{cdp:eq:C.RA.15}
\]
The simple zero of the Hessian coefficient says
$\partial_zB_0(0,0)\ne0$. Thus
$z'=B_1(y,z)x'+B_0(y,z)$ is an actual holomorphic
coordinate with $x',y$, and the reduced equation is a
unit times $x'^3+y^2z'$. The full fibre equation is a
unit times its $m$th power. \ref{cdp:st:A.5.15} excludes precisely
this Hirzebruch germ, including that unit and multiplicity.

\noindent\textbf{Step 13.} Every point of $C$ above the smooth part of $H$ where a
vertical conductor component meets it has $q=0$.
Indeed a nonzero Hessian would, by \ref{cdp:st:A.5.17}, give a constant
transverse cusp product with smooth normalization and pure
conductor $2C$; its minimal normalization resolution could
have no such vertical component. At a zero of $q$ the
cubic is holomorphic. If it also vanishes, the section
$q^3k^2$ has zero order at least five. If it does not
vanish, \eqref{cdp:eq:C.RA.14}--\eqref{cdp:eq:C.RA.15} force the zero of $q$ to
have order at least two, and the order is at least six.

\noindent\textbf{Step 14.} Each of the $r$ original ruling fibres in \eqref{cdp:eq:C.RA.1} has
a positive-boundary total-transform component meeting $C$,
at a different point of $C$. This remains true at each
blowup: a blowup at its intersection with the coefficient-two
section inserts an exceptional curve of positive coefficient.
Using \eqref{cdp:eq:C.RA.6} and the degree in \eqref{cdp:eq:C.RA.13} now gives
\[
 \begin{array}{c|c|c|c}
 \delta& s\text{ at most}&
 \text{bad points above }H_{\rm reg}\text{ at most}&r\text{ at most}\\ \hline
 0&0&2&2\\
 1&2&1&3\\
 2&3&0&3 .
 \end{array}                                    \tag{C.RA.16}\label{cdp:eq:C.RA.16}
\]
This proves the asserted bound on $r$.

\noindent\textbf{Step 15.} Finally (A.RT.12) and (A.RT.16) give
$e(Z\setminus D_{\rm red})=2-r+N_0$.
The forest argument of (A.RT.21) applies to this entire
star class, whether its minimal complement is negative
or not, and gives $e(G)\ge2$.
Together with $r\le3$, \eqref{cdp:eq:C.RA.2} now gives directly
\[
                 e(S)\ge4-r+N_0\ge1+N_0.
 \tag{C.RA.17}\label{cdp:eq:C.RA.17}
\]
\end{proof}



The two classes in \ref{cdp:st:C.2.RAT.1} are proved without a restriction
on nilpotent conductor structure. The other horizontal
coefficient-two types and contraction of that horizontal
section are not included in this statement.




\begin{proposition}[A rational double cover of a smooth tacnode image]\label{cdp:st:C.2.RAT.2}
\StatementRecord{cdp:st:C.2.RAT.2}{proved}{Historical events: Claude Fable 5.2 (Anthropic), effort high, 16 November 2026: carried as background without fault found. GPT-6 Astra (OpenAI), ChatGPT Pro mode, effort Pro, S6\_09\_PAPERS, 17 November 2026: targeted proof check. Also checked in this writing. Pre-continuation writing events: GPT-6 Astra (OpenAI), ChatGPT Pro mode, effort Pro, S6\_10\_PAPERS, 18 November 2026: targeted proof check. Continuation metadata events: GPT-6 Astra (OpenAI), ChatGPT Pro mode, effort Pro, S6\_10\_PAPERS continuation, 19 November 2026: used and rechecked in this run; status-vocabulary reconciliation; metadata reconciliation only; no new mathematical check. \emph{Pre-continuation recorded proof scope.} The actual quadratic-root conductor extension gives the two ramification points. The Hessian and quartic line degrees, together with the simple-quartic-zero exclusion, force precisely two simple Hessian zeros and two quartic zeros of order three. Each ramified fibre has one conductor preimage, so the distinct positive-boundary attachments give r<=2 and the stated Euler inequality}{the hypotheses, notation and definitions stated here; \cref{cdp:st:C.2.RAT.1}; \cref{cdp:st:A.5.17}; \cref{cdp:eq:C.RA.19}}

Retain the torus-fibre, rational-normalization and specified
ruled-boundary morphism of \ref{cdp:st:C.2.RAT.1}. Suppose that $C$ is not
contracted, that $H=h(C)_{\rm red}$ is a smooth rational curve,
that the generic transversal along $H$ is an ordinary tacnode,
and that $C\to H$ has degree two. Then $r\le2$ and
\[
                         e(S)\ge2+N_0.
\tag{C.RA.18}\label{cdp:eq:C.RA.18}
\]
The assertion allows arbitrary vertical conductor and singular
normalization away from the generic points of $C$.
\end{proposition}


\begin{proof}
\noindent\textbf{Step 1.} The two branches of a generic tacnode are both supplied by
$C\to H$; no other strict conductor component can dominate $H$.
Apply the exact normal-jet construction (A.T.5)--(A.T.7) on
$H\simeq\PP^1$. Here the curve conductor divisor $\Xi$ is zero.
The quotient $Q$ of its rank-one Hessian has degree $q_0\le0$.
The quadratic coefficient $a$ and the quartic invariant
$\Delta=ae-b^2/4$ are nonzero sections with
\[
 A:=\deg\operatorname{div}(a)=-2q_0,\qquad
 \deg\operatorname{div}(\Delta)=8+2q_0=8-A.
\tag{C.RA.19}\label{cdp:eq:C.RA.19}
\]
In particular $A$ is a nonnegative even integer. These are
sections of the actual lines $M\otimes Q^{-2}$ and
$(M\otimes Q^{-1}\otimes R^{-2})^2$, with
$\deg M=0$ and $\deg R=-2-q_0$.

\noindent\textbf{Step 2.} The normalized conductor cover is obtained by adjoining
$\sqrt{-\Delta}$, by the quadratic equation for the ratio of
the first normal coordinate to the square of the radical
coordinate. Thus its branch divisor consists exactly of the
odd-order zeros of $\Delta$, as proved in (A.T.7). Since
$C\simeq\PP^1$ and the map has degree two, Riemann--Hurwitz
gives two distinct branch points on $H$. At each of them
$a$ vanishes: a nonzero Hessian would give the constant
transverse $A_3$ product by \ref{cdp:st:A.5.17} and hence an unramified
conductor cover. The coordinate comparison
(A.T.10)--(A.T.12) excludes a simple zero of $\Delta$ at
such a point. Its odd zero order is consequently at least
three. Equations \eqref{cdp:eq:C.RA.19} now give
\[
                        A\ge2,\qquad 8-A\ge6.
\tag{C.RA.20}\label{cdp:eq:C.RA.20}
\]
Therefore $A=2$. The two branch points are exactly the
zeros of $a$, each simple, and the zeros of $\Delta$ have
order three there, with no further zeros.

\noindent\textbf{Step 3.} Every original fibre $F_j$ of (C.RA.1) gives a positive
vertical boundary component meeting $C$. The positivity
is preserved under blowup at its intersection with the
coefficient-two horizontal section, and these $r$
intersection points on $C$ are distinct. At any such point
the Hessian coefficient at its image must vanish. Otherwise
\ref{cdp:st:A.5.17} supplies the constant $A_3$ product, whose two
normalization branches are smooth and have no vertical
conductor; the minimal resolution $\pi$ is an isomorphism
near them. Each of the two zeros of $a$ is a branch point
of $C\to H$, and consequently has just one point above it
on $C$. It follows that $r\le2$.

\noindent\textbf{Step 4.} Finally the exact complement calculation and (A.RT.21) give
\[
 e(S)=2-r+N_0+e(\Gamma_{\rm red})
                  \ge4-r+N_0\ge2+N_0,
\]
which proves \eqref{cdp:eq:C.RA.18}.
\end{proof}






\begin{proposition}[Every uncontracted degree-one horizontal image has nonnegative Euler number]\label{cdp:st:C.2.RAT.3}
\StatementRecord{cdp:st:C.2.RAT.3}{proved}{Historical events: Claude Fable 5.2 (Anthropic), effort high, 16 November 2026: carried as background without fault found. GPT-6 Astra (OpenAI), ChatGPT Pro mode, effort Pro, S6\_09\_PAPERS, 17 November 2026: targeted proof check. Also checked in this writing. Pre-continuation writing events: GPT-6 Astra (OpenAI), ChatGPT Pro mode, effort Pro, S6\_10\_PAPERS, 18 November 2026: targeted proof check. Continuation metadata events: GPT-6 Astra (OpenAI), ChatGPT Pro mode, effort Pro, S6\_10\_PAPERS continuation, 19 November 2026: used and rechecked in this run; status-vocabulary reconciliation; metadata reconciliation only; no new mathematical check. \emph{Pre-continuation recorded proof scope.} The residue comparison gives conductor exponent equal to the intrinsic branch exponent plus all pairwise intersection multiplicities, bounding dominating vertical branches by two. A complete positive ruling fibre cannot lie in a negative-definite exceptional matrix, so each lost boundary requires its own coefficient-zero blowup. The surviving actual intersections give distinct normalization-incidence edges; their exact graph Betti count proves the Euler bound}{the hypotheses, notation and definitions stated here; \cref{cdp:st:C.2.RAT.1}; \cref{cdp:eq:C.RA.22}; \cref{cdp:eq:C.RA.23}; \cref{cdp:eq:C.RA.24}}

Retain the torus-fibre, rational-normalization and ruled-boundary
morphism of \ref{cdp:st:C.2.RAT.1}. Suppose that $C$ is not contracted and
that its actual map to the normalization of
$H=h(C)_{\rm red}$ has degree one. Let $r_{\rm dom}$ be the
number of original ruling fibres $F_j$ whose positive vertical
boundary contains a strict conductor component mapping onto
$H$. Then
\[
       0\le r_{\rm dom}\le2,\qquad
       e(S)\ge2-r_{\rm dom}\ge0.
\tag{C.RA.21}\label{cdp:eq:C.RA.21}
\]
There is no assumption on the generic transverse type, on the
local complete intersection property of $H$, or on the
nilpotent structure of the conductor image.
\end{proposition}


\begin{proof}
\noindent\textbf{Step 1.} The conductor coefficient on $C$ is two, because $C$ is
the strict transform of the coefficient-two section in
(C.RA.1) and is not $\pi$-exceptional. At a general point
of $H$, its degree-one map supplies one branch of the reduced
transverse plane curve. If that branch has intrinsic
conductor exponent $c_i^{\rm int}$, its exponent in the
whole plane curve is
\[
             2=c_i^{\rm int}+\sum_{j\ne i}I(C_i,C_j).
\tag{C.RA.22}\label{cdp:eq:C.RA.22}
\]
Here $I(C_i,C_j)$ is the positive intersection multiplicity
of two different plane branches. To verify the displayed
comparison, write the reduced equation as $u\prod f_j$.
On the normalization of the $i$th branch the residue
description of its dualizing generator gains the factor
$\prod_{j\ne i}f_j|_{C_i}$; its valuation is the sum in
\eqref{cdp:eq:C.RA.22}. The remaining factor is exactly the conductor
of the individual branch. This gives the actual inclusion
of its conductor ideal into its normalization, with the
stated exponent. Thus there are at most two other generic
branches. Each vertical strict component dominating $H$
supplies at least one such branch, counted with its covering
degree. This proves $r_{\rm dom}\le2$.

\noindent\textbf{Step 2.} For an original fibre $F_j$, let $B_j$ be the positive
vertical boundary in its total transform. Its reduced
support is connected and meets $C$ at a single point.
These assertions follow at each point blowup: a centre
at the intersection of positive components inserts a
positive exceptional component; a centre at a smooth
coefficient-one point creates only a coefficient-zero
leaf; and a centre at a smooth point of coefficient
at least two creates a positive leaf. These are all the
possibilities in the exact divisor formula
$D_{\rm new}=\sigma^*D_{\rm old}-E$.
The $r$ points of $C$ so obtained are distinct.

\noindent\textbf{Step 3.} If every component of $B_j$ is $\pi$-exceptional, there
has been at least one multiplicity-one blowup centre
above $F_j$. Indeed, in the absence of such a centre,
every component of the entire total transform of $F_j$
has positive conductor coefficient. That complete
ruling fibre, with its positive multiplicities, has
self-intersection zero. It cannot be supported on the
negative-definite exceptional matrix of a normal surface
resolution. Distinct original fibres require distinct
blowup events. Hence at most $N_0$ of the $B_j$ can be
entirely $\pi$-exceptional.

\noindent\textbf{Step 4.} Discard those fibres and the $r_{\rm dom}$ fibres counted
above. There remain at least
\[
                        k\ge r-N_0-r_{\rm dom}
\tag{C.RA.23}\label{cdp:eq:C.RA.23}
\]
fibres with a surviving strict conductor component and
none dominating $H$. Follow the connected positive
boundary from its point on $C$ to the first component
which is not $\pi$-exceptional. Its image is a curve
different from $H$ meeting $H$ at that point's image.
Every intervening connected chain is contracted by
$\pi$, so this is an actual intersection in $S$, not
an identification imposed on an auxiliary graph.

\noindent\textbf{Step 5.} Use the normalization-incidence graph of
$G=\Gamma_{\rm red}$: its vertices are the rational
normalization components and the singular points with
more than one normalization branch, and its edges are
the individual branch preimages. The preceding $k$
intersections give $k$ distinct edges adjacent to the
vertex belonging to $H$. They remain distinct even
when two of their images are the same singular point:
$C\to H^\nu$ is an isomorphism, and the points on $C$
were distinct. Remove that vertex and all its adjacent
edges. Each of the $c$ remaining connected components
which receives one of these edges contains a
normalization vertex of a curve other than $H$.
In particular $c\le n_G-1$, where $n_G$ is the number
of irreducible components of $G$.

\noindent\textbf{Step 6.} Contract a spanning tree in each of those $c$ graph
components. The $k$ retained edges then give a connected
multistar with $c+1$ vertices and $k$ edges. Its first
Betti number is $k-c$, so the first Betti number
$\beta$ of the full incidence graph satisfies
\[
             k-c\le\beta\le1,
          \qquad n_G\ge k+1-\beta\ge k.
\tag{C.RA.24}\label{cdp:eq:C.RA.24}
\]
The upper bound on $\beta$ is the actual local torus-group
quotient (A.RT.18)--(A.RT.20). Adding genus-zero
normalization components changes no graph first Betti
number, and the normalization quotient gives
$e(G)=n_G+1-\beta\ge n_G$. Combining \eqref{cdp:eq:C.RA.23}--\eqref{cdp:eq:C.RA.24}
with the exact surface Euler formula yields
\[
\begin{split}
 e(S)&=2-r+N_0+e(G)\\
     &\ge2-r+N_0+n_G
      \ge2-r+N_0+(r-N_0-r_{\rm dom})
       =2-r_{\rm dom}.
\end{split}
\tag{C.RA.25}\label{cdp:eq:C.RA.25}
\]
This proves the assertion.
\end{proof}



The degree-one result accounts for every generic transverse
type allowed by conductor exponent two on that branch. It
does not replace the other branches by a presumed ordinary
triple: those branches may have larger conductor exponents.
The cases in which the horizontal section is contracted, or
maps with higher degree to a singular conductor image, require
the further arguments identified separately below.




\begin{proposition}[Every rational degree-two tacnode cover has positive Euler number]\label{cdp:st:C.2.RAT.4}
\StatementRecord{cdp:st:C.2.RAT.4}{proved}{Historical events: Claude Fable 5.2 (Anthropic), effort high, 16 November 2026: carried as background without fault found. GPT-6 Astra (OpenAI), ChatGPT Pro mode, effort Pro, S6\_09\_PAPERS, 17 November 2026: targeted proof check. Also checked in this writing. Pre-continuation writing events: GPT-6 Astra (OpenAI), ChatGPT Pro mode, effort Pro, S6\_10\_PAPERS, 18 November 2026: targeted proof check. Continuation metadata events: GPT-6 Astra (OpenAI), ChatGPT Pro mode, effort Pro, S6\_10\_PAPERS continuation, 19 November 2026: used and rechecked in this run; status-vocabulary reconciliation; metadata reconciliation only; no new mathematical check. \emph{Pre-continuation recorded proof scope.} The prior non-lci jet exclusions leave the actual lci conductor lines. A unit Hessian entry would remove the conductor poles and exceed the quartic degree budget, forcing its vanishing at every singular preimage. The exact x=y\textasciicircum{}2*lambda factorization shows that a unit quartic leaves one branch with pure conductor. Counting ramification, conductor multiplicities and duplicate attachments gives at most one lost incidence edge and hence positive Euler number}{the hypotheses, notation and definitions stated here; \cref{cdp:st:C.2.RAT.1}; \cref{cdp:st:A.5.18}; \cref{cdp:eq:C.RA.26}; \cref{cdp:st:A.5.17}; \cref{cdp:eq:C.RA.28}; \cref{cdp:eq:C.RA.30}; \cref{cdp:eq:C.RA.31}; \cref{cdp:st:C.2.RAT.3}}

Retain the torus-fibre, rational-normalization and specified
ruled-boundary morphism of \ref{cdp:st:C.2.RAT.1}. Suppose $C$ is not
contracted, $H=h(C)_{\rm red}$ has generic transverse type
$A_3$, and the actual map $C\to H^\nu$ has degree two.
Then $e(S)\ge1$. No smoothness or local complete intersection
assumption on the image curve is required.
\end{proposition}


\begin{proof}
\noindent\textbf{Step 1.} Write $\Sigma=H^\nu\simeq\PP^1$, $G=\Gamma_{\rm red}$,
$\delta=p_a(H)\le2$, and
$s=\#\nu^{-1}(\operatorname{Sing}H)\le3$ as in
(C.RA.5)--(C.RA.6). There are exactly two generic
normalization branches above $H$, both supplied by the
degree-two map from $C$. Thus no vertical strict conductor
component dominates $H$.

\noindent\textbf{Step 2.} We first justify using the lci jet construction without
assuming that property. The low-delta classification in the
proof of \ref{cdp:st:C.2.RAT.1} applies to this same reduced image curve.
Its three-independent-axis germ is excluded by
$b_1(H)\le b_1(G)\le1$. The remaining non-lci possibilities
are the space cusp with semigroup $\langle3,4,5\rangle$
and a cusp with a transverse smooth branch. The
genus-zero double-cover calculation at the end of \ref{cdp:st:A.5.18}
excludes both: the space-cusp inequalities force
$q_0=-1$ and two simple zeros of the cleared quartic,
one over $H_{\rm reg}$, contradicting (A.T.12); for the
cusp-plus-line ring its cleared quartic has nonpositive
degree and cannot have two branch points. Therefore $H$
is lci at all its points.

\noindent\textbf{Step 3.} Let $\Xi$ be the normalization conductor divisor of $H$,
and use the actual jet lines and sections in
(A.T.5)--(A.T.7). Put
\[
 A=\deg\operatorname{div}(a)=-2\deg Q,\qquad
 \sigma\in H^0(\Sigma,L^2(4\Xi)),\qquad
 \deg\operatorname{div}(\sigma)=8-A.
\tag{C.RA.26}\label{cdp:eq:C.RA.26}
\]
Here $A$ is a nonnegative even integer and $\sigma$ is
the conductor-cleared version of $\Delta=ae-b^2/4$.
The degree-two cover $C\to\Sigma$ has two branch points,
exactly the odd-order zeros of $\sigma$. Let $B_s$ be
the number of these branch points lying over singular
points of $H$; there are $2-B_s$ over its smooth locus.
Every smooth branch point has zero order at least three
in $\sigma$, by the actual coordinate exclusion
(A.T.10)--(A.T.12).

\noindent\textbf{Step 4.} The section $a$ must vanish at every point of $\Sigma$
over a singular point of $H$. To see this, the Hessian
matrix has entries in $\OO_H$. If its normalized scalar
$a$ is nonzero on one branch above $p$, its primitive
rank-one factor has a unit component there. A diagonal
entry consequently has a nonzero value at $p$. That
same entry has the same nonzero value on every branch,
so $a$ is nonzero at all normalization points above $p$.
In that event (A.T.13) removes the conductor denominator:
$\Delta$ is holomorphic in $L^2$ there. The orders of
$\sigma$ above $p$ are at least those of $4\Xi$,
which are even. Its total zero divisor must additionally
contain the two parity contributions of the two branch
points, whether or not they lie over $p$. Thus
\[
 \deg\operatorname{div}(\sigma)
       \ge4\deg(\Xi|_{\nu^{-1}(p)})+2
       =8\delta_p+2\ge10,
\tag{C.RA.27}\label{cdp:eq:C.RA.27}
\]
contrary to \eqref{cdp:eq:C.RA.26}. This proves the claimed vanishing.

\noindent\textbf{Step 5.} It follows in particular that
\[
                              A\ge2s-2.
\tag{C.RA.28}\label{cdp:eq:C.RA.28}
\]
For $s\le1$ this is automatic; for $s=2$ the two
zeros give $A\ge2$. If $s=3$, the bounds
$\delta\le2$ and $b_1(H)\le1$ force exactly two
singular points, each of delta invariant one: one node
and one cusp. The two node preimages contribute at
least one each. At the cusp the ring is
$\mathbb C\{t^2,t^3\}$, so any vanishing Hessian entry
has order at least two. Dividing such a diagonal entry
by the square of a unit component of the primitive
normalized factor shows that $a$ has order at least
two there. Thus $A\ge1+1+2=4=2s-2$.

\noindent\textbf{Step 6.} We now check exactly which smooth image points can
receive two vertical boundary attachments on their two
conductor preimages. At a smooth point with $a\ne0$,
\ref{cdp:st:A.5.17} gives the constant $A_3$ product and no vertical
attachment. Suppose instead that $a=0$ and $\Delta$
is a unit. In the actual coordinates (A.T.10), the
coefficient $b$ is then a unit. Substitute $x=y^2\lambda$
in that equation and divide by $y^4$. The result is
holomorphic in $(\lambda,y,z)$ and has, at the point
in question, the simple root
$\lambda=-e(0)/b(0)$. The implicit-function theorem
gives a holomorphic $\lambda(y,z)$ and an actual
factorization
\[
 f_0=F_1F_2,\qquad
 F_1=x-y^2\lambda(y,z),\qquad
 F_2|_{F_1=0}=y^2\,u(y,z),\quad u(0)\ne0.
\tag{C.RA.29}\label{cdp:eq:C.RA.29}
\]
The last equality is the restriction of
$\partial f_0/\partial x$ to $F_1=0$; its quotient
by $y^2$ has constant value $b(0)$.
The branch $F_1=0$ is smooth, and its conductor in the
whole reduced hypersurface is exactly the divisor
$2\{y=0\}$, by the intersection factor in
(C.RA.22). There is therefore no vertical conductor
near this normalization point. Its minimal resolution
is an isomorphism there as well. Since $\Delta$ is a
unit, this is one of the two distinct conductor
preimages over the given point of $\Sigma$.
Consequently a smooth nonbranch value with two vertical
attachments has a positive even zero order of $\sigma$,
at least two.

\noindent\textbf{Step 7.} Let $d_{\rm reg}$ count distinct smooth values with
two such attachments. The two branch points and these
distinct nonbranch zeros give
\[
             3(2-B_s)+B_s+2d_{\rm reg}\le8-A.
\tag{C.RA.30}\label{cdp:eq:C.RA.30}
\]
Among the $s$ points above the singular locus, the
$B_s$ branch points each have only one preimage on
$C$. Thus at most $s-B_s$ of these values can have
two attachments. If $d$ counts all values on
$\Sigma$ with two attachments, \eqref{cdp:eq:C.RA.28}--\eqref{cdp:eq:C.RA.30}
give
\[
 d\le s-B_s+d_{\rm reg}
       \le s+1-\frac A2\le2.
\tag{C.RA.31}\label{cdp:eq:C.RA.31}
\]

\noindent\textbf{Step 8.} In fact $d\le1$. If $d=2$, \eqref{cdp:eq:C.RA.31} requires
$A\le2s-2$. For $s=0$ this is impossible, and for
$s=1$ it contradicts the vanishing of $a$ over the
singular point. For $s=2$, it forces $A=2$ and
the two singular preimages exhaust the zero divisor
of $a$. For $s=3$, it forces $A=4$, exhausted by
the node and cusp orders $1+1+2$ established above.
In both remaining cases $a$ is nowhere zero over
$H_{\rm reg}$. The constant $A_3$ product there
has neither a cover branch point nor a vertical
attachment. Thus both cover branch points lie
over $\operatorname{Sing}H$, and only the other
$s-2\le1$ singular preimages can have two
attachments. This again contradicts $d=2$.
Consequently
\[
                            d\le1.
\tag{C.RA.31a}\label{cdp:eq:C.RA.31a}
\]

\noindent\textbf{Step 9.} It remains to connect this local count to the exact
surface Euler number. As in \ref{cdp:st:C.2.RAT.3}, at least
$K\ge r-N_0$ original positive vertical boundaries
have a component which is not $\pi$-exceptional.
None maps onto $H$, so they give $K$ distinct points
of $C$ anchored to image curves other than $H$.
After projection to $\Sigma$ their number of distinct
normalization-incidence edges is $k=K-d'$ where
$d'\le d\le1$: a degree-two fibre has at most two
points, and the loss of one occurs precisely at a
value with two selected attachments. Apply the graph
construction of (C.RA.24) to these $k$ distinct
edges. It gives $n_G\ge k$ and $e(G)\ge n_G$.
The actual Euler decomposition therefore yields
\[
\begin{split}
 e(S)&=2-r+N_0+e(G)
       \ge2-K+e(G)\\
     &\ge2-K+k=2-d'\ge1.
\end{split}
\tag{C.RA.32}\label{cdp:eq:C.RA.32}
\]
Every contraction, covering fibre and graph edge in
this count comes from the specified maps
$Z\to T\to S$ and $C\to\Sigma\to H$.
\end{proof}






\begin{proposition}[A locally unibranch tangent-root component gives positive Euler number]\label{cdp:st:C.2.RAT.5}
\StatementRecord{cdp:st:C.2.RAT.5}{proved}{Historical events: Claude Fable 5.2 (Anthropic), effort high, 16 November 2026: carried as background without fault found. GPT-6 Astra (OpenAI), ChatGPT Pro mode, effort Pro, S6\_09\_PAPERS, 17 November 2026: targeted proof check. Also checked in this writing. Pre-continuation writing events: GPT-6 Astra (OpenAI), ChatGPT Pro mode, effort Pro, S6\_10\_PAPERS, 18 November 2026: targeted proof check. Continuation metadata events: GPT-6 Astra (OpenAI), ChatGPT Pro mode, effort Pro, S6\_10\_PAPERS continuation, 19 November 2026: used and rechecked in this run; status-vocabulary reconciliation; metadata reconciliation only; no new mathematical check. \emph{Pre-continuation recorded proof scope.} The root-field identification extends by integral closure to the actual normalized conductor cover. At a simple cubic root, holomorphic division gives a smooth branch whose complementary factor restricts to y squared times a unit, so its conductor is exactly twice the horizontal curve. Local unibranchness prevents two attached conductor points sharing a root; a cubic cannot have two different multiple roots. The remaining generic branches and the exact graph inequality give e(S)>=d-1}{the hypotheses, notation and definitions stated here; \cref{cdp:st:C.2.RAT.1}; \cref{cdp:eq:C.RD.2}; \cref{cdp:eq:C.RD.4}; \cref{cdp:st:C.2.RAT.3}}

Retain the torus-fibre, rational-normalization and specified
ruled-boundary morphism of \ref{cdp:st:C.2.RAT.1}. Suppose $C$ is not
contracted and maps with degree $d=2$ or $3$ onto a smooth
rational curve $H$ with generic ordinary-triple transversal.
Let $q_3\in H^0(H,\operatorname{Sym}^3N^*\otimes M)$ be
its cubic normal tensor, where $N=N_{H/Y}$ and
$M=\OO_Y(S)|_H$. Suppose $q_3$ is nonzero as a tensor at
every point of $H$. Its tangent-root divisor
$R_3\subset\PP(N)$ is then finite of degree three over $H$.
Let $R_d$ be the component corresponding to $C$ over the
generic point. If $R_d$ is locally unibranch at every point,
then
\[
                           e(S)\ge d-1>0.
\tag{C.RD.1}\label{cdp:eq:C.RD.1}
\]
In particular this applies when $d=2$ and the only singularities
of $R_2$ are an $A_4$ cusp or two $A_2$ cusps.
\end{proposition}


\begin{proof}
\noindent\textbf{Step 1.} The nonzero cubic has three distinct roots over the generic
point and is never identically zero on a projective fibre.
It therefore defines a reduced finite divisor $R_3$.
The function-field algebra of its roots agrees with that of
the normalized conductor curves dominating $H$: a branch
of an ordinary triple is distinguished by its tangent line.
Taking the integral closure of $\OO_H$ in that algebra gives
the exact finite-cover identification
\[
                            C\simeq R_d^\nu
                            \quad\text{over }H.
\tag{C.RD.2}\label{cdp:eq:C.RD.2}
\]
Thus two distinct points of $C$ over one point $p\in H$
give distinct root points of $R_d$, because normalization
of a locally unibranch reduced curve is bijective on its
underlying points.

\noindent\textbf{Step 2.} We check what a simple root implies for the actual
hypersurface, including its conductor. In coordinates
$H=(x,y)=0$, with parameter $z$ on $H$, the local equation
$f_0$ of $S$ belongs to $(x,y)^3$. The lower normal jets
vanish on the reduced smooth curve, and the graded pieces
of its regular-sequence ideal are locally free, which
justifies this membership. Choose the normal coordinates
so a specified simple tangent root is finite in $x/y$.
Then
\[
 f_0(y\lambda,y,z)/y^3
       =q_3(\lambda,1,z)+y\,A(\lambda,y,z).
\tag{C.RD.3}\label{cdp:eq:C.RD.3}
\]
The simple root gives a holomorphic solution
$\lambda=\lambda(y,z)$ by the implicit-function theorem.
Division by the monic linear factor yields an actual
factorization
\[
 f_0=(x-y\lambda(y,z))\,G(x,y,z),\qquad
 G|_{x=y\lambda(y,z)}=y^2 u(y,z),\quad u(0)\ne0.
\tag{C.RD.4}\label{cdp:eq:C.RD.4}
\]
The unit is the derivative of the cubic with respect to
its root coordinate, at the chosen simple root. Hence
this normalization branch is smooth and its conductor
in $S$ is exactly twice the curve lying over $H$.
There is no vertical conductor component there, and
the minimal resolution $\pi$ is an isomorphism nearby.

\noindent\textbf{Step 3.} Two distinct $C$-points over $p$ cannot both receive a
vertical boundary attachment. Their root points are
distinct by \eqref{cdp:eq:C.RD.2}. If both could receive attachments,
\eqref{cdp:eq:C.RD.4} would make both roots multiple, requiring total
cubic degree at least four. Thus each point of $H$
receives at most one attachment from $C$.

\noindent\textbf{Step 4.} There are at most $3-d$ generic normalization branches
left to be supplied by vertical strict conductor
components. In the notation of \ref{cdp:st:C.2.RAT.3} this gives
$r_{\rm dom}\le3-d$. At least
$K\ge r-N_0-r_{\rm dom}$ original positive vertical
boundaries therefore survive with an image different
from $H$. Their $K$ points of $C$ give $K$ distinct
normalization-incidence edges at $H$, by the preceding
paragraph. Equation (C.RA.24) gives $n_G\ge K$, and
the exact Euler decomposition gives
\[
 e(S)\ge2-r+N_0+K
       \ge2-r_{\rm dom}\ge d-1.
\]
This proves the assertion. The final examples are locally
unibranch by their displayed cusp equations.
\end{proof}






\begin{proposition}[Every primitive rational double ordinary-triple cover has nonnegative Euler number]\label{cdp:st:C.2.RAT.7}
\StatementRecord{cdp:st:C.2.RAT.7}{proved}{Historical events: Claude Fable 5.2 (Anthropic), effort high, 16 November 2026: carried as background without fault found. GPT-6 Astra (OpenAI), ChatGPT Pro mode, effort Pro, S6\_09\_PAPERS, 17 November 2026: targeted proof check. Also checked in this writing. Pre-continuation writing events: GPT-6 Astra (OpenAI), ChatGPT Pro mode, effort Pro, S6\_10\_PAPERS, 18 November 2026: targeted proof check. Continuation metadata events: GPT-6 Astra (OpenAI), ChatGPT Pro mode, effort Pro, S6\_10\_PAPERS continuation, 19 November 2026: used and rechecked in this run; status-vocabulary reconciliation; metadata reconciliation only; no new mathematical check. \emph{Pre-continuation recorded proof scope.} The cyclic character algebra and relative duality give H1(X0,O)=0 for every normal-bundle splitting. The resolved-boundary incidence formula retains disconnected root components and the exact class-relation number. Hodge index and Gysin identify the compact pairing rank as 2*nu+2*b1(X)-lambda+epsilon; actual compact-core transport into the torus forces nu<=1. Each duplicate attachment consumes a distinct multibranch root point, giving the final nonnegative Euler bound}{the hypotheses, notation and definitions stated here; \cref{cdp:st:C.2.RAT.1}; \cref{cdp:st:A.5.19}; \cref{cdp:st:C.2.RAT.6}; \cref{cdp:eq:C.RD.14}; \cref{cdp:eq:C.RD.15}}

Retain the torus-fibre, rational-normalization and specified
ruled-boundary morphism of \ref{cdp:st:C.2.RAT.1}. Suppose $C$ is not
contracted, its image $H$ has generic ordinary-triple
transversal, and $C\to H^\nu$ has degree two. Suppose
the cubic normal tensor along $H$ has no common zero of
its coefficients. Then $e(S)\ge0$. This holds for every
splitting type of the normal bundle $N_{H/Y}$.
\end{proposition}


\begin{proof}
\noindent\textbf{Step 1.} Statement \ref{cdp:st:A.5.19} makes $H$ smooth: over a singular
rational ordinary-triple image every connected normalized
conductor cover has degree one. Thus $H\simeq\PP^1$.
Set $N=N_{H/Y}$, $W=\PP(N)$, and
$\xi=c_1(\OO_W(1))$, where $\PP(N)$ parameterizes lines.
The ambient Chern-class hypothesis gives
$\deg N=-2$ and $\OO_Y(S)|_H\simeq\OO_H$.
The cubic is an actual section of $\OO_W(3)$.
The assumed absence of a common coefficient zero makes
its root divisor finite of degree three over $H$:
\[
 R=R_1+R_2\sim3\xi,\qquad
 \xi^2=2,\qquad K_W=-2\xi,\qquad p_a(R)=4.
\tag{C.RD.11}\label{cdp:eq:C.RD.11}
\]
Here $R_1$ is a section and hence is smooth rational;
$R_2$ has normalization $C\simeq\PP^1$, by the exact
integral-closure comparison (C.RD.2). Both are reduced
and they are different. The normalization sequence gives
\[
                  \sum_p\delta_p(R)=5.
\tag{C.RD.12}\label{cdp:eq:C.RD.12}
\]
This equality also holds if the two components are
disjoint: the arithmetic-genus convention is
$p_a(R)=1-\chi(\OO_R)$ and the normalization has Euler
characteristic two.

\noindent\textbf{Step 2.} Form the cyclic triple cover
$p_0:X_0\to W$ with algebra
$\OO_W\oplus\OO_W(-1)\oplus\OO_W(-2)$ and multiplication
given by the cubic. It is normal because its branch
divisor $R$ is reduced, and is connected because a
meridian around a smooth branch point permutes its
three sheets cyclically. Resolve its isolated
singularities and its reduced boundary to obtain
$X$, with $D=X\setminus U$, where
$U=\{v\in N:q_3(v)=1\}$ is the smooth normal-cone
level surface. The isomorphism over $W\setminus R$
is the exact radial-vector/dual-tautological-coordinate
comparison (C.D4.6).

\noindent\textbf{Step 3.} The projective-bundle cohomology calculation gives
\[
 H^1(W,\OO_W)=H^1(W,\OO_W(-1))
             =H^1(W,\OO_W(-2))=0.
\tag{C.RD.13}\label{cdp:eq:C.RD.13}
\]
Indeed both relative cohomologies of $\OO_W(-1)$
vanish, while $p_*\OO_W(-2)=0$ and
$R^1p_*\OO_W(-2)=\det N=\OO_{\PP^1}(-2)$.
Thus $H^1(X_0,\OO_{X_0})=0$ for every normal-bundle
splitting. Resolution Leray and \ref{cdp:st:C.2.RAT.6} give
$q(X)\le\sum p_g(X_0,x)\le1$, so $b_1(X)\le2$.

\noindent\textbf{Step 4.} We spell out the boundary calculation with its number
of connected components retained. Write
\[
 \nu=\sum_p(r_p(R)-1),\qquad
 c=b_0(R),\qquad
 \lambda=\sum_x b_1(\operatorname{link}(X_0,x))\le2.
\]
Normalization of $R$ gives $e(R)=4-\nu$.
Since $U\to W\setminus R$ is etale of degree three,
$e(U)=3\nu$. The resolved boundary has the same
number $c$ of connected components. Replacing each
singular point of $R$ by the connected exceptional
divisor of its cubic suspension gives, by the
Mayer--Vietoris sequence with its individual branch
attachment points,
\[
                 b_1(D)=\nu-2+c+\lambda.
\tag{C.RD.14}\label{cdp:eq:C.RD.14}
\]
The exceptional first Betti number is the link's
first Betti number: the exceptional intersection
matrix in its plumbed neighbourhood is invertible
over $\Q$. This is precisely the local comparison
proved with (C.D4.7). Additional boundary-resolution
blowups add rational trees and leave \eqref{cdp:eq:C.RD.14}
unchanged.

\noindent\textbf{Step 5.} Let $d_D$ be the number of irreducible components
of $D$, and let $V\subset H_2(X,\Q)$ be the span
of their classes. The exceptional classes are
independent. Pushing a relation to $W$ and pulling
back a relation from $R_1,R_2$ shows that
\[
 \varepsilon:=d_D-\dim V
      =2-\dim\operatorname{span}\{[R_1],[R_2]\}
      \in\{0,1\}.
\]
This proves the relation count rather than
assuming the two branch classes are independent.
The space $V$ contains the pullback of $R$,
whose square is $3R^2=54>0$. The Hodge index
theorem therefore makes its intersection pairing
nondegenerate. The divisor Gysin sequence gives
$b_1(U)=b_1(X)+\varepsilon\le3$; its degree-two
part identifies the image of $H_2(U)$ in
$H_2(X)$ with $V^\perp$. Consequently the rank
of the compact intersection pairing on $U$ is
$b_2(X)-d_D+\varepsilon$.
Now use $e(X)=2-2b_1(X)+b_2(X)$ and
$e(D)=c-b_1(D)+d_D$. Equations \eqref{cdp:eq:C.RD.14} and
$e(U)=3\nu$ yield the exact identity
\[
\begin{split}
 \operatorname{rank}I_U
   &=e(U)-2+2b_1(X)+c-b_1(D)+\varepsilon\\
   &=2\nu+2b_1(X)-\lambda+\varepsilon
     \ge2\nu-2.
\end{split}
\tag{C.RD.15}\label{cdp:eq:C.RD.15}
\]
In particular this formula includes disjoint root
components; no connected-boundary assumption has
been introduced.

\noindent\textbf{Step 6.} If $\nu\ge2$, \eqref{cdp:eq:C.RD.15} gives a nonzero compact
intersection pairing on a surface with $b_1(U)\le3$.
The compact-core transport (C.D4.2) embeds that
pairing into an actual nearby torus. Its induced
first-homology image has dimension at most three,
so its second-homology image has zero wedge pairing,
as established immediately after (C.D4.2). This
contradiction proves $\nu\le1$.

\noindent\textbf{Step 7.} Finally suppose a value of $H$ receives vertical
boundary attachments at both its $C$-preimages.
If their tangent-root points in $R_2$ were distinct,
at least one would be a simple root of the cubic,
and (C.RD.3)--(C.RD.4) would give pure conductor
there, a contradiction. They therefore map to a
single multibranch point of $R_2$. Distinct values
give distinct such points, each contributing at
least one to $\nu$. Hence there is at most one
value with two attachments. With
$K\ge r-N_0-r_{\rm dom}$ surviving original
vertical boundaries, at least $K-1$ distinct
normalization-incidence edges remain at $H$.
There is only one remaining generic triple branch,
so $r_{\rm dom}\le1$. The graph comparison
(C.RA.24) and exact Euler decomposition now give
\[
 e(S)\ge2-r+N_0+(K-1)
                         \ge1-r_{\rm dom}\ge0.
\tag{C.RD.16}\label{cdp:eq:C.RD.16}
\]
This completes the proof.
\end{proof}






\begin{proposition}[A single common cubic zero in the rational double-cover case]\label{cdp:st:C.2.RAT.8}
\StatementRecord{cdp:st:C.2.RAT.8}{proved}{Historical events: Claude Fable 5.2 (Anthropic), effort high, 16 November 2026: carried as background without fault found. GPT-6 Astra (OpenAI), ChatGPT Pro mode, effort Pro, S6\_09\_PAPERS, 17 November 2026: targeted proof check. Also checked in this writing. Pre-continuation writing events: GPT-6 Astra (OpenAI), ChatGPT Pro mode, effort Pro, S6\_10\_PAPERS, 18 November 2026: targeted proof check. Continuation metadata events: GPT-6 Astra (OpenAI), ChatGPT Pro mode, effort Pro, S6\_10\_PAPERS continuation, 19 November 2026: used and rechecked in this run; status-vocabulary reconciliation; metadata reconciliation only; no new mathematical check. \emph{Pre-continuation recorded proof scope.} Every possible horizontal intersection with the common ruling fibre was enumerated from intersection total three and horizontal delta three. The two Newton facets of the cusp with two transverse branches have exactly one positive lattice point under their boundary; all higher terms lie above the faces. The geometric-genus and link bounds give pairing rank 2*nu-4+2*b1(X)-lambda, forcing nu<=3. The exact common-fibre branch count 1+a and the outside multibranch count bound the total lost attachment value by one}{the hypotheses, notation and definitions stated here; \cref{cdp:st:C.2.RAT.1}; \cref{cdp:st:C.2.RAT.7}; \cref{cdp:eq:C.RD.17}; \cref{cdp:st:C.2.RAT.6}; \cref{cdp:eq:C.RD.18}; \cite[Theorem 2.1.1 and Corollary 2.1.3]{MerleTeissier1977}; \cref{cdp:eq:C.RD.20}; \cref{cdp:eq:C.RD.21}; \cref{cdp:eq:C.RD.22}}

Retain the torus-fibre, rational-normalization and specified
ruled-boundary morphism of \ref{cdp:st:C.2.RAT.1}. Suppose $C$ is not
contracted, its image has generic ordinary-triple transversal,
and $C\to H^\nu$ has degree two. If the common zero divisor
of the cubic normal coefficients has degree one, then
$e(S)\ge0$.
\end{proposition}


\begin{proof}
\noindent\textbf{Step 1.} As in \ref{cdp:st:C.2.RAT.7}, $H$ is smooth rational. On $W=\PP(N)$
write the complete root divisor as
\[
 B=R_1+R_2+F\sim3\xi,
 \qquad B_h=R_1+R_2\sim3\xi-F.
\tag{C.RD.17}\label{cdp:eq:C.RD.17}
\]
Here $F$ is the ruling fibre over the unique simple
common zero. The horizontal divisor $B_h$ is finite
over $H$, $R_1$ is a section, and the normalization of
$R_2$ is the rational degree-two cover $C$.
Since $K_W=-2\xi$, $\xi^2=2$, and $\xi F=1$,
\[
 p_a(B_h)=2,
 \qquad \sum_p\delta_p(B_h)=3,
 \qquad I(F,B_h)=3.
\tag{C.RD.18}\label{cdp:eq:C.RD.18}
\]
The full divisor $B$ has total delta six. We will
use its particular decomposition \eqref{cdp:eq:C.RD.17}, not a
statement about arbitrary delta-six plane germs.

\paragraph{The local cubic-suspension genus along $F$.}
At a point of $F$, put $i=I_p(F,B_h)$. If $i=1$,
the germ of $B$ is an ordinary node. If $i=2$,
the horizontal germ is either smooth with contact
two with $F$, two smooth branches transverse to
$F$, or a multiplicity-two cusp transverse to $F$.
The bound $\delta_p(B_h)\le3$ makes the resulting
plane types respectively
\[
 A_3,\qquad D_{2k+2}\ (1\le k\le3),\qquad
                 D_{2k+3}\ (1\le k\le3).
\tag{C.RD.19}\label{cdp:eq:C.RD.19}
\]
For two smooth branches $k$ is their contact;
for the cusp $k$ is its delta invariant.
The exact analytic identifications are the graph
straightening and square-completion comparisons
in \ref{cdp:st:C.2.RAT.6}. Their cubic suspensions have
geometric genus zero or one by (C.RD.10).

\noindent\textbf{Step 2.} If $i=3$, the point contains $R_1$ and all the
intersection of $R_2$ with $F$. The section $R_1$
is transverse to $F$. A smooth branch of $R_2$
with intersection two is tangent to $F$ and
transverse to $R_1$, giving $D_6$. If $R_2$ has
two local branches, both are smooth and transverse
to $F$. The horizontal delta bound three forces
all three horizontal branches to be pairwise
transverse. The full germ is therefore four
distinct transverse smooth branches.
The last possibility is one singular branch
of $R_2$. Its multiplicity is two, since its
intersection with $F$ is two. Its intersection
with $R_1$ is at least two; hence \eqref{cdp:eq:C.RD.18}
forces it to be an ordinary cusp, transverse
to both smooth branches $F$ and $R_1$.
These alternatives exhaust $i=3$.

\noindent\textbf{Step 3.} For four transverse branches, the degree-four
initial plane polynomial has four distinct
linear factors. Its cubic suspension is Newton
nondegenerate with weights $(1/4,1/4,1/3)$;
the only positive lattice point on or below
its compact Newton face is $(1,1,1)$. Its
geometric genus is therefore one by
\cite[Theorem 2.1.1 and Corollary 2.1.3]{MerleTeissier1977}.

\noindent\textbf{Step 4.} We give the precise same calculation for the cusp
and two transverse branches, retaining higher
terms of the actual equation. Straighten $F,R_1$
to coordinate axes and make a linear change so
the cusp tangent is $u=0$. Up to units, the
equation is then
\[
 r(u,v)=v(u-v)\bigl(u^2+c v^3+
           u\,O((u,v)^2)+O((u,v)^4)\bigr),
                  \qquad c\ne0.
\tag{C.RD.20}\label{cdp:eq:C.RD.20}
\]
The nonzero radical cubic is the ordinary-cusp
condition. The compact Newton facets of
$z^3-r(u,v)$ have normalized inward vectors
\[
             (1/4,1/4,1/3),\qquad
             (3/10,1/5,1/3).
\tag{C.RD.21}\label{cdp:eq:C.RD.21}
\]
Their plane edges have vertex sets
$\{(3,1),(2,2)\}$ and
$\{(2,2),(0,5)\}$. The higher terms in
\eqref{cdp:eq:C.RD.20} are strictly above these faces.
The first edge polynomial is a nonzero multiple
of $u^2v(u-v)$; the second is a nonzero multiple
of $v^2(u^2+c v^3)$. Both have smooth zero sets
in $(\C^*)^2$. Each two-dimensional compact
face additionally contains the term $z^3$,
whose $z$-derivative is nonzero on $(\C^*)^3$.
This proves the required Newton nondegeneracy.
Every noncompact face inequality is strictly
satisfied by positive integral points; the only
additional one beyond coordinate positivity is
$v+z/3\ge1$.
For either vector in \eqref{cdp:eq:C.RD.21}, the point
$(1,1,1)$ has weight $5/6$, and increasing any
one coordinate makes its weight greater than
one. Thus it is the only positive lattice
point outside the interior of the Newton
polyhedron. The same adjunction theorem gives
$p_g(z^3-r)=1$.

\noindent\textbf{Step 5.} This proves that every suspension singularity
over $F$ has geometric genus at most one.
A nonzero value requires $i\ge2$ and
$\delta_p(B_h)\ge1$: the cases with horizontal
delta zero were the ordinary node and $A_3$.
Since the total intersection with $F$ is three,
there is at most one such nonzero value over
$F$. Away from $F$, \ref{cdp:st:C.2.RAT.6} shows that a
nonzero geometric genus costs horizontal delta
at least three. It cannot coexist with a
nonzero value over $F$, by \eqref{cdp:eq:C.RD.18}.
Consequently the cyclic cover $X_0$ branched
over $B$ satisfies
\[
                        \sum_x p_g(X_0,x)\le1.
\tag{C.RD.22}\label{cdp:eq:C.RD.22}
\]

\noindent\textbf{Step 6.} We also obtain the required link bound without
assuming a delta-six monodromy table. For a
normal surface singularity, let $E$ be the
reduced exceptional divisor of a good resolution.
The restriction maps from infinitesimal
neighbourhoods of $E$ onto $E$ are surjective
on $H^1$, since the kernels have one-dimensional
support and no $H^2$. Formal functions therefore
gives $h^1(E,\OO_E)\le p_g$.
For an SNC exceptional curve,
$b_1(E)=2\sum g_i+\beta\le
2(\sum g_i+\beta)=2h^1(E,\OO_E)$.
The invertible exceptional intersection matrix
identifies this $b_1(E)$ with the link's first
Betti number, as in the proof of \ref{cdp:st:C.2.RAT.6}.
Equation \eqref{cdp:eq:C.RD.22} yields
\[
              \lambda:=\sum_x b_1(\operatorname{link}(X_0,x))
                            \le2.
\tag{C.RD.23}\label{cdp:eq:C.RD.23}
\]

\paragraph{The compact pairing and the attachment count.}
The cyclic algebra and projective-bundle cohomology
are unchanged from (C.RD.13), so its projective
resolution $X$ has $b_1(X)\le2$.
The divisor $B$ is connected, since $F$ meets
both horizontal components. Its three classes
span a two-dimensional space: the section
$R_1$ and the fibre $F$ are independent.
Thus the boundary-relation number is exactly
$\varepsilon=1$, and $b_1(U)=b_1(X)+1\le3$.
Put $\nu=\sum_p(r_p(B)-1)$.
Normalization of its three rational components
and the resolved-boundary comparison give
\[
 e(U)=3(\nu-2),\qquad
 b_1(D)=\nu-2+\lambda,
 \qquad
 \operatorname{rank}I_U
       =2\nu-4+2b_1(X)-\lambda\ge2\nu-6.
\tag{C.RD.24}\label{cdp:eq:C.RD.24}
\]
The rank identity is the same Gysin, Hodge-index
and Euler calculation as (C.RD.15), with all
three boundary components and their one relation
included. Compact-core transport into the
nearby torus therefore forces $\nu\le3$.

\noindent\textbf{Step 7.} Let $p\in H$ be the common zero value and
$a=\#(C\to H)^{-1}(p)\in\{1,2\}$.
At every point of $F\cap B_h$, the full divisor
has the branch of $F$ together with the
individual horizontal normalization branches.
It follows exactly that
\[
        \nu_F:=\sum_{x\in F}(r_x(B)-1)=1+a.
\tag{C.RD.25}\label{cdp:eq:C.RD.25}
\]
The contribution one is the section $R_1$;
the other $a$ branches are the normalization
points of $R_2$ above $p$. This identity retains
coinciding tangent roots and ramification.

\noindent\textbf{Step 8.} The loss of distinct horizontal anchor values
at $p$ is at most $a-1=\nu_F-2$. Away from
$p$, the cubic is primitive. The argument at
the end of \ref{cdp:st:C.2.RAT.7} makes each duplicate
attachment there a multibranch point of $R_2$,
costing at least one in
$\nu_{\rm out}=\nu-\nu_F$. Hence the total
loss $\ell$ of distinct normalization-incidence
edges satisfies
\[
                    \ell\le\nu_F-2+\nu_{\rm out}
                              =\nu-2\le1.
\tag{C.RD.26}\label{cdp:eq:C.RD.26}
\]
Finally $r_{\rm dom}\le1$ and
$K\ge r-N_0-r_{\rm dom}$. Projecting these
$K$ surviving attachments to $H$ retains at
least $K-\ell$ distinct edges. Equation
(C.RA.24) and the Euler decomposition give
\[
                 e(S)\ge2-r_{\rm dom}-\ell\ge0.
\tag{C.RD.27}\label{cdp:eq:C.RD.27}
\]
The two Newton facets, their exact lower
lattice-point counts and the attachment
budgets are recorded in
\path{code/A_cubic_scalar_one.py} and its output.
\end{proof}






\begin{proposition}[Every rational double ordinary-triple cover has nonnegative Euler number]\label{cdp:st:C.2.RAT.9}
\StatementRecord{cdp:st:C.2.RAT.9}{proved}{Historical events: Claude Fable 5.2 (Anthropic), effort high, 16 November 2026: carried as background without fault found. GPT-6 Astra (OpenAI), ChatGPT Pro mode, effort Pro, S6\_09\_PAPERS, 17 November 2026: targeted proof check. Also checked in this writing. Pre-continuation writing events: GPT-6 Astra (OpenAI), ChatGPT Pro mode, effort Pro, S6\_10\_PAPERS, 18 November 2026: targeted proof check. Continuation metadata events: GPT-6 Astra (OpenAI), ChatGPT Pro mode, effort Pro, S6\_10\_PAPERS continuation, 19 November 2026: used and rechecked in this run; status-vocabulary reconciliation; metadata reconciliation only; no new mathematical check. \emph{Pre-continuation recorded proof scope.} The exact normal-bundle coefficient degrees bound the common divisor by two. With two simple common values, canonical duality forces q(X)=0 and the boundary pairing gives nu<=5. With one double value, the normalized character generator w squared over s and every weighted node/cusp blowup have character increments zero or one. Their proper birational direct-image comparison and final cyclic quotient resolutions prove rationality. The two exact boundary-rank and attachment counts establish the claimed Euler bounds in both cases}{the hypotheses, notation and definitions stated here; \cref{cdp:st:C.2.RAT.1}; \cref{cdp:st:A.5.19}; \cref{cdp:st:C.2.RAT.7}; \cref{cdp:st:C.2.RAT.8}; \cref{cdp:eq:C.RD.28}; \cref{cdp:eq:C.RD.34}}

Retain the torus-fibre, rational-normalization and specified
ruled-boundary morphism of \ref{cdp:st:C.2.RAT.1}. Suppose $C$ is not
contracted, its image has generic ordinary-triple transversal,
and $C\to H^\nu$ has degree two. Then $e(S)\ge0$,
including all common zeros and their multiplicities in
the cubic normal tensor.
\end{proposition}


\begin{proof}
\noindent\textbf{Step 1.} Again \ref{cdp:st:A.5.19} makes $H$ smooth rational. Let $A$ be the
common zero divisor of the cubic coefficients and put
$v=\deg A$. The cases $v=0,1$ are \ref{cdp:st:C.2.RAT.7}--8.
We first verify that $v\le2$. Write
$N=\OO(a)\oplus\OO(-2-a)$ with $a\ge-1$.
For $a\ge0$ and $v>0$, the coefficient of the cube
of the first normal coordinate has degree $-3a-v$
and vanishes. Thus the other coordinate is a factor.
The remaining quadratic is irreducible over the
function field, because it describes the connected
degree-two cover. Its first-coordinate-square
coefficient must be nonzero and has degree
$2-a-v$. This gives $a+v\le2$.
For $a=-1$, on $\PP^1\times\PP^1$ the section
$R_1$ and quadratic component have bidegrees
$(1,d)$ and $(2,3-v-d)$, with $d\ge0$ and the
first entry denoting degree over $H$. An integral
quadratic component requires $3-v-d\ge1$:
a divisor of bidegree $(2,0)$ is a union of
constant sections over $\C$. Thus $v\le2$ here
as well.

\noindent\textbf{Step 2.} It remains to consider $v=2$. With the notation
of \ref{cdp:st:C.2.RAT.7}, the horizontal divisor is
\[
 B_h=R_1+R_2\sim3\xi-2F,
       \qquad p_a(B_h)=0,
       \qquad \sum_p\delta_p(B_h)=1.
\tag{C.RD.28}\label{cdp:eq:C.RD.28}
\]
Both normalization components are rational.
Consequently its only possible singularity is
one ordinary node or one ordinary cusp. An
intersection between $R_1$ and $R_2$, if present,
is the node. Otherwise a node or cusp belongs
to $R_2$. Every other point of the horizontal
divisor is smooth.

\paragraph{Two distinct common zero values.}
First let $A=p+q$ with $p\ne q$. The full branch
divisor is reduced:
$B=R_1+R_2+F_p+F_q\sim3\xi$.
Use the ordinary cyclic algebra of (C.RD.13).
The local analysis in the proof of \ref{cdp:st:C.2.RAT.8}
applies at each of the two fibres. A nonzero
local geometric genus there requires positive
horizontal delta. The total horizontal delta
in \eqref{cdp:eq:C.RD.28} is one, and the two fibres are
disjoint. Away from these fibres the local
plane delta is at most one and its cubic
suspension is rational. Therefore
\[
 P:=\sum_x p_g(X_0,x)\le1,\qquad
 \lambda:=\sum_x b_1(\operatorname{link}(X_0,x))\le2.
\tag{C.RD.29}\label{cdp:eq:C.RD.29}
\]

\noindent\textbf{Step 3.} In this case the irregularity is zero, not
merely at most one. The finite-cover dualizing
formula gives
$\omega_{X_0}=p_0^*(\omega_W\otimes\OO_W(2))
\simeq\OO_{X_0}$, and (C.RD.13) gives
$H^1(X_0,\OO)=0$. Thus $\chi(\OO_{X_0})=2$.
For a resolution $\rho:X\to X_0$, local duality
and Grauert--Riemenschneider vanishing identify
$\rho_*\omega_X\subset\omega_{X_0}$ with an
ideal of colength $p_g(X_0,x)$ at each singular
point. Explicitly this is the dual of the
finite-length $R^1\rho_*\OO_X$ in the surface
duality sequence. If $P>0$, the global generator
of $\omega_{X_0}$ does not belong to that ideal
at a point with positive colength. No nonzero
global canonical form on $X$ can then exist,
because its pushforward would be a nonzero
constant multiple of that generator.
Resolution Leray now gives
\[
 \chi(\OO_X)=2-P,
 \qquad q(X)=P-1\quad(P>0),
 \qquad q(X)=0\quad(P=0).
\tag{C.RD.30}\label{cdp:eq:C.RD.30}
\]
Since $P\le1$, this proves $b_1(X)=0$.

\noindent\textbf{Step 4.} The four branch classes span the two-dimensional
space generated by a section and a fibre, so
$\varepsilon=2$ and $b_1(U)=2$.
The branch divisor and resolved boundary are
connected. For
$\nu=\sum_x(r_x(B)-1)$, normalization and the
exact pairing calculation give
\[
 e(U)=3(\nu-4),\qquad b_1(D)=\nu-3+\lambda,
 \qquad \operatorname{rank}I_U=2\nu-8-\lambda.
\tag{C.RD.31}\label{cdp:eq:C.RD.31}
\]
The compact torus comparison forces this rank
to be zero; otherwise its nondegenerate compact
pairing would embed into a rank-at-most-two
first-homology image in the torus. Hence
$\nu=4+\lambda/2\le5$.

\noindent\textbf{Step 5.} Let $a_p,a_q\in\{1,2\}$ be the numbers of
$C$-preimages of these two common values.
The exact branch count (C.RD.25) gives
$\nu_{F_p}+\nu_{F_q}=2+a_p+a_q$.
The loss of distinct attachment values on
these fibres is at most $a_p+a_q-2$.
Away from them each loss costs a multibranch
point of $R_2$. Consequently
\[
                    \ell\le\nu-4\le1.
\tag{C.RD.32}\label{cdp:eq:C.RD.32}
\]
The graph count gives
$e(S)\ge2-r_{\rm dom}-\ell\ge0$.

\paragraph{One double common zero value.}
Now let $A=2p$. Normalize the cyclic cover
with branch multiplicities $R_1+R_2+2F_p$.
Its exact character algebra is
\[
 (p_0)_*\OO_{X_0}
       =\OO_W\oplus L^{-1}\oplus L^{-2}(F_p),
                  \qquad L=\OO_W(1).
\tag{C.RD.33}\label{cdp:eq:C.RD.33}
\]
The additional integral generator at $F_p$
is $w^2/s$ in $w^3=s^2u$. Its divisorwise
character line extends reflexively over the
smooth base, which proves the sheaf identity
at the finitely many intersections as well;
this is also (C.DD.16). Relative duality gives
$R^1p_*L^{-2}(F_p)=\OO_H(-1)$.
Thus $H^1(X_0,\OO_{X_0})=0$.

\noindent\textbf{Step 6.} We check that every singularity of this
normalized cover is rational, including the
horizontal node and cusp allowed by \eqref{cdp:eq:C.RD.28}.
For this purpose write a cyclic branch divisor
as $D_1+2D_2$, with root line $L$.
At a point blowup $\sigma$ let
$m_i=\operatorname{mult}D_i$,
$n=m_1+2m_2=3k+j$, $0\le j<3$, and let $E$
be the new exceptional curve. Keeping the
same function field gives root line
$L'=\sigma^*L(-kE)$ and new exceptional
branch weight $j$. Its two character sheaves
are exactly
\[
\begin{split}
 (L')^{-1}&=\sigma^*(L^{-1})(kE),\\
 (L')^{-2}(D'_2)
   &=\sigma^*(L^{-2}(D_2))
       \bigl((2k-m_2+\mathbf1_{j=2})E\bigr).
\end{split}
\tag{C.RD.34}\label{cdp:eq:C.RD.34}
\]
Here $D'_2=\sigma^*D_2-m_2E+\mathbf1_{j=2}E$.
These identities follow by replacing the
cyclic root by its quotient by the $k$th
power of a local exceptional equation.

\noindent\textbf{Step 7.} At a smooth horizontal branch the contact
with $F_p$ is one or two. The weight-$(2,1)$
contact calculation (C.DD.18)--(C.DD.20)
separates the branches with character
increments $(1,1)$ at each blowup.
If the horizontal node lies on $F_p$, there
are three smooth local branches of weights
$(2,1,1)$. They are transverse except possibly
for contact two between $F_p$ and the
quadratic branch in the case of an intersection
$R_1\cap R_2$. A blowup has $n=4,m_2=1$,
so both character increments in \eqref{cdp:eq:C.RD.34}
are one. In the transverse case it separates
the three branches. In the contact-two case
one remaining ordinary triple has the same
weights $(2,1,1)$; the same blowup separates it.

\noindent\textbf{Step 8.} If the horizontal cusp lies on $F_p$, its
local form is a weight-one cusp transverse
to the weight-two fibre:
\[
                        w^3=s^2(x^2-s^3).
\tag{C.RD.35}\label{cdp:eq:C.RD.35}
\]
The first blowup has $n=4,m_2=1$ and
increments $(1,1)$. The strict cusp becomes
smooth with contact two with the new
weight-one exceptional curve. Blow up that
weight-$(1,1)$ tangency; here $n=2,m_2=0$
and the increments are $(0,1)$.
The remaining ordinary triple has weights
$(1,1,2)$, so a final blowup has increments
$(1,1)$. The resulting branch has normal
crossings. At a crossing of weights $a,b$
in $\{1,2\}$, normalization of
$w^3=x^a y^b$ is the quotient of $\C^2$
under the subgroup
\[
 \{(\zeta^i,\zeta^j):ai+bj=0\pmod3\},
 \qquad x=u^3,\quad y=v^3,\quad w=u^av^b.
\]
These are the rational cyclic quotient
singularities $1/3(1,1)$ or $1/3(1,2)$.
Their standard resolutions have respectively
one $(-3)$-curve or a two-curve $(-2)$ chain,
and $R^1\OO=0$. The horizontal node or cusp
away from $F_p$ likewise has a rational
cubic suspension by (C.RD.10).

\noindent\textbf{Step 9.} For completeness, none of the preceding
base blowups adds a geometric-genus defect.
Every exceptional character increment in
\eqref{cdp:eq:C.RD.34} was zero or one. A point blowup
has $\sigma_*\OO(E)=\OO$ and
$R^1\sigma_*\OO(E)=0$, from
$\OO(E)|_E=\OO_{\PP^1}(-1)$ and the
structure-sheaf direct images; the same
holds for increment zero. Projection formula
therefore makes the first direct image of
each normalized character algebra vanish.
The normalized cyclic covers over successive
base blowups map properly and birationally
to the preceding cover through their same
function field. Finite direct image and
Leray identify their $R^1\OO$ with the
just computed eigensheaf direct images;
faithfulness of finite pushforward gives
$R^1\OO=0$. Resolving the final rational
cyclic quotients preserves this property.
This proves the claimed local rationality.

\noindent\textbf{Step 10.} We consequently have $b_1(X)=0$ and
$\lambda=0$. The reduced branch support
$B=R_1+R_2+F_p$ has three rational
normalization components, is connected,
and has class rank two. Thus
\[
 b_1(U)=1,\qquad e(U)=3(\nu-2),\qquad
 b_1(D)=\nu-2,\qquad
          \operatorname{rank}I_U=2\nu-4.
\tag{C.RD.36}\label{cdp:eq:C.RD.36}
\]
Again compact transport forces rank zero,
hence $\nu=2$.
Equation (C.RD.25) then gives exactly one
$C$-preimage over the common value and no
multibranch point of $R_2$ away from it.
There is no loss of distinct attachment
values. The graph inequality yields
\[
                    e(S)\ge2-r_{\rm dom}\ge1.
\tag{C.RD.37}\label{cdp:eq:C.RD.37}
\]
This proves the degree-two-common-divisor
case and, together with \ref{cdp:st:C.2.RAT.7}--8 and
the degree bound at the start, every case
in the statement.
The additional weighted node/cusp charts,
their exact character increments and both
attachment budgets are checked in
\path{code/A_cubic_scalar_two.py} and its output.
\end{proof}






\begin{proposition}[Global exclusion of connected positive-genus triple covers over a rational image]\label{cdp:st:C.2.D4.1}
\StatementRecord{cdp:st:C.2.D4.1}{proved}{Historical events: Claude Fable 5.2 (Anthropic), effort high, 16 November 2026: carried as background without fault found. GPT-6 Astra (OpenAI), ChatGPT Pro mode, effort Pro, S6\_09\_PAPERS, 17 November 2026: targeted proof check. Also checked in this writing. Pre-continuation writing events: GPT-6 Astra (OpenAI), ChatGPT Pro mode, effort Pro, S6\_10\_PAPERS, 18 November 2026: targeted proof check; used and rechecked in this run. Continuation metadata events: GPT-6 Astra (OpenAI), ChatGPT Pro mode, effort Pro, S6\_10\_PAPERS continuation, 19 November 2026: used and rechecked in this run; status-vocabulary reconciliation; metadata reconciliation only; no new mathematical check. \emph{Pre-continuation recorded proof scope.} The actual normal-cone transport embeds the compact Milnor core into a nearby torus. Splitting degrees, connected cyclic triple-cover construction, boundary intersection rank, local genus and link budgets, and every positive-genus case were re-derived. The Newton and weighted-spectrum source loci were opened separately. Arnold1976 section1 I.1.0 p92 and section2 weighted-jet criteria4-9 p101 were subsequently opened and visually matched to the exact analytic ADE normal forms used here}{the hypotheses, notation and definitions stated here; \cref{cdp:eq:C.D4.2}; \cite[Theorem 2.1.1 and Corollary 2.1.3]{MerleTeissier1977}; \cite[pp. 165--167]{Steenbrink1989Spectrum}; \cite[\S1 and Example 2.4(a)]{Nemethi1991Join}; \cref{cdp:eq:C.D4.7}; \cref{cdp:eq:C.D4.6}; \cref{cdp:eq:C.D4.5}; \cref{cdp:eq:C.D4.9}; \cref{cdp:eq:C.D4.12}; \cref{cdp:eq:C.D4.4}}

Let $Y$ be a smooth compact complex threefold with
$b_2(Y;\Q)=0$, and let $f:Y\to B$ be a holomorphic map to a
smooth curve whose smooth nearby fibres are complex two-tori.
Fix a fibre $mS$, with $S$ reduced and irreducible, and restrict
the base to a small coordinate disc about it. Let
$H\subset S$ be a smooth rational curve along which the generic
transversal of $S$ is an ordinary triple point. Suppose that the
three generic normalization branches form one connected degree-three
cover $C\to H$, where $C$ denotes its smooth complete model.
Then $C$ is rational, and the cubic normal jet of $S$ vanishes
identically on at least one normal fibre over $H$.
In particular, a positive-genus connected degree-three cover is
impossible, and a primitive cubic normal jet is impossible for
any connected degree-three cover.
\end{proposition}


\begin{proof}
\noindent\textbf{Step 1.} We construct all comparison maps, including the map into the
actual nearby fibre. Put $N=N_{H/Y}$ and
$M=\OO_Y(S)|_H$. The ambient Chern-class hypothesis gives
$M\simeq\OO_H$ and $\deg N=-2$. The order-three normal jet is
therefore an actual homogeneous function
\[
 P\in H^0(H,\operatorname{Sym}^3N^*),\qquad
 U=\{v\in N:P(v)=1\}.
\tag{C.D4.1}\label{cdp:eq:C.D4.1}
\]
The generic cubic has three distinct roots. The fibrewise Euler
identity $dP(v)(v)=3P(v)$ proves that $U$ is smooth everywhere,
including above a base point where the cubic degenerates. Above
a point where the entire cubic is zero there is no point of $U$.

\paragraph{Comparison with the actual fibre.}
Let $t$ be the base coordinate, so $f$ is a holomorphic function
on a neighbourhood of $H$. Locally $f=u s^m$, where $s=0$
defines $S$ and $u$ is a unit. The degree-$3m$ normal jet of
$f$ is $P^m$, after rescaling the chosen trivialization of $M$
by a constant. Indeed the ratio of the two jets is an invertible
holomorphic function on $H\simeq\PP^1$, hence constant.
Take a smooth tubular map $\tau$ from a neighbourhood of the
zero section in $N$ to $Y$, whose derivative on normal fibres
is the specified complex-linear identification. For positive
real $s\to0$, Taylor's formula gives
\[
 s^{-3m}f\bigl(\tau(sv)\bigr)\longrightarrow P(v)^m
                 \quad\hbox{in }C^\infty\hbox{ on compact subsets of }N.
\tag{C.D4.2}\label{cdp:eq:C.D4.2}
\]
The order of vanishing is at least $3m$ along every point of
$H$: its ideal is a regular codimension-two ideal, and its
ordinary powers are primary along the smooth curve, so the
order verified at its generic point extends over $H$.

\noindent\textbf{Step 2.} Choose any compact codimension-zero smooth submanifold
$K\subset U$. On a neighbourhood of $K$, the differential of
$P^m$ has a smooth right inverse as a real map to $\C$.
The implicit-function theorem, uniformly on $K$, or the
corresponding right-inverse vector field, transports this
neighbourhood to the level
$s^{-3m}f(\tau(sv))=1$. Composing with $v\mapsto\tau(sv)$
gives an orientation-preserving embedding into the actual
fibre $f^{-1}(s^{3m})$. In particular, intersections of compact
two-cycles, computed with parallel perturbations inside $K$,
are preserved. This construction uses the actual function
$f$ and its leading normal jet; it makes no assertion that
arbitrary changes of local equations preserve a Milnor fibre.

\noindent\textbf{Step 3.} The algebraic surface $U$ admits a compact smooth core whose
inclusion is a homotopy equivalence. For example, resolve the
boundary in a projective compactification, triangulate the
resulting pair, take the complement of a sufficiently small
regular neighbourhood of the boundary and round its corners.
The radial collars of this regular neighbourhood give the
retraction. We choose this core for $K$. Thus the preceding
embedding induces one map carrying all of $H_1(U)$ and
$H_2(U)$, together with their compact intersection pairing.

\noindent\textbf{Step 4.} Here is the precise torus obstruction we will use. For a map
$K\to T^4$, let $W$ be the rational image of $H_1(K)$ in
$H_1(T^4)$. The map lifts to the covering corresponding to
its fundamental-group image, a space homotopy equivalent to
$T^r$, where $r=\dim W$. Its image on $H_2$ consequently lies
in $\Lambda^2W\subset\Lambda^2\Q^4$. For $r\le3$ the wedge
pairing on this image is zero, because $\Lambda^4W=0$.
Consequently a compact core with $b_1(U)\le3$ and nonzero
intersection pairing cannot embed by \eqref{cdp:eq:C.D4.2} into a nearby
torus. This is a comparison of the actual homology maps.

\paragraph{The normal bundle and the root divisor.}
Write
\[
                  N=\OO_H(a)\oplus\OO_H(-2-a),\qquad a\ge-1.
\tag{C.D4.3}\label{cdp:eq:C.D4.3}
\]
If $a>0$, the coefficient of the cube of the first fibre
coordinate vanishes, since it has degree $-3a$. The other
coordinate is then a global linear factor of the cubic.
This contradicts connectedness of the degree-three root
cover. Thus $a=-1$ or $a=0$.

\noindent\textbf{Step 5.} For $a=0$, irreducibility of the generic cubic forces the
coefficient of the first coordinate cubed to be a nonzero
constant. Along the trivial summand of $N$, the equation
$P=1$ therefore has three constant sections. Each is a
compact $\PP^1$ in $U$ with normal bundle $\OO_{\PP^1}(-2)$:
the derivative in the trivial summand is nonzero, so the
implicit-function theorem identifies its normal direction
with the other summand. Its homology square is $-2$.
Transporting its tubular neighbourhood by \eqref{cdp:eq:C.D4.2} gives a
sphere of square $-2$ in a torus, whereas the torus has
$\pi_2=0$. This is impossible.

\noindent\textbf{Step 6.} We are left with $N=\OO(-1)^{\oplus2}$ and
$E=\PP(N)\simeq\PP^1\times\PP^1$. Let $A$ be the common
zero divisor of the cubic coefficients on the base, of degree
$v$. Removing that scalar factor gives an integral horizontal
root divisor $R$ of bidegree $(3,3-v)$, where the first entry
is its degree over $H$. Its normalization is the given
cover $C$: at a general point the maps identify a normalization
branch with its distinct tangent line, and both complete
normal covers are the integral closure in this same function
field. Thus, writing $g=g(C)$,
\[
                       0\le g\le p_a(R)=4-2v.
\tag{C.D4.4}\label{cdp:eq:C.D4.4}
\]
For a primitive cubic $v=0$. For a positive-genus cover with
$v>0$, the bound forces $v=1$. We treat these two cases.
Their total root divisors are respectively
\[
 B=R_{(3,3)},\qquad
 B=R_{(3,2)}+F_{(0,1)}.
\tag{C.D4.5}\label{cdp:eq:C.D4.5}
\]
Both are reduced divisors of class $(3,3)$; the second has
one simple vertical component. They are connected. Let
$c=1$ or $2$ be their number of components, and set
$\nu=\sum_p(r_p(B)-1)$, where $r_p(B)$ is the number of
analytic branches at a singular point. Their total delta
invariants are $4-g$ and $5-g$, respectively, so in the
cases under consideration they are at most four.

\paragraph{The cyclic compactification and its local invariants.}
Put $L=\OO_E(1,1)$ and form the finite cyclic triple cover
\[
 X_0=\operatorname{Spec}_E
       (\OO_E\oplus L^{-1}\oplus L^{-2}),\qquad w^3=b,
\tag{C.D4.6}\label{cdp:eq:C.D4.6}
\]
where $b$ is the section of $L^3$ defining $B$. The algebra
multiplication is given by $b$. Because $B$ is reduced,
$X_0$ is normal: its hypersurface local rings are Cohen--Macaulay
and regular in codimension one. The cover is connected,
as a meridian around any smooth point of $B$ acts by a
three-cycle. Above $E\setminus B$ it is exactly $U$.
Indeed a line $[v]\in E\setminus B$ has three radial vectors
with $P(v)=1$; equivalently, the dual tautological coordinate
has cube $b$. This gives the claimed isomorphism, including
the transition functions.

\noindent\textbf{Step 7.} Let $\rho:X\to X_0$ resolve the surface singularities and
then the reduced boundary $D=X\setminus U$. Let $Q_p$ denote
the reduced exceptional divisor above a singular point of
$X_0$. We need the following finite computation:
\[
  \sum_p p_g(X_0,p)\le1,\qquad
  \sum_p b_1(Q_p;\Q)\le2.
\tag{C.D4.7}\label{cdp:eq:C.D4.7}
\]
Here are its equations and derivation. A reduced plane germ
with delta invariant at most four has one of the forms
\[
\begin{array}{c|c|c}
 \hbox{type}&r(x,y)&\hbox{range}\\ \hline
 A_n&x^2+y^{n+1}&1\le n\le8\\
 D_n&x^2y+y^{n-1}&4\le n\le7\\
 E_6&x^3+y^4&\\
 E_7&x^3+xy^3&\\
 E_8&x^3+y^5&
\end{array}
\]
For completeness, multiplicity at least four gives
$\delta\ge\binom42=6$. At multiplicity two, Weierstrass
preparation and completing the square give $A_n$, with
$\delta=\lfloor(n+1)/2\rfloor$. At multiplicity three,
three different tangent lines give $D_4$. A double tangent
and a transverse line give that smooth line times an
$A_k$ branch or pair, adding intersection multiplicity two;
the delta bound gives $D_5,D_6,D_7$. With a single tangent
line, an irreducible branch has initial Puiseux pair $(3,4)$
or $(3,5)$, giving $E_6,E_8$; a smooth branch and a cusp
have total delta at least $1+3=4$, and equality gives $E_7$.
Three smooth branches with one tangent line already give
delta at least six. These reductions give the displayed
analytic forms; the higher terms in the last three forms
are eliminated by their positive weighted Jacobian terms.

\noindent\textbf{Step 8.} The singularity of $X_0$ is $w^3=r(x,y)$. Normalize the
weights of $r$ to degree one and give $w$ weight $1/3$.
The adjunction calculation of
\cite[Theorem 2.1.1 and Corollary 2.1.3]{MerleTeissier1977}
counts $p_g$ by the positive lattice points under the
weight-one plane. Equivalently, it counts monomials
$x^iy^jw^k$ with
$(i+1)w_x+(j+1)w_y+(k+1)/3\le1$.
The normalized monodromy weights are those of the Jacobian
monomials, shifted by the sum of the variable weights
\cite[pp. 165--167]{Steenbrink1989Spectrum}.
The following table records the result; all entries are
also computed with rational arithmetic in
\path{code/root_cubic_normal_cone.py}.
\[
\begin{array}{c|c|c|c}
 \hbox{plane type}&\delta&p_g(w^3-r)&b_1(Q_p;\Q)\\ \hline
 A_1,A_2&1&0&0\\
 A_3,A_4&2&0&0\\
 A_5&3&1&2\\
 A_6&3&1&0\\
 A_7,A_8&4&1&0\\
 D_4&3&1&2\\
 D_5&3&1&0\\
 D_6,D_7&4&1&0\\
 E_6&3&1&0\\
 E_7,E_8&4&1&0
\end{array}
\]
To verify the last column directly, use the Jacobian bases
$\{y^j:0\le j<n\}$ for $A_n$;
$\{1,y,\ldots,y^{n-2},x\}$ for $D_n$;
$\{x^iy^j:0\le i<2,0\le j<3\}$ for $E_6$;
$\{1,y,\ldots,y^4,x,xy\}$ for $E_7$; and
$\{x^iy^j:0\le i<2,0\le j<4\}$ for $E_8$.
Count weights $(i+1)w_x+(j+1)w_y$ congruent to $1/3$ or
$2/3$ modulo one. The tensor-product monodromy for adding
$w^3$ identifies this count with the eigenvalue-one
multiplicity of the surface monodromy
\cite[\S1 and Example 2.4(a)]{Nemethi1991Join}.
The variation sequence of its simply connected Milnor fibre
identifies that multiplicity with the rational first Betti
number of its link. The negative-definite exceptional
intersection matrix then identifies $H_1$ of the link with
$H_1(Q_p)$: in the homology sequence of its plumbed
neighbourhood, the map on second homology is that invertible
matrix. This proves the column without replacing the actual
exceptional divisor by an unspecified graph.
Every nonzero entry in either last column costs delta at
least three. Since the total delta is at most four,
\eqref{cdp:eq:C.D4.7} follows.

\paragraph{Homology and the compact intersection pairing.}
The finite algebra in \eqref{cdp:eq:C.D4.6} and the K\"unneth formula give
$H^1(X_0,\OO_{X_0})=0$: the three summands on $E$ are
$\OO$, $\OO(-1,-1)$ and $\OO(-2,-2)$, all with $H^1=0$.
The resolution Leray sequence injects $H^1(X,\OO_X)$ into
$H^0(X_0,R^1\rho_*\OO_X)$, of dimension
$\sum p_g(X_0,p)$. Since $X$ is projective, \eqref{cdp:eq:C.D4.7} gives
\[
                         b_1(X;\Q)=2q(X)\le2.
\tag{C.D4.8}\label{cdp:eq:C.D4.8}
\]
The irreducible components of $D$ have independent homology
classes. The exceptional classes have a negative-definite
intersection matrix. In the first case of \eqref{cdp:eq:C.D4.5} the
remaining class pushes forward to the nonzero class of $R$;
in the second, the two remaining classes push forward to
$(3,2)$ and $(0,1)$, which are independent on $E$.
Moreover their span contains the pullback of the ample
class $B=(3,3)$, whose square is positive. The Hodge index
theorem implies that the intersection form on this span is
nondegenerate: its part orthogonal to the positive pullback
class is negative definite.

\noindent\textbf{Step 9.} Write $n_D$ for the number of boundary components. The exact
homology sequence of the pair $(X,U)$ identifies the map
$H_2(X;\Q)\to H_2(X,U;\Q)=\Q^{n_D}$ with intersection
against these components. Its surjectivity gives
$H_1(U;\Q)\simeq H_1(X;\Q)$. Its kernel is the orthogonal
complement of their span, and is the image of $H_2(U)$.
The pairing on that image is nondegenerate, of dimension
$b_2(X)-n_D$. As $D$ is connected,
$e(D)=n_D+1-b_1(D)$, and consequently
\[
\begin{split}
 b_1(U)&=b_1(X)\le2,\\
 \operatorname{rank}\bigl(\hbox{intersection pairing on }H_2(U)\bigr)
     &=b_2(X)-n_D\\
     &=e(U)-1+2b_1(X)-b_1(D).
\end{split}
\tag{C.D4.9}\label{cdp:eq:C.D4.9}
\]
Here the equality of ranks follows also for a degenerate
pairing on $H_2(U)$: its kernel includes the kernel of the
map to $H_2(X)$, and the induced form on its image is the
nondegenerate form just identified.

\noindent\textbf{Step 10.} There is a useful exact version allowing relations among
the strict boundary classes. If there are $\varepsilon$
independent relations, while their span still contains a
positive-square class, the same pair sequence and the same
Hodge-index argument give
\[
\begin{split}
 b_1(U)&=b_1(X)+\varepsilon,\\
 \operatorname{rank}\bigl(\hbox{intersection pairing on }H_2(U)\bigr)
 &=e(U)-1+2b_1(X)-b_1(D)+\varepsilon.
\end{split}
\tag{C.D4.10}\label{cdp:eq:C.D4.10}
\]
These identities concern the actual classes in the indicated
cyclic cover. They will also be used for the degree-two
root divisor below.

\noindent\textbf{Step 11.} To compute $b_1(D)$, start with the normalizations of the
$c$ components of $B$, of total genus $g$, and attach the
connected exceptional divisor over each singular point.
There are $r_p(B)$ attaching branches there. The incidence
graph contributes $\nu-c+1$ cycles. Subsequent point
blowups resolving the boundary subdivide edges or add trees
and preserve this first Betti number. Thus \eqref{cdp:eq:C.D4.7} gives
\[
                         b_1(D)\le2g+\nu-c+3.
\tag{C.D4.11}\label{cdp:eq:C.D4.11}
\]
The unramified cover $U\to E\setminus B$ has degree three,
so additivity and multiplicativity of Euler characteristic
give the exact value
\[
                         e(U)=3(4-2c+2g+\nu).
\tag{C.D4.12}\label{cdp:eq:C.D4.12}
\]
For $c=1$, \eqref{cdp:eq:C.D4.9}--\eqref{cdp:eq:C.D4.12} give intersection rank at
least $3+4g+2\nu>0$. For $c=2$ they give rank at least
$4g+2\nu-2$. In this second case $g\ge1$, and $R$ meets
$F$ with intersection number three, so $\nu\ge1$ and the
rank is again positive. Both cases have $b_1(U)\le2$.
They contradict the compact-core comparison with $T^4$
proved immediately after \eqref{cdp:eq:C.D4.2}.

\noindent\textbf{Step 12.} This excludes $v=0$ for every connected triple cover and
$v=1$ for every positive-genus one. By \eqref{cdp:eq:C.D4.4} these are
all possibilities for a positive-genus cover, proving the
statement. The genus-zero cases with $v=1$ or $2$ are not
asserted to be excluded by this argument.
\end{proof}






\begin{proposition}[A rational ordinary-triple image has no positive-genus double branch cover]\label{cdp:st:C.2.D4.2}
\StatementRecord{cdp:st:C.2.D4.2}{proved}{Historical events: Claude Fable 5.2 (Anthropic), effort high, 16 November 2026: carried as background without fault found. GPT-6 Astra (OpenAI), ChatGPT Pro mode, effort Pro, S6\_09\_PAPERS, 17 November 2026: targeted proof check. Also checked in this writing. Pre-continuation writing events: GPT-6 Astra (OpenAI), ChatGPT Pro mode, effort Pro, S6\_10\_PAPERS, 18 November 2026: targeted proof check; used and rechecked in this run. Continuation metadata events: GPT-6 Astra (OpenAI), ChatGPT Pro mode, effort Pro, S6\_10\_PAPERS continuation, 19 November 2026: used and rechecked in this run; status-vocabulary reconciliation; metadata reconciliation only; no new mathematical check. \emph{Pre-continuation recorded proof scope.} The factor degrees, normalized eigensheaves, boundary relation defect and intersection-rank formula were re-derived. The degree-two common-fibre case was checked separately for distinct fibres and a double fibre, including the exact w-squared-over-s generator and each normalized-cover blowup. All positive-genus cases produce a nonzero intersection pairing incompatible with the transported torus core}{the hypotheses, notation and definitions stated here; \cref{cdp:st:C.2.D4.1}; \cref{cdp:st:A.5.19}; \cref{cdp:eq:C.DD.3}; \cref{cdp:eq:C.DD.6}; \cref{cdp:eq:C.DD.2}; \cref{cdp:st:C.2.RAT.6}; \cref{cdp:eq:C.DD.9}; \cref{cdp:eq:C.DD.10}; \cref{cdp:eq:C.DD.8}; \cref{cdp:eq:C.DD.4}; \cref{cdp:eq:C.DD.16}; \cref{cdp:eq:C.DD.7}; \cref{cdp:eq:C.DD.20}; \cref{cdp:eq:C.DD.5}}

Let $f:Y\to C_0$ have the torus-fibre and ambient hypotheses
of \ref{cdp:st:C.2.D4.1}. Let $H\subset S$ be a smooth rational curve
with generic ordinary-triple transversal. Suppose two of
the three generic normalization branches form a connected
degree-two cover $C\to H$, and the remaining branch forms
a degree-one cover. Then the smooth complete curve $C$
is rational. The assertion includes all zeros, with their
multiplicities, of the entire cubic normal tensor.

Consequently, in the irrational ruled setup of \ref{cdp:st:A.5.19},
a noncontracted coefficient-two horizontal conductor
with generic ordinary-triple transversal has smooth
elliptic image. Its normalization is smooth minimal
elliptic ruled, the horizontal cover is unramified of
degree three, and $e(S)=0$.
\end{proposition}


\begin{proof}
\noindent\textbf{Step 1.} Suppose $g=g(C)\ge1$. Use the actual normal tensor and
normal-cone comparison (C.D4.1)--(C.D4.2), and write
\[
 N=N_{H/Y}=\OO(a)\oplus\OO(-2-a),\quad a\ge-1,
 \qquad P\in H^0(H,\operatorname{Sym}^3N^*),
 \qquad U=\{P=1\}\subset N.
\tag{C.DD.1}\label{cdp:eq:C.DD.1}
\]
The comparison embeds a compact core of $U$, including
its compact intersection pairing, into an actual nearby
torus. As proved there, $b_1(U)\le3$ and a nonzero
intersection pairing contradict this embedding.

\noindent\textbf{Step 2.} Put $E=\PP(N)$, using lines in $N$, and let
$L=\OO_E(1)$, $\xi=c_1(L)$ and $F$ be the fibre class.
These conventions give
\[
 \xi^2=-\deg N=2,\qquad \xi F=1,\quad F^2=0,
 \qquad K_E=-2\xi.
\tag{C.DD.2}\label{cdp:eq:C.DD.2}
\]
Let $A$ be the common zero divisor of the coefficients
of $P$, and put $v=\deg A$. The divisor of the induced
section of $L^3$ is
\[
                    B_{\rm tot}=R_1+R_2+p^*A.
\tag{C.DD.3}\label{cdp:eq:C.DD.3}
\]
Here $R_1$ is a section and $R_2$ is an integral
degree-two curve over $H$. The section is smooth:
its rational tangent-root map extends over the smooth
complete base into the proper bundle $E$.
The normalization of $R_2$ is precisely $C$. Over the
generic point a normalization branch is identified
with its distinct tangent root; taking the integral
closure in this same quadratic function field identifies
the two complete normal covers, as in (C.RD.2).

\paragraph{All possible common divisors.}
For $a\ge0$ and $v>0$, the pure cube of the coordinate
in $\OO(a)$ has coefficient of degree $-3a-v$, and
therefore vanishes. The other coordinate is a factor.
The quadratic factor is irreducible over the function
field. Its coefficient of the first coordinate squared
must be nonzero; that coefficient has degree
\[
                              2-a-v.
\tag{C.DD.4}\label{cdp:eq:C.DD.4}
\]
Thus $a+v\le2$. In particular, $v\le2$, and $v=2$
forces $a=0$.

\noindent\textbf{Step 3.} For $a=-1$, identify $E$ with the product of the base
and the projective direction line. Write $d\ge0$ for
the degree of the map defining the section $R_1$.
The bidegrees, with degree over $H$ first, are
\[
 R_1:(1,d),\qquad R_2:(2,3-v-d),\qquad
 p_a(R_2)=2-v-d.
\tag{C.DD.5}\label{cdp:eq:C.DD.5}
\]
The inequality $g\le p_a(R_2)$ gives $v+d\le1$.
Consequently the only common-divisor possibilities
for our positive-genus cover are $v=0$, $v=1$, and
$v=2$ with $N=\OO\oplus\OO(-2)$.

\paragraph{The exact cyclic-cover homology formulas.}
We record the comparison with relations and disconnected
boundary retained. First suppose that \eqref{cdp:eq:C.DD.3} is reduced,
and let $B=B_{\rm tot}$. The cyclic cover
\[
 \pi:X_0=\operatorname{Spec}_E
       (\OO_E\oplus L^{-1}\oplus L^{-2})\longrightarrow E
\tag{C.DD.6}\label{cdp:eq:C.DD.6}
\]
has multiplication given by $P$. It is normal, since
the branch divisor is reduced. Its complement over
$B$ is exactly $U$: if $w\in\ell^*$ is a nonzero
cube root of $P|_\ell$, its corresponding vector in
the line $\ell\subset N$ is the unique vector $z$
with $w(z)=1$. This description specifies inverse
maps and their bundle transitions.

\noindent\textbf{Step 4.} All three summands in \eqref{cdp:eq:C.DD.6} have vanishing $H^1$.
Indeed the relative fibre cohomology of $L^{-1}$
vanishes, while relative duality gives
\[
 p_*L^{-2}=0,\qquad R^1p_*L^{-2}=\det N=\OO(-2).
\tag{C.DD.7}\label{cdp:eq:C.DD.7}
\]
Thus $H^1(E,L^{-2})=H^0(H,\OO(-2))=0$;
$H^1(E,\OO_E)=0$ as well. For a resolution
$\rho:X\to X_0$, including resolution of the reduced
boundary $D=X\setminus U$, Leray therefore gives
\[
 b_1(X)=2q(X)\le2\sum_x p_g(X_0,x).
\tag{C.DD.8}\label{cdp:eq:C.DD.8}
\]
The resolution is projective, so the equality here
is Hodge decomposition on the actual surface $X$.

\noindent\textbf{Step 5.} Write $c$ for the number of components of $B$,
$b_0=b_0(B)$, and
$\nu=\sum_x(r_x(B)-1)$. Its normalized components
have total genus $g$. Let $r$ be the rank of their
classes in $H_2(E,\Q)$ and put $\varepsilon=c-r$.
The boundary classes on $X$ have exactly
$\varepsilon$ relations. To see the equality, modulo
the independent exceptional classes, the strict
component over $B_i$ is $\pi^*[B_i]/3$.
Indeed its ramification index is three, and the
remaining terms in the divisor pullback are exceptional.
The pullback of $H_2(E,\Q)$ is injective and its
intersection pairing is multiplied by three.
This proves both directions of the asserted relation
count. The span of the boundary classes is
nondegenerate: it contains the positive-square
pullback of $B_{\rm tot}=3\xi$, and the Hodge index
theorem makes the part orthogonal to that class
negative definite.

\noindent\textbf{Step 6.} The pair sequence, exactly as in (C.D4.9)--(C.D4.10),
now gives the version with $b_0$ retained:
\[
\begin{split}
 b_1(U)&=b_1(X)+\varepsilon,\\
 \operatorname{rank}\langle\ ,\ \rangle_U
 &=e(U)+b_0-2+2b_1(X)-b_1(D)+\varepsilon.
\end{split}
\tag{C.DD.9}\label{cdp:eq:C.DD.9}
\]
The second identity follows from
$e(D)=b_0+n_D-b_1(D)$ and the nondegenerate
orthogonal complement of the boundary span. Here
$\langle\ ,\ \rangle_U$ is the compact intersection
pairing, so it is precisely the form preserved by
the normal-cone embedding.

\noindent\textbf{Step 7.} If $\lambda$ is the sum of the first Betti numbers
of the reduced exceptional divisors over the original
surface singularities, normalization of $B$ and the
attaching-branch graph give
\[
 b_1(D)\le2g+\nu-c+b_0+\lambda,
 \qquad e(U)=3(4-2c+2g+\nu).
\tag{C.DD.10}\label{cdp:eq:C.DD.10}
\]
Each exceptional divisor is connected. Its attachments
number $r_x(B)$ at a singular point, giving graph
cycle rank $\nu-c+b_0$. Further boundary blowups
subdivide edges or add trees. The Euler identity
uses the actual unramified degree-three cover
$U\to E\setminus B$ and
$e(B)=2c-2g-\nu$.

\paragraph{No common zero.}
For $v=0$, the divisor $B=R_1+R_2$ is reduced,
has $c=2$, and has arithmetic genus four by
\eqref{cdp:eq:C.DD.2}. Its total delta invariant is therefore
\[
                         \delta(B)=5-g\le4.
\tag{C.DD.11}\label{cdp:eq:C.DD.11}
\]
The local calculation (C.D4.7), or \ref{cdp:st:C.2.RAT.6},
gives $\sum p_g\le1$ and $\lambda\le2$.
The rank $r$ is one or two, so
$\varepsilon=2-r\le1$. Thus $b_1(U)\le3$.
Substitution in \eqref{cdp:eq:C.DD.9}--\eqref{cdp:eq:C.DD.10} gives
\[
                  \operatorname{rank}\langle\ ,\ \rangle_U
                        \ge4g+2\nu-r>0.
\tag{C.DD.12}\label{cdp:eq:C.DD.12}
\]
This calculation includes a disconnected $B$ and
dependent section classes. It contradicts the
normal-cone comparison.

\paragraph{One simple common zero.}
For $v=1$, one has $B=R_1+R_2+F_0$, still reduced.
It is connected because the fibre meets both
horizontal components. Here $c=3$, $b_0=1$,
$\nu\ge2$, and
$\delta(B)=6-g\le5$.
Statement \ref{cdp:st:C.2.RAT.6} supplies
$\sum p_g\le1$ and $\lambda\le2$.
The section and the fibre have independent classes,
so $r=2$ and $\varepsilon=1$.
Equations \eqref{cdp:eq:C.DD.8}--\eqref{cdp:eq:C.DD.10} yield
\[
 b_1(U)\le3,\qquad
 \operatorname{rank}\langle\ ,\ \rangle_U
                         \ge4g+2\nu-6\ge2.
\tag{C.DD.13}\label{cdp:eq:C.DD.13}
\]
This is the same contradiction.

\paragraph{The exact form when the common divisor has degree two.}
Now $N=\OO\oplus\OO(-2)$ by \eqref{cdp:eq:C.DD.4}.
Let $x$ and $y$ be its projections, valued in
$\OO$ and $\OO(-2)$ respectively. The coefficient
of $x^2$ in the irreducible quadratic is a nonzero
constant. Completing its square is a global bundle
automorphism: the mixed coefficient is a section
of $\OO(2)$ and therefore defines an actual map
$\OO(-2)\to\OO$. After a nonzero constant rescaling,
the tensor has the form
\[
 P=a\,y\,(x^2-b y^2),\qquad
 a\in H^0(H,\OO(2)),\quad b\in H^0(H,\OO(4)).
\tag{C.DD.14}\label{cdp:eq:C.DD.14}
\]
Its horizontal root divisors are
$R_1=\{y=0\}$ and $R_2=\{x^2-b y^2=0\}$.
They are disjoint. The quadratic cover has branch
degree at most four. Riemann--Hurwitz and $g\ge1$
force $g=1$, branch degree four, and four distinct
simple zeros of $b$. Thus $R_2$ is already a smooth
elliptic curve. A ruling fibre meets it either at
two transverse points or at a single point of
contact two. This verifies the local models used
in both remaining cases.

\noindent\textbf{Step 8.} If $p^*A=F_1+F_2$ with distinct fibres, the reduced
branch divisor has four components. Its only plane
singularities are ordinary nodes and tangencies of
order two, of types $A_1$ and $A_3$.
Their cubic suspensions have geometric genus zero
and exceptional first Betti number zero, by
(C.RD.10). Consequently $b_1(X)=0$.
The branch divisor is connected, its class rank
is two, and $4\le\nu\le6$.
Thus $\varepsilon=2$, $b_1(D)=\nu-1$, and
\[
 b_1(U)=2,\qquad
 \operatorname{rank}\langle\ ,\ \rangle_U
                       =2\nu-4\ge4.
\tag{C.DD.15}\label{cdp:eq:C.DD.15}
\]
Again the torus comparison is impossible.

\paragraph{A double common zero: normalization of the cyclic cover.}
It remains to take $p^*A=2F_0$. The branch divisor
with multiplicities is $R_1+R_2+2F_0$; its reduced
support $B$ has three components. We use the
normalization of the cyclic cover. Its exact
eigensheaf algebra has underlying module
\[
       \pi_*\OO_{X_0}
       =\OO_E\oplus L^{-1}\oplus L^{-2}(F_0).
\tag{C.DD.16}\label{cdp:eq:C.DD.16}
\]
At the generic point of $F_0$, write the equation
as $w^3=s^2u$ with $u$ a unit. The extra integral
generator is $w^2/s$. This gives the last summand
in \eqref{cdp:eq:C.DD.16}. At the other branch components the
usual generators $1,w,w^2$ apply. These comparisons
extend across the finitely many crossings: the
normal finite surface is Cohen--Macaulay, so its
finite direct image over the smooth surface $E$
is locally free, and each character summand is
the rank-one reflexive extension of its divisorwise
description. This proves \eqref{cdp:eq:C.DD.16} as a sheaf
identity, not only over the nonsingular locus.

\noindent\textbf{Step 9.} Relative duality as in \eqref{cdp:eq:C.DD.7} now gives
\[
 R^1p_*\bigl(L^{-2}(F_0)\bigr)
                       =\det N\otimes\OO(1)=\OO(-1).
\tag{C.DD.17}\label{cdp:eq:C.DD.17}
\]
Thus every summand in \eqref{cdp:eq:C.DD.16} has vanishing
$H^1$, and $H^1(X_0,\OO_{X_0})=0$.
We next prove that all singularities of this
normalized cover are rational.

\noindent\textbf{Step 10.} The only intersections are between the weight-two
branch $F_0$ and a weight-one smooth branch. Their
contact is one or two. The following explicit
resolution works for every contact $k\ge1$.
In coordinates the branch equations are
\[
                F_0=(s=0),\qquad R=(s-u^k=0),
                \qquad w^3=s^2(s-u^k).
\tag{C.DD.18}\label{cdp:eq:C.DD.18}
\]
Blow up their common point until the strict
branches separate, requiring $k$ blowups. At
each step their total branch multiplicity is
$2+1=3$. The exceptional divisor is therefore
unbranched after normalization, and the new
root line and weight-two divisor are exactly
\[
 L'=\sigma^*L(-E_0),\qquad
 F'=\sigma^*F-E_0.
\tag{C.DD.19}\label{cdp:eq:C.DD.19}
\]
For example $s=u s_1$, $w=u w_1$ changes the
equation to
$w_1^3=s_1^2(s_1-u^{k-1})$; the extra normalized
generator changes by
$w^2/s=u(w_1^2/s_1)$.
Both nontrivial eigensheaves consequently satisfy
the exact identities
\[
 (L')^{-1}=\sigma^*(L^{-1})(E_0),\qquad
 (L')^{-2}(F')=\sigma^*(L^{-2}(F))(E_0).
\tag{C.DD.20}\label{cdp:eq:C.DD.20}
\]
For a point blowup of a smooth surface,
$\sigma_*\OO(E_0)=\OO$ and
$R^1\sigma_*\OO(E_0)=0$.
The first follows by extension of functions across
the centre; the second follows from
$0\to\OO\to\OO(E_0)\to\OO_{\PP^1}(-1)\to0$
and the structure-sheaf direct images of the
point blowup. Projection formula applied to
\eqref{cdp:eq:C.DD.20} therefore makes $R^1$ of the finite
cyclic algebra vanish at every step.

\noindent\textbf{Step 11.} For precision, the normalized covers over the
successive blown-up bases have proper birational
maps to $X_0$, induced by their common function
field. Finite direct image and Leray identify the
pushforward of the $R^1$ of each such map with
the just computed $R^1$ of its eigensheaf algebra.
The finite pushforward is faithful, so that $R^1$
is zero. At the end the two branch curves are
smooth and disjoint, and the normalized cover
is smooth: a single branch of multiplicity one
or two is parameterized by $s=t^3$, with $w=t$
or $w=t^2$, up to a unit root. This is a resolution
with $R^1\OO=0$ and proves the claimed rationality.
Its connected exceptional divisors have rational
components and tree incidence graphs: their
$H^1(\OO)$ is a quotient of that $R^1$ and is zero.
Hence their first rational Betti numbers are zero.

\noindent\textbf{Step 12.} In our global double-fibre case it follows that
$b_1(X)=0$ and $\lambda=0$.
The comparison $U\to E\setminus B$ is still an
unramified degree-three cover, since normalization
does not change it off the branch divisor.
Each reduced branch component still has one
strict transform with ramification index three,
so the relation and intersection calculations
\eqref{cdp:eq:C.DD.9}--\eqref{cdp:eq:C.DD.10} apply to this reduced support.
It is connected, $c=3$, its class rank is two,
and $\nu=2$ or $3$. Thus
\[
 b_1(U)=1,\qquad b_1(D)=\nu,\qquad e(U)=3\nu,
 \qquad\operatorname{rank}\langle\ ,\ \rangle_U
                                      =2\nu>0.
\tag{C.DD.21}\label{cdp:eq:C.DD.21}
\]
This is the final contradiction. Every common
divisor allowed by \eqref{cdp:eq:C.DD.4}--\eqref{cdp:eq:C.DD.5} has now
been treated, proving that $C$ cannot have
positive genus.

\noindent\textbf{Step 13.} For the stated irrational ruled consequence,
\ref{cdp:st:A.5.19} leaves a smooth rational or smooth elliptic
image. A rational image would have horizontal
degree two or three. The preceding proof excludes
the first, and \ref{cdp:st:C.2.D4.1} excludes the second.
The elliptic conclusion, its exact minimality
comparison, and its Euler number are those already
proved in \ref{cdp:st:A.5.19}.
\end{proof}



The splitting bounds, all intersection-rank lower
bounds, the exact scalar-zero coordinate example,
and the normalized cyclic-cover blowup substitutions
are retained in \path{code/E_cubic_double.py} and
its saved output. The local delta-five table is
independently computed in \path{code/A_cubic_delta5.py}.




\begin{proposition}[Rational connected triple covers have positive fibre Euler number]\label{cdp:st:C.2.D4.3}
\StatementRecord{cdp:st:C.2.D4.3}{proved}{Historical events: Claude Fable 5.2 (Anthropic), effort high, 16 November 2026: carried as background without fault found. GPT-6 Astra (OpenAI), ChatGPT Pro mode, effort Pro, S6\_09\_PAPERS, 17 November 2026: targeted proof check. Also checked in this writing. Pre-continuation writing events: GPT-6 Astra (OpenAI), ChatGPT Pro mode, effort Pro, S6\_10\_PAPERS, 18 November 2026: targeted proof check. Continuation metadata events: GPT-6 Astra (OpenAI), ChatGPT Pro mode, effort Pro, S6\_10\_PAPERS continuation, 19 November 2026: used and rechecked in this run; status-vocabulary reconciliation; metadata reconciliation only; no new mathematical check. \emph{Pre-continuation recorded proof scope.} The one-simple-zero case has independent branch classes and forces nu=1. The two-simple-zero case has exactly one boundary relation; canonical duality gives q=max(P-1,0), and the full contact partitions leave at most one duplicate attachment. For the double zero, the normalized character algebra and all weight-(2,1) contact blowups have vanishing R1; its positive compact pairing contradicts the actual torus embedding. The three cases exhaust the common cubic divisor and prove positive Euler number}{the hypotheses, notation and definitions stated here; \cref{cdp:st:C.2.RAT.1}; \cref{cdp:st:C.2.D4.1}; \cref{cdp:st:C.2.RAT.3}; \cref{cdp:st:C.2.RAT.5}; \cref{cdp:st:C.2.RAT.6}; \cref{cdp:eq:C.D4.13}; \cref{cdp:eq:C.D4.19}}

Retain the torus-fibre, rational-normalization and specified
Hirzebruch-boundary morphism of \ref{cdp:st:C.2.RAT.1}. Suppose the horizontal
section $C$ is not contracted, its image $H$ is smooth rational,
the generic transversal along $H$ is an ordinary triple point,
and $C\to H$ has degree three. Then $e(S)\ge1$.
\end{proposition}


\begin{proof}
\noindent\textbf{Step 1.} Use the normal-cone construction of \ref{cdp:st:C.2.D4.1}. That result forces
$N=\OO(-1)^{\oplus2}$ and a common cubic-coefficient divisor
of degree $v=1$ or $2$. The horizontal root curve is integral
of bidegree $(3,3-v)$, with rational normalization $C$.
No vertical strict conductor component dominates $H$, since
all three generic branches are supplied by $C$.

\noindent\textbf{Step 2.} We first recall exactly how a finite attachment bound gives
the Euler inequality. In the notation of \ref{cdp:st:C.2.RAT.1} let $r$
be the number of original positive ruling fibres and let
$N_0$ count coefficient-zero blowup events. At least
$K\ge r-N_0$ of their positive boundary trees contain a
strict conductor curve. The negative-definite argument of
\ref{cdp:st:C.2.RAT.3} proves this bound without an assumption about
normalization singularities. They yield $K$ distinct
attachment points on $C$. Their images on $H$ give $k$
distinct normalization-incidence edges. Put $\ell=K-k$.
The incidence graph has first Betti number at most one,
so the exact graph count (C.RA.24) gives
$e(\Gamma_{\rm red})\ge k$. Consequently
\[
             e(S)=2-r+N_0+e(\Gamma_{\rm red})\ge2-\ell.
\tag{C.D4.13}\label{cdp:eq:C.D4.13}
\]

\noindent\textbf{Step 3.} At a point where the cubic is primitive, a simple tangent
root gives a smooth normalization branch with conductor
exactly twice the horizontal curve. Indeed the ambient
blowup chart divides the reduced equation by the third
power of its normal coordinate; the nonzero root derivative
solves the remaining coordinate holomorphically. The
remaining quadratic factor has nonzero value at the chosen
simple root. Its restriction therefore gives the conductor
factor $y^2$ times a unit, even when its own two roots coincide. This is the explicit
factorization and conductor comparison of \ref{cdp:st:C.2.RAT.5}.
Such a point cannot carry a vertical boundary attachment.
If the horizontal root curve is locally unibranch, a
multiple root has just one point on its normalization.
A cubic has at most one multiple root unless all its roots
are simple. Thus off the common divisor each image value
receives at most one attachment when that root curve is
locally unibranch.

\paragraph{One simple common zero.}
Here the total reduced branch divisor is
$B=R_{(3,2)}+F_{(0,1)}$, with two rational normalizations
and total delta five. The complete computation \ref{cdp:st:C.2.RAT.6}
gives $P:=\sum p_g\le1$ and
$J:=\sum b_1(Q_p)\le2$. The classes of $R,F$ are
independent. As in (C.D4.8)--(C.D4.12), writing
$\nu=\sum(r_p(B)-1)$, we obtain
\[
 b_1(U)\le2,\qquad e(U)=3\nu,\qquad
 \operatorname{rank}(\hbox{intersection pairing})
                         =2\nu+2b_1(X)-J.
\tag{C.D4.14}\label{cdp:eq:C.D4.14}
\]
The torus comparison excludes $\nu\ge2$. Since $R\cdot F=3$,
we have $\nu\ge1$, so $\nu=1$. All branches of $R$ are
then locally unibranch: a multibranch point would contribute
to $\nu$ in addition to the intersection with $F$, or would
raise its contribution at that intersection to at least two.
There is exactly one point of $R\cap F$, and it has just
one preimage on $C$. Thus the common value also receives
at most one attachment. The preceding simple-root argument
handles every other value. Hence $\ell=0$ and \eqref{cdp:eq:C.D4.13}
gives $e(S)\ge2$.

\paragraph{Two distinct common zeros.}
Now $B=R_{(3,1)}+F_1+F_2$. The curve $R$ is smooth rational,
and each $R\cap F_i$ has total intersection number three.
A contact of order $k=1,2,3$ is the plane type $A_{2k-1}$.
Let $P$ be the number of order-three contacts. The local
table gives precisely
\[
              \sum p_g=P,\qquad \sum b_1(Q_p)=2P,
                         \qquad 0\le P\le2.
\tag{C.D4.15}\label{cdp:eq:C.D4.15}
\]
The cyclic cover $X_0$ has $\omega_{X_0}\simeq\OO_{X_0}$
and $\chi(\OO_{X_0})=2$, directly from its cyclic algebra
and adjunction. Proper duality identifies the quotient
$\omega_{X_0}/\rho_*\omega_X$ with a finite-length dual
module of length $P$. If $P>0$, a generator of the trivial
canonical line does not belong to the proper adjunction
ideal at at least one point. Thus no nonzero global
canonical section extends to $X$, and $p_g(X)=0$.
Leray gives $\chi(\OO_X)=2-P$. If $P=0$, the canonical
section extends and the same calculation gives $q(X)=0$.
Together these identities yield
\[
                         q(X)=\max(P-1,0).
\tag{C.D4.16}\label{cdp:eq:C.D4.16}
\]
There is exactly one relation among the three strict
boundary classes: $F_1,F_2$ have the same class, while
$R$ and a fibre are independent. The exceptional matrix
is negative definite, so no further relation occurs.
Thus (C.D4.10) gives $b_1(U)=2q(X)+1\le3$. With
$\nu=\#(R\cap(F_1\cup F_2))$, the boundary topology
and Euler number give
\[
 \operatorname{rank}(\hbox{intersection pairing})
          =2\nu-4-2P+4\max(P-1,0).
\tag{C.D4.17}\label{cdp:eq:C.D4.17}
\]
We check every possibility. For $P=0$ each fibre has at
least two distinct intersections, so $\nu\ge4$ and the
rank is positive, which is excluded. For $P=1$, one fibre
has a single order-three contact. The other has either
three simple intersections or one contact of order two
and one simple intersection. The first has $\nu=4$ and
positive rank; the only remaining value is $\nu=3$.
For $P=2$, each fibre has one point of intersection and
$\nu=2$. These are all the cases.

\noindent\textbf{Step 4.} Since $R$ is smooth, the simple-root argument gives at
most one attachment at any noncommon value. A common
value has as many points on $C$ as the corresponding
intersection with $F_i$. In the remaining $P=1$ case
these numbers are one and two, so $\ell\le1$. In the
$P=2$ case they are both one, so $\ell=0$.
Equation \eqref{cdp:eq:C.D4.13} gives $e(S)\ge1$ and $e(S)\ge2$,
respectively.

\paragraph{A double common zero.}
Finally the cubic divisor is $R_{(3,1)}+2F$. The normalized
cyclic cover has the exact eigensheaf algebra
\[
 \pi_*\OO_{X_0}
       =\OO_E\oplus L^{-1}\oplus L^{-2}(F),
       \qquad L=\OO_E(1,1).
\tag{C.D4.18}\label{cdp:eq:C.D4.18}
\]
One obtains the last summand by dividing the square of the
cyclic coordinate by a local equation of the weight-two
branch; this integral element satisfies its monic cubic
relation. These summands describe the normalization in
codimension one and hence everywhere by normality. Their
$H^1$ groups vanish, since their bidegrees are $(0,0)$,
$(-1,-1)$ and $(-2,-1)$.

\noindent\textbf{Step 5.} We verify that the normalization singularities are rational.
At a contact of order $k\le3$ the local cover is the
normalization of
\[
                            w^3=s^2(s-u^k).
\tag{C.D4.19}\label{cdp:eq:C.D4.19}
\]
Blow up the common point of the two smooth branch curves
repeatedly, $k$ times, until their strict transforms are
disjoint. At each step $\sigma$, the branch multiplicity
is three, so the new cyclic line and weight-two branch
are $L'=\sigma^*L(-E)$ and $F'=\sigma^*F-E$; the exceptional
curve has branch weight zero. Both nontrivial normalized
eigensheaves on the blown-up base are therefore the
pullbacks of their predecessors tensored with $\OO(E)$.
For a point blowup,
$\sigma_*\OO(E)=\OO$ and $R^1\sigma_*\OO(E)=0$, as follows
from $0\to\OO\to\OO(E)\to\OO_E(-1)\to0$.
Finite-cover Leray then proves that the induced birational
map of the normalized covers has $R^1\OO=0$ at each step.
The final cover is smooth: each remaining branch is smooth,
the branches are disjoint, and a single branch of either
weight one or two has smooth normalization. Composing
these maps proves rationality of every singularity in
\eqref{cdp:eq:C.D4.19}. In particular $q(X)=0$, and their exceptional
configurations have $b_1=0$ (rational singularities have
rational-tree exceptional divisors, also seen here from
the indicated successive cyclic covers of $E$).

\noindent\textbf{Step 6.} The reduced boundary has two components with independent
classes, and contains the pullback of the ample class
$(3,2)$. If $\nu=\#(R\cap F)$, the same exact homology
calculation now gives
\[
        b_1(U)=0,\qquad e(U)=3\nu,\qquad
        b_1(D)=\nu-1,\qquad
        \operatorname{rank}(\hbox{intersection pairing})=2\nu>0.
\tag{C.D4.20}\label{cdp:eq:C.D4.20}
\]
This contradicts (C.D4.2), and excludes the double common
zero. The three cases exhaust $v=1,2$, proving $e(S)\ge1$.
\end{proof}






\begin{proposition}[The irreducible local comparison and its rational and ruled consequences]\label{cdp:st:C.2.UN.7}
\StatementRecord{cdp:st:C.2.UN.7}{proved}{Historical events: Claude Fable 5.2 (Anthropic), effort high, 16 November 2026: carried as background without fault found. GPT-6 Astra (OpenAI), ChatGPT Pro mode, effort Pro, S6\_09\_PAPERS, 17 November 2026: targeted proof check. Also checked in this writing. Pre-continuation writing events: GPT-6 Astra (OpenAI), ChatGPT Pro mode, effort Pro, S6\_10\_PAPERS, 18 November 2026: targeted proof check. Continuation metadata events: GPT-6 Astra (OpenAI), ChatGPT Pro mode, effort Pro, S6\_10\_PAPERS continuation, 19 November 2026: used and rechecked in this run; status-vocabulary reconciliation; metadata reconciliation only; no new mathematical check. \emph{Pre-continuation recorded proof scope.} Canonical torsion, coherent Euler zero, full-multiplicity trace and finite-character Euler reductions were re-derived. The rational, reduced-horizontal, cusp, tacnode and ordinary-triple cases were followed through A.5.9,16,18-20, IRR.1 and D4.1-2. In the contracted case the exact C.CR.5-6 boundary/support steps and van Kampen map give a free group of rank at most one and the asserted positive Euler alternatives. Those scoped inputs are explicit in the revision; outstanding primary-source access gates are recorded separately and no wider catalogue certification is asserted}{the hypotheses, notation and definitions stated here; \cref{cdp:st:D.1.9}; \cref{cdp:eq:C.UI.2}; \cref{cdp:eq:C.UI.3}; \cref{cdp:st:A.5.3}; \cref{cdp:eq:C.UI.4}; \cref{cdp:eq:C.UI.5}; \cite[\S2 and \S3]{ACampo1975}; \cref{cdp:st:C.2.NC.1}; \cref{cdp:eq:C.UI.6}; \cref{cdp:eq:C.UI.7}; \cref{cdp:st:C.2.MH.2}; \cref{cdp:st:A.5.11}; \cref{cdp:st:C.2.MH.1}; \cref{cdp:st:C.2.UN.5}; \cref{cdp:st:A.5.16}; \cref{cdp:st:A.5.17}; \cref{cdp:eq:C.UI.12}; \cref{cdp:eq:C.UI.11}; \cref{cdp:st:C.2.UN.1}; \cref{cdp:st:C.2.IRR.1}; \cref{cdp:st:A.5.14}; \cref{cdp:st:A.5.18}; \cref{cdp:st:A.5.20}; \cref{cdp:st:A.5.5}; \cref{cdp:st:A.5.19}; \cref{cdp:st:C.2.D4.1,cdp:st:C.2.D4.2}; \eqref{cdp:eq:C.CR.5}--\eqref{cdp:eq:C.CR.6} (only the divisor and support comparisons, proved independently of the later consequence of \cref{cdp:st:C.2.CR.1}); \cref{cdp:st:A.5.10}; \cref{cdp:eq:C.UI.15}}

Let $f:Y\to B$ be a proper holomorphic map from a smooth
complex threefold to a smooth curve, with connected fibres
and smooth general fibre a complex two-torus. Let $mS$ be a full
fibre with irreducible reduced support, and let $F$ be a nearby
smooth fibre. Write $T_1$ for the local monodromy on
$H^1(F,\mathbb Q)$. Then
\[
 m\,e(S_{\rm reg})=\det(1-T_1^m)\ge0.
 \tag{C.UI.1}\label{cdp:eq:C.UI.1}
\]
For its normalization and minimal normalization resolution
$h=\eta\pi:Z\to T\to S$, the following statements hold.

\begin{enumerate}
\item If $Z$ is rational, then $e(S)\ge2$. More precisely,
      whenever $H^1(S,\mathbb Q)\ne0$, one has $e(S)=2$.
\item Suppose $Z$ is irrational ruled. A reduced horizontal
      conductor gives a smooth minimal elliptic ruled
      normalization, with disjoint smooth elliptic conductor
      curves and their finite etale double-cover maps to the
      conductor image. The support has ordinary double-curve
      normal crossings and Euler number zero.
\item Suppose the horizontal conductor is twice a noncontracted
      section. The normalization is again smooth minimal
      elliptic ruled. Its conductor is twice that section;
      its image is smooth elliptic, and the map from the
      section is etale of degree $1$, $2$, or $3$.
      The respective constant transverse types are the
      ordinary cusp, tacnode, and ordinary triple point.
      Every such support has Euler number zero.
\item Suppose the horizontal conductor is twice a section
      contracted by $\pi$. Then $e(S)\ge2$ and
      $\pi_1(S)=1$. In particular this case cannot occur
      under the full CDP hypotheses, since \ref{cdp:st:D.1.9} supplies
      a nonzero integral degree-one class on every fibre
      support there.
\end{enumerate}
The alternatives in the last three items exhaust the
irrational ruled case. No vanishing of $b_1(Y)$ or $b_2(Y)$
is used in the first four conclusions.
\end{proposition}


\begin{proof}
\noindent\textbf{Step 1.} \leavevmode\par\medskip\noindent\textbf{The actual canonical and Chern-class comparisons.}\quad
Choose a coordinate $t$ on a small disc $\Delta$ about the
chosen value and put $U=f^{-1}(\Delta)$. The torsion-free
rank-one sheaf $f_*\omega_{U/\Delta}$ is a line bundle on
$\Delta$. Its rank is one by base change on the punctured
disc, where a two-torus has a one-dimensional space of
holomorphic top forms. A generator evaluates to a
holomorphic section of $\omega_{U/\Delta}$ with no zeros
on any smooth nearby fibre: on that fibre it is a nonzero
holomorphic two-form. Its divisor is therefore $kS$ for
some integer $k\ge0$. The coordinate $dt$ identifies the
relative canonical line with $K_U$. The actual fibre
equation and adjunction give
\[
 K_U\simeq\OO_U(kS),\qquad
 \OO_U(mS)\simeq\OO_U,\qquad
 \omega_S^{\otimes m}\simeq\OO_S.
 \tag{C.UI.2}\label{cdp:eq:C.UI.2}
\]
This restricts a section of the relative canonical line;
it does not pull a three-form back as a differential form
to a surface.

\noindent\textbf{Step 2.} Here is the Euler calculation on this actual tube.
The map is flat: its local rings are torsion-free
over the regular one-dimensional base ring.
Proper flatness makes the holomorphic Euler
characteristic of its full fibres constant, so
$\chi(\OO_{mS})=\chi(\OO_F)=0$.
The Cartier filtration of $\OO_{mS}$ has successive
quotients $\OO_S(-jS)$, $0\le j<m$; by \eqref{cdp:eq:C.UI.2}
these are torsion line-bundle twists.
Each such twist has the same Euler characteristic
as $\OO_S$. One can check this without a compact
ambient manifold: use the normalization sequence
and the resolution Leray sequence to compute the
Euler characteristic on the compact smooth surface
$Z$, the conductor curve, and the finite-support
sheaf $R^1\pi_*\OO_Z$. Smooth-surface
Riemann--Roch leaves the first term unchanged
under a torsion twist; curve Riemann--Roch,
applied to the generic-rank filtration of the
conductor quotient, leaves the second unchanged;
the finite lengths in the third are unchanged.
These exact sequences prove the assertion on $S$.
Consequently, writing $c_i=c_i(TU)$, one has
\[
 \begin{split}
 0&=\chi(\OO_{mS})=m\chi(\OO_S),\\
 c_1(h^*TU)&=0\quad\hbox{in }H^2(Z,\mathbb Q),
 \qquad c_2(h^*TU)=0\quad\hbox{in }H^4(Z,\mathbb Q).
 \end{split}
 \tag{C.UI.3}\label{cdp:eq:C.UI.3}
\]
The first Chern-class assertion follows from \eqref{cdp:eq:C.UI.2}.
For the second we specify the local cycle comparison.
The inverse image of a radial segment $[0,t_0]$
is a compact triangulable real five-dimensional
space. Over the open segment it has the orientation
of the complex fibres followed by the radial
orientation. At a smooth point of $S$, the exact
local equation $t=u z^m$ gives $m$ radial sheets
with the same boundary orientation on $S$.
The remaining central singular locus has real
dimension at most two and contributes no
four-dimensional boundary. The resulting oriented
fundamental chain therefore has boundary
$[F]-m[S]$. This proves $[F]=m[S]$ in
$H_4(U,\mathbb Q)$, with its multiplicity retained.
The proper map $h$ has degree one onto $S$, hence
\[
 \int_Z h^*c_2=\langle c_2,[S]\rangle
       =\frac1m\int_F c_2(TU|_F)=0.
\]
The last equality follows from
$0\to T_F\to TU|_F\to\OO_F\to0$ and the
triviality of $T_F$. Top rational cohomology of
the connected closed oriented four-manifold $Z$
is one-dimensional, proving the second class
equality in \eqref{cdp:eq:C.UI.3}.

\noindent\textbf{Step 3.} Finite duality and the minimal-resolution calculation
\ref{cdp:st:A.5.3} now give the actual effective integral divisor $D$
with
\[
                    K_Z(D)\simeq h^*\omega_S.
 \tag{C.UI.4}\label{cdp:eq:C.UI.4}
\]
Thus that line is torsion. Equations \eqref{cdp:eq:C.UI.2}--\eqref{cdp:eq:C.UI.4}
also give degree zero for both $K_Y$ and $\OO_Y(S)$
on every compact curve in this fibre. These are precisely
the ambient numerical inputs used below.

\par\medskip\noindent\textbf{The proper analytic Lefschetz formula.}\quad
Take an embedded resolution
$r:\widetilde U\to U$ of the pair, isomorphic over
$U\setminus\operatorname{Sing}S$, with a simple
normal-crossings central divisor
$\sum_i a_iE_i$. Its strict transform $E_0$ has
multiplicity $a_0=m$. Every exceptional multiplicity
satisfies
\[
                    a_i\ge2m\quad(i\ne0).
 \tag{C.UI.5}\label{cdp:eq:C.UI.5}
\]
Indeed a first centre outside the existing exceptional
locus lies in the singular locus of the reduced strict
transform, where its order is at least two. A centre
inside the existing exceptional locus already receives
the coefficient of an exceptional component, at least
$2m$. The coefficient of a new exceptional divisor is
the order of the current pullback of $f$ along its centre,
so the other contributions are nonnegative. Induction
proves \eqref{cdp:eq:C.UI.5}. The resolution can be chosen with all
centres above $\operatorname{Sing}S$. Its strict open
stratum is consequently
$E_0^\circ\simeq S_{\rm reg}$.

\noindent\textbf{Step 4.} We specify the proper version of the trace formula used
here. A'Campo's original Theorem~1 is stated for a local
hypersurface germ. His \S2 begins with a proper analytic
map from a smooth analytic space to a disc, smooth off
the central simple-normal-crossings divisor, with no
Kahler assumption. The global boundary fibration and
its cohomology spectral sequence are constructed there;
the stalk-trace computation is on pp.~243--244
\cite[\S2 and \S3]{ACampo1975}. In the present notation
that construction gives
\[
 L(T^j):=\sum_q(-1)^q\operatorname{Tr}
       (T^j|H^q(F,\mathbb Q))
       =\sum_{a_i\mid j}a_i e(E_i^\circ).
 \tag{C.UI.6}\label{cdp:eq:C.UI.6}
\]
For clarity, the local and global maps producing this
identity are as follows. The real oriented blowup of
the divisor, followed by the angular map to $S^1$,
has a boundary fibre homeomorphic to $F$, as in
\ref{cdp:st:C.2.NC.1}. Above an $r$-branch stratum its fibre is
\[
 \{(u_1,\ldots,u_r)\in(S^1)^r:
                 u_1^{a_1}\cdots u_r^{a_r}=e^{i\theta}\}.
 \tag{C.UI.7}\label{cdp:eq:C.UI.7}
\]
An integral change of angular basis identifies this with
$\gcd(a_1,\ldots,a_r)$ copies of $(S^1)^{r-1}$.
Monodromy permutes these copies and acts by translations
on them. For $r\ge2$ the alternating trace is zero,
whether the permutation fixes a copy or not. For $r=1$
it is $a_1$ when $a_1\mid j$ and zero otherwise.
The constructible higher-direct-image sheaves of this
boundary projection carry those stalk actions. Their
cohomology spectral sequence computes $H^*(F)$ with
its actual monodromy. Triangulate the compact central
divisor compatibly with the strata and these local
systems. On every open cell the alternating trace of
compactly supported cochains is its signed cell
contribution times the stalk trace. Additivity and
the spectral sequence prove \eqref{cdp:eq:C.UI.6}. The proper
modification is an isomorphism over $\Delta^*$, so
these are the nearby fibre and monodromy of the
original map.

\noindent\textbf{Step 5.} This description also proves quasi-unipotence without
an extra hypothesis. If $N$ is a common multiple of
the finitely many $a_i$, then $T^N$ is the identity
on every stalk cohomology group in \eqref{cdp:eq:C.UI.7}, since
the finite permutation disappears and translations
act trivially on cohomology. It is the identity on
the corresponding constructible cohomology sheaves,
hence on the second page of their cohomology spectral
sequence. The finite filtration on $H^q(F)$ has
identity action on its associated graded. Thus
$T^N$ is unipotent on $H^q(F)$.

\noindent\textbf{Step 6.} Put $j=m$ in \eqref{cdp:eq:C.UI.6}. By \eqref{cdp:eq:C.UI.5} only $E_0$
contributes. The cohomology ring of the real
four-torus is the exterior algebra of its first
cohomology, and a monodromy diffeomorphism acts by
the exterior powers of $T_1$. Therefore
\[
 m e(S_{\rm reg})
   =\sum_{q=0}^4(-1)^q\operatorname{Tr}(\Lambda^qT_1^m)
   =\det(1-T_1^m).
 \tag{C.UI.8}\label{cdp:eq:C.UI.8}
\]
Every eigenvalue of $T_1^m$ is a root of unity by
the preceding quasi-unipotence. A real eigenvalue
is $1$ or $-1$ and contributes $0$ or $2$; each
nonreal conjugate pair contributes $|1-\lambda|^2$.
Their product is nonnegative. This proves \eqref{cdp:eq:C.UI.1}.

\par\medskip\noindent\textbf{The local degree-one class and the finite-character comparison.}\quad
We record the exact local replacement for the global
transgression used in \ref{cdp:st:C.2.MH.2}. Every nonzero class
$\gamma_S\in H^1(S,\mathbb Z)$ extends to a sufficiently
small tube and restricts nontrivially over $\mathbb Q$
to $F$. To see the extension, triangulate the analytic
pair $(Y,S)$ as in \ref{cdp:st:A.5.11} and take a regular
neighbourhood $N$ with a retraction $N\to S$.
Properness puts $f^{-1}(\Delta)$ inside $N$ after
shrinking $\Delta$. Pullback under that retraction
gives the specified class $\gamma_U$ on the tube,
whose restriction to $S$ is $\gamma_S$.

\noindent\textbf{Step 7.} For injectivity upon restriction to $F$, the inclusion
$U^*=U\setminus S\hookrightarrow U$ surjects on
fundamental groups: loops can be perturbed off the
real codimension-two divisor, avoiding its higher
codimension singular strata. If a rational
degree-one class on $U$ restricts to zero on $F$,
the exact homotopy sequence of the punctured-disc
bundle makes its restriction to $U^*$ a scalar
multiple $a$ of the base-circle class. A small
meridian at a smooth point of $S$ is null-homologous
in $U$ and maps to $m$ times that base circle.
It follows that $ma=0$, hence $a=0$. Surjectivity
on fundamental groups now makes the original
class zero. In particular $\gamma_U|_F\ne0$.

\noindent\textbf{Step 8.} Choose a prime $\ell$ for which both restrictions
to $S$ and $F$ stay nonzero modulo $\ell$.
Their values on one loop in each space exclude
only finitely many primes. The reduction of
$\gamma_U$ defines the actual cyclic etale cover
of $U$, with rational augmentation local system
$V$ of rank $d=\ell-1$. It is self-dual,
$H^0(S,V)=0$, and every nontrivial complex character
in it is nontrivial on $F$. The Koszul complex
of $\mathbb Z^4$ is then contracted by the inverse
of a nonzero character difference. Hence
\[
 H^*(F,V)=H^*(U^*,V)=0.
 \tag{C.UI.9}\label{cdp:eq:C.UI.9}
\]
The second equality is the bundle Leray comparison.

\noindent\textbf{Step 9.} When $Z$ is rational or ruled over a curve of genus
$g$, all strata of an embedded resolution of the
central fibre are projective: the strict transform
is a blowup of $Z$, and every other surface arises
by blowups from a projective bundle over a compact
curve or from $\mathbb P^2$. \ref{cdp:st:C.2.MH.1} thus applies
to this actual local system. Its support comparison
makes $H^3(S,V)$ pure of weight three. The exact
normalization-conductor descent sequence has a
kernel which is a quotient of the degree-two
cohomology of a compact curve, hence pure of weight
two; strictness makes that kernel zero. The
negative-definite exceptional intersection matrices
then give the actual isomorphisms
\[
 H^3(S,V)\simeq H^3(Z,V),\qquad
 H^4(S,V)\simeq H^4(Z,V).
 \tag{C.UI.10}\label{cdp:eq:C.UI.10}
\]
These are precisely the maps (C.MH.10)--(C.MH.12);
none of their proofs uses ambient cohomology.
The acyclic punctured tube and local support duality
give, in the direct limit of shrinking tubes,
$H^j(S,V)\simeq H_{6-j}(S,V)$. Thus
\[
 d e(S)=2h^4(Z,V)-h^3(Z,V)
 =\begin{cases}
 2d,&Z\text{ rational},\\
 (2-2g)d,&Z\text{ ruled over genus }g.
 \end{cases}
 \tag{C.UI.11}\label{cdp:eq:C.UI.11}
\]
For the second line, the local system pulls back
from the ruling base. If its character is trivial
there, the two dimensions on the left are $d$
and $2gd$; otherwise they are $0$ and $(2g-2)d$.
This proves the precise local character version
of \ref{cdp:st:C.2.MH.2} and the additional assertion in item~1.

\par\medskip\noindent\textbf{Strict positivity for a rational normalization.}\quad
If $S$ is normal, \ref{cdp:st:C.2.UN.5} already makes it a
smooth torus or bielliptic surface, so its
resolution is neither rational nor irrational
ruled. Thus the remaining surface cases under
consideration are nonnormal.
The proof of \ref{cdp:st:A.5.16} uses the ambient assumption
only for $\omega_S$ torsion and $\chi(\OO_S)=0$,
now supplied by \eqref{cdp:eq:C.UI.2}--\eqref{cdp:eq:C.UI.3}. Consequently
its exact support and genus conclusions apply:
\[
 \begin{gathered}
 K_Z+D\simeq\OO_Z,\qquad D_{\rm red}\text{ connected},
 \qquad p_a(D_{\rm red})\le1,\\
 \Gamma_{\rm red}\text{ connected},\qquad
 p_a(\Gamma_{\rm red})\le2,\qquad
 e(\Gamma_{\rm red})\ge0,\\
 D_{\rm red}=h^{-1}(\Gamma_{\rm red}),\qquad
 S_{\rm reg}\simeq Z\setminus D_{\rm red}.
 \end{gathered}
 \tag{C.UI.12}\label{cdp:eq:C.UI.12}
\]
The final isomorphism retains all exceptional curves:
any exceptional fibre outside the conductor would
give an isolated normal hypersurface singularity,
excluded by the actual Milnor-sphere comparison
in \ref{cdp:st:A.5.17}. Thus
$e(S)=e(S_{\rm reg})+e(\Gamma_{\rm red})\ge0$.

\noindent\textbf{Step 10.} Suppose it were zero. Both summands would be zero.
If $D_{\rm red}$ has an elliptic-normalized component,
the bound on its arithmetic genus makes that component
smooth and unique. Follow a sequence of rational
surface blowdowns to $\mathbb P^2$ or a Hirzebruch
surface. The divisor comparison at each step is
$D_{\rm new}=\sigma^*D_{\rm old}-E$; a centre of
multiplicity one increases the Euler number of the
complement by one, and a centre of multiplicity at
least two preserves it. A centre outside the
boundary is impossible by effectiveness. The
elliptic component survives these blowdowns.
On $\mathbb P^2$ the resulting anticanonical cubic
is a smooth elliptic cubic. On a Hirzebruch
surface the elliptic component is a smooth
bisection: a section is rational, and for a
coefficient-one bisection $H$ with vertical part
$V$ adjunction gives
$2p_a(H)-2+H\cdot V=0$, forcing $V=0$.
The minimal complement Euler numbers are therefore
$3$ and $4$, respectively. Their blowup comparison
contradicts $e(S_{\rm reg})=0$.

\noindent\textbf{Step 11.} All normalized components of $D_{\rm red}$ are
therefore rational. Its first Betti number is at
most one, since the normalization sequence gives
$b_1(D_{\rm red})\le p_a(D_{\rm red})\le1$.
Every component of $\Gamma_{\rm red}$ is also
rational-normalized, being the finite image of a
strict conductor curve. Let their number be
$n\ge1$. For a connected compact curve with
rational normalizations the exact normalization
topology gives
$e(\Gamma_{\rm red})=1+n-b_1(\Gamma_{\rm red})$.
The supposed equality therefore gives
$b_1(\Gamma_{\rm red})=n+1\ge2$.

\noindent\textbf{Step 12.} Here the required comparison with $S$ is an
actual pushout, not a presumed identification
of cohomology groups. By \eqref{cdp:eq:C.UI.12}, the natural
map from
$Z\cup_{D_{\rm red}}\Gamma_{\rm red}$ to $S$
is a continuous bijection from a compact space
to a Hausdorff space, hence a homeomorphism.
Triangulation makes $D_{\rm red}\hookrightarrow Z$
a cofibration, so its double mapping cylinder
gives the cohomology sequence of that pushout.
Connectedness and $H^1(Z,\mathbb Q)=0$ give
\[
 H^1(S,\mathbb Q)
  \simeq\ker\bigl[H^1(\Gamma_{\rm red},\mathbb Q)
                       \to H^1(D_{\rm red},\mathbb Q)\bigr].
 \tag{C.UI.13}\label{cdp:eq:C.UI.13}
\]
Its dimension is at least $n>0$. The proved
local finite-character calculation \eqref{cdp:eq:C.UI.11}
now gives $e(S)=2$, contradicting zero.
Thus $e(S)>0$, and \ref{cdp:st:C.2.UN.1} makes it an even
integer. This proves item~1.

\par\medskip\noindent\textbf{All noncontracted irrational ruled alternatives.}\quad
On a general ruling fibre $R$, \eqref{cdp:eq:C.UI.4} gives
$D\cdot R=2$. Thus its horizontal part is reduced
of total degree two, or is $2C$ for a smooth
section $C$. The proof of \ref{cdp:st:C.2.IRR.1} uses the
ambient condition only for the zero degrees
in (C.IR.2), (C.IR.7)--(C.IR.8) and for the
torsion-canonical, zero-$\chi$ genus bound \ref{cdp:st:A.5.14}.
All are furnished by \eqref{cdp:eq:C.UI.2}--\eqref{cdp:eq:C.UI.4}.
Its exact branch-reception, pinch-point and
blowdown comparisons therefore give item~2,
including its assertion that no horizontal
component is contracted.

\noindent\textbf{Step 13.} Consider $D=2C+V$ with $C$ not contracted.
Every resolution-exceptional curve is vertical.
Indeed a positive-coefficient horizontal one
would have to be $C$. A coefficient-zero
exceptional curve $E$ satisfies
$K_Z\cdot E=D\cdot E=0$ by minimality and
adjunction. A horizontal exceptional curve over
a positive-genus base has $p_a(E)\ge1$ and
$E^2<0$, so $K_Z\cdot E>0$, a contradiction.
These are exactly the exceptional-curve
hypotheses of \ref{cdp:st:A.5.18}--\ref{cdp:st:A.5.20}.

\noindent\textbf{Step 14.} Let $H$ be the reduced image of $C$ and
$d=\deg(C\to H^\nu)$. The branchwise conductor
formula is its intrinsic plane-branch conductor
plus the sum of its intersections with the
other branches. Its value on each horizontal
branch is two. If $d=1$, $H^\nu$ has the
positive genus of $C$, and no vertical rational
curve can dominate it. The sole branch with
conductor two is the ordinary cusp. If $d=2$,
the two horizontal branches must be smooth:
a singular branch would contribute at least two
intrinsically and at least one from the other
branch. Their allowed intersections give either
contact two between those branches, or simple
contact between them and with one additional
smooth branch. These are exactly a tacnode and
an ordinary triple point. For $d=3$ all three
branches must be smooth and pairwise transverse,
and there is no further branch. For $d\ge4$
the three other intersection contributions
already exceed two. This proves exhaustion
of the transverse cases used in the following
comparisons.

\noindent\textbf{Step 15.} The local numerical hypotheses of \ref{cdp:st:A.5.14} are
exactly $\chi(\OO_S)=0$ and torsion $\omega_S$.
Its genus bound and delta-one complete-intersection
comparison consequently remain valid. The
conductor-cleared jets \ref{cdp:st:A.5.5} use degree zero for
$K_Y$ and $\OO_Y(S)$ on the image curve.
The cusp proof \ref{cdp:st:A.5.20} and the complete tacnode
proof \ref{cdp:st:A.5.18} therefore apply. The latter's
rational-image fundamental-group restriction
uses only the local quotient of
$\mathbb Z^4\rtimes\mathbb Z$, not an ambient
degree-one class. For the triple case, \ref{cdp:st:A.5.19}
excludes singular rational images with a
positive-genus covering component. \cref{cdp:st:C.2.D4.1,cdp:st:C.2.D4.2}
exclude the connected degree-three and degree-two
covers of a smooth rational image. Their normal
cone is the actual cubic tensor in
$\operatorname{Sym}^3N_{H/Y}^*$; on
$H\simeq\mathbb P^1$ the needed identities
$\OO_Y(S)|_H=\OO_H$ and $\deg N_{H/Y}=-2$
follow from \eqref{cdp:eq:C.UI.2}. Their compact-core
transport is into the original nearby torus,
so no global Picard condition enters. The
independent ruled Chern comparison \ref{cdp:st:A.5.10} also
extends, since its additional input
$c_2(h^*TY)=0$ is \eqref{cdp:eq:C.UI.3}.

\noindent\textbf{Step 16.} Consequently the image in each surviving case
is smooth elliptic and the transverse analytic
type is constant. No vertical rational curve
dominates that image, so all its generic
branches belong to $C$; the actual local
normalizations give etale degrees $1$, $2$,
or $3$, respectively. The normal line of $C$
has degree zero. The divisor comparison
(A.T.22), together with this square-zero
section and the absence of further conductor
along it, contracts no point of $C$ and makes
the minimal ruled boundary exactly $2C$.
A first reverse blowup away from $C$ would
create coefficient $-1$; a blowup on $C$
would create a conductor component meeting
it, contrary to the constant local type.
Thus the normalization is smooth minimal
elliptic ruled, exactly as proved in
\ref{cdp:st:A.5.18}--\ref{cdp:st:A.5.20}. The normalization Euler
comparison $e(S)=e(T)-e(C)+e(H)$ gives zero.
This proves item~3.

\par\medskip\noindent\textbf{A contracted section and its fundamental group.}\quad
Finally let $D=2C+V$ with $C$ contracted and
let $g\ge1$ be the ruling genus. The exact
boundary blowdown relation \eqref{cdp:eq:C.CR.5}, now
supplied by \eqref{cdp:eq:C.UI.4}, shows that $D_{\rm red}$
is the smooth section $C$ with trees of smooth
rational curves attached. All its crossings
are transverse nodes. Its inclusion in $Z$
induces an isomorphism on fundamental groups,
both identified with $\pi_1(C)$.
The proof of \eqref{cdp:eq:C.CR.6} uses only this
connectedness, minimality and the isolated
singularity exclusion. It gives
\[
 D_{\rm red}=h^{-1}(\Gamma_{\rm red}),\qquad
 S_{\rm reg}\simeq Z\setminus D_{\rm red}.
 \tag{C.UI.14}\label{cdp:eq:C.UI.14}
\]
Every component of the connected curve
$\Gamma_{\rm red}$ is dominated by a vertical
rational curve and hence has normalization
$\mathbb P^1$. The same actual topological
pushout as above and van Kampen give
\[
 \pi_1(S)\simeq
 \pi_1(Z)*_{\pi_1(D_{\rm red})}\pi_1(\Gamma_{\rm red})
 \simeq\pi_1(\Gamma_{\rm red}).
 \tag{C.UI.15}\label{cdp:eq:C.UI.15}
\]
The second amalgamating map is trivial because
$C$ generates $\pi_1(D_{\rm red})$ and is
contracted. This group is free of rank
$\beta=b_1(\Gamma_{\rm red})$.

\noindent\textbf{Step 17.} For its comparison with the nearby torus,
take the same regular neighbourhood and
small tube used above. A loop in $S$ lies
in the tube, and the neighbourhood retraction
is the identity on $S$. Thus the tube group
surjects onto $\pi_1(S)$. Perturbing loops
off $S$ makes it a quotient of
$\pi_1(U^*)=\mathbb Z^4\rtimes\mathbb Z$.
The latter is metabelian, whereas a free
group of rank at least two surjects onto
the nonmetabelian group $S_4$. Hence
$\beta\le1$. If $n\ge1$ is the number
of conductor-image components, its exact
normalization topology gives
\[
 e(\Gamma_{\rm red})=n+1-\beta\ge n\ge1,
 \qquad e(S)=e(S_{\rm reg})+e(\Gamma_{\rm red})>0.
 \tag{C.UI.16}\label{cdp:eq:C.UI.16}
\]
Here \eqref{cdp:eq:C.UI.1} proves the last inequality.
Parity makes $e(S)\ge2$. If $\beta=1$,
\eqref{cdp:eq:C.UI.15} supplies a nonzero integral
degree-one class on $S$; \eqref{cdp:eq:C.UI.11} then
gives $e(S)=2-2g\le0$, a contradiction.
Thus $\beta=0$, proving $\pi_1(S)=1$.
The transgression of \ref{cdp:st:D.1.9} excludes this
case under the full CDP hypotheses. This
proves item~4 and completes the proof.
\end{proof}



The assertion $e(S)=2$ in item~1 is deliberately tied
to the proved nonzero local-character comparison.
\ref{cdp:st:C.2.UN.8} constructs an irreducible rational-normalized
support with $\pi_1(S)=1$ and $e(S)=6$, so that
local character cannot be inserted without its
existence argument. The classification of all
irreducible zero-Euler supports is assembled in
\ref{cdp:st:C.2.UN.12} using this result, \ref{cdp:st:C.2.UN.5}, and the
class VII conductor comparisons. Reducible
unrestricted supports are not covered by this
irreducible statement.




\begin{proposition}[The unrestricted irreducible zero-Euler criterion]\label{cdp:st:C.2.UN.12}
\StatementRecord{cdp:st:C.2.UN.12}{proved}{Historical events: Claude Fable 5.2 (Anthropic), effort high, 16 November 2026: carried as background without fault found. GPT-6 Astra (OpenAI), ChatGPT Pro mode, effort Pro, S6\_09\_PAPERS, 17 November 2026: targeted proof check. Also checked in this writing. Pre-continuation writing events: GPT-6 Astra (OpenAI), ChatGPT Pro mode, effort Pro, S6\_10\_PAPERS, 18 November 2026: targeted proof check. Continuation metadata events: GPT-6 Astra (OpenAI), ChatGPT Pro mode, effort Pro, S6\_10\_PAPERS continuation, 19 November 2026: used and rechecked in this run; status-vocabulary reconciliation; metadata reconciliation only; no new mathematical check. \emph{Pre-continuation recorded proof scope.} The canonical tube, all ruled conductor maps, the class VII cyclic-quotient reduction and the exact complement isomorphism were re-derived. The Hopf criterion is exactly b2 minus the rational-component count plus the actual conductor-image Euler number; the two zero patterns follow from its branch and delta sequences. Both D4 cover inputs and IRR.1 are now explicit references in the revision}{the hypotheses, notation and definitions stated here; \cref{cdp:st:C.2.UN.5}; \cref{cdp:st:C.2.UN.7}; \cref{cdp:st:C.2.UN.19}; \cref{cdp:st:A.5.20}; \cref{cdp:st:A.5.18}; \cref{cdp:st:C.2.D4.1,cdp:st:C.2.D4.2}; \cref{cdp:st:C.2.IRR.1}; \cref{cdp:st:C.2.VII.1}; \cref{cdp:st:C.2.VII.2}; \cref{cdp:st:A.5.17}; \cref{cdp:eq:C.IZ.6}; \cref{cdp:eq:C.IZ.5}; \cref{cdp:eq:C.IZ.8}; \cref{cdp:eq:C.IZ.7}}

Let $f:Y\to B$ be a proper holomorphic map from a smooth
complex threefold to a smooth complex curve, with connected
fibres and smooth general fibre a complex two-torus. Fix an
irreducible full fibre $mS$, and write
$h=\eta\pi:Z\to T\to S$ for normalization followed by its
minimal resolution. Every such support satisfies $e(S)\ge0$.
Its Euler number is zero precisely in the following geometric
cases, stated for these actual maps.

A normal support is a smooth torus or a smooth bielliptic
surface, with multiplicities as in \ref{cdp:st:C.2.UN.5}.

A nonnormal support with irrational ruled normalization has
smooth minimal normalization ruled over an elliptic curve.
Its conductor image and reduced inverse image are disjoint
unions of smooth elliptic curves. The transversal types and
their finite conductor maps are exactly those in the table:
\[
\begin{array}{c|c|c}
 \text{normalization conductor}&\text{transversal type}
       &\text{degree over the image}\\ \hline
 H\text{ reduced, ruling degree }2&xy=0&2\\
 2C,\ C\text{ a section}&x^2-y^3=0&1\\
 2C,\ C\text{ a section}&x^2-y^4=0&2\\
 2C,\ C\text{ a section}&xy(x-y)=0&3
\end{array}                                                   \tag{C.IZ.1}\label{cdp:eq:C.IZ.1}
\]
Every displayed map of elliptic curves is unramified; the
degree-one map is an isomorphism. The transversal type is
constant along each conductor component. In the first row,
degree two means the total degree over each image component;
its inverse image may have two connected components. There are no
contracted horizontal sections in the zero-Euler case.

For a nonnormal support whose normalization resolution
has a Hopf surface as its minimal model, put $R=D_{\rm red}$,
where $D$ is the effective
resolution conductor with $K_Z+D=h^*\omega_S$. Let $r$ be
the number of rational components of $R$. Put
$C=\Gamma_{\rm red}$ for the conductor image, let $n_0$ be
the number of its rational-normalized components, and set
$v(C)=\sum_x(s_x-1)$, using the actual number $s_x$ of
normalization branches at each singular point. The exact
zero criterion is
\[
 b_2(Z)=r,\qquad
 \begin{cases}
 n_0=0,\quad v(C)=0,\quad\sum_x\delta_x\le1,\quad\text{or}\\
 n_0=1,\quad v(C)=2,\quad\sum_x\delta_x=2.
 \end{cases}                                                   \tag{C.IZ.2}\label{cdp:eq:C.IZ.2}
\]
Every other normalized conductor component is elliptic.
In the second line every singular point satisfies
$\delta_x=s_x-1$. No normalization singularity lies outside
the conductor in either line. An irreducible support with
rational normalization resolution always has $e(S)\ge2$.

This is an equivalence for fibres of the specified map.
The displayed curve data retain their normalization and
resolution morphisms; no assertion of smoothability of
arbitrary abstract surface gluings is included.
\end{proposition}


\begin{proof}
\noindent\textbf{Step 1.} We first check the classification's scope without an ambient
Chern-class assumption. On a sufficiently small fibre tube
$U=f^{-1}(\Delta)$, the evaluation of a generator of the
rank-one torsion-free sheaf $f_*K_U$ gives
$K_U\simeq\OO_U(kS)$, as proved with its divisor comparison
in \ref{cdp:st:C.2.UN.7}. Since $mS$ is principal, $K_U$ and
$\omega_S=\OO_U((k+1)S)|_S$ are torsion. The actual
$\OO_{mS}$ filtration and flatness give
$\chi(\OO_S)=0$. For nonnormal $S$, finite duality and
resolution adjunction give the nonempty effective divisor
$D$ and the stated torsion line $K_Z+D$.

\noindent\textbf{Step 2.} Every positive pluricanonical section on $Z$ would, after
taking a power trivializing that torsion line, give a
nonzero section of $\OO_Z(-aD)$ for some $a>0$.
Its inclusion in $\OO_Z$ would be a nonzero holomorphic
function on the compact connected surface vanishing along
$D$, a contradiction. Thus $\kappa(Z)=-\infty$.
The surface classification, together with the actual
minimal-model blowdown, makes $Z$ rational, irrational
ruled, or class VII. In the last case \ref{cdp:st:C.2.UN.19} proves,
through the specified cyclic cover and quotient maps,
that its minimal model is Hopf. Thus the formulation
above includes every class VII possibility for the
actual degeneration. Normal supports have already been
classified with their finite root covers in \ref{cdp:st:C.2.UN.5}.

\ref{cdp:st:C.2.UN.7} proves strict positivity for rational $Z$.
It also proves strict positivity when an irrational
horizontal section is contracted, and extends every
noncontracted irrational ruled comparison to this tube.
The precise replacements for the earlier ambient $b_2=0$
uses are the torsion canonical and normal lines, the
vanishing of $\chi(\OO_S)$, and
$\int_Zh^*c_2(T_U)=0$, all proved there through restriction
to the nearby torus.

\noindent\textbf{Step 3.} For clarity, the horizontal possibilities exhaust the
ruling degree two of $D$. A reduced horizontal part is
covered by C.2.IRR.1: its adjunction comparison makes its
curves smooth elliptic and disjoint, its normalization
smooth minimal ruled, and its conductor map an unramified
double cover with product-node germs. A nonreduced
horizontal part is $2C$ for a section $C$. At a generic
conductor point, the plane-curve conductor exponent of a
normalization branch equals its intrinsic conductor plus
its intersections with the other branches. Exponent two
permits one cusp branch, two smooth branches meeting with
multiplicity two, or three pairwise transverse smooth
branches. These are the last three rows of \eqref{cdp:eq:C.IZ.1}.
An additional branch would give exponent at least three.
The exact cusp, tacnode and cubic-cover comparisons in
\ref{cdp:st:A.5.20}, \ref{cdp:st:A.5.18} and \cref{cdp:st:C.2.D4.1,cdp:st:C.2.D4.2}, extended in \ref{cdp:st:C.2.UN.7},
give the asserted elliptic image, constant type and
degrees one, two and three. A vertical rational curve
cannot dominate that elliptic image, so every generic
branch is accounted for by the specified horizontal map.
All these comparisons remove vertical conductor and
normalization exceptional curves. Their Euler formula is
\[
 e(S)=e(T)-e((\eta^{-1}C)_{\rm red})+e(C)=0,             \tag{C.IZ.3}\label{cdp:eq:C.IZ.3}
\]
since each term vanishes. This proves both directions of
the ruled part of the criterion.

\noindent\textbf{Step 4.} It remains to check and extend the class VII part. The
proof of \ref{cdp:st:C.2.VII.1} uses only the actual torsion relation
$K_Z+D$; it gives
\[
 e(R)=r,\qquad r\le b_2(Z).                              \tag{C.IZ.4}\label{cdp:eq:C.IZ.4}
\]
The proof of \ref{cdp:st:C.2.VII.2} uses only this relation,
$\omega_S$ torsion and $\chi(\OO_S)=0$. They have just
been established on $U$, so its finite-duality and
curve-normalization maps give
\[
 \chi(\OO_C)\ge-1,\qquad e(C)\ge0.                      \tag{C.IZ.5}\label{cdp:eq:C.IZ.5}
\]
Its normalizations are rational or elliptic, because they
are finite images of the rational or elliptic curves in
$R$. These statements have no remaining ambient Betti
assumption.

\noindent\textbf{Step 5.} We sharpen the exceptional terms in the earlier comparison.
At a point of $T$ outside the conductor, $\eta$ is an
isomorphism. A singular point there would consequently be
an isolated normal hypersurface singularity of $S$ in
the smooth tube. The isolated Milnor-fibre argument
\ref{cdp:st:A.5.17}, retaining the local unit and multiplicity of
the full fibre, excludes that point. Thus $T$ is smooth
outside the conductor. Every resolution-exceptional curve
of coefficient zero in $D$ is disjoint from $R$: minimal
resolution adjunction gives $K_Z\cdot E\ge0$, while
$D\cdot E\ge0$ and $(K_Z+D)\cdot E=0$ force both
numbers to vanish. Each exceptional fibre over the
conductor meets its strict transform and is connected.
Hence it cannot contain such a disjoint curve, and its
whole inverse image lies in $R$. There are no exceptional
fibres outside the conductor, by the preceding local
argument. We have therefore identified the actual
complement isomorphism
\[
 Z\setminus R\xrightarrow{\ \simeq\ }
 T\setminus(\eta^{-1}C)_{\rm red}
 \xrightarrow{\ \simeq\ }S\setminus C.                \tag{C.IZ.6}\label{cdp:eq:C.IZ.6}
\]
Its first map is the resolution restricted to this open
set, and its second is the finite normalization restricted
where its conductor ideal is the unit ideal.

\noindent\textbf{Step 6.} A class VII surface and its point blowups have
$b_1=b_3=1$, so $e(Z)=b_2(Z)$. Euler additivity in
\eqref{cdp:eq:C.IZ.6} now yields the exact formula
\[
 e(S)=b_2(Z)-r+e(C),\qquad
 b_2(Z)-r\ge0,\quad e(C)\ge0.                           \tag{C.IZ.7}\label{cdp:eq:C.IZ.7}
\]
The curve normalization sequence identifies the remaining
term, rather than replacing its singular points by a
model:
\[
 e(C)=2n_0-v(C),\qquad
 \chi(\OO_C)=n_0-\sum_x\delta_x,\qquad
 \delta_x\ge s_x-1.                                    \tag{C.IZ.8}\label{cdp:eq:C.IZ.8}
\]
Equations \eqref{cdp:eq:C.IZ.5} and \eqref{cdp:eq:C.IZ.8} imply $e(C)\ge n_0-1$.
Consequently $e(C)=0$ allows exactly $n_0=0,v=0$ or
$n_0=1,v=2$. In the first case the coherent bound gives
$\sum\delta_x\le1$. In the second, that bound and
$\delta_x\ge s_x-1$ give $\sum\delta_x=2$, with equality
at every singular point. Conversely either of these
two curve patterns makes $e(C)=0$ in the exact
topological formula. Equation \eqref{cdp:eq:C.IZ.7} proves precisely
\eqref{cdp:eq:C.IZ.2}, and proves the sign for every class VII
support. All three nonnormal surface classes have now
been included, which proves the assertion.
\end{proof}



The class VII surface input is the proved anticanonical
configuration theorem checked in \ref{cdp:st:C.2.VII.1}, with its exact
references to \cite{Dloussky2006,DOT2003}. \ref{cdp:st:C.2.UN.19} further
identifies the minimal model in this branch as Hopf;
the existing conductor comparison applies through its
actual point-blowup map. The result here
transfers that input to arbitrary irreducible torus fibres
through the canonical evaluation, isolated-singularity
exclusion and the complement map \eqref{cdp:eq:C.IZ.6}. The unreduced
fibre multiplicity is retained throughout. The remaining
unrestricted classification problem in C.2 concerns
reducible supports.




\begin{proposition}[Proved elliptic-image restriction]\label{cdp:st:C.2.INT.1}
\StatementRecord{cdp:st:C.2.INT.1}{proved}{Historical events: Claude Fable 5.2 (Anthropic), effort high, 16 November 2026: carried as background without fault found. GPT-6 Astra (OpenAI), ChatGPT Pro mode, effort Pro, S6\_09\_PAPERS, 17 November 2026: targeted proof check. Also checked in this writing. Pre-continuation writing events: GPT-6 Astra (OpenAI), ChatGPT Pro mode, effort Pro, S6\_10\_PAPERS, 18 November 2026: targeted proof check; used and rechecked in this run. Continuation mathematical and source checks: GPT-6 Astra (OpenAI), ChatGPT Pro mode, effort Pro, S6\_10\_PAPERS continuation, 19 November 2026: targeted proof check; used and rechecked in this run. Continuation metadata events: GPT-6 Astra (OpenAI), ChatGPT Pro mode, effort Pro, S6\_10\_PAPERS continuation, 19 November 2026: used and rechecked in this run; status-vocabulary reconciliation; metadata reconciliation only; no new mathematical check. \emph{Pre-continuation recorded proof scope.} The adjacent-divisor restriction and nonconstant elliptic Picard map were re-derived in the Hopf and elliptic cases. The precise source lemmas for the covering and effective-divisor classification were opened}{the hypotheses, notation and definitions stated here; \cref{cdp:eq:C.INT.3}; \cref{cdp:st:D.1.7}; \cref{cdp:eq:C.INT.1}; \cref{cdp:eq:C.INT.2}; \cite[Lemma~2.2 and \S2]{BrinzanescuMoraru2005}; \cite[Remark~2.7]{DlousskyTeleman2012}}


Let $Y$ be a smooth compact complex threefold with
$H^2(Y,\Q)=0$, let $S_i$ be an irreducible component of
a connected reducible divisor $S$, and let
$h:Z\to S_i\hookrightarrow Y$ be the minimal resolution
of its normalization. Suppose that the minimal model
$\sigma:Z\to Z_0$ is a Hopf surface or a non-K\"ahler
elliptic surface. Every strict transform $E\subset Z$ of
an elliptic curve on $Z_0$ has a one-dimensional image
under $h$, and the normalization of that image is elliptic.
The actual restriction homomorphism
$\Pic^0(Y)\to\Pic^0(E)$ is nonconstant for each such $E$.
In particular such a component cannot meet a component
whose minimal model is an Enoki surface.
\end{proposition}


\begin{proof}
\noindent\textbf{Step 1.} Choose an adjacent component $S_j$. The actual pullback
of its canonical divisor section gives a nonzero effective
Cartier divisor
\[
 A=h^*S_j>0,\qquad \OO_Z(A)=h^*\OO_Y(S_j).
 \tag{C.INT.1}\label{cdp:eq:C.INT.1}
\]
Its first Chern class is torsion. For some $q>0$,
$\OO_Y(qS_j)\in\Pic^0(Y)$, so the restriction homomorphism
$\rho:\Pic^0(Y)\to\Pic^0(Z)$ has in its image the line
$\OO_Z(qA)$. That line has strictly positive degree for
every Gauduchon metric on $Z$.

\noindent\textbf{Step 2.} We spell out the blowdown comparison. Put $A_0=\sigma_*A$.
The difference $A-\sigma^*A_0$ is exceptional and has
zero real Chern class: both $A$ and its pushforward have
zero real class. The negative definite exceptional
intersection matrix makes every coefficient of this
difference zero. Thus
\[
 A=\sigma^*A_0,\qquad A_0>0.
 \tag{C.INT.2}\label{cdp:eq:C.INT.2}
\]
Nonzeroness follows by the same negative definiteness if
$A_0$ were zero. Pullback is an isomorphism
$\Pic^0(Z_0)\simeq\Pic^0(Z)$; the integral blowup
formula and $R\sigma_*\OO_Z=\OO_{Z_0}$ identify both
the lattice and the coefficient space defining it.
We therefore regard $\rho$ as a map into $\Pic^0(Z_0)$.
Its image contains $\OO_{Z_0}(qA_0)$.

\noindent\textbf{Step 3.} We first prove the stronger assertion that the composite
restriction to $E$ cannot be zero. Suppose that it is zero.
This situation would in particular arise if $h(E)$ were
a point or had rational normalization.
In the latter case the map from the normal curve $E$
factors as $E\to\PP^1\to Y$. In both cases every line
from $\Pic^0(Y)$ restricts trivially to $E$. Here the
pullback to $\PP^1$ is in $\Pic^0(\PP^1)=0$; this is
the specified map giving the vanishing. The restriction
of $\sigma$ to $E$ is an isomorphism onto its smooth
elliptic image $E_0$.

\noindent\textbf{Step 4.} For a Hopf $Z_0$ the homomorphism
\[
 H^1(Z_0,\R)\longrightarrow H^1(E_0,\R)
 \tag{C.INT.3}\label{cdp:eq:C.INT.3}
\]
is nonzero. On a primary Hopf surface this follows from
its universal cover $\C^2\setminus\{0\}$: a finite image
of $\pi_1(E_0)$ in $\pi_1(Z_0)=\Z$ would lift a finite
unramified cover of $E_0$ to that cover. Its two coordinate
functions would be constant, contradicting the nonconstant
inclusion of the curve. For a secondary Hopf surface,
pull back along a finite unramified primary cover. A
component above $E_0$ is an elliptic curve. Transfer makes
$H^1(Z_0,\R)\to H^1(Z'_0,\R)$ an isomorphism between
one-dimensional spaces, and the preceding primary
argument proves \eqref{cdp:eq:C.INT.3} after this pullback, hence
before it. The real-to-holomorphic coefficient injection
of \ref{cdp:st:D.1.7}, applied on $E_0$, shows that
$H^1(Z_0,\OO)\to H^1(E_0,\OO)$ is nonzero. Since
$\Pic^0(Z_0)=\C^*$, the identity component of the
kernel of $\Pic^0(Z_0)\to\Pic^0(E_0)$ is zero.
Connectedness of $\Pic^0(Y)$ would therefore make
$\rho$ zero, contradicting \eqref{cdp:eq:C.INT.1}--\eqref{cdp:eq:C.INT.2}.

\noindent\textbf{Step 5.} For a non-K\"ahler elliptic $p:Z_0\to B$, every curve
on its minimal model is a reduced elliptic fibre.
The actual covering construction of
\cite[Lemma~2.2 and \S2]{BrinzanescuMoraru2005}
provides a finite cyclic map $\psi:Z'_0\to Z_0$ and a
principal elliptic bundle $p':Z'_0\to B'$ such that a
reduced fibre $E'$ above $E_0$ maps to $E_0$ by a finite
isogeny. To include multiple fibres, take a base cover
whose branching orders are divisible by their
multiplicities and normalize the fibre product. In the
local cyclic presentation of a smooth multiple fibre,
this removes the transverse multiplicity and leaves a
translation quotient on the reduced fibre, giving exactly
the stated isogeny. Extra branch points over ordinary
fibres make the total map branched there, but its map on
the reduced elliptic fibre is an isomorphism.

\noindent\textbf{Step 6.} On this principal bundle the source's explicit
presentation is
\[
 Z'_0=\Theta^*/\langle\tau\rangle,
 \quad |\tau|>1,\quad \deg\Theta>0,
 \qquad
 \Pic^0(Z'_0)=p'^*\Pic^0(B')\times\C^*.
 \tag{C.INT.4}\label{cdp:eq:C.INT.4}
\]
In the last factor a scalar $\alpha$ is the automorphy
factor of a line bundle. Its restriction to
$E'=\C^*/\langle\tau\rangle$ is that same automorphy
line; its kernel is the discrete subgroup
$\langle\tau\rangle$. Consequently the identity
component of the restriction kernel is precisely the
compact group $p'^*\Pic^0(B')$.
Our assumed trivial restriction on $E$ makes
$\psi^*\rho(\Pic^0(Y))$ lie in that identity component.
Gauduchon degree is a continuous additive homomorphism
to $\R$, and is zero on any compact subgroup. On the
other hand this image contains
$\OO_{Z'_0}(q\psi^*A_0)$, the line of a nonzero effective
divisor, which has strictly positive Gauduchon degree.
This contradiction proves the first assertion in both
cases, and proves the stated nonconstant restriction.
Riemann--Hurwitz gives genus at most one for the
normalization of the nonconstant image of $E$, so it is
elliptic.

\noindent\textbf{Step 7.} For the last assertion, an effective divisor of zero
real Chern class on a minimal class VII surface is,
by \cite[Remark~2.7]{DlousskyTeleman2012}, an elliptic
divisor on a Hopf surface or a positive multiple of the
rational cycle on an Enoki surface. If the source states
integral homological triviality, first multiply the
divisor by an integer killing its torsion integral
class and use Poincar\'e duality; the resulting support
and coefficient conclusion descends to the original
divisor. Formula \eqref{cdp:eq:C.INT.2} also applies to an Enoki
minimal model. Thus the adjacent divisor on its
resolution is the full pullback of a rational cycle;
all its curve images have rational normalization.
On a Hopf or non-K\"ahler elliptic neighbouring
resolution, the adjacent divisor is the full pullback
of a nonzero divisor consisting of elliptic curves.
At least one of their strict transforms must map
into the same intersection. It cannot map to a point
or a rational-normalized curve by the result just
proved. This is the required contradiction.
\end{proof}






\begin{proposition}[Proved exclusion of Enoki components]\label{cdp:st:C.2.INT.2}
\StatementRecord{cdp:st:C.2.INT.2}{proved}{Historical events: Claude Fable 5.2 (Anthropic), effort high, 16 November 2026: carried as background without fault found. GPT-6 Astra (OpenAI), ChatGPT Pro mode, effort Pro, S6\_09\_PAPERS, 17 November 2026: targeted proof check. Also checked in this writing. Pre-continuation writing events: GPT-6 Astra (OpenAI), ChatGPT Pro mode, effort Pro, S6\_10\_PAPERS, 18 November 2026: targeted proof check; used and rechecked in this run. Continuation mathematical and source checks: GPT-6 Astra (OpenAI), ChatGPT Pro mode, effort Pro, S6\_10\_PAPERS continuation, 19 November 2026: computation or formal certificate; targeted proof check; used and rechecked in this run. Continuation metadata events: GPT-6 Astra (OpenAI), ChatGPT Pro mode, effort Pro, S6\_10\_PAPERS continuation, 19 November 2026: used and rechecked in this run; status-vocabulary reconciliation; metadata reconciliation only; no new mathematical check. \emph{Pre-continuation recorded proof scope.} The coefficient maximum principle on Enoki components, parabolic Inoue canonical relation, the preserved external cycle, and the local Lagrangian sphere of square -2 in a torus were re-derived}{the hypotheses, notation and definitions stated here; \cref{cdp:st:C.3.4}; \cref{cdp:st:C.2.INT.1}; \cref{cdp:st:A.5.3}; \cite[Lemma~2.22 and the following paragraph]{Dloussky2006}; \cref{cdp:st:A.4.4}; \cref{cdp:eq:C.INT.5}; \cref{cdp:eq:C.INT.7}; \cref{cdp:eq:C.INT.6}}


Let $f:Y\to B$ be a proper connected-fibre map from a
smooth compact complex threefold with $H^2(Y,\Q)=0$
and general fibre a complex two-torus. No component
of a reducible fibre has normalization resolution
with Enoki minimal model.
\end{proposition}


\begin{proof}
\noindent\textbf{Step 1.} Write the full fibre as $F=\sum_i m_iS_i$. Its
normalization resolutions $Z_i$ are non-K\"ahler by
\ref{cdp:st:C.3.4}. A neighbouring component gives on each $Z_i$
a nonzero effective divisor $E_{ij}=h_i^*S_j$ of
zero real Chern class. On a class VII minimal model
the precise classification in \ref{cdp:st:C.2.INT.1} makes that
model Hopf or Enoki. The other non-K\"ahler minimal
models are elliptic surfaces, including both types
of Kodaira surfaces. \ref{cdp:st:C.2.INT.1} rules out an interface
between either such type and an Enoki component.
Thus connectedness of the fibre shows that, if an
Enoki component occurs, every component has Enoki
minimal model.

\noindent\textbf{Step 2.} Let $D_i^{\rm self}\ge0$ be the effective resolution
conductor for $S_i$ from \ref{cdp:st:A.5.3}, and set
$D_i=D_i^{\rm self}+\sum_{j\ne i}E_{ij}$. This is
also the resolution conductor for the total reduced
support. Adjunction specifies the exact comparisons
\[
 K_{Z_i}+D_i=h_i^*(K_Y+\sum_j S_j),\qquad
 K_{Z_i}+D_i^{\rm self}=h_i^*(K_Y+S_i).
 \tag{C.INT.5}\label{cdp:eq:C.INT.5}
\]
The effective numerically anticanonical divisor on
an Enoki minimal model makes it parabolic Inoue;
this is the Enoki case in
\cite[Lemma~2.22 and the following paragraph]{Dloussky2006}.
Write $E_i^0$ for its elliptic curve and $C_i$ for
its reduced rational cycle. They are disjoint and
\[
 K_{Z_i^0}+E_i^0+C_i=0.
 \tag{C.INT.6}\label{cdp:eq:C.INT.6}
\]
The pushforward of $D_i$ has the form
$E_i^0+k_i' C_i$, with $k_i'\ge0$, by intersection
with $E_i^0$ and the kernel of the cycle intersection
matrix. Each nonzero adjacent divisor has the exact
form $E_{ij}=n_{ij}\sigma_i^*C_i$, $n_{ij}\ge1$.
In particular $k_i'\ge\sum_{j\ne i}n_{ij}\ge1$.

\noindent\textbf{Step 3.} The canonical evaluation used in \ref{cdp:st:A.4.4} gives, on a
neighbourhood of this fibre,
$K_Y=\OO_Y(\sum_i k_iS_i)$ for integers $k_i\ge0$;
the line from the base is trivial on the chosen disc.
Put $r_i=(k_i+1)/m_i$. The principal full-fibre
relation yields the following equality in
$\Pic(Z_i)\otimes\Q$:
\[
 K_{Z_i}+D_i
   =\sum_{j\ne i}m_j(r_j-r_i)E_{ij}.
 \tag{C.INT.7}\label{cdp:eq:C.INT.7}
\]
Indeed subtract $r_i h_i^*F=0$ from the first
formula \eqref{cdp:eq:C.INT.5}. Choose $i$ for which $r_i$ is
maximal. The negative of the right side of
\eqref{cdp:eq:C.INT.7} is an effective rational divisor $A_i$.
Thus $K_{Z_i}+D_i+A_i$ is torsion after clearing
denominators. Push to its minimal model and use
\eqref{cdp:eq:C.INT.6}. Taking Gauduchon degree gives
\[
 k_i'+\sum_{j\ne i}m_j(r_i-r_j)n_{ij}=1,
 \tag{C.INT.8}\label{cdp:eq:C.INT.8}
\]
since $\deg\OO(C_i)>0$. All terms after $k_i'$
are nonnegative, and $k_i'\ge\sum n_{ij}\ge1$.
Equality forces $k_i'=1$, exactly one neighbour,
$n_{ij}=1$, and $r_j=r_i$. That neighbour is again
a maximum. Propagating along the connected
intersection graph proves that there are exactly
two components $S_1,S_2$, all ratios are equal,
and on both minimal models
\[
 (D_i^{\rm self})_0=E_i^0,
 \qquad (E_{ij})_0=C_i.
 \tag{C.INT.9}\label{cdp:eq:C.INT.9}
\]

\noindent\textbf{Step 4.} We retain the geometric information in the actual
blowups and normalization maps. Every topologically
torsion line descends through a point blowup, by the
integral Picard blowup formula. Applying the
canonical divisor formula to the second equality
\eqref{cdp:eq:C.INT.5} therefore gives at every blowup
\[
 (D_i^{\rm self})_{\rm new}
   =\sigma^*(D_i^{\rm self})_{\rm old}-E_{\rm exc}.
 \tag{C.INT.10}\label{cdp:eq:C.INT.10}
\]
Starting with the smooth multiplicity-one elliptic
curve $E_i^0$, effectivity permits centres only on
its current strict transform, and the new
exceptional coefficient is zero. Hence
$D_i^{\rm self}$ is that elliptic strict transform,
and all these centres are disjoint from $C_i$.
The external cycle $\sigma_i^*C_i=C_i$ is therefore
unchanged, and disjoint from the self conductor.

\noindent\textbf{Step 5.} It follows that $S_i$ is normal near the external
intersection. Here is the needed resolution
comparison at possible exceptional points. Any
resolution-exceptional curve outside $D_i$ is
disjoint from it: adjunction and minimality give
$K_{Z_i}\cdot E\ge0$, while $D_i\cdot E\ge0$,
so $(K_{Z_i}+D_i)\cdot E=0$ forces $D_i\cdot E=0$.
Connected resolution fibres cannot join the
elliptic self conductor to the external rational
cycle. The conductor on the normal surface is a
reflexive ideal, by the duality calculation in
\ref{cdp:st:A.5.3}, so it has no isolated additional support
on that cycle. Thus normalization is an isomorphism
near its image. Resolution over this image is
crepant by the second formula \eqref{cdp:eq:C.INT.5}, since
$D_i^{\rm self}$ is disjoint from it. Its
exceptional curves are consequently $(-2)$
curves in $C_i$. A whole cycle cannot be contracted,
because its intersection matrix is not negative
definite. The connected fibres which are
contracted are proper chains in that cycle.

\noindent\textbf{Step 6.} The intersection scheme $D=S_1\cap S_2$ is a
complete-intersection curve in $Y$. Its pullback
has reduced divisor $C_i$, so it is generically
reduced; its Cohen--Macaulay property then makes
it reduced everywhere. It is connected, since it
is the image of $C_i$. The equality of all $r_i$
means $K_Y+S_1+S_2=rF$ in $\Pic(Y)\otimes\Q$
on the fibre neighbourhood. Adjunction gives
$\omega_D$ torsion, hence $p_a(D)=1$.
Topologically the map $C_i\to D$ contracts just
the proper chains described above and is an
isomorphism elsewhere. It preserves the graph
circle and all its remaining components are
rational. Thus its normalized components are
$\PP^1$, its graph has exactly one cycle, and
$\sum\delta_x=\sum(r_x-1)$. Each contracted
chain has only two ends; every singular point
therefore has two branches and $\delta=1$.
The colength-one ring comparison (C.VII.6)
proves that these are ordinary nodes. At least
one node occurs.

\noindent\textbf{Step 7.} At such a node the complete-intersection curve
has embedding dimension two. At least one of
the two surface equations has a nonzero linear
part; take it to be $z$. On $z=0$ the other
equation has an ordinary node. The holomorphic
Morse lemma with parameter $z$ therefore writes
it as $xy+h(z)$. Normality of that surface near
$D$ excludes $h\equiv0$, and a change of the
$z$ coordinate writes $h(z)=z^n$, $n\ge1$.
Consequently the original fibration has local
equation
\[
 f=u(x,y,z)\,z^a(xy+z^n)^b,
 \qquad a,b,n\ge1,
 \tag{C.INT.11}\label{cdp:eq:C.INT.11}
\]
where $u$ is a unit and $a,b$ are the two full
fibre multiplicities. The unit can be removed
by an actual biholomorphic change of variables.
Put $d=2(a+nb)$ and choose a local logarithm of
$u$. For $f_0=z^a(xy+z^n)^b$, weighted homogeneity
gives
\[
 \Phi(x,y,z)=(x u^{n/d},y u^{n/d},z u^{2/d}),
 \qquad f_0\circ\Phi=u f_0.
 \tag{C.INT.12}\label{cdp:eq:C.INT.12}
\]
The derivative of $\Phi$ at the origin is
invertible, proving the required equivalence
of the actual local fibres.

\noindent\textbf{Step 8.} For $t\ne0$ small, take a simply connected
sector in the punctured $z$-plane and one
branch of $(t z^{-a})^{1/b}$. An open part
of the actual smooth fibre is
\[
 xy=g(z):=(t z^{-a})^{1/b}-z^n.
 \tag{C.INT.13}\label{cdp:eq:C.INT.13}
\]
Consecutive zeros on this logarithmic branch
have angular separation $2\pi b/(a+nb)<2\pi$;
two such zeros fit in one sector. They are
simple, since at either one
$g'(z)=-(a/b+n)z^{n-1}\ne0$.
Join them by a smooth embedded path with no
other zero and take the circle $|x|=|y|$
over its interior. Explicitly, for a continuous
choice of $\arg g$ along that interior, it is
\[
 x=|g|^{1/2}e^{i(\arg g/2+\theta)},\qquad
 y=|g|^{1/2}e^{i(\arg g/2-\theta)}.
 \tag{C.INT.14}\label{cdp:eq:C.INT.14}
\]
The circle collapses to a smooth disc at each
simple zero, giving an embedded sphere.
The two $x,y$ terms of the standard K\"ahler
form cancel on this sphere, and the $z$ term
vanishes on the one-dimensional path. The
sphere is Lagrangian; its normal bundle is
$JTS^2$ with the opposite Euler orientation,
so its self-intersection is $-2$.
The whole sphere lies in the local chart:
$|z|=O(|t|^{1/(a+nb)})$ and
$|x|,|y|=O(|t|^{n/(2(a+nb))})$.
Its image in the compact smooth torus fibre
would be null-homotopic, because a torus has
contractible universal cover. Its integral
homology class would therefore be zero and
its self-intersection zero, contradicting
$-2$. This contradiction excludes the Enoki
alternative.
\end{proof}






\begin{proposition}[Proved sign and equality for every reducible fibre]\label{cdp:st:C.2.INT.3}
\StatementRecord{cdp:st:C.2.INT.3}{proved}{Historical events: Claude Fable 5.2 (Anthropic), effort high, 16 November 2026: carried as background without fault found. Claude Opus 5.5 (Anthropic), effort max, sweep 1, scope s6, 18 November 2026: used and rechecked in this run. GPT-6 Astra (OpenAI), ChatGPT Pro mode, effort Pro, S6\_09\_PAPERS, 17 November 2026: targeted proof check. Also checked in this writing. Pre-continuation writing events: GPT-6 Astra (OpenAI), ChatGPT Pro mode, effort Pro, S6\_10\_PAPERS, 18 November 2026: targeted proof check; used and rechecked in this run. Continuation mathematical and source checks: GPT-6 Astra (OpenAI), ChatGPT Pro mode, effort Pro, S6\_10\_PAPERS continuation, 19 November 2026: targeted proof check. Continuation metadata events: GPT-6 Astra (OpenAI), ChatGPT Pro mode, effort Pro, S6\_10\_PAPERS continuation, 19 November 2026: used and rechecked in this run; status-vocabulary reconciliation; metadata reconciliation only; no new mathematical check. \emph{Pre-continuation recorded proof scope.} The factorization through the minimal model, conductor direct images, generic elliptic twists, curve Euler identity, component complement inequality, and every equality implication were re-derived}{the hypotheses, notation and definitions stated here; \cref{cdp:st:C.2.INT.2}; \cref{cdp:st:C.2.INT.1}; \cref{cdp:st:A.5.3}; \cref{cdp:eq:C.INT.18}; \cref{cdp:eq:C.INT.21}; \cref{cdp:eq:C.INT.20}; \cref{cdp:eq:C.INT.22}; \cref{cdp:eq:C.INT.26}; \cref{cdp:eq:C.INT.27}}


Let $f:Y\to B$ satisfy the hypotheses of \ref{cdp:st:C.2.INT.2},
and let $S$ be a connected reducible fibre support.
Then
\[
 e(S)\ge0.
 \tag{C.INT.15}\label{cdp:eq:C.INT.15}
\]
Equality holds precisely when the normalization of $S$
is a disjoint union of smooth minimal Hopf or
non-K\"ahler elliptic surfaces, and the reduced
conductor image and its reduced inverse image are
disjoint unions of smooth elliptic curves. The maps
between these elliptic curves are finite unramified
covers. This criterion allows arbitrary conductor
scheme structure and arbitrary unequal multiplicities
of the full fibre.
\end{proposition}


\begin{proof}
\noindent\textbf{Step 1.} Let $\nu:T=\coprod T_i\to S$ be normalization,
$\pi_i:Z_i\to T_i$ its minimal resolutions, and
$\sigma_i:Z_i\to Z_i^0$ the minimal-model maps.
By \ref{cdp:st:C.2.INT.1}--\ref{cdp:st:C.2.INT.2} the $Z_i^0$ are Hopf or
non-K\"ahler elliptic surfaces. They have no rational
curves and have both Euler characteristics
$e(Z_i^0)=\chi(\OO_{Z_i^0})=0$.
Every rational curve of $Z_i$ is exceptional for
$\sigma_i$; their total number is exactly the number
$N_i$ of point blowups, and $e(Z_i)=N_i$.
The elliptic strict transforms cannot be contracted
by $\pi_i$, by \ref{cdp:st:C.2.INT.1}. Thus every
$\pi_i$-exceptional curve is $\sigma_i$-exceptional.

\noindent\textbf{Step 2.} The connected fibres of $\pi_i$ consequently map to
points under $\sigma_i$. Normality and
$(\pi_i)_*\OO_{Z_i}=\OO_{T_i}$ descend the latter
map to a proper birational morphism
$\tau_i:T_i\to Z_i^0$, with
\[
 \sigma_i=\tau_i\pi_i.
 \tag{C.INT.16}\label{cdp:eq:C.INT.16}
\]
This is an actual factorization, not a replacement
of the normalization by its minimal model.
Since point blowups have $R\sigma_{i*}\OO_{Z_i}
=\OO_{Z_i^0}$, the low-degree Leray sequence gives
\[
 0\longrightarrow R^1\tau_{i*}\OO_{T_i}
 \longrightarrow 0
 \longrightarrow \tau_{i*}R^1\pi_{i*}\OO_{Z_i}
 \longrightarrow R^2\tau_{i*}\OO_{T_i}=0.
 \tag{C.INT.17}\label{cdp:eq:C.INT.17}
\]
The last vanishing follows from fibre dimension at
most one. The middle sheaf $R^1\pi_{i*}\OO_{Z_i}$
is supported at finitely many points, and its
pushforward under $\tau_i$ detects it. Hence it is
zero. This proves that $T_i$ has rational
singularities, and in particular
$\chi(\OO_{T_i})=\chi(\OO_{Z_i})=0$.

\noindent\textbf{Step 3.} Let $\Gamma$ be the full conductor image scheme,
$A_i\subset T_i$ its full inverse-image conductor,
and let $D_i\ge0$ be the resolution conductor from
\ref{cdp:st:A.5.3}. Its line relation
$K_{Z_i}+D_i=h_i^*(K_Y+S)$ has torsion topological
first Chern class. That line descends through
$\sigma_i$. Consequently the exact canonical
divisor formula gives
\[
 D_i=\sigma_i^*D_i^0-K_{Z_i/Z_i^0},
 \qquad D_i^0=(\sigma_i)_*D_i.
 \tag{C.INT.18}\label{cdp:eq:C.INT.18}
\]
The divisor $D_i^0$ is nonzero, since $D_i$ contains
a nonzero adjacent component divisor whose
minimal-model pushforward is nonzero by (C.INT.2).
It is a sum of positive multiples of disjoint
smooth elliptic curves. At each point blowup,
\eqref{cdp:eq:C.INT.18} has coefficient
$\operatorname{mult}_pD_{\rm old}-1$ on the new
exceptional curve. Effectivity makes the centre
lie on the old support. Starting from the disjoint
elliptic curves, the reduced support consequently
acquires only rational leaves and subdivisions of
edges. Each of its connected components is one
elliptic curve with rational trees attached.
Writing $r_i$ for its number of rational curves,
we obtain
\[
 e((D_i)_{\rm red})=r_i.
 \tag{C.INT.19}\label{cdp:eq:C.INT.19}
\]

\noindent\textbf{Step 4.} We next prove a simultaneous cohomological
comparison on the entire, possibly nonreduced,
conductor. Relative duality applied to
$R\pi_{i*}\OO_{Z_i}=\OO_{T_i}$ gives
$R\pi_{i*}\omega_{Z_i}=\omega_{T_i}$.
Finite conductor duality identifies
$\omega_{T_i}\otimes\nu^*\omega_S^{-1}$ with
the conductor ideal $\mathcal I_{A_i/T_i}$.
The divisor sequence on $Z_i$, using
$\OO_{Z_i}(-D_i)=\omega_{Z_i}\otimes h_i^*\omega_S^{-1}$,
therefore gives
\[
 R\pi_{i*}\OO_{D_i}=\OO_{A_i}.
 \tag{C.INT.20}\label{cdp:eq:C.INT.20}
\]
The identical calculation for $\sigma_i$, using
\eqref{cdp:eq:C.INT.18}, gives
\[
 R\sigma_{i*}\OO_{D_i}=\OO_{D_i^0},
 \qquad R\tau_{i*}\OO_{A_i}=\OO_{D_i^0}.
 \tag{C.INT.21}\label{cdp:eq:C.INT.21}
\]
For completeness the canonical pushforward for a
point blowup is elementary: its relative canonical
line is $\OO(E)$, and
$0\to\OO\to\OO(E)\to\OO_E(-1)\to0$
has zero cohomology in its last term. It gives
$R\sigma_*\omega=\omega$; composition proves the
version used in \eqref{cdp:eq:C.INT.21}.

\noindent\textbf{Step 5.} For $M\in\Pic^0(Y)$ its pullback to $Z_i$ descends
to a line $M_i^0\in\Pic^0(Z_i^0)$, by the Picard
blowup isomorphism. Projection formula and
\eqref{cdp:eq:C.INT.20}--\eqref{cdp:eq:C.INT.21} identify the actual spaces
\[
 H^0(A_i,\nu^*M|_{A_i})
   =H^0(D_i^0,M_i^0|_{D_i^0}).
 \tag{C.INT.22}\label{cdp:eq:C.INT.22}
\]
For a component $nE$ of $D_i^0$, the filtration by
its powers of the defining ideal has graded lines
$M_i^0|_E\otimes\OO_E(-jE)$, $0\le j<n$.
They have degree zero, since $E^2=0$ on the
minimal model. \ref{cdp:st:C.2.INT.1} proves that the map from
$\Pic^0(Y)$ to $\Pic^0(E)$ is nonconstant. We can
therefore choose one $M$ avoiding, for each of
the finitely many such curves and integers $j$,
the prescribed value $\OO_E(jE)$.
Each excluded locus is a proper closed analytic
subset of the connected complex Lie group
$\Pic^0(Y)$; their finite union does not fill it.
All the graded lines are nontrivial of degree
zero on an elliptic curve and have no sections.
The filtration and \eqref{cdp:eq:C.INT.22} prove
\[
 H^0(A_i,\nu^*M)=0\quad\hbox{for every }i.
 \tag{C.INT.23}\label{cdp:eq:C.INT.23}
\]

\noindent\textbf{Step 6.} The conductor inclusion is
$\OO_\Gamma\hookrightarrow\nu_*\OO_A$, where
$A=\coprod A_i$. It now gives $H^0(\Gamma,M)=0$.
The normalization sequence and the preceding
rational-singularity computation give
$\chi(\mathcal Q)=0$ for
$\mathcal Q=\nu_*\OO_T/\OO_S$: here
$\chi(\OO_S)=0$ is the Cartier-divisor
Riemann--Roch calculation under $H^2(Y,\Q)=0$.
Conductor duality is
$\mathcal Q=\omega_\Gamma\otimes\omega_S^{-1}|_\Gamma$.
Twisting on this compact Cohen--Macaulay curve by
an ambient topologically torsion line preserves
Euler characteristic; its degree on every
component is zero. Serre duality thus gives
\[
 \chi(\OO_\Gamma)=-\chi(\mathcal Q)=0.
 \tag{C.INT.24}\label{cdp:eq:C.INT.24}
\]
The same holds after twisting by $M$, so
$H^0(\Gamma,M)=0$ implies $H^1(\Gamma,M)=0$.
Restriction to $C=\Gamma_{\rm red}$ surjects
on $H^1$, since its kernel is supported in
dimension at most one and has no $H^2$.
Consequently
\[
 \chi(\OO_C)=\chi(M|_C)=h^0(C,M|_C)\ge0.
 \tag{C.INT.25}\label{cdp:eq:C.INT.25}
\]
This is where the full conductor scheme is used;
no passage to its reduction is assumed harmless.

\noindent\textbf{Step 7.} Each normalized component of $C$ has genus zero
or one, since it is dominated by a curve of
$(D_i)_{\rm red}$. Let $n_0$ be the number of
rational-normalized components, and let
$\delta_x$ and $s_x$ be the delta invariant and
branch number of a singular point of $C$.
The exact curve normalization formulas give
\[
 e(C)=\chi(\OO_C)+n_0+
                \sum_x\bigl(\delta_x-(s_x-1)\bigr)\ge0.
 \tag{C.INT.26}\label{cdp:eq:C.INT.26}
\]
Every summand is nonnegative, because
$\delta_x\ge s_x-1$. Equality holds precisely
when $C$ is a disjoint union of smooth elliptic
curves. Indeed equality first forces $n_0=0$;
then $\chi(\OO_C)=-\sum\delta_x\ge0$
forces every $\delta_x=0$. The converse is immediate.

\noindent\textbf{Step 8.} We finish the topological comparison on each
normalization component. A $\pi_i$-exceptional
curve outside $(D_i)_{\rm red}$ is a $(-2)$ curve
disjoint from that support, by the adjunction
argument following (C.INT.10). Each connected
configuration of such curves is a tree: a cycle
or a repeated intersection would contradict
negative definiteness of its matrix with
all diagonal entries $-2$. Such fibres lie
off $A_i$, because a fibre over its conductor
meets the strict conductor and is connected.
Let $a_i$ count their irreducible curves, and
let $b_i$ count exceptional fibres contained
in $(D_i)_{\rm red}$ whose image lies outside
$A_i$. Removing the conductor and contracting
these specified fibres gives the exact equality
\[
 e(T_i\setminus(A_i)_{\rm red})
       =N_i-r_i-a_i+b_i\ge0.
 \tag{C.INT.27}\label{cdp:eq:C.INT.27}
\]
A tree of $a$ rational curves has Euler number
$a+1$, so its contraction removes $a$; a
positive-conductor fibre removed upstairs but
represented by a point downstairs adds one.
These are all changes, since normalization is
an isomorphism off the conductor. Nonnegativity
follows because the $r_i$ support curves and
the $a_i$ disjoint exceptional curves are
distinct rational curves among the $N_i$
rational curves of $Z_i$.
The normalization is an isomorphism
$T\setminus A\simeq S\setminus C$. Additivity
with \eqref{cdp:eq:C.INT.26}--\eqref{cdp:eq:C.INT.27} proves \eqref{cdp:eq:C.INT.15}.

\noindent\textbf{Step 9.} If $e(S)=0$, every nonnegative summand in these
formulas is zero. In particular $C$ is a
union of smooth elliptic curves, $b_i=0$,
and $r_i+a_i=N_i$. Every rational curve in
$(D_i)_{\rm red}$ must then be $\pi_i$-exceptional:
a rational curve cannot map nonconstantly to
one of these elliptic curves. The $a_i$ curves
are exceptional by definition. Thus all $N_i$
rational curves of $Z_i$ are $\pi_i$-exceptional.
If $N_i>0$, the last point blowup in
$\sigma_i$ supplies a $(-1)$ curve on $Z_i$.
Its being $\pi_i$-exceptional contradicts the
minimality of that resolution. Hence $N_i=0$.
There are now no rational curves to contract,
and \ref{cdp:st:C.2.INT.1} excludes contracting elliptic
curves. Therefore $Z_i=T_i=Z_i^0$ is smooth
and minimal. Its reduced conductor is a
disjoint union of elliptic curves. A finite
map between smooth elliptic curves is
unramified by Riemann--Hurwitz, proving the
stated map comparison.
Conversely these normalization and conductor
conditions give $e(T_i)=e((A_i)_{\rm red})=e(C)=0$;
the normalization formula gives $e(S)=0$.
The underlying spaces, and all comparisons
just used, are independent of the positive
full-fibre multiplicities. This completes the
proof including the equality criterion.
\end{proof}






\begin{proposition}[The analytic support comparison used here]\label{cdp:st:C.2.MH.1}
\StatementRecord{cdp:st:C.2.MH.1}{proved}{Historical events: Claude Fable 5.2 (Anthropic), effort high, 16 November 2026: carried as background without fault found. GPT-6 Astra (OpenAI), ChatGPT Pro mode, effort Pro, S6\_09\_PAPERS, 17 November 2026: targeted proof check. Also checked in this writing. Pre-continuation writing events: GPT-6 Astra (OpenAI), ChatGPT Pro mode, effort Pro, S6\_10\_PAPERS, 18 November 2026: targeted proof check; used and rechecked in this run. Continuation mathematical and source checks: GPT-6 Astra (OpenAI), ChatGPT Pro mode, effort Pro, S6\_10\_PAPERS continuation, 19 November 2026: targeted proof check; used and rechecked in this run. Continuation metadata events: GPT-6 Astra (OpenAI), ChatGPT Pro mode, effort Pro, S6\_10\_PAPERS continuation, 19 November 2026: used and rechecked in this run; status-vocabulary reconciliation; metadata reconciliation only; no new mathematical check. \emph{Pre-continuation recorded proof scope.} The quotient logarithmic complex, local residues, support cone, and support-to-restriction map were re-derived. Steenbrink 1999 section 3 Theorem 1 and Corollary 2 and section 4 equation (3) were opened. Its algebraic-ambient hypothesis is replaced here by an explicit construction on compact projective strata; finite characters and degree-one trace were checked separately}{the hypotheses, notation and definitions stated here; \cref{cdp:eq:C.MH.2}; \cite[Section~3, Theorem~1 and
Corollary~2]{Steenbrink1999}; \cite[Section~4, equation~(3)]{Steenbrink1999}; \cref{cdp:eq:C.MH.3}; \cref{cdp:eq:C.MH.5}}

Let $U$ be a smooth complex threefold, let $S\subset U$ be a
compact divisor, and suppose that an embedded resolution
$r:U'\to U$ is a sequence of blowups with smooth compact
centres in $S$, with $D=(r^{-1}S)_{\rm red}$ a simple
normal crossings divisor whose components are projective.
Let $V$ be a finite-monodromy rational local system on $U$.
Suppose that the cohomology of the punctured sufficiently
small fibre tubes around $S$, with coefficients in $V$,
vanishes in every degree. Then the mixed Hodge structure
on $H^k(S,V)$ obtained from a projective hyperresolution
is pure of weight $k$.

In this assertion the support-to-restriction morphism is
the actual morphism
\[
 R\Gamma_D(U',r^*V)\longrightarrow
 \underset{D\subset W}{\mathop{\rm colim}}\,
 R\Gamma(W,r^*V)=R\Gamma(D,r^*V).
 \tag{C.MH.1}\label{cdp:eq:C.MH.1}
\]
Neither a K\"ahler metric on $U$ nor a projective
compactification of $U$ is needed for this comparison.
\end{proposition}


\begin{proof}
\noindent\textbf{Step 1.} First take constant rational coefficients. Write $D_I$
for the intersection of the components indexed by a
nonempty set $I$, and $D^{[p]}=\coprod_{|I|=p+1}D_I$.
All these spaces are smooth and projective. The ordinary
\v Cech de Rham complex on $D^{[\bullet]}$, with its
\v Cech weight filtration and its differential-form
filtration, is therefore a cohomological mixed Hodge
complex $K(D)$ for $\mathbb Q_D$. Its cohomology in
degree $k$ has weights at most $k$.

\noindent\textbf{Step 2.} Here is the precise analytic construction of the other
complex and of the arrow. For each stratum $C=D_I$, put
\[
 K(C\log D)_{\mathbb C}
   =\Omega_{U'}^\bullet(\log D)\otimes\OO_C.
 \tag{C.MH.2}\label{cdp:eq:C.MH.2}
\]
The tensor product in \eqref{cdp:eq:C.MH.2} is a complex: in local
normal crossings coordinates, the ideal of $C$ is
preserved by the logarithmic differential. Endow it with
the quotient residue filtration and the filtration by
form degree. The quotient map from the restriction of
$\Omega_{U'}^\bullet(\log D)$ is a filtered
quasi-isomorphism to \eqref{cdp:eq:C.MH.2}. Its rational realization is
$i_C^*Rj_*\mathbb Q_{U'\setminus D}$, where $j$ is the
open inclusion. To see the comparison, apply the
iterated Poincar\'e residue: its residue-degree-$a$
quotient is a direct sum of the de Rham complexes of
$D_J\cap C$, shifted by $[-a]$ and with Tate twist
$(-a)$. The ideal and its differentials restrict these
forms to precisely $D_J\cap C$. This is the local
coordinate proof of \cite[Section~3, Theorem~1 and
Corollary~2]{Steenbrink1999}; it uses holomorphic
coordinates and the projectivity, or compact K\"ahler
property, of the displayed strata. It does not use a
compactification of their ambient neighbourhood.

\noindent\textbf{Step 3.} Totalizing \eqref{cdp:eq:C.MH.2} over $D^{[\bullet]}$ gives the
cohomological mixed Hodge complex $K(D\log D)$, with
rational realization $i_D^*Rj_*\mathbb Q_{U'\setminus D}$.
The natural inclusion $K(D)\to K(D\log D)$ is a
morphism of mixed Hodge complexes. Its mixed cone,
shifted by $[-1]$, computes cohomology with support in
$D$; the cone projection is exactly \eqref{cdp:eq:C.MH.1}. This is
the excision comparison of
\cite[Section~4, equation~(3)]{Steenbrink1999}, applied
on the compact divisor. All complexes just used are
defined on that divisor, and all purity checks are on
its projective strata.

\noindent\textbf{Step 4.} For completeness, the weight bound on the source can
be checked without invoking an invariant-cycle theorem.
The closed-support Mayer--Vietoris complex, together
with the holomorphic Thom isomorphisms, gives
\[
 E_1^{-p,q}=
 \bigoplus_{|I|=p+1}
 H^{q-2p-2}(D_I,\mathbb Q)(-p-1)
 \ \Longrightarrow\ H_D^{q-p}(U',\mathbb Q),
 \qquad p\ge0.
 \tag{C.MH.3}\label{cdp:eq:C.MH.3}
\]
The first differential is the alternating sum of Gysin
maps for $D_I\hookrightarrow D_{I\setminus\{i\}}$.
The comparison with the preceding mixed cone can be
verified by induction on the number of components.
For one smooth divisor the residue exact sequence
identifies
\[
 \operatorname{Cone}\bigl[\Omega_{U'}^\bullet
       \longrightarrow\Omega_{U'}^\bullet(\log D_i)\bigr][-1]
       \simeq\Omega_{D_i}^\bullet[-2](-1),
\]
where $(-1)$ records both the shift of the Hodge
filtration by one and the rational $2\pi i$ residue.
For a union $A\cup B$, use the actual support triangle
$R\Gamma_{A\cap B}\to R\Gamma_A\oplus R\Gamma_B
\to R\Gamma_{A\cup B}$. Iterated residues commute
with its inclusion maps, with the alternating signs
fixed by the ordering of the omitted coordinates.
Induction identifies its totalization with the same
mixed cone. These comparisons agree on overlaps,
since they are the residue maps for the holomorphic
divisor line bundles. Thus \eqref{cdp:eq:C.MH.3} compares the actual support
complex, including its arrow to $K(D)$, rather than
only its unfiltered vector spaces. Its displayed term
has pure weight $q$. In total degree $k=q-p$ this is
weight $k+p\ge k$. It follows that $H_D^k(U',\mathbb Q)$
has weights at least $k$.

\noindent\textbf{Step 5.} Finite coefficients are handled by taking a finite
\'{e}tale cover on which $V$ is constant and then a
rational direct summand of a finite direct sum of its
regular representation.
The inverse images of the $D_I$ are still projective:
a finite holomorphic \'{e}tale cover of a
projective manifold is projective, since the pullback
of an ample bundle is positive. The deck action acts
on the displayed mixed Hodge complexes. Rational
idempotents consequently give exactly the same
comparison and weight bounds with coefficients in $V$.

\noindent\textbf{Step 6.} By the assumed punctured-tube acyclicity, the support
long exact sequence makes \eqref{cdp:eq:C.MH.1} an isomorphism on
cohomology. Properness of $r$ identifies the punctured
tubes upstairs and downstairs, since its centres lie
in $S$. The map is a morphism of mixed Hodge structures,
whose source has weights at least $k$ and whose target
has weights at most $k$. Strictness therefore proves
that $H^k(D,r^*V)$ is pure of weight $k$.

\noindent\textbf{Step 7.} We finally make the comparison with the original
support explicit. The proper degree-one map between
the smooth ambient manifolds has the trace map, defined
by oriented Poincar\'e duality and proper pushforward,
and the projection formula gives
\[
 r_*r^*=\operatorname{id}:H^k(U,V)\longrightarrow H^k(U,V).
 \tag{C.MH.4}\label{cdp:eq:C.MH.4}
\]
The same identity holds on every smaller ambient
neighbourhood. Passing to the direct limit of these
identities gives an injection
\[
                  r^*:H^k(S,V)\hookrightarrow H^k(D,r^*V).
 \tag{C.MH.5}\label{cdp:eq:C.MH.5}
\]
The identifications in this limit are the usual
continuity identifications for a compact triangulable
subset. They can equally be made using cofinal
triangulated regular neighbourhoods retracting onto
$S$ and $D$. Properness ensures that inverse images of
neighbourhoods are cofinal; thus the map in \eqref{cdp:eq:C.MH.5}
is the pullback by $D\to S$.

\noindent\textbf{Step 8.} This last map is a mixed Hodge morphism as well. One
may construct the projective hyperresolution of $S$
from its projective normalization resolution and
resolutions of the conductor and exceptional curves.
To compare it with $D^{[\bullet]}$, resolve their
fibre products over $S$. Each such fibre product is a
closed analytic subspace of a product of projective
spaces already occurring in the two resolutions;
Chow's theorem makes it algebraic and projective,
and its resolution is projective. The resulting
common hyperresolution realizes \eqref{cdp:eq:C.MH.5} by pullbacks
between projective strata. Cohomological descent then
identifies this with the stated mixed Hodge structures.
Injectivity into the pure target proves the assertion.
\end{proof}






\begin{proposition}[Euler number from a ruled normalization]\label{cdp:st:C.2.MH.2}
\StatementRecord{cdp:st:C.2.MH.2}{proved}{Historical events: Claude Fable 5.2 (Anthropic), effort high, 16 November 2026: carried as background without fault found. GPT-6 Astra (OpenAI), ChatGPT Pro mode, effort Pro, S6\_09\_PAPERS, 17 November 2026: targeted proof check. Also checked in this writing. Pre-continuation writing events: GPT-6 Astra (OpenAI), ChatGPT Pro mode, effort Pro, S6\_10\_PAPERS, 18 November 2026: targeted proof check. Continuation mathematical and source checks: GPT-6 Astra (OpenAI), ChatGPT Pro mode, effort Pro, S6\_10\_PAPERS continuation, 19 November 2026: computation or formal certificate; targeted proof check. Continuation metadata events: GPT-6 Astra (OpenAI), ChatGPT Pro mode, effort Pro, S6\_10\_PAPERS continuation, 19 November 2026: used and rechecked in this run; status-vocabulary reconciliation; metadata reconciliation only; no new mathematical check. \emph{Pre-continuation recorded proof scope.} The self-dual finite-character augmentation system, torus Koszul acyclicity, support duality, conductor weight separation, negative-definite exceptional Gysin map, and formula d e(S)=2h4(Z,V)-h3(Z,V) were re-derived. This gives 2 in the rational case and 2-2g in the ruled case}{the hypotheses, notation and definitions stated here; \cref{cdp:st:D.1.9}; \cref{cdp:st:D.1.8}; \cref{cdp:st:C.2.MH.1}; \cref{cdp:eq:C.MH.9}; \cref{cdp:eq:C.MH.8}; \cref{cdp:eq:C.MH.13}; \cref{cdp:eq:C.MH.11}; \cref{cdp:eq:C.MH.14}}

Let $f:Y\to\PP^1$ satisfy the full CDP hypotheses
$b_1(Y)=b_2(Y)=0$ with smooth torus general fibre.
Let $mS$ be a full fibre with irreducible reduced support,
and let $Z\to T\to S$ be a resolution of its
normalization. If $Z$ is rational, then
\[
                              e(S)=2.
 \tag{C.MH.6}\label{cdp:eq:C.MH.6}
\]
If $Z$ is ruled over a curve of genus $g\ge1$, then
\[
                              e(S)=2-2g.
 \tag{C.MH.7}\label{cdp:eq:C.MH.7}
\]
The calculation retains all multiplicities in the full
fibre and places no normal crossings assumption on $S$.
\end{proposition}


\begin{proof}
\noindent\textbf{Step 1.} The integral class of \ref{cdp:st:D.1.9} restricts nontrivially to
$S$ and, by \ref{cdp:st:D.1.8}, nontrivially to the nearby torus.
More precisely, restrict the section
$\gamma\in H^0(\PP^1,R^1f_*\mathbb Z)$ to a small
disc $\Delta$ about this fibre. The integral Leray edge
isomorphism
$H^1(f^{-1}\Delta,\mathbb Z)
 \simeq H^0(\Delta,R^1f_*\mathbb Z)$ supplies the
actual integral class on the tube: the two adjacent
base groups $H^1(\Delta,\mathbb Z)$ and
$H^2(\Delta,\mathbb Z)$ vanish.
Choose a prime $\ell$ for which its restriction to that
torus is nonzero modulo $\ell$. On a sufficiently small
fibre tube $U$, its reduction defines a connected cyclic
\'{e}tale cover with group $\mathbb Z/\ell$. Let $V$ be
the rational augmentation local system, that is, the
kernel of the trace from the regular representation
to the constant representation. Its rank is
$d=\ell-1$, it is self-dual, and $H^0(S,V)=0$.
Over $\mathbb C$ it is the sum of the nontrivial
characters of this cyclic group. Since $\ell$ is prime,
each of these characters restricts nontrivially to
the nearby torus $F$. The group-cochain complex for
$\mathbb Z^4$ with such a character is acyclic: one of
the four commuting operators $\rho(e_i)-1$ is
invertible, and its inverse gives a contraction of the
Koszul complex. Thus
\[
 H^j(F,V)=0\quad\hbox{for every }j,
 \qquad H^j(U\setminus S,V)=0\quad\hbox{for every }j.
 \tag{C.MH.8}\label{cdp:eq:C.MH.8}
\]
The second assertion is the Leray sequence of the
punctured-disc torus bundle. This argument uses the
smooth fibres of the original map, so it also applies
when the full central fibre is $mS$.

\noindent\textbf{Step 2.} Choose an embedded resolution of $(U,S)$. Its strict
transform is a blowup of $Z$, hence projective. Every
other exceptional surface is obtained by blowups from
a projective bundle over a compact smooth curve, or
from $\PP^2$. It is therefore projective. All centres
are compact and lie in the central fibre. \ref{cdp:st:C.2.MH.1}
applies and gives
\[
                         H^3(S,V)\text{ pure of weight }3.
 \tag{C.MH.9}\label{cdp:eq:C.MH.9}
\]

\noindent\textbf{Step 3.} Let $A\subset T$ and $\Gamma\subset S$ be the reduced
conductor inverse image and image. The finite
normalization square is an actual topological pushout.
Its cohomological descent sequence, with the pulled
back local system, gives an exact sequence of mixed
Hodge structures
\[
\begin{gathered}
0\longrightarrow C_V\longrightarrow H^3(S,V)
  \longrightarrow H^3(T,V)\longrightarrow0,\\
C_V=\operatorname{coker}\bigl[H^2(T,V)\oplus H^2(\Gamma,V)
                         \longrightarrow H^2(A,V)\bigr].
\end{gathered}
 \tag{C.MH.10}\label{cdp:eq:C.MH.10}
\]
This also covers the normal case, with the conductor
terms zero. The left term is pure of weight two:
$H^2(A,V)$ is the top cohomology of a compact curve
with a finite local system, and hence pure of weight
two. Strictness and \eqref{cdp:eq:C.MH.9} make that term zero.
Consequently
\[
                           H^3(S,V)\simeq H^3(T,V).
 \tag{C.MH.11}\label{cdp:eq:C.MH.11}
\]

\noindent\textbf{Step 4.} For the resolution $\pi:Z\to T$, the restriction
$H^2(Z,V)\to H^2(E,V)$ to its exceptional curve is
surjective. Indeed $V$ is constant on each exceptional
fibre, since it is pulled back from $T$; the Gysin
classes of its irreducible exceptional curves restrict
by their negative-definite intersection matrix. That
matrix is invertible over $\mathbb Q$, separately
for each exceptional fibre and each coefficient
vector. The resolution descent sequence therefore
gives
\[
 H^3(T,V)\simeq H^3(Z,V),\qquad
 H^4(S,V)\simeq H^4(T,V)\simeq H^4(Z,V).
 \tag{C.MH.12}\label{cdp:eq:C.MH.12}
\]
The last two identifications also follow directly
because all the omitted loci have real dimension two.

\noindent\textbf{Step 5.} There is a second, purely topological consequence of
\eqref{cdp:eq:C.MH.8}. The support sequence and oriented local
Poincar\'e duality give
\[
 H^k(S,V)\simeq H_S^k(U,V)
             \simeq H_{6-k}(S,V).
 \tag{C.MH.13}\label{cdp:eq:C.MH.13}
\]
Again these are the natural direct-limit maps on
sufficiently small tubes. Since $S$ has real
dimension four, \eqref{cdp:eq:C.MH.13} gives $H^1(S,V)=0$ and,
using the self-duality of $V$,
$h^2(S,V)=h^4(S,V)$. Therefore \eqref{cdp:eq:C.MH.11}--\eqref{cdp:eq:C.MH.13}
give the exact numerical identity
\[
 d\,e(S)=\chi(S,V)=2h^4(Z,V)-h^3(Z,V).
 \tag{C.MH.14}\label{cdp:eq:C.MH.14}
\]

\noindent\textbf{Step 6.} If $Z$ is rational, the pullback of $V$ is constant,
$h^4(Z,V)=d$, and $h^3(Z,V)=0$; this proves \eqref{cdp:eq:C.MH.6}.
If $Z$ is ruled over $B$, blowups leave its fundamental
group and its third cohomology unchanged, and the
pulled-back local system comes from $B$. The
projective-bundle formula gives
$H^3(Z,V)=H^1(B,V)(-1)$. If the cyclic character is
trivial on $B$, its dimensions are
$h^4(Z,V)=d$ and $h^3(Z,V)=2gd$. Otherwise every
nontrivial character of the prime-order group remains
nontrivial on $B$, so $h^4(Z,V)=0$ and
$h^3(Z,V)=(2g-2)d$. Both substitutions in \eqref{cdp:eq:C.MH.14}
give \eqref{cdp:eq:C.MH.7}.
\end{proof}






\begin{proposition}[The contracted irrational section contributes zero]\label{cdp:st:C.2.MH.3}
\StatementRecord{cdp:st:C.2.MH.3}{proved}{Historical events: Claude Fable 5.2 (Anthropic), effort high, 16 November 2026: carried as background without fault found. GPT-6 Astra (OpenAI), ChatGPT Pro mode, effort Pro, S6\_09\_PAPERS, 17 November 2026: targeted proof check. Also checked in this writing. Pre-continuation writing events: GPT-6 Astra (OpenAI), ChatGPT Pro mode, effort Pro, S6\_10\_PAPERS, 18 November 2026: targeted proof check. Continuation mathematical and source checks: GPT-6 Astra (OpenAI), ChatGPT Pro mode, effort Pro, S6\_10\_PAPERS continuation, 19 November 2026: independent re-proof. Continuation metadata events: GPT-6 Astra (OpenAI), ChatGPT Pro mode, effort Pro, S6\_10\_PAPERS continuation, 19 November 2026: used and rechecked in this run; status-vocabulary reconciliation; metadata reconciliation only; no new mathematical check. \emph{Pre-continuation recorded proof scope.} CR.5 gives genus one, MH.2 gives the actual-fibre zero Euler equation, and CR.1 gives n=v-N0. The independent unconditional calculation UN.7 has Euler at least two and trivial fundamental group for that contracted configuration, excluding it under the stated CDP hypotheses. This record asserts no occurrence example}{the hypotheses, notation and definitions stated here; \cref{cdp:st:C.2.CR.1}; \cref{cdp:st:C.2.CR.5}; \cref{cdp:st:C.2.MH.2}; \cref{cdp:st:C.2.UN.7}}

In the contracted-horizontal-section situation of
\ref{cdp:st:C.2.CR.1}, the normalization resolution is ruled over
an elliptic curve and
\[
                    e(S)=0,\qquad n=v-N_0.
 \tag{C.MH.15}\label{cdp:eq:C.MH.15}
\]
\end{proposition}


\begin{proof}
\noindent\textbf{Step 1.} \ref{cdp:st:C.2.CR.5} proves that the ruling genus is one.
\ref{cdp:st:C.2.MH.2} then gives $e(S)=0$ for the actual fibre.
The exact equality (C.CR.3) gives $n=v-N_0$.
The later unconditional calculation \ref{cdp:st:C.2.UN.7}
gives $e(S)\ge2$ and $\pi_1(S)=1$ for this
contracted irrational configuration. Thus its
comparison with the necessary equations here
excludes the case entirely under the CDP
hypotheses; this reduction supplies no example
of an actual CDP fibre.
\end{proof}






\begin{proposition}[Every irreducible class VII normalization contributes zero]\label{cdp:st:C.2.MH.4}
\StatementRecord{cdp:st:C.2.MH.4}{proved}{Historical events: Claude Fable 5.2 (Anthropic), effort high, 16 November 2026: carried as background without fault found. GPT-6 Astra (OpenAI), ChatGPT Pro mode, effort Pro, S6\_09\_PAPERS, 17 November 2026: targeted proof check. Also checked in this writing. Pre-continuation writing events: GPT-6 Astra (OpenAI), ChatGPT Pro mode, effort Pro, S6\_10\_PAPERS, 18 November 2026: targeted proof check. Continuation mathematical and source checks: GPT-6 Astra (OpenAI), ChatGPT Pro mode, effort Pro, S6\_10\_PAPERS continuation, 19 November 2026: computation or formal certificate; independent re-proof. Continuation metadata events: GPT-6 Astra (OpenAI), ChatGPT Pro mode, effort Pro, S6\_10\_PAPERS continuation, 19 November 2026: used and rechecked in this run; status-vocabulary reconciliation; metadata reconciliation only; no new mathematical check. \emph{Pre-continuation recorded proof scope.} The class-VII calculation was re-derived without mixed Hodge theory: h3(S,V)>=h3(Z,V), h4 is preserved, the trivial and nontrivial restrictions give e(S)<=1 or <=0, and the independently proved parity and nonnegativity force e(S)=0}{the hypotheses, notation and definitions stated here; \cref{cdp:st:C.2.VII.4}; \cref{cdp:st:C.2.MH.2}; \cref{cdp:eq:C.MH.17}; \cref{cdp:st:C.2.UN.1}; \cref{cdp:st:C.2.VII.3}}

Under the full CDP torus-fibration hypotheses, an
irreducible fibre support whose normalization resolution
is a surface of class VII has
\[
                              e(S)=0.
 \tag{C.MH.16}\label{cdp:eq:C.MH.16}
\]
Consequently every such actual fibre satisfies the
equality data of \ref{cdp:st:C.2.VII.4}.
\end{proposition}


\begin{proof}
\noindent\textbf{Step 1.} Use the prime-order rational augmentation local system
$V$ of \ref{cdp:st:C.2.MH.2}, of rank $d=\ell-1$. Its construction,
the torus acyclicity (C.MH.8), and the support duality
(C.MH.13) do not require a projective normalization.
They give $h^0(S,V)=h^1(S,V)=0$ and
$h^2(S,V)=h^4(S,V)$. Conductor descent still gives the
surjection $H^3(S,V)\to H^3(T,V)$ and the isomorphism
$H^4(S,V)\simeq H^4(T,V)$. The negative-definite
exceptional Gysin matrix in the proof of (C.MH.12)
is valid for the resolution of any normal complex
surface. It gives
\[
 \begin{gathered}
 h^3(S,V)\ge h^3(Z,V),\qquad h^4(S,V)=h^4(Z,V),\\
 d\,e(S)\le 2h^4(Z,V)-h^3(Z,V).
 \end{gathered}
 \tag{C.MH.17}\label{cdp:eq:C.MH.17}
\]
No mixed Hodge assertion is used in these comparisons.

\noindent\textbf{Step 2.} If the cyclic character restricts nontrivially to $Z$,
all its nontrivial powers remain nontrivial because
$\ell$ is prime. Poincar\'e duality on the compact
oriented surface gives $h^4(Z,V)=0$, and \eqref{cdp:eq:C.MH.17}
gives $e(S)\le0$. If its restriction is trivial,
$h^4(Z,V)=d$ and $h^3(Z,V)=d$: a class VII surface
has $b_1=b_3=1$, and point blowups preserve $b_3$.
Thus $e(S)\le1$ in this second case. \ref{cdp:st:C.2.UN.1} proves
that the Euler number of every actual torus-degeneration
fibre is even, so the second case also gives
$e(S)\le0$. The nonnegative lower bound of \ref{cdp:st:C.2.VII.3}
now gives \eqref{cdp:eq:C.MH.16} in both cases. \ref{cdp:st:C.2.VII.4} supplies
its exact normalization-and-conductor equality data.
\end{proof}






\section{Normal-crossing topology and canonical-flat component comparisons}\label{sec:II:08}


\begin{proposition}[Euler localization for torus fibrations]\label{cdp:st:C2a}
\StatementRecord{cdp:st:C2a}{proved}{Historical events: Claude Fable 5.2 (Anthropic), effort high, 16 November 2026: carried as background without fault found. GPT-6 Astra (OpenAI), ChatGPT Pro mode, effort Pro, S6\_09\_PAPERS, 17 November 2026: targeted proof check. Also checked in this writing. Pre-continuation writing events: GPT-6 Astra (OpenAI), ChatGPT Pro mode, effort Pro, S6\_10\_PAPERS, 18 November 2026: targeted proof check. Continuation mathematical and source checks: GPT-6 Astra (OpenAI), ChatGPT Pro mode, effort Pro, S6\_10\_PAPERS continuation, 19 November 2026: independent re-proof. Continuation metadata events: GPT-6 Astra (OpenAI), ChatGPT Pro mode, effort Pro, S6\_10\_PAPERS continuation, 19 November 2026: used and rechecked in this run; status-vocabulary reconciliation; metadata reconciliation only; no new mathematical check. \emph{Pre-continuation recorded proof scope.} The finite critical-value set, proper submersion trivializations over the complement, compact-support Euler additivity and the top Chern/Euler comparison give the stated sum over reduced supports. Scheme multiplicities leave the underlying topological support unchanged}{the hypotheses, notation and definitions stated here; the explicit constructions and calculations in the proof}

Let $g:Y\to C$ be a proper surjective holomorphic map from a connected compact complex manifold to a compact Riemann surface, with connected general fibre $F$. Let $\Sigma\subset C$ contain every critical value. Then $\Sigma$ can be chosen finite, and
\[
 e(Y)=e(C\setminus\Sigma)e(F)+\sum_{c\in\Sigma}e((Y_c)_{\rm red}). \tag{C.1}\label{cdp:eq:C.1}
\]
For a complex threefold with general fibre a complex two-torus this becomes
$\int_Yc_3(T_Y)=e(Y)=\sum_{c\in\Sigma}e((Y_c)_{\rm red})$.
\end{proposition}


\begin{proof}
\noindent\textbf{Step 1.} The critical locus is a closed analytic subset of $Y$. Properness and the proper mapping theorem make its image analytic in $C$; generic smoothness over characteristic zero makes this image proper. Thus it is finite. Over $C\setminus\Sigma$, the underlying differentiable map is a proper submersion. Ehresmann's theorem makes it a smooth locally trivial fibre bundle. Euler characteristic with compact supports is additive for the finite stratification by this open part and the fibres; triangulating the base and trivializing over its cells shows that the open part contributes $e_c(C\setminus\Sigma)e(F)$. The real dimensions of a complex curve and its punctured strata are even, so $e_c(C\setminus\Sigma)=e(C\setminus\Sigma)$. Nilpotents change no underlying topological space. This proves \eqref{cdp:eq:C.1}. A two-torus has the cohomology of the real four-torus and Euler number $(1-1)^4=0$. Finally the top Chern class of the complex tangent bundle is the Euler class for the complex orientation, giving the last identity.
\end{proof}






\begin{proposition}[Counterexample: Nonnegativity of every fibre in unrestricted torus fibrations]\label{cdp:st:C2b}
\StatementRecord{cdp:st:C2b}{proved}{Historical events: Claude Fable 5.2 (Anthropic), effort high, 16 November 2026: carried as background without fault found. GPT-6 Astra (OpenAI), ChatGPT Pro mode, effort Pro, S6\_09\_PAPERS, 17 November 2026: targeted proof check. Also checked in this writing. Pre-continuation writing events: GPT-6 Astra (OpenAI), ChatGPT Pro mode, effort Pro, S6\_10\_PAPERS, 18 November 2026: targeted proof check. Continuation mathematical and source checks: GPT-6 Astra (OpenAI), ChatGPT Pro mode, effort Pro, S6\_10\_PAPERS continuation, 19 November 2026: computation or formal certificate; targeted proof check. Continuation metadata events: GPT-6 Astra (OpenAI), ChatGPT Pro mode, effort Pro, S6\_10\_PAPERS continuation, 19 November 2026: used and rechecked in this run; status-vocabulary reconciliation; metadata reconciliation only; no new mathematical check. \emph{Pre-continuation recorded proof scope.} The genus-two Abel-Jacobi map is an embedding by the exact degree-one divisor and basepoint-free differential arguments. Blowing up its curve in the torus product gives central Euler number minus two; the subsequent point blowup gives zero. The local t=u*tau chart, projectivized normal bundle and both integral blowup maps give every printed fibre and ambient Betti number}{the hypotheses, notation and definitions stated here; the explicit constructions and calculations in the proof}

There is a smooth projective threefold $Y$ with a proper connected-fibre map $Y\to\PP^1$, smooth general fibre an abelian surface, and a reduced normal-crossing fibre of Euler number $-2$. There is also such a threefold with a singular reduced fibre of Euler number zero whose three irreducible components are a blown-up abelian surface, a ruled surface over a genus-two curve, and $\PP^2$.
\end{proposition}


\begin{proof}
\noindent\textbf{Step 1.} Let $C_2$ be the smooth projective completion of $w^2=z^6-1$. Its six simple branch points and the unramified points at infinity give genus two by Riemann--Hurwitz. Let $A=\operatorname{Jac}(C_2)$ and embed $C_2$ in $A$ by $p\mapsto\OO(p-p_*)$. This is injective because a principal divisor $p-q$ for distinct points would define a degree-one map $C_2\to\PP^1$. It is an immersion because the canonical system of a genus-two curve has no base point; the explicit differentials $dz/w,z\,dz/w$ verify that property, including at infinity. A proper injective immersion is an embedding. Adjunction on the torus gives $N_{C_2/A}=K_{C_2}$.

\noindent\textbf{Step 2.} Blow up $A\times\PP^1$ along $C_2\times\{0\}$, and denote the resulting smooth projective threefold by $Y$. Compose its blowdown with projection to $\PP^1$. The exceptional divisor is $E=\PP(K_{C_2}\oplus\OO_{C_2})$. The strict transform of $A\times\{0\}$ is $A$ itself, since blowing up a Cartier divisor in a smooth surface is an isomorphism. These two components meet along a section isomorphic to $C_2$. In local coordinates $(u,v,t)$ with centre $(u,t)=(0,0)$, a blowup chart has $t=u\tau$, so the fibre is the reduced normal-crossing divisor $u\tau=0$. The other chart and the complement give the claimed global map and connectedness of its fibres. Euler additivity gives
\[
 e((Y_0)_{\rm red})=e(A)+e(E)-e(C_2)=0+2(-2)-(-2)=-2.
\]
Every other fibre is $A$.

\noindent\textbf{Step 3.} Now blow up one point of the strict-transform component $A$ away from $C_2$. The new central fibre has components $\operatorname{Bl}_pA$, $E$, and $\PP^2$. Its pairwise intersections are $C_2$ and $\PP^1$, with no triple intersection. Its Euler number is
\[
1-4+3-(-2)-2=0.
\]
Both modifications have smooth total space and reduced normal-crossing central fibre. For completeness the integral cohomology is torsion-free. The Betti numbers of $A\times\PP^1$ are $(1,4,7,8,7,4,1)$; the blowup formula adds $H^{k-2}(C_2,\Z)$ in degree $k$, giving $(1,4,8,12,8,4,1)$. The point blowup adds one generator in degrees two and four, giving $(1,4,9,12,9,4,1)$. These formulas follow from excision for the exceptional projective bundle and its hyperplane class, whose restriction identifies the new relative summands. Their Euler numbers are respectively $-2$ and zero, independently checking the fibre computation. The exact arithmetic is in \texttt{code/C\_euler.py}.
\end{proof}






\begin{proposition}[Euler numbers in semistable toric rank-two models]\label{cdp:st:C2c}
\StatementRecord{cdp:st:C2c}{proved}{Historical events: Claude Fable 5.2 (Anthropic), effort high, 16 November 2026: carried as background without fault found. GPT-6 Astra (OpenAI), ChatGPT Pro mode, effort Pro, S6\_09\_PAPERS, 17 November 2026: targeted proof check. Also checked in this writing. Pre-continuation writing events: GPT-6 Astra (OpenAI), ChatGPT Pro mode, effort Pro, S6\_10\_PAPERS, 18 November 2026: targeted proof check; used and rechecked in this run. Continuation mathematical and source checks: GPT-6 Astra (OpenAI), ChatGPT Pro mode, effort Pro, S6\_10\_PAPERS continuation, 19 November 2026: computation or formal certificate; targeted proof check; used and rechecked in this run. Continuation metadata events: GPT-6 Astra (OpenAI), ChatGPT Pro mode, effort Pro, S6\_10\_PAPERS continuation, 19 November 2026: used and rechecked in this run; status-vocabulary reconciliation; metadata reconciliation only; no new mathematical check. \emph{Pre-continuation recorded proof scope.} The cone fan over the explicit unimodular triangles gives t=x1*x2*x3. The lattice shear acts by the displayed t-periods; logarithmic barycentric coordinates prove proper discontinuity and compact representatives, giving a proper analytic quotient with the stated product torus nearby. Its vertex and triangle counts give every positive even Euler number. Mumford section6 pp265-268 was opened with its complete-DVR hypotheses; the analytic quotient is established by the receiving proof's own maps}{the hypotheses, notation and definitions stated here; \cite[Section~6, especially (6.3)--(6.7)]{Mumford1972}}

Consider a rank-two Mumford model with smooth three-dimensional total space obtained from a unimodular periodic triangulation of $\R^2$, divided by a rank-two translation lattice acting freely on the total model. Let $V,E,F$ be the numbers of vertices, edges and triangles in the quotient triangulation. Its reduced central fibre has Euler number
\[
 e(S)=F=2V\in2\Z_{>0}.                                \tag{C.2}\label{cdp:eq:C.2}
\]
Every positive even number occurs in this class. A rank-one model assembled from $\PP^1$-bundles over elliptic curves with their boundary elliptic sections identified has Euler number zero.
\end{proposition}


\begin{proof}
\noindent\textbf{Step 1.} The toric orbit strata of a rank-two central fibre have dimensions two, one and zero. Its two-dimensional strata are copies of $(\C^*)^2$, its one-dimensional strata copies of $\C^*$, and its zero-dimensional strata correspond to triangles of the quotient triangulation. The first two have compactly supported Euler number zero. Hence the Euler number is $F$. The quotient complex is a real two-torus, so $V-E+F=0$, and every edge has two incident triangles, giving $3F=2E$, with incidences counted even when a face has identified edges. Thus $F=2V$.

\noindent\textbf{Step 2.} For existence, use the lattice $N=\mathbb Z^2\oplus\mathbb Z$
and the standard triangulation of $\mathbb R^2$ whose triangles are
\[
 [u,u+e_1,u+e_1+e_2],\qquad
 [u,u+e_2,u+e_1+e_2],\qquad u\in\mathbb Z^2.
\]
These are the Delaunay triangles for the positive integral quadratic
form $q(u_1,u_2)=u_1^2-u_1u_2+u_2^2$.
Take the cones over these triangles at height one and all their
faces. They form a fan. Every three-dimensional cone is generated
by a basis of $N$, because its triangle is unimodular. Gluing its
affine toric charts therefore gives a smooth complex threefold
$\mathcal X$ with a holomorphic map to the affine line. On each
maximal chart there are coordinates $x_1,x_2,x_3$ in which the map
is exactly $t=x_1x_2x_3$. The central support is consequently
reduced and has normal crossings. Over $t\ne0$ the fibre is the
common algebraic torus $(\mathbb C^*)^2$.

\noindent\textbf{Step 3.} Let $\Lambda_N=\mathbb Z(Ne_1)\oplus\mathbb Z e_2$.
For $\ell\in\Lambda_N$, the lattice shear
\[
 (u,k)\longmapsto(u+k\ell,k)
\]
preserves the fan and induces an automorphism over the base.
On the punctured torus chart it is
\[
 (z_1,z_2,t)\longmapsto
 (t^{\ell_1}z_1,t^{\ell_2}z_2,t).
\]
Restrict to $|t|<\epsilon<1$. This action is free: away from the
central fibre, a fixed point would imply $t^{\ell_i}=1$ for both
coordinates and hence $\ell=0$; on the central fibre, a fixed
point would require its bounded indexing face of the triangulation
to be invariant under translation by $\ell$, which again forces
$\ell=0$.

\noindent\textbf{Step 4.} Here is also the required analytic properness check. For $t\ne0$,
put $v_i=\log|z_i|/\log|t|$. In the chart of a triangle with
vertices $u_1,u_2,u_3$, write
\[
 v=\sum_{j=1}^3\lambda_j u_j,\qquad
 \lambda_j=\frac{\log|x_j|}{\log|t|},\qquad
 \sum_j\lambda_j=1.
\]
On a compact subset of finitely many charts the $|x_j|$ are
bounded. As $t$ tends to zero, this gives a uniform lower bound
for every $\lambda_j$, and their sum then gives a uniform upper
bound. Thus $v$ remains bounded there. Away from zero the same
conclusion follows from compactness in the torus charts.
Translation by $\ell$ sends $v$ to $v+\ell$. Consequently only
finitely many elements can carry a fixed compact set into itself,
or make two fixed compact sets intersect. On the central fibre
this is also immediate from the finitely many indexing faces.
The action is properly discontinuous, so its quotient is a smooth
complex threefold.

\noindent\textbf{Step 5.} The quotient map to the disc is proper. Indeed, translate any
punctured-fibre point so that $v$ belongs to a fixed compact
fundamental parallelogram for $\Lambda_N$. Only finitely many
triangles meet that parallelogram. Choose one containing $v$;
then all its barycentric coordinates are nonnegative, so
$|x_j|=|t|^{\lambda_j}\le1$. These finitely many closed unit
polydiscs provide compact representatives for every sequence
approaching the central fibre. The central fibre itself has
finitely many toric components modulo $\Lambda_N$, and each is
compact: the star of a vertex is a complete two-dimensional fan.
These facts prove compactness over every closed smaller disc.
For $t\ne0$ the quotient fibre is explicitly
\[
 \mathbb C^*/t^{N\mathbb Z}\ \times\
 \mathbb C^*/t^{\mathbb Z},
\]
a product of elliptic curves. This supplies the actual convergent
analytic degeneration, including its smooth total space and
properness, without inferring analytic convergence from a formal
model. The quadratic form supplies the usual polarization data:
its difference under every lattice translation is an integral
affine function. Thus the construction is the analytic realization
of the polyhedral model described in
\cite[Section~6, especially (6.3)--(6.7)]{Mumford1972}.
That source formulates its construction over a complete discrete
valuation ring; the preceding chart and compactness argument is
the complex-disc comparison used here.

\noindent\textbf{Step 6.} The quotient triangulation has $V=N$ and $F=2N$, as required.
For rank one every normalization component is a $\PP^1$-bundle over an elliptic curve, with Euler number zero, and the identified boundary sections also have Euler number zero. Additivity gives the result.
\end{proof}




\begin{proposition}[All rank-four monodromy stratum signatures]\label{cdp:st:C.2.UN.11}
\StatementRecord{cdp:st:C.2.UN.11}{proved}{Historical events: Claude Fable 5.2 (Anthropic), effort high, 16 November 2026: carried as background without fault found. Claude Fable 5.2 (Anthropic), effort high, 17 November 2026: carried as background without fault found. GPT-6 Astra (OpenAI), ChatGPT Pro mode, effort Pro, S6\_09\_PAPERS, 17 November 2026: targeted proof check. Also checked in this writing. Pre-continuation writing events: GPT-6 Astra (OpenAI), ChatGPT Pro mode, effort Pro, S6\_10\_PAPERS, 18 November 2026: targeted proof check; used and rechecked in this run. Continuation metadata events: GPT-6 Astra (OpenAI), ChatGPT Pro mode, effort Pro, S6\_10\_PAPERS continuation, 19 November 2026: used and rechecked in this run; status-vocabulary reconciliation; metadata reconciliation only; no new mathematical check. \emph{Pre-continuation recorded proof scope.} The actual nearby-cycle stalk action and two-step degree-one spectral-sequence filtration give quasi-unipotence and the square-zero unipotent bound. Alternating traces, exterior powers, Mobius inversion, minimal exceptional multiplicities and the full finite-order torus quotient constructions were re-derived. All nineteen signature rows and exterior-power zeta identities were independently replayed over the integers. The point/quartic charts preserve SNC and realize zero Euler with exactly two singular values. A'Campo's primary-host indexed text was opened at section2 pp237-243 and section3 pp243-245; direct visual-scan retrieval remains separately unavailable}{the hypotheses, notation and definitions stated here; \cite[\S2, pp.~237--243, and \S3]{ACampo1975}; \cite[Theorem~1, p.~235,
and proof on pp.~243--244]{ACampo1975}; \cref{cdp:eq:C.MS.6}; \cite[Theorem~1.6]{BierstoneMilman1997}; \cref{cdp:st:C.2.UN.1}; \cref{cdp:st:C.2.UN.4}}

Let $f:U\to\Delta$ be a proper holomorphic map from a
smooth complex threefold, with connected fibres and
smooth fibres complex two-tori away from zero. Take
an embedded resolution, unchanged away from the
central fibre, and write its full central divisor as
$\sum_jN_jE_j$, with simple normal crossings. Put
\[
 E_j^\circ=E_j\setminus\bigcup_{i\ne j}E_i,
 \qquad a_n=\sum_{N_j=n}e(E_j^\circ).
 \tag{C.MS.1}\label{cdp:eq:C.MS.1}
\]
Let $T$ be the actual monodromy on $H^1(F,\Z)$,
and let $M$ be the order of its semisimple part.
Then $T\in\operatorname{SL}_4(\Z)$ is
quasi-unipotent, and the following exact maps
determine the entire sequence $(a_n)$:
\[
 \begin{aligned}
 L_k&=\det(I-T^k)
       =\sum_{n\mid k}n a_n,\\
 a_n&=\frac1n\sum_{d\mid n}\mu(n/d)\det(I-T^d).
 \end{aligned}                                  \tag{C.MS.2}\label{cdp:eq:C.MS.2}
\]
Here $\mu$ is the integer M\"obius function.
The complete table follows; $\Phi_d$ denotes
the $d$th cyclotomic polynomial and every
unlisted $a_n$ is zero.

\begin{center}
\small
\begin{tabular}{ccl}
$P_T$ & $M$ & Nonzero $a_n$\\\hline
$\Phi_{1}^{4}$ & 1 & $\text{all zero}$ \\
$\Phi_{1}^{2}\Phi_{2}^{2}$ & 2 & $\text{all zero}$ \\
$\Phi_{1}^{2}\Phi_{3}$ & 3 & $\text{all zero}$ \\
$\Phi_{1}^{2}\Phi_{4}$ & 4 & $\text{all zero}$ \\
$\Phi_{1}^{2}\Phi_{6}$ & 6 & $\text{all zero}$ \\
$\Phi_{2}^{4}$ & 2 & $a_{1}=16,\quad a_{2}=-8$ \\
$\Phi_{3}^{2}$ & 3 & $a_{1}=9,\quad a_{3}=-3$ \\
$\Phi_{2}^{2}\Phi_{4}$ & 4 & $a_{1}=8,\quad a_{2}=-4$ \\
$\Phi_{4}^{2}$ & 4 & $a_{1}=4,\quad a_{2}=6,\quad a_{4}=-4$ \\
$\Phi_{5}$ & 5 & $a_{1}=5,\quad a_{5}=-1$ \\
$\Phi_{2}^{2}\Phi_{3}$ & 6 & $a_{1}=12,\quad a_{2}=-6,\quad a_{3}=-4,\quad a_{6}=2$ \\
$\Phi_{2}^{2}\Phi_{6}$ & 6 & $a_{1}=4,\quad a_{2}=-2,\quad a_{3}=4,\quad a_{6}=-2$ \\
$\Phi_{3}\Phi_{6}$ & 6 & $a_{1}=3,\quad a_{2}=3,\quad a_{3}=-1,\quad a_{6}=-1$ \\
$\Phi_{6}^{2}$ & 6 & $a_{1}=1,\quad a_{2}=4,\quad a_{3}=5,\quad a_{6}=-4$ \\
$\Phi_{8}$ & 8 & $a_{1}=2,\quad a_{2}=1,\quad a_{4}=3,\quad a_{8}=-2$ \\
$\Phi_{10}$ & 10 & $a_{1}=1,\quad a_{2}=2,\quad a_{5}=3,\quad a_{10}=-2$ \\
$\Phi_{3}\Phi_{4}$ & 12 & $a_{1}=6,\quad a_{2}=3,\quad a_{3}=-2,\quad a_{4}=-3,\quad a_{6}=-1,\quad a_{12}=1$ \\
$\Phi_{4}\Phi_{6}$ & 12 & $a_{1}=2,\quad a_{2}=5,\quad a_{3}=2,\quad a_{4}=-3,\quad a_{6}=-3,\quad a_{12}=1$ \\
$\Phi_{12}$ & 12 & $a_{1}=1,\quad a_{3}=1,\quad a_{4}=2,\quad a_{6}=2,\quad a_{12}=-2$ \\
\end{tabular}
\end{center}


In particular $a_n=0$ unless $n\mid M$, where
\[
                 M\in\{1,2,3,4,5,6,8,10,12\}.
 \tag{C.MS.3}\label{cdp:eq:C.MS.3}
\]
On $H^1(F,\Q)$ one also has
\[
                  (T^M-I)^2=0,
                  \qquad\operatorname{rank}(T^M-I)\le2.
 \tag{C.MS.10}\label{cdp:eq:C.MS.10}
\]
All $a_n$ vanish exactly when $1$ is an eigenvalue
of $T$. The table is the complete range of \eqref{cdp:eq:C.MS.2}
on orientation-preserving quasi-unipotent integral
rank-four matrices. Every actual filling has one
of these signatures. Every row is realized by a
map from a smooth projective threefold to $\PP^1$
whose general fibre is an abelian surface and
whose every fibre has Euler number zero. It
can be chosen to have exactly two singular
fibres, both with reducible simple-normal-crossings
support.

There are two useful consequences for the original
fibre before resolution. Write it as $\sum_i m_iS_i$,
put $m_0=\min_i m_i$, and let $S_i^\circ$ be the
smooth locus of $S_i$ with all its intersections
with the other components removed. Then
\[
 m_0\sum_{m_i=m_0}e(S_i^\circ)
                    =\det(I-T^{m_0})\ge0.        \tag{C.MS.4}\label{cdp:eq:C.MS.4}
\]
If every component has multiplicity greater than
one, the entire signature is zero. For an
irreducible full fibre $mS$ one has
\[
 \begin{gathered}
 m>1\quad\Longrightarrow\quad e(S_{\rm reg})=0,\\
 e(S_{\rm reg})\in
       \{0,1,2,3,4,5,6,8,9,12,16\},\\
 e(S_{\rm reg})>0\quad\Longrightarrow\quad m=1.
 \end{gathered}                                  \tag{C.MS.5}\label{cdp:eq:C.MS.5}
\]
These constraints apply in particular to every
reducible or irreducible zero-Euler filling.
\end{proposition}


\begin{proof}
\noindent\textbf{Step 1.} We specify first the comparison with the analytic
family. The resolution is an isomorphism over
$\Delta^*$, so it identifies the actual nearby
torus and its monodromy before and after resolution.
The proper global construction in
\cite[\S2, pp.~237--243, and \S3]{ACampo1975}
gives the nearby-cycle cohomology sheaves on the
resolved central support and their spectral sequence
to $H^*(F,\Q)$. In a chart
$z_1^{N_1}\cdots z_r^{N_r}=t$, a local nearby
fibre has $\gcd(N_1,\ldots,N_r)$ components,
each with the homotopy type of $(S^1)^{r-1}$.
Monodromy permutes these components cyclically;
after their permutation is killed its action
is a translation on the torus and induces the
identity on cohomology. A common multiple of all
the $N_j$ therefore acts identically on every
stalk of every nearby-cycle cohomology sheaf.
It acts identically on the corresponding
cohomology groups of those sheaves and hence
on the associated graded of the finite
spectral-sequence filtration of $H^*(F,\Q)$.
Its difference from the identity is nilpotent
on that cohomology. This proves quasi-unipotence
using the actual analytic normal-crossings
comparison; no ambient K\"ahler assumption enters.
In degree one this filtration has only the
two possible terms $H^1(E,\mathcal H^0\psi)$
and $H^0(E,\mathcal H^1\psi)$, where
$E=\bigcup_jE_j$ and $\mathcal H^q\psi$
denotes the nearby-cycle cohomology sheaf
already described. Thus for a
common multiplicity multiple $N$ one has
$(T^N-I)^2=0$. The semisimple order $M$
divides $N$, and $U=T^M$ is unipotent.
Writing
$U^{N/M}-I=(U-I)(I+U+\cdots+U^{N/M-1})$,
the second factor is invertible over $\Q$
and commutes with $U$. Hence $(U-I)^2=0$.
Its image is contained in its kernel on
the four-dimensional space, proving
the rank bound in \eqref{cdp:eq:C.MS.10}.

\noindent\textbf{Step 2.} The same local model computes the alternating
stalk trace of the $k$th monodromy power. At a
point of $E_j^\circ$ it is $N_j$ if $N_j\mid k$
and zero otherwise. At a point lying on two or
more components it is zero, since the torus
factor has positive dimension and Euler number
zero. Taking the alternating trace through the
constructible spectral sequence, or its finite
cell filtration, gives
\[
 \sum_q(-1)^q\operatorname{Tr}
             (T^k;H^q(F,\Q))
       =\sum_{N_j\mid k}N_j e(E_j^\circ).
 \tag{C.MS.6}\label{cdp:eq:C.MS.6}
\]
This is the global form of the explicit local
calculation in \cite[Theorem~1, p.~235,
and proof on pp.~243--244]{ACampo1975}.
The cup-product map
$\Lambda^qH^1(F,\Z)\to H^q(F,\Z)$ is an
isomorphism for a torus and commutes with the
actual monodromy. Thus the left side of
\eqref{cdp:eq:C.MS.6} is exactly $\det(I-T^k)$.
The complex orientation of the smooth fibre
is preserved by transport. The map on its
top cohomology is consequently the identity;
under the fourth exterior power this says
$\det T=1$. Integer M\"obius inversion now
gives both equations of \eqref{cdp:eq:C.MS.2}.

\noindent\textbf{Step 3.} We compute their complete range. The characteristic
polynomial of a quasi-unipotent integral rank-four
matrix is a product of cyclotomic polynomials
whose degrees sum to four. The possible individual
orders are
$1,2,3,4,5,6,8,10,12$. Since $\det T=1$,
the multiplicity of the eigenvalue $-1$ is even.
There are exactly nineteen resulting products:
the five containing $\Phi_1$ and the fourteen
remaining products displayed in the table.
This enumeration follows directly from
$\deg\Phi_1=\deg\Phi_2=1$,
$\deg\Phi_3=\deg\Phi_4=\deg\Phi_6=2$,
and degree four for $\Phi_5,\Phi_8,\Phi_{10},
\Phi_{12}$. The sum of the degree-one
multiplicities is even, so the $\Phi_1$
multiplicity is even as well.

\noindent\textbf{Step 4.} The determinants $\det(I-T^k)$ depend only on
these eigenvalues, including their multiplicities;
Jordan blocks make no change to them. Each listed
product is represented by the block diagonal
matrix of its integral cyclotomic companion
matrices, which has determinant one. Its traces
therefore compute the exact full range on
matrices. Applying \eqref{cdp:eq:C.MS.2} gives the displayed
integers.

\noindent\textbf{Step 5.} There is also a direct verification that no
unlisted index can be concealed by cancellation.
The actual exterior-power identity gives
\[
 \begin{aligned}
 \exp\!\left(\sum_{k\ge1}\frac{L_k}{k}t^k\right)
 &=\prod_{q=0}^4
       \det\bigl(I-t\Lambda^qT\bigr)^{(-1)^{q+1}}\\
 &=\prod_{n\ge1}(1-t^n)^{-a_n}.
 \end{aligned}                                  \tag{C.MS.7}\label{cdp:eq:C.MS.7}
\]
Every eigenvalue of every exterior power has
order dividing $M$. Hence the first product
factors into cyclotomic polynomials $\Phi_d(t)$
with $d\mid M$. The identity
$\Phi_d(t)=\prod_{n\mid d}(1-t^n)^{\mu(d/n)}$
for $d>1$, with the harmless sign at $d=1$,
expresses it using only indices dividing $M$.
Finite products of the factors $1-t^n$ have
unique integer exponent vectors: the coefficient
at the smallest differing degree in their
formal logarithms detects any difference.
Thus $a_n=0$ for every $n\nmid M$, including
all possible cancellations. If $T$ has an
eigenvalue one, every $L_k$ is zero and so is
the signature. Conversely, without that
eigenvalue, the real eigenvalues are $-1$
and the others occur in complex-conjugate
root-of-unity pairs. Their factors in
$\det(I-T)$ are respectively $2$ and
$|1-\lambda|^2$, all positive. Thus $a_1=L_1>0$.
The same factorization proves $L_k\ge0$ for
every $k$.

\noindent\textbf{Step 6.} For \eqref{cdp:eq:C.MS.4}, choose the resolution with centres
over the singular locus of the original reduced
support. Its first exceptional multiplicities
are at least $2m_0$: at such a centre either
one support component has order at least two
or at least two components occur. All later
exceptional multiplicities are sums of the
orders of the components containing their
centres. If the generic point of a centre is
contained in an existing exceptional component,
that component already contributes at least
$2m_0$. Otherwise the map is locally an
isomorphism at its generic point and the
original reduced support still has order at
least two there. Thus this case has the same
lower bound, including successive resolutions
over disjoint original singular points.
The open stratum of a strict
transform is exactly the specified
$S_i^\circ$. In \eqref{cdp:eq:C.MS.6} for $k=m_0$, only
strict transforms of multiplicity $m_0$
contribute; this proves \eqref{cdp:eq:C.MS.4}.

\noindent\textbf{Step 7.} If $m_0>1$, the resolved divisor has no
multiplicity-one component. Hence $a_1=0$,
so $T$ has eigenvalue one and the entire
signature vanishes. For an irreducible
fibre $mS$, \eqref{cdp:eq:C.MS.4} is
$m e(S_{\rm reg})=L_m$. When $m>1$ it
is zero by the preceding argument. When
$m=1$, its possible values are the
$a_1$ entries of the table, exactly the
set in \eqref{cdp:eq:C.MS.5}.

\noindent\textbf{Step 8.} We now realize every row by an actual compact
family. The fifteen products of two degree-two
blocks are obtained from products of elliptic
curves. The five blocks
\[
                 \Phi_1^2,\quad\Phi_2^2,\quad
                 \Phi_3,\quad\Phi_4,\quad\Phi_6
\]
come respectively from the identity, negation,
and multiplication by a primitive third, fourth
or sixth root of unity on an elliptic curve
with the corresponding triangular or square
lattice. Taking the product automorphism gives
exactly the indicated product on integral
first cohomology.

\noindent\textbf{Step 9.} For $\Phi_5$, let $A$ be the Jacobian of
the smooth genus-two curve
\[
                           C:\ y^2=x^5-1.
\]
Its double cover of $\PP^1$ has six simple
branch points, including infinity, so the
genus assertion follows from Riemann--Hurwitz.
The automorphism $(x,y)\mapsto(\zeta_5x,y)$
fixes infinity and acts on the basis
$dx/y,x\,dx/y$ of holomorphic differentials
with eigenvalues $\zeta_5,\zeta_5^2$.
The Abel--Jacobi map based at infinity
intertwines this automorphism with its induced
automorphism of $A$ and induces the actual
isomorphism on integral $H^1$. Including the
conjugate eigenvalues gives $\Phi_5$.
Composing the Jacobian automorphism with
negation gives order ten and characteristic
polynomial $\Phi_{10}$. Here negation is also
the map induced by the curve's hyperelliptic
involution: the divisor of $x-x(p)$ gives
$[p-\infty]+[\iota(p)-\infty]=0$ on the
Jacobian, verifying that comparison explicitly.

\noindent\textbf{Step 10.} For the remaining two rows take
$A=E_i\times E_i$, where
$E_i=\C/(\Z+i\Z)$. The Gaussian-integer
matrices
\[
 G_8=\begin{pmatrix}0&i\\1&0\end{pmatrix},
 \qquad
 G_{12}=\begin{pmatrix}0&1\\1&i\end{pmatrix}
 \tag{C.MS.8}\label{cdp:eq:C.MS.8}
\]
define actual automorphisms of this torus.
One has $G_8^2=iI$, and its holomorphic
eigenvalues are $\zeta_8,\zeta_8^5$.
The characteristic polynomial of $G_{12}$
over $\C$ is $x^2-ix-1$, with roots
$(i+\sqrt3)/2=\zeta_{12}$ and
$(i-\sqrt3)/2=\zeta_{12}^5$.
Together with the conjugate eigenvalues,
these give the integral characteristic
polynomials $\Phi_8$ and $\Phi_{12}$ and
the exact orders eight and twelve.
Both matrices preserve the displayed lattice;
their determinants are Gaussian units.
This accounts for all nineteen rows with
finite-order automorphisms $g$ of projective
two-tori $A$.

\noindent\textbf{Step 11.} Let $M$ be the order of the selected $g$.
Take the projective finite quotient
\[
 X=(A\times\PP^1)/
       \langle(a,z)\mapsto(g(a),\zeta_Mz)\rangle,
 \qquad f([a,z])=z^M.                           \tag{C.MS.9}\label{cdp:eq:C.MS.9}
\]
The action is free over $z\ne0,\infty$.
Over a nonzero finite base value the $M$
slices are permuted transitively, and their
quotient is the original connected torus $A$.
The actual positive meridian lifts from $z$
to $\zeta_Mz$, so under this identification
the monodromy is $g^{-1}$. Inverting all
the eigenvalues preserves each cyclotomic
polynomial and hence the displayed signature.
Resolve the projective quotient and then its
two central supports by projective blowups
over $0$ and $\infty$, using
\cite[Theorem~1.6]{BierstoneMilman1997}.
The resulting threefold is smooth and projective,
with SNC supports, and the punctured family is
unchanged. Its Stein factor has generic
degree one over the normal curve $\PP^1$,
so every fibre is connected. When $M=1$
the construction is simply the product,
which already has zero-Euler fibres.

\noindent\textbf{Step 12.} For the remaining cases, each of the two
possibly singular supports has even Euler
number by \ref{cdp:st:C.2.UN.1}. Apply the exact vertical
blowups of \ref{cdp:st:C.2.UN.4} to make each number
zero. A point blowup changes it by $+2$;
a point blowup followed by the blowup of
a smooth plane quartic in its exceptional
$\PP^2$ changes it by $-2$. These
operations are projective and are performed
only over the selected value, so they
preserve the actual punctured family and
the signature \eqref{cdp:eq:C.MS.2}.

\noindent\textbf{Step 13.} They preserve SNC supports as well. Choose
the point on a single component away from
the others. Its exceptional plane meets
the old strict transform along a line,
and choose the quartic transverse to that
line. Near any of their four intersection
points write the map as $t=(ru)^m$,
with exceptional component $r=0$, old
component $u=0$, and quartic centre
$(r,v)=(0,0)$. Its blowup charts
$v=rw$ and $r=vw$ give respectively
$t=(ru)^m$ and $t=(vwu)^m$.
Thus every resulting component is still
a coordinate hyperplane; away from these
four points the same assertion uses only
two coordinates. These explicit charts
verify the claimed SNC zero-Euler
realization of every row.
After the Euler number is zero, one additional
$+2$ operation and one $-2$ operation over the
same value leave it zero and introduce
exceptional components. Performing this at
both $0$ and $\infty$ ensures precisely two
reducible singular fibres, including the
identity-monodromy row. All other fibres
remain the original smooth tori.

\noindent\textbf{Step 14.} The exact computation is saved in
\path{code/D_monodromy_signatures.py}, its
output, and
\path{records/D_monodromy_signatures.csv}.
It verifies every determinant and exterior-power
zeta identity over the integers or $\Q(t)$,
and the M\"obius inversions use exact integers.
The matrix-to-signature map itself forgets
information: for example $I_4$ and
$I_4+E_{12}$ have the same characteristic
polynomial and zero signature, while their
minimal polynomials differ. The table therefore
records precisely the stated invariant and
does not supply an analytic isomorphism
classification of fillings.
\end{proof}






\begin{proposition}[K3 components are elliptic in positive algebraic dimension; algebraic dimension zero has an exact occurrence criterion]\label{cdp:st:C.2.UN.24}
\StatementRecord{cdp:st:C.2.UN.24}{proved}{Historical events: Claude Fable 5.2 (Anthropic), effort high, 16 November 2026: carried as background without fault found. GPT-6 Astra (OpenAI), ChatGPT Pro mode, effort Pro, S6\_09\_PAPERS, 17 November 2026: targeted proof check. Also checked in this writing. Pre-continuation writing events: GPT-6 Astra (OpenAI), ChatGPT Pro mode, effort Pro, S6\_10\_PAPERS, 18 November 2026: targeted proof check; used and rechecked in this run. Continuation metadata events: GPT-6 Astra (OpenAI), ChatGPT Pro mode, effort Pro, S6\_10\_PAPERS continuation, 19 November 2026: used and rechecked in this run; status-vocabulary reconciliation; metadata reconciliation only; no new mathematical check. \emph{Pre-continuation recorded proof scope.} The exact covering quotient and fibration descent, cyclic symplectic eigenvalue restriction, fixed-point orbit counts and resolution Euler numbers were re-derived and the counts independently replayed. The revision proves Kodaira dimension at most one by the elliptic-fibre canonical restriction before classification. The positive rational four-square step was matched to Lagrange's original text in a modern transcription, pp5-6. Piatetski-Shapiro-Shafarevich section3 Corollary3 p560 and section6 Theorem1(3) pp573-574 supply the exact isotropic-class reflection and elliptic-pencil construction. The quotient family and actual local Euler adjustments prove occurrence}{the hypotheses, notation and definitions stated here; \cref{cdp:st:C.2.UN.16}; \cref{cdp:eq:C.K3.2}; \cref{cdp:st:C.2.UN.20}; \cref{cdp:st:C.2.UN.22}; \cref{cdp:eq:C.K3.3}; \cref{cdp:st:C.3.1}; \cite[Chapter~1, Theorem~3.1, p.~10]{Peters1990}; \cref{cdp:st:C.2.UN.23}; \cref{cdp:eq:C.K3.7}; \cite[Theorem~II, Corollary~II]{Lagrange1770}; \cite[\S3, Corollary~3, with the reflection
argument of \S6]{PiatetskiShapiroShafarevich1971}; \cref{cdp:eq:C.K3.4}; \cref{cdp:eq:C.K3.5}; \cref{cdp:eq:C.K3.6}; \cref{cdp:st:C.2.UN.1}; \cref{cdp:st:C.2.UN.4}}

Let $f:U\to\Delta$ be a proper holomorphic map from a
smooth complex threefold, with connected two-torus fibres
over $\Delta^*$. Let $S_i$ be an irreducible component of
its central support, and suppose that a resolution of
$S_i^\nu$ has a K3 minimal model $M$.
If $a(M)>0$, there is an actual elliptic fibration
$M\to\mathbb P^1$.

Among K3 surfaces of algebraic dimension zero, the minimal
birational classes which occur in this way are exactly the
minimal resolutions of cyclic quotients of complex two-tori
by nontrivial symplectic automorphisms without eigenvalue
one on the tangent space. After translating the origin, the generator
is linear of order $n=2,3,4$, or $6$. The quotient singularities
and exceptional configurations are exactly
\[
\begin{array}{c|c|c|c}
n&\text{quotient singularities}&
 \text{number of exceptional curves}&e(A/\langle g\rangle)\\ \hline
2&16A_1&16&8\\
3&9A_2&18&6\\
4&4A_3+6A_1&18&6\\
6&A_5+4A_2+5A_1&18&6.
\end{array}                                                    \tag{C.K3.1}\label{cdp:eq:C.K3.1}
\]
Here the cyclic quotient in the criterion is required to
have algebraic dimension zero; the table applies to the
cyclic K3 quotient models in every algebraic dimension.

Every cyclic quotient K3 minimal class just described is
realized by a compact smooth threefold over $\mathbb P^1$
with connected torus general fibre and with both special
reduced supports of Euler characteristic zero. The central
support has a component whose normalization resolution has
exactly that K3 minimal model. The total space and construction
are projective when the given torus is projective.
\end{proposition}



\begin{proof}
\noindent\textbf{Step 1.} Use the actual normalized root base change and equivariant
resolution of \ref{cdp:st:C.2.UN.16}. For the chosen component, let $E$
be a smooth strict covering component and $G$ its cyclic
stabilizer. The comparison is
\[
 E\longrightarrow S_i^\nu,
 \qquad E/G\dashrightarrow S_i^\nu\dashrightarrow M,
                                                               \tag{C.K3.2}\label{cdp:eq:C.K3.2}
\]
where the first map is generically finite, the quotient
comparison is proper birational, and the remaining arrows
are resolved by their proper graphs. In particular, a
nonzero holomorphic two-form on $M$ pulls back, by its
actual Jacobian map, to a nonzero holomorphic two-form on
a resolution of the graph in \eqref{cdp:eq:C.K3.2}. Point-blowup
invariance then gives
\[
                  p_g(E)\ge1,\qquad \kappa(E)\ge0.
                                                               \tag{C.K3.3}\label{cdp:eq:C.K3.3}
\]
Also $a(E)=a(M)$. To check this equality through the
specified map, a meromorphic function on $E$ has elementary
symmetric functions along the finite generic sheets.
They descend meromorphically to $M$, so that the original
function satisfies a monic polynomial over $\mathcal M(M)$.
Thus the two meromorphic function fields have the same
transcendence degree.

\noindent\textbf{Step 2.} The two-dimensional space of fields on $E$ in \ref{cdp:st:C.2.UN.16}
has generic rank one or two. In rank one, \ref{cdp:st:C.2.UN.20} gives
an actual rational or elliptic fibration on a birational
model of $M$. The rational alternative would make all
positive pluricanonical sections vanish, by restricting
them to its smooth general fibre. The pulled-back K3
form excludes it. Hence the quotient fibration is elliptic.

\noindent\textbf{Step 3.} In rank two, \ref{cdp:st:C.2.UN.22} gives $\chi(\mathcal O_E)=0$ and
its precise torus or elliptic-fibration alternatives.
First suppose $a(M)=a(E)=1$. In the elliptic alternative,
the connected-fibre map $q:E'\to B$ has
\[
                     \mathcal M(E')=q^*\mathcal M(B).
                                                               \tag{C.K3.4}\label{cdp:eq:C.K3.4}
\]
Indeed every meromorphic function is algebraic over any
one nonconstant function on $B$ and hence constant on
the connected general fibres; the proper graph descends
it, as in (C.RF.9). The equality makes this field
intrinsic and therefore $G$-invariant. Resolve the graph
of $q$ equivariantly and take its Stein factor. The group
acts on that factor, and the actual finite quotient is a
map to $B/G$ whose general fibre is $F/K$, with
$K=\ker(G\to\operatorname{Aut}B)$. This is the same
quotient and Stein comparison proved in \ref{cdp:st:C.2.UN.20}.
Its general fibre is connected and rational or elliptic;
the K3 form again excludes the rational case. In the
torus alternative the minimal model of $E$ is a torus
and the cyclic quotient argument below applies.

\noindent\textbf{Step 4.} Now suppose $a(M)=a(E)=2$. We work on the smooth
projective minimal model $E_0$, to which the finite
action descends. In rank two, \eqref{cdp:eq:C.K3.3} and
$\chi(\mathcal O_{E_0})=0$, together with the compact
surface classification already used in \ref{cdp:st:C.3.1}
\cite[Chapter~1, Theorem~3.1, p.~10]{Peters1990}, leave the following
possibilities. The precise models in \ref{cdp:st:C.2.UN.22}
already give $\kappa(E_0)\le1$: on an elliptic general fibre
$F$ of its stated model, adjunction gives
$K|_F=K_F=\OO_F$, so every pluricanonical map is constant
on $F$ and has image of dimension at most one; a torus
has trivial canonical line. If $\kappa(E_0)=0$, it is an Abelian
surface: K3 and Enriques surfaces have coherent Euler
characteristic $2$ and $1$, and bielliptic surfaces have
$p_g=0$. If $\kappa(E_0)=1$, its elliptic Iitaka
fibration is intrinsic. More explicitly, $G$ acts on
each space of pluricanonical sections, so a sufficiently
large pluricanonical map, its equivariant graph resolution,
and its Stein factor are all $G$-equivariant. The
connected general fibres are elliptic. Taking the
quotient gives the actual fibration with general fibre
$F/K$ just described. It is elliptic because its total
quotient has the K3 minimal model. Thus the only
remaining positive-dimensional case is a cyclic quotient
of a two-torus.

\noindent\textbf{Step 5.} For use in all these branches, an elliptic fibration on
a smooth birational model of $M$ descends to $M$ itself.
Take an actual common resolution
\[
                  M\xleftarrow{\rho}X\xrightarrow{q}B
                                                               \tag{C.K3.5}\label{cdp:eq:C.K3.5}
\]
with connected general elliptic fibre $F$. Since $K_M=0$,
the point-blowup formula gives
$K_X=\sum_j a_jE_j$ with $a_j>0$ on every
$\rho$-exceptional prime curve. Adjunction gives
$K_X\cdot F=0$. Each $E_j\cdot F\ge0$, so all these
intersections are zero and every exceptional curve is
vertical for $q$. Every connected fibre of $\rho$ is
therefore sent to a point. The proper graph, together
with $\rho_*\mathcal O_X=\mathcal O_M$, descends $q$
to a holomorphic map $M\to B$ with connected fibres.
The Leray injection
$H^1(B,\mathcal O_B)\hookrightarrow H^1(M,\mathcal O_M)=0$
gives $B\simeq\mathbb P^1$.

\noindent\textbf{Step 6.} We next compute the cyclic quotient branch. By
\ref{cdp:st:C.2.UN.23}, a torus minimal model $A$ of $E$ has an
actual holomorphic cyclic action whose quotient is
birational to $M$. The K3 alternative there says that
its generator $g(z)=Lz+b$ has no eigenvalue one and
$\det L=1$. Since $I-L:A\to A$ is an isogeny, the
equation $(I-L)p=b$ has a solution. Translation by
$p$ conjugates $g$ to the linear automorphism $L$.
Write its eigenvalues as $\lambda,\lambda^{-1}$,
with $\lambda$ a primitive $n$th root. The integral
action on $H_1(A,\mathbb Z)$ has characteristic polynomial
\[
 (x^2-sx+1)^2
   =x^4-2sx^3+(s^2+2)x^2-2sx+1,
 \qquad s=\lambda+\lambda^{-1}\in[-2,2].
                                                               \tag{C.K3.6}\label{cdp:eq:C.K3.6}
\]
Its integral coefficients imply $2s\in\mathbb Z$
and $s^2\in\mathbb Z$. Writing $s=k/2$ shows that
$k$ is even, hence $s\in\mathbb Z$. The value $s=2$
is excluded by the absence of eigenvalue one.
Thus $s=-2,-1,0,1$, giving precisely
$n=2,3,4,6$.

\noindent\textbf{Step 7.} For every nonidentity power, its fixed points are the
kernel of an isogeny and their exact number is
\[
 \#\operatorname{Fix}(g^k)
 =\det_{\mathbb R}(I-L^k)
 =|\det_{\mathbb C}(I-L^k)|^2
 =\bigl(2-2\cos(2\pi k/n)\bigr)^2.
                                                               \tag{C.K3.7}\label{cdp:eq:C.K3.7}
\]
At a point with stabilizer of order $d$, the local
action is $(u,v)\mapsto(\zeta_du,\zeta_d^{-1}v)$.
Its invariant coordinates $x=u^d,y=v^d,z=uv$ give
exactly $xy=z^d$. The height-one toric resolution
in (C.CQ.11) gives its chain of $d-1$ rational
$(-2)$-curves. Formula \eqref{cdp:eq:C.K3.7} now gives the table
by subtracting points with larger stabilizers and
dividing by the remaining orbit sizes. For $n=4$
there are $4$ full-stabilizer points and
$(16-4)/2=6$ order-two orbits. For $n=6$ the
counts are $1$, $(9-1)/2=4$, and
$(16-1)/3=5$ for stabilizer orders $6,3,2$.
The order-two and order-three cases have $16$ and
$9$ fixed points respectively. The finite quotient
Euler formula gives the last column; adding the
exceptional ranks gives $24$ in every row.
The computation is reproduced with exact integral
matrices in \texttt{code/D\_cyclic\_k3.py}.

\noindent\textbf{Step 8.} If this quotient K3 $M$ is projective, the table gives
an explicit elliptic pencil criterion. Choose one
exceptional $(-2)$-curve from each of four different
singular points, and call them $C_1,\ldots,C_4$.
They are pairwise disjoint. For an ample integral
class $H$ on $M$, put
\[
 H'=H+\sum_{i=1}^4\frac{H\cdot C_i}{2}C_i,
 \quad H'\cdot C_i=0,
 \quad H'^2=H^2+\frac12\sum_{i=1}^4(H\cdot C_i)^2>0.
                                                               \tag{C.K3.8}\label{cdp:eq:C.K3.8}
\]
Lagrange's four-square theorem
\cite[Theorem~II, Corollary~II]{Lagrange1770}
expresses the positive rational number $H'^2/2$ as
$\sum_i x_i^2$ with $x_i\in\mathbb Q$: for a
positive fraction $a/b$, apply the integral theorem
to $ab$ and divide the four summands by $b$.
Consequently
\[
 D=H'-\sum_{i=1}^4x_iC_i,
 \qquad D^2=0,
 \qquad D\cdot H'=H'^2>0.
                                                               \tag{C.K3.9}\label{cdp:eq:C.K3.9}
\]
Clear denominators to obtain a nonzero isotropic
integral divisor class. The original elliptic-pencil
criterion \cite[\S3, Corollary~3, with the reflection
argument of \S6]{PiatetskiShapiroShafarevich1971}
then produces an actual base-point-free elliptic
pencil on $M$. This proves the remaining projective
branch. If the quotient has algebraic dimension
one, then $a(A)=1$ as well. Choose a nonconstant
meromorphic function on $A$, and let $D$ be a
general member of its pencil after removing fixed
components. Its self-intersection is nonnegative,
by intersecting two distinct members. In fact
$D^2=0$. To verify this, average the positive
integration current $[D]$ under translations of
the compact torus. The result is the constant
semipositive $(1,1)$-form representing $c_1(D)$.
If $D^2>0$, this form is positive definite, so
the integral class is a polarization and the
torus is projective, contrary to $a(A)=1$.
The equality $D^2=0$ shows that two distinct
members have no common point. The pencil thus
defines a holomorphic map, and its Stein factor
$q:A\to B$ has connected smooth general fibre
with trivial normal line. Adjunction with
$K_A=0$ gives genus one. The field identity
\eqref{cdp:eq:C.K3.4} makes this fibration $G$-equivariant;
its quotient and \eqref{cdp:eq:C.K3.5} give the actual
elliptic pencil on $M$. This completes the
assertion for $a(M)>0$.

\noindent\textbf{Step 9.} If $a(M)=0$, equality of algebraic dimensions in
\eqref{cdp:eq:C.K3.2} gives $a(E)=0$. The final assertion of
\ref{cdp:st:C.2.UN.16} identifies its minimal model as a torus
or a Hopf surface. The nonzero canonical form
in \eqref{cdp:eq:C.K3.3} excludes Hopf. Thus \eqref{cdp:eq:C.K3.2} and
\ref{cdp:st:C.2.UN.23} give exactly the cyclic symplectic torus
quotient of \eqref{cdp:eq:C.K3.6}--\eqref{cdp:eq:C.K3.7}. This proves necessity
of the stated algebraic-dimension-zero criterion.

\noindent\textbf{Step 10.} Finally we prove the occurrence assertion with its
specified compact comparison. Start with any torus
$A$ and one of the actual cyclic symplectic actions
above, of order $n$. Form the finite analytic quotient
\[
 W=(A\times\mathbb P^1_s)/
       \langle(a,s)\mapsto(g(a),\zeta_ns)\rangle,
 \qquad W\longrightarrow\mathbb P^1_t,
 \qquad t=s^n.
                                                               \tag{C.K3.10}\label{cdp:eq:C.K3.10}
\]
Over $t\ne0,\infty$, the $n$ sheets of the base
change are permuted transitively, and their quotient
fibre is exactly $A$. The reduced fibres over zero
and infinity are exactly $A/\langle g\rangle$.
Choose an analytic resolution of the total quotient
and its two central supports, isomorphic over the
punctured base and at a generic point of each of
these two original components. The strict transform
of the latter maps properly birationally to the
specified torus quotient. Its normalization resolution
therefore has the K3 minimal model already proved
in \ref{cdp:st:C.2.UN.23}. The resulting smooth compact total
space has connected fibres: its Stein factor over
$\mathbb P^1$ has generic degree one and hence is
the identity finite map to the normal base.

\noindent\textbf{Step 11.} Each of its two reduced special supports has even
Euler characteristic by \ref{cdp:st:C.2.UN.1}. Apply the actual
vertical point and quartic-curve blowup operations
of \ref{cdp:st:C.2.UN.4} separately over zero and infinity.
They change the reduced Euler number by $+2$ or
$-2$ until it is zero, preserve the punctured torus
family, and leave the original strict component
birational to the same torus quotient. Resolution
and these blowups are projective morphisms when
$A$ is projective. This supplies all asserted
compact realizations, including every class of
algebraic dimension zero in the criterion.

\noindent\textbf{Step 12.} The algebraic-dimension-zero occurrence criterion
is thus an equivalence. In positive algebraic
dimension the theorem gives the elliptic-fibration
restriction and realizes its cyclic quotient
subfamily by the explicit compact construction.
\end{proof}






\begin{proposition}[K3 components in arbitrary torus fillings are exactly the cyclic generalized Kummer surfaces]\label{cdp:st:C.2.UN.32}
\StatementRecord{cdp:st:C.2.UN.32}{proved}{Historical events: Claude Fable 5.2 (Anthropic), effort high, 16 November 2026: carried as background without fault found. GPT-6 Astra (OpenAI), ChatGPT Pro mode, effort Pro, S6\_09\_PAPERS, 17 November 2026: targeted proof check. Also checked in this writing. Pre-continuation writing events: GPT-6 Astra (OpenAI), ChatGPT Pro mode, effort Pro, S6\_10\_PAPERS, 18 November 2026: targeted proof check. Continuation metadata events: GPT-6 Astra (OpenAI), ChatGPT Pro mode, effort Pro, S6\_10\_PAPERS continuation, 19 November 2026: used and rechecked in this run; status-vocabulary reconciliation; metadata reconciliation only; no new mathematical check. \emph{Pre-continuation recorded proof scope.} The fibrewise kernel quotient is smooth by its actual pseudoreflection coordinates. The canonical character is lambda times beta. At a branched base value the free root-deck lift gives both m dividing s and gcd(m,s)=1; at a free orbit, the uncontracted elliptic curve and pulled K3 form give m=1. The explicit Cech Chern cocycle is beta-invariant and hence zero, and its additive torsor is the displayed rank-four translation lattice. This constructs the actual torus quotient in every algebraic dimension}{the hypotheses, notation and definitions stated here; \cref{cdp:st:C.2.UN.24}; \cref{cdp:st:C.2.UN.16}; \cref{cdp:st:C.2.UN.19}; \cref{cdp:st:C.2.UN.28}; \cref{cdp:st:C.2.UN.23}; \cref{cdp:st:C.2.UN.22}; \cref{cdp:st:C.2.UN.29}; \cref{cdp:st:C.2.UN.20}; \cref{cdp:eq:C.KX.4}; \cref{cdp:eq:C.KX.8}; \cref{cdp:eq:C.KX.10}; \cref{cdp:eq:C.KX.7}; \cref{cdp:eq:C.KX.13}; \cref{cdp:eq:C.KX.1}}

The K3 minimal birational classes occurring as
normalization components of proper two-torus
degenerations are exactly the minimal resolutions
of cyclic quotients of complex two-tori by
nontrivial symplectic automorphisms.
The generator has no eigenvalue one on the tangent space
and, after an actual translation of the origin,
is linear of order $2,3,4$, or $6$.
The exceptional configurations are those of
(C.K3.1). This criterion holds in every
algebraic dimension.

Every such class occurs in a smooth compact
filling with both special reduced supports
of Euler number zero, by the explicit quotient
and balancing construction of \ref{cdp:st:C.2.UN.24}.
In particular the algebraic-dimension-zero
equivalence of that statement extends to
all K3 components.
\end{proposition}


\begin{proof}
\noindent\textbf{Step 1.} Let $M$ be the K3 minimal model of an original
component. Take the actual multiplicity-one
component cover $E$ and its cyclic stabilizer
$G$ from \ref{cdp:st:C.2.UN.16}. The proper graph comparisons
of \ref{cdp:st:C.2.UN.19} give
\[
                         E/G\dashrightarrow M
                                                               \tag{C.KX.1}\label{cdp:eq:C.KX.1}
\]
as a proper birational comparison after
resolving the graph. Pullback of the K3
two-form is injective, whereas the
compact-core bound \ref{cdp:st:C.2.UN.28} gives $p_g(E)\le1$.
Thus
\[
                       p_g(E)=1,\qquad\kappa(E)\ge0.
                                                               \tag{C.KX.2}\label{cdp:eq:C.KX.2}
\]
Pass to its unique smooth minimal model
$E_0$; the finite action descends through
the actual point blowdowns.
If $E_0$ is a torus, \ref{cdp:st:C.2.UN.23} already gives
the required cyclic K3 quotient.

\noindent\textbf{Step 2.} In the other cases, \ref{cdp:st:C.2.UN.22} and \ref{cdp:st:C.2.UN.29} give
an actual relatively minimal elliptic map
\[
                   q:E_0\longrightarrow B
                                                               \tag{C.KX.3}\label{cdp:eq:C.KX.3}
\]
with identity integral monodromy and
only smooth multiple fibres. Its
relative Jacobian is one fixed elliptic
curve $F$. The local invariant
direct-image calculation (C.RG.13)
applies and gives
$p_g(E_0)=g(B)$, so $B$ is elliptic.
The map can be taken $G$-equivariant.
In the rank-one branch this is \ref{cdp:st:C.2.UN.20}.
In algebraic dimension one its
connected-fibre function field is
intrinsic, by (C.RG.13) and
(C.K3.4). In algebraic dimension
two a non-torus minimal model
with $p_g=1$ and $\chi(\OO)=0$
has Kodaira dimension one, so
its elliptic map is its intrinsic
Iitaka fibration, as in \ref{cdp:st:C.2.UN.24}.
Algebraic dimension zero was
already the torus branch of \ref{cdp:st:C.2.UN.16}.
These are the same equivariant
graph and Stein comparisons
used there.

\noindent\textbf{Step 3.} We first describe every multiple
fibre with its actual cyclic
uniformization. For a fibre $mC$,
the normalized root change
$t=u^m$ is \'etale over the
original smooth surface and
has connected smooth elliptic
central fibre. The identity
monodromy and constant period
ratio make this root family
a torsor under the fixed $F$
over a disc. Its torsor is
trivial: the sheaf sequence
$0\to\Lambda_F\to\OO\to F\to0$
on a disc has $H^1(\OO)=0$
and $H^2(\Lambda_F)=0$.
Thus it is an actual product.
The generator of its root
deck group has translation
linear part and can be
written
$(u,z)\mapsto(\zeta_m u,z+\tau+a(u))$,
where $\tau=a_{\rm total}(0)$
and $a(0)=0$.
Lift $a$ holomorphically to
$\C$. The finite-order
identity gives
$\sum_{j=0}^{m-1}a(\zeta_m^ju)=0$.
Replacing $z$ by $z+b(u)$,
with
$b(u)=-m^{-1}\sum_{j=0}^{m-1}j\,a(\zeta_m^ju)$,
removes $a(u)$.
The local model is therefore
\[
 (F\times\Delta_u)/\langle\delta\rangle,
 \qquad
 \delta(u,z)=(\zeta_m u,z+\tau),\quad
 \operatorname{ord}(\tau)=m,\quad t=u^m .
                                                               \tag{C.KX.4}\label{cdp:eq:C.KX.4}
\]
The exact order of $\tau$
follows from freeness of
the \'etale root covering
on its central fibre.

\noindent\textbf{Step 4.} Let $H$ be the kernel of the
action of $G$ on $B$.
Every element of $H$ acts
on a general elliptic fibre
by a translation. Indeed
a nonidentity linear part
would make the general
finite quotient fibre
rational. The actual
quotient and Stein map
would then give a rational
fibration on a birational
model of $M$, contradicting
its nonzero canonical form
by restriction to a general
rational fibre.
The translation subgroup
on the fixed $F$ is
constant: its points have
bounded finite order and
belong to a discrete
finite torsion set.

\noindent\textbf{Step 5.} The quotient
\[
              E_1=E_0/H,\qquad Q=G/H
                                                               \tag{C.KX.5}\label{cdp:eq:C.KX.5}
\]
is smooth. To check this
also at multiple fibres,
use \eqref{cdp:eq:C.KX.4}. A fibrewise
finite automorphism has
linear part one on the
reduced elliptic fibre,
so its restriction is
a translation. If it
fixes one point, it
fixes that curve pointwise.
Finite holomorphic
linearization then gives
$(x,z)\mapsto(\zeta_e x,z)$;
its quotient coordinates
are $(x^e,z)$. Every
point stabilizer is
cyclic, so these charts
prove smoothness of the
full quotient. They also
show exactly how the
fibre multiplicity can
decrease: $t=x^m$ becomes
$t=(x^e)^{m/e}$, with
$e\mid m$. The reduced
support remains a smooth
elliptic curve. The
general fibre is the
fixed isogenous elliptic
curve obtained by the
constant translation
subgroup. Its integral
linear monodromy remains
the identity.

\noindent\textbf{Step 6.} Thus $q_1:E_1\to B$ is
again a smooth-multiple
elliptic fibration with
constant Jacobian.
The same direct-image
calculation gives
\[
 p_g(E_1)=1,\quad q(E_1)=2,\quad
                 \chi(\OO_{E_1})=0 .
                                                               \tag{C.KX.6}\label{cdp:eq:C.KX.6}
\]
The actual quotient
$E_1/Q=E_0/G$ remains
birational to $M$.
The group $Q$ acts
faithfully on the base.
Its quotient fibration
has connected elliptic
general fibre. The
K3 descent argument
(C.K3.5) makes its
base $B/Q\simeq\PP^1$.
Consequently $Q$ is
nontrivial and its
action on $B$ is not
a group of translations.
After translating the
origin on $B$, its
generator is a linear
elliptic automorphism.
The derivative character
$\lambda:Q\to\C^*$ is
faithful.

\noindent\textbf{Step 7.} We record the exact
canonical character.
Choose a nonzero
one-form $\omega_B$
on the elliptic base
and a time form $dz$
for the fixed Jacobian.
Their wedge glues
over the translation
trivializations. In
\eqref{cdp:eq:C.KX.4} it is
\[
 \Omega=\omega_B\wedge dz
        =m\,c(u^m)u^{m-1}\,du\wedge dz ,
                  \qquad c(0)\ne0 .
                                                               \tag{C.KX.7}\label{cdp:eq:C.KX.7}
\]
It is invariant under
$\delta$ and extends
holomorphically. Its
divisor is exactly
$\sum_b(m_b-1)C_b$.
The group acts on
the fixed elliptic
Jacobian with a
constant linear
character $\beta$;
translation terms
add only base
differentials to $dz$.
Hence
$g^*\Omega=\lambda(g)\beta(g)\Omega$.
The nonzero K3 form
pulls back through
the actual graph to
the unique canonical
form on $E_1$.
It is invariant, so
\[
                   \lambda(g)\beta(g)=1
                       \quad(g\in Q).                \tag{C.KX.8}\label{cdp:eq:C.KX.8}
\]

\noindent\textbf{Step 8.} Here is the local
multiple-fibre
obstruction, with
all covering maps
retained. Let a
base point have
stabilizer generated
by $g$ of order
$s>1$, and suppose
its fibre has
multiplicity $m$.
In \eqref{cdp:eq:C.KX.4} choose
a lift of $g$:
\[
 \widetilde g(u,z)
       =(\rho u,\beta z+a(u)),\qquad
       \rho^m=\lambda .
                                                               \tag{C.KX.9}\label{cdp:eq:C.KX.9}
\]
Such a lift is the
functorial lift of
the automorphism
to the normalized
root change, after
choosing an $m$th
root of $\lambda$.
Conjugation preserves
the root deck group.
Its action on $u$
shows that
$\widetilde g\delta
 \widetilde g^{-1}=\delta$.
Restricting this
identity to $u=0$
gives
\[
                         \beta\tau=\tau .
                                                               \tag{C.KX.10}\label{cdp:eq:C.KX.10}
\]
By \eqref{cdp:eq:C.KX.8}, $\beta$
is nonidentity and
has order $s$.
The sum of its
endomorphisms
$1+\beta+\cdots+\beta^{s-1}$
is zero. Applying
it to $\tau$ in
\eqref{cdp:eq:C.KX.10} gives
$s\tau=0$.
Since $\tau$ has
exact order $m$,
we obtain $m\mid s$.

\noindent\textbf{Step 9.} On the reduced
elliptic fibre
$F/\langle\tau\rangle$,
the automorphism
$g$ has nonidentity
linear part $\beta$.
It has a fixed
point, because
$1-\beta$ is an
isogeny of that
elliptic curve.
Choose a lift of
this point and
multiply the
chosen $\widetilde g$
by a root deck
transformation so
that it fixes
that lift.
Then $\widetilde g^s$
is a root deck
transformation
fixing that point.
The deck action
is free, so
$\widetilde g^s=1$.
In particular
$\rho^s=1$, while
$\rho^m=\lambda$
has exact order
$s$. This implies
$\gcd(m,s)=1$.
Together with
$m\mid s$, it
proves
\[
          \text{a branched base value has }m=1 .
                                                               \tag{C.KX.11}\label{cdp:eq:C.KX.11}
\]

\noindent\textbf{Step 10.} At a base point
with trivial
stabilizer, the
quotient is locally
an isomorphism
on a neighbourhood
of one fibre in
the orbit. Its
smooth elliptic
support therefore
retains multiplicity
$m$ and the zero
order $m-1$ of
\eqref{cdp:eq:C.KX.7}.
On a common smooth
resolution of
$E_1/Q$ and $M$,
the strict transform
of this curve is
elliptic and cannot
be contracted by
a birational map
of smooth surfaces:
all curves
exceptional for
successive point
blowdowns are
rational.
At its generic
point the pullback
of the nowhere-zero
K3 form must
therefore have
zero order zero.
Its equality with
the descended form
\eqref{cdp:eq:C.KX.7} forces
$m-1=0$.
This treats every
free base orbit.
We have proved
that $q_1$ has
no multiple fibre
at all.

\noindent\textbf{Step 11.} We finish by an
explicit Chern-class
comparison, without
assuming that a
finite image or
preimage is
K\"ahler.
The smooth map
$q_1:E_1\to B$
is now a principal
bundle under its
fixed elliptic
curve $F_1$.
Its Chern vector is
\[
        c\in H^2(B,\Lambda_{F_1})
                  \simeq\Lambda_{F_1}.
                                                               \tag{C.KX.12}\label{cdp:eq:C.KX.12}
\]
A bundle
automorphism
covering $g_B$
and having
fibre linear
part $\beta$
gives
$g_B^*c=\beta_*c$.
For the precise
map, lift the
translation
transition
functions on
a base cover
to $\C$.
Their triple
differences are
the lattice
cocycle for $c$.
Conjugating the
trivializations
multiplies that
cocycle by $\beta$,
up to a lattice
coboundary; this
is the displayed
equality.
The holomorphic
automorphism of
the elliptic
base has degree
one, so it acts
as the identity
on $H^2(B,\Z)$.
Thus $c=\beta c$.
For a generator,
$\beta\ne1$, and
$\beta-I$ is
invertible on
$\Lambda_{F_1}
\otimes\R$.
The lattice is
torsion-free, so
$c=0$.

\noindent\textbf{Step 12.} The sheaf sequence
\[
 0\longrightarrow\Lambda_{F_1}
 \longrightarrow\OO_B
 \longrightarrow\mathcal F_1
 \longrightarrow0
                                                               \tag{C.KX.13}\label{cdp:eq:C.KX.13}
\]
now lifts the
principal-bundle
class to
$H^1(B,\OO_B)$.
Write $B=\C/\Gamma$.
Every class in
that one-dimensional
complex space
is represented
by a constant
$(0,1)$-form
$a\,d\overline z$.
On the universal
cover its
primitive is
$a\overline z$,
and its changes
under $\gamma\in\Gamma$
are the constants
$a\overline\gamma$.
Consequently its
additive torsor,
and its image
under \eqref{cdp:eq:C.KX.13},
have the exact
quotient presentation
\[
 E_1\simeq
 \C^2\big/
 \left\langle
 (z,w)\mapsto
       (z+\gamma,w+a\overline\gamma),\
 (z,w)\mapsto(z,w+\lambda)
 \right\rangle_{
       \gamma\in\Gamma,\,
       \lambda\in\Lambda_{F_1}} .
                                                               \tag{C.KX.14}\label{cdp:eq:C.KX.14}
\]
Equality of the
Cech classes gives
this actual
isomorphism:
the difference
of two transition
systems is a
holomorphic
zero-cochain,
and translating
the fibre
coordinate by
that cochain
glues the local
maps.
The displayed
group is a
discrete rank-four
lattice of
translations,
with compact
fundamental
parallelepiped.
Thus $E_1$
is a complex
two-torus.

\noindent\textbf{Step 13.} Finally $Q$ is
cyclic and
$E_1/Q=E_0/G$
has the actual
birational
comparison
\eqref{cdp:eq:C.KX.1} to $M$.
\ref{cdp:st:C.2.UN.23} applies
to this torus,
and gives
exactly its
symplectic K3
quotient case.
The order and
exceptional
configuration
calculation
(C.K3.6)--(C.K3.7)
gives $2,3,4,6$
and (C.K3.1).
Conversely every
one of those
actual torus
quotients is
realized by
(C.K3.10), its
proper resolution,
and the two
explicit balancing
operations of
\ref{cdp:st:C.2.UN.24}. This
proves both
directions in
every algebraic
dimension.
\end{proof}






\begin{proposition}[Terminal ruled pairs, their actual contractions and the full normal extension]\label{cdp:st:C.2.UN.64}
\StatementRecord{cdp:st:C.2.UN.64}{proved}{Historical events: Claude Fable 5.2 (Anthropic), effort high, 16 November 2026: carried as background without fault found. GPT-6 Astra (OpenAI), ChatGPT Pro mode, effort Pro, S6\_09\_PAPERS, 17 November 2026: targeted proof check. Also checked in this writing. Pre-continuation writing events: GPT-6 Astra (OpenAI), ChatGPT Pro mode, effort Pro, S6\_10\_PAPERS, 18 November 2026: targeted proof check; used and rechecked in this run. Continuation metadata events: GPT-6 Astra (OpenAI), ChatGPT Pro mode, effort Pro, S6\_10\_PAPERS continuation, 19 November 2026: used and rechecked in this run; status-vocabulary reconciliation; metadata reconciliation only; no new mathematical check. \emph{Pre-continuation recorded proof scope.} The actual canonical restriction and dual normal maps give the complete section, bisection, line and cubic alternatives; the residue map dual to finite pullback excludes an elliptic-base bisection. Nakano's Main Theorem p484 and Fujiki-Nakano's corrected Theorem2-prime p640 and Proposition9-prime p642 were visually opened with their analytic hypotheses. They apply to the exact fibrewise O(-1) normal. The descended function has nonzero normal derivative at every centre, and the ordinary blowup retains the full K\_B-by-O extension and exact Euler change. The actual twistor period derivatives give the split and nonzero elliptic normal cocycles; their Laurent constant coefficient proves nonsplitting and scalar variation realizes the prescribed extension class}{the hypotheses, notation and definitions stated here; \cref{cdp:eq:C.EL.6}; \cref{cdp:eq:C.EL.7}; \cite[p.~484]{Nakano1970}; \cite[\S2, Theorem~2' and Proposition~9']{FujikiNakano1971}; \cref{cdp:eq:C.EL.9}; \cref{cdp:eq:C.EL.11}}

Let $f:U\to\Delta$ be a proper degeneration of connected complex
two-tori from a smooth threefold. Assume that the full central fibre
$S$ is reduced and SNC. Let $T$ be one component, meeting the rest
of $S$ in exactly one smooth connected curve $C$, contained in a
single other component $P$. Let $k_T,k_P$ be the orders of the
canonical evaluation along these components. Suppose that either
$q:T\to B$ is a specified holomorphic $\PP^1$-bundle over a smooth
compact curve, or $T\simeq\PP^2$.

For a ruled $T$, precisely one of the following two numerical and
marked-pair descriptions holds. Either $C$ is a section,
\[
 K_T=-2C,\qquad k_T-k_P=1,\qquad
 N_{C/T}\simeq K_B^{-1},\qquad N_{C/P}\simeq K_B,
                                                        \tag{C.EL.1}\label{cdp:eq:C.EL.1}
\]
where the section identifies $C$ with $B$; or $B\simeq\PP^1$,
$C$ is a smooth anticanonical elliptic bisection and
\[
 K_T=-C,\qquad k_T=k_P,\qquad
 T\simeq\mathbb F_e\quad\hbox{with }e\in\{0,1,2\}.
                                                        \tag{C.EL.2}\label{cdp:eq:C.EL.2}
\]
For $T\simeq\PP^2$, either $C$ is a line and $k_T-k_P=2$,
or $C$ is a smooth cubic and $k_T=k_P$.
These assertions classify the stated terminal pairs in actual
degenerations; they make no occurrence assertion for the
anticanonical alternatives.

In the section case \eqref{cdp:eq:C.EL.1} there is an actual proper holomorphic
map $\pi:U\to U'$ onto a smooth threefold, equal to an isomorphism
off $T$, which contracts the specified ruling to a smooth curve
$B\subset U'$. It identifies $U$ with the ordinary blowup of $U'$
along $B$. The map $f$ descends to a proper holomorphic torus
degeneration $f':U'\to\Delta$, smooth along $B$; the restriction
$P\to P'$ is an isomorphism and identifies $C$ with $B$.
The full new fibre $S'$ is again reduced and SNC. Its exact normal
bundle and Euler comparisons are
\[
\begin{gathered}
 0\longrightarrow K_B\longrightarrow V:=N_{B/U'}
        \xrightarrow{\,df'\,}\OO_B\longrightarrow0,\qquad
 (T,C)\simeq\bigl(\PP_B(V),\PP_B(K_B)\bigr),\\
 N_{T/U}\simeq\OO_{\PP_B(V)}(-1)\simeq\OO_T(-C),\qquad
 e(S)=e(S')+2-2g(B).
\end{gathered}                                                   \tag{C.EL.3}\label{cdp:eq:C.EL.3}
\]
Here $\PP_B(V)$ parametrizes lines. The extension class belongs
to the actual one-dimensional space $H^1(B,K_B)$ and is retained,
with its endpoint identifications, by the contraction record.

In the plane-line case the corresponding map contracts $T$ to a
smooth point of a smooth reduced central component, identifies
$U$ with its point blowup, and gives
\[
       P=\operatorname{Bl}_{p}P',\qquad e(S)=e(S')+2.
                                                        \tag{C.EL.4}\label{cdp:eq:C.EL.4}
\]
These contractions also exist on a compact total threefold over a
compact curve. Iteration on eligible leaves terminates, retaining
the punctured torus family, with exact Euler comparison
\[
 e(S)=e(S_{\rm rem})+2a+\sum_{j=1}^{b}\bigl(2-2g(B_j)\bigr),
                                                        \tag{C.EL.5}\label{cdp:eq:C.EL.5}
\]
where $a$ is the number of recorded point contractions and the
$B_j$ are the recorded curve centres. Reversing those actual
ordinary blowups restores the original marked support and all
its neighbourhood maps.

For every specified elliptic curve $B$, both the split and the
nonzero-extension pairs in \eqref{cdp:eq:C.EL.3} occur in a compact torus
degeneration over $\PP^1$ with exactly one singular fibre. That
fibre has two components: a two-torus and the prescribed ruled
pair. Its Euler number is zero, and every other fibre is a smooth
two-torus.
\end{proposition}



\begin{proof}
\noindent\textbf{Step 1.} \par\medskip\noindent\textbf{The actual canonical and normal restrictions.}
A local generator of the rank-one direct image of the canonical
bundle evaluates to a holomorphic three-form on $U$, using the
disc coordinate to trivialize $K_\Delta$. Its divisor is vertical:
on every smooth torus fibre the corresponding nonzero canonical
form has no zeros. Thus, after shrinking the disc, its divisor
is $\sum_v k_vT_v$. Adding a common integer to the $k_v$ does not
change the following comparisons. Since $S$ is the full reduced
principal fibre, restriction to $T$ gives actual line isomorphisms
\[
 \OO_U(T)|_T\simeq\OO_T(-C),\qquad
 K_T\simeq\OO_T\bigl((k_P-k_T-1)C\bigr).
                                                        \tag{C.EL.6}\label{cdp:eq:C.EL.6}
\]
The second map is adjunction applied to the evaluation section
after dividing by a local equation of $T$ to its order $k_T$.
It also specifies a meromorphic canonical section with exactly
the indicated divisor. At $C$ the reduced equation is $xy=t$;
its product trivialization gives
$N_{C/T}\otimes N_{C/P}\simeq\OO_C$.

\noindent\textbf{Step 2.} For a ruling fibre $L\simeq\PP^1$, put $d=C\cdot L$ and
$a=k_P-k_T-1$. Intersecting \eqref{cdp:eq:C.EL.6} with $L$ gives
$-2=ad$. Thus $d>0$, $a<0$, and the integer pairs are exactly
$(d,a)=(1,-2)$ and $(2,-1)$. The map $C\to B$ is finite:
it is a nonconstant proper map of smooth compact curves. In the
degree-one case it is an isomorphism and hence a section.
Adjunction along this section gives
$K_C=\OO_C(-C)$, proving both normal identifications in
\eqref{cdp:eq:C.EL.1}, including their actual dual-normal map.

\noindent\textbf{Step 3.} In the degree-two case $K_T+C$ is trivial, so adjunction gives
$K_C\simeq\OO_C$ and $g(C)=1$. Riemann--Hurwitz restricts
$g(B)$ to zero or one. We exclude the latter by the actual
cohomology maps. The projection satisfies
$q_*\OO_T=\OO_B$ and $R^1q_*\OO_T=0$, by the cohomology of
$\PP^1$ in a local bundle chart. The restriction
\[
 H^1(T,\OO_T)\longrightarrow H^1(C,\OO_C)
                                                        \tag{C.EL.7}\label{cdp:eq:C.EL.7}
\]
is, under this Leray isomorphism, pullback by $C\to B$.
For $g(B)=1$ that degree-two map is etale; its trace composed
with pullback is multiplication by two. Hence \eqref{cdp:eq:C.EL.7} is
nonzero. On the other hand the residue exact sequence
\[
 0\longrightarrow K_T\longrightarrow K_T(C)
       \longrightarrow K_C\longrightarrow0
                                                        \tag{C.EL.8}\label{cdp:eq:C.EL.8}
\]
has $H^0(K_T)=0$, $H^0(K_T(C))=\C$ and $H^0(K_C)=\C$.
Its residue map is consequently an isomorphism and its
connecting map is zero. Serre duality identifies that connecting
map with the dual of \eqref{cdp:eq:C.EL.7}, by the residue pairing. This
contradiction excludes $g(B)=1$. Thus $B=\PP^1$ and
$T=\mathbb F_e$. The negative section $C_0$ has
$(-K_T)\cdot C_0=2-e$. A distinct irreducible effective
anticanonical curve has nonnegative intersection with $C_0$,
so $e\le2$. This proves \eqref{cdp:eq:C.EL.2}.

\noindent\textbf{Step 4.} On the plane, write $C\sim dH$. Equation \eqref{cdp:eq:C.EL.6} instead
gives $ad=-3$. Thus $(d,a)$ is $(1,-3)$ or $(3,-1)$,
giving exactly the stated line and cubic alternatives.

\par\medskip\noindent\textbf{The corrected analytic contraction theorem and its application.}
We use the Main Theorem of \cite[p.~484]{Nakano1970}, in its
corrected and strengthened form in
\cite[\S2, Theorem~2' and Proposition~9']{FujikiNakano1971}.
For a smooth divisor which is a holomorphic projective-space
bundle, the condition that its ambient normal line restricts to
$\OO(-1)$ on every projective fibre produces a smooth analytic
blowdown with that bundle as its exceptional divisor. The
supplement repairs the cohomology argument and proves that this
normal condition suffices. These are the versions read and used
here; their local blowdown maps glue by the actual isomorphisms
off the exceptional divisor.

\noindent\textbf{Step 5.} In the section case, \eqref{cdp:eq:C.EL.6} restricts to $\OO_L(-1)$ on
every ruling fibre. The theorem therefore supplies the stated
proper holomorphic blowdown $\pi:U\to U'$ along the given
bundle projection. It is the ordinary blowup of a smooth curve
$B$ in a smooth threefold. The equality
$\pi_*\OO_U=\OO_{U'}$ gives a unique holomorphic $f'$ with
$f=f'\circ\pi$. This equality can also be checked in the two
ordinary blowup charts by extending holomorphic functions
across the smooth codimension-two centre. Properness follows
from $f'^{-1}(K)=\pi(f^{-1}(K))$ for every compact
$K\subset\Delta$.

\noindent\textbf{Step 6.} The coefficient of $T$ in $f^{-1}(0)$ is one, so the first
normal term of $f'$ along $B$ is a section of $N_{B/U'}^*$.
Its zero set on $\PP_B(N_{B/U'})$ is exactly $C$.
It is nonzero at every point of $B$: a zero at $b$ would
make that zero set contain the entire ruling fibre over $b$,
whereas $C$ is a section. Its normal differential is therefore
surjective along $B$. Thus $f'$ is a submersion there and
$P'=f'^{-1}(0)$ is smooth near the centre. The restriction of
the blowup to $P'$ is the blowup of its Cartier curve $B$,
which is an isomorphism. This proves $P\simeq P'$ with
the specified curve identification. Every other component
misses the exceptional divisor, so no other central component
meets this centre. The new full fibre remains reduced and SNC.

\noindent\textbf{Step 7.} The normal exact sequence of $B\subset P'\subset U'$ is
precisely the first row of \eqref{cdp:eq:C.EL.3}: its quotient is trivialized
by $df'$, and its subline is $N_{C/P}\simeq K_B$ from
\eqref{cdp:eq:C.EL.1}. The exceptional divisor of the ordinary blowup is
the lines projectivization of this very normal bundle. Its
intersection with the strict central component is
$\PP_B(K_B)$. The quotient $V\to\OO_B$ defines this section
as the zero divisor of the corresponding section of
$\OO_{\PP_B(V)}(1)$, which gives the second exact line
identification in \eqref{cdp:eq:C.EL.3}. Serre duality gives
$H^1(B,K_B)\simeq H^0(B,\OO_B)^*\simeq\C$.
In local normal frames the transition matrices have the form
\[
        \begin{pmatrix}k_{ij}&a_{ij}\\0&1\end{pmatrix}.
                                                        \tag{C.EL.9}\label{cdp:eq:C.EL.9}
\]
Changing a lift of the quotient changes $a_{ij}$ by its
actual Cech coboundary; rescaling the subline multiplies the
extension class by that scalar. Thus there are the zero and
nonzero orbits when endpoint frames are allowed to change.
The contraction retains the matrices \eqref{cdp:eq:C.EL.9}, rather than
discarding this extension class.

\noindent\textbf{Step 8.} In the plane-line case the normal line is $\OO_{\PP^2}(-1)$.
The same corrected theorem, with centre a point, produces an
ordinary point blowdown onto a smooth threefold. The exceptional
coefficient one says that the descended $f'$ has a nonzero
linear term at that point. It is consequently a submersion
there. Its central surface pulls back to its point blowup,
with exceptional curve the specified line $C$, proving
\eqref{cdp:eq:C.EL.4}. The normal-dual map gives $C^2_P=-1$, consistently
with this actual surface blowdown.

\par\medskip\noindent\textbf{Euler and neighbourhood comparisons.}
For a curve contraction, the map is an isomorphism off $T$,
with fibre $\PP^1$ over each point of $B$. Additivity gives
$e(S)-e(S')=e(T)-e(B)=e(B)=2-2g(B)$.
For a point contraction it gives $e(T)-1=3-1=2$.
The same maps on a compact total space give the identical
Euler changes there. The descended family is unchanged over
the punctured base, and its fibres there are the original
marked two-tori. Each contraction removes one central
component; every iteration therefore terminates. Summing
these equalities gives \eqref{cdp:eq:C.EL.5}. The record consists of the
actual smooth centre, its normal extension, the local ordinary
blowup charts and their unchanged exterior identifications.
Reversing the maps restores the initial threefold near its
support, so this comparison retains all neighbourhood data.

\par\medskip\noindent\textbf{Both elliptic normal extensions occur compactly.}
Fix the specified elliptic $B$, choose another elliptic curve
$E$, and set $A=B\times E$. Use the compact smooth twistor
torus family (C.PT.13) with period order $R=1$ or $R=2$.
Its central torus is exactly $A$. Fix the curve $B\times\{0\}$
in that central member. In the actual primitive period cover
near this curve, (C.PR.13) writes the deck transformation as
\[
 (p,w,t)\longmapsto
       (p+\delta(t),\gamma(t)w,t),\qquad
 \delta(t)=t^R d(t),\quad d(0)\ne0,
                                                        \tag{C.EL.10}\label{cdp:eq:C.EL.10}
\]
where $B=\C^*/\gamma(0)^{\Z}$ and $0<|\gamma(0)|<1$.
The curve is $p=t=0$. Differentiating these actual deck
coordinates gives its normal transition over this elliptic
cover:
\[
 V:\quad
 \begin{pmatrix}p\\t\end{pmatrix}\longmapsto
 \begin{pmatrix}1&\delta'(0)\\0&1\end{pmatrix}
 \begin{pmatrix}p\\t\end{pmatrix}.
                                                        \tag{C.EL.11}\label{cdp:eq:C.EL.11}
\]
The term $\gamma'(0)w$ is tangent to the central elliptic
curve and hence vanishes in this normal quotient. The other
primitive period has identity normal transition.

\noindent\textbf{Step 9.} For $R=2$, \eqref{cdp:eq:C.EL.11} is the identity, so the normal extension
is split. For $R=1$, its upper entry is nonzero. A splitting
would supply a holomorphic $b$ on $\C^*$ satisfying
$b(\gamma(0)w)-b(w)=\delta'(0)$, with one possible opposite
sign from changing the frame convention. The Laurent
expansion on $\C^*$ has zero constant coefficient on the
left and a nonzero one on the right. No such splitting
exists. Thus its extension class is nonzero in $H^1(B,\OO_B)$.
To prescribe the nonzero class while keeping the endpoint
frames and the quotient $df'$ fixed, use instead the actual
period family $V_t(w)=w+a tJ\overline w$, with $a\in\C^*$.
Its central marked torus is unchanged, and differentiation
multiplies the cocycle in \eqref{cdp:eq:C.EL.11} by $a$. Since
$H^1(B,K_B)$ is one-dimensional and the cocycle just proved
is nonzero, this realizes the prescribed class exactly.
The real period determinant is $(1+|at|^2)^2$. At infinity,
dividing the period frame by $at$ gives
$J\overline w+(s/a)w$, with $s=1/t$; this is the same
holomorphic compact completion used in (C.PT.13).

\noindent\textbf{Step 10.} Blow up this actual central curve in the compact smooth
twistor family. The exceptional marked pair is exactly
$\bigl(\PP_B(V),\PP_B(\OO_B)\bigr)$, by the ordinary
normal-bundle description just proved. The strict central
torus is unchanged because its centre is Cartier there.
Ordinary coordinates for the ideal $(p,t)$ have $t=xy$,
so the full new fibre has its two components with
multiplicity one. Its Euler number is zero by \eqref{cdp:eq:C.EL.3},
and all other fibres are the unchanged smooth two-tori.

\noindent\textbf{Step 11.} The exact degree alternatives, line-frame comparisons,
ordinary blowup Jacobians, normal-extension matrices and
Euler identities are computed in
\path{code/root_exceptional_leaves.py}.
\end{proof}






\begin{proposition}[Trivial local canonical class, triple points and the complete zero case]\label{cdp:st:C.2.UN.65}
\StatementRecord{cdp:st:C.2.UN.65}{proved}{Historical events: Claude Fable 5.2 (Anthropic), effort high, 16 November 2026: carried as background without fault found. GPT-6 Astra (OpenAI), ChatGPT Pro mode, effort Pro, S6\_09\_PAPERS, 17 November 2026: targeted proof check. Also checked in this writing. Pre-continuation writing events: GPT-6 Astra (OpenAI), ChatGPT Pro mode, effort Pro, S6\_10\_PAPERS, 18 November 2026: targeted proof check; used and rechecked in this run. Continuation metadata events: GPT-6 Astra (OpenAI), ChatGPT Pro mode, effort Pro, S6\_10\_PAPERS continuation, 19 November 2026: used and rechecked in this run; status-vocabulary reconciliation; metadata reconciliation only; no new mathematical check. \emph{Pre-continuation recorded proof scope.} Dividing the actual evaluation by its common fibre power proves the canonical trivialization. Adjunction gives the double-curve line maps, 2r=3t, total canonical square 3t, and the exact coherent and topological inclusion-exclusion identities e(S)=t. Surface Euler and coherent Euler sums force minimal nonrational components. The residue/trace map excludes the ruled bisection, and the Hopf eigenfunction and deck-residue signs give the individually marked diagonal axes. The actual annulus-bundle vertex groups inject into the nearby torus group, forcing each ruled degree zero and complement product identity. The displayed dt-wedge logarithmic form extends as a nowhere-zero form through xy=t. Kulikov Lemmas2.1-2.2 and their proofs were opened in the original Russian edition, printed p1010; their known-result credit remains explicit}{the hypotheses, notation and definitions stated here; \cite[Lemma 2.1]{Kulikov1977}; \cref{cdp:eq:C.CZ.3}; \cref{cdp:eq:C.CZ.5}; \cref{cdp:eq:C.CZ.7}; \cite[Lemma 2.2]{Kulikov1977}; \cref{cdp:eq:C.CZ.4}; \cref{cdp:st:C.2.UN.5}; \cref{cdp:st:C.3.1}; \cref{cdp:eq:C.CZ.8}; \cref{cdp:eq:C.CZ.9}; \cref{cdp:st:C.2.UN.19}; \cref{cdp:st:C.2.UN.43}; \cref{cdp:st:C.2.UN.59}; \cref{cdp:st:C.2.UN.57}; \cref{cdp:eq:C.CZ.12}; \cref{cdp:st:C.2.UN.46}}

Let $f:X\to\Delta$ be a proper degeneration of connected complex
two-tori, with smooth three-dimensional total space and full
reduced simple-normal-crossing central fibre
$S=\sum_iV_i$. Suppose that the orders of the actual canonical
evaluation along all $V_i$ are equal. Equivalently, the canonical
line of a neighbourhood of $S$ is trivial. Let $t$ be the number
of triple points and $r$ the number of rational double curves.
Then
\[
                         e(S)=t,\qquad 2r=3t.
\tag{C.CZ.1}\label{cdp:eq:C.CZ.1}
\]
In particular $e(S)$ is a nonnegative even integer.

The complete zero-Euler list in this class is as follows.
The fibre is either a smooth two-torus, or is a cycle of smooth
unblown components, each a marked diagonal primary or secondary
Hopf surface, or a marked surface
\[
                   \PP(\OO_E\oplus L),
          \qquad E\text{ elliptic},\quad L\in\operatorname{Pic}^0(E).
\tag{C.CZ.2}\label{cdp:eq:C.CZ.2}
\]
The boundary curves are exactly the two Hopf axes or the two
ruled sections. Their actual normal lines are dual across
every edge, and the actual integral complement-lattice product
is $M=I$.

Conversely every specified marked cycle in this list with
the stated dual-normal data and $M=I$ has a compact realization
over $\PP^1$, retaining all components and elliptic gluing
maps, with smooth total space and trivial canonical line
near the specified fibre. Its other singular fibre is the
conjugate marked cycle and all reduced fibre supports have
Euler number zero. The construction chooses normal-product
constants. A smooth two-torus is realized by its product
with $\PP^1$.
\end{proposition}



\begin{proof}
\noindent\textbf{Step 1.} \par\medskip\noindent\textbf{The volume form and the actual double-curve line maps.}
Shrink the disc to remove every other critical value.
The rank-one canonical direct image has an evaluation
generator $\Omega$ whose divisor is supported on $S$.
If all its orders are $k$, the quotient
$\Omega/f^k$ is a nowhere-zero holomorphic three-form.
This is the actual trivialization of $K_X$ used below.
Conversely a nowhere-zero three-form generates this
direct image: every other section is its product with
a holomorphic function constant on the compact connected
fibres, hence pulled back from the disc. Its evaluation
orders are all zero. This proves the equivalence in the
statement by the displayed multiplication and division maps.

\noindent\textbf{Step 2.} Put $D_i=\sum_{j\ne i}(V_i\cap V_j)$, with each connected
double curve included once. Adjunction and the principal
divisor $S=(f)$ give the actual isomorphism
\[
                           K_{V_i}\simeq\OO_{V_i}(-D_i).
\tag{C.CZ.3}\label{cdp:eq:C.CZ.3}
\]
For a connected double curve $C\subset V_i\cap V_j$,
let $T_C$ be its reduced divisor of triple points.
Adjunction on $V_i$ gives
\[
             K_C=(K_{V_i}+C)|_C\simeq\OO_C(-T_C),
             \qquad 2g(C)-2=-\#T_C.
\tag{C.CZ.4}\label{cdp:eq:C.CZ.4}
\]
Every double curve is therefore elliptic and has no
triple points, or is rational and has exactly two.
Each triple point lies on exactly three double curves.
Counting these actual incidences proves $2r=3t$.

\noindent\textbf{Step 3.} Restriction of the trivial line $\OO_X(S)$ to $C$
also gives the exact triple-point comparison
\[
 N_{C/V_i}\otimes N_{C/V_j}\otimes\OO_C(T_C)
                                  \simeq\OO_C.
\tag{C.CZ.5}\label{cdp:eq:C.CZ.5}
\]
Indeed $\OO_X(V_i)|_C=N_{C/V_j}$,
$\OO_X(V_j)|_C=N_{C/V_i}$, and every remaining
component contributes its actual transverse triple
point. In degrees this is
$C_{V_i}^2+C_{V_j}^2=-\#T_C$.
This is also the direct line-bundle proof of
\cite[Lemma 2.1]{Kulikov1977} in the present notation.

\par\medskip\noindent\textbf{The exact Euler identity.}
By \eqref{cdp:eq:C.CZ.3}, $K_{V_i}^2=D_i^2$.
The self-intersections of the double curves sum,
by \eqref{cdp:eq:C.CZ.5}, to $-3t$. The cross terms in the
boundary squares contribute $2$ on each of the
three surfaces through a triple point. Consequently
\[
                         \sum_iK_{V_i}^2=3t.
\tag{C.CZ.6}\label{cdp:eq:C.CZ.6}
\]
The alternating restriction sequence on the actual
SNC strata is
\[
 0\longrightarrow\OO_S\longrightarrow\bigoplus_i\OO_{V_i}
  \longrightarrow\bigoplus_C\OO_C
  \longrightarrow\bigoplus_p\C_p\longrightarrow0.
\tag{C.CZ.7}\label{cdp:eq:C.CZ.7}
\]
Proper flatness gives $\chi(\OO_S)=\chi(\OO_{F_t})=0$.
Only the rational double curves contribute to
$\sum_C\chi(\OO_C)$, so \eqref{cdp:eq:C.CZ.7} gives
\[
                         \sum_i\chi(\OO_{V_i})=r-t.
\tag{C.CZ.8}\label{cdp:eq:C.CZ.8}
\]
The sequence and equality also agree with
\cite[Lemma 2.2]{Kulikov1977}.

\noindent\textbf{Step 4.} Hirzebruch--Riemann--Roch and complex Gauss--Bonnet
on each compact complex surface give
$e(V_i)=12\chi(\OO_{V_i})-K_{V_i}^2$.
Topological inclusion-exclusion on the same strata
then yields
\[
 \begin{aligned}
 e(S)&=\sum_i e(V_i)-\sum_C e(C)+t\\
     &=12(r-t)-3t-2r+t\\
     &=10r-14t=t,
 \end{aligned}
\tag{C.CZ.9}\label{cdp:eq:C.CZ.9}
\]
where the last equality uses $2r=3t$.
These formulas require no K\"ahler assumption.
The count $2r=3t$ also makes $t$ even.

\par\medskip\noindent\textbf{The zero-case surfaces are minimal.}
Assume $e(S)=0$. Equations \eqref{cdp:eq:C.CZ.1} and \eqref{cdp:eq:C.CZ.4}
give no triple points and only elliptic double
curves. If there is just one component, the
central fibre is smooth and reduced. The
actual smooth-family comparison in \ref{cdp:st:C.2.UN.5} makes
it a complex two-torus. We therefore treat
the connected reducible case, where every
$D_i$ is nonempty.

\noindent\textbf{Step 5.} Equation \eqref{cdp:eq:C.CZ.3} gives
$H^0(V_i,mK_{V_i})=H^0(V_i,\OO(-mD_i))=0$
for every $m>0$: its image in $H^0(V_i,\OO)=\C$
would be a constant vanishing on $D_i$.
Thus every component has Kodaira dimension
$-\infty$. The compact-surface classification
in \ref{cdp:st:C.3.1} makes its minimal model ruled or
of class VII.

\noindent\textbf{Step 6.} For a ruled minimal model over a curve $B$
of genus at least two, every elliptic curve
on $V_i$ maps constantly to $B$, by
Riemann--Hurwitz. A general ruling fibre
$F\simeq\PP^1$ avoids their vertical images
and every blowup centre. It would satisfy
$D_i\cdot F=0$, whereas \eqref{cdp:eq:C.CZ.3} and
adjunction give $D_i\cdot F=-K_{V_i}\cdot F=2$.
This contradiction leaves ruled genus zero
or one, and class VII.

\noindent\textbf{Step 7.} Their coherent Euler characteristics are,
respectively, one, zero and zero. Equation
\eqref{cdp:eq:C.CZ.8} is now $\sum_i\chi(\OO_{V_i})=0$,
so no rational component occurs. A ruled
component over an elliptic curve has Euler
number equal to its number of point blowups.
A class VII component has $b_1=b_3=1$ and
$e(V_i)=b_2(V_i)\ge0$. Since all double
curves are elliptic, \eqref{cdp:eq:C.CZ.9} is also
$0=e(S)=\sum_i e(V_i)$. Every term is
therefore zero. The ruled components have
no blowups. Every class VII component is
minimal with $b_2=0$, and its actual
minimal-model restriction in \ref{cdp:st:C.2.UN.19} makes
it a Hopf surface. This proves the claimed
absence of blowups and identifies both
remaining kinds of surface.

\par\medskip\noindent\textbf{The ruled boundary consists of two actual sections.}
Let $\pi:V\to E$ be one of the minimal ruled
surfaces over an elliptic curve. Each component
of $D$ is elliptic and maps finitely to $E$:
it cannot be contained in a rational ruling
fibre. The map is unramified by
Riemann--Hurwitz. Since $D\cdot F=2$, the
only possibilities are two degree-one
sections or one connected elliptic
bisection $C$.

\noindent\textbf{Step 8.} The latter is excluded by an exact residue
map. Suppose $K_V+C\simeq\OO_V$.
The sequence
\[
 0\longrightarrow K_V\longrightarrow K_V(C)
      \xrightarrow{\operatorname{res}}K_C\longrightarrow0
\tag{C.CZ.10}\label{cdp:eq:C.CZ.10}
\]
has a nonzero residue of the nowhere-zero
section of $K_V(C)$: in a local equation
$x=0$ it is the residue of a unit times
$dx\wedge dz/x$. Thus
$H^0(K_V(C))\to H^0(K_C)$ is onto and
its connecting map to $H^1(K_V)$ is zero.
Under Serre duality that connecting map
is dual to the actual restriction
$H^1(V,\OO_V)\to H^1(C,\OO_C)$.
But $\pi^*:H^1(E,\OO_E)\to H^1(V,\OO_V)$
is an isomorphism, and its composite with
restriction is pullback by the degree-two
unramified map $C\to E$. The trace of
that pullback is multiplication by two,
so it is nonzero. This contradicts the
zero connecting map and rules out the
bisection.

\noindent\textbf{Step 9.} The two disjoint sections give two
quotient lines of the rank-two ruling
bundle. Their combined quotient map is
an isomorphism on every fibre, because
the sections are distinct. It is therefore
an actual bundle isomorphism with the
sum of those two lines. Tensoring by
one inverse line gives
$V\simeq\PP(\OO_E\oplus L)$ with the
two specified sections. At this point
$L$ may have any degree; that degree
is determined below by the actual
nearby fundamental group.

\par\medskip\noindent\textbf{The Hopf boundary is the individually marked diagonal pair.}
Pull $K_H+D\simeq\OO_H$ back to an actual
central primary Hopf cover, as in \ref{cdp:st:C.2.UN.43}.
For a nonlinear contraction the only
elliptic curve is its axis $C$, with
$K=-(a+1)C$, $a\ge1$. The effective
axis line has infinite order. Thus a
reduced nonempty pullback boundary
cannot satisfy $K+D=0$ there.

\noindent\textbf{Step 10.} For a diagonal contraction, the lifted
anticanonical section is
$g(z,w)\partial_z\wedge\partial_w$,
where Hartogs extends $g$ across zero
and
\[
                        g(\alpha z,\beta w)
                                      =\alpha\beta g(z,w).
\tag{C.CZ.11}\label{cdp:eq:C.CZ.11}
\]
In the nonresonant case its only
possible monomial is $zw$.
In the resonant case a further central
power has multipliers $(\gamma^p,\gamma^q)$,
with coprime positive $p,q$.
Its monomials have weight $p+q$.
For $p,q>1$ this again gives only
$zw$. For $p>1,q=1$ the equation is
$g=w(az+bw^p)$, and reducedness
requires $a\ne0$. The triangular
coordinate change $z\mapsto z+(b/a)w^p$
commutes with the contraction and
gives the two axes. The symmetric
case is identical. When $p=q=1$,
the reduced quadratic is a product
of two distinct linear forms, and
an invertible linear change commuting
with the scalar contraction gives
the axes. Thus the pulled boundary
is exactly two axes in every case.

\noindent\textbf{Step 11.} Their unique logarithmic two-form,
up to a constant, is
$dz/z\wedge dw/w$. A deck lift
preserving both axes individually
is diagonal by the contraction-unit
argument of \ref{cdp:st:C.2.UN.59} and fixes this
form. A lift interchanging the axes
is a diagonal map followed by the
interchange and negates the form.
The nonzero section of $K_H(D)$
pulls back to a deck-invariant
section, so no interchange occurs.
Consequently the two images are
distinct elliptic curves, the full
deck group is diagonal, and the
original marked pair is precisely
one of the diagonal primary or
secondary Hopf pairs of \ref{cdp:st:C.2.UN.57}.

\par\medskip\noindent\textbf{The integral obstruction to nonzero ruled degree.}
Every component now has exactly two
boundary curves. The connected
incidence graph is a cycle. At a
ruled vertex its complement is the
actual nonzero line $L^\times$ over
$E$. For $d=\deg L$ its group is
\[
 \pi_1(L^\times)=
 \langle a,b,c\mid [a,c]=[b,c]=1,
                               \ [a,b]=c^d\rangle .
\tag{C.CZ.12}\label{cdp:eq:C.CZ.12}
\]
This is the actual circle-bundle
extension: trivialize over an
oriented elliptic parallelogram;
the boundary commutator is the
degree-$d$ clutching loop of its
fibre. The fibre generator $c$
has infinite order, since the
homotopy sequence injects
$\pi_1(\C^*)$ when $\pi_2(E)=0$.

\noindent\textbf{Step 12.} For a given smoothing take the
component cores and the actual
$xy=t$ collars. A Hermitian radial
retraction identifies a ruled
core with an annulus bundle in
$L^\times$; either boundary
inclusion induces the isomorphism
onto \eqref{cdp:eq:C.CZ.12}. For a diagonal
Hopf core, the same statement is
its actual rank-three logarithmic
lattice comparison. The dual-normal
maps identify the intervening
circle bundles. The canonical
divisor sections and a small smooth
unit root give these product-preserving
tubular maps for the arbitrary
given smoothing, exactly as in
\ref{cdp:st:C.2.UN.46}. Flatness of $L$ is not used
for this topological comparison.

\noindent\textbf{Step 13.} The resulting cycle of group
isomorphisms is a graph of groups
whose vertex groups inject into
the total group. Explicitly cut
one edge, compose the remaining
isomorphisms, and obtain a single
vertex group with one stable
letter implementing the closing
automorphism. Its normal form
contains that vertex group as
a subgroup. Since the actual
nearby group is $\Z^4$, every
injected vertex group is abelian.
Equation \eqref{cdp:eq:C.CZ.12} and the
infinite order of $c$ force
$d=0$ at every ruled vertex.
This proves \eqref{cdp:eq:C.CZ.2} by the
actual collar and clutching maps.

\noindent\textbf{Step 14.} The remaining component groups
are all their actual rank-three
abelian lattices. The same cycle
calculation gives
$\pi_1(F_t)=\Lambda\rtimes_M\Z$.
Its injected lattice and the
abelian group $\Z^4$ force
$M=I$. We have obtained exactly
the marked data of \ref{cdp:st:C.2.UN.57}, with
no omitted unblown component
or boundary alternative.

\par\medskip\noindent\textbf{Realization with a nowhere-zero local volume.}
For every such marked cycle,
\ref{cdp:st:C.2.UN.57} gives the actual proper
plumbing, the compact two-ended
period family, and its exact
punctured-family gluing. It
retains all marked components
and elliptic maps. We check
the required local canonical
condition on this construction.

\noindent\textbf{Step 15.} On its common period charts
put
\[
                       \Omega=dt\wedge d\eta_1\wedge d\eta_2.
\tag{C.CZ.13}\label{cdp:eq:C.CZ.13}
\]
Here $\eta$ is the common
vector coordinate obtained
from each component chart
by its actual moving frame
$P_i(t)$. Thus every overlap
has identity linear part in
this common frame.
Every period and affine edge
translation depends holomorphically
on $t$; its differential changes
the vector differentials only by
multiples of $dt$. Thus this
three-form descends through the
actual lattice and collar maps.
It is nonzero on component cores.
At a double curve the relative
two-form is a holomorphic unit
times $dx/x\wedge dz$, where
$z$ is an elliptic coordinate
and $xy=t$. Hence
\[
             dt\wedge(dx/x\wedge dz)
                               =dy\wedge dx\wedge dz.
\tag{C.CZ.14}\label{cdp:eq:C.CZ.14}
\]
The right side is holomorphic
and nowhere zero. The moving
Hopf and ruled frames have
nonzero determinant at the
central value, so the displayed
unit never vanishes after
shrinking the disc. The form
therefore extends across all
of $S$ as a nowhere-zero
three-form. Its evaluation
orders are equal to zero.
The compact realization in
the statement now follows
from \ref{cdp:st:C.2.UN.57}, with the same
verification at its conjugate
fibre. The smooth torus case
is its literal product family.
\end{proof}



\par\medskip\noindent\textbf{Exact relation to the primary modification literature.}
Kulikov's definition of a good degeneration in
\cite[original Russian ed., pp.~1008--1009]{Kulikov1977} is the smooth-total,
reduced SNC condition used here. Its Lemmas 2.1--2.2
give the same principal-divisor and coherent-stratum
maps as \eqref{cdp:eq:C.CZ.5}--\eqref{cdp:eq:C.CZ.8}. The general canonical
modification statement is
\cite[Theorem 1$'$]{Kulikov1981}: a good degeneration
of surfaces with trivial canonical class admits a
good modification with trivial total canonical
class precisely under its stated minimum-order
condition on class VII components. The present
equal-order hypothesis gives the actual identity
comparison to such a target through $\Omega/f^k$.
The zero-case classification above is proved on
that given neighbourhood and supplies the explicit
marked-cycle maps and compact occurrence. The
1981 appendix also records the correction to the
1977 Enriques modification assertion; the cited
line and Euler comparisons are the displayed
1977 lemmas.

Kamenova's author-hosted primary text
\cite[Theorem 3.1]{KamenovaTori} assumes a
K\"ahler total space and trivial $K_X$.
Its necessary list contains a smooth torus,
elliptic-ruled paths and cycles, and a rational
case with a triangulated torus as dual complex.
The residue and finite-trace maps in
\eqref{cdp:eq:C.CZ.10} exclude a terminal elliptic bisection,
so the path alternative contributes no
additional zero-Euler support in the stated
smooth-component setting. This sharpens that
necessary list. In the K\"ahler case Hopf
components are also excluded by restriction
of the K\"ahler form: their vanishing real
$H^2$ would make its positive square integrate
to zero by Stokes. The present proof keeps
the actual component and collar maps without
a K\"ahler hypothesis and supplies compact
occurrence of every marked zero-Euler cycle
specified in the statement.




\begin{proposition}[Actual leaf contractions, elliptic cores and the intrinsic Euler formula]\label{cdp:st:C.2.UN.66}
\StatementRecord{cdp:st:C.2.UN.66}{proved}{Historical events: Claude Fable 5.2 (Anthropic), effort high, 16 November 2026: carried as background without fault found. GPT-6 Astra (OpenAI), ChatGPT Pro mode, effort Pro, S6\_09\_PAPERS, 17 November 2026: targeted proof check. Also checked in this writing. Pre-continuation writing events: GPT-6 Astra (OpenAI), ChatGPT Pro mode, effort Pro, S6\_10\_PAPERS, 18 November 2026: targeted proof check; used and rechecked in this run. Continuation metadata events: GPT-6 Astra (OpenAI), ChatGPT Pro mode, effort Pro, S6\_10\_PAPERS continuation, 19 November 2026: used and rechecked in this run; status-vocabulary reconciliation; metadata reconciliation only; no new mathematical check. \emph{Pre-continuation recorded proof scope.} The ordinary terminal contractions preserve the actual family and the stated surface class. The finite maximum-order partitions leave either an equal-order core or a torsion-normal Hopf leaf. The actual meromorphic-field projections, period comparison, ruled finite pullback with trace, and torus or Euclidean root-cover quotient align every adjacent boundary with a fixed-Jacobian elliptic fibration. Propagation and the component-group/cohomology maps give exactly the rooted or unicyclic cores. Rational boundaries cannot meet the nonrational alternatives, so no rational component can disappear in the finite contraction sequence. Each higher-genus ruled component is removed once, yielding the exact negative Euler sum. For arbitrary elliptic double curves, the actual minimal-model blowup maps and triple-point inclusion-exclusion give the nonnegative sum and its precise zero case}{the hypotheses, notation and definitions stated here; \cref{cdp:st:C.2.UN.64}; \cref{cdp:st:C.2.UN.65}; \cref{cdp:st:C.2.UN.61}; \cref{cdp:st:C.2.UN.62}; \cref{cdp:eq:C.MR.3}; \cref{cdp:st:C.2.UN.5}; \cref{cdp:eq:C.MR.5}; \cref{cdp:st:C.2.UN.19}; \cref{cdp:st:C.2.UN.58}; \cref{cdp:st:C.2.UN.50}; \cref{cdp:st:C.2.UN.29}; \cref{cdp:st:C.2.UN.22}; \cref{cdp:st:C.2.UN.28}; \cref{cdp:eq:C.MR.10}; \cref{cdp:eq:C.MR.11}; \cref{cdp:st:C.2.UN.38}; \cref{cdp:st:C.2.UN.31}; \cref{cdp:eq:C.MR.12}; \cref{cdp:st:C.2.UN.33}; \cref{cdp:eq:C.MR.8}; \cref{cdp:eq:C.MR.4}; \cref{cdp:eq:C.MR.9}; \cref{cdp:st:C.2.UN.47}; \cref{cdp:st:C.2.UN.54}; \cref{cdp:st:C.2.UN.44}; \cref{cdp:eq:C.MR.15}; \cref{cdp:eq:C.MR.16}}

Let $f:U\to\Delta$ be a proper degeneration of connected complex
two-tori from a smooth threefold. Suppose that its full central
fibre $S=\sum_vT_v$ is reduced and SNC, has no triple points,
and that every component is either a smooth minimal compact
complex surface or a holomorphic $\PP^1$-bundle over a smooth
compact curve. The latter convention includes $\mathbb F_1$.
All components, intersection curves and neighbourhood maps refer
to this actual degeneration.

There is a finite sequence of the ordinary holomorphic
contractions of \ref{cdp:st:C.2.UN.64}, each along a terminal component, which
retains the punctured torus family and ends in one of the
following marked supports:
\begin{enumerate}
\item a smooth two-torus;
\item a marked diagonal Hopf and degree-zero elliptic ruled
cycle of \ref{cdp:st:C.2.UN.65}, with its actual dual normal maps and
complement-lattice product $M=I$;
\item a rooted elliptic tree passing the complete test of
\ref{cdp:st:C.2.UN.61}, including its actual zero-Chern inverse torsor;
\item a unicyclic genus-zero elliptic quasi-bundle support
passing the complete test of \ref{cdp:st:C.2.UN.62}.
\end{enumerate}
Every original surface and neighbourhood is recovered by
reversing the recorded ordinary blowups. In particular the
normal extension at every curve centre is retained.

No rational component occurs in $S$. Let $\mathcal R_{\ge2}$
be the set of its ruled components whose ruling bases $B_v$
have genus at least two. Then
\[
 e(S)=-2\sum_{v\in\mathcal R_{\ge2}}\bigl(g(B_v)-1\bigr)
                                                        \le0.
\tag{C.MR.1}\label{cdp:eq:C.MR.1}
\]
Consequently $e(S)=0$ exactly when $\mathcal R_{\ge2}$ is
empty. In that case all the contractions in the preceding
reduction have elliptic curve centres. Their exceptional
marked pairs are precisely
\[
 \left(\PP_B(V),\PP_B(\OO_B)\right),
 \qquad
 0\longrightarrow\OO_B\longrightarrow V
             \longrightarrow\OO_B\longrightarrow0,
 \qquad g(B)=1,
\tag{C.MR.2}\label{cdp:eq:C.MR.2}
\]
with their actual split or nonzero extension and endpoint
identifications.

Conversely, start with an actual torus degeneration having
one of the displayed cores and perform any finite sequence of
ordinary blowups along smooth compact elliptic curves in the
smooth locus of its full reduced central fibre, avoiding its
double curves at each stage. The result is again a full
reduced SNC torus fibre without triple points, its components
are in the stated minimal-or-bundle class, and its Euler
number is zero. These are actual blowup maps on the chosen
neighbourhood. The reduction does not identify an arbitrary
other realization of the same core support with that
neighbourhood or discard the recorded normal extensions.

There is also a complete extension with no initial minimality
or no-triple assumption. For any full reduced SNC torus fibre
whose connected double curves are all elliptic, $e(S)\ge0$.
Its Euler number is zero exactly when it has no triple
points and every component is minimal and nonrational.
Every such zero-Euler fibre is therefore covered by the
preceding reduction and exact reconstruction.
\end{proposition}



\begin{proof}
\noindent\textbf{Step 1.} \par\medskip\noindent\textbf{Canonical orders and actual contractions.}
For the connected reducible case let $k_v$ be the orders of
a canonical evaluation generator. Its divisor is vertical,
and adjunction with the full principal fibre gives a
meromorphic canonical section on each component with exact
divisor
\[
 K_{T_v}=
       \sum_{e=(v,w)}(k_w-k_v-1)C_e.
\tag{C.MR.3}\label{cdp:eq:C.MR.3}
\]
The notation denotes this section and its line isomorphism.
The curves in the sum are pairwise disjoint because there
are no triple points. Along each edge the actual equation
$xy=t$ gives
\[
                 N_{C_e/T_v}\otimes N_{C_e/T_w}
                                        \simeq\OO_{C_e}.
\tag{C.MR.4}\label{cdp:eq:C.MR.4}
\]

\noindent\textbf{Step 2.} Apply \ref{cdp:st:C.2.UN.64} whenever a terminal ruled component meets its
parent in a section, or a terminal plane meets its parent
in a line. Each is the actual ordinary blowdown of the
specified smooth curve or point, and the new family is
proper, smooth-total, reduced and SNC. Its other components
and marked curves are unchanged except for the parent
comparison proved there.

\noindent\textbf{Step 3.} This procedure preserves the surface convention in the
statement. At a curve contraction the parent surface is
unchanged. At a plane-line contraction the parent contains
the boundary $(-1)$-curve. A minimal surface has none.
Among the allowed projective-line bundles, a rational
curve on a bundle over a positive-genus base is a ruling
fibre and has square zero. On $\mathbb F_n$ the only
$(-1)$-curve occurs on $\mathbb F_1$ and is its negative
section: an irreducible curve different from the negative
section has class $aC_0+bf$ with $b\ge na$ when $a>0$,
and hence nonnegative square; the fibres have square zero.
Thus that parent is $\mathbb F_1$, and its specified
surface blowdown is $\PP^2$. It remains in the allowed
class. Every other boundary was disjoint from the
contracted curve, so no new intersection or triple point
is introduced.

\noindent\textbf{Step 4.} Each operation removes one component, so it terminates.
The canonical order differences on surviving components
are unchanged: the ordinary curve and point blowup
Jacobians have exceptional orders one and two,
respectively, as in (C.EL.1) and (C.EL.4). The remaining
family therefore still has \eqref{cdp:eq:C.MR.3}, and no eligible
\ref{cdp:st:C.2.UN.64} leaf. We call its support $S_0$.
For one smooth reduced component, \ref{cdp:st:C.2.UN.5} already identifies
the fibre as a two-torus. Assume below that $S_0$ is
reducible.

\par\medskip\noindent\textbf{A nonconstant maximum is a Hopf leaf.}
At a component $T$ of maximum canonical order, put
$a_e=k_T-k_w+1\ge1$. Equation \eqref{cdp:eq:C.MR.3} reads
\[
                         -K_T=\sum_e a_eC_e.
\tag{C.MR.5}\label{cdp:eq:C.MR.5}
\]
Its boundary is nonempty. If a positive pluricanonical
power had a nonzero section, its product with the
corresponding power of the nonzero section in \eqref{cdp:eq:C.MR.5}
would be a nonzero holomorphic function with a nonempty
zero divisor on the compact connected surface. This is
impossible. Thus $\kappa(T)=-\infty$. The surface
classification and the actual class-VII restriction \ref{cdp:st:C.2.UN.19}
leave a projective-line bundle, $\PP^2$, or a Hopf
surface.

\noindent\textbf{Step 5.} For a ruled $T\to B$, intersect \eqref{cdp:eq:C.MR.5} with a general
ruling fibre. No boundary curve is vertical: a vertical
curve would be a ruling fibre and would meet every
horizontal boundary, whereas some horizontal boundary
is required by the following equality. Hence
\[
              2=\sum_e a_ed_e,\qquad
              d_e=C_e\cdot(\text{ruling fibre})>0.
\tag{C.MR.6}\label{cdp:eq:C.MR.6}
\]
The integer partitions are exactly one section of
coefficient two, one bisection of coefficient one,
or two sections of coefficient one each. The first
is an eligible contraction. The terminal-pair
classification in \ref{cdp:st:C.2.UN.64} makes the second a rational
ruled surface with its smooth anticanonical elliptic
bisection; its neighbour has the same canonical order.
In the third case adjunction along either disjoint
section gives
\[
       2g(B)-2=(K_T+C_e)\cdot C_e=0.
\]
Thus $B$ is elliptic, and both neighbours again have
the same order as $T$. On $\PP^2$ disjoint effective
boundary curves can have only one component, by
Bezout. The two possibilities are the eligible
plane-line contraction or the anticanonical cubic
with an equal-order neighbour.

\noindent\textbf{Step 6.} For a Hopf surface of algebraic dimension one use
its actual elliptic map over $\PP^1$, the connected
normal-root chart and the invariant relative
differential of \ref{cdp:st:C.2.UN.58}. All its curves are fibre
supports. The canonical order formula gives
\[
 d_{T,e}:=\frac{k_w-k_T}{m_e}\in\Z_{\le0},
 \qquad
              \sum_{e\ni T}d_{T,e}
                                  =\operatorname{val}(T)-2.
\tag{C.MR.7}\label{cdp:eq:C.MR.7}
\]
An unmarked multiple fibre is excluded by the same
root formula $m\operatorname{ord}\theta+m-1=0$.
Thus the valence is one or two. Valence one gives
$k_T-k_w=m_e>0$; valence two gives two equal-order
neighbours. A nonlinear Hopf surface has its single
elliptic axis, so is a leaf; its actual canonical
axis form gives the positive drop equal to its
resonance, including the secondary quotient in \ref{cdp:st:C.2.UN.50}.

\noindent\textbf{Step 7.} For completeness, a diagonal algebraic-dimension-zero
Hopf surface cannot give another maximum case. Pull
\eqref{cdp:eq:C.MR.5} to its primary contraction cover. Its only
curves are the two axes. If both occur, the
anticanonical eigenfunction equation gives
$\alpha^{a_1-1}\beta^{a_2-1}=1$.
Since $|\alpha|,|\beta|<1$ and $a_i\ge1$, both
exponents are zero. If only one occurs, the same
equation gives a positive integral resonance
between $\alpha$ and $\beta$, contradicting
algebraic dimension zero. Thus this case has
exactly the two equal-order boundaries.

\noindent\textbf{Step 8.} It follows that every noneligible maximum component
which is not a Hopf leaf has all of its neighbours
at that same maximum. An equal-order edge cannot
lead to one of the strict-drop leaves just listed.
Following these edges therefore gives an entire
connected component of the whole incidence graph.
Connectedness of $S_0$ forces all orders to be equal.
In that event \ref{cdp:st:C.2.UN.65} applies; since there are no triple
points, its complete zero-case classification is
exactly the displayed cycle or smooth torus.
If the orders are not all equal, a maximum Hopf
leaf exists. Its boundary is elliptic and its normal
line has finite exact order.

\par\medskip\noindent\textbf{The actual field space aligns a vertical elliptic fibration.}
We record the comparison needed to propagate from that
leaf. On any component $T$ of the current reduced smooth
SNC family, the actual multiplicity-one construction
(C.CF.3)--(C.CF.7) gives two independent commuting
meromorphic fields, all of whose constant combinations
are semicomplete. Their polar divisors are confined to
the boundary of $T$. Indeed the original module consists
of holomorphic fields on $U$; multiplying its basis by
meromorphic functions of $t$ introduces poles only in
the central fibre. Cancelling the generic order on
$T$ makes the restriction holomorphic at every point
of $T$ away from its intersections with the other
components. This is a restriction on the original
component, with no different covering surface substituted.

\noindent\textbf{Step 9.} Let $q:T\to B$ be an actual connected-fibre elliptic
map for which every boundary is vertical. Over the
complement of finitely many values the fields are
holomorphic on the entire compact fibres. Their
coefficients in a base direction are holomorphic
functions on each such fibre and hence constant.
Proper Stein descent gives commuting meromorphic
fields on $B$.

\noindent\textbf{Step 10.} The projected fields have no poles, by the exact
continuation argument in (C.RF.1). At a pole of
order $r$, the primitive of the one-dimensional
time form vanishes to order $r+1$. Continue one
inverse branch around its time value, using a
small loop whose lifted base paths avoid all
other bad values. The preimage under $q$ of
the compact union of these paths is compact
and contained in the holomorphic locus.
The upstairs solution therefore continues
along that whole loop. Its endpoint has a
different base value, contradicting
semicompleteness. This proves holomorphy of
the projected fields without an assumed
descent theorem for semicompleteness.

\noindent\textbf{Step 11.} If a projection is nonzero, any other commuting
projection is a constant multiple of it:
the derivative of their meromorphic ratio
along the nonzero field is zero. Subtracting
that constant multiple upstairs gives a
nonzero vertical commuting field. The
actual period-preservation map (C.RF.2)
for this transverse/vertical pair shows
that the integral elliptic periods are
constant and monodromy is the identity.
If both projections are zero, the two
fields are translation fields on each
regular elliptic fibre. Their nonconstant
meromorphic ratio is constant there,
so its connected Stein fibres are
exactly those of $q$. Applying \ref{cdp:st:C.2.UN.29}
to this actual first integral gives
the same conclusion.

\noindent\textbf{Step 12.} Thus the fibration just constructed has
one fixed Jacobian and identity integral
monodromy. Its multiple-root comparison
in \ref{cdp:st:C.2.UN.22}/29 gives only smooth multiple
supports, and its invariant cohomology
frame gives
\[
             R^1q_*\OO_T=\OO_B,\qquad
                         p_g(T)=g(B)\le1.
\tag{C.MR.8}\label{cdp:eq:C.MR.8}
\]
The last inequality is the actual simultaneous
compact-core bound \ref{cdp:st:C.2.UN.28}.

\par\medskip\noindent\textbf{A torsion-normal elliptic boundary on a ruled component.}
Consider a component $T$ meeting a known elliptic
boundary $C$ whose normal line is torsion. It has
$C^2=0$. A minimal rational surface contains no
such elliptic curve. On $\PP^2$ a smooth elliptic
curve has degree three and square nine. On a
Hirzebruch surface write an irreducible elliptic
curve as $aC_0+bf$, with $a>0$. Its arithmetic
genus equation is
\[
           (a-1)(2b-na-2)=2.
\tag{C.MR.9}\label{cdp:eq:C.MR.9}
\]
It is distinct from $C_0$ and has $b\ge na$.
The only possibilities have $a=2$, $b=n+2$,
$n\le2$, with square eight, or $a=3,n=1,b=3$,
with square nine. These include $\mathbb F_1$.
None has square zero. On a ruled surface over
a base of genus at least two an elliptic curve
cannot be horizontal by Riemann--Hurwitz, and
cannot be contained in a rational ruling fibre.
Those components are excluded as well.

\noindent\textbf{Step 13.} Suppose that $T\to E$ is ruled over an elliptic
curve. Its finite map $C\to E$ is unramified.
Pull the bundle back along this actual map.
The diagonal copy of $C$ is a section whose
normal line is the original torsion line $L$.
Its relative degree-one line $\OO(C)$ has
rank-two direct image, and the actual
fibrewise evaluation identifies this
pulled-back surface with the corresponding
projective bundle. After twisting that
rank-two bundle by its section subline,
it has the exact form
\[
     0\longrightarrow\OO_C\longrightarrow V
                         \longrightarrow L\longrightarrow0.
\tag{C.MR.10}\label{cdp:eq:C.MR.10}
\]
If $L\ne\OO_C$, its extension group is
$H^1(C,L^{-1})=H^0(C,L)^*=0$, so it splits.
If $L=\OO_C$, the complement of this section
is the additive torsor of the actual extension
class in $H^1(C,\OO_C)$.

\noindent\textbf{Step 14.} There is a second boundary curve on $T$.
Otherwise \ref{cdp:st:C.2.UN.64} would make this terminal pair
a section leaf: its connected anticanonical
bisection over an elliptic ruling base is
excluded by the residue/finite-trace map.
It would already have been contracted.
Pull a second boundary back and normalize
one component. It is a compact curve mapping
finitely onto $C$ and lying in the complement
of the section. It supplies a section of
the pulled-back additive torsor. Finite
trace composed with pullback on $H^1(\OO)$
is multiplication by its positive degree,
so the original additive class is zero.
Thus \eqref{cdp:eq:C.MR.10} splits in this case also.

\noindent\textbf{Step 15.} Kill the torsion $L$ by a further finite
etale cover. The resulting actual surface
is $E'\times\PP^1$, with its finite Galois
covering map to $T$. The base cover is
an isogeny after selecting origins; its
deck group acts by translations on $E'$.
A deck map on the product is a translation
on $E'$ and a holomorphic map
$E'\to\operatorname{PGL}_2(\C)$ on the
projective fibre. The latter is constant:
$\operatorname{PGL}_2$ is affine, so its
coordinate functions are constant on the
compact elliptic curve. These constants
give an actual finite action on $\PP^1$.
Taking its quotient constructs
\[
        T=(E'\times\PP^1)/G
                         \longrightarrow\PP^1/G.
\tag{C.MR.11}\label{cdp:eq:C.MR.11}
\]
Its regular connected fibre is $E'/K$,
where $K$ is the kernel of the projective
action. All fibre actions are translations.
At a projective stabilizer the base action
on $E'$ is free, so the local quotient
is exactly the smooth multiple chart
of (C.QT.2). This proves fixed Jacobian
and identity integral monodromy directly.

\noindent\textbf{Step 16.} Every other boundary pulls back to a
compact curve disjoint from the fixed
section of $E'\times\PP^1$. Its second
coordinate is a holomorphic function
with values in the affine complement
of that section, and is constant.
It is therefore a constant section.
Consequently every original boundary
is a fibre support of \eqref{cdp:eq:C.MR.11}, with
its actual normal-root cover retained.

\par\medskip\noindent\textbf{A torsion-normal elliptic boundary on a nonruled component.}
A Hopf component with a torsion-normal
elliptic curve is either a nonlinear
leaf, or is an elliptic Hopf component
whose curves are fibres of its fixed
Jacobian map over $\PP^1$. The
algebraic-dimension-zero diagonal case
has nontorsion axis normals and cannot
occur here, by the exact characters
in \ref{cdp:st:C.2.UN.38} and \ref{cdp:st:C.2.UN.58}.

\noindent\textbf{Step 17.} For a minimal component of nonnegative
Kodaira dimension, the common-root
comparison in \ref{cdp:st:C.2.UN.31} has degree one in
this full reduced SNC family. Thus
the component itself is a two-torus
or has the actual smooth-multiple,
fixed-Jacobian elliptic map of that
theorem. There is no further surface
quotient in this use of \ref{cdp:st:C.2.UN.31}.
On a torus the elliptic curve $C$
is a translate of an elliptic subtorus.
Its actual quotient makes $C$ a fibre,
and every disjoint boundary is vertical.

\noindent\textbf{Step 18.} Let $q_0:T\to B_0$ be the elliptic
alternative. For $N$ divisible by
all its fibre multiplicities, its
translation-root differential gives
the actual canonical line formula
\[
 K_T^{\otimes N}
  =q_0^*\left(K_{B_0}^{\otimes N}
       \otimes\OO_{B_0}
          \left(\sum_iN(1-1/m_i)b_i\right)\right).
\tag{C.MR.12}\label{cdp:eq:C.MR.12}
\]
Put $\delta=2g(B_0)-2+\sum_i(1-1/m_i)$.
In Kodaira dimension one $\delta>0$,
as follows from the exact base
pluricanonical pushdown. Since
$C$ is elliptic and $C^2=0$,
adjunction gives $K_T\cdot C=0$.
Equation \eqref{cdp:eq:C.MR.12} then makes $C$
vertical. Every disjoint boundary
is also vertical, since a horizontal
curve meets this entire fibre support.

\noindent\textbf{Step 19.} In Kodaira dimension zero, $\delta=0$.
For an elliptic $B_0$ there are no
multiple fibres, so $q_0$ is a
principal elliptic torsor. A horizontal
$C$ maps finitely and unramifiedly
onto $B_0$ and gives a section after
that actual base pullback. The integral
Chern vector therefore pulls back to
zero. Pullback is multiplication by
the positive covering degree on the
free lattice group $H^2(B_0,\Lambda_F)$,
so this vector is zero. The additive
torsor comparison (C.PT.7)--(C.PT.8)
identifies this very surface with
a complex two-torus. This is already
the preceding torus case. If $C$ is
vertical there is nothing to change.

\noindent\textbf{Step 20.} For $B_0=\PP^1$, $\delta=0$ gives
exactly the four multiplicity lists
\[
       (2,2,2,2),\quad(3,3,3),\quad
                       (2,4,4),\quad(2,3,6).
\tag{C.MR.13}\label{cdp:eq:C.MR.13}
\]
There is an actual elliptic covering
$E\to\PP^1$ with precisely these
ramification indices. For four
specified double points it is the
double cover branched at those
points. For the three-point lists
it is the elliptic quotient by
the cyclic group of order three,
four or six, with the three
branch values moved to the
specified points by a projective
coordinate. These are the actual
quotient maps and stabilizer lists
in \ref{cdp:st:C.2.UN.33}.

\noindent\textbf{Step 21.} Normalize the pullback of $T$ by
this map. In each chart (C.QT.2)
it is the finite etale surface
map $F\times\Delta_s\to
(F\times\Delta_s)/\mu_{m_i}$;
the free translation removes
the multiple fibre. Thus the
global pullback $T'\to T$ is
an actual finite etale Galois
cover and $T'\to E$ is a
principal fixed-elliptic torsor.
If $C$ was horizontal, its
pullback components are elliptic
and map finitely to $E$.
The same trace/Chern argument
and (C.PT.8) make $T'$ an
actual two-torus.

\noindent\textbf{Step 22.} The preimage of $C$ in $T'$ is
a finite disjoint union of
parallel elliptic translates.
The deck group preserves their
elliptic subtorus and therefore
acts on its elliptic quotient.
The actual descended map
\[
 T=T'/G\longrightarrow
                  (T'/F')/G
\tag{C.MR.14}\label{cdp:eq:C.MR.14}
\]
has connected elliptic fibres.
Indeed the kernel on the base
acts by translations on $F'$:
a nontranslation finite elliptic
automorphism would have a fixed
point on a fibre, contradicting
the etale surface action.
The same argument at a base
stabilizer gives precisely
smooth multiple fibres.
The curve $C$ is one such
fibre support, and every
disjoint boundary is vertical.
The field-space comparison
leading to \eqref{cdp:eq:C.MR.8} now applies
to this very descended map,
and removes any possible
nontrivial linear monodromy.
This completes the nonruled
alignment by actual finite
cover and quotient morphisms.

\par\medskip\noindent\textbf{Propagation and the exhaustive remaining cores.}
Start with the maximum Hopf
leaf when the canonical orders
are not constant. Its boundary
has finite normal order, and
\eqref{cdp:eq:C.MR.4} gives the same property
on its neighbour. The preceding
comparisons put that neighbour
in the fixed-Jacobian elliptic
class, with all its boundaries
vertical, unless it is itself
a nonlinear leaf. Every vertical
boundary again has finite exact
normal order by (C.QT.2).
Repeat across the connected
incidence graph. No rational
or higher-genus ruled component
can occur in this propagation,
by \eqref{cdp:eq:C.MR.9} and
Riemann--Hurwitz. A one-boundary
elliptic ruled surface would be
an eligible \ref{cdp:st:C.2.UN.64} section leaf
and is absent from $S_0$.
Consequently the entire $S_0$
consists of the displayed
elliptic quasi-bundles and
nonlinear Hopf leaves.

\noindent\textbf{Step 23.} Equation \eqref{cdp:eq:C.MR.8} and the
simultaneous bound \ref{cdp:st:C.2.UN.28} give
at most one genus-one base.
If there is one, its component
has $p_g=1$. The actual
component-group and dual-complex
comparison \ref{cdp:st:C.2.UN.47} makes the
incidence graph simply connected.
It is therefore a tree. The
entire support is in the exact
scope of the extended \ref{cdp:st:C.2.UN.61}
and satisfies all its necessary
tests, including its integral
inverse-torsor obstruction.
If there is no genus-one base,
all vertices have coherent
dimensions $(1,1,0)$ and the
nonzero boundary restrictions
of (C.QT.3) or \ref{cdp:st:C.2.UN.54}.
The actual graph quotient of
\ref{cdp:st:C.2.UN.44} has at most one cycle.
Leaf elimination would give
$H^2(\OO_{S_0})=0$ on a tree,
contradicting upper
semicontinuity from a torus.
Thus it is unicyclic and the
whole support is in the exact
scope of \ref{cdp:st:C.2.UN.62}, with all of
its necessary tests satisfied.
This proves that the
contraction procedure reaches
one of the stated cores.
No support-only inverse
logarithm has been used to
assert a new smoothing:
the intermediate degenerations
are the actual \ref{cdp:st:C.2.UN.64} blowdowns,
and \ref{cdp:st:C.2.UN.61}/62 are applied to
their actual final family.

\par\medskip\noindent\textbf{Rational components cannot persist or disappear.}
We now use the proved endpoint
description to exclude rational
components in the original
support. A nonrational
component in the allowed class
has no rational boundary.
For a Hopf surface a rational
curve would lift to
$\C^2\setminus\{0\}$ and its
coordinate functions would be
constant. On a ruled surface
over a positive-genus base a
rational curve is a ruling
fibre. If it were a boundary,
disjointness would leave no
horizontal boundary. A general
ruling fibre would then have
intersection zero with the
right side of \eqref{cdp:eq:C.MR.3}, while
its canonical degree is $-2$.
This is impossible.

\noindent\textbf{Step 24.} For a minimal component with
$\kappa\ge0$, use its degree-one
\ref{cdp:st:C.2.UN.31} description. A torus has
no rational curve by its
universal covering. In the
elliptic alternative \eqref{cdp:eq:C.MR.12}
has $\delta\ge0$. A rational
curve cannot be vertical in
a smooth elliptic support.
If its map to the base had
degree $d>0$, each local
ramification order over $b_i$
would be divisible by $m_i$,
from the actual equation
$t=u x^{m_i}$. Riemann--Hurwitz
would therefore give
\[
 -2\ \ge\
 d\left(2g(B_0)-2+
              \sum_i(1-1/m_i)\right)
                              =d\delta\ge0,
\tag{C.MR.15}\label{cdp:eq:C.MR.15}
\]
a contradiction. This proves
the asserted absence of
rational boundaries at every
stage of the reduction.

\noindent\textbf{Step 25.} Whenever a rational component
is removed by \ref{cdp:st:C.2.UN.64}, its
section or line boundary is
rational. Its parent must
therefore be rational too.
The operation cannot remove
the last rational component.
Every final core, however,
has no rational component:
this is explicit in \ref{cdp:st:C.2.UN.65},
and in \ref{cdp:st:C.2.UN.61}/62 every
component has coherent
Euler zero. Thus no rational
component existed initially.
In particular no plane-line
contraction occurred.

\par\medskip\noindent\textbf{The intrinsic Euler identity and the complete zero description.}
All contractions are now
curve contractions, with
base genus at least one.
All original ruled components
of base genus at least two
must be contracted, since
none occurs in the final
cores. A curve contraction
does not modify its parent
surface, and there are no
point contractions. Hence
these are exactly the
original higher-genus ruled
components, each counted once
with its original ruling base.
Every other contracted
component has elliptic base.
The exact Euler comparison
(C.EL.5), and the zero Euler
number of each final core,
give precisely \eqref{cdp:eq:C.MR.1}.
This also proves the stated
zero criterion independently
of any choice of contraction
order.

\noindent\textbf{Step 26.} In the zero case every
centre is elliptic.
Equation (C.EL.3) then
gives \eqref{cdp:eq:C.MR.2}, including
its full normal extension.
The recorded ordinary
blowups and unchanged
exterior identifications
recover the original
neighbourhood exactly.

\noindent\textbf{Step 27.} For the converse construction,
the displayed cores are in
the stated surface class.
For the quasi-bundle vertices
their actual smooth-multiple
fibration gives $e=\chi(\OO)=0$.
In nonnegative Kodaira dimension
the argument \eqref{cdp:eq:C.MR.15} excludes
exceptional rational curves.
In Kodaira dimension $-\infty$,
the surface classification,
$\chi(\OO)=0$ and $e=0$
give an unblown elliptic ruled
surface or a Hopf surface.
The root and the canonical-flat
cores are already minimal by
\ref{cdp:st:C.2.UN.61} and \ref{cdp:st:C.2.UN.65}.

\noindent\textbf{Step 28.} Now
an elliptic centre in the
smooth central locus lies
in exactly one component.
The full fibre is reduced,
so the map is a submersion
near it. Blowing up its
two-generator ideal gives
the ordinary local equation
$t=xy$, with two full
multiplicity-one components.
The strict parent is the
blowup of a Cartier curve
and hence is unchanged;
the new component is the
projective normal bundle
over the elliptic centre.
No other boundary meets the
centre, so there are no
new triple points.
Properness and the punctured
torus family are preserved
by this actual blowup.
Its Euler change is
$2-2g(B)=0$. Iterating
proves the converse in
the statement and retains
every actual neighbourhood
map and extension.

\noindent\textbf{Step 29.} The finite maximum partitions,
the ruled elliptic-curve
intersection alternatives,
the Euclidean multiplicity
lists and the Euler
contraction identities are
checked in
\path{code/A_minimal_reduction.py}
and its saved output.

\par\medskip\noindent\textbf{All elliptic double curves force minimality in the zero case.}
Finally drop the component-minimality and no-triple
assumptions, retaining a full reduced SNC torus fibre
whose connected double curves are all elliptic.
The single-component case is the smooth two-torus
case of \ref{cdp:st:C.2.UN.5}. In the reducible case every component
has a nonempty boundary, and \eqref{cdp:eq:C.MR.3} remains the
actual adjunction formula even at triple points.

\noindent\textbf{Step 30.} No component can be birationally ruled over a
curve of genus at least two. Its map through
the actual contraction to that minimal ruled
surface and then to the base sends each
elliptic boundary to a point, by
Riemann--Hurwitz. A general ruling fibre
avoids those finitely many base values and
all point-blowup centres. Its strict
transform is a rational curve of
canonical degree $-2$ and has zero
intersection with every term of
\eqref{cdp:eq:C.MR.3}, a contradiction.

\noindent\textbf{Step 31.} Write $Z_v$ for the actual smooth minimal
model of $T_v$, and $n_v$ for the number
of point blowups in that contraction.
If its Kodaira dimension is nonnegative,
the degree-one comparison of \ref{cdp:st:C.2.UN.31} gives
a torus or the smooth-multiple elliptic
minimal model, so $e(Z_v)=0$. If it has
Kodaira dimension $-\infty$, the preceding
paragraph leaves rational, elliptic ruled
or class-VII minimal models. Their Euler
numbers are respectively strictly positive,
zero and zero; the last assertion uses
the actual Hopf restriction \ref{cdp:st:C.2.UN.19}.
Thus every $e(Z_v)$ is nonnegative,
and it is positive exactly for a
rational component. The actual
point-blowup comparisons give
$e(T_v)=e(Z_v)+n_v$.

\noindent\textbf{Step 32.} Let $t$ denote the number of triple
points. Topological inclusion-exclusion
on the actual SNC strata gives
\[
 \begin{aligned}
 e(S)&=\sum_v e(T_v)-\sum_C e(C)+t\\
     &=\sum_v e(Z_v)+\sum_v n_v+t
                                      \ge0.
 \end{aligned}
\tag{C.MR.16}\label{cdp:eq:C.MR.16}
\]
The triple term is retained until this
equality is used. If $e(S)=0$, all
nonnegative terms vanish: $t=0$,
every $n_v=0$, and no component is
rational. These are exactly the
initial hypotheses of the proved
reduction, with no higher-genus ruled
component. Conversely those same
conditions make every term in
\eqref{cdp:eq:C.MR.16} zero. This proves the
extension and its asserted exact
zero description.
\end{proof}






\begin{proposition}[A torsion-normal elliptic boundary cannot lie on a rational component]\label{cdp:st:C.2.UN.67}
\StatementRecord{cdp:st:C.2.UN.67}{proved}{Historical events: Claude Fable 5.2 (Anthropic), effort high, 16 November 2026: carried as background without fault found. GPT-6 Astra (OpenAI), ChatGPT Pro mode, effort Pro, S6\_09\_PAPERS, 17 November 2026: targeted proof check. Also checked in this writing. Pre-continuation writing events: GPT-6 Astra (OpenAI), ChatGPT Pro mode, effort Pro, S6\_10\_PAPERS, 18 November 2026: targeted proof check. Continuation metadata events: GPT-6 Astra (OpenAI), ChatGPT Pro mode, effort Pro, S6\_10\_PAPERS continuation, 19 November 2026: used and rechecked in this run; status-vocabulary reconciliation; metadata reconciliation only; no new mathematical check. \emph{Pre-continuation recorded proof scope.} The exact normal-order filtration constructs a basepoint-free pencil and its Stein map has degree one. The actual meromorphic field projections and return germs force identity integral monodromy. Vertical point contractions, the canonical intersection formula on the normalized unit-root cover, and the matched Kodaira fibre table preserve relative minimality and leave only smooth multiple elliptic supports. Their exact direct images give coherent Euler zero, contradicting rational coherent Euler one. The full Hopf-axis normal and canonical character maps then exclude every stated Hopf-rational interface}{the hypotheses, notation and definitions stated here; \cref{cdp:eq:C.RH.1}; \cref{cdp:eq:C.RH.4}; \cref{cdp:st:C.2.UN.16}; \cite[Theorem~6.2; Section~8(iv)(11), p.~585]{Kodaira1963II}; \citet[Section~2.3, pp.~5--6]{UlmerUrzua2021}; \cref{cdp:eq:C.RH.8}; \cref{cdp:eq:C.RH.7}; \cref{cdp:st:C.2.UN.37}; \cref{cdp:st:C.2.UN.50}; \cref{cdp:st:C.2.UN.53}}

Let \(f:U\to\Delta\) be a proper degeneration of connected
complex two-tori, with smooth total threefold and full
reduced SNC central fibre. Let \(P\) be a rational
component, with no minimality assumption. It cannot
have a smooth elliptic boundary curve \(J\) whose normal
line \(N_{J/P}\) is torsion and which is disjoint from
every other boundary curve on \(P\).

In particular a rational component, including a rational
elliptic surface and any point blowups of a rational
surface, cannot be adjacent to an elliptic Hopf component
or a nonlinear Hopf component. No Hopf leaf can have a
rational parent. The assertion permits rational double
curves and triple points elsewhere in the central fibre.
It imposes neither a zero-Euler hypothesis nor a
minimality hypothesis on the other components.

The obstruction is obtained on the actual component.
Its torsion normal produces the connected elliptic
pencil \(|mJ|\), where \(m=\operatorname{ord}N_{J/P}\).
The actual limiting translation fields force this
pencil to have identity integral monodromy. After its
vertical exceptional curves are contracted, it has
only smooth multiple fibres and coherent Euler
characteristic zero, contradicting
\(\chi(\OO_P)=1\).
\end{proposition}



\begin{proof}
\noindent\textbf{Step 1.} \par\medskip\noindent\textbf{The exact pencil made by the torsion normal.}
Suppose that such a component and curve exist, and put
\(N=N_{J/P}\), with exact order \(m\ge1\).
The degree of \(N\) is zero, hence \(J^2=0\).
For every \(k\ge1\) use the actual divisor restriction
sequence
\[
 0\longrightarrow\OO_P((k-1)J)
 \longrightarrow\OO_P(kJ)
 \longrightarrow N^{\otimes k}\longrightarrow0.
 \tag{C.RH.1}\label{cdp:eq:C.RH.1}
\]
For \(1\le k<m\), the last line is a nontrivial
degree-zero line on the elliptic curve. It has
both \(H^0\) and \(H^1\) zero: a nonzero section
would be nowhere zero, and Serre duality gives
the \(H^1\) assertion. Since \(P\) is rational,
\(H^1(P,\OO_P)=0\) and \(H^0(P,\OO_P)=\C\).
Induction in the actual long exact sequences gives
\[
 h^0(P,\OO_P(kJ))=1,\qquad
 H^1(P,\OO_P(kJ))=0\quad(0\le k<m).
 \tag{C.RH.2}\label{cdp:eq:C.RH.2}
\]
At \(k=m\), the restriction in \eqref{cdp:eq:C.RH.1} is
onto \(H^0(J,\OO_J)=\C\). Thus
\[
 h^0(P,\OO_P(mJ))=2,\qquad
 H^0(P,\OO_P(mJ))\longrightarrow H^0(J,\OO_J)
                         \text{ is onto}.
 \tag{C.RH.3}\label{cdp:eq:C.RH.3}
\]
One section has divisor \(mJ\); a second restricts
to a nonzero constant on \(J\). Their common
zero locus is empty: on \(J\) the second
does not vanish, and off \(J\) the first
does not vanish. They define the actual
holomorphic map
\[
                   q_0:P\longrightarrow\PP^1,\qquad
                             q_0^{-1}(\infty)=mJ .
 \tag{C.RH.4}\label{cdp:eq:C.RH.4}
\]

\noindent\textbf{Step 2.} We check its connectedness through the Stein map,
rather than assuming that a pencil has connected
fibres. Write \(P\xrightarrow q B\xrightarrow\nu\PP^1\)
for its Stein factor, and let \(d=\deg\nu\).
The injection \(H^1(B,\OO_B)\to H^1(P,\OO_P)\)
gives \(B\simeq\PP^1\).
Since the support of \eqref{cdp:eq:C.RH.4} is connected,
\(\nu^{-1}(\infty)\) has just one point.
It is totally ramified of order \(d\), and
the actual fibre of \(q\) there is
\((m/d)J\). Its principal local equation
trivializes \(\OO_P((m/d)J)|_J\), so
\[
                   N^{\otimes(m/d)}=\OO_J.
 \tag{C.RH.5}\label{cdp:eq:C.RH.5}
\]
Exact order \(m\) forces \(d=1\).
Thus \(q=q_0\) has connected fibres.
Adjunction gives \(K_P\cdot J=0\);
a general member \(F\sim mJ\) is smooth
by generic smoothness and satisfies
\[
                 2g(F)-2=F\cdot(K_P+F)=0.
 \tag{C.RH.6}\label{cdp:eq:C.RH.6}
\]
It is a smooth connected elliptic curve.
This constructs the required elliptic
map on the original nonminimal surface.

\noindent\textbf{Step 3.} Every other boundary component \(C\) on
\(P\) is disjoint from \(J\) by hypothesis.
Consequently \(C\cdot q^*(\infty)=mC\cdot J=0\).
If \(q|_C\) were nonconstant, that intersection
would be its positive degree. Hence every
boundary is vertical for this very map \(q\).

\par\medskip\noindent\textbf{The actual translation fields align this pencil.}
The component has full multiplicity one.
The saturated relative translation module
of \ref{cdp:st:C.2.UN.16} supplies two commuting meromorphic
fields \(V_1,V_2\) on \(P\), independent over
\(\C\), all of whose constant combinations
are semicomplete. Their polar divisors
are contained in the boundary of \(P\).
Indeed they are obtained from holomorphic
ambient relative fields by meromorphic
base coefficients; cancelling the generic
order on \(P\) leaves a holomorphic
restriction at every point of \(P\)
away from the other central components.
This is the original surface and its
actual fields, not a replacement cover.

\noindent\textbf{Step 4.} Remove the finitely many base values
supporting a boundary, a singular fibre
or a vertical zero of the relevant field.
On each remaining compact elliptic fibre
the fields are holomorphic. Their base
coefficients are holomorphic functions
on that fibre and hence are constant.
The connected-fibre map therefore
descends them to commuting meromorphic
fields on \(B=\PP^1\).

\noindent\textbf{Step 5.} These projected fields have no poles.
For clarity, at a pole of order \(a>0\),
the primitive of the one-dimensional
time form has a zero of order \(a+1\).
Continue one inverse branch around
its corresponding time value. Choose
the small loop so that its lifted
base paths avoid every other bad
value. Their compact union has a
compact inverse image under the proper
map \(q\), contained in the holomorphic
locus of the upstairs field. The
upstairs solution therefore continues
along the entire loop. Its final
base value is a different inverse
branch, contradicting the
single-valued maximal solution of
a semicomplete field. This is the
actual continuation map used in
(C.RF.1) and (C.MR.8); it does
not require a minimal surface.

\noindent\textbf{Step 6.} If one projection is nonzero,
the other is a constant multiple,
because the derivative of their
meromorphic ratio along a nonzero
commuting curve field is zero.
Subtract that constant multiple.
The resulting nonzero vertical
field is a translation field on
each compact regular elliptic fibre.
Its commuting transverse field
identifies neighbouring fibres and
preserves its time form. Its two
marked integral periods are
therefore locally constant and,
by analytic continuation on the
connected smooth base, globally
single-valued. Integration of a
nonzero elliptic time form is
injective on \(H_1(F,\Z)\);
the integral monodromy is the
identity.

\noindent\textbf{Step 7.} If both projections vanish,
the fields are vertical translation
fields on each regular fibre.
Their meromorphic ratio is a
nonconstant function pulled
back from \(B\), since the two
fields are independent over
\(\C\). Its connected Stein
fibres are exactly those of
the map \(q\) just constructed.
The actual ambient return germs
(C.RG.4)--(C.RG.12) apply to
this same first integral:
\(\lambda_i=d_i-k_iH\) for its
two continued integral periods.
The right-hand side is
single-valued, and injectivity
of elliptic period integration
again gives identity integral
monodromy. No elliptic pencil
on an ambient neighbourhood
has been assumed.

\par\medskip\noindent\textbf{Relative minimization and the exact singular-fibre comparison.}
Contract only the vertical
\((-1)\)-curves of \(q\).
Each contraction is the
actual smooth surface
blowdown of a rational
exceptional curve, and \(q\)
descends because that curve
is one fibre of the blowdown
and maps to a point of \(B\).
There are finitely many such
steps, since each lowers
\(b_2\). The resulting map
\[
               q_{\min}:P_{\min}\longrightarrow\PP^1
 \tag{C.RH.7}\label{cdp:eq:C.RH.7}
\]
is relatively minimal. It
has the same marked smooth
fibres and the same identity
monodromy. Moreover
\(\chi(\OO_{P_{\min}})=\chi(\OO_P)=1\),
by the direct-image equality
for every ordinary point
blowup.

\noindent\textbf{Step 8.} We now use the relatively
minimal singular-fibre
classification, precisely in
its stated scope:
\cite[Theorem~6.2; Section~8(iv)(11), p.~585]{Kodaira1963II} and \citet[Section~2.3, pp.~5--6]{UlmerUrzua2021}.
Its nonmultiple non-smooth
types have monodromies
conjugate to nontrivial
unipotent matrices, their
negatives, or the nonidentity
finite-order matrices of
orders \(2,3,4,6\).
Only the smooth type \(I_0\)
has identity monodromy.

\noindent\textbf{Step 9.} For a full multiple fibre
write its local function as
\(t=u h^a\), with \(u\) a
unit and the divisor of
\(h\) primitive. Normalize
the actual base change
\(t=s^a\) by
\[
                         w^a=u,\qquad s=wh .
 \tag{C.RH.8}\label{cdp:eq:C.RH.8}
\]
The resulting surface map
is \'{e}tale over the
original surface, because
\(w\) is a root of a unit.
Its monodromy is the
\(a\)-th power of the
original one and is
therefore the identity.
It is still relatively
minimal by the following
canonical intersection
comparison. For every
irreducible component \(C\)
of a fibre \(F\) below,
\(F\cdot C=0\) gives
\(C^2\le0\). Adjunction
gives
\(K_{P_{\min}}\cdot C
=2p_a(C)-2-C^2\ge0\),
except possibly for a
smooth rational curve of
square zero or minus one.
The latter is excluded
by relative minimality.
In the square-zero case,
all intersections with
the other fibre components
vanish. Connectedness
would give \(F=bC\),
contradicting
\(K_{P_{\min}}\cdot F=0\).
Thus every component has
nonnegative canonical
intersection. Their
positive weighted sum
is \(K_{P_{\min}}\cdot F=0\),
so each intersection
is zero. For the actual
\'{e}tale covering map
\(\pi\) in \eqref{cdp:eq:C.RH.8},
\(K_{\rm up}=\pi^*K_{\rm down}\).
Any vertical \((-1)\)-curve
\(E\) upstairs would then
satisfy
\(K_{\rm up}\cdot E
=\deg(E\to C)K_{\rm down}\cdot C=0\),
contrary to adjunction,
which gives minus one.
This uses the cycle
pushforward and does
not assume that the
image curve is embedded
or nonsingular.
Its primitive Kodaira
fibre must consequently
be \(I_0\). Thus every
fibre of \(q_{\min}\)
is a positive multiple
of a smooth elliptic
curve.

\noindent\textbf{Step 10.} The oriented period ratio
is now a single-valued
holomorphic map from a
punctured \(\PP^1\) to
the upper half-plane.
Its bounded Cayley
transform extends across
the missing points and
is constant on \(\PP^1\).
All regular fibres have
one fixed Jacobian.
Their transition maps
are translations, since
the integral monodromy
is the identity. At a
multiple value use
\eqref{cdp:eq:C.RH.8}. The free
central elliptic inertia
is a translation, and
its action on \(H^1(\OO)\)
is trivial. Taking the
actual invariant coefficient
functions gives
\[
 (q_{\min})_*\OO_{P_{\min}}=\OO_{\PP^1},
 \qquad R^1(q_{\min})_*\OO_{P_{\min}}=\OO_{\PP^1}.
 \tag{C.RH.9}\label{cdp:eq:C.RH.9}
\]
The coherent Leray
sequence now yields
\(\chi(\OO_{P_{\min}})=1-1=0\),
contradicting the preserved
value one in \eqref{cdp:eq:C.RH.7}.
This proves the obstruction
without applying a
minimal-only assertion
to the original \(P\).

\par\medskip\noindent\textbf{The Hopf interface.}
On an elliptic Hopf
component every curve
is a fibre support and
has normal of finite
exact order, by the
actual normalized-root
map in \ref{cdp:st:C.2.UN.37}. A nonlinear
Hopf has its one axis,
whose normal has exact
resonance order, including
the secondary comparison
of \ref{cdp:st:C.2.UN.50}. Distinct Hopf
curves are disjoint.
There can therefore be
no triple point on an
interface with a Hopf
component: such a point
would make two different
boundary curves meet
inside that Hopf surface.
The actual \(xy=t\)
normal-product map gives
\[
                     N_{J/P}\otimes N_{J/H}=\OO_J.
 \tag{C.RH.10}\label{cdp:eq:C.RH.10}
\]
It makes \(N_{J/P}\)
torsion, and the same
absence of triple points
makes \(J\) disjoint
from every other
boundary curve on \(P\).
The obstruction just
proved applies.

\noindent\textbf{Step 11.} Every terminal Hopf
pair has one of these
torsion-normal axes,
by the exact
one-boundary canonical
and character comparison
in \ref{cdp:st:C.2.UN.53}. The two
nontorsion axes of
an algebraic-dimension-zero
diagonal Hopf cannot
form such a leaf.
This proves all the
asserted parent and
interface exclusions.
\end{proof}






\begin{proposition}[Rational toric component pairs and the exact dual torus]\label{cdp:st:C.2.UN.68}
\StatementRecord{cdp:st:C.2.UN.68}{proved}{Historical events: Claude Fable 5.2 (Anthropic), effort high, 16 November 2026: carried as background without fault found. GPT-6 Astra (OpenAI), ChatGPT Pro mode, effort Pro, S6\_09\_PAPERS, 17 November 2026: targeted proof check. Also checked in this writing. Pre-continuation writing events: GPT-6 Astra (OpenAI), ChatGPT Pro mode, effort Pro, S6\_10\_PAPERS, 18 November 2026: targeted proof check; used and rechecked in this run. Continuation metadata events: GPT-6 Astra (OpenAI), ChatGPT Pro mode, effort Pro, S6\_10\_PAPERS continuation, 19 November 2026: used and rechecked in this run; status-vocabulary reconciliation; metadata reconciliation only; no new mathematical check. \emph{Pre-continuation recorded proof scope.} The actual surface blowdowns preserve the nodal anticanonical cycle. Rationality propagates across every boundary, and the coherent SNC complex identifies the dual complex with an oriented closed Euler-zero surface. The direct charge e(V)-number-of-boundary-components decreases exactly at an interior blowup and is unchanged by a node blowup. Equality on the plane or Hirzebruch minimal model gives the explicit coordinate torus action; reversing node ideals lifts it. The global charge sum is zero, so every original marked pair is toric and e(S)=2 times the component count. GHK2015 Definition1.2 and Proposition1.3 pp76-77 were opened for the credited toric-model comparison}{the hypotheses, notation and definitions stated here; \cref{cdp:st:C.2.UN.65}; \cref{cdp:eq:C.TC.3}; \cref{cdp:st:C.2.UN.19}; \cref{cdp:eq:C.TC.2}; \cref{cdp:eq:C.TC.5}; \cref{cdp:eq:C.TC.6}; \cref{cdp:eq:C.TC.9}; \cref{cdp:eq:C.TC.8}; \cite[Definition 1.2 and Proposition 1.3]{GHK2015}; \cref{cdp:eq:C.TC.11}; \cref{cdp:eq:C.TC.12}}

Let $f:X\to\Delta$ have the full reduced simple-normal-crossing
central fibre $S=\sum_{i=1}^{v}V_i$ of \ref{cdp:st:C.2.UN.65}, and let its
canonical evaluation orders be equal. Suppose that $e(S)>0$.
Then every $V_i$ is rational and its actual boundary
$D_i=\sum_{j\ne i}V_i\cap V_j$ is the full toric boundary
for an actual algebraic two-torus action on $V_i$.
The dual complex of this given fibre is a closed real
two-torus. If $r$ counts its connected double curves and
$t$ its triple points, then
\[
                  r=3v,\qquad t=2v,\qquad e(S)=2v.
\tag{C.TC.1}\label{cdp:eq:C.TC.1}
\]
In particular $v\ge3$ and $e(S)\ge6$ in the embedded
simple-normal-crossing setting of the statement.

Together with \ref{cdp:st:C.2.UN.65} this identifies every component
pair and the dual topological space of a canonical-flat
torus degeneration. The assertion concerns the actual
given degeneration and its actual component pairs.
No smoothing of arbitrary proposed toric gluing data
is inferred from these necessary comparisons.
\end{proposition}



\begin{proof}
\noindent\textbf{Step 1.} \par\medskip\noindent\textbf{A rational boundary cycle forces a rational component.}
Use the actual three-form $\Omega/f^k$ of \ref{cdp:st:C.2.UN.65}.
It gives
\[
               K_{V_i}+D_i=0,\qquad
               K_C=\OO_C(-T_C)
\tag{C.TC.2}\label{cdp:eq:C.TC.2}
\]
by its component residue and the principal-fibre line maps.
Thus a boundary component is either elliptic and disjoint
from the remaining boundary, or rational and meets it
at exactly two transverse points. Consequently the
rational boundary components form disjoint nodal cycles.
The positive Euler identity $e(S)=t$ of \ref{cdp:st:C.2.UN.65} supplies
at least one such cycle.

\noindent\textbf{Step 2.} We check how an actual point blowdown on a component
acts on this divisor. Let $E$ be a smooth rational
$(-1)$-curve and $\pi:V\to V'$ its ordinary blowdown.
Adjunction gives
\[
                         D\cdot E=1.
\tag{C.TC.3}\label{cdp:eq:C.TC.3}
\]
If $E$ is not in $D$, it meets exactly one smooth
boundary point transversely. It cannot meet a node,
which would contribute at least two to \eqref{cdp:eq:C.TC.3}.
The image $D'=\pi_*D$ is again reduced nodal and
$K_{V'}+D'=0$. Locally this is the ordinary blowup
at a smooth point of $D'$ with $D$ its strict transform.

\noindent\textbf{Step 3.} If $E$ is in $D$, then
$(D-E)\cdot E=2$. Its two intersection points are
distinct and transverse. The blowdown identifies
their two transverse tangent directions and creates
one ordinary node on $D'$. The two branches may
belong to the same irreducible curve. This is the
actual inverse of the blowup of a boundary node;
again $K_{V'}+D'=0$ by the ordinary canonical
Jacobian formula. In particular a rational cycle
persists. When it has become an irreducible nodal
rational curve, it cannot be the smooth exceptional
curve of a further ordinary point blowdown.

\noindent\textbf{Step 4.} These two maps preserve the separation between
the rational cycles and the elliptic boundary
components. Contracting the point blowups to
a minimal model therefore retains a rational
cycle whenever one was present. Each component
has $\kappa=-\infty$, because $-K=D$ is
nonempty effective. A nonrational ruled minimal
model has rational curves only in its ruling
fibres: a map from $\PP^1$ to the positive-genus
base is constant. Those fibres are mutually
disjoint smooth rational curves and cannot
contain a nodal rational cycle. The other
nonrational possibility is the Hopf minimal
model of the actual class-VII component,
by \ref{cdp:st:C.2.UN.19}. A rational curve on a Hopf surface
would lift to $\C^2\setminus\{0\}$, where
its coordinate functions are constant. It
has no rational cycle either. Hence a
component whose boundary contains a rational
cycle is rational.

\noindent\textbf{Step 5.} On a rational component the exact sequence
\[
 0\longrightarrow K_V\longrightarrow\OO_V
                 \longrightarrow\OO_D\longrightarrow0
\tag{C.TC.4}\label{cdp:eq:C.TC.4}
\]
has $H^0(K_V)=H^1(K_V)=0$, the second vanishing
being Serre dual to $H^1(\OO_V)=0$.
It follows that $H^0(\OO_D)=\C$, so $D$ is
connected. If $D$ contains a rational curve,
it is therefore exactly one rational nodal
cycle. A neighbouring component contains
that same rational double curve and hence
is rational by the preceding argument.
Propagation through the connected fibre
shows that every $V_i$ is rational and
every $D_i$ is one rational cycle.

\par\medskip\noindent\textbf{The actual dual complex and its cohomology.}
Let $\mathcal B$ be the dual complex with
one vertex per $V_i$, one edge per connected
double curve and one triangle per triple
point. Multiple edges and triangles are
retained. Every edge belongs to two
triangles by \eqref{cdp:eq:C.TC.2}, and the link of
each vertex is exactly the boundary cycle
on that component. Thus the actual
topological realization of $\mathcal B$
is a connected closed real surface.

\noindent\textbf{Step 6.} All positive-degree structure-sheaf
cohomology of the rational components
and rational double curves vanishes.
The SNC restriction sequence therefore
identifies, by its alternating restriction
maps, the full complex computing
$H^*(S,\OO_S)$ with
\[
           \C^{v}\longrightarrow\C^{r}
                                  \longrightarrow\C^{t},
\tag{C.TC.5}\label{cdp:eq:C.TC.5}
\]
the cellular cochain complex of
$\mathcal B$. The principal-fibre
adjunction and the actual volume form
give $\omega_S\simeq\OO_S$.
Serre duality on the compact Gorenstein
surface $S$ consequently gives
\[
             H^2(S,\OO_S)\simeq
                         H^0(S,\OO_S)^*\simeq\C.
\tag{C.TC.6}\label{cdp:eq:C.TC.6}
\]
Here $H^0(S,\OO_S)=\C$ follows from
compactness, connectedness and reducedness.
Equations \eqref{cdp:eq:C.TC.5}--\eqref{cdp:eq:C.TC.6} make
$\mathcal B$ orientable. Proper flatness
gives $\chi(\OO_S)=0$, hence
$\chi(\mathcal B)=0$. A connected closed
oriented surface with this Euler
characteristic is a real two-torus.

\noindent\textbf{Step 7.} The same cochain counts give
$v-r+t=0$. Together with the exact
incidence count $2r=3t$ of \ref{cdp:st:C.2.UN.65},
this proves \eqref{cdp:eq:C.TC.1}. A triple point
in a simple-NC divisor lies on three
distinct embedded components, so
$t>0$ implies $v\ge3$.

\par\medskip\noindent\textbf{An elementary charge test with its actual equality morphisms.}
For a smooth rational surface with an
anticanonical nodal rational cycle
$D=D_1+\cdots+D_n$, allowing an
irreducible nodal rational curve when
$n=1$, put
\[
                        Q(V,D)=e(V)-n=e(V\setminus D).
\tag{C.TC.7}\label{cdp:eq:C.TC.7}
\]
Indeed a rational nodal cycle has
Euler characteristic $n$, including
the one-component case. We prove
\[
 Q(V,D)\ge0,\qquad Q(V,D)=0
       \quad\Longleftrightarrow\quad
       (V,D)\text{ is an actual toric pair}.
\tag{C.TC.8}\label{cdp:eq:C.TC.8}
\]

\noindent\textbf{Step 8.} Follow any sequence of ordinary
$(-1)$-curve contractions to a
minimal rational surface. The two
maps already proved after \eqref{cdp:eq:C.TC.3}
keep an anticanonical nodal cycle.
At a contraction outside the boundary,
$e(V)$ drops by one and $n$ is
unchanged. At a boundary contraction,
both drop by one. Thus
\[
 Q(V,D)=Q(V_{\min},D_{\min})+a,
\tag{C.TC.9}\label{cdp:eq:C.TC.9}
\]
where $a$ is the number of actual
contracted exceptional curves outside
the current boundary.

\noindent\textbf{Step 9.} On $\PP^2$ the boundary has total
degree three, so $n\le3=e(\PP^2)$.
Equality means three distinct lines;
their nodal cycle condition makes
them nonconcurrent. The invertible
matrix of their linear equations is
an actual projective transformation
to the coordinate triangle.

\noindent\textbf{Step 10.} On a minimal rational ruled surface
$q:\mathbb F_b\to\PP^1$, the total
horizontal degree of $D$ is two.
There are at most two horizontal
components. Each vertical prime is
a complete ruling fibre and meets
the horizontal part in two points
counted with multiplicity. The cycle
has only two incidences at each
horizontal component, so there are
at most two vertical components.
Hence $n\le4=e(\mathbb F_b)$.
Equality forces two degree-one
sections and two ruling fibres,
with the sections disjoint.
Their two quotient line maps split
the ruling bundle. A base coordinate
puts the two marked fibres at zero
and infinity. The resulting actual
pair is
\[
 \left(\PP_{\PP^1}(\OO\oplus\OO(b)),
        \text{two sections plus the fibres over }0,\infty\right).
\tag{C.TC.10}\label{cdp:eq:C.TC.10}
\]
In a line coordinate the transitions
are $w_\infty=z^{-b}w_0$.
The maps $(z,w_0)\mapsto(\lambda z,\mu w_0)$
and $(z^{-1},w_\infty)\mapsto
(\lambda^{-1}z^{-1},\lambda^{-b}\mu w_\infty)$
give its actual two-torus action
and precisely the four specified
boundary components.

\noindent\textbf{Step 11.} The two minimal computations prove
$Q(V_{\min},D_{\min})\ge0$.
If the initial charge is zero,
\eqref{cdp:eq:C.TC.9} says that $a=0$ and
the minimal pair is one of the
two explicit toric pairs just
identified. Reversing the recorded
contractions therefore blows up
only boundary nodes. Each such
node is fixed by the torus action;
the invariant ideal $(x,y)$ gives
the actual lifted action on its
blowup charts $y=xu$ and $x=yv$.
The new full boundary is the
toric boundary. Induction gives
an actual torus action on the
original marked pair $(V,D)$.
Conversely a smooth complete
toric surface has one fixed
point per boundary node; its
orbit decomposition gives
$e(V)=n$, proving the reverse
implication in \eqref{cdp:eq:C.TC.8}.

\noindent\textbf{Step 12.} This direct argument agrees with
the actual toric-model comparison
of \cite[Definition 1.2 and Proposition 1.3]{GHK2015},
whose hypotheses are exactly a
smooth rational projective surface
and an anticanonical nodal rational
cycle. The displayed contractions
prove the equality case without
discarding the original pair.

\par\medskip\noindent\textbf{The global charge is zero.}
Write $n_i$ for the number of
curves in $D_i$. Every double
curve belongs to two component
boundaries, so
\[
                         \sum_i n_i=2r=3t.
\tag{C.TC.11}\label{cdp:eq:C.TC.11}
\]
The exact HRR and boundary-square
calculation in \ref{cdp:st:C.2.UN.65} gives
\[
            \sum_i e(V_i)
             =12(r-t)-3t=3t.
\tag{C.TC.12}\label{cdp:eq:C.TC.12}
\]
Subtracting \eqref{cdp:eq:C.TC.11} from
\eqref{cdp:eq:C.TC.12} yields
$\sum_i Q(V_i,D_i)=0$.
Every summand is nonnegative by
\eqref{cdp:eq:C.TC.8}, so each is zero.
The actual equality morphisms
above make every specified
$(V_i,D_i)$ a toric pair.
This proves the asserted
classification on the given
degeneration.
\end{proof}






\section{Compact occurrence of prescribed component classes}\label{sec:II:09}


\begin{proposition}[Zero-Euler product fillings with essential Kodaira-dimension-one components]\label{cdp:st:C.2.UN.6}
\StatementRecord{cdp:st:C.2.UN.6}{proved}{Historical events: Claude Fable 5.2 (Anthropic), effort high, 16 November 2026: carried as background without fault found. GPT-6 Astra (OpenAI), ChatGPT Pro mode, effort Pro, S6\_09\_PAPERS, 17 November 2026: targeted proof check. Also checked in this writing. Pre-continuation writing events: GPT-6 Astra (OpenAI), ChatGPT Pro mode, effort Pro, S6\_10\_PAPERS, 18 November 2026: targeted proof check; used and rechecked in this run. Continuation metadata events: GPT-6 Astra (OpenAI), ChatGPT Pro mode, effort Pro, S6\_10\_PAPERS continuation, 19 November 2026: used and rechecked in this run; status-vocabulary reconciliation; metadata reconciliation only; no new mathematical check. \emph{Pre-continuation recorded proof scope.} The hypersurface sigma=a\textasciicircum{}m*ell, free q-action, normalized finite quotient charts and exact punctured product were re-derived. The canonical form, floor-function plurigenera, solid-torus meridian and longitude maps give the stated line, all plurigenera and integral H1. The actual multiplicity-one base-change chart and non-uniruled divisor comparison give persistence and separation. The following unnumbered audit was also checked: the compact diagonal Hopf threefold maps to P2 by [a,b,c] to [a\textasciicircum{}m,b,c], with connected elliptic fibres and full pullback m times a smooth divisor over a line. It meets the exact Fujimoto pp91-92 definition and refutes the printed Theorem2.8 divisibility and fundamental-group assertions; the covering equation splits into its m global graphs}{the hypotheses, notation and definitions stated here; \cref{cdp:eq:C.ER.4}; \cref{cdp:eq:C.ER.5}; \cref{cdp:eq:C.ER.7}; \cref{cdp:eq:C.ER.8}; \cref{cdp:eq:C.ER.14}; \cref{cdp:eq:C.ER.9}; \cref{cdp:eq:C.ER.3}}

Put $q=e^{-2\pi}$ and $E=\C^*/q^{\Z}\simeq\C/(\Z+\mathrm i\Z)$.
For every integer $m\ge2$, the fixed punctured product
$(E\times E)\times\Delta^*$ has a proper filling
$f_m:U_m\to\Delta$ with smooth total space and full central fibre
\[
                    (U_m)_0=P_m+mQ_m.             \tag{C.ER.1}\label{cdp:eq:C.ER.1}
\]
Here $Q_m$ is a primary Hopf surface, and $P_m$ is a smooth compact
non-K\"ahler elliptic surface over $E$ with precisely one multiple
fibre $mJ$. The surfaces meet transversely along the smooth elliptic
curve $J$. One has
\[
 \begin{gathered}
 e((U_m)_0{}_{\rm red})=0,\qquad H_1(P_m,\Z)=\Z^3,\\
 K_{P_m}\simeq\OO_{P_m}((m-1)J),\qquad \kappa(P_m)=1,\\
 h^0(P_m,K_{P_m}^{\otimes r})=
 \begin{cases}
 1,&d_r=0,\\
 d_r,&d_r>0,
 \end{cases}
 \qquad d_r=\left\lfloor\frac{r(m-1)}m\right\rfloor,
 \quad r\ge1.
 \end{gathered}                                   \tag{C.ER.2}\label{cdp:eq:C.ER.2}
\]
Neither $P_m$ nor $Q_m$ can be contracted by a bimeromorphic map
between smooth threefold fillings. The $U_m$ are pairwise
non-bimeromorphic over the disc. After every finite base change and
resolution of its normalization, a central component remains
birational to $P_m$. Thus a properly elliptic component of Kodaira
dimension one is forced to survive in these examples, including
after finite base change. All these fillings have compact smooth
threefold realizations over $\PP^1$.
\end{proposition}


\begin{proof}
\noindent\textbf{Step 1.} We give the construction as an actual hypersurface and a free
quotient. Let
\[
 V=\operatorname{Bl}_{(0,0)}(E\times\Delta),\quad
 R=\operatorname{Exc}(V\to E\times\Delta),\quad
 L=\OO_V(R),\quad \sigma\in H^0(V,L)
                                                        \tag{C.ER.3}\label{cdp:eq:C.ER.3}
\]
with $\sigma$ the canonical section. Define
\[
 W_m=\{(v,\ell,a):v\in V,\ \ell\in L_v^\times,\ a\in\C,
                         \ \sigma(v)=a^m\ell\}.
                                                        \tag{C.ER.4}\label{cdp:eq:C.ER.4}
\]
The generator of a $\Z$-action is
\[
             (v,\ell,a)\longmapsto(v,q^m\ell,q^{-1}a),
 \qquad U_m=W_m/\Z.                               \tag{C.ER.5}\label{cdp:eq:C.ER.5}
\]
The action preserves \eqref{cdp:eq:C.ER.4}. Near $R$, choose a holomorphic frame
of $L$ in which $\sigma=xe$ and write $\ell=be$, $b\ne0$.
Then \eqref{cdp:eq:C.ER.4} is $x=a^m b$ and eliminates the coordinate $x$;
thus $W_m$ is smooth. It is also smooth off $R$, where $a\ne0$
and $\ell=\sigma/a^m$. A Hermitian norm on $L$ is multiplied by
$q^m$ under \eqref{cdp:eq:C.ER.5}. It is bounded above and bounded away from zero
on each compact subset of $W_m$, proving proper discontinuity.
No nonzero power fixes a point, since $\ell\ne0$. Consequently
$U_m$ is a smooth complex threefold.

\noindent\textbf{Step 2.} We verify properness and the local maps explicitly. With $y$ a
second base coordinate, the map
\[
 (z,w,y)\longmapsto
 (x=z^m,y;\ b=w^m,\ a=z/w),\qquad w\in\C^*,
                                                        \tag{C.ER.6}\label{cdp:eq:C.ER.6}
\]
identifies the local $W_m$ with the quotient by
$(z,w,y)\mapsto(\zeta z,\zeta w,y)$, $\zeta^m=1$.
It is onto because an $m$th root of $b$ determines $w$ and then
$z=aw$; its fibres are exactly the displayed orbits. The action is
free. The generator \eqref{cdp:eq:C.ER.5} lifts to $(z,w,y)\mapsto(z,qw,y)$.
Hence the local projection $\rho_m:U_m\to V$ is exactly
\[
 (\Delta_z\times E\times\Delta_y)/\boldsymbol\mu_m
        \longrightarrow\Delta_x\times\Delta_y,
 \quad [z,[w],y]\longmapsto(z^m,y),
 \quad \zeta:(z,[w],y)\longmapsto(\zeta z,[\zeta w],y).
                                                        \tag{C.ER.7}\label{cdp:eq:C.ER.7}
\]
The last finite action is free: $[\zeta w]=[w]$ would imply
$\zeta=q^j$ for an integer $j$, forcing $j=0$ and $\zeta=1$.
The map before this finite quotient is the product of a finite
disc map with a compact elliptic curve. It is proper, so \eqref{cdp:eq:C.ER.7}
is proper. These charts prove that $\rho_m$ is proper everywhere.
Its fibres are connected elliptic curves, including the reduced
fibre over $R$. Thus composition with $V\to\Delta$ is proper
with connected fibres.

\noindent\textbf{Step 3.} The comparison with the punctured product is particularly simple.
For $v\notin R$ and $c\in\C^*$ set
\[
       (v,[c])\longmapsto[(v,\sigma(v)c^m,c^{-1})].
                                                        \tag{C.ER.8}\label{cdp:eq:C.ER.8}
\]
Changing $c$ to $qc$ changes the representative by \eqref{cdp:eq:C.ER.5}.
Conversely $[c]=[a^{-1}]$ is its inverse. Thus \eqref{cdp:eq:C.ER.8} is an
actual isomorphism $(V\setminus R)\times E\simeq
\rho_m^{-1}(V\setminus R)$, for every $m$, and over $\Delta^*$
it gives the specified fixed product $(E\times E)\times\Delta^*$.

\noindent\textbf{Step 4.} Let $E_0\simeq E$ be the strict transform of $E\times\{0\}$.
It meets $R\simeq\PP^1$ at one point, and
$L|_{E_0}=\OO_E([0])$, $L|_R=\OO_{\PP^1}(-1)$.
The inverse image of $R$, with its reduction, is $a=0$ in
\eqref{cdp:eq:C.ER.4}, so it is
\[
 Q_m=L|_R^\times/(q^m)^{\Z}
     \simeq(\C^2\setminus\{0\})/\langle q^m\operatorname{id}\rangle.
                                                        \tag{C.ER.9}\label{cdp:eq:C.ER.9}
\]
The first isomorphism sends the nonzero tautological vector to
the same vector in $\C^2\setminus\{0\}$. It identifies $Q_m$
as a primary Hopf surface. The reduced inverse image $P_m$ of
$E_0$ is smooth by \eqref{cdp:eq:C.ER.7}, since $E_0$ meets $R$ transversely.
Its map $p:P_m\to E_0$ has smooth elliptic fibres off $[0]$ and
one full fibre $mJ$ over $[0]$. The reduced fibre is
\[
                    J=\C^*/(q^m)^{\Z}.
                                                        \tag{C.ER.10}\label{cdp:eq:C.ER.10}
\]
For example the isomorphism $E/\boldsymbol\mu_m\to J$ in
\eqref{cdp:eq:C.ER.7} is $[w]\mapsto[w^m]$, which also specifies the finite
covering between the two elliptic curves.

\noindent\textbf{Step 5.} Near $E_0\cap R$, choose coordinates with $R=(x=0)$,
$E_0=(y=0)$ and base map $t=xy$. In the smooth chart before
the free quotient the equation is $t=z^m y$; in \eqref{cdp:eq:C.ER.4} it is
\[
                         t=a^m b y,\qquad b\ne0.
                                                        \tag{C.ER.11}\label{cdp:eq:C.ER.11}
\]
It follows that the full central fibre is exactly $P_m+mQ_m$,
and its reduction has the two transverse components stated in
\eqref{cdp:eq:C.ER.1}. This calculation also retains the ramification at $J$.

\noindent\textbf{Step 6.} We compute the canonical line without a classification assertion.
Choose a nowhere-zero holomorphic one-form $\omega$ on $E_0$.
On a local cyclic chart for $P_m$ as in \eqref{cdp:eq:C.ER.7}, the form
\[
                         p^*\omega\wedge\frac{dw}{w}
                                                        \tag{C.ER.12}\label{cdp:eq:C.ER.12}
\]
is invariant under both $w\mapsto qw$ and the finite action.
Changing a local frame changes $dw/w$ by a form from the
one-dimensional base, which has zero wedge with $p^*\omega$.
It therefore glues to a global holomorphic two-form on $P_m$.
If $\omega=h(x)\,dx$ near $[0]$, its pullback is
\[
 m h(z^m)z^{m-1}\,dz\wedge\frac{dw}{w}.
                                                        \tag{C.ER.13}\label{cdp:eq:C.ER.13}
\]
It has precisely the divisor $(m-1)J$, proving the canonical
isomorphism in \eqref{cdp:eq:C.ER.2}.

\noindent\textbf{Step 7.} For every integer $k\ge0$ there is a precise pushforward equality
\[
        p_*\OO_{P_m}(kJ)
          =\OO_E\bigl(\lfloor k/m\rfloor[0]\bigr).
                                                        \tag{C.ER.14}\label{cdp:eq:C.ER.14}
\]
A section on the left is a meromorphic function whose only allowed
poles are along $J$. Off $[0]$ it is holomorphic on every compact
elliptic fibre, so it is constant on that fibre and descends to a
holomorphic function on the base. On the cyclic chart at $[0]$,
its Laurent coefficients in $z$ are holomorphic functions on the
compact elliptic curve and hence constant. Finite-group invariance
forces the powers of $z$ to be multiples of $m$. The pole bound
$k$ is consequently exactly a pole bound $\lfloor k/m\rfloor$
in $x=z^m$. This proves both inclusions in \eqref{cdp:eq:C.ER.14}.
Applying it to $k=r(m-1)$ and applying Riemann--Roch on the elliptic
curve proves the displayed plurigenera. Their growth is linear
with positive slope $(m-1)/m$, so $\kappa(P_m)=1$.

\noindent\textbf{Step 8.} Here is also an integral topological check that this surface is
non-K\"ahler. Choose a Hermitian metric on $L|_{E_0}$ and put
\[
 M_m=\{(v,\ell,a)\in W_m|_{E_0}:\|\ell\|=1\}.
                                                        \tag{C.ER.15}\label{cdp:eq:C.ER.15}
\]
The maps $\ell=e^{mu}\ell_0$, $a=e^{-u}a_0$ identify
$W_m|_{E_0}=M_m\times\R$ smoothly. The generator \eqref{cdp:eq:C.ER.5}
translates $u$ by $\log q$, so
$P_m\simeq M_m\times S^1$ smoothly. Over the complement of a
small disc about $[0]$, the canonical section trivializes $M_m$
as the product with a circle, whose coordinate is the phase of $a$.
On the disc choose the metric to be constant in the frame used
above. Then $M_m$ has the solid-torus coordinates
$a=\rho e^{\mathrm i\theta}$, $b=e^{\mathrm i\varphi}$, and
\[
              x=a^m b,\qquad \arg x=m\theta+\varphi.
                                                        \tag{C.ER.16}\label{cdp:eq:C.ER.16}
\]
Let $\delta$ denote the boundary of the punctured elliptic base
and let $h$ denote its product circle. The meridian of this solid
torus maps to $m\delta+h$; its longitude maps to $\delta$.
The class $\delta$ is zero in the first homology of the once-punctured
elliptic curve. Mayer--Vietoris therefore kills $h$ and the solid-torus
longitude, leaving exactly the two free base classes. All spaces
and their boundary torus are connected, so the degree-zero part
adds no further class. Hence
$H_1(M_m,\Z)=\Z^2$ and $H_1(P_m,\Z)=\Z^3$.
The odd first Betti number excludes a K\"ahler metric, by Hodge
decomposition on a compact K\"ahler manifold.

\noindent\textbf{Step 9.} The Euler calculation retains the multiple fibre. Over
$E_0\setminus\{[0]\}$, $P_m$ is locally a product with an elliptic
curve; its remaining reduced fibre $J$ is also elliptic.
Stratified Euler additivity gives $e(P_m)=0$.
The Hopf surface in \eqref{cdp:eq:C.ER.9} is a principal elliptic bundle over
$\PP^1$, so $e(Q_m)=0$. Finally $e(J)=0$, and consequently
$e(P_m\cup_JQ_m)=0$.

\noindent\textbf{Step 10.} Every map $\PP^1\to P_m$ is constant: it first maps constantly
to the elliptic base $E_0$, and then constantly to the reduced
elliptic fibre containing its image. Every map $\PP^1\to Q_m$
lifts to $\C^2\setminus\{0\}$ and is constant there by its
coordinate functions. Thus neither component is covered by
rational curves. The inverse-map resolution in (C.UC.1) applies
and proves that neither can be contracted by a bimeromorphic
comparison of smooth fillings. The Kodaira dimension of $Q_m$
is $-\infty$: a lifted pluricanonical form extends across the
origin by Hartogs, and its invariance under the contraction
$q^m\operatorname{id}$ forces all its power-series coefficients
to vanish. It cannot correspond birationally to $P_{m'}$.
Therefore any bimeromorphic comparison $U_m\dashrightarrow U_{m'}$
would identify $P_m$ birationally with $P_{m'}$.
Plurigenera are unchanged by a point blowup of a smooth compact
surface: a pluricanonical form descends off the centre and its
coefficients extend across that point by Hartogs. Resolution of
a surface bimeromorphism proves their birational invariance.
The different slopes $(m-1)/m$ in \eqref{cdp:eq:C.ER.2} force $m=m'$.

\noindent\textbf{Step 11.} For the persistence assertion, take a finite base change and choose
a local parameter in which it is $t=s^d$; a unit is absorbed by
its local $d$th root. At every point of the nonempty open subset
$P_m\setminus J$, the map $f_m$ is a submersion. Its base change
is the smooth chart $t=s^d$ with the two coordinates along $P_m$
unchanged. Thus a unique central component maps with degree one
onto this open subset. Its closure, followed by normalization and
any resolution, is birational to $P_m$. The preceding
non-uniruled-divisor comparison prevents its contraction under
any further bimeromorphic modification between smooth fillings.
This proves the asserted survival without imposing a description
on the other components of the base change.

\noindent\textbf{Step 12.} Finally replace $\Delta$ throughout \eqref{cdp:eq:C.ER.3} by $\PP^1$ and use
the same exceptional divisor, line bundle, and canonical section.
The local proofs of smoothness and properness still apply. Since
the base $V$ is now compact, the quotient is a compact smooth
threefold, with the same central fibre and with every other fibre
the product $E\times E$. This gives the claimed compact realizations.
The exact quotient-chart, canonical, homology, and plurigenus
calculations are retained in \path{code/E_elliptic_roots.py} and
its output.
\end{proof}



\paragraph{Literature comparison and an explicit source audit.}
The occurrence of a Kodaira-dimension-one component in a degeneration
of abelian surfaces already appears in
\cite[Example~3.4(1), p.~99]{Fujimoto1988Log}.
The preceding statement supplies its own quotient construction,
plurigenus distinctions, and finite-base-change comparison.

There is a precise reason not to use the divisibility assertion in
\cite[Theorem~2.8 and Corollary~2.10, pp.~96--97]{Fujimoto1988Log}
as an obstruction to \eqref{cdp:eq:C.ER.4}. For every $m\ge2$ consider the
diagonal Hopf threefold and the displayed map
\[
 \begin{split}
 H_m&=(\C^3\setminus\{0\})/
       \langle(a,b,c)\longmapsto(qa,q^m b,q^m c)\rangle,\\
 \Phi_m:H_m&\longrightarrow\PP^2,
       \qquad[a,b,c]\longmapsto[a^m:b:c].
 \end{split}                                             \tag{C.ER.17}\label{cdp:eq:C.ER.17}
\]
The action is free and properly discontinuous. A compact annulus
between a small sphere and its image under the contraction meets every
orbit, so $H_m$ is compact and smooth. The three projective coordinates
in \eqref{cdp:eq:C.ER.17} are multiplied by the same factor $q^m$, proving descent.
The map is holomorphic, surjective, and proper. Its fibre over
$[1:Y:Z]$ is parametrized by
$a\mapsto(a,Ya^m,Za^m)$ and is $\C^*/q^{\Z}$; its fibre over
$[0:Y:Z]$ is parametrized by
$s\mapsto(0,sY,sZ)$ and is $\C^*/(q^m)^{\Z}$.
In particular every fibre is connected and elliptic. These are all
the hypotheses in the definition of an elliptic threefold in
\cite[pp.~91--92]{Fujimoto1988Log}.
Near a point of the line $B=(X=0)$, use a chart with $b\ne0$ or
$c\ne0$. In the first chart $X/Y=a^m/b$, where $b$ is a unit.
It follows that
\[
             \Phi_m^*B=m\{a=0\},\qquad B^2=1,
             \qquad\pi_1(\PP^2\setminus B)=0.             \tag{C.ER.18}\label{cdp:eq:C.ER.18}
\]
Thus \eqref{cdp:eq:C.ER.17} is an explicit counterexample to both the asserted
divisibility and the asserted surjection to $\Z/m$ in that theorem
as printed. This conclusion follows from the verified map, rather
than from an inference that an incomplete proof makes a statement
false.

The problematic covering step can also be located exactly. For an
effective divisor $D$ with canonical section $s_D$, the construction
in that proof takes the equation $\xi^m=s_D^m$ in the total space of
$\OO(D)$. Its normalization is the disjoint union of the $m$ global
graphs $\xi=\zeta s_D$, $\zeta^m=1$, each isomorphic to the original
space. It therefore supplies no connected branched covering of the
base. Finally, \eqref{cdp:eq:C.ER.17} is itself a special case of \eqref{cdp:eq:C.ER.4}: take
$V=\PP^2$, $L=\OO(1)$ and $\sigma=X$. A nonzero functional
$\ell$ on a projective line determines its unique vector $v$ with
$\ell(v)=1$; the equation $\sigma=a^m\ell$ becomes
$v=(a^m,b,c)$. Under this identification \eqref{cdp:eq:C.ER.5} is the inverse
of the contraction in \eqref{cdp:eq:C.ER.17}, giving an actual isomorphism of
the two quotient constructions.




\paragraph{Credit and coordinate comparison.}
The cuspidal Weierstrass self-product and its projective small resolutions
are the construction of Namikawa described in
\cite[Proposition~4, pp.~10--11]{Rossi2017}; the ambient Hodge numbers
$h^{1,1}=21$ and $h^{2,1}=3$ are also recorded in
\cite[Remark~5, p.~20]{Rossi2017}.
The coordinate exchange $(x_{\rm Rossi},y_{\rm Rossi})=(y,x)$
and the coefficient $B(t)=t^6-1$ identify the underlying Weierstrass
surface and its fibre self-product with the equations below.
We retain the explicit two blowup ideals, the full support, and its
integral collapse and normalization maps in the proof.

\begin{proposition}[A simply connected irreducible rational-normalized fibre]\label{cdp:st:C.2.UN.8}
\StatementRecord{cdp:st:C.2.UN.8}{proved}{Historical events: Claude Fable 5.2 (Anthropic), effort high, 16 November 2026: carried as background without fault found. GPT-6 Astra (OpenAI), ChatGPT Pro mode, effort Pro, S6\_09\_PAPERS, 17 November 2026: targeted proof check. Also checked in this writing. Pre-continuation writing events: GPT-6 Astra (OpenAI), ChatGPT Pro mode, effort Pro, S6\_10\_PAPERS, 18 November 2026: targeted proof check; used and rechecked in this run. Continuation metadata events: GPT-6 Astra (OpenAI), ChatGPT Pro mode, effort Pro, S6\_10\_PAPERS continuation, 19 November 2026: used and rechecked in this run; status-vocabulary reconciliation; metadata reconciliation only; no new mathematical check. \emph{Pre-continuation recorded proof scope.} Both projective small-resolution charts of uv=a*b*(alpha*a+beta*b), the actual exceptional chain collapse and integral cofibration sequence were re-derived. Adjunction, the unimodular type-II monodromy relations, both generic divisor-restriction maps and the four explicit Neron-Severi classes give all support and ambient groups. The intersection determinant was independently recomputed as minus three. Rossi2017 Proposition4 pp10-11 and Remark5 p20 were opened: the compact cuspidal self-product construction and ambient Hodge numbers are known, while the displayed proof retains the specific support and comparison maps}{the hypotheses, notation and definitions stated here; \cite[Proposition~4, pp.~10--11, and Remark~5, p.~20]{Rossi2017}; \cref{cdp:st:C.2.MH.2}; \cref{cdp:eq:C.CP.5}; \cref{cdp:eq:C.CP.7}; \cref{cdp:eq:C.CP.6}; \cref{cdp:eq:C.CP.4}; \cref{cdp:eq:C.CP.10}}

There is a smooth projective threefold $Y$ with $K_Y\simeq\OO_Y$
and a connected-fibre map $f:Y\to\PP^1$ whose general fibre is a
product of elliptic curves. It has six singular fibres. Each is
reduced and irreducible, has rational normalization resolution, and
satisfies
\[
 \begin{gathered}
 \pi_1(S)=1,\qquad e(S)=6,\\
 (H^i(S,\Z))_{i=0}^4=(\Z,0,\Z^4,0,\Z),\qquad
 (h^i(S,\OO_S))_{i=0}^2=(1,2,1).
 \end{gathered}
 \tag{C.CP.1}\label{cdp:eq:C.CP.1}
\]
The ambient cohomology is
\[
 \begin{gathered}
 (b_i(Y))_{i=0}^6=(1,0,21,8,21,0,1),\qquad e(Y)=36,\\
 (h^i(Y,\OO_Y))_{i=0}^3=(1,0,0,1),\qquad
 h^{1,1}(Y)=21,\quad h^{2,1}(Y)=3,
 \end{gathered}
 \tag{C.CP.2}\label{cdp:eq:C.CP.2}
\]
where the Betti numbers use rational coefficients. In particular
the nonzero local character used in \ref{cdp:st:C.2.MH.2} need not exist after
the global CDP hypotheses are removed, and its Euler conclusion
$e(S)=2$ does not extend to every irreducible rational-normalized
torus degeneration.
\end{proposition}


\begin{proof}
\noindent\textbf{Step 1.} Let $B=\PP^1$, let $L=\OO_B(1)$, and take the Weierstrass
surface $g:R\to B$ with zero section and equation
\[
                  y^2z=x^3+p(t)z^3,
             \qquad p(t)=t^6-1.                 \tag{C.CP.3}\label{cdp:eq:C.CP.3}
\]
Here $x,y$ have the usual transition weights $L^2,L^3$;
thus $p$ is a section of $L^6$. The projective Weierstrass
bundle supplies the compactification. At infinity the transformed
coefficient is $1-t'^6$, so that fibre is smooth. The six
zeros of $p$ are simple. The only singular fibres are the six
cuspidal cubics, and the surface is smooth: in the affine
Weierstrass chart a possible singular point has $x=y=0$, while
the derivative $p'(t)$ there is nonzero. The zero-section chart
is smooth directly from its equation. Adjunction in the
Weierstrass bundle gives
\[
                         K_R=g^*\OO_B(-1).
 \tag{C.CP.4}\label{cdp:eq:C.CP.4}
\]
This formula can also be read from the local canonical form
$dt\wedge dx/y$ and the transition weights just specified.

\noindent\textbf{Step 2.} Form the projective fibre product $X=R\times_B R$. Its only
singularities are the six points where both factors are cusp
points. At one of them use the base coordinate $\tau=p(t)$,
and write
\[
 y_1^2-x_1^3=\tau=y_2^2-x_2^3.
\]
Fix a primitive cube root $\rho$. With
\[
 \begin{gathered}
 u=y_1-y_2,\quad v=y_1+y_2,\quad
 a=x_1-x_2,\quad b=x_1-\rho x_2,\\
 c=x_1-\rho^2x_2=\alpha a+\beta b,
 \qquad \alpha=-\rho,\quad\beta=-\rho^2,
 \end{gathered}
\]
the exact threefold equation is
\[
                         uv=abc.                \tag{C.CP.5}\label{cdp:eq:C.CP.5}
\]
Both $\alpha$ and $\beta$ are nonzero. We resolve \eqref{cdp:eq:C.CP.5}
by two specified blowups. First blow up the ideal $(u,a)$.
In its $u$-chart, $a=ur$ and
\[
                  v=rb(\alpha ur+\beta b),
 \tag{C.CP.6}\label{cdp:eq:C.CP.6}
\]
so $(u,r,b)$ are smooth coordinates. In its $a$-chart put
$u=as$. The equation is $sv=b(\alpha a+\beta b)$; its
only singular point is the origin, an ordinary threefold
node. The strict transform of $(u,b)$ has ideal $(s,b)$
there. Blowing up that ideal gives the two smooth charts
\[
 \begin{array}{c|c}
 \text{coordinates}&\text{blowdown formulas}\\ \hline
 (a,s,h)&b=sh,\quad u=as,\quad
                    v=h(\alpha a+\beta sh),\\
 (b,j,v)&s=bj,\quad a=(jv-\beta b)/\alpha,\quad u=abj.
 \end{array}                                    \tag{C.CP.7}\label{cdp:eq:C.CP.7}
\]
The inverse image of the original singular point is a chain
of two projective lines. In the first chart of \eqref{cdp:eq:C.CP.7}
the old and new curves have equations $(a,h)=(0,0)$
and $(a,s)=(0,0)$; they meet once transversely. Charts
\eqref{cdp:eq:C.CP.6}--\eqref{cdp:eq:C.CP.7} complete each to a projective line.
No exceptional divisor is introduced. Outside the origin
the respective ideals are Cartier, so these maps are
isomorphisms there.

\noindent\textbf{Step 3.} The blowups globalize without an analytic gluing choice.
The ideal $(u,a)$ is the ideal of the Weil diagonal
$R\hookrightarrow X$. The ideal $(u,b)$ is the Weil
graph of the global automorphism
$(x,y,t)\mapsto(\rho x,y,t)$ of $R$; after the first
blowup take its strict transform. Each of these divisors
is Cartier off the six specified singular points.
Their projective blowups therefore give a smooth
projective threefold $Y$ and a proper small map
$\pi:Y\to X$, with precisely the local charts above.
Set $f=g\pi$, using the common projection of the fibre
product. It agrees with the original product of elliptic
curves away from the six values. Its fibres are connected,
because the product fibres and every fibre of the two
blowdowns are connected.

\noindent\textbf{Step 4.} The original central support at each of the six values
is the product $S_0=C\times C$ of a cuspidal cubic with
itself. It is reduced and irreducible. Since $\pi$ has
no exceptional divisors, the full fibre upstairs is its
single strict transform, still reduced. Its normalization
resolution is rational: the function field is unchanged
from that of the finite birational parametrization
$\PP^1\times\PP^1\to S_0$. This assertion is birational;
we do not discard the exceptional curves of the actual
normalization resolution.

\noindent\textbf{Step 5.} Topologically a cuspidal cubic is $S^2$. Its finite
normalization is bijective, so compactness makes it a
homeomorphism. Hence $S_0$ is homeomorphic to
$S^2\times S^2$. The restriction $S\to S_0$ is an
isomorphism off the cusp-pair point and replaces that
point by the chain $E$ of two projective lines. The
quotient $S/E\to S_0$ is a continuous bijection from a
compact space to a Hausdorff space, hence a homeomorphism.
Choose a finite triangulation of $(S,E)$. Its cofibration
sequence identifies the reduced homology of $S/E$ with
the relative homology of $(S,E)$. Since
$H_2(E,\Z)=\Z^2$ and $E$ is simply connected, the
long exact sequence gives
\[
 \begin{gathered}
 H_4(S,\Z)=\Z,\qquad H_3(S,\Z)=H_1(S,\Z)=0,\\
 0\longrightarrow\Z^2\longrightarrow H_2(S,\Z)
                 \longrightarrow\Z^2\longrightarrow0.
 \end{gathered}                                  \tag{C.CP.8}\label{cdp:eq:C.CP.8}
\]
The last sequence splits as a sequence of abelian groups.
Van Kampen for the same cofibration shows
$\pi_1(S/E)=\pi_1(S)/\langle\!\langle\pi_1(E)\rangle\!\rangle$.
Thus $\pi_1(S)=1$. Universal coefficients now give the
integral cohomology in \eqref{cdp:eq:C.CP.1}, and its Euler number
is $1+4+1=6$.

\noindent\textbf{Step 6.} The fibre product is a local complete intersection in
$R\times R$, cut out by the inverse image of the base
diagonal. Adjunction and \eqref{cdp:eq:C.CP.4} give
\[
 \omega_X=p_1^*K_R\otimes p_2^*K_R\otimes f^*T_B
                         \simeq\OO_X.
 \tag{C.CP.9}\label{cdp:eq:C.CP.9}
\]
Its generator pulls back to a nowhere-vanishing canonical
form on $Y$ away from the small exceptional curves.
Its coefficients and the coefficients of its inverse
extend across that codimension-two set by Hartogs'
theorem in a local canonical frame. Hence $K_Y=\OO_Y$.
Adjunction makes $\omega_S=\OO_S$. Proper flatness over
the smooth base preserves the holomorphic Euler
characteristic, so $\chi(\OO_S)=\chi(\OO_{E_t^2})=0$.
Connectedness and Serre duality then give
$(h^0,h^1,h^2)(\OO_S)=(1,2,1)$.

\noindent\textbf{Step 7.} We include the ambient groups, with their comparison
maps. A type-II local elliptic monodromy has characteristic
polynomial $x^2-x+1$, so $T-I$ is an integral isomorphism.
This also follows directly from the cyclic sixth-root
trivialization of $y^2=x^3+\tau$. The smooth part of
the fibration has fundamental group generated by its
fibre lattice and base meridians. Its inclusion into
$Y$ is surjective on fundamental groups, by perturbing
loops off the central divisor. The pair of zero sections
extends over every special value and kills each base
meridian in $Y$. Conjugation by one such meridian
therefore kills the image of $T-I$ on each of the two
elliptic lattices. Since those two maps are unimodular,
the whole fibre lattice is killed. Thus $\pi_1(Y)=1$.
Projectivity and $K_Y=\OO_Y$ give
$(h^0,h^1,h^2,h^3)(\OO_Y)=(1,0,0,1)$ by Hodge theory
and Serre duality. The identical argument for $R$ gives
$b_1(R)=0$. Its six cuspidal fibres give $e(R)=12$,
hence $b_2(R)=10$. Moreover $K_R=g^*\OO(-1)$ makes
every positive plurigenus zero. The surface classification
therefore makes $R$ rational, so its Picard group has
rank ten.

\noindent\textbf{Step 8.} Let $K=\C(t)$ and let $E_K$ be the generic elliptic
curve of $R$. Restriction of divisors gives a surjection
$\Pic(R)\to\Pic(E_K)$: close a generic divisor in the
smooth surface. Its kernel is generated, up to linear
equivalence, by the full fibre, because every fibre is
irreducible. The zero section splits the degree map on
$\Pic(E_K)$, whose degree-zero part is $E_K(K)$.
Consequently $E_K(K)$ has rank eight.

\noindent\textbf{Step 9.} Apply the same divisor restriction to $Y$ and its
generic fibre $A_K=E_K\times E_K$. It is surjective
and its vertical kernel has rank one, since all six
singular fibres of $Y$ are irreducible. The degree-zero
Picard group of $A_K$ is its dual abelian variety's
$K$-points, which have rank sixteen under the product
principal polarization. Its N\'eron--Severi group has
rank four. To verify this last number explicitly, use
the two factors, the diagonal, and the graph of the
automorphism $\rho$. Their intersection matrix is
\[
 \begin{pmatrix}
 0&1&1&1\\1&0&1&1\\1&1&0&3\\1&1&3&0
 \end{pmatrix}.                                \tag{C.CP.10}\label{cdp:eq:C.CP.10}
\]
Its determinant is $-3$. The entry three is the degree of
$1-\rho$, whose kernel on the triangular elliptic
lattice has order three. Here is the exact comparison
for the upper bound on the generic N\'eron--Severi rank.
Over $K'=K(u)$ with $u^6=p(t)$, the algebraic
coordinate substitution $x=u^2X$, $y=u^3Y$ identifies
$E_{K'}$ with the constant elliptic curve
$E_0:\ Y^2=X^3+1$. Base change to an algebraic closure
identifies the geometric N\'eron--Severi group with that
of $(E_0\times E_0)_{\overline K}$. For the constant
complex abelian surface, the first Chern class embeds
this group into the integral $(1,1)$ classes and
$h^{1,1}(E_0^2)=4$. Algebraically closed extension
preserves the N\'eron--Severi group: the Picard scheme
and its identity component commute with field extension,
and each of its connected components over $\C$ remains
geometrically connected. The four classes
in \eqref{cdp:eq:C.CP.10} are already defined over $K$, so its rank
over $K$ is exactly four. Thus
$\operatorname{rank}\Pic(Y)=1+16+4=21$.
Since $h^1(\OO_Y)=h^2(\OO_Y)=0$, the exponential
sequence identifies this rank with $b_2(Y)$.

\noindent\textbf{Step 10.} Euler localization now gives $e(Y)=6\cdot6=36$.
Poincar\'e duality and $b_1=0$ yield $b_3=8$;
the remaining Betti numbers and Hodge numbers follow
as displayed in \eqref{cdp:eq:C.CP.2}. These rational Betti assertions
make no assertion about integral torsion in degrees
not already computed.

\noindent\textbf{Step 11.} Finally a sufficiently small compact fibre tube retracts
onto $S$. Its fundamental group is therefore trivial.
There is no nontrivial finite local system there, and
no nonzero integral degree-one class to restrict to
the nearby torus. The example consequently disproves
the proposed unrestricted replacement for the character
construction itself, with its actual comparison map.
The exact coordinate identities, intersection matrix,
and all displayed integer ranks are checked in
\path{code/root_cusp_product.py} and its saved output.
\end{proof}



The compact construction in this proof is the Namikawa
cuspidal self-product. Rossi~\cite[\S3, Proposition~4]{Rossi2017}
gives the same two graph blowups and exceptional chains;
his Remark~5 gives $h^{1,1}=21$ and $h^{2,1}=3$.
We opened and checked those passages. The coordinate,
integral fibre-cohomology and Picard calculations above
derive the precise local-character obstruction used here.





\begin{proposition}[A projective zero-Euler filling with an essential Kummer component]\label{cdp:st:C.2.UN.13}
\StatementRecord{cdp:st:C.2.UN.13}{proved}{Historical events: Claude Fable 5.2 (Anthropic), effort high, 16 November 2026: carried as background without fault found. GPT-6 Astra (OpenAI), ChatGPT Pro mode, effort Pro, S6\_09\_PAPERS, 17 November 2026: targeted proof check. Also checked in this writing. Pre-continuation writing events: GPT-6 Astra (OpenAI), ChatGPT Pro mode, effort Pro, S6\_10\_PAPERS, 18 November 2026: targeted proof check; used and rechecked in this run. Continuation metadata events: GPT-6 Astra (OpenAI), ChatGPT Pro mode, effort Pro, S6\_10\_PAPERS continuation, 19 November 2026: used and rechecked in this run; status-vocabulary reconciliation; metadata reconciliation only; no new mathematical check. \emph{Pre-continuation recorded proof scope.} The invariant quadratic cone blowup gives smooth coordinates (R,u,v) and the exact base equation t=R*v squared. The descended two-form and finite trace identify the Kummer component. The actual collapse pairs, exceptional Gysin sections and coherent formal-function filtration give every ambient and support group. Forty specified negative-two blocks preserve both K3 components, with all support multiplicities and coherent pushforwards checked. The finite normal map from A times P1 is exactly the normalization after quadratic base change. Huybrechts's Kummer construction was opened at its stated locus}{the hypotheses, notation and definitions stated here; \cref{cdp:eq:C.KM.4}; \cref{cdp:eq:C.KM.5}; \cite[Chapter~1, Example~1.3(iii)]{HuybrechtsK3Draft}; \cref{cdp:st:C.2.UN.4}; \cref{cdp:eq:C.KM.S2}; \cref{cdp:st:C.2.UN.3}}

Fix an abelian surface $A$. There is a smooth projective threefold
$W$ with a proper map $f_W:W\to\PP^1$ whose fibres over
$\C^*$ are isomorphic to $A$. It has precisely two special fibres,
over $0$ and $\infty$. Each has full divisor
\[
               2P+\sum_{i=1}^{16}E_i,\qquad
 P=\operatorname{Km}(A),\quad E_i\simeq\PP^2.       \tag{C.KM.1}\label{cdp:eq:C.KM.1}
\]
Here $P$ is the smooth Kummer K3 surface, the $E_i$ are disjoint,
and $P\cap E_i$ is a projective line on $E_i$ and a curve of
self-intersection $-2$ on $P$. The reduced divisor is simple
normal crossing and has Euler number $40$.

There is a composition of smooth vertical blowups $Y\to W$,
disjoint from both K3 components, for which both special reduced
fibres have Euler number zero. Each has $57$ irreducible
components, including its original unchanged Kummer surface.
The maps to the base and the nearby tori are retained. Their
rational Betti numbers and coherent dimensions are
\[
 \begin{array}{c|c|c|c}
 & (b_i)_{i=0}^6 & e & (h^i(\OO))_{i=0}^3\\ \hline
 W&(1,0,39,0,39,0,1)&80&(1,0,1,0)\\
 Y&(1,0,119,240,119,0,1)&0&(1,0,1,0).
 \end{array}                                             \tag{C.KM.2}\label{cdp:eq:C.KM.2}
\]
For either special value, let $S_W$ and $S_Y$ denote its
reduced supports in these two models. Their cohomology is
\[
 \begin{gathered}
 (H^i(S_W,\Q))_{i=0}^4=(\Q,0,\Q^{22},0,\Q^{17}),\\
 (H^i(S_Y,\Q))_{i=0}^4=(\Q,0,\Q^{62},\Q^{120},\Q^{57}),\\
 (h^i(S_W,\OO_{S_W}))_{i=0}^2
   =(h^i(S_Y,\OO_{S_Y}))_{i=0}^2=(1,0,1).
 \end{gathered}                                          \tag{C.KM.S1}\label{cdp:eq:C.KM.S1}
\]
Neither K3 component can be contracted by a bimeromorphic map
between smooth fillings. For comparison, after the quadratic
base change the normalization of the original quotient is
exactly $A\times\PP^1$, and the corresponding base changes
of $W$ and $Y$ are birational modifications of that product.
Thus the asserted persistence concerns bimeromorphic comparison
over the original base; its relation to base change is given
by this specified map.
\end{proposition}


\begin{proof}
\noindent\textbf{Step 1.} On $A\times\PP^1_s$ take the involution
\[
       \iota(a,s)=(-a,-s),\qquad
 X=(A\times\PP^1_s)/\langle\iota\rangle,
       \qquad t=s^2.                                  \tag{C.KM.3}\label{cdp:eq:C.KM.3}
\]
The quotient is projective and normal. Its fixed points are
exactly $A[2]\times\{0,\infty\}$, so there are $32$ of them.
Away from them the action is free. In local coordinates centred
at any fixed point, including the coordinate $s^{-1}$ at
infinity, the action is scalar multiplication by $-1$ on
$\C^3$. Thus the only singularities of $X$ are those $32$
quotient points.

\noindent\textbf{Step 2.} Blow up these points in $X$, and call the resulting map
$\pi:W\to X$. We verify its smoothness and its exact relation
to the quotient before using it. Write the coordinates of a
covering chart as $(a,b,s)$. The invariant local ring is
generated by their six quadratic monomials. In the blowup
chart corresponding to $a^2$, put
\[
             a=r,\qquad b=ru,\qquad s=rv,
             \qquad R=r^2.                            \tag{C.KM.4}\label{cdp:eq:C.KM.4}
\]
All invariant generators pull back as
\[
 (a^2,ab,as,b^2,bs,s^2)
           =(R,Ru,Rv,Ru^2,Ruv,Rv^2).                  \tag{C.KM.5}\label{cdp:eq:C.KM.5}
\]
Consequently $(R,u,v)$ are smooth coordinates on the blowup.
The analogous charts for $b^2$ and $s^2$ cover it: the
exceptional projectivized invariant cone is the quadratic
Veronese image of $\PP^2$, on which at least one of the
three squared coordinates is nonzero. Equivalently, blow up
the origin of $\C^3$ first. The lifted involution fixes its
exceptional $\PP^2$ and sends its normal coordinate $r$ to
$-r$; its quotient has precisely \eqref{cdp:eq:C.KM.4}--\eqref{cdp:eq:C.KM.5}.
The transitions for $R=r^2$ identify its normal line with
$\OO_{\PP^2}(-2)$. This proves that $W$ is smooth and
projective, with exceptional divisors $E_i\simeq\PP^2$.

\noindent\textbf{Step 3.} The base map in \eqref{cdp:eq:C.KM.4} is exactly
\[
                             t=R v^2.                 \tag{C.KM.6}\label{cdp:eq:C.KM.6}
\]
Here $E_i=(R=0)$ and the strict central component is
$P=(v=0)$. In the chart where $s$ is the exceptional
coordinate the map is $t=R$ and only $E_i$ occurs.
These charts prove the multiplicities and the transverse
intersection in \eqref{cdp:eq:C.KM.1}. They also show that
$N_{E_i/W}|_{P\cap E_i}=\OO_{\PP^1}(-2)$, which is
the normal line of that intersection in $P$.
The same calculation applies over infinity. Off those two
values, choose either square root $s$ of $t$; the map
$a\mapsto[a,s]$ identifies the fibre with $A$, and switching
the square root composes this identification with $a\mapsto-a$.
In particular the map has connected fibres and its monodromy
on $H^1(A,\Z)$ is $-I_4$.

\noindent\textbf{Step 4.} The restriction $P\to A/\{\pm1\}$ resolves exactly the
sixteen quotient surface nodes. It is the classical Kummer
construction; see \cite[Chapter~1, Example~1.3(iii)]{HuybrechtsK3Draft}.
Here the two required properties can also be checked in the
same charts. A translation-invariant two-form on $A$ is
invariant under $-1$. In a surface chart $a=r$, $b=ru$,
its pullback is
\[
                da\wedge db=r\,dr\wedge du
                            =\tfrac12\,dR\wedge du.   \tag{C.KM.7}\label{cdp:eq:C.KM.7}
\]
It therefore descends to a nowhere-zero form on $P$.
The finite covering $\operatorname{Bl}_{A[2]}A\to P$
injects $H^1(P,\OO_P)$ into the invariant part of
$H^1(A,\OO_A)$, by the averaging trace. This invariant
part is zero because $-1$ acts as $-1$ on $H^1(A,\OO_A)$.
Thus $K_P=\OO_P$ and $h^1(P,\OO_P)=0$, verifying that
$P$ is a K3 surface.

\noindent\textbf{Step 5.} The finite quotient map identifies rational cohomology with
invariants. Explicitly, its higher direct images of the
constant rational sheaf vanish by proper base change and
the finiteness of its fibres; averaging identifies the
constant sheaf downstairs with the invariant direct image.
Applied to $A/\{\pm1\}$, this gives Betti numbers
$(1,0,6,0,1)$ and Euler number eight. Resolving each of its
sixteen points replaces that point by $\PP^1$, so
$e(P)=8+16=24$. Hence the actual inclusion--exclusion
comparison in \eqref{cdp:eq:C.KM.1} gives
\[
          e\left(P\cup\bigcup_{i=1}^{16}E_i\right)
                    =24+16\cdot3-16\cdot2=40.          \tag{C.KM.8}\label{cdp:eq:C.KM.8}
\]
Mayer--Vietoris gives the support groups as well. The restriction
map from the direct sum of $H^2(P,\Q)$ and the sixteen
$H^2(E_i,\Q)$ onto the sixteen $H^2(P\cap E_i,\Q)$ is
surjective: the hyperplane class on each $E_i$ restricts to
the degree-one class of its intersection line. Its kernel
has dimension $22+16-16=22$, and the same exact sequence
gives $H^1(S_W,\Q)=H^3(S_W,\Q)=0$ and
$H^4(S_W,\Q)=\Q^{17}$. The coherent sequence
\[
 0\longrightarrow\OO_{S_W}\longrightarrow
 \OO_P\oplus\bigoplus_i\OO_{E_i}
 \longrightarrow\bigoplus_i\OO_{P\cap E_i}
 \longrightarrow0
\]
has a surjective difference map on $H^0$. The lines and planes
have no positive coherent cohomology, so it gives
$(h^0,h^1,h^2)(\OO_{S_W})=(1,0,1)$.

\noindent\textbf{Step 6.} We give the ambient cohomology maps as well. For the finite
quotient $A\times\PP^1\to X$, the same invariant-sheaf
argument gives
\[
 H^k(X,\Q)\simeq H^k(A\times\PP^1,\Q)^{\iota},
 \qquad (b_k(X))_{k=0}^6=(1,0,7,0,7,0,1).              \tag{C.KM.9}\label{cdp:eq:C.KM.9}
\]
Indeed $-1$ acts by $(-1)^j$ on $H^j(A)$, while the
base involution acts trivially on both nonzero cohomology
groups of $\PP^1$.

\noindent\textbf{Step 7.} Let $j_i:E_i\hookrightarrow W$ and let
$h_i=c_1(\OO_{E_i}(1))$. Collapsing the $32$ exceptional
divisors separately gives exactly $X$. The relative
cohomology sequence of this quotient identifies the
kernel of restriction to the $E_i$ with the pullback
of $H^k(X,\Q)$. The self-intersection formula gives
\[
 j_i^*j_{i*}(1)=-2h_i,
       \qquad j_i^*j_{i*}(h_i)=-2h_i^2.                 \tag{C.KM.10}\label{cdp:eq:C.KM.10}
\]
Thus restriction is surjective in degrees two and four,
with explicit rational sections $-j_{i*}(1)/2$ and
$-j_{i*}(h_i)/2$. In these two degrees pullback is
injective because the preceding odd cohomology of the
exceptional divisors vanishes. The same relative
sequence then leaves the other positive-degree groups
unchanged. Each exceptional point consequently adds one
class in degrees two and four and none in any other
degree. This proves the row for $W$ in \eqref{cdp:eq:C.KM.2}.

\noindent\textbf{Step 8.} For coherent cohomology the finite quotient and averaging
give
$H^k(X,\OO_X)=H^k(A\times\PP^1,\OO)^{\iota}$,
of dimensions $(1,0,1,0)$. The local resolution just
used has $R^q\pi_*\OO_W=0$ for $q>0$. To verify this
without suppressing the quotient singularity, filter the
infinitesimal neighbourhoods of $E_i$ by its ideal.
Their successive quotients are $\OO_{\PP^2}(2r)$,
$r\ge0$, whose positive-degree coherent cohomology
vanishes. Formal functions therefore makes the completed
stalks of those higher direct images zero, and coherence
makes the stalks themselves zero. Normality gives
$\pi_*\OO_W=\OO_X$. Leray proves the coherent row for
$W$.

\noindent\textbf{Step 9.} We now perform exactly twenty of the negative-two
operations of \ref{cdp:st:C.2.UN.4} over each of $0$ and $\infty$.
Every operation starts at a fresh smooth point on the
strict transform of one fixed $E_i$, away from its
intersection with $P$ and from all other components.
The first blowup creates a new $\PP^2$, and the second
blows up a smooth plane quartic of genus three inside
it, transverse to its old intersection line. These
centres and their entire modifications are disjoint
from $P$. The exact charts of \ref{cdp:st:C.2.UN.4} retain simple
normal crossings and show that both new components
have multiplicity one. Each operation changes the
reduced-fibre Euler number by $-2$ and adds two
components. Thus both final supports have Euler number
$40-20\cdot2=0$ and $17+20\cdot2=57$ components,
with their unchanged multiplicity-two K3 surfaces.

\noindent\textbf{Step 10.} We retain the support comparison under each of these blowups,
including the four boundary intersections of the quartic.
Let $\sigma:\widehat W\to W'$ be one blowup, with smooth
centre $C$ of ambient codimension $c$, let $S$ be the current
reduced support, and put $S'=\sigma^{-1}(S)_{\rm red}$ and
$E=\PP(N_{C/W'})$. Collapsing $E$ in $S'$ and $C$ in $S$
gives the actual homeomorphism $S'/E\simeq S/C$.
Hence pullback identifies the relative cohomology of the
pairs $(S,C)$ and $(S',E)$. With $p_E:E\to C$ and
$\xi=c_1(\OO_E(1))$, the projective-bundle decomposition is
\[
 H^k(E,\Q)=\bigoplus_{i=0}^{c-1}
                \xi^i p_E^*H^{k-2i}(C,\Q).
\]
For $i\ge1$, the restriction to $S'$ of the ambient Gysin
class $j_*(\xi^{i-1}p_E^*\alpha)$ restricts further to $E$
as $-\xi^i p_E^*\alpha$, because $N_{E/\widehat W}=\OO_E(-1)$.
These classes split every positive projective-bundle summand.
The two relative exact sequences now identify the remaining
summand with the pullback of $H^k(S,\Q)$ and give
\[
 H^k(S',\Q)\simeq H^k(S,\Q)\oplus
                \bigoplus_{i=1}^{c-1}H^{k-2i}(C,\Q).
                                                        \tag{C.KM.S2}\label{cdp:eq:C.KM.S2}
\]
This argument uses the full inverse image and does not require
the centre to avoid the other support components at every point.
A point step therefore adds one class in degrees two and four.
The quartic step adds ranks one, six and one in degrees two,
three and four. After twenty blocks, \eqref{cdp:eq:C.KM.S2} gives exactly
the second support row of \eqref{cdp:eq:C.KM.S1}.

\noindent\textbf{Step 11.} Each chosen centre is generically contained in just one
reduced component, of multiplicity one in the reduced support.
Consequently the equality of Cartier divisors
$\sigma^*S=S'$ holds, also at the four crossing points as
the explicit equations show. Apply derived pushforward to
$0\to\OO_{\widehat W}(-S')\to\OO_{\widehat W}\to
\OO_{S'}\to0$. The projection formula and
$R\sigma_*\OO_{\widehat W}=\OO_{W'}$ give
\[
                         R\sigma_*\OO_{S'}=\OO_S.       \tag{C.KM.S3}\label{cdp:eq:C.KM.S3}
\]
Thus all support coherent cohomology groups are unchanged,
proving the remaining assertion of \eqref{cdp:eq:C.KM.S1}.

\noindent\textbf{Step 12.} The pullback and exceptional Gysin maps in \ref{cdp:st:C.2.UN.4}
give cohomology increments $(\Delta b_2,\Delta b_3,
\Delta b_4)=(2,6,2)$ per operation. Forty operations
give the second Betti row in \eqref{cdp:eq:C.KM.2}. Smooth-centre
blowups preserve coherent cohomology, by their direct
image $\OO$ and vanishing higher direct images, so
the second coherent row is unchanged. This also checks
$e(Y)=0$ directly from its displayed Betti numbers.

\noindent\textbf{Step 13.} Finally, $P$ is not uniruled. A dominating ruled surface
would pull its nonzero holomorphic two-form back to a
nonzero two-form, whereas a ruled surface has no such
form. The inverse-map resolution in \ref{cdp:st:C.2.UN.3} therefore
prevents its contraction by any bimeromorphic comparison
between smooth threefold fillings. The finite-base-change
comparison is different and explicit: the map
\[
 A\times\PP^1_s\longrightarrow
 X\times_{\PP^1_t}\PP^1_s,
         \qquad(a,s)\longmapsto([a,s],s)               \tag{C.KM.11}\label{cdp:eq:C.KM.11}
\]
is finite and generically one-to-one. Its source is
normal, so it is the normalization of the reduced
base change. The normalized base changes of $W$ and
$Y$ factor properly and birationally through this
normalization. This specifies exactly how the Kummer
component is related to a torus after that base change.
The construction is based on the classical Kummer
quotient. Its chart,
trace, Euler and cohomology calculations are retained
in \path{code/E_kummer_filling.py} and its output.
\end{proof}






\begin{proposition}[Opposite logarithmic data produce a projective essential component]\label{cdp:st:C.2.UN.18}
\StatementRecord{cdp:st:C.2.UN.18}{proved}{Historical events: Claude Fable 5.2 (Anthropic), effort high, 16 November 2026: carried as background without fault found. GPT-6 Astra (OpenAI), ChatGPT Pro mode, effort Pro, S6\_09\_PAPERS, 17 November 2026: targeted proof check. Also checked in this writing. Pre-continuation writing events: GPT-6 Astra (OpenAI), ChatGPT Pro mode, effort Pro, S6\_10\_PAPERS, 18 November 2026: targeted proof check. Continuation metadata events: GPT-6 Astra (OpenAI), ChatGPT Pro mode, effort Pro, S6\_10\_PAPERS continuation, 19 November 2026: used and rechecked in this run; status-vocabulary reconciliation; metadata reconciliation only; no new mathematical check. \emph{Pre-continuation recorded proof scope.} Both root charts, their signed deck actions, the proper finite quotients and the exact punctured product were re-derived. The cyclic curve cover and nonzero Jacobians prove the finite etale projective product cover. The canonical form, floor-function plurigenera, primitive Mayer-Vietoris maps, all integral Betti ranks, coherent direct images and birational persistence were checked with the actual boundary and cap maps}{the hypotheses, notation and definitions stated here; \cref{cdp:eq:C.PE.5}; \cref{cdp:eq:C.PE.6}; \cref{cdp:eq:C.PE.7}; \cref{cdp:eq:C.PE.12}; \cref{cdp:eq:C.PE.8}; \cref{cdp:eq:C.PE.9}; \cref{cdp:eq:C.PE.16}; \cref{cdp:eq:C.PE.18}; \cref{cdp:st:C.2.UN.3}; \cref{cdp:st:C.2.UN.6}; \cite[Example~3.4(1)]{Fujimoto1988Log}}

Put $q=e^{-2\pi}$ and $E=\C^*/q^{\Z}$. Choose distinct
points $p_+,p_-\in E$ whose difference is torsion.
For every $m\ge2$ the fixed punctured family
$(E\times E)\times\Delta^*$ has a smooth proper filling
with full central fibre
\[
                         P_m+mQ_++mQ_-.                 \tag{C.PE.1}\label{cdp:eq:C.PE.1}
\]
Here $Q_+$ and $Q_-$ are disjoint primary Hopf surfaces,
$P_m$ is a smooth projective surface of Kodaira dimension
one, and $J_\pm=P_m\cap Q_\pm$ are disjoint smooth
elliptic curves. The support is a simple normal crossing
surface of Euler number zero. The map $P_m\to E$ is an
elliptic fibration with exactly the two multiple fibres
$mJ_+$ and $mJ_-$, and
\[
 \begin{gathered}
 K_{P_m}\simeq\OO_{P_m}((m-1)(J_++J_-)),\qquad
 H_1(P_m,\Z)\simeq\Z^4,\\
 h^0(P_m,K_{P_m}^{\otimes r})=
 \begin{cases}1,&d_r=0,\\d_r,&d_r>0,\end{cases}
 \qquad d_r=2\left\lfloor\frac{r(m-1)}m\right\rfloor.
 \end{gathered}                                         \tag{C.PE.2}\label{cdp:eq:C.PE.2}
\]
The surface $P_m$ admits a finite \'etale cover by a
product of two projective curves, one elliptic and the
other of genus at least two. All three components contain
no rational curve and survive birationally in every
bimeromorphic comparison between smooth fillings. The
multiplicity-one component $P_m$ also survives birationally
after every finite base change and resolution. The
fillings for different $m$ are not bimeromorphic over the
disc.

The construction extends to a smooth compact threefold
$Y_m\to\PP^1$ with just this special fibre. All integral
cohomology groups below are free. Writing $S_m$ for its
reduced special support, their ranks and coherent
cohomology dimensions are
\[
 \begin{array}{c|c|c}
 &\operatorname{rank}H^k(-,\Z)&h^k(\OO)\\ \hline
 P_m&(1,4,6,4,1)&(1,2,1)\\
 Q_\pm&(1,1,0,1,1)&(1,1,0)\\
 S_m&(1,3,6,7,3)&(1,2,1)\\
 Y_m&(1,3,7,10,7,3,1)&(1,2,1,0).
 \end{array}                                           \tag{C.PE.3}\label{cdp:eq:C.PE.3}
\]
\end{proposition}


\begin{proof}
\noindent\textbf{Step 1.} Let $V$ be the blowup of $E\times\Delta$ at
$(p_+,0)$ and $(p_-,0)$, and let $R_+,R_-$ be its
disjoint exceptional curves. Put
\[
 D=R_+-R_-,\qquad L=\OO_V(D),\qquad
 \operatorname{div}(\sigma)=D,                         \tag{C.PE.4}\label{cdp:eq:C.PE.4}
\]
where $\sigma$ is the canonical meromorphic section of
$L$. We define a smooth covering space $W_m$ using two
actual charts. On $V\setminus R_-$ use
\[
 \sigma=a^m\ell,\qquad \ell\in L^\times,\quad a\in\C;
                                                               \tag{C.PE.5}\label{cdp:eq:C.PE.5}
\]
on $V\setminus R_+$ use
\[
 \sigma^{-1}=(a')^m\ell',\qquad
 \ell'\in(L^{-1})^\times,\quad a'\in\C.                 \tag{C.PE.6}\label{cdp:eq:C.PE.6}
\]
On the common open set both $a,a'$ are nonzero, and the
gluing is $a'=a^{-1}$, $\ell'=\ell^{-1}$. Thus $a$
is globally a $\PP^1$-valued coordinate. The generator
of the free $\Z$-action is
\[
 (\ell,a)\longmapsto(q^m\ell,q^{-1}a),\qquad
 (\ell',a')\longmapsto(q^{-m}\ell',qa').                 \tag{C.PE.7}\label{cdp:eq:C.PE.7}
\]
Near $R_+$, a frame of $L$ makes \eqref{cdp:eq:C.PE.5} the equation
$x=a^m b$, $b\ne0$. Near $R_-$, a frame of $L^{-1}$
makes \eqref{cdp:eq:C.PE.6} the equation $x=(a')^m b'$, $b'\ne0$.
Each eliminates $x$ and proves smoothness. A Hermitian
norm of the globally defined nonzero vector $\ell$ is
multiplied by $q^m$ in \eqref{cdp:eq:C.PE.7}; it proves proper
discontinuity and the absence of stabilizers. The
quotient $U_m=W_m/\Z$ is therefore a smooth threefold.

\noindent\textbf{Step 2.} The local proper maps are explicit. In the positive
chart write $x=z^m$, $b=w^m$, $a=z/w$ and quotient
by $(z,w)\mapsto(\zeta z,\zeta w)$. The generator
\eqref{cdp:eq:C.PE.7} lifts to $w\mapsto qw$. In the negative
chart use $x=z^m$, $b'=(w')^m$, $a'=z/w'$ and the
same diagonal finite action; the generator lifts to
$w'\mapsto q^{-1}w'$. In each case the local map to
$V$ is a free finite quotient of
\[
 \Delta_z\times E\times\Delta_y
   \longrightarrow\Delta_x\times\Delta_y,\qquad
 (z,[w],y)\longmapsto(z^m,y).                           \tag{C.PE.8}\label{cdp:eq:C.PE.8}
\]
Multiplication by a nontrivial $m$th root of unity is a
nontrivial translation of $E$, so the finite action is
free. The map before this quotient is proper; hence
the resulting $\rho:U_m\to V$ is proper. Away from
$R_+\cup R_-$ the actual product identification is
\[
 (v,[c])\longmapsto[(v,\ell=\sigma(v)c^m,a=c^{-1})].     \tag{C.PE.9}\label{cdp:eq:C.PE.9}
\]
It gives exactly the stated fixed punctured product.

\noindent\textbf{Step 3.} Write $B\simeq E$ for the strict central curve of $V$.
The base function has divisor $B+R_++R_-$. At the two
intersections it has equation $t=xy$, so its pullbacks
are respectively
\[
                   t=a^mby,\qquad t=(a')^m b'y.         \tag{C.PE.10}\label{cdp:eq:C.PE.10}
\]
These are transverse products with multiplicities
$m,1$. They prove \eqref{cdp:eq:C.PE.1}, smoothness of $P_m$ and
the stated intersections. Over $R_+$ the component is
$L|_{R_+}^{\times}/q^{m\Z}$; over $R_-$ it is
$(L^{-1}|_{R_-})^{\times}/q^{-m\Z}$. Both line bundles
are $\OO_{\PP^1}(-1)$. Sending a nonzero tautological
vector to the same vector in $\C^2\setminus0$ identifies
both components with
$(\C^2\setminus0)/q^{m\Z}$, a primary Hopf surface.
The two reduced multiple fibres $J_\pm$ are each
$\C^*/q^{m\Z}$.

\noindent\textbf{Step 4.} We construct the finite projective cover of $P_m$.
The restriction of $L$ to $B$ is
$L_B=\OO_E(p_+-p_-)$, a torsion line. Choose a torsion
line $N$ with $N^{\otimes m}\simeq L_B$; multiplication
by $m$ is onto on the torsion subgroup of $\Pic^0(E)$.
Normalize the cyclic cover described in the projective
bundle $\PP(\OO_E\oplus N)$ by
\[
                             y^m=\sigma|_B.              \tag{C.PE.11}\label{cdp:eq:C.PE.11}
\]
Here $y$ is the tautological meromorphic $N$-coordinate.
At $p_+$ this is $y^m=z$, and at $p_-$ its reciprocal
is $(y^{-1})^m=z$, after local frame choices. The
normalization $C$ is a connected smooth projective
curve: the order-one zero of $\sigma$ prevents the
cyclic equation from splitting, and the local maps
are the stated fully ramified ones. Riemann--Hurwitz
gives $2g(C)-2=2(m-1)$, so $g(C)=m$.

\noindent\textbf{Step 5.} Pull back a finite \'etale cover of $E$ trivializing
$N$, and choose a connected component of its fibre
product with $C$. This gives a finite \'etale cover
$C'\to C$ and a map $\tau:C'\to E$ such that
$\tau^*N$ has a nowhere-zero frame $n$. Write
$y=y_0n$. There is a holomorphic map
\[
 \begin{split}
 \Phi:C'\times E&\longrightarrow P_m,\\
 (z,[u])&\longmapsto
       [\ell=u^m n(z)^m,\ a=y_0(z)/u].
 \end{split}                                             \tag{C.PE.12}\label{cdp:eq:C.PE.12}
\]
At a pole of $y_0$ this formula is interpreted in
\eqref{cdp:eq:C.PE.6}, so $a'=u/y_0$ and $\ell'=u^{-m}n^{-m}$.
It is well defined under $u\mapsto qu$ by \eqref{cdp:eq:C.PE.7}.
It is onto and has finite fibres: away from the two
branch values its restriction to each elliptic fibre
is an isomorphism; at a branch value it is the finite
map $[u]\mapsto[u^m]$ from $E$ to
$\C^*/q^{m\Z}$. Thus it is finite and proper.

\noindent\textbf{Step 6.} In fact $\Phi$ is \'etale everywhere. Local nonvanishing
frame factors are absorbed into the coordinate $u$.
At a zero branch choose the root parameter $\zeta$;
in the actual surface coordinates
$(a,b)$ its map is
\[
             (\zeta,u)\longmapsto(\zeta/u,u^m),
             \qquad \det d\Phi=m u^{m-2}\ne0.           \tag{C.PE.13}\label{cdp:eq:C.PE.13}
\]
At a pole use $y_0=\zeta^{-1}$ and the map
$(\zeta,u)\mapsto(a'=\zeta u,b'=u^{-m})$, whose
Jacobian is $-m u^{-m}\ne0$. Away from those branches
the assertion follows from the \'etale base map and
the fibre isomorphism. This proves the claimed finite
\'etale product cover. Its source is projective, and
a finite holomorphic image of a projective surface is
Moishezon; a smooth compact Moishezon surface is
projective. The comparison can also be made by taking
the finite \'etale Galois closure of \eqref{cdp:eq:C.PE.12}, a
projective surface, and its finite algebraic quotient.

\noindent\textbf{Step 7.} Let $\omega$ be a nowhere-zero one-form on $E$.
Off the multiple fibres, the form
$p^*\omega\wedge d\log\ell$ is well defined: a change
of frame adds a base one-form to $d\log\ell$, whose
wedge with $p^*\omega$ is zero. Near $J_+$, write
$\omega=h(x)dx$ and $x=a^m b$; then this two-form is
\[
                m h(a^m b)a^{m-1}\,da\wedge db.         \tag{C.PE.14}\label{cdp:eq:C.PE.14}
\]
Near $J_-$ the dual formula has a minus sign and the
same vanishing order $m-1$. It has no further zero.
This proves the canonical line in \eqref{cdp:eq:C.PE.2}. Constancy
on compact elliptic fibres and the invariant local
root charts give the exact direct image
\[
 p_*\OO_{P_m}(k_+J_++k_-J_-)
 =\OO_E\!\left(
  \left\lfloor\frac{k_+}m\right\rfloor p_+
 +\left\lfloor\frac{k_-}m\right\rfloor p_-\right).         \tag{C.PE.15}\label{cdp:eq:C.PE.15}
\]
A permitted Laurent term in a root parameter descends
precisely when its order is divisible by $m$.
Substitute $k_+=k_-=r(m-1)$ and use Riemann--Roch
on the elliptic curve. This proves the plurigenera
and $\kappa(P_m)=1$.

\noindent\textbf{Step 8.} We give the integral comparisons in \eqref{cdp:eq:C.PE.3}. The
real action
$(\ell,a)\mapsto(\lambda^m\ell,\lambda^{-1}a)$,
$\lambda>0$, preserves the equations and translates
$\log\|\ell\|$. Since $q$ is positive real, its
quotient gives actual differentiable products
\[
                   P_m\simeq M_3\times S^1,\qquad
                   Y_m\simeq M_5\times S^1,             \tag{C.PE.16}\label{cdp:eq:C.PE.16}
\]
where the $M$'s are the level sets $\|\ell\|=1$ in
the corresponding covering spaces. For $Y_m$ use
the construction over
$V=\operatorname{Bl}_{(p_+,0),(p_-,0)}(E\times\PP^1)$.
The local proper maps \eqref{cdp:eq:C.PE.8} make this quotient
compact and smooth. Off $0$ the product comparison
\eqref{cdp:eq:C.PE.9} applies, including at infinity.

\noindent\textbf{Step 9.} For $M_3$, remove small discs around $p_+,p_-$.
Over their complement the angular coordinate of
$a$ trivializes its circle fibration. Let $h$ be
the positive angular $a$-loop and let $\delta_\pm$
be positive base boundary loops with $a$ fixed.
In the positive solid torus the phases satisfy
$\delta_+=m\arg(a)+\arg(b)$; in the negative one
they satisfy $\delta_-=m\arg(a')+\arg(b')$ and
$\arg(a)=-\arg(a')$. The two meridians are
\[
                         m\delta_++h,\qquad
                         m\delta_- -h.                 \tag{C.PE.17}\label{cdp:eq:C.PE.17}
\]
In the abelianization of the punctured elliptic base,
$\delta_++\delta_-=0$. These relations give
$h=-m\delta_+$ and $\delta_-=-\delta_+$, leaving
$H_1(M_3,\Z)=\Z^3$, generated by the two
elliptic-base loops and $\delta_+$. There is no
torsion. Poincar\'e duality on this closed oriented
three-manifold and \eqref{cdp:eq:C.PE.16} give all the asserted
integral groups of $P_m$.

\noindent\textbf{Step 10.} For the compact five-manifold put
$Q=(E\times\PP^1)\setminus(B^4_+\sqcup B^4_-)$,
with two disjoint open balls. The exterior is
$Q\times S^1$, while each filling is
$S^3\times D^2$. Near an exceptional curve regard
a nonzero tautological vector as
$w\in\C^2\setminus0$. On the unit norm slice
$w\in S^3$, and the blowdown is
$u=a^m w$ or $u=(a')^m w$. The boundary maps are
\[
 \begin{split}
 (w,e^{\mathrm i\theta})&\longmapsto
       (e^{\mathrm i m\theta}w,e^{\mathrm i\theta}),\\
 (w,e^{\mathrm i\theta})&\longmapsto
       (e^{\mathrm i m\theta}w,e^{-\mathrm i\theta}).
 \end{split}                                             \tag{C.PE.18}\label{cdp:eq:C.PE.18}
\]
The pair sequence for the two removed balls gives
free homology of $Q$ of ranks $(1,2,2,3,0)$; its
two boundary sphere classes are opposite
primitive classes $\delta,-\delta$ in $H_3(Q,\Z)$.
In Mayer--Vietoris for \eqref{cdp:eq:C.PE.18}, the map on
$H_1(S^3\times S^1)^{\oplus2}$ has primitive
rank-one image, the circle generator of the
exterior. The map on $H_3$ is injective with
primitive image, since each of the two filling
$S^3$ summands receives its own identity map.
The map on $H_4$ has primitive rank-one image
generated by $\delta\times[S^1]$. Every kernel
and cokernel used in the sequence is therefore
free, and it gives
\[
       \operatorname{rank}H_k(M_5,\Z)=(1,2,5,5,2,1).
                                                               \tag{C.PE.19}\label{cdp:eq:C.PE.19}
\]
K\"unneth and \eqref{cdp:eq:C.PE.16} give the $Y_m$ row of
\eqref{cdp:eq:C.PE.3}, including integral torsion-freeness.

\noindent\textbf{Step 11.} Each Hopf surface is differentiably $S^3\times S^1$.
For the support use the actual union
$S_m=P_m\cup Q_+\cup Q_-$ with intersections
$J_+\sqcup J_-$. In the preceding $M_3$ presentation,
the core angular loop of $J_+$ is $\delta_+$ and
that of $J_-$ is $-\delta_+$. Their radial loops
are the $S^1$ factor of \eqref{cdp:eq:C.PE.16}. Inclusion in
$Q_\pm$ kills the angular loop and takes the
radial loop to its Hopf generator. Consequently
the integral map
\[
 \begin{split}
 H_1(J_+\sqcup J_-,\Z)\longrightarrow
 H_1(P_m,\Z)\oplus H_1(Q_+,\Z)\oplus H_1(Q_-,\Z)
 \end{split}                                             \tag{C.PE.20}\label{cdp:eq:C.PE.20}
\]
has primitive rank three: its columns are, up to
orientation signs, $\delta_+$, $z-q_+$,
$-\delta_+$, $z-q_-$. On $H_2$, the two elliptic
fundamental classes map to the same primitive
fibre class of $P_m$ and to zero on the Hopf
surfaces. Primitivity follows directly from
$[J_\pm]=\pm\delta_+\times[S^1]$ with the appropriate
product orientations. More explicitly, the dual chart at
$J_-$ reverses both its angular loop and its radial loop:
its complex orientation is
$(-\delta_+)\times(-z)=\delta_+\times z$.
Using the same global radial generator $z$ at both
ends instead reverses the displayed product orientation
at $J_-$. The restriction map on
cohomology thus has rank one in degree two and
rank three in degree one. Mayer--Vietoris gives
the free support ranks $(1,3,6,7,3)$ and Euler zero.

\noindent\textbf{Step 12.} The coherent groups also retain exact maps.
On every chart \eqref{cdp:eq:C.PE.8}, finite direct image and
the exact finite-group invariant functor give
\[
 \rho_*\OO_{Y_m}=\OO_V,\qquad
 R^1\rho_*\OO_{Y_m}=\OO_V,\qquad
 R^j\rho_*\OO_{Y_m}=0\quad(j\ge2).                       \tag{C.PE.21}\label{cdp:eq:C.PE.21}
\]
Translations act trivially on $H^1(E,\OO_E)$.
The trivialization of $R^1$ is global: the
constant cocycle assigning $1$ to the generator
\eqref{cdp:eq:C.PE.7} gives a class restricting to the
nonzero logarithmic cocycle on each elliptic
fibre, including the multiple fibres in the
finite-quotient chart. It is a nowhere-zero
frame. Since $h^k(V,\OO_V)=(1,1,0)$, Leray gives
$(1,2,1,0)$ for $Y_m$. The same local calculation
on $P_m\to E$ gives $(1,2,1)$ there, consistent
with \eqref{cdp:eq:C.PE.2}.

\noindent\textbf{Step 13.} For a primary Hopf surface, a holomorphic
two-form lifts and extends to $\C^2$; invariance
under its strict scalar contraction forces
every power-series coefficient to vanish.
Hence $p_g=0$. Its topological $e=0$ and
$H^2(-,\Z)=0$ give $\chi(\OO)=0$ by Noether's
formula, so its coherent dimensions are
$(1,1,0)$. Restriction
$H^1(Q_\pm,\OO)\to H^1(J_\pm,\OO)$ is nonzero:
the deck cocycle $1$ restricts to the same
cocycle on $\C^*/q^{m\Z}$. It is not a
coboundary, since the constant Laurent
coefficient of $g(q^m u)-g(u)$ is zero for
every holomorphic $g$ on $\C^*$.
The coherent normalization sequence for
$S_m$ therefore has a surjective difference
map in degrees zero and one. Its remaining
terms give $h^k(S_m,\OO)=(1,2,1)$.

\noindent\textbf{Step 14.} Finally, a rational curve in $P_m$ would lift
through the finite \'etale map \eqref{cdp:eq:C.PE.12}, since
$\PP^1$ is simply connected. Its two projections
to the positive-genus curves would be constant.
The Hopf components have no rational curve
either: a map from $\PP^1$ lifts to
$\C^2\setminus0$, and its coordinate functions
are constant. The inverse-map resolution
comparison in \ref{cdp:st:C.2.UN.3} prevents contraction
of any of these surfaces in a bimeromorphic
map between smooth threefold fillings. A
finite base change $t=s^d$ is generically,
along $P_m$, the graph $x=s^d$ for the local
transverse coordinate $x=t$. Its normalized
strict component maps with degree one to
$P_m$; resolution preserves that birational
component. The distinct plurigenus slopes
$2(1-1/m)$ distinguish the unique properly
elliptic components, and hence the fillings
for different $m$.

\noindent\textbf{Step 15.} This is an explicit signed logarithmic-transform
example. The prior Kodaira-dimension-one
occurrences discussed after \ref{cdp:st:C.2.UN.6}, including
\cite[Example~3.4(1)]{Fujimoto1988Log}, are retained
as the literature comparison.
\end{proof}






\begin{proposition}[A free half-period quotient realizes secondary Kodaira normalization]\label{cdp:st:C.2.UN.21}
\StatementRecord{cdp:st:C.2.UN.21}{proved}{Historical events: Claude Fable 5.2 (Anthropic), effort high, 16 November 2026: carried as background without fault found. GPT-6 Astra (OpenAI), ChatGPT Pro mode, effort Pro, S6\_09\_PAPERS, 17 November 2026: targeted proof check. Also checked in this writing. Pre-continuation writing events: GPT-6 Astra (OpenAI), ChatGPT Pro mode, effort Pro, S6\_10\_PAPERS, 18 November 2026: targeted proof check. Continuation metadata events: GPT-6 Astra (OpenAI), ChatGPT Pro mode, effort Pro, S6\_10\_PAPERS continuation, 19 November 2026: used and rechecked in this run; status-vocabulary reconciliation; metadata reconciliation only; no new mathematical check. \emph{Pre-continuation recorded proof scope.} The lifted involution and half-period translation give the exact free quotient and multiplicity-two fibres. The mapping-torus and integral Gysin maps, invariant bounding chain for the double curve, all torsion groups, coherent eigensheaf calculation and the actual normalized quadratic base change were re-derived. The latter recovers the primary Kodaira cover; no persistence of the secondary quotient type is asserted}{the hypotheses, notation and definitions stated here; \cref{cdp:st:C.2.UN.3}; \cref{cdp:eq:C.SKD.4}; \cref{cdp:eq:C.SKD.5}; \cref{cdp:eq:C.SKD.7}; \cref{cdp:eq:C.SKD.9}; \cref{cdp:eq:C.SKD.6}; \cref{cdp:st:C.2.UN.18}; \cref{cdp:eq:C.SKD.12}}

Let $q=e^{-2\pi}$ and $E=\C^*/q^{\Z}$, and use the
origin of its group structure. There is a smooth compact
threefold $Y\to\PP^1$ whose smooth general fibre is
$E\times E$ and whose two special full fibres are
\[
                        2P+2Q\quad\hbox{and}\quad 2B.
                                                            \tag{C.SKD.1}\label{cdp:eq:C.SKD.1}
\]
The surface $P$ is a secondary Kodaira surface, $Q$ is
a primary Hopf surface, and they meet transversely along
a smooth elliptic curve $J$. The surface $B$ is smooth
bielliptic. Both reduced special supports have Euler
number zero. In particular $P$ occurs as an essential
component in a zero-Euler torus filling.

The canonical bundle of $P$ has exact order two. For
$S=P\cup Q$, the rational cohomology ranks and coherent
dimensions are
\[
 \begin{array}{c|c|c}
 &\dim H^k(-,\Q)&h^k(\OO)\\ \hline
 P&(1,1,0,1,1)&(1,1,0)\\
 Q&(1,1,0,1,1)&(1,1,0)\\
 J&(1,2,1)&(1,1)\\
 S&(1,1,1,3,2)&(1,1,0)\\
 B&(1,2,2,2,1)&(1,1,0)\\
 Y&(1,1,2,4,2,1,1)&(1,1,0,0).
 \end{array}                                             \tag{C.SKD.2}\label{cdp:eq:C.SKD.2}
\]
The full schemes in \eqref{cdp:eq:C.SKD.1}, in contrast to their
reductions, each have coherent dimensions $(1,2,1)$.
The primary Kodaira double cover of $P$ is recovered
by the actual quadratic base change at the first
special value.
\end{proposition}


\begin{proof}
\noindent\textbf{Step 1.} Start with the compact version of \ref{cdp:st:C.2.UN.3} at $n=1$:
\[
 V=\operatorname{Bl}_{(0,0)}(E\times\PP^1_t),\qquad
 L=\OO_V(R),\qquad
 \widetilde Y=L^\times/q^{\Z},
                                                            \tag{C.SKD.3}\label{cdp:eq:C.SKD.3}
\]
where $R$ is the exceptional curve. Let $\sigma$ be
the canonical section of $L$. The projection
$\widetilde Y\to V\to\PP^1_t$ is proper and smooth
off its reduced central divisor
$\widetilde P+\widetilde Q$. The first component is
the primary Kodaira surface
$\OO_E([0])^\times/q^{\Z}$, the second is the Hopf
surface $\OO_{\PP^1}(-1)^\times/q^{\Z}$, and their
intersection is an elliptic curve. These are the
actual line-bundle quotient maps of \ref{cdp:st:C.2.UN.3}.

\noindent\textbf{Step 2.} The involution $(p,t)\mapsto(-p,-t)$ fixes the blowup
centre and lifts to an involution $\iota$ of $V$.
It has the canonical order-two lift $\widehat\iota$
to $L$ preserving $\sigma$. Define an automorphism
of $\widetilde Y$ by
\[
             g[\ell]=[q^{1/2}\widehat\iota(\ell)].
                                                            \tag{C.SKD.4}\label{cdp:eq:C.SKD.4}
\]
Its square is multiplication by $q$, so $g^2=1$
on $\widetilde Y$. It acts freely. Indeed at a
fixed base point the order-two linearization on
the line $L_v$ has eigenvalue $+1$ or $-1$.
Neither $q^{1/2}$ nor $-q^{1/2}$ belongs to
$q^{\Z}$. Thus \eqref{cdp:eq:C.SKD.4} has no fixed point on
any elliptic fibre, including those over $R$.
The quotient
\[
              \pi:\widetilde Y\longrightarrow Y
                        =\widetilde Y/\langle g\rangle
                                                            \tag{C.SKD.5}\label{cdp:eq:C.SKD.5}
\]
is finite \'etale of degree two, and $Y$ is smooth
and compact. The function $T=t^2$ defines its map
to $\PP^1_T$. For $T\ne0,\infty$, the two values
of $t$ are interchanged by $g$, so each quotient
fibre is exactly one of the original product
tori.

\noindent\textbf{Step 3.} The strict central base curve and $R$ are
invariant. The local equation upstairs at their
intersection is $t=xy$, hence $T=x^2y^2$.
Since \eqref{cdp:eq:C.SKD.5} is \'etale, this proves the first
full fibre and the normal-crossing assertion
in \eqref{cdp:eq:C.SKD.1}. At infinity the original smooth
fibre is $E\times E$, and the action is base
inversion on the first elliptic factor and
translation by the class of $q^{1/2}$ on the
second. It is free. Its quotient is a
bielliptic surface $B$, and the transverse
coordinate is squared, giving the full fibre
$2B$.

\noindent\textbf{Step 4.} We identify the central quotients more precisely.
The canonical linearization on
$L|_R=\OO_{\PP^1}(-1)$ acts by $-1$ on its
tautological vector: the derivative at the
blown-up point is the scalar $-I$ on $\C^2$.
Thus the quotient of $\widetilde Q$ is
\[
                 Q=(\C^2\setminus0)/
                         (-q^{1/2})^{\Z},              \tag{C.SKD.6}\label{cdp:eq:C.SKD.6}
\]
again a primary Hopf surface. On $\widetilde P$
the involution covers inversion on $E$.
The nowhere-zero canonical form of the primary
Kodaira surface is the wedge of the base
one-form and the fibre logarithmic differential.
The first changes sign; the second does not.
Consequently $g$ acts by $-1$ on that form.
Its square descends to a nowhere-zero section
of $K_P^{\otimes2}$, while no canonical form
descends. Thus $K_P$ has exact order two.
The finite \'etale cover $\widetilde P\to P$
and the computed $b_1(P)=1$ below identify $P$
as a secondary Kodaira surface.

\noindent\textbf{Step 5.} We give the topology through actual mapping
tori. Choose an $\iota$-invariant Hermitian
metric on $L$, by averaging. Let $M_5$ be its
unit circle bundle. The positive-real norm
coordinate gives
$\widetilde Y\simeq M_5\times S^1$. In that
description \eqref{cdp:eq:C.SKD.4} is
$(x,u)\mapsto(\iota_Mx,u+\tfrac12)$, where
the period of $u$ is one. Hence the quotient
is the mapping torus
\[
          Y\simeq T(\iota_M)
            =(M_5\times\R)/
               ((x,u+1)\sim(\iota_Mx,u)).               \tag{C.SKD.7}\label{cdp:eq:C.SKD.7}
\]
For example the map $[x,u]\mapsto[x,2u]$
identifies the two quotient descriptions.

\noindent\textbf{Step 6.} The integral circle-bundle Gysin sequence
over $V$ has Euler class $[R]$. The base has
free cohomology of ranks $(1,2,3,2,1)$.
Cup product with $[R]$ has ranks $1,0,1$
in degrees zero, one and two; the nonzero
images are primitive because $R^2=-1$.
Thus $M_5$ has free homology of ranks
\[
                         (1,2,4,4,2,1).                \tag{C.SKD.8}\label{cdp:eq:C.SKD.8}
\]
Its involution acts as $-I$ on $H_1$ and,
by oriented Poincar\'e duality, on $H_4$.
On $H_2$ it has an integral diagonal
decomposition with two $+1$ and two $-1$
summands. To specify that decomposition,
lift $E\times\{\infty\}$ and
$\{p_*\}\times\PP^1$, where $p_*\ne0$ is
a two-torsion point, using the normalized
canonical section of $L$. These are
invariant cycles, avoid $R$, and generate
the two base summands. The preimages of
the two elliptic-base loops at infinity
are circle tori; they generate the two
remaining summands and are sent to their
negatives. The homological Gysin sequence
identifies these four cycles with an
integral basis. Duality gives the same
two-positive, two-negative decomposition
on $H_3$. The actions on $H_0,H_5$ are
the identity. The Wang sequence of
\eqref{cdp:eq:C.SKD.7} therefore gives the rational
$Y$ row of \eqref{cdp:eq:C.SKD.2}.

\noindent\textbf{Step 7.} This also determines its integral groups,
not only their ranks. Writing $T_2=(\Z/2)^2$,
the integral cohomology of $Y$ is
\[
 \begin{array}{c|ccccccc}
 k&0&1&2&3&4&5&6\\ \hline
 H^k(Y,\Z)&
 \Z&\Z&\Z^2\oplus T_2&\Z^4\oplus T_2&
 \Z^2\oplus T_2&\Z\oplus T_2&\Z .
 \end{array}                                             \tag{C.SKD.9}\label{cdp:eq:C.SKD.9}
\]
Indeed the Wang maps are zero or multiplication
by $-2$ on the stated integral summands.
Their kernel terms are free, so the short
exact sequences split as abelian groups;
the universal coefficient sequence gives
\eqref{cdp:eq:C.SKD.9}.

\noindent\textbf{Step 8.} On $\widetilde P$, the unit circle bundle
$M_3\to E$ has Euler number one and free
homology ranks $(1,2,2,1)$. Its involution
acts by $-I$ in degrees one and two and
by $+1$ in degrees zero and three.
The quotient $P$ is its mapping torus.
The same Wang maps give
\[
 \begin{array}{c|ccccc}
 k&0&1&2&3&4\\ \hline
 H^k(P,\Z)&
 \Z&\Z&T_2&\Z\oplus T_2&\Z .
 \end{array}                                             \tag{C.SKD.10}\label{cdp:eq:C.SKD.10}
\]
In particular its rational ranks are those
of \eqref{cdp:eq:C.SKD.2}. The Hopf ranks follow from
\eqref{cdp:eq:C.SKD.6}, which is differentiably
$S^3\times S^1$.

\noindent\textbf{Step 9.} For the support use the exact union
$S=P\cup Q$ with intersection $J$. The
radial deck loop of $J$ maps to the
mapping-torus generator on $P$ and the
Hopf generator on $Q$. Its angular loop
maps to zero in rational first homology:
upstairs it is the circle of an
Euler-number-one bundle, hence is already
zero in the abelianization. Therefore
the map
\[
           H^1(P,\Q)\oplus H^1(Q,\Q)
                         \longrightarrow H^1(J,\Q)
                                                            \tag{C.SKD.11}\label{cdp:eq:C.SKD.11}
\]
has rank one. Both surfaces have zero
rational $H^2$, so the next restriction
map has rank zero. Mayer--Vietoris gives
the support ranks $(1,1,1,3,2)$.
There is also an exact integral comparison at the only
potential torsion restriction. In the unit circle bundle
$M_3\to E$, the core fibre over zero bounds an
$\iota_M$-invariant oriented surface. To construct it,
take the normalized canonical section over the complement
of a small invariant disc at zero. Choose a local
coordinate and an equivariant frame with
$\sigma=z e$; the lift is
$(z,u)\mapsto(-z,-u)$ in the unit-circle coordinates.
The section on the disc boundary is
$(z,u)=(\varepsilon e^{i\theta},e^{i\theta})$.
Attach the annulus
\[
 (r,\theta)\longmapsto(re^{i\theta},e^{i\theta}),
                 \qquad 0\le r\le\varepsilon .
\]
Its inner boundary is the core circle. The involution
is $\theta\mapsto\theta+\pi$ on this annulus and
preserves its orientation; it also preserves the
orientation of the section over the punctured elliptic
curve. The joined surface, or its invariant rounded
chain, therefore has an oriented mapping torus whose
boundary is $J$. Thus
$H_2(J,\Z)\to H_2(P,\Z)$ is zero integrally,
including its possible torsion target. The angular
first-homology class is already zero in $M_3$;
the other generator maps to a primitive mapping-torus
class on $P$ and to a primitive Hopf class on $Q$.
The integral Mayer--Vietoris sequence now gives
\[
 \begin{array}{c|ccccc}
 k&0&1&2&3&4\\ \hline
 H_k(S,\Z)&
 \Z&\Z\oplus T_2&\Z\oplus T_2&\Z^3&\Z^2\\
 H^k(S,\Z)&
 \Z&\Z&\Z\oplus T_2&\Z^3\oplus T_2&\Z^2 .
 \end{array}                                             \tag{C.SKD.14}\label{cdp:eq:C.SKD.14}
\]
The free kernel terms split as abelian groups.
Any torsion summand in the radial image is absorbed
by eliminating the primitive Hopf generator, so
no choice of radial lift changes these groups.
The invariant cohomology of the free
quotient of $E\times E$ at infinity
gives the row for $B$. In both cases
the displayed Euler number is zero.
For completeness, $B$ is also the product of an
angular circle with the mapping torus of inversion
on the first elliptic curve. The Wang maps on
that curve are $0,-2I_2,0$. Its integral
cohomology is therefore
\[
 H^k(B,\Z)=
 \bigl(\Z,\ \Z^2,\ \Z^2\oplus T_2,\
                  \Z^2\oplus T_2,\ \Z\bigr)_k .
                                                            \tag{C.SKD.15}\label{cdp:eq:C.SKD.15}
\]
The integral groups of $J$ and $Q$ are free with
the ranks already displayed.

\noindent\textbf{Step 10.} For coherent cohomology, the elliptic
line-bundle quotient $\widetilde Y\to V$
has direct images $\OO_V$ in degrees
zero and one and zero above one. The
degree-one frame is the constant
logarithmic deck cocycle, as in
(C.PE.21). Its involution action is
$+1$, while the action on
$H^1(V,\OO_V)=H^1(E,\OO_E)$ is $-1$.
The finite \'etale invariant-sheaf
comparison for \eqref{cdp:eq:C.SKD.5}, applied to
Leray's filtration, gives
$h^k(Y,\OO_Y)=(1,1,0,0)$. The identical
calculation for $\widetilde P\to E$
gives $h^k(P,\OO_P)=(1,1,0)$, consistent
with the canonical character.
The Hopf and elliptic rows were
computed from their actual deck cocycles
in \ref{cdp:st:C.2.UN.18}. In particular
$H^1(Q,\OO_Q)\to H^1(J,\OO_J)$ is
nonzero and hence onto. The coherent
normalization sequence gives
$h^k(S,\OO_S)=(1,1,0)$. The quotient
at infinity has the same base-negative,
fibre-positive action, giving
$h^k(B,\OO_B)=(1,1,0)$.

\noindent\textbf{Step 11.} We finally verify the full scheme fibres.
Upstairs let $\widetilde S$ be the
reduced central fibre of the original
map to $t$. Its coherent normalization
sequence gives $(1,2,1)$, the same
dimensions as a nearby product torus.
Grauert's theorem therefore makes all
$R^qf_*\OO_{\widetilde Y}$ locally free
near $t=0$, of ranks $1,2,1$, with
base change. The actual inverse image
of the full fibre $2S$ under the
\'etale map \eqref{cdp:eq:C.SKD.5} is $2\widetilde S$,
defined by $t^2=0$. Thus
\[
 H^q(2S,\OO_{2S})
  \simeq
 \left((R^qf_*\OO_{\widetilde Y})/
                   t^2(R^qf_*\OO_{\widetilde Y})\right)^g .
                                                            \tag{C.SKD.12}\label{cdp:eq:C.SKD.12}
\]
Choose a basis of eigenvectors on the
central fibre of each direct-image
bundle. Lift each vector and average
with $g$ with its specified sign; the
resulting holomorphic sections remain
a frame near zero and satisfy
$g(e)=\pm e$. Since $g(t)=-t$, the
invariants in each summand
$\C[t]/(t^2)e$ have dimension one:
they are generated by $e$ for sign
$+1$ and by $te$ for sign $-1$.
Equation \eqref{cdp:eq:C.SKD.12} gives precisely
$(1,2,1)$. The same argument at
infinity proves the assertion for
$2B$. This comparison retains both
the nilpotent direction and its sign.

\noindent\textbf{Step 12.} No rational curve lies in $P$:
its map from $\PP^1$ would lift to
the primary Kodaira cover, whose
elliptic base and elliptic fibres
force that map to be constant.
The same lifting argument to
$\C^2\setminus0$ treats $Q$.
The inverse-map resolution comparison
of \ref{cdp:st:C.2.UN.3} therefore proves their
essentiality. At $T=0$, the finite
map
\[
 \widetilde Y\longrightarrow
             Y\times_{\PP^1_T}\PP^1_t,\qquad
 y\longmapsto(\pi(y),f(y)),\quad T=t^2,
                                                            \tag{C.SKD.13}\label{cdp:eq:C.SKD.13}
\]
is generically one-to-one and has
normal source. It identifies
$\widetilde Y$ with the normalization
of the reduced fibre product.
Consequently the primary Kodaira
component is recovered by this
actual quadratic base change.
The construction supplies an
explicit secondary Kodaira component;
no birational persistence of its
secondary type under that base change
is asserted.
\end{proof}






\begin{proposition}[An essential Enriques surface occurs in a compact zero-Euler torus filling]\label{cdp:st:C.2.UN.25}
\StatementRecord{cdp:st:C.2.UN.25}{proved}{Historical events: Claude Fable 5.2 (Anthropic), effort high, 16 November 2026: carried as background without fault found. GPT-6 Astra (OpenAI), ChatGPT Pro mode, effort Pro, S6\_09\_PAPERS, 17 November 2026: targeted proof check. Also checked in this writing. Pre-continuation writing events: GPT-6 Astra (OpenAI), ChatGPT Pro mode, effort Pro, S6\_10\_PAPERS, 18 November 2026: targeted proof check; used and rechecked in this run. Continuation metadata events: GPT-6 Astra (OpenAI), ChatGPT Pro mode, effort Pro, S6\_10\_PAPERS continuation, 19 November 2026: used and rechecked in this run; status-vocabulary reconciliation; metadata reconciliation only; no new mathematical check. \emph{Pre-continuation recorded proof scope.} The signed root cover, commuting product involutions, eight surface A1 points, four ambient fixed curves and both resolution charts were re-derived. The full multiplicities, Enriques double cover, rational-curve relations and all restriction ranks give the printed cohomology. Six explicit balancing blocks preserve the Enriques component. The discriminant and its gcd were recomputed exactly in analysis/UN25\_discriminant.json, and Yoshihara's comparison loci were opened}{the hypotheses, notation and definitions stated here; \cref{cdp:st:C.2.UN.14}; \cref{cdp:st:C.2.UN.18}; \cref{cdp:eq:C.EN.7}; \cref{cdp:eq:C.EN.8}; \cref{cdp:st:C.2.UN.23}; \cite[Theorem~2.1 and Example~2.3]{Yoshihara1995}; \cref{cdp:eq:C.EN.15}; \cref{cdp:st:C.2.UN.4}; \cref{cdp:eq:C.EN.S2}; \cref{cdp:eq:C.EN.S1}; \cref{cdp:eq:C.EN.S4}}

There is a smooth compact complex threefold over
$\PP^1$ with connected two-torus general fibre and
a simple-normal-crossing special support of Euler
number zero containing an unchanged smooth Enriques
surface as an irreducible component.
Before six explicit vertical modifications, its full
central fibre is
\[
 2M+2A+4Q_1+4Q_2+E_1+E_2+E_3+E_4,                 \tag{C.EN.1}\label{cdp:eq:C.EN.1}
\]
where $M$ is Enriques, $A\simeq\PP^1\times\PP^1$,
$Q_1,Q_2$ are primary Hopf surfaces, and every
$E_i$ is the Hirzebruch surface $\mathbb F_2$.
That reduced support has Euler number $12$.
All six modifications can be made disjoint from $M$.
The Enriques component is therefore essential in the
sense of \ref{cdp:st:C.2.UN.14}.
The reduced supports before and after balancing have
respectively
\[
 \begin{array}{c|c|c}
 & (b_0,b_1,b_2,b_3,b_4)&(h^0(\OO),h^1(\OO),h^2(\OO))\\ \hline
 S&(1,1,8,4,8)&(1,1,0)\\
 S_{\rm bal}&(1,1,20,40,20)&(1,1,0).
 \end{array}                                             \tag{C.EN.18}\label{cdp:eq:C.EN.18}
\]
\end{proposition}


\begin{proof}
\noindent\textbf{Step 1.} We first give a compact base with specified local
coordinates. Take the Weierstrass elliptic surface
over $\PP^1_T$
\[
 y^2=x^3+(T^4-3)x+(T^6+T+2).                       \tag{C.EN.2}\label{cdp:eq:C.EN.2}
\]
The coefficients are sections of $\OO(4)$ and
$\OO(6)$. Its discriminant, up to the usual nonzero
constant factor, is
\[
 \begin{split}
 \Delta(T)={}&31T^{12}-36T^8+54T^7+108T^6\\
             &+108T^4+27T^2+108T .
 \end{split}                                       \tag{C.EN.3}\label{cdp:eq:C.EN.3}
\]
The exact Euclidean algorithm gives
$\gcd(\Delta,\Delta')=1$; also
$\Delta'(0)=108$ and its leading coefficient is
$31$. These checks are saved in
\texttt{code/E\_enriques\_filling.py}.
Thus every singular fibre is nodal of type $I_1$,
the fibre at zero is nodal, and infinity is smooth.
The total Weierstrass surface is smooth.

\noindent\textbf{Step 2.} Here is an explicit chart at the node of the
zero fibre. Put $\xi=x-1$ and
$U=1+T^3(1+\xi)+T^5$. On the surface,
\[
 X=\frac{y-\xi\sqrt{3+\xi}}{U},\qquad
 Y=y+\xi\sqrt{3+\xi}
 \quad\Longrightarrow\quad XY=T .                  \tag{C.EN.4}\label{cdp:eq:C.EN.4}
\]
The coordinate Jacobian at $\xi=y=T=0$ is
nonzero, so these are actual local coordinates.
Make the quadratic base change $T=t^2$ and
resolve $XY=t^2$. One resolution chart is
\[
       (X,Y,t)=(u,uv^2,uv),\qquad
       h(u,v)=(u,-v),                              \tag{C.EN.5}\label{cdp:eq:C.EN.5}
\]
where $h$ is the deck involution. The other chart
interchanges $X,Y$ and replaces $v$ by $v^{-1}$.
The central fibre of the smooth base surface $V_0$
is $R_0+R_1$ of type $I_2$. Its strict nodal
normalization $R_0\simeq\PP^1$ is fixed pointwise
by $h$. The exceptional curve $R_1\simeq\PP^1$
has coordinate $z=v$, with $h(z)=-z$, and
meets $R_0$ exactly at $z=0,\infty$.

\noindent\textbf{Step 3.} Blow up the four points $z=1,2,-1,-2$ of $R_1$.
Write their exceptional curves as
$D_1,D_2,D_{-1},D_{-2}$, and denote the resulting
compact surface by $V$. The involution lifts
and exchanges $D_j,D_{-j}$. Set
\[
 D=D_1+D_2-D_{-1}-D_{-2},\qquad L=\OO_V(D).
                                                               \tag{C.EN.6}\label{cdp:eq:C.EN.6}
\]
The canonical meromorphic section $\sigma$ of
$L$ supplies the section-preserving identification
$h^*L=L^{-1}$. Its square is the identity.

\noindent\textbf{Step 4.} Apply the explicit signed root construction of
\ref{cdp:st:C.2.UN.18} with order $m=2$ at these four disjoint
curves. Thus on the complement of the negative
curves use $\sigma=a^2\ell$, with $\ell\in L^\times$;
on the complement of the positive curves use
$\sigma^{-1}=(a')^2\ell'$, where
$a'=a^{-1}$ and $\ell'=\ell^{-1}$. Divide by
\[
       (\ell,a)\longmapsto(q^2\ell,q^{-1}a),
                  \qquad 0<q<1 .                  \tag{C.EN.7}\label{cdp:eq:C.EN.7}
\]
Every local root chart is the already verified
free finite quotient of an elliptic product in
(C.PE.8), so the resulting compact threefold
$\widetilde Y$ is smooth and proper over $V$.
Away from the four exceptional curves the
actual product coordinate is
$\ell=\sigma c^2,\ a=c^{-1}$, with
$[c]\in E=\C^*/q^{\Z}$.

\noindent\textbf{Step 5.} Dualization gives an involution of
$\widetilde Y$: it takes the global meromorphic
coordinates $(\ell,a)$ to $(\ell^{-1},a^{-1})$
over $h$. Equivalently, a positive chart maps
to the target negative chart with unchanged
numerical coordinates $(\ell',a')=(\ell,a)$.
It conjugates \eqref{cdp:eq:C.EN.7} to its inverse, so
descends to the quotient and still squares
to one. In the product coordinate it is
\[
                       (v,[c])\longmapsto(hv,[c^{-1}]).
                                                               \tag{C.EN.8}\label{cdp:eq:C.EN.8}
\]
It covers $t\mapsto-t$.

\noindent\textbf{Step 6.} We identify the component over the strict
transform of $R_1$. The restriction of $L$
there is trivial; in a compatible frame,
\[
 \sigma|_{R_1}=\phi(z)
   =\frac{(z-1)(z-2)}{(z+1)(z+2)},\qquad
                    \phi(-z)=\phi(z)^{-1}.          \tag{C.EN.9}\label{cdp:eq:C.EN.9}
\]
Let $C$ be the smooth elliptic curve
\[
 C:\quad Y^2=(z^2-1)(z^2-4),\qquad
 y=\frac{Y}{(z+1)(z+2)},\quad y^2=\phi(z).
                                                               \tag{C.EN.10}\label{cdp:eq:C.EN.10}
\]
Its four simple branch values are $1,2,-1,-2$.
The exact finite product-cover map from \ref{cdp:st:C.2.UN.18} is
\[
 C\times E\longrightarrow B,\qquad
 (z,y,[u])\longmapsto[\ell=u^2,\ a=y/u].
                                                               \tag{C.EN.11}\label{cdp:eq:C.EN.11}
\]
It is \'etale, including at the four root values,
by the two Jacobians in (C.PE.13). Its deck
involution is
\[
 \eta:(z,y,[u])\longmapsto(z,-y,[-u]).
                                                               \tag{C.EN.12}\label{cdp:eq:C.EN.12}
\]
The translation $[u]\mapsto[-u]$ is a nonzero
two-torsion translation, so $\eta$ is free.
Thus $B=(C\times E)/\langle\eta\rangle$ is
bielliptic. The involution \eqref{cdp:eq:C.EN.8} lifts to
\[
 k:(z,y,[u])\longmapsto(-z,y^{-1},[u^{-1}])
       \quad\text{or}\quad
       (z,Y,[u])\longmapsto(-z,Y,[u^{-1}]).
                                                               \tag{C.EN.13}\label{cdp:eq:C.EN.13}
\]
These formulas commute with $\eta$.

\noindent\textbf{Step 7.} On $C$, the involution $j:(z,Y)\mapsto(z,-Y)$
and $k_C:(z,Y)\mapsto(-z,Y)$ each act by $-1$
on $dz/Y$. The latter has four fixed points,
the two points over each of $0,\infty$.
Their product $jk_C$ has no fixed point:
over zero its equation would require $Y=0$,
and at infinity it exchanges the two points.
It is consequently a nonzero two-torsion
translation of the elliptic curve $C$.
It follows that $\eta$ and $\eta k$ are
free on $C\times E$, whereas $k$ has exactly
$4\cdot4=16$ fixed points. They occur in
eight free $\eta$-orbits. At each, the
derivative of $k$ is $(-1,-1)$.

\noindent\textbf{Step 8.} The quotient $B/\langle h\rangle$ therefore
has exactly eight $A_1$ singularities and
no branch curve. Its minimal resolution
$M$ is Enriques. One exact verification
uses the intermediate quotient
$(C\times E)/\langle k\rangle$. Its
minimal resolution is K3 by \ref{cdp:st:C.2.UN.23}.
The involution $\eta$ extends over its
sixteen exceptional curves, permuting
them in pairs. It is free away from
them because both $\eta$ and $\eta k$
are free upstairs, and free on them
because they are paired. Its quotient
is precisely $M$. Equivalently, the
canonical form has character $-1$
under $\eta$ and $+1$ under $k$;
its square descends through the
explicit crepant $A_1$ resolutions,
while $p_g(M)=q(M)=0$. Thus $K_M$
has exact order two and $e(M)=12$.
This is an actual noncyclic product
quotient, consistent with the two
Enriques holonomy cases of
\cite[Theorem~2.1 and Example~2.3]{Yoshihara1995}.

\noindent\textbf{Step 9.} It remains to realize this surface resolution
inside the threefold, not just on an abstract
component. Near $R_0$ the four root curves
are absent, so \eqref{cdp:eq:C.EN.8} is a product.
Its fixed locus over zero consists of four
disjoint copies of $R_0$, indexed by $E[2]$.
Near either node of $R_0+R_1$, choose a
local coordinate $w$ at the corresponding
two-torsion point of $E$, so inversion is
$w\mapsto-w$. The ambient action is
\[
 (u,v,w)\longmapsto(u,-v,-w),\qquad
                         T=u^2v^2 .
                                                               \tag{C.EN.14}\label{cdp:eq:C.EN.14}
\]
Its quotient has invariant coordinates
$A_0=v^2,B_0=vw,C_0=w^2$, with
$A_0C_0=B_0^2$ and $T=u^2A_0$.
The two smooth resolution charts are
\[
 \begin{array}{lll}
 A_0=r,&B_0=rs,&C_0=rs^2,\quad T=u^2r,\\
 A_0=r's'^2,&B_0=r's',&C_0=r',\quad T=u^2r's'^2 .
 \end{array}                                             \tag{C.EN.15}\label{cdp:eq:C.EN.15}
\]
They glue by $s'=s^{-1}$ and $r'=rs^2$.
Globally they resolve the four fixed
curves. Restricting to $u=0$ is exactly
the minimal $A_1$ resolution of the
eight surface singularities just found.
In particular the strict central component
is the smooth Enriques surface $M$.

\noindent\textbf{Step 10.} The component over $R_0$ becomes
$A=R_0\times(E/\{\pm1\})\simeq\PP^1\times\PP^1$.
Each exceptional surface of \eqref{cdp:eq:C.EN.15}
is the projective normal-direction bundle
of a fixed curve. That normal bundle is
$\OO_{R_0}(-2)\oplus\OO_{R_0}$: the
product trivialization is valid on an
entire neighbourhood of $R_0$, and
$R_0^2=-2$. Hence these surfaces are
$E_i\simeq\mathbb F_2$.
The four root Hopf components are
exchanged in two pairs and give
$Q_1,Q_2$. Squaring the base parameter,
and the exact monomials in \eqref{cdp:eq:C.EN.15},
give every coefficient in \eqref{cdp:eq:C.EN.1}.
Those same monomials show that its
reduced support is simple normal
crossing.

\noindent\textbf{Step 11.} The full incidence calculation is
\[
 \begin{array}{c|c}
 \text{intersection}&\text{reduced space}\\ \hline
 M\cap A&\PP^1\sqcup\PP^1\\
 M\cap Q_j&\text{one elliptic curve}\quad(j=1,2)\\
 A\cap E_i&\PP^1\quad(i=1,\ldots,4)\\
 M\cap E_i&\PP^1\sqcup\PP^1\quad(i=1,\ldots,4)\\
 M\cap A\cap E_i&\text{two points}\quad(i=1,\ldots,4).
 \end{array}                                             \tag{C.EN.16}\label{cdp:eq:C.EN.16}
\]
There are no other nonempty intersections.
Thus inclusion--exclusion gives, with
every nonzero term displayed,
\[
 e(S)=12+4+4\cdot4-(4+4\cdot2+4\cdot4)+8=12 .
                                                               \tag{C.EN.17}\label{cdp:eq:C.EN.17}
\]

\noindent\textbf{Step 12.} For the compact construction, the base involution
also fixes the smooth elliptic fibre at infinity.
The ambient fixed loci there are four elliptic
curves, and the same transverse $A_1$ charts
resolve them. Their exceptional surfaces and
all intersections are elliptic ruled and
have Euler number zero. At the remaining
nodal base values the elliptic-product
coordinate gives a nodal curve times $E$,
also of Euler number zero. The quotient
and the resolutions are proper and compact.
Over a general $T$, the two $t$-sheets
are exchanged by $h$, so their quotient
is exactly one product of two elliptic
curves. The Stein factor has generic
degree one over the normal curve
$\PP^1$, proving connectedness.

\noindent\textbf{Step 13.} Finally perform six of the point-then-quartic
operations of \ref{cdp:st:C.2.UN.4}, all over $T=0$ and
away from $M$. They can begin at fresh
smooth points of an $E_i$ outside its
intersections. Their exceptional plane
quartics and subsequent centres remain
disjoint from $M$. Each changes the
reduced Euler number by $-2$, preserves
the smooth total space and the normal
crossing support, and leaves the
punctured family unchanged. The resulting
support has Euler number zero, with
the same smooth Enriques component.
Its nonzero even pluricanonical form
precludes a covering family of rational
curves, so it is non-uniruled.
The exact inverse-map resolution argument
of \ref{cdp:st:C.2.UN.14} proves essentiality. This completes
the compact occurrence construction.
\par\medskip\noindent\textbf{The reduced support and its exact cohomology.}
Let $S$ denote the reduced support before the six balancing
operations. Label the Enriques component $M$, the component over
$R_0$ by $A\simeq\PP^1\times\PP^1$, the two Hopf components by
$Q_a,Q_b$, and the four exceptional ruled surfaces by
$E_1,\ldots,E_4\simeq\mathbb F_2$.
The actual quotient and resolution charts give
\[
\begin{array}{c|c|c}
\text{components}&\text{double curve}&\text{number}\\ \hline
M,A&D_0,D_\infty\simeq\PP^1&2\\
M,Q_j&J_j\text{ elliptic},\ j=a,b&2\\
A,E_i&Z_i\simeq\PP^1&4\\
M,E_i&U_{i0},U_{i\infty}\simeq\PP^1&8.
\end{array}                                                   \tag{C.EN.S1}\label{cdp:eq:C.EN.S1}
\]
There are precisely eight triple points
$D_\epsilon\cap Z_i\cap U_{i\epsilon}$, one for
$1\le i\le4$ and $\epsilon=0,\infty$, and no higher
intersection. On $A$, the curves $D_0,D_\infty$ are two
fibres of one projection and the $Z_i$ are sections of that
projection, all with the same divisor class. On $E_i$, $Z_i$
is the negative section and $U_{i0},U_{i\infty}$ are two
ruling fibres. These assertions retain the actual intersections:
the local quotient equation is $AC=B^2$, with base function
$T=u^2A$, and its $C$-chart after resolution has
$A=C\delta^2$, $B=C\delta$, hence $T=u^2C\delta^2$.
The three reduced components are $u=0$, $C=0$, and
$\delta=0$. Thus this support has simple normal crossings.

\noindent\textbf{Step 14.} The elliptic map $M\to\PP^1$, induced by $z\mapsto z^2$,
has at $0,\infty$ the actual fibre divisors
\[
 \begin{gathered}
 F_0=2D_0+\sum_{i=1}^4U_{i0},\qquad
 F_\infty=2D_\infty+\sum_{i=1}^4U_{i\infty},\\
 [F_0]=[F_\infty]=2[J_a]=2[J_b]
       \quad\text{in }H_2(M,\Q).
 \end{gathered}                                          \tag{C.EN.S2}\label{cdp:eq:C.EN.S2}
\]
Indeed near a fixed base value the elliptic involution has
four fixed points; the displayed resolution gives one central
rational curve of multiplicity two and four exceptional arms
of multiplicity one. At the two root values the chart
$x=a^2b$ gives the two smooth double fibres $2J_a,2J_b$.
All these are fibres of the specified proper elliptic map,
which proves the equality of their rational fundamental classes.

\noindent\textbf{Step 15.} The ten rational curves in \eqref{cdp:eq:C.EN.S2} have precisely one rational
homology relation, namely $F_0-F_\infty=0$. To check this,
their intersection matrix is the direct sum of two affine
$D_4$ matrices: each central curve and its four arms has
self-intersection $-2$, each arm meets the central curve
once, and all remaining intersections are zero. Each matrix
has rank four and kernel generated by $(2,1,1,1,1)$.
Any homology relation lies in this two-dimensional kernel.
The common fibre class is nonzero, since its intersection
with an ample divisor of the projective surface $M$ is
positive. Thus the only kernel combination which is zero in
homology is the difference of the two fibres. In particular
these ten curves span a nine-dimensional subspace of
$H_2(M,\Q)$.

\noindent\textbf{Step 16.} Here is the complete constant-sheaf comparison. If $S^{[p]}$
denotes the disjoint union of the $(p+1)$-fold intersections
listed in \eqref{cdp:eq:C.EN.S1}, the alternating restriction maps give
an exact complex
\[
 0\longrightarrow\Q_S\longrightarrow\Q_{S^{[0]}}
 \longrightarrow\Q_{S^{[1]}}
 \longrightarrow\Q_{S^{[2]}}\longrightarrow0.           \tag{C.EN.S3}\label{cdp:eq:C.EN.S3}
\]
At a point this is the augmented cochain complex of the
simplex formed by the branches through that point, proving
its exactness. Its cohomology spectral sequence has
$E_1^{p,q}=H^q(S^{[p]},\Q)$.
In degree zero the dimensions are $8,16,8$. The first
map has rank seven, because the incidence graph is connected.
The second has rank eight: the boundary of each triple
point contains its own edge $U_{i\epsilon}$, which occurs
in no other triangle. Consequently its degree-zero cohomology
has dimensions $(1,1,0)$ in columns $0,1,2$.

\noindent\textbf{Step 17.} Only the two Hopf components contribute to $H^1(S^{[0]},\Q)$,
of total dimension two; only $J_a,J_b$ contribute to
$H^1(S^{[1]},\Q)$, of total dimension four.
Each Hopf restriction has rank one. In the actual presentation
$(\C^2\setminus0)/q^{2\Z}$, inclusion of its elliptic
attaching curve preserves the radial deck generator and
kills its angular generator. The two images lie in separate
elliptic summands. This map therefore has rank two.

\noindent\textbf{Step 18.} The degree-two restriction map has dimensions $20\to16$
and rank fourteen. We verify its rank through the dual map
on fundamental curve classes. Give coefficients
$x_0,x_\infty$ to $D_0,D_\infty$, coefficients $y_a,y_b$
to $J_a,J_b$, coefficients $z_i$ to $Z_i$, and coefficients
$u_{i0},u_{i\infty}$ to $U_{i0},U_{i\infty}$, orienting
edges from the first component to the second in \eqref{cdp:eq:C.EN.S1}.
The $E_i$ equations force
$z_i=0$ and $u_{i\infty}=-u_{i0}$; the $A$ equation
forces $x_\infty=-x_0$. Equation \eqref{cdp:eq:C.EN.S2} and the unique
rational-curve relation then force
\[
 y_a+y_b=0,\qquad u_{i0}=x_0/2\quad(1\le i\le4).
                                                               \tag{C.EN.S4}\label{cdp:eq:C.EN.S4}
\]
For completeness, the common half-fibre class is independent
of the four-dimensional span of the differences
$D_0-D_\infty,U_{i0}-U_{i\infty}$: in the curve vector
space modulo the sole relation $F_0-F_\infty$, a multiple
of $F_0$ cannot be a difference vector unless that multiple
is zero. Thus \eqref{cdp:eq:C.EN.S4} lists all kernel elements. Its two
parameters are the relation $J_a-J_b$ and the relation
$2D_0-2D_\infty+\sum_i(U_{i0}-U_{i\infty})$.
This proves rank fourteen.

\noindent\textbf{Step 19.} Finally the components contribute dimension two in degree
three, from the Hopf components, and dimension eight in
degree four. Intersections have no cohomology in those
degrees. The sole possible higher differential with nonzero
target would land in $E_2^{2,0}$, which was proved zero;
all groups $E_2^{2,q}$ for $q>0$ vanish because triple
intersections are points. The preceding actual maps give
\[
 \bigl(b_i(S)\bigr)_{i=0}^4=(1,1,8,4,8),\qquad e(S)=12.
                                                               \tag{C.EN.S5}\label{cdp:eq:C.EN.S5}
\]

\noindent\textbf{Step 20.} The identical alternating complex with structure sheaves
is exact in the same normal crossing charts. All components
have $H^2(\OO)=0$; the only nonzero component $H^1(\OO)$
groups are those of the two Hopf surfaces, one-dimensional
each. Their restriction to $H^1(J_j,\OO_{J_j})$ is an
isomorphism. One can see this through the actual Hopf
elliptic fibration over $\PP^1$: its transition maps on the
fixed elliptic curve are translations, so its $R^1\OO$
is the trivial line; the Leray edge map is an isomorphism
and evaluation at any fibre is an isomorphism. The degree-zero
incidence calculation is unchanged. Hence
\[
 \bigl(h^i(S,\OO_S)\bigr)_{i=0}^2=(1,1,0).
                                                               \tag{C.EN.S6}\label{cdp:eq:C.EN.S6}
\]

\noindent\textbf{Step 21.} The six negative balancing blocks may be placed away from
$M$, keeping the Enriques surface itself untouched. For
each point or curve blowup, the support is the full inverse
image. The relative-pair comparison and the restrictions
of the ambient exceptional Gysin classes give exactly
(C.KM.S2); each negative block adds $(2,6,2)$ in degrees
two, three, and four. For coherent cohomology the chosen
centre has generic reduced-support multiplicity one, and
the displayed blowup charts give $\sigma^*S=S'$ as
Cartier divisors, including the four crossing points of
the quartic step. Applying the projection formula to the
Cartier sequence, with $R\sigma_*\OO=\OO$, gives
$R\sigma_*\OO_{S'}=\OO_S$ as in (C.KM.S3).
Thus the balanced support $S_{\rm bal}$ has
\[
 \bigl(b_i(S_{\rm bal})\bigr)_{i=0}^4=(1,1,20,40,20),
 \qquad
 \bigl(h^i(S_{\rm bal},\OO_{S_{\rm bal}})\bigr)_{i=0}^2
       =(1,1,0),\qquad e(S_{\rm bal})=0.                 \tag{C.EN.S7}\label{cdp:eq:C.EN.S7}
\]
The complete incidence matrices, curve-class map, and
their exact ranks are executed in
\path{code/A_enriques_support.py} and retained in its
\texttt{.out} file.
\end{proof}






\begin{proposition}[A fixed punctured torus family has infinitely many zero-Euler fillings with essential projective elliptic components of coherent Euler characteristic one]\label{cdp:st:C.2.UN.27}
\StatementRecord{cdp:st:C.2.UN.27}{proved}{Historical events: Claude Fable 5.2 (Anthropic), effort high, 16 November 2026: carried as background without fault found. GPT-6 Astra (OpenAI), ChatGPT Pro mode, effort Pro, S6\_09\_PAPERS, 17 November 2026: targeted proof check. Also checked in this writing. Pre-continuation writing events: GPT-6 Astra (OpenAI), ChatGPT Pro mode, effort Pro, S6\_10\_PAPERS, 18 November 2026: targeted proof check. Continuation metadata events: GPT-6 Astra (OpenAI), ChatGPT Pro mode, effort Pro, S6\_10\_PAPERS continuation, 19 November 2026: used and rechecked in this run; status-vocabulary reconciliation; metadata reconciliation only; no new mathematical check. \emph{Pre-continuation recorded proof scope.} The odd and even lift signs, all sixteen fixed points in eight quotient orbits, the exact invariant two-form and its resolution divisor were re-derived. The direct-image floor formula proves every plurigenus and minimality. The parameterized incidence ranks, half-fibre relation and six balancing operations give all support-cohomology rows. The product coordinate proves equality of the specified punctured families for all r}{the hypotheses, notation and definitions stated here; \cref{cdp:st:C.2.UN.25}; \cref{cdp:eq:C.EN.R7}; \cref{cdp:eq:C.EN.R10}; \cref{cdp:eq:C.EN.R11}; \cref{cdp:eq:C.EN.R8}; \cref{cdp:st:C.2.UN.4}; \cref{cdp:eq:C.EN.G1}; \cref{cdp:st:C.2.UN.14}}

For every integer $r\ge2$, the compact construction
of \ref{cdp:st:C.2.UN.25} has a version whose distinguished smooth
component $M_r$ is a minimal projective elliptic
surface with exactly two fibres of type $I_0^*$
and $r$ double smooth elliptic fibres $2J_1,\ldots,2J_r$.
Writing $F$ for a full fibre, its actual canonical
line and plurigenera are
\[
 \begin{gathered}
 K_{M_r}\simeq\OO_{M_r}\left(\sum_{j=1}^rJ_j-F\right),
 \qquad
 2K_{M_r}\simeq\OO_{M_r}((r-2)F),
 \\
 P_m(M_r)=\max\{0,\,1-m+r\lfloor m/2\rfloor\}.
 \end{gathered}
                                                               \tag{C.EN.R1}\label{cdp:eq:C.EN.R1}
\]
It has
\[
 (b_0,b_1,b_2,b_3,b_4)=(1,0,10,0,1),\qquad
 (h^0(\OO),h^1(\OO),h^2(\OO))=(1,0,0).
                                                               \tag{C.EN.R2}\label{cdp:eq:C.EN.R2}
\]
The case $r=2$ is the Enriques component of \ref{cdp:st:C.2.UN.25}.
Every $r\ge3$ gives Kodaira dimension one and
coherent Euler characteristic one.

Before balancing, the full central fibre is
\[
 2M_r+2A+4\sum_{j=1}^r Q_j+\sum_{i=1}^4E_i,
                                                               \tag{C.EN.R3}\label{cdp:eq:C.EN.R3}
\]
with $A\simeq\PP^1\times\PP^1$, primary Hopf
$Q_j$, and $E_i\simeq\mathbb F_2$.
Its reduced support has Euler number $12$.
Six of the explicit negative-two operations,
all disjoint from $M_r$, make that number zero.
The reduced supports before and after these
operations have respectively
\[
 \begin{array}{c|c|c}
 &(b_0,b_1,b_2,b_3,b_4)&(h^0(\OO),h^1(\OO),h^2(\OO))\\ \hline
 S_r&(1,1,r+6,2r,r+6)&(1,1,0)\\
 S_{r,\mathrm{bal}}&(1,1,r+18,2r+36,r+18)&(1,1,0).
 \end{array}                                             \tag{C.EN.R4}\label{cdp:eq:C.EN.R4}
\]
All these fillings have the same specified punctured
torus family. Their distinguished essential components
are pairwise non-bimeromorphic, with canonical growth
slope $r/2-1$.
\end{proposition}


\begin{proof}
\noindent\textbf{Step 1.} Use exactly the compact Weierstrass base, quadratic
base change, resolution, and involution of
(C.EN.2)--(C.EN.5). On its exceptional $R_1$,
blow up the $2r$ points
$z=1,\ldots,r,-1,\ldots,-r$.
Replace (C.EN.6) by
\[
 D_r=\sum_{j=1}^r(D_j-D_{-j}),\qquad L_r=\OO(D_r).
                                                               \tag{C.EN.R5}\label{cdp:eq:C.EN.R5}
\]
The canonical section and the involution give
$h^*L_r=L_r^{-1}$ with its actual order-two
section-preserving dualization. The two signed
root charts still have order two, and their
free $q$-quotient and properness are exactly
(C.EN.7)--(C.EN.8). The formula
$\ell=\sigma c^2,\ a=c^{-1}$ supplies an
actual product identification away from the
central divisor, and the involution there
is always $(v,[c])\mapsto(hv,[c^{-1}])$.
It is independent of $r$. This proves the
assertion about the fixed punctured family.

\noindent\textbf{Step 2.} On the strict transform of $R_1$, use
\[
 \phi_r(z)=\prod_{j=1}^r\frac{z-j}{z+j},\qquad
 C_r:\quad Y^2=\prod_{j=1}^r(z^2-j^2),\qquad
 y=\frac{Y}{\prod_{j=1}^r(z+j)}.
                                                               \tag{C.EN.R6}\label{cdp:eq:C.EN.R6}
\]
The $2r$ branch points are distinct, so $C_r$
is smooth and has genus $r-1$, by the
degree-two Riemann--Hurwitz formula.
The actual finite \'etale product cover
of the corresponding component is again
$\ell=u^2,\ a=y/u$ from $C_r\times E$,
with $E=\C^*/q^{\Z}$. Its deck involution
and the lifted dualization are
\[
 \begin{split}
 \eta(z,Y,[u])&=(z,-Y,[-u]),\\
 k(z,Y,[u])&=(-z,(-1)^rY,[u^{-1}]).
 \end{split}                                             \tag{C.EN.R7}\label{cdp:eq:C.EN.R7}
\]
They commute and generate a Klein four group.
The first action is free. The second formula
is forced by the actual identity
$\phi_r(-z)=\phi_r(z)^{-1}$, so it includes
the sign for odd $r$.

\noindent\textbf{Step 3.} We check both parities at the fixed base values.
At infinity, $Y/z^r=\pm1$, and $k_C$
fixes the two points while $jk_C$ exchanges
them. At zero, the two values of $Y$ are
nonzero. When $r$ is even, $k_C$ fixes
them and $jk_C$ exchanges them; when $r$
is odd, the roles are reversed. Therefore
\[
 \begin{array}{c|cc}
 &\#\operatorname{Fix}(k_C)&
       \#\operatorname{Fix}(jk_C)\\ \hline
 r\text{ even}&4&0\\
 r\text{ odd}&2&2 .
 \end{array}                                             \tag{C.EN.R8}\label{cdp:eq:C.EN.R8}
\]
Each nontrivial fixed-point action on the
smooth curve has derivative $-1$.
On $E$, both $[u]\mapsto[u^{-1}]$
and $[u]\mapsto[-u^{-1}]$ have four
fixed points and derivative $-1$.
Consequently the two elements $k,\eta k$
have a total of sixteen fixed points,
all with local representation $(-1,-1)$.
Their sets are disjoint because $\eta$
is free. Dividing by the remaining free
involution pairs them into eight quotient
points. Thus
\[
 X_r=(C_r\times E)/\langle\eta,k\rangle
                                                               \tag{C.EN.R9}\label{cdp:eq:C.EN.R9}
\]
has precisely eight $A_1$ singularities,
and no branch curve. Let $M_r$ be its
minimal $A_1$ resolution. The same
ambient charts (C.EN.14)--(C.EN.15)
realize that exact resolution as a
central component. For odd $r$ the
fixed lifts over zero are the
$\eta k$ lifts; downstairs the local
involution, the four fixed curves,
and all monomials of those charts
are unchanged.

\noindent\textbf{Step 4.} The map $w=z^2$ gives the elliptic map
$p:M_r\to\PP^1_w$. Above a general $w$,
the four points of $C_r$ are acted on
transitively by $\langle\eta,k\rangle$,
so its fibre is exactly $E$. At $w=j^2$,
the local double-root chart and the free
translation in $\eta$ give a smooth
double fibre $2J_j$. At zero and infinity,
the involution has four fixed points on
the elliptic fibre; its $A_1$ resolutions
give the full divisor
$2D+\sum_{i=1}^4 U_i$ of type $I_0^*$,
as in (C.EN.S2). These are all special
values. In particular the morphism $p$
has connected fibres.

\noindent\textbf{Step 5.} We compute its canonical line using an
actual invariant meromorphic form,
instead of deducing it from numerical
invariants. On $C_r\times E$ use
\[
                       \Omega=dz\wedge\frac{du}{u}.
                                                               \tag{C.EN.R10}\label{cdp:eq:C.EN.R10}
\]
It is fixed by both generators in
\eqref{cdp:eq:C.EN.R7}. On $C_r$, $dz$ has a simple
zero at each of the $2r$ branch points
and a double pole at each of the two
points over infinity, and no other
zero or pole. The quotient is free
in codimension one. The root charts
therefore give one zero along every
$J_j$. At infinity, the local invariant
coordinate is $w^{-1}$, whose divisor
pulls back with order two at both
points of $C_r$.

\noindent\textbf{Step 6.} The resolution comparison at the isolated
quotient points retains these orders.
In local coordinates $(s,v)$ where the
involution is $(-s,-v)$, its crepant
volume form satisfies
$ds\wedge dv=\tfrac12\,dr\wedge d\delta$
in the chart $s^2=r,\ sv=r\delta$.
Thus the double pole $s^{-2}$ of
\eqref{cdp:eq:C.EN.R10} becomes a simple pole
on the exceptional curve, and in the
other chart has orders one and two
on the exceptional and strict central
curves. These are exactly the
multiplicities of the full $I_0^*$
fibre. At zero the form is regular
and nowhere zero, and its crepant
extension is likewise nowhere zero.
Consequently the descended form has
the exact divisor
\[
                \operatorname{div}(\Omega)
                         =\sum_{j=1}^r J_j-F_\infty.
                                                               \tag{C.EN.R11}\label{cdp:eq:C.EN.R11}
\]
This proves the first formula of
\eqref{cdp:eq:C.EN.R1}. Since $2J_j$ is linearly
equivalent to a full fibre, it also
proves its second formula.

\noindent\textbf{Step 7.} Every meromorphic section of $mK_{M_r}$
is a meromorphic function times
$\Omega^{\otimes m}$. A holomorphic
such section restricts to a constant
multiple on a general elliptic fibre,
because that fibre's canonical line
is trivial and $\Omega$ has no zero
or pole there. Its ratio is therefore
a function from the connected-fibre
base. The exact allowed orders at
the double fibres in \eqref{cdp:eq:C.EN.R11}
give the sheaf equality
\[
 p_*\OO_{M_r}(mK_{M_r})
  \simeq
 \OO_{\PP^1}\left(
      \sum_{j=1}^r\lfloor m/2\rfloor[j^2]
                       -m[\infty]\right).
                                                               \tag{C.EN.R12}\label{cdp:eq:C.EN.R12}
\]
The ordinary full fibre at infinity,
including all its components, is
Cartier and contributes exactly the
displayed $-m$. Taking sections gives
the plurigenus formula in \eqref{cdp:eq:C.EN.R1}.
For $r=2$, it gives $P_m=1$ for even
$m$ and zero for odd $m$, recovering
\ref{cdp:st:C.2.UN.25}. For $r\ge3$, its linear growth
has positive slope $r/2-1$, so
$\kappa(M_r)=1$. The line $2K_{M_r}$
is a nonnegative multiple of a fibre
and is nef. Adjunction therefore
excludes a $(-1)$-curve, proving that
$M_r$ is minimal for all $r\ge2$.

\noindent\textbf{Step 8.} The component is projective because
it is the resolution of the finite
quotient of the projective surface
$C_r\times E$. Its coherent groups
can also be read through that exact
quotient. On $H^1(C_r,\OO)$ the
hyperelliptic involution $j$ is
$-I$; this follows directly from
the basis $z^a dz/Y$, $0\le a\le r-2$,
of canonical forms and duality.
The action $\eta$ therefore removes
this summand, whereas $k$ removes
the remaining elliptic summand.
Thus $q(M_r)=0$. The $A_1$ resolution
does not change coherent cohomology:
the explicit exceptional $(-2)$-curve
has positive successive conormal
powers, whose $H^1$ vanish, so
formal functions gives $R^1\pi_*\OO=0$.
The formula at $m=1$ gives $p_g=0$.
The finite quotient Euler formula is
$(0+0+16)/4=4$, for both parity
patterns in \eqref{cdp:eq:C.EN.R8}. Resolving
eight $A_1$ points adds eight,
giving $e(M_r)=12$. Hodge symmetry
and Poincar\'e duality on this
projective surface now yield
\eqref{cdp:eq:C.EN.R2}.

\noindent\textbf{Step 9.} The ambient fixed curves over zero
are still the four copies of $R_0$,
and their normal bundles are still
$\OO(-2)\oplus\OO$. The $2r$
root Hopf components are paired
by the involution into $r$ components.
The exact quotient and resolution
monomials give \eqref{cdp:eq:C.EN.R3}, with no
parity change in its coefficients.
The incidence data of \ref{cdp:st:C.2.UN.25} are
unchanged except for the $r$
elliptic attaching curves $M_r\cap Q_j$.
Those curves all have Euler number
zero. Hence the same calculation
gives $e(S_r)=12$, and six \ref{cdp:st:C.2.UN.4}
operations away from $M_r$ give
zero while preserving $M_r$ itself.
The normal crossing and compactness
arguments are the actual charts
of \ref{cdp:st:C.2.UN.25}; no new analytic gluing
is required.

\par\medskip\noindent\textbf{Cohomology of all the paired-root supports.}\par\noindent
The support calculation (C.EN.S1)--(C.EN.S7) extends through
the same maps, with $r$ Hopf components instead of two.
We give the dimensions and the changed maps explicitly.
There are $r+6$ components: the projective elliptic surface
$M$, $A\simeq\PP^1\times\PP^1$, four surfaces
$E_i\simeq\mathbb F_2$, and $Q_1,\ldots,Q_r$.
The double curves are the fourteen rational curves
$D_0,D_\infty,Z_i,U_{i0},U_{i\infty}$ already listed
in (C.EN.S1), and the $r$ elliptic curves
$J_j=M\cap Q_j$. The same eight triple points remain.
On $M$, the actual elliptic fibration gives
\[
 [2D_0+\textstyle\sum_iU_{i0}]
 =[2D_\infty+\textstyle\sum_iU_{i\infty}]
 =2[J_j]\quad(1\le j\le r).
                                                               \tag{C.EN.G1}\label{cdp:eq:C.EN.G1}
\]
Its two $I_0^*$ configurations still span dimension nine
in $H_2(M,\Q)$, by the same affine-$D_4$ intersection
calculation and nonzero fibre class. Here $b_2(M)=10$,
$b_1(M)=b_3(M)=0$, and $h^1(M,\OO_M)=h^2(M,\OO_M)=0$
by the invariant and resolution computations above.

\noindent\textbf{Step 10.} For the exact constant-sheaf complex (C.EN.S3), the
degree-zero row now has dimensions
$r+6,r+14,8$. Its two ranks are $r+5$ and eight:
each extra Hopf vertex and attaching edge is a leaf,
and each of the eight triangles retains its private
edge $U_{i\epsilon}$. Thus its cohomology is still
$(1,1,0)$. In degree one the map has dimensions
$r\to2r$ and rank $r$, since each Hopf radial
generator restricts nontrivially to its own elliptic
attaching curve. In particular this row contributes
dimension $r$ to $H^2$ of the support.

\noindent\textbf{Step 11.} The degree-two restriction map has dimensions
$20\to r+14$ and rank fourteen. The $E_i$ and $A$
equations in (C.EN.S4) are unchanged. The $M$
equation, using \eqref{cdp:eq:C.EN.G1}, becomes
\[
 \sum_{j=1}^r y_j=0,\qquad
 u_{i0}=x_0/2,\quad u_{i\infty}=-x_0/2,\quad
 x_\infty=-x_0,\quad z_i=0.
                                                               \tag{C.EN.G2}\label{cdp:eq:C.EN.G2}
\]
These are all solutions by the already proved
independence of the common half-fibre and the
rational-curve difference subspace. The kernel
therefore has dimension $(r-1)+1=r$, proving
rank fourteen. The component degree-three and
degree-four groups have dimensions $r$ and $r+6$,
respectively. As before, every possible higher
differential has zero target. Consequently
\[
 \bigl(b_i(S_r)\bigr)_{i=0}^4=(1,1,r+6,2r,r+6),
 \qquad e(S_r)=12.
                                                               \tag{C.EN.G3}\label{cdp:eq:C.EN.G3}
\]
For structure sheaves the $r$ Hopf-to-elliptic
restrictions are isomorphisms, and every component
has zero $H^2(\OO)$. The same exact complex gives
$(h^i(S_r,\OO_{S_r}))_{i=0}^2=(1,1,0)$.
Applying six negative blocks away from $M$, with
the support pullback/Gysin maps and derived Cartier
comparison used in (C.EN.S7), gives
\[
 \begin{gathered}
 \bigl(b_i(S_{r,\mathrm{bal}})\bigr)_{i=0}^4
      =(1,1,r+18,2r+36,r+18),\qquad e(S_{r,\mathrm{bal}})=0,\\
 \bigl(h^i(S_{r,\mathrm{bal}},\OO_{S_{r,\mathrm{bal}}})\bigr)_{i=0}^2
      =(1,1,0).
 \end{gathered}                                                \tag{C.EN.G4}\label{cdp:eq:C.EN.G4}
\]
The parameterized incidence and curve-class matrices
are included in \path{code/A_enriques_support.py}; its
executed output checks both the original case and
the first two properly elliptic cases.

\noindent\textbf{Step 12.} Finally a covering family of rational
curves would kill every positive
pluricanonical section by its actual
dominant ruled parametrization.
Thus $M_r$ is non-uniruled, and
\ref{cdp:st:C.2.UN.14} gives its essentiality under
proper birational comparison of
smooth fillings. The canonical
growth slopes in \eqref{cdp:eq:C.EN.R1} are
distinct, so these components
are pairwise non-bimeromorphic.
For $r\ge3$ each is the unique
component of Kodaira dimension
one; the other original components
are rational or Hopf, and the
additional blowup components are
ruled. The same inverse-map
comparison therefore also
distinguishes the corresponding
fillings over the fixed punctured
family.
\end{proof}






\begin{proposition}[The fixed Kummer punctured family has zero-Euler fillings with arbitrarily large canonical growth on a projective elliptic component]\label{cdp:st:C.2.UN.30}
\StatementRecord{cdp:st:C.2.UN.30}{proved}{Historical events: Claude Fable 5.2 (Anthropic), effort high, 16 November 2026: carried as background without fault found. GPT-6 Astra (OpenAI), ChatGPT Pro mode, effort Pro, S6\_09\_PAPERS, 17 November 2026: targeted proof check. Also checked in this writing. Pre-continuation writing events: GPT-6 Astra (OpenAI), ChatGPT Pro mode, effort Pro, S6\_10\_PAPERS, 18 November 2026: targeted proof check. Continuation metadata events: GPT-6 Astra (OpenAI), ChatGPT Pro mode, effort Pro, S6\_10\_PAPERS continuation, 19 November 2026: used and rechecked in this run; status-vocabulary reconciliation; metadata reconciliation only; no new mathematical check. \emph{Pre-continuation recorded proof scope.} The torsion square-root product cover, sixteen surface fixed points, thirty-two ambient quotient points and their exact T=R*u\textasciicircum{}2 resolution were re-derived. The invariant two-form has exactly the stated multiple-fibre divisor, giving its line identities and plurigenera. The star-support curve-class map has rank seventeen, giving all Betti and coherent rows. Twenty actual balancing blocks at each special value preserve the distinguished surfaces and punctured family}{the hypotheses, notation and definitions stated here; \cref{cdp:st:C.2.UN.18}; \cref{cdp:st:C.2.UN.23}; \cref{cdp:st:C.2.UN.13}; \cref{cdp:st:C.2.UN.25}; \cref{cdp:eq:C.ET.11}; \cref{cdp:st:C.2.UN.14}}

Fix the square elliptic curve
$E_0=\C/(\Z+i\Z)$ and $E=\C^*/q^{\Z}$,
where $q=e^{-2\pi}$. For every $r\ge1$ there is
a smooth compact threefold over $\PP^1$,
with the same punctured product-torus family
as the diagonal involution quotient
\[
 (E_0\times E\times\PP^1_t)/
      \langle(p,[c],t)\mapsto(-p,[c^{-1}],-t)\rangle,
                                                               \tag{C.ET.1}\label{cdp:eq:C.ET.1}
\]
whose two special reduced supports have Euler
number zero. The support over zero contains an
unchanged smooth minimal projective properly
elliptic surface $M_r$ with four $I_0^*$ fibres
and $r$ double smooth elliptic fibres $2J_j$.
For its connected elliptic map $p_r:M_r\to\PP^1$
and a full fibre $F$,
\[
 \begin{aligned}
 K_{M_r}&\simeq\OO_{M_r}\left(\sum_{j=1}^rJ_j\right),
 \qquad 2K_{M_r}\simeq\OO_{M_r}(rF),\\
 (p_r)_*\OO(mK_{M_r})
   &\simeq\OO_{\PP^1}\left(
                \sum_{j=1}^r\lfloor m/2\rfloor[b_j]\right),
 \\
 P_m(M_r)&=1+r\lfloor m/2\rfloor .
 \end{aligned}                                             \tag{C.ET.2}\label{cdp:eq:C.ET.2}
\]
Here $b_j$ is the value of $2J_j$. In particular
$\chi(\OO_{M_r})=2$, $p_g(M_r)=1$, $q(M_r)=0$,
and
\[
            (b_0,b_1,b_2,b_3,b_4)(M_r)
                         =(1,0,22,0,1).              \tag{C.ET.3}\label{cdp:eq:C.ET.3}
\]
These essential components, and the associated
fillings over the fixed punctured family, are
pairwise non-bimeromorphic.

Before balancing, the full fibre over zero is
\[
                  2M_r+4\sum_{j=1}^rQ_j
                           +\sum_{\nu=1}^{16}P_\nu,
                                                               \tag{C.ET.4}\label{cdp:eq:C.ET.4}
\]
where the $Q_j$ are primary Hopf surfaces and
$P_\nu\simeq\PP^2$. It has Euler number $40$.
Its reduced support and its balanced version
have respectively
\[
 \begin{array}{c|c|c}
 &(b_0,b_1,b_2,b_3,b_4)&
                 (h^0(\OO),h^1(\OO),h^2(\OO))\\ \hline
 S_r&(1,0,21+r,2r-1,17+r)&(1,0,1)\\
 S_{r,\mathrm{bal}}&(1,0,61+r,2r+119,57+r)&(1,0,1).
 \end{array}                                             \tag{C.ET.5}\label{cdp:eq:C.ET.5}
\]
\end{proposition}


\begin{proof}
\noindent\textbf{Step 1.} Choose an odd prime $N>2r$ and the torsion
points $p_j=j/N$ of $E_0$, for $1\le j\le r$.
Their pairs $\{p_j,-p_j\}$ are disjoint and
avoid the two-torsion points. Blow up
$E_0\times\PP^1_t$ at $(p_j,0),(-p_j,0)$,
and call the exceptional curves $R_j,R_{-j}$.
The involution $(p,t)\mapsto(-p,-t)$ lifts
to the resulting surface $V$, exchanging
the paired curves. Put
\[
       D=\sum_{j=1}^r(R_j-R_{-j}),\qquad L=\OO_V(D).
                                                               \tag{C.ET.6}\label{cdp:eq:C.ET.6}
\]
The canonical meromorphic section $\sigma$
gives $h^*L=L^{-1}$ with its order-two
section-preserving dualization.

\noindent\textbf{Step 2.} Perform the signed order-two root construction
of \ref{cdp:st:C.2.UN.18} along these disjoint curves:
$\sigma=a^2\ell$ on positive charts and
$\sigma^{-1}=(a')^2\ell'$ on negative charts,
with $a'=a^{-1},\ell'=\ell^{-1}$.
Divide by $(\ell,a)\mapsto(q^2\ell,q^{-1}a)$.
The explicit finite free charts (C.PE.8)
prove that the resulting compact threefold
$\widetilde Y$ is smooth and proper over $V$.
Dualization over $h$ takes the global
meromorphic pair $(\ell,a)$ to
$(\ell^{-1},a^{-1})$ and conjugates the
$q$-action to its inverse. It therefore
defines an actual involution of
$\widetilde Y$.

\noindent\textbf{Step 3.} Away from the root curves the formula
$\ell=\sigma c^2,\ a=c^{-1}$ identifies the
family with $V\times E$. Under that actual
identification the involution is
$(v,[c])\mapsto(hv,[c^{-1}])$.
All modifications lie over $t=0$.
Consequently the quotient family over
$T=t^2$, on the complement of zero and
infinity, is exactly the punctured family
in \eqref{cdp:eq:C.ET.1}, independently of $r$.
The quotient and every resolution are proper over
$\PP^1$; the Stein factor has generic degree one
and is therefore the identity finite map to the
normal base. All fibres are connected.

\noindent\textbf{Step 4.} Let $B_r$ be the component above the
strict central elliptic curve $E_0$.
It is an elliptic surface over $E_0$
with $2r$ smooth double fibres, at
$p_j,-p_j$. It is projective; we
give the actual finite product cover
which proves this. The restricted line
\[
 L|_{E_0}\simeq
       \OO_{E_0}\left(\sum_{j=1}^r[p_j]-[-p_j]\right)
                                                               \tag{C.ET.7}\label{cdp:eq:C.ET.7}
\]
is torsion because all $p_j$ are torsion.
Choose a torsion square root $\mathcal N$
of this line, and form the normalization
$C$ of $y^2=\sigma|_{E_0}$ in the
corresponding line-bundle compactification.
It is a connected double cover of $E_0$
branched simply at the $2r$ specified
points, so $g(C)=r+1$.
Choose a finite \'etale cover $C'\to C$
trivializing the pulled-back torsion
line $\mathcal N$, and a nowhere-zero
frame $n$ there. Write $y=y_0n$.
Then
\[
 C'\times E\longrightarrow B_r,\qquad
 (z,[u])\longmapsto
             [\,\ell=u^2n^2,\ a=y_0/u\,]             \tag{C.ET.8}\label{cdp:eq:C.ET.8}
\]
is the exact finite \'etale map of \ref{cdp:st:C.2.UN.18}.
At a positive root its coordinate Jacobian
is $2$; in the negative chart it is
$-2u^{-2}$. Thus it is \'etale also
over the multiple fibres. The finite
product-cover comparison proves
projectivity of $B_r$, and hence of
the finite quotient and its surface
resolution used below.

\noindent\textbf{Step 5.} The involution on $B_r$ covers inversion
on $E_0$. At each of the four two-torsion
base points, the root curves are absent,
and the actual product formula gives
inversion on the elliptic fibre.
There are exactly sixteen fixed points
on $B_r$, each with derivative $(-1,-1)$.
There is no fixed curve. Its quotient
has sixteen $A_1$ singularities; let
$M_r$ be their minimal resolution.
The induced base is
$E_0/\{\pm1\}\simeq\PP^1$.
The four fixed base values give $I_0^*$
fibres by the same explicit resolution
as (C.EN.S2). The paired root values
give the $r$ smooth double fibres
$2J_j$. Away from those values the
quotient is a smooth elliptic fibration,
so this is its complete fibre list.

\noindent\textbf{Step 6.} Let $\omega_0$ be a nonzero holomorphic
one-form on $E_0$. On $B_r$ the
meromorphic expression
\[
                        \Omega=\omega_0\wedge d\log\ell
                                                               \tag{C.ET.9}\label{cdp:eq:C.ET.9}
\]
is global. A change of frame of $L$
adds a base logarithmic differential
to $d\log\ell$; its wedge with
$\omega_0$ is zero. Multiplication
by the constant $q^2$ also leaves
the expression unchanged. In a
positive root chart $z=a^2b$ it is
a unit times $2a\,da\wedge db$;
the negative dual chart has the
opposite sign and the same simple
zero. Thus its divisor on $B_r$
is precisely the sum of the
$2r$ reduced multiple fibres.
The involution changes both
$\omega_0$ and $d\log\ell$ by a
minus sign, up to the harmless
base differential already described.
Hence $\Omega$ is invariant.

\noindent\textbf{Step 7.} At its sixteen fixed points the form
is regular and nowhere zero. The
height-one $A_1$ charts of \ref{cdp:st:C.2.UN.23}
extend it without a zero or pole.
Paired root fibres are exchanged
freely and become $J_j$. The
descended form consequently has
the exact divisor
\[
                      \operatorname{div}(\Omega)
                                      =\sum_{j=1}^rJ_j .
                                                               \tag{C.ET.10}\label{cdp:eq:C.ET.10}
\]
This proves the first two line-bundle
identities in \eqref{cdp:eq:C.ET.2}.
The ratio of any holomorphic
$m$-canonical section to
$\Omega^{\otimes m}$ is constant
on a connected general elliptic
fibre. It descends meromorphically
to the base. Its allowed pole
order at $b_j$ is exactly
$\lfloor m/2\rfloor$, because
pullback of a base parameter
has divisor $2J_j$; there are
no other allowed poles.
This proves the actual direct-image
equality in \eqref{cdp:eq:C.ET.2}, and hence
all its plurigenera.
Their growth proves $\kappa(M_r)=1$.
The line $2K_{M_r}=rF$ is nef,
so adjunction excludes a rational
$(-1)$-curve; $M_r$ is minimal.

\noindent\textbf{Step 8.} The surface $B_r$ has Euler number
zero: every fibre over its elliptic
base has elliptic reduced support,
and the map is locally a smooth
elliptic bundle off finitely many
values. The finite involution
quotient therefore has Euler
number $(0+16)/2=8$. Resolving
sixteen $A_1$ points gives
$e(M_r)=24$. Its canonical
divisor is a sum of disjoint
square-zero fibres, so
$K_{M_r}^2=0$. Surface
Riemann--Roch gives
$\chi(\OO_{M_r})=24/12=2$.
The first plurigenus in \eqref{cdp:eq:C.ET.2}
is one; hence $q(M_r)=0$.
Projectivity, Hodge symmetry and
Poincar\'e duality give \eqref{cdp:eq:C.ET.3}.

\noindent\textbf{Step 9.} We now specify the ambient resolution.
The only fixed base points of
$h$ lie at $E_0[2]\times\{0,\infty\}$,
and none was blown up. Each has
four fixed elliptic-fibre points.
The thirty-two ambient fixed
points therefore have local
linear action $(-1,-1,-1)$.
The actual invariant quadratic
blowup charts of \ref{cdp:st:C.2.UN.13} resolve
each quotient point smoothly.
For example, write upstairs
$(x,t,w)=(s,su,sv)$ and put
$R=s^2$. The resolved quotient
chart has coordinates $(R,u,v)$
and
\[
                              T=Ru^2 .                \tag{C.ET.11}\label{cdp:eq:C.ET.11}
\]
The exceptional surface is $\PP^2$.
Its intersection with the strict
central component is a line;
restriction is exactly the
minimal $A_1$ surface resolution.
Thus the strict component is
the same smooth $M_r$ proved
above, and the monomial gives
its multiplicity two and the
exceptional multiplicity one.
The paired root Hopf surfaces
give $r$ components of
multiplicity four, exactly as
in \ref{cdp:st:C.2.UN.25}. This proves \eqref{cdp:eq:C.ET.4}
with all full-fibre coefficients.

\noindent\textbf{Step 10.} The reduced support is a star:
each $P_\nu$ meets $M_r$ along
one $\PP^1$, each $Q_j$ meets
$M_r$ along the elliptic curve
$J_j$, and there are no other
intersections or triple points.
It is simple normal crossing
by \eqref{cdp:eq:C.ET.11} and the root charts.
Therefore
\[
                 e(S_r)=24+16\cdot3-16\cdot2=40.
                                                               \tag{C.ET.12}\label{cdp:eq:C.ET.12}
\]
At infinity the same ambient
resolution gives the original
product-Kummer support
$2K+16\PP^2$ of \ref{cdp:st:C.2.UN.13}, also
of Euler number forty.
Twenty negative-two blocks
on each special support,
all beginning at fresh
smooth points of its plane
components away from the
K3 or $M_r$, make both
Euler numbers zero and
leave those distinguished
surfaces unchanged.

\noindent\textbf{Step 11.} For completeness the support
cohomology in \eqref{cdp:eq:C.ET.5} follows
through exact restriction maps.
The constant-sheaf union
sequence has component term
$M_r\sqcup\bigsqcup Q_j
\sqcup\bigsqcup P_\nu$
and curve term
$\bigsqcup J_j\sqcup
\bigsqcup\PP^1_\nu$.
Its degree-zero incidence
graph is a tree. In degree
one the $r$ Hopf generators
map injectively, each onto
the radial one-dimensional
subspace of its elliptic
attaching curve. The cokernel
has dimension $r$.
The degree-two restriction
has source dimension $38$
and target dimension $16+r$.
Its dual curve-class map
has kernel of dimension
$r-1$: the plane components
force each of their sixteen
edge coefficients to vanish;
the remaining elliptic
classes all equal the
nonzero class $\tfrac12F$
on $M_r$ (a full fibre meets
an ample divisor positively), leaving exactly
the relations whose
coefficients sum to zero.
The map thus has rank
seventeen. Its kernel
and cokernel have dimensions
$21$ and $r-1$.
The components contribute
$r$ in degree three and
$17+r$ in degree four.
The resulting long exact
sequence gives the first
Betti row of \eqref{cdp:eq:C.ET.5}.

\noindent\textbf{Step 12.} The analogous structure-sheaf
sequence has isomorphisms
$H^1(Q_j,\OO)\to H^1(J_j,\OO)$,
proved by the actual Hopf
elliptic Leray map in \ref{cdp:st:C.2.UN.25}.
The sole component
$H^2(\OO)$ is the one
on $M_r$. Hence the
coherent dimensions are
$(1,0,1)$.
Each of the twenty
balancing blocks adds
$(2,6,2)$ to support
cohomology in degrees
two, three and four,
by the restricted
exceptional Gysin maps
of (C.KM.S2), and
preserves coherent
cohomology by the
actual Cartier pullback
in (C.KM.S3). This
proves the second row.

\noindent\textbf{Step 13.} The canonical growth slopes
$r/2$ are distinct and
positive. Therefore these
components are non-uniruled,
essential by \ref{cdp:st:C.2.UN.14}, and
pairwise non-bimeromorphic.
They are the unique components
of Kodaira dimension one
in their respective supports;
the other components are
rational, Hopf, K3, or
the ruled blowup components.
The inverse-map resolution
comparison consequently
also distinguishes the
fillings over their same
specified punctured family.
\end{proof}






\begin{proposition}[Every fixed-modulus elliptic quasi-bundle over a rational or elliptic curve occurs unchanged]\label{cdp:st:C.2.UN.34}
\StatementRecord{cdp:st:C.2.UN.34}{proved}{Historical events: Claude Fable 5.2 (Anthropic), effort high, 16 November 2026: carried as background without fault found. GPT-6 Astra (OpenAI), ChatGPT Pro mode, effort Pro, S6\_09\_PAPERS, 17 November 2026: targeted proof check. Also checked in this writing. Pre-continuation writing events: GPT-6 Astra (OpenAI), ChatGPT Pro mode, effort Pro, S6\_10\_PAPERS, 18 November 2026: targeted proof check. Continuation metadata events: GPT-6 Astra (OpenAI), ChatGPT Pro mode, effort Pro, S6\_10\_PAPERS continuation, 19 November 2026: used and rechecked in this run; status-vocabulary reconciliation; metadata reconciliation only; no new mathematical check. \emph{Pre-continuation recorded proof scope.} The exact exponential two-torus quotient has additive kernel O, and H2 of that coherent sheaf on the root curve vanishes by finite averaging and the coarse-curve calculation. The lifted pair of line bundles represents the given full torsor class, including Picard-zero data. Actual signed vertical divisors on the blown-up auxiliary surface extend those bundles and their root characters. The inertia translations have exact order, so the coarse threefold is smooth, the specified component is unchanged, and the central support has Euler number zero}{the hypotheses, notation and definitions stated here; \cref{cdp:st:C.2.UN.22}; \cref{cdp:st:C.2.UN.29}; \cref{cdp:eq:C.SF.4}; \cref{cdp:eq:C.SF.3}; \cref{cdp:eq:C.SF.5}; \cref{cdp:eq:C.SF.7}; \cref{cdp:eq:C.SF.8}; \cref{cdp:eq:C.SF.1}; \cref{cdp:eq:C.SF.9}; \cref{cdp:eq:C.SF.10}; \cref{cdp:eq:C.SF.6}; \cref{cdp:eq:C.SF.12}; \cref{cdp:eq:C.SF.2}; \cref{cdp:eq:C.SF.11}; \cref{cdp:eq:C.SF.13}}

Let $q:S\to B$ be a proper holomorphic map from a
smooth compact connected complex surface to a smooth
compact curve with $g(B)\le1$. Assume that every
fibre is a positive multiple of a smooth elliptic
curve, that its smooth fibres have one fixed elliptic
modulus, and that its integral $H_1$ monodromy is
the identity. Then there are a smooth compact complex
threefold $Y$, a proper holomorphic map
$f:Y\to\mathbb P^1$, and a specified isomorphism
from $S$ to a component of its reduced zero fibre,
with the following properties:

\begin{itemize}
\item the component $S$ is unchanged and has
multiplicity one in the full fibre;
\item over $\mathbb P^1\setminus\{0\}$ the map
is an actual product with a fixed product of two
elliptic curves, and hence has connected torus
fibres and identity monodromy;
\item the reduced zero fibre has simple normal
crossings, all its components are smooth elliptic
quasi-bundles, all its double curves are smooth
elliptic, and it has no triple points;
\item every reduced fibre, including the zero
fibre, has Euler characteristic zero.
\end{itemize}
This realizes every elliptic covering model in
\ref{cdp:st:C.2.UN.22} and \ref{cdp:st:C.2.UN.29}. The prescribed surface
is retained through its full holomorphic torsor
class, including its degree-zero parameters and
its individual multiple-fibre translations.
\end{proposition}



\begin{proof}
\noindent\textbf{Step 1.} Fix the elliptic curve $A=\mathbb C/\Lambda$ of
the smooth fibres, and choose an ordered lattice
basis $\lambda_1,\lambda_2$. Let the multiple
values be $p_i\in B$, with full fibre
$m_iC_i$ and $m_i>1$. Introduce the actual root
curve
\[
                 \mathcal B=B[(p_i,m_i)_i].
                                                               \tag{C.SF.1}\label{cdp:eq:C.SF.1}
\]
Its local chart at $p_i$ is
$[\Delta_w/\mu_{m_i}]\to\Delta_z$,
$z=w^{m_i}$, with the indicated standard action
on $w$. We will use these charts to specify
all coarse-space maps in the construction.
The line and section $\mathcal O_S(C_i),s_{C_i}$,
whose $m_i$th power is the pullback of
$\mathcal O_B(p_i),s_{p_i}$, give the actual
lift of $q$ to this root curve.

\noindent\textbf{Step 2.} First the given surface is the coarse space
of a principal $A$-torsor $\mathcal T$ over
$\mathcal B$. Over the smooth base, choose
local holomorphic sections and integrate a
marked differential. Constant modulus and
identity integral monodromy make every
overlap map a translation of $A$, giving
the actual principal torsor. At a multiple
value, normalize the actual base change
$z=w^{m_i}$ of $q$. The local equations
$z=u h^{m_i}$ become $v^{m_i}=u$, $w=vh$;
the normalized cover is \emph{\'etale} over
$S$ and is a smooth elliptic family. Its
fibres are connected: the proper Stein
factor over the root disc has generic
degree one and is the identity over
that normal disc. Its
deck action is free on the central elliptic
curve and hence is a translation of exact
order $m_i$, denoted $\tau_i\in A[m_i]$.
The translation can be made constant
on the whole root chart. If its product
expression is $a\mapsto a+b(w)$, lift
$b(w)-\tau_i$ to a holomorphic function
$\widetilde b(w)$ with value zero at
the origin. The order-$m_i$ relation
gives $\sum_{k=0}^{m_i-1}\widetilde b(\zeta^kw)=0$.
The coordinate change $a'=a+c(w)$, where
$c(w)=-m_i^{-1}\sum_{k=0}^{m_i-1}
 k\widetilde b(\zeta^kw)$, changes this
translation to exactly $\tau_i$.
The marked product description on the root
disc therefore has the local form
\[
 (A\times\Delta_w)/
 \langle(a,w)\mapsto(a+\tau_i,\zeta_{m_i}w)\rangle.
                                                               \tag{C.SF.2}\label{cdp:eq:C.SF.2}
\]
This is precisely the original multiple-fibre
chart, including its torsor gluing on the
punctured disc, by the normalized-root and
translation comparison in \ref{cdp:st:C.2.UN.29}.
Thus the assertion identifies the given
$S$, rather than only another surface with
the same multiplicities.

\noindent\textbf{Step 3.} We next lift the entire torsor to two
line bundles. The holomorphic group map
\[
 \begin{split}
 \Phi:(\mathbb C^*)^2&\longrightarrow A,\\
 (u,v)&\longmapsto
 \left[\frac{\lambda_1\log u+\lambda_2\log v}{2\pi\mathrm i}
                                                    \right]
 \end{split}
                                                               \tag{C.SF.3}\label{cdp:eq:C.SF.3}
\]
is independent of both logarithm choices.
It is surjective, since its derivative
is surjective and its image is therefore
an open subgroup of the connected group
$A$. Its kernel is the closed copy of
$\mathbb C$ given by
\[
 w\longmapsto
 \bigl(\exp(2\pi\mathrm i\lambda_2w),
                  \exp(-2\pi\mathrm i\lambda_1w)\bigr).
                                                               \tag{C.SF.4}\label{cdp:eq:C.SF.4}
\]
To verify exactness, take logarithms of
a point in the kernel. Their linear
combination lies in $\Lambda$; subtract
the corresponding integral logarithm
periods to make that combination zero.
The resulting pair lies in the line
parametrized by \eqref{cdp:eq:C.SF.4}. Its exponential
map is injective, since an integral
relation between $\lambda_1$ and
$\lambda_2$ is zero. The image is the
kernel of \eqref{cdp:eq:C.SF.3}, hence is closed;
local logarithms give its holomorphic
inverse. This proves the exact group
sequence and the corresponding exact
sequence of analytic sheaves on
$\mathcal B$,
\[
 0\longrightarrow\mathcal O_{\mathcal B}
   \longrightarrow(\mathcal O_{\mathcal B}^*)^2
   \xrightarrow{\Phi}\mathcal A_{\mathcal B}
   \longrightarrow0,
                                                               \tag{C.SF.5}\label{cdp:eq:C.SF.5}
\]
where $\mathcal A$ is the sheaf of
holomorphic maps to the fixed group $A$.

\noindent\textbf{Step 4.} The coarse root-curve map
$\pi:\mathcal B\to B$ has
$\pi_*\mathcal O=\mathcal O_B$ and
no higher coherent direct images.
Locally this is exactly the invariants
functor of a finite group in characteristic
zero; averaging makes it exact. Hence
$H^2(\mathcal B,\mathcal O)=H^2(B,\mathcal O_B)=0$.
The long exact sequence of the abelian
sheaves in \eqref{cdp:eq:C.SF.5} now gives a surjection
\[
 \operatorname{Pic}(\mathcal B)^2
                 \longrightarrow H^1(\mathcal B,\mathcal A).
                                                               \tag{C.SF.6}\label{cdp:eq:C.SF.6}
\]
Choose a pair of line bundles
$\mathcal L_1,\mathcal L_2$ whose
image is the \emph{given} torsor
$\mathcal T$. This equality is an
equality of its analytic $H^1$ class,
so fixes its full holomorphic gluing.

\noindent\textbf{Step 5.} Let $\mathcal R_i$ be the tautological
root line with
$\mathcal R_i^{m_i}=\pi^*\mathcal O_B(p_i)$.
Each line bundle on $\mathcal B$ has
an expression
\[
 \mathcal L_a\simeq\pi^*M_a\otimes
                  \bigotimes_i\mathcal R_i^{e_{ai}},
                  \qquad a=1,2.
                                                               \tag{C.SF.7}\label{cdp:eq:C.SF.7}
\]
Indeed choose $e_{ai}$ to be its
inertia character modulo $m_i$.
After tensoring by the inverse root
characters it has trivial inertia;
its invariant local frames descend
to the ordinary line bundle $M_a$
on the coarse curve. This also
proves \eqref{cdp:eq:C.SF.7} by actual local
frame identifications.

\noindent\textbf{Step 6.} Represent each $M_a$ by an ordinary
signed divisor
\[
                        M_a\simeq\mathcal O_B\left(
                                    \sum_p d_{ap}p\right)
                                                               \tag{C.SF.8}\label{cdp:eq:C.SF.8}
\]
whose support avoids the finitely
many marked points $p_i$. It may
also avoid any one further chosen
point. For an exact choice, take
an effective divisor away from
that finite set, twist $M_a$ by
a sufficiently large multiple,
and use a section which is nonzero
at all points of the finite set;
curve Riemann--Roch and global
generation provide it. Subtracting
the chosen twist gives \eqref{cdp:eq:C.SF.8}.
This represents degree-zero classes
as well as classes of nonzero
degree.

\noindent\textbf{Step 7.} We construct the auxiliary compact
surface. If $B$ is elliptic, take
$V_0=B\times\mathbb P^1$ with its
projection and selected component
$B=B\times\{0\}$. If $B=\mathbb P^1$,
start with any product
$A_0\times\mathbb P^1$ for an
elliptic curve $A_0$, blow up one
point over zero, and identify its
exceptional curve with the prescribed
$B$. Its single intersection with
the strict central elliptic curve
may be chosen away from all $p_i$;
choose \eqref{cdp:eq:C.SF.8} away from it as well.
In both cases the selected component
has multiplicity one and the
auxiliary central divisor has
simple normal crossings.

\noindent\textbf{Step 8.} Blow up the distinct points of
the selected $B$ consisting of
all $p_i$ and all supports in
\eqref{cdp:eq:C.SF.8}. Call the resulting
surface $V$, its map to
$\mathbb P^1$ $h$, and the
exceptional curves $T_i,T_p$.
They are pairwise disjoint and
vertical, each meets the strict
transform of $B$ transversely
at its specified point, and
\[
 \mathcal O_V(T_i)|_B=\mathcal O_B(p_i),
 \qquad\mathcal O_V(T_p)|_B=\mathcal O_B(p).
                                                               \tag{C.SF.9}\label{cdp:eq:C.SF.9}
\]
These are the actual restriction
maps of the exceptional divisors.
Every component of $h^{-1}(0)$
still has multiplicity one; its
reduced support is a tree of
smooth curves with no triple
points. The strict $B$ remains
isomorphic to the given curve.

\noindent\textbf{Step 9.} Root $V$ along the disjoint
$T_i$ with their specified
orders, obtaining
$\mathcal V=V[(T_i,m_i)_i]$.
The inverse image of the strict
$B$ is exactly \eqref{cdp:eq:C.SF.1}, since
the intersections are transverse.
Write $\widetilde{\mathcal R}_i$
for the corresponding root lines
on $\mathcal V$. The vertical
extensions
\[
 \widetilde{\mathcal L}_a
   =\mathcal O_V\left(\sum_p d_{ap}T_p\right)
       \otimes\bigotimes_i
                       \widetilde{\mathcal R}_i^{e_{ai}}
                          \qquad(a=1,2)
                                                               \tag{C.SF.10}\label{cdp:eq:C.SF.10}
\]
restrict, through \eqref{cdp:eq:C.SF.9}, to
exactly $\mathcal L_a$. Here
ordinary line bundles are pulled
back to $\mathcal V$ in the
displayed formula.

\noindent\textbf{Step 10.} Push the pair in \eqref{cdp:eq:C.SF.10}
out by the actual group map
\eqref{cdp:eq:C.SF.3}. Explicitly, if their
nonvanishing local transition
functions are $g^1_{\alpha\beta}$
and $g^2_{\alpha\beta}$, glue
the products with $A$ by
\[
 (x,a)_\beta=(x,a+
     \Phi(g^1_{\alpha\beta}(x),
          g^2_{\alpha\beta}(x)))_\alpha.
                                                               \tag{C.SF.11}\label{cdp:eq:C.SF.11}
\]
The triple-overlap equality follows
from the line-bundle cocycles and
the homomorphism property of $\Phi$.
Local logarithm choices change the
translation only by an element
of $\Lambda$. This produces a
principal $A$-torsor over
$\mathcal V$ whose restriction
is isomorphic to $\mathcal T$
by \eqref{cdp:eq:C.SF.6}--\eqref{cdp:eq:C.SF.10}.

\noindent\textbf{Step 11.} Its coarse space $Y$ is smooth.
In a root chart transverse to
$T_i$ it is precisely
\[
 (A\times\Delta_w\times\Delta_x)/
 \langle(a,w,x)\mapsto
             (a+\tau_i,\zeta_{m_i}w,x)\rangle,
 \qquad z=w^{m_i}.
                                                               \tag{C.SF.12}\label{cdp:eq:C.SF.12}
\]
The translation is the same
$\tau_i$ as in the original
torsor, because the inertia
characters in \eqref{cdp:eq:C.SF.7} are
preserved. Its exact order
is $m_i$, so the finite action
in \eqref{cdp:eq:C.SF.12} is free, including
on $w=0$. Thus the quotient
is a complex manifold. There
are no simultaneous root
intersections. Elsewhere it
is an ordinary elliptic bundle.
These charts also prove
properness over $V$: locally
the map is a finite quotient
of a product with the compact
curve $A$. Since $V$ is
compact, $Y$ is compact.

\noindent\textbf{Step 12.} Compose with $h$ to obtain
$f:Y\to\mathbb P^1$.
The component over the strict
$B$ is exactly the coarse
torsor $S$ in \eqref{cdp:eq:C.SF.2}, with
the specified isomorphism
coming from its identical
$H^1$ class. At a generic
point no root divisor is
present and the auxiliary
component has multiplicity
one, so $S$ also has
multiplicity one in the
full fibre of $f$.

\noindent\textbf{Step 13.} Every ordinary vertical line
in \eqref{cdp:eq:C.SF.10} has its canonical
meromorphic divisor section,
nonvanishing off the central
fibre. Every root line has
its tautological section,
also nonvanishing there.
Thus both line bundles are
canonically trivial on the
entire complement of zero,
and \eqref{cdp:eq:C.SF.11} gives the
actual product identification
\[
 f^{-1}(\mathbb P^1\setminus\{0\})
   \simeq A\times A_0\times
                         (\mathbb P^1\setminus\{0\}),
                                                               \tag{C.SF.13}\label{cdp:eq:C.SF.13}
\]
where $A_0=B$ in the elliptic
base case. The fibres of
$Y\to V$ are connected
elliptic curves, including
the finite translation
quotients over root divisors.
Since $h$ has connected
fibres, so does $f$.

\noindent\textbf{Step 14.} Finally the auxiliary central
support is an SNC curve
with no triple points.
Its inverse image in the
free quotient charts
\eqref{cdp:eq:C.SF.12} has coordinate
normal crossings. For
example a crossing whose
root coordinate is $z$
has local full equation
$w^{m_i}x=0$. Every
component of the reduced
support is a smooth
elliptic quasi-bundle over
one auxiliary curve; its
special fibres still have
smooth elliptic reduction.
Its Euler characteristic
is therefore zero by
additivity over the smooth
base stratum and the
finitely many multiple
values. Every double
curve is one smooth
elliptic translation
quotient, and there are
no triple points.
Inclusion--exclusion gives
$e(f^{-1}(0)_{\mathrm{red}})=0$.
All other fibres in
\eqref{cdp:eq:C.SF.13} are smooth tori,
so have the same zero
Euler characteristic.
No balancing modification
of the given $S$ is needed.
\end{proof}






\begin{proposition}[Every prescribed cyclic smooth-multiple elliptic quotient occurs in a zero-Euler torus filling]\label{cdp:st:C.2.UN.35}
\StatementRecord{cdp:st:C.2.UN.35}{proved}{Historical events: Claude Fable 5.2 (Anthropic), effort high, 16 November 2026: carried as background without fault found. GPT-6 Astra (OpenAI), ChatGPT Pro mode, effort Pro, S6\_09\_PAPERS, 17 November 2026: targeted proof check. Also checked in this writing. Pre-continuation writing events: GPT-6 Astra (OpenAI), ChatGPT Pro mode, effort Pro, S6\_10\_PAPERS, 18 November 2026: targeted proof check; used and rechecked in this run. Continuation metadata events: GPT-6 Astra (OpenAI), ChatGPT Pro mode, effort Pro, S6\_10\_PAPERS continuation, 19 November 2026: used and rechecked in this run; status-vocabulary reconciliation; metadata reconciliation only; no new mathematical check. \emph{Pre-continuation recorded proof scope.} The equivariant exponential kernel and vanishing H2 lift the full given torsor and action. The displayed Riemann-Roch C1 argument, norm surjectivity, all five integral cyclic character modules and exact restriction-transfer calculation prove generic torus-torsor vanishing. Steinberg1965 Theorem1.9 p51 and its proof pp78-80 were also opened. The invariant meromorphic frame identifies every character line and linearisation with its actual equivariant divisor. Orbit blowups, root charts, free inertia translations, the faithful diagonal quotient and resolution realize the prescribed quotient minimal model and product general fibre}{the hypotheses, notation and definitions stated here; \cref{cdp:st:C.2.UN.23}; \cref{cdp:st:C.2.UN.32}; \cref{cdp:st:C.2.UN.29}; \cref{cdp:eq:C.EO.1}; \cref{cdp:eq:C.EO.3}; \cref{cdp:eq:C.EO.5}; \cite[Theorem~1.9]{Steinberg1965}; \cref{cdp:eq:C.EO.H3}; \cref{cdp:eq:C.EO.H1}; \cref{cdp:eq:C.EO.H4}; \cref{cdp:eq:C.EO.H2}; \cref{cdp:eq:C.EO.H7}; \cref{cdp:eq:C.EO.H5}; \cref{cdp:eq:C.EO.H6}; \cref{cdp:eq:C.EO.7}; \cref{cdp:eq:C.EO.9}; \cref{cdp:eq:C.EO.10}; \cref{cdp:eq:C.EO.4}; \cref{cdp:eq:C.EO.11}; \cref{cdp:eq:C.EO.14}; \cref{cdp:eq:C.EO.6}; \cref{cdp:eq:C.EO.12}; \cref{cdp:eq:C.EO.13}; \cref{cdp:st:C.2.UN.1}; \cref{cdp:st:C.2.UN.4}; \cref{cdp:st:C.2.UN.31}; \cref{cdp:st:C.2.UN.22}; \cref{cdp:st:C.2.UN.20}; \cref{cdp:st:C.2.UN.24}}

Let $q:E\to B$ be a connected compact smooth elliptic
surface with a fixed elliptic Jacobian $F$, identity
integral $H_1$ monodromy, and only smooth multiple
fibres. Suppose $g(B)\le1$. Let $Q$ be a finite
cyclic group acting holomorphically on $E$, preserving
$q$, and acting faithfully on $B$. The trivial group
is allowed. The smooth minimal model of $E/Q$ occurs
as the minimal normalization model of a component of
a fibre in a smooth compact threefold over $\PP^1$
whose general fibre is a complex two-torus. It also
occurs in such a fibre with reduced-support Euler
number zero. The construction retains the given
surface and its given action before taking the
quotient: the distinguished component of the
constructed covering threefold is $Q$-equivariantly
isomorphic to $E$ over $B$.

Consequently the minimal normalization models of
nonnegative Kodaira dimension occurring in arbitrary
two-torus degenerations are exactly the following
two explicitly specified classes:
\begin{itemize}
\item the minimal resolutions of finite cyclic
quotients of complex two-tori that have nonnegative
Kodaira dimension;
\item the minimal models of $E/Q$ in the preceding
construction that have nonnegative Kodaira dimension.
\end{itemize}
The same two classes occur among zero-Euler fillings.
The torus-quotient class is described completely in
\ref{cdp:st:C.2.UN.23} and \ref{cdp:st:C.2.UN.32}. In the second class, all multiplicities,
translation gluing, and the cyclic action are part
of the specified input; none is discarded by the
occurrence construction.
\end{proposition}



\begin{proof}
\noindent\textbf{Step 1.} We give the extension on actual finite root charts,
and then prove the exhaustive assertion.

\par\medskip\noindent\textbf{The actual orbicurve torsor.}
Write the multiple fibres as $m_iC_i$ over $p_i\in B$,
with $m_i>1$. The normalized root chart proved in
\ref{cdp:st:C.2.UN.29} and (C.KX.3)--(C.KX.5) is
\[
 (F\times\Delta_u)/\langle\delta_i\rangle,
 \qquad
 \delta_i(z,u)=(z+\tau_i,\zeta_{m_i}u),
 \qquad \operatorname{ord}(\tau_i)=m_i,
 \qquad t=u^{m_i}.
                                                               \tag{C.EO.1}\label{cdp:eq:C.EO.1}
\]
The free translation on $F$ is essential here.
Let $\mathcal B=B[(p_i,m_i)]$ denote the root
orbicurve, whose chart at $p_i$ is
$[\Delta_u/\mu_{m_i}]$ with $t=u^{m_i}$.
Changes of local origins of the fixed elliptic
curve are translations: the linear part is the
identity because the integral monodromy is the
identity. Thus these charts and the ordinary
product charts exhibit the given $E$ as a
principal $F$-torsor over $\mathcal B$.
Its total space is the original smooth surface,
because all the root inertia actions in \eqref{cdp:eq:C.EO.1}
are free. The action of $Q$ induces the functorial
action on the root orbicurve, and its linear action
on $F$ is a constant finite character
\[
                 \beta:Q\longrightarrow
                            \operatorname{Aut}(F,0).
                                                               \tag{C.EO.2}\label{cdp:eq:C.EO.2}
\]
The constancy follows from the discreteness of
$\operatorname{Aut}(F,0)$ and connectedness of
the regular base. The actual action on the
torsor includes its translation part.

\par\medskip\noindent\textbf{A lift that retains the whole action.}
Put $F=\C/\Lambda$ and let
$\ell:\Lambda\otimes_{\Z}\C\to\C$
be the complex-linear extension of the inclusion
of the period lattice. There is an exact sequence
of complex Lie groups
\[
 0\longrightarrow K\longrightarrow
 T=(\Lambda\otimes_{\Z}\C)/\Lambda
 \xrightarrow{\,\Phi\,}F\longrightarrow0,
 \qquad K=\ker\ell\simeq(\C,+),
 \qquad T\simeq(\C^*)^2.                              \tag{C.EO.3}\label{cdp:eq:C.EO.3}
\]
For a basis $(\lambda_1,\lambda_2)$ of $\Lambda$,
the last map is explicitly
\[
 \Phi(z_1,z_2)=
 \frac{\lambda_1\log z_1+\lambda_2\log z_2}{2\pi i}
                    \pmod\Lambda .                 \tag{C.EO.4}\label{cdp:eq:C.EO.4}
\]
Changing either logarithm changes the numerator
by the corresponding period. The map is locally
split and surjective. Its kernel is exactly
the image of the complex line $\ker\ell$:
a vector mapping to a period can first be
changed by that integral lattice vector.
The intersection $\ker\ell\cap\Lambda$ is zero,
and the kernel is closed because $F$ is Hausdorff.
This proves the exactness, including the actual
additive kernel, without a choice of a cyclic
subgroup of the elliptic periods.

\noindent\textbf{Step 2.} Every $\beta(g)$ acts on $\Lambda$ by an integral
matrix and hence on $T$ by the corresponding
monomial automorphism. It preserves $K$; on
that line its character is $\beta(g)^{-1}$,
because its determinant on $\Lambda$ is one.
Sequence \eqref{cdp:eq:C.EO.3} is therefore $Q$-equivariant.
On the quotient orbicurve $\mathcal C=[\mathcal B/Q]$
it gives an exact sequence of sheaves of
holomorphic sections
\[
 0\longrightarrow\mathcal K\longrightarrow
 \mathcal T\longrightarrow\mathcal F\longrightarrow0.
                                                               \tag{C.EO.5}\label{cdp:eq:C.EO.5}
\]
Here $\mathcal K$ is a holomorphic line bundle
with its additive group law; the other sheaves
retain the indicated finite monodromy. The given
equivariant elliptic torsor is an element of
$H^1(\mathcal C,\mathcal F)$.

\noindent\textbf{Step 3.} One has $H^2(\mathcal C,\mathcal K)=0$.
Indeed finite-group invariants on complex vector
spaces are exact by averaging. On each finite
root chart the higher invariant direct images
therefore vanish, and the invariant direct image
of a coherent sheaf is coherent on the coarse
curve. Coherent cohomology in degree two on a
compact complex curve is zero. The same averaging
for $Q$ proves the asserted vanishing on
$\mathcal C$. The long exact sequence of
\eqref{cdp:eq:C.EO.5} consequently supplies an actual
$\mathcal T$-torsor $P$ and an isomorphism
\[
                  P\mathbin{\times^{T}}F
                           \simeq E                 \tag{C.EO.6}\label{cdp:eq:C.EO.6}
\]
in the equivariant orbicurve category. In
particular the descent isomorphisms of $P$
satisfy the full cocycle identity; they are
not merely isomorphisms of underlying line
bundle classes.

\par\medskip\noindent\textbf{An invariant meromorphic frame.}
The faithful action on $B$ makes
$L=\C(B)$ a cyclic extension of
$k=\C(B/Q)$. At the generic point the root
inertia is absent. The monomial descent of
$(\C^*)^2_L$ defines an algebraic $k$-torus
$T_k$, and the generic fibre of $P$ is an
algebraic torsor under it. To see algebraicity
directly, pull back to $\mathcal B$ and choose
meromorphic sections of its two associated
orbifold line bundles. Such sections exist:
after separating the finitely many root weights,
this is the ordinary meromorphic-section theorem
for line bundles on a compact curve. The
equivariant descent isomorphisms are then
monomials multiplied by meromorphic functions
on $B$, which belong to $L$. Their existing
cocycle identities give precisely this algebraic
generic torsor.

\noindent\textbf{Step 4.} It has a $k$-rational point. This is the
rank-two torus case of Steinberg's theorem
\cite[Theorem~1.9]{Steinberg1965}; the independent
curve-field and cyclic-lattice proof below
establishes the required vanishing for exactly
the representations used here.
\par\medskip\noindent\textbf{An independent calculation of the generic torus torsor.}
Here is a direct verification of the required vanishing
for the actual cyclic rank-two torus. Put
$K=\mathbb C(B/Q)$. The splitting action on its
character lattice has order $d\in\{1,2,3,4,6\}$;
if the base group is larger, pass first to the fixed
field of the kernel of this lattice action. This gives
the actual cyclic splitting extension $L/K$ of degree $d$.
We prove
\[
                         H^1(K,T)=0.                 \tag{C.EO.H1}\label{cdp:eq:C.EO.H1}
\]
The cohomology here is the Galois cohomology of the
algebraic torus obtained by descent of the two-dimensional
split multiplicative group, and classifies its actual
torsors over $K$.

\noindent\textbf{Step 5.} First every function field $K'$ of a smooth compact
complex curve has the following elementary property:
a homogeneous polynomial of degree $a$ in $b>a$
variables has a nonzero zero over $K'$. To verify it,
let $g$ be the genus of the curve, choose a point $P$,
and choose an effective divisor $D$ bounding the poles
of all coefficients. Seek each variable in
$H^0(\OO(rP))$. For large $r$ there are
$b(r+1-g)$ scalar unknowns, whereas the polynomial
value belongs to $H^0(\OO(arP+D))$, of dimension
$ar+\deg D+1-g$. Its vanishing imposes at most that
many homogeneous equations of positive degree.
For large $r$ the number of unknowns is strictly
larger. The height bound for an ideal generated by
that many elements gives an affine homogeneous zero
set of dimension at least one. Over $\mathbb C$ it
contains a point other than the origin. This yields
the asserted nonzero zero over $K'$.

\noindent\textbf{Step 6.} Consequently, for every finite field extension
$L'/K'$ the norm map is surjective:
\[
               N_{L'/K'}:(L')^*\longrightarrow(K')^*.
                                                               \tag{C.EO.H2}\label{cdp:eq:C.EO.H2}
\]
Indeed, for $c\in(K')^*$, choose a $K'$-basis of
$L'$ of size $a=[L':K']$ and apply the preceding
property to
$N_{L'/K'}(x_1e_1+\cdots+x_ae_a)-cz^a$.
This has degree $a$ in $a+1$ variables.
A nonzero zero cannot have $z=0$, because the norm
of a field element vanishes only for the zero element.
Dividing by $z$ gives an element of norm $c$.
Every intermediate field used below is again the
function field of a smooth compact complex curve,
so the same proof applies to it.

\noindent\textbf{Step 7.} The elliptic lattice action has, after an integral
choice of basis, the following matrices:
\[
\begin{array}{c|c}
d&M_d\\ \hline
1&I_2\\
2&-I_2\\
3&\begin{pmatrix}0&-1\\1&-1\end{pmatrix}\\
4&\begin{pmatrix}0&-1\\1&0\end{pmatrix}\\
6&\begin{pmatrix}0&-1\\1&1\end{pmatrix}.
\end{array}                                                   \tag{C.EO.H3}\label{cdp:eq:C.EO.H3}
\]
For completeness, a rank-two lattice preserved by a
primitive automorphism of order $3,4$, or $6$ is a
rank-one module over the corresponding ring
$\mathbb Z[\zeta_d]$. After tensoring with $\mathbb Q$
and choosing a nonzero lattice vector, it becomes
a fractional ideal of that ring. The Gaussian and
Eisenstein rings are Euclidean for their usual norm,
so this ideal is principal. The same argument applies
to the dual character lattice. It gives the matrices
in \eqref{cdp:eq:C.EO.H3}; replacing a generator by its inverse
gives an integrally conjugate matrix. The order-two
case is scalar multiplication by $-1$.

\noindent\textbf{Step 8.} If $d=1$, the torus is split and Hilbert's Theorem 90
gives \eqref{cdp:eq:C.EO.H1}. For $d=2$, it is the product of two
quadratic norm-one tori, since each character lattice
is the sign representation. For $d=3$, it is the
cubic norm-one torus. To specify this identification,
the exact sequence
\[
 1\longrightarrow R^1_{L/K}\mathbb G_m
 \longrightarrow R_{L/K}\mathbb G_m
 \xrightarrow{\ N\ }\mathbb G_m\longrightarrow1
                                                               \tag{C.EO.H4}\label{cdp:eq:C.EO.H4}
\]
has character quotient
$\mathbb Z[C_3]/\mathbb Z(1+g+g^2)$.
In the basis $[1],[g]$, multiplication by $g$ is
exactly $M_3$. In degree two the same quotient is
the sign lattice. Galois descent of characters
therefore gives the stated isomorphisms of tori.
The long exact sequence of \eqref{cdp:eq:C.EO.H4} and Hilbert 90
identify its $H^1$ with $K^*/N_{L/K}L^*$, which is
zero by \eqref{cdp:eq:C.EO.H2}.

\noindent\textbf{Step 9.} For $d=4$, let $K_2$ be the degree-two intermediate
field. The torus is
\[
        T\simeq R_{K_2/K}
                  \bigl(R^1_{L/K_2}\mathbb G_m\bigr).
                                                               \tag{C.EO.H5}\label{cdp:eq:C.EO.H5}
\]
This is an exact character-lattice identification:
induce the sign module of $\langle g^2\rangle$
to $\langle g\rangle$; with basis $e,ge$, the
relations $g(e)=ge$ and $g(ge)=-e$ give $M_4$.
Weil restriction and the corresponding Shapiro
map identify its $H^1$ over $K$ with the $H^1$
of the norm-one torus over $K_2$. That group is
$K_2^*/N_{L/K_2}L^*$ and vanishes by \eqref{cdp:eq:C.EO.H2}.

\noindent\textbf{Step 10.} For $d=6$, let $K_2,K_3$ be the degree-two and
degree-three intermediate fields. On restriction
to $K_2$, a generator acts by $M_6^2$; on restriction
to $K_3$ it acts by $M_6^3=-I_2$. The integral matrix
\[
 P=\begin{pmatrix}1&-1\\0&1\end{pmatrix},\qquad
                  P^{-1}M_6^2P=M_3                  \tag{C.EO.H6}\label{cdp:eq:C.EO.H6}
\]
identifies the first restricted torus with the
cubic norm-one torus, and the second with a product
of quadratic norm-one tori. Thus both restricted
$H^1$ groups vanish.

\noindent\textbf{Step 11.} We spell out the transfer used to return to $K$.
For a finite extension $K'/K$, put
$\mathcal G=\operatorname{Gal}(K^{\rm sep}/K)$,
$\mathcal H=\operatorname{Gal}(K^{\rm sep}/K')$,
and $A=T(K^{\rm sep})$, written additively for this
calculation. Restriction is restriction of continuous
one-cocycles to $\mathcal H$. Choose left coset
representatives $R$ for $\mathcal G/\mathcal H$;
write $\sigma r=r_\sigma h_{\sigma,r}$.
For a one-cocycle $c$ on $\mathcal H$ the transfer
is represented by
\[
       (\operatorname{cor}c)(\sigma)
            =\sum_{r\in R}r_\sigma\,
                         c(h_{\sigma,r}).             \tag{C.EO.H7}\label{cdp:eq:C.EO.H7}
\]
The cocycle identity verifies that this is a cocycle.
For $c=\operatorname{res}b$, the identity
$b(\sigma r)=b(\sigma)+\sigma b(r)
=b(r_\sigma)+r_\sigma b(h_{\sigma,r})$
shows that \eqref{cdp:eq:C.EO.H7} differs from
$[K':K]b(\sigma)$ by the coboundary of
$\sum_{r\in R}b(r)$. Thus
$\operatorname{cor}\circ\operatorname{res}$ is
multiplication by $[K':K]$ on $H^1(K,T)$.
The restrictions to $K_2$ and $K_3$ just proved
to vanish therefore kill every class by both
two and three. Their coprimality gives \eqref{cdp:eq:C.EO.H1}.

\noindent\textbf{Step 12.} This proves the required generic torsor vanishing
in every possible elliptic character, with the
actual restriction and norm maps specified.
The integral identities \eqref{cdp:eq:C.EO.H3}, \eqref{cdp:eq:C.EO.H5},
and \eqref{cdp:eq:C.EO.H6}, together with the compatible
elliptic quotient and additive-kernel characters,
are independently executed in
\path{code/A_equivariant_torus.py} and its saved output.

\noindent\textbf{Step 13.} The curve-field proof in that calculation also
checks the precise field hypothesis in the
printed statement of Steinberg's theorem.
It applies to every finite extension of $k$,
since that extension is another complex curve
function field. A central division algebra
of degree $d>1$ over any such field has a
reduced norm of degree $d$ in $d^2$ variables.
The proved $C_1$ property would give a nonzero
element of reduced norm zero, impossible in
a division algebra. Thus all the indicated
division algebras are commutative, which is
the cohomological-dimension-at-most-one
condition used in \cite[Theorem~1.9]{Steinberg1965}.
Pull back that rational point to $L$.
It is a meromorphic section $s$ of $P$ on
$\mathcal B$ satisfying the exact equality
\[
                         g(s(b))=s(gb)
                              \qquad(g\in Q).        \tag{C.EO.7}\label{cdp:eq:C.EO.7}
\]
Thus the translation cocycle has been retained
in a genuine invariant meromorphic frame.

\noindent\textbf{Step 14.} Let $M=\operatorname{Hom}(T,\C^*)$ be the
character lattice. For $\chi\in M$, write
$L_\chi$ for the associated line bundle and
$s_\chi$ for its meromorphic section induced
by $s$. Its divisor defines an additive map
\[
 D:M\longrightarrow\operatorname{Div}(\mathcal B),
       \qquad D(\chi)=\operatorname{div}(s_\chi).
                                                               \tag{C.EO.8}\label{cdp:eq:C.EO.8}
\]
This map is $Q$-equivariant with its actual
action on characters and divisors, by \eqref{cdp:eq:C.EO.7}.
Choose two basis characters. Their supports,
together with the root points, lie in one
finite $Q$-invariant subset $\Sigma\subset B$.
In terms of coarse points, every coefficient
of $D(\chi)$ is integral away from the root
points and belongs to $m_i^{-1}\Z$ at $p_i$.
We have the canonical multiplicative
identifications
\[
          (L_\chi,s_\chi)
                   \simeq(\OO_{\mathcal B}(D(\chi)),1),
                                                               \tag{C.EO.9}\label{cdp:eq:C.EO.9}
\]
where $1$ denotes the canonical meromorphic
divisor section. They commute with $Q$ and
with addition of characters. There is no
remaining scalar linearization or group
two-cocycle in \eqref{cdp:eq:C.EO.9}.

\par\medskip\noindent\textbf{An equivariant vertical extension.}
We construct a smooth projective surface
$V_0\to\PP^1_t$ containing $B$ in its
central fibre and extending its given
$Q$-action. Let $n=|Q|$.
For elliptic $B$, choose any faithful
character $\xi:Q\to\C^*$ when $n>1$
(and $\xi=1$ for the trivial group),
and take $V_0=B\times\PP^1_t$
with action $(b,t)\mapsto(gb,\xi(g)t)$,
and use $B\times\{0\}$.
For $B\simeq\PP^1$, first diagonalize
its given cyclic action as
$z\mapsto\lambda(g)z$ and define
$\xi=\lambda^{-1}$. Faithfulness makes
this a faithful disc character too.
For the trivial group use both
characters equal to one. Choose any
elliptic curve $A$ and point $o\in A$.
Blow up $(o,0)$ in $A\times\PP^1_t$,
where the action is
$(a,t)\mapsto(a,\xi(g)t)$.
The exceptional curve is identified with
the given $B$: in a local coordinate $x$
at $o$, its affine coordinate $z=x/t$
has exactly the required action. The
other chart supplies its point at infinity.

\noindent\textbf{Step 15.} Blow up every point of $\Sigma$ on the
distinguished strict curve $B$, in complete
$Q$-orbits. Call the resulting disjoint
exceptional curves $T_p$, $p\in\Sigma$,
and call the new surface $V$. The strict
curve is still isomorphic to $B$, and
\[
 \OO_V(T_p)|_B=\OO_B(p),\qquad
                     g(T_p)=T_{gp}.                \tag{C.EO.10}\label{cdp:eq:C.EO.10}
\]
These identities include the canonical
sections and their equivariant structures.
There is no restriction at a fixed point
or at the node of the original central
fibre. At a node with coordinates $xy=0$
and $B=(x=0)$, its blowup has chart
$x=uv,y=v$. Here the strict $B$ is
$u=0$, $T_p$ is $v=0$, and the
restriction of its canonical section
is the simple parameter $v$ on $B$.
This proves \eqref{cdp:eq:C.EO.10} also in that case.
The other central branch meets $T_p$
at the other point of this exceptional
curve; it does not meet the strict $B$.

\noindent\textbf{Step 16.} Root $V$ along each $T_{p_i}$ to order
$m_i$, obtaining the orbifold surface
$\mathcal V$. Its local chart at a
point of $T_{p_i}$ is
\[
 [\Delta_u\times\Delta_v/\mu_{m_i}],
           \qquad x=u^{m_i},\qquad
           \zeta\cdot(u,v)=(\zeta u,v).
                                                               \tag{C.EO.11}\label{cdp:eq:C.EO.11}
\]
The $T_{p_i}$ are disjoint, so there
is only this one inertia group in
each such chart. Its restriction to
$B$ is precisely the previous root
orbicurve $\mathcal B$, including the
functorial $Q$-action. The transverse
intersection in \eqref{cdp:eq:C.EO.10} is what
identifies its root line and root
section with those on $\mathcal B$.

\noindent\textbf{Step 17.} Replace each coarse point $p$ in
$D(\chi)$ by $T_p$, with the same
coefficient. Denote the resulting
orbifold Cartier divisor by
$\widetilde D(\chi)$. All its
coefficients are legitimate integral
root-divisor coefficients on
$\mathcal V$. The map
$\chi\mapsto\widetilde D(\chi)$ is
additive and $Q$-equivariant because
the replacement $p\mapsto T_p$ is.
The character lines
$\OO_{\mathcal V}(\widetilde D(\chi))$
therefore define a principal $T$-torsor
$\widetilde P$ with its exact monomial
$Q$-action. Define the principal
elliptic bundle
\[
              \widetilde Y
                =\widetilde P\mathbin{\times^T}F
                     \longrightarrow\mathcal V .   \tag{C.EO.12}\label{cdp:eq:C.EO.12}
\]
Concretely, on every line-bundle
trivialization its transition map is
translation on $F$ by the value under
\eqref{cdp:eq:C.EO.4} of the two multiplicative
transition functions. Equations
\eqref{cdp:eq:C.EO.9}--\eqref{cdp:eq:C.EO.10} give an actual
$Q$-equivariant isomorphism
\[
                   \widetilde Y|_{\mathcal B}
                                   \simeq E,         \tag{C.EO.13}\label{cdp:eq:C.EO.13}
\]
over the given base, not merely
an equality of topological Chern
classes. Locally the meromorphic
frame identification extends
holomorphically because the two
divisors agree: their quotient is
a nowhere-zero holomorphic unit.
The same assertion holds in the
root coordinates. Since both
frames in \eqref{cdp:eq:C.EO.9} are exactly
equivariant, these local isomorphisms
also retain the given $Q$-action.
More explicitly, a lift of a base
map to $x=u^{m_i}$ involves a
choice of an $m_i$-th root of
its nonzero local multiplier.
Two choices differ by the
root inertia in \eqref{cdp:eq:C.EO.11}.
The character lines simultaneously
apply the corresponding elliptic
translation in \eqref{cdp:eq:C.EO.14}.
Hence the two lifts induce the
same map on the total coarse
quotient. Compositions of lifts
agree there for the same reason.
This is the finite-chart form
of the canonical section-preserving
action, and proves that no local
root choice changes its group law.

\par\medskip\noindent\textbf{Smoothness, properness and the generic torus.}
Let $(a_1,a_2)\pmod {m_i}$ be the
two weights of $P$ at $p_i$.
In \eqref{cdp:eq:C.EO.11} the total elliptic
bundle has the finite chart
\[
 (\Delta_u\times\Delta_v\times F)/\mu_{m_i},
 \quad
 \zeta_{m_i}(u,v,z)
     =(\zeta_{m_i}u,v,z+\tau_i),\quad
 \tau_i=\frac{a_1\lambda_1+a_2\lambda_2}{m_i}
                                           \pmod\Lambda .
                                                               \tag{C.EO.14}\label{cdp:eq:C.EO.14}
\]
The order of $\tau_i$ is exactly
$m_i$ by \eqref{cdp:eq:C.EO.1} and \eqref{cdp:eq:C.EO.6}.
Every nonidentity inertia element
therefore acts freely on the total
chart, including $u=0$ and every
central crossing $v=0$. Thus the
coarse space in \eqref{cdp:eq:C.EO.12} is a
smooth compact complex threefold;
we continue to call it
$\widetilde Y$. Properness follows
on these finite charts from
properness of the compact factor
$F$ and glues over the compact
surface $V$. The principal-bundle
transition functions in \eqref{cdp:eq:C.EO.12}
give the gluing maps explicitly.

\noindent\textbf{Step 18.} All divisors used to construct
$\widetilde P$ are vertical over
$t=0$. Their canonical meromorphic
sections are therefore nowhere
zero on the complement of that
fibre. There \eqref{cdp:eq:C.EO.12} is the
actual product
\[
 \widetilde Y|_{t\ne0}
       \simeq A\times F\times(\PP^1\setminus\{0\}),
                                                               \tag{C.EO.15}\label{cdp:eq:C.EO.15}
\]
where $A=B$ in the elliptic-base
case and $A$ is the chosen auxiliary
elliptic curve in the rational-base
case. The action is the prescribed
action on $A$ (the identity in the
rational-base construction),
the linear action $\beta$ on $F$,
and multiplication by $\xi$ on
$t$. This formula follows from
the invariant divisor frames,
so no additional translation
cocycle survives off zero.

\noindent\textbf{Step 19.} Take the finite quotient by $Q$
and resolve its singularities.
Finite holomorphic actions can
be linearized locally, and these
cyclic quotients have compatible
proper resolutions; a resolution
of the compact quotient, chosen
over its singular locus, is
enough. Its map to
$\PP^1_T$, $T=t^n$, has general
fibre exactly $A\times F$:
the $n$ covering fibres are
permuted simply transitively
for $T\ne0,\infty$. The quotient
is smooth there, so the
resolution leaves those fibres
unchanged. The Stein factor is
a finite map of degree one
to the normal curve $\PP^1$,
and is the identity. Thus
the resulting proper map has
connected fibres.

\noindent\textbf{Step 20.} The distinguished component
before the quotient is exactly
the given $E$ by \eqref{cdp:eq:C.EO.13}.
Its image in the finite quotient
is $E/Q$: taking invariants
is exact in characteristic zero,
so the restriction of the
quotient structure sheaf to
this image is
$(\OO_E)^Q$. It is normal
because $E$ is normal. The
strict component on the ambient
resolution maps properly and
birationally to $E/Q$.
Resolving that component and
contracting its surface
$(-1)$-curves gives exactly
the prescribed minimal model.
Finally \ref{cdp:st:C.2.UN.1} gives even Euler
number for the reduced support.
The actual vertical point and
quartic blowups of \ref{cdp:st:C.2.UN.4} change
that number by $+2$ and $-2$
respectively. Finitely many
such operations make it zero
and preserve the minimal
normalization model through
their explicit proper
birational maps. This proves
the occurrence assertion.

\par\medskip\noindent\textbf{Exhaustion of the nonnegative classes.}
Let $Z$ be a minimal normalization
model of nonnegative Kodaira
dimension in an arbitrary
two-torus degeneration. \ref{cdp:st:C.2.UN.31}
provides its actual cyclic
component cover and equivariant
minimal contraction. A torus
cover gives exactly the first
class in the statement. In
the remaining case \ref{cdp:st:C.2.UN.22} and
\ref{cdp:st:C.2.UN.29} give a fixed-Jacobian,
identity-monodromy fibration
with only smooth multiples
over a rational or elliptic
curve. We specify why an
appropriate such fibration
is preserved by the cyclic
action.

\noindent\textbf{Step 21.} For a cover of Kodaira
dimension one it is the
intrinsic connected Iitaka
fibration. In algebraic
dimension one its base is
the intrinsic meromorphic
function field, and the
equivariant Stein factor
of \ref{cdp:st:C.2.UN.20} applies. The
non-torus projective case
of Kodaira dimension zero
is bielliptic. Use its
actual finite free product
description
$(F_1\times F_2)/\Gamma$,
with $\Gamma$ acting by
translations on $F_1$
and with a nontrivial
linear action on $F_2$.
The invariant holomorphic
vector fields on the
product are exactly the
constant fields in the
$F_1$ direction. These
integrate to the connected
automorphism group of
the quotient, an elliptic
quotient of $F_1$.
Its orbit map is
\[
 (F_1\times F_2)/\Gamma
                   \longrightarrow F_2/\Gamma\simeq\PP^1.
                                                               \tag{C.EO.16}\label{cdp:eq:C.EO.16}
\]
Every automorphism normalizes
the connected automorphism
group and hence preserves
this orbit fibration. Its
local stabilizers act by
free translations on
$F_1$; its root charts
are \eqref{cdp:eq:C.EO.1}. Therefore
it has the required
identity monodromy and
only smooth multiples.
This is the preserved
fibration used in the
bielliptic case.

\noindent\textbf{Step 22.} Let $H$ be the kernel of
the cyclic action on the
chosen base. Its action
on a general elliptic
fibre is by translations:
a nontrivial linear
elliptic automorphism
has rational quotient,
which would make the
quotient surface ruled
and annihilate all
positive pluricanonical
sections, contrary to
$\kappa(Z)\ge0$.
The exact translation
and reflection charts
(C.EQ.3)--(C.EQ.4)
show that quotienting
the cover by $H$ is
again a smooth elliptic
surface with fixed
Jacobian, identity
monodromy, and only
smooth multiples.
The residual cyclic
group acts faithfully
on the base. Its
quotient is the same
actual surface quotient
as before, through
the finite quotient
maps. It is therefore
one of the inputs
already constructed
above.

\noindent\textbf{Step 23.} Conversely, the
construction proves
occurrence for every
member of the second
class. For the first,
take the compact
diagonal quotient
$(A\times\PP^1_t)/G$
using a faithful
character on $t$,
resolve it, and use
the same \ref{cdp:st:C.2.UN.4} balancing;
the exact component
comparison is \ref{cdp:st:C.2.UN.24}.
All occurrences may
therefore be chosen
with zero-Euler
support, proving
both asserted
equalities of classes.
\end{proof}






\begin{proposition}[Occurrence of every Hopf surface of algebraic dimension one]\label{cdp:st:C.2.UN.37}
\StatementRecord{cdp:st:C.2.UN.37}{proved}{Historical events: Claude Fable 5.2 (Anthropic), effort high, 16 November 2026: carried as background without fault found. GPT-6 Astra (OpenAI), ChatGPT Pro mode, effort Pro, S6\_09\_PAPERS, 17 November 2026: targeted proof check. Also checked in this writing. Pre-continuation writing events: GPT-6 Astra (OpenAI), ChatGPT Pro mode, effort Pro, S6\_10\_PAPERS, 18 November 2026: targeted proof check. Continuation metadata events: GPT-6 Astra (OpenAI), ChatGPT Pro mode, effort Pro, S6\_10\_PAPERS continuation, 19 November 2026: used and rechecked in this run; status-vocabulary reconciliation; metadata reconciliation only; no new mathematical check. \emph{Pre-continuation recorded proof scope.} The primary finite-cover transfer and determinant character give the actual rational Betti groups for secondary surfaces as well. Positivity of local curve intersections removes meromorphic base points; adjunction and the universal cover show every curve is smooth elliptic and every full connected fibre is a smooth multiple. The normalized root cover is etale and its free elliptic deck action is by translations, proving identity integral monodromy. Its invariant j-function extends and is constant. Applying UN.34 to that exact root-curve torsor preserves the specified Hopf surface itself}{the hypotheses, notation and definitions stated here; \cref{cdp:st:C.2.UN.34}; \cref{cdp:eq:C.EH.2}; \cref{cdp:eq:C.EH.4}; \cref{cdp:eq:C.EH.6}; \cref{cdp:eq:C.EH.7}}

Let $H$ be any primary or secondary Hopf surface with $a(H)=1$.
There is a proper holomorphic map from a smooth compact complex
threefold to $\mathbb P^1$, with connected fibres and smooth
two-tori away from one value, whose fibre over that value contains
$H$ itself as a smooth component of multiplicity one. Its reduced
special fibre has simple normal crossings and Euler number zero.
Over the complement of the special value the family is the
actual product of two fixed elliptic curves with the base.

The map used to apply \ref{cdp:st:C.2.UN.34} is the actual connected elliptic
reduction
\[
                         q:H\longrightarrow\mathbb P^1.
\tag{C.EH.1}\label{cdp:eq:C.EH.1}
\]
Every fibre of $q$ is a multiple of a smooth elliptic curve,
the modulus is constant, and the integral $H^1$ monodromy is
the identity. The assertion includes secondary Hopf surfaces
without imposing a cyclic condition on their finite primary cover.
\end{proposition}



\begin{proof}
\noindent\textbf{Step 1.} We first construct \eqref{cdp:eq:C.EH.1} and check its full fibres directly
on the given surface. For any Hopf surface one has
\[
                    b_1(H)=1,\qquad H^2(H,\mathbb Q)=0.
\tag{C.EH.2}\label{cdp:eq:C.EH.2}
\]
These statements follow through the actual finite unramified
primary cover: the primary cover has the rational cohomology
of $S^1\times S^3$, transfer gives $b_1\leq1$ and $b_2=0$,
and the nonzero deck character
$\gamma\mapsto\log|\det(d\gamma)_0|$ gives $b_1\geq1$.
This is the comparison proved in (C.HQ.5)--(C.HQ.7).
Thus it applies to secondary surfaces as well.

\noindent\textbf{Step 2.} Choose a nonconstant meromorphic function $h$ on $H$.
Its effective zero and pole divisors have no common
irreducible component and have intersection number zero
by \eqref{cdp:eq:C.EH.2}. Local intersection multiplicities of distinct
effective curves on a smooth surface are positive.
Consequently those two divisors are disjoint. In a local
coprime expression $h=f/g$, an indeterminacy point would
belong to both $(f=0)$ and $(g=0)$, so no such point exists.
The function therefore defines an actual holomorphic map
$h:H\to\mathbb P^1$. Its Stein factorization is
\[
               H\xrightarrow{q}B\xrightarrow{\nu}\mathbb P^1,
\tag{C.EH.3}\label{cdp:eq:C.EH.3}
\]
where $B$ is a smooth compact connected curve, $q$ has
connected fibres, and $\nu$ is finite. The degree-one
Leray edge map injects $H^1(B,\mathbb Q)$ into
$H^1(H,\mathbb Q)$, since $q_*\mathbb Q=\mathbb Q$.
Thus $2g(B)\leq1$, and $B\simeq\mathbb P^1$.

\noindent\textbf{Step 3.} Every irreducible compact curve $C\subset H$ is smooth
elliptic. Indeed, adjunction and \eqref{cdp:eq:C.EH.2} give
\[
                 2p_a(C)-2=(K_H+C)\cdot C=0.
\tag{C.EH.4}\label{cdp:eq:C.EH.4}
\]
Its normalization cannot be rational: a holomorphic map
$\mathbb P^1\to H$ lifts through the universal cover
$\mathbb C^2\setminus\{0\}\to H$, and both coordinate
functions of the lift are constant. If $g$ is the genus
of the normalization and $\delta$ is its normalization
defect, \eqref{cdp:eq:C.EH.4} says $g+\delta=1$. Hence $g=1$ and
$\delta=0$. Furthermore two distinct irreducible curves
are disjoint, again by their zero intersection number
and local positivity. Since $q$ has connected fibres,
each full fibre is therefore
\[
                              q^{-1}(b)=m_b C_b,
\tag{C.EH.5}\label{cdp:eq:C.EH.5}
\]
with $C_b$ a smooth connected elliptic curve. The scheme
fibre is the displayed Cartier divisor: a map from the
smooth surface to a curve has no embedded fibre component.

\noindent\textbf{Step 4.} We check the monodromy by an actual finite map, including
at a multiple fibre. Take a small disc with coordinate $t$
at $b$, write $m=m_b$, and normalize the base change
$t=s^m$. If $x=0$ defines $C_b$ locally, then
$t=u x^m$ for a holomorphic unit $u$. The normalized
cover is given locally by
\[
                     w^m=u,\qquad s=wx.
\tag{C.EH.6}\label{cdp:eq:C.EH.6}
\]
It is a finite unramified cover of the original surface
tube: locally it is the disjoint union of the $m$ choices
of a holomorphic root of $u$. Denote it by $V\to U$.
The deck group $\mu_m$ acts freely, and $V/\mu_m=U$.
The punctured base-changed family has connected fibres
and connected total space. The normalization has no
component supported over zero. Its Stein factor over
$\Delta_s$ is finite and generically degree one and
therefore is the disc itself. In particular the central
fibre $C'$ is connected. It is smooth by \eqref{cdp:eq:C.EH.6}, and
its multiplicity as the fibre of $V\to\Delta_s$ is one.
Thus this map is a submersion on a neighbourhood of $C'$.
The actual central comparison is the finite unramified
quotient
\[
                         C'\longrightarrow C'/\mu_m=C_b.
\tag{C.EH.7}\label{cdp:eq:C.EH.7}
\]

\noindent\textbf{Step 5.} A fixed-point-free finite-order automorphism of an elliptic
curve is a translation. To see this, choose an origin and
write it as $z\mapsto\beta z+a$. If $\beta\ne1$, the
nonzero elliptic endomorphism $1-\beta$ is an isogeny,
so $(1-\beta)z=a$ has a solution. Hence every nonidentity
element of the deck group in \eqref{cdp:eq:C.EH.7} is a translation.
Choose a $\mu_m$-invariant smooth horizontal connection
for the submersion $V\to\Delta_s$, by finite averaging.
Radial parallel transport then gives an equivariant
smooth trivialization. A positive meridian in the
$t$-disc acts on $H^1$ by the corresponding deck
translation or its inverse, according to the transport
convention; both induce the identity. We have proved
\[
                  T_b=I\quad\hbox{on }H^1(q^{-1}(t),\mathbb Z).
\tag{C.EH.8}\label{cdp:eq:C.EH.8}
\]
The finitely many meridians generate the fundamental
group of the punctured base $\mathbb P^1$, so the entire
integral monodromy representation is trivial.

\noindent\textbf{Step 6.} The holomorphic elliptic $j$-function extends across
each omitted value because the actual cover
$V\to\Delta_s$ is a smooth elliptic family there;
its $j$-function is deck invariant and descends to
the $t$-disc. Thus $j$ is holomorphic on $\mathbb P^1$
with values in $\mathbb C$ and is constant. The map
\eqref{cdp:eq:C.EH.1} has exactly the fixed-modulus, identity-monodromy,
smooth-multiple-fibre structure required in \ref{cdp:st:C.2.UN.34}.
Apply that theorem to the given surface $H$ and this
specific map $q$. Its root-curve torsor is the one
defined by \eqref{cdp:eq:C.EH.6}--\eqref{cdp:eq:C.EH.7}, and its complete analytic
torsor class is retained by the two-line-bundle lift.
Consequently the distinguished component in that
construction is isomorphic to $H$, with multiplicity
one, rather than merely having its topological or
birational invariants. The proper compact map, exact
punctured product and Euler-zero simple-normal-crossing
support are those constructed in (C.SF.8)--(C.SF.13).
\end{proof}






\begin{proposition}[Every diagonal primary Hopf surface occurs unchanged in a compact torus filling]\label{cdp:st:C.2.UN.41}
\StatementRecord{cdp:st:C.2.UN.41}{proved}{Historical events: Claude Fable 5.2 (Anthropic), effort high, 16 November 2026: carried as background without fault found. GPT-6 Astra (OpenAI), ChatGPT Pro mode, effort Pro, S6\_09\_PAPERS, 17 November 2026: targeted proof check. Also checked in this writing. Pre-continuation writing events: GPT-6 Astra (OpenAI), ChatGPT Pro mode, effort Pro, S6\_10\_PAPERS, 18 November 2026: targeted proof check; used and rechecked in this run. Continuation metadata events: GPT-6 Astra (OpenAI), ChatGPT Pro mode, effort Pro, S6\_10\_PAPERS continuation, 19 November 2026: used and rechecked in this run; status-vocabulary reconciliation; metadata reconciliation only; no new mathematical check. \emph{Pre-continuation recorded proof scope.} The two dual-normal bundles, their inversion gluing, Hausdorff separation and compact-core bounds give the actual proper smooth filling. The four-period lattice and height-circle comparison prove degree one of the punctured torus map. The nodal cubic pencil supplies the specified section value; the explicit Tate coordinate isomorphism identifies the compact suspension with the plumbing. Primitive integral restriction matrices and constant coherent deck cocycles give all stated support groups. Schutt-Shioda section7.6 pp33-34 and Papikian Theorem2.1 p2 were opened for translation extension and the exact complex uniformization series}{the hypotheses, notation and definitions stated here; \cref{cdp:eq:C.DH.5}; \cref{cdp:eq:C.DH.6}; \cref{cdp:eq:C.DH.7}; \cite[\S7.6, pp.~33--34]{SchuttShioda2010}; \cref{cdp:eq:C.DH.13}; \cref{cdp:eq:C.DH.14b}; \cref{cdp:eq:C.DH.12}; \cref{cdp:st:C.2.UN.40}; \cite[Theorem~2.1, p.~2]{Papikian2004Tate}}

For every pair $\alpha,\beta\in\C$ with
$0<|\alpha|,|\beta|<1$, the diagonal Hopf surface
\[
 H_{\alpha,\beta}=(\C^2\setminus\{0\})/
       \langle(z,w)\mapsto(\alpha z,\beta w)\rangle
                                                               \tag{C.DH.1}\label{cdp:eq:C.DH.1}
\]
occurs as an unchanged smooth component of a
reduced normal-crossing fibre of a proper map
from a smooth compact complex threefold to
$\PP^1$, with general fibre a two-torus.
The distinguished fibre consists of two copies
of \eqref{cdp:eq:C.DH.1}, joined along both axis elliptic
curves by inversion. Every reduced fibre
support in this compact construction has
Euler number zero. This includes all
multiplicatively independent pairs, whose
Hopf surfaces have algebraic dimension zero.

The distinguished support $S$ has torsion-free
integral cohomology of ranks
\[
                  (1,2,3,4,2),\qquad
                  h^\bullet(S,\OO_S)=(1,2,1).        \tag{C.DH.2}\label{cdp:eq:C.DH.2}
\]
Each normalization component has integral
cohomology of ranks $(1,1,0,1,1)$ and
$h^\bullet(\OO)=(1,1,0)$.
Locally the nearby torus has the exact
multiplicative period presentation
\[
 A_t=(\C^*)^2/
        \langle(\alpha,\beta),(t,t^{-1})\rangle
                    \qquad(0<|t|\ll1).             \tag{C.DH.3}\label{cdp:eq:C.DH.3}
\]
In the compact construction the constants
$\alpha,\beta$ in this formula vary
holomorphically with unchanged limiting
values, and $t$ is replaced by a holomorphic
parameter with a simple zero. The explicit
comparison below retains the torus fibres.
\end{proposition}



\begin{proof}
\noindent\textbf{Step 1.} We first construct the proper local filling
and its actual torus comparison, then glue
it into a compact suspension.

\par\medskip\noindent\textbf{Exact neighborhoods and the two plumbing maps.}
The two axes of $H=H_{\alpha,\beta}$ are
$C_\alpha=\C^*/\alpha^{\Z}$ and
$C_\beta=\C^*/\beta^{\Z}$. The open
neighborhood $z\ne0$ of the first is
the entire normal line bundle
\[
 L_\alpha=(\C^*\times\C)/
         ((z,x)\sim(\alpha z,\beta x));
                                                               \tag{C.DH.4}\label{cdp:eq:C.DH.4}
\]
the second has the analogous presentation
with $\alpha,\beta$ interchanged.
Inversion on $C_\alpha$ pulls this
normal line to its inverse. Explicitly
the rank-two plumbing bundle over that
elliptic curve is
\[
 \mathcal N_\alpha=(\C^*\times\C^2)/
       ((z,x,y)\sim(\alpha z,\beta x,\beta^{-1}y)),
                         \qquad t=xy.               \tag{C.DH.5}\label{cdp:eq:C.DH.5}
\]
Its branch $y=0$ maps to the first
Hopf copy by $(z,x)$, and its branch
$x=0$ maps to the second by
$(Z,W)=(z^{-1},y)$.
For the other axis use
\[
 \mathcal N_\beta=(\C^*\times\C^2)/
       ((w,x,y)\sim(\beta w,\alpha x,\alpha^{-1}y)),
                         \qquad t=xy,                \tag{C.DH.6}\label{cdp:eq:C.DH.6}
\]
whose branches map to the first
copy by $(z,w)=(x,w)$ and to the
second by $(Z,W)=(y,w^{-1})$.
These formulas prove the two dual-normal
identifications, including their
base and fibre transition functions.

\noindent\textbf{Step 2.} Put $A=-\log|\alpha|>0$ and
$B=-\log|\beta|>0$. On the common
open torus of a Hopf copy define
\[
 h=B\log|z|-A\log|w|,\qquad
 \nu_\alpha=|w||z|^{-B/A}=e^{-h/A},\qquad
 \nu_\beta=|z||w|^{-A/B}=e^{h/B}.        \tag{C.DH.7}\label{cdp:eq:C.DH.7}
\]
All three expressions are invariant
under the Hopf deck generator.
The two $\nu$ are the norms of
flat Hermitian metrics on the
normal lines. Neighborhoods
$\nu_\alpha<R$ and $\nu_\beta<R$
are disjoint for $R<1$.
Fix $0<r<R<1$ and $0<\delta<r^2$.
For each copy take the product
of the open compact core
$U_i=\{\nu_\alpha>r,\nu_\beta>r\}$
with $\Delta_\delta$.
In \eqref{cdp:eq:C.DH.5} take
\[
 \|x\|=|x||z|^{-B/A}<R,
 \qquad \|y\|=|y||z|^{B/A}<R,
 \qquad |xy|<\delta;                                \tag{C.DH.8}\label{cdp:eq:C.DH.8}
\]
take the analogous open set in
\eqref{cdp:eq:C.DH.6}. Glue it to the first
core where $\|x\|>r$, using the
first branch coordinate and
$t=xy$, and to the second core
where $\|y\|>r$, using the second
branch coordinate. The product
of the two norms is exactly
$|t|$. Thus the two core
overlaps in any one neck are
disjoint, since $|t|<r^2$.
The coordinate maps are
biholomorphisms on the stated
overlaps; their inverses insert
the coordinate $t/x$ or $t/y$.

\noindent\textbf{Step 3.} This finite gluing is Hausdorff.
Indeed at the boundary $\|x\|=r$
of a core overlap the core
point is absent, and at its
other boundary $\|x\|=R$ the
neck point is absent. Thus
the graph of each identification
is closed relative to the
two pieces. The same holds
for $y$; distinct necks have
disjoint overlaps in a given
core, and the preceding
$r^2$ inequality excludes
an additional identification
through the opposite core.
This also proves separation
for the resulting finite
equivalence relation. The
quotient is a smooth complex
threefold $\mathcal Y$ with
the holomorphic map $t$.
At the two double curves it
has the exact coordinates
$xy=t$; at zero its fibre
is the two unchanged Hopf
surfaces with the asserted
two inversion identifications.

\noindent\textbf{Step 4.} The map is proper after the
indicated restriction to the
disc. Choose
$r<r'<R'<R$.
Over a closed subdisc,
the two compact sets
$\{\nu_\alpha\ge r',\nu_\beta\ge r'\}$
in the Hopf copies, together
with the two closed neck sets
$\|x\|,\|y\|\le R'$, cover
the entire inverse image.
For a core point missing its
compact set, the small normal
coordinate is less than $r'$
and the other neck coordinate
has norm at most $\delta/r<r$.
For a neck point with one
norm greater than $R'$,
the corresponding core point
belongs to its compact set.
All these closed sets are
compact: their bases are
compact elliptic curves and
their vector-bundle coordinates
are bounded. They cover the
closed inverse image, proving
properness. This argument
also explicitly verifies
compactness of the fibres.

\par\medskip\noindent\textbf{The exact nearby torus.}
For $t\ne0$, the punctured
identifications are
\[
 \phi_\alpha(z,w)=(z^{-1},t/w),\qquad
 \phi_\beta(z,w)=(t/z,w^{-1}).                       \tag{C.DH.9}\label{cdp:eq:C.DH.9}
\]
Each conjugates the Hopf deck
generator to its inverse.
Their difference is the
additional period
$(z,w)\mapsto(tz,t^{-1}w)$.
The four additive period
vectors of \eqref{cdp:eq:C.DH.3} are
\[
 (2\pi i,0),\ (0,2\pi i),\
 (\log\alpha,\log\beta),\
                         (\log t,-\log t).          \tag{C.DH.10}\label{cdp:eq:C.DH.10}
\]
They are real-linearly
independent: after quotienting
the two imaginary directions,
the real determinant is
\[
       -\log|t|\bigl(\log|\alpha|+\log|\beta|\bigr)
                                      \ne0.          \tag{C.DH.11}\label{cdp:eq:C.DH.11}
\]
Their integral span is
therefore a discrete
cocompact lattice in $\C^2$.
This proves directly that
$A_t$ is a compact complex
two-torus. A turn of $t$
adds $(2\pi i,-2\pi i)$
to the last vector, which
is already in the lattice,
so these quotients form
the stated family over
the punctured disc.

\noindent\textbf{Step 5.} There is an actual
holomorphic comparison
$\Psi_t:\mathcal Y_t\to A_t$.
On the first core use
$[z,w]$, and on the
second use $[Z^{-1},t/W]$.
On the first neck use
$[z,x]$; on the second
use $[x,w]$. On its
second-core overlap
the latter gives instead
$[x/t,tw]$, differing by
the inverse of the period
$(t,t^{-1})$. Thus the
maps agree exactly on
all the overlaps. They
are local biholomorphisms.
The proper source makes
$\Psi_t$ a finite
unramified covering
onto the connected
torus $A_t$.

\noindent\textbf{Step 6.} Its degree is one.
For completeness the
height in \eqref{cdp:eq:C.DH.7}
gives a direct check.
Put $L=-\log|t|$ and
$c_r=-\log r$, $c_R=-\log R$.
The first core has
height interval
$(-Bc_r,Ac_r)$.
Under the chosen comparison
the second core has interval
$(AL-Ac_r,AL+Bc_r)$.
The first neck fills
$(Ac_R,AL-Ac_R)$,
and the other fills,
modulo the period
$(A+B)L$,
$(AL+Bc_R,(A+B)L-Bc_R)$.
These intervals cover
one height circle;
their consecutive
overlaps are exactly
the gluing annuli.
Since $L>2c_r$,
height zero has one
preimage in the first
core and none in the
other three pieces,
apart from the already
identified overlaps.
On that core the only
remaining identification
is the original Hopf
deck generator. Thus
the covering degree
is one, proving the
actual biholomorphism
in \eqref{cdp:eq:C.DH.3}.

\noindent\textbf{Step 7.} The entire construction
also works for holomorphic
$a(t),b(t)$ with limiting
values $\alpha,\beta$
and for a holomorphic
smoothing parameter $Q(t)$
with a simple zero.
Use the smooth family
of Hopf quotients with
deck multipliers $a(t),b(t)$
and put $xy=Q(t)$ in
the two necks. Their
quotient and gluing maps
are the same holomorphic
formulas. The norms and
height vary continuously,
and the preceding bounds
hold uniformly on a
smaller disc. Since
$Q'(0)\ne0$, the total
space is still smooth.
The comparison is now
\[
 \mathcal Y_t\simeq
 (\C^*)^2/
       \langle(a(t),b(t)),(Q(t),Q(t)^{-1})\rangle.
                                                               \tag{C.DH.12}\label{cdp:eq:C.DH.12}
\]

\par\medskip\noindent\textbf{A compact suspension with these exact punctured fibres.}
We construct a compact elliptic
surface $J\to\PP^1$ with a
section $P$ whose value on a
nodal fibre is the specified
$\alpha$. Start with the
nodal plane cubic
$y^2z=x^2(x+z)$, with origin
$O=[0:1:0]$. On its
normalization put
$v=y/x$; its smooth group
coordinate is
$\zeta=(v-1)/(v+1)$,
with $O$ corresponding to
$\zeta=1$. Choose the
smooth point $P_0$ with
$\zeta=\alpha$.
Take a second cubic
through $O,P_0$, otherwise
general. Its intersection
with the nodal cubic
consists of nine distinct
smooth points, with
transverse intersections,
and misses the node.
These requirements are
nonempty open conditions
on its restriction to
the cubic and its first
jets; the restriction
map of cubic polynomials
is surjective, so a
general such choice
satisfies them.
Blow up the nine
base points of the
pencil. The resulting
smooth rational surface
$J$ has a connected
elliptic pencil, zero
section from $O$, and
section $P$ from $P_0$.
Its central member is
the original nodal
cubic, a fibre of type
$I_1$. At the node the
second cubic is nonzero,
so the pencil parameter
smooths the node to
first order. The full
fibre class is
$F=3H-\sum E_i=-K_J$.
Consequently there is
no vertical rational
$(-1)$-curve, by
adjunction, and this
is the smooth relatively
minimal elliptic model.

\noindent\textbf{Step 8.} Translation by $P$ is an
automorphism $T_P$ of
this entire surface.
Here is the required
extension argument.
The group law first
defines a birational
map over the base,
with inverse translation
by $-P$. Resolve its
graph minimally,
$J\xleftarrow{p}W\xrightarrow{q}J$.
Every exceptional curve
is vertical. If an
exceptional $(-1)$-curve
$D$ of $p$ is not
exceptional for $q$,
then $q^*K_J\cdot D=0$
because $K_J=-F$.
The discrepancy formula
$K_W=q^*K_J+\sum a_iD_i$
has $a_i>0$ and
$D_i\cdot D\ge0$,
contradicting $K_W\cdot D=-1$.
Thus $D$ is exceptional
for both maps and can
be contracted in their
common graph, contrary
to its minimality.
There are no exceptional
curves, so the translation
and its inverse are
holomorphic automorphisms.
This agrees with the explicit
section-translation extension
in \cite[\S7.6, pp.~33--34]{SchuttShioda2010}.

\noindent\textbf{Step 9.} Put $q_0=\alpha\beta$.
The quotient
\[
 Y^{\rm std}=(\C^*_u\times J)/
        \langle(u,y)\mapsto(q_0u,T_P(y))\rangle
                                                               \tag{C.DH.13}\label{cdp:eq:C.DH.13}
\]
is a smooth compact
threefold. Projection to
$E_{q_0}=\C^*/q_0^{\Z}$
is a locally trivial
holomorphic bundle with
compact fibre $J$:
local choices of a
logarithm give its
trivializations, and
their transition is
the indicated integral
power of $T_P$.
The deck action is
free and properly
discontinuous because
its first coordinate
is multiplication by
$q_0$, with $|q_0|<1$.
It preserves the
elliptic-pencil parameter,
so \eqref{cdp:eq:C.DH.13} maps
properly to $\PP^1$.

\noindent\textbf{Step 10.} On a small punctured
disc at its $I_1$ value
there is the actual
analytic Tate presentation
\[
 J_t\simeq\C^*/Q(t)^{\Z},\qquad
 Q(t)=t\,u(t),\quad u(0)\ne0,
 \qquad P(t)=[a(t)],\quad a(0)=\alpha.
                                                               \tag{C.DH.14}\label{cdp:eq:C.DH.14}
\]
The presentation can be
obtained by normalizing
the nonvanishing cycle
period to one. Its
other period has the
form $(2\pi i)^{-1}\log t$
plus a holomorphic
function; exponentiation
gives $Q(t)$ with a
simple zero. The
standard nodal Tate
chart identifies its
smooth central group
with $\C^*$, so the
section through $P_0$
has the holomorphic
nonzero lift $a(t)$
with that limiting
value. This is an
isomorphism of the
elliptic groups and
their smooth central
charts, not only an
identity of their
$j$-invariants.
For clarity, the central
chart and its sign can
also be checked directly.
The complex uniformization in \cite[Theorem~2.1 and equation~(2.1), p.~2]{Papikian2004Tate} gives
\[
 \begin{gathered}
 y^2+xy=x^3+a_4(q)x+a_6(q),\\
 a_4(q)=-5q+O(q^2),\quad
 a_6(q)=-q+O(q^2),\quad
 \Delta(q)=q+O(q^2).
 \end{gathered}                                 \tag{C.DH.14a}\label{cdp:eq:C.DH.14a}
\]
Here $a_4=-5\sum_{n\ge1}n^3q^n/(1-q^n)$ and
$a_6=-\sum_{n\ge1}(5n^3+7n^5)q^n/(12(1-q^n))$.
The multiplicative
uniformization has the
central values
\[
 x=\frac{u}{(1-u)^2},\qquad
 y=\frac{u^2}{(1-u)^3},\qquad
 X=4x,\quad Y=8y+4x,\qquad
 \frac{Y/X-1}{Y/X+1}=u.                          \tag{C.DH.14b}\label{cdp:eq:C.DH.14b}
\]
Its equation becomes
$Y^2=X^2(X+1)$, precisely
the chosen nodal cubic.
The series converge
normally on compact
subsets of the smooth
central chart. Invert
the holomorphic series
$1/j(q)=q+O(q^2)$ at
the given simple zero
of $1/j(t)$ to obtain
$Q(t)$. After completing
the square and cube,
the two short Weierstrass
coefficients are units
at zero. Their equal
$j$-invariants identify
them by a holomorphic
unit fourth root; choose
that root at zero to
match the sixth-weight
coefficient. The identity
of the squared ratios
then matches that
coefficient throughout
the disc. This supplies
the actual Weierstrass
coordinate isomorphism.
Choose its central value
as in \eqref{cdp:eq:C.DH.14b}.
The inverse function
theorem at the smooth
point $u=\alpha$ gives
the holomorphic lift
$a(t)$ with exactly
$a(0)=\alpha$.

\noindent\textbf{Step 11.} Let $b(t)=q_0/a(t)$.
It tends to $\beta$,
and both multipliers
have modulus less than
one on a sufficiently
small disc. The actual
coordinate change
\[
                        z=y,\qquad w=u/y
                                                               \tag{C.DH.15}\label{cdp:eq:C.DH.15}
\]
on the multiplicative
cover of \eqref{cdp:eq:C.DH.13}
turns its two deck
generators into
$(a(t),b(t))$ and
$(Q(t),Q(t)^{-1})$.
Thus its punctured
family is exactly
\eqref{cdp:eq:C.DH.12}. Replace
the local filling
of \eqref{cdp:eq:C.DH.13} by
the proper plumbing
just constructed,
using that explicit
isomorphism on the
punctured disc.
This is gluing over
the open cover
$\PP^1\setminus\{0\}$
and a disc at zero.
The total spaces
are identified by
biholomorphisms over
their common base,
so the result is
Hausdorff and smooth;
properness is local
on the base and
holds on each member
of this cover.
The compact base
therefore gives a
compact total space.
The distinguished
central components
are exactly the
prescribed Hopf
surfaces. General
fibres are the
actual tori of
\eqref{cdp:eq:C.DH.12}, or,
elsewhere, the same
translation suspensions
of smooth elliptic
fibres of $J$.
The latter are tori
by their additive
four-period lattice,
or by the direct
constant-translation
bundle presentation
in \eqref{cdp:eq:C.DH.13}.

\noindent\textbf{Step 12.} The support over zero
has Euler number
$2e(H)-2e(C)=0$.
At every other value
the support is a
locally trivial bundle
over $E_{q_0}$ with
fibre the corresponding
reduced elliptic-pencil
support; its monodromy
is $T_P$. Its Euler
number is the product
with $e(E_{q_0})=0$.
Thus every reduced
support has Euler
number zero.

\par\medskip\noindent\textbf{The support cohomology and its maps.}
For the two-component
support there is the
actual normalization
Mayer--Vietoris sequence
with middle map the
difference of the two
restrictions to
$C_\alpha\sqcup C_\beta$.
In degree zero its
matrix is
$\left(\begin{smallmatrix}1&-1\\1&-1\end{smallmatrix}\right)$,
of rank one and
primitive image.
Choose on each axis
the circle of its
Hopf deck period and
its angular circle.
The integral $H^1$
generator of each
Hopf surface restricts
as $(1,0)$ on each
axis. The inversion
gluing reverses these
two generators on
the second side.
Consequently the
degree-one restriction
matrix is
\[
                     \begin{pmatrix}
                       1&1\\0&0\\1&1\\0&0
                     \end{pmatrix},                \tag{C.DH.16}\label{cdp:eq:C.DH.16}
\]
again of rank one
with primitive image.
The Hopf components
have $H^2=0$, and
the two elliptic
curves have $H^2=\Z^2$
and no higher
cohomology. Exactness
therefore gives
\eqref{cdp:eq:C.DH.2}, including
torsion-freeness.
For coherent cohomology,
$H^1(H,\OO_H)=\C$
is represented by
the constant deck
cocycle one. Its
restriction to each
axis is the nonzero
constant cocycle on
that elliptic curve;
inversion changes
its sign. The
degree-one coherent
restriction matrix
has both rows $(1,1)$.
Together with
$h^\bullet(\OO_H)=(1,1,0)$
and the degree-zero
matrix already given,
the coherent
normalization sequence
gives $(1,2,1)$.
All maps in these
calculations are
the stated axis
restrictions and
inversion pullbacks.
The curve calculation in \ref{cdp:st:C.2.UN.40} leaves only the two axes when
the multipliers are multiplicatively independent. A nonconstant meromorphic
function would have infinitely many distinct nonempty generic level divisors,
which is impossible when the only compact curves are those two axes.
Consequently the Hopf component in this case has algebraic dimension zero.
\end{proof}






\begin{proposition}[Every nonlinear primary Hopf surface occurs unchanged over a rational base]\label{cdp:st:C.2.UN.42}
\StatementRecord{cdp:st:C.2.UN.42}{proved}{Historical events: Claude Fable 5.2 (Anthropic), effort high, 16 November 2026: carried as background without fault found. GPT-6 Astra (OpenAI), ChatGPT Pro mode, effort Pro, S6\_09\_PAPERS, 17 November 2026: targeted proof check. Also checked in this writing. Pre-continuation writing events: GPT-6 Astra (OpenAI), ChatGPT Pro mode, effort Pro, S6\_10\_PAPERS, 18 November 2026: targeted proof check. Continuation metadata events: GPT-6 Astra (OpenAI), ChatGPT Pro mode, effort Pro, S6\_10\_PAPERS continuation, 19 November 2026: used and rechecked in this run; status-vocabulary reconciliation; metadata reconciliation only; no new mathematical check. \emph{Pre-continuation recorded proof scope.} The explicit hyperelliptic base gives the exact translation jet at its two infinity points. The weighted quotient charts make the nonzero tautological line smooth. Polynomial bounds for every iterate of the lifted translation, combined with the prescribed contraction, prove freeness, proper discontinuity and compactness for every beta in the stated range. The commuting half-translation involution is free by its line modulus and gives the rational base. The exact central function xw, Seifert filling, canonical form and integral inclusion maps give every multiplicity, cohomology rank and Euler assertion}{the hypotheses, notation and definitions stated here; \cref{cdp:eq:C.NH.6}; \cref{cdp:eq:C.NH.5}; \cref{cdp:eq:C.NH.4}; \cref{cdp:eq:C.NH.10}; \cref{cdp:eq:C.NH.12}; \cref{cdp:eq:C.NH.11}; \cref{cdp:eq:C.NH.8}; \cref{cdp:st:C.2.UN.40}; \cref{cdp:eq:C.NH.14}; \cref{cdp:eq:C.NH.13}}

For every integer $r\ge1$ and every $\beta$ with $0<|\beta|<1$,
the primary Hopf surface
\[
 H_r(\beta)=(\C^2\setminus\{0\})/
   \langle(z,w)\mapsto(\beta^r z+w^r,\beta w)\rangle
                                                        \tag{C.NH.1}\label{cdp:eq:C.NH.1}
\]
occurs unchanged as a smooth multiplicity-one component of a
reduced SNC fibre of a proper map $f:Y\to\PP^1$, where $Y$ is
a smooth compact complex threefold and the general fibre is a
complex two-torus. Every reduced fibre support in this
construction has Euler number zero.

The distinguished full fibre is $S=H_r(\beta)+P_r$, with one
smooth elliptic intersection $J$. The second component is a
smooth elliptic surface over an elliptic curve, with exactly
one multiple fibre $rJ$ when $r>1$, and
\[
 \begin{gathered}
 K_{P_r}=\OO_{P_r}((r-1)J),\qquad
 b(P_r)=(1,3,4,3,1),\qquad h^\bullet(\OO_{P_r})=(1,2,1),\\
 H^\bullet(S,\Z)\text{ is torsion-free of ranks }(1,3,5,5,2),
 \qquad h^\bullet(S,\OO_S)=(1,2,1).
 \end{gathered}                                         \tag{C.NH.2}\label{cdp:eq:C.NH.2}
\]
The Hopf component has integral ranks $(1,1,0,1,1)$ and
coherent dimensions $(1,1,0)$. The construction works for
every prescribed $\beta$, without a smallness condition
beyond $|\beta|<1$.
\end{proposition}



\begin{proof}
\noindent\textbf{Step 1.} \par\medskip\noindent\textbf{A compact base with the required translation jet.}
Let $B$ be the elliptic cubic $v^2=u^3-u+1$, with origin the
point at infinity. Its cubic has distinct nonzero roots. Let
$C$ be the smooth projective completion of
\[
 v^2=t^{6r}-t^{2r}+1,
 \qquad \psi:C\longrightarrow B,\quad
          (t,v)\longmapsto(t^{2r},v).
                                                        \tag{C.NH.3}\label{cdp:eq:C.NH.3}
\]
The polynomial has distinct roots: a repeated root would be
nonzero and would give a common root of $u^3-u+1$ and its
derivative. At infinity the two charts
$s=t^{-1}$, $v s^{3r}=\pm\sqrt{1-s^{4r}+s^{6r}}$ give two
distinct smooth points $c_+,c_-$. Thus $C$ is a smooth compact
curve of genus $3r-1$. The involution $h_C(t,v)=(t,-v)$ has
quotient $\PP^1_t$, interchanges $c_+,c_-$, and satisfies
$\psi h_C=-\psi$. Its fixed points are the $6r$ finite roots
of the displayed polynomial.

\noindent\textbf{Step 2.} The elliptic uniformizer $u/v$ has order $r$ at both
$c_\pm$. Choose a local additive coordinate $p$ on $B$ at
zero. At each of the two points a holomorphic coordinate
$s$ on $C$ can therefore be chosen so that
\[
 2\psi(s)=s^r\quad\hbox{in the additive coordinate on }B.
                                                        \tag{C.NH.4}\label{cdp:eq:C.NH.4}
\]
Taking the $r$-th root of its nonzero holomorphic unit proves
this coordinate equality. On the compact surface $B\times C$
consider the commuting automorphisms
\[
 A(p,c)=(p+2\psi(c),c),\qquad
 A_{1/2}(p,c)=(p+\psi(c),c),\qquad
 h(p,c)=(-p,h_Cc).
                                                        \tag{C.NH.5}\label{cdp:eq:C.NH.5}
\]
Here $A_{1/2}^2=A$, $h^2=1$, and the equality
$\psi h_C=-\psi$ proves that $h$ commutes with both
translations. All three preserve the pair of points
$(0,c_+),(0,c_-)$.

\par\medskip\noindent\textbf{The exact weighted blowup and its smooth nonzero line.}
Perform the weighted blowup of weights $(r,1)$ at both
points, retaining its smooth orbifold charts. Denote the
resulting compact orbifold surface by $\mathcal V$, its
exceptional divisors by $R_+,R_-$, and put
$L=\OO_{\mathcal V}(R_++R_-)$. Locally its weighted blowup
has the explicit presentation
\[
 \left[((\C^2_{z,w}\setminus\{0\})\times\C_x)/\C^*\right],
 \quad \lambda(z,w,x)=(\lambda^r z,\lambda w,\lambda^{-1}x),
 \quad p=x^r z,\quad s=xw.
                                                        \tag{C.NH.6}\label{cdp:eq:C.NH.6}
\]
The exceptional divisor is $x=0$. The line $L$ has fibre
coordinate $\ell$ of weight $-1$ and canonical section
$\ell=x$. The quotient chart $w\ne0$ is the smooth chart
$p=s^r(z/w^r)$; the chart $z\ne0$ has the residual finite
action
\[
 (w,x,\ell)\longmapsto
      (\zeta w,\zeta^{-1}x,\zeta^{-1}\ell),\qquad \zeta^r=1.
                                                        \tag{C.NH.7}\label{cdp:eq:C.NH.7}
\]
It is free when $\ell\ne0$. More explicitly, on $L^\times$
use $\lambda=\ell$ in \eqref{cdp:eq:C.NH.6} to set $\ell=1$. This gives
the actual smooth coordinates
$(z,w,x)\in(\C^2\setminus\{0\})\times\C$, after replacing
the coordinates by $(\ell^rz,\ell w,\ell^{-1}x)$.
Thus the coarse nonzero line space $L^\times$ is a smooth
complex threefold, including the orbifold points.

\noindent\textbf{Step 3.} All three maps in \eqref{cdp:eq:C.NH.5} lift canonically. In the
coordinates \eqref{cdp:eq:C.NH.4}, $A$ is $p\mapsto p+s^r$ and its lift
in \eqref{cdp:eq:C.NH.6} is $(z,w,x)\mapsto(z+w^r,w,x)$.
The weighted ideals are therefore preserved. The half
translation replaces the coefficient by $1/2$; $h$
interchanges the two centres and preserves their weighted
orders. The line lifts preserve its canonical divisor
section, so their group laws hold exactly, including
$\widetilde A_{1/2}^{\,2}=\widetilde A$ and
$\widetilde h^{\,2}=1$.

\noindent\textbf{Step 4.} Multiply the line fibre by $\beta$ and compose with
$\widetilde A$. On the smooth coordinates just constructed
the resulting automorphism is
\[
 g(z,w,x)=
       (\beta^r(z+w^r),\beta w,\beta^{-1}x),
 \qquad s=xw.
                                                        \tag{C.NH.8}\label{cdp:eq:C.NH.8}
\]
It preserves the map to $C$. On $x=0$, the change
$Z=\beta^{-r}z$ turns its first coordinate into
$\beta^r Z+w^r$, proving the exact identification with
\eqref{cdp:eq:C.NH.1}.

\par\medskip\noindent\textbf{Uniform properness for every contraction.}
We verify the global quotient, rather than inferring its
compactness from a local contraction. Choose a Hermitian
orbifold metric on $L$. Write $\sigma$ for its canonical
section. Near either centre, its norm off the exceptional
divisor is comparable, with constants independent of the
point, to
\[
       \rho(p,s)=\bigl(d_B(p,0)^2+|s|^{2r}\bigr)^{1/(2r)}.
                                                        \tag{C.NH.9}\label{cdp:eq:C.NH.9}
\]
Here $d_B$ is a fixed flat elliptic distance. To see the
comparison, normalize the weighted coordinates so that
$|z|^2+|w|^{2r}=1$. Then
$(|p|^2+|s|^{2r})^{1/(2r)}=|x|$, which is the norm of the
canonical section up to a smooth positive bounded factor
on the compact weighted exceptional curve. Away from the
two centres all these norms are bounded above and below.
More precisely this comparison holds on $B$ times a small
$c_\pm$-disc: outside the chosen $p$-coordinate neighborhood,
both $\rho$ and $\|\sigma\|$ are bounded above and away from
zero. It therefore also covers an iterate $A^n$ that moves
$p$ out of that local coordinate chart.

\noindent\textbf{Step 5.} For every integer $n$, the triangle inequality on the flat
elliptic curve gives
\[
 d_B(p+2n\psi(s),0)
       \le d_B(p,0)+C_0|n||s|^r.
\]
The same inequality applied to the inverse translation
gives the reverse comparison. Since the canonical lift
sends $\sigma(p,s)$ to $\sigma(p+2n\psi(s),s)$, there is
a global constant $C_1\ge1$ such that
\[
 C_1^{-1}(1+|n|)^{-1/r}\|\ell\|
 \le \|\widetilde A^{\,n}\ell\|
 \le C_1(1+|n|)^{1/r}\|\ell\|.
                                                        \tag{C.NH.10}\label{cdp:eq:C.NH.10}
\]
The comparisons first hold off the exceptional divisors.
The ratio of line norms extends continuously and positively
on the entire orbifold line, so they hold there as well.
On the compact part of $C$ away from both centres the
canonical section is nowhere zero and supplies the same
bounds with a fixed constant. This proves the uniform
global assertion.

\noindent\textbf{Step 6.} Consequently $\|g^n\ell\|$ tends uniformly to zero on the
unit circle bundle as $n\to+\infty$, and to infinity as
$n\to-\infty$: multiply \eqref{cdp:eq:C.NH.10} by $|\beta|^n$.
The action is properly discontinuous because any two
compact subsets have line norms bounded above and away
from zero. It is free, since a fixed point of a nonzero
power would also be fixed by arbitrarily large positive
powers, contradicting the norm limit.

\noindent\textbf{Step 7.} Choose $N$ such that $g^N$ decreases every norm by a factor
strictly less than one. Its minimum norm ratio is a positive
number $c_N$. Every $g^{N\Z}$-orbit meets the compact annulus
$c_N\le\|\ell\|\le1$: take the first iterate whose norm is
at most one. This annulus is compact over the compact
orbifold base. It follows that
\[
                 \widetilde Y=L^\times/\langle g\rangle
                                                        \tag{C.NH.11}\label{cdp:eq:C.NH.11}
\]
is a smooth compact threefold, with a proper map to $C$.
This proof applies to every $0<|\beta|<1$.

\noindent\textbf{Step 8.} Off the two exceptional fibres, the canonical section
trivializes $L$ and gives the actual quotient
\[
 (B\times\C^*)/
   \langle(p,\ell)\mapsto(p+2\psi(c),\beta\ell)\rangle.
                                                        \tag{C.NH.12}\label{cdp:eq:C.NH.12}
\]
If $B=\C/\Lambda_B$, its additive covering lattice is
generated by $(\lambda,0)$ for $\lambda\in\Lambda_B$,
$(0,2\pi i)$, and $(a_c,\log\beta)$, where $a_c$ is any
lift of $2\psi(c)$. These four real generators are
independent because $\log|\beta|\ne0$. Changing the lift
of $a_c$ changes a generator by $\Lambda_B$. This proves
that \eqref{cdp:eq:C.NH.12} is precisely a compact two-torus.

\noindent\textbf{Step 9.} At either distinguished value the exact function is
$s=xw$. Its full fibre is reduced SNC, with components
$H:x=0$ and $P:w=0$, meeting along
\[
          J=(\C^*_z/\beta^{r\Z}),\qquad x=w=0.
                                                        \tag{C.NH.13}\label{cdp:eq:C.NH.13}
\]
Both components and their intersection are compact since
they are closed in \eqref{cdp:eq:C.NH.11}. On $P$ the map to the
unchanged elliptic base $B$ has the exact local form
\[
  p=x^r z,\qquad (z,x)\sim(\beta^r z,\beta^{-1}x).
                                                        \tag{C.NH.14}\label{cdp:eq:C.NH.14}
\]
Away from $p=0$ its canonical frame is the product
$B\setminus\{0\}\times E_\beta$; on the overlap the
frame coordinate is $\ell=x^{-1}$ and $z=p\ell^r$.
These equalities identify the surface with the one-fibre
logarithmic model in (C.ER.10)--(C.ER.16), with $m=r$.
They also give the same model for $r=1$ directly.
In particular every fibre support of $P\to B$ is a
smooth elliptic curve, and the special full fibre is
$rJ$. Stratified Euler additivity gives $e(P)=0$.
The Hopf surface and $J$ also have Euler number zero,
so the entire distinguished support has Euler number zero.
The local holomorphic canonical evaluation form is
$x^r\,dz\wedge dw\wedge dx$. The Jacobian of \eqref{cdp:eq:C.NH.8}
is $\beta^r$, and the multiplier of $x^r$ is $\beta^{-r}$,
so this form is invariant. Its holomorphic vanishing orders
along $(H,P)$ are exactly $(r,0)$. This realizes the
strict canonical-order difference in \ref{cdp:st:C.2.UN.40} without
replacing a holomorphic generator by a meromorphic
full-fibre shift.

\par\medskip\noindent\textbf{The free quotient giving the rational base.}
Fix either square root $\beta^{1/2}$. The commuting
canonical lifts define
\[
           G=\beta^{1/2}\widetilde A_{1/2}\widetilde h,
                    \qquad G^2=g.
                                                        \tag{C.NH.15}\label{cdp:eq:C.NH.15}
\]
Thus $G$ induces an involution on $\widetilde Y$. A fixed
point there would have $h_Cc=c$. Such a point is not one
of the two centres. Moreover $\psi(c)$ is a two-torsion
point of $B$, so $2\psi(c)=0$. In the canonical line frame,
\[
 G(p,c,\ell)=(-p-\psi(c),c,\beta^{1/2}\ell),
 \qquad g^n(p,c,\ell)=(p,c,\beta^n\ell).
                                                        \tag{C.NH.16}\label{cdp:eq:C.NH.16}
\]
The equality $G(p,c,\ell)=g^n(p,c,\ell)$ is impossible by
the absolute values of the line coordinates: it would
give $|\beta|^{1/2}=|\beta|^n$. The involution is therefore
free, and
\[
                 Y=\widetilde Y/\langle G\rangle
                  \longrightarrow C/\langle h_C\rangle=\PP^1
                                                        \tag{C.NH.17}\label{cdp:eq:C.NH.17}
\]
is a smooth compact threefold with the required map.
The two Hopf components over $c_+,c_-$ are exchanged.
Restriction of the quotient to either one is an
isomorphism onto its image, so the prescribed Hopf
surface is unchanged. The entire distinguished fibre is
isomorphic to either $H+P$ above, and remains reduced.
At a fixed value of $h_C$ the full fibre is twice a free
quotient of a two-torus; its reduced support has Euler
number zero. At every other ordinary value it is a
two-torus. All fibres are connected, as is seen directly
from these quotient descriptions and from $H\cap P=J$.
This proves all the stated occurrence and Euler assertions.

\par\medskip\noindent\textbf{The component and support groups.}
The exact local frame in \eqref{cdp:eq:C.NH.14} gives a global canonical
form $dp\wedge d\log\ell$ off $J$. In that chart it is
$-x^{r-1}dz\wedge dx$, and it is invariant under the
displayed deck map. Hence its divisor is $(r-1)J$.
The direct-image equality proved from this same chart is
$p_*\OO_P(kJ)=\OO_B(\lfloor k/r\rfloor[0])$.
It gives $h^0(K_P)=1$. Since $J^2=0$ and $e(P)=0$,
surface Riemann--Roch gives $\chi(\OO_P)=0$ and
$h^\bullet(\OO_P)=(1,2,1)$.

\noindent\textbf{Step 10.} For the integral groups use the exact Seifert filling in
(C.ER.16). Its meridian is $r\delta+h$ and its longitude
is $\delta$, where $\delta$ is the boundary of the
once-punctured elliptic base and $h$ the product circle.
The base relation kills $\delta$ in first homology, and
the meridian then kills $h$. Thus the closed oriented
three-manifold $M_r$ has $H_1(M_r,\Z)=\Z^2$; all its
homology is free, by integral Poincare duality and the
universal coefficient sequence. The radial line coordinate
gives $P\simeq M_r\times S^1$ smoothly. For a nonreal
$\beta$ its angular rotation is isotopic to the identity
through the same circle action, so the mapping torus has
this same product identification. This proves the ranks
of $P$ in \eqref{cdp:eq:C.NH.2}.
That circle isotopy preserves the Seifert core. Thus it
identifies the pair $(P,J)$ with
$(M_r\times S^1,\mathrm{core}\times S^1)$, not merely
the ambient cohomology groups.

\noindent\textbf{Step 11.} The angular loop of $J$ is the solid-torus longitude and
is zero in $H_1(M_r,\Z)$; its radial loop maps primitively
to the last $S^1$. Therefore $[J]=0$ in $H_2(P,\Z)$,
including integrally. In the Hopf component its angular
loop is zero and its radial loop is the primitive deck
generator. The integral normalization sequence of
$S=P\cup_JH$ consequently has the restriction maps
\[
 H^0(P\sqcup H,\Z)\to H^0(J,\Z):\ (a,b)\mapsto a-b,
 \qquad
 H^1(P\sqcup H,\Z)\to H^1(J,\Z):
 \begin{pmatrix}0&0&1&-1\\0&0&0&0\end{pmatrix},
                                                        \tag{C.NH.18}\label{cdp:eq:C.NH.18}
\]
up to simultaneous orientation choices for the radial
loops, and its degree-two map is zero. Both displayed
images are primitive. The exact sequence gives precisely
the free groups of ranks $(1,3,5,5,2)$.

\noindent\textbf{Step 12.} Finally the constant deck cocycle represents the nonzero
class in $H^1(H,\OO_H)$. It is not a coboundary: a
holomorphic solution of $a\circ F-a=1$ on $\C^2\setminus0$
extends across zero and gives $0=1$ there. Its restriction
to $J$ is still nonzero, since the constant Laurent
coefficient of $a(\beta^r z)-a(z)$ is zero. Thus the
coherent degree-one restriction onto $H^1(J,\OO_J)$
is surjective. Together with the degree-zero surjection
and the component groups already computed, the coherent
normalization sequence gives $h^\bullet(S,\OO_S)=(1,2,1)$.
All support maps are the actual inclusions in \eqref{cdp:eq:C.NH.13}--\eqref{cdp:eq:C.NH.14}.
\end{proof}






\begin{proposition}[Every secondary Hopf surface of algebraic dimension zero occurs unchanged]\label{cdp:st:C.2.UN.43}
\StatementRecord{cdp:st:C.2.UN.43}{proved}{Historical events: Claude Fable 5.2 (Anthropic), effort high, 16 November 2026: carried as background without fault found. GPT-6 Astra (OpenAI), ChatGPT Pro mode, effort Pro, S6\_09\_PAPERS, 17 November 2026: targeted proof check. Also checked in this writing. Pre-continuation writing events: GPT-6 Astra (OpenAI), ChatGPT Pro mode, effort Pro, S6\_10\_PAPERS, 18 November 2026: targeted proof check; used and rechecked in this run. Continuation metadata events: GPT-6 Astra (OpenAI), ChatGPT Pro mode, effort Pro, S6\_10\_PAPERS continuation, 19 November 2026: used and rechecked in this run; status-vocabulary reconciliation; metadata reconciliation only; no new mathematical check. \emph{Pre-continuation recorded proof scope.} The normal core and determinant character give a central primary contraction. Its explicit centralizer makes the deck group abelian; the free finite characters give exactly the two stated cyclic presentations. The torsion-section cover is unramified at the prescribed nodal sheet, and the suspended finite action matches both plumbing charts. Quotient resolution and Euler balancing leave that whole tube unchanged. The nonlinear construction uses the intermediate compact weighted-line family over C, before the additional rational-base involution; its scalar quotient is free by the exact coordinate and modulus equations. Kato section3 pp230-232 and the 1989 erratum were opened; the erratum's nonabelian topological corrections are not used}{the hypotheses, notation and definitions stated here; \cite[Theorem~2.1]{Khaled2024Hopf}; \cref{cdp:st:C.2.UN.40}; \cref{cdp:eq:C.SH.3}; \cite[\S3, pp.~230--232]{Kato1975Hopf}; \cref{cdp:st:C.2.UN.41}; \cref{cdp:eq:C.SH.9}; \cref{cdp:eq:C.SH.10}; \cref{cdp:st:C.2.UN.1}; \cref{cdp:st:C.2.UN.4}; \cref{cdp:st:C.2.UN.42}}

Every secondary Hopf surface $H$ of algebraic dimension zero
has an actual finite quotient presentation of one of the forms
\[
\begin{array}{ll}
 H_{\alpha,\beta}/\langle(z,w)\mapsto
          (\zeta_n^a z,\zeta_n^b w)\rangle,
 &\gcd(a,n)=\gcd(b,n)=1,\\[2pt]
 H_r(\beta)/\langle(z,w)\mapsto
          (\zeta_n^r z,\zeta_n w)\rangle,
 &\gcd(r,n)=1.
\end{array}
\tag{C.SH.1}\label{cdp:eq:C.SH.1}
\]
The first primary contraction is $(\alpha z,\beta w)$ with no
multiplicative resonance, and the second is
$(\beta^r z+w^r,\beta w)$, where $0<|\alpha|,|\beta|<1$.
This is an exhaustive presentation, without a uniqueness
assertion about the displayed parameters.

Every prescribed such surface occurs itself as a smooth
multiplicity-one component of a reduced SNC fibre of a proper
map from a smooth compact complex threefold to a smooth
compact curve, with connected fibres and general fibre a
complex two-torus. All reduced fibre supports can have Euler
number zero. In the diagonal construction the base is the
finite torsion-cover curve $B'$ constructed below; no claim
that this curve is rational is needed.
\end{proposition}



\begin{proof}
\noindent\textbf{Step 1.} \par\medskip\noindent\textbf{The actual deck group.}
Write $W=\mathbb C^2\setminus\{0\}$ and $H=W/\Gamma$.
The deck group has an infinite cyclic contraction subgroup of
finite index. Its normal core, the intersection of its finitely
many conjugates, is a normal infinite cyclic subgroup
$\langle F\rangle$ of finite index. Conjugation acts on it by
$1$ or $-1$. The nonzero conjugation-invariant character
$\gamma\mapsto\log|\det(d\gamma)_0|$ excludes the second
possibility. Thus $F$ is central and the actual primary cover is
\[
                 P=W/\langle F\rangle\longrightarrow H.
\tag{C.SH.2}\label{cdp:eq:C.SH.2}
\]
Taking elementary symmetric functions on its finite generic
sheets gives $a(P)=a(H)=0$.

\noindent\textbf{Step 2.} The primary normal forms and their full centralizers are
\cite[Theorem~2.1]{Khaled2024Hopf}. In the nonresonant diagonal
case a commuting power series has the form $(Az,Bw)$:
the eigenmonomial argument of \ref{cdp:st:C.2.UN.40} leaves only $z$ and
$w$ in its respective coordinate functions. In the
non-diagonal case the full centralizer is
\[
       (z,w)\longmapsto(A^r z+B w^r,Aw),\qquad A\ne0.
\tag{C.SH.3}\label{cdp:eq:C.SH.3}
\]
Indeed write a commuting map as $(P,Q)$. Its second
coordinate satisfies $Q\circ F=\beta Q$. The weighted
eigenfunction calculation of \ref{cdp:st:C.2.UN.40} forces $Q=Aw$.
The first coordinate then satisfies
$P\circ F=\beta^rP+A^r w^r$. On a weighted-degree-$d$
piece with $d\ne r$, the operator $F^*-\beta^r$ is
invertible, since its only eigenvalue is
$\beta^d-\beta^r\ne0$. Thus $P=Cz+Bw^r$.
Comparing its $w^r$ coefficients gives $C=A^r$;
invertibility requires $A\ne0$. This proves the full
formula \eqref{cdp:eq:C.SH.3} directly.
The parameters $(A,B/A^r)$ in this formula have product
$(A_1,c_1)(A_2,c_2)=(A_1A_2,c_1+c_2)$. Both centralizers
are abelian. Since every element of $\Gamma$ commutes
with $F$, the actual deck group is abelian.

\noindent\textbf{Step 3.} A finitely generated abelian group containing $\mathbb Z$
with finite index has the form $\mathbb Z\oplus T$, with $T$
its finite torsion subgroup. Choose a generator $f$ of the
infinite factor with the contracting sign. Some positive power
is $F^k$ times an element of $T$, $k>0$. Thus both diagonal
eigenvalue moduli of $f$ are less than one. In \eqref{cdp:eq:C.SH.3} it has
$|A|<1$ and nonzero additive parameter, since the parameter
of $F$ is nonzero. Rescaling $z$ makes its resonant coefficient
one. We obtain the actual holomorphic quotient maps
\[
       H=(W/\langle f\rangle)/T,\qquad
       \Gamma=\langle f\rangle\times T.
\tag{C.SH.4}\label{cdp:eq:C.SH.4}
\]
For diagonal $T\subset(\mathbb C^*)^2$, each coordinate
character is injective: its kernel would fix a nonzero point
of the corresponding axis. Hence $T$ is cyclic and both
weights are units modulo its order. In \eqref{cdp:eq:C.SH.3}, a finite-order
element has additive parameter zero and acts as
$(A^r z,Aw)$. Its finite subgroup is cyclic, and freeness
on the first axis forces $\gcd(r,n)=1$. This proves \eqref{cdp:eq:C.SH.1}
and the finite-primary-cover comparison, including the
exclusion of a nonabelian group in algebraic dimension zero.
The abelian calculation also appears in
\cite[\S3, pp.~230--232]{Kato1975Hopf}; no topological
classification is used here.

\noindent\textbf{Step 4.} Conversely the indicated finite maps are unramified.
An equation $\tau^k(z,w)=f^\ell(z,w)$ at a nonzero coordinate
forces $\ell=0$ by the modulus of the contraction multiplier;
the coprimality conditions then force $k=0$ modulo $n$.
In the non-diagonal case use $w$ when it is nonzero and
$z$ on its complement. Thus these are the exact free
finite actions to be extended.

\par\medskip\noindent\textbf{A prescribed torsion section at the nodal fibre.}
For the diagonal case start with the compact elliptic surface
$J\to\mathbb P^1$, zero section, nodal fibre at zero and
section $P$ of \ref{cdp:st:C.2.UN.41}. Its multiplicative value at that
fibre is the prescribed $\alpha$. Explicitly, on
$y^2=x^2(x+1)$ the point with group coordinate $\alpha$ is
\[
 \left(\frac{4\alpha}{(1-\alpha)^2},
       \frac{4\alpha(1+\alpha)}{(1-\alpha)^3}\right).
\tag{C.SH.5}\label{cdp:eq:C.SH.5}
\]
Blowing up the nine distinct smooth base points of a general
cubic pencil through this point and the origin gives that
surface and its section, with an unchanged $I_1$ fibre.

\noindent\textbf{Step 5.} Let $B^\circ$ be the smooth elliptic-fibre locus.
The exact-order-$n$ part of $J[n]|_{B^\circ}$ is a finite
etale algebraic cover. Select a component by the following
actual local section. The smooth group locus of the nodal
family has central group $\mathbb C^*$; multiplication by
$n$ has nonzero derivative at $\zeta_n$. The holomorphic
implicit-function theorem gives a unique local $n$-torsion
section specializing to $\zeta_n$, necessarily of exact order.
Take the connected cover containing its punctured restriction,
and normalize $\mathbb P^1$ in its function field. This gives
a smooth connected projective curve and finite map
\[
                             B'\longrightarrow\mathbb P^1.
\tag{C.SH.6}\label{cdp:eq:C.SH.6}
\]
At the selected point $b_0$ this map is unramified: the local
sheet is the graph of the just constructed section and has
degree one over the disc. The actual $I_1$ fibre, its simple
smoothing parameter and $P(b_0)=\alpha$ remain unchanged.

\noindent\textbf{Step 6.} Resolve the compact algebraic base-changed elliptic surface
and contract vertical $(-1)$-curves, away from this unchanged
tube, to obtain a smooth relatively minimal surface $J'\to B'$.
The pulled-back section and the tautological torsion point
extend as sections $P',R$: rational maps from the smooth
projective base curve to the proper regular surface extend by
the valuative criterion. They satisfy
\[
                 P'(b_0)=\alpha,\qquad R(b_0)=\zeta_n,
                 \qquad \operatorname{ord}(R)=n.
\tag{C.SH.7}\label{cdp:eq:C.SH.7}
\]
Their translations extend to automorphisms of $J'$.
The minimal-graph argument of \ref{cdp:st:C.2.UN.41} applies because
$K_{J'}\cdot C\ge0$ for every vertical curve, by adjunction
and relative minimality, or by Kodaira's fibre components.
A $(-1)$-curve exceptional for one graph projection but not
the other would have nonnegative intersection with its
canonical pullback and all other effective discrepancy
curves, contradicting canonical intersection $-1$.
Both translations and inverses are therefore holomorphic;
their commutation on the smooth fibres proves commutation
everywhere.

\par\medskip\noindent\textbf{The diagonal action and the actual compact quotient.}
Put $q_0=\alpha\beta$ and form the compact suspension
\[
 Y^{\rm std}=(\mathbb C^*_u\times J')/
       \langle(u,y)\mapsto(q_0u,y+P')\rangle.
\tag{C.SH.8}\label{cdp:eq:C.SH.8}
\]
Its proper compact bundle comparison is the one in
\ref{cdp:st:C.2.UN.41}. The commuting order-$n$ automorphism
\[
               (u,y)\longmapsto
                  (\zeta_n^{a+b}u,y+aR)
\tag{C.SH.9}\label{cdp:eq:C.SH.9}
\]
descends. On the selected Tate tube $R=[\zeta_n]$ exactly,
by uniqueness of the torsion section in the smooth group
locus. In the multiplicative coordinate change $z=y$, $w=u/y$,
\eqref{cdp:eq:C.SH.9} is precisely
\[
                    (z,w)\longmapsto
                    (\zeta_n^a z,\zeta_n^b w).
\tag{C.SH.10}\label{cdp:eq:C.SH.10}
\]
Replace the tube by \ref{cdp:st:C.2.UN.41}'s two-Hopf plumbing using
its exact punctured-family isomorphism. The action extends
on the first Hopf copy by \eqref{cdp:eq:C.SH.10}, on the second by the
inverse weights, and on the first neck by
\[
 (z,x,y)\longmapsto
     (\zeta_n^a z,\zeta_n^b x,\zeta_n^{-b}y),
                         \qquad xy=t.
\tag{C.SH.11}\label{cdp:eq:C.SH.11}
\]
Interchange $a,b$ and the axes for the other neck.
Substitution in $(z,w)\mapsto(z^{-1},t/w)$ and
$(z,w)\mapsto(t/z,w^{-1})$ proves equivariance of both
gluing maps. The action is free on the whole central
support: on a double curve the equation
$\zeta_n^{ak}z=\alpha^\ell z$, or its counterpart, forces
$\ell=0$ by modulus and then $k=0$ by coprimality.
Off the double curves use the already verified Hopf
freeness. On a nearby torus it is a finite translation,
also free. Hence the quotient is smooth on this entire
tube and its two central components are the prescribed
secondary Hopf surface itself. The inverse weights
generate the same subgroup.

\noindent\textbf{Step 7.} Take the finite quotient of the entire compact threefold
and resolve it over the other finitely many bad values,
leaving this smooth tube unchanged. General fibres are
quotients of tori by translations, hence tori. Connectedness
passes through the finite quotient and the connected fibres
of the birational resolution. The distinguished reduced
SNC fibre has two Hopf components and two elliptic double
curves, and therefore Euler number zero. Every other
reduced fibre has even Euler number by \ref{cdp:st:C.2.UN.1}.
Apply \ref{cdp:st:C.2.UN.4} only over those other values to make their
Euler numbers zero; none of these operations changes
the specified Hopf surface or its distinguished tube.
This proves compact diagonal secondary occurrence over
the actual curve $B'$.

\par\medskip\noindent\textbf{The nonlinear scalar action.}
Use the intermediate compact family $\widetilde Y\to C$ in
\eqref{cdp:eq:C.NH.11} for the prescribed $H_r(\beta)$, before the
additional involution in \eqref{cdp:eq:C.NH.15}. Its smooth nonzero
tautological line space has coordinates
\[
 (z,w,x)\in(\mathbb C^2\setminus\{0\})\times\mathbb C,
                   \qquad p-p_0=x^r z,\quad s=xw.
\tag{C.SH.12}\label{cdp:eq:C.SH.12}
\]
Multiplication of that line by $\zeta_n$ is the actual map
\[
 (z,w,x)\longmapsto
       (\zeta_n^r z,\zeta_n w,\zeta_n^{-1}x).
\tag{C.SH.13}\label{cdp:eq:C.SH.13}
\]
It commutes with the lifted base translations and the
contracting deck map. It is free when $\gcd(r,n)=1$:
a nonzero $w$ or $x$ forces the corresponding root to be
one; if both vanish then $z\ne0$ and its $r$-th power
must be one. It stays free after the infinite cyclic
quotient. Away from the exceptional component the
canonical line trivialization turns equality with a
deck power into $\zeta_n^k=\beta^\ell$, so $\ell=0$.
On the Hopf component use its nonzero $w$ coordinate,
or its $z$ coordinate if $w=0$, for the same modulus
argument. Coprimality then finishes the proof.

\noindent\textbf{Step 8.} The resulting quotient is a smooth compact threefold.
Its desired component is exactly the second quotient
in \eqref{cdp:eq:C.SH.1}, with no change to that surface. On every
general torus fibre the action is a line-coordinate
translation, so the quotient is a torus. The distinguished
full fibre remains reduced SNC, and its Euler number
zero is preserved by the free finite quotient. Any
remaining values can be balanced by \ref{cdp:st:C.2.UN.4} away
from that fibre.

\noindent\textbf{Step 9.} The base for this construction is the specified compact
curve in \ref{cdp:st:C.2.UN.42}. No rational-base globalization is
needed for this occurrence assertion.
\end{proof}






\begin{proposition}[Every compact curve is the ruling base of an actual normalization component]\label{cdp:st:C.2.UN.45}
\StatementRecord{cdp:st:C.2.UN.45}{proved}{Historical events: Claude Fable 5.2 (Anthropic), effort high, 16 November 2026: carried as background without fault found. GPT-6 Astra (OpenAI), ChatGPT Pro mode, effort Pro, S6\_09\_PAPERS, 17 November 2026: targeted proof check. Also checked in this writing. Pre-continuation writing events: GPT-6 Astra (OpenAI), ChatGPT Pro mode, effort Pro, S6\_10\_PAPERS, 18 November 2026: targeted proof check. Continuation metadata events: GPT-6 Astra (OpenAI), ChatGPT Pro mode, effort Pro, S6\_10\_PAPERS continuation, 19 November 2026: used and rechecked in this run; status-vocabulary reconciliation; metadata reconciliation only; no new mathematical check. \emph{Pre-continuation recorded proof scope.} The primitive meromorphic function-field map gives the actual plane normalization. Each singular-point blowup lowers its arithmetic-genus defect by m(m-1)/2, proving termination and embedding of the specified curve. Its generic position in the multiplicity-one exceptional plane gives the claimed coefficient of the projectivized normal bundle. The blowup maps add exactly two or e(C), and balancing disjoint from the new component preserves it. The meromorphic frame gives the explicit ruled birational map; the integral projective-bundle map and coherent direct images give all ranks. The all-genus branch calculation was re-derived directly; historical code execution is not assigned new credit}{the hypotheses, notation and definitions stated here; \cref{cdp:st:C.2.UN.4}; \cref{cdp:eq:C.RU.7}; \cref{cdp:eq:C.RU.8}; \cref{cdp:eq:C.RU.5}}

For every smooth compact connected complex curve $C$, there is a smooth
compact complex threefold $Y$ and a proper map $f:Y\to\PP^1$ with connected
fibres, whose restriction over $\PP^1\setminus\{0\}$ is the product with a
fixed two-torus, and whose reduced zero fibre has Euler number zero and
contains a smooth multiplicity-one component
\[
                         R=\PP(N_{C/U})\longrightarrow C.
\tag{C.RU.1}\label{cdp:eq:C.RU.1}
\]
Here $C\hookrightarrow U$ is the explicit vertical embedding constructed
below. The component $R$ is retained unchanged by the Euler-balancing
operations. Its birational class is that of $C\times\PP^1$.
Consequently every ruled surface birational class occurs in an actual
compact zero-Euler torus fibre. A rational component is also obtained
unchanged as the exceptional plane of a point blowup followed by balancing
away from that plane.

Writing $g=g(C)$, the integral cohomology of $R$ is torsion-free and has ranks
\[
              (b_0,b_1,b_2,b_3,b_4)(R)=(1,2g,2,2g,1),
              \qquad h^\bullet(R,\OO_R)=(1,g,0).
\tag{C.RU.2}\label{cdp:eq:C.RU.2}
\]
The construction proves occurrence of every ruled birational class; it
does not prescribe an arbitrary rank-two bundle as the normal bundle in
\eqref{cdp:eq:C.RU.1}.
\end{proposition}



\begin{proof}
\noindent\textbf{Step 1.} Choose a nonconstant meromorphic function $x$ on $C$, supplied by
Riemann--Roch. It gives a finite map $C\to\PP^1$, so the meromorphic
function field of $C$ is a finite separable extension of $\C(x)$.
Choose a primitive element $y$. The meromorphic map
\[
                    C\longrightarrow\PP^2,\qquad [1:x:y]
\tag{C.RU.3}\label{cdp:eq:C.RU.3}
\]
extends at every point: multiply its three local meromorphic coordinates
by a power of a uniformizer so that their minimum vanishing order is zero.
At least one resulting holomorphic coordinate is nonzero. The image is an
integral projective plane curve $D$, and the map is finite and birational,
because $\C(x,y)=\mathcal M(C)$. Thus $C$ is the actual normalization of
$D$. For $C=\PP^1$ one may choose $x=s,y=s^2$, giving a smooth conic.
In every case choose projective coordinates so that $D$ is not the fixed
line used below.

\noindent\textbf{Step 2.} Start with $A\times\PP^1$, where $A$ is any fixed smooth complex two-torus,
and blow up a point of $A\times\{0\}$. Its exceptional divisor is an
actual $\PP^2$. Place $D$ in this plane. Successively blow up singular
points of its strict transform, performing each point blowup in the
ambient smooth threefold. The strict transform of the plane is precisely
its surface point blowup: in local coordinates this is the strict
transform of the coordinate plane in the three ordinary blowup charts.
The full map to $\PP^1$ is unchanged over its punctured base.

\noindent\textbf{Step 3.} Here the point-resolution process is finite without an unproved
termination assertion. For an integral curve $D_0$ on a smooth surface,
normalization gives
\[
  p_a(D_0)-g(C)=\operatorname{length}
                (\nu_*\OO_C/\OO_{D_0})\ge0.
\tag{C.RU.4}\label{cdp:eq:C.RU.4}
\]
At a singular point the multiplicity $m$ is at least two. Under a surface
point blowup, $D_1=\pi^*D_0-mE$ and
$K_{V_1}=\pi^*K_{V_0}+E$. Adjunction gives
\[
                        p_a(D_1)=p_a(D_0)-\frac{m(m-1)}2.
\tag{C.RU.5}\label{cdp:eq:C.RU.5}
\]
The normalization is unchanged through this proper birational map.
Therefore its total nonnegative normalization length decreases by at
least one at every singular-point blowup. After finitely many such
blowups the strict transform is smooth and is the specified $C$ itself.
This gives a genuine closed embedding $C\hookrightarrow U$ into the
resulting smooth compact threefold.

\noindent\textbf{Step 4.} The strict transform of the original exceptional plane has full-fibre
multiplicity one, since that was its coefficient after the first point
blowup and later strict transforms retain the coefficient. The curve
$C$ is not any boundary curve of this plane: it dominates the original
$D$, which is different from the line of intersection with the original
torus component, whereas every subsequently exceptional boundary curve
maps to a point of that plane. Consequently a generic point of $C$ lies
only in this multiplicity-one component. Blow up the smooth curve $C$ in
$U$. The exceptional divisor is exactly the projectivized normal bundle
in \eqref{cdp:eq:C.RU.1}, with coefficient one in the full fibre. This follows also
from the local charts of the blowup of $(x,y)$: its exceptional fibre at
each point of $C$ is the projective space of normal directions.

\noindent\textbf{Step 5.} Choose two meromorphic sections of $N_{C/U}$ which form a basis at the
generic point. Such sections exist by twisting by a sufficiently positive
line bundle on the curve and applying Riemann--Roch. They give an actual
bundle trivialization away from a finite set, and hence a birational map
\[
                         \PP(N_{C/U})\dashrightarrow C\times\PP^1.
\tag{C.RU.6}\label{cdp:eq:C.RU.6}
\]
The same meromorphic-frame map applies to every rank-two bundle on $C$.
It therefore identifies this birational class with that of every ruled
surface over this specified curve. In genus zero it also identifies the
rational class, for example through the affine-coordinate birational map
$\PP^2\dashrightarrow\PP^1\times\PP^1$.

\noindent\textbf{Step 6.} We retain the exact Euler changes, including all singularities along the
centres. The reduced inverse image of a point centre replaces that point
by $\PP^2$, adding two to the support Euler number. The reduced inverse
image of the smooth curve $C$ replaces each point by $\PP^1$, adding
$e(C)=2-2g$. These are the proper-map comparisons of (C.ZC.8), which do
not impose a normal-crossing or reduced full-fibre hypothesis. If $p$ is
the total number of point blowups, including the initial one, the support
after blowing up $C$ has
\[
                              e=2p+2-2g.
\tag{C.RU.7}\label{cdp:eq:C.RU.7}
\]
It is even. There is a smooth open part of the original torus component
disjoint from $R$, since only finitely many proper curves and points have
been affected there. Perform the positive point blocks or negative
point--quartic blocks of \ref{cdp:st:C.2.UN.4} in such an open part and in the new
exceptional components they create. Each changes Euler number by exactly
$2$ or $-2$. Choosing the indicated finite number makes \eqref{cdp:eq:C.RU.7} zero.
All these centres are disjoint from $R$, so the displayed smooth component,
its normal projection to $C$, and its coefficient one are unchanged.
All other fibres remain the original smooth tori. Every blowup has
connected fibres, so the composite map retains connected fibres as well.
Compactness and smoothness follow at each step from the actual point or
smooth-curve blowup of the compact smooth threefold.

\noindent\textbf{Step 7.} For the cohomology comparison use the lines convention in \eqref{cdp:eq:C.RU.1} and
put $\xi=c_1(\OO_{\PP(N)}(1))$. Pullback and cup product give the integral
projective-bundle map
\[
 H^k(C,\Z)\oplus H^{k-2}(C,\Z)
       \longrightarrow H^k(R,\Z),\qquad
 (a,b)\longmapsto\pi^*a+\xi\smile\pi^*b.
\tag{C.RU.8}\label{cdp:eq:C.RU.8}
\]
It is an isomorphism: on every fibre the classes $1,\xi$ are the standard
integral basis of $H^*(\PP^1,\Z)$, so the fibrewise basis glues and the
Leray--Hirsch comparison has these exact maps. The relation for this lines
convention is $\xi^2+\pi^*c_1(N)\xi=0$. Since the curve cohomology is
torsion-free, \eqref{cdp:eq:C.RU.8} proves the integral assertion. The fibrewise
calculation $H^0(\PP^1,\OO)=\C$, $H^1(\PP^1,\OO)=0$ gives
$\pi_*\OO_R=\OO_C$, $R^1\pi_*\OO_R=0$, with their restriction and
evaluation maps. Leray therefore gives the coherent tuple in \eqref{cdp:eq:C.RU.2}.

\noindent\textbf{Step 8.} There is an explicit all-genus computation supporting the construction.
For $g\ge1$ take the plane curve
\[
                 y^2z^{2g-1}=x^{2g+1}-z^{2g+1}.
\tag{C.RU.9}\label{cdp:eq:C.RU.9}
\]
It is smooth on $z\ne0$. At its single point at infinity, use
$u=x/y,v=z/y$; the equation is $v^{2g-1}(1+v^2)=u^{2g+1}$.
The holomorphic change $v'=v(1+v^2)^{1/(2g-1)}$ identifies its germ with
the plane branch $v'^{2g-1}=u^{2g+1}$. The two exponents are coprime.
Successive point blowups subtract the smaller exponent from the larger,
and the strict curve is smooth when the smaller exponent becomes one.
Summing \eqref{cdp:eq:C.RU.5} through these actual charts gives total normalization
defect $2g(g-1)$, so its normalization has genus $g$.
The executed code \path{code/root_ruled_occurrence.py} checks all exponent
charts, genus drops, Euler changes, balancing counts and the integral
projective-bundle ranks for $g=1,\ldots,50$. The displayed subtraction
and adjunction formulas prove the same identities for every $g$.
\end{proof}






\section{Marked Hopf, ruled and elliptic supports with all multiplicities}\label{sec:II:10}


\begin{proposition}[Hopf-axis adjacency with arbitrary multiplicities and point blowups]\label{cdp:st:C.2.UN.40}
\StatementRecord{cdp:st:C.2.UN.40}{proved}{Historical events: Claude Fable 5.2 (Anthropic), effort high, 16 November 2026: carried as background without fault found. GPT-6 Astra (OpenAI), ChatGPT Pro mode, effort Pro, S6\_09\_PAPERS, 17 November 2026: targeted proof check. Also checked in this writing. Pre-continuation writing events: GPT-6 Astra (OpenAI), ChatGPT Pro mode, effort Pro, S6\_10\_PAPERS, 18 November 2026: targeted proof check; used and rechecked in this run. Continuation metadata events: GPT-6 Astra (OpenAI), ChatGPT Pro mode, effort Pro, S6\_10\_PAPERS continuation, 19 November 2026: used and rechecked in this run; status-vocabulary reconciliation; metadata reconciliation only; no new mathematical check. \emph{Pre-continuation recorded proof scope.} Adjunction and the principal full fibre give the integral coefficient m\_i(k\_j+1)-m\_j(k\_i+1). The common function field defines its exact divisor-class pushforward through every point blowdown. Pullback through the actual unramified primary cover and the injective character-to-line-bundle map force both diagonal axes with equal ratios, or the unique nonlinear axis with the stated strict decrease. The shared elliptic curve survives every neighbouring point blowdown, proving the claimed adjacency exclusion}{the hypotheses, notation and definitions stated here; \cref{cdp:eq:C.HA.3}; \cref{cdp:eq:C.HA.4}; \cite[Theorem~2.1]{Khaled2024Hopf}}

Let $f:U\to\Delta$ be a proper degeneration of smooth complex
two-tori with smooth total space, and suppose that its reduced
central support has simple normal crossings. Write its full
fibre as $\sum_i m_iE_i$. Choose a generator of the rank-one
sheaf $f_*K_U$ and write its divisor as $\sum_i k_iE_i$.
Set $r_i=(k_i+1)/m_i$.

Suppose that the minimal model $H_i$ of a component $E_i$
is a Hopf surface of algebraic dimension zero. The contraction
$\rho_i:E_i\to H_i$ may consist of arbitrarily many point
blowdowns. Choose an actual finite unramified primary Hopf
cover $p:P\to H_i$.

If $P$ has diagonal contraction $(z,w)\mapsto(\alpha z,\beta w)$,
the images of both coordinate axes are boundary curves on
$H_i$: their strict transforms in $E_i$ meet another central
component. For each such strict transform, let $E_j$ be its
unique adjacent component at the generic point. Then
\[
                 m_i(k_j+1)-m_j(k_i+1)=0,
                 \qquad r_j=r_i.
\tag{C.HA.1}\label{cdp:eq:C.HA.1}
\]
If $P$ has non-diagonal contraction
$(z,w)\mapsto(\beta^r z+w^r,\beta w)$, $r\geq1$, the image
of its unique elliptic curve is a boundary curve. Its unique
generic neighbour satisfies
\[
       m_i(k_j+1)-m_j(k_i+1)=-m_i r,
       \qquad m_j(r_j-r_i)=-r<0.
\tag{C.HA.2}\label{cdp:eq:C.HA.2}
\]
These are integral divisor-class consequences on the actual
components. In particular a non-diagonal algebraic-dimension-zero
Hopf component cannot meet another algebraic-dimension-zero
Hopf component along this elliptic curve. Diagonal such
components can meet along their axes only at equal ratios.
\end{proposition}



\begin{proof}
\noindent\textbf{Step 1.} The direct image $f_*K_U$ is coherent, torsion-free and of
rank one on the disc, because a smooth torus fibre has one
holomorphic two-form. Shrink the disc and choose a generator.
Its evaluation is an actual holomorphic three-form on $U$.
On every punctured fibre it is a nonzero translation volume
form, so its zero divisor is vertical and gives
$K_U\simeq\mathcal O_U(\sum_i k_iE_i)$, with $k_i\geq0$.
For $j\ne i$ put $D_{ij}=E_j|_{E_i}$ and
$D_i=\sum_{j\ne i}D_{ij}$. Each is the reduced intersection
divisor specified by the simple-normal-crossing coordinates.
The full fibre is principal. Adjunction and its restriction
to $E_i$ therefore give the following integral identity:
\[
 m_i(K_{E_i}+D_i)
 \sim\sum_{j\ne i}
       \bigl[m_i(k_j+1)-m_j(k_i+1)\bigr]D_{ij}.
\tag{C.HA.3}\label{cdp:eq:C.HA.3}
\]
Indeed, first use
$K_{E_i}=((k_i+1)E_i+\sum_{j\ne i}k_jE_j)|_{E_i}$,
then replace $m_i E_i|_{E_i}$ by
$-\sum_{j\ne i}m_jD_{ij}$, and finally add $m_iD_i$.
Thus \eqref{cdp:eq:C.HA.3} has no rational divisor or divisibility convention.

\noindent\textbf{Step 2.} We specify its pushforward through the point blowdowns.
For a proper birational map of smooth surfaces the common
meromorphic function field sends the divisor of a function
to the divisor of the same function. Sending every prime
curve dominating a curve to that curve and every exceptional
prime curve to zero consequently induces a homomorphism
of divisor-class groups. It sends $K_{E_i}$ to $K_{H_i}$:
pull back a meromorphic two-form on $H_i$ and observe that
its additional Jacobian divisor is exceptional. This is a
divisor-class homomorphism; no equality with the coherent
direct image of an arbitrary line bundle is asserted.
Applying it to \eqref{cdp:eq:C.HA.3} gives
\[
 m_i(K_{H_i}+B_i)
 \sim\sum_{j\ne i}
       \bigl[m_i(k_j+1)-m_j(k_i+1)\bigr]\rho_{i*}D_{ij},
 \qquad B_i=\rho_{i*}D_i.
\tag{C.HA.4}\label{cdp:eq:C.HA.4}
\]
The divisor $B_i$ is reduced. A point-blowup morphism is an
isomorphism at the generic point of every curve on its target,
so there is exactly one strict transform dominating that curve.
At its generic point the SNC equation has exactly two branches,
$x_i x_j=0$; an additional branch can meet it only at isolated
triple points. Hence this strict transform occurs in exactly
one $D_{ij}$, with coefficient one.

\noindent\textbf{Step 3.} The pullback of \eqref{cdp:eq:C.HA.4} through the specified unramified
cover $p:P\to H_i$ is an identity of integral divisor classes,
and $p^*K_{H_i}=K_P$. We record precisely the curve and line
classes on $P$. The primary normal forms are those in
\cite[Theorem~2.1]{Khaled2024Hopf}. The finite meromorphic
field comparison proves $a(P)=a(H_i)=0$, by taking elementary
symmetric functions on the finite generic sheets, as in
(C.CF.9).

\noindent\textbf{Step 4.} Here is also a direct check of the curves needed in that
comparison. The inverse image of a compact curve in
$\mathbb C^2\setminus\{0\}$ has analytic closure across
the origin by the Remmert--Stein extension theorem. Its
divisor on $\mathbb C^2$ is principal, say $(F=0)$.
For the contraction $\gamma$ one has
$F\circ\gamma=uF$ with $u$ a holomorphic unit. Write
$u=u(0)e^h$, $h(0)=0$. The locally uniformly convergent sum
$a=\sum_{n\geq0}h\circ\gamma^n$ satisfies
$a-a\circ\gamma=h$. Replacing $F$ by $e^aF$ makes it an
eigenfunction of $\gamma^*$ with constant eigenvalue.

\noindent\textbf{Step 5.} For a diagonal contraction, distinct monomials can have
the same eigenvalue only when $\alpha^a\beta^b=1$ for
some nonzero integer pair. Because $|\alpha|,|\beta|<1$,
such a pair has opposite signs and gives a nonconstant
invariant meromorphic monomial, contrary to $a(P)=0$.
Thus each eigenfunction is one monomial, and the only
curves are the two axes $C_1,C_2$. In the non-diagonal
case use weighted degree $\deg z=r$, $\deg w=1$.
On the finite-dimensional space of weighted degree $d$
the contraction has eigenvalue $\beta^d$ times the
translation operator on polynomials in $z/w^r$.
Its eigenvectors are precisely the constants in that
variable. The eigenvalues $\beta^d$ for different $d$
are distinct, so every eigenfunction is a power of $w$.
There is exactly the curve $C=(w=0)$.

\noindent\textbf{Step 6.} For completeness, denote by $L_\lambda$ the line bundle
whose lifted fibre coordinate is multiplied by $\lambda$
under the contraction. The map
$\lambda\mapsto L_\lambda$ is injective. A trivialization
would be a nowhere-zero holomorphic eigenfunction on
$\mathbb C^2\setminus\{0\}$. Hartogs extends it and its
inverse across zero, and evaluation there forces
$\lambda=1$. Coordinate functions and the Jacobian
therefore give the exact classes
\[
\begin{array}{ll}
 \mathcal O(C_1)=L_\alpha,
 \quad\mathcal O(C_2)=L_\beta,
 &K_P=\mathcal O(-C_1-C_2),\\[2pt]
 \mathcal O(C)=L_\beta,
 &K_P=\mathcal O(-(r+1)C)
 \quad\hbox{in the non-diagonal case}.
\end{array}
\tag{C.HA.5}\label{cdp:eq:C.HA.5}
\]
The first two classes are integrally independent: a relation,
or even a torsion relation, would give the forbidden
positive-power resonance after raising to a power.
The class of $C$ in the second line has infinite order
because $0<|\beta|<1$.

Every curve on $H_i$ has nonempty reduced inverse image
on $P$, since $p$ is unramified and surjective. Thus the
pullback of $B_i$ is a reduced subset of the indicated
axis curves. In the diagonal case, compare coefficients
of their independent classes in the pullback of \eqref{cdp:eq:C.HA.4}.
If an axis were absent, its coefficient on the left
would be $-m_i$ and on the right zero. Each axis is
therefore present. Its coefficient on the left is then
zero, and the coefficient on the right is exactly that
of its unique generic neighbour. This proves \eqref{cdp:eq:C.HA.1},
including when one boundary curve has more than one
component in its finite inverse image.

In the non-diagonal case, absence of $C$ would instead
give the impossible nonzero class $-m_i(r+1)C\sim0$.
It is present, with coefficient one. Comparing its
infinite-order class gives
\[
       -m_i r=m_i(k_j+1)-m_j(k_i+1),
\tag{C.HA.6}\label{cdp:eq:C.HA.6}
\]
which is \eqref{cdp:eq:C.HA.2}. This proves both assertions on the
original component after all its point blowups, with
the actual finite-primary-cover comparison included.

Finally, the shared strict axis is a smooth elliptic
curve. It cannot be contracted by a point-blowdown
morphism on the neighbouring smooth component, since
all exceptional curves of such a morphism are rational.
If that neighbour also has an algebraic-dimension-zero
Hopf minimal model, this same curve is one of its
minimal axes. Its diagonal rule would force equal
ratios, contradicting \eqref{cdp:eq:C.HA.2}; its non-diagonal rule
would force the opposite strict inequality. Neither
is possible. The diagonal--diagonal case merely repeats
the equality \eqref{cdp:eq:C.HA.1}. The theorem gives these exact
adjacency restrictions and does not infer a global
nonexistence assertion from a local decrease.
\end{proof}






\begin{proposition}[The diagonal Hopf cycle criterion and the complete two-component case]\label{cdp:st:C.2.UN.46}
\StatementRecord{cdp:st:C.2.UN.46}{proved}{Historical events: Claude Fable 5.2 (Anthropic), effort high, 16 November 2026: carried as background without fault found. GPT-6 Astra (OpenAI), ChatGPT Pro mode, effort Pro, S6\_09\_PAPERS, 17 November 2026: targeted proof check. Also checked in this writing. Pre-continuation writing events: GPT-6 Astra (OpenAI), ChatGPT Pro mode, effort Pro, S6\_10\_PAPERS, 18 November 2026: targeted proof check. Continuation metadata events: GPT-6 Astra (OpenAI), ChatGPT Pro mode, effort Pro, S6\_10\_PAPERS continuation, 19 November 2026: used and rechecked in this run; status-vocabulary reconciliation; metadata reconciliation only; no new mathematical check. \emph{Pre-continuation recorded proof scope.} The actual entire axis line bundles give affine logarithmic edge lifts and integral lattice isomorphisms. Finite normal-product plumbing is Hausdorff and proper by the disjoint collars and compact cores. The product-preserving smooth tubular maps give the same injected semidirect-product fundamental group for every given smoothing. Product identity glues a closed holomorphic coframe, whose complete commuting flows produce the actual torus quotient. The two-matrix product forces precisely coordinate exchange and the two inversion gluings; a single Hopf-axis interchange cannot have identity product}{the hypotheses, notation and definitions stated here; \cref{cdp:st:C.2.UN.41}; \cref{cdp:st:C.2.UN.43}; \cref{cdp:st:C.2.UN.40}; \cref{cdp:eq:C.HC.3}; \cref{cdp:eq:C.HC.4}; \cref{cdp:eq:C.HC.7}; \cref{cdp:eq:C.HC.5}; \cref{cdp:eq:C.HC.8}; \cref{cdp:eq:C.HC.10}; \cref{cdp:eq:C.HC.9}}

Consider a reduced SNC fibre of an actual smooth-total
degeneration of two-tori, all of whose irreducible components
are unblown diagonal Hopf surfaces of algebraic dimension zero.
Primary and secondary Hopf surfaces are both allowed.
Its components form a cycle: each of the two axis elliptic
curves of each component is a double curve, and there are
no triple points.

For an oriented cycle of specified diagonal Hopf surfaces
$H_1,\ldots,H_\ell$, specify the actual elliptic gluing
isomorphisms and dual-normal isomorphisms along its edges.
Let
\[
 G_i=H_i\setminus(C_i^-\cup C_i^+)
       =\mathbb C^2/\Lambda_i,\qquad \Lambda_i\simeq\mathbb Z^3,
\tag{C.HC.1}\label{cdp:eq:C.HC.1}
\]
where $\Lambda_i$ is the inverse image of its diagonal
Hopf deck group under coordinatewise exponentiation.
The edge data determine isomorphisms of these integral
lattices and a cycle product $M:\Lambda_1\to\Lambda_1$.
For the finite proper local plumbing constructed below,
the exact criterion is
\[
 \text{$Y_t$ is a two-torus}\quad(0<|t|\ll1)
       \quad\Longleftrightarrow\quad M=I.
\tag{C.HC.2}\label{cdp:eq:C.HC.2}
\]
Every actual torus degeneration in the stated SNC branch
has $M=I$. The definition of $M$, its independence from
the choices of logarithms, and the actual nearby quotient
are given in the proof.

For two components, the criterion is more explicit.
The two Hopf surfaces are isomorphic. After an actual
diagonal automorphism of one of them, both gluing maps
are inversion on the corresponding axis elliptic curves.
Thus the reduced fibre is precisely two copies of the
same diagonal Hopf surface joined by the two inversion
maps. Every such reduced fibre occurs in a compact
zero-Euler torus filling, by \ref{cdp:st:C.2.UN.41} and \ref{cdp:st:C.2.UN.43}.
This completely determines the reduced supports in
this two-component branch. The constants in the
dual-normal isomorphisms remain explicit local
smoothing parameters in the quotient below.

The same local lattice criterion applies to the
specified operation of identifying the two distinct
axes of one Hopf surface. Its single transition
cannot have $M=I$, so that particular self-gluing
admits no smooth-total torus smoothing.
\end{proposition}



\begin{proof}
\noindent\textbf{Step 1.} \par\medskip\noindent\textbf{The cycle and its actual edge data.}
\ref{cdp:st:C.2.UN.40} forces both axis curves into the boundary of
every component of the given SNC fibre. An
algebraic-dimension-zero diagonal primary Hopf surface
has exactly these two curves, and its free diagonal
finite quotients have their two distinct images, by
\ref{cdp:st:C.2.UN.43}. They are disjoint. Consequently there is
no triple point, and at the generic point of either
axis the SNC equation has exactly one other branch.
Every vertex of the connected incidence graph has
degree two, so the graph is a cycle, with a pair
of edges allowed between two vertices.

\noindent\textbf{Step 2.} Orient this cycle. On $H_i$ choose diagonal coordinates
so that the incoming axis $C_i^-$ has coordinate $z_i$
and the outgoing axis $C_i^+$ has coordinate $w_i$.
Thus the outgoing normal coordinate is $z_i$, and
the incoming normal coordinate is $w_i$.
For the edge from $i$ to $i+1$ write its data as
\[
 \phi_i:C_i^+\xrightarrow{\sim}C_{i+1}^-,
 \qquad
 \theta_i:N_{C_i^+/H_i}\otimes
       \phi_i^*N_{C_{i+1}^-/H_{i+1}}
           \xrightarrow{\sim}\mathcal O_{C_i^+}.
\tag{C.HC.3}\label{cdp:eq:C.HC.3}
\]
These are holomorphic isomorphisms. In an actual
reduced fibre the second comes from the defining
function: locally $f=xy$ times a unit, and its
transverse bilinear term is precisely \eqref{cdp:eq:C.HC.3}.
There is no extra divisor in this isomorphism,
because there are no triple points.

\noindent\textbf{Step 3.} Each axis neighbourhood is an entire holomorphic
normal line bundle. For a primary diagonal Hopf
surface the outgoing example is
\[
 (\mathbb C^*_{w}\times\mathbb C_z)/
       ((w,z)\sim(\beta w,\alpha z)).
\tag{C.HC.4}\label{cdp:eq:C.HC.4}
\]
For a secondary surface take also the specified
finite diagonal quotient. Its action on the base
elliptic curve is free, since both finite characters
are faithful; the quotient of \eqref{cdp:eq:C.HC.4} is therefore
an actual line bundle on the quotient elliptic curve.
Removing its zero section gives exactly $G_i$.
Its universal cover is $\mathbb C^2$ in logarithmic
coordinates. The lattice $\Lambda_i$ includes
$2\pi i\mathbb Z^2$ and one logarithm of a contracting
deck transformation, together with the finite
diagonal roots. It is a discrete rank-three lattice,
and it spans $\mathbb C^2$ over $\mathbb C$.

\par\medskip\noindent\textbf{The affine edge maps and the integral product.}
For $t\ne0$, the normal-product equation and
\eqref{cdp:eq:C.HC.3} define a global biholomorphism
$\Phi_i(t):G_i\to G_{i+1}$: send the base by $\phi_i$
and send a nonzero normal vector to its dual vector
whose product is $t$. Lift it to the logarithmic
universal covers. An elliptic isomorphism lifts to
an affine map, so its base coordinate is
$c_i\log w_i+d_i$, with $c_i\ne0$.
The normal coordinate has the form
$-\log z_i+h_i(\log w_i)+\log t$.
The nonzero holomorphic fibre multiplier has a
holomorphic logarithm $h_i$ on the simply connected
base cover. For every elliptic period $\lambda$,
$h_i(u+\lambda)-h_i(u)$ is constant, since both
normal bundles in \eqref{cdp:eq:C.HC.4} have constant factors of
automorphy. Hence $h_i'$ descends to a holomorphic
function on a compact elliptic curve and is constant.
We obtain the exact lift
\[
 \binom{\log z_{i+1}}{\log w_{i+1}}
 =
 L_i\binom{\log z_i}{\log w_i}
 +\binom{d_i}{e_i+\log t},
 \qquad
 L_i=\begin{pmatrix}0&c_i\\-1&k_i\end{pmatrix}.
\tag{C.HC.5}\label{cdp:eq:C.HC.5}
\]
This derivation uses the actual entire line bundles
and their isomorphism; it does not replace an
analytic neighbourhood by a formal normal bundle.

\noindent\textbf{Step 4.} Since $\Phi_i(t)$ is a global biholomorphism,
$L_i\Lambda_i=\Lambda_{i+1}$. Choosing integral bases
gives actual matrices $M_i\in\operatorname{GL}_3(\mathbb Z)$.
Put
\[
 M=M_\ell\cdots M_1,\qquad L=L_\ell\cdots L_1.
\tag{C.HC.6}\label{cdp:eq:C.HC.6}
\]
Changing the initial basis conjugates $M$.
Changing any logarithm changes only its translation
by a lattice vector. Multiplying $\theta_i$ by a
nonzero constant changes $e_i$, not $L_i$.
Thus the condition $M=I$ is well defined.
It is equivalent to $L=I$, because $\Lambda_1$
spans the complex vector space.

\par\medskip\noindent\textbf{Actual finite plumbing and properness.}
Use \eqref{cdp:eq:C.HC.4} to choose disjoint tubular neighbourhoods
of the two axes in each Hopf component. They have
flat Hermitian normal metrics. The dual-normal
isomorphism in \eqref{cdp:eq:C.HC.3} identifies the two flat
metrics up to a positive constant: the ratio of
compatible flat metrics is the exponential of a
harmonic function on the compact elliptic curve.
Adjust that constant so that the product of the
two normal norms is $|t|$ in each neck.

\noindent\textbf{Step 5.} For each edge take the open part of
$N_{C_i^+/H_i}\oplus N_{C_i^+/H_i}^{-1}$
with bounded fibre coordinates $x,y$, and the
actual holomorphic function $xy=t$.
Its two branches are identified with the two
specified axis neighbourhoods using $\phi_i,\theta_i$.
Glue these finitely many necks to the products
of the compact component cores with a disc.
Explicitly choose $0<r_0<R_0$ and $|t|<r_0^2$,
use $\|x\|,\|y\|<R_0$, and attach where one of
these norms is greater than $r_0$.
The inverse map inserts $t/x$ or $t/y$.
The two ends of a neck cannot overlap, because
$\|x\|\|y\|=|t|<r_0^2$; the two axis collars in
a component are disjoint as well.

\noindent\textbf{Step 6.} These are actual biholomorphic gluing maps.
Their graphs are closed relative to the pieces:
at one boundary of an attaching annulus its core
point is absent, and at the other boundary its
neck point is absent. Each equivalence class
meets at most one neck and one core. The finite
gluing is therefore Hausdorff. The intermediate
closed cores and closed smaller neck disc bundles
used in \ref{cdp:st:C.2.UN.41} cover the inverse image of every
closed subdisc; each is compact over a compact
elliptic curve or compact Hopf core. This proves
properness. The result is a smooth complex
threefold over a disc, with exactly the prescribed
reduced normal-crossing central surface.

\par\medskip\noindent\textbf{The exact fundamental group, including actual given smoothings.}
The open group $G_i=\mathbb C^2/\Lambda_i$ is
smoothly $T^3\times\mathbb R$.
For example, a linear height annihilating
$\operatorname{span}_{\mathbb R}\Lambda_i$ has
compact three-torus levels; the usual axis norms
are exponentials of this height. A component
core and each annular neck consequently retract
onto a three-torus. Every collar inclusion
induces the indicated isomorphism on its
rank-three fundamental group. Van Kampen,
first along a spanning tree of the cycle and
then at its last edge, gives the actual group
\[
                       \pi_1(Y_t)=\Lambda_1\rtimes_M\mathbb Z.
\tag{C.HC.7}\label{cdp:eq:C.HC.7}
\]
The copy of $\Lambda_1$ is injected, as in the
displayed semidirect product presentation.

\noindent\textbf{Step 7.} For an arbitrary actual smoothing of the same
reduced SNC data, this calculation still uses
actual smooth maps. Near a double curve $C$,
the divisor line bundles of its two branches
have canonical sections $s_1,s_2$, and $f$
trivializes their tensor product. Choose a
smooth retraction to $C$ and smooth line-bundle
identifications with the pullbacks of their
restrictions, choosing the second identification
so that the product trivialization is preserved.
The map consisting of the retraction and
$(s_1,s_2)$ has invertible derivative along $C$.
After shrinking the neighbourhood it is a
smooth tubular diffeomorphism to the two
normal line bundles, and its actual function
is $xy$. These are the normal bundles and
the product isomorphism \eqref{cdp:eq:C.HC.3}.
On a compact core away from these neighbourhoods
$f$ is a submersion; horizontal transport
identifies its nearby cores. The attaching
maps are isotopic to the normal maps just
specified, by radial contraction of the
tubular identifications. Thus their maps on
$\pi_1$ are exactly $M_i$, proving \eqref{cdp:eq:C.HC.7}
for the given smoothing as well. No holomorphic
linearization of its full neighbourhood is
asserted or needed.

\noindent\textbf{Step 8.} If the given $Y_t$ is a two-torus, its fundamental
group is abelian. The injected copy of $\Lambda_1$
in \eqref{cdp:eq:C.HC.7} then forces $M=I$.
This proves the necessary half of \eqref{cdp:eq:C.HC.2}.

\par\medskip\noindent\textbf{Closed one-forms give the sufficient torus comparison.}
Suppose $M=I$, equivalently $L=I$.
The vector of closed holomorphic one-forms
$(d\log z_i,d\log w_i)$ on $G_i$ transforms
by the constant matrix $L_i$ on each overlap.
Choose constant invertible matrices $B_i$
successively so that $B_{i+1}L_i=B_i$.
The identity of the cycle product makes this
choice consistent on its last edge.
These forms therefore glue on each actual
nearby compact fibre, including all its necks,
to two everywhere independent closed holomorphic
one-forms.

\noindent\textbf{Step 9.} Their dual holomorphic vector fields commute,
because the two dual one-forms are closed.
On a compact complex manifold holomorphic
vector fields have entire complex-time flows:
a finite coordinate cover supplies a uniform
local existence interval, which can be iterated.
The two commuting flows give a transitive
holomorphic $\mathbb C^2$-action; every orbit
is open, and the fibre is connected.
Its stabilizer is discrete and cocompact,
since the infinitesimal action is an
isomorphism and the fibre is compact.
The orbit map is therefore the actual
biholomorphism
\[
                  Y_t\simeq\mathbb C^2/\Lambda(t)
\tag{C.HC.8}\label{cdp:eq:C.HC.8}
\]
for a rank-four real lattice.

\noindent\textbf{Step 10.} One may read this lattice without discarding
the translation data. Compose the affine
maps \eqref{cdp:eq:C.HC.5} on a chosen logarithm branch:
\[
                    \Phi_\ell(t)\cdots\Phi_1(t)(v)
                          =v+b_0+v_0\log t.
\tag{C.HC.9}\label{cdp:eq:C.HC.9}
\]
In the normalization $B_1=I$, the period
integrals on the core generators are
$\Lambda_1$ and the cycle generator has
period $b_0+v_0\log t$. Equation \eqref{cdp:eq:C.HC.7}
lists all generators, so the exact lattice is
\[
                 \Lambda(t)=
                    \Lambda_1+\mathbb Z(b_0+v_0\log t).
\tag{C.HC.10}\label{cdp:eq:C.HC.10}
\]
Its discreteness and full real rank follow
from the preceding transitive-action proof,
rather than being assumed. A turn of $\log t$
changes the last vector by $2\pi i v_0\in\Lambda_1$,
as follows directly from the globally defined
edge maps. Hence \eqref{cdp:eq:C.HC.8}--\eqref{cdp:eq:C.HC.10} descend
to the actual punctured-disc family. This proves
the sufficient half and records every smoothing
translation parameter.

\par\medskip\noindent\textbf{One and two components.}
A single transition in \eqref{cdp:eq:C.HC.5} cannot be the
identity, because its upper-left entry is zero.
Thus the specified self-identification of the
two distinct axes cannot give a torus smoothing.

\noindent\textbf{Step 11.} For two components multiply the two matrices:
\[
 L_2L_1=
 \begin{pmatrix}
   -c_2&c_2k_1\\
   -k_2&-c_1+k_2k_1
 \end{pmatrix}.
\tag{C.HC.11}\label{cdp:eq:C.HC.11}
\]
It is the identity precisely when
$c_1=c_2=-1$ and $k_1=k_2=0$.
Consequently $L_1=L_2=-S$, where $S$ exchanges
the two coordinates, and
$\Lambda_2=-S\Lambda_1$.
Exponentiation gives
$\Gamma_2=(S\Gamma_1)^{-1}=S\Gamma_1$,
since these are groups. The actual coordinate
exchange $(z,w)\mapsto(w,z)$ therefore descends
to a biholomorphism $H_1\to H_2$, including
their finite diagonal subgroups.

\noindent\textbf{Step 12.} After this coordinate exchange, the two
elliptic identifications have the form
$z\mapsto a/z$ and $w\mapsto b/w$.
The actual diagonal automorphism
$(Z,W)\mapsto(a^{-1}Z,b^{-1}W)$ of the second
Hopf copy commutes with its complete deck group
and makes both identifications inversion.
It follows that the reduced central surface
is isomorphic to exactly the support in
\ref{cdp:st:C.2.UN.41}, or its free finite quotient in
\ref{cdp:st:C.2.UN.43}. Those constructions give the asserted
compact realizations and retain this entire
reduced surface. The constants in $\theta_i$
are recorded by \eqref{cdp:eq:C.HC.9}--\eqref{cdp:eq:C.HC.10}; the
compact realization assertion here identifies
the reduced support, while the explicit finite
plumbing realizes each specified local choice.
\end{proof}






\begin{proposition}[The full compatible cycle is realized over a compact rational base]\label{cdp:st:C.2.UN.48}
\StatementRecord{cdp:st:C.2.UN.48}{proved}{Historical events: Claude Fable 5.2 (Anthropic), effort high, 16 November 2026: carried as background without fault found. GPT-6 Astra (OpenAI), ChatGPT Pro mode, effort Pro, S6\_09\_PAPERS, 17 November 2026: targeted proof check. Also checked in this writing. Pre-continuation writing events: GPT-6 Astra (OpenAI), ChatGPT Pro mode, effort Pro, S6\_10\_PAPERS, 18 November 2026: targeted proof check. Continuation metadata events: GPT-6 Astra (OpenAI), ChatGPT Pro mode, effort Pro, S6\_10\_PAPERS continuation, 19 November 2026: used and rechecked in this run; status-vocabulary reconciliation; metadata reconciliation only; no new mathematical check. \emph{Pre-continuation recorded proof scope.} The contracting logarithms put all normal meridians on one signed half-space, so their actual sum has a primitive elliptic quotient. The original open group excludes trivial line parameter by its explicit embedded elliptic curve. The nodal-pencil suspension, exact Tate coordinate and moving integral lattice retain every primary or secondary component. Two independent normal constants solve the complete affine period equation, giving an actual punctured-family isomorphism and proper compact annular replacement. The three-component matrices, contraction signs, period and secondary divisibility formula were re-derived directly; historical grid-execution counts retain their historical scope}{the hypotheses, notation and definitions stated here; \cref{cdp:st:C.2.UN.46}; \cref{cdp:eq:C.HG.3}; \cref{cdp:eq:C.HG.5}; \cref{cdp:st:C.2.UN.40}; \cref{cdp:st:C.2.UN.43}; \cref{cdp:eq:C.HG.7}; \cref{cdp:st:C.2.UN.41}; \cref{cdp:eq:C.HG.6}; \cref{cdp:eq:C.HG.12}; \cref{cdp:eq:C.HG.10}; \cref{cdp:eq:C.HG.13}; \cref{cdp:eq:C.HG.14}; \cref{cdp:eq:C.HG.11}}

Take the specified cycle of primary or secondary diagonal
Hopf surfaces of algebraic dimension zero in \ref{cdp:st:C.2.UN.46},
with its specified elliptic gluing isomorphisms.
Suppose its dual-normal bundles are isomorphic to the
trivial bundle, and compute the integral edge product
from those actual isomorphisms as in that statement.
Every cycle whose product is the identity occurs,
with precisely its given components and gluing maps,
as a reduced SNC fibre of a holomorphic map
\[
                 f:Y\longrightarrow\PP^1
\]
from a smooth compact complex threefold with connected
two-torus general fibre. Every reduced fibre of this
map has Euler characteristic zero.

In particular the necessary criterion of \ref{cdp:st:C.2.UN.46}
is sufficient for compact occurrence of the whole
reduced support, for arbitrary cycle length and for
secondary components as well. The construction makes
specified choices of the normal-product constants;
it preserves the entire reduced surface, while
recording those choices and the actual punctured-family
comparison below.

For example there is a three-component compact
realization with primary Hopf multipliers
\[
 \begin{gathered}
 (e^{-1},e^{-\sqrt2}),\qquad
 (e^{-\sqrt2},e^{-1-\sqrt2}),\qquad
 (e^{-1-\sqrt2},e^{-1}),\\
 L_1=\begin{pmatrix}0&-1\\-1&-1\end{pmatrix},\qquad
 L_2=\begin{pmatrix}0&1\\-1&1\end{pmatrix},\qquad
 L_3=\begin{pmatrix}0&-1\\-1&1\end{pmatrix}.
 \end{gathered}
 \tag{C.HG.1}\label{cdp:eq:C.HG.1}
\]
Here $L_3L_2L_1=I$ and the actual monodromy period
on the common open group is $(t^{-2},t^2)$.
Thus the longer-cycle occurrence assertion includes
an explicit case beyond the two-component inversion
construction.
\end{proposition}



\begin{proof}
\noindent\textbf{Step 1.} \par\medskip\noindent\textbf{The common lattice and the sign of the normal periods.}
Retain the exact logarithmic coordinates of \ref{cdp:st:C.2.UN.46}.
The product identity permits constant matrices $P_i$
identifying every open group with one common group
\[
                  G=\C^2/\Lambda,\qquad \Lambda\simeq\Z^3,
                  \qquad P_{i+1}L_i=P_i.
\tag{C.HG.2}\label{cdp:eq:C.HG.2}
\]
Take indices cyclically. Write $n_i=P_i(2\pi i,0)$
for the outgoing normal meridian. Since
$L_i(1,0)=(0,-1)$, the incoming meridian on the
$i$th component is $-n_{i-1}$, and
\[
 P_i=\left(\frac{n_i}{2\pi i},
                   -\frac{n_{i-1}}{2\pi i}\right),
 \qquad
 2\pi i v_0=-\sum_i n_i=:u\in\Lambda.
\tag{C.HG.3}\label{cdp:eq:C.HG.3}
\]
The last equality follows by adding the actual
edge translations $P_{i+1}(d_i,e_i+\log t)$.
It retains their signs, rather than choosing a
meridian sign afterwards.

\noindent\textbf{Step 2.} Put $W=\operatorname{span}_{\R}\Lambda$, and choose
a nonzero real linear functional $\ell:\C^2\to\R$
with kernel $W$. On the $i$th Hopf surface choose
its actual contracting generator with multipliers
$\alpha_i,\beta_i$, including when a finite diagonal
subgroup is also present. Its common logarithm
\[
 c_i=P_i(\log\alpha_i,\log\beta_i)
       =\frac{\log\alpha_i}{2\pi i}n_i
          -\frac{\log\beta_i}{2\pi i}n_{i-1}
\]
belongs to $\Lambda$. Applying $\ell$ and using
$\ell(n_i)=\ell(n_{i-1})=0$ gives the exact real
identity
\[
 \log|\alpha_i|\,\ell(i n_i)
       =\log|\beta_i|\,\ell(i n_{i-1}).
\tag{C.HG.4}\label{cdp:eq:C.HG.4}
\]
Both logarithms are strictly negative.
Neither displayed value of $\ell$ is zero: if
both vanished, $\ell$ would vanish on the complex
span of $n_i,n_{i-1}$, which is $\C^2$ by
\eqref{cdp:eq:C.HG.3}. The identity would make both zero if
one were zero. It follows successively around
the cycle that all $\ell(i n_i)$ have the same
nonzero sign. Consequently
\[
                       \ell(iu)=-\sum_i\ell(i n_i)\ne0.
\tag{C.HG.5}\label{cdp:eq:C.HG.5}
\]
In particular $u\ne0$, and $\C u\cap W=\R u$.

\noindent\textbf{Step 3.} Write $u=d p$, where $p$ is a primitive element
of the free lattice $\Lambda$ and $d>0$.
The real quotient map
$W/\R p\to\C^2/\C p$ is an isomorphism, by
\eqref{cdp:eq:C.HG.5}. Therefore the image of $\Lambda$ in
$\C^2/\C p$ is a discrete rank-two lattice.
Its quotient is an actual elliptic curve.
Extend $p$ to an integral basis $p,\lambda,c$
of $\Lambda$. Choose a linear quotient
coordinate $X$ with $X(\lambda)=2\pi i$.
The two projected periods are real-linearly
independent, so $\Re X(c)\ne0$; reverse $c$
if necessary to make it negative. Choose a
second complex linear coordinate $Z$ with
$Z(p)=2\pi i$ and $Z(\lambda)=0$. These
conditions determine it because $p,\lambda$
are complex-linearly independent. In these
actual linear coordinates the common lattice
is
\[
 \Lambda=\left\langle
     (0,2\pi i),\ (2\pi i,0),\
     (\log q_0,\log a)\right\rangle_{\Z},
 \qquad 0<|q_0|<1,\qquad a\in\C^*,
 \qquad u=d(0,2\pi i).
\tag{C.HG.6}\label{cdp:eq:C.HG.6}
\]
The logarithms in this formula are the actual
chosen coordinates of $c$.

\noindent\textbf{Step 4.} Here $a\ne1$. Indeed, if $a=1$, subtract an
integral multiple of the first period from $c$
so that its second coordinate is zero. The map
\[
 \C/\langle2\pi i,\log q_0\rangle_{\Z}
          \longrightarrow G,\qquad [X]\longmapsto[(X,0)]
\tag{C.HG.7}\label{cdp:eq:C.HG.7}
\]
is well defined and injective: equality in $G$
forces the coefficient of $(0,2\pi i)$ to vanish,
leaving exactly the two periods in its source.
It is an immersion and, by compactness, an embedded
elliptic curve. But $G$ is the complement of the
two axes in any of the prescribed Hopf surfaces.
The curve calculation of \ref{cdp:st:C.2.UN.40} and \ref{cdp:st:C.2.UN.43}
says that these are all its compact curves.
The image of \eqref{cdp:eq:C.HG.7} contradicts that calculation.
This proves the precise specialization needed
for the pencil construction below.

\par\medskip\noindent\textbf{A compact elliptic suspension with the required limiting periods.}
The cubic construction in \ref{cdp:st:C.2.UN.41} works for
every specified $a\in\C^*\setminus\{1\}$:
its choice of the smooth nodal point uses no
inequality on $|a|$. Explicitly use
\[
 \begin{gathered}
 y^2z=x^2(x+z),\qquad O=[0:1:0],\\
 P_0=\left(\frac{4a}{(1-a)^2},
                  \frac{4a(1+a)}{(1-a)^3}\right)
 \end{gathered}
\tag{C.HG.8}\label{cdp:eq:C.HG.8}
\]
in the affine chart $z=1$, whose smooth group
coordinate is $a$. The point is distinct from
$O$ and from the node. Choose a second cubic
through $O,P_0$, transverse at nine distinct
smooth points of the nodal cubic and missing
its node. The jet-surjectivity and nonempty
open conditions in \ref{cdp:st:C.2.UN.41} give this choice.
Blowing up the nine actual base points gives
a compact rational elliptic surface with
sections $O,P$ and the selected $I_1$ fibre.

\noindent\textbf{Step 5.} Make the global base change $\PP^1_s\to\PP^1$
which is $s\mapsto s^d$ at that value.
Normalize, resolve and take the relatively
minimal elliptic model $J\to\PP^1_s$.
Near the chosen node the equation before
resolution is $xy=s^d$. Its resolution has
charts $(u_j,v_j)$, $0\le j<d$, with exact maps
\[
 x=u_j^{j+1}v_j^j,\qquad
 y=u_j^{d-j-1}v_j^{d-j},\qquad s=u_jv_j.
\]
Adjacent charts glue by
$u_{j+1}=v_j^{-1}$, $v_{j+1}=u_jv_j^2$.
These maps are isomorphisms off the node and
insert the chain of $d-1$ projective lines
between its two branches; each displayed
central equation is $u_jv_j=0$.
Together with the old component
this gives an $I_d$ fibre. The sections avoid
these charts and retain their given values.
At all other values these operations are
proper vertical modifications of a compact
elliptic surface. The section translation by
$P$ extends to an automorphism of this
relatively minimal model, by the exact graph
and relative-canonical-nef argument of
\ref{cdp:st:C.2.UN.41} and \ref{cdp:st:C.2.UN.43}.

\noindent\textbf{Step 6.} The analytic Tate comparison at this fibre
gives, on a sufficiently small punctured disc,
\[
       J_s=\C^*/Q(s)^{\Z},\qquad
       \operatorname{ord}_0Q=d,\qquad P_s=[a(s)],
       \qquad a(0)=a.
\tag{C.HG.9}\label{cdp:eq:C.HG.9}
\]
The lift $a(s)$ is a holomorphic unit, since
the pulled-back section meets the smooth
identity component. Choose the holomorphic
$d$th root of the unit $Q(s)/s^d$ and use the
actual coordinate
$t=s\bigl(Q(s)/s^d\bigr)^{1/d}$.
It has nonzero derivative at zero and makes
$Q=t^d$ exactly. Denote the section lift in
this coordinate by $a(t)$, and fix the
holomorphic logarithm with value $\log a$.

\noindent\textbf{Step 7.} Now take the actual suspension
\[
 Y^{\rm std}=
  (\C^*_b\times J)/
      \langle(b,y)\longmapsto(q_0 b,y+P)\rangle
          \longrightarrow\PP^1_s .
\tag{C.HG.10}\label{cdp:eq:C.HG.10}
\]
The action is free and properly discontinuous
because its first coordinate has modulus
multiplier $|q_0|<1$. A closed annulus in
$\C^*$ times compact $J$ meets every orbit,
so the quotient is compact and smooth.
Its general connected fibres are two-tori.
Near the distinguished value their exact
logarithmic periods are
\[
 \Lambda(t)+\Z(0,d\log t),\qquad
 \Lambda(t)=
 \left\langle(0,2\pi i),(2\pi i,0),
                  (\log q_0,\log a(t))\right\rangle_{\Z}.
\tag{C.HG.11}\label{cdp:eq:C.HG.11}
\]
This is the displayed suspension quotient,
not just equality of characteristic polynomials.

\par\medskip\noindent\textbf{Transporting the prescribed Hopf surfaces and the whole cycle.}
View each original $n_i$ and contracting
$c_i$ as its actual integral combination
of the three basis periods in \eqref{cdp:eq:C.HG.6}.
Use those same integer coefficients in
$\Lambda(t)$ to define $n_i(t),c_i(t)$.
All integral relations are retained exactly.
In particular $\sum_i n_i(t)=-d(0,2\pi i)$.
Set
\[
 \begin{gathered}
 P_i(t)=\left(\frac{n_i(t)}{2\pi i},
                    -\frac{n_{i-1}(t)}{2\pi i}\right),\\
 \Lambda_i(t)=P_i(t)^{-1}\Lambda(t),
 \quad L_i(t)=P_{i+1}(t)^{-1}P_i(t).
 \end{gathered}
\tag{C.HG.12}\label{cdp:eq:C.HG.12}
\]
These matrices are holomorphic and invertible
after shrinking the disc. The first column
of $L_i(t)$ is $(0,-1)$, so it has exactly the
form
$\left(\begin{smallmatrix}0&c_i^*(t)\\-1&k_i^*(t)\end{smallmatrix}\right)$
required for the normal-product map.
Its cycle product is the identity.

\noindent\textbf{Step 8.} The group $\Lambda_i(t)/(2\pi i\Z^2)$ is
the same abstract finitely generated group
as at zero, by the fixed integral relations.
Its selected free generator is the actual
diagonal contraction with logarithms
$P_i(t)^{-1}c_i(t)$. Both moduli remain
strictly between zero and one on a smaller
disc. Every finite torsion generator has
a relation $n v=A(2\pi i,0)+B(0,2\pi i)$
with the same integers as at zero. Therefore
its two roots of unity and their coprimality
conditions remain unchanged. These are
actual holomorphic families of the prescribed
primary or secondary Hopf quotients. The
contractions are uniformly contracting over
a closed smaller disc, and their fixed
finite groups remain free. Thus the
quotients of $(\C^2\setminus\{0\})\times\Delta$
are smooth proper families, with precisely
the original surfaces at zero.

\noindent\textbf{Step 9.} Let $d_i$ be the constants specifying the
original elliptic gluing maps in the
coordinates of \ref{cdp:st:C.2.UN.46}. Keep these constants.
Their normal multipliers have freely
adjustable nonzero constants, whose logarithms
we denote by $e_i(t)$. The full affine cycle
translation in common coordinates is
\[
 \sum_iP_{i+1}(t)
        \binom{d_i}{e_i(t)+\log t}
       =b_0(t)+(0,d\log t).
\tag{C.HG.13}\label{cdp:eq:C.HG.13}
\]
Two adjacent coefficient vectors of the
unknowns $e_i(t)$ are
$-n_i(t)/(2\pi i)$ and $-n_{i+1}(t)/(2\pi i)$.
They are complex-linearly independent by
\eqref{cdp:eq:C.HG.12}. Fix all the other $e_i(t)$, and
solve the resulting invertible two-by-two
linear system holomorphically to obtain
\[
                              b_0(t)=0.
\tag{C.HG.14}\label{cdp:eq:C.HG.14}
\]
Only the normal-product constants change.
The central surfaces and every given
elliptic gluing map remain exactly the
prescribed ones.

\noindent\textbf{Step 10.} Use the actual varying normal line bundles
and these affine edge maps to perform the
finite plumbing of \ref{cdp:st:C.2.UN.46}, with $xy=t$
in each neck. The flat metrics, disjoint
axis neighbourhoods and compact cores may
be chosen uniformly on a smaller parameter
disc; the strict inequalities used in that
proof persist by continuity. The result
$\mathcal Y\to\Delta_t$ is a smooth proper
threefold whose reduced central fibre is
the given cycle, with every component
of multiplicity one.

\noindent\textbf{Step 11.} There is an actual punctured-family
biholomorphism to \eqref{cdp:eq:C.HG.10}. To see the map
without suppressing affine data, choose one
core as origin in the common logarithmic
coordinates. On consecutive cores use
$P_i(t)$ and subtract the accumulated edge
translations from \eqref{cdp:eq:C.HG.13}. Use the same
formulas on the necks, where $xy=t$.
They agree on successive overlaps. On the
last overlap their difference is the
period $(0,d\log t)$ by \eqref{cdp:eq:C.HG.14};
changing logarithm branches changes them
by $\Lambda(t)$. Hence they define a
holomorphic map to the quotient with
periods \eqref{cdp:eq:C.HG.11}. Fibrewise this is the
closed-coframe comparison of \ref{cdp:st:C.2.UN.46},
which is a biholomorphism: its periods
are all the generators of the actual
fundamental group. Its differential is
invertible in the fibre directions and
it is the identity on the base, so it
is a local biholomorphism of total
spaces as well. The fibrewise inverses
therefore give the actual holomorphic
inverse over the punctured disc.

\par\medskip\noindent\textbf{Compact replacement and Euler numbers.}
Replace the distinguished tube in
$Y^{\rm std}$ by $\mathcal Y$, using this
biholomorphism on a punctured annulus.
More precisely remove a closed smaller
disc from the base of the first family
and glue the second over a larger disc.
Their annular overlap is the displayed
isomorphism. Properness of both families
and two nested radii give a Hausdorff
gluing: at either boundary of the annulus
only one of the two pieces contains the
limiting base point. Closed smaller
pieces in these two proper families
cover the resulting space, proving
compactness. Smoothness and the
holomorphic map to the same $\PP^1_s$
are preserved. Its fibres are connected
and its general fibre is the specified
suspension two-torus.

\noindent\textbf{Step 12.} The new central support has Euler number
zero, since every Hopf component and every
double elliptic curve has Euler number
zero and there are no triple points.
For every other base value the support
is unchanged from \eqref{cdp:eq:C.HG.10}. Its map to
the actual elliptic curve $\C^*/q_0^{\Z}$
is a locally trivial bundle with fibre
the reduced curve $J_s^{\rm red}$ and
transition the automorphism induced by
$P$. A finite triangulation, or the
Euler characteristic of its finite
local-system cellular complex, gives
Euler number
$e(\C^*/q_0^{\Z})e(J_s^{\rm red})=0$.
Thus every reduced fibre has Euler
number zero, as asserted.

\noindent\textbf{Step 13.} Finally \eqref{cdp:eq:C.HG.1} is verified by taking,
in the common angular lattice,
\[
 n_1=2\pi i(1,0),\quad
 n_2=2\pi i(1,-1),\quad
 n_3=2\pi i(0,-1),\quad
 c=(-1,-\sqrt2).
\tag{C.HG.15}\label{cdp:eq:C.HG.15}
\]
The three chosen contracting logarithms
are $P_1^{-1}c$, $-P_2^{-1}c$ and
$-P_3^{-1}c$, exactly the pairs in
\eqref{cdp:eq:C.HG.1}. Their ratios are irrational,
so the three Hopf surfaces have
algebraic dimension zero. The edge
maps in multiplicative coordinates,
with constant factors one, are
\[
 \begin{gathered}
 (z,w)\mapsto(w^{-1},t z^{-1}w^{-1}),\\
 (z,w)\mapsto(w,t z^{-1}w),\\
 (z,w)\mapsto(w^{-1},t z^{-1}w).
 \end{gathered}
\tag{C.HG.16}\label{cdp:eq:C.HG.16}
\]
Their composite is $(z,w)\mapsto(t^{-2}z,t^2w)$.
The exact integral matrices, negative
contraction logarithms, primitive
period divisor $d=2$ and the two-constant
translation solution are checked in
\path{code/D_hopf_cycle_matrices.py}
and its saved output.

\noindent\textbf{Step 14.} The finite-secondary divisibility is also
explicit. For a repeated two-component
cycle of length $2k$ with finite weights
$(a,b)$ modulo $n$, the common angular
lattice is
$\Z^2+\Z(a,b)/n$, and the period vector is
$k(-1,1)$ in units of $2\pi i$.
Its intersection with $\R(-1,1)$ is
$\Z(-1,1)/\gcd(n,a+b)$: the sum of the
coordinates requires $n\mid j(a+b)$ for
a representative $j(a,b)/n$, and the
coprimality of $a$ to $n$ then gives
exactly those fractions on this line.
Consequently the precise positive divisor is
\[
                    d=k\gcd(n,a+b).
\tag{C.HG.17}\label{cdp:eq:C.HG.17}
\]
The sign is fixed by the actual outgoing
vectors $(1,0)$ and $(0,-1)$.
For example weights $(1,4)$ modulo $5$
give $d=5$ for two components.
In the three-component example, the
finite weights $(1,2)$ modulo $5$ become
$(1,2),(3,2),(2,4)$ on its three actual
Hopf quotients. They are all free, the
common period still has divisor $2$,
and the same construction gives this
compact secondary three-cycle.
The saved computation checks 12,180
repeated-cycle cases and 1,376 free
secondary three-cycles by exact
Hermite normal forms.
\end{proof}






\begin{proposition}[Exact nonlinear Hopf star criterion and compact realization]\label{cdp:st:C.2.UN.49}
\StatementRecord{cdp:st:C.2.UN.49}{proved}{Historical events: Claude Fable 5.2 (Anthropic), effort high, 16 November 2026: carried as background without fault found. GPT-6 Astra (OpenAI), ChatGPT Pro mode, effort Pro, S6\_09\_PAPERS, 17 November 2026: targeted proof check. Also checked in this writing. Pre-continuation writing events: GPT-6 Astra (OpenAI), ChatGPT Pro mode, effort Pro, S6\_10\_PAPERS, 18 November 2026: targeted proof check. Continuation metadata events: GPT-6 Astra (OpenAI), ChatGPT Pro mode, effort Pro, S6\_10\_PAPERS continuation, 19 November 2026: used and rechecked in this run; status-vocabulary reconciliation; metadata reconciliation only; no new mathematical check. \emph{Pre-continuation recorded proof scope.} The normal-root charts define the exact inverse logarithmic transform, including its full analytic torsor. The determinant-one meridian matrix identifies every Hopf core with the same integral fibre class. Canonical adjunction forces the elliptic base and removes unmarked multiples; the nonzero integral of the leading two-form on that common class forces equal resonances. The zero Chern vector gives the actual rank-four translation lattice. Its compact twistor family, exact four-period presentation, weighted line insertion and logarithmic cluster estimate provide simultaneous proper replacements preserving every component and gluing map}{the hypotheses, notation and definitions stated here; \cref{cdp:eq:C.NS.1}; \cref{cdp:eq:C.NS.3}; \cref{cdp:eq:C.NS.2}; \cref{cdp:eq:C.NS.5}; \cref{cdp:eq:C.NS.4}; \cref{cdp:eq:C.NS.6}; \cref{cdp:eq:C.NS.9}; \cref{cdp:eq:C.NS.10}; \cref{cdp:eq:C.NS.8}; \cref{cdp:eq:C.NS.11}; \cref{cdp:eq:C.NS.15}; \cref{cdp:eq:C.NS.18}; \cref{cdp:eq:C.NS.20}; \cref{cdp:eq:C.NS.19}}

Let $S$ be a specified compact SNC surface with components
$P,H_1,\ldots,H_n$, $n\ge1$, and no triple points. Suppose
$q:P\to B$ is a connected fixed-Jacobian elliptic quasi-bundle:
its only nonsmooth scheme fibres are multiples of smooth
elliptic curves, and its integral $H^1$ monodromy is the
identity. Each $H_i=H_{r_i}(\beta_i)$ is a nonlinear primary
Hopf surface as in (C.NH.1). It meets $P$ along its unique
axis, identified by a specified holomorphic isomorphism
with a distinct reduced fibre $J_i$ of $q$. There are no
other intersections.

The following assertions are equivalent.

\begin{enumerate}
\item $S$ is the reduced full central fibre of a proper
one-parameter degeneration of two-tori with smooth total space.
\item $S$ is the reduced full special fibre of a smooth compact
threefold over $\PP^1$, all of whose other fibres are smooth
two-tori.
\item The normal lines on every double curve are dual;
$B$ is elliptic; all $r_i$ equal one integer $r$; the full
marked fibres are $rJ_i$ and there are no other multiple
fibres; and the actual inverse logarithmic transform
$T_0\to B$ defined in \eqref{cdp:eq:C.NS.1}--\eqref{cdp:eq:C.NS.3} has zero integral
Chern vector in $H^2(B,\Lambda_F)$.
\end{enumerate}

Here $F$ is the fixed elliptic Jacobian and $\Lambda_F$ its
integral period lattice. The inverse transform retains
the entire analytic torsor of $P$ and the given axis
identifications. It is unique up to the explicit product
translations described below. Thus the criterion concerns
the specified surface $S$, including its component
moduli and gluing maps. Every fibre support in assertion
(2) has Euler number zero. No extension of the original
smoothing's elliptic quotient is assumed.
\end{proposition}



\begin{proof}
\noindent\textbf{Step 1.} \par\medskip\noindent\textbf{The actual inverse transform.}
We first define the operation whenever the marked
multiplicity is $r_i$ and the two normal lines are dual.
Write the Hopf axis as $\C^*/\beta_i^{r_i\Z}$, with
coordinate $z$. Its normal has multiplier $\beta_i$.
The matching neighborhood in $P$ has the precise model
\[
 p=x^{r_i}z,\qquad
 (z,x)\sim(\beta_i^{r_i}z,\beta_i^{-1}x).
                                                        \tag{C.NS.1}\label{cdp:eq:C.NS.1}
\]
Indeed, normalize the $r_i$-th root of $q$ at this fibre.
The result is a smooth principal $F$-bundle over a disc,
hence is trivial there. The free root action is an
order-$r_i$ translation on $F$, together with the
primitive normal character. A finite averaging gauge
makes that translation independent of the root
parameter. The dual-normal isomorphism identifies this
connected cyclic cover with the one coming from the
Hopf axis. In multiplicative coordinates the cover is
$z=w^{r_i}$, $x=u/w$, with deck $w\mapsto\beta_i w$;
its quotient is exactly \eqref{cdp:eq:C.NS.1}, with $p=u^{r_i}$.
This also proves that the normal of a smooth
multiple fibre has exact order its multiplicity.

\noindent\textbf{Step 2.} On the punctured base disc the map
\[
       (z,x)\longmapsto(p,w)=(x^{r_i}z,x^{-1}),
       \qquad z=pw^{r_i},
       \qquad [w]\in F_i=\C^*/\beta_i^{\Z}
                                                        \tag{C.NS.2}\label{cdp:eq:C.NS.2}
\]
is an actual product trivialization. Replace \eqref{cdp:eq:C.NS.1}
by $F_i\times\Delta_p$ through \eqref{cdp:eq:C.NS.2}, at every
marked point. The original torsor on the complement
and these exact transition maps define
\[
                             q_0:T_0\longrightarrow B.
                                                        \tag{C.NS.3}\label{cdp:eq:C.NS.3}
\]
When there are no unmarked multiple fibres, this is a
smooth principal $F$-bundle. All identifications of
$F_i$ with $F$ are the ones furnished by the given
normal covers and the original torsor.

\noindent\textbf{Step 3.} The construction is independent of the permitted local
choices. Two dual-normal isomorphisms differ by a
nonzero holomorphic function on the compact elliptic
axis and hence by a constant. This changes $w$ by a
constant translation of $F_i$, which extends over the
product disc. An automorphism of a root-trivialized
principal $F$-bundle that extends across its central
curve is a holomorphic translation there, followed by
a fixed linear automorphism. If it preserves the
specified normal cover, that translation is invariant
under the root action and descends to a holomorphic
translation on the ordinary disc. It again extends
over the product filling. The specified axis map
fixes the remaining linear identification. Thus
\eqref{cdp:eq:C.NS.3} retains the full holomorphic torsor class;
it is not defined merely by a rational Chern number.

\par\medskip\noindent\textbf{The integral boundary matrix and the smoothing topology.}
We compute all meridians before applying a smoothing
argument. In \eqref{cdp:eq:C.NS.1}, let $a$ be the positive angular
loop of $z$, let $b$ be its deck loop, and let
$\mu_P$ be the positive angular loop of $x$. On the
punctured product \eqref{cdp:eq:C.NS.2}, let $\delta$ be the
positive base loop with $w$ fixed, $e$ the angular
loop of $w$, and $h$ its $\beta_i$-deck loop. Then
\[
 \begin{pmatrix}a&b&\mu_P\end{pmatrix}_{(\delta,e,h)}
       =
 \begin{pmatrix}
  1&0&r_i\\
  0&0&-1\\
  0&1&0
 \end{pmatrix};
 \qquad
 a=\delta,\quad b=h,\quad \mu_P=r_i\delta-e.
                                                        \tag{C.NS.4}\label{cdp:eq:C.NS.4}
\]
The matrix is integral of determinant one. For example,
holding $p$ fixed while $w$ winds once gives
$x=w^{-1}$ and $z=pw^{r_i}$, proving
$e=r_i a-\mu_P$. The formulas retain the normal
character, not just its order.

\noindent\textbf{Step 4.} The nonlinear Hopf pair is smoothly equivalent, for
this boundary calculation, to the diagonal pair with
contraction $(\beta_i^{r_i}z,\beta_i w)$. One exact
comparison on the complement is
\[
 U=\frac{z}{w^{r_i}},\qquad
 V=U-\frac{\beta_i^{-r_i}}{\log|\beta_i|}\log|w|.
                                                        \tag{C.NS.5}\label{cdp:eq:C.NS.5}
\]
The deck map sends $U$ to $U+\beta_i^{-r_i}$ and
$w$ to $\beta_iw$, so $V$ is invariant. This gives
$H_i\setminus J_i\simeq\C_V\times F_i$ smoothly.
The map
$z\mapsto z-\beta_i^{-r_i}w^{r_i}\log|w|/\log|\beta_i|$
extends continuously as the identity on the axis and
preserves orientation; it is smooth on every boundary
torus. Equivalently, vary the resonant coefficient from
one to zero in the contraction and use a uniform
compact fundamental annulus. This is a smooth proper
family of pairs, so the same integral boundary
identification follows by relative Ehresmann.

\noindent\textbf{Step 5.} Let $\mu_H$ be the Hopf normal meridian. The
smoothing equation is $xw=t$, so
$\mu_H=-\mu_P$. The regular fibre cycles on the
Hopf side are
\[
                   e_H=r_i a+\mu_H=e,\qquad h_H=b=h.
                                                        \tag{C.NS.6}\label{cdp:eq:C.NS.6}
\]
The complementary solid torus $F_i\times D^2$ in
\eqref{cdp:eq:C.NS.5} kills $a=\delta$. Consequently replacing
the tubular neighborhood in $P$ by the Hopf
complement is precisely the inverse logarithmic
transform \eqref{cdp:eq:C.NS.2}, with its indicated integral
framing.

\noindent\textbf{Step 6.} For completeness, this calculation applies to every
actual SNC smoothing of the specified support.
Remove small tubular neighborhoods of the disjoint
double curves. On each remaining compact component
core the original map is a submersion through zero;
Ehresmann with boundary identifies that core with
its nearby copy. On the collars the local equations
are exactly $xw=t$ after absorbing a holomorphic
unit. The product of the two normal lines is
trivialized by this same function. Two such
trivializations along a compact double curve differ
by a constant in $\C^*$, whose rotation is isotopic
to the identity. Thus the nearby four-manifold is
the indicated tubular plumbing, and \eqref{cdp:eq:C.NS.4}--\eqref{cdp:eq:C.NS.6}
identify it with the underlying four-manifold of
$T_0$. Under this identification, the oriented
generator of $H_2(H_i\setminus J_i,\Z)$ is the
ordinary fibre class of $q_0$, for every $i$.
The identifications between the normal covers are
holomorphic elliptic isomorphisms, hence have degree
one and preserve that orientation. In particular all
these transported core cycles are the same integral
two-cycle in the nearby fibre, even when their
primitive angular directions in $\Lambda_F$ differ.

\noindent\textbf{Step 7.} Here is an explicit global version of the Chern test.
Trivialize the principal $F$-bundle on the punctured
base and compare that trivialization with \eqref{cdp:eq:C.NS.2}
around the $i$-th disc. Let $\nu_i\in\Lambda_F$ be
the winding of its translation transition, with
positive disc-boundary orientation. Then
\[
 c(T_0)=\sum_i\nu_i,\qquad
 \mu_{P,i}=r_i\delta_{\rm ext}
                    +(r_i\nu_i-e_i).
                                                        \tag{C.NS.7}\label{cdp:eq:C.NS.7}
\]
The first equality is the integral obstruction to
extending these boundary trivializations over the
discs. It can alternatively be used as the sign
convention for the Chern vector. The second equality
is \eqref{cdp:eq:C.NS.4} after $\delta_i=\delta_{\rm ext}+\nu_i$.
Thus the rational Seifert vector in this convention
is $\sum_i(\nu_i-e_i/r_i)$, and the inverse transform
adds exactly $\sum_i e_i/r_i$. The holomorphic
transition functions themselves, rather than only
their windings, remain in \eqref{cdp:eq:C.NS.3}.

\par\medskip\noindent\textbf{Canonical orders, the elliptic base, and equal resonance.}
Suppose first that $S$ occurs as in assertion (1).
The equation $xw=t$ proves duality of the normal
lines. Their exact torsion orders then give the
marked multiplicities $r_i$.

\noindent\textbf{Step 8.} Choose a generator of the rank-one direct image of
the relative canonical line over the disc. Its
evaluation has a vertical divisor
$k_PP+\sum_i k_iH_i$. It is nowhere zero on every
nonzero torus fibre. Adjunction on the Hopf component
and $P|_{H_i}=J_i$, $H_i|_{H_i}=-J_i$ give
\[
 (k_P-k_i)J_i=K_{H_i}+J_i=-r_iJ_i,
       \qquad k_i-k_P=r_i.
                                                        \tag{C.NS.8}\label{cdp:eq:C.NS.8}
\]
The last equality is an equality of integers:
$\OO_{H_i}(J_i)$ has infinite order, since a
positive effective divisor on a compact complex
surface cannot have a trivial positive tensor power.
Adjunction on $P$ now gives the actual line equality
\[
                         K_P=\OO_P\Bigl(\sum_i(r_i-1)J_i\Bigr).
                                                        \tag{C.NS.9}\label{cdp:eq:C.NS.9}
\]

\noindent\textbf{Step 9.} For the given quasi-bundle the root charts give
$q_*K_P=K_B$ and
$K_P=q^*K_B+\sum_{\text{all multiple fibres}}(m-1)J$.
Indeed, the local form is the base differential
wedged with the invariant elliptic differential;
on the root chart its exact zero is $(m-1)J$,
and taking the fibrewise constant Laurent
coefficients gives the displayed pushforward.
Equation \eqref{cdp:eq:C.NS.9} implies $h^0(K_P)\ge1$, so
$g(B)\ge1$. Subtracting the marked terms from
the two canonical formulas gives
\[
 q^*K_B+\sum_{\text{unmarked multiple fibres}}(m-1)J=0
                    \quad\hbox{as a line bundle}.       \tag{C.NS.10}\label{cdp:eq:C.NS.10}
\]
For $g(B)\ge2$ a nonzero canonical section on $B$
has a nonzero effective divisor, and its pullback
together with the other terms would be a nonzero
effective divisor linearly equivalent to zero on
the compact surface. This is impossible. Therefore
$B$ is elliptic, and the same argument in \eqref{cdp:eq:C.NS.10}
removes every unmarked multiple fibre.

\noindent\textbf{Step 10.} Let $\omega_t$ be the evaluated relative two-form
divided by $t^{k_P}$. On a compact core of $H_i\setminus J_i$
its leading term is $t^{r_i}\eta_i$, where
\[
 \eta_i=c_i\,\frac{dz\wedge dw}{w^{r_i+1}},
                         \qquad c_i\ne0.                \tag{C.NS.11}\label{cdp:eq:C.NS.11}
\]
To verify this exact pole order, write the total
canonical evaluation in a double chart as
$x^{k_i}w^{k_P}$ times a holomorphic unit and a
coordinate three-form, then divide by $dt$ with
$t=xw$. Equation \eqref{cdp:eq:C.NS.8} gives \eqref{cdp:eq:C.NS.11}.
The resulting coefficient on the entire Hopf
component is a nowhere-zero section of
$K_{H_i}((r_i+1)J_i)=\OO_{H_i}$, so it is the
nonzero constant $c_i$.

\noindent\textbf{Step 11.} On the logarithmic cover of $H_i\setminus J_i$,
put $U=z/w^{r_i}$ and $W=\log w$. Its two deck
periods are
$(0,2\pi i)$ and $(\beta_i^{-r_i},\log\beta_i)$.
Thus the integral of $\eta_i=c_i\,dU\wedge dW$
on its oriented core torus is
$\pm2\pi i c_i\beta_i^{-r_i}\ne0$.
Transport that compact torus through the
submersion near the core. By \eqref{cdp:eq:C.NS.6} all such
cycles are the same class $\gamma$ in the
nearby fibre. Therefore the single holomorphic
period satisfies, for every $i$,
\[
        \int_\gamma\omega_t
             =t^{r_i}\bigl(\pm2\pi i c_i\beta_i^{-r_i}
                              +O(t)\bigr).              \tag{C.NS.12}\label{cdp:eq:C.NS.12}
\]
Its vanishing order is simultaneously every
$r_i$. Hence all resonances equal $r$.
This uses the transported two-cycle; it makes
no assumption about parallel primitive angular
directions or an elliptic quotient of the
nearby complex torus.

\noindent\textbf{Step 12.} By the plumbing comparison the underlying
four-manifold of $T_0$ is a real four-torus.
For a principal elliptic bundle over the
elliptic curve, the homotopy sequence injects
$\Lambda_F$ into its fundamental group, and
the exact presentation is
\[
 \langle a,b,e,h:\ e,h\text{ central},\
             [a,b]=e^{c_1}h^{c_2}\rangle,
                         \quad(c_1,c_2)=c(T_0).          \tag{C.NS.13}\label{cdp:eq:C.NS.13}
\]
Its fundamental group is abelian only if
$c(T_0)=0$. This proves every necessity in (3).

\par\medskip\noindent\textbf{The inverse transform is then an actual torus.}
Assume (3). The sheaf sequence
$0\to\Lambda_F\to\OO_B\to F\to0$ lifts the
principal torsor class of $T_0$ to $H^1(B,\OO_B)$
because its connecting class is precisely
$c(T_0)=0$. On $B=\C/\Gamma_B$ represent this
class by a constant $(0,1)$-form. Integrating it
on the universal cover gives the actual
translation lattice
\[
 (p,z)\mapsto(p+\gamma,z+a\bar\gamma),\quad
 (p,z)\mapsto(p,z+\lambda),
       \qquad \gamma\in\Gamma_B,\quad\lambda\in\Lambda_F.
                                                        \tag{C.NS.14}\label{cdp:eq:C.NS.14}
\]
Changes of the chosen Cech cocycle are fibre
translations and give an isomorphism with the
original torsor, exactly as in (C.KX.14).
The lattice is discrete of real rank four.
Consequently $T_0$ is a complex two-torus,
with its specified principal elliptic quotient
and all the marked fibres retained.

\par\medskip\noindent\textbf{A compact twistor family with the exact local period presentations.}
Write $T_0=\C^2/\Lambda$ in coordinates $(p,z)$
so that the common elliptic fibre direction is
$p=0$. On the affine $t$-chart of $\PP^1$ put
\[
 \ell_t(p,z)=(p+t\bar z,\ z-t\bar p).
                                                        \tag{C.NS.15}\label{cdp:eq:C.NS.15}
\]
This is real-linear and invertible for every
$t$, since its inverse has denominator
$1+|t|^2$. The values $\ell_t(v)$ for
$v\in\Lambda$ are holomorphic sections of
$\OO_{\PP^1}(1)^{\oplus2}$: on the other chart
they are $(s p+\bar z,s z-\bar p)$, $s=t^{-1}$.
Translate the vector bundle by these sections
and take the quotient. Each fibre lattice has
real rank four. The inverse real-linear maps
give local uniform fundamental parallelepipeds
over the compact base, proving a free proper
action and a smooth compact quotient
$\mathcal T\to\PP^1$. Its fibres are two-tori
and its fibre at zero is the actual $T_0$.
Pull this family back by $t=s^r$; the pullback
remains smooth and compact because the
original map is a submersion.

\noindent\textbf{Step 13.} For the $i$-th Hopf identification choose the
oriented fibre-lattice basis
$\lambda_i=(0,a_i)$, $\mu_i=(0,b_i)$ so that
$\beta_i=\exp(2\pi i b_i/a_i)$. This is the
basis supplied by \eqref{cdp:eq:C.NS.2}, not a replacement
of the given Hopf modulus. On the vector
space with coordinates $(P,Z)$ after
\eqref{cdp:eq:C.NS.15}, define
\[
 p_i=P-s^r\frac{\bar a_i}{a_i}Z,\qquad
 z_i=\frac{2\pi i}{a_i}Z.
                                                        \tag{C.NS.16}\label{cdp:eq:C.NS.16}
\]
The $\lambda_i$-period is $(0,2\pi i)$ and
the $\mu_i$-period is
\[
   \bigl(\delta_i(s),\log\beta_i\bigr),\qquad
   \delta_i(s)=s^r
        \left(\bar b_i-\frac{\bar a_i b_i}{a_i}\right).
                                                        \tag{C.NS.17}\label{cdp:eq:C.NS.17}
\]
The coefficient is nonzero because $a_i,b_i$
are an oriented real basis of a complex line.

\noindent\textbf{Step 14.} Complete this pair to an integral basis of
$\Lambda$ by two lifts of the base lattice.
Their $p_i$-periods define a smooth local
family of elliptic curves $B_{i,s}$, and
their exponentiated $z_i$-periods define a
holomorphic family of flat degree-zero line
bundles $M_i$ on $B_{i,s}$. Thus, on a
sufficiently small disc, the actual torus
family has the exact presentation
\[
 \mathcal T_s
   =M_{i,s}^{\times}/\langle g_i\rangle,\qquad
 g_i(p,\ell)=(p+\delta_i(s),\beta_i\ell).
                                                        \tag{C.NS.18}\label{cdp:eq:C.NS.18}
\]
This equality follows by quotienting the
same four displayed lattice generators,
first by the three excluding $\mu_i$.
It retains the flat part of the torsor,
even when that part is not torsion.

\par\medskip\noindent\textbf{The local weighted insertion with the varying flat line.}
At the marked point of $B_{i,0}$ perform
the weight-$(r,1)$ blowup of the relative
elliptic surface $B_i\to\Delta_s$, and replace
$M_i$ by its pullback tensored with
$\OO(R_i)$. The nonzero line has the same
smooth coarse chart as (C.NH.6)--(C.NH.7).
If $\delta_i(s)=d_i s^r$, the lifted deck
map in a flat local frame is exactly
\[
 (z,w,x)\mapsto
   \bigl(\beta_i^r(z+d_iw^r),
              \beta_iw,\beta_i^{-1}x\bigr),
           \qquad p-p_i=x^rz,\quad s=xw.
                                                        \tag{C.NS.19}\label{cdp:eq:C.NS.19}
\]
The change $Z=z/(\beta_i^r d_i)$ gives
the specified $H_r(\beta_i)$ on $x=0$.
On $w=0$ the chart is precisely \eqref{cdp:eq:C.NS.1};
hence the other component is obtained
from $T_0$ by the inverse of \eqref{cdp:eq:C.NS.2}.

\noindent\textbf{Step 15.} We check properness with the varying
flat line. On a small closed base disc,
choose smoothly varying unitary flat
metrics on $M_{i,s}$. If its period
characters are $e^{Z_j(s)}$ and base
periods $P_j(s)$, solve the two real
equations
$\operatorname{Re}(Z_j-c(s)P_j)=0$.
Their determinant is nonzero, so $c(s)$
is smooth and bounded. In the resulting
metric, the norm multiplier of the flat
lift in \eqref{cdp:eq:C.NS.18} is
$|\beta_i|\exp(-\operatorname{Re}(c(s)\delta_i(s)))$.
It is uniformly less than one after
shrinking the disc. The additional
exceptional line has the two-sided
polynomial distortion (C.NH.9)--(C.NH.10).
Exponential contraction still dominates
that polynomial uniformly on the closed
disc. The compact-annulus proof of
(C.NH.11) therefore gives a smooth
proper filling $\mathcal Y_i\to\Delta_s$.
Its infinite cyclic action is free,
and every prescribed $\beta_i$ is
retained.

\noindent\textbf{Step 16.} Off the exceptional divisor the canonical
section of $\OO(R_i)$ identifies the new
line with $M_i$, equivariantly for the
two deck maps. Consequently there is
the actual isomorphism
\[
       \mathcal Y_i\setminus H_i
            \simeq \mathcal T|_{\Delta_s}\setminus F_i,
                                                        \tag{C.NS.20}\label{cdp:eq:C.NS.20}
\]
where $F_i$ is the marked central
elliptic fibre in $T_0$. On the remaining
central component this is exactly the
punctured product map \eqref{cdp:eq:C.NS.2}.

\par\medskip\noindent\textbf{Simultaneous replacement and compactness.}
The central curves $F_i$ are disjoint.
We justify that \eqref{cdp:eq:C.NS.20} allows
independent replacements near them,
although its generic tori need not
have an elliptic quotient.

\noindent\textbf{Step 17.} If a sequence approaches $H_i$ in
$\mathcal Y_i$, use a compact Hopf
fundamental domain in \eqref{cdp:eq:C.NS.19}.
Then $z,w$ stay bounded, do not
simultaneously approach zero, and
$x\to0$. Its image in the original
line presentation has
$p-p_i=x^rz$, $s=xw$, and flat
line coordinate proportional to $x^{-1}$.
To put that coordinate in a fixed
compact norm annulus requires
$O(1+|\log|x||)$ deck iterates.
Their total base displacement is
$O(|s|^r(1+|\log|x||))$, which
tends to zero. Hence every cluster
point of the images lies in $F_i$.
Conversely properness of $\mathcal Y_i$
shows that images approaching a point
outside $F_i$ have only their unique
inverse-image limit outside $H_i$.

\noindent\textbf{Step 18.} Choose disjoint relatively compact
neighborhoods $W_i$ of $F_i$ in
$\mathcal T$ over a small base disc.
The set consisting of $H_i$ and the
inverse image of $W_i\setminus F_i$
under \eqref{cdp:eq:C.NS.20} is open, by the
cluster-set calculation. Glue these
open replacements to
$\mathcal T\setminus\bigcup_i F_i$.
Their identification graphs are
relatively closed by that same
calculation, and the replacements
have disjoint central neighborhoods.
Thus the finite gluing is Hausdorff
and smooth. Nested compact
neighborhoods inside the $W_i$,
properness of each $\mathcal Y_i$,
and the compact original complement
cover every closed inverse image
of a sufficiently small base disc.
The resulting map is proper.

\noindent\textbf{Step 19.} The central main component is exactly
$P$: outside the marked fibres its
torsor is unchanged, and at every
marked fibre \eqref{cdp:eq:C.NS.19} reverses
the precise map \eqref{cdp:eq:C.NS.2}.
The Hopf components are the prescribed
$H_i$. Any constant axis translation
in their identification is realized
by the Hopf automorphism
$(z,w)\mapsto(a^rz,aw)$; the linear
part has already been retained in
the chosen fibre-lattice basis.
Thus all specified elliptic gluing
maps, not just their degrees, are
the original ones. The full central
equations are $xw=s$, so the fibre
is reduced SNC.

\noindent\textbf{Step 20.} Finally carry out these replacements
inside the compact pulled-back
twistor family over $\PP^1$ and
leave its complement unchanged.
The preceding proper local gluing
and the original compact complement
give a smooth compact total space.
Every fibre away from zero remains
the original two-torus. The central
support has Euler number zero since
each elliptic quasi-bundle, Hopf
component, and double elliptic curve
has Euler number zero. We have
proved (3)$\Rightarrow$(2), including
the full compact-total-space and
torsor comparisons. Assertion
(2)$\Rightarrow$(1) is restriction
to a disc, completing the equivalence.
\end{proof}






\begin{proposition}[Complete primary and secondary nonlinear Hopf star criterion]\label{cdp:st:C.2.UN.50}
\StatementRecord{cdp:st:C.2.UN.50}{proved}{Historical events: Claude Fable 5.2 (Anthropic), effort high, 16 November 2026: carried as background without fault found. GPT-6 Astra (OpenAI), ChatGPT Pro mode, effort Pro, S6\_09\_PAPERS, 17 November 2026: targeted proof check. Also checked in this writing. Pre-continuation writing events: GPT-6 Astra (OpenAI), ChatGPT Pro mode, effort Pro, S6\_10\_PAPERS, 18 November 2026: targeted proof check. Continuation metadata events: GPT-6 Astra (OpenAI), ChatGPT Pro mode, effort Pro, S6\_10\_PAPERS continuation, 19 November 2026: used and rechecked in this run; status-vocabulary reconciliation; metadata reconciliation only; no new mathematical check. \emph{Pre-continuation recorded proof scope.} The Bezout coordinates u=z\textasciicircum{}n and X=x*z\textasciicircum{}k are the actual free coarse quotient. Their determinant-one punctured product map retains the full integer n in the meridian r*delta-n*e. The finite-invariant Hopf complement gives the same core fibre class and its nonzero period is divided by n. Tensoring the weighted exceptional line by its nth power produces exactly the specified free finite action, infinite contraction and main-component chart. The polynomial norm and cluster bounds give the compact simultaneous construction with the original torsor}{the hypotheses, notation and definitions stated here; \cref{cdp:st:C.2.UN.49}; \cref{cdp:st:C.2.UN.43}; \cref{cdp:eq:C.SN.4}; \cref{cdp:eq:C.SN.5}; \cref{cdp:eq:C.SN.6}; \cref{cdp:eq:C.SN.9}}

In \ref{cdp:st:C.2.UN.49} allow each prescribed leaf to be
\[
 H_i=H_{r_i}(\beta_i)/\mu_{n_i},\qquad
 \zeta\cdot(z,w)=(\zeta^{r_i}z,\zeta w),\qquad
                         \gcd(r_i,n_i)=1,               \tag{C.SN.1}\label{cdp:eq:C.SN.1}
\]
including $n_i=1$. These are all nonlinear primary and
secondary Hopf surfaces of algebraic dimension zero,
by \ref{cdp:st:C.2.UN.43}. The same equivalence between an actual
smooth-total torus smoothing and compact realization
over $\PP^1$ holds. Its necessary-and-sufficient tests
are dual normal lines, an elliptic central base with
no unmarked multiple fibres, one common resonance
$r_i=r$, and zero integral Chern vector of the
inverse logarithmic transform defined below.

The marked full fibres in $P$ still have multiplicity
$r_i$, while the finite integer $n_i$ enters the
exact inverse-transform framing. In a global
punctured-base trivialization its meridian is
\[
       \mu_{P,i}=r_i\delta_{\rm ext}
                         +(r_i\nu_i-n_i e_i),\qquad
                         c(T_0)=\sum_i\nu_i.             \tag{C.SN.2}\label{cdp:eq:C.SN.2}
\]
Thus its rational Seifert vector is
$\sum_i(\nu_i-n_i e_i/r_i)$ in this convention.
The full integer $n_i$ is retained; its residue
modulo $r_i$ alone does not determine the criterion.
All original component moduli and elliptic gluing
maps are preserved by the compact realization.
\end{proposition}



\begin{proof}
\noindent\textbf{Step 1.} We give every change to the exact maps of \ref{cdp:st:C.2.UN.49}.
Fix one leaf and suppress its index. Choose integers
$k,l$ with
\[
                         rk-nl=1,\qquad \gamma=\beta^n.
                                                        \tag{C.SN.3}\label{cdp:eq:C.SN.3}
\]
On the primary cover of the paired $P$ neighborhood
use $p=x^rz$, the infinite deck
$(z,x)\mapsto(\beta^rz,\beta^{-1}x)$, and the finite
action $(z,x)\mapsto(\zeta^rz,\zeta^{-1}x)$.
The following are genuine invariant coarse coordinates:
\[
 u=z^n,\qquad X=xz^k,\qquad
 p=X^ru^{-l},\qquad
 (u,X)\sim(\gamma^ru,\gamma^lX).
                                                        \tag{C.SN.4}\label{cdp:eq:C.SN.4}
\]
Here $u\ne0$. Since $rk-1=nl$, both coordinates
are finite-group invariant, and the finite quotient
map has no ramification. Conversely, a root of $u$
recovers $z$ and $x=Xz^{-k}$; all choices are exactly
the original finite action. Thus \eqref{cdp:eq:C.SN.4} is the
actual coarse neighborhood, not a formal coordinate
substitution.

\noindent\textbf{Step 2.} The normal of $J$ in $P$ has multiplier $\gamma^l$
over $\C^*/\gamma^{r\Z}$. Its order is exactly $r$,
since $\gcd(l,r)=1$. On the normal root cover write
$u=y^r$, $X=s y^l$; then $p=s^r$.
The generator $(y,s)\mapsto(\zeta y,\zeta^{-l}s)$
preserves these coordinates. Normalize its normal
character to $\zeta$: the fibre translation is then
$y\mapsto\zeta^n y$, because $-ln\equiv1\pmod r$.
This proves the exact torsion-character comparison,
including its dependence on $n$.
The dual normal on the secondary Hopf axis is
the inverse character, so its normal-root
identification with the given $P$ gives precisely
\eqref{cdp:eq:C.SN.4}.

\noindent\textbf{Step 3.} The punctured product map used for the inverse
transform is
\[
 (u,X)\longmapsto(p,W)
       =(X^ru^{-l},\,X^{-n}u^k),\qquad
 u=p^nW^r,\quad X=p^kW^l,\qquad
                       F=\C^*/\gamma^{\Z}.               \tag{C.SN.5}\label{cdp:eq:C.SN.5}
\]
The exponent matrix has determinant $rk-nl=1$.
Consequently these are inverse holomorphic maps
on the punctured chart. On the primary cover
they are $W=x^{-n}$ and $u=z^n=p^nW^r$.
Replacing the chart by $F\times\Delta_p$ through
\eqref{cdp:eq:C.SN.5} defines the actual principal torsor $T_0$.
It retains the original analytic transitions
exactly as in (C.NS.3).

\noindent\textbf{Step 4.} Let $a$ be the primitive positive angular loop of
$u$ with $X$ fixed, $b$ the infinite deck loop,
and $\mu_P$ the normal loop of $X$. In the
punctured product basis $(\delta,e,h)$,
\eqref{cdp:eq:C.SN.5} gives the full integral matrix
\[
 \begin{pmatrix}a&b&\mu_P\end{pmatrix}_{(\delta,e,h)}
   =
 \begin{pmatrix}
   -l&0&r\\
    k&0&-n\\
    0&1&0
 \end{pmatrix},
 \qquad
 \delta=n a+k\mu_P,\quad e=r a+l\mu_P,\quad h=b.
                                                        \tag{C.SN.6}\label{cdp:eq:C.SN.6}
\]
Its determinant is one. In particular the normal
meridian is $r\delta-ne$ and is primitive even
when $n>1$.

\noindent\textbf{Step 5.} On the Hopf side the invariant normal coordinate
is $Y=wz^{-k}$. Its normal-line multiplier is $\gamma^{-l}$,
and the diagonal reference pair has reciprocal
base coordinate $Y^ru^l=w^r/z$ and fibre coordinate
$w^n=Y^nu^k$. The complementary solid torus
therefore has meridian $n a-k\mu_H$ and fibre
cycles $(r a-l\mu_H,b)$. The actual smoothing
identifies $\mu_H=-\mu_P$. Formula \eqref{cdp:eq:C.SN.6}
then identifies its meridian with $\delta$
and its fibre cycles with $(e,h)$. This proves
the inverse-transform topology and the common
integral fibre-cycle identification for every
secondary leaf. Substituting
$\delta_i=\delta_{\rm ext}+\nu_i$ proves \eqref{cdp:eq:C.SN.2}.
The nonlinear-to-diagonal comparison (C.NS.5)
is finite-group invariant, so its boundary
calculation applies to this actual quotient.

\noindent\textbf{Step 6.} The canonical form
$dz\wedge dw/w^{r+1}$ is invariant under both
the infinite contraction and $\mu_n$. It gives
$K_H=-(r+1)J$ on the quotient. Thus the
canonical-order calculation, the effectivity
argument forcing elliptic $B$ and no unmarked
multiple fibres, and the Chern-zero necessity
are unchanged. The only change to the core
period is explicit. In $U=z/w^r$, $V=\log w$,
the two primitive periods are
\[
             (0,2\pi i/n),\qquad
             (\beta^{-r},\log\beta).
                                                        \tag{C.SN.7}\label{cdp:eq:C.SN.7}
\]
The leading two-form has core integral
$\pm2\pi i c\beta^{-r}/n$, still nonzero.
All cores are the same integral fibre class
by \eqref{cdp:eq:C.SN.6}, so their one global period
again forces all resonances $r_i$ to coincide.
The finite orders $n_i$ need not coincide.

\noindent\textbf{Step 7.} For compact sufficiency perform the exact
inverse transformations and use their
Chern-zero torsor $T_0$, which is an actual
torus by (C.NS.14). In the twistor family
(C.NS.15), choose the $i$-th fibre basis
so that its multiplicative modulus is
$\gamma_i=\beta_i^{n_i}$. After $t=s^r$,
the same four-period calculation gives
\[
 \mathcal T_s=M_{i,s}^{\times}/
   \langle(p,\ell)\mapsto
        (p+d_i s^r,\gamma_i\ell)\rangle,\qquad d_i\ne0.
                                                        \tag{C.SN.8}\label{cdp:eq:C.SN.8}
\]
At the marked point perform the weight-$(r,1)$
blowup, but now tensor the flat line by
$\OO(n_iR_i)$. Its local line coordinate
has weight $-n_i$. Setting it equal to one
leaves the finite action
\[
 (z,w,x)\mapsto
      (\zeta^{r}z,\zeta w,\zeta^{-1}x),
                    \qquad \zeta^{n_i}=1.                \tag{C.SN.9}\label{cdp:eq:C.SN.9}
\]
It is free: a nonzero $w$ or $x$ forces
$\zeta=1$, and if both vanish then $z\ne0$
and $\gcd(r,n_i)=1$ forces the same conclusion.
This proves directly that the coarse nonzero
line is smooth. Equivalently the weighted
orbifold inertia has order $r$ and acts
on the line with the faithful weight $n_i$.

\noindent\textbf{Step 8.} Choose the given $n_i$-th root $\beta_i$
of $\gamma_i$. The infinite line deck
lifts to the exact map
\[
 (z,w,x)\mapsto
   \bigl(\beta_i^r(z+d_iw^r),\beta_iw,\beta_i^{-1}x\bigr)
                                                        \tag{C.SN.10}\label{cdp:eq:C.SN.10}
\]
on \eqref{cdp:eq:C.SN.9}. It commutes with the finite
action. Its exceptional component is
exactly $H_r(\beta_i)/\mu_{n_i}$ after
$Z=z/(\beta_i^r d_i)$.
Off that component the flat-line
coordinate is $W=x^{-n_i}$; consequently
the remaining central chart and its
outer product are exactly
\eqref{cdp:eq:C.SN.4}--\eqref{cdp:eq:C.SN.5}.

\noindent\textbf{Step 9.} The polynomial norm distortion is now
$C(1+|j|)^{n_i/r}$ for the $j$-th
translation lift. The flat-line
contraction multiplier is
$|\gamma_i|\exp(-\operatorname{Re}(c_i(s)d_i s^r))$,
uniformly less than one on a small disc.
It still dominates that polynomial.
The same compact norm annulus proves
properness, and comparison with arbitrarily
large powers proves freeness of the
infinite cyclic quotient. Thus there
is again an actual isomorphism of the
complement of the new Hopf component
with the original twistor tube minus
its marked central elliptic curve.

\noindent\textbf{Step 10.} In the cluster calculation the original
flat coordinate is $x^{-n_i}$, so only
$O(1+|\log|x||)$ deck iterates are
needed. Their base displacement remains
$O(|s|^r(1+|\log|x||))$. All cluster
points therefore lie on the prescribed
central fibre curve. The simultaneous
disjoint replacements and compact gluing
of \ref{cdp:st:C.2.UN.49} apply with these exact
coarse maps. The main component is the
original $P$, and the finite diagonal
actions commute with the axis-translation
automorphisms used there. Every original
gluing map is retained. The full
central equations are still $xw=s$;
the finite action preserves this
function and is free, so the full
fibre remains reduced SNC.
This proves the equivalence and
compact realization asserted here.
\end{proof}






\begin{proposition}[The mixed-cycle criterion and compact occurrence with Hopf components]\label{cdp:st:C.2.UN.51}
\StatementRecord{cdp:st:C.2.UN.51}{proved}{Historical events: Claude Fable 5.2 (Anthropic), effort high, 16 November 2026: carried as background without fault found. GPT-6 Astra (OpenAI), ChatGPT Pro mode, effort Pro, S6\_09\_PAPERS, 17 November 2026: targeted proof check. Also checked in this writing. Pre-continuation writing events: GPT-6 Astra (OpenAI), ChatGPT Pro mode, effort Pro, S6\_10\_PAPERS, 18 November 2026: targeted proof check. Continuation metadata events: GPT-6 Astra (OpenAI), ChatGPT Pro mode, effort Pro, S6\_10\_PAPERS continuation, 19 November 2026: used and rechecked in this run; status-vocabulary reconciliation; metadata reconciliation only; no new mathematical check. \emph{Pre-continuation recorded proof scope.} The actual ruled nonzero-line complement and Hopf axis complements have the displayed rank-three lattices and affine reciprocal-normal maps. Van Kampen and the closed coframe give the exact local torus criterion. Ruled vertices preserve the outgoing normal period, so a compatible cycle cannot have just one Hopf vertex. The separate moving frame at each ruled vertex retains its actual degree-zero line, while the two independent Hopf normal parameters solve the full affine period equation. The compact suspension comparison preserves all marked surfaces; in the all-ruled case the actual horizontal translation remains in the period lattice}{the hypotheses, notation and definitions stated here; \cref{cdp:eq:C.MX.3}; \cref{cdp:st:C.2.UN.46}; \cref{cdp:eq:C.MX.4}; \cref{cdp:eq:C.MX.5}; \cref{cdp:st:C.2.UN.48}; \cref{cdp:eq:C.MX.8}; \cref{cdp:eq:C.MX.9}}

Specify an oriented cycle of surfaces, each of which is either
a diagonal primary or secondary Hopf surface of algebraic dimension
zero, marked by its two axis curves, or
\[
                  R_i=\PP(\OO_{E_i}\oplus L_i),
                  \qquad E_i\ \hbox{elliptic},\quad
                         L_i\in\operatorname{Pic}^0(E_i),
                                                        \tag{C.MX.1}\label{cdp:eq:C.MX.1}
\]
marked by its zero and infinity sections. Specify the elliptic
isomorphism joining the outgoing curve to the next incoming curve,
and a trivialization of their tensor product of normal lines.
The associated proper normal-crossing plumbing has two-torus
general fibre exactly when the cycle product $M$ of its actual
rank-three integral complement lattices is the identity.
Every actual reduced smooth-total two-torus smoothing of the
specified support satisfies this criterion.

For every such cycle containing a Hopf component, $M=I$ is also
equivalent to occurrence of the entire specified reduced support
in a smooth compact threefold over $\PP^1$, with two-torus
general fibre and Euler number zero for every reduced fibre
support. This retains all component surfaces and every elliptic
gluing map. The construction chooses the normal-product constants.
A compatible cycle containing Hopf components contains at least
two of them.

The local criterion also includes all-ruled cycles and the
identification of the two marked sections of one ruled surface.
In that case the exact punctured period lattice has the form
\[
                 \Lambda+\Z\bigl(b_0+v_0\log t\bigr),
           \qquad \Lambda\simeq\Z^3\subset\C^2,          \tag{C.MX.2}\label{cdp:eq:C.MX.2}
\]
where the projection of $b_0$ to the elliptic quotient of the
common complement is the actual total translation of the base
gluing maps. It is retained in this formula.
\end{proposition}



\begin{proof}
\noindent\textbf{Step 1.} \par\medskip\noindent\textbf{The complements and their actual line-bundle maps.}
For a ruled component, removing the two sections gives the
nonzero line bundle $L_i^\times$. Choose an origin on $E_i$
and a rigidification of $L_i$ there. Since $\deg L_i=0$, it
has constant factors of automorphy on the universal cover
of $E_i$. Taking a logarithm of its fibre coordinate gives
the exact presentation
\[
 L_i^\times=\C^2/\Lambda_i,\qquad
 \Lambda_i=\left\langle
       (\gamma_1,\log\chi_i(\gamma_1)),
       (\gamma_2,\log\chi_i(\gamma_2)),(0,2\pi i)
                         \right\rangle_{\Z} .             \tag{C.MX.3}\label{cdp:eq:C.MX.3}
\]
Here $\gamma_1,\gamma_2$ are an oriented integral period
basis of $E_i$, and $\chi_i$ is the chosen constant
factor character. The lattice is discrete of real rank
three and spans $\C^2$ over $\C$. Each marked-section
neighbourhood is the entire normal line bundle obtained
by removing the other section. Inverting the line
coordinate identifies the two descriptions on their
common complement. Their positive normal meridians are
opposite primitive elements of $\Lambda_i$.

\noindent\textbf{Step 2.} For Hopf components use the actual entire normal
line-bundle neighbourhoods and logarithmic lattices
of (C.HC.4)--(C.HC.5), including their finite diagonal
deck groups. Every edge is the same explicit operation:
apply its elliptic base isomorphism, and send a nonzero
normal vector to its dual vector whose product is $t$.
It is a global biholomorphism of the two complements.
On universal covers it is affine. Indeed its base lift
is affine, and the derivative of its logarithmic normal
multiplier is doubly periodic and holomorphic, hence
constant. Thus it has a constant complex linear part
which maps $\Lambda_i$ isomorphically to $\Lambda_{i+1}$.
The product defines the actual integral automorphism
$M$ up to conjugacy. Changes of logarithms change only
translations by lattice vectors.

\noindent\textbf{Step 3.} Every complement is smoothly $T^3\times\R$.
For \eqref{cdp:eq:C.MX.3}, a unitary flat metric on $L_i$ makes
the logarithm of its norm a linear height, with compact
three-torus level sets. The same description for a
Hopf complement was proved in \ref{cdp:st:C.2.UN.46}. Hence all component
cores and annular necks have the same rank-three
fundamental group, with collar maps exactly the
indicated integral maps. The finite normal-product
plumbing of \ref{cdp:st:C.2.UN.46} applies without a changed local map:
choose disjoint section neighbourhoods, use $xy=t$
in each dual-line neck, and glue where one normal
norm exceeds a fixed inner radius. The two ends
are disjoint for sufficiently small $|t|$.
Closed intermediate cores and necks give properness;
the opposite open ends give the same Hausdorff
comparison as in that proof.

\noindent\textbf{Step 4.} Van Kampen consequently gives
\[
                        \pi_1(Y_t)=\Lambda_1\rtimes_M\Z .
                                                        \tag{C.MX.4}\label{cdp:eq:C.MX.4}
\]
The rank-three lattice injects in this presentation.
For an actual given smoothing, the divisor sections
and the product-preserving smooth tubular map used
in \ref{cdp:st:C.2.UN.46} identify its collars with these very normal
products. Radial contraction of the tubular maps
gives the same maps on $\pi_1$. Thus \eqref{cdp:eq:C.MX.4} holds
also for that smoothing. A two-torus has abelian
fundamental group, forcing $M=I$.

\noindent\textbf{Step 5.} Conversely, when $M=I$, the complex linear cycle
product is the identity because $\Lambda_1$
spans $\C^2$. Constant changes of the two logarithmic
one-forms on successive pieces therefore glue them
to a closed holomorphic coframe on each compact
nearby fibre. The commuting dual fields have
complete complex flows and give the actual
biholomorphism with $\C^2$ divided by its period
lattice, exactly as in (C.HC.8). The core periods
and the one cycle period are all the fundamental
group generators, giving \eqref{cdp:eq:C.MX.2}. This proves
the complete local criterion for every stated
combination of components.

\par\medskip\noindent\textbf{Common normal periods and the Hopf restriction.}
Assume $M=I$ and identify all complements with
$G=\C^2/\Lambda$ using their actual linear maps.
Let $n_i\in\Lambda$ be the positive outgoing normal
meridian. The incoming normal meridian of component
$i$ is $-n_{i-1}$. On a ruled component the two
meridians are opposite, so
\[
                      n_i=n_{i-1}
                  \quad\hbox{at every ruled vertex}.   \tag{C.MX.5}\label{cdp:eq:C.MX.5}
\]
At a Hopf vertex $n_i,n_{i-1}$ are complex-linearly
independent: its two logarithmic coordinate
meridians are a complex basis. A cycle with
exactly one Hopf vertex would have equality of
these two meridians by following \eqref{cdp:eq:C.MX.5} through
all the other vertices. This proves its impossibility.

\noindent\textbf{Step 6.} Put $W=\operatorname{span}_{\R}\Lambda$ and choose
a nonzero real linear functional $\ell$ with
kernel $W$. At a Hopf vertex, its contracting
logarithm gives
\[
 \log|\alpha_i|\,\ell(i n_i)
        =\log|\beta_i|\,\ell(i n_{i-1}).
                                                        \tag{C.MX.6}\label{cdp:eq:C.MX.6}
\]
Both coefficients are negative, and neither value
can be zero because the two meridians span the
complex plane. At a ruled vertex their values
are equal by \eqref{cdp:eq:C.MX.5}. Thus, as soon as there is
a Hopf vertex, all $\ell(i n_i)$ have one nonzero
sign. The actual logarithmic cycle translation has
\[
                  2\pi i v_0=u=-\sum_i n_i,\qquad
                              \ell(iu)\ne0 .            \tag{C.MX.7}\label{cdp:eq:C.MX.7}
\]
Write $u=d p$ with $p$ primitive in $\Lambda$ and
$d>0$. The projection along $\C p$ takes $\Lambda$
to a genuine elliptic lattice, since
$\C p\cap W=\R p$. Exactly as in (C.HG.6), actual
linear coordinates give
\[
 \Lambda=\left\langle
 (0,2\pi i),(2\pi i,0),(\log q_0,\log a)
                     \right\rangle_{\Z},\qquad
              |q_0|<1,\quad u=d(0,2\pi i).             \tag{C.MX.8}\label{cdp:eq:C.MX.8}
\]
Here $a\ne1$: otherwise the explicit embedded
elliptic curve (C.HG.7) would lie in the common
complement of an algebraic-dimension-zero Hopf
surface, contradicting its two-curve calculation.

\par\medskip\noindent\textbf{Transporting both sorts of component.}
Use the compact elliptic surface and suspension
constructed in \ref{cdp:st:C.2.UN.48} with the exact periods
\[
 \Lambda(t)+\Z(0,d\log t),\qquad
 \Lambda(t)=\left\langle
 (0,2\pi i),(2\pi i,0),(\log q_0,\log a(t))
                         \right\rangle_{\Z},             \tag{C.MX.9}\label{cdp:eq:C.MX.9}
\]
where $a(0)=a$ and the actual Tate parameter is
$Q=t^d$. All the prescribed normal meridians and
contracting logarithms have fixed integral
coordinates in \eqref{cdp:eq:C.MX.8}. Retain those coordinates
in \eqref{cdp:eq:C.MX.9}.

\noindent\textbf{Step 7.} At a Hopf vertex use the moving coordinate basis
$n_i(t)/(2\pi i),-n_{i-1}(t)/(2\pi i)$.
It is invertible near zero, retains both
contracting multipliers and all finite diagonal
characters, and gives the actual Hopf family
of \ref{cdp:st:C.2.UN.48}.

\noindent\textbf{Step 8.} At a ruled vertex those two columns are
dependent, so a different actual frame is
required. Its primitive normal period $n_i$
extends to a lattice basis $n_i,\lambda_i,c_i$.
Projection along its complex line gives the
elliptic base lattice of this particular
ruled component. Choose a base coordinate
with the projected $\lambda_i$-period $2\pi i$,
and a logarithmic fibre coordinate having
normal period $2\pi i$ and zero $\lambda_i$-period.
Transport $n_i,\lambda_i,c_i$ by their fixed
integer coordinates. The first two remain
complex-linearly independent. They give
holomorphic coordinates in which the complement
lattice is
\[
        \left\langle
          (0,2\pi i),(2\pi i,0),
                      (\gamma_i(t),b_i(t))
                                  \right\rangle_{\Z} .   \tag{C.MX.10}\label{cdp:eq:C.MX.10}
\]
Its projected periods remain real-linearly
independent on a smaller disc. They define
a proper elliptic family $E_i(t)$, and the
factor $\exp b_i(t)$ defines an actual
degree-zero line bundle $L_i(t)$ on it.
The resulting
$\PP(\OO_{E_i(t)}\oplus L_i(t))$ is a smooth
proper family of the prescribed ruled
components, with their marked sections.
At zero its line-bundle presentation is
the original one in \eqref{cdp:eq:C.MX.3}, so it is
an isomorphism of marked surfaces.

\noindent\textbf{Step 9.} Every moving edge linear map sends the
outgoing normal vector to minus the next
incoming normal vector. Its induced base
map is therefore a holomorphic elliptic
isomorphism, and its nonzero normal map
is the reciprocal linear-fibre map with
the prescribed constant factors. Thus it
extends to exactly the required dual
line-bundle data. Keep the constants of
all original elliptic gluing maps.
The coefficients of the adjustable
normal-product logarithms in the total
common translation are
$-n_i(t)/(2\pi i)$. At a Hopf vertex
the two adjacent such vectors are
complex-linearly independent. Fix the
other constants and solve the resulting
two-by-two holomorphic system to make
the constant cycle translation zero.
The total period is then exactly
$(0,d\log t)$ in \eqref{cdp:eq:C.MX.9}, while every
central component and elliptic gluing
map is unchanged.

\noindent\textbf{Step 10.} The uniformly chosen flat metrics, cores
and normal-product necks now give a
smooth proper local filling. On its
punctured disc the piecewise affine
maps into the common period quotient
agree modulo precisely \eqref{cdp:eq:C.MX.9};
their changes of logarithm are integral
normal periods. Their fibrewise inverses
are the closed-coframe inverses already
proved above, so they give an actual
biholomorphism of the punctured
families with the compact suspension.
The nested-annulus replacement of
\ref{cdp:st:C.2.UN.48} produces a smooth compact total
space over $\PP^1$ with exactly the
specified mixed central support.
Each of its components and double
curves has Euler number zero. All
other reduced supports are the
elliptic-base bundles in that
suspension and also have Euler
number zero. This proves the
compact assertion.

\par\medskip\noindent\textbf{The retained all-ruled translation.}
When every vertex is ruled, \eqref{cdp:eq:C.MX.5}
makes all normal periods equal to
one primitive $n$. The common
elliptic quotient is
$\C^2/(\C n+\Lambda)$. The coefficient
vectors of all normal-product constants
lie in the single line $\C n$,
whereas the projection of the full
affine cycle constant $b_0$ is the
composition of the actual base
translations. Consequently varying
normal-product constants changes
only its vertical part; the period
formula \eqref{cdp:eq:C.MX.2} retains the projected
translation. The local torus proof
above made no assumption that it
vanishes. In particular the
single ruled self-identification
is included by its actual
normal-line inversion, rather than
the Hopf coordinate interchange.
The exact examples and moving-frame comparisons in
\path{code/A_mixed_cycles.py}
check 169 mixed cycles with two Hopf vertices and
arbitrary numbers of ruled insertions in the
displayed range. Their primitive cycle divisor is
$\gcd(a,b)$, and their real period determinant is
$\log|t|\bigl(a\log|\beta|+b\log|\alpha|\bigr)>0$.
The same calculation retains the all-ruled
horizontal translation exactly.
\end{proof}






\begin{proposition}[The weighted root criterion, exact inertia and compact occurrence]\label{cdp:st:C.2.UN.52}
\StatementRecord{cdp:st:C.2.UN.52}{proved}{Historical events: Claude Fable 5.2 (Anthropic), effort high, 16 November 2026: carried as background without fault found. GPT-6 Astra (OpenAI), ChatGPT Pro mode, effort Pro, S6\_09\_PAPERS, 17 November 2026: targeted proof check. Also checked in this writing. Pre-continuation writing events: GPT-6 Astra (OpenAI), ChatGPT Pro mode, effort Pro, S6\_10\_PAPERS, 18 November 2026: targeted proof check. Continuation metadata events: GPT-6 Astra (OpenAI), ChatGPT Pro mode, effort Pro, S6\_10\_PAPERS continuation, 19 November 2026: used and rechecked in this run; status-vocabulary reconciliation; metadata reconciliation only; no new mathematical check. \emph{Pre-continuation recorded proof scope.} The cyclic component covers and the exact normalization uv=s\textasciicircum{}h, x=u\textasciicircum{}B, y=v\textasciicircum{}A give every cover degree, bridge length and inertia subgroup. Their connected graph and injected complement lattice give the full torus criterion. The finite root action preserves orientation, fixes the primitive period, and has identity linear part because a nontrivial elliptic automorphism fixing the non-torsion line would force torsion. The exact affine compatibility equations force u/N integral and remain solvable under lattice transport. The local invariant ring is exactly C\{x,y\}, with t=x\textasciicircum{}a*y\textasciicircum{}b; the punctured translation torsor is trivial by its explicit sheaf sequence. These maps give the compact smooth quotient with the original multiplicities and both unequal-multiplicity examples}{the hypotheses, notation and definitions stated here; \cref{cdp:eq:C.WH.1}; \cref{cdp:eq:C.WH.2}; \cref{cdp:eq:C.WH.3}; \cref{cdp:eq:C.WH.4}; \cref{cdp:st:C.2.UN.40}; \cref{cdp:st:C.2.UN.46}; \cref{cdp:st:C.2.UN.51}; \cref{cdp:st:C.2.UN.43}; \cref{cdp:eq:C.WH.9}; \cref{cdp:eq:C.WH.8}; \cref{cdp:st:C.2.UN.48}; \cref{cdp:eq:C.WH.12}; \cref{cdp:eq:C.WH.11}; \cref{cdp:eq:C.WH.15}}

Let $S$ be a simple normal-crossing oriented cycle
with at least two embedded components, which are specified unblown
diagonal primary or secondary Hopf surfaces $H_i$
of algebraic dimension zero, with their two axes
joined by specified elliptic isomorphisms between distinct components.
Assign positive integers $m_i$ to its components.
The compatible weighted normal data defined in
\eqref{cdp:eq:C.WH.1}--\eqref{cdp:eq:C.WH.2} give an explicit smooth proper
local plumbing with full fibre $\sum_i m_iH_i$.

There is the following complete test for its
general fibre to be a connected two-torus.
Put $N=\operatorname{lcm}_i(m_i)$, form the finite
cyclic component covers \eqref{cdp:eq:C.WH.3}, and join their
axes by the normalized root necks \eqref{cdp:eq:C.WH.4}.
Their graph must be connected, and the product
$M$ of their actual rank-three integral complement
lattice maps must be the identity. Every actual
smooth-total torus degeneration with the given
weighted support supplies these data and satisfies
this test. Conversely every weighted support
passing this test occurs, with its specified
reduced surface and component multiplicities,
over a compact rational base, with a smooth
compact total threefold and Euler number zero
for every reduced fibre support. The compact
construction chooses normal-product constants;
it does not identify different full analytic
neighbourhoods without a comparison.

The finite-root comparison also determines the
inertia. There is a common number $r$ of connected
strict covers of every $H_i$, with
$r\mid\gcd_i m_i$. Write
\[
 g_i=\gcd(m_i,m_{i+1}),\qquad
 h_i=\frac{N}{\operatorname{lcm}(m_i,m_{i+1})}.
\]
Each strict cover of $H_i$ has degree $m_i/r$;
each connected cover of its outgoing axis has
degree $g_i/r$ and is \'{e}tale. If $K$ is the
stabilizer of a strict component and $I_i$
its generic inertia subgroup, then
\[
 \begin{gathered}
 |K|=N/r,\qquad |I_i|=N/m_i,\\
 |I_iI_{i+1}|=N/g_i,\qquad
 |I_i\cap I_{i+1}|=h_i.
 \end{gathered}
\tag{C.WH.0}\label{cdp:eq:C.WH.0}
\]
The root resolution inserts exactly $h_i-1$
ruled elliptic components along each lifted edge.

There are compact examples with multiplicities
$(1,n)$ for every $n\ge1$, described explicitly
below. There is also a compact three-component
example with multiplicities $(1,2,2)$ in which
all three components are primary Hopf surfaces.
\end{proposition}



\begin{proof}
\noindent\textbf{Step 1.} \par\medskip\noindent\textbf{The actual weighted data and smooth plumbing.}
Let $C_i^-$ and $C_i^+$ be the incoming and
outgoing axes on $H_i$, and let
$\phi_i:C_i^+\to C_{i+1}^-$ be the specified map.
The data consist of a holomorphic line
$\mathcal L_i$ and a section
\[
 \sigma_i\in H^0(H_i,\mathcal L_i^{-m_i}),
 \qquad
 \operatorname{div}(\sigma_i)=
       m_{i-1}C_i^-+m_{i+1}C_i^+,
\tag{C.WH.1}\label{cdp:eq:C.WH.1}
\]
together with the actual edge identifications
\[
 \begin{aligned}
 \mathcal L_i|_{C_i^+}
   &\simeq \phi_i^*N_{C_{i+1}^-/H_{i+1}},\\
 \phi_i^*(\mathcal L_{i+1}|_{C_{i+1}^-})
   &\simeq N_{C_i^+/H_i}.
 \end{aligned}
\tag{C.WH.2}\label{cdp:eq:C.WH.2}
\]
The leading coefficients of $\sigma_i$ and
$\sigma_{i+1}$ at this axis must agree under
these identifications. Equivalently they
give the same trivialization of
$N_{C_i^+/H_i}^{m_{i+1}}\otimes
 \phi_i^*N_{C_{i+1}^-/H_{i+1}}^{m_i}$.
This is an equality of specified holomorphic
line maps, including their nonzero constants.

\noindent\textbf{Step 2.} These are precisely the maps furnished by
an actual defining function $f$:
take $\mathcal L_i=\OO_Y(H_i)|_{H_i}$,
divide $f$ by the $m_i$th power of the
canonical divisor section, and restrict
to $H_i$. The divisor of the resulting
section is \eqref{cdp:eq:C.WH.1}. On a double curve
the two canonical sections give \eqref{cdp:eq:C.WH.2}
and the single equation $f=x^{m_i}y^{m_{i+1}}$
after absorbing a unit. \ref{cdp:st:C.2.UN.40} forces
both axes into the boundary in every
actual case; their disjointness gives
the stated cycle and no triple points.

\noindent\textbf{Step 3.} The line $\mathcal L_i$ is flat. Its
$m_i$th tensor power is the inverse of
the indicated axis divisor line, which
has the constant factors of automorphy
calculated in \ref{cdp:st:C.2.UN.40}. Transfer that flat
connection through \eqref{cdp:eq:C.WH.1}. In a local
frame $e$, divide the connection form
on $e^{m_i}$ by $m_i$; the frame-change
formula divides in the same way, giving
an actual flat holomorphic connection
on $\mathcal L_i$.
On either entire axis line-bundle chart
of the Hopf surface it is therefore the
pullback of its axis restriction:
the projection is a homotopy equivalence,
and the two flat holomorphic lines have
the same monodromy character. This gives
an actual holomorphic line isomorphism,
not a formal neighbourhood assertion.

\noindent\textbf{Step 4.} On a component core use the total line
$\mathcal L_i$ and the function
$\sigma_i(v^{m_i})$. On the edge use the
two normal lines in \eqref{cdp:eq:C.WH.2}, with function
$x^{m_{i+1}}y^{m_i}$ in the indicated
weighted trivialization. The preceding
entire line-bundle maps glue these
functions exactly. Choose bounded normal
norms and disjoint axis collars as in
\ref{cdp:st:C.2.UN.46}. If both attaching norms exceed
the inner radius $\epsilon$, their
weighted product has modulus at least
$\epsilon^{m_i+m_{i+1}}$; take the base
disc smaller than all these bounds.
Thus the two ends of each neck are
disjoint. The same closed-graph and
compact-core comparison proves a
Hausdorff smooth proper threefold.
Its central divisor has exactly the
assigned multiplicities. Away from
zero its function is a submersion.

\par\medskip\noindent\textbf{The exact component covers and normalized necks.}
For each component form the finite cyclic
cover
\[
 \mathcal E_i=
 \operatorname{Norm}
 \{\xi\in\mathcal L_i^{-1}:
                         \xi^{m_i}=\sigma_i\}.
\tag{C.WH.3}\label{cdp:eq:C.WH.3}
\]
It has degree $m_i$, possibly with more
than one connected component. It is
the actual strict cover in the
normalization of $s^N=f$:
on the component core the formula
$\xi=s^{N/m_i}/v$ gives \eqref{cdp:eq:C.WH.3}.
The root generator acts on it by
$\xi\mapsto\zeta_N^{N/m_i}\xi$.

\noindent\textbf{Step 5.} For an edge put $a=m_i$, $b=m_{i+1}$,
$g=\gcd(a,b)$, $A=a/g$, $B=b/g$ and
$h=N/\operatorname{lcm}(a,b)$.
The normalization of $x^a y^b=s^N$
has $g$ local branches. After choosing
one of their constant roots of unity,
each is exactly
\[
      uv=s^h,\qquad x=u^B,\quad y=v^A.
\tag{C.WH.4}\label{cdp:eq:C.WH.4}
\]
Indeed $(x^Ay^B)^g=s^{gABh}$ first
gives its $g$ branches. In one branch
the displayed finite map is
birational, because $\gcd(A,B)=1$
makes the choices of the roots
$u^B=x$, $v^A=y$, $uv=s^h$ unique
off $xy=0$. Its source is the normal
$A_{h-1}$ hypersurface, so it is the
normalization. Units in the original
equation are absorbed in local
holomorphic roots before this map.
The maps glue by uniqueness of
normalization.

\noindent\textbf{Step 6.} Every connected component of
$\mathcal E_i$ is a smooth diagonal
Hopf surface of algebraic dimension
zero. Here is the actual quotient
comparison. Its open part is a
connected finite \'{e}tale cover of
$G_i=\C^2/\Lambda_i$, hence is
$\C^2/\Lambda_i'$ for a finite-index
sublattice $\Lambda_i'\subset\Lambda_i$.
Let $q_+,q_-$ be the smallest positive
multiples of the two angular meridians
belonging to $\Lambda_i'$.
The local normalization
$\xi^{m_i}=z^{m_{i+1}}$ shows
$q_+=m_i/\gcd(m_i,m_{i+1})$;
the incoming formula is identical.
Divide the two logarithmic coordinates
by $q_+,q_-$. The resulting lattice
contains the two usual angular periods,
and exponentiating it gives a diagonal
deck group containing a contraction.
A finite element fixing a nonzero axis
point would give a smaller positive
angular multiple in $\Lambda_i'$,
contrary to the definition of $q_\pm$.
Thus its quotient of $\C^2\setminus0$
is a smooth primary or secondary
diagonal Hopf surface.

\noindent\textbf{Step 7.} The coordinate-power map
$(z,w)\mapsto(z^{q_+},w^{q_-})$
is equivariant for these actual deck
groups. It is finite on the compact
quotients, agrees with the open
lattice covering, and has precisely
the local branch maps above. Uniqueness
of the finite normalization therefore
identifies it with the given component
of \eqref{cdp:eq:C.WH.3}, including its two axes.
Finite pullback preserves algebraic
dimension, so it is still of algebraic
dimension zero. Each preimage axis is
one connected elliptic curve.

\noindent\textbf{Step 8.} The local branches \eqref{cdp:eq:C.WH.4} join these
covers by their actual root maps.
Every vertex has degree two. Resolving
$uv=s^h$ inserts the explicit $h-1$
projective-line bundles over the
corresponding elliptic axis cover,
with the resolution charts of
(C.HG.9)'s proof. They have two
disjoint sections and no triple
points. Thus the connected components
of the root support graph are cycles;
the compressed graph retains only
the Hopf vertices and labels each
edge by its positive integer $h$.
On the resolved family every central
component has multiplicity one.
Its Stein factor therefore has degree
one over each connected root disc.
It follows that the original general
fibre is connected exactly when this
root graph is connected.

\par\medskip\noindent\textbf{The integral criterion and the exact root inertia.}
On a compressed edge the two normal
lines are dual and the product equation
is $uv=s^h$. Its logarithmic map is
the one in (C.HC.5), with $\log s$
replaced by $h\log s$. It identifies
the actual rank-three lattices of
the two open Hopf covers. Let $M$
be their cycle product. The resolved
bridges give the same linear product:
their two normal meridians are
opposite, and their successive open
maps compose to this exact compressed
edge map. \ref{cdp:st:C.2.UN.46} and \ref{cdp:st:C.2.UN.51}, or their
same core-and-neck Van Kampen
calculation, give
\[
                  \pi_1(F_s)=\Lambda\rtimes_M\Z.
\tag{C.WH.5}\label{cdp:eq:C.WH.5}
\]
Consequently a two-torus forces $M=I$.
Conversely product identity glues
the closed logarithmic coframes.
Their complete commuting flows give
an actual complex two-torus, as in
(C.HC.8). The root fibre over a
nonzero $s$ maps isomorphically to
the original fibre over $s^N$.
This proves the stated local criterion.

\noindent\textbf{Step 9.} For an actual smoothing the same
calculation uses its actual normal
data. On cores, the canonical divisor
section and a smooth tubular
identification make its function
$v^{m_i}\sigma_i$: the ratio to this
expression is one on the zero
section and has a unique small smooth
$m_i$th root. On double curves use
both divisor sections and their
weighted product trivialization.
The resulting tubular maps give the
same covering and collar maps on
fundamental groups. Equivalently
resolve its explicit root charts
\eqref{cdp:eq:C.WH.4} and use the actual reduced
SNC comparison of \ref{cdp:st:C.2.UN.51}. No holomorphic
linearization of an arbitrary given
neighbourhood is inferred.

\noindent\textbf{Step 10.} Assume the root graph is connected.
Orient the original cycle and use
the actual factorization of each
strict root cover through its
original normalization $H_i$.
Its two axes map to the two distinct
marked axis branches on that $H_i$.
The root deck group lies over the
identity on the original space,
so it preserves these labels and
the lifted cyclic orientation.
It acts on the compressed graph
by rotations. This excludes a
reflection by the actual covering
map; an unrelated automorphism of
a Hopf cycle can reverse its
orientation.

\noindent\textbf{Step 11.} The cyclic root group acts transitively
on the strict covers of each original
component. In a connected oriented
cycle the stabilizer of any vertex
is the kernel of the rotation action
and is the same subgroup $K$ for all
vertices. If $r=N/|K|$, there are
exactly $r$ strict covers of every
$H_i$. At a generic point its normal
ramification index is $N/m_i$.
Counting the $N$ sheets gives
\[
               \deg(E_i\to H_i)=m_i/r,
                    \qquad r\mid\gcd_i m_i.
\tag{C.WH.6}\label{cdp:eq:C.WH.6}
\]
At a generic point of the original
double curve \eqref{cdp:eq:C.WH.4} has $g_i$
points above it. There are $r$
connected lifted edges, so their
axis covers have degree $g_i/r$.
They are \'{e}tale by that same
local chart, including at every
point of the compact elliptic axis.

\noindent\textbf{Step 12.} The pointwise axis stabilizer $J_i$
has order $N/g_i$. It contains the
two component inertia subgroups,
whose orders are $N/m_i$ and
$N/m_{i+1}$. In a cyclic group their
product and intersection have orders
the least common multiple and
greatest common divisor of these
numbers. These are $N/g_i$ and $h_i$,
respectively. This proves \eqref{cdp:eq:C.WH.0}
and the entire degree comparison.

\par\medskip\noindent\textbf{The finite action on the common open group is by translations.}
Assume now $M=I$. Identify all open
Hopf covers with $G=\C^2/\Lambda$
using the actual linear edge maps.
Let $n_j$ be the outgoing normal
meridian of its $j$th vertex, and
let $h_j$ be the compressed edge
length. The common period is
\[
                u=-\sum_j h_j n_j\in\Lambda.
\tag{C.WH.7}\label{cdp:eq:C.WH.7}
\]
At every vertex the contracting
generator proves (C.HG.4). The
numbers $\ell(i n_j)$ therefore
have one nonzero sign. All $h_j$
are positive, so $\ell(iu)\ne0$.
Writing $u=\delta p$ with $p$
primitive gives the actual elliptic
quotient $G\to E$ along $\C p$.
It is the nonzero bundle of a
degree-zero line $L$ on $E$.

\noindent\textbf{Step 13.} The line $L$ is not torsion. If
$L^k=\OO_E$, its pullback by the
multiplication-$k$ elliptic isogeny
is trivial. The actual finite map
$E'\times\C^*\to L^\times$ would
then send $E'\times\{1\}$ to a
compact elliptic curve in $G$.
But $G$ is the axis complement of
an algebraic-dimension-zero Hopf
surface, and its proved curve list
excludes such a curve.

\noindent\textbf{Step 14.} The deck generator is affine on
each common open group. Indeed
isomorphisms of these diagonal
Hopf covers preserving the two
labelled axes lift to diagonal
linear maps on their universal
$\C^2\setminus0$ covers, by the
nonresonant calculation of \ref{cdp:st:C.2.UN.43}.
Their logarithmic maps are affine.
Compatibility with the actual
normal edge maps makes their
common linear part a single map
$R$ preserving $\Lambda$. It has
finite order. The already proved
rotation property gives
$R n_j=n_{\sigma j}$ and
$h_{\sigma j}=h_j$, hence $Ru=u$.
In particular $Rp=p$.

\noindent\textbf{Step 15.} Thus $R$ acts as the identity on
the kernel $\C^*$ of $G\to E$
and induces a finite origin-preserving
elliptic automorphism $\beta$ with
$\beta^*L\simeq L$. If $\beta\ne1$,
the map $\beta^*-1$ on
$\operatorname{Pic}^0(E)$ is an
isogeny, so its finite kernel
would make $L$ torsion. If
$\beta=1$, $R$ is a shear with
both eigenvalues one; a finite
order shear is the identity.
We have therefore proved
\[
                              R=I.
\tag{C.WH.8}\label{cdp:eq:C.WH.8}
\]

\par\medskip\noindent\textbf{Integral descent of the period and holomorphic finite transport.}
Write the compressed edge maps in
the common coordinates as
\[
 v\longmapsto v+b_j+l_j\log s,
 \qquad l_j=-h_jn_j/(2\pi i),
\]
and the root generator as
$(j,v,s)\mapsto(\sigma j,v+\tau_j,\zeta_Ns)$.
Its exact compatibility and order give
lattice elements $\lambda_j,\mu_j$ with
\[
 \begin{aligned}
 \tau_{j+1}-\tau_j+b_j-b_{\sigma j}
                   -l_j\,2\pi i/N&=\lambda_j\in\Lambda,\\
 \sum_{a=0}^{N-1}\tau_{\sigma^a j}&=\mu_j\in\Lambda.
 \end{aligned}
\tag{C.WH.9}\label{cdp:eq:C.WH.9}
\]
These identities are obtained by composing
the actual affine maps, with
$\log\zeta_N=2\pi i/N$.
Summing the first identity around the
cycle gives
\[
                   u/N\in\Lambda,\qquad
                 \delta=Nq\quad\hbox{for an integer }q>0.
\tag{C.WH.10}\label{cdp:eq:C.WH.10}
\]
This is an integral assertion in the
actual rank-three lattice.

\noindent\textbf{Step 16.} The constants $b_j$ can first be
adjusted to have sum zero, retaining
the entire central surface and its
deck action. Adding the same multiple
of $n_j$ to $b_j$ on its whole
$\sigma$-orbit changes only a
normal-product constant and does
not change either equation in
\eqref{cdp:eq:C.WH.9}. At a Hopf vertex the
adjacent normal vectors are
complex-linearly independent.
They belong to different rotation
orbits, because \eqref{cdp:eq:C.WH.8} makes
the normal vector constant on
each orbit. These two adjustable
orbit constants span $\C^2$,
so they set $\sum_j b_j=0$.

\noindent\textbf{Step 17.} Normalize the common lattice as in
(C.HG.6). Then $u=\delta(0,2\pi i)$,
$|q_0|<1$ and the section value
$a$ is different from one by the
same explicit embedded-elliptic
argument (C.HG.7).
Construct the compact $I_q$
elliptic suspension of \ref{cdp:st:C.2.UN.48},
where $q=\delta/N$. Its periods are
\[
 \begin{gathered}
 \Lambda(t)+\Z(0,q\log t),\\
 \Lambda(t)=
 \langle(0,2\pi i),(2\pi i,0),
                       (\log q_0,\log a(t))\rangle_{\Z} .
 \end{gathered}
\tag{C.WH.11}\label{cdp:eq:C.WH.11}
\]
Put $t=s^N$. Transport every normal
period, contracting vector and
lattice element in \eqref{cdp:eq:C.WH.9} using
its same integral coefficients in
$\Lambda(t)$. This gives the actual
Hopf families of \ref{cdp:st:C.2.UN.48}, with all
finite characters retained.

\noindent\textbf{Step 18.} There are holomorphic $b_j(t),\tau_j(t)$
through their central values satisfying
\eqref{cdp:eq:C.WH.9} and $\sum b_j(t)=0$.
Here is the precise linear-algebra
reason. In common coordinates the
unknowns occur in a matrix with
constant integer coefficients.
Its right side is a rational linear
combination of the three transported
lattice vectors, because
$l_j\,2\pi i/N=-h_jn_j/N$.
Every rational left-kernel relation
on that matrix vanishes on the
right side at zero. The three
original lattice vectors are
$\Q$-linearly independent, so each
of its rational coefficients is
zero separately. The relation
therefore remains zero after
transport. A fixed linear right
inverse on the matrix image gives
a holomorphic solution, and adding
the original kernel vector makes
its value at zero the prescribed
one. The extra equation $\sum b_j=0$
is included in this matrix.

\noindent\textbf{Step 19.} Use these maps in the finite
compressed plumbing, retaining the
literal hypersurfaces $uv=s^{h_j}$
in its necks. The holomorphic
finite action exists on this entire
normal threefold: \eqref{cdp:eq:C.WH.9} is
exactly its overlap equation and
its $N$th power is a deck period.
On component cores it is the
diagonal translation isomorphism
between the transported Hopf
quotients. On necks it is the
induced line map preserving the
product up to $\zeta_N^{h_j}$.
Thus it extends across their
zero sections and respects the
actual equation. All central
components, axis maps and finite
characters are the original ones.

\par\medskip\noindent\textbf{The quotient is smooth with the original multiplicities.}
This assertion is checked on the
actual local invariant rings.
At an axis put $a=gA$, $b=gB$,
$\gcd(A,B)=1$, $h=N/(gAB)$.
The pointwise axis stabilizer
$J$ has order $hAB$. In the chart
$uv=s^h$, its subgroup of order
$h$ fixes $u,v$ and rotates $s$.
The invariant ring of this
subgroup is $\C\{u,v\}$.
Its remaining quotient acts on
$u,v$ as the independent groups
of orders $B,A$. Therefore
\[
 \bigl(\C\{u,v,s\}/(uv-s^h)\bigr)^J
           =\C\{u^B,v^A\}=\C\{x,y\},
 \qquad t=s^N=x^a y^b .
\tag{C.WH.12}\label{cdp:eq:C.WH.12}
\]
For an explicit generator choose
integers $c,d$ with $cA+dB=1$.
It has characters
$(\zeta_B^c,\zeta_A^d,\zeta_{hAB})$
on $(u,v,s)$, so their product
indeed matches the character of
$s^h$. These characters are
unchanged in the transported
finite family: they are finite
roots of unity and have the
specified central values.
The stabilizer action along the
axis is trivial; the residual
axis covering group acts freely.
Thus \eqref{cdp:eq:C.WH.12}, followed by that
\'{e}tale quotient, is a smooth
threefold chart. At a general
component point the inertia is
the normal reflection of order
$N/m_i$, so the quotient is
smooth there as well. Away from
the central support the action
moves $s$, hence is free.
This proves smoothness everywhere
near the fibre and retains
precisely the multiplicities
$m_i$. Properness follows from
the finite quotient of the
proper compressed plumbing.

\noindent\textbf{Step 20.} The central quotient is the
given reduced surface. The
transport left its cover surfaces,
axis maps and all finite action
maps unchanged at zero; these
are the actual normalization
and quotient maps constructed
from \eqref{cdp:eq:C.WH.1}--\eqref{cdp:eq:C.WH.4}.
The adjusted orbit constants
change only the normal-product
maps, not this quotient surface.

\noindent\textbf{Step 21.} The punctured quotient has the
periods \eqref{cdp:eq:C.WH.11}, since
$\log s=(\log t)/N$ and the
change of branch adds $u/N\in\Lambda$.
The logarithmic coframe descends
because of \eqref{cdp:eq:C.WH.8}.
For clarity, a possible translation
torsor does not obstruct the
actual comparison. The quotient
is a holomorphic torsor under
the marked family of tori with
these periods: local origins
differ by translations, since
the coframe is fixed. On a
punctured disc the sheaf sequence
\[
 0\longrightarrow\mathscr L
   \longrightarrow\OO^{\oplus2}
   \longrightarrow\mathscr A
   \longrightarrow0
\tag{C.WH.13}\label{cdp:eq:C.WH.13}
\]
has $H^1(\OO^{\oplus2})=0$ by
the holomorphic Cousin theorem
on the Stein punctured disc,
and $H^2(\mathscr L)=0$ because
it retracts to a circle.
Thus $H^1(\mathscr A)=0$.
Lifting the translation cocycle
through this sequence produces
an actual holomorphic section
and an isomorphism with the
period family. This identifies
the punctured quotient with
the compact suspension in
\eqref{cdp:eq:C.WH.11}, as a holomorphic
family, including all affine
translation data.

\noindent\textbf{Step 22.} The proper nested-annulus
replacement of \ref{cdp:st:C.2.UN.48} now
globalizes this smooth quotient
over the same $\PP^1$.
Every central Hopf component
and double elliptic curve
has Euler number zero.
Every other reduced support
is unchanged from the elliptic
suspension and also has Euler
number zero. This completes
the compact assertion and
the full criterion.

\par\medskip\noindent\textbf{Explicit unequal multiplicities.}
Fix nonresonant contractions
$\alpha,\beta$ and put
\[
 H=(\C^2\setminus0)/
  \langle(\alpha z,\beta w),
              (\zeta_n z,\zeta_n^{-1}w)\rangle .
\tag{C.WH.14}\label{cdp:eq:C.WH.14}
\]
Take the two-copy inversion
plumbing of this exact secondary
Hopf surface. Let $g(s)=\zeta_n s$,
let $g$ be the identity on the
first copy, and on the second
put $g(Z,W)=(Z,\zeta_n W)$.
The first neck map
$(z,w)\mapsto(z^{-1},s/w)$
intertwines these actions.
The other neck gives instead
$(Z,W)\mapsto(\zeta_n Z,W)$;
the two differ by the actual
anti-diagonal deck element.
Thus they define a genuine
order-$n$ action.

\noindent\textbf{Step 23.} At every double curve it is
$(x,y)\mapsto(x,\zeta_n y)$,
in the actual line-bundle chart.
Its invariant coordinates are
$x,Y=y^n$ and the full quotient
equation is
\[
                              t=s^n=x^nY .
\tag{C.WH.15}\label{cdp:eq:C.WH.15}
\]
The quotient is smooth and its
two multiplicities are $(1,n)$.
The first component is exactly
$H$. On the second, the finite
anti-diagonal group and this
extra action generate the two
independent scalar groups.
The coordinate-power map
$(Z,W)\mapsto(Z^n,W^n)$
therefore identifies its
quotient with the primary
Hopf surface
$H(\alpha^n,\beta^n)$.
The axis identifications are
the actual maps $z\mapsto z^{-n}$
from the corresponding axis
quotients; they are isomorphisms,
since their kernel is exactly
the given finite angular group.

\noindent\textbf{Step 24.} This example also has an
independent compact comparison.
Its common coordinates are
$(b,y)=(zw,w^n)$, with periods
$q_0=\alpha\beta$ and
$a=\beta^n$. The normal cycle
has primitive divisor $n$.
Start with the compact $I_1$
suspension with section $a(t)$,
pull back by $t=s^n$, and use
$\beta(s)=a(s^n)^{1/n}$ with
the given value at zero,
$\alpha(s)=q_0/\beta(s)$.
The two-copy replacement and
the displayed finite action
are equivariant. On the
punctured common quotient this
action is exactly the base
deck action. Its smooth
quotient consequently replaces
a tube in the original compact
suspension, giving \eqref{cdp:eq:C.WH.15}
over a compact rational base.

\noindent\textbf{Step 25.} For the primary triangle
(C.HG.1), let the base action
be $s\mapsto-s$ and the three
component actions be
\[
            (z,w),\qquad(z,-w),\qquad(-z,w).
\tag{C.WH.16}\label{cdp:eq:C.WH.16}
\]
The three actual maps
(C.HG.16) intertwine them.
The action on the middle
double curve is a free
elliptic translation; on
the other two curves it
is a normal reflection.
The smooth quotient has
multiplicities $(1,2,2)$
and primary multipliers
\[
 (e^{-1},e^{-\sqrt2}),\quad
 (e^{-\sqrt2},e^{-2-2\sqrt2}),\quad
 (e^{-2-2\sqrt2},e^{-1}).
\tag{C.WH.17}\label{cdp:eq:C.WH.17}
\]
Its axis maps are inversion,
identity and inversion.
The root primitive divisor
is two, so the same quadratic
base change of a compact
$I_1$ suspension gives its
compact replacement.

\noindent\textbf{Step 26.} The exact normalization,
inertia and invariant-ring
exponents, 300 resolution
charts, sixty anti-diagonal
secondary indices, and the
three equivariance identities
are checked in
\path{code/D_weighted_hopf_cycles.py}
and its saved output.
\end{proof}






\begin{proposition}[Complete reduced Hopf-star support criterion]\label{cdp:st:C.2.UN.53}
\StatementRecord{cdp:st:C.2.UN.53}{proved}{Historical events: Claude Fable 5.2 (Anthropic), effort high, 16 November 2026: carried as background without fault found. GPT-6 Astra (OpenAI), ChatGPT Pro mode, effort Pro, S6\_09\_PAPERS, 17 November 2026: targeted proof check. Also checked in this writing. Pre-continuation writing events: GPT-6 Astra (OpenAI), ChatGPT Pro mode, effort Pro, S6\_10\_PAPERS, 18 November 2026: targeted proof check. Continuation metadata events: GPT-6 Astra (OpenAI), ChatGPT Pro mode, effort Pro, S6\_10\_PAPERS continuation, 19 November 2026: used and rechecked in this run; status-vocabulary reconciliation; metadata reconciliation only; no new mathematical check. \emph{Pre-continuation recorded proof scope.} The canonical character equality on the actual primary cover classifies every possible marked Hopf leaf, including the scalar secondary group and the excluded axis-interchange residue sign. Adjunction and the degree-one root-component comparison force the central component to be the stated minimal elliptic quasi-bundle, with exactly its marked fibres. The full inverse-transform framing gives zero Chern vector. All leaves carry one transported core class: its nonzero nonlinear period fixes R, while its zero diagonal leading coefficient forces r<R. The explicit x\textasciicircum{}(R-r)*w\textasciicircum{}R insertion, finite characters, uniform norm bounds and cluster estimate realize every specified datum compactly}{the hypotheses, notation and definitions stated here; \cref{cdp:st:C.2.UN.38}; \cref{cdp:st:C.2.UN.43}; \cref{cdp:eq:C.AS.6}; \cref{cdp:eq:C.AS.7}; \cref{cdp:eq:C.AS.9}; \cref{cdp:st:C.2.UN.28}; \cref{cdp:eq:C.AS.13}; \cref{cdp:st:C.2.UN.31}; \cref{cdp:st:C.2.UN.36}; \cref{cdp:st:C.2.UN.49}; \cref{cdp:eq:C.AS.12}; \cref{cdp:eq:C.AS.15}; \cref{cdp:eq:C.AS.16}; \cref{cdp:eq:C.AS.18}; \cref{cdp:eq:C.AS.17}; \cref{cdp:st:C.2.UN.50}}

Let $S$ be a specified compact connected SNC surface with smooth
components $P,H_1,\ldots,H_n$, $n\ge1$. Each $H_i$ is a primary
or secondary Hopf surface. Its intersection with $P$ is one
connected smooth curve $J_i$, the $J_i$ are disjoint on $P$,
and there are no other intersections. We impose no initial
classification or elliptic-fibration hypothesis on $P$.

The following two occurrence assertions are equivalent: $S$
is the full reduced central fibre of a proper two-torus
degeneration with smooth total space; and $S$ is such a fibre
of a smooth compact threefold over $\PP^1$, with every other
fibre a smooth two-torus. They hold precisely for the following
explicit data.

Every leaf pair is isomorphic to one of
\[
 \begin{gathered}
 (H^{\rm diag}_{r_i}(\beta_i)/\mu_{n_i},J_i),\qquad
 (H^{\rm nl}_{r_i}(\beta_i)/\mu_{n_i},J_i),\\
 r_i\ge1,\qquad 0<|\beta_i|<1,\qquad \gcd(r_i,n_i)=1 ,
 \end{gathered}                                             \tag{C.AS.1}\label{cdp:eq:C.AS.1}
\]
where the two contractions on $\C^2\setminus\{0\}$ are,
respectively,
\[
 (z,w)\longmapsto(\beta_i^{r_i}z,\beta_iw),\qquad
 (z,w)\longmapsto(\beta_i^{r_i}z+w^{r_i},\beta_iw),
 \quad
 \zeta(z,w)=(\zeta^{r_i}z,\zeta w),                         \tag{C.AS.2}\label{cdp:eq:C.AS.2}
\]
and $J_i$ is the image of $w=0$. The normal lines on every
double curve are dual. The central component has a
fixed-Jacobian, identity-integral-monodromy elliptic
quasi-bundle
\[
                         q:P\longrightarrow B             \tag{C.AS.3}\label{cdp:eq:C.AS.3}
\]
over an elliptic curve. Its marked full fibres are $r_iJ_i$;
there are no other multiple fibres and every fibre support
is smooth. The actual inverse logarithmic transform $T_0$
defined by (C.SN.3)--(C.SN.6), retaining all the specified
axis identifications and the entire analytic torsor,
satisfies
\[
                         c(T_0)=0\in H^2(B,\Lambda_F).     \tag{C.AS.4}\label{cdp:eq:C.AS.4}
\]
Finally there is a positive integer $R$ such that
\[
 \begin{cases}
 r_i=R,&H_i\text{ is of the nonlinear form in \eqref{cdp:eq:C.AS.2}},\\
 r_i<R,&H_i\text{ is of the diagonal form in \eqref{cdp:eq:C.AS.2}}.
 \end{cases}                                             \tag{C.AS.5}\label{cdp:eq:C.AS.5}
\]
When every leaf is diagonal, the last test is satisfied by
any $R>\max_i r_i$. When a nonlinear leaf occurs, it
determines $R$, and every other nonlinear leaf has that
same resonance.

The conclusion retains the full isomorphism class of $S$.
In particular it retains arbitrary flat torsor parameters,
finite secondary orders, independent primitive angular
directions, and the given elliptic gluing maps. Every support
in the compact realization has Euler characteristic zero.
\end{proposition}



\begin{proof}
\noindent\textbf{Step 1.} \par\medskip\noindent\textbf{The canonical identity on an arbitrary Hopf leaf.}
We first use an actual reduced smoothing. Let the divisor
of a local canonical evaluation generator be
$k_PP+\sum_i k_iH_i$. Its existence and the fact that it
has only vertical zeroes are the exact direct-image
calculation in \ref{cdp:st:C.2.UN.38}. Adjunction on one leaf gives
\[
                  K_H=\OO_H(-aJ),\qquad
                  a=k_H-k_P+1\in\Z.                      \tag{C.AS.6}\label{cdp:eq:C.AS.6}
\]
We classify this equality on the actual pair $(H,J)$,
including the full secondary deck group.

\noindent\textbf{Step 2.} Take a normal primary cyclic subgroup of the Hopf deck
group $\Gamma$. Its generator $F$ is central: conjugation
can send it only to $F$ or $F^{-1}$, while the nonzero
logarithm of the modulus of its derivative determinant
excludes the latter. This is the exact normal-core
argument of \ref{cdp:st:C.2.UN.43}. We therefore have a finite
unramified primary cover $H'\to H$, and the full inverse
image $D$ of $J$ is a disjoint union of smooth elliptic
curves. The finite quotient acts transitively on these
curves, since $J$ is connected. Pullback gives
\[
                           K_{H'}=-aD.                  \tag{C.AS.7}\label{cdp:eq:C.AS.7}
\]

\noindent\textbf{Step 3.} For a nonlinear primary contraction the curve computation
of \ref{cdp:st:C.2.UN.38} gives its unique axis, and the full
centralizer computation of \ref{cdp:st:C.2.UN.43} gives exactly the
nonlinear quotient in \eqref{cdp:eq:C.AS.1}--\eqref{cdp:eq:C.AS.2}.
The invariant form $dz\wedge dw/w^{r+1}$ has divisor
$-(r+1)J$, on the quotient as well. Since $\OO_H(J)$ has
infinite order, \eqref{cdp:eq:C.AS.6} forces $a=r+1$.

\noindent\textbf{Step 4.} Consider a diagonal primary contraction. Its invariant
curve computation can be read directly from an equation
on $\C^2$. Closure across the origin and removal of the
nonconstant unit in its transformation give an eigenfunction
of $F$. Comparison of its power-series monomials gives
the axes in the absence of a multiplicative resonance.
In that case the complete centralizer is diagonal, so
it cannot interchange the axes. Thus $D$ is one axis.
If, for example, it is $w=0$, the exact automorphy
characters in \eqref{cdp:eq:C.AS.7} give
$\alpha\beta=\beta^a$, which is the resonance
$\alpha=\beta^{a-1}$. The contraction moduli make
$a-1>0$. This contradicts the assumed absence of a
resonance and accounts for that case through an actual
character equality.

\noindent\textbf{Step 5.} In the resonant case pass to a further central cyclic
subgroup to remove the finite root-of-unity discrepancy.
Its generator has the form
\[
                 F(z,w)=(\gamma^p z,\gamma^q w),
       \qquad p,q\ge1,\quad\gcd(p,q)=1,\quad |\gamma|<1 .
                                                               \tag{C.AS.8}\label{cdp:eq:C.AS.8}
\]
Here the axes have defining weights $p,q$, and every
other irreducible curve has equation $z^q=cw^p$ of
weight $pq$. This follows from the same eigenfunction
calculation: after an axis monomial is removed, a
weighted homogeneous polynomial factors into these
binomials over $\C$. Its irreducible factor is unique
because $p,q$ are coprime. The curves are smooth away
from the origin, as required for the elliptic curves
on $H'$.

\noindent\textbf{Step 6.} The canonical weight is $p+q$. If $D$ is an orbit of
$m$ non-axis curves, \eqref{cdp:eq:C.AS.7} therefore gives
\[
                              ampq=p+q.                 \tag{C.AS.9}\label{cdp:eq:C.AS.9}
\]
The same calculation with a single axis replaces $mpq$
by $p$ or $q$. These are genuine character equalities:
a purported extra holomorphic unit and its inverse
extend across the origin, and their nonzero value there
forces the identical character. In particular $a>0$.

\noindent\textbf{Step 7.} For $p,q>1$, coprimality gives $pq>p+q$, so \eqref{cdp:eq:C.AS.9}
is impossible. A single axis would require $p\mid q$
or $q\mid p$, also impossible. In this case the
centralizer preserves each axis, so an orbit containing
both axes cannot occur. For $p=r>1,q=1$, the curves
$z=cw^r$, including $z=0$, all have weight $r$.
The triangular centralizer can permute these curves;
an orbit of $m$ of them has total weight $mr$ and
would require $amr=r+1$, hence $r\mid1$.
The only possibility is the weight-one axis $w=0$,
and then $a=r+1$.

\noindent\textbf{Step 8.} We retain the full deck group in this last calculation.
Commuting with $(\gamma^rz,\gamma w)$ gives exactly
\[
              (z,w)\longmapsto(Az+Cw^r,Bw).             \tag{C.AS.10}\label{cdp:eq:C.AS.10}
\]
Its induced action on $z/w^r$ is affine. The central
cyclic subgroup acts trivially there; hence the image
is a finite affine group. Averaging the orbit of zero
gives an actual common fixed point. The corresponding
coordinate change $z\mapsto z-cw^r$ makes every element
diagonal, while retaining the marked axis. Equation
\eqref{cdp:eq:C.AS.6} then says, for every deck element,
$AB=B^{r+1}$, so $A=B^r$. Thus $\Gamma$ is the
weighted scalar group with an infinite contracting
generator $\beta$ and finite roots of unity
$\mu_n$. Freeness on the marked axis gives
$\gcd(r,n)=1$. This is exactly the diagonal pair in
\eqref{cdp:eq:C.AS.1}--\eqref{cdp:eq:C.AS.2}.

\noindent\textbf{Step 9.} It remains to check the scalar case $p=q=1$,
where the full secondary group need not initially
be diagonal. Every component of $D$ lifts to a
line, and \eqref{cdp:eq:C.AS.7} gives $am=2$. If $m=1$,
the finite projective group fixes that line.
On its affine complement, averaging gives a
second fixed point, so the linear group becomes
diagonal in coordinates retaining the first line.
The character equality $AB=B^2$ forces $A=B$
for every element. We obtain the case $r=1$
of \eqref{cdp:eq:C.AS.2}, with scalar finite action.
If $m=2$, then $a=1$ and an element of the full
deck group interchanges the two lines. Choose
them as the two coordinate axes. The unique
log-canonical form up to a constant is
\[
                             \frac{dz}{z}\wedge\frac{dw}{w}.
                                                               \tag{C.AS.11}\label{cdp:eq:C.AS.11}
\]
Every such interchange multiplies this form
by $-1$. A form with exactly the same divisor
differs by a nowhere-zero holomorphic function
on the compact primary Hopf surface, hence by
a constant. It therefore cannot descend.
This contradicts the actual equality
$K_H+J=0$ in \eqref{cdp:eq:C.AS.6}. Thus this last orbit
does not satisfy \eqref{cdp:eq:C.AS.6}. We have proved,
without an a priori restriction on the secondary
group, precisely the leaf-pair list \eqref{cdp:eq:C.AS.1},
and in every surviving case
\[
                              k_i-k_P=r_i\ge1.          \tag{C.AS.12}\label{cdp:eq:C.AS.12}
\]

\par\medskip\noindent\textbf{The central component is forced into the stated class.}
Adjunction on $P$ now gives the exact effective
canonical divisor
\[
                         K_P=\sum_i(r_i-1)J_i .         \tag{C.AS.13}\label{cdp:eq:C.AS.13}
\]
Thus $\kappa(P)\ge0$ and $p_g(P)\ge1$.
\ref{cdp:st:C.2.UN.28} gives $p_g(P)=1$. The surface $P$
is minimal. Indeed, a $(-1)$-curve is rational,
so is distinct from every elliptic $J_i$.
Its intersections with all of them are
nonnegative, contradicting $K_P\cdot E=-1$
and \eqref{cdp:eq:C.AS.13}.

\noindent\textbf{Step 10.} We now specify the covering comparison needed
to apply \ref{cdp:st:C.2.UN.31} to this original component.
Its full-fibre multiplicity is one. In the
normalized root base change $t=s^N$, a generic
chart is $x=s^N$ after a holomorphic unit is
absorbed. There is exactly one strict covering
component and its map to $P$ has degree one.
Equivalently, the total generic degree of the
strict component covers is the original
multiplicity, as calculated in \ref{cdp:st:C.2.UN.36}.
After resolution it remains birational to $P$.
Both minimal models are therefore the same
smooth surface. The actual root-cover conclusion
of \ref{cdp:st:C.2.UN.31} thus applies to $P$ itself:
it is a two-torus, or a fixed-Jacobian,
identity-monodromy elliptic quasi-bundle over
a curve of genus at most one, with only smooth
multiple fibres. In the latter description
$p_g(P)=g(B)$, so the base is elliptic.
No assertion about an arbitrary quotient
component has been substituted for this
degree-one comparison.

\noindent\textbf{Step 11.} We also verify that all the given curves
$J_i$ are fibres of one such fibration.
On a torus, an embedded elliptic curve is a
translate of an elliptic subgroup: lift its
map to universal covers and differentiate;
the derivative is periodic and holomorphic,
hence constant. Two elliptic subgroups with
distinct tangent lines give a surjective
isogeny from their product to the two-torus.
Every two translates then meet. The disjoint
$J_i$ consequently have one common tangent
line and are fibres of the actual elliptic
quotient.

\noindent\textbf{Step 12.} In the quasi-bundle description over an
elliptic base the canonical divisor is
$\sum(m-1)F_{\rm red}$, by the exact root
chart calculation in \ref{cdp:st:C.2.UN.49}.
If a multiple fibre occurs, a horizontal
$J_i$ would have positive intersection with
this canonical divisor. But its torsion
normal line gives $J_i^2=0$, and adjunction
gives $K_P\cdot J_i=0$. Thus every $J_i$
is vertical. If no multiple fibre occurs,
we have a principal elliptic bundle.
A nonzero integral Chern vector excludes
a horizontal elliptic curve: its finite
unramified map to the base would pull the
bundle back to one with a section, while
its Chern vector would be multiplied by
the nonzero covering degree. If that Chern
vector is zero, the exact Cech lattice
calculation (C.NS.14) identifies $P$ with
a torus, reducing to the preceding argument.
This proves \eqref{cdp:eq:C.AS.3} with every marked
curve retained.

\noindent\textbf{Step 13.} The normal of the leaf axis has exact
torsion order $r_i$, including its finite
secondary character by (C.SN.3)--(C.SN.4).
Normal duality in the reduced smoothing
therefore gives full marked fibre $r_iJ_i$
on $P$. Subtracting the marked terms of
its canonical root-chart formula from
\eqref{cdp:eq:C.AS.13} leaves an effective sum of
unmarked multiple fibres linearly
equivalent to zero. Such a nonzero
effective divisor on a compact connected
complex surface cannot be principal:
its defining meromorphic function would
have no poles and hence be constant.
Consequently there are no unmarked
multiple fibres.

\par\medskip\noindent\textbf{The exact topology and the one period shared by all leaves.}
The primary and secondary inverse-transform
charts (C.SN.3)--(C.SN.6) apply to both
contractions in \eqref{cdp:eq:C.AS.2}: they use their
identical axis normal character. For a
nonlinear leaf the explicit deformation
of its resonant coefficient to zero gives
the relative smooth-pair comparison
proved in \ref{cdp:st:C.2.UN.49}. For a diagonal leaf
the product complement is already exact.
Thus every Hopf complement kills the
same displayed complementary meridian,
and each oriented core two-torus becomes
the ordinary fibre class of the actual
inverse transform $T_0\to B$.
The nearby smoothing is its underlying
four-manifold. Its fundamental group is
the principal-bundle presentation
(C.NS.13), with the unaltered integer
Chern vector. A two-torus has abelian
fundamental group; the injected fibre
lattice in that presentation therefore
forces \eqref{cdp:eq:C.AS.4}.

\noindent\textbf{Step 14.} Let $\gamma$ be this one transported
integral core class. Divide the relative
canonical generator by $t^{k_P}$.
On the $i$th Hopf core its leading
term is $t^{r_i}\eta_i$, with
\[
                   \eta_i=c_i\,\frac{dz\wedge dw}{w^{r_i+1}},
                          \qquad c_i\ne0.               \tag{C.AS.14}\label{cdp:eq:C.AS.14}
\]
This is the actual coefficient calculation
of (C.NS.11), together with \eqref{cdp:eq:C.AS.12}.
Put $U=z/w^{r_i}$ and $V=\log w$.
The primitive angular period of the
secondary quotient is $(0,2\pi i/n_i)$.
For a nonlinear leaf the other period is
$(\beta_i^{-r_i},\log\beta_i)$, so
\[
       \int_{\gamma}\omega_t
          =t^{r_i}\left(\frac{\pm2\pi i c_i\beta_i^{-r_i}}{n_i}
                         +O(t)\right).                 \tag{C.AS.15}\label{cdp:eq:C.AS.15}
\]
Its coefficient is nonzero. For a diagonal
leaf the other period is
$(0,\log\beta_i)$: the core lies at
constant $U$, and $c_i\,dU\wedge dV$
has zero integral there. Consequently
\[
                         \int_\gamma\omega_t=O(t^{r_i+1})
                  \quad\hbox{at a diagonal leaf}.        \tag{C.AS.16}\label{cdp:eq:C.AS.16}
\]
These are periods of the same actual
class in the same nearby fibre. They
are holomorphic functions: integration
of a closed relative holomorphic form
over the transported compact cycle is
holomorphic by differentiation and
Stokes' theorem. The displayed leading
terms also give their holomorphic
extension at zero.

\noindent\textbf{Step 15.} If a nonlinear leaf occurs, \eqref{cdp:eq:C.AS.15}
gives one finite vanishing order $R$.
Every nonlinear $r_i$ equals it, and
\eqref{cdp:eq:C.AS.16} gives $r_i<R$ for every
diagonal leaf. If all leaves are
diagonal, no finite vanishing order
is asserted for the given smoothing;
the explicit numerical test \eqref{cdp:eq:C.AS.5}
is satisfied by choosing
$R=1+\max_i r_i$. This proves its
necessity without imposing an unproved
nonzero period on that family.

\par\medskip\noindent\textbf{One compact construction realizes all the specified data.}
We now use data \eqref{cdp:eq:C.AS.1}--\eqref{cdp:eq:C.AS.5}.
The zero-Chern inverse transform is
the actual torus $T_0$ with its marked
fibres, by (C.NS.14). Take its explicit
compact $\OO_{\PP^1}(1)^{\oplus2}$
twistor family (C.NS.15), and pull it
back by $t=s^R$. For the $i$th marked
curve the exact four-period coordinate
change (C.NS.16)--(C.NS.18), with
$\gamma_i=\beta_i^{n_i}$, gives
\[
 \mathcal T_s=M_{i,s}^{\times}/
   \langle(p,\ell)\mapsto(p+d_i s^R,\gamma_i\ell)\rangle,
                    \qquad d_i\ne0 .                  \tag{C.AS.17}\label{cdp:eq:C.AS.17}
\]
The non-torsion flat line and every
original primitive fibre direction
are retained by these actual period
maps.

\noindent\textbf{Step 16.} At this marked point use weight
$(r_i,1)$ and tensor the flat line
by $\OO(n_iE_i)$, where $E_i$ is
the weighted exceptional divisor.
The same free coarse quotient as
(C.SN.9) has coordinates
$p-p_i=x^{r_i}z$, $s=xw$.
The exact infinite deck map is now
\[
 \begin{split}
 (z,w,x)\longmapsto\bigl(&
   \beta_i^{r_i}(z+d_i x^{R-r_i}w^R),\\
   &\beta_iw,\ \beta_i^{-1}x\bigr).
 \end{split}                                             \tag{C.AS.18}\label{cdp:eq:C.AS.18}
\]
The exponent $R-r_i$ is nonnegative.
The added monomial has weighted degree
$r_i$, so this map commutes with the
finite action of weights
$(r_i,1,-1)$ modulo $n_i$. It preserves
$s=xw$ and changes the original base
coordinate by exactly $d_i s^R$.
Thus \eqref{cdp:eq:C.AS.18} is the actual lift
of \eqref{cdp:eq:C.AS.17}.

\noindent\textbf{Step 17.} On $x=0$, equality $R=r_i$ gives the
specified nonlinear contraction after
$Z=z/(\beta_i^{r_i}d_i)$. Strict
inequality $R>r_i$ gives the specified
diagonal contraction. On $w=0$ the
map is always
$(z,x)\mapsto(\beta_i^{r_i}z,\beta_i^{-1}x)$,
so the main component and its outer
trivialization are exactly
(C.SN.4)--(C.SN.5), independently of
$R$. Every given central component
and every marked normal cover is
therefore the one prescribed.

\noindent\textbf{Step 18.} We record the analytic bounds to
verify the simultaneous compact
replacement in this varying-weight
case. On a small closed disc,
$|s|^R\le C|s|^{r_i}$ for every $i$.
The weighted norm comparison of
(C.NH.9)--(C.NH.10) consequently
gives at most
$C(1+|j|)^{n_i/r_i}$ distortion
for the $j$th translation lift.
The unitary flat metric used in
\ref{cdp:st:C.2.UN.50} has scalar multiplier
$|\gamma_i|\exp(-\operatorname{Re}(c_i(s)d_i s^R))$,
uniformly less than one after
shrinking the disc. Its exponential
contraction dominates that polynomial.
The same compact norm annulus proves
properness and freeness of the
infinite quotient. Coprimality in
\eqref{cdp:eq:C.AS.1} makes its finite coarse
quotient smooth.

\noindent\textbf{Step 19.} Outside the new Hopf component the
canonical exceptional section gives
the actual isomorphism with the
original twistor family minus its
marked central elliptic curve.
In a compact Hopf fundamental domain,
$z,w$ remain bounded and $x\to0$.
The old flat coordinate is proportional
to $x^{-n_i}$. Putting it in a compact
norm annulus uses
$O(1+|\log|x||)$ iterates, with total
base displacement
\[
                   O\bigl(|s|^R(1+|\log|x||)\bigr)
                          \longrightarrow0,
                         \qquad s=xw.                   \tag{C.AS.19}\label{cdp:eq:C.AS.19}
\]
All cluster points consequently lie
on the specified original central
elliptic curve. The disjoint-curve
gluing proof of (C.NS.20) therefore
applies to all these insertions
simultaneously. It proves a smooth
Hausdorff proper local family, and
the unchanged compact twistor
complement makes the total space
compact over $\PP^1$.

\noindent\textbf{Step 20.} The central main component is
precisely $P$, because all its
analytic torsor transitions have
been reversed by the exact inverse
charts. The diagonal automorphisms
$(z,w)\mapsto(a^{r_i}z,aw)$ commute
with either contraction and with
its finite action; they retain
every specified constant in the
axis gluing. The linear parts were
already retained by the chosen
normal-cover lattice bases.
Thus the full central support is
the specified $S$. Its equations
are $xw=s$ in the free coarse
charts, so every component has
multiplicity one. All other fibres
are the original smooth two-tori.
The Euler number of the central
support is zero by additivity:
each quasi-bundle, Hopf leaf and
double elliptic curve has Euler
number zero.

\noindent\textbf{Step 21.} The exact symbolic checks in
\path{code/root_mixed_hopf_stars.py}
retain the weighted base map, its
Jacobian, the finite-action weights
and the common period exponent.
The finite integer search there
checks the degree exclusions in
\eqref{cdp:eq:C.AS.9}; the preceding integer
argument proves them for all weights.
\end{proof}






\begin{proposition}[Every reduced all-Hopf support has one cycle, with an exact algebraic-dimension-zero classification]\label{cdp:st:C.2.UN.54}
\StatementRecord{cdp:st:C.2.UN.54}{proved}{Historical events: Claude Fable 5.2 (Anthropic), effort high, 16 November 2026: carried as background without fault found. GPT-6 Astra (OpenAI), ChatGPT Pro mode, effort Pro, S6\_09\_PAPERS, 17 November 2026: targeted proof check. Also checked in this writing. Pre-continuation writing events: GPT-6 Astra (OpenAI), ChatGPT Pro mode, effort Pro, S6\_10\_PAPERS, 18 November 2026: targeted proof check. Continuation metadata events: GPT-6 Astra (OpenAI), ChatGPT Pro mode, effort Pro, S6\_10\_PAPERS continuation, 19 November 2026: used and rechecked in this run; status-vocabulary reconciliation; metadata reconciliation only; no new mathematical check. \emph{Pre-continuation recorded proof scope.} Adjunction and the universal-cover map make every Hopf curve smooth elliptic and all such curves disjoint. The constant deck cocycle restricts nontrivially to each elliptic curve, including finite secondary covers. The normalization exact sequence, reduced-fibre semicontinuity and the actual graph-group quotient force one cycle and rank |V|-1 of the weighted restriction matrix. Duality identifies the unique dualizing form with the residue kernel, supported exactly on cycle edges and components. The explicit local canonical restriction therefore identifies the canonical-order minima with the cycle, and the actual a-zero curve list removes all trees in that branch}{the hypotheses, notation and definitions stated here; \cref{cdp:st:C.2.UN.48}; \cref{cdp:st:C.2.UN.37}; \cref{cdp:st:C.2.UN.38}; \cref{cdp:st:C.2.UN.53}; \cref{cdp:st:C.2.UN.43}; \cref{cdp:st:C.2.UN.44}; \cref{cdp:eq:C.HR.6}; \cref{cdp:eq:C.HR.7}; \cref{cdp:st:C.2.UN.46}}

Let $f:U\to\Delta$ be a proper degeneration of connected complex
two-tori with smooth total space, whose full central fibre
$S=\bigcup_{v\in V}H_v$ is reduced and simple normal crossing.
Suppose that every $H_v$ is a primary or secondary Hopf surface.
Then every connected double curve is elliptic, there are no triple
points, and the connected incidence graph has exactly one independent
cycle.

For an oriented edge $e$ from $v$ to $w$, let $C_e$ be its actual
double curve and let
\[
 r_{v,e}:H^1(H_v,\OO_{H_v})\longrightarrow H^1(C_e,\OO_{C_e})
                                                        \tag{C.HR.1}\label{cdp:eq:C.HR.1}
\]
be restriction through the specified inclusion (and through the
specified gluing map at the other endpoint). Every such map is a
nonzero map between one-dimensional complex vector spaces. The
weighted incidence map
\[
 d_1:\bigoplus_vH^1(H_v,\OO_{H_v})
       \longrightarrow\bigoplus_eH^1(C_e,\OO_{C_e}),\qquad
 (d_1a)_e=r_{v,e}a_v-r_{w,e}a_w
                                                        \tag{C.HR.2}\label{cdp:eq:C.HR.2}
\]
has rank $|V|-1$. In any choices of nonzero bases, the product of
$r_{v,e}/r_{w,e}$ around the oriented cycle is one. This product
retains the actual restriction maps; it is unchanged by changing
the chosen bases. The full coherent dimensions are
\[
                    h^q(S,\OO_S)=(1,2,1),\qquad q=0,1,2.
                                                        \tag{C.HR.3}\label{cdp:eq:C.HR.3}
\]

The one-dimensional space $H^0(S,\omega_S)$ has a more precise
description. Its nonzero section is represented by logarithmic
two-forms on the components. Their residues are nonzero on exactly
the cycle edges and zero on every edge of an attached tree. The
forms themselves are nonzero on exactly the cycle components.
For a canonical evaluation generator with vertical orders $k_v$,
the minimum of the $k_v$ is therefore attained exactly at the
cycle components.

In particular, if every $H_v$ has algebraic dimension zero, the
graph has no attached trees and every component is a diagonal Hopf
surface. There are at least two components. This gives the following
complete criterion for this entire reduced branch, without initially
assuming either diagonality or a cycle. A specified compact reduced
SNC surface all of whose components are algebraic-dimension-zero
Hopf surfaces occurs as the full fibre of a smooth-total torus
degeneration precisely when its components are diagonal, its graph
is a cycle, its actual axis normal lines are dual, and the integral
edge product $M$ of (C.HC.3)--(C.HC.6) is the identity. Every such
specified support has the compact realization of \ref{cdp:st:C.2.UN.48}, with
all support Euler characteristics zero.
\end{proposition}



\begin{proof}
\noindent\textbf{Step 1.} \par\medskip\noindent\textbf{The curves and the actual restriction maps.}
A Hopf surface has $H^2(H,\Q)=0$. For an irreducible compact curve
$C\subset H$, adjunction gives $p_a(C)=1$. A map from $\PP^1$ to
$H$ lifts to its universal cover $\C^2\setminus\{0\}$, and its two
coordinate functions are constant. The normalization of $C$ hence
cannot be rational. The equality
$p_a(C)=g(C^\nu)+\delta(C)$ now gives $g(C^\nu)=1$ and
$\delta(C)=0$. Thus $C$ is smooth elliptic. Two distinct such
curves have intersection number zero, since $H^2(H,\Q)=0$;
positivity of every local intersection forces them to be disjoint.
These are the same explicit curve comparisons used in \ref{cdp:st:C.2.UN.37}.
It follows at once that the present support has only elliptic
double curves and no triple point.

\noindent\textbf{Step 2.} We prove nonvanishing in \eqref{cdp:eq:C.HR.1}, including secondary surfaces.
First let $H'=(\C^2\setminus\{0\})/\langle F\rangle$ be a primary
Hopf surface. The constant additive deck cocycle $1$ represents a
nonzero class in $H^1(H',\OO_{H'})$. Indeed, a trivialization would
be a holomorphic function $b$ on $\C^2\setminus\{0\}$ satisfying
$b(Fz)-b(z)=1$. Hartogs extension to the origin and evaluation there
give a contradiction. Since $h^1(H',\OO_{H'})=1$, this class spans
the group.

\noindent\textbf{Step 3.} For every elliptic curve $C'\subset H'$, its lift can have finitely
many branches permuted by $F$. Choose a positive power $F^d$
preserving one branch and removing any root-of-unity discrepancy
in the diagonal multipliers. The actual curve list used in
\ref{cdp:st:C.2.UN.38} and \ref{cdp:st:C.2.UN.53} gives its normalized parameter $\C^*$,
with $F^d$ acting by $u\mapsto\lambda u$, where
$0<|\lambda|<1$. For an axis this is its displayed diagonal
multiplier; for a resonant binomial curve $z^q=cw^p$ it follows
from $(z,w)=(a u^p,b u^q)$. The scalar and nonlinear-axis cases
are included. The actual map
$\C^*/\langle F^d\rangle\to C'$ induced by this inclusion is a
finite unramified elliptic covering. Pulling the constant deck
cocycle $1$ to the subgroup generated by $F^d$ gives the cocycle
$d$. Thus, if the original restriction were zero, its pullback
would give a holomorphic function on $\C^*$ satisfying
\[
                         b(\lambda u)-b(u)=d.          \tag{C.HR.4}\label{cdp:eq:C.HR.4}
\]
In its Laurent expansion on $\C^*$ the constant coefficient on
the left is zero, whereas $d>0$. This proves nonvanishing without
assuming that $F$ fixes every lifted branch individually.

\noindent\textbf{Step 4.} For a secondary $H$, take the actual finite unramified primary
cover $\pi:H'\to H$ of \ref{cdp:st:C.2.UN.43}. A component $C'$ of the inverse
image of $C$ maps by a finite unramified elliptic covering to $C$.
The trace identities
$\operatorname{Tr}\circ\pi^*=(\deg\pi)\operatorname{id}$, on $H$
and on this curve covering, show that both pullbacks on $H^1(\OO)$
are injective. The surface pullback is an isomorphism of
one-dimensional spaces. Naturality of the restriction square,
together with the primary calculation, proves that the original
$H^1(H,\OO_H)\to H^1(C,\OO_C)$ is nonzero. Thus every coefficient
in \eqref{cdp:eq:C.HR.2} is nonzero, with the actual gluing maps included.

\par\medskip\noindent\textbf{Normalization, the graph, and its exact cycle product.}
The no-triple SNC equation gives the exact normalization sequence
\[
 0\longrightarrow\OO_S
 \longrightarrow\bigoplus_v\OO_{H_v}
 \xrightarrow{\ d\ }\bigoplus_e\OO_{C_e}
 \longrightarrow0,                                    \tag{C.HR.5}\label{cdp:eq:C.HR.5}
\]
where $d$ is the difference of the actual two restrictions. On
$H^0$ it is the ordinary signed incidence matrix $d_0$. The
connectedness of $S$ gives $\operatorname{rank}d_0=|V|-1$.
Since $H^2(H_v,\OO_{H_v})=0$, its long exact sequence gives
\[
 \begin{gathered}
 0\longrightarrow\operatorname{coker}d_0
 \longrightarrow H^1(S,\OO_S)
 \longrightarrow\ker d_1\longrightarrow0,\\
             H^2(S,\OO_S)\simeq\operatorname{coker}d_1.
 \end{gathered}                                        \tag{C.HR.6}\label{cdp:eq:C.HR.6}
\]
Here $\dim\operatorname{coker}d_0$ is the first Betti number of
the incidence graph.

\noindent\textbf{Step 5.} Choose a spanning tree and a nonzero value at one vertex.
Its edge equations $r_{v,e}a_v=r_{w,e}a_w$ successively determine
nonzero values at all other vertices. Thus the tree rows of $d_1$
are independent and have a one-dimensional kernel. In particular
if the whole graph were a tree, $d_1$ would be surjective and
$H^2(S,\OO_S)=0$.

\noindent\textbf{Step 6.} The map $f$ is flat: multiplication by its nonzero local parameter
is injective in each regular local ring of the smooth total space,
and torsion-free modules over the base discrete valuation ring
are flat. Upper semicontinuity for this proper flat family gives
$h^2(S,\OO_S)\ge h^2(F,\OO_F)=1$. The central fibre is the
reduced $S$ by hypothesis, so no thickened-fibre cohomology has
been substituted here. The tree case is impossible.

\ref{cdp:st:C.2.UN.44} gives a surjection from $\pi_1(S)$ to the incidence
graph group, and $\pi_1(S)$ is a quotient of $\Z^4$ because the
full fibre is reduced. The free graph group is therefore abelian
and has rank at most one. Together with the preceding paragraph
this proves that the graph is unicyclic and $|E|=|V|$.

\noindent\textbf{Step 7.} There is now just one row left after the spanning-tree equations.
Writing its endpoints using the successive nonzero values shows
that it vanishes precisely when the product of the actual ratios
$r_{v,e}/r_{w,e}$ around the cycle is one. If it did not vanish,
$d_1$ would have rank $|V|$ and again $H^2(S,\OO_S)=0$.
The product is consequently one, the rank is $|V|-1$, and
\eqref{cdp:eq:C.HR.6} gives exactly \eqref{cdp:eq:C.HR.3}. A vertex-basis change cancels
between the incoming and outgoing ratios; an edge-basis change
cancels inside its ratio. This proves the asserted independence
without discarding the analytic restriction coefficients.

\par\medskip\noindent\textbf{The dualizing form and the canonical-order minimum.}
Let $D_v$ be the disjoint union of the double curves on $H_v$.
The residue exact sequence is
\[
 0\longrightarrow K_{H_v}
 \longrightarrow K_{H_v}(D_v)
 \xrightarrow{\operatorname{Res}}\bigoplus_{e\ni v}K_{C_e}
 \longrightarrow0.                                    \tag{C.HR.7}\label{cdp:eq:C.HR.7}
\]
Since $H^0(K_{H_v})=0$, a logarithmic two-form is determined by
its residues. Serre duality identifies the connecting map from
these residues to $H^1(K_{H_v})$ with the dual of the actual
restriction maps \eqref{cdp:eq:C.HR.1}. The signs may be fixed explicitly:
in a double chart $xy=0$ the two residues of a dualizing form
sum to zero. Orienting the edge takes its residue on the tail
and the negative on the head. Thus the global residue vector of
a section of $\omega_S$ belongs exactly to $\ker d_1^\vee$.
Conversely a vector in that kernel lifts uniquely by \eqref{cdp:eq:C.HR.7} on
each component, and the opposite-residue rule glues these forms
to a section of $\omega_S$. These are inverse maps.

\noindent\textbf{Step 8.} At a leaf of an attached tree, the corresponding row of
$d_1^\vee$ contains only its one nonzero coefficient. The
residue on that edge is zero. Delete that leaf and repeat.
Every tree-edge residue is zero. On the remaining cycle each
vertex relates its two adjacent residues by nonzero coefficients;
its one-dimensional kernel has a nonzero residue on every cycle
edge. On an off-cycle component all residues vanish, so the
form is an ordinary holomorphic two-form and is zero. On a cycle
component its two nonzero residues make it nonzero. Vanishing
residue along any other boundary curve removes the logarithmic
pole there, directly from the local expression
$a(y,z)\,dy\wedge dz/y$.

\noindent\textbf{Step 9.} Finally take the canonical evaluation section of \ref{cdp:st:C.2.UN.38} and
write its vertical divisor as $\sum k_vH_v$. Put $k_0=\min k_v$.
Because $S$ is the full reduced divisor of $t$, division by
$t^{k_0}$ is an actual holomorphic section of the relative
canonical line. Its restriction to $S$ is a nonzero section of
$\omega_S$. At an edge joining $v,w$, the explicit adjunction
calculation, with $t=xy$, is
\[
 \frac{x^{k_v-k_0}y^{k_w-k_0}\,dx\wedge dy\wedge dz}{xy}
 \quad\longmapsto\quad
 y^{k_w-k_0-1}\,dy\wedge dz\quad\text{on }x=0
                                                        \tag{C.HR.8}\label{cdp:eq:C.HR.8}
\]
when $k_v=k_0$, times a holomorphic unit; when $k_v>k_0$ the
restriction on that component is zero. The local logarithmic
expression is the usual restriction of the relative dualizing
line, or equivalently the residue of the displayed meromorphic
three-form. It follows that its nonzero components are exactly
those with $k_v=k_0$. The preceding residue calculation identifies
these with the cycle components, as claimed.

\par\medskip\noindent\textbf{The full algebraic-dimension-zero branch.}
The actual primary and secondary curve computations in
\ref{cdp:st:C.2.UN.38} and \ref{cdp:st:C.2.UN.43} give two possibilities in algebraic
dimension zero: a nonlinear Hopf surface has just one elliptic
curve, and a diagonal Hopf surface has just its two distinct
axis elliptic curves. The diagonal finite action cannot exchange
them: such an exchange would force equal positive powers of the
two contracting multipliers, giving a nonconstant meromorphic
invariant and algebraic dimension one. Thus every vertex of the
present graph has degree at most two.

\noindent\textbf{Step 10.} A connected unicyclic finite graph with all degrees at most two
is a cycle and has every degree equal to two. Indeed its equality
$|E|=|V|$ gives average degree two, so no smaller degree is possible.
Consequently there is no nonlinear component and no attached
tree. There is no loop at a single vertex, since the components
of the specified SNC surface are smooth embedded components:
such a loop would be a self-identification, not a double curve
between distinct components. Hence there are at least two
components.

\noindent\textbf{Step 11.} For this now established diagonal cycle, the actual normal-product
trivialization of $t=xy$ gives the dual-normal maps
(C.HC.3). Removing the zero sections gives the actual open groups
$G_v=\C^2/\Lambda_v$ and the affine logarithmic edge maps
(C.HC.5). Their integral lattice product is precisely $M$, and
\ref{cdp:st:C.2.UN.46} proves that an actual torus smoothing forces $M=I$.
Conversely the specified axis maps and any compatible dual-normal
maps with $M=I$ give the proper plumbing of \ref{cdp:st:C.2.UN.46}; the actual
varying-period comparison in \ref{cdp:st:C.2.UN.48} extends it to a compact
smooth threefold over a compact curve. That theorem retains every
component and every specified gluing constant, so the central
support is isomorphic to the specified $S$, rather than merely
having the same graph. Its other supports have Euler zero.

\noindent\textbf{Step 12.} The exact matrix calculation in
\path{code/E_hopf_graph_reduction.py} checks the spanning-tree
elimination, the cycle product, both coherent dimensions, and the
support of the dual residue vector on several branched graphs.
The argument above proves each comparison for every finite graph.
\end{proof}






\begin{proposition}[The complete all-ruled cycle criterion with its full elliptic translation]\label{cdp:st:C.2.UN.55}
\StatementRecord{cdp:st:C.2.UN.55}{proved}{Historical events: Claude Fable 5.2 (Anthropic), effort high, 16 November 2026: carried as background without fault found. GPT-6 Astra (OpenAI), ChatGPT Pro mode, effort Pro, S6\_09\_PAPERS, 17 November 2026: targeted proof check. Also checked in this writing. Pre-continuation writing events: GPT-6 Astra (OpenAI), ChatGPT Pro mode, effort Pro, S6\_10\_PAPERS, 18 November 2026: targeted proof check. Continuation metadata events: GPT-6 Astra (OpenAI), ChatGPT Pro mode, effort Pro, S6\_10\_PAPERS continuation, 19 November 2026: used and rechecked in this run; status-vocabulary reconciliation; metadata reconciliation only; no new mathematical check. \emph{Pre-continuation recorded proof scope.} The coordinate change preserves the actual rank-three lattice and all horizontal gluing translations. The real-linear twistor formula gives the exact positive four-period determinant for every nonzero finite base parameter, independently of those translations. Near zero its moving ruled line charts have precisely the prescribed affine period, and the reciprocal base chart gives the opposite-oriented second ruled filling. The explicit punctured-family maps and a finite compact base cover prove smooth proper compact gluing, including the single-normalization self-identification case}{the hypotheses, notation and definitions stated here; \cref{cdp:st:C.2.UN.51}; \cref{cdp:eq:C.TL.3}; \cref{cdp:eq:C.TL.1}; \cref{cdp:eq:C.TL.2}; \cref{cdp:eq:C.TL.5}; \cref{cdp:eq:C.TL.4}; \cref{cdp:eq:C.TL.7}; \cref{cdp:eq:C.TL.8}}

Specify an oriented cycle of $N\geq1$ marked ruled elliptic surfaces
$R_i=\PP(\OO_{E_i}\oplus L_i)$, with
$L_i\in\operatorname{Pic}^0(E_i)$, elliptic identifications of
successive marked sections, and duality of their actual normal
lines. For $N=1$ identify the two distinct marked sections of
the same normalization. The following are equivalent:
the integral complement-lattice product $M$ of \ref{cdp:st:C.2.UN.51} is the
identity; the specified support has an actual smooth-total
two-torus smoothing; and it occurs as a full reduced fibre of
a smooth compact complex threefold over $\PP^1$ with two-torus
general fibre.

The compact realization retains every marked surface and every
elliptic gluing map, including an arbitrary total elliptic
translation around the cycle. It chooses the normal-product
constants. Its only singular fibres are the specified cycle and
another reduced cycle of $N$ ruled elliptic components. Every
reduced fibre support has Euler number zero.
\end{proposition}



\begin{proof}
\noindent\textbf{Step 1.} The necessary condition and the actual local equivalence are
\ref{cdp:st:C.2.UN.51}. We prove the compact assertion without removing the
horizontal affine translation.

\par\medskip\noindent\textbf{An actual common lattice and its retained translation.}
Use the actual logarithmic linear identifications of \ref{cdp:st:C.2.UN.51} to put
all complements in $G=\C^2/\Lambda$, where $\Lambda$ has real
rank three. All positive outgoing normal meridians are the same
primitive $n\in\Lambda$. Put $p=-n$. The normal-line quotient of
each ruled surface is an elliptic curve. Consequently
$\C p\cap W=\R p$, where $W=\operatorname{span}_{\R}\Lambda$.
The complex line $W\cap iW$ and $p$ are a complex basis of
$\C^2$. The complex linear map sending this line to the first
coordinate line and $p$ to $(0,1)$ puts $W$ exactly in
$\C\times\R$.

\noindent\textbf{Step 2.} Extend the primitive $p$ to an integral basis of $\Lambda$.
After ordering the other two vectors, its exact form is
\[
 \begin{gathered}
 w_1=(a_1,c_1),\qquad w_2=(a_2,c_2),\qquad p=(0,1),\\
 c_1,c_2\in\R,\qquad
 C=\operatorname{Im}(\overline{a_1}a_2)>0.
 \end{gathered}
 \tag{C.TL.1}\label{cdp:eq:C.TL.1}
\]
This changes coordinates, not the lattice or its integral
subgroups. In common coordinates every actual edge map is a
translation with its normal logarithmic term. Choose logarithmic
lifts of its constants $b_i$ and put $b=\sum_i b_i$.
Its first coordinate, modulo
$\Z a_1+\Z a_2$, is exactly the total elliptic gluing translation.
A normal-product constant changes $b_i$ only in the line
$\C p$. Adjust one such constant so that
\[
                         b=(\sigma,i).
 \tag{C.TL.2}\label{cdp:eq:C.TL.2}
\]
The prescribed elliptic maps, hence the horizontal class of
$\sigma$, are unchanged. No restriction on that class has been
imposed.

\par\medskip\noindent\textbf{A holomorphic period family on the entire punctured sphere.}
For $z\in\C$ and a fixed vector $w=(w^{(1)},w^{(2)})\in\C^2$
define
\[
 V_z(w)=
 \bigl(w^{(1)}+z\overline{w^{(2)}},
       w^{(2)}-z\overline{w^{(1)}}\bigr).
 \tag{C.TL.3}\label{cdp:eq:C.TL.3}
\]
For each fixed $w$ this is a holomorphic vector-valued function
of $z$. For each fixed $z$ it is a real linear isomorphism,
with real determinant $(1+|z|^2)^2$. On a branch of $\log z$,
put
\[
 \begin{gathered}
 P_1(z)=V_z(w_1),\quad P_2(z)=V_z(w_2),\quad
 P_3(z)=V_z(p),\\
 P_4(z)=V_z(b)+\frac{N}{2\pi i}\log z\,V_z(p).
 \end{gathered}
 \tag{C.TL.4}\label{cdp:eq:C.TL.4}
\]
Changing the logarithm by $2\pi i$ replaces $P_4$ by
$P_4+NP_3$. Thus the subgroup generated by these four vectors
is single-valued. We check its real rank everywhere, including
where the rank-three group alone has no elliptic normal quotient.

\noindent\textbf{Step 3.} Write $r=|z|$ and $\tau=N\log z/(2\pi i)$. The actual inverse
of \eqref{cdp:eq:C.TL.3} gives
\[
 V_z^{-1}\bigl(iV_z(p)\bigr)
       =\frac{i(2z,1-r^2)}{1+r^2}.
 \tag{C.TL.5}\label{cdp:eq:C.TL.5}
\]
The first three vectors in \eqref{cdp:eq:C.TL.1} span $\C\times\R$, so
their real determinant with any fourth vector $v$ is
$C\operatorname{Im}v^{(2)}$. Equation \eqref{cdp:eq:C.TL.2} therefore gives
$\det_{\R}(w_1,w_2,p,b)=C$. The real multiple
$\operatorname{Re}\tau\,V_z(p)$ has no effect on a determinant.
Using \eqref{cdp:eq:C.TL.5} for the remaining multiple yields the exact value
\[
 \begin{aligned}
 D(z)&=\det_{\R}(P_1(z),P_2(z),P_3(z),P_4(z))\\
 &=C(1+r^2)
       \left(1+r^2+\frac{N}{2\pi}(r^2-1)\log r\right)>0
                         \qquad (0<r<\infty).
 \end{aligned}
 \tag{C.TL.6}\label{cdp:eq:C.TL.6}
\]
Indeed $(r^2-1)\log r\geq0$, while $1+r^2>0$.
In particular the retained horizontal $\sigma$ and the real
characters $c_1,c_2$ require no inequality.

\noindent\textbf{Step 4.} The four holomorphic period sections on a simply connected
base chart act on $\C^2$ by translations. Their real determinant
is nonzero; on each compact base subset their real matrices and
inverses are uniformly bounded. Thus their integral translation
action is properly discontinuous, and a continuously varying
closed real parallelepiped gives properness of the quotient.
The integral change $P_4\mapsto P_4+NP_3$ identifies these
quotients on overlaps. This constructs a smooth proper family
$\mathcal T\to\C^*$ of connected two-tori, with these actual
periods.

\par\medskip\noindent\textbf{The prescribed filling at zero.}
For $z$ sufficiently close to zero, the line
$\C V_z(p)$ is transverse to the real rank-three lattice in
the same sense as at zero. Indeed its transverse determinant
is $C(1+r^2)(1-r^2)$, which is nonzero for $r<1$.
Use the separate ruled moving frame of (C.MX.10):
the normal vector $V_z(n)$ and a transverse integral lattice
lift define an actual elliptic family and a holomorphic
degree-zero line on it. Its projective completion gives the
marked ruled family through each prescribed $R_i$.
The original integral identifications and all finite changes
of logarithm are retained.

\noindent\textbf{Step 5.} Choose the moving edge translation to be $V_z(b_i)$ in the
common coordinates. It induces the actual translation on the
elliptic quotient and the reciprocal normal-line map; these
extend over the marked sections by their line-bundle formulas.
The normal logarithmic coefficient of each edge is
$-V_z(n)/(2\pi i)=V_z(p)/(2\pi i)$. Consequently the entire
cycle has precisely the fourth period in \eqref{cdp:eq:C.TL.4}, not only
its leading asymptotic term. The central value of each edge
is its prescribed elliptic gluing map.

\noindent\textbf{Step 6.} Apply the proper normal-product plumbing of \ref{cdp:st:C.2.UN.51} to these
moving line charts, with actual equation $xy=z$ in each neck.
The moving component cores are smooth over the disc; the
normal-product equation is smooth at $x=y=z=0$.
Uniform flat norms and finitely many bounded cores and necks
give the same Hausdorff proper local filling as in that proof.
Its full central divisor is reduced, and its support is the
specified marked cycle.

\noindent\textbf{Step 7.} On a punctured base chart, send each core's logarithmic
coordinate to the common vector space, adding the successive
edge translations. Across each neck these maps agree modulo
the periods in \eqref{cdp:eq:C.TL.4}. Across a change of branch of $\log z$
their changes are integral multiples of $P_3$. Going around
the cycle adds exactly $P_4$. They therefore give a holomorphic
map to $\mathcal T$. Its fibre maps are the closed-coframe
isomorphisms of \ref{cdp:st:C.2.UN.51}: all four fundamental-group generators
are exactly the displayed periods. Their local inverses
depend holomorphically on the base, so this is an actual
biholomorphism of punctured families.

\par\medskip\noindent\textbf{The second proper filling at infinity.}
Put $s=z^{-1}$ and change the two-dimensional vector-bundle
frame by $\eta_\infty=\eta/z$. In this frame \eqref{cdp:eq:C.TL.3} becomes
\[
 V^\infty_s(w)
       =\bigl(\overline{w^{(2)}},-\overline{w^{(1)}}\bigr)
                          +s w.
 \tag{C.TL.7}\label{cdp:eq:C.TL.7}
\]
This is an actual holomorphic frame change on $\C^*$, and
$V^\infty_0$ is an invertible antilinear map. The periods
in the new frame are exactly
\[
 \begin{gathered}
 V^\infty_s(w_1),\quad V^\infty_s(w_2),\quad V^\infty_s(p),\\
 V^\infty_s(b)
          -\frac{N}{2\pi i}\log s\,V^\infty_s(p).
 \end{gathered}
 \tag{C.TL.8}\label{cdp:eq:C.TL.8}
\]
The ambiguity of the logarithm again adds an integral period.

\noindent\textbf{Step 8.} Near $s=0$ the rank-three lattice and the complex line through
$V^\infty_s(p)$ give a genuine elliptic quotient; the
transversality at zero follows by applying the invertible
antilinear map in \eqref{cdp:eq:C.TL.7} to the original quotient.
Construct $N$ marked ruled components from this actual
nonzero line bundle. Give each positive outgoing meridian
the vector $n_\infty(s)=V^\infty_s(p)$, and each incoming
meridian its negative. Choose the sum of their full affine
edge constants to be $V^\infty_s(b)$, for example by putting
this constant on one edge and zero on the others.
Every such constant induces an actual elliptic translation
and its compatible normal-line map.

\noindent\textbf{Step 9.} The normal-period sum of this opposite-oriented cycle is
$-N V^\infty_s(p)$. Its normal-product plumbing with equation
$xy=s$ therefore has precisely the periods \eqref{cdp:eq:C.TL.8}.
The same finite-core construction is a smooth proper local
filling with reduced central divisor. Its exact logarithmic
maps identify the punctured family with the frame change of
$\mathcal T$ just displayed. This explicitly constructs the
second filling; no contraction of a global auxiliary line
bundle is involved.

\par\medskip\noindent\textbf{Compact gluing and the Euler comparison.}
Choose small disjoint discs about $0$ and $\infty$ and a middle
annulus covering the rest of $\PP^1$. Glue their two proper
fillings to $\mathcal T$ by the actual punctured-family
biholomorphisms above. Over every point of the base all pieces
are identified with one Hausdorff proper local model. Hence
the glued space is Hausdorff and the map is proper locally
on the base. A finite closed refinement of the three base
charts makes its inverse image a finite union of compact
sets, proving properness over the compact $\PP^1$.
The total space is smooth by its displayed local models.
All other fibres are the tori of \eqref{cdp:eq:C.TL.4}.

\noindent\textbf{Step 10.} Each ruled elliptic component has Euler number
$e(\PP^1)e(E)=0$, and each double curve has Euler number zero.
There are no triple points, so both singular supports have
Euler number zero. This proves every asserted compact and
marked-data comparison, including the one-component
self-identification.

\noindent\textbf{Step 11.} The symbolic determinant, frame-change, inverse-map and
monodromy identities are certified in
\path{code/A_twistor_log_cycles.py} and its saved output.
\end{proof}






\begin{proposition}[The complete weighted mixed-cycle criterion]\label{cdp:st:C.2.UN.56}
\StatementRecord{cdp:st:C.2.UN.56}{proved}{Historical events: Claude Fable 5.2 (Anthropic), effort high, 16 November 2026: carried as background without fault found. GPT-6 Astra (OpenAI), ChatGPT Pro mode, effort Pro, S6\_09\_PAPERS, 17 November 2026: targeted proof check. Also checked in this writing. Pre-continuation writing events: GPT-6 Astra (OpenAI), ChatGPT Pro mode, effort Pro, S6\_10\_PAPERS, 18 November 2026: targeted proof check. Continuation metadata events: GPT-6 Astra (OpenAI), ChatGPT Pro mode, effort Pro, S6\_10\_PAPERS continuation, 19 November 2026: used and rechecked in this run; status-vocabulary reconciliation; metadata reconciliation only; no new mathematical check. \emph{Pre-continuation recorded proof scope.} The actual ruled normal bundle has degree minus (a+b)/m, and its finite-index complement lattice gives the exact degree-q fibre map and extended ruled root cover. The connected root graph, inertia groups and Hopf normal periods exclude a one-Hopf cycle and force every ruled line to be non-torsion. The finite linear action is then identity. Orbitwise moving ruled frames, the exact affine equations, local invariant rings and the punctured translation-torsor comparison give the compact descent with every original multiplicity and gluing map}{the hypotheses, notation and definitions stated here; \cref{cdp:st:C.2.UN.52}; \cref{cdp:eq:C.WM.4}; \cref{cdp:eq:C.WM.5}; \cref{cdp:eq:C.WM.6}; \cref{cdp:st:C.2.UN.51}}

The criterion and compact realization of \ref{cdp:st:C.2.UN.52}
remain valid when some vertices are marked ruled surfaces
\[
        R_i=\PP(\OO_{E_i}\oplus L_i),\qquad
        E_i\ \hbox{elliptic},\qquad L_i\in\operatorname{Pic}^0(E_i),
\tag{C.WM.1}\label{cdp:eq:C.WM.1}
\]
with their zero and infinity sections, provided at least
one vertex is a diagonal primary or secondary Hopf surface
of algebraic dimension zero. The support is a simple-NC
cycle of embedded components, with edges between distinct
components. All marked surfaces, elliptic gluing maps and
positive component multiplicities are retained.

More precisely choose the actual compatible weighted normal
lines, sections and boundary maps (C.WH.1)--(C.WH.2) on
these marked surfaces. Form their normalized common root
covers and the compressed necks $uv=s^{h_i}$ as in \ref{cdp:st:C.2.UN.52}.
The resulting root graph must be connected and its actual
rank-three integral lattice product must be the identity.
These conditions are necessary and sufficient for connected
two-torus fibres in the explicit smooth weighted plumbing.
Every weighted support passing them occurs over $\PP^1$
with smooth compact total space and Euler number zero for
every reduced fibre support. Every actual torus degeneration
in this specified branch gives these data and passes the test.

There are explicit additional necessary relations at a ruled
vertex of multiplicity $m$, with neighbouring multiplicities
$a,b$:
\[
 \begin{gathered}
      m\mid a+b,\qquad
      \deg(\mathcal L|_{\PP^1})=-\frac{a+b}{m},\\
      \gcd(m,a)=\gcd(m,b),\qquad
      q=\frac{m}{\gcd(m,a)}
         =\frac{m}{\gcd(m,b)}.
 \end{gathered}
\tag{C.WM.2}\label{cdp:eq:C.WM.2}
\]
Every connected strict root cover of this ruled component
is exactly a marked surface
$\PP(\OO_{E'}\oplus L')$, with an \'{e}tale elliptic map
$\pi:E'\to E$ and an actual line isomorphism
\[
                   (L')^q\simeq\pi^*L.
\tag{C.WM.3}\label{cdp:eq:C.WM.3}
\]
Its finite map has fibre degree $q$, ramification order
$q$ at each marked section, and base degree
$\gcd(m,a)/r$, where $r$ is the common number of strict
covers in \ref{cdp:st:C.2.UN.52}. In particular its total degree is $m/r$.
All the component, axis and inertia formulas (C.WH.0)
hold unchanged.

A compatible cycle in this branch has at least two
original Hopf vertices. Every ruled line $L_i$ in it
has infinite order in $\operatorname{Pic}^0(E_i)$.
\end{proposition}



\begin{proof}
\noindent\textbf{Step 1.} \par\medskip\noindent\textbf{The weighted normal line on an actual ruled component.}
Let $C^-,C^+$ be its two marked sections and write
\[
  \mathcal L^{m}
       \simeq \OO_R(-aC^- -bC^+),\qquad
       \operatorname{div}(\sigma)=aC^-+bC^+.
\tag{C.WM.4}\label{cdp:eq:C.WM.4}
\]
Restriction to any ruling fibre gives
$m\deg(\mathcal L|_{\PP^1})=-(a+b)$,
proving the first two assertions of \eqref{cdp:eq:C.WM.2}.
The congruence $b\equiv-a\pmod m$ gives
the equality of the two gcd's and the displayed
value of $q$.

\noindent\textbf{Step 2.} The local line-bundle comparison needed for the weighted
plumbing is holomorphic, including when $\mathcal L$
has nonzero degree on ruling fibres.
Each section neighbourhood obtained by removing the
opposite section is the entire total space of its
normal degree-zero line on $E$.
On this chart, the right side of \eqref{cdp:eq:C.WM.4} is the
pullback of a degree-zero line on $E$:
the removed divisor is trivial there, and the divisor
of the zero section is the pullback of its normal line.
Thus $\mathcal L^m$ is flat on this chart.
Dividing its transferred flat connection by $m$
gives an actual flat connection on $\mathcal L$ there.
The chart retracts to its zero section, so its
monodromy character agrees with the pullback of
the restricted character. This gives the actual
holomorphic isomorphism with the pulled-back
axis restriction. The weighted component cores
and normal-product necks consequently glue exactly
as in \ref{cdp:st:C.2.UN.52}. No assertion that $\mathcal L$ is
globally flat on the ruled surface is used.

\noindent\textbf{Step 3.} For an alternative global description of this
normal line, take $\OO_R(C^+)$ as a relative
degree-one line and put
$d=\deg(\mathcal L|_{\PP^1})$.
The line $\mathcal L\otimes\OO_R(-dC^+)$
has trivial restriction to every ruling fibre.
Its direct image is a line on $E$, and evaluation
is an isomorphism on every fibre and hence everywhere,
by proper base change and the local Nakayama lemma.
Therefore
$\mathcal L=\OO_R(dC^+)\otimes\pi_R^*A$.
Equation \eqref{cdp:eq:C.WM.4} and the fact that
$C^--C^+$ is pulled back from $\operatorname{Pic}^0(E)$
give $\deg A=0$. This describes the same local
flat pullbacks by an explicit bundle map.

\par\medskip\noindent\textbf{The complete ruled root-cover comparison.}
Remove the two sections. Its actual open group is
$G=L^\times=\C^2/\Lambda$, with the rank-three
period lattice (C.MX.3). Its positive normal
meridian $n$ is primitive in $\Lambda$.
Let $\Lambda'\subset\Lambda$ be the finite-index
sublattice of one connected cyclic root component
over this open set, and let $q$ be the smallest
positive integer for which $qn\in\Lambda'$.
Projection along the complex line $\C n$
gives the actual elliptic base covering
\[
 E'=\C^2/(\C n+\Lambda')
      \longrightarrow
 E=\C^2/(\C n+\Lambda).
\tag{C.WM.5}\label{cdp:eq:C.WM.5}
\]
This is finite and \'{e}tale: the two projected
rank-two lattices are discrete, and one has
finite index in the other. The source open
group maps to $E'$ with kernel
$\C n/\Z qn\simeq\C^*$.
Taking logarithms of its fibre coordinate
gives constant factors of automorphy, hence
an actual degree-zero line $L'$ on $E'$.

\noindent\textbf{Step 4.} The homomorphism of open groups induced by
$\Lambda'\subset\Lambda$ acts on this
$\C^*$ kernel by the power $z\mapsto z^q$.
Equivalently it gives exactly \eqref{cdp:eq:C.WM.3}.
For example choose lifts of the two base
periods to the source logarithmic lattice.
If their normal coordinates are $b_j'$,
the image coordinates are $qb_j'$ modulo
integral angular periods; exponentiation
is the line isomorphism in \eqref{cdp:eq:C.WM.3}.
Thus the open finite map extends to the
actual projective-bundle map
\[
 \PP(\OO_{E'}\oplus L')
          \longrightarrow \PP(\OO_E\oplus L),
\tag{C.WM.6}\label{cdp:eq:C.WM.6}
\]
using the given base map and the homogeneous
power map on each ruling fibre.
This map is finite, has ramification order
$q$ at both marked sections, and agrees
with the original cover on the dense open
group. Its source is smooth and normal.
Uniqueness of finite normalization therefore
identifies it with the actual strict root
component, including both sections.

\noindent\textbf{Step 5.} This also proves directly that the two
ramification orders are equal. Applying
the local normalization
$\xi^m=z^a$ and $\xi^m=z^b$ to the two
sections identifies their common value
with \eqref{cdp:eq:C.WM.2}. The lattice index splits as
\[
 [\Lambda:\Lambda']
     =q\,[\Lambda:\Lambda'+\Z n].
\tag{C.WM.7}\label{cdp:eq:C.WM.7}
\]
The second factor is the degree of \eqref{cdp:eq:C.WM.5}.
Since \ref{cdp:st:C.2.UN.52}'s component count gives total
degree $m/r$, it equals $\gcd(m,a)/r$.
The two section covers have that same
degree, proving all the asserted maps
and degrees without replacing the root
surface by an unspecified ruled model.

\par\medskip\noindent\textbf{The connected root cycle and its finite action.}
Each strict root component is now either
one of the actual diagonal Hopf covers of
\ref{cdp:st:C.2.UN.52} or the marked ruled surface \eqref{cdp:eq:C.WM.6}.
Every one has two connected marked elliptic
curves. The local root charts are still
exactly (C.WH.4), and their resolutions
insert the same two-section ruled bridges.
Thus the root graph is a union of cycles.
Connectedness, reduced-root Stein factorization
and the actual core-and-neck group calculation
give the same criterion
\[
                         M=I
\tag{C.WM.8}\label{cdp:eq:C.WM.8}
\]
for a connected two-torus general fibre.
This is also the mixed-cycle criterion
of \ref{cdp:st:C.2.UN.51}, applied to the resolved root,
with the bridges compressed by their
exact $uv=s^h$ maps.

\noindent\textbf{Step 6.} The root deck group preserves the actual
map to each original normalization and
its two labelled marked branches. It
therefore preserves the lifted cyclic
orientation and acts by rotations.
Its common vertex stabilizer, connected
cover counts, generic inertia and axis
stabilizers are exactly those computed
in (C.WH.6) and (C.WH.0). The local
degrees at a ruled vertex just proved
are consistent with, and refine, those
same actual sheet counts.

\noindent\textbf{Step 7.} When $M=I$, identify all open complements
with one $G=\C^2/\Lambda$.
On a strict ruled component the two
normal meridians are opposite, so
the common outgoing periods satisfy
$n_j=n_{j-1}$ there. At every Hopf
component they are complex-linearly
independent and satisfy the same
contracting-height identity (C.HG.4).
The positive compressed edge lengths
therefore give
\[
                  u=-\sum_j h_j n_j,\qquad \ell(iu)\ne0.
\tag{C.WM.9}\label{cdp:eq:C.WM.9}
\]
The common group has a Hopf open chart,
so it contains no compact elliptic curve.
Its quotient along the primitive part
of $u$ is the same nontorsion degree-zero
elliptic line used in \ref{cdp:st:C.2.UN.52}.

\noindent\textbf{Step 8.} The finite root action is affine also
on the ruled open charts. A marked
projective-bundle isomorphism has an
affine elliptic base map and a
nonzero fibre multiplier. The logarithmic
derivative of that multiplier is
holomorphic and doubly periodic,
hence constant, exactly as in \ref{cdp:st:C.2.UN.51}.
The common linear part $R$ fixes $u$
because the actual action is a
rotation. The nontorsion line argument
of (C.WH.8) consequently proves
$R=I$ without any change.
Summing the exact affine equations
(C.WH.9) again gives
$u/N\in\Lambda$.

\noindent\textbf{Step 9.} Every ruled line $L_i$ has infinite
order. Otherwise its finite pullback
and $q$th root \eqref{cdp:eq:C.WM.3} would make
$L'$ torsion as well. Its nonzero
bundle would contain the image of
an elliptic curve under a finite
trivializing isogeny. This would be
a compact elliptic curve in the
common Hopf open group, a contradiction.

\noindent\textbf{Step 10.} There must be at least two original
Hopf vertices. If there were only one,
start at a lifted Hopf vertex and
follow the cycle to the next lifted
vertex above that same original Hopf.
Every intermediate strict vertex is
ruled, so its incoming and outgoing
normal periods agree; compressing
or resolving a neck has the same
property. Hence the common period
immediately preceding the next Hopf
is the outgoing normal period of
the first Hopf, in the notation of
(C.MX.5). A root deck rotation takes
the first marked Hopf to the second.
Its linear part is the identity,
so their respective outgoing normal
periods agree as well. At that
second Hopf the two consecutive
normal periods would therefore be
equal, contradicting their complex
linear independence. This proves
the final component-count assertion.

\par\medskip\noindent\textbf{Compact realization preserves every marked surface.}
The finite transport of \ref{cdp:st:C.2.UN.52} applies
to this mixed root with one precise
additional frame choice. At a Hopf
vertex transport its two independent
normal periods. At a ruled vertex,
transport instead its primitive
normal period, a transverse base
lattice lift and a third lattice
basis vector, as in (C.MX.10).
This gives an actual family
$\PP(\OO_{E'(t)}\oplus L'(t))$
of the given marked ruled surface,
rather than attempting to use its
dependent pair of normal meridians.

\noindent\textbf{Step 11.} These frames can be chosen identically
on each root-deck orbit: the common
linear part is the identity and
the corresponding normal periods
and lattice subgroups agree.
Translations of the common open
group extend to the marked ruled
compactification. Indeed its
projection to $E'$ is a base
translation, and the group law on
the rigidified degree-zero line
gives the fibre isomorphism;
this extends to its zero and
infinity sections. Hence the
transported affine deck maps
extend to these ruled families
as actual holomorphic maps.

\noindent\textbf{Step 12.} All equations for the edge
translations and deck translations
are still the constant-integer
system (C.WH.9).
The two adjacent independent
directions at a Hopf vertex belong
to distinct rotation orbits, so
the two orbitwise normal constants
still set the total translation
to zero without changing any
central elliptic gluing map.
The rational-left-kernel argument
and fixed right inverse of \ref{cdp:st:C.2.UN.52}
give the full holomorphic finite
action on the moving compressed
root plumbing.

\noindent\textbf{Step 13.} On every neck retain $uv=s^{h_i}$.
Its actual pointwise stabilizer
has the same invariant ring
$\C\{u^B,v^A\}$ as (C.WH.12),
and its residual axis action is
\'{e}tale. Component inertia is
again a normal reflection.
Thus the quotient total space is
smooth and has exactly the
original full multiplicities.
At zero the components and all
their marked maps and finite
quotients are the original ones,
by \eqref{cdp:eq:C.WM.6} and the Hopf comparison.
The punctured quotient is the
marked torus torsor with period
index $\operatorname{div}(u)/N$;
the exact sequence (C.WH.13)
gives its actual holomorphic
isomorphism to the compact
elliptic suspension.

\noindent\textbf{Step 14.} Performing the same proper
annular replacement therefore
gives the stated compact family
over $\PP^1$. Ruled elliptic
components, Hopf components and
their double elliptic curves all
have Euler number zero. Every
other reduced support is unchanged
from the compact suspension and
also has Euler number zero.
This proves both sufficiency and
the asserted exact retention.
Necessity for every actual
smoothing uses its actual weighted
normal data and the tubular
comparison of \ref{cdp:st:C.2.UN.52}, now with
the explicit ruled maps above.
\end{proof}






\begin{proposition}[The exact marked-cycle criterion without an algebraic-dimension restriction]\label{cdp:st:C.2.UN.57}
\StatementRecord{cdp:st:C.2.UN.57}{proved}{Historical events: Claude Fable 5.2 (Anthropic), effort high, 16 November 2026: carried as background without fault found. GPT-6 Astra (OpenAI), ChatGPT Pro mode, effort Pro, S6\_09\_PAPERS, 17 November 2026: targeted proof check. Also checked in this writing. Pre-continuation writing events: GPT-6 Astra (OpenAI), ChatGPT Pro mode, effort Pro, S6\_10\_PAPERS, 18 November 2026: targeted proof check. Continuation metadata events: GPT-6 Astra (OpenAI), ChatGPT Pro mode, effort Pro, S6\_10\_PAPERS continuation, 19 November 2026: used and rechecked in this run; status-vocabulary reconciliation; metadata reconciliation only; no new mathematical check. \emph{Pre-continuation recorded proof scope.} The integral coframe and normal-meridian sum use only the actual two marked Hopf axes and ruled sections. The real-linear period deformation preserves every finite character and contraction inequality; its explicit determinant stays positive on the whole punctured line. Separate component frames identify the zero filling with the prescribed marked cycle. Dividing the infinity frame by the base parameter gives the exact conjugate contractions, lattice maps and translations, hence an actual second filling and compact smooth family}{the hypotheses, notation and definitions stated here; \cref{cdp:st:C.2.UN.46}; \cref{cdp:st:C.2.UN.51}; \cref{cdp:eq:C.AD.1}; \cref{cdp:eq:C.AD.2}; \cref{cdp:eq:C.AD.3}; \cref{cdp:st:C.2.UN.55}; \cref{cdp:eq:C.AD.6}; \cref{cdp:eq:C.AD.10}; \cref{cdp:eq:C.AD.8}; \cref{cdp:eq:C.AD.9}}

Specify a connected oriented cycle of marked surfaces, each either a
diagonal primary or secondary Hopf surface with its two distinct axes,
or $\PP(\OO_E\oplus L)$ with $L\in\operatorname{Pic}^0(E)$ and its
zero and infinity sections. Specify the elliptic gluing maps and
duality of the actual axis or section normal lines. Components are
smooth embedded components, except that a single ruled normalization
with its two sections identified is also allowed.

The actual rank-three complement-lattice product $M$ is the identity
precisely for those specified supports which admit a smooth-total
two-torus smoothing. Every such support occurs as a full reduced fibre
of a smooth compact complex threefold over $\PP^1$. The construction
retains each component and every specified elliptic gluing map. Its
only other singular fibre is the complex conjugate marked support,
and every reduced fibre support has Euler number zero. No restriction
on the algebraic dimension of the diagonal Hopf components is imposed.
The construction chooses normal-product constants.
\end{proposition}



\begin{proof}
\noindent\textbf{Step 1.} \par\medskip\noindent\textbf{The actual complement maps and the local criterion.}
For a diagonal Hopf component the exponential preimage of its full
diagonal deck group is an actual rank-three lattice $\Lambda_i$ in
$\C^2$. Its complement of both axes is $\C^2/\Lambda_i$. This
description retains the finite secondary group, including its two
angular characters. Both small normal meridians are primitive: a
smaller angular lattice vector would give a nonidentity finite deck
element fixing a nonzero point of the other axis, contrary to the
freeness of this Hopf quotient. For a ruled component the complement
is exactly $L^\times$; its degree-zero line-bundle factors give the
rank-three lattice and primitive normal meridian in (C.MX.3).

\noindent\textbf{Step 2.} An elliptic gluing map lifts to an affine map of universal covers.
A dual-normal line isomorphism lifts to a nowhere-zero function;
the derivative of its logarithm is periodic under the two elliptic
periods, hence constant. Its logarithm is therefore affine.
Consequently the actual normal-product neck has exactly the affine
logarithmic maps of (C.HC.5) and (C.MX.2), including their constants.
These statements use the displayed axes and sections; they use no
absence of further curves on a component.

\noindent\textbf{Step 3.} The component cores and the $xy=t$ necks give the actual fundamental
group
\[
                         \pi_1(F_t)=\Lambda\rtimes_M\Z.
                                                        \tag{C.AD.1}\label{cdp:eq:C.AD.1}
\]
For an arbitrary specified smoothing, this is the same comparison
as in \ref{cdp:st:C.2.UN.46} and \ref{cdp:st:C.2.UN.51}: its canonical divisor sections and a smooth
tubular retraction identify the component cores with their actual
normal circle bundles, and its two divisor sections identify each
double collar with $xy=t$. The changes of these identifications
give the specified elliptic and normal-line maps on fundamental
groups. Thus \eqref{cdp:eq:C.AD.1} also holds for the actual smoothing, without
asserting that its holomorphic neighbourhood is linear. The
abelian group of a two-torus forces $M=I$.

\noindent\textbf{Step 4.} Conversely the actual flat normal lines and affine necks give the
proper normal-product plumbing of those proofs. Their argument
requires only the entire line-bundle axis charts, disjoint marked
curves and the displayed constant factors, all present here.
When $M=I$, the common logarithmic coframe glues on every nearby
compact fibre. Its two dual holomorphic fields commute and are
complete. Their $\C^2$ action has discrete cocompact stabilizer
and identifies that fibre with a complex two-torus. This proves
the full local equivalence in the present generality.

\par\medskip\noindent\textbf{The positive total normal period.}
Assume $M=I$ and identify the complement groups with
$G=\C^2/\Lambda$ by the actual linear edge maps. Write $n_i$
for the outgoing normal meridian at vertex $i$; the incoming
one is $-n_{i-1}$. At a Hopf vertex these two vectors are
complex-linearly independent, and its actual logarithmic frame is
\[
       P_i=\left(\frac{n_i}{2\pi i},
                         -\frac{n_{i-1}}{2\pi i}\right).
                                                        \tag{C.AD.2}\label{cdp:eq:C.AD.2}
\]
At a ruled vertex $n_i=n_{i-1}$, and its separate frame is the
one in (C.MX.10).

\noindent\textbf{Step 5.} Put $W=\operatorname{span}_{\R}\Lambda$ and choose a nonzero
real linear functional $\ell$ with kernel $W$. A contracting
Hopf period with logarithms $\log\alpha_i,\log\beta_i$ lies
in $W$. Applying $\ell$ to \eqref{cdp:eq:C.AD.2} times that period gives
\[
 \log|\alpha_i|\,\ell(i n_i)
       =\log|\beta_i|\,\ell(i n_{i-1}).                 \tag{C.AD.3}\label{cdp:eq:C.AD.3}
\]
Both logarithms are negative. Neither transverse value can
vanish: their simultaneous vanishing would put both complex
normal lines, hence all $\C^2$, inside the real three-plane
$W$. Thus they have the same nonzero sign. At a ruled vertex
the two values coincide. In an all-ruled cycle their common
value is nonzero because the actual elliptic quotient satisfies
$\C n\cap W=\R n$. Hence all values around every present
cycle have the same nonzero sign.

\noindent\textbf{Step 6.} The exact fourth period is
\[
 b+\frac{\log t}{2\pi i}u,\qquad
             u=-\sum_i n_i\in\Lambda,
                                                        \tag{C.AD.4}\label{cdp:eq:C.AD.4}
\]
where $b$ is the sum of the actual affine edge constants.
Equation \eqref{cdp:eq:C.AD.3} gives $\ell(iu)\ne0$. Write $u=d p$ with
$p$ primitive and $d>0$. The actual complex linear coordinate
change of (C.TL.1) therefore gives
\[
 \begin{gathered}
 W=\C\times\R,\qquad p=(0,1),\\
 \Lambda=\langle (a_1,c_1),(a_2,c_2),p\rangle_{\Z},
 \quad c_1,c_2\in\R,\quad
 C=\operatorname{Im}(\overline{a_1}a_2)>0.
 \end{gathered}                                        \tag{C.AD.5}\label{cdp:eq:C.AD.5}
\]
Every $n_i$ has nonzero second coordinate, by the transverse
calculation. Changing one normal-product constant changes
its affine edge constant by an arbitrary multiple of $n_i$.
Choose it to make $b=(\sigma,i)$. This changes no elliptic
gluing map. In the all-ruled case it leaves the full horizontal
translation $\sigma$ unchanged, exactly as in \ref{cdp:st:C.2.UN.55}.

\par\medskip\noindent\textbf{The entire punctured sphere of actual tori.}
Use $V_z$ from (C.TL.3), and set
\[
 \begin{gathered}
 \Lambda(z)=V_z\Lambda,\qquad
 B(z)=V_z(b)+\frac{d}{2\pi i}\log z\,V_z(p),\\
                   T_z=\C^2/(\Lambda(z)+\Z B(z)).
 \end{gathered}                                        \tag{C.AD.6}\label{cdp:eq:C.AD.6}
\]
These are holomorphic period sections for every $z\in\C^*$.
The change of a logarithm adds the integral period $dV_z(p)$.
The exact real determinant proved in (C.TL.5)--(C.TL.6), with
$N$ replaced by the positive integer $d$, is
\[
 C(1+|z|^2)\left(1+|z|^2+
          \frac{d}{2\pi}(|z|^2-1)\log|z|\right)>0.
                                                        \tag{C.AD.7}\label{cdp:eq:C.AD.7}
\]
Thus \eqref{cdp:eq:C.AD.6} gives an actual smooth proper family of two-tori
on all $\C^*$, by the locally uniformly bounded lattice
matrices and parallelepipeds in \ref{cdp:st:C.2.UN.55}. No elliptic quotient
of this entire middle family is required.

\par\medskip\noindent\textbf{The exact moving components at zero.}
Transport every lattice vector by $V_z$, retaining its integral
coefficients. At a Hopf vertex put
\[
 P_i(z)=\left(\frac{V_z(n_i)}{2\pi i},
                       -\frac{V_z(n_{i-1})}{2\pi i}\right),
 \qquad \Lambda_i(z)=P_i(z)^{-1}\Lambda(z).
                                                        \tag{C.AD.8}\label{cdp:eq:C.AD.8}
\]
The frame is invertible near zero. Its two distinguished
angular periods are literally $(2\pi i,0)$ and $(0,2\pi i)$;
the finite angular enlargement is retained. The finite quotient characters
are unchanged, since all their integer relations were
transported. An original contracting generator becomes a
holomorphic pair of logarithms, still with negative real
parts on a smaller disc. Exponentiation therefore gives an
actual smooth family of the marked primary or secondary
diagonal Hopf quotients. Its central member is the original
marked surface, including every finite character.

\noindent\textbf{Step 7.} At a ruled vertex use its primitive normal vector and a
transverse integral lattice lift, as in (C.MX.10). Transporting
both gives the actual moving elliptic base, degree-zero line
and projective completion. This is valid on a smaller disc
because its two base periods remain real-linearly independent.
Its two marked sections are retained.

\noindent\textbf{Step 8.} If $b_i$ is the original common affine edge constant after
the one permitted normal-product change, choose its moving
value to be $V_z(b_i)$. The common lattice and the paired
normal meridians give the required linear edge maps. Their
affine constants induce actual elliptic translations and
normal-line isomorphisms, extending over the marked curves.
The sum of the edge translations is $V_z(b)$; the sum of
their logarithmic coefficients is $dV_z(p)/(2\pi i)$.
Thus the proper moving plumbing with equation $xy=z$ has
exactly the periods \eqref{cdp:eq:C.AD.6}. Its logarithmic chart maps
identify its entire punctured family with that family:
each collar change is the stated affine map, every branch
change is an integral period, and the complete circuit is
exactly $B(z)$. The inverse is the closed-coframe inverse
already used in the local criterion. These maps and their
inverses vary holomorphically with $z$.

\par\medskip\noindent\textbf{The conjugate Hopf and ruled filling at infinity.}
Put $s=z^{-1}$ and divide the common vector frame by $z$.
Write
\[
 W_s(w)=J\overline w+s w,\qquad
              J=\begin{pmatrix}0&1\\-1&0\end{pmatrix}.
                                                        \tag{C.AD.9}\label{cdp:eq:C.AD.9}
\]
The rank-three periods become $W_s\Lambda$, and the fourth
period is
\[
                 W_s(b)-\frac{d}{2\pi i}\log s\,W_s(p).
                                                        \tag{C.AD.10}\label{cdp:eq:C.AD.10}
\]
Choose the outgoing normal at vertex $i$ to be
$n_i^\infty(s)=-W_s(n_i)$. Its incoming normal is
$-n_{i-1}^\infty(s)$ and its total normal period is
$-dW_s(p)$, exactly as required by \eqref{cdp:eq:C.AD.10}.

\noindent\textbf{Step 9.} For a Hopf vertex the corresponding frame satisfies the
following exact identities at zero:
\[
 \begin{gathered}
 P_i^\infty(0)=J\overline{P_i},\\
 (P_i^\infty(0))^{-1}W_0\Lambda
                         =\overline{\Lambda_i}.
 \end{gathered}                                        \tag{C.AD.11}\label{cdp:eq:C.AD.11}
\]
The signs follow from $\overline{2\pi i}=-2\pi i$.
Consequently the actual contracting logarithms are the
complex conjugates of the original ones, whose real parts
are still negative. The finite angular characters are
their conjugates and remain free. For small $s$, the same
holomorphic lattice construction as \eqref{cdp:eq:C.AD.8} gives the
smooth marked Hopf family through this conjugate surface.

\noindent\textbf{Step 10.} For a ruled vertex the separate frame transforms by the
same $J$ and conjugation, with its angular normal changed
by the displayed minus sign. The resulting central ruled
surface is exactly the conjugate marked projective bundle;
the actual line-bundle lattice proves this identification.
Choose the affine constant on edge $i$ to be $W_s(b_i)$.
At zero its component-coordinate expression is the complex
conjugate of the original expression, so it induces the
conjugate elliptic gluing map. Its logarithmic coefficient
is $W_s(n_i)/(2\pi i)$, and their sum is exactly the one
in \eqref{cdp:eq:C.AD.10}.

\noindent\textbf{Step 11.} The same proper normal-product plumbing, now with $xy=s$,
therefore has central support equal to the specified
complex conjugate support. Its actual punctured-family
logarithmic map identifies it with \eqref{cdp:eq:C.AD.10}. This proves
both the component and the entire-family comparisons at
infinity; complex conjugation is used only in the fixed
central coefficients of holomorphic period sections.

\noindent\textbf{Step 12.} Finally glue these two proper fillings to \eqref{cdp:eq:C.AD.6} on
their punctured discs. The exact frame change \eqref{cdp:eq:C.AD.9}
and the two logarithmic biholomorphisms give the transition
maps. The finite closed-base refinement argument of \ref{cdp:st:C.2.UN.55}
proves the resulting total space Hausdorff and proper
over $\PP^1$, hence compact. Its displayed local models
are smooth. All middle fibres are smooth two-tori.
Every marked component and double elliptic curve has
Euler number zero, so both reduced singular supports
have Euler number zero as well.

\noindent\textbf{Step 13.} The symbolic real determinant and the two conjugated
Hopf-frame identities are checked in
\path{code/root_all_diagonal_cycles.py} and its output.
\end{proof}






\begin{proposition}[The complete canonical and period reduction for all-Hopf trees]\label{cdp:st:C.2.UN.58}
\StatementRecord{cdp:st:C.2.UN.58}{proved}{Historical events: Claude Fable 5.2 (Anthropic), effort high, 16 November 2026: carried as background without fault found. GPT-6 Astra (OpenAI), ChatGPT Pro mode, effort Pro, S6\_09\_PAPERS, 17 November 2026: targeted proof check. Also checked in this writing. Pre-continuation writing events: GPT-6 Astra (OpenAI), ChatGPT Pro mode, effort Pro, S6\_10\_PAPERS, 18 November 2026: targeted proof check. Continuation metadata events: GPT-6 Astra (OpenAI), ChatGPT Pro mode, effort Pro, S6\_10\_PAPERS continuation, 19 November 2026: used and rechecked in this run; status-vocabulary reconciliation; metadata reconciliation only; no new mathematical check. \emph{Pre-continuation recorded proof scope.} The finite normal-root charts and averaged translation gauge identify every elliptic component with its actual fixed-Jacobian quasi-bundle. The exact order m*ord(theta)+m-1 excludes unmarked multiples and gives the integral leaf height equation. Normal-root collars and the determinant-one secondary matrices identify the same integral fibre class through the tree. Its leading period is zero on each elliptic component and nonzero on a nonlinear leaf, proving the precise equal-height and strict-height tests. The residue restriction also gives identity Jacobian linear holonomy around the cycle}{the hypotheses, notation and definitions stated here; \cref{cdp:st:C.2.UN.46}; \cref{cdp:st:C.2.UN.48}; \cref{cdp:st:C.2.UN.54}; \cref{cdp:st:C.2.UN.38}; \cref{cdp:st:C.2.UN.37}; \cref{cdp:eq:C.HT.4}; \cref{cdp:st:C.2.UN.50}; \cref{cdp:eq:C.HT.5}; \cref{cdp:eq:C.HT.7}; \cref{cdp:eq:C.HT.6}; \cref{cdp:eq:C.HT.9}; \cref{cdp:st:C.2.UN.49}; \cref{cdp:eq:C.HT.10}; \cref{cdp:eq:C.HT.11}}

Let $S$ be a full reduced SNC fibre of a smooth-total degeneration
of connected complex two-tori, with every irreducible component a
primary or secondary Hopf surface. If one component is a diagonal
Hopf surface of algebraic dimension zero, the whole support is
one of the diagonal cycles classified by \ref{cdp:st:C.2.UN.46} and \ref{cdp:st:C.2.UN.48}.

In every remaining case the incidence graph is the unicyclic graph
of \ref{cdp:st:C.2.UN.54}. Every cycle vertex and every nonterminal tree vertex
has algebraic dimension one. Its actual elliptic reduction is
\[
                    q_v:H_v\longrightarrow\PP^1,        \tag{C.HT.1}\label{cdp:eq:C.HT.1}
\]
with fixed Jacobian, identity integral monodromy and smooth multiple
supports. The other possible vertices are nonlinear
algebraic-dimension-zero Hopf leaves. Every multiple fibre of
\eqref{cdp:eq:C.HT.1} is a double curve of $S$.

For each edge $e$ the normal line has a finite exact order $m_e$.
It is the full fibre multiplicity on both elliptic Hopf endpoints,
and is the resonance at a nonlinear endpoint. The connected
normal-root covers of the two endpoints are identified by the
actual dual-normal map. Their Jacobians therefore identify along
every edge, and the product of their linear identifications around
the cycle is the identity. The actual elliptic cycle product can
still be a translation, whose value is retained by these maps.

Let $k_v$ be the canonical evaluation orders, and let $k_0$ be their
common minimum on the cycle. Root each attached tree at that cycle.
The heights are determined by the weighted graph:
\[
 h_v:=k_v-k_0=
 \begin{cases}
 0,&v\text{ is on the cycle},\\
 \displaystyle\sum_{e\text{ on the path from the cycle to }v}m_e,
       &v\text{ is in an attached tree}.
 \end{cases}                                           \tag{C.HT.2}\label{cdp:eq:C.HT.2}
\]
More precisely, after its invariant relative elliptic one-form is
factored out, the canonical form on a cycle vertex is a nonzero
constant multiple of $dz/z$, where the two cycle points are
$0,\infty$. On a tree vertex it is a nonzero constant multiple of
$dz$, where its parent point is $\infty$.

All regular elliptic-fibre two-tori on the algebraic-dimension-one
vertices, together with the real core two-tori on the nonlinear
leaves, represent one transported integral class $\gamma$ in the
nearby fibre. If a nonlinear leaf occurs, the single period
\[
                     P(t)=\int_{\gamma_t}\omega_t       \tag{C.HT.3}\label{cdp:eq:C.HT.3}
\]
of the relative canonical evaluation divided by $t^{k_0}$ has a
finite positive vanishing order $R$. Every nonlinear leaf has
height $R$, and every algebraic-dimension-one vertex has height
strictly smaller than $R$. These are assertions about the actual
common period. When there is no nonlinear leaf its period is
allowed to vanish identically.
\end{proposition}



\begin{proof}
\noindent\textbf{Step 1.} \par\medskip\noindent\textbf{Separation by the actual normal lines.}
The axes of an algebraic-dimension-zero diagonal primary Hopf
surface have nontorsion normal lines. For example, torsion of the
normal of $w=0$ for $(z,w)\mapsto(\alpha z,\beta w)$ would give
$\beta^n=\alpha^a$ for some integers $n>0,a$. The corresponding
monomial quotient is a nonconstant meromorphic invariant, contrary
to algebraic dimension zero. This is the exact automorphy
calculation of \ref{cdp:st:C.2.UN.38}. Pullback to the finite primary cover
proves nontorsion on a secondary quotient as well.

\noindent\textbf{Step 2.} Every curve on an algebraic-dimension-one Hopf surface is a fibre
support of \eqref{cdp:eq:C.HT.1}: a horizontal curve would meet a regular elliptic
fibre, contradicting $H^2(H,\Q)=0$ and positivity of intersection.
Its normal has finite order, as checked explicitly below. The
nonlinear Hopf axis also has finite order equal to its resonance,
by (C.SN.3)--(C.SN.4). Normal duality at an actual reduced SNC edge
preserves torsion and nontorsion. Thus a diagonal
algebraic-dimension-zero vertex can meet only another such vertex.
Connectedness then puts every component in that class, and
\ref{cdp:st:C.2.UN.54} gives precisely the already classified diagonal cycle.

\noindent\textbf{Step 3.} In the remaining case every algebraic-dimension-zero vertex is
nonlinear and has just its one elliptic curve. It is a leaf.
All other vertices have \eqref{cdp:eq:C.HT.1} by \ref{cdp:st:C.2.UN.37}. In particular every
cycle or nonterminal tree vertex has that fibration.

\par\medskip\noindent\textbf{The root charts and the invariant relative form.}
At a full fibre $mJ$ of \eqref{cdp:eq:C.HT.1}, its actual normalized $m$th root
cover is a connected smooth elliptic family, finite unramified over
the original surface tube. The fixed modulus and identity linear
monodromy give a local principal bundle with fixed elliptic fibre
$F$. It is holomorphically trivial on a disc: the exponential
sequence $0\to\Lambda_F\to\OO\to F\to0$, together with
$H^1(\Delta,\OO)=H^2(\Delta,\Lambda_F)=0$, kills its translation
torsor. Thus the original tube is
\[
       (F\times\Delta_s)/\mu_m,\qquad
       (z,s)\longmapsto(z+\tau,\zeta_m s),\qquad
       q=s^m,\quad \operatorname{ord}(\tau)=m.          \tag{C.HT.4}\label{cdp:eq:C.HT.4}
\]
Here the translation can be made constant by an actual gauge.
If its initially holomorphic expression is $\tau+A(s)$, choose a
small complex lift with $\sum_{j=0}^{m-1}A(\zeta_m^j s)=0$ and put
$b(s)=-m^{-1}\sum_{j=0}^{m-1}jA(\zeta_m^j s)$.
Then $b(\zeta_ms)-b(s)=-A(s)$, so $z\mapsto z+b(s)$ gives \eqref{cdp:eq:C.HT.4}.
The exact order follows from the free deck action on the connected
central elliptic curve.

\noindent\textbf{Step 4.} The normal of $J=F/\langle\tau\rangle$ is the associated line for
the primitive character $\zeta_m$. It has exact order $m$:
if a smaller tensor power were trivial, its pullback trivialization
would be a nowhere-zero holomorphic function on the compact $F$
with a nontrivial character, which is impossible. The connected
normal-root cover is exactly $F\to J$. Duality of the two normal
lines at an edge identifies their exact orders and their connected
root covers. Hence the number $m_e$ in the statement is well defined,
including at a nonlinear leaf via the exact secondary charts of
\ref{cdp:st:C.2.UN.50}.

\noindent\textbf{Step 5.} In every root chart write $dz$ for a fixed nonzero invariant
elliptic one-form. Transition maps between principal elliptic
charts are translations; their differentials differ by a multiple
of the base differential. Consequently
\[
                         q^*\theta\wedge dz             \tag{C.HT.5}\label{cdp:eq:C.HT.5}
\]
is a globally defined meromorphic two-form whenever $\theta$ is a
meromorphic differential on $\PP^1$. The finite translations in
\eqref{cdp:eq:C.HT.4} preserve it. This constructs the needed form rather than
postulating a global absolute elliptic one-form.

\noindent\textbf{Step 6.} Every meromorphic two-form whose divisor is vertical has precisely
the form \eqref{cdp:eq:C.HT.5}. Its ratio to one such reference form is a
meromorphic function with vertical divisor. On a regular compact
elliptic fibre away from that divisor the ratio is a nonzero
holomorphic function and is constant. It therefore descends on
the regular base. In \eqref{cdp:eq:C.HT.4} its descent is a meromorphic function
of $s$ invariant under $s\mapsto\zeta_ms$, and hence is a
meromorphic function of $s^m$. This proves descent across every
multiple fibre as well. In this identification the exact order is
\[
 \operatorname{ord}_J(q^*\theta\wedge dz)
               =m\,\operatorname{ord}_{q(J)}\theta+m-1 .
                                                        \tag{C.HT.6}\label{cdp:eq:C.HT.6}
\]

\par\medskip\noindent\textbf{The canonical equation and the leaf induction.}
The actual canonical evaluation gives on every component a
nonzero meromorphic two-form $\eta_v$ with divisor
\[
 \operatorname{div}(\eta_v)
       =\sum_{e=(v,w)}(k_w-k_v-1)J_e .                 \tag{C.HT.7}\label{cdp:eq:C.HT.7}
\]
This is adjunction with the specified vertical orders. It can also
be read in the double chart $t=xy$ by dividing the evaluation by
$t^{k_v}$ and taking the residue after division by $t$. Away from
the boundary there is no zero or pole, since the evaluation has
only its central component divisors. Equation \eqref{cdp:eq:C.HT.7} is a divisor
of the actual section, not just numerical equivalence.

\noindent\textbf{Step 7.} Apply \eqref{cdp:eq:C.HT.6} at an algebraic-dimension-one vertex. An unmarked
multiple fibre would have order zero in \eqref{cdp:eq:C.HT.7}, requiring
$m-1\equiv0\pmod m$ and hence $m=1$. There are therefore no
unmarked multiple fibres. At a marked fibre \eqref{cdp:eq:C.HT.6} says
\[
 d_{v,e}:=\frac{k_w-k_v}{m_e}
       =\operatorname{ord}_{q_v(J_e)}\theta_v+1\in\Z.
                                                        \tag{C.HT.8}\label{cdp:eq:C.HT.8}
\]
Every other point of $\PP^1$ has order zero. Taking the degree of
the divisor of this rational differential proves
\[
                    \sum_{e\ni v}d_{v,e}
                           =\operatorname{val}(v)-2.   \tag{C.HT.9}\label{cdp:eq:C.HT.9}
\]
Opposite orientations give opposite $d$ because the same exact
normal order $m_e$ occurs at the two endpoints.

\ref{cdp:st:C.2.UN.54} identifies the cycle vertices exactly as the minimum
canonical-order vertices. Their $k_v$ are one common $k_0$,
so both cycle-edge differences are zero. A terminal
algebraic-dimension-one vertex has just its parent edge, and
\eqref{cdp:eq:C.HT.9} gives $k_{\rm leaf}-k_{\rm parent}=m_e$.
For a nonlinear leaf the invariant form
$dz\wedge dw/w^{r+1}$ and the infinite order of its effective
axis line give exactly the same equality, with $m_e=r$,
as in (C.AS.12).

\noindent\textbf{Step 8.} Now work inward along any finite attached tree. Suppose all the
$c$ child edges of a vertex have outgoing $d=1$. At that vertex
the one parent term is the negative of the parent's outgoing
$d_{\rm parent}$, and \eqref{cdp:eq:C.HT.9} gives
$-d_{\rm parent}+c=c-1$. Thus $d_{\rm parent}=1$.
Every nonterminal tree vertex has algebraic dimension one, so this
induction reaches the cycle. It proves \eqref{cdp:eq:C.HT.2} on every edge.

\noindent\textbf{Step 9.} At a cycle vertex the two cycle terms have $d=0$, giving simple
poles of $\theta_v$ there; every child term has $d=1$ and order
zero. There are no other zeros or poles. A coordinate taking the
cycle points to $0,\infty$ therefore gives $\theta_v=c\,dz/z$.
At a tree vertex the parent term has $d=-1$ and order $-2$;
all child terms have order zero. A coordinate with its parent
point at infinity gives $\theta_v=c\,dz$. In each case $c\ne0$.
This proves the more precise differential assertion.

\par\medskip\noindent\textbf{The actual Jacobian comparison around the cycle.}
The same root charts show $R^1q_{v*}\OO_{H_v}=\OO_{\PP^1}$.
On $F\times\Delta_s$ the cohomology frame is the fixed
$H^1(F,\OO_F)$; finite deck translations act trivially, so its
invariant functions are exactly functions of $s^m$. On overlaps
the identity linear monodromy and translations give the same
frame. Leray on $\PP^1$ therefore identifies
$H^1(H_v,\OO_{H_v})$ with this one-dimensional constant fibre
cohomology.

\noindent\textbf{Step 10.} Restriction to $J_e$ followed by pullback along the actual finite
cover $F\to J_e$ agrees with evaluation on the regular elliptic
fibre. This follows from the commuting pullback diagram of the
surface tube and \eqref{cdp:eq:C.HT.4}, where both maps are restriction of the
same constant cohomology frame. Thus the scalar cycle product
in (C.HR.2) is exactly the action on $H^1(F,\OO_F)$ of the product
of the lifted fibre isomorphisms. It is one by \ref{cdp:st:C.2.UN.54}.
An elliptic automorphism whose action on $H^1(\OO)$ is the identity
has linear part one: after choosing an origin, its lift is
$z\mapsto az+b$, and the nonzero invariant one-form detects $a$
(or its conjugate on $H^1(\OO)$). Hence that product is a
translation. Different choices of a lift through a finite
normal-root cover alter it only by deck translations. The
Jacobian product is therefore the identity with all maps retained.

\par\medskip\noindent\textbf{One integral core class and its nonzero nonlinear periods.}
We describe the cycle comparison before integrating. Near an
edge the two surface tubes are the disk bundles of the actual
torsion normal lines in \eqref{cdp:eq:C.HT.4}. Their dual-normal isomorphism
identifies their connected root covers. On that cover the
normal-product collar is
\[
                  xy=t,\qquad q_v=x^{m_e},\quad
                  q_w=y^{m_e}.                         \tag{C.HT.10}\label{cdp:eq:C.HT.10}
\]
Multiplying one base coordinate by a nonzero constant makes these
equalities exact: the two trivializations of the $m_e$th normal
power differ by a holomorphic unit on a compact elliptic curve,
and hence by a constant. The regular fibre $x^{m_e}=c$ on the
first side is sent to $y^{m_e}=t^{m_e}/c$ on the other. Both are
the same connected elliptic cover of $J_e$, with its complex
orientation. Varying the nonzero normal level supplies the
oriented isotopy in the neck.

\noindent\textbf{Step 11.} For an arbitrary given smoothing, the product-preserving smooth
tubular comparison used in \ref{cdp:st:C.2.UN.49} carries these normal collars
to the actual $xy=t$ collars. Compact cores are compared by
Ehresmann transport. Consequently \eqref{cdp:eq:C.HT.10} is an equality of
the transported integral two-cycle classes in the actual nearby
fibre, without a holomorphic linearization of its neighbourhood.
Regular fibres on one component are isotopic through its regular
base, so every algebraic-dimension-one vertex gives the same
class $\gamma$. At a nonlinear leaf the exact unimodular
boundary maps (C.SN.6)--(C.SN.8), together with the displayed
smooth-pair deformation, identify its real core with this same
oriented class. They retain the full secondary order.

\noindent\textbf{Step 12.} Write $\omega_t$ for the evaluation divided by $t^{k_0}$.
On a compact regular fibre of $H_v$ its leading coefficient is
$t^{h_v}\eta_v$. The restriction of the meromorphic two-form
$\eta_v$ to this holomorphic curve is zero. Therefore
\[
                     P(t)=O(t^{h_v+1})
       \quad\text{at every algebraic-dimension-one vertex}.
                                                        \tag{C.HT.11}\label{cdp:eq:C.HT.11}
\]
These are actual periods of a closed relative holomorphic form
over the same compact class. They are holomorphic by
differentiation and Stokes' theorem; the local compact-core
transport gives the displayed extension and leading coefficient.

\noindent\textbf{Step 13.} At a nonlinear leaf
$H_r(\beta)/\mu_n$ its unique-axis form is
$\eta=c\,dz\wedge dw/w^{r+1}$, $c\ne0$. In
$U=z/w^r$, $V=\log w$, the two primitive real core periods are
$(0,2\pi i/n)$ and $(\beta^{-r},\log\beta)$. Hence
\[
 P(t)=t^{h_{\rm leaf}}
       \left(\frac{\pm2\pi i c\beta^{-r}}{n}+O(t)\right),
                    \qquad c\ne0.                     \tag{C.HT.12}\label{cdp:eq:C.HT.12}
\]
The coefficient is nonzero. One such leaf fixes the finite
vanishing order $R=h_{\rm leaf}>0$ of the single period.
Every other nonlinear leaf has that same height, while
\eqref{cdp:eq:C.HT.11} gives $h_v<R$ at every elliptic Hopf vertex.
No nonvanishing conclusion is imposed when there is no nonlinear
leaf.

\noindent\textbf{Step 14.} The exact root congruence, tree equations, height sums, collar
monomials and common-period inequalities are checked in
\path{code/E_hopf_tree_heights.py} and its saved output. The
argument proves these comparisons for every actual support in
the stated branch. It does not assert that these necessary
equations alone smooth arbitrary abstract attached-tree data.
\end{proof}






\begin{proposition}[All multiplicities and all diagonal Hopf and ruled vertices]\label{cdp:st:C.2.UN.59}
\StatementRecord{cdp:st:C.2.UN.59}{proved}{Historical events: Claude Fable 5.2 (Anthropic), effort high, 16 November 2026: carried as background without fault found. GPT-6 Astra (OpenAI), ChatGPT Pro mode, effort Pro, S6\_09\_PAPERS, 17 November 2026: targeted proof check. Also checked in this writing. Pre-continuation writing events: GPT-6 Astra (OpenAI), ChatGPT Pro mode, effort Pro, S6\_10\_PAPERS, 18 November 2026: targeted proof check. Continuation metadata events: GPT-6 Astra (OpenAI), ChatGPT Pro mode, effort Pro, S6\_10\_PAPERS continuation, 19 November 2026: used and rechecked in this run; status-vocabulary reconciliation; metadata reconciliation only; no new mathematical check. \emph{Pre-continuation recorded proof scope.} The actual finite root action has linear part diag(rho,1) with rho of order 1,2,3,4 or 6. Choosing k with rho=zeta\_N\textasciicircum{}k makes the explicit real-linear deformation intertwine the base rotation. Every affine edge and order equation is preserved in the target frame; the full four-period determinant remains positive. Eigenfunction averaging gives actual product-preserving local coordinates and the smooth invariant ring with t=x\textasciicircum{}a*y\textasciicircum{}b. The reciprocal chart gives exactly the conjugate action and marked cycle. Both compact equivariant fillings, the free punctured action and all four explicit finite-order examples were re-derived}{the hypotheses, notation and definitions stated here; \cref{cdp:st:C.2.UN.40}; \cref{cdp:st:C.2.UN.43}; \cref{cdp:st:C.2.UN.52}; \cref{cdp:st:C.2.UN.56}; \cref{cdp:st:C.2.UN.51}; \cref{cdp:st:C.2.UN.55}; \cref{cdp:eq:C.EW.3}; \cref{cdp:eq:C.EW.6}; \cref{cdp:eq:C.EW.9}; \cref{cdp:eq:C.EW.10}; \cref{cdp:eq:C.EW.13}; \cref{cdp:eq:C.EW.15}; \cref{cdp:eq:C.EW.14}; \cref{cdp:st:C.2.UN.57}; \cref{cdp:eq:C.EW.12}; \cref{cdp:eq:C.EW.1}; \cref{cdp:eq:C.EW.19}; \cref{cdp:eq:C.EW.20}; \cref{cdp:eq:C.EW.17}; \cref{cdp:eq:C.EW.18}}

Specify a simple-NC cycle of smooth embedded marked surfaces, with
edges between distinct components. Each vertex is a diagonal primary
or secondary Hopf surface with its two axes, or
$\PP(\OO_E\oplus L)$ with $E$ elliptic and
$L\in\operatorname{Pic}^0(E)$, marked by its two sections.
Specify their elliptic gluing maps, positive multiplicities $m_i$,
and the actual compatible weighted normal lines, sections and
boundary maps of (C.WH.1)--(C.WH.2).

Put $N=\operatorname{lcm}_i m_i$. Form the actual normalized
$N$th-root component covers and the compressed necks $uv=s^{h_i}$.
The root graph is connected and its rank-three integral
complement-lattice product is the identity precisely when the
specified weighted normal plumbing has connected two-torus fibres.
Every actual smooth-total torus degeneration with the specified
marked support yields these data and passes this test.

Every weighted support passing the test occurs over $\PP^1$ in a
smooth compact complex threefold. The construction retains all
marked components, elliptic gluing maps and multiplicities. Its
only other singular fibre is the complex conjugate weighted
support. All reduced fibre supports have Euler number zero.
Normal-product constants are chosen in this compact realization.
There is no algebraic-dimension restriction on the Hopf vertices,
and an all-ruled cycle is included.

The complete component, axis and inertia comparisons in
(C.WH.0), and the ruled-cover identities (C.WM.2)--(C.WM.7),
hold in this scope. In common root-complement coordinates the
finite deck linear part fixes the primitive total normal period.
Its other eigenvalue is a root of unity of order $1,2,3,4$ or $6$.
The compact construction retains this linear part and every
affine deck constant.
\end{proposition}



\begin{proof}
\noindent\textbf{Step 1.} \par\medskip\noindent\textbf{Actual root maps and the local criterion.}
The weighted line on a diagonal Hopf vertex has a flat tensor
power, since both axis divisors have the explicit constant
automorphy factors of \ref{cdp:st:C.2.UN.40} and \ref{cdp:st:C.2.UN.43}. Dividing the transferred
connection by its tensor exponent gives the actual flat line
used in \ref{cdp:st:C.2.UN.52}. Its root cover on the axis complement is the
finite-index logarithmic lattice cover constructed there.
Coordinate powers extend it across both axes, and the actual
contracting generator and free finite angular characters identify
its finite normalization with a smooth diagonal Hopf surface.
This comparison does not require nonresonant multipliers.

\noindent\textbf{Step 2.} On a ruled vertex the proof (C.WM.4)--(C.WM.7) applies without
any Hopf assumption. It gives the actual finite normalization
$\PP(\OO_{E'}\oplus L')\to\PP(\OO_E\oplus L)$, with
an \'{e}tale base map $\pi$, fibre power $q$, and
$(L')^q\simeq\pi^*L$. In particular
$m\mid a+b$ and $q=m/\gcd(m,a)=m/\gcd(m,b)$.
The weighted normal line is pulled back on each entire section
chart by its transferred flat connection; its nonzero global
degree on a ruling fibre is retained as in \ref{cdp:st:C.2.UN.56}.

\noindent\textbf{Step 3.} At an original double curve put
$a=gA$, $b=gB$, $\gcd(A,B)=1$ and
$h=N/(gAB)$. The finite normalization is still exactly
\[
             uv=s^h,\qquad x=u^B,\quad y=v^A,
             \qquad t=s^N=x^a y^b.
\tag{C.EW.1}\label{cdp:eq:C.EW.1}
\]
The actual root component and axis covers, their connected
degrees, and all stabilizers therefore have precisely the
maps and orders proved in \ref{cdp:st:C.2.UN.52} and \ref{cdp:st:C.2.UN.56}. Each strict component
has two connected marked elliptic curves, so the compressed
root graph is a union of cycles. Resolving a neck inserts the
specified degree-zero ruled bridges and changes no complement
lattice product.

\noindent\textbf{Step 4.} Connectedness of this graph is equivalent to connectedness of
the root family, by the proper Stein factor and its reduced
central fibre. On a connected graph the actual core-and-neck
calculation gives $\pi_1(F_s)=\Lambda\rtimes_M\Z$.
Thus a torus forces $M=I$. Conversely $M=I$ glues the closed
logarithmic coframe; its complete commuting fields identify
the nearby compact fibre with $\C^2$ modulo its rank-four
period lattice. For an arbitrary actual smoothing the
canonical divisor sections and their small smooth unit roots
give the same weighted tubular and collar maps, exactly as
in \ref{cdp:st:C.2.UN.52}. This proves the asserted local criterion for the
actual maps, as well as for the explicit plumbing.

\par\medskip\noindent\textbf{The exact finite linear and affine data.}
Assume that the test passes and identify all open root
complements with $G=\C^2/\Lambda$ by the actual linear
edge maps. Let $n_j$ be the outgoing normal meridian at
vertex $j$, and let $h_j$ be its outgoing compressed
neck exponent. The incoming meridian is $-n_{j-1}$.
The root generator lies over the identity on each original
marked normalization and preserves its two marked branches.
It consequently induces an orientation-preserving rotation
$\sigma$ of the connected root cycle. The common number
$r$ of strict covers and every inertia order are therefore
exactly those of (C.WH.0).

\noindent\textbf{Step 5.} The generator is affine on the common complement. For clarity,
an isomorphism of diagonal Hopf surfaces preserving both axes
lifts to a biholomorphism of $\C^2\setminus0$, extends with
its inverse across zero, and has coordinates $zU(z,w)$ and
$wV(z,w)$, with $U,V$ units. Conjugating a contracting deck
generator and evaluating at zero makes the corresponding
multiplier of each unit equal to one. Iteration of the
contraction then makes each unit constant. Thus this lift
is diagonal even for resonant multipliers. For a ruled
vertex the actual base isomorphism and marked line-bundle
isomorphism have affine logarithmic lifts by the constant
factors of automorphy, as in \ref{cdp:st:C.2.UN.51}. Compatibility with the
actual neck maps makes the linear part a single map $R$.
Write the generator and the edge maps as
\[
 \begin{aligned}
 g(j,v,s)&=(\sigma j,Rv+\tau_j,\zeta_Ns),\\
 e_j(v,s)&=v+b_j+l_j\log s,\\
 l_j&=-\frac{h_j n_j}{2\pi i}.
 \end{aligned}
\tag{C.EW.2}\label{cdp:eq:C.EW.2}
\]
Here $R\Lambda=\Lambda$, $R^N=I$,
$Rn_j=n_{\sigma j}$ and $h_{\sigma j}=h_j$.
Actual composition of these maps, with
$\log\zeta_N=2\pi i/N$, gives lattice elements
$\lambda_j,\mu_j\in\Lambda$ satisfying
\[
 \begin{aligned}
 \tau_{j+1}+Rb_j-\tau_j-b_{\sigma j}
             +\frac{h_j n_{\sigma j}}N&=\lambda_j,\\
 \sum_{a=0}^{N-1}R^{N-1-a}\tau_{\sigma^a j}&=\mu_j.
 \end{aligned}
\tag{C.EW.3}\label{cdp:eq:C.EW.3}
\]
These are equalities in the actual logarithmic vector space.

\noindent\textbf{Step 6.} Put
\[
             b=\sum_jb_j,\qquad
             u=-\sum_jh_jn_j=d p\in\Lambda,
\tag{C.EW.4}\label{cdp:eq:C.EW.4}
\]
where $p$ is primitive and $d>0$. The contraction-height
calculation (C.AD.3) applies at every diagonal Hopf vertex;
at a ruled vertex the consecutive outgoing normals agree.
Thus for $W=\operatorname{span}_{\R}\Lambda$ and
$\ker\ell=W$, all numbers $\ell(i n_j)$ have one
nonzero sign. In the all-ruled case the same assertion
is the actual elliptic quotient calculation of \ref{cdp:st:C.2.UN.55}.
Consequently $\ell(iu)\ne0$. Rotation and the displayed
normal maps give $Ru=u$, hence $Rp=p$.

\noindent\textbf{Step 7.} The complex line $W\cap iW$ is $R$-invariant and is
complementary to $\C p$. Choose the complex coordinates
given by these two lines, with $p=(0,1)$. Then
\[
 \begin{gathered}
 W=\C\times\R,\qquad R=\operatorname{diag}(\rho,1),\\
 \Lambda=\langle w_1=(a_1,c_1),\ w_2=(a_2,c_2),\ p\rangle_{\Z},\\
 c_1,c_2\in\R,\qquad
 C=\operatorname{Im}(\overline{a_1}a_2)>0.
 \end{gathered}
\tag{C.EW.5}\label{cdp:eq:C.EW.5}
\]
The projection of $\Lambda$ to the first coordinate is
the actual elliptic lattice. Since $R$ preserves it,
$\rho$ has order $1,2,3,4$ or $6$: its integral real
two-dimensional matrix has determinant one, trace
$2\operatorname{Re}\rho\in\Z$, and finite order.
Also $\rho^N=1$. The stabilizer of a root vertex acts
as the deck translations of its finite open-group cover,
so in fact $R^r=I$.

\noindent\textbf{Step 8.} Summing \eqref{cdp:eq:C.EW.3} retains the full affine descent relation
\[
                         (R-I)b-u/N\in\Lambda.
\tag{C.EW.6}\label{cdp:eq:C.EW.6}
\]
When $R=I$ this recovers $u/N\in\Lambda$ from \ref{cdp:st:C.2.UN.52}.
Equation \eqref{cdp:eq:C.EW.6}, with its possibly nontrivial $R$, is
the relation used below.

\noindent\textbf{Step 9.} Adding $c n_j$ to $b_j$ on a whole $\sigma$-orbit,
with the same $c\in\C$, changes only the normal-product
constants. It preserves both equations \eqref{cdp:eq:C.EW.3}, since
$Rn_j=n_{\sigma j}$. The second coordinates of all
$n_j$ are real and have the same nonzero sign, by the
height calculation. The sum on any such orbit therefore
has nonzero second coordinate. Choose $c$ on one orbit
so that
\[
                              b=(\sigma_0,i).
\tag{C.EW.7}\label{cdp:eq:C.EW.7}
\]
Every original elliptic gluing map and every finite
action map on the central support is retained.
The orbit-normal sum is also $R$-fixed; when
$\rho\ne1$ its first coordinate is zero. Thus the
entire common horizontal affine constant is retained
in that case as well.

\par\medskip\noindent\textbf{A finite-equivariant period family on the punctured sphere.}
Choose $k\in\{1,\ldots,N\}$ with
$\rho=\zeta_N^k$, using $k=N$ for $\rho=1$.
Let
\[
 J=\begin{pmatrix}0&1\\-1&0\end{pmatrix},\qquad
 V_s(w)=w+s^kJ\overline w.
\tag{C.EW.8}\label{cdp:eq:C.EW.8}
\]
For each fixed vector $w$ this is holomorphic in $s$.
The exact relation $RJ\overline R^{-1}=\rho J$
gives
\[
                           R V_s=V_{\zeta_Ns}R.
\tag{C.EW.9}\label{cdp:eq:C.EW.9}
\]
Transport all integral vectors by $V_s$, and define
\[
 \begin{gathered}
 \Lambda_s=V_s\Lambda,\qquad
 B_s=V_s b+\frac{d\log s}{2\pi i}V_s p,\\
                         T_s=\C^2/(\Lambda_s+\Z B_s).
 \end{gathered}
\tag{C.EW.10}\label{cdp:eq:C.EW.10}
\]
A logarithm change adds $dV_s p$, an actual integral
period. The determinant calculation of \ref{cdp:st:C.2.UN.55} with
$z=s^k$ and logarithmic coefficient $d\log s/(2\pi i)$
gives, for $r_s=|s|>0$,
\[
 C(1+r_s^{2k})
 \left(1+r_s^{2k}
       +\frac d{2\pi}(r_s^{2k}-1)\log r_s\right)>0.
\tag{C.EW.11}\label{cdp:eq:C.EW.11}
\]
Thus these are actual compact two-tori for every
$s\in\C^*$, and form a smooth proper holomorphic
family there. Properness follows from the uniformly
bounded period matrices and inverses on compact
base charts, exactly as in \ref{cdp:st:C.2.UN.55}.

\noindent\textbf{Step 10.} The transported finite action and edge data are explicitly
\[
 \begin{aligned}
 b_j(s)&=V_s b_j,& n_j(s)&=V_s n_j,\\
 l_j(s)&=-h_jn_j(s)/(2\pi i),&
 \tau_j(s)&=V_{\zeta_Ns}\tau_j.
 \end{aligned}
\tag{C.EW.12}\label{cdp:eq:C.EW.12}
\]
Notice that the translation is expressed in the
\emph{target} frame. Applying \eqref{cdp:eq:C.EW.9} to the first
equation of \eqref{cdp:eq:C.EW.3} gives exactly
\[
 \begin{aligned}
 &\tau_{j+1}(s)+Rb_j(s)-\tau_j(s)-b_{\sigma j}(\zeta_Ns)\\
 &\qquad-l_{\sigma j}(\zeta_Ns)\,2\pi i/N
       =V_{\zeta_Ns}\lambda_j.
 \end{aligned}
\tag{C.EW.13}\label{cdp:eq:C.EW.13}
\]
The logarithm correction transports because
$h_jn_{\sigma j}/N$ is a real rational multiple
of an integral vector. For the order equation,
\[
 R^{N-1-a}V_{\zeta_N^{a+1}s}
                         =V_sR^{N-1-a},
\]
so the $N$th iterate translates by $V_s\mu_j$.
This proves the full finite action, including all
overlap and order equalities. Likewise
\[
                    R B_s-B_{\zeta_Ns}
                   =V_{\zeta_Ns}\bigl((R-I)b-u/N\bigr)
                   \in\Lambda_{\zeta_Ns}.
\tag{C.EW.14}\label{cdp:eq:C.EW.14}
\]

\noindent\textbf{Step 11.} Here is an explicit global comparison with \eqref{cdp:eq:C.EW.10}.
On a logarithm chart put
$\theta_j(s)=\sum_{a<j}(b_a(s)+l_a(s)\log s)$.
The common coordinate of the $j$th component chart is
$v_j-\theta_j(s)$. The generator on the period quotient
therefore has the affine translation
\[
 R\theta_j(s)+\tau_j(s)-\theta_{\sigma j}(\zeta_Ns).
\tag{C.EW.15}\label{cdp:eq:C.EW.15}
\]
Changing $j$ changes this by an integral period, by
\eqref{cdp:eq:C.EW.13}; passing the last edge adds exactly $B_s$.
Changing the logarithm changes each partial sum by
an integral vector. Thus \eqref{cdp:eq:C.EW.15} defines a single
holomorphic finite action on the entire punctured
period family. Its order is the order already proved
in \eqref{cdp:eq:C.EW.13}, and its linear and period maps are
\eqref{cdp:eq:C.EW.9} and \eqref{cdp:eq:C.EW.14}.

\par\medskip\noindent\textbf{The actual finite filling at zero.}
Near zero construct the moving strict components from
their transported lattices. At a Hopf vertex use
\[
 P_j(s)=\left(\frac{V_s n_j}{2\pi i},
                        -\frac{V_s n_{j-1}}{2\pi i}\right).
\tag{C.EW.16}\label{cdp:eq:C.EW.16}
\]
It is invertible near zero, retains the distinguished
angular periods and their finite angular enlargement,
and transports the original contracting logarithms
holomorphically. Their real parts remain negative.
It therefore defines the actual marked Hopf family.
At a ruled vertex use its primitive normal and an
independent projected elliptic lattice lift, as in
\ref{cdp:st:C.2.UN.56} and \ref{cdp:st:C.2.UN.57}. This gives the actual moving
degree-zero line and its projective completion.

\noindent\textbf{Step 12.} Use the edge constants \eqref{cdp:eq:C.EW.12} and keep the literal
compressed hypersurfaces $uv=s^{h_j}$. They give the
proper moving plumbing, whose exact periods are
\eqref{cdp:eq:C.EW.10}. The component logarithmic maps and their
closed-coframe inverses give a holomorphic
biholomorphism of its entire punctured family with
that period family. Formula \eqref{cdp:eq:C.EW.15} makes this
biholomorphism equivariant.

\noindent\textbf{Step 13.} The finite maps extend across all marked curves.
At a Hopf vertex
$RP_j(s)=P_{\sigma j}(\zeta_Ns)$, so the logarithmic
action in these component coordinates is a translation;
its exponential is the actual diagonal map on the
marked Hopf family. At a ruled vertex the same
normal-period map and its elliptic base map give the
actual line-bundle map and homogeneous projective
extension from \ref{cdp:st:C.2.UN.56}. The overlap equalities
\eqref{cdp:eq:C.EW.13} extend them across the compressed necks.

\noindent\textbf{Step 14.} The smoothness of the finite quotient is the same
local ring calculation as (C.WH.12), which uses
the actual inertia characters rather than $R=I$.
At \eqref{cdp:eq:C.EW.1} the pointwise axis stabilizer has order
$hAB$. Its order-$h$ subgroup fixes $u,v$ and rotates
$s$; the remaining independent characters have
orders $B,A$. These finite characters retain their
specified central values throughout the moving
families. We give the actual character coordinates.
The stabilizer preserves each adjacent component,
so it sends each of $u,v$ to itself times a unit.
Reynolds-average these functions to their specified
central eigencharacters. The results are $u'=uU$
and $v'=vV$, with $U,V$ units. Average the coordinate
along the pointwise fixed axis to an invariant
coordinate; its derivative there is unchanged.
The character identity
$\chi_u\chi_v=\chi_s^h$ and the equation
$u'v'=s^hUV$ show that $A=UV$ is invariant.
Replace $u'$ by $u'/A$. The product is then exactly
$s^h$, both eigencharacters are retained, and the
actual base coordinate $s$ has not changed. The
coordinate change has invertible differential at
the axis and is an analytic local isomorphism.
In these exact coordinates the ring is
\[
 \bigl(\C\{u,v,s\}/(uv-s^h)\bigr)^{J_{\rm axis}}
                         =\C\{u^B,v^A\}
                         =\C\{x,y\},
 \qquad t=x^a y^b.
\tag{C.EW.17}\label{cdp:eq:C.EW.17}
\]
The residual axis cover is \'{e}tale. At a general
component point the inertia is the normal reflection
of order $N/m_i$. Thus the quotient is smooth at
every point and has exactly the original full
multiplicities. At zero the covers, central maps and
finite action are the original ones; the quotient
support is therefore the specified marked surface.

\par\medskip\noindent\textbf{The conjugate finite filling at infinity.}
Put $w=s^{-1}$ and divide the vector frame by $s^k$.
For a fixed vector $c$ write
\[
 W_w(c)=J\overline c+w^k c.
\tag{C.EW.18}\label{cdp:eq:C.EW.18}
\]
The periods in this frame are $W_w\Lambda$ and
\[
                    W_wb-\frac{d\log w}{2\pi i}W_wp.
\tag{C.EW.19}\label{cdp:eq:C.EW.19}
\]
Choose outgoing normals $-W_wn_j$, the same neck
exponents $h_j$, and affine edge constants $W_wb_j$.
Their logarithmic coefficients are
$h_jW_wn_j/(2\pi i)$, so their complete circuit is
exactly \eqref{cdp:eq:C.EW.19}.

\noindent\textbf{Step 15.} The actual transformed deck generator is
\[
 \begin{aligned}
 (j,v,w)&\longmapsto
     (\sigma j,R_\infty v+\tau_{j,\infty}(w),\zeta_N^{-1}w),\\
 R_\infty&=\rho^{-1}R=J\overline R J^{-1},\\
 \tau_{j,\infty}(w)
     &=J\overline{\tau_j}+\rho^{-1}w^k\tau_j
       =W_{\zeta_N^{-1}w}\tau_j.
 \end{aligned}
\tag{C.EW.20}\label{cdp:eq:C.EW.20}
\]
Indeed these formulas follow by multiplying the image
vector $Rv+\tau_j(s)$ by $(\zeta_Ns)^{-k}$.
They retain all finite equations, since they are the
actual frame change of the action already constructed.

\noindent\textbf{Step 16.} At $w=0$ the Hopf frame is exactly
$J\overline{P_j(0)}$, as in (C.AD.11).
The angular normals have the displayed minus sign
because conjugation changes $2\pi i$ to $-2\pi i$.
Thus the contracting logarithms and every finite
angular character are their actual complex
conjugates. The separate ruled frame gives the
conjugate line bundle and marked sections by the
same logarithmic calculation. Each affine edge
constant, in its component coordinates, is the
conjugate of the original one. Equation \eqref{cdp:eq:C.EW.20}
is precisely the conjugate central action under
the anti-linear comparison $v\mapsto J\overline v$,
with the conjugate base rotation.

\noindent\textbf{Step 17.} The resulting proper compressed filling at infinity
is therefore the complex conjugate root plumbing
with its actual finite action. Its punctured-family
map has exactly \eqref{cdp:eq:C.EW.19} and \eqref{cdp:eq:C.EW.20}.
Its quotient rings are the conjugates of
\eqref{cdp:eq:C.EW.17}, so the quotient is smooth and its full
central fibre has the conjugate support with the
same multiplicities.

\noindent\textbf{Step 18.} Finally glue the two finite-equivariant proper
fillings to the entire period family \eqref{cdp:eq:C.EW.10},
using their explicit logarithmic biholomorphisms
and the frame change \eqref{cdp:eq:C.EW.18}. The closed-base
refinement argument of \ref{cdp:st:C.2.UN.55} proves this normal
root total space Hausdorff and proper over
$\PP^1_s$. Take the actual finite quotient by
the constructed action. The base quotient is
$\PP^1_t$ with $t=s^N$. Its total space is smooth
by \eqref{cdp:eq:C.EW.17} at the two special values and by
freeness of the base action elsewhere. It is
compact and proper. A general quotient fibre is
identified with any one of its $N$ root fibres,
so is the specified connected two-torus.
Every component and every double curve of either
special support has Euler number zero. Additivity
therefore proves the final Euler assertion.

\par\medskip\noindent\textbf{Explicit nontrivial finite linear actions.}
For $N=2,3,4,6$, choose an elliptic curve $E$ with
an origin-preserving automorphism $\rho$ of exact
order $N$. Start with a root cycle of $2N$ copies
of $E\times\PP^1$, glued by the identity on their
elliptic sections. In the coordinates above take
\[
 \begin{gathered}
 \Lambda=\Lambda_E\oplus\Z p,\quad
 n_j=-p,\quad h_j=1,\quad
 b_j=ip/(2N),\quad \tau_j=j p/N,\\
 \sigma(j)=j+2\pmod{2N},\qquad
 R=\operatorname{diag}(\rho,1).
 \end{gathered}
\tag{C.EW.21}\label{cdp:eq:C.EW.21}
\]
The first equation \eqref{cdp:eq:C.EW.3} is zero except at the
cyclic endpoint, where it is $-2p$.
The $N$th translation sum is $(N-1)p$ on even
vertices and $Np$ on odd vertices. These are
actual integral periods, so this is the specified
finite action. It is free on the central graph.
The quotient has two components $E\times\PP^1$,
both with full multiplicity $N$. One elliptic
gluing is the identity; the other is $\rho^{-1}$,
as follows by identifying vertex $2$ with vertex
$0$ using the inverse generator. The root cycle
has identity product, while this latter gluing
retains the nontrivial finite linear part.
The compact construction above realizes these
exact weighted supports and their conjugates.

\noindent\textbf{Step 19.} The equivariance, target-frame affine equations,
infinity identities, positive determinant and
all four finite-action examples are checked
exactly in \path{code/D_all_weighted_cycles.py}
and its saved output.
\end{proof}






\begin{proposition}[Exact pruning and compact reinsertion classify every reduced all-Hopf support]\label{cdp:st:C.2.UN.60}
\StatementRecord{cdp:st:C.2.UN.60}{proved}{Historical events: Claude Fable 5.2 (Anthropic), effort high, 16 November 2026: carried as background without fault found. GPT-6 Astra (OpenAI), ChatGPT Pro mode, effort Pro, S6\_09\_PAPERS, 17 November 2026: targeted proof check. Also checked in this writing. Pre-continuation writing events: GPT-6 Astra (OpenAI), ChatGPT Pro mode, effort Pro, S6\_10\_PAPERS, 18 November 2026: targeted proof check. Continuation metadata events: GPT-6 Astra (OpenAI), ChatGPT Pro mode, effort Pro, S6\_10\_PAPERS continuation, 19 November 2026: used and rechecked in this run; status-vocabulary reconciliation; metadata reconciliation only; no new mathematical check. \emph{Pre-continuation recorded proof scope.} The root-line exponential sequence classifies the complete one-point and two-point elliptic torsors, retaining full secondary integers and primitive bases. Their explicit inverse product charts perform finite leaf pruning without changing exterior analytic transitions. The actual boundary homomorphisms preserve the Van Kampen pushouts and force identity core monodromy. The compact core's primitive fibre period has exact order R. Each weighted line insertion or ordinary cap divides its displacement by the recorded height, restoring all component torsors and gluings. Exponential contraction dominates the polynomial line-norm distortion; the exact logarithmic cluster bound proves Hausdorff proper replacement at every stage}{the hypotheses, notation and definitions stated here; \cref{cdp:st:C.2.UN.46}; \cref{cdp:st:C.2.UN.50}; \cref{cdp:eq:C.PR.2}; \cref{cdp:eq:C.PR.4}; \cref{cdp:st:C.2.UN.43}; \cref{cdp:eq:C.PR.6}; \cref{cdp:eq:C.PR.3}; \cref{cdp:st:C.2.UN.49}; \cref{cdp:st:C.2.UN.58}; \cref{cdp:eq:C.PR.5}; \cref{cdp:st:C.2.UN.57}; \cref{cdp:st:C.2.UN.54}; \cref{cdp:eq:C.PR.9}; \cref{cdp:st:C.2.UN.48}; \cref{cdp:eq:C.PR.10}; \cref{cdp:eq:C.PR.12}; \cref{cdp:eq:C.PR.13}; \cref{cdp:eq:C.PR.14}; \cref{cdp:eq:C.PR.16}; \cref{cdp:st:C.2.UN.42}}

Let $S$ be a specified compact connected reduced SNC surface whose
smooth components are primary or secondary Hopf surfaces. Its
elliptic gluing maps are part of the specified surface. The following
finite geometric test is necessary and sufficient for $S$ to occur
as the full central fibre of a proper degeneration of connected
complex two-tori with smooth total space.

\begin{enumerate}
\item The actual normal lines on every double curve are dual, and
the incidence graph has exactly one independent cycle.
\item If a diagonal algebraic-dimension-zero component occurs, the
whole support is its diagonal cycle and its integral lattice product
$M$ is the identity, as in \ref{cdp:st:C.2.UN.46}.
\item In the remaining case every algebraic-dimension-zero vertex
is a nonlinear leaf. Every other vertex has its actual elliptic
reduction to $\PP^1$, and every multiple fibre of that reduction
is one of its marked double curves.
\item Let $m_e$ be the exact order of either normal line at an edge,
and let $h_v$ be zero on the cycle and the sum of the $m_e$ along
the path from the cycle to $v$. There is an integer $R>0$ such that
every nonlinear leaf has $h_v=R$ and every other vertex has $h_v<R$.
If there is no nonlinear leaf, any $R>\max h_v$ passes this test.
\item Perform the exact leaf-pruning operations specified below,
including the full integer secondary framing and all torsor maps.
The resulting marked Hopf/ruled core cycle has integral
complement-lattice product $M=I$.
\end{enumerate}

In the second case the tests referring to elliptic tree vertices are
unneeded. In the remaining case the pruning operations always exist
under the preceding geometric requirements. Each terminal step is
one of two explicit operations: the inverse logarithmic chart of
\ref{cdp:st:C.2.UN.50} for a primary or secondary Hopf pair; or removal of the
product pair $(F\times\PP^1,F\times\{\infty\})$. The latter leaves
the parent surface unchanged. After its children have been removed,
every vertex has exactly one of these two forms. The pruned core
and its test are independent, up to the stated actual marked
isomorphisms, of the order of pruning disjoint branches.

Every support passing the test has a realization over $\PP^1$
with smooth compact total threefold. It retains all the original
components, their analytic torsors and every specified elliptic
gluing map. Its other singular support can be taken to be the
complex conjugate of its pruned core cycle. All reduced fibre
supports have Euler characteristic zero.
\end{proposition}



\begin{proof}
\noindent\textbf{Step 1.} \par\medskip\noindent\textbf{The exact one-point torsor and its marked compact surface.}
We first prove the classification used at a terminal pruning step.
Let $q:T\to\PP^1$ be a fixed-Jacobian, identity-monodromy elliptic
quasi-bundle, with just one possible multiple fibre $mJ$, marked
over infinity. Its actual normal-root charts make it a principal
$F$-torsor on
\[
                 \mathcal B=\PP^1[(\infty,m)].          \tag{C.PR.1}\label{cdp:eq:C.PR.1}
\]
Choose an oriented period basis $\lambda_1,\lambda_2$ of
$F=\C/\Lambda_F$. The exact sheaf sequence used here is
\[
 \begin{gathered}
 0\longrightarrow\OO_{\mathcal B}
 \xrightarrow{\ a\mapsto
 (e^{2\pi i\lambda_2a},e^{-2\pi i\lambda_1a})\ }
 (\OO_{\mathcal B}^{*})^2
 \xrightarrow{\ \Phi\ }F_{\mathcal B}\longrightarrow0,\\
 \Phi(u,v)=
 \frac{\lambda_1\log u+\lambda_2\log v}{2\pi i}
                          \pmod{\Lambda_F}.            \tag{C.PR.2}\label{cdp:eq:C.PR.2}
 \end{gathered}
\]
Changing a logarithm changes the numerator by a lattice period.
Conversely a pair in the kernel can be lifted to logarithms,
adjusted by integral periods, and then lies on the displayed
complex line. The intersection of that line with the integral
two-period lattice is zero because
$\lambda_2/\lambda_1$ is not real. Thus the additive kernel and
every map in \eqref{cdp:eq:C.PR.2} are exact.

\noindent\textbf{Step 2.} Finite stabilizer invariants are exact in characteristic zero, so
$H^j(\mathcal B,\OO)=H^j(\PP^1,\OO)=0$ for $j=1,2$.
Moreover
\[
 \operatorname{Pic}(\mathcal B)=\Z[\mathcal R],
                 \qquad \mathcal R^m=\OO_{\PP^1}(\infty).
                                                        \tag{C.PR.3}\label{cdp:eq:C.PR.3}
\]
Indeed every line has a unique stabilizer weight modulo $m$;
subtract that power of the canonical root line and it descends
to $\PP^1$. The displayed relation then gives the asserted
infinite cyclic group, with its canonical sections retained.
The long exact sequence of \eqref{cdp:eq:C.PR.2} therefore identifies the
actual $F$-torsor class with a unique vector in
$\operatorname{Pic}(\mathcal B)^2=\Z^2$.

\noindent\textbf{Step 3.} If this vector is zero, the torsor is the product
$F\times\mathcal B$. The inertia must act freely on its total
space because the actual connected root cover of $mJ$ is
unramified over the original surface tube. A nontrivial inertia
would act trivially on its product fibre. Thus $m=1$ and
$T=F\times\PP^1$, with $J=F\times\{\infty\}$.

\noindent\textbf{Step 4.} For a nonzero vector, choose its primitive direction and an
oriented period basis so that it is $(-n,0)$, $n>0$. This is
an integral change of the actual fibre basis, not removal of
its direction. The torsor is the pushout of
$(\mathcal R^{-n},\OO)$ by \eqref{cdp:eq:C.PR.2}. Gauge the second nonzero
line coordinate to one. The remaining ambiguity is multiplication
of the first coordinate by
$\gamma^{\Z}$, where
$\gamma=\exp(2\pi i\lambda_2/\lambda_1)$ and $|\gamma|<1$.
Consequently the actual total surface is
\[
                         (\mathcal R^{-n})^\times/\gamma^{\Z}.
                                                        \tag{C.PR.4}\label{cdp:eq:C.PR.4}
\]
Write $\mathcal B=\PP(m,1)$ with coordinates $(Z,W)$ of
weights $(m,1)$ and coarse coordinate $U=Z/W^m$.
The nonzero total line of $\mathcal R^{-n}$ is
\[
 [((\C^2\setminus0)\times\C^*_\ell)/
              \lambda\cdot(Z,W,\ell)
                =(\lambda^mZ,\lambda W,\lambda^{-n}\ell)].
                                                        \tag{C.PR.5}\label{cdp:eq:C.PR.5}
\]
Setting $\ell=1$ leaves the finite group $\mu_n$ with weights
$(m,1)$. A scalar multiplication by $\gamma$ lifts, after choosing
$\beta^n=\gamma$, to $(Z,W)\mapsto(\beta^mZ,\beta W)$.
Thus \eqref{cdp:eq:C.PR.4} is exactly
$H_m^{\rm diag}(\beta)/\mu_n$, with the marked axis $W=0$.
Freeness of the original inertia is exactly
$\gcd(m,n)=1$. Changing the chosen root of $\gamma$ changes
the infinite generator by the already included finite group.
Every primitive fibre basis, the full integer $n$, and the marked
axis identification have been retained.

\noindent\textbf{Step 5.} This also explains why reduction of $n$ modulo $m$ is impermissible:
the entire root-line degree in \eqref{cdp:eq:C.PR.4}, rather than only its
stabilizer character, is the torsor class.

\par\medskip\noindent\textbf{The actual inverse operation at a Hopf leaf.}
For a nonlinear leaf use its full normal form from \ref{cdp:st:C.2.UN.43};
for an elliptic terminal surface use \eqref{cdp:eq:C.PR.4}. Both have the
same marked normal-line data. Put $r=m$, choose $k,l$ with
$mk-nl=1$, and write $\gamma=\beta^n$. On the parent the
actual coarse root chart and its inverse product chart are
\[
 \begin{gathered}
 (u,X)\sim(\gamma^m u,\gamma^lX),\qquad p=X^m u^{-l},\\
 W=X^{-n}u^k,\qquad
                u=p^nW^m,\quad X=p^kW^l.
 \end{gathered}                                        \tag{C.PR.6}\label{cdp:eq:C.PR.6}
\]
These are precisely (C.SN.3)--(C.SN.5), with the given dual-normal
map and elliptic gluing used to choose the same connected root
cover. Replace this marked multiple-fibre tube by
$(p,W)\in\Delta\times(\C^*/\gamma^{\Z})$ through the displayed
isomorphism on $p\ne0$. This is the actual inverse logarithmic
transform. Its base point becomes an ordinary elliptic fibre.
Every other chart and transition of the parent is unchanged.
In particular the parent remains the same fixed-Jacobian
identity-monodromy quasi-bundle off this one point, with its full
analytic torsor retained by the displayed transition map.

\noindent\textbf{Step 6.} The choices do not discard gluing data. A change of a root
trivialization or a normal-product constant changes \eqref{cdp:eq:C.PR.6} by
a translation gauge extending over the inserted product disc.
The primitive nonzero direction in \eqref{cdp:eq:C.PR.3} fixes the angular
direction; the other period can change by an integral multiple,
which changes the generator by its same finite deck group.
Translations of the marked axis are retained in the transition.
This is the exact choice comparison already proved for the
full inverse torsor in \ref{cdp:st:C.2.UN.49} and \ref{cdp:st:C.2.UN.50}.

\noindent\textbf{Step 7.} Remove terminal vertices inward from every attached tree.
All original nonterminal vertices have the elliptic reduction
of \ref{cdp:st:C.2.UN.58}. Removing a child removes exactly its marked
multiple fibre by \eqref{cdp:eq:C.PR.6}, or leaves an ordinary fibre unchanged
in the product-cap case. Hence, after all children of a vertex
are removed, its only possible multiple fibre is its parent
mark. The one-point classification above applies to that
actual modified surface. This proves that every terminal step
has one of the two asserted forms.

\noindent\textbf{Step 8.} Operations in distinct discs of one base keep the same transition
atlas on their complement and hence commute. Operations in disjoint
subtrees do likewise. Every allowable pruning order differs by
interchanging such operations; a parent is processed only after
its children. This proves independence through actual marked
torsor isomorphisms, not merely through the same degree vector.
Keep the entire maps \eqref{cdp:eq:C.PR.6}, the primitive periods, the integer
$n$ and all remaining marked points as the record of each step.

\par\medskip\noindent\textbf{The exact two-point core models.}
At the end each cycle vertex is still a connected fixed-$F$
quasi-bundle over $\PP^1$, with only its two marked possible
multiple fibres, of orders $a,b$. Put $g=\gcd(a,b)$,
$A=a/g$, $B=b/g$. There is the actual representable finite
etale orbicurve map
\[
 \PP^1[(0,A),(\infty,B)]
       \longrightarrow \PP^1[(0,a),(\infty,b)],
                         \qquad t=z^g .                \tag{C.PR.7}\label{cdp:eq:C.PR.7}
\]
At zero its root uniformizers have map $v=u$ and inertia
$\mu_A\subset\mu_a$; the infinity chart has the analogous
inclusion. Thus it is etale as a map of these root curves.
The source, since $A,B$ are coprime, is the weighted projective
line $\PP(B,A)$, with coarse coordinate
$z=X^A/Y^B$. Its Picard group is generated by $\OO(1)$:
the two root divisors have classes $B$ and $A$, and Bezout
expresses their generator. Its higher coherent groups vanish.

\noindent\textbf{Step 9.} Pull back the actual $F$-torsor. The same sequence \eqref{cdp:eq:C.PR.2}
therefore describes it by an integral vector. A nonzero vector
$(-n,0)$ gives the explicit nonzero line $\OO(-n)$, and the
gauge \eqref{cdp:eq:C.PR.5}, now of weights $(B,A)$, gives the actual diagonal
Hopf cover with contraction $(\beta^BX,\beta^AY)$ and finite
action $\mu_n(B,A)$. The inertia is free precisely when
$\gcd(n,A)=\gcd(n,B)=1$. Its cyclic descent preserves both
marked axes. Lifts to a central primary cover commute with a
contraction and preserve each axis, so they have the form
$(Xf(X,Y),Yh(X,Y))$; contraction equivariance forces the
nonvanishing units $f,h$ to be constant. Thus the full quotient
is the actual marked diagonal primary or secondary Hopf surface.

\noindent\textbf{Step 10.} For the zero vector, free inertia forces $A=B=1$ and the cover
is $F\times\PP^1$. Its cyclic descent is
$(y,z)\mapsto(y+\tau,\zeta_g z)$: a holomorphic map
$\PP^1\to F$ is constant, and the principal-$F$ action excludes
a linear part. Freeness at $0,\infty$ makes $\tau$ have exact
order $g$. The quotient is exactly
\[
        \PP(\OO_E\oplus L),\quad
        E=F/\langle\tau\rangle,\quad
        L=(F\times\C)/((y,v)\sim(y+\tau,\zeta_gv)).
                                                        \tag{C.PR.8}\label{cdp:eq:C.PR.8}
\]
This exhibits its two marked sections, the exact torsion order
of $L$, the full fibre map $q=z^g$ on the cover, and both
elliptic identifications. It also gives the actual finite
projective-bundle map from $F\times\PP^1$.
Thus the pruned core is precisely one of the marked diagonal
Hopf/ruled cycles of \ref{cdp:st:C.2.UN.57}, with all its data retained.

\par\medskip\noindent\textbf{Necessity of the core lattice test through actual group maps.}
For an actual torus degeneration, the graph, normal and height
requirements are \ref{cdp:st:C.2.UN.54} and \ref{cdp:st:C.2.UN.58}. It remains to compare
the nearby group to that of the pruned core; no contraction of
the given holomorphic total space is required.

\noindent\textbf{Step 11.} The complement of the marked axis in a terminal diagonal Hopf
pair is exactly $\C\times F$. For a nonlinear pair the explicit
proper smooth-pair deformation of \ref{cdp:st:C.2.UN.49} identifies its compact
core and boundary with the diagonal pair. The product ruled cap
has the same complement $F\times\C$. In the primitive boundary
basis of \ref{cdp:st:C.2.UN.50} the actual homomorphism of this compact
piece is
\[
        \Z^3_{\delta,e,h}\longrightarrow\Z^2_{e,h},
                  \delta\longmapsto0,\quad
                  e\longmapsto e,\quad h\longmapsto h. \tag{C.PR.9}\label{cdp:eq:C.PR.9}
\]
For the Hopf case the exact matrix
$[(-l,k,0),(0,0,1),(m,-n,0)]$ has determinant one;
the normal reversal in $xy=t$ turns the killed Hopf meridian
into precisely $\delta$. The inverse product filling
$F\times D^2$ has the identical homomorphism \eqref{cdp:eq:C.PR.9},
with the same two retained generators. For a product cap this
is the literal product-disc boundary map.

\noindent\textbf{Step 12.} Van Kampen applied to the unchanged exterior and this boundary
gives two pushout diagrams with identical maps. Their resulting
groups are canonically isomorphic by the identity on the exterior
and the two displayed fibre generators. Repeat the comparison,
retaining every remaining collar. The actual tubular comparison
of \ref{cdp:st:C.2.UN.49} carries these maps to an arbitrary original smoothing.
Consequently the nearby group of the core plumbing is the
original $\Z^4$. The exact core calculation of \ref{cdp:st:C.2.UN.57} gives
$\Lambda\rtimes_M\Z$, with $\Lambda$ injected. Its abelianness
forces $M=I$. This proves necessity using the actual maps needed
for the test.

\par\medskip\noindent\textbf{The compact core family with a prescribed remaining period order.}
For sufficiency first handle the all-algebraic-dimension-zero
diagonal branch by \ref{cdp:st:C.2.UN.48}. Otherwise prune as above and assume
the core test passes. Its components have retained elliptic
quasi-bundles. In the common complement
$G=\C^2/\Lambda$ of \ref{cdp:st:C.2.UN.57} their regular fibres are translates
of the compact elliptic subgroup
\[
       F=(W\cap iW)/(\Lambda\cap W\cap iW),
                     \qquad W=\operatorname{span}_{\R}\Lambda.
                                                        \tag{C.PR.10}\label{cdp:eq:C.PR.10}
\]
Their lattice is primitive of rank two: their actual quotient
in the open two-boundary component is $\C^*$.
The root-cover maps already recorded identify every later
inserted regular fibre with this same elliptic group.

\noindent\textbf{Step 13.} In \ref{cdp:st:C.2.UN.57} replace $V_z(w)=w+zJ\overline w$ by
$V_{t^R}(w)=w+t^RJ\overline w$, retaining the term
$d\log t$ and the normal neck equation $xy=t$.
The exact real determinant becomes
\[
 C(1+|t|^{2R})
 \left(1+|t|^{2R}
       +\frac{d}{2\pi}(|t|^{2R}-1)\log|t|\right)>0.
                                                        \tag{C.PR.11}\label{cdp:eq:C.PR.11}
\]
The product $(|t|^{2R}-1)\log|t|$ is nonnegative.
Thus the same actual period parallelepipeds and both endpoint
plumbings give a smooth compact torus family on $\PP^1$ with
the prescribed core at zero and its conjugate at infinity.
At infinity divide the frame by $t^R$; the displayed frame is
$J\overline w+s^Rw$, and the logarithmic term still has exactly
the sign and coefficient for the reduced neck $xy=s$.

\noindent\textbf{Step 14.} Let $\lambda,\mu=c\lambda$ be an oriented primitive period basis
of the compact group in \eqref{cdp:eq:C.PR.10}, with $\operatorname{Im}c>0$.
The ordinary determinant in the actual vector frame is
\[
 \det_{\C}(V_{t^R}\lambda,V_{t^R}\mu)
       =2i\operatorname{Im}(c)\,\|\lambda\|^2 t^R\ne0
                      \quad(t\ne0).                    \tag{C.PR.12}\label{cdp:eq:C.PR.12}
\]
Here $\det(\lambda,J\overline\lambda)=-\|\lambda\|^2$.
This is the exact common core period, including its integral
class, rather than a chosen formal order.

\par\medskip\noindent\textbf{The recursive local model and the exact order drops.}
The actual logarithmic charts near a regular compact fibre give
a neighbourhood in the following rank-two period quotient:
\[
 \mathcal G=
 (\C_p\times\C_w^*\times\Delta_t)/
 \langle(p,w,t)\mapsto(p+\delta(t),\gamma(t)w,t)\rangle,
 \quad
 \delta(t)=t^\ell d(t),\quad d(0)\ne0,\quad|\gamma(0)|<1 .
                                                        \tag{C.PR.13}\label{cdp:eq:C.PR.13}
\]
Indeed choose the primitive first period as $(0,2\pi i)$ in
holomorphic linear coordinates; the second is then
$(\delta(t),\log\gamma(t))$. Its first coordinate is a unit
times their determinant. The charts of the core plumbing select
a bounded transverse disc in this cover and identify it with
the actual neighbourhood of the chosen central elliptic fibre.
Initially \eqref{cdp:eq:C.PR.12} gives $\ell=R$.

\noindent\textbf{Step 15.} Every prescribed new primitive fibre basis is allowed. If
$\lambda'=a\lambda+b\mu$, $\mu'=c\lambda+d\mu$ with
$ad-bc=1$, the coordinate
$p'=p-b\delta V/(2\pi ia+b\log\gamma)$ and the rescaled
logarithmic coordinate
$V'=2\pi iV/(2\pi ia+b\log\gamma)$ make $\lambda'=(0,2\pi i)$.
The second displacement is
$2\pi i\,\delta/(2\pi ia+b\log\gamma)$.
The denominator is a holomorphic unit near zero.
Thus its order is unchanged and the full new fibre basis is
retained by an actual linear period map.

\noindent\textbf{Step 16.} For a recorded Hopf step of weight $m$ and secondary integer $n$,
perform the weight-$(m,1)$ blowup of the base coordinates $(p,t)$
and tensor its nonzero line by $\OO(nE)$. In the free coarse
chart of \ref{cdp:st:C.2.UN.50} write $p=x^mZ$, $t=xW$.
A chosen holomorphic root $\beta(t)^n=\gamma(t)$ gives the exact
infinite map
\[
 (Z,W,x)\longmapsto
 \left(\beta(t)^m
    \bigl(Z+d(t)x^{\ell-m}W^\ell\bigr),
       \beta(t)W,\beta(t)^{-1}x\right).
                                                        \tag{C.PR.14}\label{cdp:eq:C.PR.14}
\]
The finite action has weights $(m,1,-1)$ modulo $n$ and is free,
because $\gcd(m,n)=1$. The added monomial has weight $m$,
so these maps commute. The inequality $\ell\ge m$ follows
from the height test. At $\ell=m$ the exceptional component
is exactly the recorded nonlinear Hopf surface after a constant
rescaling of $Z$. At $\ell>m$ it is exactly the recorded
diagonal pair. On the new child's complement of its parent curve
put $U=Z/W^m$. The actual deck map becomes
\[
        (U,W^n)\longmapsto
       \left(U+\frac{\delta(t)}{t^m},\,\gamma(t)W^n\right).
                                                        \tag{C.PR.15}\label{cdp:eq:C.PR.15}
\]
It has remaining displacement order $\ell-m$.
At any other marked point use $U-U_0$ as its base coordinate.
This retains the original marked point and every later child
transition. No extension of an arbitrary original smoothing's
elliptic quotient is being asserted.

\noindent\textbf{Step 17.} For a recorded product cap, $m=1$ and $\ell\ge2$, because this
vertex was an elliptic vertex of height strictly smaller than $R$.
Blow up the actual ideal $(p,t)$ along its compact elliptic
curve in \eqref{cdp:eq:C.PR.13}. The deck map preserves this ideal.
On its exceptional normal directions it is the identity,
since $\delta'(0)=0$; hence the exceptional is precisely
$F\times\PP^1$ with its parent section. The chart $p=tU$
has the exact map
\[
                    (U,w)\longmapsto
          (U+\delta(t)/t,\gamma(t)w).                   \tag{C.PR.16}\label{cdp:eq:C.PR.16}
\]
Thus this step has exactly the same order drop one. The strict
parent surface is unchanged, since its centre is a Cartier
curve on that surface. The two other blowup coordinates give
the reduced normal-crossing equation $t=xy$.

\noindent\textbf{Step 18.} Inductively, the model at a vertex of height $h$ has remaining
order $R-h$. All vertices with descendants are elliptic and
have positive remaining order. Only a nonlinear terminal step
has order zero after insertion. The height test is therefore
precisely what makes every recorded step in \eqref{cdp:eq:C.PR.14}--\eqref{cdp:eq:C.PR.16}
holomorphic with the required central type.

\par\medskip\noindent\textbf{Local properness and simultaneous compact replacement.}
We give the needed bounds also for the affine, noncompact base
of \eqref{cdp:eq:C.PR.13}. On its weight-$(m,1)$ blowup use the actual
canonical-section norm
\[
                  \rho(p,t)=
                  (|p|^2+|t|^{2m})^{1/(2m)}.           \tag{C.PR.17}\label{cdp:eq:C.PR.17}
\]
In the chart $p=x^mZ$, $t=xW$ this is
$|x|(|Z|^2+|W|^{2m})^{1/(2m)}$, the nonvanishing smooth
metric factor on the weighted exceptional charts. Translation
by $j\delta(t)$ gives two-sided norm distortion at most
$C(1+|j|)^{1/m}$, because $|\delta(t)|\le C|t|^m$.
For $\OO(nE)$ the exponent is $n/m$. On a sufficiently small
closed disc the scalar $|\gamma(t)|$ is uniformly smaller
than one. Its exponential contraction dominates this polynomial.
The infinite action is free and properly discontinuous, and
a high fixed iterate gives a norm-annulus representative for
every orbit, exactly as in \ref{cdp:st:C.2.UN.42}.

\noindent\textbf{Step 19.} The old line coordinate is the new line divided by the $n$th
power of the canonical exceptional section. If an old image
stays in a compact old annular chart, putting both line norms
in annuli changes the base by at most a constant times
$|\delta(t)|(1+|\log\rho|)$. For $\rho<1$ use $\rho\ge|t|$;
this quantity is bounded and tends to zero as $t\to0$.
For $\rho\ge1$, a new base representative consequently obeys
\[
 |p_{\rm new}|\le B+C|\delta(t)|\log(2+|p_{\rm new}|).
                                                        \tag{C.PR.18}\label{cdp:eq:C.PR.18}
\]
It is uniformly bounded: a linear function eventually dominates
the logarithm. Thus the new annulus over the required base set
is compact, including the compact exceptional Hopf component.
This proves the local compactness comparison without assuming
that the affine base itself is compact.

\noindent\textbf{Step 20.} Off the new Hopf component the canonical exceptional section
gives the actual isomorphism with the old model minus its
marked central elliptic curve. In a compact Hopf fundamental
domain, $x\to0$ gives $p=x^mZ$, $t=xW$ and old line coordinate
proportional to $x^{-n}$. Normalizing that coordinate takes
$O(1+|\log|x||)$ iterates. Every old base cluster point is zero,
because
$|x|^m+|t|^\ell(1+|\log|x||)\to0$.
The normal-annulus coordinate has a compact limit on the old
elliptic curve. Hence every old cluster point lies on that
curve and on no other point of the old model.

\noindent\textbf{Step 21.} Choose disjoint small neighbourhoods of the finitely many
curves at each insertion stage. The new open piece is the
inserted compact component together with the inverse image
of the punctured old neighbourhood. The preceding cluster
comparison proves it is open and that all gluing graphs are
relatively closed. The resulting complex space is Hausdorff.
The bounded-annulus comparison over smaller closed
neighbourhoods and the unchanged compact exterior prove
properness. Product-cap steps are ordinary proper blowups.
Finite iteration therefore gives a smooth proper family over
a disc, glued by actual biholomorphisms to the unchanged
compact core family away from zero.

\par\medskip\noindent\textbf{Restoration of the prescribed support.}
On the strict parent of every Hopf insertion the chart is
exactly \eqref{cdp:eq:C.PR.6}, in the opposite direction: the old coordinate
is $W=x^{-n}$ and the full secondary integer is unchanged.
It restores the complete recorded elliptic torsor, not just
its rational Chern class. All untouched charts and all
other marked fibres remain their recorded charts. On a cap
step the actual product exceptional pair and its parent
section are the recorded ones. Fibre-basis maps retain
the linear parts of every elliptic gluing. The constant
translations are retained in the recorded transition maps;
equivalently the automorphisms
$(Z,W)\mapsto(a^mZ,aW)$ of either Hopf contraction realize
every required marked-axis translation. They commute with
the finite action. Product caps have the corresponding
elliptic translations and projective automorphisms fixing
the parent section.

\noindent\textbf{Step 22.} Reverse the complete pruning record. At each step these
actual maps recover exactly the surface and marked gluing
removed at that step. Induction therefore restores the
specified $S$ in its full analytic isomorphism class.
The normal equations are $t=xW$ in the free weighted
coarse charts and $t=xy$ for an ordinary blowup, so every
component has full multiplicity one. The other fibres
are unchanged from the compact period family, and the
filling at infinity is its conjugate pruned core.
All these reduced supports have Euler zero: their
components are Hopf or elliptic ruled surfaces, their
double curves are elliptic, and there are no triple points.
This proves sufficiency and the compact assertion.

\noindent\textbf{Step 23.} The exact comparisons are reproduced in
\path{code/E_hopf_tree_pruning.py} and its saved output:
primitive integral fibre-basis changes, both-point root Picard
lattices, the complete weighted deck and finite-action maps,
the nested remaining-order identities, and the product-cap
blowup. These checks retain the full integers in the pruning
record; the proof establishes them for arbitrary finite trees.
\end{proof}






\begin{proposition}[The exact rooted elliptic-tree criterion and compact occurrence]\label{cdp:st:C.2.UN.61}
\StatementRecord{cdp:st:C.2.UN.61}{proved}{Historical events: Claude Fable 5.2 (Anthropic), effort high, 16 November 2026: carried as background without fault found. GPT-6 Astra (OpenAI), ChatGPT Pro mode, effort Pro, S6\_09\_PAPERS, 17 November 2026: targeted proof check. Also checked in this writing. Pre-continuation writing events: GPT-6 Astra (OpenAI), ChatGPT Pro mode, effort Pro, S6\_10\_PAPERS, 18 November 2026: targeted proof check. Continuation metadata events: GPT-6 Astra (OpenAI), ChatGPT Pro mode, effort Pro, S6\_10\_PAPERS continuation, 19 November 2026: used and rechecked in this run; status-vocabulary reconciliation; metadata reconciliation only; no new mathematical check. \emph{Pre-continuation recorded proof scope.} The rooted normalization incidence map is surjective by its nonzero leaf pivots, giving the exact root Hodge dimensions. The canonical residue divisor makes the root minimal; the degree-one root-component comparison yields its actual elliptic fibration and all marked fibre supports. Normal-root collars give the exact height increments and common integral period. Pruning retains the complete principal inverse torsor; its unchanged boundary group forces the Chern commutator to vanish in the injected fibre lattice. The resulting actual torus has an everywhere invertible twistor period family. Reversing the full recorded weighted and ordinary charts restores every specified torsor and gluing, with the proved proper cluster estimates and reduced local equations}{the hypotheses, notation and definitions stated here; \cref{cdp:st:C.2.UN.60}; \cref{cdp:st:C.2.UN.5}; \cref{cdp:st:C.2.UN.54}; \cref{cdp:eq:C.PT.3}; \cref{cdp:st:C.2.UN.28}; \cref{cdp:eq:C.PT.5}; \cref{cdp:eq:C.PT.6}; \cref{cdp:st:C.2.UN.31}; \cref{cdp:eq:C.PT.7}; \cref{cdp:eq:C.PT.8}; \cref{cdp:st:C.2.UN.58}; \cref{cdp:eq:C.PT.9}; \cref{cdp:st:C.2.UN.50}; \cref{cdp:eq:C.PT.11}}

Let $S$ be a specified compact connected reduced SNC surface whose
incidence graph is a nontrivial tree. Distinguish one smooth component
$P$, with no initial restriction on its geometric genus or surface
type. Every other component is either a primary or secondary Hopf
surface, or is equipped with a specified fixed-Jacobian elliptic
quasi-bundle over $\PP^1$, with identity integral monodromy,
connected fibres and only smooth multiple fibre supports.
In the latter case every boundary curve is a fibre support of
the specified map. Retain the actual elliptic gluing maps
and the complete analytic torsors. The following finite
geometric test is necessary and sufficient for $S$ to occur
as a full fibre of a smooth-total degeneration of connected complex
two-tori.

\begin{enumerate}
\item The actual normal lines at every double curve are dual.
\item The root $P$ has a fixed-Jacobian elliptic quasi-bundle
$q_P:P\to B$ over a smooth elliptic curve, with identity integral
monodromy and only smooth multiple supports. Its marked boundary
curves are fibre supports, and every multiple fibre is marked.
Every other component is either a nonlinear algebraic-dimension-zero
Hopf leaf or an elliptic quasi-bundle over $\PP^1$ of the stated
type, whose multiple fibres are all marked double curves.
\item For each edge $e$, let $m_e$ be the exact finite order of
its normal line. It is the full elliptic multiplicity at each
elliptic endpoint and the resonance at a nonlinear endpoint.
Set $h_P=0$ and
\[
                  h_v=\sum_{e\in[P,v]}m_e.
                  \tag{C.PT.1}\label{cdp:eq:C.PT.1}
\]
There is an integer $R>0$ such that every nonlinear leaf has
$h_v=R$ and every elliptic nonroot vertex has $h_v<R$.
In the absence of nonlinear leaves any $R>\max_vh_v$ passes.
\item Perform the exact leaf-pruning operations of \ref{cdp:st:C.2.UN.60},
retaining the complete inverse logarithmic charts, primitive fibre
directions, secondary integers and torsor transitions. The final
surface is a principal elliptic torsor $T_0\to B$. Its integral
Chern vector in $H^2(B,\Lambda_F)$ is zero.
\end{enumerate}

The normal-root maps identify the one fixed Jacobian $F$ along
the tree. The last condition refers to the actual analytic
inverse torsor; it does not replace it by its rational Chern
data. The pruning result is independent, up to the actual
marked isomorphisms of \ref{cdp:st:C.2.UN.60}, of the order of disjoint operations.

Every support passing the test has a realization over $\PP^1$
with smooth compact total threefold. Its only singular fibre
is the given reduced $S$; every other fibre is a smooth
two-torus. All reduced fibre supports have Euler number zero.
Every component and every specified elliptic gluing map is
retained, including the complete torsor of every original
properly elliptic nonroot component.

In every actual occurrence the root is minimal,
$h^q(P,\OO_P)=h^q(S,\OO_S)=(1,2,1)$, and its actual canonical
divisor is
\[
                 K_P=\sum_{e\ni P}(m_e-1)J_e.
                 \tag{C.PT.2}\label{cdp:eq:C.PT.2}
\]
It is the unique minimum-order component for canonical evaluation;
the differences of the other orders are exactly \eqref{cdp:eq:C.PT.1}.
The zero-edge case, excluded from the nontrivial tree here,
is the smooth torus fibre already treated in \ref{cdp:st:C.2.UN.5}.
\end{proposition}



\begin{proof}
\noindent\textbf{Step 1.} \par\medskip\noindent\textbf{The root cohomology is forced by the actual incidence maps.}
Every double curve has a nonroot endpoint. At a Hopf
endpoint it is smooth elliptic by the curve argument of
\ref{cdp:st:C.2.UN.54}; at a supplied quasi-bundle endpoint it is a smooth
elliptic fibre support by hypothesis. Distinct boundary
curves on either such endpoint are disjoint, so there
are no triple points. In particular the curves $J_e$
on $P$ are also pairwise disjoint. Write $T_v$ for a
nonroot component. The exact normalization sequence is
\[
 0\longrightarrow\OO_S
 \longrightarrow\OO_P\oplus\bigoplus_{v\ne P}\OO_{T_v}
 \longrightarrow\bigoplus_e\OO_{J_e}\longrightarrow0.
 \tag{C.PT.3}\label{cdp:eq:C.PT.3}
\]
For a Hopf component, $h^q(\OO_{T_v})=(1,1,0)$ and
every boundary restriction on $H^1$ is nonzero by \ref{cdp:st:C.2.UN.54}.
For a supplied quasi-bundle, (C.QT.2)--(C.QT.3) give
the same dimensions and the actual commuting square:
restriction to $J_e$ followed by its connected normal-root
pullback is evaluation of the fixed $H^1(F,\OO_F)$
frame. Thus its boundary restriction is also nonzero.
These are the actual restriction and gluing maps.
Root the tree at $P$ and order the nonroot columns from the
leaves inward. The submatrix of the $H^1$ incidence map
on those columns is square, with one column and row per
nonroot vertex. Elimination at a leaf has its nonzero
restriction as pivot. Repeating gives an invertible
matrix. The full $H^1$ incidence map is therefore onto,
regardless of the root restriction maps.

\noindent\textbf{Step 2.} The $H^0$ incidence map of a tree is also onto. Since
$H^2(\OO_{T_v})=0$ on every nonroot component, the actual long exact
sequence of \eqref{cdp:eq:C.PT.3} gives
\[
 H^2(S,\OO_S)\simeq H^2(P,\OO_P),\qquad
                 h^1(S,\OO_S)=h^1(P,\OO_P).
 \tag{C.PT.4}\label{cdp:eq:C.PT.4}
\]
The smooth total map is flat over the disc. Upper
semicontinuity gives $h^2(S,\OO_S)\ge1$, while the
compact-core bound \ref{cdp:st:C.2.UN.28} gives $p_g(P)\le1$. Thus
$p_g(P)=1$. Coherent Euler characteristic in this proper
flat family is zero. Additivity in \eqref{cdp:eq:C.PT.3}, since every
nonroot component and every elliptic curve has coherent Euler zero,
then gives $\chi(\OO_P)=0$ and $q(P)=2$. This proves
the asserted coherent dimensions without an initial
geometric-genus hypothesis on $P$.

\par\medskip\noindent\textbf{The actual dualizing form and the unique minimum.}
Use the residue sequence (C.HR.7), or the identical
sequence of the actual boundary divisor, at every
nonroot vertex. At a terminal vertex its one nonzero restriction
coefficient forces the sole residue of a global
dualizing form to be zero. Since this component has no
holomorphic two-form, the form on that component is zero.
Remove that leaf and repeat. All nonroot forms and every
edge residue are zero. The remaining form on $P$ is
an ordinary holomorphic two-form. Conversely such a form,
with zero on all nonroot components, glues by the zero-residue
rule. These are inverse maps, giving the actual isomorphism
\[
                  H^0(S,\omega_S)\simeq H^0(P,K_P).
 \tag{C.PT.5}\label{cdp:eq:C.PT.5}
\]

\noindent\textbf{Step 3.} Take a canonical evaluation generator, with vertical orders
$k_v$, and divide it by $t^{\min k_v}$. Its nonzero restriction
to the reduced fibre spans \eqref{cdp:eq:C.PT.5}. The local adjunction map
(C.HR.8) identifies its nonzero components precisely with the
minimum-order components. Hence the minimum is attained
only at $P$. Writing $k_P=k_0$, the same actual residue
calculation on $P$ gives a holomorphic two-form with divisor
\[
          K_P=\sum_{e=(P,w)}(k_w-k_0-1)J_e,\qquad
                          k_w>k_0.
 \tag{C.PT.6}\label{cdp:eq:C.PT.6}
\]
All coefficients are nonnegative integers; there is no
additional zero away from the boundary, because the
evaluation has only its vertical component divisors.
If $E$ were a $(-1)$-curve, adjunction would give
$K_P\cdot E=-1$. The effective divisor in \eqref{cdp:eq:C.PT.6}
is supported on elliptic curves, none equal to $E$,
and has nonnegative intersection with $E$. This
contradiction proves that $P$ is minimal.

\par\medskip\noindent\textbf{The actual elliptic fibration on the root.}
Its nonzero canonical form gives $\kappa(P)\ge0$.
Every component has full multiplicity one, and the
given total space and central support are already
smooth and SNC. Thus the actual common-root map in
\ref{cdp:st:C.2.UN.31} has degree one and its component over $P$ is
$P$ itself. That theorem makes $P$ a two-torus or
gives on it a fixed-modulus, identity-monodromy
elliptic quasi-bundle with only smooth multiples.
Its base has genus $p_g(P)=1$.

\noindent\textbf{Step 4.} We record an exact principal-torsor comparison used twice
below. For a principal $F=\C/\Lambda_F$-torsor over a
smooth elliptic curve $B$, the sequence
\[
 0\longrightarrow\Lambda_F
 \longrightarrow\OO_B\longrightarrow F_B
 \longrightarrow0
 \tag{C.PT.7}\label{cdp:eq:C.PT.7}
\]
maps its actual analytic class to its integral Chern
vector. If that vector is zero, the class comes from
$H^1(B,\OO_B)$. On $B=\C/\Lambda_B$, a representative
$a\,d\overline z$ has additive transition functions
$a\overline\gamma$ for $\gamma\in\Lambda_B$.
Its actual total torsor is therefore
\[
 \C^2/
 \left\langle
  (z,w)\mapsto(z+\gamma,w+a\overline\gamma),\
  (z,w)\mapsto(z,w+\lambda):
   \gamma\in\Lambda_B,\ \lambda\in\Lambda_F
 \right\rangle .
 \tag{C.PT.8}\label{cdp:eq:C.PT.8}
\]
The four real periods are independent by projection
to $\Lambda_B$ and then to $\Lambda_F$; they give
a discrete cocompact lattice. This is a complex
two-torus, with the original torsor class represented
by the same additive transitions. A different lift
through \eqref{cdp:eq:C.PT.7} gives its isomorphic original
$F$-torsor. No section of the original torsor has
been assumed.

\noindent\textbf{Step 5.} We now identify every marked elliptic curve on $P$
with a fibre of one such elliptic fibration.
On a torus an elliptic curve is a translate of an
elliptic subtorus: lifting its map to universal
covers makes its derivative a periodic holomorphic
function, hence constant. Two different elliptic
directions give an isogeny from their product to
the two-torus, so their translates intersect.
The pairwise disjoint boundary curves therefore
have the same direction and are fibres of the
actual quotient by that elliptic subtorus.

\noindent\textbf{Step 6.} In the other alternative let $q:P\to B$ be the
fibration supplied by \ref{cdp:st:C.2.UN.31}. If it has a multiple
fibre, any nonvertical marked elliptic curve would
map finitely to the elliptic $B$. Riemann--Hurwitz
makes that map unramified. But at a multiple fibre
the actual local function is $s^m$, $m>1$,
so the restriction to that curve has positive
ramification, a contradiction.
If there is no multiple fibre, $q$ is a principal
$F$-torsor. A nonvertical elliptic curve would give
a section after its finite base pullback. Pullback
on $H^2(B,\Lambda_F)$ multiplies its Chern vector
by the positive covering degree. The vector group
is torsion-free, so the original vector would be
zero. Formula \eqref{cdp:eq:C.PT.8} then makes $P$ a torus,
already handled above. Otherwise all marked
curves are vertical. Since every reduced fibre
is one smooth elliptic curve, each is exactly
a marked fibre support. This proves the required
actual fibration $q_P$ in all cases.

\par\medskip\noindent\textbf{Finite normal roots and the exact height equations.}
The normal-root chart (C.HT.4) gives each root
boundary normal a finite exact order equal to its
full elliptic multiplicity. Normal duality carries
that order to the adjacent component.
An algebraic-dimension-zero diagonal Hopf has only
nontorsion axis normals, so it cannot be adjacent
to $P$. At an elliptic Hopf all curve normals have
finite order, and a nonlinear Hopf has only its
single torsion-normal axis. A supplied quasi-bundle
has finite boundary normal orders by (C.QT.2).
Following any path from $P$ therefore excludes
its first algebraic-dimension-zero diagonal
vertex, and hence all such vertices. A nonlinear
vertex is terminal because it has only one curve.
Every remaining nonroot vertex is an elliptic
quasi-bundle of the required genus-zero type.

\noindent\textbf{Step 7.} For each elliptic component, the actual root
charts give the global meromorphic form
$q^*\theta\wedge dz$ and its exact order
$m\operatorname{ord}\theta+m-1$, as in
(C.HT.5)--(C.HT.6) and (C.QT.6). These formulas
use translation charts, not the Kodaira dimension
of the component. The canonical residue form
on that component has divisor only on its
marked curves. At an unmarked multiple fibre
its order would be zero, which is incompatible
with $m\operatorname{ord}\theta+m-1=0$ for
$m>1$. This proves the absence of all unmarked
multiples on $P$ and on every elliptic nonroot
vertex.

\noindent\textbf{Step 8.} The dual-normal maps identify the connected
normal-root elliptic covers. Beginning with
the fixed Jacobian of $q_P$ and following the
unique path in the tree gives its actual
identification with every later Jacobian and
with each nonlinear axis root cover. In
particular the exact normal order $m_e$ is
the same at both endpoints.

\noindent\textbf{Step 9.} At every elliptic nonroot vertex the residue
calculation gives the integer equation
\[
 \sum_{e=(v,w)}
       \frac{k_w-k_v}{m_e}
                        =\operatorname{val}(v)-2.
 \tag{C.PT.9}\label{cdp:eq:C.PT.9}
\]
A terminal elliptic vertex gives the outward
increment $m_e$. A terminal nonlinear vertex
gives the same increment from its exact
axis form $dz\wedge dw/w^{m_e+1}$ and its
infinite-order effective axis line.
The finite leaf induction of \ref{cdp:st:C.2.UN.58} applies
without an equation at the root: once every
child term is one, \eqref{cdp:eq:C.PT.9} makes its parent
term minus one. Therefore every outward
increment is $m_e$, proving \eqref{cdp:eq:C.PT.1}.
Substituting the root increments into
\eqref{cdp:eq:C.PT.6} proves \eqref{cdp:eq:C.PT.2}.

\noindent\textbf{Step 10.} The actual root collars $xy=t$ and the full
secondary boundary matrices of \ref{cdp:st:C.2.UN.50} identify
every regular elliptic fibre and every
nonlinear real core with one transported
integral two-cycle $\gamma_t$ in the nearby
torus. They are exactly the maps
(C.HT.10)--(C.HT.12); the genus of the root
base does not change those collar maps.
For the canonical evaluation divided by
$t^{k_0}$, its leading form on an elliptic
vertex restricts to zero on the elliptic
curve. Its period is therefore
$O(t^{h_v+1})$, including at $P$.
At a nonlinear leaf it is
\[
 t^{h_v}
       \left(\frac{\pm2\pi i c\beta^{-m_e}}{n}
                         +O(t)\right),\qquad c\ne0.
 \tag{C.PT.10}\label{cdp:eq:C.PT.10}
\]
These are the same holomorphic period, by the
actual integral collar identifications.
One nonlinear leaf fixes its finite positive
order $R$; all nonlinear heights equal $R$,
and every elliptic vertex has smaller height.
This proves the stated height test.

\par\medskip\noindent\textbf{The exact inverse torsor and its integral obstruction.}
Perform the operations (C.PR.6) from the leaves
inward. After its children have been removed,
an elliptic nonroot vertex retains its actual
fixed-Jacobian map over $\PP^1$: each operation
replaced exactly the specified multiple-fibre
tube by the ordinary product tube and left
all other transitions unchanged. Its only
possible multiple fibre is now the parent
mark. It is therefore the actual one-mark
torsor of (C.PR.1)--(C.PR.5), hence is a
recorded diagonal Hopf pair or the product
ruled cap. The root-line classification retains
its entire torsor, even when the original
unpruned surface had Kodaira dimension one.
Nonlinear leaves use the same
inverse product map, with their full
secondary integer and primitive fibre
direction. The root always remains its
fixed-Jacobian quasi-bundle over the original
elliptic $B$. Each operation replaces exactly
one marked multiple-fibre tube by its
displayed ordinary product tube; a product
cap leaves that ordinary tube unchanged.
There are no unmarked multiples. After all
nonroot vertices have been removed, the final
root is therefore a principal $F$-torsor
$T_0\to B$, with its full analytic transition
class retained.

\noindent\textbf{Step 11.} The exact boundary homomorphism of every
removed compact piece is (C.PR.9): it kills
the one meridian $\delta$ and retains the
two generators $e,h$. The inverse product
filling has those identical maps. Applying
the actual Van Kampen pushouts successively,
as in \ref{cdp:st:C.2.UN.60}, gives
\[
                         \pi_1(T_0)\simeq\pi_1(F_t)
                                             \simeq\Z^4.
 \tag{C.PT.11}\label{cdp:eq:C.PT.11}
\]
Only these smooth collar maps are used for
an arbitrary original degeneration; no
holomorphic contraction of its neighbourhood
is asserted.

\noindent\textbf{Step 12.} For a principal elliptic bundle over an
elliptic curve the homotopy sequence has
an injected fibre lattice, since
$\pi_2(B)=0$. The group is the central
extension
\[
 1\longrightarrow\Lambda_F
 \longrightarrow\pi_1(T_0)
 \longrightarrow\Lambda_B\longrightarrow1.
 \tag{C.PT.12}\label{cdp:eq:C.PT.12}
\]
The commutator of lifts of an oriented
base basis is its integral Chern vector.
This follows directly by trivializing over
the usual base parallelogram: its boundary
commutator is the clutching loop of the
fibre torus. Equation \eqref{cdp:eq:C.PT.11} makes that
commutator zero in the injected lattice.
Thus the Chern vector is zero. Formula
\eqref{cdp:eq:C.PT.8} identifies this very inverse
torsor, including its analytic class,
with a complex two-torus. This proves
necessity of the last test. Independence
of disjoint pruning orders is the actual
chart comparison of \ref{cdp:st:C.2.UN.60}.

\par\medskip\noindent\textbf{Compact reconstruction of the complete prescribed tree.}
Assume all the tests pass and retain the
complete pruning record. Formula \eqref{cdp:eq:C.PT.8}
gives the actual initial torus
$T_0=\C^2/\Gamma$. Its elliptic fibre
lattice is a primitive rank-two subgroup
of $\Gamma$, because the quotient base
lattice is free. Choose $R$ as in the
height test and put
\[
 V_t(w)=w+t^R J\overline w,\qquad
              J=\begin{pmatrix}0&1\\-1&0\end{pmatrix}.
 \tag{C.PT.13}\label{cdp:eq:C.PT.13}
\]
For any real basis of $\Gamma$ its image
has determinant $(1+|t|^{2R})^2$ times
the original nonzero determinant.
These holomorphic period sections
therefore give a proper smooth torus
family on the finite chart. At infinity
the frame divided by $t^R$ gives
$J\overline w+s^Rw$, $s=t^{-1}$.
It is again real-linearly invertible
at every point, including $s=0$.
The two charts consequently give an
actual compact smooth torus family
over $\PP^1$, with central member
exactly $T_0$.

\noindent\textbf{Step 13.} For an oriented primitive fibre basis
$\lambda,\mu=c\lambda$,
$\operatorname{Im}c>0$, its exact period
in this frame is
\[
 \det_{\C}(V_t\lambda,V_t\mu)
       =2i\operatorname{Im}(c)\|\lambda\|^2t^R.
 \tag{C.PT.14}\label{cdp:eq:C.PT.14}
\]
The coefficient is nonzero. Near each
recorded central elliptic fibre use
the actual rank-two period cover and
the primitive-basis change (C.PR.13).
This gives
$(p,w)\mapsto(p+\delta(t),\gamma(t)w)$
with $\delta(t)=t^R d(t)$, $d(0)\ne0$,
and the original marked Jacobian and
primitive fibre direction.

\noindent\textbf{Step 14.} Reverse the pruning record. At a
recorded Hopf step of weight $m$ use
the free weighted chart (C.PR.14),
including its full secondary integer.
Its child's affine pencil has the
exact displacement $\delta(t)/t^m$.
At a product-cap step use the ordinary
blowup (C.PR.16), with displacement
$\delta(t)/t$. Thus the remaining
order at every vertex is $R-h_v$.
The height test makes each recorded
elliptic child linear and every
recorded nonlinear terminal child
nonlinear, with its prescribed
resonance and contraction. The
recorded primitive basis changes
retain all linear gluing maps,
and the retained transition charts
restore all elliptic translations
and every analytic torsor. In particular an
original properly elliptic nonroot surface
is restored by reversing its successive
root-line and ordinary-tube changes. No
replacement by a surface with only the same
degrees or normal characters occurs.

\noindent\textbf{Step 15.} All local compactness comparisons
are precisely (C.PR.17)--(C.PR.18).
Exponential contraction dominates
the polynomial exceptional norm
distortion, and the old base
displacement is
$O(|\delta(t)|\,|\log|x||)$.
Its cluster set lies in the old
marked elliptic curve. These exact
bounds give disjoint Hausdorff
proper replacements at each stage;
ordinary caps are proper blowups.
There are finitely many stages.
Each is the identity on the
punctured family and outside its
chosen small central neighbourhood.
They therefore preserve the
compact family over the rest of
$\PP^1$, including infinity.

\noindent\textbf{Step 16.} Induction through the recorded
inverse charts restores precisely
the specified $P$, every nonroot
component and each prescribed
elliptic gluing map. Their normal
equations are the free coarse
$t=xW$ charts or the ordinary
$t=xy$ charts, so the full fibre
is reduced and the total space
is smooth. Every other fibre is
the unchanged smooth torus.
The root quasi-bundle and every supplied
genus-zero quasi-bundle have Euler number
zero by additivity over their smooth
elliptic fibre supports. Every remaining
Hopf component and every double elliptic
curve also has Euler number zero. Additivity proves
the stated zero Euler number
for $S$ and for all other supports.

\noindent\textbf{Step 17.} The rooted incidence matrices,
the integral Chern commutator,
the exact torsor period matrix
and the height equations are
checked in
\path{code/A_rooted_hopf_trees.py};
the primitive period and nested
order identities are also
certified in
\path{code/A_tree_periods.py}.
\end{proof}






\begin{proposition}[Unicyclic quasi-bundle supports with all torsors retained]\label{cdp:st:C.2.UN.62}
\StatementRecord{cdp:st:C.2.UN.62}{proved}{Historical events: Claude Fable 5.2 (Anthropic), effort high, 16 November 2026: carried as background without fault found. GPT-6 Astra (OpenAI), ChatGPT Pro mode, effort Pro, S6\_09\_PAPERS, 17 November 2026: targeted proof check. Also checked in this writing. Pre-continuation writing events: GPT-6 Astra (OpenAI), ChatGPT Pro mode, effort Pro, S6\_10\_PAPERS, 18 November 2026: targeted proof check. Continuation metadata events: GPT-6 Astra (OpenAI), ChatGPT Pro mode, effort Pro, S6\_10\_PAPERS continuation, 19 November 2026: used and rechecked in this run; status-vocabulary reconciliation; metadata reconciliation only; no new mathematical check. \emph{Pre-continuation recorded proof scope.} The actual finite root chart gives the exact normal order, trivial R1 direct image and nonzero boundary restriction. The normalization incidence matrix and dual residue maps force a single cycle and identify the canonical minima. Its meromorphic differential order gives all height increments and removes unmarked multiples; the common integral period gives the strict height test. Full root-line torsor pruning and the unchanged boundary maps yield identity core monodromy. Reversing every recorded insertion restores the specified quasi-bundles, torsors and gluings compactly. The four triple-root line pushout was identified with its actual transition y=sW, and the invariant-coordinate pole calculation gives the printed plurigenera and K-cubed pullback}{the hypotheses, notation and definitions stated here; \cref{cdp:st:C.2.UN.60}; \cref{cdp:eq:C.QT.2}; \cref{cdp:st:C.2.UN.54}; \cref{cdp:st:C.2.UN.50}; \cref{cdp:st:C.2.UN.44}; \cref{cdp:eq:C.QT.3}; \cref{cdp:eq:C.QT.7}; \cref{cdp:eq:C.QT.6}; \cref{cdp:eq:C.QT.8}; \cref{cdp:eq:C.QT.5}; \cref{cdp:eq:C.QT.9}; \cref{cdp:st:C.2.UN.57}; \cref{cdp:eq:C.QT.11}}

Let $S$ be a specified compact connected SNC surface with smooth
embedded components. Each component is either a nonlinear primary
or secondary Hopf surface of algebraic dimension zero, or is
equipped with an actual map
\[
                         q_v:T_v\longrightarrow\PP^1
\tag{C.QT.1}\label{cdp:eq:C.QT.1}
\]
with connected fibres, every fibre a multiple of a smooth elliptic
curve, fixed elliptic Jacobian, and identity monodromy on the
integral first cohomology of the regular fibre. Every boundary
curve on $T_v$ is a fibre support of this specified map. The
elliptic gluing maps, the surfaces and their complete analytic
torsors are part of the data.

The following geometric test is necessary and sufficient for
$S$ to occur as a full reduced fibre of a smooth-total
degeneration of connected complex two-tori.
\begin{enumerate}
\item The actual boundary normal lines are dual, and the incidence
graph has exactly one independent cycle. Nonlinear Hopf vertices
are leaves.
\item Every multiple fibre of every map \eqref{cdp:eq:C.QT.1} is marked by
a boundary curve. Let $m_e$ be the exact order of its normal
line, also the normal order at the other endpoint. At a
nonlinear endpoint this is its resonance.
\item Set $h_v=0$ on the cycle and set $h_v$ equal to the sum
of the $m_e$ along its path from the cycle. There is an integer
$R>0$ for which every nonlinear leaf has height $R$ and every
elliptic quasi-bundle vertex has height strictly smaller than
$R$. With no nonlinear leaves one may choose any
$R>\max_v h_v$.
\item Apply the exact inverse-logarithmic and product-cap
pruning maps of \ref{cdp:st:C.2.UN.60}, retaining their full integer
framings and all torsor identifications. The resulting
marked diagonal Hopf or degree-zero ruled core cycle has
integral complement-lattice product $M=I$.
\end{enumerate}
All pruning operations in this test exist under its preceding
requirements. Their marked result is independent of the order
of disjoint branch operations. Every support passing the test
has a realization over $\PP^1$ with smooth compact total space,
retaining every original component, analytic torsor and
specified elliptic gluing map. The only other singular support
is the complex conjugate of its pruned core. Every reduced
fibre support has Euler number zero.

In particular the statement includes elliptic quasi-bundles
of Kodaira dimension one and coherent Euler characteristic
zero, together with the genus-zero quasi-bundle components
of Kodaira dimension zero and the Hopf and ruled components
in the same explicitly specified class.
\end{proposition}



\begin{proof}
\noindent\textbf{Step 1.} \par\medskip\noindent\textbf{The actual root chart and its cohomology maps.}
Let $mJ$ be a fibre of \eqref{cdp:eq:C.QT.1}. The normalized local root
$t=s^m$ is finite unramified over the surface tube, by
the exact chart $s=wx$, $w^m=u$ for $t=ux^m$.
Its generic fibre is connected. Its finite Stein factor
over the root disc has degree one, so the central
elliptic cover is connected as well. The fixed Jacobian
and identity linear monodromy make the root family a
translation torsor under the fixed elliptic curve $F$.
The sequence $0\to\Lambda_F\to\OO\to F\to0$ and
$H^1(\Delta,\OO)=H^2(\Delta,\Lambda_F)=0$ trivialize
this actual torsor. The finite translation gauge in
(C.HT.4) therefore gives the surface tube exactly as
\[
 (F\times\Delta_s)/\mu_m,
 \qquad(z,s)\longmapsto(z+\tau,\zeta_ms),
 \qquad\operatorname{ord}(\tau)=m.
\tag{C.QT.2}\label{cdp:eq:C.QT.2}
\]
The last order is exact because the central action is
free. Its associated normal character is primitive.
Thus $N_{J/T_v}$ has exact order $m$, and its connected
normal-root cover is precisely $F\to J$.

\noindent\textbf{Step 2.} The fixed cohomology frame $H^1(F,\OO_F)$ on this
root chart is unchanged by the finite translations.
Taking invariant functions gives the ordinary
coordinate $s^m$. The regular frames glue by the
specified identity monodromy and by translations.
Consequently the actual direct images and their
cohomology are
\[
 q_{v*}\OO_{T_v}=\OO_{\PP^1},\qquad
 R^1q_{v*}\OO_{T_v}=\OO_{\PP^1},\qquad
 (h^0,h^1,h^2)(T_v,\OO)=(1,1,0).
\tag{C.QT.3}\label{cdp:eq:C.QT.3}
\]
At a marked fibre $J$, restriction to $J$ followed
by the finite pullback $F\to J$ is evaluation on this
same cohomology frame. The resulting map
$H^1(T_v,\OO)\to H^1(F,\OO_F)$ is an isomorphism.
In particular the original restriction to $J$ is
nonzero. The commuting square is the actual root
chart \eqref{cdp:eq:C.QT.2}; no identification of different
cohomology groups by dimensions alone is used.

\noindent\textbf{Step 3.} A nonlinear Hopf leaf has the same values $(1,1,0)$
and the nonzero restriction map proved in \ref{cdp:st:C.2.UN.54}.
Duality of boundary normal lines identifies their
exact orders and their connected normal-root covers.
It thus identifies the endpoint Jacobians by actual
finite-cover maps, including the secondary character
of a nonlinear leaf from \ref{cdp:st:C.2.UN.50}.

\par\medskip\noindent\textbf{Graph, residues and the canonical minimum.}
Distinct fibre supports on each $T_v$ are disjoint,
and a nonlinear Hopf has only its one elliptic curve.
Hence the specified SNC surface has no triple points.
For an actual degeneration its full fibre is reduced,
so the meridian quotient of \ref{cdp:st:C.2.UN.44} makes $\pi_1(S)$
abelian. The actual map onto the incidence graph
then makes that graph's free fundamental group
abelian. It has at most one independent cycle.

\noindent\textbf{Step 4.} The normalization sequence is exactly
\[
 0\longrightarrow\OO_S\longrightarrow
   \bigoplus_v\OO_{T_v}\longrightarrow
   \bigoplus_e\OO_{J_e}\longrightarrow0,
\tag{C.QT.4}\label{cdp:eq:C.QT.4}
\]
with nonlinear Hopf terms included in the vertex
sum. The induced $H^1$ matrix is a weighted
incidence matrix whose every endpoint coefficient
is nonzero, by \eqref{cdp:eq:C.QT.3}. A tree gives a surjective
matrix by leaf elimination, and therefore
$H^2(S,\OO_S)=0$. Flatness over the smooth curve
and upper semicontinuity from a torus instead
give $h^2(S,\OO_S)\ge1$. Thus the graph is
unicyclic. Its spanning-tree minor has rank
$\#V-1$; the remaining row must vanish on that
minor's kernel. This is exactly product-one
holonomy for the actual endpoint restriction
maps, and
\[
                    (h^0,h^1,h^2)(S,\OO_S)=(1,2,1).
\tag{C.QT.5}\label{cdp:eq:C.QT.5}
\]

\noindent\textbf{Step 5.} Serre duality identifies the transpose incidence
map with the component residue obstruction maps.
Its kernel has zero residues on every attached
tree, by leaf elimination, and nonzero residues
on every cycle edge. Since every vertex has
$H^0(K_{T_v})=0$, a component section with all
residues zero is zero. The unique global section
of $\omega_S$ is consequently nonzero exactly
on the cycle components. Divide the actual
canonical evaluation by its minimum full-fibre
power. Its nonzero reduction is this section.
Its minimum vertical order $k_0$ is therefore
attained exactly on the cycle, as in the explicit
evaluation comparison of \ref{cdp:st:C.2.UN.54}.

\par\medskip\noindent\textbf{The exact differential and height calculation.}
On a vertex \eqref{cdp:eq:C.QT.1}, invariant fibre differentials
give actual meromorphic two-forms
$q_v^*\theta\wedge dz$. Translation charts
change $dz$ only by a multiple of the base
differential, which disappears in this wedge.
The forms therefore glue. Every meromorphic
two-form with vertical divisor has this form:
its ratio to one reference form is constant
on the regular compact fibres and descends;
\eqref{cdp:eq:C.QT.2} makes its descended meromorphic
function invariant under $s\mapsto\zeta_ms$,
so it is meromorphic in $s^m$ across a multiple
fibre. The local order is exactly
\[
 \operatorname{ord}_{J}(q_v^*\theta\wedge dz)
             =m\operatorname{ord}_{q_v(J)}\theta+m-1.
\tag{C.QT.6}\label{cdp:eq:C.QT.6}
\]

\noindent\textbf{Step 6.} Restriction of the canonical evaluation divided
by $t^{k_v}$ gives the actual meromorphic
component form with divisor
\[
           \sum_{e=(v,w)}(k_w-k_v-1)J_e.
\tag{C.QT.7}\label{cdp:eq:C.QT.7}
\]
An unmarked multiple fibre would have order zero
in \eqref{cdp:eq:C.QT.7}, and \eqref{cdp:eq:C.QT.6} would require
$m\operatorname{ord}\theta+m-1=0$, hence $m=1$.
At a marked fibre these same equations give
\[
 d_{v,e}:=\frac{k_w-k_v}{m_e}
       =\operatorname{ord}_{q_v(J_e)}\theta+1\in\Z,
 \qquad\sum_{e\ni v}d_{v,e}=\operatorname{val}(v)-2.
\tag{C.QT.8}\label{cdp:eq:C.QT.8}
\]
The last equality is the degree $-2$ of an
actual rational differential on $\PP^1$.
Opposite orientations negate $d$ because their
normal-root covers have the same exact order.

\noindent\textbf{Step 7.} At a terminal elliptic vertex \eqref{cdp:eq:C.QT.8} gives
height difference $m_e$ from its parent. A
nonlinear terminal vertex gives precisely the
same difference by its invariant axis form
$dz\wedge dw/w^{m_e+1}$ and the exact secondary
normal comparison in \ref{cdp:st:C.2.UN.50}. Induct inward along
each tree. If the $c$ child terms are one,
\eqref{cdp:eq:C.QT.8} gives $-d_{\rm parent}+c=c-1$,
so the parent term is also one. The cycle
terms are zero by the canonical minimum
already identified. Hence $k_v-k_0$ is exactly
the stated path sum $h_v$.

\noindent\textbf{Step 8.} This proof also gives the actual base differentials:
on a cycle vertex it is $c\,dz/z$, with the two
cycle points $0,\infty$; on a tree vertex it is
$c\,dz$, with its parent point at infinity.
The other marked points have order zero.
No Kodaira-dimension assumption was used in
this calculation.

\par\medskip\noindent\textbf{The common integral period.}
Dual normal-root charts give on an edge
$xy=t$, $q_v=x^{m_e}$ and $q_w=y^{m_e}$,
with the same connected elliptic cover and
the actual normal-product trivialization.
The level $x^{m_e}=c$ corresponds to
$y^{m_e}=t^{m_e}/c$. Varying its nonzero
normal level gives the oriented isotopy
between the two regular-fibre two-tori.
The product-preserving tubular comparison
transfers this equality to the actual
nearby torus, as in (C.HT.10).
At a nonlinear leaf the unimodular boundary
map of \ref{cdp:st:C.2.UN.50} identifies its real core with
this same integral class. All vertices
therefore give one class $\gamma$.

\noindent\textbf{Step 9.} The finite-root restriction square in
\eqref{cdp:eq:C.QT.3} also shows that the coherent
cycle holonomy in \eqref{cdp:eq:C.QT.5} is the linear
action of the actual product of elliptic
root-cover isomorphisms. Its value one
forces that product to be a translation.
This retains its possible translation
and all Jacobian identifications.

\noindent\textbf{Step 10.} The period of the evaluation divided by
$t^{k_0}$ is holomorphic. On a regular
holomorphic fibre curve of $T_v$ its
leading two-form restricts to zero, so
\[
                  P(t)=O(t^{h_v+1}).
\tag{C.QT.9}\label{cdp:eq:C.QT.9}
\]
At a nonlinear leaf of secondary order
$n$ and contraction parameter $\beta$,
the exact primitive core computation is
\[
 P(t)=t^{h_v}
       \left(\frac{\pm2\pi i c\beta^{-m_e}}n+O(t)\right),
                         \qquad c\ne0.
\tag{C.QT.10}\label{cdp:eq:C.QT.10}
\]
Thus every nonlinear leaf has the one
finite order $R$ of this actual common
period, and \eqref{cdp:eq:C.QT.9} makes every elliptic
vertex's height smaller. Without such a
leaf this condition places no further
restriction: an integer above the finite
set of heights exists. This proves all
the height requirements for actual fibres.

\par\medskip\noindent\textbf{Pruning with the full torsors.}
The pruning proof of \ref{cdp:st:C.2.UN.60} applies to the
maps \eqref{cdp:eq:C.QT.1} themselves. We specify the
comparisons that retain the present
broader surface class. Delete a terminal
nonlinear Hopf by the inverse chart
(C.PR.6), with its full primitive period
basis and integer $n$. After the children
of an elliptic vertex have been removed,
its actual modified map still has fixed
Jacobian and identity monodromy over
$\PP^1$; each inverse operation replaced
one multiple fibre by the displayed
ordinary product tube and changed no
other transition. Its only possible
multiple fibre is now its parent mark.

\noindent\textbf{Step 11.} It is therefore the actual $F$-torsor
on $\PP^1[(\infty,m)]$ classified by
(C.PR.1)--(C.PR.5). The exact logarithm
sequence and $H^1(\OO)=H^2(\OO)=0$
identify the entire torsor with its
two root-line classes. A nonzero
vector gives the literal weighted
Hopf quotient, including $n$ and
the primitive fibre direction;
a zero vector forces $m=1$ and
the product ruled cap. This conclusion
concerns the actual modified surface,
even when its original unmodified
surface had Kodaira dimension one.
Remove that pair by its corresponding
map and retain the entire inverse
record.

\noindent\textbf{Step 12.} At the end the two-mark root-stack
cover (C.PR.7) and its actual line
pushout identify each core vertex
with a marked diagonal Hopf surface
or the torsion ruled bundle (C.PR.8).
These are exact maps on the original
transition atlases. Disjoint branch
operations commute on those atlases,
so their marked result has the
independence stated in \ref{cdp:st:C.2.UN.60}.

\noindent\textbf{Step 13.} For an actual degeneration, every
terminal complement and its inverse
product cap have the identical
boundary homomorphism (C.PR.9),
retaining both elliptic generators
and killing exactly the specified
primitive meridian. The van Kampen
pushouts are identical on the
unchanged exterior. Their iteration
identifies the core plumbing group
with the original $\Z^4$. The
injected-lattice semidirect product
for that core consequently has
$M=I$. This proves necessity of
the last test by actual group maps.

\par\medskip\noindent\textbf{Compact restoration of every prescribed component.}
Conversely suppose the test passes.
Use the compact marked core of \ref{cdp:st:C.2.UN.57}
with the power $V_{t^R}$ and the
unchanged logarithmic coefficient
of (C.PR.11). Its necks remain
$xy=t$ and its opposite filling
is the conjugate pruned core.
The compact elliptic subgroup of
its common open group is the
actual regular-fibre group, and
(C.PR.12) makes its transverse
period determinant a nonzero
constant times $t^R$.

\noindent\textbf{Step 14.} Reverse the entire pruning record.
At a vertex of height $h$, the
local quotient (C.PR.13) has
displacement of order $R-h$.
The recorded primitive basis is
implemented by the displayed
unimodular linear period map
following (C.PR.13). At a Hopf
step, the weight-$(m,1)$ insertion
with $\OO(nE)$ has exactly the
free finite chart (C.PR.14).
The new child's coordinate
$U=Z/W^m$ gives the actual
displacement $\delta(t)/t^m$,
as in (C.PR.15). At a product
cap the ordinary blowup has
the exact displacement
$\delta(t)/t$, as in (C.PR.16).
The height test makes all these
maps holomorphic and gives
precisely the prescribed diagonal
or nonlinear central type.

\noindent\textbf{Step 15.} All operations use the entire
recorded inverse transitions.
In particular an original
properly elliptic vertex is
restored by reversing the
successive changes of its
root lines and ordinary tubes;
its complete analytic torsor,
all multiple fibres and all
marked points reappear, rather
than only its degree or normal
characters. The strict-parent
map is exactly (C.PR.6) in the
opposite direction. The same
argument applies to every
other original quasi-bundle.
Fibre translations and the
recorded linear maps retain
each specified elliptic gluing.

\noindent\textbf{Step 16.} The norm and affine-base bounds
(C.PR.17)--(C.PR.18) prove the
actual local properness and
the old-cluster comparison for
each replacement. Their constants
depend on the local displacement
order and contraction, not on the
Kodaira dimension of the surface
being restored. The finitely many
disjoint replacements therefore
give the same Hausdorff compact
gluing as \ref{cdp:st:C.2.UN.60}. The weighted
free coordinates have $t=xW$;
ordinary cap charts have $t=xy$.
Every final component consequently
has full multiplicity one.

\noindent\textbf{Step 17.} Each component \eqref{cdp:eq:C.QT.1} has Euler
number zero by additivity over
its elliptic fibre supports;
each nonlinear Hopf also has
Euler number zero. Every double
curve is elliptic and there are
no triple points. Both special
supports thus have Euler zero,
as do the general tori. This
proves the complete criterion
and the compact assertion.

\par\medskip\noindent\textbf{An explicit properly elliptic component.}
Fix $F=\C^*/\beta^{\Z}$, $0<|\beta|<1$.
Start with the two-component
ruled core $F\times\PP^1$
glued along its two marked
sections by the identity.
Its lattice product is the
identity. On one component
choose four distinct ordinary
fibre points away from those
sections. Perform four recorded
weight-three, secondary-order-one
insertions, with the common
period order $R=3$.
They add four nonlinear primary
Hopf leaves $H_3(\beta)$.

\noindent\textbf{Step 18.} The modified central component
has an explicit independent
description. Put
\[
 \mathcal B=\PP^1[(p_1,3),\ldots,(p_4,3)],\qquad
 L=\OO_{\mathcal B}(D_1+\cdots+D_4),
 \qquad 3D_i=p_i.
\tag{C.QT.11}\label{cdp:eq:C.QT.11}
\]
It is the smooth coarse surface
of $L^\times/\beta^{\Z}$.
Each inertia acts by a primitive
fibre translation, so its coarse
surface is smooth and its four
fibres are exactly triple smooth
elliptic curves. In a root chart
$p=s^3$ its line coordinate is
$y=sW$, where $W$ is the old
product coordinate. The inverse
chart (C.PR.6) with $m=3,n=1$
has precisely $W=y/s$.
Thus this given line pushout
is the component produced by
the four insertions, including
its full torsor.
Explicitly the positive divisor
line has frames $e_{\rm in}=s^{-1}e_{\rm out}$.
The equality of the actual vector
$W e_{\rm out}=y e_{\rm in}$
is precisely $y=sW$; its canonical
section has local coefficient $s$.
This fixes the sign of each $D_i$
in \eqref{cdp:eq:C.QT.11}.

\noindent\textbf{Step 19.} For any map of the class
\eqref{cdp:eq:C.QT.1}, the actual invariant
fibre differential and the
root chart give the plurigenus
formula
\[
 h^0(T_v,\omega_{T_v}^{\otimes a})
 =\max\left\{0,\ 1-2a+
       \sum_i\left\lfloor
                    a(1-1/m_i)\right\rfloor\right\}.
\tag{C.QT.12}\label{cdp:eq:C.QT.12}
\]
Indeed a pluricanonical section
is a base meromorphic
$a$-differential times the
$a$th invariant fibre factor;
holomorphy at a root chart
permits precisely the displayed
pole order. Descent and its
inverse are those actual
invariant-coordinate maps.
For \eqref{cdp:eq:C.QT.11} this gives
\[
 h^0(T,\omega_T^{\otimes a})
    =\max\{0,1-2a+4\lfloor2a/3\rfloor\},
 \qquad \omega_T^{\otimes3}\simeq q^*\OO_{\PP^1}(2).
\tag{C.QT.13}\label{cdp:eq:C.QT.13}
\]
In particular $p_g(T)=0$,
$q(T)=1$, $\chi(\OO_T)=0$
and $h^0(T,\omega_T^{\otimes3b})=2b+1$.
Thus this is an actual properly
elliptic component. The complete
compact reduced fibre consists
of it, the unchanged ruled core
component, and the four Hopf
leaves, with every gluing given
by the preceding maps.
The root-line maps, the incidence
matrix, the root-coordinate pole
bounds and these exact plurigenera
are checked in
\path{code/D_quasibundle_trees.py}
and its saved output.
\end{proof}






\begin{proposition}[The weighted-root criterion for elliptic and Hopf trees]\label{cdp:st:C.2.UN.63}
\StatementRecord{cdp:st:C.2.UN.63}{proved}{Historical events: Claude Fable 5.2 (Anthropic), effort high, 16 November 2026: carried as background without fault found. GPT-6 Astra (OpenAI), ChatGPT Pro mode, effort Pro, S6\_09\_PAPERS, 17 November 2026: targeted proof check. Also checked in this writing. Pre-continuation writing events: GPT-6 Astra (OpenAI), ChatGPT Pro mode, effort Pro, S6\_10\_PAPERS, 18 November 2026: targeted proof check. Continuation metadata events: GPT-6 Astra (OpenAI), ChatGPT Pro mode, effort Pro, S6\_10\_PAPERS continuation, 19 November 2026: used and rechecked in this run; status-vocabulary reconciliation; metadata reconciliation only; no new mathematical check. \emph{Pre-continuation recorded proof scope.} The literal cyclic normalizations and their isogeny transition maps give smooth strict covers and their exact Jacobians. Compressed necks multiply each canonical height increment by h, and the transported integral period has the actual determinant character of the pruned finite action. Signed affine cycle equations retain reflections and all translations; the equivariant twistor frame has an everywhere positive four-period determinant. The finite norm-one and polynomial jet gauges retain the nonlinear resonant translation and make each actual weighted ideal invariant. The pulled tautological line has coefficient nh, and the literal insertion xW=s\textasciicircum{}h divides the displacement by s\textasciicircum{}(hm). Its exact norm and cluster bounds give proper replacements. Character coordinates yield the smooth invariant ring C\{z,u\textasciicircum{}B,v\textasciicircum{}A\}, retain every weighted eigensheaf, and restore the complete original and conjugate supports}{the hypotheses, notation and definitions stated here; \cref{cdp:st:C.2.UN.59}; \cref{cdp:eq:C.WT.7}; \cref{cdp:st:C.2.UN.37}; \cref{cdp:eq:C.WT.10}; \cref{cdp:st:C.2.UN.40}; \cref{cdp:st:C.2.UN.28}; \cref{cdp:eq:C.WT.9}; \cref{cdp:st:C.2.UN.61}; \cref{cdp:st:C.2.UN.62}; \cref{cdp:st:C.2.UN.47}; \cref{cdp:st:C.2.UN.58}; \cref{cdp:st:C.2.UN.50}; \cref{cdp:eq:C.WT.14}; \cref{cdp:eq:C.WT.8}; \cref{cdp:eq:C.WT.15}; \cref{cdp:st:C.2.UN.51}; \cref{cdp:st:C.2.UN.55}; \cref{cdp:eq:C.WT.18}; \cref{cdp:eq:C.WT.19}; \cref{cdp:eq:C.WT.23}; \cref{cdp:st:C.2.UN.57}; \cref{cdp:eq:C.WT.30}; \cref{cdp:eq:C.WT.32}; \cref{cdp:st:C.2.UN.60}; \cref{cdp:eq:C.WT.34}; \cref{cdp:eq:C.WT.31}; \cref{cdp:eq:C.WT.28}; \cref{cdp:eq:C.WT.36}; \cref{cdp:eq:C.WT.33}; \cref{cdp:eq:C.WT.38}; \cref{cdp:eq:C.WT.42}; \cref{cdp:eq:C.WT.25}}

Specify a connected SNC surface with embedded smooth components,
edges between distinct components, and positive component
multiplicities \(m_i\). Each component is a primary or secondary
Hopf surface, or is supplied with its actual fixed-Jacobian,
identity-integral-monodromy elliptic quasi-bundle over a smooth
compact curve, with connected fibres and only smooth multiple
supports. At a supplied quasi-bundle every marked double curve
is a fibre support. Retain every component, analytic torsor and
elliptic gluing map.

Specify the actual weighted line data
\[
 \begin{gathered}
 \sigma_i\in H^0(S_i,\mathcal L_i^{-m_i}),\qquad
 \operatorname{div}\sigma_i=\sum_{j\sim i}m_jJ_{ij},\\
 \mathcal L_i|_{J_{ij}}\simeq
       \phi_{ij}^{*}N_{J_{ji}/S_j},\qquad
 \phi_{ij}^{*}\mathcal L_j|_{J_{ji}}
       \simeq N_{J_{ij}/S_i},
 \end{gathered}
 \tag{C.WT.1}\label{cdp:eq:C.WT.1}
\]
with the equal leading-coefficient line maps at each edge.
Form the finite normalizations
\[
 \mathcal E_i=\operatorname{Norm}
       \{\xi\in\mathcal L_i^{-1}:\xi^{m_i}=\sigma_i\},
 \qquad N=\operatorname{lcm}_i m_i,
 \tag{C.WT.2}\label{cdp:eq:C.WT.2}
\]
their actual elliptic boundary maps and their cyclic
\(\mu_N\)-action. At an edge of multiplicities \(a=gA,b=gB\),
where \(\gcd(A,B)=1\), retain the literal normalized neck
\[
       uv=s^h,\qquad x=u^B,\quad y=v^A,\qquad
       h=\frac{N}{gAB},\qquad t=s^N=x^a y^b.
 \tag{C.WT.3}\label{cdp:eq:C.WT.3}
\]
These are compressed necks; no unspecified contraction of a
reconstructed ruled chain is part of the test.

The following test is necessary and sufficient for the given
weighted support to occur as a fibre of a smooth-total
degeneration of connected complex two-tori.

\begin{enumerate}
\item The root graph is connected. The branch containing an
algebraic-dimension-zero diagonal Hopf component is the complete
two-axis cycle and passes \ref{cdp:st:C.2.UN.59}. In every remaining branch,
each strict elliptic root cover has base genus at most one, and
at most one such cover has genus one.
\item If all elliptic root bases have genus zero, the compressed
root graph is unicyclic. If there is a genus-one root component,
the graph is a tree rooted there. All other root components are
the supplied genus-zero elliptic quasi-bundles or nonlinear
Hopf leaves. There are no unmarked multiple elliptic fibres.
The actual normal lines at each root edge are dual.
\item Let \(m_e\) be their common exact torsion order, and let
\(h_e\) be the exponent in \eqref{cdp:eq:C.WT.3}. Give the cycle, or the
genus-one root, height zero, and set
\[
       H_v=\sum_{e\text{ on the rooted path to }v}h_e m_e.
 \tag{C.WT.4}\label{cdp:eq:C.WT.4}
\]
All nonlinear leaves have one common height \(R>0\);
every elliptic vertex has height strictly less than \(R\).
When there is no nonlinear leaf, \(R\) may be chosen larger
than all the heights.
\item Perform the exact inverse operations below in whole
cyclic-group orbits, retaining the complete torsors, the full
secondary integers, primitive fibre bases and all marked maps.
The final cycle has actual rank-three complement product
\(M=I\), or the final principal elliptic torsor over its
elliptic base has integral Chern vector zero.
\item Let \(D\) be the finite linear part on this actual pruned
core. For a cycle, it is the linear part on its common
two-dimensional logarithmic vector space; for a torus core,
it is the linear part of its actual affine torus action.
For the root generator with base action \(s\mapsto\zeta_Ns\),
the height satisfies
\[
                         \zeta_N^R=\det_{\C}D.
 \tag{C.WT.5}\label{cdp:eq:C.WT.5}
\]
In the absence of nonlinear leaves one chooses any sufficiently
large positive \(R\) in this residue class.
\end{enumerate}

Every step uses finite covers, actual line and torsor maps,
integral lattices, and the displayed inverse charts. In
particular \eqref{cdp:eq:C.WT.5} is an equality of the actual canonical
character. A cycle action may reverse orientation through
vertices; such reflections are included.

Every support passing the test has a realization over
\(\PP^1\) with smooth compact total threefold. The construction
retains its original components, gluing maps and all
multiplicities. It can have exactly the prescribed weighted
support at zero and its complex conjugate at infinity, with
all other fibres smooth connected two-tori. Every reduced
fibre support has Euler number zero. Normal-product constants
are chosen in the realization; no equivalence of different
full analytic neighbourhoods is asserted.
\end{proposition}



\begin{proof}
\noindent\textbf{Step 1.} \par\medskip\noindent\textbf{The actual component covers and their local inertia.}
An actual map supplies \eqref{cdp:eq:C.WT.1} by dividing its defining
function by the \(m_i\)-th power of the canonical divisor
section and restricting to \(S_i\). On its strict root
component the exact formula is
\(\xi=s^{N/m_i}/v_i\), where \(v_i\) is that normal
divisor coordinate. Thus the root generator is
\[
                  \xi\longmapsto
                         \zeta_N^{N/m_i}\xi.
 \tag{C.WT.6}\label{cdp:eq:C.WT.6}
\]
At a double curve, absorbing a unit in one normal
coordinate gives \(x^a y^b=s^N\). Its \(g\) local
branches have precisely the finite birational
normal maps \eqref{cdp:eq:C.WT.3}, as proved in (C.WH.4).
The equal leading-coefficient maps in \eqref{cdp:eq:C.WT.1}
identify their actual curve covers and normal
characters. This constructs the compressed root
data without replacing a nonlinear Hopf
neighbourhood by its normal line.

\noindent\textbf{Step 2.} Every strict component cover is smooth. Locally along
a marked curve its equation is
\(\xi^m=z^b\) times a unit. Write
\(d=\gcd(m,b)\), \(q=m/d\), \(B=b/d\).
On each branch its normalization is
\[
                         z=w^q,\qquad \xi=cw^B,
 \tag{C.WT.7}\label{cdp:eq:C.WT.7}
\]
with \(c\) a fixed root of unity. This is a smooth
coordinate chart with the unchanged curve
coordinate. Away from these curves the map is
finite \'{e}tale. The induced finite maps of
reduced boundary curves are therefore \'{e}tale.
Every stabilizer preserving such a connected
elliptic curve acts on it by deck translations.
The same is true on its connected normal-root
elliptic cover.

\noindent\textbf{Step 3.} For a supplied elliptic component, remove the finitely
many marked base values. A connected component of
its finite surface cover has a fixed finite-index
fibre subgroup \(\Lambda'\subset\Lambda_F\).
Take the actual connected-fibre Stein base cover.
On it the fibre is the fixed elliptic curve
\(F'=\C/\Lambda'\), with the actual isogeny
\(\pi:F'\to F\). A lifted transition \(a\)
satisfies
\[
                            \pi a=T_b\pi.
 \tag{C.WT.8}\label{cdp:eq:C.WT.8}
\]
Its derivative is therefore the identity, so it
is a translation. This proves identity integral
monodromy on the connected-fibre base cover,
including loops which permuted the original
components of a fibre cover.

\noindent\textbf{Step 4.} Near a deleted value lift the actual chart
\((F\times\Delta)/\langle(\tau,\zeta_m)\rangle\)
and apply \eqref{cdp:eq:C.WT.7} to its normal power. Quotient
first by the subgroup trivial on \(F'\); it
acts by normal pseudoreflections. The remaining
inertia acts freely by translations on \(F'\).
The result is again a smooth multiple-fibre
chart
\((F'\times\Delta)/\langle(\tau',\zeta_{m'})\rangle\).
The finite translation gauge of (C.HT.4) makes
\(\tau'\) constant. Thus the compact strict cover
has an actual fixed-Jacobian quasi-bundle over
the smooth compact Stein base \(B'\), with only
smooth multiple supports. Its invariant
cohomology frame gives
\[
                  R^1q_*\OO=\OO_{B'},\qquad
                         p_g=g(B').
 \tag{C.WT.9}\label{cdp:eq:C.WT.9}
\]
This argument also applies to every
algebraic-dimension-one Hopf component, using
its actual fibration from \ref{cdp:st:C.2.UN.37}.

\noindent\textbf{Step 5.} For a nonlinear Hopf component there is a
different explicit map. On its primary cover
write the contraction
\[
                    (z,w)\longmapsto
                     (\beta^r z+c w^r,\beta w),
                 \qquad c\ne0 .
 \tag{C.WT.10}\label{cdp:eq:C.WT.10}
\]
Its only curve is \(w=0\). The power of the
weighted normal line is the constant-automorphy
axis line; dividing its transferred connection
gives a flat connection on the line itself.
On \(\C^2\setminus0\) the root section is a
constant times \(w^b\). Indeed its quotient
by \(w^b\) extends as a holomorphic unit at
zero; its automorphy factor is one by
evaluation there, and contraction makes that
unit constant. Formula \eqref{cdp:eq:C.WT.7} is consequently
the actual map \((z,v)\mapsto(z,v^q)\).
A positive power of \eqref{cdp:eq:C.WT.10} preserving one
branch lifts to
\[
                 (z,v)\longmapsto
                    (\alpha^{qr}z+c'v^{qr},\alpha v),
                     \qquad c'\ne0 .
 \tag{C.WT.11}\label{cdp:eq:C.WT.11}
\]
Its branch stabilizer is free on
\(\C^2\setminus0\). A fixed point would project
to a fixed point of the original free deck
group; an element over its identity fixing
both \(v^q\) and \(v^B\) fixes \(v\), since
\(\gcd(q,B)=1\). Hence the finite normalization
is exactly a primary or secondary nonlinear
Hopf surface, with the displayed map and
resonance \(qr\).

\noindent\textbf{Step 6.} The nontorsion axis normals of a diagonal
algebraic-dimension-zero Hopf cannot satisfy
the weighted normal equality with a torsion
normal. Its axis covers remain nonresonant:
a multiplicative relation after taking finite
coordinate roots would give a multiplicative
relation for the original multipliers.
\ref{cdp:st:C.2.UN.40} forces both its axes into the boundary.
Thus that connected branch is the complete
cycle already covered by \ref{cdp:st:C.2.UN.59}. We henceforth
treat the torsion-normal branch.

\noindent\textbf{Step 7.} Resolve the \(A_{h-1}\) necks only to apply
the compact-core inequality \ref{cdp:st:C.2.UN.28}. It gives
\(\sum p_g\le1\) for all strict components.
By \eqref{cdp:eq:C.WT.9}, their base genera are at most one
and at most one is one. No assertion that
a many-marked Hopf root cover is still Hopf
has been made.

\par\medskip\noindent\textbf{Compressed canonical heights and the actual period character.}
The central fibre of \(uv=s^h\) is reduced,
and its abstract two-component crossing is
SNC. Resolving the neck inserts the actual
elliptic ruled chain, with the same normal
root maps. Either description gives the
same component and edge cohomology maps.
The exact restriction square (C.QT.2)--(C.QT.3)
and the normalization and residue arguments
of \ref{cdp:st:C.2.UN.61}--\ref{cdp:st:C.2.UN.62} therefore give the following.
With no genus-one vertex the root graph is
unicyclic, and the minimum canonical order
is attained exactly on its cycle. With one
genus-one vertex the graph is a tree and
that vertex is the unique minimum. One can
also apply \ref{cdp:st:C.2.UN.47} to the latter assertion.
The connectedness here follows from the
proper Stein factor of the actual root
family, whose nearby fibres are connected.

\noindent\textbf{Step 8.} Here is the precise change caused by a
compressed neck. Use a curve coordinate
\(z\) and the relative canonical frame
\(du/u\wedge dz=-dv/v\wedge dz\).
If the leading form on the \(u\)-side has
curve order \(d\), then
\[
 s^k u^{d+1}\frac{du}{u}\wedge dz
      =-s^{k+h(d+1)}v^{-d-1}
                         \frac{dv}{v}\wedge dz .
 \tag{C.WT.12}\label{cdp:eq:C.WT.12}
\]
Thus the canonical-order difference is
\(h(d+1)\), rather than \(d+1\).
At an elliptic multiple fibre of order \(m\),
the exact wedge order is
\(d=m\operatorname{ord}\theta+m-1\).
Consequently the edge integer is
\[
       \frac{k_w-k_v}{h_e m_e}
                        =\operatorname{ord}_{p_e}\theta+1.
 \tag{C.WT.13}\label{cdp:eq:C.WT.13}
\]
At an unmarked multiple fibre the required
order is zero; the formula excludes every
unmarked \(m>1\). The degree of the actual
base differential gives the same
valence-minus-two equation as \ref{cdp:st:C.2.UN.58}.
The finite leaf induction now has outward
increment \(h_e m_e\). This proves exactly
\eqref{cdp:eq:C.WT.4}, including product caps with
normal order one. Nonlinear terminal
forms give the same increment, by \eqref{cdp:eq:C.WT.10}.

\noindent\textbf{Step 9.} The root collar sends the connected cover
\(u^{m_e}=c\) to
\(v^{m_e}=s^{h_e m_e}/c\). These are the
same oriented integral two-cycle. Together
with the full secondary boundary matrix
of \ref{cdp:st:C.2.UN.50}, they identify every regular
elliptic fibre and every nonlinear core
with one integral class \(\gamma_s\).
The canonical period has order greater
than \(H_v\) on an elliptic vertex, and
has order exactly \(H_v\), with the
nonzero coefficient of (C.PT.10), at a
nonlinear leaf. This proves the height
test even on an arbitrary actual
smoothing.

\noindent\textbf{Step 10.} Choose a finite-equivariant relative
canonical generator: a nonconstant unit
in its finite character is removed by
the finite averaging of its logarithm.
Divide by the minimum power of \(s\).
The inverse charts used below carry this
form to the nonzero canonical frame on
the pruned core. Its finite character
there is exactly \(\chi=\det_{\C}D\).
The deck maps preserve the oriented
regular-fibre class, so
\[
 P(\zeta_Ns)=\chi P(s),\qquad
 P(s)=\int_{\gamma_s}\Omega_s .
 \tag{C.WT.14}\label{cdp:eq:C.WT.14}
\]
If a nonlinear leaf is present, its
nonzero leading coefficient proves
\(\zeta_N^R=\chi\). Without such a leaf
one may choose a larger positive exponent
in this character class. This proves
\eqref{cdp:eq:C.WT.5} with the actual minimum-form
character, without inserting an
unproved parity assumption.

\par\medskip\noindent\textbf{Equivariant inverse pruning and the exact core obstruction.}
Use the one-mark root-torsor classification
(C.PR.1)--(C.PR.8), or its identical application
to the supplied quasi-bundles in \ref{cdp:st:C.2.UN.61}--\ref{cdp:st:C.2.UN.62}.
After all children of a vertex have been
removed, its only possible multiple fibre
is its parent mark. Its actual torsor is
therefore a recorded diagonal Hopf pair,
with its full integer \(n\), or the product
ruled cap. At a Hopf step choose
\(mk-nl=1\). The actual inverse maps are
\[
 \begin{gathered}
 (u,X)\sim(\gamma^m u,\gamma^lX),\qquad
 p=X^m u^{-l},\qquad W=X^{-n}u^k,\\
                       u=p^nW^m,\qquad X=p^kW^l.
 \end{gathered}
 \tag{C.WT.15}\label{cdp:eq:C.WT.15}
\]
They identify the punctured parent tube
with the specified product tube. The
parameter \(p\) is its actual base
coordinate, since its pullback is the
displayed multiple-fibre function.
The tensor and torsor maps therefore
extend across the new ordinary fibre,
with all other base charts unchanged.
The identity
\[
             d\log p\wedge d\log W
                         =d\log X\wedge d\log u
 \tag{C.WT.16}\label{cdp:eq:C.WT.16}
\]
also proves the canonical-frame comparison
used in \eqref{cdp:eq:C.WT.14}: the old order \(m-1\)
becomes the nonzero frame \(dp\wedge
d\log W\). The normal-product neck
exponent \(h\) changes the order in the
family parameter, but does not change
this central inverse map.

\noindent\textbf{Step 11.} Carry out these operations simultaneously
on whole finite-group orbits. A stabilizer
of a connected boundary curve acts there
by a deck translation, as proved after
\eqref{cdp:eq:C.WT.7}. On the connected normal-root
elliptic cover its linear part is also
one, by the isogeny derivative argument
of \eqref{cdp:eq:C.WT.8}. It consequently fixes the
entire primitive fibre lattice. Choose
one primitive basis and the full integer
\(n\) on each orbit and transport it.
The maps \eqref{cdp:eq:C.WT.15}, their canonical
divisor sections and their inverses
then give the finite maps on the
replacement tubes. Their group law
holds on the punctured tubes and hence
on their holomorphic extensions.
This retains the stabilizer characters;
no independent choice of roots of
transition multipliers is made.
Disjoint replacements commute by their
actual atlases, so a group-invariant
postorder gives the required exact
equivariant pruning record.

\noindent\textbf{Step 12.} The root action cannot invert an edge:
its two endpoints lie over two distinct
original components, each fixed by the
root generator. It can, however, reflect
a cycle through vertices, interchanging
two different preimage curves over one
original boundary. Tree parent
orientations are preserved.

\noindent\textbf{Step 13.} For necessity no analytic contraction
of the given family is used.
At a fixed \(s\ne0\) the neck
\(uv=s^h\) has the same boundary
homomorphisms as \(uv=s\).
The full secondary matrix in \ref{cdp:st:C.2.UN.50}
shows that removal of the compact
leaf piece kills precisely the
primitive meridian \(\delta\) and
retains the two fibre generators.
The inverse product filling has
those identical maps. The actual
Van Kampen pushouts of (C.PR.9)
therefore preserve the nearby
fundamental group throughout.
For a cycle the final injected
semidirect product is
\(\Lambda\rtimes_M\Z\); its equality
with \(\Z^4\) forces \(M=I\).
For a genus-one core the injected
fibre lattice in (C.PT.12) has
commutator its integral Chern
vector, so that vector is zero.
The explicit Cech lattice
(C.PT.8) identifies this very
principal torsor with a complex
two-torus. This proves the fourth
test while retaining the full
analytic class.

\noindent\textbf{Step 14.} In the cycle case the exact two-mark
classification gives marked diagonal
Hopf or degree-zero ruled components.
Their regular fibres are the common
compact elliptic subgroup of their
open group \(\C^2/\Lambda\).
The finite maps on that group are
affine. For a diagonal Hopf pair,
an axis-preserving lift has form
\((zU,wV)\); contraction makes
the units \(U,V\) constant, exactly
as in \ref{cdp:st:C.2.UN.59}. An axis-interchanging
lift is the composition of such
a map with the coordinate
exchange. Marked ruled maps have
the actual affine logarithmic
lifts of \ref{cdp:st:C.2.UN.51}. Hence the finite
linear part \(D\) is a single
map on the common vector space.
For a torus core its derivative
is constant by periodicity,
so its actual finite action is
affine as well.

\par\medskip\noindent\textbf{Signed cycle equations, including reflections.}
Orient the pruned cycle and let \(n_j\)
be the outgoing normal meridian,
\(h_j\) its retained neck exponent.
Let \(\epsilon=1\) or \(-1\) be the
orientation character of the root
generator, and \(\sigma\) its
permutation of vertices. The image
edge is \(k=\sigma(j)\) for
\(\epsilon=1\), and
\(k=\sigma(j+1)\) for \(\epsilon=-1\).
In actual logarithmic coordinates
write
\[
 \begin{gathered}
 g(j,v,s)=(\sigma j,Dv+\tau_j,\zeta_Ns),\\
 e_j(v,s)=v+b_j+l_j\log s,\qquad
       l_j=-\frac{h_jn_j}{2\pi i},\\
 Dn_j=\epsilon n_k,\qquad
 Dl_j=\epsilon l_k.
 \end{gathered}
 \tag{C.WT.17}\label{cdp:eq:C.WT.17}
\]
For a reverse edge the comparison
is with \(e_k^{-1}\), not \(e_k\).
Taking \(\log\zeta_N=2\pi i/N\),
the actual compositions give
\[
 \begin{aligned}
 \tau_{j+1}+Db_j-\tau_j
       -\epsilon b_k-\epsilon l_k\,2\pi i/N
                                      &\in\Lambda,\\
 \sum_{a=0}^{N-1}D^{N-1-a}
                       \tau_{\sigma^a j}&\in\Lambda.
 \end{aligned}
 \tag{C.WT.18}\label{cdp:eq:C.WT.18}
\]
For example when \(\epsilon=-1\),
\(g_{j+1}e_j=e_k^{-1}g_j\) gives
the first line directly. These
equalities retain every affine
constant of the actual central
action.

\noindent\textbf{Step 15.} Put \(b=\sum b_j\) and
\(u=-\sum h_jn_j=dp\), with \(p\)
primitive and \(d>0\). The exact
contraction-height calculation
(C.AD.3) gives
\(\ell(in_j)\) of one nonzero
sign, where
\(\ker\ell=\operatorname{span}_{\R}\Lambda\).
For ruled vertices this is the
actual elliptic quotient
calculation of \ref{cdp:st:C.2.UN.55}.
Thus \(\ell(iu)\ne0\) and
\(Du=\epsilon u\).
Choose the common elliptic
line and \(\C p\) as coordinates.
Then
\[
 \begin{gathered}
 W=\operatorname{span}_{\R}\Lambda=\C\times\R,\quad
 p=(0,1),\quad D=\operatorname{diag}(\rho,\epsilon),\\
 \Lambda=\langle(a_1,c_1),(a_2,c_2),p\rangle_{\Z},
 \qquad c_1,c_2\in\R,\qquad
        \operatorname{Im}(\overline{a_1}a_2)>0,\\
              (D-\epsilon I)b-\epsilon u/N\in\Lambda.
 \end{gathered}
 \tag{C.WT.19}\label{cdp:eq:C.WT.19}
\]
The last line is the sum of
\eqref{cdp:eq:C.WT.18}. The finite elliptic
linear part \(\rho\) has order
\(1,2,3,4\) or \(6\), by its
integral trace, as in \ref{cdp:st:C.2.UN.59}.
In particular
\(\det D=\rho\epsilon\), including
the sign at a reflected core.

\noindent\textbf{Step 16.} Adding \(c n_j\) to \(b_j\) on
a whole edge orbit, with the
same \(c\in\C\), preserves
\eqref{cdp:eq:C.WT.18}, since
\(Dn_j=\epsilon n_k\).
The second coordinates of all
the \(n_j\) have the same
nonzero sign. Such an orbit
therefore allows the second
coordinate of \(b\) to be
set to \(i\). This changes
only a normal-product constant:
the translation along the
normal meridian restricts
to zero on the double curve.
It preserves the specified
elliptic gluing maps and the
entire finite action.

\par\medskip\noindent\textbf{The finite-equivariant compact core.}
For a cycle use the coordinates
in \eqref{cdp:eq:C.WT.19}; for a torus choose
a finite-invariant Hermitian
metric and an orthonormal
complex basis. In either case
\(D\) is unitary and
\[
 D J\overline D^{-1}=(\det D)J,
       \qquad J=\begin{pmatrix}0&1\\-1&0\end{pmatrix}.
 \tag{C.WT.20}\label{cdp:eq:C.WT.20}
\]
This follows by multiplying
the two-by-two matrices, or
from the alternating form
on a two-dimensional space.
With \(R\) as in \eqref{cdp:eq:C.WT.5} set
\[
 V_s(w)=w+s^R J\overline w,\qquad
             D V_s=V_{\zeta_Ns}D .
 \tag{C.WT.21}\label{cdp:eq:C.WT.21}
\]
For a torus core the four
real period vectors \(V_s\Gamma\)
remain independent because
\(V_s\) has real determinant
\((1+|s|^{2R})^2\).
Transport an affine constant
to the target frame:
\(\tau(s)=V_{\zeta_Ns}\tau\).
All finite lattice and order
equalities are then preserved
exactly, so these are an actual
finite action on the torus
family.

\noindent\textbf{Step 17.} For a cycle transport
\[
 b_j(s)=V_s b_j,\qquad
 \tau_j(s)=V_{\zeta_Ns}\tau_j,\qquad
 B_s=V_s b+\frac{d\log s}{2\pi i}V_s p .
 \tag{C.WT.22}\label{cdp:eq:C.WT.22}
\]
The correction
\(l_j\,2\pi i/N=-h_jn_j/N\)
in \eqref{cdp:eq:C.WT.18} is a real-rational
period vector. Applying \(V_s\)
therefore carries the signed
equations to the identical
equations in the target frame.
Summing gives
\[
                 D B_s-\epsilon B_{\zeta_Ns}
                               \in V_{\zeta_Ns}\Lambda .
 \tag{C.WT.23}\label{cdp:eq:C.WT.23}
\]
The finite action consequently
preserves the actual rank-four
nearby period lattice. The
global map itself can be
written on the actual middle
charts. Choose a cut and put
\[
 \begin{gathered}
 \theta_j(s)=\sum_{a<j}
      \bigl(b_a(s)+l_a(s)\log s\bigr),\qquad
 \eta=v_j-\theta_j(s),\\
 \eta\longmapsto
 D\eta+D\theta_j(s)+\tau_j(s)
                         -\theta_{\sigma j}(\zeta_Ns),
 \qquad l_a(s)=-h_aV_s(n_a)/(2\pi i).
 \end{gathered}
 \tag{C.WT.23a}\label{cdp:eq:C.WT.23a}
\]
Changing an adjacent chart
changes this expression by
the actual lattice element
in \eqref{cdp:eq:C.WT.18}. Changing the
cut adds \(B_s\) on the source
and \(\epsilon B_{\zeta_Ns}\)
on the target, so \eqref{cdp:eq:C.WT.23}
proves that the expression
is unchanged as a map of
the middle torus. A logarithm
change adds integral normal
periods. The second equation
of \eqref{cdp:eq:C.WT.18} proves its
\(N\)-th-order relation.
Thus this is an actual
global finite holomorphic
action, including reflections,
not only an invariant lattice.
The
normal maps are the literal
\(uv=s^{h_j}\) maps; reversing
an edge exchanges their two
coordinates. The moving
Hopf and ruled frame maps
of \ref{cdp:st:C.2.UN.57}--\ref{cdp:st:C.2.UN.59} give these
actual charts also for a
reflection through a vertex.

\noindent\textbf{Step 18.} The core is compact over the
whole parameter sphere. In
the cycle case its determinant
is a positive constant times
\[
 (1+r^{2R})
 \left[1+r^{2R}
    +\frac d{2\pi}(r^{2R}-1)\log r\right],
                \qquad r=|s|>0 .
 \tag{C.WT.24}\label{cdp:eq:C.WT.24}
\]
Each factor is positive.
This is the exact computation
of \ref{cdp:st:C.2.UN.57} with \(V_s\) replaced
by \(V_{s^R}\), keeping
\(d\log s\) unchanged.
It gives compact connected
two-tori on every punctured
fibre. Finite core and neck
collars give the proper
filling at zero. No torus
rank is inferred merely
from the fundamental group.

\noindent\textbf{Step 19.} At infinity put \(z=s^{-1}\)
and divide the vector frame
by \(s^R\). The periods become
\[
              W_z(w)=J\overline w+z^Rw,\qquad
              D_{\infty}=(\det D)^{-1}D
                       =J\overline D J^{-1}.
 \tag{C.WT.25}\label{cdp:eq:C.WT.25}
\]
The base character is
\(\zeta_N^{-1}\). An affine
constant becomes
\(J\overline\tau+(\det D)^{-1}z^R\tau
=W_{\zeta_N^{-1}z}(\tau)\).
The normal meridians are
\(-J\overline{n_j}\); changing
\(\log s\) to \(-\log z\)
gives the conjugate core
with the same positive neck
exponents. Thus the infinity
maps, including every affine
constant and reflection, are
the actual complex conjugates
of the central core maps.
This provides the second
proper compact end.

\par\medskip\noindent\textbf{The precise finite jet gauge at a surgery orbit.}
The common regular elliptic fibre has a
primitive period basis \(\lambda,\mu=c\lambda\),
\(\operatorname{Im}c>0\). In the core frame
its exact determinant is
\[
 \det_{\C}(V_s\lambda,V_s\mu)
              =2i\operatorname{Im}(c)\|\lambda\|^2s^R .
 \tag{C.WT.26}\label{cdp:eq:C.WT.26}
\]
Consequently the actual rank-two period
cover near a marked fibre has the
semi-torus presentation
\[
 (p,w,s)\longmapsto(p+\delta(s),\gamma(s)w,s),
 \qquad \delta(s)=s^{\ell} d(s),\quad d(0)\ne0,
                     \quad|\gamma(0)|<1 .
 \tag{C.WT.27}\label{cdp:eq:C.WT.27}
\]
Initially \(\ell=R\). The exact primitive
basis change (C.PR.13) multiplies \(\delta\)
by a holomorphic unit and retains its
order. Every later chart below has this
same presentation.

\noindent\textbf{Step 20.} Consider the stabilizer of one surgery
curve, of order \(d_0\), with base
character \(\xi\). This is the stabilizer
order, not the order of the whole
graph group when that curve moves.
As proved after \eqref{cdp:eq:C.WT.15}, its linear
action on the actual fibre lattice
is the identity. On the universal
logarithmic cover its finite lift
therefore has an upper-triangular
affine form
\[
 (p,V,s)\longmapsto
    (a(s)p+b(s),\,V+c(s)p+e(s),\,\xi s),
                     \qquad V=\log w .
 \tag{C.WT.28}\label{cdp:eq:C.WT.28}
\]
The lift may have \(d_0\)-th power
equal to the \(k\)-th infinite deck
translation, with an integral
angular deck translation as well.
Both integers are retained from the
actual root action. The \(p\)-part
of that equality is
\[
 \prod_{j=0}^{d_0-1}a(\xi^j s)=1,\qquad
                 g^{d_0}(p)=p+k\delta(s).
 \tag{C.WT.29}\label{cdp:eq:C.WT.29}
\]
The norm-one unit cocycle for \(a\)
is made constant \(a_0=a(0)\) by
a holomorphic unit gauge, equal
to one at zero. Explicitly take
a logarithm of \(a/a_0\), whose
finite orbit sum is zero, and
use the same finite weighted
sum as in (C.HT.4). Its
coboundary is that logarithm.
Commutation with the infinite
period then gives
\[
 \delta(\xi s)=a_0\delta(s),\qquad
 \sum_{j=0}^{d_0-1}a_0^{d_0-1-j}b(\xi^j s)
                                      =k\delta(s).
 \tag{C.WT.30}\label{cdp:eq:C.WT.30}
\]

\noindent\textbf{Step 21.} For any coefficient \(b_j s^j\)
of order \(j<\ell\), the sum
in \eqref{cdp:eq:C.WT.30} is zero unless
\(\xi^j=a_0\); in that resonant
case it is
\(d_0 a_0^{d_0-1}b_j s^j\),
so \(b_j=0\). Every nonresonant
coefficient is removed by
\(p\mapsto p+c_j s^j\), since
its change is
\((\xi^j-a_0)c_j s^j\).
Thus, whenever \(\ell\ge hm\),
one makes
\[
                               b(s)=O(s^{hm})
 \tag{C.WT.31}\label{cdp:eq:C.WT.31}
\]
by an actual finite polynomial
jet gauge. This preserves the
central marked point and every
resonant coefficient. For an
intermediate diagonal Hopf
component a constant nontrivial
affine action is first centered
at its finite fixed point;
this is its actual triangular
Hopf automorphism and is
recorded on all remaining
marks.

\noindent\textbf{Step 22.} In particular the nonlinear
translation is never discarded.
On a nonlinear Hopf universal
cover put \(U=z/w^m,V=\log w\).
Its infinite generator is
\((U,V)\mapsto(U+c,V+\log\beta)\),
with \(c\ne0\). A parent-axis-
preserving lift commuting with
a positive central contraction
has \(U\mapsto U+b\):
an affine coefficient \(a\)
must satisfy \(ac=c\).
If its \(d_0\)-th power is the
\(k\)-th infinite generator
times a finite angular element,
that element fixes \(U\), so
\[
                                d_0b=kc.
 \tag{C.WT.32}\label{cdp:eq:C.WT.32}
\]
This is exactly the term which
survives \eqref{cdp:eq:C.WT.30} at a terminal
insertion. Indeed if \(\ell=hm\),
then \(a_0=\xi^{\ell}\) and the
leading coefficient there is
\(b_{\ell}=k a_0d(0)/d_0\).
After dividing by \((\xi s)^{hm}\)
it is \(k d(0)/d_0\), as required
by \eqref{cdp:eq:C.WT.32}. For a diagonal
component the corresponding
equation is
\(b(1+a+\cdots+a^{d_0-1})=0\);
it forces \(b=0\) when \(a=1\)
and permits exactly the finite
fixed-point gauge otherwise.
This proves the full central
triangular comparison.

\par\medskip\noindent\textbf{The literal compressed weighted insertion.}
For a recorded Hopf step of
normal order \(m\), full secondary
integer \(n\), and neck exponent
\(h\), use the actual normalized
pullback, by \(t=s^h\), of the
weight-\((m,1)\) base modification
in \ref{cdp:st:C.2.UN.60}. Its root chart is
\[
       p=x^m Z,\qquad xW=s^h,\qquad (Z,W)\ne(0,0).
 \tag{C.WT.33}\label{cdp:eq:C.WT.33}
\]
Retain the pullback of the old
tautological line \(\OO(nE_{\rm old})\).
Its canonical local section is
\(x^n\); since \(x=s^h/W\) at a
generic point of the new component,
\(\operatorname{div}(x)=hE_{\rm new}\).
Thus the pulled line has coefficient
\(nh\), not \(n\), on this generic
divisor. Replacing it by
\(\OO(nE_{\rm new})\) would change
the actual root action and is not
performed here.

\noindent\textbf{Step 23.} Choose the recorded \(\beta(s)\)
with \(\beta(s)^n=\gamma(s)\).
The nonzero line has the free
finite chart
\[
       (Z,W,x)\longmapsto
                   (\eta^mZ,\eta W,\eta^{-1}x),
                           \qquad\eta\in\mu_n.
 \tag{C.WT.34}\label{cdp:eq:C.WT.34}
\]
It is free because
\(\gcd(m,n)=1\), including
where one of \(Z,W\) vanishes.
The infinite deck map is exactly
\[
 \begin{split}
 (Z,W,x,s)\longmapsto\bigl(
   &\beta(s)^m[Z+d(s)s^{\ell-hm}W^m],\\
   &\beta(s)W,\ \beta(s)^{-1}x,\ s\bigr).
 \end{split}
 \tag{C.WT.35}\label{cdp:eq:C.WT.35}
\]
It preserves \(xW=s^h\) and
changes \(p=x^mZ\) by precisely
\(d(s)s^{\ell}=\delta(s)\).
Every term has the same finite
character as \(Z\) in
\eqref{cdp:eq:C.WT.34}. The central child
is nonlinear exactly when
\(\ell=hm\), and diagonal
when \(\ell>hm\). Its recorded
nonzero nonlinear coefficient
is matched by the actual
rescaling of \(Z\); the finite
translation in \eqref{cdp:eq:C.WT.32}
scales with it.

\noindent\textbf{Step 24.} Away from its parent curve put
\(U=Z/W^m\) and use the finite
coarse fibre coordinate \(W^n\).
The exact new deck is
\[
            (U,W^n)\longmapsto
            (U+\delta(s)/s^{hm},\,\gamma(s)W^n).
 \tag{C.WT.36}\label{cdp:eq:C.WT.36}
\]
This is the required remaining
order \(\ell-hm\), at every
other marked child fibre.
There is no change of the
given analytic torsor outside
the one recorded tube.

\noindent\textbf{Step 25.} The finite action also lifts
exactly. The gauge \eqref{cdp:eq:C.WT.31}
makes its base map preserve
the weighted ideal
\((p,s^{hm})\), and hence its
normalized weighted modification.
The canonical divisor section
gives the action on the pulled
tautological line. In a chart,
an \(n\)-th root of a nonzero
multiplier may be chosen;
different choices are precisely
the \(\mu_n\)-action \eqref{cdp:eq:C.WT.34}.
Thus the coarse maps are
independent of the choice.
Their group law holds on the
unchanged open set, and then
on the normal chart. This
proves the actual finite action,
including \(g^{d_0}=G^k\)
on a period lift, without
a new linearization cocycle.

\noindent\textbf{Step 26.} More explicitly the new
transverse coordinate is
\(U=p/s^{hm}\). Its finite
map is
\[
             U\longmapsto
       \frac{a_0}{\xi^{hm}}U
                 +\frac{b(s)}{\xi^{hm}s^{hm}} .
 \tag{C.WT.37}\label{cdp:eq:C.WT.37}
\]
On logarithmic fibre coordinates,
the substitution adds \(nh\log s\);
the new off-diagonal coefficient
is \(c(s)s^{hm}\). Both are
holomorphic in the new chart.
Thus subsequent stabilizers
again have the exact affine
form \eqref{cdp:eq:C.WT.28}, with the
remaining displacement in
\eqref{cdp:eq:C.WT.36}. This proves the
induction for the complete
triangular action, not just
its action on homology.

\noindent\textbf{Step 27.} A product ruled cap uses the
actual normalized blowup of
\((p,s^h)\). Its two charts are
\[
                  p=s^hU,\qquad\text{and}\qquad xW=s^h .
 \tag{C.WT.38}\label{cdp:eq:C.WT.38}
\]
The height test has \(\ell>h\);
the central projective normal
action is therefore the identity
for the infinite deck map.
The cap is the specified
product \(F\times\PP^1\), and
its child displacement is
\(\delta(s)/s^h\).
The same finite jet gauge,
now through order \(h-1\),
gives its actual equivariant
blowup; its central affine
map is the recorded finite
projectivized normal map.

\noindent\textbf{Step 28.} Each step subtracts exactly
\(h_em_e\) from the period
order, with \(m_e=1\) for
a product cap. The height
test consequently supplies
the required holomorphic
inequalities at every step
and stops with zero remaining
order only at a nonlinear
terminal leaf. Reverse the
complete group-invariant
pruning record. Equations
\eqref{cdp:eq:C.WT.15}, \eqref{cdp:eq:C.WT.33}--\eqref{cdp:eq:C.WT.38}
restore every original root
component, all its child marks,
and its full torsor and
linear and translation
gluing maps.

\par\medskip\noindent\textbf{Properness of the compressed replacements.}
The affine-base estimate must
retain \(h\). Put
\[
 \rho(p,s)=
       (|p|^2+|s|^{2hm})^{1/(2m)}
       =|x|(|Z|^2+|W|^{2m})^{1/(2m)} .
 \tag{C.WT.39}\label{cdp:eq:C.WT.39}
\]
It satisfies \(\rho\ge|s|^h\).
Since \(\ell\ge hm\), the
translation by \(j\delta(s)\)
distorts this norm by at most
a uniform polynomial in
\(1+|j|\), in both directions.
The pulled line has the same
local fibre \(x^{-n}\) as
before base change. Its norm
distortion is a polynomial
of exponent at most \(n/m\).
Uniform exponential contraction
by \(\gamma(s)^j\) dominates
that polynomial, so a compact
norm annulus meets every orbit
and the action is proper
and free. The finite-group
orbit of a chosen norm gives
the same estimates uniformly.

\noindent\textbf{Step 29.} For comparison with the old
annular chart, the required
normalizing deck exponent is
\(O(1+|\log\rho|)\). Its base
displacement is bounded by
\[
              C|\delta(s)|(1+|\log\rho|).
 \tag{C.WT.40}\label{cdp:eq:C.WT.40}
\]
If \(\rho<1\), use
\(\rho\ge|s|^h\). If \(\rho\ge1\),
an old compact base bound
gives
\(|p_{\rm new}|\le B+
C|\delta(s)|\log(2+|p_{\rm new}|)\),
so \(p_{\rm new}\) is uniformly
bounded. The required new
annular domain is consequently
compact over every old compact
domain, even though the
universal affine base is
not compact.

\noindent\textbf{Step 30.} Finally, on a compact Hopf
fundamental domain, \(x\to0\)
implies \(|s|^h=|xW|=O(|x|)\).
The old base coordinate after
normalization is
\[
          O(|x|^m+|s|^\ell|\log|x||)\longrightarrow0.
 \tag{C.WT.41}\label{cdp:eq:C.WT.41}
\]
Every cluster point is therefore
on the single removed compact
elliptic curve. The closed
overlap graphs and nested
compact annuli used in \ref{cdp:st:C.2.UN.60}
now give the same Hausdorff
proper replacement, retaining
the literal \(A_{h-1}\) neck.
Product caps are proper
normalized blowups. Distinct
central marks have disjoint
finite-group-invariant tubes;
their finite simultaneous
replacements therefore glue.
There are only finitely many
stages.

\par\medskip\noindent\textbf{The smooth weighted quotient and the compact second end.}
The constructed root family
has only its recorded neck
singularities. At a generic
point of a strict component
over \(S_i\), the inertia has
order \(N/m_i\), fixes the
component and acts on its
normal parameter by that
primitive character. Its
quotient is smooth and the
full fibre multiplicity is
exactly \(m_i\).

\noindent\textbf{Step 31.} At an original double curve
of multiplicities \(a=gA,b=gB\),
the stabilizer of a point of
one lifted curve has order
\(hAB=N/g\). Its subgroup
\(\mu_h\) fixes \(u,v\) and
acts primitively on \(s\).
Quotienting \(uv=s^h\) first
by it gives the smooth
\((u,v)\)-chart. The residual
cyclic group of order \(AB\)
has the two independent
characters of orders \(B,A\).
Since \(\gcd(A,B)=1\), its
invariant ring is exactly
\[
 \bigl(\C\{z,u,v,s\}/(uv-s^h)\bigr)^{\,I}
       =\C\{z,u^B,v^A\},\qquad
                  t=(u^B)^a(v^A)^b .
 \tag{C.WT.42}\label{cdp:eq:C.WT.42}
\]
Here \(I\) is that inertia
group. The curve coordinate
\(z\) is invariant, because
the actual curve map is
\'{e}tale. Finite character
averaging gives coordinates
of the form \(uU,vV,z+\cdots\)
with their original characters.
Their product is \(s^h\)
times an invariant unit;
dividing one coordinate by
that unit keeps both the
characters and the literal
equation \(uv=s^h\).
This proves \eqref{cdp:eq:C.WT.42} for the
actual finite action, not
only its tangent representation.

\noindent\textbf{Step 32.} There is no edge inversion
to add an unrecorded quotient
of this chart. At interior
points the strict component
maps are \'{e}tale after the
generic normal inertia, so
there are no other quotient
singularities. The resulting
total space near zero is
smooth and its weighted
central support is the
original one.

\noindent\textbf{Step 33.} For completeness the weighted
normal line is retained as
well. On the root complement
the original normal coordinate
is \(v_i=s^{N/m_i}/\xi\).
The character eigensheaves
of the finite component map
extend by \eqref{cdp:eq:C.WT.7}; explicitly
their \(k\)-th pieces are
\[
 \mathcal L_i^{\,k}\otimes
  \OO_{S_i}\!\left(
       \sum_j\left\lfloor\frac{k m_j}{m_i}\right\rfloor J_{ij}
             \right),\qquad 0\le k<m_i .
 \tag{C.WT.43}\label{cdp:eq:C.WT.43}
\]
This follows by dividing
\(\xi^k\) by its integral
normal zero order in that
chart. The retained full
section and its orders
recover \(\mathcal L_i\)
and the specified edge line
maps; no reduction of an
integer modulo \(m_i\)
has occurred. The chosen
normal-product constants
are the only freedom stated
in the theorem.

\noindent\textbf{Step 34.} Carry out the conjugate
pruning record at infinity
as well. Formula \eqref{cdp:eq:C.WT.25}
has base character
\(\zeta_N^{-1}\) and canonical
character \((\det D)^{-1}\);
its primitive elliptic period
again has exact order \(R\).
Thus the conjugate height
and finite-jet comparisons
hold with precisely the
conjugate characters.
All charts, including the
compressed necks and
\eqref{cdp:eq:C.WT.42}, are the actual
conjugates of those at zero.
The quotient is therefore
smooth at that end too and
has the conjugate prescribed
weighted support. The two
sets of replacements have
disjoint parameter
neighbourhoods.

\noindent\textbf{Step 35.} The compact core and the
proper replacements give
a compact normal root
threefold. Its finite
quotient is smooth by the
charts just proved. Away
from zero and infinity
the base action is free;
one fibre over each orbit
maps biholomorphically to
the quotient fibre. They
are consequently connected
complex two-tori. Each
original component and
each double elliptic curve
has Euler number zero,
so both reduced central
supports have Euler zero.
This completes necessity,
sufficiency and the compact
comparison, including all
finite triangular actions.
\end{proof}






\section{Actual contractions, canonical reduction and finite generation}\label{sec:II:11}


\begin{proposition}[Maximum components without a minimality assumption]\label{cdp:st:C.2.UN.69}
\StatementRecord{cdp:st:C.2.UN.69}{proved}{Historical events: Claude Fable 5.2 (Anthropic), effort high, 16 November 2026: carried as background without fault found. GPT-6 Astra (OpenAI), ChatGPT Pro mode, effort Pro, S6\_09\_PAPERS, 17 November 2026: targeted proof check. Also checked in this writing. Pre-continuation writing events: GPT-6 Astra (OpenAI), ChatGPT Pro mode, effort Pro, S6\_10\_PAPERS, 18 November 2026: targeted proof check. Continuation metadata events: GPT-6 Astra (OpenAI), ChatGPT Pro mode, effort Pro, S6\_10\_PAPERS continuation, 19 November 2026: used and rechecked in this run; status-vocabulary reconciliation; metadata reconciliation only; no new mathematical check. \emph{Pre-continuation recorded proof scope.} The exact canonical and dual-normal restrictions give every edge square. The actual minimal-model map and square-zero pencil first prove the rationality and verticality assertions. Integer intersections exclude all hidden point blowups at a strict-drop maximum. A ruled rational-section centre would transport its square-minus-two sphere into the actual nearby torus, contradicting its contractible universal cover. Strict height ascent along rational boundaries therefore ends in a genuine plane-line contraction. Each such map removes exactly one original rational edge, preserves all remaining boundary neighbourhoods and subtracts two from the support Euler number}{the hypotheses, notation and definitions stated here; \cref{cdp:st:C.2.UN.64}; \cref{cdp:st:C.2.UN.50}; \cref{cdp:st:C.2.UN.58}; \cref{cdp:st:C.2.UN.19}; \cref{cdp:st:C.2.UN.15}; \cref{cdp:st:C.2.UN.31}; \cref{cdp:eq:C.NM.4}; \cref{cdp:eq:C.NM.7}; \cref{cdp:eq:C.NM.6}; \cref{cdp:st:C.2.UN.66}; \cref{cdp:eq:C.NM.9}; \cref{cdp:eq:C.NM.3}; \cref{cdp:eq:C.NM.8}; \cref{cdp:eq:C.NM.10}}

Let $f:U\to\Delta$ be a proper degeneration of connected complex
two-tori from a smooth threefold. Let its full central fibre
$S=\sum_vT_v$ be reduced and simple normal crossing, without triple
points. No minimality assumption is made on any $T_v$.
Write $k_v$ for the actual canonical evaluation orders.

At a local maximum $T_v$ of these orders, precisely one of the
following descriptions applies:
\begin{enumerate}
\item Every adjacent component has order $k_v$.
\item The component is an unblown terminal ruled surface meeting
its parent in a section, with order drop one, or an unblown
terminal plane meeting its parent in a line, with order drop two.
These are the actual contractions of \ref{cdp:st:C.2.UN.64}.
\item The component is an unblown terminal Hopf surface, with
the exact axis normal order and positive canonical drop of
\ref{cdp:st:C.2.UN.50} and \ref{cdp:st:C.2.UN.58}.
\end{enumerate}
In the second alternative the ruled base has positive genus.
Thus, when the orders are not all equal, a global maximum is
one of the two kinds of strict-drop leaves.

Moreover every rational double curve can be removed by a finite
sequence of the actual plane-line contractions of \ref{cdp:st:C.2.UN.64}.
At each step the other double curves and their neighbourhoods
are unchanged. If $r_0$ is the number of connected rational
double curves of the original support, the resulting actual
torus degeneration has no rational double curve and satisfies
\[
             e(S)=e(S_{\rm nr})+2r_0 .
\tag{C.NM.1}\label{cdp:eq:C.NM.1}
\]
There are exactly $r_0$ point contractions. Reversing their
recorded ordinary blowups recovers the original support,
including its original smooth threefold neighbourhood.
In particular the rational double edges are removed as terminal
edges of the actual incidence graph; none belongs to a cycle.
\end{proposition}



\begin{proof}
\noindent\textbf{Step 1.} \par\medskip\noindent\textbf{The exact restrictions at an edge.}
Canonical evaluation, adjunction and the actual equation $xy=t$
give the same line maps as (C.MR.3)--(C.MR.4), without a
minimality assumption:
\[
 \begin{split}
 K_{T_v}&=\operatorname{div}(\omega_v)
       =\sum_{e=(v,w)}(k_w-k_v-1)C_e,\\
 N_{C_e/T_v}\otimes N_{C_e/T_w}&\simeq\OO_{C_e},\\
 K_{C_e}&\simeq N_{C_e/T_v}^{\,k_w-k_v}.
 \end{split}
\tag{C.NM.2}\label{cdp:eq:C.NM.2}
\]
The last equality uses that the other boundaries are disjoint
from $C_e$. In particular, for $g_e=g(C_e)$,
\[
                  (k_w-k_v)C_e^2=2g_e-2.
\tag{C.NM.3}\label{cdp:eq:C.NM.3}
\]
These formulas retain the actual meromorphic section and
dual-normal trivialization.

\par\medskip\noindent\textbf{A positive-square rational curve on a component.}
We shall use the following implication for any component, not
only a maximum:
\[
 \text{a smooth rational curve of positive square on }T
                         \quad\Longrightarrow\quad T
                         \text{ is rational}.
\tag{C.NM.4}\label{cdp:eq:C.NM.4}
\]
Here is the comparison with its actual minimal-model map.
For a nonrational ruled minimal model, every rational curve
maps to a point of the positive-genus ruling base. A surviving
such curve is a ruling fibre, of square zero; its strict
transform has nonpositive square. Every exceptional curve
also has negative square. On the Hopf minimal model supplied
by \ref{cdp:st:C.2.UN.19} there are no rational curves, since a map from
$\PP^1$ lifts to $\C^2\setminus\{0\}$ and has constant
coordinate functions. Hence a rational curve on a blowup of
that model is exceptional and has negative square.

\noindent\textbf{Step 2.} The remaining actual minimal models have Kodaira dimension
zero or one by \ref{cdp:st:C.2.UN.15}. In this multiplicity-one setting \ref{cdp:st:C.2.UN.31}
gives a torus or the smooth-multiple elliptic model. Its
canonical line is numerically nonnegative on every curve:
on a torus it is trivial; in the elliptic case its positive
power is the pullback of the actual base line of degree
\[
                   2g(B)-2+\sum_b(1-1/m_b)\ge0.
\tag{C.NM.5}\label{cdp:eq:C.NM.5}
\]
The inequality follows from the pluricanonical pushdown and
the nonnegative Kodaira dimension. For its point-blowup map
$\pi:T\to T_{\min}$, one has
$K_T=\pi^*K_{T_{\min}}+\sum b_EE$, with $b_E>0$.
A nonexceptional rational curve $C$ would therefore have
$K_T\cdot C\ge0$. Positive square instead gives
$K_T\cdot C=-2-C^2<0$ by adjunction. An exceptional curve
has negative square. This proves \eqref{cdp:eq:C.NM.4} through the
actual minimal-model morphism.

\noindent\textbf{Step 3.} If $C\subset T$ is rational with square $s=1$ or $2$,
blow up $s$ distinct smooth points of $C$. The strict
transform $C'$ has square zero. Since $T$ is rational,
$H^1(\OO)=0$ on this blown-up surface. The restriction
sequence for $\OO(C')$, with
$\OO_{C'}(C')\simeq\OO_{\PP^1}$, gives a second section
nonzero everywhere on $C'$. Together with the canonical
section cutting out $C'$ it gives an actual basepoint-free
pencil
\[
                         q:T'\longrightarrow\PP^1.
\tag{C.NM.6}\label{cdp:eq:C.NM.6}
\]
There are no basepoints outside $C'$ because the canonical
section is nonzero there. The fibre $C'$ is reduced and
connected. In the Stein factorization a finite map of
degree greater than one would either give more than one
point above this value or a multiplicity greater than one
on $C'$. Thus $q$ has connected fibres. Adjunction gives
genus zero on a smooth fibre. The ordinary relative
minimal-model contractions identify this rational fibration
with a projective-line bundle followed by point blowups
in fibres; in particular every irreducible vertical curve
is rational. Every original curve disjoint from $C$
is untouched by the selected blowups and has intersection
zero with $q^*\OO_{\PP^1}(1)=\OO(C')$, so it is vertical.
This supplies the precise pencil used below.

\par\medskip\noindent\textbf{A maximum with a rational boundary is already minimal.}
At a local maximum put $a_e=k_v-k_w+1\ge1$. Then
\[
             -K_T=\sum_e a_eC_e,\qquad
             (1-a_e)C_e^2=2g(C_e)-2.
\tag{C.NM.7}\label{cdp:eq:C.NM.7}
\]
The boundary is nonempty in the reducible case. A nonzero
pluricanonical section multiplied by the corresponding
power of this anticanonical section would be a nonzero
holomorphic function having zeros on a compact connected
surface. Consequently $\kappa(T)=-\infty$.

\noindent\textbf{Step 4.} For a rational boundary, \eqref{cdp:eq:C.NM.7} gives exactly
\[
                     (a_e,C_e^2)=(3,1)\ \text{or}\ (2,2).
\tag{C.NM.8}\label{cdp:eq:C.NM.8}
\]
Equation \eqref{cdp:eq:C.NM.4} first makes $T$ rational; no rational
pencil is used before this conclusion. The actual pencil
\eqref{cdp:eq:C.NM.6} then makes every other disjoint boundary rational.
Every such boundary would again have positive square by
\eqref{cdp:eq:C.NM.7}, whereas the Hodge index theorem permits at most
one positive-square class among mutually disjoint curves.
Thus there is only one boundary $C$, and $K_T=-aC$.
Any smooth rational $(-1)$-curve $E$ would satisfy
$a(C\cdot E)=1$, which is impossible for $a=2$ or $3$.
The surface is therefore minimal. The terminal-pair
classification and actual maps in \ref{cdp:st:C.2.UN.64} give its ruled
section or plane-line contraction.

\par\medskip\noindent\textbf{All other blowdowns occur only at coefficient-one elliptic boundaries.}
Suppose now that no boundary is rational. A $(-1)$-curve
$E$ is not in the boundary, and adjunction gives
\[
                     1=(-K_T)\cdot E
                              =\sum_e a_e(C_e\cdot E).
\tag{C.NM.9}\label{cdp:eq:C.NM.9}
\]
Every summand is a nonnegative integer. Exactly one is
one: $E$ meets one boundary transversely once, that
boundary has coefficient one, and $E$ misses all others.
The coefficient-one boundary is elliptic by \eqref{cdp:eq:C.NM.7}.
Its ordinary surface blowdown keeps this boundary smooth
and keeps all boundary curves disjoint. Pushforward of
the actual meromorphic canonical section gives the same
canonical divisor with the same coefficients on the
new surface. This is a statement about the surface
morphism; no ambient contraction of $E$ is invoked.
Repeat to the minimal model.

\noindent\textbf{Step 5.} The minimal possibilities with \eqref{cdp:eq:C.NM.7} are exactly those
computed in (C.MR.6)--(C.MR.7) and \ref{cdp:st:C.2.UN.64}. On a ruled
surface, intersection with a ruling fibre gives the three
partitions of two: one coefficient-two section, one
coefficient-one bisection, or two coefficient-one sections.
The connected elliptic bisection over an elliptic base is
excluded by the actual residue/trace map (C.EL.7)--(C.EL.8).
On the plane the possibilities are a coefficient-three
line or a coefficient-one cubic. On an elliptic Hopf
surface the invariant relative differential gives
$\sum (k_w-k_T)/m_e=\operatorname{val}(T)-2$ with every
summand nonpositive; thus it is either a strict-drop
one-boundary leaf or a two-boundary equal-height component.
The nonlinear Hopf axis formula gives the same strict
leaf alternative, including its exact secondary framing.
For an algebraic-dimension-zero diagonal Hopf surface,
the two-axis character equation and $|\alpha|,|\beta|<1$
make both coefficients one; a single axis would impose
an integral resonance. These are actual marked pairs
by \ref{cdp:st:C.2.UN.50}, \ref{cdp:st:C.2.UN.58} and the proof of \ref{cdp:st:C.2.UN.66}.

\noindent\textbf{Step 6.} In particular whenever a coefficient is greater than
one, there is just one boundary, with that coefficient.
No coefficient-one boundary then exists. Since the
surface blowdowns above preserve every boundary and
its coefficient, \eqref{cdp:eq:C.NM.9} shows that no such blowdown
could have taken place. The original strict-drop maximum
was already minimal. In all remaining cases every
coefficient is one and every neighbour has the same
height. A component at the global maximum with this
property cannot be adjacent to a strict-drop maximum
leaf. Its equal-height neighbours have the same property,
so connectedness propagates it to the whole support.
This proves the maximum assertion, except for the
rational ruled-section exclusion proved next.

\par\medskip\noindent\textbf{A rational section centre cannot lie in the smooth torus-fibre locus.}
In a ruled-section contraction of \ref{cdp:st:C.2.UN.64} with base
$B\simeq\PP^1$, the actual descended family is a
submersion near the compact centre $B$. Its containing
central component $P'$ is smooth there and
\[
                    N_{B/P'}\simeq K_B\simeq\OO_{\PP^1}(-2).
\tag{C.NM.10}\label{cdp:eq:C.NM.10}
\]
Choose a relatively compact tubular neighbourhood of
$B$ inside that smooth central surface and a slightly
larger compact neighbourhood in the submersion locus.
A smooth metric gives horizontal lifts of the two
coordinate vector fields on the disc. For sufficiently
small time their flows exist on the whole smaller
neighbourhood and embed it in an actual nearby fibre.
The induced diffeomorphism preserves the oriented
normal bundle and its Euler number. Thus the nearby
two-torus would contain an embedded sphere of
self-intersection minus two. Its universal cover is
the contractible space $\C^2$, so this sphere map is
nullhomotopic, its homology class is zero, and its
self-intersection is zero. This contradiction excludes
the rational ruled-section alternative.

\par\medskip\noindent\textbf{Following rational edges gives actual point contractions.}
Take a rational double curve and orient it toward its
higher-height endpoint. Equation \eqref{cdp:eq:C.NM.3} makes the
height difference nonzero; at that higher endpoint
its square is positive, equal to one or two. By
\eqref{cdp:eq:C.NM.4} the endpoint surface is rational. The pencil
\eqref{cdp:eq:C.NM.6} makes every other boundary on it rational.
Hodge index leaves at most one positive-square
boundary, the incoming one. Every other boundary
therefore has negative square by \eqref{cdp:eq:C.NM.3}, and its
neighbour has strictly higher height. Follow one
of those edges when it exists.

\noindent\textbf{Step 7.} The heights strictly increase and there are finitely
many components. The path ends at a rational
component having exactly its one incoming boundary.
It has $K=-aC$, with $(a,C^2)$ as in \eqref{cdp:eq:C.NM.8}.
The same integer-intersection argument excludes every
$(-1)$-curve, so it is an actual \ref{cdp:st:C.2.UN.64} leaf. Its ruled
section possibility has just been excluded by
\eqref{cdp:eq:C.NM.10}. It is therefore a plane-line leaf, and
\ref{cdp:st:C.2.UN.64} contracts it to a smooth point of the smooth
central locus of its parent.

\noindent\textbf{Step 8.} Every other boundary misses that line, so the actual
blowdown is an isomorphism on its neighbourhood.
It removes exactly one rational double curve, creates
no new double curve or triple point, retains the
punctured torus family, and decreases $e(S)$ by two.
Apply the argument to the new actual degeneration
until no rational double curve remains. Exactly
$r_0$ steps have been made; \eqref{cdp:eq:C.NM.1} and the recorded
ordinary-blowup reconstruction follow.
\end{proof}






\begin{proposition}[Torsion-normal parent alignment with all vertical blowups retained]\label{cdp:st:C.2.UN.70}
\StatementRecord{cdp:st:C.2.UN.70}{proved}{Historical events: Claude Fable 5.2 (Anthropic), effort high, 16 November 2026: carried as background without fault found. GPT-6 Astra (OpenAI), ChatGPT Pro mode, effort Pro, S6\_09\_PAPERS, 17 November 2026: targeted proof check. Also checked in this writing. Pre-continuation writing events: GPT-6 Astra (OpenAI), ChatGPT Pro mode, effort Pro, S6\_10\_PAPERS, 18 November 2026: targeted proof check. Continuation metadata events: GPT-6 Astra (OpenAI), ChatGPT Pro mode, effort Pro, S6\_10\_PAPERS continuation, 19 November 2026: used and rechecked in this run; status-vocabulary reconciliation; metadata reconciliation only; no new mathematical check. \emph{Pre-continuation recorded proof scope.} The actual minimal-model canonical intersection formula forces every blowdown to miss the marked elliptic curve. In the ruled case its finite elliptic pullback gives disjoint square-zero sections; the exact ruling-fibre intersection sum supplies contractions missing one section, and Hodge index proves they miss every lift, yielding the equivariant descended contraction. The retained rank-two extension splits by degree-zero cohomology or finite trace. The resulting finite product quotient gives the actual elliptic map, while the original fields and normalized unit-root comparison force identity monodromy and smooth multiple supports. All original boundary curves and recorded exceptional curves are vertical}{the hypotheses, notation and definitions stated here; \cref{cdp:st:C.2.UN.67}; \cref{cdp:st:C.2.UN.15}; \cref{cdp:st:C.2.UN.31}; \cref{cdp:eq:C.NP.3}; \cref{cdp:st:C.2.UN.19}; \cref{cdp:st:C.2.UN.37}; \cref{cdp:eq:C.NP.7}; \cref{cdp:eq:C.NP.8}; \cref{cdp:eq:C.NP.9}}

Let \(P\) be a component of a full reduced SNC fibre of
a smooth-total degeneration of connected complex two-tori.
Suppose it has a smooth elliptic boundary curve \(J\)
disjoint from its other boundary curves, with torsion
normal line. Suppose its coefficient in the actual
canonical divisor on \(P\) is nonnegative:
\[
 K_P=\operatorname{div}\omega_P
       =a_JJ+\sum_C a_C C,\qquad a_J\ge0 ,
 \tag{C.NP.1}\label{cdp:eq:C.NP.1}
\]
where the sum is over the other boundary curves and
the equality includes the actual meromorphic canonical
section induced by the ambient evaluation.

Then \(P\) is not rational. There is an actual proper
birational morphism
\[
 \pi:P\longrightarrow Z,\qquad q_Z:Z\longrightarrow B,
                    \qquad q=q_Z\pi ,
 \tag{C.NP.2}\label{cdp:eq:C.NP.2}
\]
with the following properties. The morphism \(\pi\)
is a finite sequence of point blowdowns and is an
isomorphism on a neighbourhood of \(J\). The map
\(q_Z\) is a relatively minimal fixed-Jacobian elliptic
quasi-bundle with identity integral monodromy, only
smooth multiple supports, and base genus at most one.
The image of \(J\) is one of these fibre supports,
of full multiplicity exactly
\(\operatorname{ord}N_{J/P}\).
Every original boundary curve on \(P\) is vertical
for \(q\). Every exceptional curve of \(\pi\) is
vertical; the complete remaining nonminimal data
are the recorded point blowups of fibres of \(q_Z\),
with their actual positions and normal maps.
In particular every rational boundary curve on
\(P\) is contracted by \(\pi\).

This applies, in particular, to the parent of a
Hopf leaf: its canonical coefficient is
\(a_J=m-1\ge0\), with \(m\) the exact normal
order of the Hopf axis. No contraction of any
of these surface exceptional curves in the
ambient threefold is asserted.
\end{proposition}



\begin{proof}
\noindent\textbf{Step 1.} Rational \(P\) is excluded by \ref{cdp:st:C.2.UN.67}. A component
of general type is excluded by \ref{cdp:st:C.2.UN.15}. Its remaining
minimal model is therefore ruled, Hopf, or has
Kodaira dimension zero or one, by the actual
component comparisons already established.
We construct the maps separately and then
apply the actual fields on \(P\).

\par\medskip\noindent\textbf{Nonnegative Kodaira dimension: every blowdown misses \(J\).}
Let \(\pi:P\to Z\) be its ordinary minimal-model
map and let \(J_0=\pi(J)\). The surface \(Z\)
is the actual minimal model supplied by
\ref{cdp:st:C.2.UN.31} in this multiplicity-one case. It is
a torus, or has the smooth-multiple elliptic
map with the canonical line formula
(C.MR.12). In the latter case that formula
has base degree
\[
       \delta=2g(B_0)-2+\sum_i(1-1/m_i)\ge0 .
 \tag{C.NP.3}\label{cdp:eq:C.NP.3}
\]
The inequality follows from its exact
pluricanonical pushdown and
\(\kappa(Z)\ge0\). Thus \(K_Z\) has
nonnegative intersection with every
curve: it is zero on a vertical
curve, and its intersection with
a horizontal curve is its positive
base degree times \(\delta\).
For a torus \(K_Z\) is trivial.
This proves the needed intersection
property by the explicit base-line
map, without a new nefness assumption.

\noindent\textbf{Step 2.} For the actual sequence of point
blowups write
\(K_P=\pi^*K_Z+\sum_E b_EE\),
where every exceptional coefficient
\(b_E\) is a positive integer.
Adjunction and \(J^2=0\) give
\[
       0=K_P\cdot J
          =K_Z\cdot J_0+\sum_E b_E(E\cdot J).
 \tag{C.NP.4}\label{cdp:eq:C.NP.4}
\]
Every term is nonnegative; \(J\)
is not an exceptional rational
curve. Thus each \(E\cdot J=0\).
The exceptional locus is disjoint
from \(J\), so \(\pi\) is an
actual isomorphism of its
neighbourhood. In particular
\(J_0\) is smooth, elliptic,
of square zero, with the
same torsion normal.

\noindent\textbf{Step 3.} The exact nonruled alignment
(C.MR.12)--(C.MR.14) applies
to this marked pair \((Z,J_0)\).
In Kodaira dimension one,
\eqref{cdp:eq:C.NP.3} is positive, so
adjunction makes \(J_0\)
vertical for the given map.
On a torus it is a translate
of an elliptic subtorus and
its quotient gives the
required map. In the remaining
Kodaira-dimension-zero case,
the finite translation-root
cover in that alignment is a
principal elliptic torsor.
A horizontal lift of \(J_0\)
trivializes its integral Chern
vector after a finite base
pullback; the actual Cech
period lattice then makes
the cover a torus. Its
parallel lifted elliptic
curves define the quotient
pencil, and the finite
group preserving them
descends it to \(Z\).
These are the actual
finite-cover and quotient
maps of (C.MR.14).
They make \(J_0\) a full
reduced fibre support.
Every other noncontracted
boundary image is disjoint
from \(J_0\) and is therefore
vertical. Contracted curves
map to points and are
vertical as well.

\par\medskip\noindent\textbf{Hopf minimal model.}
Suppose the actual minimal
model \(Z\) is Hopf, as in
\ref{cdp:st:C.2.UN.19}. Every irreducible
curve on it is a smooth
elliptic curve with square
zero. The ordinary blowup
square formula gives
\[
                 0=J^2=J_0^2-\sum_\nu r_\nu^2.
 \tag{C.NP.5}\label{cdp:eq:C.NP.5}
\]
Thus no centre lies on
any successive image of
\(J\); again the map
preserves its entire
neighbourhood. Its normal
is torsion. This excludes
the nontorsion axes of an
algebraic-dimension-zero
diagonal Hopf.

\noindent\textbf{Step 4.} If \(Z\) were nonlinear,
its only curve would be
its axis \(J_0\). All
other boundaries would
be contracted. Push the
actual meromorphic form
in \eqref{cdp:eq:C.NP.1} through the
common open isomorphism.
Its divisor on \(Z\)
would be \(a_JJ_0\).
But \(K_Z=-(r+1)J_0\)
with \(r\ge1\). This
would trivialize the
line of the positive
effective divisor
\((a_J+r+1)J_0\), contrary
to constancy of holomorphic
functions on the compact
surface. Thus \(Z\) is
an elliptic Hopf surface.
Its actual map from \ref{cdp:st:C.2.UN.37}
makes every surviving
boundary a smooth multiple
fibre support. Composing
with \(\pi\) gives \eqref{cdp:eq:C.NP.2},
with every exceptional
curve vertical.

\par\medskip\noindent\textbf{Ruled minimal model: the contraction preserving all lifts of \(J\).}
A ruled minimal base of
genus at least two cannot
occur. The normalization
of the image of \(J\)
would map nontrivially
from an elliptic curve
to that base, contradicting
Riemann--Hurwitz; a
vertical curve in a
rational ruling fibre
cannot have elliptic
normalization. The
rational case was
excluded already.
Thus the original ruling
extends to a rational
fibration
\[
                             p:P\longrightarrow E
 \tag{C.NP.6}\label{cdp:eq:C.NP.6}
\]
over an elliptic curve.
Its restriction \(J\to E\)
is a finite unramified
map. Pull back this
actual fibration along
that cover. After origins
are chosen, it is an
isogeny and is Galois.
On the smooth pulled-back
surface \(P'\), the
preimage of \(J\) consists
of disjoint sections
\(J_0',\ldots,J_{d-1}'\),
each of square zero and
with the original torsion
normal. Its deck group
permutes these sections
and acts freely on the
elliptic base.

\noindent\textbf{Step 5.} There is a vertical
\((-1)\)-curve disjoint
from \(J_0'\) whenever
a ruling fibre is reducible.
To prove this, let
\(F=\sum_\alpha m_\alpha E_\alpha\)
be such a fibre. The
section meets just one
component, with multiplicity
one, and misses all nodes.
All \(E_\alpha\) are
smooth rational curves;
adjunction and \(K\cdot F=-2\)
give
\[
             \sum_\alpha
                    m_\alpha(2+E_\alpha^2)=2 .
 \tag{C.NP.7}\label{cdp:eq:C.NP.7}
\]
If the only possible
\((-1)\)-curve were the
component meeting the
section, its contribution
would be one and all
others would be nonpositive.
If it were not a
\((-1)\)-curve, all terms
would be nonpositive.
Both alternatives contradict
\eqref{cdp:eq:C.NP.7}. A vertical
\((-1)\)-curve disjoint
from the section therefore
exists.

\noindent\textbf{Step 6.} Contract such curves
successively. This gives
a \(\PP^1\)-bundle \(Z'\)
over the same elliptic
base, preserving a
neighbourhood of \(J_0'\).
Consider any other
lifted section \(J_i'\).
Its final image is
disjoint from \(J_0'\).
The latter is a nonzero
nef square-zero curve,
so the Hodge index
theorem gives square
at most zero for that
image. The blowdown
square formula gives
square at least zero.
Equality holds; hence
every contracted curve
also missed every
\(J_i'\).

\noindent\textbf{Step 7.} This proves equivariance
of the contraction, not
just existence of a
choice after base change.
In each original ruling
fibre, the unique component
meeting \(J_0'\) survives
and all the others are
contracted. Since all the
other sections were
preserved, they meet
that same surviving
component. It is
therefore the intrinsic
component meeting the
full set of lifted
sections. The exceptional
locus, its connected
fibres and the resulting
contraction are preserved
by the whole deck group.
An isomorphism on their
common open complement
extends uniquely through
the proper contraction,
so the action descends
to \(Z'\).
Its action on the base
is free, hence its
action on \(Z'\) is
free. Taking the actual
quotients gives a
birational morphism
\(\pi:P\to Z=Z'/G\),
where \(Z\) is a smooth
\(\PP^1\)-bundle over
\(E\). This morphism
is a sequence of point
blowdowns and is an
isomorphism near \(J\).
These facts also follow
by contracting the
disjoint deck orbits
of the indicated
exceptional curves.

\noindent\textbf{Step 8.} At least one other
boundary survives on
\(Z\). Otherwise the
pushforward of the
actual canonical section
would give \(K_Z=a_JJ_0\).
Its intersection with
a rational ruling fibre
would be nonnegative,
whereas that canonical
intersection is minus two.
This uses the
nonnegative coefficient
in the statement.
Every surviving boundary
is disjoint from \(J_0\).

\par\medskip\noindent\textbf{The actual split bundle and descended elliptic map.}
On \(Z'\) the diagonal
section has torsion
normal \(L\). Its actual
section subline gives,
after tensoring, an
extension
\[
           0\longrightarrow\OO
             \longrightarrow V
             \longrightarrow L\longrightarrow0,
                  \qquad Z'=\PP(V).
 \tag{C.NP.8}\label{cdp:eq:C.NP.8}
\]
If \(L\ne\OO\), its
extension group
\(H^1(L^{-1})\) is zero.
If \(L=\OO\), the
complement of the section
is the additive torsor
of this extension class.
Normalize a component
of the pullback of the
second surviving boundary.
It maps finitely to the
elliptic base and lies
in that complement,
so it supplies a section
of the pulled additive
torsor. Trace composed
with pullback on
\(H^1(\OO)\) is its
positive degree.
The original extension
class is therefore zero.
Thus \eqref{cdp:eq:C.NP.8} splits
in both cases.

\noindent\textbf{Step 9.} Take a further finite
etale elliptic base cover
which kills \(L\) and
is Galois over \(E\).
The actual pulled-back
surface is \(E'\times\PP^1\).
Its finite deck maps have
the form
\((b,z)\mapsto(b+\tau,A(b)z)\).
The holomorphic map
\(A:E'\to\operatorname{PGL}_2(\C)\)
is constant, since the
target is affine and
the source compact.
Consequently the
actual quotient has
the map
\[
        Z=(E'\times\PP^1)/G'
                   \longrightarrow\PP^1/G'.
 \tag{C.NP.9}\label{cdp:eq:C.NP.9}
\]
The regular connected
fibre is \(E'/K\), where
\(K\) is the kernel of
the projective action.
At a projective stabilizer
the elliptic translation
is free. Thus the chart
is precisely a smooth
multiple-fibre chart,
with identity integral
monodromy and one fixed
Jacobian. Every boundary
lift disjoint from the
fixed section has constant
projective coordinate:
on its compact
normalization that
coordinate is a
holomorphic function
with values in the
affine complement.
Thus every surviving
boundary is a fibre
support of \eqref{cdp:eq:C.NP.9}.
Every contracted boundary
maps to a point, so
all original boundaries
are vertical for the
actual composite \(q\).

\par\medskip\noindent\textbf{The original fields and the remaining vertical blowups.}
In all cases we have
constructed an actual
connected elliptic map
on \(P\), with every
boundary vertical, and
the specified isomorphism
near \(J\). Apply the
field alignment proof
of (C.MR.8), in the
explicit nonminimal
form of \ref{cdp:st:C.2.UN.67}, to this
map on \(P\). Its
saturated fields are
holomorphic on entire
general elliptic fibres.
The projected-field
continuation and the
actual return germs
give identity integral
monodromy. No field
pushforward through a
surface contraction is
needed for this step.

\noindent\textbf{Step 10.} The constructed \(Z\)
is already relatively
minimal for its map.
In the nonruled case
it is the minimal
surface; Hopf surfaces
contain no rational
curves, and an elliptic
ruled minimal surface
contains no rational
\((-1)\)-curves. The
exact relatively minimal
unit-root comparison in
\ref{cdp:st:C.2.UN.67} consequently gives
only smooth multiple
fibres and
\(R^1(q_Z)_*\OO_Z=\OO_B\).
The compact-core bound
then gives
\(g(B)=p_g(Z)=p_g(P)\le1\).
Its connected normal-root
chart identifies the full
multiplicity at \(J_0\)
with the exact order of
its normal line. Since
\(\pi\) preserves that
neighbourhood, this is
the original order.

\noindent\textbf{Step 11.} All exceptional curves
of \(\pi\) map to
points of \(Z\), hence
are vertical for \(q\).
Conversely a rational
boundary curve cannot
map onto a fibre support
of \(q_Z\), which is
smooth elliptic; it
must be contracted.
Factor the proper
birational surface map
into its actual point
blowdowns and retain
the inverse centres
and charts. This is
exactly the claimed
record of vertical
point blowups. It
does not contract
anything in the
ambient threefold.
For a Hopf leaf,
the exact terminal
canonical relation
gives \(a_J=m-1\),
so all hypotheses
hold. This proves
the marked parent
alignment.
\end{proof}






\begin{proposition}[All positive canonical-flat central log spaces over a disc]\label{cdp:st:C.2.UN.71}
\StatementRecord{cdp:st:C.2.UN.71}{proved}{Historical events: Claude Fable 5.2 (Anthropic), effort high, 16 November 2026: carried as background without fault found. GPT-6 Astra (OpenAI), ChatGPT Pro mode, effort Pro, S6\_09\_PAPERS, 17 November 2026: targeted proof check. Also checked in this writing. Pre-continuation writing events: GPT-6 Astra (OpenAI), ChatGPT Pro mode, effort Pro, S6\_10\_PAPERS, 18 November 2026: targeted proof check; used and rechecked in this run. Continuation metadata events: GPT-6 Astra (OpenAI), ChatGPT Pro mode, effort Pro, S6\_10\_PAPERS continuation, 19 November 2026: used and rechecked in this run; status-vocabulary reconciliation; metadata reconciliation only; no new mathematical check. \emph{Pre-continuation recorded proof scope.} The triple-chart angular maps and nonzero residue integral inject each rank-two lattice into the actual nearby Z4 and kill its linear holonomy. The complete affine developing map gives the periodic unimodular triangulation. Gross-Siebert v3 Theorem4.14 pp114-116 supplies the exact units-to-residues squares; Theorem3.22 pp95-97 and Example3.30 p106 match the normal-product log torsor, and Theorems5.2-5.4 pp128-134 supply the normalized chart changes and H1 classification with the checked positive/simple hypotheses. The receiving proof explicitly transports these maps through convergent local rings. The actual shear quotient has free action, bounded logarithmic fundamental domain, compact barycentric charts, exact periods and nowhere-zero canonical form}{the hypotheses, notation and definitions stated here; \cref{cdp:st:C.2.UN.68}; \cref{cdp:eq:C.TG.4}; \cref{cdp:st:C.2.UN.65}; \cref{cdp:eq:C.TG.8}; \cref{cdp:eq:C.TG.3}; \cite[Theorems 4.14, 5.2 and 5.4]{GrossSiebert2006I}; \cref{cdp:eq:C.TG.6}; \cite[Example 3.30 and Theorem 3.22]{GrossSiebert2006I}; \cref{cdp:eq:C.TG.10}; \cref{cdp:eq:C.TG.12}; \cref{cdp:eq:C.TG.11}; \cref{cdp:eq:C.TG.13}}

The positive canonical-flat central fibres of \ref{cdp:st:C.2.UN.68},
with their actual divisorial log structure and marked
dual complex, have the following complete description.
Take a rank-two lattice $N$, a finite-index sublattice
$\Gamma\subset N$, a $\Gamma$-periodic unimodular
triangulation $\mathscr P$ of $N_{\R}$ with vertex set
$N$, and a homomorphism
\[
       c:\Gamma\longrightarrow T_N
          :=\operatorname{Hom}(N^*,\C^*).
\tag{C.TG.1}\label{cdp:eq:C.TG.1}
\]
Require that each triangle have three distinct vertices
in $N/\Gamma$; this is the embedded simple-NC condition.
The central component at $p\in N/\Gamma$ is the actual
toric surface of the star fan
$\{\R_{\ge0}(\sigma-p):p\in\sigma\in\mathscr P\}$.
The homomorphism in \eqref{cdp:eq:C.TG.1} supplies all the remaining
logarithmic gluing constants, modulo the explicit torus
changes of the local charts described below.

Every such datum gives a smooth proper analytic
degeneration of connected complex two-tori over a
sufficiently small disc, with trivial canonical line,
and with exactly the specified marked central log space.
Every actual fibre in \ref{cdp:st:C.2.UN.68} is obtained in this way.
In a basis of $N$, let the columns of the integral
matrix $Q$ be a basis of $\Gamma$, and choose logarithms
of its two values under $c$ as the columns of $C$.
The actual nearby torus is
\[
 F_t=\C^2\big/
       \left(2\pi i\Z^2+(\log t\,Q+C)\Z^2\right),
       \qquad 0<|t|<\epsilon .
\tag{C.TG.2}\label{cdp:eq:C.TG.2}
\]
The central Euler number is $2[N:\Gamma]$.
This is a proper local occurrence and central-log-space
classification. It does not identify the full analytic
germ of an arbitrary given degeneration with this
chosen quotient, and it makes no compact-base
globalization assertion.
\end{proposition}



\begin{proof}
\noindent\textbf{Step 1.} \par\medskip\noindent\textbf{The actual integral-affine charts.}
On each toric component $V_i$ of \ref{cdp:st:C.2.UN.68} write
$N_i=\pi_1(V_i\setminus D_i)$, and let $n_{ij}$
be the primitive angular meridian of its boundary
curve $C_{ij}$. A triple point $ijk$ gives the
unimodular triangle
\[
                 0,\ n_{ij},\ n_{ik}\quad\hbox{in }(N_i)_{\R}.
\tag{C.TG.3}\label{cdp:eq:C.TG.3}
\]
To pass from the chart at $j$ to that at $i$,
use the unique integral-affine map which sends
$0$ to $n_{ij}$ and whose linear part $P_{ij}$ satisfies
\[
 P_{ij}n_{ji}=-n_{ij},\qquad
 P_{ij}n_{jk}=n_{ik}-n_{ij}.
\tag{C.TG.4}\label{cdp:eq:C.TG.4}
\]
These are the actual angular maps. Indeed in a
triple chart $f=xyz$, the coordinates on $V_i$
are $(y,z)$ and those on $V_j$ are $(x,z)$.
On the nearby fibre $x=t/(yz)$, a loop varying
$x$ with $z$ fixed has angular vector $-n_{ij}$,
and a loop varying $z$ with $x$ fixed has vector
$n_{ik}-n_{ij}$. Multiplication of these coordinates
by holomorphic units leaves these integral maps
unchanged after the chart is shrunk.

\noindent\textbf{Step 2.} There are exactly two triple points on $C_{ij}$.
Write their third vertices as $k,l$, retaining
their two distinct incidences if the same global
component occurs twice. The toric fan relation on
$V_i$ is
\[
 n_{ik}+n_{il}=-(C_{ij}^2)_{V_i}n_{ij}.
\tag{C.TG.5}\label{cdp:eq:C.TG.5}
\]
The corresponding relation at $j$ agrees with
\eqref{cdp:eq:C.TG.4} at the other endpoint precisely because
\[
             (C_{ij}^2)_{V_i}+(C_{ij}^2)_{V_j}=-2,
\tag{C.TG.6}\label{cdp:eq:C.TG.6}
\]
the actual normal-product and two-triple-point
formula of \ref{cdp:st:C.2.UN.65}. Thus there is no affine singularity
in an edge. The complete smooth toric fan at each
vertex makes its star an affine neighbourhood of
the origin. The triangle charts and \eqref{cdp:eq:C.TG.4}
therefore give an integral-affine structure on
the entire actual dual two-torus $B$.

\par\medskip\noindent\textbf{The angular lattices inject into the actual nearby torus.}
The nowhere-zero three-form of \ref{cdp:st:C.2.UN.65} is, in an
actual triple chart,
\[
       \Omega=a(x,y,z)\,dx\wedge dy\wedge dz,
          \quad a(0)\ne0,
          \qquad f=xyz.
\tag{C.TG.7}\label{cdp:eq:C.TG.7}
\]
The induced relative two-form on $F_t$ is
$a(x,y,t/(xy))\,d\log x\wedge d\log y$.
Its integral over the real angular two-torus
$|x|=|y|=|z|=|t|^{1/3}$ is exactly
\[
 (2\pi i)^2\sum_{q\ge0}a_{q,q,q}t^q
                  =(2\pi i)^2(a(0)+O(t))\ne0.
\tag{C.TG.8}\label{cdp:eq:C.TG.8}
\]
This follows by the convergent power-series expansion
and the two circle integrals; only the diagonal
monomials survive.

\noindent\textbf{Step 3.} For a fixed component, its toric complement retracts
onto such an angular two-torus. Any finite collection
of its loops, their homotopies and the connecting
paths to the corner lies in a compact subset away
from the boundary. The submersion $f$ near this
compact subset gives actual smooth horizontal
transports into $F_t$. Consequently its angular
lattice has a well-defined homomorphism
$N_i\to\pi_1(F_t)$, up to conjugation, and the
corner homomorphism is this same map. Since
$\pi_1(F_t)=\Z^4$ is abelian, conjugation makes
no ambiguity. If this map had rank at most one,
the fundamental class of the angular torus would
map to zero in $H_2(F_t,\Z)$: the actual complex
torus is a $K(\Z^4,1)$ and its degree-two homology
map is the exterior square of its degree-one map.
This contradicts \eqref{cdp:eq:C.TG.8}. Thus every $N_i$
injects into the actual $\Z^4$.

\noindent\textbf{Step 4.} The corner maps and the compact double-curve collars
identify the injected lattices by exactly \eqref{cdp:eq:C.TG.4}.
Transport around any loop in $B$ gives a linear
holonomy $A$ with $\iota_iA=\iota_i$ in $\Z^4$.
Injectivity gives $A=I$. This proves trivial linear
holonomy by maps into the given smooth fibre; no
invariant-cycle surjectivity is being assumed.

\noindent\textbf{Step 5.} The affine charts now differ by translations.
Their Euclidean metrics give a complete flat metric
on compact $B$. Extending straight lines on its
universal cover proves that the developing map is
an actual affine isomorphism with $N_{\R}$.
The deck group is a full lattice of translations
$\Gamma\subset N$, because all the transition
translations in \eqref{cdp:eq:C.TG.3}--\eqref{cdp:eq:C.TG.4} are integral.
Every developed triangle is unimodular. Conversely
every lattice point of $N$ is a vertex, since a
unimodular triangle contains no other lattice point
and the triangles cover $N_{\R}$. Thus the actual
dual complex is $(N_{\R},\mathscr P)/\Gamma$,
with exactly $[N:\Gamma]$ vertices.

\par\medskip\noindent\textbf{The exact logarithmic gluing comparison.}
We spell out the primary-source comparison used to
retain the given gluing constants. The divisorial
log structure of $(X,S)$ restricted to $S$ is locally
generated by the actual factors of $f$. At a triple
point its chart is $\mathbb N^3$, and the base generator
is $(1,1,1)$. At a double curve it is the analogous
two-factor chart. The normalization is the disjoint
union of the actual toric $V_i$; its conductor is
the reduced full toric boundary, with the exact
cartesian and cocartesian SNC conductor square.
The central surface is Gorenstein and this log
map is smooth at every point, directly from
$xyz=t$, $xy=t$ and the submersion chart. Thus
the log-singular set is empty.

\noindent\textbf{Step 6.} The precise gluing result is
\cite[Theorems 4.14, 5.2 and 5.4]{GrossSiebert2006I}.
Its use here does not import a polarization or an
algebraic total space. The following description
checks the analytic-to-algebraic comparison in the
proof. Fix the toric identifications of every
normalization stratum with its actual counterpart.
Every boundary identification is an automorphism
of $\PP^1$ preserving its two specified corner
points, and hence, in the specified coordinates,
is multiplication by a nonzero constant. These
are the closed gluing maps of that source.

\noindent\textbf{Step 7.} At a triple point, restrict an actual log chart
to each of the three toric affine planes. A
monomial $z^p$ is multiplied by a unit $h_p$
in its analytic local ring, with
$h_{p+q}=h_ph_q$. The map
\[
       p\longmapsto h_pz^p
       \quad\longmapsto\quad
       p\longmapsto h_p(0)z^p
\tag{C.TG.9}\label{cdp:eq:C.TG.9}
\]
is exactly the units-to-residues map in the proof
of Theorem 4.14. Restricting to a shared coordinate
axis and evaluating at its endpoint commute.
The resulting constant monoid maps therefore
intertwine the actual closed gluing diagrams.
Their colimits give the actual open toric charts
and the isomorphism with the given central space,
as in the commutative squares of that proof.
All these maps are defined by the same convergent
local rings in the analytic setting. The constant
residues are algebraic monomial maps; thus the
resulting finite open atlas and its actual
equivalence relation also give an algebraic
space whose analytification is the given $S$.

\noindent\textbf{Step 8.} The remaining log data are the nowhere-zero
sections of the lines
$N_{C/V_i}\otimes N_{C/V_j}\otimes\OO_C(p+q)$
on the double $\PP^1$'s, where $p,q$ are their
two triple points, with the actual triple
compatibilities. These corrected normal-product
lines have degree zero by \eqref{cdp:eq:C.TG.6}, and
their sections are nonzero constants in their
trivializations. They are algebraic as well.
This is the exact log-smooth-structure torsor
in \cite[Example 3.30 and Theorem 3.22]{GrossSiebert2006I};
the constant maps just described retain its
given section, not merely the underlying surface.
Consequently the log central space satisfies
Definition 4.3 of that source over $\C$ with
empty singular set. Positivity says these
sections have no poles and is immediate.
The affine monodromy polytopes are points;
in Definition 1.60 simplicity is the case
$p=0$, whose remaining conditions are vacuous.

\noindent\textbf{Step 9.} Theorems 5.2 and 5.4 now supply actual torus
changes of the open charts, preserving the
marked dual complex, which turn these data
into a Cech cocycle for the constant sheaf
$N\otimes\C^*=T_N$ on $B$.
The maps of a cohomologous change are explicitly
$s'_{\alpha\beta}=a_\alpha^{-1}
s_{\alpha\beta}a_\beta$, acting on each chart
by its torus element $a_\alpha$; their pullbacks
also preserve the normalized log section.
Thus they are actual isomorphisms of central
log spaces. On the universal cover
$N_{\R}$ a cocycle is trivialized by path
products, and its only remaining values are
the two commuting deck holonomies. Therefore
\[
 H^1(B,N\otimes\C^*)
       =\operatorname{Hom}(\Gamma,T_N),
\tag{C.TG.10}\label{cdp:eq:C.TG.10}
\]
with precisely the marked gluing datum \eqref{cdp:eq:C.TG.1}.
This is the actual comparison to be used in
the construction, rather than an assertion
that arbitrary analytic neighbourhoods are toric.

\par\medskip\noindent\textbf{The smooth toric threefold and its exact deck maps.}
For every triangle $\sigma$ of $\mathscr P$, take
the cone over $\sigma\times\{1\}$ in
$N_{\R}\oplus\R$. These cones and their faces
form a fan $\Sigma$. Its three primitive
generators $(p_0,1),(p_1,1),(p_2,1)$ are an
integral basis because the triangle is
unimodular. Its toric chart is therefore
literally $\C^3$ with coordinates $x_0,x_1,x_2$,
and the character of the last dual coordinate is
\[
                  t=x_0x_1x_2,
       \qquad z^m=\prod_{a=0}^2x_a^{\langle m,p_a\rangle}
                       \quad(m\in N^*).
\tag{C.TG.11}\label{cdp:eq:C.TG.11}
\]
The charts glue by their actual Laurent
monomial localizations. The fan intersection
condition makes this a Hausdorff smooth toric
analytic threefold. Although the fan has
infinitely many cones, each central orbit has
the ordinary finite star of its cone, and the
open torus is a single common open chart.
These are actual analytic neighbourhoods.

\noindent\textbf{Step 10.} A translation $\lambda\in\Gamma$ acts on the
lattice $N\oplus\Z$ by the integral shear
$(n,h)\mapsto(n+h\lambda,h)$. It preserves
the fan and induces the actual monomial map
between its charts. Multiply this map by
$c_\lambda$ on the torus factor. On the
punctured open torus the result is exactly
\[
              g_\lambda(z,t)=(c_\lambda t^\lambda z,t).
\tag{C.TG.12}\label{cdp:eq:C.TG.12}
\]
The $c_\lambda$ are a homomorphism, and
constant torus multiplication commutes with
the shears. Thus $g_\lambda g_\mu=g_{\lambda+\mu}$
on every chart. On the central fibre these
maps identify the component and boundary
charts by the cocycle in \eqref{cdp:eq:C.TG.10}. The
quotient central log space is consequently
the exact space specified earlier, including
its given log section.

\par\medskip\noindent\textbf{Freeness, properness and the actual compact fibres.}
Extend the homomorphism
$\lambda\mapsto\log|c_\lambda|$ real-linearly
to a map $K:N_{\R}\to N_{\R}$.
For $0<|t|<\epsilon$ put $L=\log|t|$ and
\[
                 q=\frac{\log|z|}{L}.
\tag{C.TG.13}\label{cdp:eq:C.TG.13}
\]
Then \eqref{cdp:eq:C.TG.12} sends $q$ to
$q+(I+K/L)\lambda$.
Choose $\epsilon<1$ so small that $I+K/L$
and its inverse are uniformly bounded and
invertible for $0<|t|<\epsilon$.
The action on every nonzero fibre is
therefore free and properly discontinuous;
its real logarithmic quotient has a compact
fundamental parallelogram. Its angular
directions already form a compact two-torus.
Taking logarithms gives exactly \eqref{cdp:eq:C.TG.2},
whose four real periods are independent
because $LQ+\operatorname{Re}C$ is invertible.

\noindent\textbf{Step 11.} At $t=0$ a nonzero translation cannot fix
any orbit: it would preserve a bounded cell
of the triangulation, which is impossible.
The central action is thus free. To verify
proper discontinuity also at the central
fibre, use a finite collection of fixed
toric charts about a proposed compact set.
For points in those charts tending to zero,
the coordinates in \eqref{cdp:eq:C.TG.11} show that $q$
stays bounded: each $\log|x_a|/L$ is bounded
below by a quantity tending to zero and
the three coefficients sum to one.
If both a point and its $g_\lambda$-translate
stay in such fixed compact sets, \eqref{cdp:eq:C.TG.13}
and the uniform inverse bound force
$\lambda$ to stay bounded. Only finitely
many lattice elements occur. The same
argument at central points uses the finite
sets of their translated cells. Hence the
quotient is Hausdorff and smooth.

\noindent\textbf{Step 12.} Finally choose the representatives of $q$
in a uniformly bounded fundamental
parallelogram for $(I+K/L)\Gamma$.
Only finitely many triangles of
$\mathscr P$ meet this bounded set.
In a triangle containing $q$, its
barycentric coordinates are precisely
\[
             \ell_a=\frac{\log|x_a|}{L},
             \qquad \ell_a\ge0,\quad\sum_a\ell_a=1.
\tag{C.TG.14}\label{cdp:eq:C.TG.14}
\]
Consequently $|x_a|\le1$ in that chart.
Any sequence of fibres tending to zero
has, after its actual deck translations,
a subsequence in one of these finitely
many closed unit polydiscs, and hence a
convergent subsequence. The central fibre
itself is the finite union of the complete
toric star surfaces. This proves properness
over a smaller disc. The punctured fibres
are connected by their literal quotient
presentation, and the central dual torus
is connected as well.

\noindent\textbf{Step 13.} The triangle-distinctness condition in
the statement prevents a component from
meeting a translate of itself. The
quotient therefore retains smooth embedded
components and the full reduced
simple-NC equation in \eqref{cdp:eq:C.TG.11}.
The differential
\[
                  dt\wedge d\log z_1\wedge d\log z_2
\tag{C.TG.15}\label{cdp:eq:C.TG.15}
\]
is invariant under \eqref{cdp:eq:C.TG.12}, because all
additional terms are multiples of $dt$.
In each unimodular chart \eqref{cdp:eq:C.TG.11} it is,
up to the basis determinant $\pm1$,
$dx_0\wedge dx_1\wedge dx_2$.
It extends nowhere vanishingly across
the central fibre. Thus the constructed
total space has trivial canonical line.
The exact central-space comparison and
\ref{cdp:st:C.2.UN.68} give its Euler number and complete
both the occurrence and exhaustion
assertions in the stated local scope.
\end{proof}






\begin{proposition}[The prescribed toric fibre and its conjugate over the projective line]\label{cdp:st:C.2.UN.73}
\StatementRecord{cdp:st:C.2.UN.73}{proved}{Historical events: Claude Fable 5.2 (Anthropic), effort high, 16 November 2026: carried as background without fault found. GPT-6 Astra (OpenAI), ChatGPT Pro mode, effort Pro, S6\_09\_PAPERS, 17 November 2026: targeted proof check. Also checked in this writing. Pre-continuation writing events: GPT-6 Astra (OpenAI), ChatGPT Pro mode, effort Pro, S6\_10\_PAPERS, 18 November 2026: targeted proof check; used and rechecked in this run. Continuation metadata events: GPT-6 Astra (OpenAI), ChatGPT Pro mode, effort Pro, S6\_10\_PAPERS continuation, 19 November 2026: used and rechecked in this run; status-vocabulary reconciliation; metadata reconciliation only; no new mathematical check. \emph{Pre-continuation recorded proof scope.} The angular and nonangular periods retain the logarithm-branch integral shifts. An independent symbolic construction of the four-by-four real period matrix gives exactly the stated determinant, with zero expanded difference; see analysis/UN73\_determinant\_replay.json. Its sign is positive from the chosen positive symmetric part and the exact sign of (1-r-squared)y. The zero-end exponential coordinates and base scale restore the original log constants. At infinity the reversed generators give precisely their complex conjugates. Bounded-unit cone charts and compact fundamental parallelepipeds give the two proper fillings, and their actual punctured isomorphisms produce the claimed compact family and Euler number}{the hypotheses, notation and definitions stated here; \cref{cdp:st:C.2.UN.71}; \cref{cdp:eq:C.TX.4}; \cref{cdp:eq:C.TX.8}; \cref{cdp:eq:C.TX.10}; \cref{cdp:eq:C.TX.12}}

Every marked central log space classified in \ref{cdp:st:C.2.UN.71} occurs
unchanged in a smooth compact complex threefold mapping
properly to $\PP^1$, with connected complex two-torus
fibres away from two values. Its two full singular fibres
are the prescribed reduced simple-NC log space and its
complex conjugate. Both neighbourhoods have a nowhere-zero
canonical form. If its lattice index is $d=[N:\Gamma]$,
each singular support has Euler number $2d$ and the
compact total space has Euler number $4d$.

The local base coordinates at the two values, and the
maps restoring the prescribed log sections, are part
of the construction below. The assertion preserves
the specified central log spaces; it does not assert
that an arbitrary given full degeneration germ is
the germ of this compact construction.
\end{proposition}



\begin{proof}
\noindent\textbf{Step 1.} \par\medskip\noindent\textbf{The exact initial periods and a positive symmetric part.}
Use the actual data $(N,\Gamma,\mathscr P,c)$ of \ref{cdp:st:C.2.UN.71}.
Choose $N=\Z^2$, write $\Gamma=Q\Z^2$ with
$Q\in\operatorname{Mat}_2(\Z)$ invertible over $\R$,
and choose the logarithms $C$ of its two torus
holonomies. In the vector coordinate
$(2\pi i)^{-1}Q^{-1}\log z$, the angular lattice is
\[
        N_0=Q^{-1}\Z^2\supset\Z^2,
             \qquad b=(2\pi i)^{-1}Q^{-1}C.
\tag{C.TX.1}\label{cdp:eq:C.TX.1}
\]
The two other periods of the local model are
$b_j+(\log u)/(2\pi i)e_j$. Write $b=U+iA$
with $U,A$ real matrices. Choose $R>0$ large enough
that the symmetric part $S$ of $A+RI$ is positive
definite, and put
\[
 b'=b+iRI,
 \qquad \kappa=e^{-2\pi R}>0,
 \qquad A+RI=S+aJ,
 \qquad J=\begin{pmatrix}0&1\\-1&0\end{pmatrix}.
\tag{C.TX.2}\label{cdp:eq:C.TX.2}
\]
In particular
$\exp(2\pi iQb'_j)=\kappa^{Qe_j}c_{Qe_j}$.
This changes only a recorded base scale; the
actual coordinate maps undo it at both ends.
We do not commute this imaginary shift through
the real-linear operation used next.

\par\medskip\noindent\textbf{A holomorphic family on the entire punctured line.}
Fix any integer $k\ge1$. For each constant
vector $w\in\C^2$, let
\[
 V_t(w)=w+t^kJ\overline w,
 \qquad B_t=I+t^kJ,
 \qquad 0\ne t\in\C.
\tag{C.TX.3}\label{cdp:eq:C.TX.3}
\]
This is holomorphic in $t$ for every fixed $w$.
Define the period lattice by its actual generators
\[
 \Lambda_t=V_t(N_0)
   +\sum_{j=1}^2\Z
       \left(V_t(b'_j)+\frac{\log t}{2\pi i}B_te_j\right).
\tag{C.TX.4}\label{cdp:eq:C.TX.4}
\]
Changing the logarithm by $2\pi i$ adds $B_te_j$,
which is an angular period because $\Z^2\subset N_0$.
Thus these lattices are globally single-valued
over $\C^*$, with the displayed integral change
of basis on overlaps of logarithm branches.

\noindent\textbf{Step 2.} We prove their real rank is four everywhere.
It suffices to use $\Z^2$ in place of its finite
angular enlargement $N_0$. Subtract the real
angular combinations supplied by $U$ and by
$\operatorname{Re}(\log t/(2\pi i))$ from the
last two columns. Set
\[
 x=\operatorname{Re}(t^k),\quad
 w=\operatorname{Im}(t^k),\quad
 r=|t|^k,\quad
 y=-\frac{\log|t|}{2\pi}.
\tag{C.TX.5}\label{cdp:eq:C.TX.5}
\]
With $A'=S+aJ=A+RI$, the resulting real
period matrix is exactly
\[
 \begin{pmatrix}
 I+xJ & wJ(A'-yI)\\
 wJ   & (I-xJ)A'+y(I+xJ)
 \end{pmatrix}.
\tag{C.TX.6}\label{cdp:eq:C.TX.6}
\]
Its determinant, by the displayed block
multiplication, is
\[
 \begin{aligned}
 \Delta(t)={}&(1+r^2)^2\det S
       +(1-r^2)(1+r^2)y\operatorname{tr}S\\
       &+y^2(1-r^2)^2
       +\bigl((1+r^2)a+2xy\bigr)^2>0.
 \end{aligned}
\tag{C.TX.7}\label{cdp:eq:C.TX.7}
\]
Indeed $(1-r^2)y\ge0$, the first term is
strictly positive, and the other three
terms are nonnegative. The identity uses
$r^2=x^2+w^2$ and holds for every argument
of $t$; it also holds where $B_t$ fails
to be complex-invertible. Thus \eqref{cdp:eq:C.TX.4}
is a discrete rank-four real lattice
for every $t\in\C^*$.

\noindent\textbf{Step 3.} The actual quotient of the holomorphic
vector bundle $\C^2\times\C^*$ by these
locally holomorphic period sections is
therefore a smooth family of connected
complex two-tori. On every compact
annulus, the real period matrices and
their inverses are uniformly bounded.
Their real fundamental parallelepipeds
give compact representatives. This
proves properness of this middle family
over the annulus, and hence over $\C^*$.

\par\medskip\noindent\textbf{The exact chart and prescribed log section at zero.}
Near zero $B_t$ is complex-invertible.
In the middle vector coordinate $v$ use
the actual exponential map
\[
                   z=\exp(2\pi iQ B_t^{-1}v).
\tag{C.TX.8}\label{cdp:eq:C.TX.8}
\]
It quotients exactly the angular lattice
$B_tN_0$, since $QB_t^{-1}B_tN_0=\Z^2$.
The other two period translations act by
\[
 z\longmapsto c^0_{Qe_j}(t)\,t^{Qe_j}z,
 \qquad
 c^0_{Qe_j}(t)=
       \exp\bigl(2\pi iQ B_t^{-1}V_t(b'_j)\bigr).
\tag{C.TX.9}\label{cdp:eq:C.TX.9}
\]
These are holomorphic nonzero vectors
near zero; extend them multiplicatively
to a homomorphism $c^0(t):\Gamma\to T_N$.
Their exact central values are
$c^0_\lambda(0)=\kappa^\lambda c_\lambda$.

\noindent\textbf{Step 4.} Use the base coordinate $u=\kappa t$.
The same action is then
\[
 z\longmapsto d^0_\lambda(u)u^\lambda z,
 \qquad
 d^0_\lambda(u)=\kappa^{-\lambda}
                    c^0_\lambda(u/\kappa),
 \qquad d^0_\lambda(0)=c_\lambda.
\tag{C.TX.10}\label{cdp:eq:C.TX.10}
\]
This is an actual equality of holomorphic
maps, not a replacement of its limiting
constant by a numerical invariant.

\noindent\textbf{Step 5.} Apply the cone-chart construction of
\ref{cdp:st:C.2.UN.71} to these variable holomorphic
units $d^0_\lambda(u)$. Multiplication
by them is a holomorphic torus action
on every toric chart over the same
base. Their homomorphism identity
makes the shear action exact. The
proof of freeness and properness is
unchanged: the real-linear maps
$\lambda\mapsto\log|d^0_\lambda(u)|$
are uniformly bounded as $u\to0$,
so the perturbation of the tropical
translation lattice tends uniformly
to the identity. The same finitely
many barycentric unit polydiscs give
properness at zero. Thus \eqref{cdp:eq:C.TX.8}
identifies this filling with the
middle family on a punctured disc.
Its actual central log constants in
the coordinate $u$ are exactly $c$,
by \eqref{cdp:eq:C.TX.10}, so the specified
marked log space is retained.

\par\medskip\noindent\textbf{The exact conjugate filling at infinity.}
Put $s=t^{-1}$ and change the vector
frame by $v_\infty=s^kv$. The angular
periods become
\[
 B^\infty_sN_0,
       \qquad B^\infty_s=s^kI+J.
\tag{C.TX.11}\label{cdp:eq:C.TX.11}
\]
Reverse the two nonangular period
generators. Their resulting exact
expressions are
\[
 -(s^kb'_j+J\overline{b'_j})
                +\frac{\log s}{2\pi i}B^\infty_se_j.
\tag{C.TX.12}\label{cdp:eq:C.TX.12}
\]
The integer generator reversal is
permitted because it changes neither
the subgroup nor the actual quotient.
The matrix $B^\infty_s$ is invertible
near $s=0$. The exponential coordinates
$z_\infty=\exp(2\pi iQ(B^\infty_s)^{-1}v_\infty)$
therefore give the exact deck constants
\[
 c^\infty_{Qe_j}(s)=
  \exp\bigl(-2\pi iQ(B^\infty_s)^{-1}
                  (s^kb'_j+J\overline{b'_j})\bigr).
\tag{C.TX.13}\label{cdp:eq:C.TX.13}
\]
At zero this is
\[
 c^\infty_\lambda(0)
       =\overline{c^0_\lambda(0)}
       =\kappa^\lambda\overline{c_\lambda}.
\tag{C.TX.14}\label{cdp:eq:C.TX.14}
\]
The sign in \eqref{cdp:eq:C.TX.12} is essential:
$\exp(-2\pi i\overline{Qb'})$
is the conjugate of $\exp(2\pi iQb')$.

\noindent\textbf{Step 6.} Choose the base coordinate $v=\kappa s$
at infinity. Define
$d^\infty_\lambda(v)=\kappa^{-\lambda}
c^\infty_\lambda(v/\kappa)$.
Then $d^\infty_\lambda(0)=\overline{c_\lambda}$,
and the exact action is
$z_\infty\mapsto d^\infty_\lambda(v)v^\lambda z_\infty$.
The two base charts satisfy $uv=\kappa^2$;
they are literal holomorphic coordinates
at the two points of one projective line.
Use the same integral fan $\Sigma$ and
its variable-constant cone quotient
at this end. All its monomial chart
maps are defined over $\Z$, while its
central torus gluing constants are
the complex conjugates of the given
ones. Hence its marked central log
space is exactly the conjugate of
the prescribed one, with the stated
local base coordinate.

\par\medskip\noindent\textbf{Compact gluing and canonical neighbourhoods.}
The exponential maps \eqref{cdp:eq:C.TX.8} and
their infinity counterparts are
actual isomorphisms of the
punctured torus families, since
they first divide by precisely
the angular lattice and then
by the displayed remaining
period group. Glue the two
proper cone fillings to the
middle family along these
isomorphisms. Disjoint smaller
base discs and a compact
intermediate annulus show
directly that the glued total
space is Hausdorff and proper
over $\PP^1$. It is smooth,
and its fibres are connected.
Properness over the compact
base makes the total space
compact. There are no other
singular fibres, because the
middle family is smooth.

\noindent\textbf{Step 7.} In both end fillings the form
$du\wedge d\log z_1\wedge d\log z_2$
(respectively with $v,z_\infty$)
is invariant under the actual
variable-constant shear action:
the additional logarithmic
derivatives are multiples of
the base differential and
vanish in this wedge. The
unimodular cone calculation
of \ref{cdp:st:C.2.UN.71} extends it to a
nowhere-zero three-form.
Thus both end neighbourhoods
are canonical-flat with the
full reduced simple-NC
fibres specified in the
statement. Finally Euler
additivity over the base,
with smooth torus Euler
number zero, gives
$e(Y)=2d+2d=4d$.
This completes the compact
occurrence construction.
\end{proof}






In this subsection the marked weighted data include the reduced
central analytic space, its component multiplicities, the central
log space of the normalized cyclic root, and the old cone cells
before their semistable subdivision.  A coarsening is an actual
proper toric morphism between the displayed monoid charts.  These
data retain more than the dual incidence complex.

\begin{proposition}[Exact weighted toric descent and compact occurrence]\label{cdp:st:C.2.UN.76}
\StatementRecord{cdp:st:C.2.UN.76}{proved}{Historical events: Claude Fable 5.2 (Anthropic), effort high, 16 November 2026: carried as background without fault found. GPT-6 Astra (OpenAI), ChatGPT Pro mode, effort Pro, S6\_09\_PAPERS, 17 November 2026: targeted proof check. Also checked in this writing. Pre-continuation writing events: GPT-6 Astra (OpenAI), ChatGPT Pro mode, effort Pro, S6\_10\_PAPERS, 18 November 2026: targeted proof check; used and rechecked in this run. Continuation metadata events: GPT-6 Astra (OpenAI), ChatGPT Pro mode, effort Pro, S6\_10\_PAPERS continuation, 19 November 2026: used and rechecked in this run; status-vocabulary reconciliation; metadata reconciliation only; no new mathematical check. \emph{Pre-continuation recorded proof scope.} The saturated monoid lattice, all normalized root branches and equivariant height-one star subdivisions give the actual semistable root and retained coarse quotient. The canonical valuation is M times the common ratio minus one; its root character equals the integral affine determinant, forcing precisely the two stated ratios. Finite covers of the original open strata prove their Euler numbers vanish. The exact affine order equations, simultaneous base shifts and real-linear period frame retain the full finite action with an everywhere positive determinant. At infinity the reflected fan, generator reversal and torus inversion give the conjugate data. A finite-stabilizer logarithmic norm gauge makes its coefficients constant in the original coarse monoid algebra, yielding exactly the smooth original weighted invariant chart and compact quotient}{the hypotheses, notation and definitions stated here; \cref{cdp:st:C.2.UN.71}; \cref{cdp:eq:C.WF.9}; \cref{cdp:eq:C.WF.11}; \cref{cdp:eq:C.WF.28}; \cref{cdp:eq:C.WF.5}; \cref{cdp:eq:C.WF.6}; \cref{cdp:eq:C.WF.4}; \cref{cdp:eq:C.WF.7}; \cref{cdp:st:C.2.UN.68}; \cref{cdp:eq:C.WF.8}; \cref{cdp:st:C.2.UN.73}; \cref{cdp:eq:C.WF.10}; \cref{cdp:eq:C.WF.14}; \cref{cdp:eq:C.WF.15}; \cref{cdp:eq:C.WF.17}; \cref{cdp:eq:C.WF.19}; \cref{cdp:eq:C.WF.22}; \cref{cdp:eq:C.WF.25}; \cref{cdp:eq:C.WF.24}; \cref{cdp:eq:C.WF.26}; \cref{cdp:eq:C.WF.18}}

Let \(f:X\to\Delta_t\) be a proper degeneration of connected complex
two-tori, with smooth total space and full fibre
\(F=\sum_i m_iV_i\) whose reduced support is simple normal crossing.
Suppose there is a triple point and that the actual orders of a
generator of \(f_*K_X\) satisfy
\[
 \operatorname{div}(\Omega)=\sum_i k_iV_i,\qquad
 \frac{k_i+1}{m_i}=\lambda\quad\hbox{for every }i.
\tag{C.WF.1}\label{cdp:eq:C.WF.1}
\]
Then the following conclusions hold.

The normalized cyclic root, followed by the explicit height-one
subdivision below, is a positive canonical-flat toric degeneration
of \ref{cdp:st:C.2.UN.71}.  The original marked weighted support is exactly its
retained coarse-cone quotient by the root group.  Thus this branch
is described by the toric data \((N,\Gamma,\mathscr P,c)\) of \ref{cdp:st:C.2.UN.71},
the retained periodic coarse triangles, and the actual cyclic
monomial maps \eqref{cdp:eq:C.WF.9}--\eqref{cdp:eq:C.WF.11}, with their specified smooth
coarse invariant charts \eqref{cdp:eq:C.WF.28}.

There are precisely two possible common ratios:
\[
 \lambda=1,\quad k_i=m_i-1;
 \qquad\hbox{or}\qquad
 \lambda=\tfrac12,\quad k_i=\tfrac{m_i}{2}-1.
\tag{C.WF.2}\label{cdp:eq:C.WF.2}
\]
In the second case every \(m_i\) is even.  Every positive-dimensional
open stratum of the original support has Euler number zero.
Consequently
\[
             e(F_{\mathrm{red}})
               =\#\{\hbox{original triple points}\}>0.
\tag{C.WF.3}\label{cdp:eq:C.WF.3}
\]
In particular this entire weighted branch contains no zero-Euler
support.

Every such marked weighted log support has a realization in a
smooth compact complex threefold over \(\PP^1\), with connected
two-torus fibres away from two values.  The two marked weighted
supports are the given one and its complex conjugate.  The actual
old-cone contractions and finite invariant-ring maps are retained
at both ends, so the reduced analytic support, all its gluing maps,
and every full-fibre multiplicity are unchanged.  The total Euler
number of this compact realization is twice \eqref{cdp:eq:C.WF.3}.

The exact original nonreduced fibre and its entire neighbourhood
are recovered by the original coarsening and quotient morphisms.
The compact occurrence assertion preserves the marked reduced
log support and multiplicity cycle; it does not identify an
arbitrary original higher infinitesimal thickening or full
degeneration germ with the new compact realization.
\end{proposition}



\begin{proof}
\noindent\textbf{Step 1.} \par\medskip\noindent\textbf{The actual normalized-root and coarse-cone maps.}
Choose \(M\) divisible by all \(m_i\), and use \(t=s^M\).
Locally a holomorphic root of the nonzero unit in \(f\) gives
actual coordinates in which
\[
 t=x_1^{m_1}\cdots x_r^{m_r},\qquad 1\le r\le3.
\tag{C.WF.4}\label{cdp:eq:C.WF.4}
\]
The coordinate change is invertible because its derivative in
the changed normal direction is a unit.  Keep its choice and
all its transition units in the normalization diagram.
The normalized root has \(\gcd(m_1,\ldots,m_r)\) local components.
Each is the normal toric chart with cocharacter lattice and cone
\[
 L=\{(v,h)\in\Z^r\oplus\Z:\ \textstyle\sum_i m_iv_i=Mh\},
 \qquad \sigma=\{v_i\ge0\}.
\tag{C.WF.5}\label{cdp:eq:C.WF.5}
\]
The remaining \(3-r\) coordinates are unchanged.  Its old rays
are \((M e_i/m_i,1)\), which are primitive because their height
is one.  Formula \eqref{cdp:eq:C.WF.5}, its saturated dual monoid, and the
characters \(x_i,s\) are the actual normalization maps; no
unramified branches are discarded.

\noindent\textbf{Step 2.} The height-one slice is an interval or a lattice triangle.
Subdivide its edges at every lattice point and triangulate
using all its lattice points.  Every resulting empty lattice
triangle is unimodular: its Euclidean lattice area is \(1/2\)
by the elementary lattice-area formula.  Equivalently its
three lifted height-one vectors are an integral basis of \(L\).
Thus the resulting charts are literally
\(\C^3\), with \(s=u_1u_2u_3\), or their two-factor and
one-factor versions.

\noindent\textbf{Step 3.} For the actual original charts, perform this refinement by
successive height-one star subdivisions.  Each is the normalized
blowup of its monomial ideal.  Changing a boundary equation by a
holomorphic unit multiplies each monomial generator by a unit
and leaves that ideal unchanged.  The normalized blowup maps
therefore glue with the original tangent-dependent transition
units retained.  We have not treated those original overlaps
as constant toric maps.  After all height-one lattice points
have been inserted the same empty-triangle argument gives the
unimodular refinement.

\noindent\textbf{Step 4.} The subdivisions can be chosen equivariantly.  A stabilizer of an old
coarse cone in the actual cyclic root fixes its old rays
individually: they lie over the distinct original SNC
components at that stratum.  It therefore acts trivially on
the height-one simplex.  Choose one subdivision on each orbit
of old triangles after the common edge subdivisions, and
translate it to the other members of the orbit.  Equivalently
perform the corresponding star operations on their whole
orbits; the old component labels and the trivial simplex
stabilizers give the same restrictions on every shared face.
The proper normalized blowup maps are compatible on overlaps.  They give
an actual proper equivariant morphism
\[
 \widetilde X\longrightarrow X_M
       :=\operatorname{Norm}(X\times_{\Delta_t}\Delta_s),
 \qquad X_M/\mu_M=X.
\tag{C.WF.6}\label{cdp:eq:C.WF.6}
\]
The equality on the right is the actual finite analytic
quotient: invariant meromorphic functions are those on \(X\),
and normality of the smooth \(X\) makes its integral invariant
algebra exactly \(\OO_X\).  The retained coarsening in
\eqref{cdp:eq:C.WF.6} is essential; taking the quotient of the subdivided
space alone would usually retain additional divisors.
The semistable support is embedded SNC globally.  Its old
rays retain the distinct original component labels.  A new
ray retains its original stratum and its lattice position
in that stratum's old simplex.  Transition maps and root
inertia fix these positions on each old cone.  Rays occurring
together in a subdivided cone therefore cannot become two
branches of one global prime.  The star charts of each
prime glue as the smooth toric boundary charts of that
labelled ray.

\par\medskip\noindent\textbf{The canonical orders, with their exact normalization.}
At a height-one ray \((v,1)\), direct differentiation of
\eqref{cdp:eq:C.WF.4} gives
\[
 \operatorname{ord}_{v}(\rho^*\Omega)
       =\sum_i(k_i+1)v_i-1
       =M\lambda-1.
\tag{C.WF.7}\label{cdp:eq:C.WF.7}
\]
For an old ray this is \(M(k_i+1)/m_i-1\).
For every new ray it is the affine interpolation of those
old values.  This also proves the formula in a two-factor
chart after wedging the unchanged tangent differential.
The pulled-back nonzero coefficient is a unit throughout
each normalization chart.  Hence
\[
       \widetilde\Omega
          =s^{-(M\lambda-1)}\rho^*\Omega
\tag{C.WF.8}\label{cdp:eq:C.WF.8}
\]
is a nowhere-zero holomorphic three-form on \(\widetilde X\).
Its exponent is an integer by \eqref{cdp:eq:C.WF.7}.
An old triple cone contains at least one unimodular triangle,
so the root fibre has a triple point.  \ref{cdp:st:C.2.UN.68} and \ref{cdp:st:C.2.UN.71} apply to
this actual reduced semistable root.  In particular its marked
central log space is the toric space of
\((N,\Gamma,\mathscr P,c)\).  The old cones give a periodic
coarsening \(\mathscr P_0\) of this same developed triangulation.

\noindent\textbf{Step 5.} Notice the normalization of \eqref{cdp:eq:C.WF.1}.  In a weighted fibre it
is not necessary that some \(k_i\) be zero.  Positivity of
the \(k_i+1\) gives \(\lambda>0\).  If \(\lambda>1\), every
integer \(k_i-m_i=(\lambda-1)m_i-1\) is nonnegative, so
\(\Omega/t\) is holomorphic on the smooth \(X\).  This
contradicts that \(\Omega\) generates \(f_*K_X\).
Thus \(0<\lambda\le1\).

\par\medskip\noindent\textbf{The exact central finite action.}
Write \(\zeta=\exp(2\pi i/M)\).  Lift the root generator
to the developed affine plane and to the path-product
trivialization of the toric cocycle in \ref{cdp:st:C.2.UN.71}.  The resulting
maps on the open torus are
\[
 g(z,s)=(a\,s^{b_0}z^A,\zeta s),\qquad
 A\in\operatorname{GL}(N),\quad b_0\in N,\quad a\in T_N,
 \qquad p\longmapsto Ap+b_0.
\tag{C.WF.9}\label{cdp:eq:C.WF.9}
\]
They preserve both \(\mathscr P\) and the marked coarsening
\(\mathscr P_0\).  Here \(z^A\) denotes the map whose
pullback on characters is \(A^{\mathsf T}\).

\noindent\textbf{Step 6.} We justify the central comparison, including its constants.
Each component map preserves its full toric boundary, so on
its dense torus it is a torus translation followed by the
integral map of its fan.  This follows directly by pulling
back a character: its divisor is a specified integral
combination of boundary primes, and the quotient of two
functions with that divisor is a holomorphic unit on the
compact component, hence a constant.  In the global
characters at a vertex \(p\), its dense component coordinates
are \(z^m s^{-\langle m,p\rangle}\).  Substitution in
\eqref{cdp:eq:C.WF.9} gives the component constant
\(a\,\zeta^{-(Ap+b_0)}\).
Conversely the actual double-curve map fixes the tangential
character comparison, and its marked normal-product log
section fixes the other character comparison.  In the
three-factor chart these are exactly the two substitutions
just displayed.  Thus after the \ref{cdp:st:C.2.UN.71} path-product
trivialization the quantities
\(a_p\zeta^{Ap+b_0}\) agree across every edge and equal one
constant \(a\).  This proves \eqref{cdp:eq:C.WF.9} on the entire central
log atlas.  It is a central-log comparison; no conjugacy of
the arbitrary original whole germ has been asserted.

\noindent\textbf{Step 7.} The lifted \(g\) need not have order \(M\) on this developed
cover.  Instead it normalizes the actual period maps
\(g_\gamma(z,s)=(c_\gamma s^\gamma z,s)\) and satisfies
\[
 \begin{split}
 A\Gamma&=\Gamma,\qquad A^M=I,\\
 c_\gamma^A&=\zeta^{A\gamma}c_{A\gamma},\\
 g^M&=g_{\gamma_0}\quad\hbox{for a retained }\gamma_0\in\Gamma .
 \end{split}
\tag{C.WF.10}\label{cdp:eq:C.WF.10}
\]
All are equalities of monomial maps.  For clarity the last
one means
\[
 \gamma_0=\sum_{j=0}^{M-1}A^{M-1-j}b_0,\qquad
 \prod_{j=0}^{M-1}(a^{A^{M-1-j}})
       \ \zeta^{\sum_{j=0}^{M-1}jA^{M-1-j}b_0}
       =c_{\gamma_0}.
\tag{C.WF.11}\label{cdp:eq:C.WF.11}
\]
The central log comparison preserves these relations.
Indeed two monomial maps with the same affine map and base
character differ by a constant torus translation; their
actions on a height-one component determine that constant.
Equality of the actual central maps therefore implies the
displayed equality on every toric chart.

\noindent\textbf{Step 8.} The relative residue of \eqref{cdp:eq:C.WF.8} transforms under the root
generator by \(\zeta^{-M\lambda}\).  On the developed toric
log atlas it is a common nonzero scalar times
\(d\log z_1\wedge d\log z_2\), and \eqref{cdp:eq:C.WF.9} transforms this
residue by \(\det A\).  The scalar is common because the
developed transition maps are translations and the double
residues agree.  Consequently
\[
                  \zeta^{-M\lambda}=\det A\in\{1,-1\}.
\tag{C.WF.12}\label{cdp:eq:C.WF.12}
\]
Together with \(0<\lambda\le1\), this proves \eqref{cdp:eq:C.WF.2},
including the evenness of every \(m_i\) in the second case.

\par\medskip\noindent\textbf{Euler number on the actual original strata.}
Let \(V_i^\circ\) be an original component with all double
curves removed.  The strict root components above it have
finite unramified maps to \(V_i^\circ\); locally these are
exactly the branches of \(x_i^{m_i}=s^M\) with the
nonzero units retained.  Their total degree is \(m_i\).
The semistable subdivisions change nothing there.
On every old root component this inverse image is its full
toric complement \((\C^*)^2\).  The sum of their Euler
numbers is therefore zero, and finite-cover multiplicativity
gives \(e_c(V_i^\circ)=0\).

\noindent\textbf{Step 9.} For a connected original double curve \(C\), remove its
original triple points and write \(C^\circ\).  The normalized
root over this open curve is locally the two-dimensional
monomial root times \(C^\circ\).  Its subdivision is the
corresponding chain times that curve.  It creates no triple
point over the interior of \(C^\circ\).  Choose a boundary
curve of an old strict root component which dominates \(C\);
boundary curves supported over an original triple point
are not used here.  Every such dominating curve is a toric
\(\PP^1\), and the
preimage of the removed original endpoints is precisely
its two toric endpoints.  Thus its open part is \(\C^*\)
and maps finitely and unramifiedly to \(C^\circ\).
This remains true component by component when the root
cover is disconnected over \(C^\circ\).  Its positive total
degree gives \(e_c(C^\circ)=0\).  Additivity on the original
SNC stratification now proves \eqref{cdp:eq:C.WF.3}.  This argument
uses the retained old cells, rather than the Euler number
of the subdivided root fibre.

\par\medskip\noindent\textbf{The period equations retaining the finite action.}
Choose \(N=\Z^2\), \(\Gamma=Q\Z^2\) and logarithms \(C\)
of \(c\), as in \ref{cdp:st:C.2.UN.73}.  Initially put
\[
 \begin{gathered}
 \Lambda= \Z^2,\quad \Lambda_0=Q^{-1}\Z^2,\quad
 B=(2\pi i)^{-1}Q^{-1}C,\\
 D=Q^{-1}AQ,\quad
 d=Q^{-1}b_0,\quad
 a_0=(2\pi i)^{-1}Q^{-1}\log a .
 \end{gathered}
\tag{C.WF.13}\label{cdp:eq:C.WF.13}
\]
The additive action is \(w\mapsto Dw+d\,l+a_0\),
where \(l=\log s/(2\pi i)\).  Replace these real lattices
by their actual images under one real invertible matrix
which makes the finite-order \(D\) orthogonal for an
averaged positive metric.  Replace \(B,D\) by their
conjugates and \(d,a_0\) by their images under the same
matrix.  We retain the names
\(\Lambda\subset\Lambda_0\), \(B,D,d,a_0\) for these
actual images.  In particular no integral lattice is
replaced by the abstract standard lattice in this step.
Let \(T\) be the actual real matrix taking this vector
coordinate back to \((2\pi i)^{-1}\log z\).

\noindent\textbf{Step 10.} For \(\ell\in\Lambda\), \eqref{cdp:eq:C.WF.10} becomes
\[
       (DB-BD-D/M)\ell\in\Lambda_0.
\tag{C.WF.14}\label{cdp:eq:C.WF.14}
\]
Define
\[
 d_M=\sum_{j=0}^{M-1}D^{M-1-j}d\in\Lambda,\quad
 a_M=\sum_{j=0}^{M-1}D^{M-1-j}a_0,\quad
 e_M=\sum_{j=0}^{M-1}jD^{M-1-j}d .
\]
The order relation is precisely
\[
                   a_M+e_M/M-Bd_M\in\Lambda_0 .
\tag{C.WF.15}\label{cdp:eq:C.WF.15}
\]
These are lattice-valued identities obtained by taking
the chosen actual logarithms of \eqref{cdp:eq:C.WF.10}--\eqref{cdp:eq:C.WF.11}.

\par\medskip\noindent\textbf{An equivariant compact middle family.}
In the orthogonal coordinates put
\[
 J=\begin{pmatrix}0&1\\-1&0\end{pmatrix},\qquad
 DJ=(\det D)JD.
\]
Choose \(k\ge1\) with \(\zeta^k=\det D\).
Choose \(R>0\) so that the symmetric part \(S\) of
\(\operatorname{Im}B+RI\) is positive definite, and put
\[
 B'=B+iRI,\quad a'_0=a_0+iRd,\quad
 \kappa=e^{-2\pi R},\qquad
 V_s(q)=q+s^kJ\overline q,\quad H_s=I+s^kJ.
\tag{C.WF.16}\label{cdp:eq:C.WF.16}
\]
The simultaneous shift of \(a_0\) is required: it is
exactly the substitution \(u=\kappa s\) in the original
affine map.  Equations \eqref{cdp:eq:C.WF.14}--\eqref{cdp:eq:C.WF.15} remain true
with \(B',a'_0\), since their added terms cancel.

\noindent\textbf{Step 11.} On the whole punctured line use the actual period lattice
\[
 \mathcal L_s=V_s(\Lambda_0)+
       \{\,V_s(B'\ell)+l(s)V_s(\ell):\ell\in\Lambda\,\}.
\tag{C.WF.17}\label{cdp:eq:C.WF.17}
\]
Changing the logarithm adds an angular period.  The
determinant calculation of \ref{cdp:st:C.2.UN.73} applies in the chosen
real coordinates, with a fixed positive basis factor.
Writing \(\operatorname{Im}B'=S+aJ\),
\(r=|s|^k\), \(x=\operatorname{Re}(s^k)\) and
\(y=-\log|s|/(2\pi)\), its variable factor is
\[
 (1+r^2)^2\det S+(1-r^2)(1+r^2)y\operatorname{tr}S
 +y^2(1-r^2)^2+\bigl((1+r^2)a+2xy\bigr)^2>0 .
\tag{C.WF.18}\label{cdp:eq:C.WF.18}
\]
Thus \eqref{cdp:eq:C.WF.17} gives a proper smooth torus family over
\(\C^*\), including the points where \(H_s\) is singular.

\noindent\textbf{Step 12.} The finite action on this family is the actual map
\[
 (v,s)\longmapsto
 \bigl(Dv+V_{\zeta s}(a'_0)
                  +l(s)V_{\zeta s}(d),\ \zeta s\bigr).
\tag{C.WF.19}\label{cdp:eq:C.WF.19}
\]
The identity \(DV_s=V_{\zeta s}D\) follows from the
choice of \(k\).  Equation \eqref{cdp:eq:C.WF.14} proves that
\eqref{cdp:eq:C.WF.19} preserves the actual period subgroups.
Iterating it \(M\) times gives the period
\[
 V_s(a'_M+e_M/M)+l(s)V_s(d_M)
       =V_s(B'd_M+n)+l(s)V_s(d_M),\qquad n\in\Lambda_0,
\tag{C.WF.20}\label{cdp:eq:C.WF.20}
\]
by \eqref{cdp:eq:C.WF.15}.  Hence its induced action has order
dividing \(M\).  Its base action has order \(M\), so
the action is faithful and free over \(\C^*\).
All these computations use only real coefficients
\(j/M\) inside \(V_s\); the complex number \(l(s)\)
is kept outside.  No complex-linearity of \(V_s\)
is presumed.

\par\medskip\noindent\textbf{The exact zero-end filling and original action.}
Near zero use
\[
                    z=\exp(2\pi iT H_s^{-1}v).
\tag{C.WF.21}\label{cdp:eq:C.WF.21}
\]
This divides by exactly the angular lattice.
The other periods give the variable toric holonomies
of \ref{cdp:st:C.2.UN.73}, with central values \(\kappa^\gamma c_\gamma\).
In this coordinate \eqref{cdp:eq:C.WF.19} is
\[
 (z,s)\longmapsto
   \bigl(a(s)s^{b_0}z^A,\zeta s\bigr),\qquad
 a(s)=\exp\bigl(2\pi iT H_{\zeta s}^{-1}
                              V_{\zeta s}(a'_0)\bigr).
\tag{C.WF.22}\label{cdp:eq:C.WF.22}
\]
The coefficient is a holomorphic nonzero vector and
\(a(0)=\kappa^{b_0}a\).
Using the actual local base \(u=\kappa s\)
therefore restores exactly \(a\), \(c_\gamma\),
and the given central root action.  Both
\(\mathscr P\) and its marked old coarsening are
preserved.  Apply the variable-unit toric chart
construction of \ref{cdp:st:C.2.UN.73}, first with the subdivision
and then with its proper coarsening.  The monomial
maps \eqref{cdp:eq:C.WF.22} extend to both atlases and commute
with the actual coarsening morphism.

\par\medskip\noindent\textbf{The reflected and conjugate infinity end.}
Set \(q=s^{-1}\), replace the vector frame by
\(v_\infty=q^kv\), and put
\[
 H^\infty_q=q^kI+J,\qquad
 V^\infty_q(w)=q^kw+J\overline w .
\]
Reverse the two nonangular period generators as
in \ref{cdp:st:C.2.UN.73}.  The new period group has generators
\[
 H^\infty_q\Lambda_0,\qquad
 -V^\infty_q(B'\ell)+l(q)H^\infty_q\ell
                 \quad(\ell\in\Lambda).
\tag{C.WF.23}\label{cdp:eq:C.WF.23}
\]
The exact finite map in this frame is
\[
 (v_\infty,q)\longmapsto
 \bigl((\det D)Dv_\infty+
       V^\infty_{\zeta^{-1}q}(a'_0)
       -l(q)V^\infty_{\zeta^{-1}q}(d),\
       \zeta^{-1}q\bigr).
\tag{C.WF.24}\label{cdp:eq:C.WF.24}
\]
Since \((\det D)D H^\infty_q=
H^\infty_{\zeta^{-1}q}D\), the exponential
coordinate with frame \(H^\infty_q\) has central
affine vertex map \(p\mapsto Ap-b_0\).
Use the reflected old and subdivided fans
\(-\mathscr P_0\) and \(-\mathscr P\).
After the local base change \(v_{\mathrm{base}}=\kappa q\),
the central map is exactly
\[
 (z_\infty,v_{\mathrm{base}})
  \longmapsto
  (\overline a^{-1}v_{\mathrm{base}}^{-b_0}
                   z_\infty^A,\ \zeta^{-1}v_{\mathrm{base}}),
\tag{C.WF.25}\label{cdp:eq:C.WF.25}
\]
and its period constants are \(\overline c_\gamma\).
Here \(\overline{a'_0}=\overline{a_0}-iRd\)
and \(-d\,l(q)=-d\,l(v_{\mathrm{base}})+iRd\),
so the recorded shifts cancel exactly.

\noindent\textbf{Step 13.} Finally apply the actual torus inversion
\(z_\infty\mapsto z_\infty^{-1}\) and relabel the
period group by \(\gamma\mapsto-\gamma\).
This sends the reflected fans to the original
ones, sends the map \eqref{cdp:eq:C.WF.25} to the conjugate
of \eqref{cdp:eq:C.WF.9}, and leaves the conjugate period
constants correct because
\(\overline c_{-\gamma}^{-1}=\overline c_\gamma\).
Thus the full coarse and subdivided central
log spaces, including the finite action, are
the complex conjugates of the prescribed data.
The two base coordinates satisfy
\(u\,v_{\mathrm{base}}=\kappa^2\).

\par\medskip\noindent\textbf{Finite stabilizers and smooth coarse quotients.}
We check the local smoothness assertion rather than
inferring it from an abstract quotient.  At a point
of a retained coarse central chart, take its finite
point stabilizer after quotienting by \(\Gamma\).
The actual Kummer chart \eqref{cdp:eq:C.WF.4}--\eqref{cdp:eq:C.WF.5} shows that
its inertia fixes each normal ray and is the identity
on the tangent stratum: the latter cover is unramified.
This property is retained by the exact central log
comparison.  Its monomial lattice map is therefore
identity on the ray span and on its quotient, and,
being of finite order, is the identity on the whole
lattice.  We use a sufficiently small invariant
neighbourhood of this point.  A stabilizer of a
lower-dimensional cone which moves points on its
torus stratum is not used as a point stabilizer.
In the newly constructed chart the point stabilizer's
action on a character \(\chi^m\)
is multiplication by a holomorphic unit \(c_m(s)\).
These coefficients depend only on \(s\), by
\eqref{cdp:eq:C.WF.22} or \eqref{cdp:eq:C.WF.24}, and are multiplicative in \(m\).
Their central constants are the original stabilizer
characters.  For a generator of order \(d\), with
base character \(\xi\), let
\[
 L_m(s)=\log\!\frac{c_m(s)}{c_m(0)},\qquad
 B_m(s)=\frac1d\sum_{j=0}^{d-1}jL_m(\xi^js).
\tag{C.WF.26}\label{cdp:eq:C.WF.26}
\]
The logarithm is the unique one vanishing at zero.
The finite-order identity gives
\(\sum_jL_m(\xi^js)=0\), and hence
\[
 B_m(\xi s)-B_m(s)=L_m(s).
\tag{C.WF.27}\label{cdp:eq:C.WF.27}
\]
The coordinate change
\(\chi^m\mapsto e^{-B_m(s)}\chi^m\)
therefore makes every stabilizer coefficient constant.
The functions \(B_m\) are additive in \(m\), so this is
an actual automorphism of the saturated monoid algebra,
with the displayed inverse.  The base character has
constant coefficient, so its \(B_m\) is zero and the
change fixes \(s\).  Averaging has only been used for
this finite chart stabilizer, not for the developed
lift with \(g^M=g_{\gamma_0}\).

\noindent\textbf{Step 14.} We can now use precisely the original coarse invariant
monoid chart, retained in \eqref{cdp:eq:C.WF.5}--\eqref{cdp:eq:C.WF.6}:
\[
 \left(\OO_{X_M,\mathrm{coarse}}\right)^{\mathrm{Stab}}
        =\C\{x_1,x_2,x_3\},\qquad
       s^M=x_1^{m_1}\cdots x_r^{m_r}.
\tag{C.WF.28}\label{cdp:eq:C.WF.28}
\]
The same equation holds in the new chart after the
explicit unit change \eqref{cdp:eq:C.WF.26}.  This proves smoothness
of the quotient total space and its exact component
multiplicities.  It also proves the embedded SNC
condition for its reduced fibre, including the original
component labels.  Reduction and finite invariants
commute in characteristic zero, so the central
coarsening and quotient identify the actual original
reduced analytic support and all its gluing maps.

\par\medskip\noindent\textbf{Compactness and the precise retained scope.}
The determinant \eqref{cdp:eq:C.WF.18} and real fundamental
parallelepipeds give a proper middle torus family
on every compact annulus.  The bounded-unit
toric filling argument of \ref{cdp:st:C.2.UN.73} gives proper
fillings at both ends.  A toric coarsening has
the same fan support and is proper; the coarse
fillings therefore remain proper.  The exact
exponential maps above identify them with the
middle family on punctured discs.  They glue
to a Hausdorff compact normal root space over
\(\PP^1_s\).  Its finite action is the actual
\eqref{cdp:eq:C.WF.19}, extended by the displayed monomial
maps at both ends.  Its quotient is proper
over \(\PP^1_s/\mu_M=\PP^1\) and is smooth
by \eqref{cdp:eq:C.WF.28}; it is free over the punctured
base, where each quotient fibre is one of
the actual connected two-tori.

\noindent\textbf{Step 15.} The two marked reduced supports and their
multiplicities are the prescribed one and
its conjugate.  Euler additivity and the
zero Euler number of a smooth torus give
the stated total Euler number.  On the
original family, \eqref{cdp:eq:C.WF.6} recovers its full
nonreduced fibre by its actual invariant
analytic algebra, not just by a cycle.
The constructed compact family instead
uses the explicit variable period units
in \eqref{cdp:eq:C.WF.17}.  Their higher \(s\)-jets were
not identified with those of an arbitrary
original germ.  No such identification is
needed for the stated compact occurrence
of the marked reduced log support and
full multiplicity cycle.  This completes
all the asserted comparisons.
\end{proof}






\begin{proposition}[A section and ruling fibres give the exact transverse curve blowup]\label{cdp:st:C.2.UN.74}
\StatementRecord{cdp:st:C.2.UN.74}{proved}{Historical events: Claude Fable 5.2 (Anthropic), effort high, 16 November 2026: carried as background without fault found. GPT-6 Astra (OpenAI), ChatGPT Pro mode, effort Pro, S6\_09\_PAPERS, 17 November 2026: targeted proof check. Also checked in this writing. Pre-continuation writing events: GPT-6 Astra (OpenAI), ChatGPT Pro mode, effort Pro, S6\_10\_PAPERS, 18 November 2026: targeted proof check. Continuation metadata events: GPT-6 Astra (OpenAI), ChatGPT Pro mode, effort Pro, S6\_10\_PAPERS continuation, 19 November 2026: used and rechecked in this run; status-vocabulary reconciliation; metadata reconciliation only; no new mathematical check. \emph{Pre-continuation recorded proof scope.} The principal full-fibre line gives fibrewise normal degree minus mP/mT, so the corrected Fujiki-Nakano blowdown applies exactly at equal parent and exceptional multiplicities. The first normal term proves the parent image is smooth and unchanged; the simple zeros on the centre prove all other image components are transverse and are the specified point-blowup targets. The literal p0=xy, r0=x, q0=z chart retains all fibre weights. Adjunction gives the full root line and extension, and Serre duality makes the extension split at a reduced local maximum with marked fibres. The exact support replacement and triple charts give both Euler and triple-count changes}{the hypotheses, notation and definitions stated here; \cite[\S2, Theorem~2' and Proposition~9']{FujikiNakano1971}; \cref{cdp:st:C.2.UN.64}; \cref{cdp:eq:C.RC.8}; \cref{cdp:eq:C.RC.9}; \cref{cdp:eq:C.RC.10}; \cref{cdp:eq:C.RC.12}}

Let $f:U\to\Delta$ be a proper degeneration of connected
complex two-tori from a smooth threefold, with SNC reduced
central support and arbitrary positive full multiplicities.
Let $T$ be a component with a specified holomorphic
projective-line bundle $p:T\to B$ over a smooth compact
curve. Suppose its boundary in the reduced support is
exactly
\[
           C+F_1+\cdots+F_r,\qquad
           C\ \text{a section},\quad F_j=p^{-1}(b_j),
\tag{C.RC.1}\label{cdp:eq:C.RC.1}
\]
where the $b_j$ are distinct. The section is the
intersection with one component $P$, and $F_j$ is the
intersection with a component $Q_j$. The components
$Q_j$ may repeat at different $b_j$, but none is $P$.
All these are actual reduced transverse intersection
divisors; their intersections $C\cap F_j$ are the
specified triple points. Let the full multiplicities
of $T,P,Q_j$ be respectively $m_T,m_P,n_j$.

There is an ordinary smooth curve contraction along
the ruling, inverse to an ordinary blowup and equal
to an isomorphism off $T$, exactly when
$m_T=m_P$. In particular this always holds for a
full reduced fibre. Set their common value equal
to $m$. The resulting actual map is
\[
              \pi:U\longrightarrow U',\qquad
              f=f'\pi,\qquad \pi(T)=B\subset P'.
\tag{C.RC.2}\label{cdp:eq:C.RC.2}
\]
The threefold $U'$ is smooth and $f'$ is proper with
the same punctured torus family. Its reduced central
support is SNC. The map $P\to P'$ is an isomorphism.
Each $Q_j\to Q'_j$ is the ordinary point blowup at
$b_j\in B$, with all indicated points blown up when
the same component repeats. The curve $B$ lies in
the smooth surface $P'$, avoids every triple point
of the descended support, and meets its double
curves transversely, exactly at the $b_j$.

All full multiplicities are retained: the parent
has multiplicity $m$ and $Q'_j$ has multiplicity
$n_j$. At $b_j$ there are actual coordinates
$(p_0,q_0,r_0)$ with
\[
 \begin{gathered}
 P'=(p_0=0),\quad Q'_j=(q_0=0),\quad
 B=(p_0=r_0=0),\qquad f'=p_0^m q_0^{\,n_j},\\
 p_0=xy,\quad r_0=x,\quad q_0=z,\qquad
                         f=x^m y^m z^{\,n_j}
 \end{gathered}
\tag{C.RC.3}\label{cdp:eq:C.RC.3}
\]
in the blowup chart containing $C\cap F_j$.
Every other component and every neighbourhood away
from the indicated exceptional configurations are
unchanged.

Write
$L=N_{B/P'}$ and $M=N_{P'/U'}|_B$.
The actual normal extension and full-fibre
trivialization are
\[
 \begin{gathered}
 0\longrightarrow L\longrightarrow
       V=N_{B/U'}\longrightarrow M\longrightarrow0,
 \qquad (T,C)=(\PP_B(V),\PP_B(L)),\\
 N_{T/U}\simeq\OO_{\PP_B(V)}(-1),\qquad
 M^{\otimes m}\simeq
             \OO_B\left(-\sum_{j=1}^r n_jb_j\right).
 \end{gathered}
\tag{C.RC.4}\label{cdp:eq:C.RC.4}
\]
Thus the full fibre weights on the ruling-fibre
boundary enter the actual root line $M$; they are
not discarded. In particular $m$ divides
$\sum_jn_j$. For actual canonical orders $k_P,k_{Q_j}$
on the descended family, the section subline is
\[
 L\simeq K_B\otimes M^{-(k_P+1)}
                    \otimes\OO_B\left(-\sum_jk_{Q_j}b_j\right),
             \qquad k_T-k_P=1.
\tag{C.RC.5}\label{cdp:eq:C.RC.5}
\]
The entire extension in \eqref{cdp:eq:C.RC.4}, with its endpoint
identifications, is retained.

For a full reduced fibre, put $Z=\sum_jb_j$. Then
these maps specialize to
\[
 \begin{gathered}
 M\simeq\OO_B(-Z),\qquad
 L\simeq K_B\otimes
              \OO_B\left(\sum_j(k_P-k_{Q_j}+1)b_j\right),\\
 df'|_V:\ V\twoheadrightarrow
            \OO_B(-Z)\hookrightarrow\OO_B.
 \end{gathered}
\tag{C.RC.6}\label{cdp:eq:C.RC.6}
\]
If in addition $T$ is a local maximum of the
canonical order and $r>0$, the extension splits:
$V\simeq L\oplus M$, with its specified subline.

For the underlying reduced supports,
\[
              e(S)=e(S')+2-2g(B),\qquad
              t(S)=t(S')+r,
\tag{C.RC.7}\label{cdp:eq:C.RC.7}
\]
where $t$ counts triple points. The same contraction
exists on a compact total space. Conversely an
ordinary blowup of a smooth curve in a single
central component, transverse to its double curves
and avoiding the triple locus, has exactly the
maps, markings and multiplicities displayed above.
\end{proposition}



\begin{proof}
\noindent\textbf{Step 1.} \par\medskip\noindent\textbf{The actual normal degree and the corrected contraction theorem.}
Restrict the principal full central fibre to $T$.
For a ruling fibre $\ell$ away from the marked
values, the resulting line isomorphism gives
\[
       m_T\deg(N_{T/U}|_\ell)=-m_P,\qquad
       N_{T/U}^{\,m_T}
           \simeq\OO_T\left(-m_PC-\sum_jn_jF_j\right).
\tag{C.RC.8}\label{cdp:eq:C.RC.8}
\]
If $m_T=m_P$, the restriction is precisely
$\OO_{\PP^1}(-1)$. The corrected theorem of
\cite[\S2, Theorem~2' and Proposition~9']{FujikiNakano1971},
in the form used in \ref{cdp:st:C.2.UN.64}, supplies the actual
smooth blowdown along $p:T\to B$. It identifies
$U$ with the ordinary blowup of a smooth curve
in a smooth threefold $U'$. Conversely an
ordinary such blowup has normal degree minus
one, so \eqref{cdp:eq:C.RC.8} forces $m_T=m_P$.

\noindent\textbf{Step 2.} The equality $\pi_*\OO_U=\OO_{U'}$ gives the
unique holomorphic $f'$ in \eqref{cdp:eq:C.RC.2}. Its
properness follows because the inverse image
of a compact set in the disc is the image
under $\pi$ of the corresponding compact
set for $f$. The smooth fibres are unchanged.
It remains to prove the claimed structure
at the centre rather than infer it from a
numerical normal degree.

\par\medskip\noindent\textbf{The parent image is smooth and its whole surface is retained.}
The proper image $P'=\pi(P)$ is a hypersurface
in the smooth target. Its actual Cartier
pullback has the form
\[
                            \pi^*P'=P+aT.
\tag{C.RC.9}\label{cdp:eq:C.RC.9}
\]
Restriction to an ordinary ruling fibre
gives $0=1-a$, since $P\cap T=C$ is a
section and $T|_\ell=\OO_\ell(-1)$.
Thus $a=1$.

\noindent\textbf{Step 3.} A local defining section of $P'$ vanishes
to first order along the generic centre.
Its first normal term on $\PP_B(N_{B/U'})$
has zero divisor exactly $P\cap T=C$,
by \eqref{cdp:eq:C.RC.9}. This term cannot vanish at
any whole fibre: $C$ is a section and
contains none. Its differential in the
normal directions is therefore nonzero
at every point of $B$. Hence $P'$ is
smooth along $B$. Since the centre is
a Cartier curve in $P'$, the actual
restriction of its blowup is an
isomorphism $P\to P'$. In particular
it identifies the section $C$ with
the original curve $B$.

\par\medskip\noindent\textbf{Each ruling-fibre boundary becomes a transverse double curve.}
Consider one of the other image components
$Q'=\pi(Q)$. It does not contain $B$:
otherwise its generic first normal
divisor would meet every ruling fibre,
whereas $Q\cap T$ is a finite union
of fibres. More explicitly, writing
$\pi^*Q'=Q+aT$ and restricting to a
generic ruling fibre gives $a=0$.
Thus $\pi^*Q'=Q$.

\noindent\textbf{Step 4.} Restrict a local defining section of
$Q'$ to $B$. Its pullback to $T$
vanishes exactly on the specified
reduced fibres belonging to $Q$.
At $b_j$ that section therefore
has a simple zero on $B$. Its
differential is nonzero on the
tangent line of $B$. Consequently
$Q'$ is smooth at $b_j$ and is
transverse to $B$. The independent
normal differential of $P'$ makes
$P'$ and $Q'$ transverse as well.

\noindent\textbf{Step 5.} There is only one such other
component at each $b_j$. A second
would have the same fibre in
$T\cap Q$, contradicting the
given SNC intersections and
\eqref{cdp:eq:C.RC.1}. Any further image
component passing through $b_j$
would likewise have a boundary
component on the exceptional
fibre; hence none is missing
from this count. Thus there is
no triple point of the descended
support on $B$.

\noindent\textbf{Step 6.} Choose coordinates with $P'$ and
$Q'$ the first two coordinate
planes and $B$ the stated
coordinate curve. Transversality
just proved permits this choice.
The full fibre is a unit times
$p_0^m q_0^{n_j}$; a local root
of that unit absorbs it into
$p_0$ without changing these
marked loci. Blow up $(p_0,r_0)$.
The exact chart $p_0=xy,r_0=x$
gives \eqref{cdp:eq:C.RC.3}. On $Q'$ this
ideal is the maximal ideal
of $b_j$, so $Q$ is precisely
its ordinary point blowup,
with exceptional fibre $F_j$.
On $P'$ the ideal is Cartier,
giving the already established
isomorphism. These are the
claimed original neighbourhood
maps.

\par\medskip\noindent\textbf{The full normal extension and the weighted line.}
The normal sequence of
$B\subset P'\subset U'$
is the first row of \eqref{cdp:eq:C.RC.4}.
The exceptional divisor of
the ordinary blowup is the
lines projectivization of
this very $V$, with
ambient normal $\OO(-1)$.
The strict parent intersects
it in the specified subline
$\PP_B(L)$. In particular,
with these line conventions,
\[
                 \OO_T(C)\simeq
                      \OO_{\PP_B(V)}(1)\otimes p^*M.
\tag{C.RC.10}\label{cdp:eq:C.RC.10}
\]
The principal full fibre on
$U'$ is $mP'+\sum_Qm_QQ'$.
Restriction to $B$ gives
exactly the last isomorphism
of \eqref{cdp:eq:C.RC.4}, since each
$Q'$ meets $B$ simply at
its listed $b_j$. This
also follows by substituting
\eqref{cdp:eq:C.RC.10} into \eqref{cdp:eq:C.RC.8}.
Taking degrees proves the
stated divisibility. The
line $M$ and its actual
power trivialization are
retained, including any
torsion information.

\noindent\textbf{Step 7.} Restrict the actual vertical
canonical evaluation section
on $U'$ to $B$. Its line
isomorphism is
\[
 K_{U'}|_B\simeq
       M^{k_P}\otimes
             \OO_B\left(\sum_jk_{Q_j}b_j\right).
\tag{C.RC.11}\label{cdp:eq:C.RC.11}
\]
Other components have a
nonvanishing defining section
on $B$ and contribute their
specified trivializations.
Adjunction for $B\subset P'$
then proves the formula for
$L$ in \eqref{cdp:eq:C.RC.5}. The
ordinary curve-blowup
Jacobian contributes $T$
once to $K_U$; the generic
centre lies only in $P'$.
Thus $k_T=k_P+1$.
Changing the canonical
generator by a power of $f'$
adds the full multiplicities
to the orders, and \eqref{cdp:eq:C.RC.4}
shows that \eqref{cdp:eq:C.RC.5} remains
the same line isomorphism.

\noindent\textbf{Step 8.} In the full reduced case
$m=n_j=1$, so $M=\OO_B(-Z)$.
The derivative of
$f'=p_0q_0$ along the
normal bundle of $B$
is the quotient normal
map to $M$, followed by
its inclusion in $\OO_B$
with simple zero at $b_j$.
Away from the marked
points it is surjective.
This proves all maps in
\eqref{cdp:eq:C.RC.6}. For $m>1$ the
ordinary derivative of
$f'$ is zero along the
generic centre; the
quotient in \eqref{cdp:eq:C.RC.4}
instead is the actual
normal map of $P'$,
and has already been
specified independently
of that derivative.

\noindent\textbf{Step 9.} For the splitting assertion,
the extension group is the
actual space $H^1(B,L\otimes
M^{-1})$. Equations \eqref{cdp:eq:C.RC.6}
and Serre duality identify
its dual with
\[
 H^0\left(B,\OO_B\left(
                   \sum_j(k_{Q_j}-k_P-2)b_j\right)\right).
\tag{C.RC.12}\label{cdp:eq:C.RC.12}
\]
At a local maximum
$k_{Q_j}\le k_T=k_P+1$.
Every coefficient in
\eqref{cdp:eq:C.RC.12} is at most
minus one. When $r>0$
the line has negative
degree and no section.
The extension group is
zero, which proves the
actual splitting with
the marked subline.
When $r=0$ the extension
is precisely the one
retained in \ref{cdp:st:C.2.UN.64}.

\par\medskip\noindent\textbf{Euler, triples, the converse and compact descent.}
The actual blowup gives
$S\setminus T\simeq S'\setminus B$,
whereas $T$ is a
projective-line bundle
of Euler number $2e(B)$.
Hence $e(S)-e(S')=e(B)$,
independently of the
full multiplicities.
At each $b_j$ the chart
\eqref{cdp:eq:C.RC.3} has exactly one
triple point upstairs
and only a smooth double
curve downstairs. There
are no other changes,
which proves \eqref{cdp:eq:C.RC.7}.

\noindent\textbf{Step 10.} Conversely, blowing up the
indicated transverse curve
uses exactly these charts.
The exceptional multiplicity
is the generic parent
multiplicity $m$, and
every other weight is
unchanged. The section
and ruling fibres are
exactly the intersections
in \eqref{cdp:eq:C.RC.1}, so all
the displayed maps and
markings are recovered.
As in \ref{cdp:st:C.2.UN.64}, the actual
local blowdowns glue to
the identity off $T$.
This proves the compact
version as well.
\end{proof}






An \emph{overlap-compatible modification record} below retains the
actual analytic charts, all centres and blowup maps, and the
isomorphisms used to reglue them. In a local replacement
$X=U\cup V$, the modified chart $U'$ has a specified overlap
$W'\simeq W=U\cap V$; this isomorphism is the restriction of
the recorded composite map on $U$. The replacement is
$X'=U'\cup_W V$. Thus this datum specifies an analytic
operation, including its neighbourhoods. It is not merely a
list of component graphs or normal degrees. The records used
here are the finite local $\sigma$-records in
\cite[\S\S4.2--4.14]{Kulikov1977}, with the corrections and
extra fibre blowups in \cite[\S\S8--11]{Kulikov1981}.

\begin{proposition}[The Hopf-minimum branch is generated by canonical-flat cores and exact sigma records]\label{cdp:st:C.2.UN.75}
\StatementRecord{cdp:st:C.2.UN.75}{proved}{Historical events: Claude Fable 5.2 (Anthropic), effort high, 16 November 2026: carried as background without fault found. GPT-6 Astra (OpenAI), ChatGPT Pro mode, effort Pro, S6\_09\_PAPERS, 17 November 2026: targeted proof check. Also checked in this writing. Pre-continuation writing events: GPT-6 Astra (OpenAI), ChatGPT Pro mode, effort Pro, S6\_10\_PAPERS, 18 November 2026: targeted proof check; used and rechecked in this run. Continuation metadata events: GPT-6 Astra (OpenAI), ChatGPT Pro mode, effort Pro, S6\_10\_PAPERS continuation, 19 November 2026: used and rechecked in this run; status-vocabulary reconciliation; metadata reconciliation only; no new mathematical check. \emph{Pre-continuation recorded proof scope.} The primitive direct-image generator and Hartogs comparison preserve every surviving divisorial order. The actual proper graph excludes contraction or creation of a Hopf-birational prime. The two-ruling normal bundles, retained self-double branches, missing-fibre insertions and fixed-level termination give the canonical-flat endpoint. The restored overlap maps glue the proper local graphs, and the normal-cone calculation retains singular centres and their entire exceptional fibres. The exact original Russian1977 and1980 Kulikov passages and all stated corrections were opened; no English-edition page-opening credit is claimed}{the hypotheses, notation and definitions stated here; \cref{cdp:st:C.2.UN.65}; \cref{cdp:st:C.2.UN.71}; \cref{cdp:st:C.2.UN.74}; \cref{cdp:st:C.2.UN.73}; \cref{cdp:eq:C.KG.6}; \cref{cdp:st:C.2.UN.19}; \cite[Theorem~1']{Kulikov1981}; \cite[original Russian ed., pp.~1008--1009]{Kulikov1977}; \cref{cdp:st:C.2.UN.64}; \cite[\S4.4]{Kulikov1977}; \cite[\S\S4.10.7--4.10.8,4.14.5]{Kulikov1977}; \cite[\S8]{Kulikov1981}; \cite[\S10]{Kulikov1981}; \cite[\S4.17]{Kulikov1977}; \cref{cdp:eq:C.KG.8}; \cite[\S2]{Kulikov1981}; \cite[\S4.10.8, original Russian ed., pp.~1032--1035]{Kulikov1977}; \cref{cdp:st:C.2.UN.68}}

Let $f:X\to\Delta$ be a proper degeneration of connected
complex two-tori, with smooth total threefold and full
reduced simple-normal-crossing central fibre $S=\sum_iV_i$.
Choose a generator $\Omega$ of $f_*K_X$ near zero and write
\[
             \operatorname{div}(\Omega)=\sum_i k_iV_i,
                  \qquad k_i\ge0,\quad \min_i k_i=0.
\tag{C.KG.1}\label{cdp:eq:C.KG.1}
\]
Let $\mathcal H$ be the set of indices whose component has
a Hopf surface as its smooth minimal model. The following
two conditions are equivalent:
\begin{enumerate}
\item $k_i=0$ for every $i\in\mathcal H$.
\item There is a finite overlap-compatible modification
record from $f$ to a smooth proper torus degeneration
$f_0:X_0\to\Delta$ whose full central fibre $S_0$ is
reduced simple NC and whose canonical line is trivial.
\end{enumerate}
The constructed comparison is an actual bimeromorphic map
over the same disc and induces an isomorphism
\[
                X\setminus S\ \simeq\ X_0\setminus S_0
                         \quad\hbox{over }\Delta^*.
\tag{C.KG.2}\label{cdp:eq:C.KG.2}
\]
Intermediate charts retain their normalization branches;
in particular a type-II move may temporarily give a
self-intersection of a global component.

The endpoint has precisely one of the following forms,
with the actual markings and line maps of the indicated
statements:
\begin{enumerate}
\item a smooth two-torus;
\item a cycle of the marked diagonal Hopf and degree-zero
elliptic ruled pairs in \ref{cdp:st:C.2.UN.65}, with dual edge normals and
the actual complement-lattice product $M=I$;
\item a positive toric central log space in \ref{cdp:st:C.2.UN.71}, given by
$(N,\Gamma,\mathscr P,c)$, with Euler number
$2[N:\Gamma]$ and the specified full toric component pairs.
\end{enumerate}
Every degeneration in the stated branch is recovered by
reversing its finite record on its actual endpoint
neighbourhood. Conversely every such reverse record with
a smooth full reduced simple-NC final fibre gives a
degeneration in this branch. These are inverse analytic
constructions, retaining every centre, overlap and normal
extension.

Here is the exact Euler bookkeeping for these records.
Orient a recorded blowup arrow positively from its target
to its blowup, and negatively in the reverse direction.
For its actual centre $L$ and ideal $I$, put
\[
 \begin{split}
  E_I&=\operatorname{Proj}_L\left(\bigoplus_{a\ge0}
                                      I^a/I^{a+1}\right),\\
  \delta(I)&=e_c(E_I)-e_c(L)
       =\int_L\bigl(e((E_I)_x)-1\bigr)\,de_c(x).
 \end{split}
\tag{C.KG.3}\label{cdp:eq:C.KG.3}
\]
Local terms are evaluated on a finite constructible
subdivision subordinate to the retained charts. Restored
overlaps cancel. If $\epsilon_\nu\in\{1,-1\}$ records
the directions in the forward reduction to $S_0$, then
\[
             e(S)=e(S_0)-\sum_\nu\epsilon_\nu\delta(I_\nu).
\tag{C.KG.4}\label{cdp:eq:C.KG.4}
\]
For an ordinary smooth centre in a smooth threefold,
$\delta=2$ for a point and $\delta=2-2g(L)$ for a compact
curve. Both type-I and type-II moves have total defect
zero. An ordinary ruled-section contraction in \ref{cdp:st:C.2.UN.74} has
Euler change $-(2-2g(B))$ and removes its specified number
of triple points. Thus the exact zero-Euler criterion is
\[
 \sum_\nu\epsilon_\nu\delta(I_\nu)=
 \begin{cases}
 0,&S_0\text{ is of type 1 or 2},\\
 2[N:\Gamma],&S_0\text{ is of type 3}.
 \end{cases}
\tag{C.KG.5}\label{cdp:eq:C.KG.5}
\]
The normal-cone formula also applies to auxiliary
singular-centre charts; their contribution is not
replaced by a smooth-curve formula.

For a given compact total space over a compact curve,
the reduction and its inverse take place on that same
compact base, by the actual isomorphism \eqref{cdp:eq:C.KG.2} on an
annulus. A reverse record on an actual compact core
realization of \ref{cdp:st:C.2.UN.65} or \ref{cdp:st:C.2.UN.73} likewise gives a compact
realization. Its total Euler change is the same as
the central-fibre change in \eqref{cdp:eq:C.KG.4}.
\end{proposition}



\begin{proof}
\par\medskip\noindent\textbf{The canonical generator and its normalized orders.}
The proper direct image $f_*K_X$ is a torsion-free
rank-one coherent sheaf on the disc, hence a line.
On $\Delta^*$ its restriction and base change are the
one-dimensional spaces of top forms on the tori.
A generator therefore has no zero away from $S$.
Its divisor is the effective divisor in \eqref{cdp:eq:C.KG.1}.
If every $k_i$ were positive, $\Omega/f$ would be
holomorphic on the smooth total space, since the
\noindent\textbf{Step 1. full} fibre is reduced. This would contradict
that $\Omega$ generates the direct-image line.
Consequently its minimum is zero. Another generator
is $a(t)\Omega$ for a holomorphic unit $a$, so these
orders are independent of that choice.

For later comparison, let $\phi:X\dashrightarrow X'$
be an actual bimeromorphic map over the disc between
smooth total spaces. Pullback of a holomorphic top
form is holomorphic wherever $\phi$ is holomorphic.
The indeterminacy has codimension at least two;
Hartogs extension of the coefficient in smooth
coordinate volume frames gives a holomorphic form
on all of $X$. The inverse map gives the inverse
pullback, since their compositions agree on their
common biholomorphic open set. The construction
works over every open subdisc, and gives the actual
isomorphism
\[
                f'_*K_{X'}\xrightarrow{\ \phi^*\ } f_*K_X.
\tag{C.KG.6}\label{cdp:eq:C.KG.6}
\]
In particular normalized generators agree under
pullback up to a base unit. At the generic point
of a divisor which survives the comparison,
the map is a biholomorphism of the corresponding
divisorial local rings, so its canonical order is
unchanged.

\par\medskip\noindent\textbf{Why every Hopf-birational component is retained.}
A holomorphic map from $\PP^1$ to a Hopf surface
is constant: lift it to the universal cover
$\C^2\setminus\{0\}$ and use the two constant
coordinate functions. Every rational curve on
a point blowup of that surface is therefore
exceptional over the Hopf minimal model. A finite
sequence of point blowups has only finitely many
such curves. These surfaces cannot be covered by
rational curves. They are also not projective:
their first Betti number is one, unchanged by
point blowups, whereas a projective surface has
even first Betti number.

In the recorded local $\sigma$-maps, a divisorial
exceptional component over a general point of a
curve centre is ruled by complete rational
fibres. A component lying over an isolated
point is projective, as a component of a
projectivized normal cone. Here the contradiction
with the original compact component uses the
proper composite graphs, not a comparison of
local analytic patches. Suppose a strict
transform were removed for the first time.
For a curve centre, choose a general complete
rational exceptional fibre through a generic
strict-transform point. The proper preceding
graph maps this curve meromorphically back to
the original compact component. The map is
nonconstant, since at that generic point the
strict-transform comparison is a local
biholomorphism. Varying the fibre gives
rational curves through a dense open subset
of the original component, contradicting the
finite exceptional-curve description above.
For a point centre, restrict the proper
preceding graph to the compact projective
exceptional component containing the generic
strict-transform piece. Its image back in
the original compact component is closed
and contains that open piece, and therefore
is the whole component. The comparison has
generic degree one. Resolving this surface
graph gives a bimeromorphic comparison with
a projective surface; point blowups preserve
$b_1$, again contradicting $b_1=1$.
Thus the strict transforms through the
finite chart record retain a prime divisor
of that birational type. The restored
overlap maps identify these strict transforms
by their actual generic-point maps. Hence a
component in $\mathcal H$ survives to the
endpoint.

At a canonical-flat smooth endpoint a nowhere-zero
three-form generates $f_{0*}K_{X_0}$, so every
canonical order is zero. Equation \eqref{cdp:eq:C.KG.6} then
proves $k_i=0$ for every $i\in\mathcal H$.
This proves the necessary minimum condition for
the full Hopf birational class, including
arbitrary point blowups. No identification of
different conventions for the notation
\noindent\textbf{Step 2. VII$_0$} is used here.

\par\medskip\noindent\textbf{The finite reduction and the actual allowed charts.}
\ref{cdp:st:C.2.UN.19} identifies every actual class-VII minimal
model in a torus degeneration as Hopf. Thus the
minimum condition just stated implies the
condition in \cite[Theorem~1']{Kulikov1981}.
That theorem applies to proper analytic good
degenerations with trivial fibre canonical
class; its hypotheses do not include
projectivity. The definition of good in
\cite[original Russian ed., pp.~1008--1009]{Kulikov1977} is exactly the smooth
total space and full reduced NC fibre with
smooth components used here.

We specify the comparison maps in that proof.
For a nonboundary $(-1)$-curve $A\subset V_i$
meeting one double curve transversely once
and avoiding the remaining boundary, the
principal-fibre normal sequence is
\[
 0\longrightarrow\OO_{\PP^1}(-1)
 \longrightarrow N_{A/X}
 \longrightarrow\OO_{\PP^1}(-1)
 \longrightarrow0.
\tag{C.KG.7}\label{cdp:eq:C.KG.7}
\]
It splits because $H^1(\PP^1,\OO)=0$.
Blow up $A$; the exceptional divisor is
$\PP^1\times\PP^1$ with normal
$\OO(-1,-1)$. Contract its other ruling
using the corrected contraction theorem in
\ref{cdp:st:C.2.UN.64}. This transfers the point blowup to
the neighbour. The exact adjunction equation
gives $k_i=k_j$ for these two components.
This is the actual type-I map.

For a double curve which is $(-1)$ on both
sides, the same normal bundle calculation
gives the same two-ruling construction.
The principal-fibre line has exactly two
triple points on that curve. If the opposite
components there have orders $k_3,k_4$,
adjunction gives
\[
                         k_1+k_2=k_3+k_4.
\tag{C.KG.8}\label{cdp:eq:C.KG.8}
\]
The ordinary blowup has exceptional full
multiplicity two; the other contraction
removes the old double curve and introduces
the new one between the two opposite
branches. If those branches belong to one
global component, both normalization branches
are retained, as in
\cite[\S4.4]{Kulikov1977}. One does not
discard that self-incidence or call this
intermediate fibre embedded simple NC.

The other contraction charts are the actual
ruled section-and-fibre maps of \ref{cdp:st:C.2.UN.74} and
the plane-line map of \ref{cdp:st:C.2.UN.64}. For a reducible
ruling fibre, the local contractions and
inverse-order restorations in
\cite[\S\S4.10.7--4.10.8,4.14.5]{Kulikov1977}
give the retained finite chart record.
The corrected proof in
\cite[\S8]{Kulikov1981} also blows up the
specified missing fibre components before
contracting the chosen ruled surface.
The new ruled components have canonical
order $k+1$ when the removed component has
order $k$. For a maximum-order component
no such additions are needed.

Here termination is an actual finite
induction, rather than a decrease of an
unspecified complexity. Let $R=\max k_i$.
Every removed old component has level
$k\in\{1,\ldots,R\}$. It produces only
finitely many new components, all at level
$k+1$, and produces none when $k=R$.
There are finitely many level-one removals;
therefore finitely many level-two objects
are ever produced. Continuing through the
fixed finite set of levels proves that all
removals terminate. The components left
have order zero. The corrected endpoint
argument in \cite[\S10]{Kulikov1981}, using
\cite[\S4.17]{Kulikov1977}, gives a smooth
total space and smooth reduced components.
In particular \eqref{cdp:eq:C.KG.8} prevents a minimum
component from acquiring a self-incidence
in a type-II move whose removed edge has
an order above the minimum. The endpoint
generator is therefore nowhere zero.

\par\medskip\noindent\textbf{The source's actual overlap maps glue global graphs.}
The bare two-cover definition of modification
in \cite[\S2]{Kulikov1981} is not used to
infer a global map. The construction itself
provides the following additional data.
In \cite[\S4.10.8, original Russian ed., pp.~1032--1035]{Kulikov1977}
one keeps $V$ fixed, performs the displayed
finite contractions and restorations in
$U$, and obtains $U'$ with an overlap
$W'$ isomorphic to the original $W$ by
the restriction of that same composite.
Section4.14.5 repeats this construction.
The added local fibre blowups of
\cite[\S8]{Kulikov1981} leave this overlap
unchanged.

For a finite $\sigma$-chain, its composite
graph is obtained as the closure of the
common biholomorphic graph in the proper
image of the fibre product of its arrow
graphs. Both projections are proper.
Over $W$ this graph is precisely the
graph of the restored isomorphism.
Glue it to the identity graph over $V$.
The graphs agree on their overlap, and
properness is local on each projection
target. The glued graph is consequently
proper over both $U\cup V$ and
$U'\cup_W V$. It gives actual mutually
inverse bimeromorphic maps. All its arrows
are over the same disc and all their
centres are central. Their restriction
over $\Delta^*$ is an isomorphism.
Iteration proves \eqref{cdp:eq:C.KG.2} and justifies
the generator comparison \eqref{cdp:eq:C.KG.6}.

This argument does not assert a single
global factorization into ordinary
blowups. The finite open covers and
their induced overlap maps are part
of the stated modification record.
Auxiliary singular charts in the
restoration are also retained rather
than silently replaced by smooth ones.

\par\medskip\noindent\textbf{Endpoint classification and the converse record.}
Apply \ref{cdp:st:C.2.UN.65} to the actual endpoint.
Its zero-Euler alternatives give types1
and2, including all the marked line
and integral-lattice comparisons.
Its positive case is \ref{cdp:st:C.2.UN.68}; \ref{cdp:st:C.2.UN.71}
identifies its actual central log
space with the specified toric data
by the displayed chart and residue
maps. This proves the stated complete
endpoint list.

Reverse each recorded arrow and each
overlap isomorphism. They reconstruct
the original neighbourhood literally.
Starting instead with any actual
canonical-flat endpoint and such a
reverse record, all off-central maps
remain isomorphisms; the general
fibres are the same connected tori.
At a smooth full reduced final fibre
the argument following \eqref{cdp:eq:C.KG.6} shows
that every Hopf-birational component
has order zero. This proves both
directions of the generation statement.
In this statement the endpoint is an
actual analytic family. The maps do
not replace its neighbourhood by a
different realization merely because
their central log spaces agree.

\par\medskip\noindent\textbf{The exact Euler defect, including singular centres.}
For an actual blowup with centre ideal
$I$, its exceptional inverse image is
the projectivized normal cone in
\eqref{cdp:eq:C.KG.3}. Outside $L$ it is an
isomorphism. Since $L$ is central,
the reduced central inverse image
replaces precisely $L$ by $E_I$.
Additivity of compactly supported
Euler characteristic proves \eqref{cdp:eq:C.KG.3}
for this arrow, including its
constructible-fibre integral.
For these computations refine the
retained analytic charts by relatively
compact subanalytic coordinate
neighbourhoods, with closures inside
the original charts. Compactness of
the central fibre gives a finite
subcover. Subdivide its overlaps and
centres compatibly, and use inverse
images under the retained proper
graphs; all pieces have finite
subanalytic triangulations. This
refinement also defines the local
terms when a displayed centre is an
open part of a curve. Euler numbers
of arbitrary analytic open covers
are not used.

In one replacement, evaluated on this
finite-type refinement,
$U\cup V\rightsquigarrow
U'\cup_WV$, inclusion--exclusion gives
\[
 \begin{split}
 e(S')-e(S)
   &=e_c(S'\cap U')-e_c(S\cap U)\\
   &=\sum_{\nu\text{ in this chart}}
                         \epsilon_\nu\delta(I_\nu),
 \end{split}
\tag{C.KG.9}\label{cdp:eq:C.KG.9}
\]
because the unchanged $V$ and the
isomorphic restored overlaps cancel.
Equivalently one applies this equality
to the finite subdivision just
described and sums; no boundary
piece is counted twice. The local
arrow contributions telescope, and
then the global replacements
telescope, proving \eqref{cdp:eq:C.KG.4}.

For a smooth centre of dimension $a$
in a smooth threefold the normal
cone is its rank-$(3-a)$ normal
bundle, so its exceptional fibre
is $\PP^{2-a}$. Thus
\[
                \delta(I)=(2-a)e_c(L).
\tag{C.KG.10}\label{cdp:eq:C.KG.10}
\]
For typeI and typeII the common
blowup replaces a compact $\PP^1$
by $\PP^1\times\PP^1$ from either
side, so the two contributions
$+2$ and $-2$ cancel. The ruled
and plane changes follow from
\ref{cdp:st:C.2.UN.74} and \ref{cdp:st:C.2.UN.64}. In an auxiliary
singular-centre chart \eqref{cdp:eq:C.KG.3}
remains the formula: for example,
the blowup of the three coordinate
axes in $\C^3$, with ideal
$(xy,xz,yz)$, has fibre $\PP^1$
over a general point of each axis
and fibre $\PP^2$ over their
common origin. The last defect
is two rather than one. This
illustrates why the full normal
cone was retained.
Substitution of the three endpoint
Euler numbers proves \eqref{cdp:eq:C.KG.5}.

\par\medskip\noindent\textbf{Compact scope.}
Restrict the given compact family
to a small disc at the chosen
value. On a smaller annulus use
the actual isomorphism \eqref{cdp:eq:C.KG.2}
to glue the new proper filling
to the unchanged complement.
The two proper pieces give a
Hausdorff proper analytic family;
their smooth charts coincide on
the overlap. This produces the
claimed compact replacement,
and the inverse record recovers
the original compact family.
The same construction starts
from the explicit compact cores
of \ref{cdp:st:C.2.UN.65} and \ref{cdp:st:C.2.UN.73} and any reverse
record on their actual tubes.
Euler additivity cancels the
unchanged complement and annulus,
so the total-space Euler change
is precisely the central one.
\end{proof}






\begin{proposition}[Actual minimalization at the largest canonical order]\label{cdp:st:C.2.UN.77}
\StatementRecord{cdp:st:C.2.UN.77}{proved}{Historical events: Claude Fable 5.2 (Anthropic), effort high, 16 November 2026: carried as background without fault found. GPT-6 Astra (OpenAI), ChatGPT Pro mode, effort Pro, S6\_09\_PAPERS, 17 November 2026: targeted proof check. Also checked in this writing. Pre-continuation writing events: GPT-6 Astra (OpenAI), ChatGPT Pro mode, effort Pro, S6\_10\_PAPERS, 18 November 2026: targeted proof check; used and rechecked in this run. Continuation metadata events: GPT-6 Astra (OpenAI), ChatGPT Pro mode, effort Pro, S6\_10\_PAPERS continuation, 19 November 2026: used and rechecked in this run; status-vocabulary reconciliation; metadata reconciliation only; no new mathematical check. \emph{Pre-continuation recorded proof scope.} The actual anticanonical point-map formula makes every coefficient-one rational boundary a leaf and excludes maximal Hopf self-incidences. Each last point is either transferred to a non-Hopf receiver, removed by the two-ruling map, or gives a complete ruled fibre with section-parent order one lower. The diagnostic resolution and largest last local branch supply the corrected ruled cleanup, with restored overlaps. The fixed lexicographic vector counts whole retained component blocks and then Hopf point defects; no bound is asserted for auxiliary common blowups. The exact Kulikov constructions and correction to the old-component count were read}{the hypotheses, notation and definitions stated here; \cref{cdp:st:C.2.UN.75}; \cref{cdp:eq:C.HM.2}; \cite[\S4.11.2(III)]{Kulikov1977}; \cref{cdp:eq:C.HM.4}; \cref{cdp:eq:C.HM.3}; \cref{cdp:eq:C.HM.7}; \cite[\S4.5]{Kulikov1977}; \cite[\S\S4.10.7--4.10.8,4.14]{Kulikov1977}; \cite[\S8]{Kulikov1981}; \cref{cdp:eq:C.HM.9}; \cite[Lemma2.19]{Kulikov1977}; \cite[\S4.11.2]{Kulikov1977}; \cite[\S4.13]{Kulikov1977}; \cite[Lemma4.14]{Kulikov1977}; \cite[\S4.16.2]{Kulikov1977}; \cite[\S9]{Kulikov1981}; \cref{cdp:eq:C.HM.5}; \cref{cdp:eq:C.HM.8}}

Let \(f:X\to\Delta\) be a proper degeneration of connected
complex two-tori with smooth total space and full reduced
simple-NC central fibre.  Normalize its canonical generator
as in \ref{cdp:st:C.2.UN.75} and put
\[
        \operatorname{div}(\Omega)=\sum k_iV_i,\qquad
        \min k_i=0,\qquad M=\max k_i .
\tag{C.HM.1}\label{cdp:eq:C.HM.1}
\]
There is a finite overlap-compatible modification record,
over the same disc and unchanged over its punctured part,
whose endpoint has smooth total space and full reduced
normal-crossing fibre, such that every component with
canonical order \(M\) and Hopf minimal model is itself a
smooth minimal Hopf surface.  These Hopf components are
embedded, their boundary curves are disjoint smooth
elliptic curves, and no triple point meets them.
Other endpoint components may retain the two normalized
branches of an NC self-intersection.

All original Hopf-birational prime divisors survive, with
their original orders.  Every whole-component removal block
removes a ruled component of positive order \(q\);
its preliminary additions, when needed, are finitely many
ruled components of order \(q+1\le M\).  There are no such
additions when \(q=M\).  Thus the retained reduced degeneration stages have
no global component of order greater than \(M\).
Auxiliary common blowups in a type-I or type-II diagram
retain their actual multiplicities and orders; the bound
just stated does not apply to those auxiliary divisors.

The record retains its actual centres, local contractions,
normalization branches and restored overlap maps.  Its
reverse recovers the original family.  For a given compact
family the same operations take place on that compact
base, by the actual annular comparison.  Their Euler
changes are exactly the normal-cone terms of \ref{cdp:st:C.2.UN.75}.

The endpoint assertion concerns the maximal-order Hopf
components.  It makes no holomorphic inverse-logarithm
assertion for an arbitrary neighbourhood of a remaining
Hopf leaf.
\end{proposition}



\begin{proof}
\noindent\textbf{Step 1.} \par\medskip\noindent\textbf{The surface divisor and coefficient-one curves.}
At every smooth-total NC stage, let \(H\) be the smooth
normalization of a component of order \(M\) whose minimal
model is Hopf.  The actual adjunction map, with both
preimages of each self-double curve retained, gives
\[
       -K_H=D_H=\sum_C a_C C,\qquad
       a_C=M-k_{\mathrm{other}(C)}+1\ge1 .
\tag{C.HM.2}\label{cdp:eq:C.HM.2}
\]
A self-double branch has coefficient one.  The equality
is the restriction of the common meromorphic volume
form; it retains its line and divisor, rather than only
its numerical class.

\noindent\textbf{Step 2.} Let \(\pi:H\to H_{\min}\) be the actual sequence of point
blowdowns to its Hopf minimal surface.  A map from
\(\PP^1\) to a Hopf surface is constant: lift it to
\(\C^2\setminus\{0\}\) and use its two constant coordinate
functions.  Thus every rational curve on \(H\) is one
of the finitely many exceptional curves of \(\pi\).
Every compact curve on \(H_{\min}\) is a smooth elliptic
curve.  Indeed rational normalization is impossible,
and \(H^2(H_{\min},\Q)=0\) and adjunction give arithmetic
genus one; the normalization genus therefore equals
one and its delta invariant is zero.  Distinct such
curves are disjoint, since their intersection number
is zero and all local intersection multiplicities are
positive.

\noindent\textbf{Step 3.} Push forward \eqref{cdp:eq:C.HM.2} through the point blowdowns.  At
each reverse point blowup \(\sigma\) one has the exact
divisor formula
\[
             D_{\mathrm{new}}=\sigma^*D_{\mathrm{old}}-E.
\tag{C.HM.3}\label{cdp:eq:C.HM.3}
\]
Every intermediate divisor is effective: a negative
coefficient would persist on its strict transform in
the final effective divisor.  Starting from the
disjoint elliptic support on \(H_{\min}\), all these
blowups consequently have centres on the positive
support.  The positive support stays SNC.  At a smooth
point of a component of coefficient \(a\), the new
coefficient is \(a-1\); at a node of coefficients
\(a,b\), it is \(a+b-1\).  Coefficient-zero exceptional
curves are omitted from the positive support, but are
retained as curves of the surface.

\noindent\textbf{Step 4.} It follows by this actual point-map induction that
\[
 \begin{gathered}
 \text{two coefficient-one positive curves never meet;}\\
 \text{a rational coefficient-one curve is a leaf
 of the positive boundary,}\\
 \text{and an elliptic coefficient-one curve is
 disjoint from that boundary.}
 \end{gathered}
\tag{C.HM.4}\label{cdp:eq:C.HM.4}
\]
For the first assertion a smooth blowup creates
\((a,a-1)\), never \((1,1)\); a node blowup could
create a new \((1,1)\) only from an old \((1,1)\).
For the second assertion a new rational coefficient
one is first created at a smooth point of coefficient
two, and has valence one.  Subsequent node blowups
preserve that valence, and a smooth blowup at its
points has coefficient zero.  An original coefficient
one elliptic curve likewise gains no positive boundary
at a smooth blowup.

\noindent\textbf{Step 5.} These facts also rule out a rational self-double curve
on a maximal Hopf component at any NC stage in the
retained class.  Its two normalized preimages, as
specified by \cite[\S4.11.2(III)]{Kulikov1977}, have
coefficient one.  Each is rational and has exactly one
boundary intersection by \eqref{cdp:eq:C.HM.4}; hence the image has
one triple point.  Each rational curve has negative
self-intersection on the blown-up Hopf, whereas the
actual two-normal-bundle formula would give
\[
              C_1^2+C_2^2=-1,
\tag{C.HM.5}\label{cdp:eq:C.HM.5}
\]
an impossibility.  A triple point with three branches
of this same component would already give two
coefficient-one curves meeting, also impossible.
The modifications below can create new self-double
curves only from rational exceptional curves.
They therefore create no elliptic self-double curve
on a maximal Hopf starting from the given embedded
components.  At every subsequent smooth-total NC
stage these Hopf components are consequently embedded.

\par\medskip\noindent\textbf{The exact last-point alternatives.}
Suppose such an \(H\) is not minimal and choose a last
point blowup in the actual map \(\pi\).  Its exceptional
curve \(E\) is a smooth \((-1)\)-curve.  Formula
\eqref{cdp:eq:C.HM.3} leaves exactly three cases.

\noindent\textbf{Step 6.} First, \(E\) is not a boundary curve.  It was created
at a smooth point of a coefficient-one boundary \(C\),
and meets that curve once.  The ambient normal
sequence is
\[
 0\longrightarrow\OO_{\PP^1}(-1)
 \longrightarrow N_{E/X}
 \longrightarrow\OO_{\PP^1}(-1)\longrightarrow0 .
\tag{C.HM.6}\label{cdp:eq:C.HM.6}
\]
Its splitting and the ordinary two-ruling diagram of
\ref{cdp:st:C.2.UN.75} give the actual type-I move.  The neighbour across
\(C\) has order \(M\).  It is not Hopf-birational.
If \(C\) is elliptic, its square on \(H\) is at most
\(-1\) and it contains no triple point; its square on
the neighbour is therefore positive.  If \(C\) is
rational, it was already exceptional before this last
blowup, so its current square is at most \(-2\);
\eqref{cdp:eq:C.HM.4} gives one triple point, and its square on
the neighbour is again positive.  Such squares are
impossible for the corresponding elliptic or rational
curves on a blown-up Hopf.  The move removes one
point blowup from \(H\), adding it to a non-Hopf
component.

\noindent\textbf{Step 7.} Second, \(E\) was created at a node of two positive
boundary curves of coefficients \(a,b\).  It has
coefficient \(a+b-1\) and two triple points.  The
two-normal-bundle equation makes its square on the
other component also \(-1\), so the actual type-II
move of \ref{cdp:st:C.2.UN.75} applies.  It blows down \(E\) on both
sides and blows up the two opposite normalized
branches.  If an opposite component has order \(M\),
the associated curve \(C\) on \(H\) has coefficient
one.  It is a rational leaf by \eqref{cdp:eq:C.HM.4}, and its
square after the last node blowup is at most \(-2\).
Its square on that opposite component is
\[
                 -1-C_H^2\ge1 .
\tag{C.HM.7}\label{cdp:eq:C.HM.7}
\]
Hence that opposite component is not Hopf-birational.
This proves the same assertion when both opposite
branches belong to one global component: their actual
normalization branches are used in \eqref{cdp:eq:C.HM.7}.

\noindent\textbf{Step 8.} Let
\[
       h=\sum_{\substack{k_i=M\\
                     V_i\ \mathrm{Hopf\text{-}birational}}}
                              b_2(V_i).
\tag{C.HM.8}\label{cdp:eq:C.HM.8}
\]
At the smooth NC stages just described, this is the
total number of surface point blowups on these Hopfs.
Each of the two moves above strictly decreases \(h\).
No assertion that \(h\) decreases inside the longer
ruled-removal blocks below is needed.

\noindent\textbf{Step 9.} Third, \(E\) was created at a smooth point of a
positive curve of coefficient \(a\ge2\).  It has
coefficient \(a-1\) and one triple point.  Write \(Q\)
for its neighbour and \(P\) for the third component.
The exact normal degrees and canonical orders are
\[
 \begin{split}
 E_H^2&=-1,\qquad E_Q^2=0,\\
 q:=k_Q&=M-a+2,\qquad k_P=q-1,\\
                         &1\le q\le M .
 \end{split}
\tag{C.HM.9}\label{cdp:eq:C.HM.9}
\]
The complete rational curve of square zero gives
the actual ruling on the compact normalization of
\(Q\).  For clarity, its local complete deformations
are disjoint, since their intersection number is zero;
the resulting compact curve family gives the ruling
after the point blowdowns of its singular fibres.
Equivalently this is the ruled-surface comparison in
\cite[\S4.5]{Kulikov1977}.  Its one boundary intersection
is the section belonging to \(P\).  Any other
horizontal boundary would meet a general complete
fibre in this same local family, which it does not.
All other boundary curves on \(Q\) are vertical.
This is the precise ruled section-and-fibre
configuration, including any nonminimal fibres.

\par\medskip\noindent\textbf{Ruled removal and cleanup of its singular charts.}
We use the actual finite maps in
\cite[\S\S4.10.7--4.10.8,4.14]{Kulikov1977}, with
the missing-fibre correction in
\cite[\S8]{Kulikov1981}.  Their overlaps and proper
composite graphs are those verified in \ref{cdp:st:C.2.UN.75}.
They remove the ruled component in \eqref{cdp:eq:C.HM.9}.
If its boundary is disconnected, they first insert
the specified missing fibre components; these
exceptional ruled surfaces have order \(q+1\).
When \(q=M\), the canonical divisor on its
normalization has only negative boundary
coefficients, and \cite[Lemma2.19]{Kulikov1977}
makes that boundary connected.  No preliminary
addition is then made.  Thus all added orders are
at most the original fixed \(M\).

\noindent\textbf{Step 10.} The result of this removal need not immediately
have smooth total space.  We retain precisely the
finite sigma-atlas of
\cite[\S4.11.2]{Kulikov1977}: its singular sets are
contained in finitely many disjoint charts; in each
chart the recorded local exceptional branches
contract successively to a smooth NC chart.  We
now describe an actual cleanup, before applying
another last-point argument to a Hopf.

\noindent\textbf{Step 11.} If one of these lists is nonempty, choose a
component of largest canonical order among the
globally occurring last locally blown-up branches.
Use an auxiliary resolution \(a:\widehat X\to X'\)
as in \cite[\S4.13]{Kulikov1977} to read its
old strict component \(\widehat V\).  This
resolution is a diagnostic model for the
hypotheses of Lemma4.14; it is not inserted
as an additional degeneration stage in the
modification record.  In particular its new
exceptional divisors receive no card labels,
and no bound on their canonical orders is used.
Each branch has a complete local exceptional
\(\PP^1\)-fibre.  These complete fibres extend to
a ruling of the old smooth surface
\(\widehat V\).  In the
generic local exceptional chart there is exactly
one horizontal boundary section.  Completeness
of those fibres shows that there can be no
additional global horizontal boundary; all the
other boundary, including the diagnostic
exceptional curves over the finite exceptional
points, is vertical.  This supplies
conditions (a), (b) and (d) of
\cite[Lemma4.14]{Kulikov1977}, by their actual
normalization and section maps.

\noindent\textbf{Step 12.} If another curve on this chosen component had
been locally blown up, its exceptional branch
would have order one larger.  If that child
is not last, follow its finite local descendant
list.  The order increases at each step, so
a last descendant has strictly larger order
than the chosen component.  This contradicts
the choice of largest last-branch order.
Thus the choice gives condition
(e).  This is the exact comparison in
\cite[\S4.16.2]{Kulikov1977}.  The reordering in
its \S4.12 makes its specified branches last
without changing their restored overlaps.
The correction in \cite[\S8]{Kulikov1981}
inserts any missing vertical links in the
specified smooth open charts, leaving their
overlap fixed.  Lemma4.14 then removes this
global ruled component, and the result again
has the same kind of sigma-atlas.

\noindent\textbf{Step 13.} Every cleanup removal has positive order:
a locally exceptional component has the
actual order of its parent plus one.  It
adds only finitely many components at the
next order, and none at \(M\).  Here we
use precisely the maximal-old-order clause
of \cite[\S8]{Kulikov1981}, including its
last-locally-blown-up case: its boundary
is connected and no missing links are
inserted.  This does not assert that
every divisor on the auxiliary resolution
has order at most \(M\).  A Hopf
component cannot be one of these locally
exceptional global components: its complete
local rational fibres would give rational
curves through an open set of its compact
normalization, contrary to the finite
exceptional-curve description above.
The strict Hopf divisors survive by the
proper graph argument in \ref{cdp:st:C.2.UN.75}.

\noindent\textbf{Step 14.} This cleanup terminates without inserting an
ambient resolution into the record or estimating
orders of
arbitrary resolution divisors.  Label every
current global component by its order in
\(\{0,\ldots,M\}\).  A removal deletes a
positive label \(q\), adds finitely many
labels \(q+1\) when \(q<M\), and adds none
when \(q=M\).  There are finitely many
removals at level one, hence finitely many
objects ever created at level two, and
inductively at every level through \(M\).
This is the finite game of
\cite[\S9]{Kulikov1981}, now applied also
to the cleanup removals.

\noindent\textbf{Step 15.} When cleanup stops, no recorded locally
exceptional branch remains.  The defining
contraction lists of the sigma-atlas are
empty; their terminal smooth NC charts
are therefore the current charts.
The total space is smooth and its full
fibre is reduced NC.  Its normalization
components are smooth.  Other global
components may retain self-incidences.
At a maximal Hopf the intrinsic argument
\eqref{cdp:eq:C.HM.2}--\eqref{cdp:eq:C.HM.5} excludes rational
self-incidences.  No elliptic self-incidence
was introduced: type-I and type-II changes
use rational exceptional curves; in a
local ruled contraction the horizontal
parent has order one lower, so a
maximal-order Hopf can only meet that
contraction in vertical rational curves.
Thus the maximal Hopfs are again smooth
embedded components, and their ordinary
point maps to their Hopf minimal models
are available for the next step.

\par\medskip\noindent\textbf{Termination of the whole procedure.}
Perform a type-I or type-II step whenever
the chosen last point gives the first or
second case.  In the third case perform
the finite ruled removal and cleanup
just specified.  Across the entire
procedure there are only finitely many
divisorial removal blocks: the same
fixed-level game counts all their
positive component labels together.
Between these blocks every step strictly
decreases the finite integer \eqref{cdp:eq:C.HM.8}.
More explicitly, if \(c_j\) counts the
retained global components of order \(j\),
the block-level potential is
\[
                 (c_1,c_2,\ldots,c_M,h)
                 \quad\hbox{in lexicographic order.}
\tag{C.HM.10}\label{cdp:eq:C.HM.10}
\]
A whole-component removal decreases its
first changed \(c_j\), whatever finite
next-level additions or Hopf point
changes occur inside the block.
A type-I or type-II step chosen on a
last Hopf point leaves all the \(c_j\)
unchanged and decreases \(h\).
Lexicographic order on this fixed
finite product of nonnegative integers
has no infinite strictly descending
sequence, by the same successive-coordinate
induction.  This potential is attached
to the specified finite blocks, not to
every auxiliary arrow in their graphs.
An infinite procedure would therefore
have either infinitely many divisorial
removals, contradicting that game, or
an infinite tail of strict decreases of
\(h\).  Both are impossible.

\noindent\textbf{Step 16.} At the endpoint a nonminimal maximal
Hopf would still have a last point
blowup and one of the three alternatives,
so every such Hopf is minimal.  Its
curves are the disjoint smooth elliptic
curves proved above, and hence no
triple point meets it.  The original
Hopf primes survive with unchanged
orders by \ref{cdp:st:C.2.UN.75}, and each additional
prime used in the record is ruled.
The same actual proper graphs and
restored overlap isomorphisms give
the asserted reversible comparison.
Their annular gluing and normal-cone
Euler terms are exactly those of \ref{cdp:st:C.2.UN.75}.
This proves the statement in its
specified scope.
\end{proof}






\begin{proposition}[The \((-1,-2)\) curve with canonical orders \((0,0,1)\) occurs in a compact zero-Euler torus fibre]\label{cdp:st:C.2.UN.79}
\StatementRecord{cdp:st:C.2.UN.79}{proved}{Historical events: Claude Fable 5.2 (Anthropic), effort high, 16 November 2026: carried as background without fault found. GPT-6 Astra (OpenAI), ChatGPT Pro mode, effort Pro, S6\_09\_PAPERS, 17 November 2026: targeted proof check. Also checked in this writing. Pre-continuation writing events: GPT-6 Astra (OpenAI), ChatGPT Pro mode, effort Pro, S6\_10\_PAPERS, 18 November 2026: targeted proof check. Continuation metadata events: GPT-6 Astra (OpenAI), ChatGPT Pro mode, effort Pro, S6\_10\_PAPERS continuation, 19 November 2026: used and rechecked in this run; status-vocabulary reconciliation; metadata reconciliation only; no new mathematical check. \emph{Pre-continuation recorded proof scope.} The simple-discriminant Weierstrass model, cubic base change and three exact resolution charts give the entire O(-2) neighbourhood. The curve blowup has transition O(-1) plus O(-2), with the stated function and canonical orders, and its rational curve is horizontal for the specified elliptic map. The two further centres avoid that full germ. An independent exact computation verifies the four quadratic discriminants, all six coprimality tests and the simple Weierstrass discriminant; see analysis/UN79\_polynomial\_replay.json. Riemann-Hurwitz gives the connected genus-three centre, and the actual three curve blowups have total Euler increment zero}{the hypotheses, notation and definitions stated here; \cref{cdp:st:C.2.UN.70}; \cref{cdp:st:C.2.UN.41}; \cref{cdp:eq:C.BC.4}; \cref{cdp:st:C.2.UN.74}; \cref{cdp:eq:C.BC.5}; \cref{cdp:eq:C.BC.8}; \cref{cdp:eq:C.BC.9}; \cref{cdp:eq:C.BC.13}}

There is a smooth compact projective threefold \(X\), a proper map
\(f:X\to\PP^1\) with connected general fibre a complex two-torus,
and a full reduced SNC fibre \(S=f^{-1}(0)\) with \(e(S)=0\),
containing a smooth rational curve \(F\) with the following exact
properties. The curve belongs to just one component \(P\) generically,
is not a double curve, has \(F_P^2=-1\), and meets two other
components transversely, once each, at distinct points. Its entire
analytic neighbourhood, including \(f\), is the germ along the zero
section of
\[
 \operatorname{Tot}\bigl(\OO_{\PP^1}(-1)_a
                 \oplus\OO_{\PP^1}(-2)_b\bigr),\qquad
 f=bUV,\quad [U:V]\in\PP^1.                         \tag{C.BC.1}\label{cdp:eq:C.BC.1}
\]
Here \(P=(b=0)\), and the other two components are the fibres
\(U=0\) and \(V=0\) of the displayed vector bundle.
The actual generator of \(f_*K_X\) has orders \(0,1,0\)
on these three components in that order. Thus, after ordering the
two neighbours by canonical height, the tuple is \((0,0,1)\).
In particular
\[
                    N_{F/X}=\OO(-1)\oplus\OO(-2).            \tag{C.BC.2}\label{cdp:eq:C.BC.2}
\]

The construction retains an actual point-blowdown map from \(P\)
to \(\PP^1\times E\), where \(E\) is an elliptic curve.
It maps \(F\) isomorphically onto \(\PP^1\times\{e_0\}\).
For the elliptic map obtained by composing this blowdown with
\(\PP^1\times E\to\PP^1\), the curve \(F\) is therefore
horizontal. The specified exceptional locus of this surface map
does not contain \(F\). This records the scope relevant to \ref{cdp:st:C.2.UN.70}:
the example concerns an arbitrary nonboundary \((-1)\)-curve,
with its exact comparison to the chosen elliptic minimalization.
\end{proposition}



\begin{proof}
\noindent\textbf{Step 1.} \par\medskip\noindent\textbf{An actual compact elliptic surface with a toric \(I_3\) chart.}
Take a compact rational elliptic surface with a simple nodal fibre.
Such a surface, with an explicit cubic-pencil construction, was
used in \ref{cdp:st:C.2.UN.41}. Alternatively the Weierstrass equation
\[
       y^2=x^3+(t^4+1)x+t^6                         \tag{C.BC.3}\label{cdp:eq:C.BC.3}
\]
gives one: its discriminant, up to a nonzero constant, is
\(31t^{12}+12t^8+12t^4+4\), which is square-free and has no
zero at infinity. Choose one nodal value, put it at \(t=0\)
by a base coordinate, and take the global cubic base change
\(t=s^3\). Normalize and minimally resolve. The resulting
compact smooth projective elliptic surface \(J\to\PP^1_s\)
has an \(I_3\) fibre
\[
                         C_0+C_1+C_2 .             \tag{C.BC.4}\label{cdp:eq:C.BC.4}
\]
The components are smooth rational curves, meeting in a triangle
at three distinct nodes, with no triple point. The base change
is unramified at the other singular values; at a smooth value
it is the ordinary pullback of a smooth family. Thus these
operations give a smooth total surface after resolving the
single local \(A_2\) equation at the chosen nodal value.

\noindent\textbf{Step 2.} Here is the required local comparison, with its entire two-chart
neighbourhood retained. At the original nodal point the map is
\(xy=t\), by the holomorphic Morse coordinates. The root is
\(xy=s^3\). Two consecutive charts of its actual minimal
resolution are
\[
 \begin{aligned}
 (x,y,s)&=(b,b^2z^3,bz),\qquad (b,z)\in\C^2,\\
 (x,y,s)&=(u^2c,uc^2,uc),\qquad (u,c)\in\C^2,\\
                 u&=z^{-1},\qquad c=z^2b .
 \end{aligned}                                             \tag{C.BC.5}\label{cdp:eq:C.BC.5}
\]
The curve \(C_1\) is \(b=0\) in the first chart and
\(c=0\) in the second. Its adjacent curves are \(z=0\)
and \(u=0\). The two-chart union is the actual total space
of \(\OO(-2)\), restricted to a neighbourhood if necessary;
its map is exactly \(s=bz=cu\).
The displayed monomials satisfy \(xy=s^3\), glue on their
overlap, and are the successive point-resolution charts,
so this comparison does not identify a neighbourhood merely
from its normal degree.

\noindent\textbf{Step 3.} The two-form \(dx\wedge ds/x\) is \(db\wedge dz\)
on the first chart and \(du\wedge dc\) on the second.
An original canonical generator is a nonvanishing
holomorphic multiple of \(dx\wedge dy\) at the node.
Pull it back and divide by \(s^2\); the resulting
form is a nonvanishing holomorphic multiple of the
displayed two-form and extends through the unchanged
smooth charts of the original fibre. The three
components have the same original root order two.
This gives a nowhere-zero canonical generator
near \eqref{cdp:eq:C.BC.4}. Hence the actual
canonical evaluation orders on all three components of
\(J\) near this fibre are zero.

\par\medskip\noindent\textbf{One global curve blowup produces the required germ and orders.}
Fix an elliptic curve \(E\) and a point \(e_0\in E\).
Set \(X_0=J\times E\), with map \(f_0\) induced by \(s\).
Its general fibre is the product of two elliptic curves.
The central components are \(P_i=C_i\times E\), all
of canonical order zero. Blow up the complete smooth
rational curve
\[
                      B_0=C_0\times\{e_0\}\subset P_0 .
                                                               \tag{C.BC.6}\label{cdp:eq:C.BC.6}
\]
The centre is contained generically in just \(P_0\) and
meets each of its two double curves transversely at one
point. The inverse chart of \ref{cdp:st:C.2.UN.74}, or the following direct
calculation, shows that the new fibre is full reduced SNC.
Near any such point write the original map as \(xy\),
the centre as \((x,a)\), and the blowup charts as
\(a=xA\) and \(x=aX\). Their maps are respectively
\(xy\) and \(aXy\). All displayed factors have
multiplicity one. The new exceptional component \(R_0\)
has canonical order one, because the Jacobian of this
smooth curve blowup has a simple exceptional zero.

\noindent\textbf{Step 4.} Choose the labelling in \eqref{cdp:eq:C.BC.5} so that \(z=0\) is
the branch of \(C_0\). Let \(a\) be a local coordinate
of \(E\) at \(e_0\). Near the whole curve
\(C_1\times\{e_0\}\), the centre is \((z,a)\)
near its first endpoint and is absent elsewhere.
In the \(z\)-chart of its blowup put
\[
               a=zA,\qquad s=zb .                          \tag{C.BC.7}\label{cdp:eq:C.BC.7}
\]
The strict transform \(F\) is \(A=b=0\). In the other
chart of \eqref{cdp:eq:C.BC.5} no modification occurs near this curve;
write its elliptic coordinate as \(a_\infty\). The actual
transition is
\[
       u=z^{-1},\qquad A=u a_\infty,\qquad b=u^2c,
             \qquad s=zb=uc .                              \tag{C.BC.8}\label{cdp:eq:C.BC.8}
\]
These are exactly the two vector-bundle charts in \eqref{cdp:eq:C.BC.1}.
They prove \eqref{cdp:eq:C.BC.2} and \(F_P^2=-1\), with the entire
map retained. The two boundary intersections occur at
\(z=0\) and \(u=0\), whereas \(F\) is contained
generically only in \(b=0\).

\noindent\textbf{Step 5.} Up to its retained nonvanishing holomorphic factor, the actual
product canonical form pulls back as
\[
 db\wedge dz\wedge da
       =z\,db\wedge dz\wedge dA .                          \tag{C.BC.9}\label{cdp:eq:C.BC.9}
\]
At the other endpoint this is a nonzero unit times
\(du\wedge dc\wedge da_\infty\). Thus it has a simple
zero exactly on \(R_0=(z=0)\), with zero orders on
\(P\) and its other old neighbour. This is the requested
canonical tuple for the actual direct-image generator.
It remains a generator: the curve blowup has the actual
canonical pushforward \(\pi_*K=K\) through its Jacobian,
or directly through Hartogs across its smooth codimension-two
centre.

\noindent\textbf{Step 6.} The surface \(P\) is \(\operatorname{Bl}_p(\PP^1\times E)\),
where \(p\) is the endpoint of
\(\PP^1\times\{e_0\}\). The strict transform of that
rational curve is \(F\). The old double curve through
\(p\) has the other tangent direction in \(P_i\), so
its strict transform meets the exceptional line at a
different point from \(F\). In particular the full
strict transform of \(P_0\) is disjoint from \(F\).
The actual surface blowdown maps \(F\) isomorphically
to \(\PP^1\times\{e_0\}\); its composition with
the first projection has degree one on \(F\).

\par\medskip\noindent\textbf{Exact zero-Euler balancing away from the entire curve.}
We make the remaining centres explicit. Use
\[
 E:\ y^2=x(x-1)(x-2),\qquad e_0=(0,0),\quad e_1=(1,0).
                                                               \tag{C.BC.10}\label{cdp:eq:C.BC.10}
\]
The strict \(P_0\) is still the identical surface
\(\PP^1\times E\): restriction of the first blowup
to this component is the blowup of its Cartier divisor
\(B_0\), hence an isomorphism. Blow up the disjoint
curve \(B_1=C_0\times\{e_1\}\). Again \(P_0\)
is unchanged and all full multiplicities are one.
The second centre is disjoint from \(F\), so this
operation preserves a neighbourhood of the whole
compact curve.

\noindent\textbf{Step 7.} Choose the coordinate \(z\) on \(C_0\simeq\PP^1\)
so its two old boundary points are \(0,\infty\). Put
\[
 P(z)=z^2+1,\qquad Q(z)=3z^2+z+5,\qquad
 D=\{P(z)x=Q(z)\}\subset\PP^1\times E .
                                                               \tag{C.BC.11}\label{cdp:eq:C.BC.11}
\]
This is the divisor of a section of
\(\OO_{\PP^1}(2)\boxtimes\OO_E(2\cdot\infty_E)\).
Its map to \(\PP^1_z\) is the exact smooth double
cover
\[
          Y^2=P Q(Q-P)(Q-2P),\qquad Y=yP^2 .                  \tag{C.BC.12}\label{cdp:eq:C.BC.12}
\]
The four quadratic factors have respectively discriminants
\(-4,-59,-31,-11\), are pairwise coprime, and have
nonzero leading coefficients. Their eight simple roots
give a smooth connected genus-three curve. The chart
at \(P=0\) is the ordinary elliptic infinity chart;
since these zeros are simple it gives the same smooth
double cover, not additional components. At \(z=0\)
and \(\infty\) the values \(x=5,3\) are outside
\(\{0,1,2,\infty\}\). Thus \(D\) meets each old
elliptic boundary transversely in two points and avoids
all four existing triple points, which have elliptic
coordinate \(e_0\) or \(e_1\). The simple zeros of
\(Q\) and \(Q-P\) prove transverse intersection
with \(B_0\) and \(B_1\), respectively. These facts
also verify all its boundary intersections directly.

\noindent\textbf{Step 8.} Blow up this \(D\) in the threefold. The same two
charts \(xy\), \(aXy\) prove reduced SNC at its
boundary crossings; elsewhere it is a smooth-centre
blowup in one reduced component. It has canonical
order one generically, since its ambient centre lies
in the zero-order component \(P_0\). The centre is
disjoint from \(F\), so \eqref{cdp:eq:C.BC.8}--\eqref{cdp:eq:C.BC.9} remain
unchanged on an open neighbourhood of the entire
curve, including every analytic transverse jet.

\noindent\textbf{Step 9.} Initially \(e((C_0+C_1+C_2)\times E)=0\).
For a smooth curve centre contained in the fibre,
the actual blowup replaces that curve by its
\(\PP^1\)-bundle exceptional divisor and is an
isomorphism on the complement. Its reduced-fibre
Euler defect is therefore exactly \(e\) of the
centre. Consequently our three operations give
\[
             e(S)=0+e(B_0)+e(B_1)+e(D)
                         =0+2+2-4=0 .                       \tag{C.BC.13}\label{cdp:eq:C.BC.13}
\]
For an independent stratum check, the six final components
have Euler numbers \(0,4,4,6,6,-8\). The double curves
have total Euler characteristic \(24\), and there are
\(12\) triple points. Thus their actual SNC
inclusion--exclusion gives \(12-24+12=0\), in
agreement with \eqref{cdp:eq:C.BC.13}.

\noindent\textbf{Step 10.} All blowups are along closed smooth curves in a
compact projective threefold. Their total spaces
are smooth, compact and projective. They preserve
the map outside the single fibre; properness and
connectedness follow from their connected blowup
fibres and the original connected elliptic product.
The other fibres remain unchanged products with
\(E\). Finally, every additional point blowup
induced on \(P\) is away from \(F\), so the
same composite surface blowdown and its degree-one
map on \(F\) prove the stated horizontal
qualification on the final zero-Euler model.
\end{proof}






\begin{proposition}[Every point-exceptional prime of an aligned parent is an actual boundary]\label{cdp:st:C.2.UN.80}
\StatementRecord{cdp:st:C.2.UN.80}{proved}{Historical events: Claude Fable 5.2 (Anthropic), effort high, 16 November 2026: carried as background without fault found. GPT-6 Astra (OpenAI), ChatGPT Pro mode, effort Pro, S6\_09\_PAPERS, 17 November 2026: targeted proof check. Also checked in this writing. Pre-continuation writing events: GPT-6 Astra (OpenAI), ChatGPT Pro mode, effort Pro, S6\_10\_PAPERS, 18 November 2026: targeted proof check. Continuation metadata events: GPT-6 Astra (OpenAI), ChatGPT Pro mode, effort Pro, S6\_10\_PAPERS continuation, 19 November 2026: used and rechecked in this run; status-vocabulary reconciliation; metadata reconciliation only; no new mathematical check. \emph{Pre-continuation recorded proof scope.} The rank-one return equations and the rank-two commuting transverse flow give a fixed nonzero vertical translation field on the actual minimal elliptic model. Its constant period normalization extends across every multiple-fibre root chart. For any exceptional prime, the pulled flow-box functions are s-to-the-a times a unit and s-to-the-b times a unit. A holomorphic lift would preserve the prime ideal from the zero derivative of the second function, contradicting derivative one of the first. This proves the polar assertion for every infinitely-near exceptional valuation of the specified point map. Historical program executions are not counted as current executions}{the hypotheses, notation and definitions stated here; \cref{cdp:st:C.2.UN.70}; \cref{cdp:st:C.2.UN.16}; \cref{cdp:eq:C.EP.6}; \cref{cdp:eq:C.EP.7}; \cref{cdp:eq:C.EP.3}; \cref{cdp:eq:C.EP.10}; \cref{cdp:eq:C.EP.9}; \cref{cdp:eq:C.EP.11}}

Let \(P\), \(\pi:P\to Z\), and \(q_Z:Z\to B\) be the
actual component and maps of \ref{cdp:st:C.2.UN.70}. Write \(D_P\)
for the reduced union of its intersections with the
other central components. There is a nonzero constant
linear combination of the two restricted saturated
ambient translation fields which is the meromorphic
lift of a nowhere-zero holomorphic vertical field
\[
 W_Z\in H^0(Z,T_Z),\qquad dq_Z(W_Z)=0.
 \tag{C.EP.1}\label{cdp:eq:C.EP.1}
\]
On every smooth elliptic fibre its complex-time period
lattice is the same fixed lattice. The field extends
without a zero or a pole across every smooth multiple
fibre of \(q_Z\).

Every prime exceptional for the specified point map
\(\pi\) is a polar divisor of its meromorphic lift
\(W_P\). Consequently
\[
                 \operatorname{Exc}(\pi)\subset D_P.
 \tag{C.EP.2}\label{cdp:eq:C.EP.2}
\]
In particular, a last point-exceptional curve for
\(\pi\) is an actual boundary curve. This assertion
concerns the exceptional primes of this actual map:
it makes no assertion that every rational
\((-1)\)-curve on \(P\) is exceptional for \(\pi\).
The component and its ambient neighbourhood are not
changed.
\end{proposition}



\begin{proof}
\noindent\textbf{Step 1.} \par\medskip\noindent\textbf{The actual field space and its polar locus.}
Use the two-dimensional saturated field space of
\ref{cdp:st:C.2.UN.16} on the original multiplicity-one component
\(P\). In this full reduced situation no root cover
of \(P\) is needed. Denote a basis of its restrictions
by \(X,Y\). They are independent over \(\C\), commute,
and are meromorphic fields on \(P\).

\noindent\textbf{Step 2.} Their actual ambient representatives are linear
combinations of holomorphic relative fields with
coefficients meromorphic in the disc coordinate.
They therefore have polar divisors only in the
central fibre. Saturation removes their pole along
the generic point of \(P\). On \(P\setminus D_P\)
the central fibre has just this one smooth branch,
so restriction gives
\[
       X,Y\ \hbox{holomorphic on }P\setminus D_P .
 \tag{C.EP.3}\label{cdp:eq:C.EP.3}
\]
This is the exact restriction comparison used in
(C.MR.8) and \ref{cdp:st:C.2.UN.70}. It applies to every constant
linear combination. All boundary and exceptional
curves of the specified map \(\pi\) are vertical
for \(q=q_Z\pi\).

\noindent\textbf{Step 3.} Remove finitely many values from \(B\) so that
\(\pi\) is an isomorphism over their complement
\(B^\circ\), all fibres are smooth, and the fields
and their projected coefficients are regular.
On a compact elliptic fibre the latter coefficients
are holomorphic functions and hence constants.
Thus \(X,Y\) project to meromorphic vector fields
on \(B\), exactly as in the alignment proof.

\par\medskip\noindent\textbf{Rank one: the period identity gives a constant combination.}
Suppose the meromorphic rank of \(\langle X,Y\rangle\)
is one. Alignment identifies its connected elliptic
first integral with the actual map \(q\). Thus
\[
                       Y=(q^*H)X
 \tag{C.EP.4}\label{cdp:eq:C.EP.4}
\]
for a nonconstant meromorphic function \(H\) on
\(B\). The actual return-germ calculation
(C.RG.4)--(C.RG.12), performed on the original
ambient representatives, gives time periods
\(\lambda_1,\lambda_2\) of \(X\) and constants
\(d_i,k_i\in\C\) such that
\[
            \lambda_i=d_i-k_iH,\qquad
            \lambda_2/\lambda_1=\tau\in\mathbb H .
 \tag{C.EP.5}\label{cdp:eq:C.EP.5}
\]
The formula continues on \(B^\circ\) with these
same constants. This uses the actual global ratio
\(H\), the continued integral cycles, and the
identity theorem; it does not extend \(q\) to a
threefold neighbourhood.

\noindent\textbf{Step 4.} Take
\[
                       W=d_1X-k_1Y=\lambda_1X.
 \tag{C.EP.6}\label{cdp:eq:C.EP.6}
\]
A nonzero integral cycle on an elliptic curve has
nonzero time period, so \(W\) is nonzero on every
fibre under consideration. Multiplying a field by
\(\lambda_1\) divides its time periods by
\(\lambda_1\); the two periods of \(W\) are exactly
\(1,\tau\). Thus \eqref{cdp:eq:C.EP.6} is a constant linear
combination with a fixed nonzero period lattice.
In an integration coordinate on each fibre it
is the fixed invariant translation field.

\par\medskip\noindent\textbf{Rank two: commuting transverse transport preserves the periods.}
Suppose instead that the meromorphic rank is two.
The two projected fields on the curve \(B\)
commute. Two nonzero commuting meromorphic fields
on a connected curve have constant ratio:
writing them locally as \(a(z)\partial_z\) and
\(b(z)\partial_z\), the bracket equation is
\(ab'-ba'=0\). The ratio therefore continues as
a constant. They cannot both project to zero,
since vertical fields on an elliptic fibre have
meromorphic rank at most one. After interchanging
the basis if necessary there is consequently a
constant \(c\) for which
\[
             W=Y-cX\ne0,\qquad q_*W=0,\qquad
                         q_*X\ne0 .
 \tag{C.EP.7}\label{cdp:eq:C.EP.7}
\]
On the compact fibres over \(B^\circ\), \(W\)
is a nonzero translation field. The local flow
of the transverse field \(X\), constructed on
an entire compact fibre for sufficiently small
time, identifies neighbouring elliptic fibres.
Since \([X,W]=0\), this identification preserves
the time form and its two integral periods.
Thus their lattice is locally constant and
hence constant on connected \(B^\circ\).
This is the actual compact-fibre flow comparison
of (C.RF.1)--(C.RF.2). A constant rescaling of
\(W\) makes one nonzero period equal to \(1\).
Once again \(W\) is a constant combination of
the original two restricted fields.

\par\medskip\noindent\textbf{Extension on the specified minimal elliptic model.}
On \(q_Z^{-1}(B^\circ)\), push \(W\) through
the open isomorphism of \(\pi\). Integration of
its fixed lattice identifies each small smooth
fibre tube with a torsor under one fixed elliptic
curve \(F\). Every overlap has translation
linear part, so its derivative preserves the
same invariant field of \(F\).

\noindent\textbf{Step 5.} At a multiple value, use the actual normalized
root chart of \ref{cdp:st:C.2.UN.70} and (C.RG.13):
\[
 (F\times\Delta_z)/\langle
       (u,z)\longmapsto(u+\eta,\zeta_mz)\rangle,
       \qquad t=z^m .
 \tag{C.EP.8}\label{cdp:eq:C.EP.8}
\]
The translation \(\eta\) gives the free cyclic
action on the total space. The invariant
vertical field on \(F\times\Delta_z\) is fixed
by this action and is nowhere zero. On the
punctured root disc our field agrees with it:
its coefficient relative to the invariant
frame has the fixed nonzero time periods and
therefore is the same constant. It follows
that the agreement extends across \(z=0\)
and descends through the actual etale quotient.
These extensions agree on overlaps with the
previously fixed field. This proves \eqref{cdp:eq:C.EP.1},
including every zero or pole that might have
been hidden at an omitted base value.

\noindent\textbf{Step 6.} The meromorphic lift of \(W_Z\) through \(\pi\)
agrees with the constant combination \eqref{cdp:eq:C.EP.6}
or \eqref{cdp:eq:C.EP.7} on their common open isomorphism,
and therefore agrees everywhere as a
meromorphic vector field. In particular its
polar divisors are contained in \(D_P\), by
\eqref{cdp:eq:C.EP.3}.

\par\medskip\noindent\textbf{Every exceptional valuation is polar.}
Let \(E\) be any prime exceptional for \(\pi\),
with image \(p\in Z\). The nowhere-zero field
\(W_Z\) has a holomorphic flow box \((z,w)\)
centred at \(p\) in which
\[
                   W_Z(z)=1,\qquad W_Z(w)=0.
 \tag{C.EP.9}\label{cdp:eq:C.EP.9}
\]
Both pullbacks vanish to positive order along
\(E\). Work at a generic smooth point of \(E\)
and let \(s\) be its local equation. Write
\[
            \pi^*w=s^b u,\qquad
            \pi^*z=s^a v,\qquad a,b\ge1,
 \tag{C.EP.10}\label{cdp:eq:C.EP.10}
\]
where \(u,v\) are holomorphic units after
shrinking this generic neighbourhood.

\noindent\textbf{Step 7.} Suppose the lifted derivation \(D=W_P\) were
holomorphic there. The identity
\(D(\pi^*w)=0\), divided by \(s^{b-1}\) and
reduced modulo \(s\), gives
\[
                    b\,u\,D(s)=0\pmod{s}.
 \tag{C.EP.11}\label{cdp:eq:C.EP.11}
\]
Characteristic zero and the fact that \(u\)
is a unit imply \(D(s)\in(s)\). Thus \(D\)
preserves the ideal of \(E\), and applying it
to the second expression in \eqref{cdp:eq:C.EP.10} gives
\(D(\pi^*z)\in(s)\). This contradicts the
exact pulled-back identity \(D(\pi^*z)=1\).
Hence \(E\) is a polar divisor. The argument
applies directly to every exceptional prime,
including all infinitely-near point divisors;
it does not assume that a first pole survives
an unexamined modification.

\noindent\textbf{Step 8.} Combining this with \eqref{cdp:eq:C.EP.3} proves \eqref{cdp:eq:C.EP.2}.
The exact period rescalings, the root-chart
invariance, and finite toric valuation
certificates for \eqref{cdp:eq:C.EP.9}--\eqref{cdp:eq:C.EP.11} are
executed in \path{code/A_translation_exceptions.py}.
The general exceptional-prime assertion is
the local derivation proof above, without a
toric hypothesis on the actual point map.
\end{proof}






\begin{proposition}[Finite removal of all vertical point defects at the marked Hopf interfaces]\label{cdp:st:C.2.UN.81}
\StatementRecord{cdp:st:C.2.UN.81}{proved}{Historical events: Claude Fable 5.2 (Anthropic), effort high, 16 November 2026: carried as background without fault found. GPT-6 Astra (OpenAI), ChatGPT Pro mode, effort Pro, S6\_09\_PAPERS, 17 November 2026: targeted proof check. Also checked in this writing. Pre-continuation writing events: GPT-6 Astra (OpenAI), ChatGPT Pro mode, effort Pro, S6\_10\_PAPERS, 18 November 2026: targeted proof check; used and rechecked in this run. Continuation metadata events: GPT-6 Astra (OpenAI), ChatGPT Pro mode, effort Pro, S6\_10\_PAPERS continuation, 19 November 2026: used and rechecked in this run; status-vocabulary reconciliation; metadata reconciliation only; no new mathematical check. \emph{Pre-continuation recorded proof scope.} The complete point-map fibre graph gives the receiving-square bounds, so a direct two-ruling step decreases the designated parent defects. The zero-triple alternative yields the basepoint-free actual plane map from the plus-one curve; its canonical discrepancy divisor accounts for every other boundary. The three last-point alternatives either raise the marked plane order, lower the fixed Betti vector or remove a positive-order component. The finite corrected sigma cleanup preserves the whole maximal Hopf neighbourhood. The fixed component-card and parent-defect vector proves termination, while the retained overlap maps recover the original family}{the hypotheses, notation and definitions stated here; \cref{cdp:st:C.2.UN.77}; \cref{cdp:st:C.2.UN.75}; \cref{cdp:st:C.2.UN.70}; \cref{cdp:st:C.2.UN.80}; \cref{cdp:eq:C.AP.5}; \cref{cdp:eq:C.AP.3}; \cref{cdp:eq:C.AP.6}; \cref{cdp:st:C.2.UN.74}; \cite[\S\S4.10.7--4.10.8,4.14]{Kulikov1977}; \cite[\S8]{Kulikov1981}; \cref{cdp:st:C.2.UN.31}; \cref{cdp:eq:C.AP.9}; \cref{cdp:eq:C.AP.10}; \cref{cdp:eq:C.AP.11}; \cref{cdp:st:C.2.UN.64}; \cref{cdp:eq:C.AP.2}; \cref{cdp:eq:C.AP.4}; \cref{cdp:eq:C.AP.14}}

Let \(f:X\to\Delta\) be a proper degeneration of connected
complex two-tori with smooth total space and full reduced
SNC central fibre. We also allow as the starting recorded model
the smooth-total NC endpoint obtained from such a fibre by
\ref{cdp:st:C.2.UN.77}, retaining its two labelled normalization branches at
each self-double curve. Normalize the actual canonical generator by
\[
 \operatorname{div}(\Omega)=\sum k_iV_i,\qquad
 \min k_i=0,\qquad M=\max k_i.
 \tag{C.AP.1}\label{cdp:eq:C.AP.1}
\]
Suppose that the components of order \(M\) with Hopf minimal
model have already been made minimal by the actual record of
\ref{cdp:st:C.2.UN.77}. Mark any finite collection of their interfaces
\(J_\alpha=H_\alpha\cap P_\alpha\) for which the normal line of
\(J_\alpha\) is torsion and \(k_{P_\alpha}<M\).
Retain repeated markings on the same parent.

There is a finite overlap-compatible modification record,
unchanged over \(\Delta^*\), whose endpoint has smooth total
space and full reduced NC central fibre, such that every
designated parent is an embedded minimal fixed-Jacobian
elliptic quasi-bundle. Its full elliptic fibres are multiples
of smooth elliptic curves, its integral linear monodromy is
the identity, and every one of its boundary curves is a
smooth elliptic fibre support. In particular no triple point
meets a designated parent at the endpoint.

All maximal-order minimal Hopf components and an open
neighbourhood of their union are unchanged by the record.
In particular every marked \(J_\alpha\), its normal map and
its gluing to \(H_\alpha\) are retained. The designated parent
prime divisors survive. Other endpoint components can retain
NC self-incidences; their two normalization branches are
part of the record.

Every retained global component has canonical order between
\(0\) and the original \(M\). Whole-component removals delete
a component of positive order \(q\); their preliminary
additions, when needed, have order \(q+1\le M\), and there
are no additions at \(q=M\). The record retains its actual
centres, local maps, normalization branches and restored
overlaps, so reversing it recovers the original family.
For a given compact family the same modifications are
performed on that compact base. Their Euler changes are
the exact normal-cone terms of \ref{cdp:st:C.2.UN.75}.

No assertion is made here that every maximal Hopf component
has one of the marked interfaces. In particular an
equal-order interface with a nontorsion normal line is not
among these markings.
\end{proposition}



\begin{proof}
\noindent\textbf{Step 1.} \par\medskip\noindent\textbf{The exact marked parent maps.}
All curves on a minimal Hopf surface are disjoint smooth
elliptic curves. Thus each marked curve is isolated from
the other boundaries on both sides. Its actual normal-product
isomorphism makes the normal on the parent torsion as well.
On the normalization of a designated parent \(P\), of order
\(p\), the marked curve has coefficient
\(M-p-1\ge0\) in the actual canonical divisor. \ref{cdp:st:C.2.UN.70} therefore
gives a specified point map and elliptic map
\[
 P\xrightarrow{\pi_P} Z_P\xrightarrow{q_P}B_P,
 \qquad R^1(q_P)_*\OO_{Z_P}=\OO_{B_P},
 \qquad g(B_P)\le1.
 \tag{C.AP.2}\label{cdp:eq:C.AP.2}
\]
The point map is an isomorphism on a neighbourhood of every
marked curve. Every boundary is vertical for \(q_P\pi_P\),
and every rational boundary is exceptional for \(\pi_P\).
The exact constant-period field of \ref{cdp:st:C.2.UN.80}, equations
(C.EP.1)--(C.EP.11), proves the additional inclusion
\[
                 \operatorname{Exc}(\pi_P)\subset D_P.
 \tag{C.AP.3}\label{cdp:eq:C.AP.3}
\]
Here \(D_P\) is the full actual reduced boundary. It includes
both preimages of a self-double curve at an NC stage.

\noindent\textbf{Step 2.} These assertions also apply at the intermediate smooth-total
NC stages used below. Indeed the actual adjunction divisor
on the normalization counts those two preimages; the marked
elliptic neighbourhood stays embedded and unchanged. The
saturated ambient fields restrict as meromorphic fields to
that normalization and are holomorphic off its full boundary.
The surface maps in \ref{cdp:st:C.2.UN.70} and the direct exceptional-valuation
argument in \ref{cdp:st:C.2.UN.80} consequently apply to this actual smooth
normalization, without identifying its two boundary branches.

\noindent\textbf{Step 3.} The target \(Z_P\) is a minimal surface in its birational
class. Thus
\[
 h_P=b_2(P)-b_2(Z_P)\ge0
 \tag{C.AP.4}\label{cdp:eq:C.AP.4}
\]
is exactly the number of point blowups in \(\pi_P\). The
minimal second Betti number is unchanged under the subsequent
birational maps: in the ruled case it is two, and in the
other cases the minimal model is unique. We count each
designated parent once in \(h=\sum_P h_P\).

\par\medskip\noindent\textbf{The full fibre graph and the receiving degrees.}
The point map \(\pi_P\) is a sequence of point blowups of
smooth elliptic fibre supports. For one such fibre, retain
its entire reduced inverse image, whether or not its strict
elliptic component is an actual boundary. Its graph is an
elliptic root with rational trees. Write \(C^2=-b_C\) and let
\(n_C\) be the number of its intersections with the other
components of this entire reduced fibre. The point maps give
\[
 n_C\le b_C+1\quad\hbox{for a rational component},\qquad
 n_C\le b_C\quad\hbox{for the elliptic component}.
 \tag{C.AP.5}\label{cdp:eq:C.AP.5}
\]
The elliptic root starts with \((b,n)=(0,0)\). A new rational
curve has \((b,n)=(1,1)\) at a smooth point and \((1,2)\) at
a node. Blowing up a smooth point of an old component raises
both \(b\) and \(n\) by one; a node blowup raises \(b\) by one
and leaves its valence unchanged. These explicit updates
prove \eqref{cdp:eq:C.AP.5} at every infinitely-near point.

\noindent\textbf{Step 4.} Choose a last point blowup of \(\pi_P\), with exceptional
curve \(E\). By \eqref{cdp:eq:C.AP.3}, \(E\) is an actual boundary curve.
It meets zero, one or two other actual boundary curves.
These are exactly its triple points in the threefold.
The complete fibre graph has at most two neighbours at this
last point, so there is no additional case.

\noindent\textbf{Step 5.} Suppose first that there are two. On the other component
\(Q\), the normal-product formula gives \(E_Q^2=-1\), because
\(E_P^2=-1\). Hence the actual type-II two-ruling diagram of
\ref{cdp:st:C.2.UN.75} applies. It blows down \(E\) on both adjacent surfaces
and blows up the two opposite normalized branches at their
specified points.

\noindent\textbf{Step 6.} Neither receiving branch belongs to a designated parent.
To prove this, let \(C\) be the opposite boundary on \(P\),
and write \(-b_0\) and \(n_0\) for its square and entire-fibre
valence just before the last node blowup. Afterwards its
square is \(-b_0-1\), and its number \(t_C\) of actual triple
points is at most \(n_0\). On the opposite receiving surface
\(R\), the exact normal map therefore gives
\[
 C_R^2=b_0+1-t_C\ge b_0+1-n_0\ge0,
 \quad\hbox{and } C_R^2\ge1\hbox{ if }C\hbox{ is elliptic}.
 \tag{C.AP.6}\label{cdp:eq:C.AP.6}
\]
On a designated parent every rational boundary is exceptional
and has negative square; every elliptic boundary is the
strict transform of a fibre support and has nonpositive
square. Both alternatives contradict \eqref{cdp:eq:C.AP.6}. This argument
uses the individual normalized branch when the two receivers
are branches of one global component. Thus a type-II move
leaves the component counts unchanged and strictly decreases
\(h\); it does not create a self-double curve on a designated
parent.

\par\medskip\noindent\textbf{One triple point: an actual ruled-removal block.}
Suppose that \(E\) has one triple point. Then \(E_Q^2=0\).
Its complete rational deformations give the actual ruling
on the compact normalization of \(Q\), after its vertical
point contractions. The boundary passing through this one
triple point is its unique section; every other boundary is
vertical. This is the same complete-curve and actual ruling
comparison used in (C.HM.9), not a ruling inferred from a
numerical class alone. If \(q=k_Q\) and \(R\) is the section
component, restriction of the actual canonical divisor to
a complete ruling fibre gives
\[
                         k_R=q-1,
                         1\le q\le M.
 \tag{C.AP.7}\label{cdp:eq:C.AP.7}
\]
The actual finite ruled-removal and cleanup block in \ref{cdp:st:C.2.UN.77}
therefore applies. On a projective-line bundle its contraction
is precisely \ref{cdp:st:C.2.UN.74}; for nonminimal fibres it is the retained
sigma-atlas construction of
\cite[\S\S4.10.7--4.10.8,4.14]{Kulikov1977}, with the
missing-fibre correction of \cite[\S8]{Kulikov1981}.

\noindent\textbf{Step 7.} For clarity, the numerical and analytic properties of this
block used here are the ones already proved in \ref{cdp:st:C.2.UN.75} and
\ref{cdp:st:C.2.UN.77}. It removes a positive-order component \(q\), adds
finitely many components of order \(q+1\) when necessary,
and adds none when \(q=M\). The diagnostic resolution used
to check Lemma4.14 is not inserted into the component list.
Its actual sigma-atlas is cleaned up by choosing a largest
last local branch, following its finite descendants and
performing the same ruled removals. Empty lists give smooth
NC charts. All overlaps are restored by the restrictions of
the same proper local graph maps. The positive component
card vector strictly decreases across this whole block.
No monotonicity of \(h\) inside the block is required.

\par\medskip\noindent\textbf{Zero triple points: the exact plane comparison.}
Suppose that \(E\) has no triple point. Its square on \(Q\)
is \(+1\), and every other boundary on \(Q\) is disjoint
from \(E\). Here rationality does not require \(Q\) to be
a canonical maximum. A rational curve of positive square
cannot be exceptional for a minimal-model point map. On a
minimal model of nonnegative Kodaira dimension, the explicit
canonical lines in \ref{cdp:st:C.2.UN.31} and its quotient classification are
nef. Their pullback and the positive point discrepancies
would give \(K_Q\cdot E\ge0\), contrary to adjunction
\(K_Q\cdot E=-3\). A Hopf minimal model contains no rational
curve. On an irrational ruled minimal model every rational
curve is vertical for its ruling and has nonpositive square,
and point blowups cannot increase that square. The component
classification therefore leaves exactly the rational case.
One can specify its required plane
map without choosing a birational presentation. The sequence
\[
 0\longrightarrow\OO_Q\longrightarrow\OO_Q(E)
       \longrightarrow\OO_E(1)\longrightarrow0
 \tag{C.AP.8}\label{cdp:eq:C.AP.8}
\]
and \(H^1(Q,\OO_Q)=0\) give three sections. They have no base
point on \(E\), and the section defining \(E\) is nonzero
off \(E\). Thus \(|E|\) gives an actual morphism
\[
 \rho:Q\longrightarrow\PP^2,
 \qquad \rho^*\OO_{\PP^2}(1)=\OO_Q(E),
 \qquad \rho|_E:E\xrightarrow{\sim}\ell .
 \tag{C.AP.9}\label{cdp:eq:C.AP.9}
\]
Its image is a surface since \(E^2=1\); its degree is one
by the same intersection number. Hence it is a birational
morphism, factoring into the actual point blowdowns away
from \(E\). Every other boundary is contracted, since it
has zero intersection with \(E\).

\noindent\textbf{Step 8.} The canonical restriction to \(E\) gives
\(k_Q=k_P+2\). More generally call a current pair \((Q,E)\)
of the form \eqref{cdp:eq:C.AP.9}, with all other boundary disjoint from
\(E\), a marked plane pair, and write \(q=k_Q\). Its actual
canonical divisor has the precise form
\[
 K_Q=-3E+\sum_C a_CC,
 \qquad \sum_C a_CC=K_Q-\rho^*K_{\PP^2},
 \qquad a_C>0.
 \tag{C.AP.10}\label{cdp:eq:C.AP.10}
\]
Indeed choose on \(\PP^2\) the meromorphic two-form whose
divisor is \(-3\ell\). The ratio of the actual form on \(Q\)
to its pullback has divisor supported on the exceptional
curves. On \(\PP^2\) this same meromorphic function has no
divisorial zero or pole, so it is a nonzero constant. Thus
the equality in \eqref{cdp:eq:C.AP.10} is an equality of actual divisors,
with the positive discrepancies of the specified point map.
It shows that every exceptional prime is an actual
boundary, and that there is no additional boundary of
coefficient zero. The only coefficient on the marked curve
is \(-3\). A self-double curve away from that marking would
have coefficient \(-1\), so such a pair stays embedded at
the smooth-total NC stages under consideration.

\noindent\textbf{Step 9.} Choose a last point blowup of \(\rho\). Its exceptional
boundary curve \(L\) lies away from \(E\). At a smooth point
of a discrepancy component of coefficient \(a>0\) its new
coefficient is \(a+1\); at a node of coefficients \(a,b>0\)
it is \(a+b+1\); at a point off the discrepancy divisor it
is one. These are the actual point-map identities for
\(K_Q+3E\). They give the following exhaustive alternatives:
\[
\begin{array}{c|c|c}
 \text{triple points of }L&\text{order of its neighbour}
                       &\text{actual next operation}\\ \hline
 0&q+2&\text{another marked plane pair}\\
 1&q+a+2\ge q+3&\text{a ruled-removal block}\\
 2&q+a+b+2&\text{a type-II move}.
\end{array}
 \tag{C.AP.11}\label{cdp:eq:C.AP.11}
\]
In the last row the two receiving orders are
\(q+a+1\) and \(q+b+1\), both at least \(q+2\). The type-II
move decreases \(b_2(Q)\) by one; every other increase in
surface second Betti number occurs at a strictly higher
order. Its centre and all its points on \(Q\) are away from
the marked \(E\), so \eqref{cdp:eq:C.AP.9} persists on the modified \(Q\).
In the first row the same complete linear system constructs
the next marked plane pair, at a strictly greater order.

\noindent\textbf{Step 10.} This yields a finite progress block. Until its first whole
component removal, let
\[
                    b_j=\sum_{k_V=j}b_2(V^\nu),
                    \qquad 0\le j\le M,
 \tag{C.AP.12}\label{cdp:eq:C.AP.12}
\]
where the sum uses the actual smooth normalizations. Each
type-II step in \eqref{cdp:eq:C.AP.11} strictly decreases
\((b_0,\ldots,b_M)\) lexicographically. Passing to another
marked plane pair makes no modification and strictly raises
its order. There can be at most \(M\) such successive calls
before an actual step. Since lexicographic order on this
fixed finite product of nonnegative integers is well founded,
these steps cannot continue indefinitely without a component
removal. At a marked plane pair with no exceptional curve,
the surface is \(\PP^2\), its sole boundary is \(E=\ell\),
and the principal full-fibre restriction gives
\[
                 N_{Q/X}=\OO_{\PP^2}(-1).
 \tag{C.AP.13}\label{cdp:eq:C.AP.13}
\]
The actual point contraction of \ref{cdp:st:C.2.UN.64} removes it to a smooth
point. The adjacent surface is its ordinary point blowup,
and no other component passes through that point. In the
one-triple alternative the positive ruled-removal block
already supplies a component removal.

\noindent\textbf{Step 11.} Stop this progress block at its first whole-component
removal, and perform the finite sigma cleanup just described.
Its net positive card vector decreases. It is unnecessary
to claim that the first marked plane pair or its marking
survives an arbitrary later cleanup: after this completed
block the outer argument uses the actual designated parent
maps again. This distinction retains every map while avoiding
an unsupported global preservation assertion.

\par\medskip\noindent\textbf{Preserved maximal Hopf neighbourhoods.}
Every type-II centre used on a designated parent is vertical
for \eqref{cdp:eq:C.AP.2}, hence disjoint from its marked elliptic curves.
An opposite curve of a type-II move cannot lie on a minimal
maximal Hopf: a rational curve is impossible there, and the
elliptic alternative has the positive square \eqref{cdp:eq:C.AP.6}.
For the plane progress moves all opposite boundary curves
are rational. They likewise cannot meet a minimal Hopf.

\noindent\textbf{Step 12.} A ruled component removed in one of the actual blocks has
one boundary section, whose neighbouring order is \(q-1\),
and all other boundary is rational and vertical. It cannot
meet a minimal Hopf of order \(M\): that Hopf contains no
rational curve, while its being the section parent would
require \(q=M+1\). The same statement applies to every
cleanup removal and missing-fibre insertion. Plane point
contractions also have rational boundary. Thus these actual
closed centres and local sigma charts avoid the entire
union of maximal minimal Hopfs. Their restored overlap maps
are the identity on an open neighbourhood of that union.

\noindent\textbf{Step 13.} The designated parent prime divisors cannot be removed:
each meets one of these unchanged Hopfs. Their marked
elliptic neighbourhoods therefore persist. At every smooth
NC stage the exact maps \eqref{cdp:eq:C.AP.2}--\eqref{cdp:eq:C.AP.4} can be applied
again to their current normalizations. No holomorphic
contraction of a Hopf neighbourhood has been invoked.

\par\medskip\noindent\textbf{The fixed finite termination potential.}
Let \(c_j\) count the retained global components of order
\(j\), for \(1\le j\le M\). Use the block-level potential
\[
                  (c_1,c_2,\ldots,c_M,h).
 \tag{C.AP.14}\label{cdp:eq:C.AP.14}
\]
All coordinates are nonnegative integers. A direct type-II
step on a last designated-parent point leaves the cards
unchanged and decreases \(h\) by \eqref{cdp:eq:C.AP.6}. Either of the
other last-point alternatives supplies a finite progress
block with a strict decrease of the first changed card.
Any finite defect increase inside that whole block is
irrelevant to this lexicographic decrease. The missing-fibre
and cleanup comparisons keep all retained orders at most
the original \(M\); their auxiliary common blowups are not
additional retained cards. Hence \eqref{cdp:eq:C.AP.14} has no infinite
strictly descending sequence.

\noindent\textbf{Step 14.} When this procedure stops, a positive designated-parent
defect would still have a last point curve. \ref{cdp:st:C.2.UN.80} and the
three exhaustive alternatives would supply another strict
decrease. Thus every such defect is zero. Each designated
normalization is the actual minimal elliptic model in
\eqref{cdp:eq:C.AP.2}, and every boundary is a fibre support. All new
self-incidences created by the type-I/type-II and ruled
sigma charts are rational exceptional curves. None can
remain on these parents, since their rational boundary
would have to be exceptional for a point map of defect
zero. No elliptic self-incidence was introduced: a ruled
removal has only one horizontal boundary section, and all
its other changes are at its rational fibre curves. The
designated parents are consequently embedded, with disjoint
elliptic boundaries and no triple point.

\noindent\textbf{Step 15.} Every arrow used above is one of the actual two-ruling,
curve, point, or restored sigma-atlas maps already specified.
The inverse record, including all normalization branches
and overlap restrictions, reconstructs the given family.
The proper graph gluing and annular comparison of \ref{cdp:st:C.2.UN.75}
give the same record on a compact base, and its normal-cone
Euler formula applies at every retained step. This proves
the stated marked-interface reduction in its exact scope.
\end{proof}






\begin{proposition}[Every positive maximal Hopf after minimalization is a torsion-normal leaf]\label{cdp:st:C.2.UN.82}
\StatementRecord{cdp:st:C.2.UN.82}{proved}{Historical events: Claude Fable 5.2 (Anthropic), effort high, 16 November 2026: carried as background without fault found. GPT-6 Astra (OpenAI), ChatGPT Pro mode, effort Pro, S6\_09\_PAPERS, 17 November 2026: targeted proof check. Also checked in this writing. Pre-continuation writing events: GPT-6 Astra (OpenAI), ChatGPT Pro mode, effort Pro, S6\_10\_PAPERS, 18 November 2026: targeted proof check. Continuation metadata events: GPT-6 Astra (OpenAI), ChatGPT Pro mode, effort Pro, S6\_10\_PAPERS continuation, 19 November 2026: used and rechecked in this run; status-vocabulary reconciliation; metadata reconciliation only; no new mathematical check. \emph{Pre-continuation recorded proof scope.} The exact anticanonical divisor makes a rational component's support connected. On an elliptic ruled minimal model the residue connecting map and nonzero degree-two trace exclude a connected bisection. The primary Hopf weight equation and nonlinear translation kernel show that one simple factor forces the entire anticanonical divisor to be reduced. The actual point-map coefficient rule preserves this property and the isolated elliptic branches. Propagation through the connected incidence graph forces every order to equal the primitive minimum. At a positive maximal Hopf this excludes simple factors, leaving one elliptic axis; adjunction and the full-fibre normal product give the exact torsion-normal interface}{the hypotheses, notation and definitions stated here; \cref{cdp:st:C.2.UN.75}; \cref{cdp:st:C.2.UN.77}; \cref{cdp:st:C.2.UN.70}; \cref{cdp:st:C.2.UN.19}; \cref{cdp:eq:C.ML.5}; \cref{cdp:eq:C.ML.6}; \cref{cdp:eq:C.ML.8}; \cref{cdp:eq:C.ML.3}}

Let \(f:X\to\Delta\) be a proper degeneration of connected
complex two-tori from a smooth threefold, with full reduced
NC central fibre and smooth normalization components. Retain
every normalized branch of a self-double curve, as in \ref{cdp:st:C.2.UN.75}.
For the primitive canonical evaluation write
\[
       \operatorname{div}(\Omega)=\sum_i k_iV_i,
       \qquad \min_i k_i=0,\qquad M=\max_i k_i .
\tag{C.ML.1}\label{cdp:eq:C.ML.1}
\]
On the normalization \(T\) of a component of order \(M\),
suppose that one boundary branch \(C\) is smooth elliptic,
is disjoint from all its other boundary branches, and has
coefficient one in its induced anticanonical divisor.
Then every central component has order zero and there
are no triple points.

Consequently consider the actual endpoint of \ref{cdp:st:C.2.UN.77} for
an original full reduced SNC family with \(M>0\).
Every Hopf component \(H\) of this maximal order is a
terminal component. Its boundary is one smooth elliptic
curve \(C=H\cap P\), and its exact canonical and normal
data satisfy
\[
 \begin{gathered}
       -K_H=aC,\qquad a=M-k_P+1\ge2,\\
       N_{C/H}^{\otimes(a-1)}\simeq\OO_C,
       \qquad N_{C/P}\simeq N_{C/H}^{-1},\\
       k_P\le M-1,\qquad
       \operatorname{coeff}_C(K_P)=a-2\ge0 .
 \end{gathered}
\tag{C.ML.2}\label{cdp:eq:C.ML.2}
\]
The normal-line equality uses the actual local branches
on \(P\) when another part of that component is nonnormal.
Each such interface therefore supplies precisely the
marked torsion-normal, nonnegative-coefficient data of
\ref{cdp:st:C.2.UN.70} on that normalization. No neighbourhood is replaced
in this conclusion; the ambient comparison remains the
finite record of \ref{cdp:st:C.2.UN.77}.
\end{proposition}



\begin{proof}
\noindent\textbf{Step 1.} \par\medskip\noindent\textbf{The induced divisor at the maximal order.}
The normalization adjunction map for the reduced NC
fibre, including both branches over every self-double
curve, gives the actual meromorphic canonical section
on \(T\) with
\[
      -K_T=D_T=\sum_B (M-k_{\mathrm{other}(B)}+1)B.
\tag{C.ML.3}\label{cdp:eq:C.ML.3}
\]
Every displayed coefficient is a positive integer.
For a self-double branch its coefficient is one.
This is the restriction of the common volume form and
the principal full-fibre trivialization; it retains
the line isomorphism, rather than just intersections.
The support of \(D_T\) is exactly the full marked
boundary. The curve \(C\) in the statement is one
isolated component of this effective divisor, with
coefficient one.

\noindent\textbf{Step 2.} We first prove the following surface consequence of
these data: the whole divisor \(D_T\) is reduced and
is a disjoint union of smooth elliptic curves. The
minimal model of \(T\) is ruled or Hopf. Indeed the
nonzero effective anticanonical divisor makes every
positive pluricanonical section vanish: through
\(K_T^{\otimes n}\simeq\OO_T(-nD_T)\), such a
section would be a holomorphic function on the compact
connected surface vanishing on \(D_T\). Thus
\(\kappa(T)=-\infty\). The surface classification
and the actual component restriction \ref{cdp:st:C.2.UN.19} give the
stated ruled or Hopf alternatives.

\par\medskip\noindent\textbf{Rational surfaces: the entire divisor is the marked elliptic curve.}
For a rational \(T\), the exact divisor sequence is
\[
 0\longrightarrow K_T\longrightarrow\OO_T
          \longrightarrow\OO_{D_T}\longrightarrow0.
\tag{C.ML.4}\label{cdp:eq:C.ML.4}
\]
Here \(H^0(K_T)=0\) and
\(H^1(K_T)=H^1(\OO_T)^*=0\) by Serre duality.
Consequently \(H^0(\OO_{D_T})=\C\). An effective
divisor with disconnected support has independent
constant functions on its disjoint components, so
its support is connected. Since \(C\) is isolated
in that support and has coefficient one, the exact
divisor is \(D_T=C\).

\par\medskip\noindent\textbf{Keeping the marked curve through the surface point maps.}
Take an actual sequence of point blowdowns
\(\pi:T\to T_0\) to a smooth minimal ruled or Hopf
surface. Push the meromorphic canonical section
through their common function-field identification.
The resulting anticanonical divisors remain effective.
At a reverse point blowup \(\sigma\), their exact
comparison is
\[
                 D_{\mathrm{new}}=\sigma^*D_{\mathrm{old}}-E.
\tag{C.ML.5}\label{cdp:eq:C.ML.5}
\]
The curve \(C\) is never exceptional, since it is
elliptic. Its image remains a smooth isolated
coefficient-one component at every blowdown. To
check this assertion at a last exceptional curve
\(E\), first suppose its coefficient in the
effective divisor is positive. It is then disjoint
from \(C\), so the contraction is an isomorphism
near \(C\). If its coefficient is zero, \eqref{cdp:eq:C.ML.5}
says that the old divisor has total multiplicity
exactly one at the centre. When \(E\) meets \(C\),
this forces the image of \(C\) to be the only
positive branch there, of multiplicity one. It is
smooth and no other boundary image meets it.
Equivalently, a curve of multiplicity \(r\) at
the centre has total intersection \(r\) with
the exceptional divisor; here \(r=1\).
This proves the assertion by induction through
the specified point maps. Write \(C_0=\pi(C)\).

\par\medskip\noindent\textbf{Irrational ruled surfaces: two coefficient-one sections.}
Let \(p:T_0\to B\) be the minimal ruling. The
elliptic curve \(C_0\) cannot be contained in a
rational ruling fibre. Its map to \(B\) is finite,
so Riemann--Hurwitz gives \(g(B)\le1\).
The rational case has already been proved; hence
\(B\) is elliptic and that finite map is etale.
For a ruling fibre \(F\),
\[
                 D_{T_0}\cdot F=-K_{T_0}\cdot F=2.
\tag{C.ML.6}\label{cdp:eq:C.ML.6}
\]
Every positive vertical curve would meet \(C_0\),
contrary to its isolation. Thus there is no such
vertical part. If \(C_0\) has degree two over
\(B\), \eqref{cdp:eq:C.ML.6} gives \(D_{T_0}=C_0\).
This connected anticanonical bisection is impossible
by the following exact residue and trace comparison.
The nowhere-zero section of \(K_{T_0}(C_0)\)
has nonzero residue on \(C_0\), so in
\[
 0\longrightarrow K_{T_0}\longrightarrow K_{T_0}(C_0)
       \xrightarrow{\operatorname{res}}K_{C_0}\longrightarrow0
\tag{C.ML.7}\label{cdp:eq:C.ML.7}
\]
the connecting map \(H^0(K_{C_0})\to H^1(K_{T_0})\)
is zero. Its Serre-dual map is restriction
\(H^1(\OO_{T_0})\to H^1(\OO_{C_0})\).
Composed with the ruling isomorphism from
\(H^1(\OO_B)\), this is the pullback of the
degree-two etale map \(C_0\to B\). Trace composed
with that pullback is multiplication by two,
which is nonzero. This contradiction excludes
the bisection.

\noindent\textbf{Step 3.} The remaining possibility in \eqref{cdp:eq:C.ML.6} is
\[
                     D_{T_0}=C_0+C'_0,
\tag{C.ML.8}\label{cdp:eq:C.ML.8}
\]
with two disjoint degree-one sections, each of
coefficient one. A degree-one irreducible curve
over the smooth ruling base is the graph of
its section and is smooth. Thus both curves
are smooth elliptic.

\noindent\textbf{Step 4.} Pull \eqref{cdp:eq:C.ML.8} through the actual reverse point
blowups. Every intermediate divisor is effective:
a negative coefficient would remain negative
on its later strict transform. A centre away
from the positive support would create coefficient
\(-1\) in \eqref{cdp:eq:C.ML.5}, so no such centre occurs.
Starting with the two disjoint coefficient-one
curves, a centre on either has total positive
multiplicity one and creates coefficient zero.
The positive strict transforms remain disjoint
and smooth. Induction shows that no positive
exceptional coefficient ever appears. Therefore
\(D_T\) consists exactly of their two smooth
elliptic strict transforms, still with coefficient
one. All these are comparisons through the actual
surface morphism \(\pi\).

\par\medskip\noindent\textbf{Hopf surfaces: one simple factor makes the anticanonical divisor reduced.}
Use the actual unramified primary cover of the
Hopf minimal surface \(T_0\). The coefficient-one
curve \(C_0\) has reduced unramified inverse
image and retains that coefficient. In a primary
normal form the pulled-back anticanonical section
is \(g(z,w)\partial_z\wedge\partial_w\).
Hartogs extends \(g\) across zero, and its
functional equation is the constant-Jacobian
equation from (C.HA.5) and (C.CZ.11).

\noindent\textbf{Step 5.} For a nonlinear contraction
\((z,w)\mapsto(\beta^r z+w^r,\beta w)\), its
only possible eigenfunction is a nonzero multiple
of \(w^{r+1}\), with \(r\ge1\). Its single
curve has coefficient at least two, contradicting
the given simple factor.

\noindent\textbf{Step 6.} For a diagonal contraction without a resonance,
the section is a multiple of \(zw\). In the
resonant case pass through the actual finite
unramified central-power cover, so its multipliers
are \((\gamma^p,\gamma^q)\), with coprime
positive \(p,q\). Its monomials have weight
\(p+q\). The exact equation
\(pi+qj=p+q\) gives
\(i=1+q\ell\), \(j=1-p\ell\) for an integer
\(\ell\), by coprimality. For \(p,q>1\),
nonnegative exponents force \(\ell=0\), so the
only monomial is \(zw\). For \(p>1,q=1\), the section has form
\[
                        g=w(az+bw^p).
\tag{C.ML.9}\label{cdp:eq:C.ML.9}
\]
When \(a=0\) its only factor has order \(p+1\),
so the given coefficient one forces \(a\ne0\).
Then its two distinct smooth factors each have
order one. The case \(p=1,q>1\) is the same
with the coordinates interchanged. For \(p=q=1\),
the nonzero quadratic is either the square of
one linear form or the product of two distinct
linear forms. The simple factor excludes the
first alternative. Thus the entire pulled-back
divisor is reduced, consisting of two disjoint
elliptic curves on the primary Hopf surface.
Unramified descent retains every coefficient,
so the original \(D_{T_0}\) is reduced as well.
Its curves are disjoint smooth elliptic curves.

\noindent\textbf{Step 7.} The same point-map induction following \eqref{cdp:eq:C.ML.8}
now applies to \(T\): every positive coefficient
is one, positive branches never meet, and the
additional exceptional coefficients are zero.
It proves the asserted surface consequence for
every blown-up Hopf component too.

\par\medskip\noindent\textbf{Propagation through the actual central support.}
The preceding arguments apply to the smooth
normalization of every maximal component carrying
one isolated coefficient-one elliptic branch.
They show that all its boundary branches have
these same properties. By \eqref{cdp:eq:C.ML.3}, the component
across each one again has order \(M\). Since the
branch has no triple point, its counterpart on
that normalization is again an isolated smooth
elliptic branch of coefficient one. The argument
can therefore be repeated across every incident
edge. Self-incidences retain both normalization
branches and introduce no new component order.

\noindent\textbf{Step 8.} The finite incidence graph of the connected
central support is connected. The resulting
set of components of order \(M\) has no boundary
edge to its complement; it is consequently the
whole support. Every normalized boundary is
elliptic and isolated, so no triple point occurs.
Primitivity of the original direct-image generator
gives \(\min_i k_i=0\), hence \(M=0\).
This proves the first assertion.

\par\medskip\noindent\textbf{The remaining maximal Hopfs are the required leaves.}
At the endpoint of \ref{cdp:st:C.2.UN.77} the maximal Hopf components
are actual embedded minimal Hopf surfaces, and
their boundary curves are disjoint smooth elliptic
curves. The comparison preserves their orders and
the primitive canonical direct-image line. For
\(M>0\), the first assertion excludes coefficient
one on any such boundary.

\noindent\textbf{Step 9.} The eigenfunction classification just proved
also shows that an effective anticanonical divisor
on a minimal Hopf surface either is reduced or
is supported on one elliptic curve: the nonlinear
form is a pure power, and the only nonreduced
diagonal alternatives above are pure powers of
one smooth factor. This conclusion descends
through the finite unramified primary cover.
Thus the nonempty boundary here is exactly one
curve \(C\), with coefficient \(a\ge2\).
Its unique neighbour has order \(M-a+1\).

\noindent\textbf{Step 10.} Adjunction on this actual elliptic curve gives
\[
     \OO_C\simeq K_C=(K_H+C)|_C
                          \simeq N_{C/H}^{\otimes(1-a)}.
\tag{C.ML.10}\label{cdp:eq:C.ML.10}
\]
There are no triple points on \(C\), so the
full reduced-fibre trivialization gives
\(N_{C/H}\otimes N_{C/P}\simeq\OO_C\).
Finally the exact coefficient on \(P\) is
\(M-k_P-1=a-2\). These are precisely all the
claims in \eqref{cdp:eq:C.ML.2}. The argument uses only the
retained normalized branches, actual surface
point maps, and finite primary covers. It
does not prune a Hopf neighbourhood or assert
a different smoothing of the remaining support.
The exact weight spaces, polynomial factor cases,
nonlinear translation kernels, labelled point-map
forests and maximal-leaf order identities are
independently checked in
\path{code/root_maximal_hopf_leaf.py}; its executed
output is retained beside it. The proof above supplies
the general maps and unbounded arguments.
\end{proof}






\begin{proposition}[The negative-coefficient interfaces have exact elliptic or one-section models]\label{cdp:st:C.2.UN.83}
\StatementRecord{cdp:st:C.2.UN.83}{proved}{Historical events: Claude Fable 5.2 (Anthropic), effort high, 16 November 2026: carried as background without fault found. GPT-6 Astra (OpenAI), ChatGPT Pro mode, effort Pro, S6\_09\_PAPERS, 17 November 2026: targeted proof check. Also checked in this writing. Pre-continuation writing events: GPT-6 Astra (OpenAI), ChatGPT Pro mode, effort Pro, S6\_10\_PAPERS, 18 November 2026: targeted proof check. Continuation metadata events: GPT-6 Astra (OpenAI), ChatGPT Pro mode, effort Pro, S6\_10\_PAPERS continuation, 19 November 2026: used and rechecked in this run; status-vocabulary reconciliation; metadata reconciliation only; no new mathematical check. \emph{Pre-continuation recorded proof scope.} The canonical intersection and square formulas retain the marked elliptic neighbourhood through the actual point maps. For a multisection, the finite etale ruling base change gives disjoint sections; their exact Picard difference and torsion order produce a connected elliptic pencil, and the deck action descends its Stein base. For a section, the explicit disjoint fibre contractions retain its normal line. A surviving horizontal boundary kills the additive torsor by finite trace, giving the actual quotient elliptic map. With no such boundary, the exact canonical divisor forces K=-2J and the retained possibly nonsplit rank-two extension. The nonlinear Hopf comparison keeps its axis and every positive exceptional discrepancy}{the hypotheses, notation and definitions stated here; \cref{cdp:st:C.2.UN.74}; \cref{cdp:st:C.2.UN.77}; \cref{cdp:st:C.2.UN.81}; \cref{cdp:st:C.2.UN.67}; \cref{cdp:st:C.2.UN.19}; \cref{cdp:st:C.2.UN.31}; \cref{cdp:st:C.2.UN.70}; \cref{cdp:st:C.2.UN.80}; \cref{cdp:eq:C.TI.5}; \cref{cdp:eq:C.TI.12}; \cref{cdp:eq:C.TI.9}; \cref{cdp:eq:C.TI.13}}

Let \(T\) be a smooth normalization component of a full
reduced NC fibre of a smooth-total degeneration of connected
complex two-tori. Retain the two normalized branches of
every self-double curve. Let \(J\subset T\) be a smooth
elliptic boundary branch, disjoint from every other boundary
branch, whose exact normal line is torsion. No sign is
imposed on its coefficient in the actual meromorphic
canonical divisor. Then \(T\) is not rational, and one
of the following explicit comparisons holds.

\begin{enumerate}
\item There is an actual sequence of point blowdowns
\[
 T\xrightarrow{\pi}Z\xrightarrow{q_Z}B
 \tag{C.TI.1}\label{cdp:eq:C.TI.1}
\]
which is an isomorphism on a neighbourhood of \(J\).
The map \(q_Z\) is a relatively minimal fixed-Jacobian
elliptic quasi-bundle with identity integral monodromy
and base genus at most one. Its fibres have smooth
elliptic supports. Every original boundary and every
exceptional curve of \(\pi\) is vertical, and every
exceptional prime is an actual boundary.

\item There is an actual point map \(\pi:T\to H\),
an isomorphism near \(J\), onto a nonlinear minimal
Hopf surface. Its only curve is \(J_0=\pi(J)\).
For an integer \(a\ge2\), the actual canonical
comparison is
\[
 K_H=-aJ_0,\qquad
 K_T=-aJ+\sum_E b_EE,\qquad b_E>0.
 \tag{C.TI.2}\label{cdp:eq:C.TI.2}
\]
The other boundary support is exactly the exceptional
set of \(\pi\). Every neighbouring order along such
an exceptional curve exceeds \(k_T+1\).

\item There is an actual point map \(\pi:T\to R\),
an isomorphism near \(J\), onto a projective-line
bundle over an elliptic curve. The image \(J_0\)
is a section, and
\[
 K_R=-2J_0,\qquad N_{J_0/R}\simeq\OO_{J_0},\qquad
 K_T=-2J+\sum_E b_EE,\quad b_E>0.
 \tag{C.TI.3}\label{cdp:eq:C.TI.3}
\]
Every other boundary is exactly one of these
exceptional primes. After tensoring its defining
rank-two bundle by a line, \(R\) is
\(\PP(V)\) for an extension
\[
              0\longrightarrow\OO\longrightarrow V
                    \longrightarrow\OO\longrightarrow0.
 \tag{C.TI.4}\label{cdp:eq:C.TI.4}
\]
The extension is retained and need not split.
If \(P\) is the component across \(J\), then
\(k_T=k_P+1>0\). This is the actual one-section
ruled-removal situation of \ref{cdp:st:C.2.UN.74} and the finite
restored ruled block of \ref{cdp:st:C.2.UN.77}/\ref{cdp:st:C.2.UN.81}.
\end{enumerate}

In the ruled birational case, alternative (3) occurs
only when \(J\) is a degree-one section of the actual
ruling and no other horizontal boundary survives the
point contractions disjoint from \(J\). A higher-degree
elliptic multisection, or a second surviving horizontal
boundary, produces the actual elliptic map in (1).
All comparisons retain the original marked normal line
and the actual boundary branches. No replacement of an
arbitrary threefold neighbourhood by a prescribed
logarithmic model is used.
\end{proposition}



\begin{proof}
\noindent\textbf{Step 1.} \par\medskip\noindent\textbf{The exact canonical data and the rational exclusion.}
Write the actual restricted canonical section as
\[
              K_T=a_JJ+\sum_C a_CC.
 \tag{C.TI.5}\label{cdp:eq:C.TI.5}
\]
The sum includes every other actual boundary branch,
even when its coefficient is zero. All intersections
of \(J\) with that support are empty. Its torsion
normal gives \(J^2=0\) and \(K_T\cdot J=0\).
\ref{cdp:st:C.2.UN.67} already excludes a rational component under
exactly these hypotheses, without a coefficient
sign assumption. Its actual restriction sequences
produce \(|mJ|\), with
\(m=\operatorname{ord}N_{J/T}\); the exact order
forces connectedness, and the limiting ambient
fields give the coherent-Euler contradiction.
The same proof works on this smooth normalization:
its fields are holomorphic off the full normalized
boundary, and its marked neighbourhood is embedded.

\noindent\textbf{Step 2.} The actual component restriction \ref{cdp:st:C.2.UN.19} and \ref{cdp:st:C.2.UN.31}
leaves a Hopf, a nonrational ruled, or a
nonnegative-Kodaira-dimension minimal model.

\par\medskip\noindent\textbf{Nonnegative Kodaira dimension.}
Use the ordinary minimal-model morphism
\(\pi:T\to Z\). Its explicit canonical base-line
formula, or the trivial canonical line in the
torus case, gives \(K_Z\cdot C\ge0\) for every
curve \(C\). The exact point-map formula is
\[
 K_T=\pi^*K_Z+\sum_E b_EE,\qquad b_E>0.
 \tag{C.TI.6}\label{cdp:eq:C.TI.6}
\]
The curve \(J\) is not exceptional. Adjunction
therefore gives
\[
 0=K_T\cdot J
   =K_Z\cdot\pi(J)+\sum_E b_E(E\cdot J).
 \tag{C.TI.7}\label{cdp:eq:C.TI.7}
\]
Every summand is nonnegative, so the point map is
an isomorphism near \(J\). Its image is a smooth
elliptic curve with the same torsion normal.
The actual marked comparisons (C.NP.3)--(C.NP.4)
and (C.MR.12)--(C.MR.14) now apply without using
the sign of \(a_J\). In Kodaira dimension one,
the positive canonical base degree makes this
curve vertical. On a torus its quotient by the
parallel elliptic subtorus supplies the map.
In the remaining dimension-zero case, a
horizontal marked curve trivializes the Chern
vector after the specified finite translation
root cover; the actual period lattice is a
torus lattice, and the parallel marked curves
give the descended elliptic quotient.

\noindent\textbf{Step 3.} Thus there is an actual elliptic map making
the image of \(J\) a full fibre support.
Every other noncontracted boundary is disjoint
from that support and hence vertical; every
contracted curve is vertical too. The ambient
field alignment and the translation-root
comparison of \ref{cdp:st:C.2.UN.70} then give \eqref{cdp:eq:C.TI.1}.
The proof of \ref{cdp:st:C.2.UN.80} uses only this actual map
and its vertical boundary, so gives its final
exceptional-prime assertion.

\par\medskip\noindent\textbf{Hopf minimal model.}
For a Hopf minimal model \(H\), every curve
is a smooth elliptic curve of square zero.
The blowup square formula
\[
                  0=J^2=J_0^2-\sum_\nu r_\nu^2
 \tag{C.TI.8}\label{cdp:eq:C.TI.8}
\]
shows that no centre lies on a successive
image of \(J\). Thus \(\pi\) is again an
isomorphism near it. A diagonal Hopf axis
has torsion normal precisely in the resonant
elliptic case, through the actual finite
primary cover and its multiplier character.
In that case the actual Hopf elliptic map
makes every surviving boundary vertical;
its composition with \(\pi\), followed by
the original ambient-field comparison,
gives (1).

\noindent\textbf{Step 4.} If \(H\) is nonlinear, its unique elliptic
curve is \(J_0\), so every other boundary
is contracted. Push \eqref{cdp:eq:C.TI.5} through the
actual common open isomorphism. On a
primary cover, the nonlinear normal form
has canonical divisor \(-(r+1)J_0'\) for
an integer \(r\ge1\). Comparing divisors
of the lifted actual meromorphic section
forces \(a_J=-(r+1)\): a different exponent
would give a nonzero principal divisor
supported on just one effective curve,
impossible on a compact surface. Unramified
descent preserves this coefficient. Put
\(a=r+1\). The point-map formula now gives
\eqref{cdp:eq:C.TI.2}. Since \eqref{cdp:eq:C.TI.5} has no other
noncontracted support, every exceptional
prime has its positive discrepancy as
coefficient and is an actual boundary.
Conversely every other boundary is
exceptional. Its neighbour has order
\(k_T+b_E+1\ge k_T+2\).

\par\medskip\noindent\textbf{The actual ruling and its elliptic multisection.}
It remains to treat a ruled birational
class. Its actual ruling is a morphism
\[
                       p:T\longrightarrow B_0
 \tag{C.TI.9}\label{cdp:eq:C.TI.9}
\]
obtained by composing the specified point
map with a minimal ruling. The image of
the elliptic curve \(J\) is not a point:
all components of a ruling fibre are
rational. Riemann--Hurwitz makes \(B_0\)
elliptic and \(p|_J\) an etale isogeny,
of degree \(d\ge1\). The rational base
case was already excluded.

\noindent\textbf{Step 5.} First let \(d\ge2\). Make the actual
finite etale base change
\[
 T'=T\times_{B_0}J\longrightarrow T.
 \tag{C.TI.10}\label{cdp:eq:C.TI.10}
\]
The isogeny is Galois. Its group \(G\)
acts freely on \(T'\), and \(T'/G=T\).
The inverse image of \(J\) consists of
\(d\) disjoint sections \(J_1,\ldots,J_d\)
of \(p':T'\to J\), each with the pulled
torsion normal and square zero.

\noindent\textbf{Step 6.} For two of these sections, Hodge index
gives numerical equivalence. Indeed their
intersection and both squares are zero;
an ample class determines a linear
combination orthogonal to that ample
class, still of square zero, which is
numerically zero. Intersection with a
general ruling fibre then makes its two
coefficients equal. On this point-blown-up
elliptic ruled surface \(H^2(T',\Z)\) is
torsion-free, and a numerically trivial
line lies in \(\operatorname{Pic}^0(T')\).
The actual ruling pullback
\[
        (p')^*:\operatorname{Pic}^0(J)
                       \xrightarrow{\sim}\operatorname{Pic}^0(T')
 \tag{C.TI.11}\label{cdp:eq:C.TI.11}
\]
is an isomorphism: it is so for a
projective-line bundle by its direct
images, and each actual point blowup
preserves this exponential-sequence map.
Consequently
\(\OO(J_1-J_2)=(p')^*L\). Restriction to
the section \(J_1\) identifies \(L\)
with \(N_{J_1/T'}\), because \(J_2\)
is disjoint from \(J_1\).

\noindent\textbf{Step 7.} Let \(m\) be the exact order of that
normal line. Then the two disjoint
divisors \(mJ_1,mJ_2\) are linearly
equivalent. Their sections have no
common zero and give an actual pencil
\[
                  q':T'\longrightarrow\PP^1.
 \tag{C.TI.12}\label{cdp:eq:C.TI.12}
\]
It has connected fibres. To check this,
take its Stein factor with finite degree
\(e\). The inverse image of the value
with divisor \(mJ_1\) has just one
point on the Stein base, so that finite
map is totally ramified there of
order \(e\). The connected-fibre map
has fibre \((m/e)J_1\). Its local
principal equation gives
\(N_{J_1/T'}^{m/e}=\OO\), and exact
order \(m\) forces \(e=1\).
Adjunction, \(J_1^2=K_{T'}\cdot J_1=0\),
then makes the connected general fibre
of \eqref{cdp:eq:C.TI.12} elliptic.

\noindent\textbf{Step 8.} Every \(J_i\) is vertical for \(q'\).
It is the whole reduced support of
its fibre: since \(J_i^2=0\), the
intersection of that full fibre
with \(J_i\) forces every other
component to be disjoint from it,
and connectedness excludes such
additional components.
For \(g\in G\) and a general connected
fibre \(F\), the intersection
\((gF)\cdot J_1=F\cdot g^{-1}J_1=0\)
makes \(gF\) vertical too. Thus \(G\)
acts on the actual Stein base and
\(q'\) descends through \(T'/G=T\).
The kernel of the base action acts
freely on a general elliptic fibre;
it therefore acts by translations,
and the descended general fibre
is elliptic. This gives an actual
connected elliptic map on \(T\).
The marked \(J\) is its whole reduced
fibre support. Every other original
boundary pulls back disjointly from
all the \(J_i\), hence is vertical
for \(q'\) by intersection and is
vertical after descent. The original
ambient fields can now be applied
on \(T\) itself. The rank-one return
identities or the rank-two transverse
flow give identity monodromy, and
vertical minimization gives exactly
(1), isomorphic near \(J\).

\par\medskip\noindent\textbf{A section: the contraction retaining its actual normal line.}
Now let \(d=1\), so \(J\) is a section
of \eqref{cdp:eq:C.TI.9}. Every reducible ruling
fibre admits a rational \((-1)\)-curve
disjoint from \(J\). Here is a direct
comparison with its actual point maps.
Take a last blowup in its finite
construction. If the last exceptional
does not meet \(J\), use it. Otherwise
the centre is the smooth point where
the old section meets the old fibre:
a section never passes through a
node or a multiple component.
If the old fibre is irreducible,
its strict transform is the required
disjoint \((-1)\)-curve. If it is
reducible, induction supplies an old
disjoint \((-1)\)-curve, unchanged
by this last blowup at the section.
The chosen curve has valence at most
two; the actual point blowdown
restores the corresponding smooth
or nodal fibre chart. Repeating
decreases the finite number of
components and is an isomorphism
near the section throughout.

\noindent\textbf{Step 9.} Doing this over its finitely many
singular values gives
\[
 T\xrightarrow{\pi_0}R=\PP(V_0)
                    \xrightarrow{p_0}B_0,
 \qquad J_0=\pi_0(J),
 \qquad N_{J_0/R}=N_{J/T}.
 \tag{C.TI.13}\label{cdp:eq:C.TI.13}
\]
All centres and their actual maps
are retained. This construction
does not push two previously
intersecting marked branches
together at the section.
Every other surviving boundary
is disjoint from \(J_0\).
A vertical ruling fibre meets
the section, so any such surviving
boundary is horizontal.

\noindent\textbf{Step 10.} Suppose that one exists. Kill the
torsion normal by the actual finite
etale cover \(\widetilde B_0\to B_0\).
After a line twist, the pulled
projective bundle has section
normal \(\OO\); its complement
of that section is an additive
\(\OO_{\widetilde B_0}\)-torsor,
of class
\(\xi\in H^1(\widetilde B_0,\OO)\).
Normalize a connected component
of the pulled horizontal boundary.
It gives a finite curve map
\(a:C\to\widetilde B_0\) and an
actual section of the pulled
torsor, hence \(a^*\xi=0\).
The finite trace satisfies
\(\operatorname{tr}\circ a^*
 =(\deg a)\operatorname{id}\)
on \(H^1(\OO)\). Therefore
\(\xi=0\), including a possibly
ramified map \(a\).

\noindent\textbf{Step 11.} The pulled marked bundle is thus
\(\widetilde B_0\times\PP^1\).
The finite deck maps preserve the
section at infinity. Their affine
coefficients on the other coordinate
are holomorphic functions on the
compact base, hence constants.
The finite affine group has a
common fixed point by averaging,
so its action preserves the actual
elliptic projection to \(\PP^1\).
The quotient gives an elliptic
quasi-bundle on \(R\), with \(J_0\)
a fibre support. Each other compact
horizontal boundary lifts into
\(\widetilde B_0\times\C\); its
second coordinate is constant
on its compact normalization.
It is therefore another fibre
of this elliptic map. Compose
with \(\pi_0\). All exceptional
curves and all original boundaries
are vertical. Applying the original
ambient fields and then minimizing
this elliptic map proves (1).

\noindent\textbf{Step 12.} Finally suppose there is no other
surviving horizontal boundary.
Push the actual canonical section
\eqref{cdp:eq:C.TI.5} through \eqref{cdp:eq:C.TI.13}.
Its divisor is supported just on
\(J_0\). Intersecting a ruling
fibre gives its coefficient \(-2\).
Thus \(K_R=-2J_0\) as the exact
divisor of this meromorphic
section. Adjunction on the
elliptic section gives
\[
 \OO_{J_0}=K_{J_0}
          =(K_R+J_0)|_{J_0}=N_{J_0/R}^{-1}.
 \tag{C.TI.14}\label{cdp:eq:C.TI.14}
\]
The section normal is actually
trivial. The defining quotient
of the rank-two bundle, after
tensoring by its kernel line,
is therefore the extension
\eqref{cdp:eq:C.TI.4}; no vanishing of its
extension class has been used.

\noindent\textbf{Step 13.} The actual point-map formula is
now \eqref{cdp:eq:C.TI.3}. Every exceptional
prime has positive discrepancy,
so comparison with \eqref{cdp:eq:C.TI.5}
forces it to be a boundary.
Every other boundary is
exceptional by construction.
The coefficient \(-2\) on \(J\)
is \(k_P-k_T-1\), so
\(k_T=k_P+1>0\). The surviving
ruled model has exactly its
one horizontal boundary section;
the remaining full boundary on
\(T\) consists of vertical
exceptional configurations.
These are precisely the actual
one-section inputs to the
restored finite ruled-removal
block of \ref{cdp:st:C.2.UN.77}/\ref{cdp:st:C.2.UN.81}, whose
bundle stage is \ref{cdp:st:C.2.UN.74}.
This proves (3) with all
normal and extension data.

\noindent\textbf{Step 14.} All three cases use actual point
maps, finite covers, trace maps,
and descended Stein fibrations.
The finite section-fibre graph
checks and exact normal-order
divisibility certificates are
executed in
\path{code/A_torsion_interfaces.py}.
\end{proof}






\begin{proposition}[The highest Hopf order can be made the global maximum]\label{cdp:st:C.2.UN.84}
\StatementRecord{cdp:st:C.2.UN.84}{proved}{Historical events: Claude Fable 5.2 (Anthropic), effort high, 16 November 2026: carried as background without fault found. GPT-6 Astra (OpenAI), ChatGPT Pro mode, effort Pro, S6\_09\_PAPERS, 17 November 2026: targeted proof check. Also checked in this writing. Pre-continuation writing events: GPT-6 Astra (OpenAI), ChatGPT Pro mode, effort Pro, S6\_10\_PAPERS, 18 November 2026: targeted proof check; used and rechecked in this run. Continuation metadata events: GPT-6 Astra (OpenAI), ChatGPT Pro mode, effort Pro, S6\_10\_PAPERS continuation, 19 November 2026: used and rechecked in this run; status-vocabulary reconciliation; metadata reconciliation only; no new mathematical check. \emph{Pre-continuation recorded proof scope.} The actual canonical divisor makes every normalization above the highest Hopf order ruled or rational. The precise Kulikov selection proof, including its two-section chains and cycles and largest last local branch, was opened in the original edition. The candidate section or line gives strictly positive removed order; missing-fibre insertions add only the next order and none at the current maximum. Counting all positive levels yields a fixed well-founded vector even when a removed candidate lies below the highest Hopf order. The proper graph preserves every Hopf prime and its order, and the exact rational-branch comparison permits the final maximal Hopf minimalization}{the hypotheses, notation and definitions stated here; \cref{cdp:st:C.2.UN.82}; \cref{cdp:st:C.2.UN.75}; \cite[\S7]{Kulikov1981}; \cite[\S\S4.8.5--4.8.12,4.16.2--4.16.6]{Kulikov1977}; \cref{cdp:st:C.2.UN.77}; \cite[Lemma2.19]{Kulikov1977}; \cref{cdp:eq:C.TK.4}; \cref{cdp:eq:C.TK.5}; \cref{cdp:st:C.2.UN.19}; \cref{cdp:eq:C.TK.3}; \cref{cdp:eq:C.TK.6}}

Let \(f:X\to\Delta\) be a proper degeneration of connected
complex two-tori with smooth total space and full reduced
SNC central fibre. Use its actual primitive canonical
generator and write
\[
 \operatorname{div}(\Omega)=\sum_i k_iV_i,\qquad
 \min_i k_i=0,\qquad R=\max_i k_i.
 \tag{C.TK.1}\label{cdp:eq:C.TK.1}
\]
Let \(\mathcal H\) be the set of prime components whose
smooth minimal model is Hopf, and suppose that
\[
                  h=\max_{i\in\mathcal H}k_i>0.
 \tag{C.TK.2}\label{cdp:eq:C.TK.2}
\]
There is a finite overlap-compatible modification record,
over the same disc and isomorphic on its punctured part,
whose endpoint has smooth total space and full reduced
NC fibre, such that its global maximum canonical order
is exactly \(h\). Every original Hopf-birational prime
survives with its original order, and no new such prime
appears. Every component of order \(h\) with Hopf minimal
model can, in the same finite record, be made an embedded
minimal Hopf surface. These maximal Hopfs are the
torsion-normal elliptic leaves of \ref{cdp:st:C.2.UN.82}.

At all retained global stages the canonical orders lie
between zero and the original \(R\). A whole-component
removal has positive order \(q\); its possible added
components have order \(q+1\), bounded by the current
maximum, and there are no additions when \(q\) is that
maximum. The removed \(q\) is not required to exceed
\(h\). The termination potential therefore counts every
positive order, including orders at most \(h\).

The record retains its actual centres, normalization
branches, local sigma maps and restored overlaps.
Reversing it recovers the given analytic family. On
an original compact family it gives the same actual
modification over its compact base. Its Euler changes
are the signed normal-cone terms of \ref{cdp:st:C.2.UN.75}.
\end{proposition}



\begin{proof}
\noindent\textbf{Step 1.} \par\medskip\noindent\textbf{The precise input to the component selection.}
We use the actual selection argument, rather than
applying the conclusion of Kulikov's canonical-flat
theorem to a family that does not satisfy its
Hopf-minimum hypothesis. The relevant distinction
is explicit in \cite[\S7]{Kulikov1981}: the search
for a component whose normalization is ruled with
one boundary section uses only that the current
maximal components are ruled. Section8 removes
the earlier dual-complex hypothesis and supplies
the missing-fibre blowups. Thus a lower-order Hopf
component is permitted during this selection.

\noindent\textbf{Step 2.} Here are the actual maps and the order information
from that argument which we retain. The boundary
analysis at the maximal ruled components in
\cite[\S\S4.8.5--4.8.12,4.16.2--4.16.6]{Kulikov1977}
produces either a plane with one line or a ruled
normalization with one section and vertical boundary.
The two-section cases are converted by the specified
finite type-I/type-II moves along their section
chains or cycles; a rational tree gives the complete
one-triple ruling curve. At a singular sigma chart,
the last locally blown-up branch of largest order
supplies the same one-section configuration. These
are the actual two-ruling, complete-curve and
sigma-atlas maps recorded in \ref{cdp:st:C.2.UN.75} and \ref{cdp:st:C.2.UN.77}.

\noindent\textbf{Step 3.} The ruled candidate need not be maximal. If its
order is \(q\), restriction of the common canonical
section to a generic complete ruling fibre gives
the section coefficient \(-2\), hence
\[
       k_{\mathrm{section\ parent}}=q-1.
 \tag{C.TK.3}\label{cdp:eq:C.TK.3}
\]
For a plane candidate, the line coefficient \(-3\)
gives the parent order \(q-2\). Consequently every
removed component has positive order. When a ruled
candidate has disconnected boundary, the corrected
procedure inserts the missing fibre components.
Their actual curve-blowup Jacobians give order
\(q+1\). If \(q\) is the current maximum, the
effective anticanonical one-section comparison
of \cite[Lemma2.19]{Kulikov1977} makes the boundary
connected and no insertion is made. In every case
\[
 \begin{split}
 &\text{one component of order }q>0\text{ is removed},\\
 &\text{only finitely many components of order }q+1
       \text{ may be added},\\
 &q+1\le R_{\mathrm{current}}\text{ whenever an addition occurs}.
 \end{split}
 \tag{C.TK.4}\label{cdp:eq:C.TK.4}
\]
This is an assertion about the retained degeneration
stages. It does not assign these bounds to auxiliary
common blowups in a two-ruling diagram.

\noindent\textbf{Step 4.} The local contractions are followed by the exact
restoration of the overlap, using the same composite
map. The proper-graph comparison in (C.KG.2) and
(C.KG.6) therefore gives an actual bimeromorphic
map over the disc, not merely equal central graphs.
If singular sigma charts remain, we use the
largest last locally blown-up branch and its finite
descendant cleanup from \ref{cdp:st:C.2.UN.77}. Such branches are
ruled by their actual complete exceptional fibres.
This cleanup uses no hypothesis on the other
Hopf orders and has the same positive-level updates
\eqref{cdp:eq:C.TK.4}. It ends at a smooth-total NC stage with
the labelled normalization branches retained.

\par\medskip\noindent\textbf{Why selection remains available above \(h\).}
The proper composite-graph argument in \ref{cdp:st:C.2.UN.75} proves
that every Hopf-birational prime survives every
recorded removal. A general complete rational
exceptional fibre cannot cover a point-blown-up
Hopf surface; a compact projective point exceptional
cannot map bimeromorphically onto it. The actual
canonical direct-image isomorphism (C.KG.6) preserves
its canonical order. Every added prime is ruled or
projective over the specified centre, so it is not
Hopf-birational. It follows that the set of Hopf
orders, with its original labels, is unchanged.
In particular its highest value remains \(h\).

\noindent\textbf{Step 5.} Consider a current smooth-total NC stage with
maximum \(R'>h\). On the smooth normalization
\(T\) of a maximal component, actual adjunction is
\[
       -K_T=\sum_C (R'-k_{\mathrm{other}(C)}+1)C.
 \tag{C.TK.5}\label{cdp:eq:C.TK.5}
\]
Both preimages of a self-double curve are counted;
their coefficients are one. The divisor is nonzero:
the connected support contains a component of Hopf
order \(h<R'\), so no maximal component is an
isolated component of the whole support. Every
coefficient in \eqref{cdp:eq:C.TK.5} is positive. A positive
pluricanonical section would be a holomorphic
function on the compact connected surface vanishing
on this effective divisor, which is impossible.
Thus \(\kappa(T)=-\infty\).

\noindent\textbf{Step 6.} The actual component theorem \ref{cdp:st:C.2.UN.19} makes a
nonruled component of this Kodaira dimension
Hopf-birational. No such component has order
\(R'>h\). Hence every current maximal
normalization is ruled or rational, including
the possible plane. This is exactly the input
to the selection proved above. Controlled
singular sigma charts are first handled by
their specified last-local ruled branches;
they do not require an anticanonical formula
on a singular ambient chart.

\par\medskip\noindent\textbf{The stopped finite-level induction.}
For each \(1\le j\le R\), let \(c_j\) be the
number of retained global components of order
\(j\), and order the tuples
\[
                     (c_1,c_2,\ldots,c_R)
 \tag{C.TK.6}\label{cdp:eq:C.TK.6}
\]
lexicographically, with the smaller index first.
Each whole-component block \eqref{cdp:eq:C.TK.4} strictly
decreases this tuple: it deletes a card of
order \(q\), leaves all smaller orders unchanged,
and only increases the next order. Type-I and
type-II preparatory moves preserve the card
counts. A fixed finite product of copies of
\(\mathbb N\) in this order has no infinite
strict descending sequence. The finite cleanup
blocks obey the same comparison.

\noindent\textbf{Step 7.} We continue only while the maximum exceeds
\(h\), and finish the finite cleanup belonging
to the last selected block. If the maximum
still exceeded \(h\), \eqref{cdp:eq:C.TK.5} would supply
another selection. The decreasing tuple
therefore forces a finite stop with maximum
at most \(h\). The retained Hopf prime of
order \(h\) forces equality. No order-zero
prime has been removed, by \eqref{cdp:eq:C.TK.3} and its
plane counterpart, so the transported canonical
generator remains primitive with minimum zero.

\noindent\textbf{Step 8.} Counting only the components above \(h\)
would not prove this induction. For example,
removing an order-one ruled candidate and
adding order-two components changes none of
those upper counts when \(h\ge2\). Such a
candidate can be the ruled neighbour produced
by the one-triple step. The full tuple
\eqref{cdp:eq:C.TK.6} is strictly decreased by the exact
update and is the potential actually used.

\par\medskip\noindent\textbf{The actual maximal Hopf components at the stopping stage.}
The original central components were embedded.
The record has introduced no elliptic self-double
curve on a surviving Hopf component. Type-I and
type-II changes create incidences on their
rational exceptional curves. A removed ruled
component has one horizontal boundary section;
its generic contraction cannot identify two
elliptic horizontal branches. Its other changes
are on its rational fibre curves, and its
restored overlaps retain the original elliptic
maps. Point contractions likewise only change
the indicated rational boundary configuration.

\noindent\textbf{Step 9.} At the stopping stage a Hopf normalization
of order \(h\) has the effective divisor
\eqref{cdp:eq:C.TK.5}, now with \(R'=h\). The actual
point-map induction (C.HM.3)--(C.HM.4)
shows that a rational coefficient-one boundary
is a leaf, and that two coefficient-one
boundaries never meet. A rational self-double
curve would have two coefficient-one preimages,
each with negative self-square and exactly
one boundary intersection. Their actual normal
product would consequently give
\[
                        C_1^2+C_2^2=-1,
 \tag{C.TK.7}\label{cdp:eq:C.TK.7}
\]
contradicting the two negative integer squares.
A triple point with three branches of this same
component would already give two intersecting
coefficient-one curves. Thus these maximal
Hopfs are embedded at the stopped NC stage.

\noindent\textbf{Step 10.} All the prerequisites of the last-point and
finite-cleanup proof of \ref{cdp:st:C.2.UN.77} now hold on this
actual recorded stage. That proof uses its
smooth normalizations, the labelled NC branches
elsewhere and the embedded maximal Hopfs; it
does not require the other components to be
embedded. Apply those exact moves to remove
their point defects. The retained global orders
stay at most \(h\), every Hopf prime and order
is preserved, and the endpoint maximal Hopfs
are smooth minimal surfaces. \ref{cdp:st:C.2.UN.82} then supplies
their single elliptic boundaries and exact
torsion-normal line maps.

\noindent\textbf{Step 11.} Finally the local proper graphs glue with the
recorded overlap isomorphisms, exactly as in
\ref{cdp:st:C.2.UN.75}. They give the punctured-family isomorphism
and the inverse record on its actual endpoint
neighbourhood. Annular gluing to the unchanged
outside performs the same operation on an
original compact family. The signed normal-cone
Euler formula (C.KG.3)--(C.KG.4) applies to
every retained arrow. No chosen smoothing of
the endpoint central space replaces that
actual neighbourhood.
\end{proof}






\begin{proposition}[Exact root, birational action and coarsening records for the Hopf-minimum branch]\label{cdp:st:C.2.UN.78}
\StatementRecord{cdp:st:C.2.UN.78}{proved}{Historical events: Claude Fable 5.2 (Anthropic), effort high, 16 November 2026: carried as background without fault found. GPT-6 Astra (OpenAI), ChatGPT Pro mode, effort Pro, S6\_09\_PAPERS, 17 November 2026: targeted proof check. Also checked in this writing. Pre-continuation writing events: GPT-6 Astra (OpenAI), ChatGPT Pro mode, effort Pro, S6\_10\_PAPERS, 18 November 2026: targeted proof check. Continuation metadata events: GPT-6 Astra (OpenAI), ChatGPT Pro mode, effort Pro, S6\_10\_PAPERS continuation, 19 November 2026: used and rechecked in this run; status-vocabulary reconciliation; metadata reconciliation only; no new mathematical check. \emph{Pre-continuation recorded proof scope.} The actual embedded resolution, saturated root lattice and equivariant height-one subdivisions give the semistable model. Reversing its recorded modification restores the regular finite action before coarsening and taking invariant rings, which equal the original normal rings. The Bezout determinant gives each toric Euler volume and the fixed-point quotient calculation preserves sheet permutations. The local normal-root algebras and integral strict-cover functions account for every Hopf component; projective exceptional covers add none. The radial density integral gives the resolution-independent minimum ratio. The compact meromorphic base function and its explicit unit-root coordinate retain every additional branch value}{the hypotheses, notation and definitions stated here; \cref{cdp:st:C.2.UN.75}; \cref{cdp:st:C.2.UN.39}; \cite[Theorems~1.6 and1.10]{BierstoneMilman1997}; \cref{cdp:eq:C.RK.9}; \cref{cdp:eq:C.RK.11}; \cref{cdp:eq:C.RK.16}; \cref{cdp:eq:C.RK.21}; \cref{cdp:eq:C.RK.22}; \cref{cdp:st:C.2.UN.71}; \cref{cdp:st:C.2.UN.73}}

Let $f:X\to\Delta_t$ be a proper degeneration of connected
complex two-tori from a smooth threefold, with arbitrary
central support $S=f^{-1}(0)_{\rm red}$. Take an actual
embedded resolution $\pi:Y\to X$ by smooth central
point and curve blowups, and write
\[
 \operatorname{div}(f\pi)=\sum_i m_iD_i,
 \qquad \operatorname{div}(\pi^*\Omega)=\sum_i k_iD_i,
 \qquad r_i=\frac{k_i+1}{m_i},
 \qquad N=\operatorname{lcm}_i(m_i),
\tag{C.RK.1}\label{cdp:eq:C.RK.1}
\]
where $\Omega$ generates $f_*K_X$. The $D_i$ are the
smooth components of the resolved SNC support. Form
the actual normalization and a toric semistable
resolution
\[
 \begin{gathered}
 Y_N=\operatorname{Norm}(Y\times_{\Delta_t}\Delta_s),
                \qquad t=s^N,\\
 Z\xrightarrow{\ q\ }Y_N\xrightarrow{\ \rho\ }Y,
                 \qquad G=\mu_N.
 \end{gathered}
\tag{C.RK.2}\label{cdp:eq:C.RK.2}
\]
These maps can be chosen $G$-equivariantly, with $Z$
smooth and full reduced simple-NC central fibre.
All components added by $q$ are ruled or rational.

The following is an explicit finite-cover condition
on this resolution record: every strict root-cover
component whose smooth minimal model is Hopf lies
over an index $i$ with
\[
                            r_i=\min_j r_j.
\tag{C.RK.3}\label{cdp:eq:C.RK.3}
\]
For exactly the records satisfying \eqref{cdp:eq:C.RK.3}, the
semistable root $Z$ admits the finite compatible
canonical-flat reduction of \ref{cdp:st:C.2.UN.75}. Denote its
actual comparison by $\phi:Z\dashrightarrow X_0$.
The flat endpoint has one of the three complete
core types in \ref{cdp:st:C.2.UN.75}, and it carries the actual
finite \emph{birational} action
\[
                 \gamma_0=\phi\gamma\phi^{-1}
                 \quad (\gamma\in G).
\tag{C.RK.4}\label{cdp:eq:C.RK.4}
\]
Its group laws hold as equalities of meromorphic
maps; it acts over $s\mapsto\zeta s$ and is regular
on the punctured family. Retain the full record
$\phi$, the toric coarsening $q$, and $\pi$.
Their inverse and quotient operations reconstruct
the given original degeneration exactly:
\[
 X_0\ \dashrightarrow\ Z\ \xrightarrow{q}\ Y_N
       \ \longrightarrow\ Y_N/G=Y
       \ \xrightarrow{\pi}\ X.
\tag{C.RK.5}\label{cdp:eq:C.RK.5}
\]
In particular the inverse record is applied
\emph{before} the finite quotient. No quotient of
an undefined central birational action is used.
Conversely the same operations on any actual
flat-core record with these retained regularization,
root-coarsening and smooth blowdown maps produce
a torus degeneration in the stated branch.
This gives complete finite generation for this
explicit resolution-and-root branch, without
an equivariant canonical-flat reduction assertion.

The Euler correction for the root is explicit.
For a nonempty index set $I$ with nonempty stratum
\[
 D_I^\circ=
    \left(\bigcap_{i\in I}D_i\right)
           \setminus\bigcup_{j\notin I}D_j,
 \qquad d_I=\gcd_{i\in I}m_i,
 \qquad \nu_I=\frac{N^{|I|-1}d_I}{\prod_{i\in I}m_i},
\tag{C.RK.6}\label{cdp:eq:C.RK.6}
\]
one has $1\le|I|\le3$ and $\nu_I$ is a positive
integer. Put
\[
 \begin{split}
 R_N&=\sum_I(d_I\nu_I-1)e_c(D_I^\circ),\\
 \Delta_\pi&=2p_\pi+
                   \sum_\alpha(2-2g(C_\alpha)),
 \end{split}
\tag{C.RK.7}\label{cdp:eq:C.RK.7}
\]
where $p_\pi$ and $C_\alpha$ are the actual point
and connected curve centres of $\pi$. Let
$\Delta_K$ be the signed normal-cone defect of
the forward \ref{cdp:st:C.2.UN.75} record from $Z$ to $X_0$.
Then the exact original-support formula is
\[
                e(S)=e(S_0)-\Delta_K-R_N-\Delta_\pi.
\tag{C.RK.8}\label{cdp:eq:C.RK.8}
\]
Thus zero Euler number is equivalent to
$\Delta_K+R_N+\Delta_\pi=e(S_0)$, where the
right side is zero for the smooth or marked
cycle core, and $2[N_0:\Gamma_0]$ for the
toric core. The symbols $N_0,\Gamma_0$ here
are its two-dimensional lattices, separate
from the root exponent $N$ in \eqref{cdp:eq:C.RK.1}.
Formula \eqref{cdp:eq:C.RK.8} agrees with the universal
original-support formula of \ref{cdp:st:C.2.UN.39} by the
actual root and blowup maps.

The test \eqref{cdp:eq:C.RK.3} is in fact intrinsic to
the original degeneration. Write its original
full fibre and generator orders as
$\sum_\alpha n_\alpha S_\alpha$ and
$\sum_\alpha a_\alpha S_\alpha$.
On the normalization $T_\alpha$ of $S_\alpha$
form the finite normal cover $C_\alpha$
by the following actual local equations
and transition maps:
\[
 \begin{gathered}
 f=g_j^{\,n_\alpha}u_j,\qquad
 w_j^{\,n_\alpha}=u_j|_{T_\alpha},\\
 w_k=(g_j/g_k)|_{T_\alpha}\,w_j .
 \end{gathered}
\tag{C.RK.18}\label{cdp:eq:C.RK.18}
\]
Here $g_j$ is a local defining equation
of the Cartier component $S_\alpha$;
the polynomial covers are normalized
after gluing. Define
\[
 \tau=\sup\left\{c\in\R:
       |f|^{-2c}|\Omega|^2\text{ is locally integrable
                              near }S\right\}.
\tag{C.RK.19}\label{cdp:eq:C.RK.19}
\]
The notation $|\Omega|^2$ denotes its
positive top-degree density, so this
definition requires no metric choice.
Then \eqref{cdp:eq:C.RK.3} is equivalent to
\[
 \frac{a_\alpha+1}{n_\alpha}=\tau
 \quad\text{for every }\alpha\text{ for which a connected
 component of }C_\alpha\text{ has Hopf minimal model}.
\tag{C.RK.20}\label{cdp:eq:C.RK.20}
\]
Consequently its truth is independent
of the embedded resolution, its
semistable triangulation and the
chosen root exponent divisible by
all resolved multiplicities.

For an original compact family the record can
be made on a finite compact cover of its base,
and \eqref{cdp:eq:C.RK.5} recovers the original compact
family. Reverse generation on compact actual
core records with the displayed finite action
and coarsening likewise remains compact.
The original endpoint neighbourhood, not only
its central log space, is retained throughout.
\end{proposition}



\begin{proof}
\noindent\textbf{Step 1.} \par\medskip\noindent\textbf{Embedded resolution and the exact normalized root charts.}
Use analytic embedded resolution by central
smooth centres, as in \ref{cdp:st:C.2.UN.39} and
\cite[Theorems~1.6 and1.10]{BierstoneMilman1997}.
Divisor centres give identity blowups and are
omitted. The remaining centres in a smooth
threefold are points and compact curves.
The actual pullback of $\Omega$ is the
canonical generator on $Y$: apply the
Hartogs direct-image comparison (C.KG.6)
to the ordinary blowup maps.

\noindent\textbf{Step 2.} Near a stratum with $r$ incident components,
absorb the unit in the fibre equation into
a coordinate. The exact local equation and
volume form have the shape
\[
 t=x_1^{m_1}\cdots x_r^{m_r},\qquad
 \pi^*\Omega=u\,x_1^{k_1}\cdots x_r^{k_r}
                         dx_1\wedge dx_2\wedge dx_3,
                         \quad u\text{ a unit}.
\tag{C.RK.9}\label{cdp:eq:C.RK.9}
\]
The normalization of $t=s^N$ has $d_I$
local branches. Each is the affine toric
chart with cocharacter lattice
\[
 \mathcal L_I=\left\{(v,h)\in\Z^r\oplus\Z:
                         \sum_i m_iv_i=Nh\right\}
\tag{C.RK.10}\label{cdp:eq:C.RK.10}
\]
and cone generated by the height-one rays
$((N/m_i)e_i,1)$. The omitted $3-r$
coordinates are unchanged. Local roots of
the unit, if not absorbed in \eqref{cdp:eq:C.RK.9},
give the same normalization maps and
the actual finite transition characters.
The group $G$ permutes the $d_I$ branches
and acts torically on each stabilised
branch, over $s\mapsto\zeta s$.

\noindent\textbf{Step 3.} The height-one slice has dimension at most
two. Subdivide its edges at every lattice
point, and triangulate each two-dimensional
lattice triangle using all its lattice
points. An empty lattice triangle is
unimodular, by Pick's formula. Thus these
subdivisions give smooth toric charts.
They agree on their subdivided edges.
The refinement can be obtained by
successively inserting the height-one
lattice points as star subdivisions.
Each step is the normalized blowup
of its actual monomial ideal. On an
overlap, replacing a boundary equation
by a holomorphic unit times that
equation multiplies every monomial
generator by a unit and leaves this
ideal unchanged. Thus the normalized
blowup maps glue through the original
tangent-dependent transition units;
those transitions have not been
replaced by constant torus maps.
The stabilizer of an old coarse cone fixes
each old ray: the original SNC components
are individually labelled, and the group
acts over their identity map in $Y$.
Its action on this lattice slice is
therefore the identity. Choose the
triangulations orbitwise. The actual
toroidal transition maps, which preserve
these labelled rays, glue the subdivisions
to a proper $G$-equivariant map $q$.

\noindent\textbf{Step 4.} Every new ray still has height one,
so $s$ has order one on every central
component of $Z$. The resulting toric
boundary is simple NC with smooth
components. For the embedded assertion
retain every old component and stratum
label as well as each subdivision ray.
Transition units fix those labels and
the ray vector; distinct rays in one
old cone are never identified. The
local components in a unimodular
chart therefore belong to distinct
global components, and their smooth
star charts glue without a global
self-incidence. Over a double curve the
added surfaces are projective-line
bundles over finite covers of that
curve, with their indicated modifications;
over a triple point they are rational
toric surfaces. None has Hopf minimal
model. The normalized fibre product
over $\Delta_s^*$ is connected because
its compact torus fibres and its base
are connected. The proper Stein factor
has degree one there and hence on the
normal disc. Thus $Z\to\Delta_s$ has
connected full central fibre as well.

\par\medskip\noindent\textbf{Canonical ratios and the explicit minimum test.}
At a primitive height-one ray $(v,1)$,
the toric Jacobian of \eqref{cdp:eq:C.RK.9} gives
\[
 \operatorname{ord}_{(v,1)}(q^*\rho^*\pi^*\Omega)
                 =\sum_i(k_i+1)v_i-1.
\tag{C.RK.11}\label{cdp:eq:C.RK.11}
\]
For a strict ray this is
$N(k_i+1)/m_i-1$. Any height-one
point is a convex combination of the
old vertices, so \eqref{cdp:eq:C.RK.11} is the
same affine combination of their
orders. The minimum among all rays is
therefore
\[
                  a_{\min}=N\min_i r_i-1.
\tag{C.RK.12}\label{cdp:eq:C.RK.12}
\]
It is an integer, since a minimum is
attained at an old vertex. The form
\[
            \Omega_Z=s^{-a_{\min}}
                             q^*\rho^*\pi^*\Omega
\tag{C.RK.13}\label{cdp:eq:C.RK.13}
\]
is holomorphic and generates the
canonical direct image on the reduced
semistable family. Its component
orders have minimum zero. Since all
new components are ruled or rational,
\ref{cdp:st:C.2.UN.75}'s Hopf-minimum condition is
exactly \eqref{cdp:eq:C.RK.3} for the strict covers.
This proves both directions of the
test for the chosen resolution record.
It also retains the actual generator
character
$\gamma^*\Omega_Z=\zeta^{-a_{\min}}\Omega_Z$.

\par\medskip\noindent\textbf{The actual finite birational action and its quotient order.}
Apply \ref{cdp:st:C.2.UN.75} to $Z$. Its comparison
$\phi$ is a bimeromorphic map over
$\Delta_s$, with a specified inverse
record and an actual punctured
isomorphism. Formula \eqref{cdp:eq:C.RK.4}
therefore defines meromorphic maps
on $X_0$. Their products satisfy
the group laws on the dense
punctured set and consequently
as meromorphic maps. They have
order $N$ on the total family,
as seen on the base. Choosing
the endpoint volume by
$\phi^*\Omega_0=\Omega_Z$ also
gives exactly the character in
the preceding paragraph under
this birational action.

\noindent\textbf{Step 5.} To reconstruct, first reverse
the retained $\phi$-record.
The action is then the original
regular $G$-action on $Z$.
Apply its equivariant toric
coarsening $q$ to get the actual
normal space $Y_N$. Its quotient
is literally $Y$. Indeed on a
normal local ring $A$ of $Y$,
the root ring is the integral
closure $B$ of
$A[s]/(s^N-f\pi)$. The invariants
of its total ring of fractions are
$\operatorname{Frac}(A)$.
An invariant element of $B$ is
integral over $A$ and lies in
that fraction field, hence lies
in $A$. Thus
\[
                              B^G=A.
\tag{C.RK.14}\label{cdp:eq:C.RK.14}
\]
This argument applies to the
normalization on the connected
total space; the local split
branches are permuted by $G$
and give the same invariant
ring equality. The retained
ordinary blowdowns $\pi$ now
recover $X$. These exact maps
prove \eqref{cdp:eq:C.RK.5} and its converse.
An equivariant replacement of
$\phi$ has not been used.

\par\medskip\noindent\textbf{The root correction from actual stratum fibres.}
Over $D_I^\circ$ the reduced
central support of $Y_N$ is an
\emph{actual} finite etale cover
of degree $d_I$. The local root
equation in \eqref{cdp:eq:C.RK.9} proves the
degree; unit-root transition
maps glue the cover and may
permute its sheets. At a point
of this cover the toric fibre
of $q$ has Euler number $\nu_I$.
To see this, use its torus-orbit
stratification. All positive-
dimensional torus orbits have
Euler number zero. The zero-
dimensional orbits are exactly
the maximal unimodular cones
in the subdivision, so their
number is the normalized
volume of the height-one
simplex.

\noindent\textbf{Step 6.} Here is its exact lattice
volume. Set $a_i=m_i/d_I$ and
choose an integral vector $b$
with $\sum_i a_i b_i=1$.
The difference lattice of
the slice is
$\ker(a:\Z^r\to\Z)$.
Appending $b$ to a basis of
this kernel gives a unimodular
basis of $\Z^r$. Consequently
the normalized volume is the
absolute determinant obtained
by appending $b$ to the $r-1$
vertex differences. Expansion
gives
\[
 \left|\det\left[
 \frac{N}{m_1}e_1-\frac{N}{m_r}e_r,
 \ldots,
 \frac{N}{m_{r-1}}e_{r-1}-\frac{N}{m_r}e_r,
 b\right]\right|
               =\frac{N^{r-1}d_I}{\prod_i m_i}.
\tag{C.RK.15}\label{cdp:eq:C.RK.15}
\]
For $r=1$ both sides are one.
This proves the formula and
integrality of $\nu_I$ without
a choice of triangulation.

\noindent\textbf{Step 7.} Proper stratified Euler
pushforward now gives
\[
 \begin{split}
 e(S_Z)&=\sum_I d_I\nu_I e_c(D_I^\circ),\\
 e(S_Y)&=\sum_I e_c(D_I^\circ),\\
 e(S_Z)-e(S_Y)&=R_N.
 \end{split}
\tag{C.RK.16}\label{cdp:eq:C.RK.16}
\]
This also checks the finite
quotient directly. A deck
element $\gamma^j$ fixes the
$d_I$-sheeted central cover
over this stratum exactly
when $d_I\mid j$. Therefore
\[
 \begin{split}
 e(S_{Y_N}^{\gamma^j})
       &=\sum_{d_I\mid j}d_I e_c(D_I^\circ),\\
 \frac1N\sum_{j=0}^{N-1}
          e(S_{Y_N}^{\gamma^j})
       &=\sum_I e_c(D_I^\circ)=e(S_Y).
 \end{split}
\tag{C.RK.17}\label{cdp:eq:C.RK.17}
\]
All these equalities hold for
the actual covers even when
they are globally connected
and their sheet labels have
nontrivial monodromy.

\noindent\textbf{Step 8.} The original resolution has
$e(S_Y)-e(S)=\Delta_\pi$ by
its ordinary point and curve
maps in \ref{cdp:st:C.2.UN.39}. The forward
\ref{cdp:st:C.2.UN.75} record has
$e(S_0)-e(S_Z)=\Delta_K$.
Combine these two equalities
with \eqref{cdp:eq:C.RK.16} to obtain
\eqref{cdp:eq:C.RK.8}, including its zero
criterion. Each summand has
been compared through an
actual proper map or a
retained overlap-compatible
record.

\par\medskip\noindent\textbf{Intrinsic covers and the intrinsic minimum.}
The component $S_\alpha$ is Cartier
in the original smooth threefold.
Its local equations satisfy
$g_k=v_{kj}g_j$ for holomorphic
units, and therefore
$u_k=v_{kj}^{-n_\alpha}u_j$.
This is precisely the transition
in \eqref{cdp:eq:C.RK.18}, including its cocycle
identity. The corresponding monic
finite algebras glue before taking
their finite normalizations.
On the open smooth component
away from other components this
is the $\mu_{n_\alpha}$-torsor
of roots of the normal-line
trivialization supplied by $f$.
No global meromorphic frame for
that normal line is chosen.

\noindent\textbf{Step 9.} For any root exponent divisible
by $n_\alpha$, the exact map from
the strict central root cover is
\[
                         w_j=\frac{s^{N/n_\alpha}}{g_j}.
\tag{C.RK.21}\label{cdp:eq:C.RK.21}
\]
At its generic point numerator
and denominator have the same
order, and their quotient satisfies
\eqref{cdp:eq:C.RK.18}. On a smooth strict
component of $Z$ this meromorphic
quotient is integral over its
holomorphic local ring, by that
monic equation. Normality makes
it holomorphic. The map therefore
extends to a proper map to the
glued finite cover, and factors
through its normalization. It
has generic degree one onto a
connected component of $C_\alpha$.
Every such component is obtained,
since on the dense open both
covers have the same complete
set of $n_\alpha$ roots.
Thus the strict surfaces are
actual birational resolutions
of these fixed covers for every
choice of embedded resolution
and divisible root exponent.

\noindent\textbf{Step 10.} Every exceptional surface of
the original embedded resolution
is projective: initially it is
$\PP^2$ or a projective-line
bundle over a compact curve,
and its subsequent transforms
are projective modifications.
The ideals on these projective
surfaces are algebraic by GAGA.
A finite normal analytic cover
of a projective surface is also
projective: its finite coherent
algebra and multiplication
algebraize by GAGA, and the
finite pullback of an ample
line is ample. Its surface
resolutions remain bimeromorphic
to projective surfaces. Such
root components cannot have
Hopf minimal model. The new
semistable components have
already been shown ruled or
rational. Hence \eqref{cdp:eq:C.RK.21}
accounts for every Hopf root
component, with no exceptional
ones omitted.

\noindent\textbf{Step 11.} Finally apply the literal change
of variables under $\pi$ to the
density in \eqref{cdp:eq:C.RK.19}, outside
the exceptional sets of measure
zero. Near a resolved crossing
it is a bounded positive unit
times the product of the radial
densities
$|x_i|^{2(k_i-cm_i)}\,dA(x_i)$,
with smooth tangent density.
On finitely many relatively
compact charts about the compact
central fibre, the integral is
finite exactly when
$k_i-cm_i>-1$ for every $i$.
At equality it diverges already
at a general point of that
divisor. Thus
\[
                             \tau=\min_i\frac{k_i+1}{m_i}.
\tag{C.RK.22}\label{cdp:eq:C.RK.22}
\]
This is the original density's
integrability threshold, so it
is independent of resolution.
The strict transform of an
original component retains
its original pair
$(a_\alpha,n_\alpha)$.
Together with \eqref{cdp:eq:C.RK.21} and the
exceptional-component comparison,
\eqref{cdp:eq:C.RK.22} proves exactly
\eqref{cdp:eq:C.RK.20} and the stated
choice independence.

\par\medskip\noindent\textbf{Compact scope.}
For a compact base $B$ choose
a meromorphic function with
a simple zero at the selected
point; Riemann--Roch supplies
one with its poles away from
that point. Normalize the
cyclic cover $s^N=h$ of $B$.
It is a compact curve with
the exact local parameter
map used in \eqref{cdp:eq:C.RK.2}.
Additional branch values
are retained in this base
record. Perform the normal
base change and equivariant
resolution globally, and
the \ref{cdp:st:C.2.UN.75} replacement near
the chosen central fibre.
Its actual annular map
glues to the unchanged
outside as in \ref{cdp:st:C.2.UN.75}. Conjugate
the global finite action
through that comparison;
it is a finite birational
action on this compact
family. Reversing the
record regularizes it,
and the compact proper
coarsening and quotient
in \eqref{cdp:eq:C.RK.5} recover the
original compact family.
The converse uses those
same proper maps on an
actual compact core record.
No replacement of a full
neighbourhood follows only
from the central-log-space
classification of \ref{cdp:st:C.2.UN.71}
or its occurrence theorem
\ref{cdp:st:C.2.UN.73}.
\end{proof}






\begin{proposition}[All no-triple fibres are ordinary vertical blowups of the classified elliptic cores]\label{cdp:st:C.2.UN.72}
\StatementRecord{cdp:st:C.2.UN.72}{proved}{Historical events: Claude Fable 5.2 (Anthropic), effort high, 16 November 2026: carried as background without fault found. GPT-6 Astra (OpenAI), ChatGPT Pro mode, effort Pro, S6\_09\_PAPERS, 17 November 2026: targeted proof check. Also checked in this writing. Pre-continuation writing events: GPT-6 Astra (OpenAI), ChatGPT Pro mode, effort Pro, S6\_10\_PAPERS, 18 November 2026: targeted proof check. Continuation metadata events: GPT-6 Astra (OpenAI), ChatGPT Pro mode, effort Pro, S6\_10\_PAPERS continuation, 19 November 2026: used and rechecked in this run; status-vocabulary reconciliation; metadata reconciliation only; no new mathematical check. \emph{Pre-continuation recorded proof scope.} After the actual rational-edge contractions, a higher-genus height ascent cannot pass through an elliptic edge: the torsion-normal parent map would make the existing higher-genus curve vertical over a smooth elliptic support. It ends in an actual positive-genus section-leaf contraction. The nearby-stratum Euler formula then forces all remaining elliptic-boundary components to be minimal. Their exact earlier classification gives the four actual core neighbourhoods. Reversing the retained ordinary blowup charts recovers the original family with its full normal extensions; each original nonelliptic edge corresponds to exactly one removed centre. The resulting intrinsic Euler balance follows from those literal replacement maps}{the hypotheses, notation and definitions stated here; \cref{cdp:st:C.2.UN.65}; \cref{cdp:st:C.2.UN.57}; \cref{cdp:st:C.2.UN.61}; \cref{cdp:st:C.2.UN.62}; \cref{cdp:st:C.2.UN.69}; \cref{cdp:st:C.2.UN.70}; \cref{cdp:eq:C.NT.5}; \cref{cdp:eq:C.NT.6}; \cref{cdp:st:C.2.UN.64}; \cref{cdp:st:C.2.UN.5}; \cref{cdp:st:C.2.UN.7}; \cref{cdp:st:C.2.UN.39}; \cref{cdp:st:C.2.UN.66}}

Let $f:U\to\Delta$ be a proper degeneration of connected complex
two-tori from a smooth threefold. Suppose that its full central
fibre $S=\sum_vT_v$ is reduced and simple normal crossing and
has no triple points. The components may be nonminimal.

There is a finite sequence of actual ordinary holomorphic
blowdowns over the disc, each the inverse of a point blowup
or a blowup of a smooth compact curve of positive genus in
the smooth locus of the current central fibre, which ends
in an actual degeneration with one of these marked cores:
\begin{enumerate}
\item a smooth complex two-torus;
\item a marked diagonal Hopf and degree-zero elliptic ruled
cycle with the exact normal maps and lattice product $M=I$
of \ref{cdp:st:C.2.UN.65} and \ref{cdp:st:C.2.UN.57};
\item a rooted elliptic tree satisfying the full criterion
of \ref{cdp:st:C.2.UN.61}, including its original analytic torsors,
secondary framings and zero-Chern inverse principal torsor;
\item a unicyclic support of genus-zero fixed-Jacobian
identity-monodromy elliptic quasi-bundles and nonlinear
Hopf leaves satisfying the full criterion of \ref{cdp:st:C.2.UN.62}.
\end{enumerate}
The original neighbourhood is recovered by reversing the
recorded ordinary blowups on this actual core degeneration.
Each centre avoids every double curve then present.
For every curve centre $B$ the record retains the exact
normal extension and marked exceptional pair
\[
 \begin{gathered}
 0\longrightarrow K_B\longrightarrow V=N_{B/U'}
                 \xrightarrow{\,df'\,}\OO_B\longrightarrow0,\\
 (T,C)=\bigl(\PP_B(V),\PP_B(K_B)\bigr),\qquad
 N_{T/U}\simeq\OO_{\PP_B(V)}(-1).
 \end{gathered}
\tag{C.NT.1}\label{cdp:eq:C.NT.1}
\]
Here $\PP_B(V)$ parametrizes lines. The actual endpoint
identifications in \eqref{cdp:eq:C.NT.1} are part of the record.

Conversely, take any actual torus degeneration with one of
these cores and perform any finite sequence of ordinary
point blowups and ordinary blowups of smooth compact
positive-genus curves in the smooth locus of its current
central fibre, avoiding all current double curves. The
result is a full reduced SNC torus fibre without triple
points. This gives the complete reconstruction in the
stated class. It operates on actual neighbourhoods; a
different realization of the same abstract marked core
is not substituted for the given neighbourhood.
For a compact total space over a compact curve all the
recorded contractions and blowups have the same compact
version.

If $a$ is the number of point blowups and $B_1,\ldots,B_b$
are the recorded curve centres, then
\[
                 e(S)=2a-2\sum_{j=1}^b\bigl(g(B_j)-1\bigr).
\tag{C.NT.2}\label{cdp:eq:C.NT.2}
\]
Thus the complete zero-Euler reconstruction is obtained
exactly by the balance
\[
                       a=\sum_{j=1}^b\bigl(g(B_j)-1\bigr).
\tag{C.NT.3}\label{cdp:eq:C.NT.3}
\]
Equivalently, let $r_0$ count the connected rational
double curves of the original support and let
$\mathcal E_{\ge2}$ be its connected double curves of
genus at least two. These intrinsic data satisfy
\[
 \begin{split}
 a&=r_0,\qquad
 \{B_j:g(B_j)\ge2\}
       \ \simeq\ \{C_e:e\in\mathcal E_{\ge2}\},\\
 e(S)&=2r_0-2\sum_{e\in\mathcal E_{\ge2}}
                                      \bigl(g(C_e)-1\bigr).
 \end{split}
\tag{C.NT.4}\label{cdp:eq:C.NT.4}
\]
The bijection is the one induced by the actual contractions
on the double curves. Every nonelliptic edge is a terminal
edge at the stage at which it is removed; in particular
every cycle of the original incidence graph consists of
elliptic double curves.
\end{proposition}



\begin{proof}
\noindent\textbf{Step 1.} \par\medskip\noindent\textbf{First remove every rational double curve by actual point contractions.}
Apply \ref{cdp:st:C.2.UN.69} to the given degeneration. Its strict-height
ascent ends in actual plane-line leaves, and each ordinary
blowdown removes one rational double curve and no other
double curve. A rational ruled-section leaf was excluded
there by transporting its square-minus-two centre into
the actual nearby torus. Thus this first stage makes
exactly $r_0$ point contractions. It leaves a proper
smooth-total torus degeneration, still full reduced
SNC without triple points, and with no rational double
curve. The maps are isomorphisms on neighbourhoods of
every surviving double curve.

\par\medskip\noindent\textbf{An increasing-height path cannot change from higher genus to elliptic.}
Use the actual canonical orders of the current family.
For every remaining double curve the exact adjunction
map is (C.NM.2). In particular a curve of genus at least
two has endpoints of different heights. Orient that
edge toward its higher endpoint and call that component
$P$.

\noindent\textbf{Step 2.} Suppose $P$ has a neighbour $Q$ of strictly greater
height, and write $J=P\cap Q$ for the corresponding
connected boundary. If $J$ were elliptic, (C.NM.2)
would give the actual isomorphism
\[
      N_{J/P}^{\,k_Q-k_P}\simeq K_J\simeq\OO_J,
        \qquad a_J=k_Q-k_P-1\ge0 .
\tag{C.NT.5}\label{cdp:eq:C.NT.5}
\]
Its normal is therefore torsion, with the actual
canonical coefficient required by \ref{cdp:st:C.2.UN.70}. All boundary
curves on $P$ are disjoint. That theorem supplies
the actual map
\[
        \pi:P\longrightarrow Z,\qquad
        q_Z:Z\longrightarrow B,
\tag{C.NT.6}\label{cdp:eq:C.NT.6}
\]
which is a sequence of point blowdowns followed by a
relatively minimal elliptic quasi-bundle with only
smooth multiple supports; every original boundary
is vertical for $q_Z\pi$.

\noindent\textbf{Step 3.} No smooth boundary of genus at least two can be
vertical for this map. An exceptional prime of the
point-blowdown map is rational. A nonexceptional
boundary has the same function field as its image
under this birational morphism. A vertical image is
one of the smooth elliptic fibre supports of $q_Z$,
so its smooth strict transform has genus one.
Both alternatives contradict the higher-genus
boundary already on $P$. Hence $J$ is not elliptic.
There are no rational double curves at this stage,
so every strictly higher neighbour is reached by
another higher-genus double curve.

\noindent\textbf{Step 4.} This uses the actual maps \eqref{cdp:eq:C.NT.5}--\eqref{cdp:eq:C.NT.6}. It does
not contract any arbitrary surface exceptional curve
inside the threefold or assume an elliptic map on an
ambient neighbourhood.

\par\medskip\noindent\textbf{Every higher-genus edge leads to an actual section-leaf contraction.}
Continue along strictly increasing heights whenever
possible. There are finitely many components, so the
path terminates at a local maximum $T$ having a
strictly lower neighbour along a higher-genus curve.
The maximum theorem \ref{cdp:st:C.2.UN.69} applies without any
minimality premise. Its equal-height alternative is
impossible because of that lower neighbour. Its Hopf
leaf alternative has an elliptic boundary, and its
plane-line alternative has a rational boundary.
Thus $T$ is exactly an unblown projective-line
bundle whose single boundary is a section of genus
at least two. The actual \ref{cdp:st:C.2.UN.64} map contracts it
to the smooth curve identified with that section.

\noindent\textbf{Step 5.} The parent surface is unchanged by this curve
contraction, since its restriction is the blowup
of a Cartier curve. Every other component misses
the exceptional divisor. Consequently the new
central fibre is still full reduced SNC without
triple points; no rational double curve has appeared,
and every surviving boundary neighbourhood is
unchanged. Canonical order differences on surviving
components are retained by the ordinary blowup
Jacobian, as in \ref{cdp:st:C.2.UN.64}.

\noindent\textbf{Step 6.} Repeat this finite procedure. Each step removes
one component and precisely one higher-genus double
curve. It therefore removes all such curves and
gives the exact bijection in the first line of
\eqref{cdp:eq:C.NT.4}. Each curve centre and its normal extension
are exactly \eqref{cdp:eq:C.NT.1}. At this point every remaining
double curve is elliptic.

\par\medskip\noindent\textbf{The remaining support has zero Euler number and is already minimal.}
Denote this remaining full support by $S_{\rm ell}$.
If it is smooth, the smooth-support comparison \ref{cdp:st:C.2.UN.5}
identifies it with a two-torus. Suppose it is
reducible. The actual nearby-cycle Euler formula
for a reduced SNC fibre is
\[
        e(F)=\sum_v e(T_v^\circ),\qquad
        T_v^\circ=T_v\setminus
                          \bigcup_{w\ne v}(T_v\cap T_w).
\tag{C.NT.7}\label{cdp:eq:C.NT.7}
\]
This is the proper analytic formula of A'Campo used
in \ref{cdp:st:C.2.UN.7} and \ref{cdp:st:C.2.UN.39}. Its local comparison is especially
simple here: over a smooth component stratum there
is one nearby sheet, while a double stratum has
the annular local fibre $xy=t$, of Euler number zero.
There are no triple strata. Compact triangulations
and the actual tubular comparison give the displayed
additivity, without a Kähler assumption.

\noindent\textbf{Step 7.} Writing $C_e$ for the double curves, all of which
are elliptic, gives
\[
 \begin{split}
 0=e(F)&=\sum_v e(T_v)-2\sum_e e(C_e)
                            =\sum_v e(T_v),\\
 e(S_{\rm ell})&=\sum_v e(T_v)-\sum_e e(C_e)=0.
 \end{split}
\tag{C.NT.8}\label{cdp:eq:C.NT.8}
\]
Now use the all-elliptic extension (C.MR.16), whose
proof applies to these actual possibly nonminimal
components. Its nonnegative minimal-model Euler
terms and positive point-blowup terms force every
component of $S_{\rm ell}$ to be minimal and
nonrational. Thus the actual support has reached
the precise scope of \ref{cdp:st:C.2.UN.66}, rather than merely a
support with a prescribed numerical Euler value.

\ref{cdp:st:C.2.UN.66} now makes its further ordinary contractions.
All their curve centres are elliptic: the boundary
curves present are elliptic and no rational or
higher-genus edge is introduced. The remaining
actual support is one of the four cores in the
statement. Its full torsor, lattice, secondary and
normal identifications are those of \ref{cdp:st:C.2.UN.61}, \ref{cdp:st:C.2.UN.62}
or \ref{cdp:st:C.2.UN.65}, as appropriate. This establishes the
endpoint before any reconstruction claim.

\par\medskip\noindent\textbf{Exact reversal and the converse construction.}
Every preceding ambient map is one of the ordinary
holomorphic blowdowns of \ref{cdp:st:C.2.UN.64}. Reverse the finite
record on the actual endpoint family. A point
step replaces a smooth point of one central
component by its exceptional plane. A curve
step has the exact normal sequence \eqref{cdp:eq:C.NT.1}.
These comparisons are the actual local blowup
charts and glue with the original isomorphism
off each exceptional divisor. They recover the
given $U$ over the disc, not just the birational
classes of its components.

\noindent\textbf{Step 8.} For clarity, the same charts prove the converse
directly. At a curve centre choose coordinates
$(z,w,t)$ with the centre $(z,t)=0$ and $f=t$.
The blowup charts $t=zs$ and $z=tr$ have fibre
equations $zs=0$ and $t=0$. At a point centre
blow up the ideal $(x,y,t)$; the chart
$y=xu,\ t=xv$ has fibre equation $xv=0$,
and the other two charts give the same ordinary
normal-crossing comparison. The centre avoids
all pre-existing double curves, so no third
component meets either exceptional pair.
The full exceptional multiplicity is one, since
$t$ has order one along the centre.
Thus each step retains a smooth total space
and a full reduced SNC fibre without triple points.
Its punctured family is unchanged. Properness
is retained by the proper blowup map.

\noindent\textbf{Step 9.} On a curve centre $B$ in the smooth central
locus the actual canonical divisor is supported
on boundaries away from $B$. Its restriction
trivializes $K_{P'}|_B$. Adjunction therefore gives
$N_{B/P'}\simeq K_B$, and the normal sequence
with the quotient $df'$ is precisely \eqref{cdp:eq:C.NT.1}.
No curve of genus zero can be such a centre:
the compact submersion transport (C.NM.10)
would again put its square-minus-two sphere
into the nearby torus. Points and smooth curves
of positive genus are exactly the centre
types used in the reconstruction.

\noindent\textbf{Step 10.} When the original total space is compact, the
actual local blowdowns glue to the identity
off their exceptional divisors by \ref{cdp:st:C.2.UN.64}.
Their targets remain compact and the map to
the compact base descends. Conversely every
listed compact centre has its ordinary global
blowup. This proves the compact version without
assuming an isomorphism between unrelated
realizations of a marked core.

\par\medskip\noindent\textbf{Euler balance and the intrinsic double-curve record.}
All four cores have Euler number zero. A point
blowup replaces one point by $\PP^2$ and adds two
to the Euler number of the support. A curve
blowup replaces $B$ by its projective-line bundle
of Euler number $2e(B)$, and so adds
$e(B)=2-2g(B)$. Additivity along the actual
blowup maps gives \eqref{cdp:eq:C.NT.2}, and its zero
condition is exactly \eqref{cdp:eq:C.NT.3}.

\noindent\textbf{Step 11.} Every such blowup attaches a single new double
edge: a rational line for a point, or the actual
curve $B$ for a curve centre. It leaves all
pre-existing double curves and their
neighbourhoods unchanged. The cores have only
elliptic double curves. The rational-edge
removal stage made exactly $r_0$ point steps,
and the higher-genus removal stage identifies
each of its centres with exactly the removed
original higher-genus double curve. These
actual maps prove all of \eqref{cdp:eq:C.NT.4} and its
independence of the chosen contraction order.
They also show that every nonelliptic edge
is removed as a leaf, which proves the
assertion about cycles. This completes the
classification in the full reduced
no-triple class.
\end{proof}






\begin{proposition}[Every reduced SNC torus degeneration has a reversible finite reduction to a classified core]\label{cdp:st:C.2.UN.85}
\StatementRecord{cdp:st:C.2.UN.85}{proved}{Historical events: Claude Fable 5.2 (Anthropic), effort high, 16 November 2026: carried as background without fault found. GPT-6 Astra (OpenAI), ChatGPT Pro mode, effort Pro, S6\_09\_PAPERS, 17 November 2026: targeted proof check. Also checked in this writing. Pre-continuation writing events: GPT-6 Astra (OpenAI), ChatGPT Pro mode, effort Pro, S6\_10\_PAPERS, 18 November 2026: targeted proof check; used and rechecked in this run. Continuation metadata events: GPT-6 Astra (OpenAI), ChatGPT Pro mode, effort Pro, S6\_10\_PAPERS continuation, 19 November 2026: used and rechecked in this run; status-vocabulary reconciliation; metadata reconciliation only; no new mathematical check. \emph{Pre-continuation recorded proof scope.} The initial canonical-order alternatives give either the actual flat endpoint or a preserved maximal minimal Hopf neighbourhood. Each newly reached torsion-normal interface records an actual point map. The adjacent shifted canonical coefficients stay coprime under smooth and node point blowups, excluding the problematic self-double last curve. Receiving-square bounds preserve the already recorded models in a direct two-ruling step. Whole removal blocks may reset those models and strictly lower the earlier component-card coordinates. The fixed vector of all cards, unrecognized components and point defects proves finite termination. Exact normal-root covers, component fundamental-group maps and the prior period identities identify the resulting core, while the retained ambient atlas reconstructs every original transverse jet}{the hypotheses, notation and definitions stated here; \cref{cdp:st:C.2.UN.65}; \cref{cdp:st:C.2.UN.68}; \cref{cdp:st:C.2.UN.71}; \cref{cdp:st:C.2.UN.61}; \cref{cdp:st:C.2.UN.62}; \cref{cdp:st:C.2.UN.75}; \cref{cdp:st:C.2.UN.84}; \cref{cdp:st:C.2.UN.57}; \cref{cdp:st:C.2.UN.73}; \cref{cdp:st:C.2.UN.82}; \cref{cdp:st:C.2.UN.77}; \cref{cdp:st:C.2.UN.81}; \cref{cdp:st:C.2.UN.83}; \cref{cdp:eq:C.GP.5}; \cref{cdp:st:C.2.UN.80}; \cref{cdp:eq:C.GP.4}; \cref{cdp:eq:C.GP.8}; \cref{cdp:st:C.2.UN.74}; \cref{cdp:st:C.2.UN.50}; \cref{cdp:st:C.2.UN.28}; \cref{cdp:st:C.2.UN.47}}

Let \(f:X\to\Delta\) be a proper degeneration of connected
complex two-tori, with smooth total space and full reduced
SNC central fibre. There is a finite overlap-compatible
modification record over the same disc, isomorphic over
\(\Delta^*\), ending in one of the following actual
smooth-total reduced degeneration germs.

\begin{enumerate}
\item A canonical-flat core. Its central fibre is a smooth
two-torus, one of the exact no-triple elliptic-ruled/diagonal
Hopf cycles of \ref{cdp:st:C.2.UN.65}, or one of the marked positive toric
log spaces of \ref{cdp:st:C.2.UN.68} and \ref{cdp:st:C.2.UN.71}. The toric alternative has
Euler number equal to its positive triple count.

\item An embedded no-triple core whose components are
minimal fixed-Jacobian elliptic quasi-bundles with identity
integral monodromy, together with any minimal nonlinear
Hopf leaves. All boundary curves are smooth elliptic fibre
supports or the specified Hopf axes. Every elliptic
quasi-bundle base has genus zero or one, and at most one
has genus one. The exact support, including its original
torsors, normal characters, translations and gluing maps,
is one of the rooted-tree or unicyclic supports classified
in \ref{cdp:st:C.2.UN.61} and \ref{cdp:st:C.2.UN.62}. Its Euler number is zero.
\end{enumerate}

If every Hopf-birational prime has minimum canonical order,
the first alternative is supplied by \ref{cdp:st:C.2.UN.75}. If some such
prime has positive order, the second alternative can be
obtained while retaining a fixed maximal-order minimal
Hopf component and its entire neighbourhood after the
preliminary reduction of \ref{cdp:st:C.2.UN.84}. No nonlinear Hopf
neighbourhood is contracted or replaced during this
reduction.

The record retains all actual centres, two-ruling maps,
point and curve contractions, local sigma charts,
normalization branches and restored overlap maps.
Its inverse reconstructs the original analytic family
from this actual endpoint germ. On an originally compact
family all operations take place on that compact base.
If \(S_0\) is the original reduced support, \(S_*\) the
endpoint support, and \(\delta_j\) the signed normal-cone
Euler change of its \(j\)-th retained step as in \ref{cdp:st:C.2.UN.75},
then
\[
                 e(S_0)=e(S_*)-\sum_j\delta_j.
 \tag{C.GP.1}\label{cdp:eq:C.GP.1}
\]

The classified core central spaces have the compact
occurrence constructions of \ref{cdp:st:C.2.UN.57}, \ref{cdp:st:C.2.UN.61}, \ref{cdp:st:C.2.UN.62} and \ref{cdp:st:C.2.UN.73}.
The inverse record here is applied to its actual endpoint
neighbourhood. Identical central log data in a different
smoothing are not asserted to identify its higher analytic
jets or to carry that particular inverse record.
\end{proposition}



\begin{proof}
\noindent\textbf{Step 1.} \par\medskip\noindent\textbf{The initial dichotomy and the protected component.}
Normalize the primitive canonical evaluation as
\[
       \operatorname{div}(\Omega)=\sum_i k_iV_i,
       \qquad \min_i k_i=0.
 \tag{C.GP.2}\label{cdp:eq:C.GP.2}
\]
If every Hopf-birational prime has order zero, \ref{cdp:st:C.2.UN.75}
gives its actual canonical-flat endpoint. \ref{cdp:st:C.2.UN.65} and
\ref{cdp:st:C.2.UN.68}--\ref{cdp:st:C.2.UN.71} give exactly the alternatives in (1),
including the actual marked log maps in the positive
case. This uses the proved component restriction that
every class-VII minimal model is Hopf. There is no
unexamined additional class-VII case.

\noindent\textbf{Step 2.} Suppose instead that a Hopf-birational prime has
positive order. Apply the finite truncated record
\ref{cdp:st:C.2.UN.84}, including its maximal Hopf minimalization. At
its smooth-total NC endpoint put
\[
                  M=\max_i k_i>0.
 \tag{C.GP.3}\label{cdp:eq:C.GP.3}
\]
The minimum is still zero and all orders are at
most this fixed \(M\). Choose one of the surviving
minimal Hopf components \(H\) of order \(M\).
\ref{cdp:st:C.2.UN.82} shows that it has exactly one elliptic
boundary curve \(J\), with torsion normal, and
that its neighbour has smaller order. The marked
neighbourhood is embedded. No elliptic self-double
curve has been introduced in the preliminary
record: its new incidences are on rational
exceptional curves, as proved in \ref{cdp:st:C.2.UN.84}.

\noindent\textbf{Step 3.} All later whole-component removal blocks will be
the positive-order blocks of \ref{cdp:st:C.2.UN.77} and \ref{cdp:st:C.2.UN.81}. Such
a ruled component of order \(q\le M\) has one
horizontal boundary section, whose neighbour has
order \(q-1\), and every other boundary is rational
and vertical. It cannot meet \(H\): a rational
curve cannot lie on a minimal Hopf, and the
section case would require \(q=M+1\). A plane
contraction also has only rational boundary.
The same argument applies to the cleanup blocks
and their missing-fibre insertions. Direct moves
below will likewise avoid \(H\). Thus its entire
neighbourhood is protected and provides the same
starting interface after every completed block.

\par\medskip\noindent\textbf{The recorded models, not a numerical eligibility class.}
At a smooth-total NC stage, call a component
\emph{recognized} only when one of the actual
comparisons of \ref{cdp:st:C.2.UN.83} has been recorded for its
smooth normalization \(T\), together with its
entrance elliptic curve and its point map
\[
                         \pi_T:T\longrightarrow Z_T.
 \tag{C.GP.4}\label{cdp:eq:C.GP.4}
\]
There are three recorded types: an elliptic
quasi-bundle model; a nonlinear Hopf model with
its one axis; or a one-section elliptic ruled
model with its actual rank-two bundle extension.
The point map is an isomorphism near the entrance
curve. All boundary branches, including both
branches of a self-double, are retained.

\noindent\textbf{Step 4.} Let \(d_T\) be the exact number of point blowups
in this map. The following properties are part
of those actual comparisons:
\[
\begin{gathered}
 d_T\ge0,\qquad
 \operatorname{Exc}(\pi_T)\subset D_T,\\
 \text{every rational boundary is exceptional
 and has negative square},\\
 \text{every elliptic boundary has nonpositive square}.
\end{gathered}
 \tag{C.GP.5}\label{cdp:eq:C.GP.5}
\]
For the elliptic type, the boundary is vertical
for the composed elliptic map; \eqref{cdp:eq:C.GP.5} follows
from \ref{cdp:st:C.2.UN.80} and the fibre squares. For the other
two types the entrance curve has square zero,
all point centres are away from it, and every
other boundary is exactly a positive discrepancy
exceptional curve, by (C.TI.2)--(C.TI.3).

\noindent\textbf{Step 5.} Initially recognize \(H\) alone. \ref{cdp:st:C.2.UN.83} applied
to its actual isolated torsion-normal axis gives
the elliptic or nonlinear Hopf type, and its
point defect is zero. Let
\[
\begin{gathered}
 u=\#\{\text{current unrecognized global component primes}\},
 \\
 d=\sum_{T\ \mathrm{recognized}}d_T.
\end{gathered}
 \tag{C.GP.6}\label{cdp:eq:C.GP.6}
\]
Recognition changes no complex space: it records
the actual maps of \ref{cdp:st:C.2.UN.83} and their marked curves.
In particular the word does not assert that an
unrecognized receiving surface can never acquire
such a comparison at a later step.

\par\medskip\noindent\textbf{A direct move on a recorded elliptic model.}
Suppose that a recognized elliptic-type component
has positive defect. Choose a last exceptional
curve of \eqref{cdp:eq:C.GP.4}. It is an actual boundary by
\ref{cdp:st:C.2.UN.80}. The complete elliptic fibre graph and the
exact receiving-degree comparison (C.AP.5)--
(C.AP.6) apply without any additional sign
assumption on the entrance coefficient.

\noindent\textbf{Step 6.} Here is the exact comparison with an embedded
curve in the total space at an NC stage. On a
complete fibre graph of the point map, write
\(a_C\) for the coefficient of each component
in the actual canonical divisor on the
normalization, assigning coefficient zero
also to a component absent from the actual
boundary. Adjacent shifted coefficients
\(a_C+1\) are relatively prime. There is no
initial edge on the smooth elliptic fibre;
a blowup at a smooth point creates consecutive
shifted coefficients, and a node blowup inserts
the sum of the two old shifted coefficients.
Thus the assertion follows through the actual
point map, including negative initial coefficients.
In particular no edge has coefficients
\((-1,-1)\).

\noindent\textbf{Step 7.} A self-double branch has coefficient \(-1\).
If the image of the chosen last \((-1)\)-curve
were nodal at an actual triple point with two
local branches of its source component, its
two neighbours on the normalization would be
the two self-double branches. The last-node
coefficient formula would also give
\(a_E=(-1)+(-1)+1=-1\), contradicting that
edge invariant. A triple with three branches
of the source gives adjacent \(-1\) branches
directly. With only one local source branch,
two image branches would already make a node
on the chosen curve in that smooth source
normalization, contrary to its being the
smooth last exceptional \(\PP^1\).
The actual image curve is consequently
smooth. For a self-double its normalized
degree-two cover is unramified, by the local
NC coordinates. A connected such cover by
\(\PP^1\) is impossible by Riemann--Hurwitz;
the cover therefore splits into two distinct
normalized branches. Over the resulting
simply connected image \(\PP^1\), those
labels give actual two smooth surface germs.
When the triple count is two, the normal-product
map gives the two normal degrees \(-1,-1\).
The actual type-II neighbourhood and its two
ruling maps therefore apply even if these
two surface germs belong to one global prime.
For the positive-discrepancy models below,
all rational coefficients are positive, so
the same self-double obstruction is immediate.

\noindent\textbf{Step 8.} With two triple points, the actual type-II move
blows down this curve on its two adjacent
normalized surfaces. Each opposite receiving
boundary has square at least zero, or at least
one when it is elliptic. By \eqref{cdp:eq:C.GP.5} neither
receiver was recognized before this move. Thus
the move adds no point defect to a recognized
surface.

\noindent\textbf{Step 9.} If its second source surface was recognized,
the contracted rational boundary was exceptional
for its recorded point map, again by \eqref{cdp:eq:C.GP.5}.
That map factors through the actual blowdown:
it is constant on the contracted curve, and
the proper surface contraction and its direct
image of the structure sheaf give this descent.
The resulting birational surface morphism is
the remaining sequence of point blowups. Its
defect decreases by one. The same factorization
applies to the first source. All recorded
entrance curves are disjoint from these
exceptional curves and are unchanged. Hence
the recognized set and its models persist,
\(u\) is unchanged, and \(d\) strictly decreases.
If both source germs belong to one global
recognized component, their two normalized
curves are disjoint. The same point map
factors through both ordinary contractions;
its defect drops by two while the global
component is counted only once.

\noindent\textbf{Step 10.} If a last curve has one triple point, use the
actual positive ruled-removal block of (C.AP.7).
If it has none, use the exact plane progress
block (C.AP.8)--(C.AP.13). Each gives a strict
component-card decrease after its finite cleanup.
No nonboundary last-point case is admitted: its
exclusion is the actual field statement \ref{cdp:st:C.2.UN.80},
not an inference from nonzero canonical coefficients.

\par\medskip\noindent\textbf{A direct move on either positive-discrepancy model.}
Now let the recognized model be nonlinear Hopf
or one-section ruled, with positive defect.
The actual canonical divisor is
\[
 K_T=-aJ_T+\sum_E b_EE,
 \qquad b_E>0,
 \qquad a\ge2,
 \tag{C.GP.7}\label{cdp:eq:C.GP.7}
\]
with \(a=2\) for the ruled type. Every point
centre is away from the preserved \(J_T\),
and the full other boundary is exactly the
exceptional set. The graph of each exceptional
tree starts with a rational curve of square
\(-1\) and valence zero. If \(C^2=-b\) and
\(n\) is its valence in this entire exceptional
set, the actual smooth and node blowups give
\[
                              n\le b+1.
 \tag{C.GP.8}\label{cdp:eq:C.GP.8}
\]
At a smooth point both \(n\) and \(b\) increase
by one on an old component; at a node \(b\)
increases and its old valence is unchanged.
A newly inserted curve has \((b,n)=(1,0)\),
\((1,1)\), or \((1,2)\), according as its
centre is off the exceptional support, on
one component, or at a node. These are
exactly the updates proving \eqref{cdp:eq:C.GP.8}.
At the last node the opposite
square is \(b+1-t\ge0\), since the actual
triple count \(t\) is at most the old valence.
All opposite curves here are rational.
Thus no receiver is already recognized by
\eqref{cdp:eq:C.GP.5}, and the exact point-map factorization
in the preceding paragraph again strictly
decreases \(d\). The positive discrepancies
also give the precise height comparison:
if the two old coefficients are \(b_1,b_2>0\)
and the current component has order \(p\),
the receiving orders are
\[
                 p+b_1+1,\qquad p+b_2+1.
 \tag{C.GP.9}\label{cdp:eq:C.GP.9}
\]

\noindent\textbf{Step 11.} A last curve with no triple point has a
rational \(+1\) receiver of order \(p+2\),
and (C.AP.8)--(C.AP.13) supply the finite
plane progress block. With one triple point
its receiver is ruled with one section and
has positive order, so the same ruled block
applies. These changes concern only the
actual rational exceptional configurations,
not the nonlinear Hopf neighbourhood itself.

\noindent\textbf{Step 12.} When a one-section ruled model reaches defect
zero, it is the actual bundle \(\PP(V)\)
with sole boundary section \(J_T\), including
its possibly nonzero extension class. \ref{cdp:st:C.2.UN.74}
gives the ordinary curve contraction along
its ruling. Its order is
\(k_T=k_{\mathrm{neighbour}}+1>0\), by \ref{cdp:st:C.2.UN.83}.
This is another whole-card decrease.

\par\medskip\noindent\textbf{Recognition across the next actual interface.}
We always eliminate positive recorded defects
before recognizing another component. We also
remove a defect-zero one-section ruled model
by the preceding actual contraction. If neither
action is available, every recognized component
is a minimal elliptic quasi-bundle or a minimal
nonlinear Hopf leaf.

\noindent\textbf{Step 13.} Such a component is embedded and has disjoint
smooth elliptic boundary curves. Indeed no
rational boundary can remain when its point
defect is zero. The original fibre was SNC,
and the retained type-II and one-section
sigma records create new self-incidences only
on rational fibre or exceptional curves.
Their generic one-section contraction cannot
identify two elliptic horizontal branches.
Thus there is no new elliptic self-double
curve on one of these minimal components.

\noindent\textbf{Step 14.} If an unrecognized component remains, the
connected incidence graph of the central
support has an edge from the recognized
region to its complement. Starting with
\(H\) and recording each new entrance edge
makes that region connected; the direct
moves have retained these entrance curves.
The edge is therefore a smooth elliptic
curve, isolated on its recognized side.
There is no triple point on it, so its
other normalization branch is isolated
as well. Its normal is torsion: on a
quasi-bundle it is the exact normal root
of a smooth multiple fibre, and on a
nonlinear Hopf it is its axis character.
The actual reduced-fibre normal-product
map gives the inverse torsion line on
the other side.

\noindent\textbf{Step 15.} Apply \ref{cdp:st:C.2.UN.83} to that actual other normalization.
It supplies one of the three point models
\eqref{cdp:eq:C.GP.4}, with the entrance neighbourhood
unchanged and every original boundary
retained. Record it as recognized. Then
\(u\) decreases by one. The newly recorded
defect may be any finite nonnegative integer;
we make no assertion that this recognition
alone decreases \(d\).

\par\medskip\noindent\textbf{Whole blocks may reset the recognized region.}
A whole-component removal and its sigma
cleanup need not preserve the point models
on the previously recognized region. For
example a missing-fibre insertion can create
additional point blowups on a section parent.
We allow all such changes explicitly.
After this completed block discard every
recognition except the unchanged maximal
Hopf \(H\), and restart with its same
torsion-normal interface. This is possible
because its entire neighbourhood is protected
as proved above. There is no assertion that
the rest of the processed region was protected.

\noindent\textbf{Step 16.} Let \(c_j\), for \(1\le j\le M\), count
the retained global components of order \(j\).
The full fixed potential is
\[
                 (c_1,c_2,\ldots,c_M,u,d).
 \tag{C.GP.10}\label{cdp:eq:C.GP.10}
\]
Every direct type-II move preserves the card
vector and \(u\), and strictly decreases \(d\).
Recognition preserves the cards and strictly
decreases \(u\), regardless of its finite
new defect. Every whole-component block
strictly decreases its first changed card,
regardless of the finite reset of \(u,d\).
The actual missing-fibre and cleanup rules
add only finitely many cards at the next
order and add none at \(M\). Thus all
retained orders stay at most the fixed \(M\).

\noindent\textbf{Step 17.} Lexicographic order on this fixed finite
product of nonnegative integers is well
founded. The internal plane progress and
sigma cleanup blocks are finite by \ref{cdp:st:C.2.UN.81}
and \ref{cdp:st:C.2.UN.77}. Consequently this procedure
cannot continue indefinitely.

\noindent\textbf{Step 18.} At its endpoint no recorded defect is
positive, no one-section ruled model
remains, and no component is unrecognized.
Otherwise the preceding paragraphs give
another strict decrease. All components
are therefore the minimal elliptic
quasi-bundles and nonlinear Hopf leaves
in (2), and every boundary is smooth
elliptic and isolated. In particular the
endpoint is embedded SNC and has no
triple point.

\par\medskip\noindent\textbf{Identification of the endpoint and the retained analytic scope.}
Every endpoint component and double curve
has topological Euler number zero, so
\(e(S_*)=0\). The exact invariant relative
cohomology frame gives
\(p_g(Z)=g(B)\) for a quasi-bundle. The
common regular Jacobian is retained by
the actual normal-root maps: along an
edge joining two quasi-bundle vertices,
the dual normal lines have the same
exact order, and inversion identifies
their connected root covers of the
elliptic double curve. These are exactly
the two regular elliptic fibres, by
(C.QT.2)--(C.QT.3); their translations
are retained in the gluing record.
At a nonlinear Hopf leaf the same root
cover is instead its exact period cover
of \ref{cdp:st:C.2.UN.50}. Such leaves do not interrupt
propagation through the connected
quasi-bundle region, and their actual
leaf periods and finite characters are
retained by that comparison.
The
compact-core bound \ref{cdp:st:C.2.UN.28} therefore permits
at most one genus-one base. If there is
one, that normalization has \(p_g=1\).
\ref{cdp:st:C.2.UN.47} gives the actual surjection from its
fundamental group onto that of the full
support. The graph-of-spaces comparison
gives a surjection from the support
fundamental group to that of its incidence
graph, and kills the image of this
component's fundamental group.
The graph fundamental group is therefore
trivial, so the connected graph is a tree.
The exact coherent restriction and
fundamental-group comparisons of \ref{cdp:st:C.2.UN.61}
then give its rooted-tree support, with the
actual integral Chern and recursive
period conditions. If there is none,
\ref{cdp:st:C.2.UN.62} gives the unicyclic support and
its exact torsor, height and common-period
conditions. Its nonlinear leaves retain
the secondary finite characters of \ref{cdp:st:C.2.UN.50}.
These applications keep the original
surfaces and gluing maps; they do not
replace the actual endpoint smoothing.

\noindent\textbf{Step 19.} All prior operations are the actual
finite diagrams already used in \ref{cdp:st:C.2.UN.75},
\ref{cdp:st:C.2.UN.77}, \ref{cdp:st:C.2.UN.81} and \ref{cdp:st:C.2.UN.84}. Their proper graphs
glue by the retained overlap maps, so
their inverse recovers this given
analytic family. Equation \eqref{cdp:eq:C.GP.1}
is the telescoping normal-cone formula
on that finite record, including its
restored-overlap subdivision. Annular
gluing to an unchanged compact outside
gives the same actual modifications
when the starting total space is compact.

\noindent\textbf{Step 20.} The separate compact constructions for
the classified core support data prove
their occurrence. They do not identify
arbitrary higher transverse jets of
this actual endpoint with another core
smoothing. No inverse sigma record is
transferred through an unproved such
identification. This is the claimed
finite reduction and exact reconstruction
of every given reduced SNC torus
degeneration.
\end{proof}






The data in the following construction include analytic
neighbourhoods and their transition maps. This is essential
when an inverse operation uses a normal extension or higher
transverse terms. We first specify those data without
referring to a family to be reconstructed.

\paragraph{Finite admissible data.}
The following are the inputs and the checks on them.
\begin{enumerate}
\item An actual proper smooth-total core family
\(f_*:X_*\to\Delta_s\), with its labelled analytic
charts and overlaps and its punctured torus-family
comparison, is given. Its full reduced central space
is one of the cores in \ref{cdp:st:C.2.UN.85}: a smooth torus; a
canonical-flat marked Hopf/degree-zero-ruled cycle;
a positive periodic toric log space; or an embedded
rooted-tree or unicyclic elliptic quasi-bundle/Hopf
support of \ref{cdp:st:C.2.UN.61}--\ref{cdp:st:C.2.UN.62}. The core chart functions,
including their transverse terms, are retained.

\item A finite inverse overlap-compatible sigma record
on that actual core is given. Each arrow includes
the local blowup or contraction map, its centre
and ideal, the target charts, and the isomorphism
on the restored overlap. The maps and their inverse
comparisons are checked on these charts as in \ref{cdp:st:C.2.UN.75}.
The resulting space \(Z\) is smooth and proper over
\(\Delta_s\), with full reduced SNC central divisor.
The record is an isomorphism on the punctured family.

\item A finite cyclic birational action on the core,
over \(s\mapsto\zeta s\), is given, where \(\zeta\)
has order \(N\). Conjugating it by the specified
inverse record gives holomorphic maps on every
chart of \(Z\), including the central divisor.
Their compositions satisfy the cyclic group law.
Thus it is an actual regular action of
\(G=\mu_N\) on \(Z\), not only a central action
or an action on its meromorphic function field.

\item A labelled equivariant toric coarsening
\(q:Z\to Y_N\) is given. Its local semigroup
charts, units and maps are explicit normalized
monomial-ideal blowups in the reverse direction.
The finite quotient charts of \(Y_N\) have
invariant ring \(\OO_Y\), with smooth coordinates
in which the descended function is
\[
                  t=u\prod_{i=1}^r x_i^{m_i},
              \qquad r\le3,\quad u\in\OO_Y^*,
              \quad m_i\mid N.
 \tag{C.AR.1}\label{cdp:eq:C.AR.1}
\]
The normal charts upstairs are exactly the
normalization of \(s^N=t\) in these coordinates.
The labelled height-one subdivisions defining
\(q\) are unimodular, and their overlap maps
and finite action are retained. These ring and
coordinate conditions verify the smooth quotient
\(Y=Y_N/G\) and its full SNC divisor with
the displayed multiplicities.

\item A finite sequence \(\pi:Y\to X\) of
ordinary inverse point or smooth-curve blowups
is given, with their actual smooth target charts
and gluing maps. Their centres are in the central
fibre. The function \(t\) descends through each
map, as is checked by its displayed pullback.
The final total space \(X\) is smooth. Its
central support is allowed arbitrary singularities
and arbitrary full multiplicities.
\end{enumerate}
All quotient, contraction and overlap requirements here
are checks on the supplied finite maps. They do not
refer to the existence of an additional original family.
An ordinary arrow can equivalently be supplied by its
Rees-algebra charts and the corresponding proper map.
For compact data, the core family and base, the cyclic
base cover and finite action, and all maps outside these
central charts are supplied globally; the local overlaps
agree with them on an annulus. This is the compact
version of the same definition.

\begin{proposition}[Complete rooted generation and the zero-Euler criterion for arbitrary fibres]\label{cdp:st:C.2.UN.86}
\StatementRecord{cdp:st:C.2.UN.86}{proved}{Historical events: Claude Fable 5.2 (Anthropic), effort high, 16 November 2026: carried as background without fault found. GPT-6 Astra (OpenAI), ChatGPT Pro mode, effort Pro, S6\_09\_PAPERS, 17 November 2026: targeted proof check. Also checked in this writing. Pre-continuation writing events: GPT-6 Astra (OpenAI), ChatGPT Pro mode, effort Pro, S6\_10\_PAPERS, 18 November 2026: targeted proof check; used and rechecked in this run. Continuation metadata events: GPT-6 Astra (OpenAI), ChatGPT Pro mode, effort Pro, S6\_10\_PAPERS continuation, 19 November 2026: used and rechecked in this run; status-vocabulary reconciliation; metadata reconciliation only; no new mathematical check. \emph{Pre-continuation recorded proof scope.} The complete finite admissible input and every forward and inverse map were read. Embedded resolution and the actual normalized cyclic root give a connected smooth semistable model; conjugating and then reversing the core record restores the regular group action before the invariant-ring quotient. A nonzero quotient fibre is exactly one original torus, and smoothness plus the non-zero-divisor base function gives flatness and connected fibres. The three signed Euler changes telescope with the exact toric volume and sheet count. The original normal-root algebras and radial density threshold give h=N(H-tau), independently of the chosen resolution. The compact jet interpolation and explicit unit-root coordinate preserve the global cover and the full analytic reconstruction data}{the hypotheses, notation and definitions stated here; \cref{cdp:eq:C.AR.1}; \cref{cdp:st:C.2.UN.84}; \cref{cdp:st:C.2.UN.85}; \cref{cdp:st:C.2.UN.78}; \cite{BierstoneMilman1997}; \cref{cdp:eq:C.AR.10}; \cref{cdp:st:C.2.UN.75}; \cref{cdp:st:C.2.UN.57}; \cref{cdp:st:C.2.UN.61}; \cref{cdp:st:C.2.UN.63}; \cref{cdp:st:C.2.UN.71}; \cref{cdp:st:C.2.UN.73}}

Every proper degeneration of connected complex two-tori
from a smooth threefold, with arbitrary full central
fibre, is the output of the finite admissible construction
above. Conversely every such input produces a proper
degeneration of connected complex two-tori from a smooth
threefold. The inverse record recovers the supplied
family exactly. For compact threefolds over compact
curves, the same assertions hold for the compact data.

Here is an explicit formula selecting precisely the
zero-Euler outputs. Let \(S_*=f_*^{-1}(0)\), let
\(D_i\) be the quotient components in \eqref{cdp:eq:C.AR.1}, and put
\[
 \begin{split}
 D_I^\circ&=\bigcap_{i\in I}D_i
                    \setminus\bigcup_{j\notin I}D_j,\\
 d_I&=\gcd_{i\in I}m_i,\qquad
 \nu_I=\frac{N^{|I|-1}d_I}{\prod_{i\in I}m_i},\\
 R_N&=\sum_I(d_I\nu_I-1)e_c(D_I^\circ).
 \end{split}
 \tag{C.AR.2}\label{cdp:eq:C.AR.2}
\]
For the ordinary resolution record \(\pi\), define
\[
 \Delta_\pi=2p_\pi+\sum_\alpha(2-2g(C_\alpha)).
 \tag{C.AR.3}\label{cdp:eq:C.AR.3}
\]
Let \(\Delta_\sigma\) be the signed normal-cone Euler
change in the forward direction \(Z\dashrightarrow X_*\)
of the retained sigma record. Then its output satisfies
\[
 e\bigl(f^{-1}(0)_{\mathrm{red}}\bigr)
       =e(S_*)-\Delta_\sigma-R_N-\Delta_\pi.
 \tag{C.AR.4}\label{cdp:eq:C.AR.4}
\]
Every zero-Euler fibre is therefore obtained, and only
such fibres are obtained, by the admissible data with
\[
 \Delta_\sigma+R_N+\Delta_\pi=
 \begin{cases}
  2[N_0:\Gamma_0],&\text{for the positive toric core},\\
  0,&\text{for each other core}.
 \end{cases}
 \tag{C.AR.5}\label{cdp:eq:C.AR.5}
\]
Here \(N_0\) and \(\Gamma_0\) are the two-dimensional
affine lattices of the toric core; the root exponent
\(N\) in \eqref{cdp:eq:C.AR.1} is a separate integer. The construction
in \eqref{cdp:eq:C.AR.5} specifies spaces, bundles, finite maps,
centres and analytic gluings, as well as their Euler
numbers. Reversing a retained record recovers every
specified normal extension and every transverse jet.

The choice between a canonical-flat core and a positive
Hopf elliptic-network core is intrinsic to the original
degeneration. Write the original full fibre and the
primitive generator as
\[
        \operatorname{div}(f)=\sum_\alpha n_\alpha S_\alpha,
        \qquad \operatorname{div}(\Omega)
                       =\sum_\alpha a_\alpha S_\alpha.
 \tag{C.AR.6}\label{cdp:eq:C.AR.6}
\]
Let \(C_\alpha\to T_\alpha\) be the actual finite normal
root cover of the normalization of \(S_\alpha\) defined
by (C.RK.18), and let \(\tau\) be the density threshold
(C.RK.19). Set
\[
 H=\max\left\{\frac{a_\alpha+1}{n_\alpha}:
       C_\alpha\text{ has a component with Hopf minimal model}
       \right\},
 \tag{C.AR.7}\label{cdp:eq:C.AR.7}
\]
with \(H=\tau\) when this set is empty. Then \(H\ge\tau\).
For every actual divisible semistable root, its highest
Hopf canonical order, when such a component exists, is
\[
                         h=N(H-\tau).
 \tag{C.AR.8}\label{cdp:eq:C.AR.8}
\]
Thus \(H=\tau\) gives the canonical-flat reduction,
and \(H>\tau\) gives the positive-Hopf reduction of
\ref{cdp:st:C.2.UN.84}--\ref{cdp:st:C.2.UN.85}. The covers, threshold and this branch test
do not depend on an embedded resolution or on its
chosen divisible root exponent.
\end{proposition}



\begin{proof}
\noindent\textbf{Step 1.} \par\medskip\noindent\textbf{The actual forward construction for an arbitrary fibre.}
Shrink the base to a disc about the chosen value.
The proper direct image \(f_*K_X\) is a torsion-free
rank-one sheaf, so it has a primitive generator
\(\Omega\). Its zero divisor is vertical, giving
\eqref{cdp:eq:C.AR.6}. Take an actual embedded resolution
\(\pi:Y\to X\), by smooth point and curve blowups
in the central fibre, with smooth SNC components
\(D_i\). Retain every centre and every blowup map.
The analytic embedded-resolution comparison used
in \ref{cdp:st:C.2.UN.78}, from \cite{BierstoneMilman1997}, gives
\[
 \operatorname{div}(f\pi)=\sum_i m_iD_i,
 \qquad \operatorname{div}(\pi^*\Omega)=\sum_i k_iD_i,
 \qquad N=\operatorname{lcm}_i m_i.
 \tag{C.AR.9}\label{cdp:eq:C.AR.9}
\]

\noindent\textbf{Step 2.} Normalize the literal fibre product for \(t=s^N\)
and use the equivariant semistable maps of \ref{cdp:st:C.2.UN.78}:
\[
 Z\xrightarrow{q}Y_N
     =\operatorname{Norm}(Y\times_{\Delta_t}\Delta_s)
                   \longrightarrow Y.
 \tag{C.AR.10}\label{cdp:eq:C.AR.10}
\]
For clarity, the local root lattice is
\[
 \left\{(v,h)\in\Z^r\oplus\Z:
                      \sum_i m_iv_i=Nh\right\}.
 \tag{C.AR.11}\label{cdp:eq:C.AR.11}
\]
Its old height-one rays are \(((N/m_i)e_i,1)\).
In dimensions at most three, subdividing each
height-one segment or triangle at its lattice
points gives unimodular cones. These refinements
are implemented by normalized blowups of
monomial ideals, which are unchanged when a
boundary equation is multiplied by a unit.
Their labelled maps therefore glue. The root
group fixes the old rays within a chart and
permutes the charts, so orbitwise choices give
the actual equivariant map \eqref{cdp:eq:C.AR.10}. Retaining
the old component and stratum labels prevents
ray identifications within an old cone. Thus
\(Z\) is smooth and its full fibre is reduced
SNC, including globally embedded components.

\noindent\textbf{Step 3.} The normalized root is connected. Its generic
fibre over \(s\ne0\) is the connected original
torus; all components of the normalized fibre
product dominate the base, so there cannot be
an additional component. Equivalently its
Stein factor has degree one over the normal
disc. All the toric maps are proper and are
isomorphisms away from zero.

\noindent\textbf{Step 4.} Apply the complete reduced theorem \ref{cdp:st:C.2.UN.85} to
\(Z\), and retain its actual comparison
\(\phi:Z\dashrightarrow X_*\). This gives
one of the listed cores and all the inverse
chart maps in the input. The root action
becomes the actual birational action
\[
            \gamma_* =\phi\gamma\phi^{-1}
                            \quad(\gamma\in\mu_N).
 \tag{C.AR.12}\label{cdp:eq:C.AR.12}
\]
Its group law holds on the dense punctured
family, hence as meromorphic maps. Reversing
\(\phi\) first restores the original regular
action on \(Z\). Applying the retained
equivariant coarsening then gives \(Y_N\).
For a normal local ring \(A\) of \(Y\),
let \(B\) be the integral closure of
\(A[s]/(s^N-f\pi)\). An invariant element
of \(B\) belongs to \(\operatorname{Frac}(A)\)
and is integral over \(A\), so it belongs
to \(A\). The reverse inclusion is immediate.
This remains true when local root branches
split and are permuted by the group. Thus
\(B^{\mu_N}=A\), and the full comparison is
\[
 X_*\dashrightarrow Z\xrightarrow{q}Y_N
              \longrightarrow Y_N/\mu_N=Y
                       \xrightarrow{\pi}X.
 \tag{C.AR.13}\label{cdp:eq:C.AR.13}
\]
Every arrow and overlap in this diagram has
been retained. In particular the finite
quotient follows the inverse record; no
central quotient of an indeterminate
birational action is taken.

\par\medskip\noindent\textbf{The converse from the specified finite input.}
Begin with the data at the start of the
subsection, independently of any original
\(X\). The inverse sigma maps give the
specified smooth proper \(Z\), and the
chart checks give an actual regular finite
action. Its equivariant coarsening and
invariant-coordinate charts construct the
smooth \(Y\) and its function \eqref{cdp:eq:C.AR.1}.
The subsequent ordinary contraction maps
construct the smooth \(X\) and its descended
function. Properness follows from the
proper core and retained proper maps;
for a local overlap replacement it is the
proper-graph and annular comparison of
\ref{cdp:st:C.2.UN.75}. The ordinary target maps and finite
quotient are proper, so the final family
is proper after the common disc has been
shrunk if required.

\noindent\textbf{Step 5.} For a fixed nonzero \(t\), the \(N\)
choices of \(s\) are permuted freely and
transitively by \(\mu_N\). Inclusion of
any one \(Z_s\) into their union induces
a biholomorphism onto the quotient fibre.
The coarsening and ordinary maps are
isomorphisms there. Consequently the
final general fibre is exactly that
connected complex two-torus. Since the
total space is smooth, the nonconstant
base function is a non-zero-divisor;
the map to the smooth curve is flat.
Its Stein factor has generic degree one
and its base is normal, so all its fibres
are connected. These deductions use only
the supplied maps and their checks.

\noindent\textbf{Step 6.} The inverse elementary record is therefore
a constructive description of the resulting
analytic support, with its actual ambient
neighbourhood. Forward construction shows
that every original support has such a
record. The converse shows that every
admissible record produces such a support.
Isomorphic labelled atlases and commuting
record maps give biholomorphic outputs
by their explicit glued isomorphism.
No uniqueness of the reduction record is
required for this exact generation statement.

\par\medskip\noindent\textbf{Original support Euler numbers, including the finite quotient.}
The actual calculation (C.RK.15)--(C.RK.17)
applies to \eqref{cdp:eq:C.AR.10} without a canonical-order
restriction. Over \(D_I^\circ\), the central
normal root is a finite etale cover of degree
\(d_I\). The toric resolution fibre has Euler
number \(\nu_I\): its zero-dimensional torus
orbits are the unimodular top cones, whose
number is the normalized volume of the
height-one simplex. The exact Bezout
determinant in (C.RK.15) gives the value
in \eqref{cdp:eq:C.AR.2}. Thus
\[
       e(S_Z)-e(S_Y)=R_N.
 \tag{C.AR.14}\label{cdp:eq:C.AR.14}
\]
The fixed-point average for the actual
finite root group gives the same quotient
comparison, including globally nontrivial
sheet permutations. This is not a free-cover
Euler division at the central fibre.

\noindent\textbf{Step 7.} The ordinary blowups give
\(e(S_Y)-e(S)=\Delta_\pi\), and the
forward sigma record gives
\(e(S_*)-e(S_Z)=\Delta_\sigma\).
For an arrow with centre ideal \(I\),
its signed term is the difference between
the Euler numbers of the actual
\(\operatorname{Proj}(\bigoplus I^a/I^{a+1})\)
and its centre, on the finite constructible
subdivision of \ref{cdp:st:C.2.UN.75}. Restored overlaps
cancel. Combining these three exact
comparisons proves \eqref{cdp:eq:C.AR.4}. The core
Euler numbers are those of \ref{cdp:st:C.2.UN.85}, proving
\eqref{cdp:eq:C.AR.5}. Since both construction
directions have already been established,
this selects the full class of zero-Euler
outputs, including singular and nonreduced
original fibres.

\par\medskip\noindent\textbf{Intrinsic root data and the two endpoint branches.}
The local finite algebras defining
\(C_\alpha\) are glued from
\(w_j^{n_\alpha}=u_j|_{T_\alpha}\), where
\(f=g_j^{n_\alpha}u_j\); their transition
is \(w_k=(g_j/g_k)w_j\). Their finite
normalization is an actual compact
normal cover. No global meromorphic
frame of a nonalgebraic normal line
is required. For every divisible
root, the strict-cover comparison is
the integral local function
\(w_j=s^{N/n_\alpha}/g_j\), as proved
in (C.RK.21). Every connected component
of this fixed cover is obtained, with
its actual proper birational resolution.

\noindent\textbf{Step 8.} All exceptional surfaces of the original
embedded resolution are projective, and
their finite normal covers are projective
by the coherent-algebra GAGA comparison
in \ref{cdp:st:C.2.UN.78}. The additional semistable
components are ruled or rational.
Thus no Hopf component is omitted by
the original covers \(C_\alpha\).

\noindent\textbf{Step 9.} Proper change of variables in the actual
density gives
\[
                   \tau=\min_i\frac{k_i+1}{m_i}.
 \tag{C.AR.15}\label{cdp:eq:C.AR.15}
\]
At a height-one root ray \((v,1)\),
the pulled canonical form has order
\(\sum_i(k_i+1)v_i-1\). Its minimum
is \(N\tau-1\); the old strict ray
over \(D_i\) has order
\(N(k_i+1)/m_i-1\). Dividing by
\(s^{N\tau-1}\) gives the primitive
root generator, whose strict orders
are therefore
\[
                  N\left(\frac{k_i+1}{m_i}-\tau\right).
 \tag{C.AR.16}\label{cdp:eq:C.AR.16}
\]
Original strict transforms retain
\((k_i,m_i)=(a_\alpha,n_\alpha)\).
The fixed-cover comparison and exclusion
of exceptional Hopfs now give exactly
\eqref{cdp:eq:C.AR.8}, including its integrality.
When no Hopf occurs, \ref{cdp:st:C.2.UN.75} applies
directly. Otherwise all Hopf root
orders are zero exactly when
\(H=\tau\). If \(H>\tau\), \ref{cdp:st:C.2.UN.84}
and \ref{cdp:st:C.2.UN.85} give the positive-Hopf core.
Both the threshold and the covers
were defined on the original family,
so this branch test is independent
of the resolution and exponent.

\par\medskip\noindent\textbf{Compact data and the extent of the core parametrizations.}
For an original compact base \(B\), let
\(c\) be the selected value and let
\(\Sigma\) be the finite critical-value set.
Choose a point \(p\notin\Sigma\cup\{c\}\).
For sufficiently large \(d\),
\[
 H^1\!\left(B,\OO_B(dp-2c-
                    \sum_{v\in\Sigma\setminus\{c\}}v)\right)=0.
\]
Indeed its degree is greater than \(2g(B)-2\),
so Serre duality gives the vanishing. The
actual restriction sequence consequently
surjects from \(H^0(B,\OO_B(dp))\) onto
the prescribed first jet at \(c\) and
values at the other critical points.
Choose the meromorphic function \(a\)
with jet \(a=t+O(t^2)\) in the fixed
coordinate at \(c\), and value one at
every other point of \(\Sigma\).
Its only possible pole is \(p\), and
all its other zeros and poles are regular
values of the original family.

\noindent\textbf{Step 10.} Normalize the compact cyclic base cover
\(s^N=a\). The simple zero at \(c\)
makes it connected and totally ramified
there. It is etale over every other
critical value. Near its additional
branch values the original map is a
submersion, so its base change is locally
a product with the smooth normalized
cover curve and is smooth. Over the
other critical values the base change
is etale over the original smooth total
space. The normalized total space is
therefore smooth away from the selected
root fibre. Apply the preceding root
resolution and modification maps there,
retaining the actual identity comparisons
on its annulus and on the complement.
The local function \(a\) is a coordinate
at \(c\). Write \(a=t u(t)\), with
\(u(0)=1\), and choose its analytic
\(N\)-th root near \(c\). The map
\[
 (t,s)\longmapsto
 \left(t,\sigma=\frac{s}{u(t)^{1/N}}\right),
 \qquad \sigma^N=t,
\]
is an actual equivariant isomorphism
between the global-cover root chart
and the retained local root model;
its inverse is \(s=\sigma u(t)^{1/N}\).

\noindent\textbf{Step 11.} The finite action is conjugated by that
actual global modification and becomes
regular again upon reversal. The quotient
and original blowdowns recover the
original compact family. Conversely the
globally supplied compact admissible
data give a compact output by the same
proper maps.

\noindent\textbf{Step 12.} The central support parameters of the
flat cycles and elliptic networks are
the complete marked line, torsor,
finite-character and gluing data in
\ref{cdp:st:C.2.UN.57} and \ref{cdp:st:C.2.UN.61}--\ref{cdp:st:C.2.UN.63}. The positive
toric core is parametrized by its
periodic affine triangulation and
normalized log gluing data in \ref{cdp:st:C.2.UN.71}.
Their compact occurrence is proved
in those statements and \ref{cdp:st:C.2.UN.73}.
Those constructions supply actual
compact starting families for records
defined on their neighbourhoods.
For recovery of a different given
family, its actual endpoint atlas
and higher transverse terms remain
in the input. The proof uses its
specified overlap isomorphisms and
does not transfer a contraction
through an equality of central log
data alone. This states exactly
which parameters are classified
centrally and which analytic data
the reversible reconstruction retains.
\end{proof}






\section{Global CDP replacement and normalization resonance}\label{sec:II:12}


\begin{theorem}[The Euler sign for every degeneration and its equality criterion]
\label{cdp:st:C.2.G.1}
\StatementRecord{cdp:st:C.2.G.1}{proved}{Historical events: Claude Fable 5.2 (Anthropic), effort high, 16 November 2026: carried as background without fault found. Claude Fable 5.2 (Anthropic), effort high, 17 November 2026: carried as background without fault found. GPT-6 Astra (OpenAI), ChatGPT Pro mode, effort Pro, S6\_09\_PAPERS, 17 November 2026: targeted proof check. Also checked in this writing. Pre-continuation writing events: GPT-6 Astra (OpenAI), ChatGPT Pro mode, effort Pro, S6\_10\_PAPERS, 18 November 2026: targeted proof check; used and rechecked in this run. Continuation mathematical and source checks: GPT-6 Astra (OpenAI), ChatGPT Pro mode, effort Pro, S6\_10\_PAPERS continuation, 19 November 2026: targeted proof check; used and rechecked in this run. Continuation metadata events: GPT-6 Astra (OpenAI), ChatGPT Pro mode, effort Pro, S6\_10\_PAPERS continuation, 19 November 2026: used and rechecked in this run; status-vocabulary reconciliation; metadata reconciliation only; no new mathematical check. \emph{Pre-continuation recorded proof scope.} The base, finite-character, smooth cover, irreducible rational and ruled Euler calculations, local Hopf and positive-b2 class-VII exclusions, arbitrary reducible conductor comparison, Euler sum and parity were re-derived in the revised seven-step proof. The proof does not apply P4.58 outside its all-support irreducible-or-NC hypothesis; it uses the explicitly local P4 comparisons and MH.2}{\ref{cdp:st:D.1.9}; \ref{cdp:st:A.3.4};
\ref{cdp:st:A.5.2}; \ref{cdp:st:A.5.3};
\ref{cdp:st:C.2.MH.1}; \ref{cdp:st:C.2.MH.2};
\ref{night:P4.26}; \ref{night:P4.32};
\ref{night:P4.36}; \ref{night:P4.37};
\ref{cdp:st:C.2.INT.1}; \ref{cdp:st:C.2.INT.2};
\ref{cdp:st:C.2.INT.3}; \ref{cdp:st:C2a};
\cite[Proposition~2.1 and Theorem~2.2, pp.~842--843]{Peters2021}}

Let $Y$ be a connected smooth compact complex threefold with
$b_1(Y)=b_2(Y)=0$, and let $f:Y\to B$ be a proper surjective
holomorphic map with connected fibres to a smooth compact curve.
Assume that the smooth general fibre is a complex two-torus.
The full fibres may have arbitrary positive component multiplicities
and their reduced supports may have arbitrary singularities. Then
\[
 B\simeq\PP^1,\qquad
 c_3(Y)=e(Y)\in\{0,2\},\qquad b_3(Y)=2-c_3(Y).
 \tag{C.G.1}\label{cdp:eq:C.G.1}
\]
Every reduced fibre support has nonnegative Euler number.
An irreducible singular support with rational normalization
resolution has Euler number two; every other irreducible support
has Euler number zero.

The equality $c_3(Y)=0$ holds exactly when every support has one
of the following descriptions. A smooth irreducible support has
a finite unramified torus cover. A singular irreducible support
has smooth minimal elliptic ruled or smooth minimal Hopf
normalization. A reducible support has normalization a disjoint
union of smooth minimal Hopf or non-K\"ahler elliptic surfaces;
its reduced conductor image and inverse image are disjoint
unions of smooth elliptic curves, and the conductor branch maps
are finite unramified covers. Each description concerns the
normalization and resolution maps of the actual fibre. In
particular the full-fibre and conductor multiplicities are
retained.
\end{theorem}


\begin{proof}
\noindent\textbf{Step 1. The base is rational and the local
finite character is available on every full fibre.}
The beginning of connected-fibre Leray injects
$H^1(B,\Q)$ into $H^1(Y,\Q)=0$, so $B=\PP^1$.
Proposition~\ref{cdp:st:D.1.9} supplies an integral section
of $R^1f_*\Z$ which restricts nontrivially to each reduced
support and to each nearby torus. On a small disc its
Leray lift is an integral class on the entire tube.
Reduction modulo a prime which does not divide its fibre
restriction gives the finite cyclic character used in
\ref{cdp:st:C.2.MH.2}. The construction is on the original
full fibre and requires no change of its multiplicity.

\medskip\noindent\textbf{Step 2. A normal irreducible support
is smooth and has a finite unramified torus cover.}
Proposition~\ref{cdp:st:A.5.2} excludes a singular normal
support. For a smooth support $S$ in a full fibre $mS$,
the normalized change $t=s^m$ is, in a local equation
$t=uz^m$, the map
\[
 w^m=u,\qquad s=wz.
\]
These charts give a finite unramified central cover of $S$
and a proper holomorphic submersion over the new disc,
as in \ref{cdp:st:A.3.4}. Its nearby fibre is the
original torus. Smooth proper transport gives $b_1=4$
and rational first Chern class zero on the central cover.
It is K\"ahler by the even-$b_1$ criterion, and the
surface classification identifies it as a torus; the
precise classification facts are
\cite[Proposition~2.1 and Theorem~2.2, pp.~842--843]{Peters2021}.
Thus $e(S)=0$.

\medskip\noindent\textbf{Step 3. The irreducible nonnormal
normalization falls into three explicit resolution classes.}
Write $Z\to T\to S$ for the minimal resolution of the
normalization. The canonical evaluation and
\ref{cdp:st:A.5.3} give a nonempty effective divisor $D$
with $K_Z+D$ torsion. A positive pluricanonical section,
raised to a power killing this torsion, would give a
nonzero holomorphic function vanishing along a positive
multiple of $D$. Compactness excludes that function.
Hence $\kappa(Z)=-\infty$. The classification in
\cite[Theorem~2.2, p.~843]{Peters2021} gives a rational
resolution, a resolution ruled over a curve of positive
genus, or a resolution of class VII. The actual maps
$Z\to T\to S$ remain the maps in every subsequent
comparison.

\medskip\noindent\textbf{Step 4. Every irreducible support has
the claimed exact Euler value and normalization.}
For a rational resolution,
\ref{cdp:st:C.2.MH.2} gives $e(S)=2$. For a resolution
ruled over a curve of genus $g\ge1$, the same proposition
gives $e(S)=2-2g$. The local irreducible inequality in
\ref{night:P4.26} gives $e(S)\ge0$, and therefore $g=1$
and $e(S)=0$. Its proof uses the effective conductor,
the local trace inequality, and the fundamental group
of this support; no assertion about the other fibre
supports enters that local inequality. The zero-Euler
comparison in \ref{night:P4.32} then identifies $T$ as
a smooth minimal elliptic ruled surface.

For a class-VII resolution, \ref{night:P4.37} excludes
every minimal model with positive second Betti number
by its saturated-foliation Chern identity and the
negative-definite curve form. The curve-free Inoue
case cannot carry the nonempty effective conductor.
The remaining Hopf case is treated by
\ref{night:P4.36}: its derivative and foliation
calculations force the actual normalization to be
smooth and minimal, with $e(S)=0$. These arguments
are local comparisons for the given irreducible
support; they do not invoke the irreducible-or-normal-
crossings hypothesis on all fibres in
\ref{night:P4.58}. This proves the asserted exact
irreducible classification in the present scope.

\medskip\noindent\textbf{Step 5. Arbitrary reducible supports
have nonnegative Euler number, with the stated equality data.}
For each normalization resolution, an adjacent global
component supplies a nonzero effective divisor with
rationally trivial first Chern class. Its positive
pairing with a proposed K\"ahler form excludes a
K\"ahler resolution. Proposition~\ref{cdp:st:C.2.INT.1}
gives the actual Picard restriction to each elliptic
curve and excludes a Hopf--Enoki or elliptic--Enoki
interface. Proposition~\ref{cdp:st:C.2.INT.2} excludes
the remaining all-Enoki possibility using the full
multiplicities, the canonical coefficient maximum,
and the local vanishing sphere of square $-2$.

Proposition~\ref{cdp:st:C.2.INT.3} now applies to the
entire reduced support. With $A_i$ its normalization
conductor and $C$ its reduced image, its actual
normalization map gives
\[
 e(S)=\sum_i e(T_i\setminus A_i)+e(C),
 \qquad e(T_i\setminus A_i)\ge0,\quad e(C)\ge0.
\]
The proof retains the scheme conductor in the
coherent-duality argument and all contractions in
the topological calculation. It proves both
directions of the equality characterization:
the normalizations are the stated smooth minimal
surfaces, the conductor curves are smooth elliptic,
and their maps are finite unramified covers.

\medskip\noindent\textbf{Step 6. The global sum is zero or two.}
The finite-critical-value and Euler-additivity
calculation of \ref{cdp:st:C2a}, together with
complex-oriented Chern--Gauss--Bonnet, gives
\[
 c_3(Y)=e(Y)=\sum_{b\in B}e((Y_b)_{\rm red})\ge0.
 \tag{C.G.2}\label{cdp:eq:C.G.2}
\]
The smooth torus contributes zero over the regular
base. Poincar\'e duality gives $e(Y)=2-b_3(Y)$.
The nonsingular middle cup-product pairing is
skew-symmetric over $\Q$, so $b_3(Y)$ is even.
Its nonnegativity now proves \eqref{cdp:eq:C.G.1}.

\medskip\noindent\textbf{Step 7. Vanishing is equivalent to
the listed geometric data on every fibre.}
Every term in \eqref{cdp:eq:C.G.2} is nonnegative;
the sum vanishes precisely when every term does.
Steps 2 and 4 give the irreducible forms, and
Step 5 gives the reducible forms with their exact
conductor maps. Conversely each of these forms has
zero Euler number by the same local calculations,
so their occurrence on every support makes the
global sum zero.
\end{proof}








\end{document}
